\documentclass[10pt]{amsart}
\usepackage{amsmath,amssymb,amsthm,mathtools,mathrsfs}
\usepackage[a4paper,hmargin=0.76in,vmargin=0.60in]{geometry}
\usepackage{booktabs}
\usepackage{array}
\usepackage{longtable}
\usepackage{tabularx}
\usepackage[hypertexnames=false]{hyperref}
\hypersetup{hidelinks}
\usepackage[final,expansion=false]{microtype}
\usepackage{url}
\emergencystretch=2.5em

\newtheorem{theorem}{Theorem}[section]
\newtheorem{proposition}[theorem]{Proposition}
\newtheorem{lemma}[theorem]{Lemma}
\newtheorem{corollary}[theorem]{Corollary}
\theoremstyle{definition}
\newtheorem{definition}[theorem]{Definition}
\newtheorem{problem}[theorem]{Problem}
\theoremstyle{remark}
\newtheorem{remark}[theorem]{Remark}

\newcommand{\C}{\mathbb C}
\newcommand{\R}{\mathbb R}
\newcommand{\N}{\mathbb N}
\newcommand{\Z}{\mathbb Z}
\newcommand{\Q}{\mathbb Q}
\newcommand{\Rel}{\operatorname{Re}}
\newcommand{\Iml}{\operatorname{Im}}
\newcommand{\supp}{\operatorname{supp}}
\newcommand{\TV}{\operatorname{TV}}
\newcommand{\Kp}{\mathcal K^{\mathrm{prim}}_{2T}}
\newcommand{\Kw}{\mathcal K_{2T}}
\newcommand{\PT}{\mathcal P_T}
\newcommand{\IT}{I_T}
\newcommand{\lamP}{\lambda^{\mathrm{prim}}_T}
\newcommand{\lamW}{\lambda^{\mathrm W}_T}
\newcommand{\lamA}{\lambda_{\mathrm A}(T)}
\newcommand{\sT}{\sigma_T}
\newcommand{\gh}[1]{\widehat g_{#1}}
\newcommand{\nuh}{\widehat\nu}
\newcommand{\wh}{\widehat w}
\newcommand{\Fh}{\widehat F}
\newcommand{\cF}{c_F}
\newcommand{\Lev}{\mathfrak E}
\newcommand{\Wcl}{\mathcal W_T}
\newcommand{\NT}{\mathcal N_T}
\newcommand{\Aconst}{\mathsf A_T}
\newcommand{\Ainf}{\mathcal A_\infty}
\newcommand{\PiT}{\Pi_T}
\newcommand{\mT}{\mathfrak m_T}
\newcommand{\laminf}{\lambda^{\infty}_{\min}}
\newcommand{\laminfP}{\lambda^{\infty,\mathrm P}_{\min}}
\newcolumntype{Y}{>{\raggedright\arraybackslash}X}

% macros of the series: every macro has in all nine papers the meaning it had in the single
% paper; the definitions needed before their place in the source files are collected here
\providecommand{\anPa}{\Phi_a}
\providecommand{\anPb}{\Phi_b}
\providecommand{\auP}{\mathcal P}
\providecommand{\auPa}{\Phi_a}
\providecommand{\auS}{\mathcal S}
\providecommand{\auSram}{\mathcal S^{\mathrm{ram}}}
\providecommand{\augain}{\operatorname{gain}}
\providecommand{\aurem}{\operatorname{rem}}
\providecommand{\ausfin}{s_{\mathrm{fin}}}
\providecommand{\ausinf}{s_\infty}
\providecommand{\ausmain}{s_{\mathrm{main}}}
\providecommand{\pao}{\mathfrak o}
\providecommand{\parsgn}{\operatorname{sgn}}
\providecommand{\pbcrd}{\mathrm{rd}_E}
\providecommand{\pbdeg}{\operatorname{deg}}
\providecommand{\pbeq}{\mathsf q}
\providecommand{\pbpCT}{\operatorname{CT}}
\providecommand{\pbvol}{\operatorname{vol}}
\providecommand{\pcoB}{\mathsf B}
\providecommand{\pcom}{\mathsf m}
\providecommand{\pdsK}{\mathbf K}
\providecommand{\pexmc}{\mathfrak c}
\providecommand{\pfT}{\operatorname{Tr}}
\providecommand{\plfK}{\mathsf K}
\providecommand{\pmA}{\mathcal A}
\providecommand{\pmGa}{\Gamma}
\providecommand{\pmP}{\mathsf P}
\providecommand{\pmQ}{[G]}
\providecommand{\pmS}{\mathfrak S}
\providecommand{\pmell}{\mathrm{ell}}
\providecommand{\pmtr}{\operatorname{tr}}
\providecommand{\ppidisc}{\operatorname{disc}}
\providecommand{\psM}{\mathsf M}
\providecommand{\psrR}{\mathcal R}
\providecommand{\psrm}{\mathfrak m}
\providecommand{\psrmc}{\mathsf{mc}}
\providecommand{\psuC}{\mathcal C}
\providecommand{\pwfA}{\mathsf A}
\providecommand{\pwfL}{\mathsf L}
\providecommand{\pwfM}{\mathsf M}
\providecommand{\pwfev}{\operatorname{ev}_0}
\providecommand{\pzF}{\mathfrak F}
% statements quoted from the other papers of the series (generated appendix)
\newenvironment{quotedthm}[2]{\par\medskip\noindent\begingroup
  \textbf{#1}#2.\ \itshape\ignorespaces}{\endgroup\par\medskip}
\newenvironment{quoteddef}[2]{\par\medskip\noindent\begingroup
  \textbf{#1}#2.\ \ignorespaces}{\endgroup\par\medskip}
\begin{document}
\addtolength{\abovedisplayskip}{-3.4pt}
\addtolength{\belowdisplayskip}{-3.4pt}
\addtolength{\belowdisplayshortskip}{-2pt}
\title[The Covolume Theorem: the chain determinant]{The Covolume Theorem: exchange identities and the determinant of a chain of theta lifts}
\author{David Alejandro Trejo Pizzo}
\date{}
\begin{abstract}
We study a chain of $m$ theta lifts between unitary groups over a CM field, which carries forms of
discrete series type on an indefinite group to algebraic modular forms on a definite group, the
$j$-th lift being formed with the $j$-th power of a tamely ramified character~$\vartheta$. We prove
that the chain is injective and that the ratio $Z(1)/Z(\vartheta)$ of its Gram determinants is the
product of the corresponding ratios for $m$ single lifts in Weil's range and of explicit ratios of
values of Hecke $L$-functions. The main tool is an exchange identity between theta integrals over
two unitary groups of complementary ranks, which holds up to a constant whose absolute value is
given by Hecke $L$-values. We also construct, by lifts through two grids of hermitian spaces, a
spanning set of forms whose images under the chain are short and become integral, after every
twist, on multiplication by an explicit scalar.
\end{abstract}
\maketitle
\setcounter{tocdepth}{1}\tableofcontents

\section{Introduction}\label{C1:intro}

The main paper~\cite{P46} of this series deduces the Riemann hypothesis for Dirichlet
$L$-functions from three theorems on integral algebraic modular forms on a chain of theta lifts
between unitary groups over a CM field: the Short Forms Theorem, the Spanning Theorem and the
Covolume Theorem. The third, {\cite[Theorem~14.15]{P46}}, is a lower bound of arithmetic nature: a
family of linearly independent forms of almost full rank cannot be at the same time short and
integral after all twists. It is deduced in {\cite[\S14.10]{P46}} from a slope inequality
({\cite[Theorem~14.4]{P46}}) and from propositions of two kinds, those that concern a single
theta lift, which are proved in~\cite{P54}, \cite{P55} and~\cite{P56}, and those that concern
the chain of lifts as a whole. We prove the latter here: the injectivity of the chain
({\cite[Proposition~14.35]{P46}}), the factorisation of its determinant ({\cite[Proposition~14.37]{P46}})
and the existence of auxiliary forms ({\cite[Proposition~14.38]{P46}}).

Let us recall the objects, in the notation of {\cite[\S\S14.1--14.3]{P46}}. The
field $E$ is a CM field of degree $2f$ with maximal totally real subfield $E^+$, and
$Y_0,Y_1,\dots,Y_m$ are hermitian spaces over $E$ with unimodular lattices; $Y_j$ has rank $r+jt$,
where $t=5r$, the space $Y_0$ has signature $(k+\varrho,\varrho+1)$ at every infinite place, with
$m=2\varrho+2$, and $Y_m$ is positive definite. Consecutive unitary groups form a dual pair. The
space $\mathcal A_0$ consists of the automorphic forms on $\mathrm U(Y_0)$ of a fixed level at the
chosen places that are reproduced by the matrix coefficient of a fixed discrete series, and the
space $\mathcal T\subseteq\mathcal A_0$ of target forms is its part of prescribed type under the
finite group $\overline G$ of the level. A \emph{twist} is a unitary automorphic character
$\vartheta$ of $\mathrm U(1)$ over $E^+$, trivial at the infinite places and tamely ramified at
a finite set $Q_\vartheta$ of primes of $E$ disjoint from its conjugate; it is \emph{$m$-fold}
if the first $m$ powers of its tame parts are non-trivial. For such $\vartheta$ the \emph{lower
chain} is the composite
\[
\iota_\vartheta=\Theta^Y_{m-1,\vartheta}\circ\dots\circ\Theta^Y_{0,\vartheta}\colon
\mathcal A_0\longrightarrow\mathcal A_m(\vartheta)
\]
of the theta lifts for the pairs $(\mathrm U(Y_j),\mathrm U(Y_{j+1}))$, the lift from
$\mathrm U(Y_j)$ being formed with the character $\mu^0_{j+1}\vartheta^{\,j+1}$; its values are
algebraic modular forms on the definite group $\mathrm U(Y_m)$. The \emph{chain determinant} is
$Z(\vartheta)=\det(\Theta_\vartheta^*\Theta_\vartheta)$, where $\Theta_\vartheta$ is the restriction
of $\iota_\vartheta$ to~$\mathcal T$. In the slope inequality the integrality of a form is tested,
at each prime~$\ell$, after a twist that depends on~$\ell$, and the price is the height of the
family of numbers $Z(1)/Z(\vartheta)$. The Covolume Theorem therefore requires that $\iota_\vartheta$
be injective and that $Z(1)/Z(\vartheta)$ be expressed through quantities whose valuations can be
controlled.

Our first result is that $\iota_\vartheta$ is injective on $\mathcal A_0$ for every $m$-fold
twist ({\cite[Proposition~14.35]{P46}}). The second compares the chain with $m$ single lifts. For
$1\le j\le m$ let $W_j$ be the hermitian space of rank $r$ of {\cite[\S14.3]{P46}}, whose
unitary group is an inner form of $\mathrm U(Y_0)$, let $V_j$ be positive definite of rank
$r+t$, and let $Z_j(\vartheta)$ be the Gram determinant of the theta lift from the good forms on
$\mathrm U(W_j)$ to $\mathrm U(V_j)$ formed with the character $\mu_j^0\vartheta$. Then there is
an integer $e$ such that, for every $m$-fold twist $\vartheta$,
\[
\frac{Z(1)}{Z(\vartheta)}=\Bigl(\bigl(\operatorname{Nm}Q_\vartheta\bigr)^{e/2}\prod_{0<i<j\le m}
\rho_{ij}(\vartheta)\Bigr)^{-d_{\mathcal T}}\cdot\prod_{j=1}^m\frac{Z_j(1)}{Z_j(\vartheta^{\,j})},
\qquad
\rho_{ij}(\vartheta)=\prod_{z=1}^t\frac{|L(z,(\nu_{j-i}\vartheta^{\,j-i})^{-1})|}{|L(z,\nu_{j-i}^{-1})|},
\]
where $d_{\mathcal T}=\dim\mathcal T$ and the $L(s,\nu^{-1})$ are Hecke $L$-functions of~$E$
attached to the characters $\nu_n$ of \eqref{pc:eq-rho}, none of whose values in the formula
vanishes; this is {\cite[Proposition~14.37]{P46}}. The determinants $Z_j$ of the single lifts are
the subject of~\cite{P54}.

Both results follow from a statement on the composition of theta lifts in two different orders.
Let $W$, $V$, $V'$ be hermitian spaces of ranks $N=a+b$, $a$ and $b$ with $a=b+2e$, $e\ge0$,
whose numbers of negative signs $q$, $q_V$, $q_{V'}$ satisfy $q=q_V+q_{V'}$ and $2q+2\le N$, and
let $\nu$ be an automorphic character of $\mathrm U(1)$ whose weight $\mu$ satisfies
$N\le2\mu_\tau+2q$ at every infinite place~$\tau$, tamely ramified with non-trivial tame part
exactly at the primes of a modulus $\mathfrak m\ne(1)$. The \emph{exchange identity}
(Theorem~\ref{pcx:X}) asserts that
\[
\int_{[\mathrm U(V)]}\theta^{W,V}_{\alpha,\gamma}(g,h;x,\Phi)\,dh
=c\int_{[\mathrm U(V')]}\theta^{V',W}_{\gamma,\alpha'}(h',g;v,\Psi)\,dh',
\qquad
|c|^2=C\cdot L_*(\mathfrak m,\nu,e)\ne0 ,
\]
where
\[
L_*(\mathfrak m,\nu,e)=\Bigl(\prod_{w\mid\mathfrak m}q_w\Bigr)^{e^2}\prod_{z=1}^e\bigl|L^{\mathfrak m}(z,\nu)\bigr|^2 ,
\]
for all splitting characters with $\alpha\alpha'=\gamma=\nu^{-1}\upsilon^{N/2}$ and all pairs of
finite data $\Phi$, $\Psi$ modulo $\mathfrak m$ that are matched in the sense of
Definition~\ref{pcx:def-matched}. Here $q_w$ is the norm of $w$, $L^{\mathfrak m}(s,\nu)$ is the
Hecke $L$-function of the character $x\mapsto\nu(x/\bar x)$ of $E$ without the Euler factors at the
primes of $\mathfrak m$ and their conjugates, $v$ is a given archimedean vector, the archimedean
vector $x$ and the constant $C>0$ depend on $v$ but not on $\nu$ and $\mathfrak m$, and both
integrals converge absolutely. Thus the theta integral of the pair
$(W,V)$ over the larger group and that of the pair $(V',W)$ over the smaller group, as
functions on $\mathrm U(W)$, differ by a constant of known absolute value. The identity has the
shape of a Siegel--Weil formula on both sides at once, but neither side is in a range in which
such a formula is available, and we do not derive it from one. The proof, whose plan is given in
\S\ref{pcx:ss-arch}, has four steps. For certain standard data at the primes of $\mathfrak m$ the
identity is first proved for a larger triple, in which $W$ is enlarged by hyperbolic planes, by
unfolding along an isotropic frame, with archimedean vectors that have prescribed zero terms
(Propositions~\ref{pcx:P-top} and~\ref{pcx:P-slice}, proved in \S\ref{pcu:sec}). The planes are
then removed one at a time (Proposition~\ref{pcx:P-descent}, proved in \S\ref{pct:sec}): the
constant term of a theta integral along the isotropic line of a plane is computed by an integral
of Tate's type, and each plane contributes one value $|L^{\mathfrak m}(z,\nu)|$ at an integer
$z\ge1$ and a power of $\prod_{w\mid\mathfrak m}q_w$; no functional equation is used.
Equivariance carries the identity from standard data to all matched pairs
(Proposition~\ref{pcx:P-match}). What remains is one constant, the ratio of the two sides in equal
rank, which is shown to be the same non-zero number for all $\nu$ (Proposition~\ref{pcx:P-equal});
its non-vanishing uses the fact that the Fourier coefficients of a theta integral along the centre
of a unipotent radical vanish at the indices that are not negative definite (\S\ref{pcn:sec}).

The lifts of the chain are embedded in a \emph{triangle}: a family of spaces of forms
$\mathcal F_{i,n}$, $i+n\le m$, on unitary groups of rank $r+nt$, with lifts
$a_{i,n}\colon\mathcal F_{i,n}\to\mathcal F_{i,n+1}$ and transposed lifts
$d_{i,n}\colon\mathcal F_{i,n+1}\to\mathcal F_{i+1,n}$, whose row $i=0$ is the lower chain.
Theorem~\ref{pc:tri} states that each square commutes up to a constant,
\[
w_n\,a_{i+1,n}\,d_{i,n}=\frac{\kappa^\sharp_{i,n}}{c_{i,n}(\vartheta)}\,w_{n+2}\,d_{i,n+1}\,a_{i,n+1},
\qquad
|c_{i,n}(\vartheta)|^2=C_{i,n}\,L_*\bigl(\mathfrak m_\vartheta,(\nu_{n+1}\vartheta^{\,n+1})^{-1},t\bigr),
\]
where $\mathfrak m_\vartheta=(\pi_M)\prod_{w\in Q_\vartheta}w$, the $w_n$ are the orders of
general linear groups over $\mathcal O_E/\mathfrak m_\vartheta$ and $\kappa^\sharp_{i,n}$,
$C_{i,n}>0$ do not depend on $\vartheta$; and that an adjoint law holds whose constants have
absolute value a power of $\operatorname{Nm}Q_\vartheta$ times a constant. Each square is an instance of the exchange
identity with $e=t$, and the adjoint law comes from a second family of squares with $e=0$
(Theorem~\ref{ptw:RS}). Two points are delicate. The archimedean vector $x$ produced by the
exchange identity is not explicit; it is compared with the polynomial of the lifts by a
statement of multiplicity one for the archimedean types at the nodes of the triangle
(Theorems~\ref{ptr:square} and~\ref{ptq:H}), which is why these types and their Fock models are
studied in \S\S\ref{ptr:sec} and~\ref{ptq:sec}. And the factor $\kappa^\sharp_{i,n}$, which
this argument does not compute, must not vanish. Here the structure of the triangle is used: only
the first lift of each row is in Weil's range, where the inner product formula
(Proposition~\ref{pc:rallis}), a consequence of Weil's Siegel formula~\cite{WeilSiegel}, expresses
its Gram form through a Hecke operator given by a series that is close to the identity
(Lemma~\ref{pc:series}); so the first lifts are injective, and the injectivity of all chains, with
the non-vanishing of all the constants of the squares, follows from the commutation relations by
linear algebra (Proposition~\ref{pc:tri-inj}). The only non-vanishing of $L$-values that enters is
that of $L^{\mathfrak m}(z,\nu)$ at the integers $z\ge1$, which makes the constants of the exchange
identity non-zero (Lemma~\ref{pcx:Lpos}); for $z\ge2$ it follows from the Euler product and for
$z=1$ it is Hecke's theorem. No value at the centre or inside the critical strip occurs. The same
linear algebra, with the adjoint law, writes the Gram form of the chain as a scalar times a
product of the $m$ operators of the first lifts (Proposition~\ref{pc:tri-gram}). These operators also govern the single lifts
$W_j\to V_j$, on a different inner form; the two determinants agree by the equality of
multiplicities of~\cite{P51} ({\cite[Proposition~14.36]{P46}}, used in Lemma~\ref{pc:detcompare}),
and this gives the factorisation.

The third result supplies the forms by which a family of almost full rank is completed to a basis.
{\cite[Proposition~14.38]{P46}} gives a function $\gamma$ on the $m$-fold twists with $\gamma(1)=1$, whose
absolute value is a power of $\operatorname{Nm}Q_\vartheta$ times a product of $m^2$ ratios of
Hecke $L$-values of the same kind as the $\rho_{ij}$, and a spanning set of $\mathcal T$ consisting
of forms $y$ such that $\gamma(\vartheta)\Theta_\vartheta y$ is integral at every prime not
dividing $\operatorname{Nm}Q_\vartheta$, for every $m$-fold $\vartheta$, and such that all
conjugates of $\Theta_1y$ have squared norm at most $\exp(2x_1^2b^2)$. It rests on
Theorem~\ref{pc:auxcore}: there is a surjective map onto $\mathcal A_0$ from a space of algebraic
modular forms of deep level on a larger definite unitary group, under which integral forms remain
integral after every twist, up to the scalar $\gamma(\vartheta)$, and short forms remain short.
The map is built from transposed lifts down a second, untwisted chain, followed by restriction to
$\mathrm U(Y_0)$. To follow a form under
$\iota_\vartheta$ we place both chains in two \emph{grids} of hermitian spaces and move a path
through the grids square by square; each square is again an exchange identity, which is the
origin of the ratios in $|\gamma(\vartheta)|$, and on the diagonal the path turns by the reversal
of the theta kernel of a line. On the last path all groups are definite and the lifts are given
by Taylor coefficients of theta functions, whose denominators and sizes are bounded directly
(\S\S\ref{pgz:sec} and~\ref{pgt:sec}). The surjectivity is the delicate point: it requires
that no form of $\mathcal A_0$ be orthogonal to all restrictions of forms of deep level on the
larger group, which we prove with Poincar\'e series of deep level (Theorem~\ref{pgq:V}), together
with the injectivity of the untwisted chain of the second grid (Theorem~\ref{pgy:I}).

Finally we prove two statements used elsewhere in the series. Proposition~\ref{pc:C} carries out
in detail the deduction of the Covolume Theorem from its inputs. Theorem~\ref{pcj:CL} describes
the action of the automorphisms of $\mathbb C$ over $K$ on the norms of the theta lifts of
characters from $\mathrm U(1)$; it is used in~\cite{P54} for the archimedean size of the
conjugates of Hecke $L$-values.

The notation and the data are those of~\cite{P46}. We take from~\cite{P49} the properties of
the Weil representation and of its models, and the results on discrete series and on the
archimedean theta correspondence; from~\cite{P51} the equality of multiplicities and estimates
for coefficients of discrete series; and from~\cite{P55} the Hecke polynomial of a theta lift
({\cite[Proposition~6.7]{P55}}). This paper and~\cite{P54} use each other, in the following order.
From~\cite{P54} we take the archimedean facts on the single lifts $W_j\to V_j$
({\cite[Lemmas~4.3 and~4.7]{P54}}, {\cite[Proposition~4.8]{P54}}), the standing facts on one lift ({\cite[Lemma~2.2]{P54}}), a bound for the kernels of
lifts between definite groups ({\cite[Lemma~4.10]{P54}}) and the growth of the Siegel Eisenstein series
({\cite[Lemma~5.41]{P54}}, with an elementary comparison of heights made in the comments on {\cite[Lemma~5.40]{P54}}); these are proved there from~\cite{P49} and from~\cite{WeilSiegel},
without the exchange identity and without the inner product formula. With them we prove the
inner product formula and the exchange identity, which~\cite{P54} uses for the three exchanges
and for the norm of the lift of a character ({\cite[Proposition~3.4]{P54}}). That norm formula, with three lemmas of~\cite{P54} on the same lifts ({\cite[Lemmas~2.1, 3.3 and~7.27]{P54}}), is
used here only for the law of conjugation (Theorem~\ref{pcj:CL}), and the latter is used
in~\cite{P54} only for the archimedean size of the untwisted norms ({\cite[Theorem~7.43]{P54}}), on
which nothing in this paper depends except the final deduction of the Covolume Theorem from all its inputs (Proposition~\ref{pc:C}). The statements quoted are reproduced in
Appendix~\ref{C1:quoted}. The published results on
which we rely most are Weil's Siegel formula in its range of convergence~\cite{WeilSiegel}, the
description of the joint harmonics by Kashiwara and Vergne and by
Howe~\cite{HoweTheta,KashiwaraVergne}, Kudla's splittings~\cite{KudlaSplitting}, the
classification of hermitian lattices and forms by Jacobowitz and
Landherr~\cite{Jacobowitz,Landherr}, reduction theory~\cite{BorelFin,Godement63}, and the theory
of the discrete series of Harish-Chandra and Schmid~\cite{Schmid75,HC66}.

The paper is organised as follows. Section~\ref{pc:sec} contains the deduction of the Covolume
Theorem, the linear algebra of triangles, the inner product formula, and the proofs of
{\cite[Propositions~14.35, 14.37 and~14.38]{P46}} from the theorems of the later
sections. Section~\ref{C1:g2} proves the exchange identity. Section~\ref{C1:g3} establishes the
triangle: the archimedean types at its nodes, positivity in the Fock model, the adjoint law and
the reversed squares. Sections~\ref{C1:g4}--\ref{C1:g7} treat the two grids: the lifts through
the second grid and the proof of Theorem~\ref{pc:auxcore} (Section~\ref{C1:g4}), the grids, the
square law and the first path (Section~\ref{C1:g5}), the denominators, the vanishing at deep
level and the sizes on the last path (Section~\ref{C1:g6}), and the kernel forms, the restriction
pairing and the coefficient lemmas (Section~\ref{C1:g7}). Section~\ref{pcj:sec} proves the law of
conjugation of the lifts of characters.


% ===== C-determinant.tex =====
% proofs/C-determinant.tex
% The determinant part of Theorem C: deduction of Theorem C from its inputs, the triangle of lifts,
% the factorisation of the chain determinant (au:factor), injectivity (au:inj), auxiliary forms
% (au:aux). Labels of this file start with pc:. Only amsmath/amssymb/amsthm are used.
\providecommand{\auO}{\mathcal O}
\providecommand{\auI}{\mathcal I}
\providecommand{\auQlb}{\overline{\Q}_\ell}
\providecommand{\auAut}{\operatorname{Aut}}
\providecommand{\auNm}{\operatorname{Nm}}
\providecommand{\auU}{\mathrm U}
\providecommand{\auGL}{\mathrm{GL}}
\providecommand{\auHom}{\operatorname{Hom}}
\providecommand{\auEnd}{\operatorname{End}}
\providecommand{\auA}{\mathbb A}
\providecommand{\auF}{\mathbb F}
\providecommand{\auPl}{\Phi}
\providecommand{\auT}{\mathcal T}
\providecommand{\auAz}{\mathcal A_0}
\providecommand{\auAm}{\mathcal A_m}
\providecommand{\auG}{\overline G}
\providecommand{\auLat}{\mathcal L}
\providecommand{\audT}{d_{\auT}}
\providecommand{\auhmain}{h_{\mathrm{main}}}
\providecommand{\auhdir}{h_{\mathrm{dir}}}
\providecommand{\aueu}{\varepsilon_{\mathrm u}}
\providecommand{\aubud}{\mathrm{bud}}
\providecommand{\aubar}[1]{\overline{#1}}
\providecommand{\auf}[1]{\texttt{#1}}
\providecommand{\psHS}{\mathcal H}
\providecommand{\psrho}{\varrho}
\providecommand{\psvac}{\mathbf v}
\providecommand{\psZ}{Z}
% macros of this file
\providecommand{\pcF}{\mathcal F}
\providecommand{\pcG}{\mathsf G}
\providecommand{\pcm}{\mathfrak m}
\providecommand{\pcb}{\mathfrak b}
\providecommand{\pcR}{R}
\providecommand{\pcH}{\mathcal H^{M}}
\providecommand{\pcRal}{\mathcal R}
\providecommand{\pcmc}{\operatorname{mc}}
\providecommand{\pcvol}{\operatorname{vol}}
\providecommand{\pctr}{\operatorname{tr}}
\providecommand{\pcMat}{\mathrm M}
\providecommand{\pcW}{\mathbb W}
\providecommand{\pcFF}{\mathfrak F}
\providecommand{\pop}{\mathsf p}

\section{The determinant of the chain of lifts and the Covolume Theorem}\label{pc:sec}

This section treats the part of the Covolume Theorem ({\cite[Theorem~14.15]{P46}}) that does not concern a single lift
$W_j\to V_j$: the deduction of the theorem from its inputs, the injectivity of the lower chain
({\cite[Proposition~14.35]{P46}}), the factorisation of the chain determinant
({\cite[Proposition~14.37]{P46}}) and the auxiliary forms ({\cite[Proposition~14.38]{P46}}). The notation is
that of {\cite[\S\S14.1--14.3]{P46}} and of {\cite[\S5.1.1]{P49}}; inner products are
conjugate-linear in the first variable.

\emph{Organisation.} The section has two layers.
\begin{itemize}
\item[(I)] \emph{Proofs given here.} {\cite[Theorem~14.15]{P46}} from {\cite[Propositions~14.35, 14.37, 14.38, 14.40 and~14.43]{P46}} (Proposition~\ref{pc:C}); the
linear algebra of a \emph{triangle} of maps, which gives injectivity and the Gram form of a chain
from commutation relations in squares (Propositions~\ref{pc:tri-inj}, \ref{pc:tri-gram}); the
series of Rallis' formula, its convergence and its injectivity
(Lemma~\ref{pc:series}); the comparison of its determinants on inner forms
(Lemma~\ref{pc:detcompare}); the inner product formula for a theta lift with \emph{sharp} data in
Weil's range (Proposition~\ref{pc:rallis}), from Weil's form of the Siegel formula and the facts
(W1)--(W5), (W7) ({\cite[Lemma~5.2]{P49}} and~Lemma~\ref{pc:SW}); the finite algebra of matched functions
(Lemmas~\ref{pc:pointlaw}, \ref{pc:next}); the bookkeeping of levels and indices
(Lemmas~\ref{pc:algpet}, \ref{pc:index}); and the deductions of {\cite[Propositions~14.35, 14.37 and~14.38]{P46}} from the four statements of layer (II).
\item[(II)] \emph{Statements proved in later sections.} Lemma~\ref{pc:arch} (archimedean types;
Proposition~\ref{ptv:arch}), Theorem~\ref{pc:tri} (the lifts form a triangle whose squares commute
up to constants of known modulus; \S\S\ref{ptr:sec}--\ref{ptw:sec}), Theorem~\ref{pc:X} (the
exchange identity, on which Theorem~\ref{pc:tri} rests; Theorem~\ref{pcx:X}, \S\S\ref{pcx:sec},
\ref{pcu:sec}, \ref{pct:sec}, \ref{pcn:sec}) and Theorem~\ref{pc:auxcore} (integral and short lifts
through the second grid; \S\S\ref{pgr:sec}--\ref{pgm:sec}). In addition {\cite[Proposition~14.36]{P46}} is
needed in the form $(\mathrm M^+)$ of \S\ref{pc:ss-rallis}, which is its statement (for the full
spherical Hecke algebra away from $M$).
\end{itemize}
\S\ref{pc:ss-remains} lists where the inputs are proved and the published results on which they
rest.

\subsection{The Covolume Theorem from its inputs}\label{pc:ss-thmC}

We recall the rules for heights of {\cite[\S14.2]{P46}} in the form used.

\begin{lemma}\label{pc:heights}
\begin{enumerate}
\item[(a)] If the families $y,y'$ have heights at most $c,c'$, then $yy'$ has height at most
$c+c'$, and $y^n$ has height at most $nc$ for every integer $n\ge0$. A family of height at most $c$
has height at most $c''$ for every $c''\ge c$.
\item[(b)] Let $N_\iota$ be positive integers with $\ell(\iota)\nmid N_\iota$ and let
$e\in\tfrac12\Z$. If $y$ has height at most $c$, so has $\iota\mapsto N_\iota^{\,e}y_\iota$.
\item[(c)] Let $\theta$ be a regular family of $m$-fold twists
\textup{({\cite[Definition~14.39]{P46}})}, let $\rho_1,\dots,\rho_n$ be ramified ratios in the
sense of {\cite[\S14.10]{P46}}, and let $e\in\Z$. Then the family
\[
\iota\longmapsto\Bigl((\auNm Q_{\theta_\iota})^{e/2}\prod_{a=1}^n\rho_a(\theta_\iota)\Bigr)^{-1}
\]
is well defined and has height at most $n\,t\,x_1$.
\end{enumerate}
\end{lemma}
\begin{proof}
(a) If $(H,\alpha,h)$ and $(H',\alpha',h')$ are as in {\cite[Definition~14.1]{P46}} for $y$ and
$y'$, then $(HH',\alpha+\alpha',h+h')$ serves for $yy'$; the last assertion is clear from the
definition. (b) The positive real number $N_\iota^{\,e}$ is an algebraic number whose square is an
integral power of $N_\iota$; as $\ell(\iota)\nmid N_\iota$, $|N_\iota^{\,e}|_\iota=1$, and the
triple $(H,\alpha,h)$ of $y$ serves. (c) By definition $\rho_a(\vartheta)=\prod_{i=1}^{t'}
|L(z_i,(\nu_i\vartheta^{\pm n_i})^{-1})|/|L(z_i,\nu_i^{-1})|$ with $t'\le t$, $1\le n_i\le m$ and
ramified critical pairs $(\nu_i,z_i)$. By (R2) each family $\iota\mapsto|L(z_i,\nu_i^{-1})|/
|L(z_i,(\nu_i\theta_\iota^{\pm n_i})^{-1})|$ is defined and has height at most $x_1$. By (a),
$\iota\mapsto\rho_a(\theta_\iota)^{-1}$ has height at most $t'x_1\le tx_1$, and the product over
$a$ has height at most $ntx_1$. By (R1), $\ell(\iota)\nmid\auNm Q_{\theta_\iota}$, and (b) applies
with the exponent $-e/2$.
\end{proof}

\begin{proposition}[the Covolume Theorem from its inputs]\label{pc:C}
Assume {\cite[Propositions~14.35, 14.37, 14.38, 14.40 and~14.43]{P46}}. Then {\cite[Theorem~14.15]{P46}} holds: if $y_1,\dots,y_n\in\auT$ are linearly independent
with $n\ge(1-4/b^2)\audT$, of index at most $\delta$ and covolume at most $c$ in $\auLat_D$, then
$c+2\delta\ge\audT(\auhmain-3fa(\log b)^4)$.
\end{proposition}
\begin{proof}
\emph{Step 0: the lattice.} By {\cite[Proposition~14.35]{P46}} the maps $\Theta_\vartheta=\iota_\vartheta|
_{\auT}$ are injective for all $m$-fold twists, and {\cite[(39)]{P46}} holds because the levels
$J_\vartheta$ are free ({\cite[Lemma~14.5]{P46}}); so $\auLat_D$ is a twisted lattice in the sense of
{\cite[Definition~14.3]{P46}}, and {\cite[Theorem~14.4]{P46}} applies to it.

\emph{Step 1: completion to a basis.} Let $e'$, $\rho'_{ij}$, $\gamma$ and the spanning set
$\mathfrak A\subseteq\auT$ be as in {\cite[Proposition~14.38]{P46}}. As $y_1,\dots,y_n$ are independent
and $\mathfrak A$ spans $\auT$, the exchange lemma of Steinitz gives $z_1,\dots,z_{n'}\in
\mathfrak A$, $n'=\audT-n$, such that $y_1,\dots,y_n,z_1,\dots,z_{n'}$ is a basis of $\auT$. By
hypothesis $n'\le4\audT/b^2$.

\emph{Step 2: the twists.} Let $\theta$ be a regular family ({\cite[Proposition~14.40]{P46}}). By (R1)
the twist $\theta_\iota$ is allowed at $\ell(\iota)$ for every $\iota\in\auI$.

\emph{Step 3: the hypotheses of {\cite[Theorem~14.4]{P46}}.} (i) is the hypothesis on $y$. (ii)
For every $i$, every $\iota$ and every $m$-fold twist $\xi$ with $\ell(\iota)\nmid\auNm Q_\xi$, the
form $\gamma(\xi)\Theta_\xi z_i$ is integral at $\iota$, by {\cite[Proposition~14.38]{P46}}; and
$\gamma(\xi)\ne0$. (iii) holds with $b_{\mathrm{short}}=2x_1^2b^2$, by the same proposition. (iv)
First, $|\gamma(\theta_\iota)|^{-2}$ is the square of the inverse of $(\auNm Q_{\theta_\iota})^{e'/2}$
times $m^2$ ramified ratios; by Lemma~\ref{pc:heights}(a),(c) it has height at most
$2m^2tx_1$, and $2m^2tx_1\le b$ because $m^2\le4x$, $t=5r\le15x^5$ and $x_1^{10^6}\le b$; so
$b_{\mathrm{den}}=b$ serves. Secondly, by {\cite[Proposition~14.37]{P46}},
\[
\frac{Z(1)}{Z(\theta_\iota)}=\biggl(\Bigl((\auNm Q_{\theta_\iota})^{e/2}\prod_{0<i<j\le m}\rho_{ij}
(\theta_\iota)\Bigr)^{-1}\biggr)^{\audT}\cdot\prod_{j=1}^m\frac{Z_j(1)}{Z_j(\theta_\iota^{\,j})} .
\]
The first factor is the $\audT$-th power of a family which, by Lemma~\ref{pc:heights}(c) with
fewer than $m^2$ ratios, has height at most $m^2tx_1\le b$. By {\cite[Proposition~14.43]{P46}} each of the
$m$ remaining factors has height at most $\audT(fa(\log b)^4/m-k\auhdir)$. By
Lemma~\ref{pc:heights}(a) and $m\,k\,\auhdir=\auhmain$
({\cite[Definition~14.8(c)]{P46}}), the family $Z(1)/Z(\theta_\iota)$ has height at most
\[
b_{\mathrm{vol}}=\audT\bigl(b+fa(\log b)^4-\auhmain\bigr).
\]

\emph{Step 4: conclusion.} {\cite[Theorem~14.4]{P46}} gives
\[
0\le c+2\delta+n'(2x_1^2b^2+b)+\audT\bigl(b+fa(\log b)^4-\auhmain\bigr),
\]
and $n'(2x_1^2b^2+b)\le4\audT b^{-2}(2x_1^2b^2+b)=4\audT(2x_1^2+b^{-1})\le\audT b$. Hence
$\audT(\auhmain-fa(\log b)^4-2b)\le c+2\delta$, and $2b\le2fa(\log b)^4$.
\end{proof}

Thus, after {\cite[Proposition~14.43]{P46}} (proved in {\cite[\S14.10]{P46}} from
{\cite[Propositions~14.36, 14.41 and~14.42]{P46}}), {\cite[Theorem~14.15]{P46}} rests on
{\cite[Propositions~14.35, 14.37, 14.38 and~14.40]{P46}}. The rest of this
section concerns the first three.

\subsection{Linear algebra: volumes, moduli, triangles}\label{pc:ss-la}

\begin{lemma}[Volumes]\label{pc:vol}
Let $u\colon U\to V$ be a linear map of finite-dimensional inner product spaces, $S\in\auEnd(U)$
and $a\in\C$ with $\langle ux,uy\rangle=a\langle x,Sy\rangle$ for all $x,y\in U$. Then
$\det(u^*u)=a^{\dim U}\det S$.
\end{lemma}
\begin{proof}
The hypothesis says $u^*u=aS$.
\end{proof}

\begin{definition}\label{pc:modulus}
Let $c$ be a complex function and $f$ a real function on the set of $m$-fold twists. We say that
$c$ \emph{has modulus $f$ up to norm powers} if there are a real $C\ge0$ and $k\in\Z$ with
$|c(\vartheta)|^2=C\,(\auNm Q_\vartheta)^k f(\vartheta)^2$ for all $\vartheta$.
\end{definition}

\begin{lemma}\label{pc:mod}
\begin{enumerate}
\item[(a)] Constants have modulus $1$; if $c,c'$ have moduli $f,f'$ up to norm powers, then $cc'$
has modulus $ff'$; if moreover $c$ has no zero, $c^{-1}$ has modulus $f^{-1}$ and $\aubar c$ has
modulus $f$.
\item[(b)] If $c$ has modulus $f$ up to norm powers, $f\ge0$, $f(1)=1$ and $c(1)\ne0$, then there
is $e\in\Z$ with $|c(\vartheta)/c(1)|=(\auNm Q_\vartheta)^{e/2}f(\vartheta)$ for all $\vartheta$.
\end{enumerate}
\end{lemma}
\begin{proof}
(a) Multiply the identities. (b) At $\vartheta=1$, $Q=\emptyset$ and $\auNm Q=1$, so $C=|c(1)|^2>0$;
divide and take the non-negative square root; $e=k$.
\end{proof}

\begin{definition}[Triangle]\label{pc:triangle}
Let $m\ge1$, let $I$ be a set, $V$ a complex inner product space, and for $i,n\ge0$, $i+n\le m$,
let $\pcF_{i,n}$ be complex inner product spaces (not necessarily complete). A \emph{triangle} of
size $m$ over $I$ consists of linear isometries $j_i\colon V\to\pcF_{i,0}$ $(0\le i\le m)$ and, for
every $\vartheta\in I$, of numbers $w_n=w_n^\vartheta>0$ $(0\le n\le m)$ and linear maps
\[
a_{i,n}=a_{i,n}^\vartheta\colon\pcF_{i,n}\to\pcF_{i,n+1},\qquad
d_{i,n}=d_{i,n}^\vartheta\colon\pcF_{i,n+1}\to\pcF_{i+1,n}\qquad(i+n<m)
\]
such that:
\begin{enumerate}
\item[(T1)] \emph{(squares)} for $i+n+2\le m$ there is $\sigma_{i,n}\colon I\to\C$ with
$w_n\,a_{i+1,n}\circ d_{i,n}=\sigma_{i,n}(\vartheta)\,w_{n+2}\,d_{i,n+1}\circ a_{i,n+1}$ on
$\pcF_{i,n+1}$;
\item[(T2)] \emph{(first adjoint)} for $i<m$, $y\in V$ and $z\in\pcF_{i,1}$,
$\langle z,a_{i,0}j_iy\rangle=\langle d_{i,0}z,j_{i+1}y\rangle$;
\item[(T3)] for $i<m$, $d_{i,0}a_{i,0}j_i(V)\subseteq j_{i+1}(V)$, and $d_{i,0}a_{i,0}j_i$ is
injective.
\end{enumerate}
The triangle \emph{has the adjoint law} if moreover
\begin{enumerate}
\item[(T4)] for $i+n+2\le m$ there is $\kappa_{i,n}\colon I\to\C$ such that every
$z\in\pcF_{i,n+2}$ has an $x\in\pcF_{i,n+1}$ with $\langle a_{i,n+1}v,z\rangle=\langle v,x\rangle$
for all $v\in\pcF_{i,n+1}$ and
$\langle d_{i,n}x,v'\rangle=\kappa_{i,n}(\vartheta)\langle d_{i,n+1}z,a_{i+1,n}v'\rangle$ for all
$v'\in\pcF_{i+1,n}$.
\end{enumerate}
The \emph{chains} are $X_{i,n}=X_{i,n}^\vartheta=(w_{n-1}a_{i,n-1})\circ\dots\circ(w_0a_{i,0})\colon
\pcF_{i,0}\to\pcF_{i,n}$, with $X_{i,0}=1$.
\end{definition}

In the application $\pcF_{i,n}$ is a space of automorphic forms on the unitary group of a space
with a stem and the blocks $i+1,\dots,i+n$; $a$ adds the last block, $d$ removes the first one, and
(T1) says that these two operations commute up to a constant.

\begin{proposition}[Chains are injective]\label{pc:tri-inj}
In a triangle of size $m$, for $i+n<m$ and all $\vartheta$,
\begin{equation}\label{pc:eq-tau}
X_{i+1,n}\circ d_{i,0}\circ a_{i,0}=\tau_{i,n}\;d_{i,n}\circ X_{i,n+1},\qquad
\tau_{i,n}=\frac{w_{n+1}}{w_0w_1}\prod_{l<n}\sigma_{i,l}(\vartheta).
\end{equation}
For $i+n\le m$ the map $X_{i,n}\circ j_i\colon V\to\pcF_{i,n}$ is injective. If $V\ne0$, then
$\sigma_{i,n}(\vartheta)\ne0$ for all $i+n+2\le m$ and all $\vartheta$.
\end{proposition}
\begin{proof}
\eqref{pc:eq-tau} for $n=0$ reads $d_{i,0}a_{i,0}=w_0^{-1}d_{i,0}(w_0a_{i,0})$. If it holds for
$n$ and $i+n+2\le m$, then by (T1)
\[
X_{i+1,n+1}d_{i,0}a_{i,0}=w_na_{i+1,n}\,\tau_{i,n}d_{i,n}X_{i,n+1}=\tau_{i,n}\sigma_{i,n}w_{n+2}\,
d_{i,n+1}a_{i,n+1}X_{i,n+1}=\tau_{i,n}\sigma_{i,n}\frac{w_{n+2}}{w_{n+1}}\,d_{i,n+1}X_{i,n+2},
\]
which is \eqref{pc:eq-tau} for $n+1$. Injectivity is proved by induction on $n$, for all $i$ with
$i+n\le m$; the case $n=0$ is clear. Let $i+n+1\le m$ and $X_{i,n+1}j_iy=0$. By (T3) there is
$y'\in V$ with $j_{i+1}y'=d_{i,0}a_{i,0}j_iy$, and by \eqref{pc:eq-tau}
$X_{i+1,n}j_{i+1}y'=\tau_{i,n}d_{i,n}X_{i,n+1}j_iy=0$. By induction $y'=0$, and $y=0$ by (T3).
Finally let $V\ne0$ and suppose $\sigma_{i,n}(\vartheta)=0$ with $i+n+2\le m$. Then
$\tau_{i,n+1}=0$, so $X_{i+1,n+1}d_{i,0}a_{i,0}j_iy=0$ for all $y\in V$; as $(i+1)+(n+1)\le m$, the
injectivity just proved and (T3) give $y=0$ for all $y\in V$, a contradiction.
\end{proof}

\begin{proposition}[The Gram form of a chain]\label{pc:tri-gram}
Let a triangle of size $m$ over $I$ have the adjoint law, let $V\ne0$, and put
$\beta_{i,n}=1/\aubar{\sigma_{i,n}}$. Suppose that for $i<m$ and $\vartheta\in I$ there are
$S_i=S_i^\vartheta\in\auEnd(V)$ and $C_i(\vartheta)\in\C$ with
\begin{equation}\label{pc:eq-first}
w_0\,\langle a_{i,0}j_ix,a_{i,0}j_iy\rangle=C_i(\vartheta)\,\langle x,S_iy\rangle\qquad(x,y\in V).
\end{equation}
Then for all $\vartheta$ and $x,y\in V$
\[
\langle X_{0,m}j_0x,X_{0,m}j_0y\rangle=\frac{w_0w_1^{\,m}}{w_m}\ \prod_{i<m}C_i(\vartheta)\
\prod_{\substack{i,n\ge0\\ i+n+2\le m}}\kappa_{i,n}(\vartheta)\beta_{i,n}(\vartheta)\ \cdot\
\bigl\langle x,S_0S_1\cdots S_{m-1}\,y\bigr\rangle .
\]
\end{proposition}
\begin{proof}
\emph{Step 1.} For $i+n<m$ put $\kappa_i^{(n)}=\prod_{l<n}\kappa_{i,l}$. Then for all
$z\in\pcF_{i,n+1}$ and $y\in V$
\begin{equation}\label{pc:eq-chainadj}
\langle z,a_{i,n}X_{i,n}j_iy\rangle=\kappa_i^{(n)}\,\langle d_{i,n}z,X_{i+1,n}j_{i+1}y\rangle .
\end{equation}
For $n=0$ this is (T2). Let it hold for $n$, let $i+n+2\le m$ and $z\in\pcF_{i,n+2}$, and let $x$
be as in (T4). Since $w_n>0$,
\begin{align*}
\langle z,a_{i,n+1}X_{i,n+1}j_iy\rangle&=\langle x,X_{i,n+1}j_iy\rangle=w_n\langle x,a_{i,n}X_{i,n}
j_iy\rangle=w_n\kappa_i^{(n)}\langle d_{i,n}x,X_{i+1,n}j_{i+1}y\rangle\\
&=w_n\kappa_i^{(n)}\kappa_{i,n}\langle d_{i,n+1}z,a_{i+1,n}X_{i+1,n}j_{i+1}y\rangle=\kappa_i^{(n+1)}
\langle d_{i,n+1}z,X_{i+1,n+1}j_{i+1}y\rangle .
\end{align*}

\emph{Step 2.} Let $i+n<m$, $x,y\in V$, and let $x'\in V$ be the vector with
$j_{i+1}x'=d_{i,0}a_{i,0}j_ix$ (T3). By Proposition~\ref{pc:tri-inj}, $\tau_{i,n}\ne0$ and
$d_{i,n}X_{i,n+1}j_ix=\tau_{i,n}^{-1}X_{i+1,n}j_{i+1}x'$. With \eqref{pc:eq-chainadj} for
$z=X_{i,n+1}j_ix$,
\begin{align*}
\langle X_{i,n+1}j_ix,X_{i,n+1}j_iy\rangle&=w_n\langle X_{i,n+1}j_ix,a_{i,n}X_{i,n}j_iy\rangle
=w_n\kappa_i^{(n)}\langle d_{i,n}X_{i,n+1}j_ix,X_{i+1,n}j_{i+1}y\rangle\\
&=\frac{w_n\,\kappa_i^{(n)}}{\aubar{\tau_{i,n}}}\,\langle X_{i+1,n}j_{i+1}x',X_{i+1,n}j_{i+1}y\rangle,
\qquad\frac1{\aubar{\tau_{i,n}}}=\frac{w_0w_1}{w_{n+1}}\prod_{l<n}\beta_{i,l}.
\end{align*}
Moreover, by (T2) with $z=a_{i,0}j_ix$ and by \eqref{pc:eq-first}, for every $v\in V$
\[
\langle x',v\rangle=\langle d_{i,0}a_{i,0}j_ix,j_{i+1}v\rangle=\langle a_{i,0}j_ix,a_{i,0}j_iv
\rangle=\frac{C_i}{w_0}\langle x,S_iv\rangle .
\]

\emph{Step 3.} We prove by induction on $n$, for $i=m-n$, that
$\langle X_{i,n}j_ix,X_{i,n}j_iy\rangle=K_{i,n}\langle x,S_iS_{i+1}\cdots S_{m-1}y\rangle$ with
$K_{m,0}=1$ and
\[
K_{i,n+1}=\frac{w_1w_n}{w_{n+1}}\,C_i\,\kappa_i^{(n)}\prod_{l<n}\beta_{i,l}\cdot K_{i+1,n}.
\]
Indeed, by Step 2 and the induction hypothesis for $(i+1,n)$ applied to $x'$ and $y$,
$\langle X_{i,n+1}j_ix,X_{i,n+1}j_iy\rangle=w_n\kappa_i^{(n)}\frac{w_0w_1}{w_{n+1}}\prod_{l<n}
\beta_{i,l}\cdot K_{i+1,n}\langle x',S_{i+1}\cdots S_{m-1}y\rangle$, and $\langle x',v\rangle=
(C_i/w_0)\langle x,S_iv\rangle$ with $v=S_{i+1}\cdots S_{m-1}y$. Unfolding the recursion from
$(0,m)$ to $(m,0)$, the weights give $\prod_{p=0}^{m-1}w_1w_{m-p-1}/w_{m-p}=w_1^{\,m}w_0/w_m$, and
each $\kappa_{i,l}$ and each $\beta_{i,l}$ with $i+l+2\le m$ occurs exactly once.
\end{proof}


\subsection{Lifts with data modulo an ideal; the series of Rallis' formula}\label{pc:ss-data}

\subsubsection*{Conventions}
A \emph{modulus} is a non-zero ideal $\pcm$ of $\auO_E$ prime to $\aubar\pcm$; so every prime $w$
dividing $\pcm$ lies above a place $v=w\aubar w$ of $E^+$ that splits in $E$. Put
$\pcR_\pcm=\auO_E/\pcm$, and let $S(\pcm)$ be the set of places of $E^+$ below $\pcm$. For a twist
$\vartheta$ with set $Q=Q_\vartheta$ we write
\[
\pcm_\vartheta=(\pi_M)\prod_{w\in Q}w,\qquad\pcR_\vartheta=\auO_E/\pcm_\vartheta=\prod_{w\mid
\pcm_\vartheta}\auF_w ,
\]
a product of finite fields. If $\vartheta$ is $m$-fold then $Q_{\vartheta^j}=Q_\vartheta$ and
$\pcm_{\vartheta^j}=\pcm_\vartheta$ for $1\le j\le m$.

A \emph{space} is a hermitian space $S$ over $E$ with a basis (its \emph{frame}) whose
$\widehat\auO_E$-span $L_S$ is a unimodular lattice, and with the same signature at all infinite
places. Let $K_S\subset\auU(S)(\auA_{E^+,f})$ be the stabiliser of $L_S$. At $v\in S(\pcm)$ the
frame identifies $\auU(S)(E^+_v)$ with $\auGL_N(E_w)$, $N$ the rank; reduction gives a surjection
$K_S\to\auGL_N(\pcR_\pcm)$, with kernel $J_S(\pcm)$. The Haar measure on $\auU(S)(\auA_{E^+})$ is the
product of a fixed Haar measure on $\auU(S)_\infty=\prod_\tau\auU(S\otimes_{E^+,\tau}\R)$ (of total
mass $1$ if $S$ is definite) and of the measure on the finite adelic points that gives $K_S$ the
volume $1$; $[\auU(S)]=\auU(S)(E^+)\backslash\auU(S)(\auA_{E^+})$. For spaces of the same rank $N$ the
finite adelic groups, with $K_S$ and the reductions, are identified by the frames: the localisations
of the lattices are isometric at every finite place, because $E/E^+$ is unramified there and the
lattices are unimodular, so they have orthonormal bases (proof: the reduction of the lattice is a non-degenerate hermitian space over the quadratic extension of residue fields, so it has a vector $\bar x$ with $\bar h(\bar x,\bar x)\ne0$ (otherwise polarisation and the surjectivity of the residue trace give $\bar h=0$); a lift $x$ has $h(x,x)$ a unit of the base ring, hence the norm of a unit since the extension is unramified, so after scaling $h(x,x)=1$; then the lattice is $\mathcal O x\perp(x^\perp\cap\text{lattice})$ with unimodular second summand, and one inducts on the rank; see also Jacobowitz~\cite{Jacobowitz}; at a split place every lattice is
standard).

Let $S,T$ be spaces of ranks $N,N'$. For automorphic characters $\alpha,\beta$ of $\auU(1)$ put
\[
\omega_{\alpha,\beta}(g,h)=\alpha(\det g)^{-1}\beta(\det h)\,\omega_{\auHom(S,T)}(g,h)\qquad
(g\in\auU(S)(\auA),\ h\in\auU(T)(\auA));
\]
the representation $\omega_\mu$ of {\cite[\S5.1.1]{P49}} is $\omega_{\mu\upsilon_{N'},\mu\upsilon_N}$.
At $v\in S(\pcm)$ we use the model (W5): $\psHS_v(\auHom(S,T))$ is the space of Schwartz--Bruhat
functions on $\pcMat_{N'\times N}(E_w)$, matrices being written in the frames. For a function
$\Psi$ on $\pcMat_{N'\times N}(\pcR_\pcm)$ let $\Psi^\natural$ be the function on
$\prod_{w\mid\pcm}\pcMat_{N'\times N}(E_w)$ that is $\Psi(a^{-1}x\bmod\pcm)$ at $x\in
a\,\pcMat_{N'\times N}(\widehat\auO)$ and $0$ elsewhere, where $a=(a_w)$ is the idele of Definition~\ref{pgx:conv}: $a_w$ is the generator of the ideal
$\alpha_w\auO_w$ fixed in {\cite[\S5.1.1]{P49}} (there $\alpha_w$ denotes the component of a scalar idele,
which has nothing to do with the characters $\alpha,\beta$ of this section), equal to $1$ when
$\alpha_w$ is a unit (Definition~\ref{pgx:conv};
the coordinates are described in Lemma~\ref{pgx:dict}); so
$\mathbf 1^\natural$ is the vacuum of $a\auHom(L_S,L_T)$ at $w$ (Lemma~\ref{pgx:vac}). For a
vector $P$ of the Fock space of $\auHom(S,T)\otimes\R$ (for instance a polynomial)
the \emph{datum} of $(P,\Psi)$ is $[P,\Psi]=P\otimes\Psi^\natural\otimes\psvac^{\pcm}$, where
$\psvac^{\pcm}$ is the vacuum of the lattice $a\auHom(L_S,L_T)$ at the finite places outside
$S(\pcm)$, normalised to have norm $1$, and the norm at $v\in S(\pcm)$ is the $L^2$-norm for the
measure giving $a_w\pcMat(\auO_w)$ volume $1$. The \emph{kernel}, the \emph{lift} and the
\emph{transposed lift} are
\begin{gather*}
\theta_{\alpha,\beta}(g,h;P,\Psi)=\theta\bigl(\omega_{\alpha,\beta}(g,h)[P,\Psi]\bigr),\\
\Theta_{\alpha,\beta}(P,\Psi)\varphi\,(h)=\int_{[\auU(S)]}\theta_{\alpha,\beta}(g,h;P,\Psi)\,
\varphi(g)\,dg,\qquad
\Theta^{\mathsf t}_{\alpha,\beta}(P,\Psi)\varphi'\,(g)=\int_{[\auU(T)]}\theta_{\alpha,\beta}(g,h;
P,\Psi)\,\varphi'(h)\,dh .
\end{gather*}
The \emph{point datum} at a matrix $x$ is the indicator function $\delta_x$ of $\{x\}$; $\iota$
denotes the matrix $\binom10$ of the inclusion of frames. Thus, with the notation of
{\cite[\S5.1.3]{P49}}, $\Theta^Y_{j,\vartheta}=[K_j:J_j]\cdot\Theta_{\mu_{j+1}\upsilon_{r+(j+1)t},\,
\mu_{j+1}\upsilon_{r+jt}}(P_j,\delta_\iota)$ with the modulus $\pcm_\vartheta$: at a level place $w$ the
datum of {\cite[\S5.1.3]{P49}} is $\mathbf 1_{a_w(\iota+\varpi_w\auHom)}$, which is $\delta_\iota^\natural$
(Lemma~\ref{pgx:dict}(c)), including at the primes of $Q_\vartheta$ where $\alpha_w$ is not a unit.

\begin{lemma}\label{pc:level}
Let $\alpha,\beta$ be unramified outside $\pcm$ and tamely ramified at the primes of $\pcm$, with
tame parts giving characters $\alpha_\pcm,\beta_\pcm$ of $\pcR_\pcm^\times$. For $k\in K_S$,
$k'\in K_T$ with reductions $\bar k,\bar k'$ modulo $\pcm$,
\[
\omega_{\alpha,\beta}(k,k')[P,\Psi]=[P,\Psi'],\qquad\Psi'(x)=\alpha_\pcm(\det\bar k)^{-1}\beta_\pcm
(\det\bar k')\,\Psi(\bar k'^{-1}x\bar k).
\]
\end{lemma}
\begin{proof}
As in {\cite[Lemma~5.5]{P49}}: outside $S(\pcm)$ the pair $(k,k')$ stabilises the lattice and the
characters are unramified, so (W4) applies; at $v\in S(\pcm)$, (W5) gives
$\Psi^\natural(k'^{-1}xk)$ times the characters (as $a$ is a scalar, $k'^{-1}xk\in a\,\pcMat(\widehat\auO)$
if and only if $x$ is, and $a^{-1}k'^{-1}xk=k'^{-1}(a^{-1}x)k$), and the factor $\upsilon'_w$ is trivial
on units.
\end{proof}

\subsubsection*{Forms and types}
Let $S$ be a space, $\pi=\bigotimes_\tau\pi_\tau$ an irreducible square-integrable representation of
$\auU(S)_\infty$ and $v_0\in\pi$ a unit vector. The \emph{type} is the function
$e(g)=d_\pi\aubar{\langle\pi(g)v_0,v_0\rangle}$ on $\auU(S)_\infty$, $d_\pi$ the formal degree. Let
$\pcF(S,e)$ be the space of bounded continuous functions $\varphi$ on $[\auU(S)]$ that are
right-invariant under an open subgroup of $K_S$ and are reproduced by right convolution with $e$
over $\auU(S)_\infty$, that is, $\varphi(x)=\int_{\auU(S)_\infty}\varphi(xg^{-1})e(g)\,dg$
(inner products are conjugate-linear in the first variable, so $e(g^{-1})=d_\pi\aubar{\langle v_0,
\pi(g)v_0\rangle}$, and the condition is $\varphi(x)=\int e(g^{-1})\varphi(xg)\,dg$, in agreement
with {\cite[(9)]{P49}} and {\cite[Definition~8.1]{P49}}); it carries the Petersson product $\langle\varphi',\varphi\rangle=
\int_{[\auU(S)]}\aubar{\varphi'}\varphi$ and the unitary action $\rho$ of the finite adelic group
by right translation. The types used in this paper are integrable functions (their parameters are
far from the walls; this is part of Lemma~\ref{pc:arch} below). For an integrable type, by Schur
orthogonality, right convolution with $e$ is the orthogonal projection of $L^2([\auU(S)])$ onto the
image of $\auHom_{\auU(S)_\infty}(\pi,L^2)\otimes v_0$: every $\varphi\in\pcF(S,e)$ is
$\Phi(v_0)$ for an intertwining map $\Phi\colon\pi\to L^2([\auU(S)])$, and
$\langle\Phi'(u),\Phi(v)\rangle=\langle\varphi',\varphi\rangle\langle u,v\rangle$ by Schur's lemma.
In particular for $\varphi,\varphi'\in\pcF(S,e)$ and an integrable function $c$ on
$\auU(S)_\infty$,
\begin{equation}\label{pc:eq-schur}
\int_{\auU(S)_\infty}c(g)\langle\varphi',\rho(g)\varphi\rangle\,dg=\Bigl(\int_{\auU(S)_\infty}c(g)
\langle v_0,\pi(g)v_0\rangle\,dg\Bigr)\langle\varphi',\varphi\rangle .
\end{equation}
The space $\auAz$ of {\cite[\S14.3]{P46}} is $\pcF(Y_0,e_0)^{J_M}$ for the type $e_0$ of the
parameter $\varsigma$, where $J_M$ also denotes $J_{Y_0}((\pi_M))$; the spaces $\mathcal A(W_j)$
are $\pcF(W_j,e_{W_j})^{J_M}$.

\begin{lemma}[Archimedean types]\label{pc:arch}
Let $(S,T)$ be one of the pairs $(W_j,V_j)$, $1\le j\le m$, or one of the pairs
$(\pcG_{i,0},\pcG_{i,1})$, $0\le i<m$, of \S\ref{pc:ss-triangle}, with the types $e_S,e_T$ of the
chain and the joint harmonic polynomial $P$ of the pair. Then:
\begin{enumerate}
\item[(A1)] the types $e_S$, $e_T$ are integrable, and every element of $\pcF(S,e_S)$ and of
$\pcF(T,e_T)$ decreases faster than every power of the height on the automorphic quotient;
\item[(A2)] for every modulus $\pcm$, every $\Psi$ and all characters $\alpha,\beta$ of the weights
of the pair, tamely ramified at $\pcm$ and unramified elsewhere, $\Theta_{\alpha,\beta}(P,\Psi)$
maps $\pcF(S,e_S)$ to $\pcF(T,e_T)$;
\item[(A3)] the integral $Z_\infty=\int_{\auU(S)_\infty}\langle P,\omega_\infty(g,1)P\rangle\,\langle
v_0,\pi_S(g)v_0\rangle\,dg$ converges absolutely and is not zero; it does not depend on the twist,
which is trivial at the infinite places.
\end{enumerate}
\end{lemma}
\noindent This is Proposition~\ref{ptv:arch}: (A1), (A2) for the pairs $(\pcG_{i,0},\pcG_{i,1})$ are
Theorem~\ref{ptr:edge}, (A3) is Proposition~\ref{ptr:Z} with Theorem~\ref{ptq:F}; for $(W_j,V_j)$,
(A1) is {\cite[Lemma~4.3]{P54}}, and (A2), (A3) follow from Step~1 of the proof of
{\cite[Proposition~4.8]{P54}} (modulo (V), (D1)--(D3)).

\subsubsection*{The coefficient of the vacuum and the series}
Let $S$ be a space of rank $n$, $M'\ge1$, and $\alpha$ an automorphic character of $\auU(1)$
unramified outside a modulus $\pcm$. Let $G^\pcm$ be the group of finite adelic points of
$\auU(S)$ with trivial components at $S(\pcm)$, and $K^\pcm=K_S\cap G^\pcm$. For $\gamma\in G^\pcm$
put
\[
\pcmc_{M',\alpha}(\gamma)=\alpha(\det\gamma)^{-1}\,\bigl\langle\psvac^\pcm,\omega_{\auHom(S,T)}
(\gamma,1)\psvac^\pcm\bigr\rangle ,
\]
where $T$ is any space of rank $M'$.

\begin{lemma}[The series]\label{pc:series}
\begin{enumerate}
\item[(a)] $\pcmc_{M',\alpha}$ does not depend on the choice of $T$, is the same for $S$ and for
the space $S^-$ with the negative form, and for all spaces $S$ of rank $n$ under the identification
of their finite adelic groups; it is bi-invariant under $K^\pcm$; and
$|\pcmc_{M',\alpha}(\gamma)|=[L_S:L_S\cap\gamma L_S]^{-M'/2}$.
\item[(b)] If $n\ge2$ and $M'\ge6n$, then
$\sum_{\gamma\in G^\pcm/K^\pcm,\ \gamma\notin K^\pcm}|\pcmc_{M',\alpha}(\gamma)|\le\exp\bigl(2f(4/5)^n
\bigr)-1$.
\item[(c)] Let $n=r$ and $M'=r+t$. The \emph{series}
$\pcRal=\pcRal_{M',\alpha}=\sum_{\gamma\in G^\pcm/K^\pcm}\pcmc_{M',\alpha}(\gamma)\rho(\gamma)$
converges absolutely in the supremum norm on the bounded $K^\pcm$-invariant functions on
$\auU(S)(\auA)$, maps them to themselves, commutes with left translations and with right
translations by $\auU(S)_\infty$ and by the components at $S(\pcm)$, and satisfies
$\|\pcRal y\|_\infty\ge\tfrac12\|y\|_\infty$. In particular $\pcRal$ is injective, and it preserves
$\pcF(S,e)^J$ for every type $e$ and every open $J\supseteq K^\pcm$.
\end{enumerate}
\end{lemma}
\begin{proof}
(a) The Weil representation of $\auU(S)(E^+_v)$ on $\psHS_v(\auHom(S,T))$, with the splitting
defined by $\upsilon$, depends only on the hermitian space $\auHom(S,T)_v$ with its lattice. At a
finite place all spaces of a given rank are isometric with their lattices; and multiplication by a
unit $\epsilon_v$ of $E_v$ with $\epsilon_v\aubar{\epsilon_v}=-1$ (which exists because $E_v/E^+_v$
is unramified or split) is an isometry $S_v^-\to S_v$ that commutes with the group. This gives the
independence. Bi-invariance follows from (W4). For the absolute value put
$\mathbb L=a\auHom(L_S,L_T)$ and $\mathbb L'=\gamma\mathbb L=a\auHom(\gamma L_S,L_T)$, two
self-dual lattices, and $\Lambda=\mathbb L\cap\mathbb L'$. By (W4),
$\omega(\gamma,1)\psvac_{\mathbb L}$ is a unit vector on the line of $\psvac_{\mathbb L'}$. Since
$\epsilon_{\mathbb L}$ and $\epsilon_{\mathbb L'}$ are given by the same formula, they agree on
$\Lambda$, and the vector
$v'=\sum_{l'\in\mathbb L'/\Lambda}\epsilon_{\mathbb L'}(l')^{-1}\psrho(l')\psvac_{\mathbb L}$ is well
defined; the relation $\epsilon(l)\epsilon(l')=\beta(l,l')\epsilon(l+l')$ gives
$\psrho(l_0)v'=\epsilon_{\mathbb L'}(l_0)v'$ for $l_0\in\mathbb L'$, so $v'$ lies on the line of
$\psvac_{\mathbb L'}$. The vectors $\psrho(l')\psvac_{\mathbb L}$, for $l'$ in distinct classes
modulo $\Lambda$, hence modulo $\mathbb L$, are eigenvectors of the operators $\psrho(l)$,
$l\in\mathbb L$, for distinct characters, hence orthogonal, of norm $1$. So
$\|v'\|^2=[\mathbb L':\Lambda]$ and $\langle\psvac_{\mathbb L},v'\rangle=1$, whence
$|\langle\psvac_{\mathbb L},\omega(\gamma,1)\psvac_{\mathbb L}\rangle|=[\mathbb L':\Lambda]^{-1/2}$.
Finally $\Lambda=a\auHom(L_S+\gamma L_S,L_T)$ and $[\mathbb L':\Lambda]=[L_S+\gamma L_S:\gamma
L_S]^{M'}=[L_S:L_S\cap\gamma L_S]^{M'}$.

(b) \emph{Counting.} Let $\mathfrak o$ be the ring of integers of a local field with residue field
of $q$ elements. By the Hermite normal form, the number $s_k$ of sublattices of index $q^k$ in
$\mathfrak o^n$ is $\sum_{k_1+\dots+k_n=k}q^{\sum_i(i-1)k_i}$; hence
\begin{equation}\label{pc:eq-sk}
\sum_{k\ge0}s_kX^k=\prod_{i=0}^{n-1}(1-q^iX)^{-1},\qquad s_k\le\tbinom{k+n-1}{n-1}q^{k(n-1)}\le
2^{k+n-1}q^{k(n-1)} .
\end{equation}
The same numbers count the lattices containing a given lattice with index $q^k$. Every residue
field of $E$ has at least $5$ elements, because $E$ contains the $5040$-th roots of unity: if the
residue characteristic is not $5$ then $5\mid q_w-1$, and otherwise $7\mid q_w-1$.

\emph{One place.} The function $|\pcmc|$ is a product over the places $v\notin S(\pcm)$ of $E^+$
of the functions $\gamma_v\mapsto[L_{S,v}:L_{S,v}\cap\gamma_vL_{S,v}]^{-M'/2}$, and
$G_v/K_v$ is identified with the set of lattices $L'=\gamma_vL_{S,v}$. Let $\varepsilon_v$ be the
sum of $[L:L\cap L']^{-M'/2}$ over $L'\ne L=L_{S,v}$. If $v=w\aubar w$ splits, $L'$ runs through the
lattices of $E_w^n$ (the component at $\aubar w$ is the dual lattice), and the index is $q^{k+k'}$
with $q=q_w$, $q^k=[L_w:L_w\cap L'_w]$, $q^{k'}=[L'_w:L_w\cap L'_w]$. A lattice $L'_w$ is determined
by $L_w\cap L'_w$ and by $L'_w$ above it, so by \eqref{pc:eq-sk}, with $M'/2\ge3n$,
\[
1+\varepsilon_v\le\Bigl(\sum_ks_kq^{-kM'/2}\Bigr)^2=\prod_{i=0}^{n-1}(1-q^{i-M'/2})^{-2}\le
\exp\Bigl(4\sum_{i<n}q^{i-3n}\Bigr)\le\exp(5q^{-2n-1}),
\]
so $\varepsilon_v\le10\,q^{-2n-1}\le2(4/5)^nq_w^{-2}$, because $50\cdot25^{-n}\le2(4/5)^n$ for
$n\ge2$. If $v$ does not split, $L'$ runs through the self-dual lattices of $S_v$; such a lattice
is determined by $\Lambda=L\cap L'$ and by $L'\supset\Lambda$, both of the same index $q_w^k$ since
$L$ and $L'$ have the same volume. So by \eqref{pc:eq-sk}
\[
\varepsilon_v\le\sum_{k\ge1}s_k^2q_w^{-3nk}\le4^{n-1}\sum_{k\ge1}\bigl(4q_w^{-n-2}\bigr)^k\le
2\cdot4^nq_w^{-n-2}\le2(4/5)^nq_w^{-2}.
\]
\emph{All places.} Choose a place $w$ of $E$ above each $v$. There are at most $2f$ places of $E$
above a rational prime $p$, each with $q_w\ge p$, and $\sum_pp^{-2}<\tfrac12$. Hence
$\sum_v\varepsilon_v\le2f(4/5)^n$ and
$\sum_{\gamma\notin K^\pcm}|\pcmc(\gamma)|=\prod_v(1+\varepsilon_v)-1\le\exp(2f(4/5)^n)-1$.

(c) By (b) with $n=r\ge x^5$ and $f\le x^{12}$ {\cite[(38)]{P46}}, the sum of
$|\pcmc(\gamma)|$ over $\gamma\notin K^\pcm$ is at most $e^{1/3}-1<\tfrac12$ (comparison of sizes).
The series is well defined on $K^\pcm$-invariant functions because $\pcmc$ is right
$K^\pcm$-invariant, its values are $K^\pcm$-invariant because $\pcmc$ is left invariant, and
$\pcRal y=y+Ey$ with $\|Ey\|_\infty\le\tfrac12\|y\|_\infty$. The commutation properties are clear,
and they give the last assertion: boundedness, continuity, the reproducing property and
$J$-invariance are preserved.
\end{proof}

For $1\le j\le m$ and a twist $\vartheta$ we write
\begin{equation}\label{pc:eq-Rj}
\pcRal_j(\vartheta)=\pcRal_{r+t,\ \mu_j^0\vartheta\,\upsilon_{r+t}}\qquad\text{with the modulus }
\pcm_\vartheta ,
\end{equation}
acting on $\pcF(S,e)^{J_M}$ for every space $S$ of rank $r$ and every type $e$; here the forms are
also invariant under $K_S$ at the places of $Q_\vartheta$. Let $\pcH$ be the algebra of compactly
supported functions on $G^{(\pi_M)}$ that are bi-invariant under $K^{(\pi_M)}$, the spherical
Hecke algebra away from $M$ at \emph{all} places, split or not; it acts on the spaces
$\pcF(S,e)^{J_M}$ and commutes with the action of $\auG=K_S/J_M$.

\begin{lemma}\label{pc:heckeclosure}
Let $S,S'$ be spaces of rank $r$ with types $e,e'$ such that $\pcF(S,e)^{J_M}$ and
$\pcF(S',e')^{J_M}$ are finite-dimensional. For every $j$ and $\vartheta$ there is $h\in\pcH$ that
acts as $\pcRal_j(\vartheta)$ on both spaces. In particular $\pcRal_j(\vartheta)$ commutes with
$\auG$ and preserves the images of $\varrho(e_0)$.
\end{lemma}
\begin{proof}
Grouping the cosets into double cosets, $\pcRal_j(\vartheta)=\sum_{[\gamma]\in K^\pcm\backslash
G^\pcm/K^\pcm}\pcmc(\gamma)\,T_\gamma$ on $K^\pcm$-invariant functions, where $\pcm=\pcm_\vartheta$
and $T_\gamma$ is the action of the characteristic function of $K^\pcm\gamma K^\pcm$, an element of
$\pcH$ (extended by the unit at the places of $Q_\vartheta$). The series converges absolutely in
$\auEnd(\pcF(S,e)^{J_M})\times\auEnd(\pcF(S',e')^{J_M})$ by Lemma~\ref{pc:series}(c). The image of
$\pcH$ in this finite-dimensional space is a linear subspace, hence closed, and it contains the
partial sums.
\end{proof}

\noindent For every $M'>2n$ the series converges absolutely and is injective by a finer count
({\cite[Proposition~5.39]{P54}}); in this section only $M'=6n$ is used, where the count above suffices.

\subsection{The inner product formula for sharp data}\label{pc:ss-rallis}

\begin{lemma}[Two standard facts]\label{pc:SW}
\begin{enumerate}
\item[(W7)] \emph{(complex conjugation).} Let $S,T$ be spaces of ranks $n,M'$ and $S^-$ the space
$S$ with the negative form. There is a conjugate-linear isometry
$\mathsf C\colon\psHS(\auHom(S,T))\to\psHS(\auHom(S^-,T))$ with
\[
\mathsf C\psrho(x)=\psrho^-(x)\mathsf C,\qquad\theta^-(\mathsf CF)=\aubar{\theta(F)},\qquad
\mathsf C\,\omega_{\auHom(S,T)}(g,h)\,\mathsf C^{-1}=\lambda(g,h)\,\omega_{\auHom(S^-,T)}(g,h),
\]
for a character $\lambda$ of $\auU(S)(\auA)\times\auU(T)(\auA)$ that is trivial on the rational
points and is a product of powers of $\upsilon\circ\det$.
\item[(SW)] \emph{(Weil's Siegel formula).} Let $\pcW$ be a hermitian space of rank $2n$ over $E$
that contains an isotropic subspace of dimension $n$, and let $T$ be a hermitian space of rank
$M'$ such that (i) the theta integral converges absolutely ($T$ anisotropic over $E^+$, or
$M'-r_T>n$ with $r_T$ the Witt index of $T$) and (ii) $M'>2n$; both hold if $M'\ge6n$. Let $\auU(T)(\auA)$ act on $\psHS(\auHom(\pcW,T))$ by the Weil representation
$\omega(1,h)$ with the trivial splitting character; for $\pcW=S\oplus S^-$ and
$\psHS(\auHom(\pcW,T))=\psHS(\auHom(S,T))\otimes\psHS(\auHom(S^-,T))$ this is
$\omega_{\auHom(S,T)}(1,h)\otimes\mathsf C\,\omega_{\auHom(S,T)}(1,h)\,\mathsf C^{-1}$. For a
Lagrangian (maximal isotropic subspace) $L\subset\pcW$ let
$I_L=\{w\in\auHom(\pcW,T):w|_L=0\}$, a maximal isotropic subspace of $\auHom(\pcW,T)$, and let
$\psZ_L=\psZ_{I_L}\colon\psHS(\auHom(\pcW,T))\to\C$ be its zero term \textup{(W2)}. Then for every
smooth $\Phi$
\[
\int_{[\auU(T)]}\theta\bigl(\omega(1,h)\Phi\bigr)\,dh=\pcvol([\auU(T)])\sum_L\psZ_L(\Phi),
\]
the sum being over all Lagrangians of $\pcW$; the integral and the sum converge absolutely, and
$\sum_L|\psZ_L(\omega(g,1)\Phi)|$ is bounded by a power of the height of $g\in\auU(\pcW)(\auA)$.
Moreover, for $\pcW=S\oplus S^-$, data $F,F'$ of $\auHom(S,T)$ and characters $\alpha,\beta$ there are
$C,A$ with
\[
\int_{[\auU(T)]}\bigl|\theta(\omega_{\alpha,\beta}(g,h)F)\,\theta(\omega_{\alpha,\beta}(g',h)F')\bigr|\,dh\le
C\,(\mathrm{ht}(g)\mathrm{ht}(g'))^A\qquad(g,g'\in\auU(S)(\auA));
\]
and if $T$ is anisotropic over $E^+$, $\int|\theta(\omega(g,h)\Phi)|\,dh$ is bounded by a power of the
height of $g\in\auU(\pcW)(\auA)$.
\end{enumerate}
\end{lemma}
\begin{proof}[References and comments]
(W7): the cocycle $\beta$ of $S^-$ is the complex conjugate of that of $S$, so the complex
conjugate of the Heisenberg representation of $\auHom(S,T)$ is that of $\auHom(S^-,T)$; this gives
$\mathsf C$, unique up to a scalar of absolute value $1$, which is fixed by the condition on
$\theta$ (the functional $F\mapsto\aubar{\theta(\mathsf C^{-1}F)}$ is invariant under the rational
vectors, hence a multiple of $\theta^-$ by (W1)). By (W1) the operator $\mathsf C\omega(g,h)
\mathsf C^{-1}$ differs from $\omega^-(g,h)$ by a character; that it is a product of powers of
$\upsilon\circ\det$ follows from the form of the splitting~\cite[Theorem~3.1]{KudlaSplitting}, where
$\upsilon$ is replaced by $\aubar\upsilon=\upsilon^{-1}$. The exponents are computed in
Lemma~\ref{ptw:W7}: $\lambda(g,h)=\upsilon(\det g)^{-M'}\upsilon(\det h)^{n}$. (SW) is the Siegel formula of
Weil~\cite[Th\'eor\`eme~5]{WeilSiegel} in the range of absolute convergence, for the dual pair
$(\auU(\pcW),\auU(T))$: the Eisenstein series $\sum_{\gamma\in P(E^+)\backslash\auU(\pcW)(E^+)}
\psZ_{L_0}(\omega(\gamma g,1)\Phi)$ attached to a Lagrangian $L_0$ with stabiliser $P$ converges
for $M'>2n$ (hypothesis (ii)), the theta integral converges under (i) (when $T$ is anisotropic
$[\auU(T)]$ is compact; otherwise this is Weil's criterion that $M'$ minus the Witt index of $T$
exceeds $n$), and the two are equal up to the volume; see also
Ichino~\cite[Introduction, p.~242]{IchinoSW}. The terms are indexed by the Lagrangians $L=\gamma^{-1}L_0$, and
$\psZ_{L_0}\circ\omega(\gamma,1)=\psZ_L$ by the characterisation of the zero term in (W2) and the
invariance of $\theta$ under rational points. In Weil's notation (the hermitian case (I$_2$) of~\cite[no.~27]{WeilSiegel}, $\varepsilon=\frac12$, $m=M'$ the rank of $T$ over $E$,
$n$ the rank of $S$, $r$ the Witt index of $T$): Th\'eor\`eme~5 assumes condition (B) of~\cite[no.~38]{WeilSiegel},
$m>2n+4\varepsilon-2$, i.e.\ $M'>2n$, and the absolute convergence of the theta integral; it also assumes a place $v$ with $U(0)_v\ne\emptyset$ ($U(0)$ is the set of elements of maximal rank whose associated form vanishes) and a subgroup $G'_v$ transitive on each $U(i)_v$; we take $G'=G$, which is transitive on each non-empty $U(i)$ by~\cite[no.~22, Proposition~3]{WeilSiegel}, and $v$ a place of $E^+$ split in $E$, where $T_v$ has Witt index $\lfloor M'/2\rfloor\ge n$, so that a map of rank $n$ onto a totally isotropic subspace exists; the Eisenstein series then converges
absolutely by~\cite[Th\'eor\`eme~1]{WeilSiegel}, and the theta integral converges absolutely when $r=0$ or
$m-r>n+2\varepsilon-1$, i.e.\ $M'-r_T>n$~\cite[Proposition~8]{WeilSiegel}, which follows from $M'>2n$ since $r_T\le M'/2$. Weil's theorems give convergence uniformly on compact sets only; the bounds by a power of the height are proved as follows. For the Eisenstein series it is {\cite[Lemma~5.41]{P54}}, which needs only $M'>2n$; its proof uses of the present comments only what precedes this point, namely the translation of Weil's statements and the indexing of the terms by the Lagrangians, which rest on~\cite{WeilSiegel} and (W2) alone, so that the bounds may be deduced from it. For anisotropic $T$, Lemma~\ref{pcx:thetagrowth} for the pair $(\pcW,T)$, with $\mathrm{ht}(h)$ bounded on a compact set of representatives of $[\auU(T)]$, gives the bound for $\int|\theta(\omega(g,h)\Phi)|\,dh$. For the product integral and any $T$ as in (SW): by the Cauchy--Schwarz inequality in $h$ it suffices to bound $\int|\theta(\omega_{\alpha,\beta}(g,h)F)|^2\,dh$ by $C\,\mathrm{ht}(g)^A$. Put $G=\omega_{\alpha,\beta}(g,1)F$; then $\omega_{\alpha,\beta}(g,h)F=\omega_{\alpha,\beta}(1,h)G$ differs from $\omega_{\auHom(S,T)}(1,h)G$ by a character of $h$ of absolute value $1$, and by (W7) and the factorisation $\theta\otimes\theta^-$ of the theta functional of $\auHom(\pcW,T)$ on pure tensors ((W1)), $\theta(\omega(1,h)\Phi_g)=|\theta(\omega_{\auHom(S,T)}(1,h)G)|^2$ for $\Phi_g=G\otimes\mathsf CG$. So $\int|\theta(\omega_{\alpha,\beta}(g,h)F)|^2\,dh=\int\theta(\omega(1,h)\Phi_g)\,dh$, an integral of a non-negative function, which converges under (i) and equals $\pcvol([\auU(T)])\sum_L\psZ_L(\Phi_g)$ by the identity. Up to a factor of absolute value $1$ (the splitting characters and the character $\lambda$ of (W7) are unitary), $\Phi_g=\omega((g,g),1)\Phi_1$ with $\Phi_1=F\otimes\mathsf CF$ and $(g,g)\in\auU(S)\times\auU(S^-)\subseteq\auU(\pcW)$. Hence the integral is at most $\pcvol([\auU(T)])\sum_L|\psZ_L(\omega((g,g),1)\Phi_1)|\le C\,\mathrm{ht}(g,g)^A=C\,\mathrm{ht}(g)^A$, by {\cite[Lemma~5.41]{P54}} and the comparison of heights in the last paragraph of the comments to {\cite[Lemma~5.40]{P54}}, an elementary computation with frames that does not use (SW). (This is all that Proposition~\ref{pc:rallis} uses for isotropic $T$: absolute convergence of the theta integral, Weil's identity, the bound for $\sum_L|\psZ_L|$ in Step~1 and for the product integral in the Fubini argument of Step~2; no bound for $\int|\theta(\omega(g,h)\Phi)|\,dh$ at a general $g\in\auU(\pcW)(\auA)$ is used.) The statement is used in
this section with $M'=6n$ (Proposition~\ref{pc:rallis}) and elsewhere with a definite $T$ of rank
$M'>2n$ ({\cite[Lemma~5.40]{P54}}, {\cite[Proposition~3.4]{P54}}, \S\ref{pgy:sec}).
\end{proof}

\begin{lemma}[Clean data]\label{pc:clean}
Let $S,T$ be as in \textup{(SW)}, $\pcW=S\oplus S^-$, and let $F\in\psHS(\auHom(S,T))$ be a datum
whose component at one split place $v=w\aubar w$ is, in the model \textup{(W5)}, a Schwartz--Bruhat
function supported on injective matrices. Let $F'\in\psHS(\auHom(S,T))$ and
$\Phi=F\otimes\mathsf CF'\in\psHS(\auHom(S,T))\otimes\psHS(\auHom(S^-,T))=\psHS(\auHom(\pcW,T))$.
\begin{enumerate}
\item[(a)] If the Lagrangian $L$ of $\pcW$ is not of the form $L_\delta=\{(\delta s,-s):s\in S\}$
for an element $\delta\in\auU(S)(E^+)$, then $\psZ_L(\Phi)=0$.
\item[(b)] There is a constant $c'\ne0$, depending only on the normalisations, such that for all
$\delta\in\auU(S)(E^+)$ and all $F,F'$
\[
\psZ_{L_\delta}(F\otimes\mathsf CF')=c'\,\bigl\langle F',\omega_{\auHom(S,T)}(\delta^{-1},1)F\bigr
\rangle .
\]
\end{enumerate}
\end{lemma}
\begin{proof}
(a) If the second projection $L\to S$ is injective it is bijective, $L=\{(\delta s,-s)\}$ for a
linear map $\delta$, and the isotropy of $L$ says that $\delta$ is an isometry. Otherwise there is
$\ell\ne0$ in $S$ with $(\ell,0)\in L$. For $z\in T(\auA)$ let $i_z\in\auHom(S,T)(\auA)\subset
\auHom(\pcW,T)(\auA)$ be the map $s\mapsto\langle\ell,s\rangle z$. It vanishes on $L$, because
$\langle\ell,s\rangle=\langle(\ell,0),(s,s')\rangle_{\pcW}=0$ for $(s,s')\in L$; so $i_z\in
I_L(\auA)$ and, by (W2), $\psZ_L(\psrho(i_z)\Phi)=\psZ_L(\Phi)$. The adjoint of $i_z$ is
$t\mapsto\langle z,t\rangle\ell$, so $\langle i_z,x\rangle=\operatorname{tr}(i_z^\dagger x)=
\langle z,x\ell\rangle_T$ for $x\in\auHom(S,T)$. Take $z$ in the component $T_{\aubar w}$ of $T$ at
$v$. By (W5), $\psrho(i_z)$ multiplies the component $F_w$ of $F$ by the function
$x\mapsto[x,i_z]$, which for fixed $x$ is the character $z\mapsto\psi_v(\tfrac12\jmath(\langle z,
x\ell\rangle-\langle x\ell,z\rangle))$ of $T_{\aubar w}$; it is non-trivial when $x\ell\ne0$, since
the form pairs $T_{\aubar w}$ and $T_w$ perfectly. On the compact support of $F_w$ the matrices
are injective, so $x\ell$ stays in a compact subset of $T_w\smallsetminus\{0\}$, and there is a
compact open subgroup $C\subset T_{\aubar w}$ with $\int_C[x,i_z]\,dz=0$ for all $x$ in the support.
Then $\int_C\psrho(i_z)\Phi\,dz=0$ (a finite sum), and $\psZ_L(\Phi)=\pcvol(C)^{-1}\int_C\psZ_L(
\psrho(i_z)\Phi)\,dz=0$.

(b) For $\Delta=L_1=\{(s,-s)\}$, $I_\Delta(\auA)$ consists of the maps $(s,s')\mapsto x(s+s')$,
that is, of the pairs $(x,x)$, and the Heisenberg operator of $(x,x)$ on the tensor product is
$\psrho(x)\otimes\psrho^-(x)$. So by (W2) and by (W7) of Lemma~\ref{pc:SW},
$\psZ_\Delta(\psrho(x)F\otimes\mathsf C\psrho(x)F')=\psZ_\Delta(F\otimes\mathsf CF')$ for all $x$:
the sesquilinear form $(F',F)\mapsto\psZ_\Delta(F\otimes\mathsf CF')$ is invariant under the
Heisenberg group. As the Heisenberg representation is irreducible, it is a multiple $c'$ of the
inner product, and $c'\ne0$ because $\psZ_\Delta\ne0$. For general $\delta$,
$L_\delta=(\delta,1)\Delta$ and $\psZ_{L_\delta}=\psZ_\Delta\circ\omega((\delta,1)^{-1},1)$, as in
the proof of (SW); the element $(\delta,1)$ acts on $F\otimes\mathsf CF'$ through
$\omega_{\auHom(S,T)}(\delta,1)F$, up to a character that is trivial on rational points.
\end{proof}

\begin{definition}[Sharp functions]\label{pc:sharp}
Let $S,T$ be spaces of ranks $n,M'$, $\pcm$ a modulus, $\alpha$ a character tamely ramified at
$\pcm$, and $J$ a subgroup of $K_S$ containing $J_S(\pcm)$ and the components of $K_S$ outside
$S(\pcm)$, with image $\bar J\subseteq\auGL_n(\pcR_\pcm)$. A function $\Psi\ne0$ on
$\pcMat_{M'\times n}(\pcR_\pcm)$ is \emph{sharp for $J$} if
$\alpha_\pcm(\det\bar k)^{-1}\Psi(x\bar k)=\Psi(x)$ for $\bar k\in\bar J$ and there is a matrix
$A\in\pcMat_{n\times M'}(\pcR_\pcm)$ with $Ax\in\bar J$ for all $x$ in the support of $\Psi$. We put
$\|\Psi\|^2=\sum_x|\Psi(x)|^2$.
\end{definition}

\begin{proposition}[Inner product formula]\label{pc:rallis}
Let $(S,T)$, of ranks $r$ and $r+t$, with types $e_S,e_T$ and polynomial $P$, be a pair for which
\textup{(A1)--(A3)} of Lemma~\ref{pc:arch} hold. Let $\pcm$ be a modulus with $S(\pcm)\ne\emptyset$ \textup(i.e.\ $\pcm\ne(1)$\textup), $\alpha,\beta$ characters
as in \textup{(A2)}, and $\Psi$ sharp for $J$. Put $\Theta=\Theta_{\alpha,\beta}(P,\Psi)$ and
$\pcRal=\pcRal_{r+t,\alpha}$. Then for all $y,y'\in\pcF(S,e_S)^J$
\[
\langle\Theta y',\Theta y\rangle=A_\infty\cdot\frac{\|\Psi\|^2}{\#\pcMat_{(r+t)\times r}(\pcR_\pcm)
\cdot[K_S:J]}\cdot\langle y',\pcRal\,y\rangle,\qquad A_\infty=c'\,Z_\infty\,\pcvol([\auU(T)])\ne0,
\]
where $A_\infty$ depends only on the pair and its archimedean data. Moreover $\Theta$ is injective
on $\pcF(S,e_S)^J$.
\end{proposition}
\begin{proof}
Write $F=[P,\Psi]$ and $\omega_\alpha(g)=\omega_{\alpha,\beta}(g,1)$.

\emph{Step 1: doubling.} The component of $F$ at any $v\in S(\pcm)$ is supported on injective
matrices: if $\Psi(\bar x)\ne0$ then $A\bar x$ is invertible. So is the component of
$\omega_\alpha(g)F$ for every $g\in\auU(S)(\auA)$, since by (W5) $\omega_\alpha(g)$ acts at $v$
through $x\mapsto xg_v$ up to a scalar, and $xg_v$ is injective when $x$ is. The theta functional of
$\auHom(\pcW,T)$ is $\theta\otimes\theta^-$ on pure tensors (it is invariant under the rational
vectors; (W1)). For $g,g'\in\auU(S)(\auA)$ and $h\in\auU(T)(\auA)$, by (W7) of Lemma~\ref{pc:SW} and since
$|\beta(\det h)|=1$,
\[
\theta(\omega_{\alpha,\beta}(g,h)F)\,\aubar{\theta(\omega_{\alpha,\beta}(g',h)F)}=\theta_{\auHom(\pcW,
T)}\bigl(\omega(1,h)\,[\omega_\alpha(g)F\otimes\mathsf C\omega_\alpha(g')F]\bigr),
\]
with $\omega(1,h)$ the action of (SW). By (SW) and Lemma~\ref{pc:clean},
\begin{equation}\label{pc:eq-doubling}
\int_{[\auU(T)]}\aubar{\theta(\omega_{\alpha,\beta}(g',h)F)}\,\theta(\omega_{\alpha,\beta}(g,h)F)\,dh
=c'\pcvol([\auU(T)])\sum_{\delta\in\auU(S)(E^+)}\bigl\langle F,\omega_\alpha(g'^{-1}\delta g)F
\bigr\rangle ,
\end{equation}
with absolute convergence and a bound by a power of the heights of $g$ and $g'$.

\emph{Step 2: unfolding.} By (A1) the functions $y,y'$ decrease rapidly, so the following integrals
converge absolutely and Fubini's theorem applies:
\begin{align*}
\langle\Theta y',\Theta y\rangle&=\int_{[\auU(S)]^2}\aubar{y'(g')}y(g)\int_{[\auU(T)]}\aubar{\theta(
\omega(g',h)F)}\theta(\omega(g,h)F)\,dh\,dg\,dg'\\
&=c'\pcvol([\auU(T)])\int_{\auU(S)(\auA)}\langle F,\omega_\alpha(g)F\rangle\,\langle y',\rho(g)y
\rangle\,dg ,
\end{align*}
where the sum over $\delta$ in \eqref{pc:eq-doubling} has been combined with the integral over
$g\in[\auU(S)]$ and $g$ replaced by $g'g$.

\emph{Step 3: the three factors.} The function $g\mapsto\langle F,\omega_\alpha(g)F\rangle$ is the
product of its values on the components $g_\infty$, $g_\pcm$ (at $S(\pcm)$) and $g^\pcm$.
\emph{Outside $\pcm$}: $y$ is $K^\pcm$-invariant and $K^\pcm$ has volume $1$, so
$\int_{G^\pcm}\langle\psvac^\pcm,\omega_\alpha(g)\psvac^\pcm\rangle\rho(g)y\,dg=\pcRal y$.
\emph{At $\pcm$}: by (W5), $\omega_\alpha(g)\Psi^\natural(x)$ is a non-zero multiple of
$\Psi^\natural(xg)$. If $\langle\Psi^\natural,\omega_\alpha(g)\Psi^\natural\rangle\ne0$ there is
$x\in a\pcMat(\widehat\auO)$ with $xg\in a\pcMat(\widehat\auO)$ and the reductions of $a^{-1}x$, $a^{-1}xg$ in
the support of $\Psi$; with an integral lift $\tilde A$ of $A$, the matrices $\tilde Aa^{-1}x$ and
$\tilde Aa^{-1}xg$ are integral with reductions in $\bar J$, hence lie in $J$, and
$g=(\tilde Aa^{-1}x)^{-1}(\tilde Aa^{-1}xg)\in J$ ($a$ is a scalar). For $g\in J$,
$\omega_\alpha(g)\Psi^\natural=\Psi^\natural$ by Lemma~\ref{pc:level} and the first condition of
sharpness. So the coefficient is $\|\Psi^\natural\|^2$ times the indicator function of $J$, with
$\|\Psi^\natural\|^2=\|\Psi\|^2/\#\pcMat_{(r+t)\times r}(\pcR_\pcm)$, and the integral of
$\rho(g_\pcm)$ against it is multiplication by $\|\Psi^\natural\|^2/[K_S:J]$ on $J$-invariant
functions. \emph{At infinity}: $\pcRal y\in\pcF(S,e_S)^J$ by Lemma~\ref{pc:series}(c), and
\eqref{pc:eq-schur} with (A3) gives the factor $Z_\infty$. This proves the formula.

\emph{Step 4: injectivity.} If $\Theta y=0$, the formula with $y'=\pcRal y\in\pcF(S,e_S)^J$ gives
$\langle\pcRal y,\pcRal y\rangle=0$, so $\pcRal y=0$ and $y=0$ by Lemma~\ref{pc:series}(c).
\end{proof}

\subsubsection*{The single lifts}
Let $1\le j\le m$ and let $\vartheta$ be a twist. The lift $u=\Theta_{j,\vartheta}$ of
{\cite[Proposition~14.37]{P46}} is defined as follows; the normalisation matters, because it depends on
$\vartheta$. The pair is $(W_j,V_j)$, with $\alpha=\mu_j^0\vartheta\,\upsilon_{r+t}$,
$\beta=\mu_j^0\vartheta\,\upsilon_r$, the joint harmonic $P_j^W$ and the modulus $\pcm_\vartheta$.
The finite datum is the \emph{block function} $\Psi^{\mathrm{bl}}_\vartheta$ on
$\pcMat_{(r+t)\times r}(\pcR_\vartheta)$: at a chosen place the indicator of $\iota$; at $w\in
Q_\vartheta$ the function $x\mapsto\psi_w(\det x^{\mathrm{top}})$, where $x^{\mathrm{top}}$ is the
block of the first $r$ rows and $\psi_w$ is the tame part of $\vartheta$ (extended by $0$). Let
$J^V_\vartheta\subseteq K_{V_j}$ consist of the elements whose reduction fixes the first $r$ frame
vectors at the chosen places and preserves the span of the last $t$ frame vectors at the places of
$Q_\vartheta$. By Lemma~\ref{pc:level}, $J^V_\vartheta$ acts on $[P,\Psi^{\mathrm{bl}}_\vartheta]$
by a character, so $\Theta_{\alpha,\beta}(P^W_j,\Psi^{\mathrm{bl}}_\vartheta)$ takes values in the
functions on $\auU(V_j)$ that correspond, by $\varphi\mapsto(h\mapsto\langle\xi_0,\pi^\natural(h_\infty)^{-1}
\varphi(h_f)\rangle)$ with $\xi_0$ the vector of {\cite[Lemma~4.7]{P54}} ($\pop_{\xi_0}=P^W_j$) and
$\pi^\natural$ the action of {\cite[\S2.1]{P54}} for which $\xi\mapsto\pop_\xi$ is equivariant, to the space
$\mathcal A(V_j;\vartheta)$ of algebraic modular forms of the weight $\lambda_V$ of {\cite[\S2.1]{P54}},
level $J^V_\vartheta$ and that character, with the inner product of {\cite[\S14.3]{P46}}; here $e_{V_j}$
is the type of the representation of highest weight $\Lambda'$ generated by $P^W_j$
(Proposition~\ref{ptw:A2}). Then
\[
\Theta_{j,\vartheta}\colon e_0\mathcal A(W_j)\to\mathcal A(V_j;\vartheta),\qquad
Z_j(\vartheta)=\det(\Theta_{j,\vartheta}^*\Theta_{j,\vartheta}).
\]

\begin{remark}[The two normalisations of the single lift]\label{pc:rem-norm}
In {\cite[\S2.1]{P54}} and {\cite[\S6.1.4]{P55}} the single lift is formed with the representation
$(g,h)\mapsto a(\det h)\,b(\det g)^{-1}\omega_{\auHom(W_j,V_j)}(g,h)$, $a=\upsilon^{r_-}$,
$b=\mu\,\upsilon^{\lfloor(r+t)/2\rfloor}$, $\mu=\mu_j^0\vartheta$, where $r_-$ is the number of negative
signs of $W_j$. On $\auU(W_j)$ this is the action used here, since $\alpha=\mu\,\upsilon_{r+t}=b$. On
$\auU(V_j)$ the two differ by the character $\chi\circ\det$, $\chi=\beta a^{-1}=\mu\,
\upsilon^{\lfloor r/2\rfloor-r_-}$, a \emph{unitary} automorphic character of $\auU(1)$. Hence the
lift $\Theta'$ of {\cite[\S2.1]{P54}} and the lift $\Theta$ defined above are related by
$\Theta=\mathsf m_\chi\circ\Theta'$, where $\mathsf m_\chi$ is multiplication by $\chi(\det h)$; this
is an isometry from the forms on $\auU(V_j)$ of level $J^V_\vartheta$ with one character of the level
(and one weight) to those with the character multiplied by $\chi\circ\det$ (and the weight shifted
by that of $\chi$), for the inner products of {\cite[\S14.3]{P46}}. Therefore
$\Theta^*\Theta=\Theta'^*\Theta'$, and $Z_j(\vartheta)$, the series of Rallis' formula and
everything in {\cite[Propositions~14.35 and~14.37]{P46}} are the same for both. The
normalisation matters only where the Hecke operators on $\auU(V_j)$ or the integrality of forms on
$\auU(V_j)$ are used: these are {\cite[Proposition~6.7]{P55}}, {\cite[Lemma~2.2]{P54}} and
{\cite[\S2]{P54}}, which all use $\Theta'$; under $\mathsf m_\chi$ every root of the Hecke polynomial at
a split place $w$ is multiplied by the value of the base change of $\chi$ at $w$.
\end{remark}

\begin{lemma}[Algebraic forms and the Petersson product]\label{pc:algpet}
Let $S$ be a definite space, $\lambda$ a weight with highest weight vector $v^+$, $J\subseteq K_S$
open and $\chi$ a character of $J$ of finite order. For $\varphi\in\mathcal A_\lambda(J,\chi)$ let
$f_\varphi(h)=\langle v^+,h_\infty^{-1}\varphi(h_f)\rangle$. Then
$\langle\varphi',\varphi\rangle=[K_S:J]\cdot\dim W_\lambda\cdot\|v^+\|^{-2}\cdot\langle f_{\varphi'},
f_\varphi\rangle_{[\auU(S)]}$. The statement and proof hold for every unit vector in place of $v^+$.
\end{lemma}
\begin{proof}
$[\auU(S)]$ is the disjoint union, over $x\in\Gamma\backslash G_f/J$, of the images of
$\auU(S)_\infty\times xJ$, of volume $(\#(\Gamma\cap xJx^{-1})\,[K_S:J])^{-1}$, and
$|\chi|=1$. On each of them Schur orthogonality for the irreducible representation $W_\lambda$ of
the compact group $\auU(S)_\infty$, of total mass $1$, gives
$\int\aubar{\langle v^+,u^{-1}\xi'\rangle}\langle v^+,u^{-1}\xi\rangle\,du=\|v^+\|^2\langle\xi',\xi
\rangle/\dim W_\lambda$.
\end{proof}

\begin{proposition}[Norm of a single lift]\label{pc:single}
For $1\le j\le m$ there is $A_j\in\C^\times$ such that for every twist
$\vartheta$
\[
\Theta_{j,\vartheta}^*\Theta_{j,\vartheta}=A_j\,c_Q(\vartheta)\,\pcRal_j(\vartheta)\ \text{ on }
e_0\mathcal A(W_j),\qquad c_Q(\vartheta)=(\auNm Q_\vartheta)^{rt}\prod_{w\in Q_\vartheta}\
\prod_{i=t+1}^{t+r}\bigl(1-q_w^{-i}\bigr).
\]
Consequently $\Theta_{j,\vartheta}$ is injective, and
$Z_j(\vartheta)=\bigl(A_jc_Q(\vartheta)\bigr)^{d_j}\det\bigl(\pcRal_j(\vartheta)\,|\,e_0\mathcal A
(W_j)\bigr)>0$, $d_j=\dim e_0\mathcal A(W_j)$.
\end{proposition}
\begin{proof}
The function $\Psi^{\mathrm{bl}}_\vartheta$ is sharp for $J=J_M$ (trivial reduction at the chosen
places, all of $\auGL_r(\auF_w)$ at $w\in Q_\vartheta$): the first condition holds at $w\in Q$
because the tame part of $\alpha$ at $w$ is $\psi_w$ and
$\psi_w(\det\bar k)^{-1}\psi_w(\det(x\bar k)^{\mathrm{top}})=\psi_w(\det x^{\mathrm{top}})$; for the
second take $Ax=x^{\mathrm{top}}$. Write $g_N=\#\auGL_N(\auF_w)$ and $q=q_w$. Then
$\|\Psi^{\mathrm{bl}}_\vartheta\|^2=\prod_{w\in Q}g_rq^{rt}$, $[K_{W_j}:J_M]=\prod_{w\text{ chosen}}
g_r$, and Proposition~\ref{pc:rallis} gives, for $y,y'\in\mathcal A(W_j)$,
\[
\langle\Theta y',\Theta y\rangle_{[\auU(V_j)]}=A_\infty\prod_{w\text{ chosen}}\frac1{g_rq^{(r+t)r}}
\prod_{w\in Q}\frac{g_rq^{rt}}{q^{(r+t)r}}\cdot\langle y',\pcRal_j(\vartheta)y\rangle .
\]
By Lemma~\ref{pc:algpet} the inner product in $\mathcal A(V_j;\vartheta)$ is this times a constant
and $[K_{V_j}:J^V_\vartheta]$, which is a factor independent of $\vartheta$ at the chosen places
times $\prod_{w\in Q}g_{r+t}/(g_rg_tq^{rt})$, the index of a parabolic subgroup with blocks $r,t$.
At $w\in Q$ the product of the factors is
$g_{r+t}/(g_tq^{(r+t)r})=q^{rt}\prod_{i=t+1}^{t+r}(1-q^{-i})$. The operator $\pcRal_j(\vartheta)$
preserves $e_0\mathcal A(W_j)$ (Lemma~\ref{pc:heckeclosure}); this gives the identity. Injectivity
is that of Proposition~\ref{pc:rallis}, and the formula for $Z_j$ is Lemma~\ref{pc:vol}; $Z_j>0$
since $\Theta_{j,\vartheta}$ is injective.
\end{proof}

\subsubsection*{Inner forms}
{\cite[Proposition~14.36]{P46}} is used in the following form.
\begin{quote}
$(\mathrm M^+)$\quad For $1\le j\le m$ and every $h\in\pcH$,\quad
$\pctr\bigl(h\,|\,e_0\mathcal A(W_j)\bigr)=\pctr\bigl(h\,|\,\auT\bigr)$.
\end{quote}
With $h=1$ it gives $\dim e_0\mathcal A(W_j)=\audT$. It involves all places prime to $\pcm_\vartheta$ (as the series $\pcRal_j(\vartheta)$ does), not
only the split ones; this is the statement of {\cite[Proposition~14.36]{P46}}, proved in
{\cite[Proposition~4.6]{P51}} and {\cite[Remark~4.22]{P51}}.

\begin{lemma}\label{pc:detcompare}
Assume $(\mathrm M^+)$. Then for $1\le j\le m$ and every twist $\vartheta$,
$\det\bigl(\pcRal_j(\vartheta)\,|\,e_0\mathcal A(W_j)\bigr)=\det\bigl(\pcRal_j(\vartheta)\,|\,\auT
\bigr)$.
\end{lemma}
\begin{proof}
By Lemma~\ref{pc:heckeclosure} there is $h\in\pcH$ acting as $\pcRal_j(\vartheta)$ on $\mathcal A
(W_j)$ and on $\auAz$, hence on the images of $\varrho(e_0)$. By $(\mathrm M^+)$ for the powers
$h^k$, $k\ge0$, the two restrictions have the same dimension and the same power sums of
eigenvalues, hence (Newton's identities) the same characteristic polynomial.
\end{proof}


\subsection{The triangle of lifts; proofs of \texorpdfstring{{\cite[Propositions~14.35 and~14.37]{P46}}}{[Row D, Propositions~14.35 and~14.37]}}
\label{pc:ss-triangle}

\subsubsection*{The spaces and the lifts}
For $i,n\ge0$ with $i+n\le m$ let $\pcG_{i,n}$ be the following space of rank $r+nt$:
$\pcG_{0,n}=Y_n$ with the lattice of {\cite[\S5.1.3]{P49}}; $\pcG_{i,0}$ is $Y_0$ with its form multiplied
by $(-1)^i$; and for $i,n\ge1$, $\pcG_{i,n}=E^{r+nt}$ with the standard lattice and the diagonal
form with entries $\epsilon_1,\dots,\epsilon_r,1,\dots,1$, where
$\epsilon_p=(-1)^{i+a+\min((a-i)^+,\,n)}$, $a=(p-k)^+$. (For $i=0$ these are the signs of $Y_n$.)
Thus $\pcG_{i,n}$ has a \emph{stem} of rank $r$ and $n$ positive definite \emph{blocks} of rank
$t$, numbered $i+1,\dots,i+n$; at most $r$ of its signs are negative. Let $e_{i,n}$ be the type of
the chain at the node $(i,n)$: for $i=0$ the type of the lower chain on $Y_n$, in general the type
whose Harish-Chandra parameter at $\tau$ consists of $\varsigma_\tau$ on the stem and, for the
block $i+l$, of the $t$ numbers $-\kappa_{i+l}(\tau)+\tfrac{t-1}2-b$, $0\le b<t$, all lowered by
$\tfrac12$ if $n$ is odd; and put $\pcF_{i,n}=\pcF(\pcG_{i,n},e_{i,n})$. The groups of
$\pcG_{i,0}$ and of $Y_0$ coincide, and so do the types, so that $\pcF_{i,0}=\pcF_{0,0}\supseteq
\auAz$; let $j_i\colon\auAz\to\pcF_{i,0}$ be the inclusion.

For an $m$-fold twist $\vartheta$ put $\mu_j=\mu_j^0\vartheta^j$,
$w_n^\vartheta=\#\auGL_{r+nt}(\pcR_\vartheta)$, and for $i+n<m$
\[
a_{i,n}^\vartheta=\Theta_{\alpha,\beta}(P_{i,n},\delta_\iota)\colon\pcF_{i,n}\to\pcF_{i,n+1},\qquad
\alpha=\mu_{i+n+1}\upsilon_{r+(n+1)t},\ \ \beta=\mu_{i+n+1}\upsilon_{r+nt},
\]
the lift for the pair $(\pcG_{i,n},\pcG_{i,n+1})$ with the modulus $\pcm_\vartheta$, the point datum
at the inclusion matrix and the joint harmonic $P_{i,n}$ of the two types. For $i=0$ this is
the lift of the lower chain without its normalisation: $[K_n:J_n]=w_n^\vartheta$ and
\begin{equation}\label{pc:eq-iota}
\Theta^Y_{n,\vartheta}=w_n^\vartheta\,a_{0,n}^\vartheta,\qquad\iota_\vartheta=X^\vartheta_{0,m}
\circ j_0 .
\end{equation}
We use the following property of the data ({\cite[\S14.3]{P46}}):
$\mu_j^0=\mu^\flat\cdot\nu_1^{\,j}$ with $\nu_1=\upsilon^{-\Delta}\chi_M$ on $\auU(1)$, where
$\chi_M$ has weight $0$, is unramified outside the chosen places, and its tame parts at the chosen
places have order greater than $2m$. Put $\nu_n=\nu_1^{\,n}$, of weight $-n\Delta$. For
$1\le z\le t$ and $1\le n\le 2m$ the pair $(\nu_n,z)$ is a ramified critical pair in the sense of
{\cite[\S14.10]{P46}} when $n\Delta\le4b$, which holds for $n\le m$. For $0<i<j\le m$ let
\begin{equation}\label{pc:eq-rho}
\rho_{ij}(\vartheta)=\prod_{z=1}^t\frac{|L(z,(\nu_{j-i}\vartheta^{\,j-i})^{-1})|}{|L(z,\nu_{j-i}^{-1})
|},
\end{equation}
a ramified ratio; the $L$-functions are those of {\cite[\S14.10]{P46}}, without the Euler factors at
the primes dividing $\pcm_\vartheta\aubar{\pcm_\vartheta}$, respectively $M$. We write
$L_*(\pcm,\nu,e)=(\prod_{w\mid\pcm}q_w)^{e^2}\prod_{z=1}^e|L^{\pcm}(z,\nu)|^2$ (equal to
$(\auNm\pcm)^{e^2}\prod_{z=1}^e|L^{\pcm}(z,\nu)|^2$ for squarefree $\pcm$), with $L^\pcm$ the Hecke
$L$-function of $x\mapsto\nu(x/\aubar x)$ without the factors at $\pcm\aubar\pcm$.

\begin{theorem}[The lifts form a triangle]\label{pc:tri}
The maps $a^\vartheta_{i,n}$ are well defined and have adjoints, and
there are linear maps $d^\vartheta_{i,n}\colon\pcF_{i,n+1}\to\pcF_{i+1,n}$ $(i+n<m)$, for all
$m$-fold twists $\vartheta$, such that:
\begin{enumerate}
\item[(a)] $d^\vartheta_{i,0}$ is the adjoint of $a^\vartheta_{i,0}$, under $\pcF_{i+1,0}=\pcF_{i,0}$;
\item[(b)] for $i+n+2\le m$ there are $\kappa^\sharp_{i,n}\in\C$ and $C_{i,n}>0$, independent of
$\vartheta$, and numbers $c_{i,n}(\vartheta)\ne0$ with
\[
w_n\,a_{i+1,n}d_{i,n}=\frac{\kappa^\sharp_{i,n}}{c_{i,n}(\vartheta)}\,w_{n+2}\,d_{i,n+1}a_{i,n+1},
\qquad|c_{i,n}(\vartheta)|^2=C_{i,n}\,L_*\bigl(\pcm_\vartheta,(\nu_{n+1}\vartheta^{\,n+1})^{-1},t
\bigr);
\]
\item[(c)] for $i+n+2\le m$ the adjoint law \textup{(T4)} holds with functions $\kappa_{i,n}$ that
have modulus $1$ up to norm powers.
\end{enumerate}
\end{theorem}
\noindent Lemma~\ref{pc:arch} holds: it is Proposition~\ref{ptv:arch}. The maps are well defined with adjoints (Theorem~\ref{ptr:edge}(d), Lemma~\ref{ptv:adj});
(a) is Proposition~\ref{ptv:tri}(ii); (b) is Corollary~\ref{ptr:tri-b} with Theorem~\ref{ptq:H} and
Proposition~\ref{ptv:chars}; (c) is Theorem~\ref{ptv:adjoint}, which uses the square law (RS)
of the reversed squares, Theorem~\ref{ptw:RS} (Corollary~\ref{ptw:adjoint}).

\begin{proposition}[First lifts]\label{pc:first}
For $0\le i<m$ there is $C_i\in\C^\times$ such that for every $m$-fold
twist $\vartheta$ and all $x,y\in\auAz$
\[
w_0^\vartheta\,\langle a^\vartheta_{i,0}x,a^\vartheta_{i,0}y\rangle=\frac{C_i}{\#\pcMat_{(r+t)
\times r}(\pcR_\vartheta)}\,\bigl\langle x,\pcRal_{i+1}(\vartheta^{\,i+1})\,y\bigr\rangle ,
\]
and $a^\vartheta_{i,0}$ is injective on $\auAz$. Moreover, for every $u\in J_M$ there are
$g\in K_{\pcG_{i,1}}$ and $c\in\C$, $|c|=1$, with $a_{i,0}\rho(u)=c\,\rho(g)a_{i,0}$.
\end{proposition}
\begin{proof}
We use Lemma~\ref{pc:arch}, which is Proposition~\ref{ptv:arch}.
The point function $\delta_\iota$ on $\pcMat_{(r+t)\times r}(\pcR_\vartheta)$ is sharp for
$J=J_{\pcG_{i,0}}(\pcm_\vartheta)$, with $A$ the projection to the first $r$ rows: $A\iota=1$. The
elements of $\auAz$ are invariant under $J_M\supseteq J$, where $J_M$ is trivial at the chosen
places and maximal elsewhere. Here $\|\delta_\iota\|^2=1$ and $[K:J]=w_0^\vartheta$, the character
is $\alpha=\mu_{i+1}\upsilon_{r+t}=\mu^0_{i+1}\vartheta^{\,i+1}\upsilon_{r+t}$, and the modulus of
$\vartheta^{\,i+1}$ is $\pcm_\vartheta$; so Proposition~\ref{pc:rallis} gives the formula, with
$C_i=A_\infty$, and the injectivity. (By Lemma~\ref{pc:series}(a) the series is the one of
\eqref{pc:eq-Rj}, although $\pcG_{i,0}$ and $\pcG_{i,1}$ are not $W_{i+1}$ and $V_{i+1}$.) For the
last assertion choose $g$ with reduction $\operatorname{diag}(\bar u,1)$ modulo $\pcm_\vartheta$;
then $\bar g\iota=\iota\bar u$, and by Lemma~\ref{pc:level} the pair $(u^{-1},g^{-1})$ multiplies
the datum $[P,\delta_\iota]$ by a root of unity $c$. Hence
$\theta(g'u^{-1},h;P,\delta_\iota)=c\,\theta(g',hg;P,\delta_\iota)$, and the change of variable
$g'\mapsto g'u$ in the integral defining $a_{i,0}\rho(u)f$ gives the claim.
\end{proof}

\begin{corollary}\label{pc:istriangle}
Let $\auAz\ne0$. Then the spaces
$\pcF_{i,n}$, the maps $j_i$, $a^\vartheta_{i,n}$, $d^\vartheta_{i,n}$ and the weights
$w^\vartheta_n$ form a triangle of size $m$ over the set of $m$-fold twists, with $V=\auAz$ and
with the adjoint law. In it $\sigma_{i,n}=\kappa^\sharp_{i,n}/c_{i,n}$ has no zero,
$\kappa^\sharp_{i,n}\ne0$, and $\beta_{i,n}=1/\aubar{\sigma_{i,n}}$ has modulus $\rho_{i+1,\,i+n+2}$
up to norm powers.
\end{corollary}
\begin{proof}
(T1) and (T4) are (b) and (c). (T2) is (a) restricted to $j_i(\auAz)$. For (T3) let $x\in\auAz$ and
$z=d_{i,0}a_{i,0}x\in\pcF_{i+1,0}=\pcF_{i,0}$. By (a), $\langle z,f\rangle=\langle a_{i,0}x,a_{i,0}f
\rangle$ for all $f\in\pcF_{i,0}$. For $u\in J_M$ let $g,c$ be as in Proposition~\ref{pc:first}
for $u^{-1}$, so that $a_{i,0}\rho(u^{-1})=c\rho(g)a_{i,0}$ and hence
$\rho(g^{-1})a_{i,0}=c\,a_{i,0}\rho(u)$. As $\rho$ is unitary and $\rho(u)x=x$,
\[
\langle\rho(u)z,f\rangle=\langle a_{i,0}x,a_{i,0}\rho(u^{-1})f\rangle=c\langle\rho(g^{-1})a_{i,0}x,
a_{i,0}f\rangle=c\,\aubar c\,\langle a_{i,0}x,a_{i,0}f\rangle=\langle z,f\rangle .
\]
So $z\in\pcF_{i+1,0}^{J_M}=j_{i+1}(\auAz)$. If $z=0$ then $\langle a_{i,0}x,a_{i,0}x\rangle=0$ and
$x=0$ by Proposition~\ref{pc:first}. Thus we have a triangle, and Proposition~\ref{pc:tri-inj}
shows that $\sigma_{i,n}$ has no zero. Finally $|\beta_{i,n}|^2=|c_{i,n}|^2/|\kappa^\sharp_{i,n}|^2$,
and
\[
L_*\bigl(\pcm_\vartheta,(\nu_{n+1}\vartheta^{\,n+1})^{-1},t\bigr)=M^{ft^2}(\auNm Q_\vartheta)^{t^2}
\prod_{z=1}^t|L(z,\nu_{n+1}^{-1})|^2\cdot\rho_{i+1,\,i+n+2}(\vartheta)^2
\]
by \eqref{pc:eq-rho}, because $\auNm\pcm_\vartheta=M^f\auNm Q_\vartheta$; the values
$L(z,\nu_{n+1}^{-1})$ do not depend on $\vartheta$ and are not zero since $c_{i,n}(1)\ne0$.
\end{proof}

\begin{proof}[Proof of {\cite[Proposition~14.35]{P46}} from Lemma~\ref{pc:arch} and
Theorem~\ref{pc:tri}\textup{(a)}, \textup{(b)}]
We may assume $\auAz\ne0$. Only (T1)--(T3) are needed, that is, parts (a) and (b) of
Theorem~\ref{pc:tri} (in (b) only the existence of the constants) and Proposition~\ref{pc:first}.
By Proposition~\ref{pc:tri-inj} with $(i,n)=(0,m)$ the map $X^\vartheta_{0,m}\circ j_0$ is
injective on $\auAz$, and it is $\iota_\vartheta$ by \eqref{pc:eq-iota}.
\end{proof}

\begin{lemma}[Indices]\label{pc:index}
Let $J_\vartheta\subseteq K_m$ be the level of $\auAm(\vartheta)$ and let
$e_m=mrt+\tfrac12m(m-1)t^2$. There is a constant $C_M>0$, depending only on the chosen places, such
that for every $m$-fold twist $\vartheta$
\[
[K_m:J_\vartheta]\cdot\frac{w_0^\vartheta\,(w_1^\vartheta)^m}{w_m^\vartheta}\cdot\frac1{\bigl(\#
\pcMat_{(r+t)\times r}(\pcR_\vartheta)\bigr)^m}=C_M\,(\auNm Q_\vartheta)^{-e_m}\,c_Q(\vartheta)^m .
\]
\end{lemma}
\begin{proof}
All four quantities are products over the primes $w\mid\pcm_\vartheta$, and the factors at the
chosen places do not depend on $\vartheta$. Let $w\in Q_\vartheta$, $q=q_w$, $g_N=\#\auGL_N(
\auF_w)$. The level $J_\vartheta$ is block upper triangular at $w$, with blocks of sizes
$r,t,\dots,t$ ({\cite[\S14.3]{P46}}, {\cite[Lemma~8.5]{P49}}; it is not the principal congruence
subgroup $J_m$ of {\cite[\S5.1.3]{P49}}, see Lemma~\ref{pgr:jcirc}), so its index is $g_{r+mt}/(g_rg_t^{\,m}q^{e_m})$, where $e_m$ is the dimension of the
unipotent radical. The factor at $w$ is therefore
\[
\frac{g_{r+mt}}{g_rg_t^{\,m}q^{e_m}}\cdot\frac{g_r\,g_{r+t}^{\,m}}{g_{r+mt}}\cdot q^{-(r+t)rm}=
q^{-e_m}\Bigl(\frac{g_{r+t}}{g_t\,q^{(r+t)r}}\Bigr)^m=q^{-e_m}\Bigl(q^{rt}\prod_{i=t+1}^{t+r}(1-q^{-i})
\Bigr)^m .\qedhere
\]
\end{proof}

\begin{proof}[Proof of {\cite[Proposition~14.37]{P46}} from Lemma~\ref{pc:arch},
Theorem~\ref{pc:tri} and $(\mathrm M^+)$]
The numbers $Z_j(\vartheta)$ are those defined before Lemma~\ref{pc:algpet}; they are positive by
Proposition~\ref{pc:single}. If $\audT=0$ there is nothing to prove; let $\audT>0$.

\emph{Step 1: the Gram form of the chain.} By Corollary~\ref{pc:istriangle} and
Proposition~\ref{pc:first}, Proposition~\ref{pc:tri-gram} applies with
$S_i=\pcRal_{i+1}(\vartheta^{\,i+1})$ and
$C_i(\vartheta)=C_i/\#\pcMat_{(r+t)\times r}(\pcR_\vartheta)$. With \eqref{pc:eq-iota}, for
$x,y\in\auAz$,
\[
\langle\iota_\vartheta x,\iota_\vartheta y\rangle_{[\auU(Y_m)]}=\frac{w_0w_1^{\,m}}{w_m}\cdot
\frac{\prod_{i<m}C_i}{(\#\pcMat_{(r+t)\times r}(\pcR_\vartheta))^m}\cdot\varkappa(\vartheta)\cdot\bigl
\langle x,\pcRal_1(\vartheta)\pcRal_2(\vartheta^2)\cdots\pcRal_m(\vartheta^m)\,y\bigr\rangle,
\]
where $\varkappa=\prod_{i+n+2\le m}\kappa_{i,n}\beta_{i,n}$. By Lemma~\ref{pc:mod}(a),
Theorem~\ref{pc:tri}(c) and Corollary~\ref{pc:istriangle}, $\varkappa$ has modulus
$\prod_{0<i<j\le m}\rho_{ij}$ up to norm powers: the square $(i,n)$ contributes
$\rho_{i+1,i+n+2}$, and $(i,n)\mapsto(i+1,i+n+2)$ is a bijection onto the pairs $0<i'<j'\le m$.

\emph{Step 2: the inner product of $\auAm(\vartheta)$.} By Lemma~\ref{pc:algpet} for the definite
space $Y_m$ and the level $J_\vartheta$, the inner product of $\auAm(\vartheta)$ is
$c_\infty[K_m:J_\vartheta]$ times the Petersson product, with $c_\infty>0$ independent of
$\vartheta$. With Lemma~\ref{pc:index}, for $x,y\in\auT$,
\[
\langle\Theta_\vartheta x,\Theta_\vartheta y\rangle=\Lambda(\vartheta)\,c_Q(\vartheta)^m\,\bigl
\langle x,\pcRal_1(\vartheta)\cdots\pcRal_m(\vartheta^m)\,y\bigr\rangle,\qquad
\Lambda(\vartheta)=c_\infty C_M\prod_{i<m}C_i\cdot(\auNm Q_\vartheta)^{-e_m}\varkappa(\vartheta).
\]
The operators $\pcRal_j(\vartheta^j)$ preserve $\auT$ (Lemma~\ref{pc:heckeclosure}). By
Lemma~\ref{pc:vol},
\[
Z(\vartheta)=\bigl(\Lambda(\vartheta)c_Q(\vartheta)^m\bigr)^{\audT}\prod_{j=1}^m\det\bigl(\pcRal_j
(\vartheta^j)\,|\,\auT\bigr).
\]

\emph{Step 3: the single lifts.} By $(\mathrm M^+)$, $d_j=\audT$. By Proposition~\ref{pc:single}
for the twist $\vartheta^j$, whose set of places is $Q_\vartheta$, and by
Lemma~\ref{pc:detcompare}, $Z_j(\vartheta^j)=(A_jc_Q(\vartheta))^{\audT}\det(\pcRal_j(\vartheta^j)\,|\,
\auT)$. Dividing,
\[
\frac{Z(\vartheta)}{\prod_jZ_j(\vartheta^j)}=\Bigl(\frac{\Lambda(\vartheta)}{\prod_jA_j}\Bigr)^{\audT}
\]
for every $m$-fold $\vartheta$; the left side is positive, so $\Lambda$ has no zero and the left
side is $|\Lambda(\vartheta)/\prod_jA_j|^{\audT}$.

\emph{Step 4: the modulus.} The function $\Lambda$ has modulus $\prod_{i<j}\rho_{ij}$ up to norm
powers (Step 1 and Lemma~\ref{pc:mod}(a)), this product is non-negative and equal to $1$ at
$\vartheta=1$, and $\Lambda(1)\ne0$. By Lemma~\ref{pc:mod}(b) there is $e\in\Z$ with
$|\Lambda(\vartheta)/\Lambda(1)|=(\auNm Q_\vartheta)^{e/2}\prod_{i<j}\rho_{ij}(\vartheta)$. Dividing
the identity of Step 3 by its case $\vartheta=1$ and inverting gives
\[
\frac{Z(1)}{Z(\vartheta)}=\Bigl((\auNm Q_\vartheta)^{e/2}\prod_{0<i<j\le m}\rho_{ij}(\vartheta)
\Bigr)^{-\audT}\prod_{j=1}^m\frac{Z_j(1)}{Z_j(\vartheta^j)} .\qedhere
\]
\end{proof}


\subsection{The squares: matched functions and the exchange identity}\label{pc:ss-exchange}

This subsection contains the algebra of Theorem~\ref{pc:tri} at the places of the modulus, the
statement of the exchange identity on which the rest depends (proved in \S\ref{pcx:sec}), and an
outline of the deduction of Theorem~\ref{pc:tri}, which is carried out in
\S\S\ref{ptr:sec}--\ref{ptw:sec}.

\subsubsection*{Matched functions over a finite ring}
Let $\pcR$ be a finite commutative ring and $\chi$ a character of $\pcR^\times$, extended by $0$.
Let the finite set $\mathsf N$ be the disjoint union of $\mathsf B$ and $\mathsf A$, of
cardinalities $b$ and $a$. For $k\in\auGL_{\mathsf N}(\pcR)$ let $y_k\in\pcMat_{\mathsf N\times
\mathsf B}(\pcR)$ be the block of the columns $\mathsf B$ of $k$, and $x_k\in\pcMat_{\mathsf A
\times\mathsf N}(\pcR)$ the block of the rows $\mathsf A$ of $k^{-1}$; so $x_ky_k=0$. Functions
$\Phi$ on $\pcMat_{\mathsf A\times\mathsf N}(\pcR)$ and $\Psi$ on $\pcMat_{\mathsf N\times\mathsf B}
(\pcR)$ are \emph{matched} (for $\chi$) if $\Phi$ vanishes off the $x_k$, $\Psi$ vanishes off the
$y_k$, and $\Phi(x_k)=\chi(\det k)\Psi(y_k)$ for all $k$. Functions $L$ on $\pcMat_{\mathsf A\times
\mathsf N}(\pcR)$ and $L'$ on $\pcMat_{\mathsf N\times\mathsf B}(\pcR)$ \emph{satisfy the square law
with constant $c$} if $L(hx)=\chi(\det h)L(x)$ for $h\in\auGL_{\mathsf A}(\pcR)$,
$L'(yh)=\chi(\det h)L'(y)$ for $h\in\auGL_{\mathsf B}(\pcR)$, and
\[
\sum_y\Psi(y)L'(y)=c\sum_x\Phi(x)L(x)\qquad\text{for all matched }\Phi,\Psi .
\]

\begin{lemma}[The law at a point]\label{pc:pointlaw}
Let $L,L'$ satisfy the square law with constant $c$. For $w\in\pcMat_{\mathsf A\times\mathsf B}
(\pcR)$ let $\binom{1}{-w}\in\pcMat_{\mathsf N\times\mathsf B}$ have the blocks $1$ (rows
$\mathsf B$) and $-w$ (rows $\mathsf A$), and $(w\ 1)\in\pcMat_{\mathsf A\times\mathsf N}$ the blocks
$w$ (columns $\mathsf B$) and $1$ (columns $\mathsf A$). Then
\[
\#\auGL_{\mathsf B}(\pcR)\cdot L'\tbinom{1}{-w}=c\cdot\#\auGL_{\mathsf A}(\pcR)\cdot L(w\ 1).
\]
\end{lemma}
\begin{proof}
Let $k_w=\bigl(\begin{smallmatrix}1&0\\-w&1\end{smallmatrix}\bigr)$, so that
$y_{k_w}=\binom1{-w}$, $x_{k_w}=(w\ 1)$ and $\det k_w=1$. Define $\Psi_w(y_{k_w}h)=\chi(\det h)^{-1}$
for $h\in\auGL_{\mathsf B}(\pcR)$ and $\Psi_w=0$ elsewhere, $\Phi_w(gx_{k_w})=\chi(\det g)^{-1}$ for
$g\in\auGL_{\mathsf A}(\pcR)$ and $\Phi_w=0$ elsewhere. They are matched: if $y_k=y_{k_w}h$, then
$k=k_w\bigl(\begin{smallmatrix}h&B\\0&g\end{smallmatrix}\bigr)$ for some $B$ and some
$g\in\auGL_{\mathsf A}(\pcR)$, so $\det k=\det h\det g$ and $x_k=g^{-1}x_{k_w}$; then
$\Phi_w(x_k)=\chi(\det g)=\chi(\det k)\Psi_w(y_k)$. Conversely, $gx_{k_w}=x_k$ for
$k=k_w\bigl(\begin{smallmatrix}1&0\\0&g^{-1}\end{smallmatrix}\bigr)$, and if $x_k=gx_{k_w}$ for some
$k$, then the rows $\mathsf A$ of $k^{-1}k_w$ are $g(0\ 1)=(0\ g)$, so
$k^{-1}k_w=\bigl(\begin{smallmatrix}P_1&P_2\\0&g\end{smallmatrix}\bigr)$ with $\det P_1\det g=\det(k^{-1}k_w)$
a unit, hence $P_1\in\auGL_{\mathsf B}(\pcR)$; then $k=k_w\bigl(\begin{smallmatrix}P_1&P_2\\0&g
\end{smallmatrix}\bigr)^{-1}$ and $y_k=y_{k_w}P_1^{-1}$. So $\Phi_w$ vanishes off the $x_k$, $\Psi_w$
vanishes off the $y_k$, and $x_k\in\auGL_{\mathsf A}(\pcR)x_{k_w}$ if and only if $y_k\in
y_{k_w}\auGL_{\mathsf B}(\pcR)$; in particular $\Phi_w(x_k)=0=\Psi_w(y_k)$ for all other $k$. By the two transformation laws, $\sum_y\Psi_w(y)L'(y)=\#\auGL_{\mathsf B}(\pcR)L'(y_{k_w})$
and $\sum_x\Phi_w(x)L(x)=\#\auGL_{\mathsf A}(\pcR)L(x_{k_w})$.
\end{proof}

The next lemma is the form in which a square of the triangle is used. Let $\mathsf N_0$,
$\mathsf T$, $\mathsf T'$ be finite sets, $\mathsf N_1=\mathsf N_0\sqcup\mathsf T$ (the index set of
$U_1$), $\mathsf U_2=\mathsf N_0\sqcup\mathsf T'$ (that of $U_2$) and $\mathsf A=\mathsf N_1\sqcup
\mathsf T'$ (that of $A$); $\mathsf N_0$ is the index set of $A'$. Take $\mathsf N=\mathsf N_1
\sqcup\mathsf U_2$ and $\mathsf B=\mathsf N_0\subset\mathsf U_2$, so that $\mathsf N=\mathsf B\sqcup
\mathsf A$. For a function $p$ on $\pcMat_{\mathsf N_1\times\mathsf N_0}(\pcR)$ its \emph{extension}
$p^+$ is the function on $\pcMat_{\mathsf A\times\mathsf U_2}(\pcR)$ with
\[
p^+\begin{pmatrix}\eta&s\\0&1\end{pmatrix}=\frac{p(-\eta)}{\#\pcMat_{\mathsf N_1\times\mathsf T'}
(\pcR)}\qquad(\eta\in\pcMat_{\mathsf N_1\times\mathsf N_0},\ s\in\pcMat_{\mathsf N_1\times\mathsf
T'}),
\]
the rows being $\mathsf N_1,\mathsf T'$ and the columns $\mathsf N_0,\mathsf T'$, and $p^+=0$ at all
other matrices.

\begin{lemma}[Extension of a datum]\label{pc:next}
Let $L$ on $\pcMat_{\mathsf A\times\mathsf N}(\pcR)$ and $L'$ on $\pcMat_{\mathsf N\times\mathsf
N_0}(\pcR)$ satisfy the square law with constant $c$. Let $\iota'\in\pcMat_{\mathsf A\times\mathsf
N_1}$ and $\iota\in\pcMat_{\mathsf U_2\times\mathsf N_0}$ be the inclusion matrices. Then for every
$p$
\[
\#\auGL_{\mathsf N_0}(\pcR)\sum_{y\in\pcMat_{\mathsf N_1\times\mathsf N_0}}p(y)\,L'\binom{y}{\iota}=
c\cdot\#\auGL_{\mathsf A}(\pcR)\sum_{x\in\pcMat_{\mathsf A\times\mathsf U_2}}p^+(x)\,L(\iota'\ x),
\]
where $\binom y\iota$ has the rows $\mathsf N_1$ and $\mathsf U_2$, and $(\iota'\ x)$ the columns
$\mathsf N_1$ and $\mathsf U_2$.
\end{lemma}
\begin{proof}
Write matrices with rows and columns in the order $\mathsf N_1,\mathsf N_0,\mathsf T'$. For
$\eta\in\pcMat_{\mathsf N_1\times\mathsf N_0}$ let $w_\eta\in\pcMat_{\mathsf A\times\mathsf N_0}$
have the blocks $\eta$ (rows $\mathsf N_1$) and $0$ (rows $\mathsf T'$). Then
$\binom1{-w_\eta}=\binom{-\eta}{\iota}$, and $(w_\eta\ 1)=\bigl(\begin{smallmatrix}1&\eta&0\\0&0&1
\end{smallmatrix}\bigr)$. For $s\in\pcMat_{\mathsf N_1\times\mathsf T'}$ the matrix
$u_s=\bigl(\begin{smallmatrix}1&s\\0&1\end{smallmatrix}\bigr)\in\auGL_{\mathsf A}(\pcR)$ has
determinant $1$, and $u_s(w_\eta\ 1)=\bigl(\begin{smallmatrix}1&\eta&s\\0&0&1\end{smallmatrix}\bigr)
=(\iota'\ x)$ with $x=\bigl(\begin{smallmatrix}\eta&s\\0&1\end{smallmatrix}\bigr)$. So
$L(\iota'\ x)=L(w_\eta\ 1)$ for all $s$, and by the definition of $p^+$ and
Lemma~\ref{pc:pointlaw}
\[
c\,\#\auGL_{\mathsf A}\sum_xp^+(x)L(\iota'\ x)=c\,\#\auGL_{\mathsf A}\sum_\eta p(-\eta)L(w_\eta\ 1)=
\#\auGL_{\mathsf N_0}\sum_\eta p(-\eta)L'\tbinom{-\eta}{\iota}.\qedhere
\]
\end{proof}

\subsubsection*{The exchange identity}
\begin{theorem}[Exchange identity]\label{pc:X}
Let $W,V,V'$ be spaces of ranks $N=a+b$, $a$, $b$, with $N$ even and $a=b+2e$, $e\ge0$. Let
$q,q_V,q_{V'}$ be their numbers of negative signs, and assume $q=q_V+q_{V'}$ and $2q+2\le N$. Let
$\mu=(\mu_\tau)_{\tau\in\auPl}$ be integers with $N\le2\mu_\tau+2q$ for all $\tau$, and let $v$ be a
vector of the Fock space of $\auHom(V',W)\otimes\R$ all of whose moments $\int(1+|z|)^k|v(z)|\,dz$
are finite. Then there are $C>0$ and a vector $x$ of the Fock space of $\auHom(W,V)\otimes\R$ with
finite moments such that the following holds. Let $\pcm\ne(1)$ be a modulus, $\nu$ an automorphic
character of $\auU(1)$ of weight $\mu$, unramified outside $\pcm$ and tamely ramified with
non-trivial tame part at every prime of $\pcm$, and $\chi$ the character of $\pcR_\pcm^\times$
given by its tame parts. Then there is $c\in\C^\times$ with
\[
|c|^2=C\cdot L_*(\pcm,\nu,e)\ne0
\]
such that for all automorphic characters $\alpha,\alpha'$ of $\auU(1)$ with
$\alpha\alpha'=\gamma:=\nu^{-1}\upsilon^{N/2}$, all functions $\Phi$ on $\pcMat_{a\times N}(\pcR_\pcm)$
and $\Psi$ on $\pcMat_{N\times b}(\pcR_\pcm)$ that are matched for $\chi$, and all
$g\in\auU(W)(\auA)$,
\[
\int_{[\auU(V)]}\theta^{W,V}_{\alpha,\gamma}(g,h;x,\Phi)\,dh=c\int_{[\auU(V')]}\theta^{V',W}_{\gamma,
\alpha'}(h',g;v,\Psi)\,dh',
\]
both integrals being absolutely convergent.
\end{theorem}
\noindent Here the kernels are those of \S\ref{pc:ss-data}, in the convention
$\omega_{\alpha,\beta}(g,h)=\alpha(\det g)^{-1}\beta(\det h)\omega(g,h)$, with the data of
Definition~\ref{pgx:conv}. Theorem~\ref{pc:X} is Theorem~\ref{pcx:X}, proved in \S\S\ref{pcx:sec},
\ref{pcu:sec}, \ref{pct:sec}, \ref{pcn:sec}, with absolutely convergent integrals
(Corollary~\ref{pcx:conv}); the dictionary with the other splitting in use is the transposition of
indices \eqref{pcx:eq-dict}, without any factor, and does not change $c$. The proof unfolds a large
triple along an isotropic frame, builds real vectors with prescribed zero terms, removes hyperbolic
planes one at a time (each plane contributes a ratio of values of $L^\pcm(z,\nu)$ and a power of
$\prod_{w\mid\pcm}q_w$, which build up $L_*(\pcm,\nu,e)$), passes from standard data to all matched
functions by equivariance, and evaluates the constant in equal rank; see \S\ref{pcu:ss-final} for the
list of dependencies.

\subsubsection*{From the exchange identity to Theorem~\ref{pc:tri}}
Fix a square: $U_1=\pcG_{i,n+1}$, $U_2=\pcG_{i+1,n+1}$, $A=\pcG_{i,n+2}$, $A'=\pcG_{i+1,n}$, with
$i+n+2\le m$, and $W=U_1\oplus U_2$, of rank $N=2(r+(n+1)t)$. The hypotheses of
Theorem~\ref{pc:X} hold for $(W,V,V')=(W,A,A')$, with $e=t$ and the weight $\mu_\tau=(n+1)\Delta$ of
$(\nu_{n+1}\vartheta^{\,n+1})^{-1}$: the ranks are $a=r+(n+2)t$, $b=r+nt$; the numbers of negative
signs add up, because the signs $\epsilon_p$ of the stems satisfy
$\#\{\epsilon^{(i,n+1)}=-1\}+\#\{\epsilon^{(i+1,n+1)}=-1\}=\#\{\epsilon^{(i,n+2)}=-1\}+\#\{\epsilon^{
(i+1,n)}=-1\}$ (compare the exponents $i+a+\min((a-i)^+,n)$ entry by entry: for $a\le i$ both sides have the
exponents $\{i+a,\,i+a+1\}$; for $i<a\le i+n+1$ all four exponents are $2a$; for $a\ge i+n+2$ both
sides have $\{i+a+n+1,\,i+a+n+2\}$);
$2q+2\le N$ because $q\le2r<r+t$; and $N\le2(n+1)\Delta$ because $\Delta=2r^8$. The character
$(\nu_{n+1}\vartheta^{\,n+1})^{-1}$ is tamely ramified exactly at the primes of $\pcm_\vartheta$,
because $\vartheta$ is $m$-fold and the tame parts of $\chi_M$ have order greater than $2m$. The
deduction then has the following steps.
\begin{enumerate}
\item[(i)] \emph{Kernels multiply}. $\psHS(\auHom(W,A))=\psHS(\auHom(U_1,A))\otimes
\psHS(\auHom(U_2,A))$, the theta functional is the tensor product, and
$[P_1P_2,\delta_\xi]=[P_1,\delta_{\xi_1}]\otimes[P_2,\delta_{\xi_2}]$ for $\xi=(\xi_1\ \xi_2)$. Hence,
for forms $y$ on $U_1$ and $y_2$ on $U_2$, the pairing of $y_2$ with ``lift $U_1\to A$ with datum
$\xi_1$, then transposed lift $A\to U_2$ with datum $\xi_2$'' applied to $y$ is the integral over
$[\auU(U_1)]\times[\auU(U_2)]$ of $y\otimes y_2$ against the theta integral over $[\auU(A)]$ of the
datum $[P_1P_2,\delta_\xi]$ of the pair $(W,A)$; call it $L(\xi)$. In the same way the path through
$A'$ gives a function $L'(\eta)$ on $\pcMat_{\mathsf N\times\mathsf N_0}(\pcR_\vartheta)$. By
Lemma~\ref{pc:level}, $L$ and $L'$ transform by $\chi\circ\det$, provided the splitting characters
of the four kernels are chosen so that their products are $\gamma$ on $W$; this fixes the
characters of the descents (Definition~\ref{ptr:desc-char}).
\item[(ii)] \emph{Convergence} (Proposition~\ref{ptr:conv}). The integrals over $[\auU(A)]$ and $[\auU(A')]$,
whose groups are not compact, converge because the data vanish on the matrices that are not
injective, respectively surjective, at the primes of $\pcm$, and the forms decrease rapidly.
\item[(iii)] \emph{The infinite places} (Theorem~\ref{ptr:square}, with Theorem~\ref{ptq:H}). Theorem~\ref{pc:X}, applied with $v$ the
vector of the polynomial of the path through $A'$, gives a vector $x$ which is not known. One
needs: for every vector $x$ with finite moments there is $\kappa(x)\in\C$, independent of $\pcm$,
of the characters, of the forms and of $\xi$, such that the function $L_x$ formed with $x$ in
place of $P_1P_2$ equals $\kappa(x)L$. This is a statement of multiplicity one for the types
$e_{i,n+1}$, $e_{i+1,n+1}$ in the Fock space, extended from polynomials to vectors with finite
moments by continuity. With it,
$L,L'$ satisfy the square law with the constant $\kappa^\sharp/c(\vartheta)$,
$\kappa^\sharp=\kappa(x)$.
\item[(iv)] \emph{The descents}. Define $d_{i,n}$ as the transposed lift
along the pair $(\pcG_{i+1,n},\pcG_{i,n+1})$ with a finite datum $\varepsilon_{i,n}$, where
$\varepsilon_{i,0}$ is the datum for which $d_{i,0}=a_{i,0}^*$ (it exists by (W7): the complex
conjugate of the kernel of $a_{i,0}$ is a kernel for the space with the negative form) and
$\varepsilon_{i,n+1}=\varepsilon_{i,n}^+$. Then Lemma~\ref{pc:next}, with
$\#\auGL_{\mathsf N_0}=w_n$ and $\#\auGL_{\mathsf A}=w_{n+2}$, is exactly (b) of
Theorem~\ref{pc:tri}.
\item[(v)] \emph{The adjoint law} (Theorem~\ref{ptv:adjoint} with Theorem~\ref{ptw:RS}; the reversal of
a kernel is Proposition~\ref{ptv:rev}, the Gauss sum Proposition~\ref{ptv:fourier}). The adjoint of a descent is, up to a constant of
modulus $1$ up to norm powers and a translation by a central element, the transposed lift along
the reversed pair with the Fourier transform of the conjugate datum (this needs the reversal
$w\mapsto w^\dagger$ of a kernel, that is, the fact (W6) ({\cite[Proposition~7.17]{P49}}), and the Gauss sum of $\chi\circ\det$). It commutes with the lifts $a$ by a second
application of Theorem~\ref{pc:X}, now with $e=0$, where $L_*$ is $1$; so its constant has modulus
$1$ up to norm powers.
\end{enumerate}


\subsection{Auxiliary forms: proof of \texorpdfstring{{\cite[Proposition~14.38]{P46}}}{[Row D, Proposition~14.38]}}\label{pc:ss-aux}

For $0\le i,j<m$ let
\begin{equation}\label{pc:eq-rhoprime}
\rho'_{ij}(\vartheta)=\prod_{z=1}^t\frac{|L(z,(\mu^0_{i+1}\mu^0_{j+1}\vartheta^{\,j+1})^{-1})|}{|L(z,(
\mu^0_{i+1}\mu^0_{j+1})^{-1})|}.
\end{equation}
The character $\mu^0_{i+1}\mu^0_{j+1}$ has weight $-(\kappa_{i+1}+\kappa_{j+1})$, with
$t\le\kappa_{i+1}(\tau)+\kappa_{j+1}(\tau)\le4b$, and tame part $\chi_M^{\,i+j+2}$ at the chosen
places, which is non-trivial since $2\le i+j+2\le2m$; so $\rho'_{ij}$ is a ramified ratio.

Let $\widetilde Y$ be the positive definite standard space of rank $\tilde n=r+1+m(t-1)=n_m-m$,
let $\pcb=(\pi_M)^{b}$, and let $J(\pcb)\subset K_{\widetilde Y}$ be the principal congruence
subgroup of level $\pcb$; it is free by {\cite[Lemma~14.5]{P46}}, since $\tilde n\le2n_m$.

\begin{theorem}[Integral and short lifts through the second grid]\label{pc:auxcore}
There are a weight $\tilde\lambda$ of $\auU(\widetilde Y)$, a function $c$ on the $m$-fold twists
with values in $\C^\times$ that has modulus $\prod_{0\le i,j<m}\rho'_{ij}$ up to norm powers, a
surjective linear map $\mathsf r\colon\pcFF\to\auAz$ from
$\pcFF=\mathcal A_{\tilde\lambda}(J(\pcb),1)$, and $s\in\C^\times$, such that for every
$\varphi\in\pcFF$:
\begin{enumerate}
\item[(a)] if $\vartheta$ is an $m$-fold twist, $\iota\in\auI$ with $\ell(\iota)\nmid\auNm
Q_\vartheta$, and $\varphi$ is integral at $\iota$, then
$(c(\vartheta)/c(1))\cdot\iota_\vartheta(s\,\mathsf r\varphi)$ is integral at $\iota$;
\item[(b)] if $\sigma\in\auAut(\C)$ and $\|\sigma(\varphi(g))\|^2\le B$ for all $g$, then
$\|\sigma(\iota_1(s\,\mathsf r\varphi)(g))\|^2\le\exp(2x^{201}b^2)\,B$ for all $g$;
\item[(c)] $\pcFF$ is spanned by forms that are integral at every $\iota\in\auI$ and whose values
have all conjugates of squared norm at most $\exp(x^{99}b^2)$.
\end{enumerate}
\end{theorem}
\noindent Parts (a) and (b) are deduced in \S\ref{pgr:ss-deduce} from Propositions~\ref{pgr:first},
\ref{pgr:turn}, \ref{pgr:last}, \ref{pgr:den}, \ref{pgr:delta}, which are proved in
\S\S\ref{pgr:sec}--\ref{pgm:sec}; (c) is Proposition~\ref{pgr:short}, proved in
Proposition~\ref{pgs:short}. See \S\ref{pgr:ss-status}.

\begin{proof}[Proof of {\cite[Proposition~14.38]{P46}} from Theorem~\ref{pc:auxcore} and
{\cite[Proposition~14.28]{P46}}]
Put $\gamma(\vartheta)=c(\vartheta)/c(1)$. Then $\gamma(1)=1$, $\gamma$ has no zero, and by
Lemma~\ref{pc:mod}(b), applied to $c$ and $f=\prod_{i,j}\rho'_{ij}$ (which is non-negative and
equal to $1$ at $\vartheta=1$), there is $e'\in\Z$ with
$|\gamma(\vartheta)|=(\auNm Q_\vartheta)^{e'/2}\prod_{0\le i,j<m}\rho'_{ij}(\vartheta)$.

Let $\mathfrak B$ be the set of the forms $z=s\,\mathsf r\varphi\in\auAz$ with $\varphi$ as in
(c). It spans $\auAz$, because $\mathsf r$ is onto and $s\ne0$. For $z\in\mathfrak B$, by (a) the
form $\gamma(\vartheta)\iota_\vartheta z$ is integral at every $\iota$ with
$\ell(\iota)\nmid\auNm Q_\vartheta$, for every $m$-fold $\vartheta$; and by (b) every value of
$\sigma(\iota_1z)$ has squared norm at most $\exp(2x^{201}b^2+x^{99}b^2)$, for every $\sigma$.

Let $y=\#\auG\cdot\varrho(e_0)z\in\auT$. By {\cite[Proposition~14.28(a)]{P46}}, with $s=1$ and the
constant $\gamma(\vartheta)$, the form $\gamma(\vartheta)\Theta_\vartheta y=\gamma(\vartheta)
\iota_\vartheta y$ is integral at every $\iota$ with $\ell(\iota)\nmid\auNm Q_\vartheta$. By
{\cite[Proposition~14.28(b)]{P46}},
\[
\|\Theta_1y\|_\sigma^2\le\exp\bigl(x^{99}+2x^{201}b^2+x^{99}b^2\bigr)\le\exp(2x_1^2b^2)
\]
for all $\sigma$, because $x^{10^6}\le x_1$. Finally the forms $y$ span $\varrho(e_0)\auAz=\auT$.
\end{proof}

{\cite[Proposition~14.28]{P46}} is proved in {\cite[\S14.8]{P46}} from {\cite[Proposition~14.27]{P46}} and
from the bound $\exp(x^{30})$ for the number of classes of the level of $Y_m$
({\cite[Proposition~7.20]{P49}}).

\medskip\noindent\emph{Outline of the proof of Theorem~\ref{pc:auxcore}\textup{(a)}, \textup{(b)}}
(\S\S\ref{pgr:sec}--\ref{pgm:sec}). There are two grids of spaces, $Y^a_b$ and $\widetilde Y^b_a$
$(0\le a,b\le m)$, defined in Definition~\ref{pgv:def-spaces}: $Y^0_b=Y_b$, and $Y^a_b$ has $a$
further blocks of rank $t$; $\widetilde Y^0_0=Y_0\oplus E\tilde e$ with $\langle\tilde e,\tilde
e\rangle=1$, and $\widetilde Y^b_a$ has $a$ blocks of rank $t+1$ and $b$ blocks of rank $t-1$, so that
$\widetilde Y=\widetilde Y^m_0$; $Y^m_m$ and $\widetilde Y^m_m$ are positive definite, and
$\widetilde Y^l_l=Y^l_l\oplus E\tilde e_l$. The edges carry theta lifts with finite data modulo
$\pcb\prod_{w\in Q}w$. \emph{First path} (\S\ref{pgw:ss-first}): the transposed lifts down the first
column of the second grid, followed by restriction to $\auU(Y_0)$ at a real point and projection to
the type $e_0$, give the surjection $\mathsf r$; its surjectivity uses the injectivity of the
untwisted chain of the second grid (Theorem~\ref{pgy:I}) and a vanishing theorem at the deep level
$\pcb$ by Poincar\'e series (Theorem~\ref{pgq:V}). \emph{Turns} (\S\S\ref{pgr:ss-off},
\ref{pgu:sec}, \ref{pgv:ss-square}): the pairing of $\iota_\vartheta(\mathsf r\varphi)$ with a form on
$Y_m$ is a period along a path through the two grids; each square of a grid moves the path by
one square, at the cost of the constant of Theorem~\ref{pc:X} for that square, of modulus
$\rho'_{ij}$ up to norm powers; on the diagonal the turn uses the reversal of the theta kernel of a
line (Lemma~\ref{pgx:line}). \emph{Last path} (\S\S\ref{pgw:ss-last}, \ref{pgz:sec}, \ref{pgt:sec}):
all spaces are definite, the lifts are algebraic lifts by jets, whose kernels are integral up to a
bounded denominator, which gives (a); the size bounds give (b). Pairing with delta forms reads the
values (Proposition~\ref{pgr:delta}).

\subsection{Where the inputs of the Covolume Theorem are proved}\label{pc:ss-remains}

Proposition~\ref{pc:C} deduces {\cite[Theorem~14.15]{P46}} from {\cite[Propositions~14.35, 14.37, 14.38, 14.40 and~14.43]{P46}}. In this section, {\cite[Proposition~14.35]{P46}} is proved from
Lemma~\ref{pc:arch}, Theorem~\ref{pc:tri}(a),(b) and (SW); {\cite[Proposition~14.37]{P46}} from Lemma~\ref{pc:arch},
Theorem~\ref{pc:tri}, (SW) and $(\mathrm M^+)$; {\cite[Proposition~14.38]{P46}} from Theorem~\ref{pc:auxcore} and
{\cite[Proposition~14.28]{P46}}. Lemma~\ref{pc:arch} is Proposition~\ref{ptv:arch} (from (V), (D1)--(D3));
Theorem~\ref{pc:tri} is proved in \S\S\ref{ptr:sec}--\ref{ptw:sec} (from Theorem~\ref{pc:X}, (V), (D1)--(D5),
(W7)); Theorem~\ref{pc:X} is Theorem~\ref{pcx:X}, whose inputs are listed in \S\ref{pct:ss-status};
Theorem~\ref{pc:auxcore} is proved in \S\S\ref{pgr:sec}--\ref{pgm:sec} (\S\ref{pgr:ss-status}).
Proposition~\ref{pc:rallis} and Propositions~\ref{pc:single}, \ref{pc:first} use Weil's formula (SW), (W7) and
Lemma~\ref{pc:arch}; Lemma~\ref{pc:detcompare} uses $(\mathrm M^+)$; Lemma~\ref{pc:clean} uses (W2), (W5),
(W7); Lemmas~\ref{pc:level}, \ref{pc:series}, \ref{pc:heckeclosure} use (W4), (W5). {\cite[Propositions~14.40, 14.41 and~14.42]{P46}} and $(\mathrm M^+)$ are proved in the sections named in their proofs in
{\cite[\S14.10]{P46}}. The character of (W7) is computed in Lemma~\ref{ptw:W7}.

Two remarks on the statements of {\cite[\S14.10]{P46}}. First, $Z_j(\vartheta)$ in
{\cite[Proposition~14.37]{P46}} must be read with the lift $\Theta_{j,\vartheta}$ normalised as before
Lemma~\ref{pc:algpet} (block function, level $J^V_\vartheta$, inner product of algebraic modular
forms): another normalisation changes $Z_j(\vartheta)$ by the $\audT$-th power of a factor such as
$\prod_{w\in Q}\prod_i(1-q_w^{-i})$, which is not a power of $\auNm Q$ times ramified ratios. With
this normalisation the constant $c_Q(\vartheta)$ of Proposition~\ref{pc:single} cancels against the
index of the level of $Y_m$ (Lemma~\ref{pc:index}). Secondly, {\cite[Proposition~14.35]{P46}} is not proved through the non-vanishing of
$L$-values in the Rallis inner product formula for each of $m$ lifts: only the first lift of each row of the triangle is in Weil's range,
its injectivity comes from the fact that the series of Lemma~\ref{pc:series} is close to the
identity, and the injectivity of the chain follows from the commutation relations by
Proposition~\ref{pc:tri-inj}, without any information on the constants of the squares.



\section{The exchange identity}\label{C1:g2}
% ===== C-exchange.tex =====
% proofs/C-exchange.tex
% The exchange identity (Theorem pc:X of C-determinant.tex): statement in the convention of
% \S pc:ss-data, dictionary of splittings, non-vanishing of the constant, architecture of the proof,
% and its proof. Labels of this file start with pcx:.
% Macros: those of C-determinant.tex (\auU, \auA, \auNm, \pcm, \pcR, \pcMat, \psvac, ...).
\providecommand{\pcxH}{\mathbb H}
\providecommand{\pcxstr}{\mathrm{r}}
\providecommand{\pcxc}{\mathfrak c}
\providecommand{\pctcone}{\operatorname{cone}}

\subsection{The exchange identity}\label{pcx:sec}

This section states the exchange identity and proves it, together with
\S\S\ref{pcu:sec}, \ref{pct:sec}, \ref{pcn:sec} (see \S\ref{pcx:ss-status}); it is
Theorem~\ref{pc:X}, on which Theorems~\ref{pc:tri} and~\ref{pc:auxcore} and
{\cite[Proposition~2.15]{P54}} rest. Notation is that of \S\ref{pc:ss-data}.

\subsubsection{Statement, splitting and non-vanishing}\label{pcx:ss-statement}

\paragraph*{Conventions}
For spaces $S,T$ and automorphic characters $\alpha,\beta$ of $\auU(1)$ we use, as in
\S\ref{pc:ss-data},
\[
\omega_{\alpha,\beta}(g,h)=\alpha(\det g)^{-1}\beta(\det h)\,\omega_{\auHom(S,T)}(g,h),\qquad
\theta^{S,T}_{\alpha,\beta}(g,h;f)=\theta\bigl(\omega_{\alpha,\beta}(g,h)f\bigr),
\]
$g\in\auU(S)(\auA)$, $h\in\auU(T)(\auA)$, where $\omega_{\auHom(S,T)}(g,h)$ is the Weil representation
of $\auU(\auHom(S,T))(\auA)$, defined with $\upsilon$, composed with $x\mapsto hxg^{-1}$
({\cite[\S5.1.1]{P49}}). \emph{The first index is the character of the source group $\auU(S)$, entering
through its inverse; the second is the character of the target group $\auU(T)$.} We write
$\upsilon$ also for $\upsilon|_{\auU(1)}$.

Let $\pcm$ be a modulus and $S=S_\pcm$ the set of primes $w$ of $E$ dividing $\pcm$. A
\emph{tame character for $\pcm$ of weight $\mu$} is an automorphic character $\nu$ of $\auU(1)$ of
weight $\mu=(\mu_\tau)$ with all $\mu_\tau\ge1$, unramified outside $\pcm$, tamely ramified with
non-trivial tame part at every $w\in S$. Put $F_S=\prod_{w\in S}E_w$, $\auO_S=\prod_{w\in
S}\auO_{E_w}$, and let $\chi_S\colon F_S^\times\to\C^\times$, $t\mapsto\nu(t/\bar t)$ ($t$ read as an
idele supported on $S$), extended by $0$ to the non-invertible elements; on $\auO_S^\times$ it
factors through $\pcR_\pcm^\times$ and gives the character $\chi$ of the residue ring given by the
tame parts. Write $q_w=\#\auF_w$ and, for $e\in\Z_{\ge0}$,
\begin{equation}\label{pcx:eq-Lstar}
L_*(\pcm,\nu,e)=\Bigl(\prod_{w\in S}q_w\Bigr)^{e^2}\prod_{z=1}^e|L^S(z,\nu)|^2,\qquad
L_*(\pcm,\nu,-e)=L_*(\pcm,\nu,e)^{-1},
\end{equation}
where $L^S(s,\nu)$ is the Hecke $L$-function of $x\mapsto\nu(x/\bar x)$ with the Euler factors at
$S\cup\bar S$ removed. For squarefree $\pcm$ (e.g.\ $\pcm_\vartheta$; in general only the primes of $\pcm$ enter) $\prod_{w\in S}q_w=
\auNm\pcm$, and \eqref{pcx:eq-Lstar} is the $L_*$ of \S\ref{pc:ss-data}.

\begin{definition}[Matched functions]\label{pcx:def-matched}
Let $R$ be a commutative ring, $\chi$ a character of $R^\times$, $N=b+a$, the indices of $R^N$
ordered so that the first $b$ come first. For $k\in\auGL_N(R)$ let $k_{[b]}\in\pcMat_{N\times b}(R)$
be its first $b$ columns and $k^{[a]}\in\pcMat_{a\times N}(R)$ the last $a$ rows of $k^{-1}$.
Functions $\Phi$ on $\pcMat_{a\times N}(R)$ and $\Psi$ on $\pcMat_{N\times b}(R)$ are \emph{matched
for $\chi$} if $\Phi$ vanishes off $\{k^{[a]}\}$, $\Psi$ vanishes off $\{k_{[b]}\}$, and
$\Phi(k^{[a]})=\chi(\det k)\Psi(k_{[b]})$ for all $k\in\auGL_N(R)$. Over $F_S$ the same definition is
used with $k\in\auGL_N(\auO_S)$.
\end{definition}

\begin{theorem}[Exchange identity]\label{pcx:X}
Let $W,V,V'$ be spaces of ranks $N=a+b$, $a$, $b$, with $a=b+2e$, $e\ge0$ (so $N$ is even). Let
$q,q_V,q_{V'}$ be their numbers of negative signs; assume $q=q_V+q_{V'}$ and $2q+2\le N$. Let
$\mu=(\mu_\tau)$ be integers with $N\le2\mu_\tau+2q$ for all $\tau$, and let $v$ be a vector of the
Fock space of $\auHom(V',W)\otimes\R$ with finite moments. Then there are $C>0$ and a vector $x$ of
the Fock space of $\auHom(W,V)\otimes\R$ with finite moments, such that: for every modulus
$\pcm\ne(1)$ and every tame character $\nu$ for $\pcm$ of weight $\mu$ there is $c\in\C$ with
\[
|c|^2=C\cdot L_*(\pcm,\nu,e)\ne0,
\]
such that for all automorphic characters $\alpha,\alpha'$ of $\auU(1)$ with
\[
\alpha\alpha'=\gamma:=\nu^{-1}\upsilon^{N/2},
\]
all functions $\Phi$ on $\pcMat_{a\times N}(\pcR_\pcm)$ and $\Psi$ on $\pcMat_{N\times b}(\pcR_\pcm)$
matched for $\chi$, and all $g\in\auU(W)(\auA)$,
\begin{equation}\label{pcx:eq-X}
\int_{[\auU(V)]}\theta^{W,V}_{\alpha,\gamma}(g,h;x,\Phi)\,dh=c\int_{[\auU(V')]}
\theta^{V',W}_{\gamma,\alpha'}(h',g;v,\Psi)\,dh' .
\end{equation}
Here $\theta(\cdot\,;x,\Phi)$ is the kernel of the datum $[x,\Phi]$ of \S\ref{pc:ss-data} (matrices
in the frames), and $C$, $x$ do not depend on $\pcm,\nu,\alpha,\alpha',\Phi,\Psi,g$.
\end{theorem}
Thus on the left $\auU(W)$ acts through $\alpha(\det g)^{-1}$ and $\auU(V)$ through $\gamma(\det h)$;
on the right $\auU(V')$ acts through $\gamma(\det h')^{-1}$ and $\auU(W)$ through $\alpha'(\det g)$.
Both integrals converge absolutely (Corollary~\ref{pcx:conv}).

\paragraph*{The dictionary with the other splitting}
The other convention in use ({\cite[\S2.1]{P54}}, {\cite[\S6.1.4]{P55}}, Remark~\ref{pc:rem-norm}) writes, for characters $a,b$,
\[
\omega'_{a,b}(g,h)=a(\det h)\,b(\det g)^{-1}\,\omega_{\auHom(S,T)}(g,h),\qquad
\theta'^{S,T}_{a,b}=\theta\circ\omega'_{a,b}.
\]
The underlying operator $\omega_{\auHom(S,T)}(g,h)$ is the same in both (the Weil representation
of the unitary group of $\auHom(S,T)$ defined with $\upsilon$, composed with $x\mapsto hxg^{-1}$), so
\begin{equation}\label{pcx:eq-dict}
\omega'_{a,b}=\omega_{b,a},\qquad\theta'^{S,T}_{a,b}=\theta^{S,T}_{b,a}
\end{equation}
\emph{exactly}: the dictionary is the transposition of the two indices, with no factor. In the
primed convention \eqref{pcx:eq-X} reads
\[
\int_{[\auU(V)]}\theta'^{W,V}_{\gamma,\alpha}(g,h;x,\Phi)\,dh=c\int_{[\auU(V')]}
\theta'^{V',W}_{\alpha',\gamma}(h',g;v,\Psi)\,dh',
\]
with the \emph{same} $c$, $C$, $x$. Theorem~\ref{pc:X} is printed in the convention of
\S\ref{pc:ss-data}, with the indices of \eqref{pcx:eq-X}. The representation $\omega_\mu$ of {\cite[\S5.1.1]{P49}} is
$\omega_{\mu\upsilon_{N'},\mu\upsilon_N}=\omega'_{\mu\upsilon_N,\mu\upsilon_{N'}}$.

Changing the splitting characters multiplies kernels by characters of the determinants:
\begin{equation}\label{pcx:eq-twist}
\theta^{S,T}_{\alpha\lambda,\beta\lambda'}(g,h;f)=\lambda(\det g)^{-1}\lambda'(\det h)\,
\theta^{S,T}_{\alpha,\beta}(g,h;f),
\end{equation}
immediately from the definition, since $\theta$ is linear. In particular the passage from $\Theta'$
to $\Theta$ of Remark~\ref{pc:rem-norm} (multiplication by $\chi(\det h)$ on the target group,
$|\chi|=1$) does not change $|c|$: if \eqref{pcx:eq-X} is used with both sides multiplied by the
same unitary character of $\det g$ or integrated against forms twisted by unitary characters, the
constant $c$ is unchanged, and if the twist is put on one side only it is not an instance of
\eqref{pcx:eq-X}.

\paragraph*{Non-vanishing of the constant}

\begin{lemma}\label{pcx:Lpos}
Let $\nu$ be a tame character for $\pcm$ of weight $\mu\ne0$. Then $L^S(z,\nu)\ne0$ for every
real $z\ge1$; hence $L_*(\pcm,\nu,e)>0$ for every $e\in\Z$.
\end{lemma}
\begin{proof}
The Hecke character $\lambda(x)=\nu(x/\bar x)$ of $\auA_E^\times/E^\times$ is unitary, and it is not
trivial: its component at an infinite place $\tau$ is $x\mapsto(x/\bar x)^{\mu_\tau}$ (up to the
normalisation of the weight), which is non-trivial since $\mu_\tau\ne0$. For real $z>1$ the Euler
product of $L^S(z,\lambda)$ converges absolutely and each factor $(1-\lambda(\varpi_w)q_w^{-z})^{-1}$
is non-zero, so $L^S(z,\nu)\ne0$. For $z=1$, $L(s,\lambda)$ has no zero on $\Rel s=1$ for a
non-trivial unitary Hecke character (Hecke; see e.g.\ Lang, \emph{Algebraic Number Theory}), and removing finitely many Euler factors, each finite and non-zero
at $s=1$, keeps this. Thus every factor of \eqref{pcx:eq-Lstar} is positive.
\end{proof}

\begin{corollary}[The constant is non-zero for every $e$]\label{pcx:cne0}
In Theorem~\ref{pcx:X}, $c\ne0$ for every $e\ge0$; for $e=0$, $|c|^2=C$.
\end{corollary}
\begin{proof}
$C>0$ and $L_*(\pcm,\nu,e)>0$ by Lemma~\ref{pcx:Lpos} (the weight is non-zero because
$\mu_\tau\ge(N-2q)/2\ge1$).
\end{proof}

\begin{remark}[The exponents $e$ used downstream]\label{pcx:rem-e}
Theorem~\ref{pcx:X} is applied: (i) in the commutation of one square of the triangle
(Theorem~\ref{pc:tri}(b), \S\ref{pc:ss-exchange}), with $(W,V,V')=(\pcG_{i,n+1}\oplus\pcG_{i+1,n+1},
\pcG_{i,n+2},\pcG_{i+1,n})$ and $e=t$; (ii) in the adjoint law (Theorem~\ref{pc:tri}(c),
\S\ref{pc:ss-exchange}(v)) with $e=0$; (iii) for $V$ positive definite with explicit determinant
data at the infinite places, which is the form used for {\cite[Proposition~2.15]{P54}} and for
Theorem~\ref{pc:auxcore}, with $e=d$ where $a=b+2d$ is dictated by the ranks of the spaces of that
application, an arbitrary integer $d\ge0$ as far as the statement is concerned. Since
Corollary~\ref{pcx:cne0} holds for all $e\ge0$, no restriction on $e$ is needed in any of these
applications; the only input is Lemma~\ref{pcx:Lpos}, valid because the characters used have
non-zero weight. (The non-vanishing can also be read off from $C_1|c|^2=L_*(\pcm,\nu,e)\,C_0$ (proof of Theorem~\ref{pcx:X}, $C_1$ the constant
of Proposition~\ref{pcx:P-matches}) with $C_0\neq0$ the equal-rank constant of Proposition~\ref{pcx:P-equal} and $L_*>0$.)
\end{remark}

\subsubsection{Architecture of the proof}\label{pcx:ss-arch}

Throughout this subsection $W$ is a space of even rank $N$ with $q$ negative signs, $\nu$ runs
through the tame characters of a fixed weight $\mu$ (for all moduli), $S$ is the set of primes of
the modulus of $\nu$, $\chi_S$ is as in \S\ref{pcx:ss-statement}, and
\[
\gamma=\gamma_\nu=\nu^{-1}\upsilon^{N/2}.
\]
Data at $S$ are functions on matrices over $F_S$: for spaces $S_1,S_2$ of ranks $n_1,n_2$, a Fock
vector $y$ of $\auHom(S_1,S_2)\otimes\R$ and a locally constant compactly supported $\Phi$ on
$\pcMat_{n_2\times n_1}(F_S)$, $[y,\Phi]_S$ denotes the datum which is $y$ at the infinite places,
$\Phi$ at the places of $S$ (in the model (W5) at $w\in S$) and the vacuum $\psvac$ of the standard
lattice elsewhere. $[y,\Phi]_{\bar S}$ is the same with $\Phi$ put at the places of $\bar S$.

\begin{definition}[The two theta integrals]\label{pcx:def-large}
For spaces $W$ (rank $N$) and $T$ (rank $r$), a Fock vector $v$ of $\auHom(W,T)\otimes\R$ and
$\Phi$ on $\pcMat_{r\times N}(F_S)$, $\Psi$ on $\pcMat_{N\times r}(F_S)$, put
\[
\Theta^{\mathrm L}_{W,T}(v,\Phi)(g)=\int_{[\auU(T)]}\theta^{W,T}_{\gamma,\gamma}(g,h;[v,\Phi]_S)\,dh,
\qquad
\Theta^{\mathrm S}_{W,T}(v,\Psi)(g)=\int_{[\auU(T)]}\theta^{W,T}_{\upsilon^{r},\,\nu\upsilon^{N/2}}
(g,h;[v,\Psi^{\mathsf t}]_{\bar S})\,dh,
\]
with $\Psi^{\mathsf t}(y)=\Psi(y^{\mathsf t})$. (In the primed convention, these are the
integrals with characters $(\gamma,\gamma)$, resp.\ $(\nu\upsilon^{N/2},\upsilon^r)$.)
\end{definition}

\begin{definition}[Standard data]\label{pcx:def-std}
Let $N=b+a$. For $z\in\pcMat_{b\times N}(\auO_S)$ define on $\pcMat_{a\times N}(F_S)$ and
$\pcMat_{N\times b}(F_S)$
\[
\Phi^{\mathrm{row}}_z(x)=\chi_S\bigl(\det\tbinom zx\bigr)^{-1},\qquad
\Psi^{\mathrm{col}}_z(y)=\chi_S\bigl(\det(zy)\bigr)^{-1},
\]
where $\Phi^{\mathrm{row}}_z(x)=0$ unless $\binom zx\in\auGL_N(\auO_S)$, and
$\Psi^{\mathrm{col}}_z(y)=0$ unless $y$ is integral and $zy\in\auGL_b(\auO_S)$.
\end{definition}

\begin{definition}[Equal-rank ratio]\label{pcx:def-ratio}
For a space $X$ of even rank $2k$ and a Fock vector $v$ of $\auHom(X,X)\otimes\R$ with finite
moments let
\[
\pcxstr_X(\nu;v)=\frac{\bigl|\Theta^{\mathrm L}_{X,X}(v,\Psi^{\mathrm{col}}_1)(1)\bigr|^2}
{L_*(\pcm,\nu,k)},
\]
where $\Psi^{\mathrm{col}}_1=\chi_S^{-1}\circ\det$ on $\auGL_{2k}(\auO_S)$ and $0$ elsewhere.
\end{definition}

\begin{definition}[Exchange on standard data]\label{pcx:def-RC}
Let $X$, $T$, $T'$ be spaces of ranks $N=b+a$, $a$, $b$, and $\ell$ a function of the tame
characters of weight $\mu$. Vectors $v$ (for $(X,T)$) and $y$ (for $(X,T')$) with finite moments
satisfy $\mathrm{RC}(X,T,T';v,y;\ell)$ if there is a real $C$ such that every $\nu$ has a $c\in\C$
with
\begin{enumerate}
\item[(a)] $C|c|^2=L_*(\pcm,\nu,a-N/2)\,\ell(\nu)$;
\item[(b)] $\Theta^{\mathrm L}_{X,T}(v,\Phi^{\mathrm{row}}_z)=c\,\Theta^{\mathrm S}_{X,T'}
(y,\Psi^{\mathrm{col}}_z)$ for every $z\in\pcMat_{b\times N}(\auO_S)$.
\end{enumerate}
\end{definition}

The proof of Theorem~\ref{pcx:X} uses the following intermediate statements. In each, the
quantifiers over $\nu$ are as above.

\begin{proposition}[Exchange on standard data in even rank]\label{pcx:P-std}
Let $W$ have even rank $N=b+r$, let $U$ (rank $b$) and $V$ (rank $r$) be spaces with
$q_U+q=b+q_V$, and let $W_0=\pcxH\oplus W$, $\pcxH$ a hyperbolic plane. For Fock vectors $z$ of
$(W_0,W_0)$ and $y$ of $(W,U)$ with finite moments there is $v$ for $(W,V)$ with
$\mathrm{RC}(W,V,U;v,y;\nu\mapsto\pcxstr_{W_0}(\nu;z))$.
\end{proposition}

This is the heart of the matter. Its proof has three parts.

\begin{proposition}[Large triple]\label{pcx:P-top}
Let $W_0=\pcxH\oplus W$ be of even rank, $U$ of rank $b$, $X=\pcxH^b\oplus W_0$, and
$T=U\oplus W_0$ (no sign condition is needed; $W_0$ splits off an adapted plane). For $z$ of $(W_0,W_0)$ and
$w$ of $(W_0,U)$ there are $x$ for $(X,T)$ and $x'$ for $(X,U)$, with $x'$ having zero term $w$
along the $b$ planes (for the splitting with the two isotropic frames exchanged, \S\ref{pcu:ss-mixed}),
such that $\mathrm{RC}(X,T,U;x,x';\nu\mapsto\pcxstr_{W_0}(\nu;z))$.
\end{proposition}

\begin{proposition}[Real vectors with prescribed zero term]\label{pcx:P-slice}
Let $W'=W\oplus\pcxH^p$ be an adapted splitting. For every Fock vector $y$ of $(W,U)$ with finite
moments there is a Fock vector $w$ of $(W',U)$ with finite moments whose zero term (restriction of
the real Schr\"odinger model to the slice of matrices vanishing on the isotropic part) is $y$.
\end{proposition}

\begin{proposition}[Removing planes]\label{pcx:P-descent}
Let $W'=W\oplus\pcxH$ be an adapted splitting, $N=\operatorname{rk}W$ even, $T'$ of rank $b$,
$T$ of rank $r'$, and $T_0$ any space of rank $r'-2$ with $q_{T_0}=q_T-1$ (zero terms taken with the
orientation of Remark~\ref{pcu:rem-descent}). If
$\mathrm{RC}(W',T,T';x,x';\ell)$ holds, then there is $v$ with $\mathrm{RC}(W,T_0,T';v,y;\ell)$,
where $y$ is the zero term of $x'$. (Here $y=Z_Ax'$ is the zero term of $x'$ along the given adapted
splitting, $\ell$ is kept and the constant changes. With $d=r'-2-N/2$, $v$ is the cone integral
$\pctcone_{d+1}(x)$ if $d\ge0$ and the dual cone integral $\pctcone^\vee_{-d}(x)$ if $d<0$; the ratio of
the constant terms has modulus $C\,L^\sharp(d+1)$, resp.\ $C\,L^\sharp(-d)$, $L^\sharp(z)=|L^S(z,\nu)|
\prod_{w\in S}q_w^{z-1/2}$; only values of $L^S$ at points $\ge1$ occur and no functional equation is
used. Proof: Proposition~\ref{pct:descent}.)
\end{proposition}

\begin{proposition}[From standard data to matched functions]\label{pcx:P-match}
Let $W$ have even rank $N=b+r$. If $v$, $y$ have finite moments and
$\Theta^{\mathrm L}_{W,V}(v,\Phi^{\mathrm{row}}_z)=c\,\Theta^{\mathrm S}_{W,U}(y,\Psi^{\mathrm{col}}_z)$
for all integral $z$, then for all $\Phi$, $\Psi$ matched for $\chi_S$ over $\auO_S$ with $\Psi$
locally constant, $\Psi$ is a Schwartz--Bruhat function supported on left-invertible matrices and
$\Theta^{\mathrm L}_{W,V}(v,\Phi)=c\,\Theta^{\mathrm S}_{W,U}(y,\Psi)$.
\end{proposition}

\begin{proposition}[Passing to adjoints]\label{pcx:P-flip}
Let $W$ have even rank $N=b+r$, $U$ of rank $b$, $U^-$ the space with the negated form, and $y$ a
Fock vector of $(U,W)$. There is a Fock vector $y^{\mathrm{ad}}$ of $(W,U^-)$, with the same
moments, a number $d_0>0$, and for each $\nu$ a number $d$ with $|d|=d_0$, such that for all $\Phi,\Psi$ matched for
$\chi_S$ with $\Psi$ Schwartz--Bruhat on left-invertible matrices and all $g$,
\[
\gamma(\det g)\,\Theta^{\mathrm S}_{W,U^-}(y^{\mathrm{ad}},\Psi)(g)
=d\int_{[\auU(U)]}\theta^{U,W}_{\gamma,\gamma}(h,g;[y,\Psi]_S)\,dh .
\]
\end{proposition}


\begin{proposition}[Value of the equal-rank ratio]\label{pcx:P-equal}
Let $W$ have rank $2k+2q$ with $q$ negative signs, and $k\le\mu_\tau$ for all $\tau$. There are a
Fock vector $z$ of $(\pcxH\oplus W,\pcxH\oplus W)$ with finite moments and $C_0\ne0$ such that
$\pcxstr_{\pcxH\oplus W}(\nu;z)=C_0$ for every tame character $\nu$ of weight $\mu$.
\end{proposition}
Its proof uses Propositions~\ref{pcx:P-matches} (hence \ref{pcx:P-std}--\ref{pcx:P-flip}) and three
further statements:

\begin{proposition}[Non-vanishing for one character]\label{pcx:P-nonvan}
Let $J$ be the positive definite standard space of rank $2k$, $k\le\mu_\tau$, $Y=\pcxH\oplus J$.
For each $\nu$ there is $z$ with $\pcxstr_Y(\nu;z)\ne0$.
\end{proposition}

\begin{proposition}[Adding a plane]\label{pcx:P-plane}
Let $W'=W\oplus\pcxH$ be an adapted splitting, $\operatorname{rk}W$ even. For every $v$ of $(W,W)$
with finite moments there are $v'$ of $(W',W')$ with finite moments and $C>0$ with
$\pcxstr_{W'}(\nu;v')=C\,\pcxstr_W(\nu;v)$ for all $\nu$.
\end{proposition}

\begin{proposition}[Only the signs matter]\label{pcx:P-signs}
If $Y$, $Y_1$ are spaces of the same even rank and the same signs, and $Y$ contains a hyperbolic
plane split off adaptedly, then $\pcxstr_{Y_1}(\nu;v)=\pcxstr_Y(\nu;v)$ for all $\nu$, $v$.
\end{proposition}

Finally two elementary statements, proved in \S\ref{pcx:ss-elem}:
Lemma~\ref{pcx:std-matched} (standard data are matched) and Lemma~\ref{pcx:mod} (data modulo
$\pcm$ come from matched data on $F_S$).

\begin{proposition}[Exchange for matched functions]\label{pcx:P-matches}
Let $W,U,V$ have ranks $N=b+r$ (even), $b$, $r$, with $q=q_V+q_U$. For $z$ of $(\pcxH\oplus W,
\pcxH\oplus W)$ and $y$ of $(U,W)$ with finite moments there are $C\in\R$ and $v$ of $(W,V)$ with
finite moments such that every $\nu$ has $c$ with
\begin{enumerate}
\item[(a)] $C|c|^2=L_*(\pcm,\nu,r-N/2)\;\pcxstr_{\pcxH\oplus W}(\nu;z)$;
\item[(b)] for all $\Phi,\Psi$ matched for $\chi_S$ over $\auO_S$, $\Psi$ locally constant, and all
$g$:
$\displaystyle\int_{[\auU(V)]}\theta^{W,V}_{1,\gamma}(g,h;[v,\Phi]_S)\,dh=c\int_{[\auU(U)]}
\theta^{U,W}_{\gamma,\gamma}(h,g;[y,\Psi]_S)\,dh$.
\end{enumerate}
\end{proposition}

The logical structure is:
\[
\begin{gathered}\ref{pcx:P-slice}+\ref{pcx:P-top}+\ref{pcx:P-descent}\Rightarrow\ref{pcx:P-std};\\
\ref{pcx:P-std}+\ref{pcx:P-match}+\ref{pcx:P-flip}\Rightarrow\ref{pcx:P-matches};\\
\ref{pcx:P-matches}+\ref{pcx:P-nonvan}+\ref{pcx:P-plane}+\ref{pcx:P-signs}\Rightarrow
\ref{pcx:P-equal};\\
\ref{pcx:P-matches}+\ref{pcx:P-equal}+\ref{pcx:mod}+\eqref{pcx:eq-twist}\Rightarrow
\ref{pcx:X}.\end{gathered}
\]

\begin{lemma}[Adapted planes exist]\label{pcx:adapted}
Let $X$ be a space of rank $n$ with $q_X$ negative signs. For every sign vector of rank $n+2$ with
$q_X+1$ negative signs there are a space $Y$ with these signs and an adapted splitting
$Y=X\oplus\pcxH$ (frames orthonormal at the infinite places and integral, with unimodular sum, at
the finite places).
\end{lemma}
\begin{proof}
Let $\pcxH$ be the plane with basis $e,f$ and Gram matrix $\bigl(\begin{smallmatrix}0&1\\1&0
\end{smallmatrix}\bigr)$; $\auO_Ee\oplus\auO_Ef$ is unimodular. Let $Y=X\oplus\pcxH$ with the
lattice $L_X\oplus\widehat\auO_Ee\oplus\widehat\auO_Ef$, unimodular, and at each infinite place the
orthonormal basis consisting of that of $X$ and of $(e+f)/\sqrt2$, $(e-f)/\sqrt2$ (signs $+1,-1$).
The signature of $Y$ at every infinite place is that of $X$ plus $(1,1)$, so $Y$ has rank $n+2$ and
$q_X+1$ negative signs; ordering the frame so that the signs come in the prescribed order gives
the required sign vector. The splitting $Y=X\oplus\pcxH$, with isotropic frame $e,f$, is adapted
by construction.
\end{proof}

\begin{lemma}[Spaces with the same signs]\label{pcx:landherr}
Let $Y,Y_1$ be spaces of the same rank and signs. There are an $E$-isometry $\phi\colon Y_1(E)\to Y(E)$
and $k\in\auU(Y)(\auA)$ such that $\phi$ carries the frame of $Y_1$ (its lattice at the finite places
and its orthonormal bases at the infinite places) to $k$ applied to the frame of $Y$.
\end{lemma}
\begin{proof}
At each finite place both lattices are unimodular, hence the local hermitian spaces have
determinant a unit, which is a local norm since $E/E^+$ is unramified at finite places; at the
infinite places the signatures agree. By Landherr's theorem (hermitian spaces over $E/E^+$ are
classified by rank, the class of the determinant in $E^{+\times}/\auNm E^\times$, and the
signatures;~\cite{Landherr}) together with the Hasse norm theorem for
the cyclic extension $E/E^+$, the determinants agree in $E^{+\times}/\auNm E^\times$ and $Y_1\cong Y$.
Transport the frame of $Y_1$ by $\phi$: at each finite place it is a unimodular lattice of $Y_v$,
hence of the form $k_vL_Y$ with $k_v\in\auU(Y)(E^+_v)$, and $k_v=1$ for almost all $v$ (both lattices
are the standard ones almost everywhere) (over an unramified extension a unimodular hermitian lattice has an orthonormal basis, proof: the reduction of the lattice is a non-degenerate hermitian space over the quadratic extension of residue fields, so it has a vector $\bar x$ with $\bar h(\bar x,\bar x)\ne0$ (otherwise polarisation and the surjectivity of the residue trace give $\bar h=0$); a lift $x$ has $h(x,x)$ a unit of the base ring, hence the norm of a unit since the extension is unramified, so after scaling $h(x,x)=1$; then the lattice is $\mathcal O x\perp(x^\perp\cap\text{lattice})$ with unimodular second summand, and one inducts on the rank; see also Jacobowitz~\cite{Jacobowitz}; at a split place every
unimodular lattice is standard); at the infinite places two orthonormal bases of the same signs
differ by an element of the real unitary group.
\end{proof}


\subsubsection{Elementary lemmas}\label{pcx:ss-elem}

\begin{lemma}[Standard data are matched]\label{pcx:std-matched}
For $z\in\pcMat_{b\times N}(\auO_S)$, $\Phi^{\mathrm{row}}_z$ and $\Psi^{\mathrm{col}}_z$ are matched
for $\chi_S$ over $\auO_S$, and $\Psi^{\mathrm{col}}_z$ is locally constant.
\end{lemma}
\begin{proof}
Let $k\in\auGL_N(\auO_S)$. Then
\[
\tbinom{z}{k^{[a]}}\,k=\begin{pmatrix}zk_{[b]}&z k_{(a)}\\0&1_a\end{pmatrix},
\]
where $k_{(a)}$ are the last $a$ columns of $k$, because the last $a$ rows of $k^{-1}k$ are
$(0\ 1_a)$. Taking determinants, $\det\binom z{k^{[a]}}\det k=\det(zk_{[b]})$. The matrix
$\binom z{k^{[a]}}$ is integral, and it lies in $\auGL_N(\auO_S)$ if and only if $\det(zk_{[b]})$ is a
unit; $k_{[b]}$ is integral. Hence both sides vanish together, and otherwise
$\Phi^{\mathrm{row}}_z(k^{[a]})=\chi_S(\det(zk_{[b]}))^{-1}\chi_S(\det k)=\chi_S(\det k)
\Psi^{\mathrm{col}}_z(k_{[b]})$. Supports: if $\Phi^{\mathrm{row}}_z(x)\ne0$ then $M=\binom zx\in
\auGL_N(\auO_S)$ and $x=k^{[a]}$ for $k=M^{-1}$. If $\Psi^{\mathrm{col}}_z(y)\ne0$ then $y$ is integral
and $(zy)^{-1}z$ is an integral left inverse of $y$; over the product of complete discrete
valuation rings $\auO_S$ this means that the reduction of $y$ modulo each prime has rank $b$, so the
columns of $y$ extend to a basis of $\auO_S^N$: there is $k\in\auGL_N(\auO_S)$ with $k_{[b]}=y$.
Finally $\Psi^{\mathrm{col}}_z$ is locally constant: the set of integral $y$ with $zy$ invertible is
open and closed, and on it $\chi_S^{-1}(\det(zy))$ is locally constant because $\chi_S$ is tame,
hence trivial on $1+\varpi\auO_S$.
\end{proof}

\begin{lemma}[Data modulo $\pcm$]\label{pcx:mod}
Let $\nu$ be a tame character for $\pcm$, $S$ the primes of $\pcm$, and let $\Phi$, $\Psi$ be
functions on $\pcMat_{a\times N}(\pcR_\pcm)$, $\pcMat_{N\times b}(\pcR_\pcm)$ matched for $\chi$. Let
$\Phi'(x)=\Phi(x\bmod\pcm)$ for integral $x\in\pcMat_{a\times N}(F_S)$ and $\Phi'(x)=0$ otherwise,
and define $\Psi'$ in the same way. Then $\Phi'$, $\Psi'$ are matched for $\chi_S$ over $\auO_S$,
$\Psi'$ is locally constant, and $[x,\Phi]=[x,\Phi']_S$, $[v,\Psi]=[v,\Psi']_S$ for all Fock
vectors $x$, $v$ (the left sides being the data of \S\ref{pc:ss-data}).
\end{lemma}
\begin{proof}
The reduction $\pi\colon\auO_S\to\pcR_\pcm=\prod_{w\in S}\auO_{E_w}/w^{n_w}$ is a surjective ring
homomorphism, and $u\in\auO_S$ is a unit if and only if $\pi(u)$ is (both mean: non-zero residue at
every $w\in S$). On units $\chi_S(u)=\prod_{w\in S}\chi_w(u_w\bmod w)=\chi(\pi(u))$, by the
definition of the tame parts. The equality of data is the definition of $\Phi^\natural$ in
\S\ref{pc:ss-data}; $\Psi'$ is locally constant since it factors through reduction modulo $\pcm$
on the open and closed set of integral matrices.

Supports. If $\Phi'(x)\ne0$, then $x$ is integral and $\pi(x)=\bar k^{[a]}$ for some $\bar k\in
\auGL_N(\pcR_\pcm)$. Lift the first $b$ rows of $\bar k^{-1}$ to an integral $z$ and put
$M=\binom zx$; $\pi(M)=\bar k^{-1}$, so $\det M$ is a unit and $M\in\auGL_N(\auO_S)$; then $x=k^{[a]}$
for $k=M^{-1}$. If $\Psi'(y)\ne0$, then $\pi(y)=\bar k_{[b]}$; lift the last $a$ columns of $\bar k$ to
an integral $u$ and put $k=(y\ u)$, whose reduction is $\bar k$; so $k\in\auGL_N(\auO_S)$ and
$y=k_{[b]}$. Values. For $k\in\auGL_N(\auO_S)$, $\pi(k^{[a]})=\pi(k)^{[a]}$ and $\pi(k_{[b]})=
\pi(k)_{[b]}$, so $\Phi'(k^{[a]})=\Phi(\pi(k)^{[a]})=\chi(\det\pi(k))\Psi(\pi(k)_{[b]})=
\chi_S(\det k)\Psi'(k_{[b]})$.
\end{proof}

\noindent\emph{Coordinates at $S$.} The data of \S\ref{pc:ss-data} are those of
Definition~\ref{pgx:conv}: at $w\in S$, $\Phi^\natural(x)=\Phi(a_w^{-1}x\bmod\pcm)$ on $a_w\pcMat(\auO_w)$,
$a_w$ a generator of $\alpha_w\auO_w$. Throughout \S\S\ref{pcx:sec}--\ref{pcn:sec}, matrices at $w\in S$
are written in the coordinates $\tilde x$ of Lemma~\ref{pgx:dict}: $\tilde x=a_w^{-1}x$ on the
$w$-component and $a_{\aubar w}^{-1}x$ on the $\aubar w$-component, in which both components of the
self-dual lattice are the integral matrices and the Heisenberg pairing is
$\lambda_w(\operatorname{tr}({}^t\tilde y\tilde x))^{\pm1}$ with $\lambda_w$ of conductor $\auO_w$. As $a_w$ is a scalar this
commutes with the actions $x\mapsto k'^{-1}xk$ of (W5) and with the measure normalisations; the vacuum
lattice $a_w\pcMat(\auO_w)$ (Lemma~\ref{pgx:vac}) becomes $\pcMat(\auO_w)$, and $\Phi^\natural$ becomes
the function $\Phi'$ of the lemma. So Lemma~\ref{pcx:mod} holds with the data of
Definition~\ref{pgx:conv}, $[x,\Phi']_S$ being read in these coordinates; where $\alpha_w$ is a unit,
$a_w=1$ and nothing changes. (See \S\ref{pgx:ss-other}.)

\subsubsection{Proof of Theorem~\ref{pcx:X} from the intermediate propositions}\label{pcx:ss-assembly}

\begin{proof}[Proof of Proposition~\ref{pcx:P-matches} from Propositions~\ref{pcx:P-std},
\ref{pcx:P-match}, \ref{pcx:P-flip}]
The space $U^-$ has $b-q_U$ negative signs, so $q_{U^-}+q=b+q_V$ is equivalent to $q=q_U+q_V$.
Let $y^{\mathrm{ad}}$ be the vector of Proposition~\ref{pcx:P-flip}; it has finite moments.
Proposition~\ref{pcx:P-std} for $(W,U^-,V)$, $z$ and $y^{\mathrm{ad}}$ gives $v$ and $C$, and for each
$\nu$ a number $c_0$ with $C|c_0|^2=L_*(\pcm,\nu,r-N/2)\pcxstr_{\pcxH\oplus W}(\nu;z)$ and
$\Theta^{\mathrm L}_{W,V}(v,\Phi^{\mathrm{row}}_z)=c_0\Theta^{\mathrm S}_{W,U^-}(y^{\mathrm{ad}},
\Psi^{\mathrm{col}}_z)$ for integral $z$. By Proposition~\ref{pcx:P-match}, for $\Phi,\Psi$ matched
with $\Psi$ locally constant, $\Psi$ is Schwartz--Bruhat on left-invertible matrices and
$\Theta^{\mathrm L}_{W,V}(v,\Phi)=c_0\Theta^{\mathrm S}_{W,U^-}(y^{\mathrm{ad}},\Psi)$. By
\eqref{pcx:eq-twist}, $\theta^{W,V}_{1,\gamma}(g,h)=\gamma(\det g)\theta^{W,V}_{\gamma,\gamma}(g,h)$, so
\[
\begin{aligned}&\int_{[\auU(V)]}\theta^{W,V}_{1,\gamma}(g,h;[v,\Phi]_S)\,dh=\gamma(\det g)\Theta^{\mathrm L}_{W,V}
(v,\Phi)(g)\\&\quad=c_0\,\gamma(\det g)\Theta^{\mathrm S}_{W,U^-}(y^{\mathrm{ad}},\Psi)(g)
\\&\quad=c_0d\int_{[\auU(U)]}\theta^{U,W}_{\gamma,\gamma}(h,g;[y,\Psi]_S)\,dh\end{aligned}
\]
by Proposition~\ref{pcx:P-flip}. Put $c=c_0d$; since $|d|=d_0$ does not depend on $\nu$,
$Cd_0^{-2}|c|^2=C|c_0|^2$ gives (a) with the constant $Cd_0^{-2}$.
\end{proof}

\begin{proof}[Proof of Theorem~\ref{pcx:X} from Propositions~\ref{pcx:P-matches},
\ref{pcx:P-equal} and Lemmas~\ref{pcx:Lpos}, \ref{pcx:mod}]
Put $k=N/2-q$. From $N\le2\mu_\tau+2q$ we get $k\le\mu_\tau$, and $\operatorname{rk}W=2k+2q$.
Proposition~\ref{pcx:P-equal} gives $z$ and $C_0\ne0$ with $\pcxstr_{\pcxH\oplus W}(\nu;z)=C_0$ for
all tame $\nu$ of weight $\mu$. Proposition~\ref{pcx:P-matches} for $(W,U,V)=(W,V',V)$, this $z$
and $y=v$ gives $C_1\in\R$ and $x$ (the vector of the theorem). Let $\nu$ be tame for $\pcm$ of weight
$\mu$ (note $\mu_\tau\ge k\ge1$ since $2q+2\le N$). The number $c$ of
Proposition~\ref{pcx:P-matches} satisfies, as $a-N/2=e$,
\[
C_1|c|^2=L_*(\pcm,\nu,e)\,C_0 .
\]
The right side is non-zero by Lemma~\ref{pcx:Lpos}; hence $C_1\ne0$, $c\ne0$, and
$|c|^2=C\,L_*(\pcm,\nu,e)$ with $C=C_0/C_1$, which is then positive. ($C$ does not depend on
$\pcm,\nu$; if no tame character of weight $\mu$ exists the theorem is empty and we take $C=1$.)

Let $\Phi,\Psi$ be matched for $\chi$ modulo $\pcm$, and let $\Phi',\Psi'$ be as in
Lemma~\ref{pcx:mod}. Proposition~\ref{pcx:P-matches}(b) for $\Phi',\Psi'$ is \eqref{pcx:eq-X} for
$\alpha=1$, $\alpha'=\gamma$. For general $\alpha\alpha'=\gamma$, \eqref{pcx:eq-twist} gives
$\theta^{W,V}_{\alpha,\gamma}(g,h)=\alpha(\det g)^{-1}\theta^{W,V}_{1,\gamma}(g,h)$ and
$\theta^{V',W}_{\gamma,\alpha'}(h',g)=(\alpha'\gamma^{-1})(\det g)\,\theta^{V',W}_{\gamma,\gamma}(h',g)
=\alpha(\det g)^{-1}\theta^{V',W}_{\gamma,\gamma}(h',g)$: both sides of \eqref{pcx:eq-X} are
multiplied by $\alpha(\det g)^{-1}$, and the identity persists with the same $c$.
\end{proof}

\subsubsection{The equal-rank ratio: proof of Proposition~\ref{pcx:P-equal}}\label{pcx:ss-equal}

Let $J$ be positive definite of rank $2k$, $Y=\pcxH\oplus J$, and $v$ a Fock vector of $(Y,Y)$ with
finite moments. Write $\mathrm{vol}_0$ for the volume of $[\auU(0)]$ (a point) and $\mathrm{vol}_J$
for that of the compact space $[\auU(J)]$; both are positive. For a datum $f$ of $(J,J)$ the kernel
$\theta^{J,J}_{\gamma,\gamma}(h,g;f)$ is continuous on the compact space $[\auU(J)]^2$, so Fubini
applies.

\begin{lemma}[The identity used twice]\label{pcx:twice}
Assume Proposition~\ref{pcx:P-matches}. There is $D\ge0$ with $\pcxstr_Y(\nu;v)^2=D$ for all tame
$\nu$ of weight $\mu$.
\end{lemma}
\begin{proof}
\emph{The triple $(J,J,0)$.} Apply Proposition~\ref{pcx:P-matches} with $(W,U,V)=(J,0,J)$ (ranks
$2k$, $0$, $2k$, sign condition $0=0+0$), $z=v$ and $y$ the vacuum of the zero Fock space
$\auHom(0,J)\otimes\R=0$. It gives $B$ and $y_J$ for $(J,J)$, and for each $\nu$ a number $c_1$ with
$B|c_1|^2=L_*(\pcm,\nu,k)\pcxstr_Y(\nu;v)$. Take $\Phi=\Phi^{\mathrm{row}}_\emptyset=\chi_S^{-1}\circ
\det$ (on $\auGL_{2k}(\auO_S)$, zero elsewhere) and $\Psi=\Psi^{\mathrm{col}}_\emptyset$, the
function $1$ on the one-point space $\pcMat_{2k\times0}$; they are matched by
Lemma~\ref{pcx:std-matched} (with $b=0$). On $\auHom(0,J)=0$ the theta functional of the datum
$[\mathrm{vac},1]$ is $1$, so the right side of (b) is $c_1\gamma(\det g)\mathrm{vol}_0$, while the left
side is $\gamma(\det g)\Theta^{\mathrm L}_{J,J}(y_J,\Phi)(g)$ by \eqref{pcx:eq-twist}. With
$f=[y_J,\chi_S^{-1}\circ\det]_S$:
\begin{equation}\label{pcx:eq-A}
\int_{[\auU(J)]}\theta^{J,J}_{\gamma,\gamma}(g,h;f)\,dh=c_1\,\mathrm{vol}_0\qquad(g\in\auU(J)(\auA)).
\end{equation}
\emph{The triple $(J,0,J)$.} Apply Proposition~\ref{pcx:P-matches} with $(W,U,V)=(J,J,0)$, $z=v$,
$y=y_J$. It gives $C$ and a vector $x$ of the zero Fock space $\auHom(J,0)\otimes\R$, and for each
$\nu$ a number $c_2$ with $C|c_2|^2=L_*(\pcm,\nu,-k)\pcxstr_Y(\nu;v)$. Take $\Phi=1$ on
$\pcMat_{0\times2k}$ and $\Psi=\chi_S^{-1}\circ\det$, matched by Lemma~\ref{pcx:std-matched} with
$z=1_{2k}$. On $\auHom(J,0)=0$ the kernel $\theta^{J,0}_{1,\gamma}(g,h;[x,1]_S)$ is the number
$\xi=\theta([x,1]_S)$, which does not depend on $g$, on $h\in\auU(0)$ or on $\nu$ (the datum is $x$
tensored with the constant $1$ at every finite place). Hence
\begin{equation}\label{pcx:eq-B}
\mathrm{vol}_0\,\xi=c_2\int_{[\auU(J)]}\theta^{J,J}_{\gamma,\gamma}(h,g;f)\,dh\qquad(g\in\auU(J)(\auA)).
\end{equation}
\emph{Comparison.} Integrate \eqref{pcx:eq-B} over $g\in[\auU(J)]$ and use Fubini and
\eqref{pcx:eq-A} (with the roles of the variables renamed):
$\mathrm{vol}_J\mathrm{vol}_0\,\xi=c_2\int_{[\auU(J)]}\int_{[\auU(J)]}\theta^{J,J}_{\gamma,\gamma}(h,g;f)
\,dg\,dh=c_2c_1\mathrm{vol}_0\mathrm{vol}_J$, so $\xi=c_1c_2$. Since $L_*(k)L_*(-k)=1$,
\[
\pcxstr_Y(\nu;v)^2=\bigl(L_*(\pcm,\nu,k)\pcxstr_Y(\nu;v)\bigr)\bigl(L_*(\pcm,\nu,-k)\pcxstr_Y(\nu;v)
\bigr)=BC|c_1c_2|^2=BC|\xi|^2=:D,
\]
which does not depend on $\nu$.
\end{proof}

\begin{proof}[Proof of Proposition~\ref{pcx:P-equal} from Propositions~\ref{pcx:P-matches},
\ref{pcx:P-nonvan}, \ref{pcx:P-plane}, \ref{pcx:P-signs} and Lemma~\ref{pcx:adapted}]
If there is no tame character of weight $\mu$ there is nothing to prove. Otherwise fix one, $\nu_0$.

\emph{First floor: $W=J$ the standard positive definite space of rank $2k$.} By
Proposition~\ref{pcx:P-nonvan} there is $z$ with $\pcxstr_Y(\nu_0;z)\ne0$. By
Lemma~\ref{pcx:twice}, $\pcxstr_Y(\nu;z)^2=\pcxstr_Y(\nu_0;z)^2$ for all $\nu$; as
$\pcxstr_Y\ge0$ by definition (Lemma~\ref{pcx:Lpos}), $\pcxstr_Y(\nu;z)=\pcxstr_Y(\nu_0;z)=:C_0\ne0$.

\emph{Induction on $q$.} For $q=0$ the space $\pcxH\oplus W$ has the signs of $Y=\pcxH\oplus J$, and
$Y$ contains the adapted plane $\pcxH$; Proposition~\ref{pcx:P-signs} gives
$\pcxstr_{\pcxH\oplus W}(\nu;z)=\pcxstr_Y(\nu;z)=C_0$. For $q+1$: let $W_0=\pcxH^q\oplus J$, of rank
$2k+2q$ with $q$ negative signs; by induction there are $z_0$ and $C\ne0$ with
$\pcxstr_{\pcxH\oplus W_0}(\nu;z_0)=C$ for all $\nu$. By Lemma~\ref{pcx:adapted} there is a space $Y'$
with the signs of $\pcxH\oplus W$ and an adapted splitting $Y'=(\pcxH\oplus W_0)\oplus\pcxH$
(the ranks differ by $2$ and the numbers of negative signs by $1$). Proposition~\ref{pcx:P-plane}
gives $z'$ and $C'>0$ with $\pcxstr_{Y'}(\nu;z')=C'C$, and Proposition~\ref{pcx:P-signs} gives
$\pcxstr_{\pcxH\oplus W}(\nu;z')=\pcxstr_{Y'}(\nu;z')=C'C\ne0$.
\end{proof}

\subsubsection{Convergence}\label{pcx:ss-conv}

For a space $S_1$ of rank $n_1$, a space $S_2$ of rank $n_2$ and a set $P$ of primes of $E$ with
$P\cap\bar P=\emptyset$, consider data $[y,\Phi]_P$ of $\auHom(S_1,S_2)$, $\Phi$ locally constant
with compact support on $\pcMat_{n_2\times n_1}(F_P)$. Write $\mathrm{ht}$ for the height on the
adelic unitary groups ({\cite[\S14.2]{P46}}).

\begin{proposition}[Decay under a rank condition]\label{pcx:decay}
\textup{(a)} Suppose that every $x$ with $\Phi(x)\ne0$ is surjective at every place of $P$ (each
component has a right inverse). For every $N_0$ there are $B,M,C$ such that for all
Fock vectors $y$ with finite moment $\|y\|_M=\int(1+|z|)^M|y(z)|\,dz$, all characters
$\alpha,\beta$ and all $(g,h)$,
\[
|\theta^{S_1,S_2}_{\alpha,\beta}(g,h;[y,\Phi]_P)|\,\mathrm{ht}(h)^{N_0}\le C\,\|y\|_M\,
\mathrm{ht}(g)^B .
\]
\textup{(b)} If every $x$ with $\Phi(x)\ne0$ is injective at every place of $P$, the same holds with
the roles of $g$ and $h$ exchanged.
\end{proposition}
\noindent The proof occupies the rest of this subsection. Heights are those of
Definition~\ref{pcx:def-ht} below (not the heights of families of {\cite[\S14.2]{P46}}).

\begin{definition}[Lattice and height at a point]\label{pcx:def-ht}
Let $S_1$ be a space of rank $n$ with frame giving $S_1(\auA)\cong\auA_E^n$. For $g\in\auU(S_1)(\auA)$
the \emph{lattice at $g$} is $\Lambda_g=\{\gamma\in S_1(E): g^{-1}\gamma\text{ is integral at all
finite places}\}$, with the length $|\gamma|_g=\|(g^{-1}\gamma)_\infty\|$ (the euclidean norm of the
archimedean coordinates). Let $\lambda_1(g)=\min\{|\gamma|_g:0\ne\gamma\in\Lambda_g\}$ and
$\mathrm{ht}(g)=\max(1,\lambda_1(g)^{-1})$.
\end{definition}

We use three classical facts.
\begin{enumerate}
\item[(R1)] (Reduction theory.) For $N_0$ large, $\int_{[\auU(S_1)]}\mathrm{ht}(g)^{-N_0}dg<\infty$; and
there is $Z$ such that $\Lambda_g$ has an $E$-basis $c_1,\dots,c_n$ with $|c_j|_g\le Z\,\mathrm{ht}(g)^n$
(Minkowski's second theorem: $\Lambda_g$ is a $\Z$-lattice of rank $n[E:\Q]$ and covolume independent
of $g$, so the product of its $n[E:\Q]$ successive minima is bounded above and below. It is an
$\auO_E$-module and $|\omega\gamma|_g\le C_E|\gamma|_g$ for $\omega$ in a fixed $\Z$-basis of $\auO_E$,
so if $\mu_1\le\dots\le\mu_n$ are the successive minima over $E$ (the least $\mu$ such that the
vectors of length $\le\mu$ span an $E$-subspace of dimension $j$), the $\Z$-minima are, up to
constants depending only on $E$, the $\mu_j$ each repeated $[E:\Q]$ times. Hence
$\prod_j\mu_j^{[E:\Q]}\asymp1$, and with $\mu_j\ge\lambda_1(g)\ge\mathrm{ht}(g)^{-1}$ this gives
$\mu_n\ll\mathrm{ht}(g)^{n-1}\le\mathrm{ht}(g)^n$; a basis realising the $\mu_j$ is an $E$-basis). (Borel, \emph{Introduction aux groupes arithm\'etiques}, reduction theory.)
\item[(R2)] (Moderate growth of theta functions; Lemma~\ref{pcx:thetagrowth} below.) Let
$X=\auHom(S_1,S_2)$, $2D=\dim_\R(X\otimes\R)$, $s=D+1$, and let $\|\cdot\|_s$ and $\pzF_s$ be the norm and
the Banach space of {\cite[Lemma~8.2]{P49}} for the archimedean Fock space of $X$. A \emph{datum of class
$s$} is a finite sum of tensors $\Psi\otimes f$ with $\Psi\in\pzF_s$ and $f$ a finite datum; $\theta$ is
defined on these and continuous in the first factor ({\cite[Lemma~8.2(c)]{P49}}), and the Weil
representation preserves them ({\cite[Lemma~8.2(a)]{P49}} at infinity). For a linear functional
$\Lambda$ on the data of class $s$ and a finite datum $f_0$ put
\[
\|\Lambda\|_{f_0}=\sup\Bigl\{\frac{|\Lambda(\psrho_\infty(z)\Psi\otimes f_0)|}{\|\Psi\|_s}:
z\in X\otimes\R,\ 0\ne\Psi\in\pzF_s\Bigr\}\in[0,\infty]
\]
($\psrho_\infty(z)$ preserves $\pzF_s$ by Lemma~\ref{pcx:superpos}(a)). As $\psrho_\infty(y)\psrho_\infty(z)$
is a scalar of absolute value $1$ times $\psrho_\infty(y+z)$, $\|\Lambda\circ\psrho_\infty(y)\|_{f_0}=
\|\Lambda\|_{f_0}$ for $y\in X\otimes\R$. Lemma~\ref{pcx:thetagrowth} says: for every finite datum $f_0$
there are $Y,A$ with
\[
\|\theta\circ\omega_{\alpha,\beta}(g,h)\|_{f_0}\le Y(\mathrm{ht}(g)\mathrm{ht}(h))^A
\]
for all characters $\alpha,\beta$ and all $g,h$. Since $\mathrm{ht}\ge1$ we may and do take $A\ge2s$.
\item[(R3)] (Short vectors are isotropic.) If $\gamma,\gamma'\in\Lambda_g$ and $|\gamma|_g|\gamma'|_g<1$
then $\langle\gamma,\gamma'\rangle=0$: indeed $\langle\gamma,\gamma'\rangle\in\auO_E$ (the lattice at $g$
is unimodular, $g$ being unitary at the finite places), and each archimedean absolute value is at
most $|\gamma|_g|\gamma'|_g<1$ (Cauchy--Schwarz in the coordinates at $g$, where the form is
$\mathrm{diag}(\pm1)$); an algebraic integer all of whose conjugates have absolute value $<1$ is $0$.
\end{enumerate}
In particular, if $\mathrm{ht}(h)>2$ then $\lambda_1(h)<1/2$ and there is $0\ne d\in\Lambda_h$ with
$|d|_h<2\lambda_1(h)\le 2/\mathrm{ht}(h)<1$, isotropic by (R3).

\begin{lemma}[Reduction to a Siegel set]\label{pcx:reduce}
Let $S_1$ be a space of rank $n$ with a frame, $G=\auU(S_1)$, viewed as a connected reductive group over
$E^+$, and $d=[E^+:\Q]$. There are a compact set $C_f\subset G(\auA_f)$ and a constant $C$ such that every
$g\in G(\auA)$ can be written $g=\gamma^{-1}x$ with $\gamma\in G(E^+)$, $x=x_\infty x_f$, $x_f,x_f^{-1}\in C_f$
and $\|x_\infty\|,\|x_\infty^{-1}\|\le C\,\mathrm{ht}(g)$ (operator norms in the
coordinates of the frame).
\end{lemma}
\begin{proof}
If $S_1$ is anisotropic, $G(\auA)/G(E^+)$ is compact~\cite[Theorem~5.6(ii)]{BorelFin}, so $G(\auA)=
C\,G(E^+)$ with $C$ compact; apply this to $g^{-1}$. Otherwise let $r\ge1$ be the Witt index and
$e_1,\dots,e_r,e'_1,\dots,e'_r$ rational vectors with $\langle e_i,e'_j\rangle=\delta_{ij}$, $\langle e_i,e_j
\rangle=\langle e'_i,e'_j\rangle=0$, spanning a subspace with anisotropic orthogonal complement $S_0$. The
torus $T$ acting by $e_i\mapsto a_ie_i$, $e'_i\mapsto a_i^{-1}e'_i$ ($a_i\in\mathrm{GL}_1$ over $E^+$) and
trivially on $S_0$ is a maximal $E^+$-split torus, and the stabiliser $P$ of the flag $Ee_1\subset Ee_1+
Ee_2\subset\cdots$ is a minimal parabolic subgroup containing it; the roots of $T$ in the Lie algebra of
the unipotent radical of $P$ are $a_i/a_j$ $(i<j)$, $a_ia_j$, $a_i^2$ and, if $S_0\ne0$, $a_i$, so the simple
roots are $a_i/a_{i+1}$ $(i<r)$ and $a_r$ or $a_r^2$. By Godement's form of Borel's
theorem~\cite[Th\'eor\`eme~7]{Godement63} there are a compact $M\subset G(\auA)$ with $G(\auA)=MP(\auA)$, a
relatively compact $F\subset P(\auA)^0$ with $P(\auA)^0=F\,P(E^+)$ (here $P(\auA)^0$ is the common kernel of
the $|\chi|$, $\chi$ a rational character of $P$; it is normal in $P(\auA)$ and contains $P(E^+)$ by the
product formula) and $c\ge1$ such that $G(\auA)=M\,T_\auA(c)\,F\,G(E^+)$, where $T_\auA(c)$ is the set of $t\in
T(\auA)$ with $|\alpha(t)|<c$ for the simple roots $\alpha$ (idelic absolute value).

Write $g^{-1}=mtf\gamma_0$ accordingly. \emph{Normalising $t$.} Let $X_f\subset\auA_{E^+,f}^\times$ be a
finite set of representatives of the ideal classes of $E^+$. Every idele $a$ of $E^+$ is $\xi b$ with
$\xi\in E^{+\times}$, $b_f\in X_f\widehat\auO_{E^+}^\times$, and, by Dirichlet's unit theorem (the logarithms
of the units form a lattice in the trace-zero hyperplane), $C_1^{-1}|a|^{1/d}\le|b_\tau|\le C_1|a|^{1/d}$ at
every real place $\tau$ ($b\mapsto\xi\varepsilon b$, $\varepsilon$ a unit, $|a|$ the idelic absolute value, which
equals $|b|$). Doing this for each $a_i(t)$ gives $t=\xi t'$ with $\xi\in T(E^+)$, $t'_f$ in a fixed
compact set, and $|a_i(t'_\tau)|\in[C_1^{-1},C_1]\,|a_i(t)|^{1/d}$; and $|\alpha(t')|=|\alpha(t)|<c$. Then
$g^{-1}=mt'(\xi f\xi^{-1})\xi\gamma_0$ and $\xi f\xi^{-1}\in P(\auA)^0=FP(E^+)$, say $\xi f\xi^{-1}=f'p$.
So $g^{-1}=mt'f'\gamma'$ with $\gamma'=p\xi\gamma_0\in G(E^+)$, $f'\in F$, $t'\in T_\auA(c)$ normalised.

\emph{Bounds.} Put $\gamma=\gamma'$ and $x=\gamma g=f'^{-1}t'^{-1}m^{-1}$. Its finite part and that of its
inverse $mt'f'$ lie in a fixed compact set $C_f$ ($M$, $\bar F$ and the finite parts of the normalised
$t'$ do). From $|a_i/a_{i+1}|<c$ and $|a_r|<c$ (or $|a_r|^2<c$), with $a_i=a_i(t)$: $|a_i|<c^r$ and
$|a_i|^{-1}<c^{r}|a_1|^{-1}$ for all $i$. The archimedean eigenvalues of $t'$ in the basis
$(e_i,e'_i,\text{basis of }S_0)$ are $a_i(t'_\tau)^{\pm1}$ and $1$, so, the change of basis to the frame being
fixed, $\|x_\infty\|,\|x_\infty^{-1}\|\le C_2\max_i\max(|a_i|^{1/d},|a_i|^{-1/d})\le C_3\max(1,|a_1|^{-1/d})$.
Finally the vector $\eta=N\gamma'^{-1}e_1\in S_1(E)$ satisfies $g^{-1}\eta=Nmt'f'e_1=N\,\chi(f')a_1(t')\,me_1$,
where $f'e_1=\chi(f')e_1$ ($P$ stabilises $Ee_1$) with $\chi(f')$ in a relatively compact subset of
$\auA_E^\times$. The finite parts of $\chi(f')a_1(t')me_1$ lie in a compact subset of $S_1(\auA_f)$, hence in
$N^{-1}$ times the integral points for a fixed integer $N\ge1$; so $\eta\in\Lambda_g$, and
$|\eta|_g\le C_4|a_1|^{1/d}$. Thus $\lambda_1(g)\le C_4|a_1|^{1/d}$, i.e.\ $|a_1|^{-1/d}\le C_4\,\mathrm{ht}(g)$,
so $\|x_\infty\|,\|x_\infty^{-1}\|\le C_3\max(1,C_4\,\mathrm{ht}(g))\le C\,\mathrm{ht}(g)$, since
$\mathrm{ht}(g)\ge1$.
\end{proof}

\begin{lemma}[Growth of theta functions]\label{pcx:thetagrowth}
Let $S_1,S_2$ be spaces, $X=\auHom(S_1,S_2)$, $D=\frac12\dim_\R(X\otimes\R)$, $s=D+1$, $\pzF_s$ and
$\|\cdot\|_s$ as in {\cite[Lemma~8.2]{P49}} for $W=X$, and $f_0$ a finite datum. There are $Y,A$ such that for
all characters $\alpha,\beta$, $g\in\auU(S_1)(\auA)$, $h\in\auU(S_2)(\auA)$, $z\in X\otimes\R$ and
$\Psi\in\pzF_s$,
\[
|\theta(\omega_{\alpha,\beta}(g,h)(\psrho_\infty(z)\Psi\otimes f_0))|\le Y\,\|\Psi\|_s\,
(\mathrm{ht}(g)\mathrm{ht}(h))^{A};\qquad A=2s\text{ will do.}
\]
\end{lemma}
\begin{proof}
\emph{{\cite[Lemma~8.2]{P49}} applies to $X$.} Its statement and proof concern $W=\auHom(Y_j,Y_{j+1})$ but
use only: the Fock realisation of the archimedean Heisenberg representation of $W$ for a positive complex
structure commuting with the maximal compact subgroups, the commutation of $\omega_\infty$ with the
Heisenberg operators (W1), the continuity of the Weil representation on smooth vectors, the existence of
a non-zero theta functional continuous in the archimedean factor (W0), and the action of a lattice on
finite data by characters (W4); none of the properties of the chain $Y_j$ enters. All of these hold
for any pair of spaces, so (a), (b), (c) hold for $X$ with $D$ as above.

\emph{Uniformity in $z$.} In the proof of (c), replacing $\Psi$ by $\psrho_\infty(z)\Psi$ replaces
$v=\gamma^\natural-w_0$ by $v+z$ (up to a phase), so $|\theta(\psrho_\infty(z)\Psi\otimes f)|\le C(D+1)^{2s}\|f\|
\|f'\|\,\|\Psi\|_s\sup_{u}\sum_{\lambda}(1+|\lambda+u|)^{-2s}$, the sum over a coset of a fixed lattice of
$X\otimes\R\cong\R^{2D}$ and the supremum over $u\in\R^{2D}$; this supremum is finite as $2s>2D$. The
lattice depends on $f$ only through the integer $\mathsf N$ of that proof. Hence for every finite set
$\mathcal S$ of finite data there is $C_{\mathcal S}$ with $|\theta(\psrho_\infty(z)\Psi\otimes f)|\le
C_{\mathcal S}\|\Psi\|_s$ for $f\in\mathcal S$, all $z$, all $\Psi$.

\emph{Reduction.} By Lemma~\ref{pcx:reduce} write $g=\gamma^{-1}x$, $h=\delta^{-1}y$ with rational
$\gamma,\delta$. The theta functional is invariant under the rational points ((W1) of {\cite[Lemma~5.2]{P49}})
and the characters are trivial on them, so the left side is
$|\theta(\omega_{\alpha,\beta}(x,y)(\psrho_\infty(z)\Psi\otimes f_0))|=|\theta(\psrho_\infty(z')\,\omega_\infty
(x_\infty,y_\infty)\Psi\otimes\omega_f(x_f,y_f)f_0)|$ with $z'=y_\infty zx_\infty^{-1}$, the characters
contributing factors of absolute value $1$. As $f_0$ is fixed by an open subgroup $K'$ of the finite
group and $(x_f,y_f)$ runs over a compact set, covered by finitely many cosets of $K'$, the vectors
$\omega_f(x_f,y_f)f_0$ lie, up to factors of absolute value $1$, in a finite set $\mathcal S$. By the
previous step and {\cite[Lemma~8.2(a)]{P49}},
the left side is at most $C_{\mathcal S}(2D\|x_\infty\|\|y_\infty\|)^{2s}\|\Psi\|_s\le Y(\mathrm{ht}(g)
\mathrm{ht}(h))^{2s}\|\Psi\|_s$ by the bounds of Lemma~\ref{pcx:reduce}.
\end{proof}

\begin{lemma}[Fock vectors are superpositions of translated vacua]\label{pcx:superpos}
Let $\mathcal F$ be the Fock space of a real hermitian space $\C^m$ (real dimension $2m$), realised as in
{\cite[\S5.1.4]{P54}} as the space of functions $v\in L^1\cap L^2$ on $\C^m$ with $v(x)=\int v(y)\beta(y,x)
e(x-y)\,dy$, where $e(x)=e^{-|x|^2/2}$ is the vacuum, $|\beta|=1$, and the Heisenberg translations are
$(\rho(y)v)(x)=\beta(y,x)v(x-y)$; the measure is the normalised Lebesgue measure. For $N\ge0$ put
$\|v\|_N=\int(1+|x|)^N|v(x)|\,dx\in[0,\infty]$. Let $s\ge0$ be an integer, and let $\|\cdot\|_s$ and
$\mathcal F_s\subseteq\mathcal F$ be the norm and the Banach space of {\cite[Lemma~8.2]{P49}} (with $D=m$),
transported by the unitary isomorphism with the model {\cite[(41)]{P49}} that intertwines the
Heisenberg translations and carries vacuum to vacuum; the operators $\partial_w$, being the derivatives
of $\rho$, correspond. Let $c_\rho$ be a constant with
$\rho(w)\rho(y)=e^{\kappa(w,y)}\rho(y)\rho(w)$, $\kappa(w,y)\in i\R$ real bilinear,
$|\kappa(w,y)|\le c_\rho|w||y|$ (in the model {\cite[(41)]{P49}}, $\kappa(w,y)=2i\Iml\operatorname{tr}(y^*w)$ by the
formula for $\psrho_\infty(x)\psrho_\infty(y)$ recalled after it, so $c_\rho=2$ will do).
\begin{enumerate}
\item[\textup{(a)}] For $y\in\C^m$, $\rho(y)$ maps $\mathcal F_s$ into itself, with
$\|\rho(y)\Psi\|_s\le(1+c_\rho|y|)^{2s}\|\Psi\|_s$, and $y\mapsto\rho(y)\Psi$ is continuous into
$\mathcal F_s$ for every $\Psi\in\mathcal F_s$.
\item[\textup{(b)}] If $\|v\|_N<\infty$ then $v=\int v(y)\rho(y)e\,dy$, a Bochner integral converging for
the norm $\|\cdot\|_N$, because $\|\rho(y)e\|_N\le(1+|y|)^N\|e\|_N$. If $\|v\|_{2s}<\infty$, then
$v\in\mathcal F_s$, the same integral converges in $\mathcal F_s$, and $\|v\|_s\le c_1\|v\|_{2s}$ with
$c_1=\max(1,c_\rho)^{2s}\|e\|_s$.
\item[\textup{(c)}] Let $\ell$ be a continuous linear functional on $\mathcal F_s$ with $|\ell(\rho(z)e)|\le
S(1+|z|)^N$ for all $z$. Then $|\ell(v)|\le S\|v\|_N$ for every $v$ with $\|v\|_{\max(N,2s)}<\infty$.
\end{enumerate}
\end{lemma}
\begin{proof}
(a) Let $u$ be a smooth vector and $w\in\mathsf B$ (the real basis defining $\|\cdot\|_s$).
Differentiating $\rho(tw)\rho(y)u=e^{t\kappa(w,y)}\rho(y)\rho(tw)u$ at $t=0$ gives $\partial_w\rho(y)u=
\rho(y)(\partial_w+\kappa(w,y))u$; by induction,
\[
\partial_{w_1}\cdots\partial_{w_n}\rho(y)u=\rho(y)\,(\partial_{w_1}+\kappa_1)\cdots(\partial_{w_n}+\kappa_n)u,
\qquad\kappa_i=\kappa(w_i,y).
\]
Expanding the product, the right side is $\rho(y)$ applied to a sum over the subsets $\Sigma\subseteq
\{1,\dots,n\}$ of $\prod_{i\in\Sigma}\kappa_i$ times $\partial_{w_{j_1}}\cdots\partial_{w_{j_r}}u$, where
$j_1<\dots<j_r$ are the indices not in $\Sigma$. For $n\le2s$ each such word has norm at most
$\|u\|_s$, $|\prod_{i\in\Sigma}\kappa_i|\le(c_\rho|y|)^{|\Sigma|}$, and $\rho(y)$ is unitary; so
$\|\partial_{w_1}\cdots\partial_{w_n}\rho(y)u\|\le(1+c_\rho|y|)^n\|u\|_s$, and $\|\rho(y)u\|_s\le
(1+c_\rho|y|)^{2s}\|u\|_s$. The same identity shows that $y\mapsto\partial_{w_1}\cdots\partial_{w_n}
\rho(y)u$ is continuous into $\mathcal F$ ($\rho$ is strongly continuous on $\mathcal F$ and the $\kappa_i$
are continuous in $y$), i.e.\ $y\mapsto\rho(y)u$ is continuous into $\mathcal F_s$. Since the smooth
vectors are dense in $\mathcal F_s$ and $\rho(y)$ is bounded on them by a constant that is locally
uniform in $y$, both statements extend to $\mathcal F_s$ (consistently with $\rho(y)$ on $\mathcal F$, as
$\|\cdot\|\le\|\cdot\|_s$).

(b) The reproducing formula says $v(x)=\int v(y)(\rho(y)e)(x)\,dy$ for every $x$. From $1+|x|\le
(1+|y|)(1+|x-y|)$ and $|(\rho(y)e)(x)|=e(x-y)$ we get $\|\rho(y)e\|_N\le(1+|y|)^N\|e\|_N$; the map
$y\mapsto\rho(y)e$ is continuous for $\|\cdot\|_N$ and $\int|v(y)|\|\rho(y)e\|_N\,dy\le\|e\|_N\|v\|_N$, so
the Bochner integral exists in the space of functions with finite $\|\cdot\|_N$; evaluating it
pointwise (which is legitimate in $L^1_{\mathrm{loc}}$) gives $v$. If $\|v\|_{2s}<\infty$, the map
$y\mapsto v(y)\rho(y)e$ is continuous into $\mathcal F_s$ by (a), with norm at most
$|v(y)|(1+c_\rho|y|)^{2s}\|e\|_s\le c_1(1+|y|)^{2s}|v(y)|$, which is integrable; so the Bochner integral
exists in $\mathcal F_s$, with norm at most $c_1\|v\|_{2s}$. Its image in $\mathcal F$ is the $\mathcal
F$-valued integral, which is $v$: pairing with $\rho(x)e$ and using $c(x)=\langle\rho(x)e,c\rangle$ for
$c\in\mathcal F$ ({\cite[\S5.1.4]{P54}}) gives $\int v(y)(\rho(y)e)(x)\,dy=v(x)$.

(c) By (b) with the index $\max(N,2s)$, $v=\int v(y)\rho(y)e\,dy$ in $\mathcal F_s$, and $\ell$, being
continuous, commutes with the Bochner integral: $|\ell(v)|\le\int|v(y)||\ell(\rho(y)e)|\,dy\le
S\int(1+|y|)^N|v(y)|\,dy=S\|v\|_N$.
\end{proof}

\begin{lemma}[Short differences gain powers]\label{pcx:diffgain}
Keep the notation of Lemma~\textup{\ref{pcx:superpos}}. For a linear functional $\ell$ on $\mathcal F_s$
put
\[
\|\ell\|^*=\sup\bigl\{|\ell(\rho(z)\Psi)|/\|\Psi\|_s:\ z\in\C^m,\ 0\ne\Psi\in\mathcal F_s\bigr\}\in[0,\infty].
\]
For each $N\ge0$ there is $C$, depending only on $N$, $m$, $s$ and $c_\rho$, with the following
property. Let $t\ge0$ and let $\ell$ be a linear functional on $\mathcal F_s$ with $\|\ell\|^*<\infty$
(so $\ell$ is continuous) such that, if $t<1$: for every $f\in\mathcal F_s$ and every $B\ge0$, if
$|\ell(f)-\ell(\rho(y)f)|\le B$ for all $|y|\le t$, then $|\ell(f)|\le B$. Then
\[
|\ell(v)|\le C\,\|\ell\|^*\,t^N\,\|v\|_N\qquad\text{for every }v\in\mathcal F\text{ with }\|v\|_{\max(N,2s)}<\infty .
\]
\end{lemma}
\begin{proof}
If $t\ge1$: $|\ell(\rho(z)e)|\le\|\ell\|^*\|e\|_s\le\|\ell\|^*\|e\|_st^N$ for all $z$, and
Lemma~\ref{pcx:superpos}(c) with the index $0$ gives $|\ell(v)|\le\|\ell\|^*\|e\|_st^N\|v\|_0\le
\|\ell\|^*\|e\|_st^N\|v\|_N$.

Let $t<1$. For a polynomial $P$ on $\C^m$ let $v_P$ be the vector of $\mathcal F$ that corresponds to $P$
in the model {\cite[(41)]{P49}} (a smooth vector, $v_1=e$). For $w\in\C^m$, $\partial_wv_P=v_{P_w}$
with $P_w(z)=(z\cdot w)P(z)-(\mathsf D_{\bar w}P)(z)$ (proof of {\cite[Lemma~8.2(b)]{P49}}), a polynomial
depending $\R$-linearly on $w$; write $\mathsf B=(b_1,\dots,b_{2m})$, so that $\partial_y=\sum_ky_k
\partial_{b_k}$ for $y=\sum_ky_kb_k$, $|y_k|\le|y|$. We show by induction on $N$:

\smallskip\noindent$(H_N)$ \emph{for every polynomial $P$ there is $C_{N,P}$, depending only on
$N,P,m,s,c_\rho$, with $|\ell(\rho(z)v_P)|\le C_{N,P}\|\ell\|^*t^N(1+|z|)^N$ for all $z\in\C^m$.}

\smallskip\noindent For $N=0$ take $C_{0,P}=\|v_P\|_s$. Assume $(H_N)$, fix $P$, $z$ and $y$ with
$|y|\le t$, and put $f(\sigma)=\ell(\rho(\sigma y)\rho(z)v_P)$ for $\sigma\in[0,1]$. With
$\kappa=\kappa(y,z)$ we have $\rho(\sigma y)\rho(z)=e^{\sigma\kappa}\rho(z)\rho(\sigma y)$; the orbit map
$\sigma\mapsto\rho(\sigma y)v_P$ of the smooth vector $v_P$ is $C^1$ for the topology of smooth vectors,
hence into $\mathcal F_s$, with derivative $\rho(\sigma y)\partial_yv_P$ ($\rho$ is additive on the line
$\R y$); $\rho(z)$ is bounded on $\mathcal F_s$ and $\ell$ is continuous. Hence
\[
f'(\sigma)=e^{\sigma\kappa}\,\ell\bigl(\rho(z)\rho(\sigma y)(\kappa v_P+\partial_yv_P)\bigr),
\]
and $\rho(z)\rho(\sigma y)$ is a scalar of absolute value $1$ times $\rho(z+\sigma y)$. Since
$|\kappa|\le c_\rho|y||z|\le c_\rho t(1+|z|)$, $|y_k|\le t$ and $1+|z+\sigma y|\le2(1+|z|)$, $(H_N)$
for $P$ and for the $P_{b_k}$ gives
\[
|f'(\sigma)|\le2^N\Bigl(c_\rho C_{N,P}+\sum_kC_{N,P_{b_k}}\Bigr)\|\ell\|^*t^{N+1}(1+|z|)^{N+1}
=:C_{N+1,P}\|\ell\|^*t^{N+1}(1+|z|)^{N+1}.
\]
By the mean value inequality, $|\ell(\rho(z)v_P)-\ell(\rho(y)\rho(z)v_P)|=|f(0)-f(1)|$ is at most this
bound, for every $|y|\le t$. The hypothesis on $\ell$, applied to $f=\rho(z)v_P\in\mathcal F_s$, gives
$|\ell(\rho(z)v_P)|\le C_{N+1,P}\|\ell\|^*t^{N+1}(1+|z|)^{N+1}$, which is $(H_{N+1})$.

Finally $(H_N)$ for $P=1$ and Lemma~\ref{pcx:superpos}(c) with $S=C_{N,1}\|\ell\|^*t^N$ give the lemma
with $C=\max(C_{N,1},\|e\|_s)$.
\end{proof}

\begin{lemma}[Comparison with the mean over an isotropic subspace]\label{pcx:mean}
Let $X$ be a space (here $X=\auHom(S_1,S_2)$), $s$, the data of class $s$ and $\|\cdot\|_{f_0}$ as in
\textup{(R2)}, $f_0$ a finite datum (a finite combination of finite Heisenberg translates of the
vacuum), $l\ge0$ and $\iota$ a finite set. There is $C$ such that: if $k\in\auU(X)(\auA)$ (acting
through the Weil representation), $\delta\ge0$, $z_i\in\Lambda_k$ $(i\in\iota)$ with $|z_i|_k\le\delta$
span an isotropic subspace $Z\subseteq X(E)$, $\Lambda$ is a linear functional on the data of class $s$,
invariant under $\rho(Z(E))$, such that $x\mapsto\Lambda(\rho(x)F)$ is continuous on $Z(\auA)$ for every
datum $F$ of class $s$ and $\|\Lambda\circ\omega(k)\|_{f_0}<\infty$, and
$\Lambda_0(F)=\int_{Z(\auA)/Z(E)}\Lambda(\rho(x)F)\,dx$ (probability Haar measure on the compact group
$Z(\auA)/Z(E)$), then for every Fock vector $v$ of $X\otimes\R$ with $\|v\|_{\max(l,2s)}<\infty$,
\[
|(\Lambda-\Lambda_0)(\omega(k)(v\otimes f_0))|\le C\,\|\Lambda\circ\omega(k)\|_{f_0}\,\delta^l\,\|v\|_l .
\]
(Here $v\in\pzF_s$ by Lemma~\textup{\ref{pcx:superpos}(b)}, so $v\otimes f_0$ is a datum of class $s$.)
\end{lemma}
\begin{proof}
Since $Z$ is isotropic, $\rho$ is additive on $Z(\auA)$ (no cocycle), so for $F$ of class $s$ the
function $x\mapsto\Lambda(\rho(x)F)$ is $Z(E)$-periodic and continuous; $\Lambda_0$ is a linear
functional on the data of class $s$, invariant under $\rho(Z(\auA))$ (invariance of Haar measure), in
particular under $\rho(Z(E))$. Put $\Lambda'=\Lambda-\Lambda_0$ and $\ell(u)=\Lambda'(\omega(k)(u\otimes
f_0))$ for $u\in\pzF_s$. Then $\Lambda'$ is invariant under $\rho(Z(E))$ and has mean zero:
$\int_{Z(\auA)/Z(E)}\Lambda'(\rho(x)F)\,dx=\Lambda_0(F)-\Lambda_0(F)=0$.

\emph{Rational translates are real translations.} There is $M\ge1$ such that $\rho(x_f)f_0=f_0$ for
every finite vector $x_f$ in $M$ times the integral points (finite data are locally constant under
finite translations, and the integral points form a compact group). There is $R$ such that every
adele is congruent modulo $E$ to $M$ times an adele with integral finite part and archimedean
components of size $\le R$ (a compact fundamental domain for $E\backslash\auA_E$). Given $x=\sum
c_iz_i\in Z(\auA)$, write $c_i+a_i=Mw_i$, $a_i\in E$, and $\gamma=\sum a_iz_i\in Z(E)$; then
$k^{-1}(\gamma+x)=M\sum w_ik^{-1}z_i=:M\mathsf q$ has integral finite part and archimedean coordinates of
size $\le M\#\iota R\delta=:C_0\delta$. As $\rho$ is additive on the isotropic space $k^{-1}Z(\auA)$ and
$\rho(M\mathsf q)=\rho_\infty(M\mathsf q_\infty)\otimes\rho_f(M\mathsf q_f)$ with $\rho_f(M\mathsf q_f)f_0=f_0$, we get, for every
$u\in\pzF_s$, $\rho(k^{-1}(\gamma+x))(u\otimes f_0)=(\rho_\infty(y_x)u)\otimes f_0$ with
$y_x=M\mathsf q_\infty$, $|y_x|\le C_0\delta$. Let $\Lambda_1$ be $\Lambda$ or $\Lambda_0$ (both are invariant under
$\rho(Z(E))$). Since $\rho(\gamma)\rho(x)=\rho(\gamma+x)$ and $\rho(\gamma+x)\omega(k)=\omega(k)
\rho(k^{-1}(\gamma+x))$ by (W1),
\begin{equation}\label{pcx:eq-meantrans}
\Lambda_1\bigl(\rho(x)\omega(k)(u\otimes f_0)\bigr)=\Lambda_1\bigl(\omega(k)((\psrho_\infty(y_x)u)\otimes
f_0)\bigr)\qquad(x\in Z(\auA),\ u\in\pzF_s).
\end{equation}

\emph{The norm of $\ell$.} Take $u=\psrho_\infty(z)\Psi$ in \eqref{pcx:eq-meantrans} with $\Lambda_1=\Lambda$;
as $\psrho_\infty(y_x)\psrho_\infty(z)$ is a scalar of absolute value $1$ times $\psrho_\infty(y_x+z)$,
$|\Lambda(\rho(x)\omega(k)(\psrho_\infty(z)\Psi\otimes f_0))|\le\|\Lambda\circ\omega(k)\|_{f_0}\|\Psi\|_s$
for all $x$; integrating over $x$, the same bound holds for $\Lambda_0$. Hence, in the notation of
Lemma~\ref{pcx:diffgain} for the Fock space of $X\otimes\R$, $\|\ell\|^*\le2\|\Lambda\circ\omega(k)\|_{f_0}$.

\emph{The hypothesis of Lemma~\ref{pcx:diffgain} with $t=C_0\delta$.} Let $u\in\pzF_s$ and
$F=\omega(k)(u\otimes f_0)$. By the mean zero property, $\ell(u)=\int_{Z(\auA)/Z(E)}(\ell(u)-
\Lambda'(\rho(x)F))\,dx$, and by \eqref{pcx:eq-meantrans} for $\Lambda$ and for $\Lambda_0$,
$\Lambda'(\rho(x)F)=\ell(\psrho_\infty(y_x)u)$ with $|y_x|\le t$. So if $|\ell(u)-\ell(\psrho_\infty(y)u)|
\le B$ for all $|y|\le t$, then $|\ell(u)|\le B$.

\emph{Conclusion.} Lemma~\ref{pcx:diffgain} with $N=l$ gives, for $\|v\|_{\max(l,2s)}<\infty$,
$|\ell(v)|\le C_l\|\ell\|^*(C_0\delta)^l\|v\|_l\le2C_lC_0^l\,\|\Lambda\circ\omega(k)\|_{f_0}\,\delta^l\,
\|v\|_l$.
\end{proof}

\begin{lemma}[The mean kills data supported off the orthogonal]\label{pcx:meanzero}
Let $\Lambda_0$ be a bounded functional on the data of $X=\auHom(S_1,S_2)$ invariant under $\rho(Z(\auA))$
for an isotropic $Z\subseteq X(E)$, and $\Phi$ locally constant compactly supported at $P$. If for
every $x$ with $\Phi(x)\ne0$ there is $\zeta\in Z(\auA)$ whose finite part at $\bar P$ pairs
non-trivially with $x$ (in the pairing of {\cite[(63)]{P49}}), then $\Lambda_0([y,\Phi]_P\otimes f)=0$.
\end{lemma}
\begin{proof}
Here \emph{bounded} is used only through the following: $\Lambda_0$ is linear and
$\Lambda_0(\rho(\zeta)F)=\Lambda_0(F)$ for $\zeta\in Z(\auA)$ and all data $F=[y,\Phi_i]_P\otimes f$
with $\Phi_i$ the indicator of a coset (times $\Phi$). Since $Z(\auA)=Z\otimes_E\auA_E$, with
$\zeta\in Z(\auA)$ also every $a\zeta$, $a\in\auA_E$, lies in $Z(\auA)$, and we may take $a$
supported at the places of $\bar P$ only, so that $a\zeta$ is an element $a'$ of $X(\auA)$ supported
at $\bar P$. By {\cite[(63)]{P49}}, $\rho(a')$ acts on functions of $x$ at $P$ by multiplication by
$\lambda(\mathrm{tr}({}^ta'x))$, where $\lambda=\prod\lambda_w$ with $\lambda_w$ trivial on $\auO_w$
and non-trivial on $\varpi_w^{-1}\auO_w$. Fix $x_0$ with $\Phi(x_0)\ne0$ and $\zeta$ as in the
hypothesis, so $\mathrm{tr}({}^t\zeta x_0)\ne0$ at some place $w$ of $\bar P$. Choosing $a=t$ at $w$
(and $0$ elsewhere) with $t\in E_w$ such that $\lambda_w(t\,\mathrm{tr}({}^t\zeta x_0))\ne1$, the
character $x\mapsto\lambda(\mathrm{tr}({}^t(a\zeta)x))$ is constant, with value $c\ne1$, on a small
coset $x_0+\varpi^{n}\pcMat(O_P)$ (it is trivial on $\varpi^n\pcMat(O_P)$ for $n$ large). Then
$\rho(a\zeta)$ multiplies the datum with the indicator of that coset by $c$, and invariance gives
$\Lambda_0=c\Lambda_0$ on it, so $\Lambda_0$ kills it. By compactness finitely many such cosets cover
the support of $\Phi$ (shrinking them to be disjoint and refining so that $\Phi$ is constant on each);
by linearity $\Lambda_0([y,\Phi]_P\otimes f)=0$.
\end{proof}

\begin{proof}[Proof of Proposition~\ref{pcx:decay}]
(a) Let $A\ge2s$ be as in (R2) for the finite datum $f_0=\Phi\otimes\psvac$ (so $F=[y,\Phi]_P=
y\otimes f_0$), put $l=N_0+A\ge2s$, and take $M=l$; let $y$ have $\|y\|_M<\infty$, so that $y\in\pzF_s$ and
$\|y\|_s\le c_1\|y\|_{2s}\le c_1\|y\|_M$ (Lemma~\ref{pcx:superpos}(b)). If $\mathrm{ht}(h)\le2$, (R2) (with
$z=0$) gives $|\theta(\omega(k)F)|\le\|\theta\circ\omega(k)\|_{f_0}\|y\|_s\le Yc_12^A\|y\|_M\,
\mathrm{ht}(g)^A$, and $\mathrm{ht}(h)^{N_0}\le2^{N_0}$; this is the claim. Fix $g,h$ with
$\mathrm{ht}(h)>2$. Let $d\in\Lambda_h$ be
isotropic with $|d|_h\le2/\mathrm{ht}(h)$, and $c_1,\dots,c_{n_1}$ a basis of $S_1(E)$ in $\Lambda_g$
with $|c_j|_g\le Z\mathrm{ht}(g)^{n_1}$ (R1). The rank-one maps $z_j=d\langle c_j,\cdot\rangle\in X(E)$
lie in the lattice at $k=(g,h)$, have length $\le|d|_h|c_j|_g\le2Z\mathrm{ht}(g)^{n_1}/\mathrm{ht}(h)
=:\delta$, and span the isotropic subspace $Z=\{d\lambda:\lambda\in S_1(E)^*\}$ ($\langle d\lambda,
d\lambda'\rangle$ is a multiple of $\langle d,d\rangle=0$). $\theta$ is invariant under $\rho(Z(E))$.
Lemma~\ref{pcx:mean} with $\Lambda=\theta$ and $l$ as above (its hypotheses hold: $\theta$ is invariant
under $\rho(Z(E))$, $x\mapsto\theta(\rho(x)F')$ is continuous for $F'$ of class $s$ by
{\cite[Lemma~8.2(c)]{P49}} and~Lemma~\ref{pcx:superpos}(a) and the local constancy of finite data,
$\|\theta\circ\omega(k)\|_{f_0}<\infty$ by (R2), and $\|y\|_{\max(l,2s)}=\|y\|_l<\infty$) gives
$|\theta(\omega(k)F)-\Lambda_0(\omega(k)F)|\le C\|\theta\circ\omega(k)\|_{f_0}\delta^l\|y\|_l$. Now
$\Lambda_0(\omega(k)F)=\Lambda_0'(F)$ with $\Lambda_0'=\Lambda_0\circ\omega(k)$ invariant under
$\rho(k^{-1}Z(\auA))$, and the pairing condition of Lemma~\ref{pcx:meanzero} holds: in coordinates at
$k$, the elements of $k^{-1}Z(\auA)$ are $d'\,\lambda'$ with $d'=h^{-1}d\ne0$ and $\lambda'$ running over
all of the dual (because $g^{-1}c$ is an adelic basis), and the pairing of $x$ with $\zeta=d'\lambda'$ is
$\mathrm{tr}({}^t\zeta\,x)=\lambda'({}^td'x)$ (with ${}^td'x$ the row vector of the linear form
$\lambda'$ is evaluated on), which is non-zero for some $\lambda'$ iff ${}^td'x\ne0$; this holds
because $x$ is surjective at $P$ and the component of $d'$ at $\bar P$ is non-zero ($d$ is a
non-zero rational vector). So $\Lambda_0(\omega(k)F)=0$. With (R2),
$\|\theta\circ\omega(k)\|_{f_0}\le Y(\mathrm{ht}(g)\mathrm{ht}(h))^A$, so
\[
|\theta(\omega(k)F)|\le C'\|y\|_l\,\mathrm{ht}(g)^{A+n_1l}\,\mathrm{ht}(h)^{A-l},
\]
which is (a) since $l-A=N_0$. (b) is the same with the roles of $S_1$, $S_2$ exchanged: $d$ a short
isotropic vector of $S_1$ at $g$, $c$ a short basis of $S_2$ at $h$, and the condition $x\,d'\ne0$ from
injectivity.
\end{proof}


\begin{corollary}[Absolute convergence in the exchange identity]\label{pcx:conv}
Assume Proposition~\ref{pcx:decay}. \textup{(i)} If $\Phi$ on $\pcMat_{r\times N}(F_S)$ is supported
on matrices of rank $r$ (at every place of $S$) then $h\mapsto\theta^{W,V}_{\alpha,\beta}(g,h;[v,\Phi]_S)$
is integrable on $[\auU(V)]$, uniformly for $g$ in compact sets, and the integral has moderate
growth in $g$. \textup{(ii)} The same holds for $h'\mapsto\theta^{V',W}_{\alpha,\beta}(h',g;[v,\Psi]_S)$
when $\Psi$ on $\pcMat_{N\times b}(F_S)$ is supported on matrices of rank $b$. \textup{(iii)} In
particular, for $\Phi,\Psi$ matched for $\chi$ (over $\pcR_\pcm$, or over $\auO_S$), both integrals of
\eqref{pcx:eq-X} converge absolutely.
\end{corollary}
\begin{proof}
(i) Proposition~\ref{pcx:decay}(a) with $N_0$ large and the finiteness of $\int_{[\auU(V)]}
\mathrm{ht}(h)^{-N_0}dh$ for large $N_0$ (reduction theory; for $V$ anisotropic $[\auU(V)]$ is
compact). (ii) is (b) for the pair $(V',W)$, the integration being over the source group. (iii)
Matched functions are supported on $\{k^{[a]}\}$, resp.\ $\{k_{[b]}\}$, $k$ invertible, i.e.\ on
matrices of full rank $a$, resp.\ $b$ (Lemma~\ref{pcx:mod} reduces data modulo $\pcm$ to this case).
\end{proof}

\begin{remark}[Which integrals are used downstream]\label{pcx:rem-conv}
In all applications (Theorem~\ref{pc:tri}(b) with $e=t$, Theorem~\ref{pc:tri}(c) with $e=0$,
{\cite[Proposition~2.15]{P54}}, Theorem~\ref{pc:auxcore}) the theta integrals of
Theorem~\ref{pcx:X} are formed with matched finite data, or with point data $\delta_\xi$ at
$\xi$ of full rank; Corollary~\ref{pcx:conv} covers all of them, and no other integrals occur. So
Theorem~\ref{pcx:X} holds with absolutely convergent integrals, as stated in Theorem~\ref{pc:X}.
\end{remark}

\subsubsection{From standard data to matched functions: proof of Proposition~\ref{pcx:P-match}}
\label{pcx:ss-match}

Throughout, $F=F_S$, $O=\auO_S$, $\varpi\in O$ is the product of uniformisers of the factors, and
Haar measures on $\auGL_m(F)$ give $\auGL_m(O)$ volume $1$. A matrix $y\in\pcMat_{N\times b}(F)$ is
\emph{primitive} if it is integral and has an integral left inverse; equivalently (each factor of
$O$ being a discrete valuation ring) $y$ is integral and its reduction modulo every prime of $S$ has
rank $b$, equivalently $y=k_{[b]}$ for some $k\in\auGL_N(O)$. Write $\mathrm{Pr}_b\subset
\pcMat_{N\times b}(F)$ for the primitive matrices and $\mathrm{Pr}^a=\{x\in\pcMat_{a\times N}(F):
x^{\mathsf t}\text{ primitive}\}=\{k^{[a]}:k\in\auGL_N(O)\}$ (apply the first description to
$k^{-1}$). Both sets are open and compact. $\mathrm{SB}(X)$ denotes the locally constant compactly
supported functions with support in $X$; $\mathrm{Inj}_b$ is the set of matrices in
$\pcMat_{N\times b}(F)$ of rank $b$, $\mathrm{Surj}^a$ that of matrices in $\pcMat_{a\times N}(F)$ of
rank $a$ (at every factor).

Let $\xi(t)=\chi_S(t)|t|^{N/2}$ on $F^\times$; then $\xi=\chi_S$ on $O^\times$.

\begin{lemma}[Cosets]\label{pcx:cosets}
\textup{(a)} For $y\in\mathrm{Inj}_b$ the set of $u\in\auGL_b(F)$ with $yu$ primitive is one coset
$u_0\auGL_b(O)$. \textup{(b)} For $x\in\mathrm{Surj}^a$ the set of $u\in\auGL_a(F)$ with $ux\in
\mathrm{Pr}^a$ is one coset $\auGL_a(O)u_0$. \textup{(c)} If $y$ is primitive and $yg$ is primitive,
$g\in\auGL_b(F)$, then $g\in\auGL_b(O)$.
\end{lemma}
\begin{proof}
(a) $\Lambda=\{c\in F^b:yc\in O^N\}$ is an $O$-lattice in $F^b$ ($y$ is injective with a left
inverse $L$, so $\Lambda\subseteq L(O^N)$ is bounded, and it contains a small multiple of $O^b$). The
matrix $yu$ is integral iff $uO^b\subseteq\Lambda$, and then primitive iff $yuO^b=yF^b\cap O^N$, i.e.\
iff $uO^b=\Lambda$; the $u$ with $uO^b=\Lambda$ form one coset of $\auGL_b(O)$. (b) is (a) for
$x^{\mathsf t}$: $ux\in\mathrm{Pr}^a$ iff $x^{\mathsf t}u^{\mathsf t}$ is primitive, so the
$u^{\mathsf t}$ form a coset $u_0'\auGL_a(O)$ and the $u$ form $\auGL_a(O)u_0'^{\mathsf t}$ (a left
coset, as it must be, $\mathrm{Pr}^a$ being stable under left multiplication by $\auGL_a(O)$). (c) is (a) applied with $y$, since $u=1$ and $u=g$ both qualify.
\end{proof}

\begin{lemma}[Right invariance of matched functions]\label{pcx:rightinv}
Let $\Phi,\Psi$ be matched for $\chi_S$ over $O$. Then $\Psi(ym)=\chi_S(\det m)^{-1}\Psi(y)$ and
$\Phi(m'x)=\chi_S(\det m')^{-1}\Phi(x)$ for $m\in\auGL_b(O)$, $m'\in\auGL_a(O)$. Moreover, if $j\in
\auGL_N(F)$ is such that $j_{[b]}$ is primitive and $j^{[a]}\in\mathrm{Pr}^a$, then
$\Phi(j^{[a]})=\xi(\det j)\Psi(j_{[b]})$.
\end{lemma}
\begin{proof}
For $k\in\auGL_N(O)$ let $D=\mathrm{diag}(m,1_a)$; then $(kD)_{[b]}=k_{[b]}m$, $(kD)^{[a]}=k^{[a]}$,
$\det kD=\det k\det m$, and the matching relation for $k$ and $kD$ gives the first claim on the
support $\{k_{[b]}\}$, both sides vanishing elsewhere (the support is stable under right
multiplication by $\auGL_b(O)$). The second claim is the same with $D=\mathrm{diag}(1_b,m'^{-1})$.
For the last claim let $i\in\auGL_N(O)$ with $i^{[a]}=j^{[a]}$. The columns of $i_{[b]}$ and of
$j_{[b]}$ both form a basis of the kernel of $j^{[a]}$, so $j_{[b]}g=i_{[b]}$ with $g\in\auGL_b(F)$,
and $g\in\auGL_b(O)$ by Lemma~\ref{pcx:cosets}(c). The matrix $j^{-1}i$ has last rows $j^{[a]}i=
i^{[a]}i=(0\ 1_a)$ and first columns $j^{-1}j_{[b]}g=\binom g0$, so $\det i=\det j\det g$. Hence
$\Phi(j^{[a]})=\Phi(i^{[a]})=\chi_S(\det i)\Psi(j_{[b]}g)=\xi(\det j)\xi(\det g)\chi_S(\det g)^{-1}
\Psi(j_{[b]})=\xi(\det j)\Psi(j_{[b]})$, using $\xi=\chi_S$ on units and $|\det i|=1$.
\end{proof}

\begin{definition}[Sections and same-section pairs]\label{pcx:def-section}
For $\Phi\in\mathrm{SB}(\mathrm{Surj}^a)$, $x\in\mathrm{Surj}^a$, and $\Psi\in\mathrm{SB}(\mathrm{Inj}_b)$,
$y\in\mathrm{Inj}_b$, put
\[
\mathsf S_\Phi(x)=\int_{\auGL_a(F)}\xi(\det r)\Phi(rx)\,dr,\qquad
\mathsf T_\Psi(y)=\int_{\auGL_b(F)}\xi(\det r)\Psi(yr)\,dr
\]
(the integrands have compact support, since $r\mapsto rx$ and $r\mapsto yr$ are proper). A pair
$(\Phi,\Psi)\in\mathrm{SB}(\mathrm{Surj}^a)\times\mathrm{SB}(\mathrm{Inj}_b)$ is a
\emph{same-section pair} if $\mathsf S_\Phi(k^{[a]})=\xi(\det k)\,\mathsf T_\Psi(k_{[b]})$ for all
$k\in\auGL_N(F)$. They form a vector space $\mathcal P$; let $\mathcal M$ be its image under
$(\Phi,\Psi)\mapsto\Psi$.
\end{definition}

\begin{lemma}\label{pcx:section-laws}
\textup{(a)} If $(\Phi,\Psi)\in\mathcal P$ and $g\in\auGL_N(F)$ then $(\xi(\det g)\Phi(\cdot\,g),
\Psi(g^{-1}\cdot))\in\mathcal P$. In particular $\mathcal M$ is stable under $\Psi\mapsto\Psi(g\,\cdot)$.
\textup{(b)} If $\Phi,\Psi$ are matched for $\chi_S$ over $O$ and $\Psi$ is locally constant, then
$(\Phi,\Psi)\in\mathcal P$.
\end{lemma}
\begin{proof}
(a) The support conditions are preserved. Since $k^{[a]}g=(g^{-1}k)^{[a]}$ and $g^{-1}k_{[b]}=
(g^{-1}k)_{[b]}$,
$\mathsf S_{\xi(\det g)\Phi(\cdot g)}(k^{[a]})=\xi(\det g)\mathsf S_\Phi((g^{-1}k)^{[a]})=
\xi(\det g)\xi(\det g^{-1}k)\mathsf T_\Psi((g^{-1}k)_{[b]})=\xi(\det k)\mathsf T_{\Psi(g^{-1}\cdot)}
(k_{[b]})$.

(b) $\Psi$ is supported on $\mathrm{Pr}_b$, compact open, so $\Psi\in\mathrm{SB}(\mathrm{Inj}_b)$.
$\Phi$ is supported on $\mathrm{Pr}^a$; it is locally constant because the map $\auGL_N(O)\to
\mathrm{Pr}^a$, $k\mapsto k^{[a]}$, is continuous, open and surjective, and $\Phi(k^{[a]})=
\chi_S(\det k)\Psi(k_{[b]})$ is a locally constant function of $k$; so
$\Phi\in\mathrm{SB}(\mathrm{Surj}^a)$. Let $k\in\auGL_N(F)$. By Lemma~\ref{pcx:cosets} choose $u$ with
$k_{[b]}u$ primitive and $u'$ with $u'k^{[a]}\in\mathrm{Pr}^a$. By Lemma~\ref{pcx:rightinv} the
integrand of $\mathsf T_\Psi(k_{[b]})$ vanishes off $u\auGL_b(O)$ and is constant on it, so
$\mathsf T_\Psi(k_{[b]})=\xi(\det u)\Psi(k_{[b]}u)$; likewise the integrand of $\mathsf S_\Phi(k^{[a]})$ vanishes off the left coset
$\auGL_a(O)u'$ and is constant on it (by $\Phi(m'x)=\chi_S(\det m')^{-1}\Phi(x)$), so
$\mathsf S_\Phi(k^{[a]})=\xi(\det u')\Phi(u'k^{[a]})$. Let $j=k\,\mathrm{diag}(u,u'^{-1})$: $j_{[b]}=k_{[b]}u$, $j^{[a]}=u'k^{[a]}$,
$\det j=\det k\det u/\det u'$. By the last part of Lemma~\ref{pcx:rightinv},
$\mathsf S_\Phi(k^{[a]})=\xi(\det u')\xi(\det j)\Psi(j_{[b]})=\xi(\det k)\mathsf T_\Psi(k_{[b]})$.
\end{proof}

\begin{lemma}[Linear factor]\label{pcx:linfactor}
Let $X$ be a vector space and $\Lambda\colon\mathrm{SB}(\mathrm{Surj}^a)\to X$ linear with
$\xi(\det g)\Lambda(\Phi(g\,\cdot))=\Lambda(\Phi)$ for all $g\in\auGL_a(F)$. There is a linear map
$\mathcal T\colon\mathrm{SB}(\mathrm{Inj}_b)\to X$ with $\Lambda(\Phi)=\mathcal T(\Psi)$ for all
$(\Phi,\Psi)\in\mathcal P$.
\end{lemma}
\begin{proof}
Let $\Phi\in\mathrm{SB}(\mathrm{Surj}^a)$. By Lemma~\ref{pcx:cosets}(b), $\mathrm{Surj}^a$ is the
disjoint union of the open sets $r\,\mathrm{Pr}^a$, $r$ running through representatives of
$\auGL_a(F)/\auGL_a(O)$; finitely many of them, $r_1,\dots,r_m$, cover the support of $\Phi$. Put
$\Phi_i=\Phi\cdot\mathbf 1_{r_i\mathrm{Pr}^a}\in\mathrm{SB}(\mathrm{Surj}^a)$. For $k\in\auGL_a(O)$ the
hypothesis gives $\Lambda(\Phi_i)=\xi(\det r_ik)\Lambda(\Phi_i(r_ik\,\cdot))$. The function
$k\mapsto\xi(\det r_ik)\Phi_i(r_ik\,\cdot)$ is locally constant on the compact group $\auGL_a(O)$ with
values in a finite-dimensional space, so averaging over $k$ is a finite sum and commutes with
$\Lambda$:
\[
\Lambda(\Phi)=\sum_i\Lambda(\Phi_i)=\Lambda(\Phi^0),\qquad\Phi^0(x)=\sum_i\int_{\auGL_a(O)}
\xi(\det r_ik)\Phi_i(r_ikx)\,dk .
\]
For $x\in\mathrm{Pr}^a$ we have $r_ikx\in r_i\mathrm{Pr}^a$, so $\Phi_i(r_ikx)=\Phi(r_ikx)$, and
$\Phi(rx)=0$ unless $r$ lies in some $r_i\auGL_a(O)$; hence $\Phi^0(x)=\mathsf S_\Phi(x)$. For
$x\notin\mathrm{Pr}^a$, $r_ikx\notin r_i\mathrm{Pr}^a$ and $\Phi^0(x)=0$. So $\Lambda(\Phi)=
\Lambda(\mathbf 1_{\mathrm{Pr}^a}\mathsf S_\Phi)$. For each $x\in\mathrm{Pr}^a$ fix $k_x\in
\auGL_N(O)$ with $k_x^{[a]}=x$; let $\tilde\Lambda$ be a linear extension of $\Lambda$ to all
functions on $\pcMat_{a\times N}(F)$, and put $\mathcal T(\Psi)=\tilde\Lambda\bigl(x\mapsto\mathbf
1_{\mathrm{Pr}^a}(x)\,\xi(\det k_x)\mathsf T_\Psi((k_x)_{[b]})\bigr)$, which is linear in $\Psi$.
For $(\Phi,\Psi)\in\mathcal P$
the function is $\mathbf 1_{\mathrm{Pr}^a}\mathsf S_\Phi$, so $\mathcal T(\Psi)=\Lambda(\Phi)$.
\end{proof}

\begin{lemma}[Generation by one standard datum]\label{pcx:generate}
Let $\mathcal Q\subseteq\mathrm{SB}(\mathrm{Inj}_b)$ be a subspace stable under $\Psi\mapsto
\Psi(g\,\cdot)$ for $g\in\auGL_N(F)$, and suppose $\Psi^{\mathrm{col}}_s\in\mathcal Q$ for an integral
$s\in\pcMat_{b\times N}(O)$ with $sx_0=1$ for some integral $x_0$. Then every locally constant $\Psi$
having a matched $\Phi$ (for $\chi_S$ over $O$) lies in $\mathcal Q$.
\end{lemma}
\begin{proof}
Write $p=\Psi^{\mathrm{col}}_s$. As $\Psi$ is locally constant with compact support there is $c\ge0$
with $\Psi(\varpi^ct+y)=\Psi(y)$ for all integral $t$ and all $y$. Let $e=x_0s$, an idempotent, and
$d=\varpi^{-c}(1-e)+e\in\auGL_N(F)$ (inverse $\varpi^c(1-e)+e$). Note $sd=s$ and $dx_0=x_0$.

\emph{Claim: for $k\in\auGL_N(O)$ and all $y$, $\Psi(k^{-1}x_0)\,p(dky)=|p(dky)|\,\Psi(y)$.} If
$p(dky)=0$ both sides vanish. Otherwise $z=dky$ is integral and $g=sz=sky\in\auGL_b(O)$, and $p(z)=
\chi_S(\det g)^{-1}$. Since $z-x_0g=(d-e)ky=\varpi^{-c}(1-e)ky$, we have
$y=k^{-1}\bigl(\varpi^c(z-x_0g)+x_0g\bigr)$ with $t=k^{-1}(z-x_0g)$ integral, so
$\Psi(y)=\Psi(k^{-1}x_0g)=\chi_S(\det g)^{-1}\Psi(k^{-1}x_0)$ by Lemma~\ref{pcx:rightinv}. Thus
$\Psi(k^{-1}x_0)p(z)=\Psi(y)$, and $|p(z)|=1$.

By stability, $y\mapsto\Psi(k^{-1}x_0)p(dky)$ lies in $\mathcal Q$, hence so does $y\mapsto|p(dky)|
\Psi(y)$, for every $k\in\auGL_N(O)$; this function is locally constant in $k$ (uniformly in $y$), so
its average $I(y)\Psi(y)$, $I(y)=\int_{\auGL_N(O)}|p(dky)|\,dk$, is a finite combination and lies in
$\mathcal Q$. If $\Psi(y)\ne0$ then $y=l_{[b]}$ with $l\in\auGL_N(O)$, and $x_0=j_{[b]}$ with
$j\in\auGL_N(O)$ ($x_0$ is primitive since $sx_0=1$), so $y=lj^{-1}x_0$ and $I(y)=I(x_0)$ by right
invariance of $dk$. Finally $p(dx_0)=p(x_0)=\chi_S(\det sx_0)^{-1}=1$, so $I(x_0)>0$ by continuity.
Hence $I(x_0)\Psi\in\mathcal Q$ and $\Psi\in\mathcal Q$.
\end{proof}

\begin{lemma}[Action at $S$]\label{pcx:Sact}
Let $W$ (rank $N$, even), $V$ (rank $r$), $U$ (rank $b=N-r$) and $\eta(t)=\upsilon_w(t)|t|^{1/2}$ on
$F^\times$, where at $w\in S$ the local component $\upsilon_w$ of $\upsilon$ is evaluated on
$t_w\in E_w^\times$. For $g\in\auGL_N(F)$, $h\in\auGL_r(F)$, read via the frames as elements of
$\auU(W)(\auA)$, $\auU(V)(\auA)$ supported on $S\cup\bar S$, and for $\Phi\in\mathrm{SB}(\mathrm{Surj}^r)$,
$\Psi\in\mathrm{SB}(\mathrm{Inj}_b)$:
\begin{enumerate}
\item[(a)] $\Theta^{\mathrm L}_{W,V}(v,\xi(\det h)\Phi(h\,\cdot))=\Theta^{\mathrm L}_{W,V}(v,\Phi)$;
\item[(b)] $\Theta^{\mathrm L}_{W,V}(v,\xi(\det g)^{-1}\Phi(\cdot\,g^{-1}))(x g)=\eta(\det g)^{-b}
\Theta^{\mathrm L}_{W,V}(v,\Phi)(x)$;
\item[(c)] $\Theta^{\mathrm S}_{W,U}(y,\Psi(g\,\cdot))(xg)=\eta(\det g)^{-b}\Theta^{\mathrm S}_{W,U}(y,
\Psi)(x)$.
\end{enumerate}
\end{lemma}
\begin{proof}
By Corollary~\ref{pcx:conv} all integrals converge absolutely. All scalars are computed from one
formula. Let $v=w\aubar w$ with $w\in S$. By {\cite[Proposition~2.17]{P49}}, in the model at $w$ of
$\auHom(S_1,S_2)$ (ranks $n_1,n_2$; source group element $g$, target group element $h$) the untwisted
operator is
\[
\Psi\mapsto\upsilon_w(\det h_w)^{-n_1}\upsilon_w(\det g_w)^{n_2}|\det h_w|_w^{-n_1/2}
|\det g_w|_w^{n_2/2}\,\Psi(h_w^{-1}xg_w).
\]
At $v$, $\auU(1)(E^+_v)=\{(t,t^{-1})\}\cong E_w^\times$ through the $w$-component, and, as in the
proof of {\cite[Proposition~2.17]{P49}}, one has
$\upsilon(t,t^{-1})=\upsilon_w(t)^2$ and $\nu(t,t^{-1})=\chi_S(t)$; in particular, for $z\in
E_w^\times$, $\upsilon_w(z)^{n}=\upsilon(z)^{n/2}$ for even $n$, where $\upsilon(z)$ means
$\upsilon|_{\auU(1)}$ at $(z,z^{-1})$. Thus $\gamma(z)=\chi_S(z)^{-1}\upsilon_w(z)^{N}$, and the
untwisted operator equals $\upsilon(\det h)^{-n_1/2}\upsilon(\det g)^{n_2/2}|\det h|^{-n_1/2}
|\det g|^{n_2/2}\Psi(h^{-1}xg)$ when $n_1,n_2$ are even, and in general with $\upsilon_w^{-n_1}$,
$\upsilon_w^{n_2}$ in place of $\upsilon^{-n_1/2}$, $\upsilon^{n_2/2}$.

(a) For $\auHom(W,V)$, $n_1=N$, $n_2=r$. The operator $\omega_{\gamma,\gamma}(1,h^{-1})$ is
$\gamma(\det h)^{-1}$ times the untwisted operator for $h^{-1}$, i.e.\ it sends $\Phi$ to
\[
\chi_S(\det h)\upsilon_w(\det h)^{-N}\cdot\upsilon_w(\det h)^{N}|\det h|^{N/2}\,\Phi(h\,\cdot)
=\xi(\det h)\Phi(h\,\cdot):
\]
the power $\upsilon_w^{N}$ of the model cancels the one carried by $\gamma=\nu^{-1}\upsilon^{N/2}$.
So $\omega_{\gamma,\gamma}(1,h^{-1})[v,\Phi]_S=[v,\xi(\det h)\Phi(h\,\cdot)]_S$, using
\eqref{pcx:eq-twist} for the twist and the invariance of the infinite and unramified components
under the element $h$, which is supported on $S\cup\bar S$. Since
$\int_{[\auU(V)]}\theta(g,h h_0)\,dh=\int_{[\auU(V)]}\theta(g,h)\,dh$, (a) follows.

(b) $\Theta^{\mathrm L}(v,\Phi)(xg)$ is $\Theta^{\mathrm L}$ of the datum translated by
$\omega_{\gamma,\gamma}(g,1)$, evaluated at $x$, and $\omega_{\gamma,\gamma}(g,1)$ sends $\Phi$ to
\[
\gamma(\det g)^{-1}\upsilon_w(\det g)^{r}|\det g|^{r/2}\Phi(\cdot\,g)
=\chi_S(\det g)\upsilon_w(\det g)^{r-N}|\det g|^{r/2}\Phi(\cdot\,g).
\]
Applied to $\xi(\det g)^{-1}\Phi(\cdot\,g^{-1})$, with $\xi(\det g)^{-1}=\chi_S(\det g)^{-1}
|\det g|^{-N/2}$, this gives $\upsilon_w(\det g)^{-b}|\det g|^{-b/2}\Phi=\eta(\det g)^{-b}\Phi$,
since $r-N=-b$.

(c) For $\Theta^{\mathrm S}_{W,U}$ the characters are $(\upsilon^{b},\nu\upsilon^{N/2})$ and the
datum $\Psi^{\mathsf t}$ sits at $\bar w$, in the model at $\bar w$ of $\auHom(W,U)$ ($n_1=N$,
$n_2=b$). Apply the formula above with the roles of $w,\bar w$ exchanged. The component of $g$ at
$\bar w$ is $g_{\bar w}=\bar g^{-\mathsf t}$, so, in the identification of $E_{\bar w}$ with $E_w$ by
conjugation, $\det g_{\bar w}=(\det g)^{-1}$, and the element $\det g$ of $\auU(1)(E^+_v)$ is
$((\det g),(\det g)^{-1})$ with $\bar w$-component $(\det g)^{-1}$. The untwisted operator therefore
multiplies by $\upsilon_{\bar w}(\det g_{\bar w})^{b}|\det g_{\bar w}|^{b/2}=\upsilon_{\bar w}
((\det g)^{-1})^{b}|\det g|^{-b/2}=\upsilon_w(\det g)^{b}|\det g|^{-b/2}$, because
$\upsilon_{\bar w}=\upsilon_w^{-1}$ on $E^{+\times}_v$ (this holds for $b$ of either parity; reading
$(t,t^{-1})$ through its $\bar w$-component $t^{-1}$ gives the same value $\upsilon_{\bar
w}(t^{-1})^2=\upsilon_w(t)^2$ of $\upsilon|_{\auU(1)}$), and acts on $\Psi^{\mathsf t}$
through $X\mapsto X g_{\bar w}$, i.e.\ on $\Psi$ through $y\mapsto g^{-1}y$ up to the
identification by conjugation. The source character contributes $\upsilon^{b}(\det g)^{-1}=
\upsilon_w(\det g)^{-2b}$. Hence $\Theta^{\mathrm S}(y,\Psi)(xg)=\upsilon_w(\det g)^{-b}
|\det g|^{-b/2}\,\Theta^{\mathrm S}(y,\Psi(g^{-1}\,\cdot))(x)$; replacing $\Psi$ by $\Psi(g\,\cdot)$
gives (c), with the scalar $\eta(\det g)^{-b}$, the same as in (b).
\end{proof}

\begin{proof}[Proof of Proposition~\ref{pcx:P-match}]
Let $\Lambda(\Phi)=\Theta^{\mathrm L}_{W,V}(v,\Phi)$ on $\mathrm{SB}(\mathrm{Surj}^r)$ (convergent by
Corollary~\ref{pcx:conv}); by Lemma~\ref{pcx:Sact}(a) it satisfies the hypothesis of
Lemma~\ref{pcx:linfactor}, which gives $\mathcal T$. Put $\mathcal D(\Psi)=\mathcal T(\Psi)-c\,
\Theta^{\mathrm S}_{W,U}(y,\Psi)$ and $\mathcal Q=\{\Psi\in\mathcal M:\mathcal D(\Psi)=0\}$.

$\mathcal Q$ is stable under $\Psi\mapsto\Psi(g\,\cdot)$: if $(\Phi,\Psi)\in\mathcal P$ then
$(\Phi',\Psi(g\,\cdot))\in\mathcal P$ with $\Phi'=\xi(\det g)^{-1}\Phi(\cdot\,g^{-1})$
(Lemma~\ref{pcx:section-laws}(a) for $g^{-1}$), and by Lemma~\ref{pcx:Sact}(b),(c),
$\mathcal D(\Psi(g\,\cdot))(xg)=\eta(\det g)^{-b}\mathcal D(\Psi)(x)$ (the two parts of
$\mathcal D$ are multiplied by the same scalar: in (b) the twist $\gamma(\det g)^{-1}=
\chi_S(\det g)\upsilon_w(\det g)^{-N}$ and the model factor $\upsilon_w(\det g)^{r}|\det g|^{r/2}$
combine with $\xi(\det g)^{-1}$ to $\upsilon_w(\det g)^{-b}|\det g|^{-b/2}$, and in (c) the twist
$\upsilon^b(\det g)^{-1}=\upsilon_w(\det g)^{-2b}$ and the model factor
$\upsilon_w(\det g)^{b}|\det g|^{-b/2}$ at $\bar w$ give the same product), which vanishes if
$\mathcal D(\Psi)=0$.

Take $s$ the first $b$ rows of $1_N$ and $x_0$ the first $b$ columns. By Lemma~\ref{pcx:std-matched}
and Lemma~\ref{pcx:section-laws}(b), $(\Phi^{\mathrm{row}}_s,\Psi^{\mathrm{col}}_s)\in\mathcal P$, so
$\mathcal D(\Psi^{\mathrm{col}}_s)=\Theta^{\mathrm L}(v,\Phi^{\mathrm{row}}_s)-c\Theta^{\mathrm S}(y,
\Psi^{\mathrm{col}}_s)=0$ by hypothesis, i.e.\ $\Psi^{\mathrm{col}}_s\in\mathcal Q$. By
Lemma~\ref{pcx:generate}, every locally constant $\Psi$ with a matched $\Phi$ lies in $\mathcal Q$;
then $(\Phi,\Psi)\in\mathcal P$ (Lemma~\ref{pcx:section-laws}(b)), $\Psi\in\mathrm{SB}(\mathrm{Inj}_b)$,
and $\Theta^{\mathrm L}(v,\Phi)=\mathcal T(\Psi)=c\Theta^{\mathrm S}(y,\Psi)$.
\end{proof}

\subsubsection{Passing to adjoints: proof of Proposition~\ref{pcx:P-flip}}\label{pcx:ss-flip}

Let $U^-$ be $U$ with the negated form; $\auU(U^-)=\auU(U)$ as groups, with the same rational points,
the same lattice $L_U$ and hence the same Haar measure, so $[\auU(U^-)]=[\auU(U)]$ as measure spaces,
and $\det$ is the same. {\cite[Proposition~7.17]{P49}} (fact (W6)) gives, for the pair $(U,W)$ of ranks
$b,N$, an isomorphism $\mathrm{Fl}\colon\psHS(\auHom(U,W))\to\psHS(\auHom(W,U^-))$ induced by $w\mapsto
w^\dagger$, a constant $c_\theta\ne0$ with $\theta'\circ\mathrm{Fl}=c_\theta\theta$, and
\[
\mathrm{Fl}\circ\omega_{\auHom(U,W)}(h,g)=\upsilon(\det h)^{N}\upsilon(\det g)^{-b}\,
\omega_{\auHom(W,U^-)}(g,h)\circ\mathrm{Fl}\qquad(h\in\auU(U)(\auA),\ g\in\auU(W)(\auA)),
\]
(the factor is $\upsilon(\det\nolimits_{\auHom(U,W)}(h,g))^{-1}$, the determinant of $(h,g)$ on the
$E$-module $\auHom(U,W)$ being $\det(g)^b\det(h)^{-N}$). We
normalise $\mathrm{Fl}$ as a tensor product of local isomorphisms: at the infinite places the map of
Fock spaces induced by $w\mapsto w^\dagger$ on real points, which defines $y\mapsto y^{\mathrm{ad}}$ and
preserves the moments $\|\cdot\|_M$ (it is induced by a measure-preserving isometry of the real
vector spaces); at the finite places the map induced by $w\mapsto w^\dagger$, normalised to be unitary.
By Lemma~\ref{pgx:dict}(f) the latter carries the self-dual lattice $\alpha\auHom(L_U,L_W)$ to
$u\,\alpha\auHom(L_W,L_U)$, $u=\aubar a/a$, a central idele of $\auU(U)=\auU(U^-)$ with $u\aubar u=1$ and
$u_v=1$ at the places not above $q_1$. Put $\mathrm{Fl}^\circ=\omega(1,u)^{-1}\circ\mathrm{Fl}$ ($u$
acting on $\auHom(W,U^-)$ through $\auU(U^-)$). Then: at finite places outside $S\cup\bar S$,
$\mathrm{Fl}^\circ$ sends the vacuum of the standard lattice $\alpha\auHom(L_U,L_W)$ to a multiple of
modulus one of the vacuum of $\alpha\auHom(L_W,L_U)$ (Lemma~\ref{pgx:dict}(e)); at $w\in S$, in the
models (W5) and the coordinates $\tilde x$ of Lemma~\ref{pgx:dict}, $\mathrm{Fl}^\circ$ is, up to a
constant of modulus one, the map $\Psi\mapsto(X\mapsto\Psi(\epsilon X^{\mathsf t}))$ from functions on
$\pcMat_{N\times b}(E_w)$ to functions on $\pcMat_{b\times N}(E_{\bar w})$ (identified by conjugation),
where $\epsilon=-1$: in the frames $w^\dagger=-\bar w^{\mathsf t}$ (the Gram matrices being unimodular
at $S$), the sign coming from the negated form; indeed by Lemma~\ref{pgx:dict}(f)
$\mathrm{Fl}$ is $\tilde X=-\nu^{-1}\,{}^t\tilde x$ with $\nu=\aubar{u_{\bar w}}^{-1}$, and
$\omega(1,u)^{-1}$ multiplies the variable $\tilde X$ by $\nu^{-1}$ ((W5): $h=u^{-1}$ acts by
$\Psi\mapsto\Psi(u_{\bar w}\,\cdot)$). The tensor product $\mathrm{Fl}$ of the local maps intertwines
the Heisenberg operators, so it is an admissible $\mathrm{Fl}$, and $c_\theta$ is a constant depending
only on $U$, $W$ and the frames. Since $u$ is central, Lemma~\ref{pgx:dict}(e) gives, for the theta
integrals over $[\auU(U)]$ used below,
$\int\theta'(\omega'(g,h)\mathrm{Fl}f)\,dh=\zeta_u\int\theta'(\omega'(g,h)\mathrm{Fl}^\circ f)\,dh$ with
$\zeta_u$ a value at $\det u$ of the target character, of modulus one; so in these integrals
$\mathrm{Fl}$ may be replaced by $\mathrm{Fl}^\circ$ at the cost of a constant of modulus one.

\begin{proof}[Proof of Proposition~\ref{pcx:P-flip}]
\emph{The datum.} By the normalisation, $\mathrm{Fl}^\circ[y,\Psi]_S=a_1[y^{\mathrm{ad}},\Psi_\epsilon^{\mathsf
t}]_{\bar S}$ with $\Psi_\epsilon(X)=\Psi(\epsilon X)$ and $|a_1|=1$. For $\Psi$ matched with some $\Phi$,
Lemma~\ref{pcx:rightinv} with $m=\epsilon 1_b$ gives $\Psi(\epsilon X)=\chi_S(\epsilon)^{-b}\Psi(X)$.
So $\mathrm{Fl}^\circ[y,\Psi]_S=a\,[y^{\mathrm{ad}},\Psi^{\mathsf t}]_{\bar S}$ with $a=a_1\chi_S(\epsilon)^{-b}$,
$|a|=1$.

\emph{The kernels.} By definition and the displayed law,
\[
\theta^{U,W}_{\gamma,\gamma}(h,g;f)=\gamma(\det h)^{-1}\gamma(\det g)\,\theta(\omega(h,g)f)
=c_\theta^{-1}\gamma(\det h)^{-1}\upsilon(\det h)^{N}\gamma(\det g)\upsilon(\det g)^{-b}\,
\theta'(\omega'(g,h)\mathrm{Fl}f).
\]
In the convention of \S\ref{pcx:ss-statement} for the pair $(W,U^-)$ this is
$c_\theta^{-1}\theta^{W,U^-}_{\gamma^{-1}\upsilon^b,\ \gamma^{-1}\upsilon^N}(g,h;\mathrm{Fl}f)$, and
$\gamma^{-1}\upsilon^N=\nu\upsilon^{N/2}$; by \eqref{pcx:eq-twist},
$\theta^{W,U^-}_{\gamma^{-1}\upsilon^b,\nu\upsilon^{N/2}}=\gamma(\det g)\,\theta^{W,U^-}_{\upsilon^b,
\nu\upsilon^{N/2}}$.

\emph{Integration.} With $f=[y,\Psi]_S$, integrating over $h\in[\auU(U)]=[\auU(U^-)]$ (absolutely
convergent by Corollary~\ref{pcx:conv}, $\Psi$ being supported on injective matrices), and
replacing $\mathrm{Fl}$ by $\mathrm{Fl}^\circ$ in the integral at the cost of the constant $\zeta_u$ of
modulus one (end of \S\ref{pcx:ss-flip}; absorbed in $a$),
\[
\int_{[\auU(U)]}\theta^{U,W}_{\gamma,\gamma}(h,g;[y,\Psi]_S)\,dh=c_\theta^{-1}a\,\gamma(\det g)\,
\Theta^{\mathrm S}_{W,U^-}(y^{\mathrm{ad}},\Psi)(g).
\]
So the proposition holds with $d=c_\theta a^{-1}$, $|d|=|c_\theta|=:d_0$, independent of $\nu$.
\end{proof}

\begin{remark}
One can normalise so that $|d|=1$; only the independence of $|d|$ from $\nu$ is needed, since constants independent of $\nu$ are absorbed in $C$.
\end{remark}

\subsubsection{Only the signs matter: proof of Proposition~\ref{pcx:P-signs}}\label{pcx:ss-signs}

\begin{lemma}[Change of space]\label{pcx:move}
Let $W,W_1$ have the same rank and signs, and $V,V_1$ likewise. There are $k\in\auU(W)(\auA)$ and,
for all characters $\alpha,\beta$, a number $u$ of modulus $1$ (independent of the data and of $g$)
such that for all data $[v,\Phi]_P$ and all $g$,
$\int_{[\auU(V_1)]}\theta^{W_1,V_1}_{\alpha,\beta}(g,h;[v,\Phi]_P)\,dh=u\int_{[\auU(V)]}
\theta^{W,V}_{\alpha,\beta}(gk,h;[v,\Phi]_P)\,dh$ (absolutely convergent integrals), where on the
right $g$ stands for its transport $\phi g\phi^{-1}$ by the isometry $\phi$ of the proof.
\end{lemma}
\begin{proof}
Apply Lemma~\ref{pcx:landherr} to $(W_1,W)$ and $(V_1,V)$, giving $\phi,k$ and $\phi',k'$. A datum is a
vector of $\psHS(\auHom)$ written in the frames, and the theta functional of $\auHom(W_1,V_1)$ is
that of $\auHom(W,V)$ transported by $(\phi,\phi')$ (it only depends on the rational points).
Transporting, the frame-datum $[v,\Phi]_P$ for $(W_1,V_1)$ becomes $\omega(k,k')$-applied to the same
frame-datum for $(W,V)$, up to the characters; hence
$\theta^{W_1,V_1}_{\alpha,\beta}(g,h;\cdot)=\alpha(\det k)\beta(\det k')^{-1}\theta^{W,V}_{\alpha,\beta}
(gk\,,hk';\cdot)$, with $g,h$ transported by $\phi,\phi'$ (indeed
$\theta(\omega(g,h)\omega(k,k')D)=\theta(\omega(gk,hk')D)$). Finally $h\mapsto hk'$ carries the quotient
by $\auU(V_1)(E^+)=\phi'^{-1}\auU(V)(E^+)\phi'$ measure-preservingly onto that by $\auU(V)(E^+)$
(conjugation by the transport, right translation by $k'$, both preserving Haar measure). So the
lemma holds with $u=\alpha(\det k)\beta(\det k')^{-1}$.
\end{proof}

\begin{lemma}[Equal rank: the theta integral is constant]\label{pcx:const}
Let $Y$ have even rank $2k$ and contain an adapted hyperbolic plane. For every Fock vector $v$ of
$(Y,Y)$ with finite moments and every tame $\nu$, the function
$g\mapsto\Theta^{\mathrm L}_{Y,Y}(v,\Psi^{\mathrm{col}}_1)(g)$ is constant on $\auU(Y)(\auA)$.
\end{lemma}
\begin{proof}
See \S\ref{pct:ss-const}: the hypotheses of Lemma~\ref{pct:SA} are verified there for
$D=\Theta^{\mathrm L}_{Y,Y}(v,\Psi^{\mathrm{col}}_1)$, the torus condition coming from
Theorem~\ref{pct:plane} (independence of $u$ in equal rank).
\end{proof}

\begin{proof}[Proof of Proposition~\ref{pcx:P-signs}]
By Lemma~\ref{pcx:move} with $(W,V)=(Y,Y)$, $(W_1,V_1)=(Y_1,Y_1)$, $\alpha=\beta=\gamma$, $P=S$,
$\Phi=\Psi^{\mathrm{col}}_1$: $\Theta^{\mathrm L}_{Y_1,Y_1}(v,\Phi)(1)=u\,\Theta^{\mathrm L}_{Y,Y}(v,\Phi)(k)
=u\,\Theta^{\mathrm L}_{Y,Y}(v,\Phi)(1)$ by Lemma~\ref{pcx:const}. Taking absolute values and dividing
by $L_*$ gives $\pcxstr_{Y_1}(\nu;v)=\pcxstr_Y(\nu;v)$.
\end{proof}

\subsubsection{The inputs of the exchange identity}\label{pcx:ss-status}

\emph{(A) Proved here} (modulo the published facts cited: Landherr's theorem, the Hasse norm
theorem, Jacobowitz, reduction theory (R1), Godement's form of Borel's reduction theorem and Dirichlet's unit theorem in
Lemma~\ref{pcx:reduce}, Hecke's
non-vanishing; (R2) is Lemma~\ref{pcx:thetagrowth}): the dictionary \eqref{pcx:eq-dict} and the twist law \eqref{pcx:eq-twist};
Lemma~\ref{pcx:Lpos}, Corollary~\ref{pcx:cne0}; Lemmas~\ref{pcx:adapted}, \ref{pcx:landherr},
\ref{pcx:std-matched}, \ref{pcx:mod}, \ref{pcx:reduce}, \ref{pcx:thetagrowth}, \ref{pcx:superpos}, \ref{pcx:diffgain}, \ref{pcx:mean},
\ref{pcx:meanzero}, \ref{pcx:cosets}, \ref{pcx:rightinv}, \ref{pcx:section-laws},
\ref{pcx:linfactor}, \ref{pcx:generate}, \ref{pcx:Sact}, \ref{pcx:move};
Proposition~\ref{pcx:decay} and Corollary~\ref{pcx:conv} (absolute convergence);
Propositions~\ref{pcx:P-match} and~\ref{pcx:P-flip} (the latter with {\cite[Proposition~7.17]{P49}});
the deductions
$\ref{pcx:P-std}+\ref{pcx:P-match}+\ref{pcx:P-flip}\Rightarrow\ref{pcx:P-matches}$,
$\ref{pcx:P-matches}+\ref{pcx:P-equal}\Rightarrow\ref{pcx:X}$,
$\ref{pcx:P-matches}+\ref{pcx:P-nonvan}+\ref{pcx:P-plane}+\ref{pcx:P-signs}\Rightarrow
\ref{pcx:P-equal}$ (Lemma~\ref{pcx:twice}); Proposition~\ref{pcx:P-signs}.

\emph{(B) Proved in the following sections:} Proposition~\ref{pcx:P-slice} in \S\ref{pcu:ss-slice};
Proposition~\ref{pcx:P-top} in \S\S\ref{pcu:ss-witt}--\ref{pcu:ss-std-data} (with (U1), (U2) in
\S\ref{pcu:ss-open}); the deduction of Proposition~\ref{pcx:P-std} in \S\ref{pcu:ss-std};
Lemma~\ref{pcx:const} and Propositions~\ref{pcx:P-descent}, \ref{pcx:P-plane}, \ref{pcx:P-nonvan} in
\S\ref{pct:sec}, the last with Lemma~\ref{pct:negcoeff}, proved in \S\ref{pcn:sec}. The full list of
dependencies of Theorem~\ref{pcx:X} is \S\ref{pcu:ss-final}.



% ===== C-exchange-top.tex =====
% proofs/C-exchange-top.tex
% Continuation of C-exchange.tex: real vectors with prescribed zero term (Proposition pcx:P-slice),
% the large triple (Proposition pcx:P-top), the deduction of pcx:P-std, and the list of what
% Theorem pcx:X depends on. Labels of this file start with pcu:. Macros: those of C-exchange.tex.
\providecommand{\pcuS}{\mathsf S}
\providecommand{\pcuF}{\mathcal F}

\subsection{The large triple and the assembly of the exchange identity}\label{pcu:sec}

This section continues \S\ref{pcx:sec}; notation is that of \S\ref{pcx:ss-arch}. It proves
Proposition~\ref{pcx:P-slice}, proves Proposition~\ref{pcx:P-top} (the unfolding statements (U1),
(U2) being proved in \S\ref{pcu:ss-open}), deduces Proposition~\ref{pcx:P-std}, and lists in \S\ref{pcu:ss-final} every
statement on which Theorem~\ref{pcx:X} depends.

\subsubsection{Splittings and the mixed model at the real places}\label{pcu:ss-mixed}

Let $X$ be a space and $X=(\ell\oplus\ell')\perp X_1$ an adapted splitting
(Lemma~\ref{pcx:adapted}): $\ell=\langle e_1,\dots,e_p\rangle$, $\ell'=\langle f_1,\dots,f_p\rangle$
are isotropic, $\langle e_i,f_j\rangle=\delta_{ij}$, $X_1=(\ell\oplus\ell')^\perp$, and the frame
$(e,f,\text{frame of }X_1)$ is orthonormal (after the change $e\pm f$) at the infinite places and
integral with unimodular sum at the finite places. We write $D$ for the splitting and $D^{\mathrm
s}$ for the one with $e$ and $f$ exchanged. For a space $T$ of rank $a$ the real points of
$\auHom(X,T)$ decompose as $\auHom(\ell',T)_\R\oplus\auHom(\ell,T)_\R\oplus\auHom(X_1,T)_\R$, and
$\auHom(\ell',T)_\R\oplus\auHom(\ell,T)_\R$ is a polarised space whose Lagrangians are the two
summands. Identify $\auHom(\ell,T)_\R$ with $\prod_\tau\pcMat_{a\times p}(\C)$ through the frames.

\emph{The mixed model.} The Fock space of $\auHom(X,T)_\R$ is isomorphic, compatibly with the
Heisenberg group, to the space of functions $x\mapsto\Lambda(x)$ on $\prod_\tau\pcMat_{a\times
p}(\C)$ with values in the Fock space of $\auHom(X_1,T)_\R$ (Schr\"odinger model in the
$\auHom(\ell,T)$ variable, Fock model in the rest; this is the real form of the models (W2), (W2$'$)
of {\cite[\S5.1.1]{P49}}, obtained by applying the Bargmann transform to the $\auHom(\ell\oplus\ell',T)$
factor of the Fock space). The value $\Lambda(x)$ is the \emph{slice} of
$\Lambda$ at $x$ along $D$, written $\mathrm{sl}_D(x)\Lambda$; for an adelic matrix $x$,
$\mathrm{sl}_D(x)$ means the slice at the archimedean components of $x$. The \emph{zero term}
of $\Lambda$ along $D$ is $\mathrm{sl}_D(0)\Lambda$; it is the restriction of the real
Schr\"odinger model to the matrices vanishing on the isotropic part, as in
Proposition~\ref{pcx:P-slice}. We use four standard properties.
\begin{enumerate}
\item[(M1)] \emph{Prescribed slices} (proved in \S\ref{pcu:ss-M1}). If $\tau$ is a smooth compactly supported function on
$\prod_\tau\pcMat_{a\times p}(\C)$, $u$ a smooth map from there to the real unitary group
$\prod_\tau\auU(X_1\otimes_\tau\C)\times\prod_\tau\auU(T\otimes_\tau\C)$, and $v$ a Fock vector of
$\auHom(X_1,T)_\R$ with finite moments, there is a Fock vector $\Lambda$ of $\auHom(X,T)_\R$ with
finite moments and $\mathrm{sl}_D(x)\Lambda=\tau(x)\,\omega_\infty(u(x))v$ for all $x$. (The unitary
groups here need not be compact; the proof in \S\ref{pcu:ss-M1} uses only that $\omega_\infty(g)$ carries the
vacuum to a Gaussian vector depending smoothly on $g$.)
\item[(M2)] \emph{Change of polarisation.} There is a constant $c_D>0$, depending only on the
normalisation of measures, such that for every $\Lambda$ with finite moments
$\mathrm{sl}_{D^{\mathrm s}}(0)\Lambda=c_D\int\mathrm{sl}_D(x)\Lambda\,dx$. \emph{Proof:} the
intertwiner between the Schr\"odinger models in $\auHom(\ell,T)$ and in $\auHom(\ell',T)$ is the
partial Fourier transform $\phi\mapsto c_D\int\phi(x)\psi_\infty(\langle x,x'\rangle)dx$ (Weil's
intertwining operator for two transversal Lagrangians, unique up to a scalar by Stone--von~Neumann);
at $x'=0$ it is $c_D\int\phi$, the integral converging absolutely by Lemma~\ref{pcu:slicedecay}:
$\int\|\mathrm{sl}_D(x)\Lambda\|_0dx\le C\|\Lambda\|_{2p+1}$.
\item[(M3)] \emph{Equivariance.} The real unitary group of $T$ acts on slices by
$\mathrm{sl}_D(x)\bigl(\omega_\infty(1,h)\Lambda\bigr)=\omega_\infty(1,h)\,\mathrm{sl}_D(h^{-1}x)\Lambda$,
and $\auU(X_1)_\R\subset\auU(X)_\R$ acts on the slices pointwise; with the characters of the
representation $\omega_{\alpha,\beta}$ (\S\ref{pcx:ss-statement}) the factor $\alpha^{-1}\beta$ of
the determinants is carried along.
\item[(M4)] \emph{Stages.} If $X_1=(\ell_1\oplus\ell'_1)\perp X_2$ is adapted, then $X=((\ell\oplus\ell_1)
\oplus(\ell'\oplus\ell'_1))\perp X_2$ is adapted, and the zero term along it is the zero term along
$X_1=(\ell_1\oplus\ell_1')\perp X_2$ of the zero term along $D$.
\end{enumerate}
(M3) and (M4) are immediate from the definition of the mixed model as a tensor product.
(M2)--(M4) are proved as indicated, (M1) in \S\ref{pcu:ss-M1}.

\begin{lemma}[The mixed model is unitary up to a constant]\label{pcu:mixedmodel}
Let $V_\ell=\auHom(\ell,T)_\R$ with a Haar measure $dx$, $\mathcal F'$, $\mathcal F$ the Fock spaces of
$\auHom(X,T)_\R$ and $\auHom(X_1,T)_\R$, and $\mathsf H=L^2(V_\ell,dx;\mathcal F)$. There is a bounded operator
$\mathsf M\colon\mathcal F'\to\mathsf H$ and a constant $c_{\mathsf M}>0$, depending only on $dx$, such that
\textup{(a)} $\|\mathsf Mv\|_{\mathsf H}^2=c_{\mathsf M}\|v\|^2$ for all $v\in\mathcal F'$; \textup{(b)} $\mathsf M$
intertwines the Heisenberg operators of $\mathcal F'$ with the operators of the mixed model (translation of the
variable by $\auHom(\ell,T)_\R$, multiplication by characters for $\auHom(\ell',T)_\R$, action on the values for
$\auHom(X_1,T)_\R$), in the form \eqref{pcu:eq-mixed}; \textup{(c)} for every $v$ with finite moments,
$(\mathsf Mv)(x)=\mathrm{sl}_D(x)v$ for all $x$. The family $\mathrm{sl}_D$ is determined on vectors with finite moments by
\eqref{pcu:eq-mixed}.
\end{lemma}
\begin{proof}
\emph{Uniqueness.} Suppose two families of bounded maps $\mathrm{sl}(x)$, $\mathrm{sl}'(x)$ satisfy
\eqref{pcu:eq-mixed}. Both send the vacuum $e$ to $e^{-|x|^2}e$, and \eqref{pcu:eq-mixed} expresses
$\mathrm{sl}(x)\psrho(y)e$ through $\mathrm{sl}(\cdot)e$ alone; so $\mathrm{sl}(x)\psrho(y)e=\mathrm{sl}'(x)\psrho(y)e$
for all $x,y$. A vector with finite moments is the absolutely convergent superposition
$v=\int v(y)\psrho(y)e\,dy$ (Lemma~\ref{pcx:superpos}), and $\mathrm{sl}(x)$, $\mathrm{sl}'(x)$ are bounded, so
$\mathrm{sl}(x)v=\mathrm{sl}'(x)v$.

\emph{Existence and unitarity.} Write $A=\auHom(\ell\oplus\ell',T)_\R$ and $B=\auHom(X_1,T)_\R$, so that
$\auHom(X,T)_\R=A\oplus B$ orthogonally, and use the Fock model {\cite[(41)]{P49}}: entire functions $F$ with
$\int|F|^2e^{-|z|^2}<\infty$, in coordinates for which the frames are orthonormal. (i) \emph{Tensor factorisation.}
In these coordinates $z=(z_A,z_B)$, $|z|^2=|z_A|^2+|z_B|^2$, and the Gaussian measure is a product; the monomials
$z_A^\alpha z_B^{\alpha'}$ are orthogonal with $\|z_A^\alpha z_B^{\alpha'}\|^2=\|z_A^\alpha\|^2\|z_B^{\alpha'}\|^2$, and
the polynomials are dense. Hence $F_A\otimes F_B\mapsto F_A(z_A)F_B(z_B)$ extends to a unitary isomorphism
$\mathcal F(A)\widehat\otimes\mathcal F(B)\to\mathcal F'$ (Hilbert tensor product). By {\cite[(41)]{P49}},
$e^{z\cdot y-|y|^2/2}F(z-\bar y)$ factorises for $y=y_A+y_B$, so the isomorphism carries
$\psrho_A(y_A)\otimes\psrho_B(y_B)$ to $\psrho(y_A+y_B)$, and the vacuum $1\otimes1$ to the vacuum.
(ii) \emph{Bargmann transform.} $A=\auHom(\ell,T)_\R\oplus\auHom(\ell',T)_\R$ is a sum of two transversal
Lagrangians for the alternating form of $\beta$, and the frame of $D$ at the infinite places is
orthonormal after the change $e\pm f$ (\S\ref{pcu:ss-mixed}), so the complex structure of the Fock model of $A$
exchanges the two summands. The Bargmann transform is a unitary operator, up to a positive constant fixed by
the normalisation of $dx$, from the Schr\"odinger model $L^2(V_\ell,dx)$ of $A$ onto $\mathcal F(A)$, which
intertwines the Schr\"odinger operators (translations by $\auHom(\ell,T)_\R$, multiplications by the characters
given by $\auHom(\ell',T)_\R$) with the Fock operators and carries a Gaussian $e^{-|x|^2}$, in suitable linear
coordinates $x$ on $V_\ell$, to a non-zero multiple of the vacuum (Folland~\cite[Chapter~1, \S6, pp.~40--44, (1.71)--(1.72)]{FollandPS}, where the Gaussian is $e^{-\boldsymbol\pi|x|^2}$). Let $\mathsf B\colon\mathcal F(A)\to L^2(V_\ell,dx)$ be its inverse, multiplied by the positive
constant for which $\mathsf B1=e^{-|x|^2}$. Then $\mathsf M=(\mathsf B\otimes1)\circ(\text{isomorphism of
(i)})^{-1}$ maps $\mathcal F'$ onto $\mathsf H=L^2(V_\ell)\widehat\otimes\mathcal F$, multiplies norms by a constant
$c_{\mathsf M}^{1/2}>0$, intertwines the Heisenberg operators as stated in (b) (this defines the linear maps $L$, $K$,
$I$ of \eqref{pcu:eq-mixed}: $L$ and $K$ are the inclusions of $\auHom(\ell,T)_\R$, $\auHom(\ell',T)_\R$ read in the
coordinates $x$, $I$ that of $B$), and $\mathsf Me=e^{-|x|^2}e$. This is (a) and (b).
(iii) \emph{Pointwise values.} For $v=\psrho(y)e$ the function $\mathsf Mv$ is, by (b), the continuous function
$x\mapsto(\text{phase})\,e^{-|x+Ay|^2}\psrho(Wy)e$ (notation of \S\ref{pcu:ss-M1}); so $x\mapsto(\mathsf M\psrho(y)e)(x)$
satisfies \eqref{pcu:eq-mixed} pointwise, and by the uniqueness above it equals $\mathrm{sl}_D(x)\psrho(y)e$. For $v$
with finite moments, $\mathsf Mv=\int v(y)\mathsf M\psrho(y)e\,dy$ in $\mathsf H$ ($\mathsf M$ bounded), and the
integral also converges pointwise in $\mathcal F$, uniformly in $x$ (the integrand has norm $\le|v(y)|$); the two
limits agree almost everywhere, and the pointwise one is $\mathrm{sl}_D(x)v$. This is (c).
\end{proof}

\subsubsection{Real vectors with prescribed zero term: proof of Proposition~\ref{pcx:P-slice}}
\label{pcu:ss-slice}

\begin{lemma}[Prescribed constant slices]\label{pcu:bump}
Let $D$ be an adapted splitting $X=(\ell\oplus\ell')\perp X_1$, $T$ a space, $\tau\ge0$ smooth with
compact support on $\prod_\tau\pcMat_{a\times p}(\C)$, $\tau\not\equiv0$, and $v$ a Fock vector of
$\auHom(X_1,T)_\R$ with finite moments. There are $\kappa\in\C$ and a Fock vector $\Lambda$ of
$\auHom(X,T)_\R$ with finite moments such that $\mathrm{sl}_D(x)\Lambda=\tau(x)\kappa\,v$ for all $x$ and
$\mathrm{sl}_{D^{\mathrm s}}(0)\Lambda=v$.
\end{lemma}
\begin{proof}
$\int\tau>0$. Put $\kappa=(c_D\int\tau)^{-1}$ and take $\Lambda$ from (M1) with the section
$\tau\kappa$ and $u\equiv1$. By (M2), $\mathrm{sl}_{D^{\mathrm s}}(0)\Lambda=c_D\kappa(\int\tau)v=v$.
\end{proof}

\begin{proof}[Proof of Proposition~\ref{pcx:P-slice}]
Apply Lemma~\ref{pcu:bump} to $D^{\mathrm s}$ (the splitting $W'=W\oplus\pcxH^p$ with the two
isotropic frames exchanged), $T=U$, any bump $\tau$ and $v=y$. The vector $w=\Lambda$ has finite
moments and its zero term along $(D^{\mathrm s})^{\mathrm s}=D$ is $y$.
\end{proof}

\begin{remark}
The zero term in Proposition~\ref{pcx:P-slice} is taken along the given splitting; by the lemma the
same holds along the splitting with the frames exchanged. Proposition~\ref{pcx:P-top} below produces
$x'$ whose zero term along $D^{\mathrm s}$ is $w$; Proposition~\ref{pcx:P-descent} must use the
same orientation (see \S\ref{pcu:ss-std}).
\end{remark}

\subsubsection{The large triple: setting}\label{pcu:ss-top-setting}

Fix the data of Proposition~\ref{pcx:P-top}: $W_0=\pcxH\oplus W$ of even rank $n=2j$ with the plane
split off adaptedly (splitting $D_1$), $U$ of rank $b$ with signs $t'$, $X=\pcxH^b\oplus W_0$ with
an adapted splitting $D\colon X=(\ell\oplus\ell')\perp W_0$, $\dim\ell=b$ (Lemma~\ref{pcx:adapted}
applied $b$ times, together with (M4)), and $T=U\oplus W_0$, of rank $b+n$ with signs $t'\oplus s$
($s$ those of $W_0$). Let $J=\mathrm{diag}(t'\oplus s)$ at each infinite place and
$x_0=\binom{1_b}{0}$, the frame of $U$ in $T$; its Gram matrix is $\mathrm{diag}(t')$, and its
stabiliser in $\auU(T)$ is $\auU(W_0)$. Write $\mathrm{Fr}=\prod_\tau\pcMat_{(b+n)\times b}(\C)$ for
families of frames, one for each infinite place. Note $\operatorname{rk}T-\operatorname{rk}X/2=
b+n-(2b+n)/2=j$, so in $\mathrm{RC}(X,T,U;\cdot)$ the factor is $L_*(\pcm,\nu,j)$, and
$\pcxstr_{W_0}$ also has $L_*(\pcm,\nu,j)$ in its denominator.

The parabolic $P_D$ of $\auU(X)$ stabilising $\ell$ consists of the products $n(c,z)\,m(q,g)$, with
$m(q,g)$ acting by $q\in\auGL_b$ on $\ell$ (in the frame $e$), by $q^{*-1}$ on $\ell'$ and by $g\in
\auU(W_0)$ on $W_0$, and $n(c,z)$ in the unipotent radical. Precisely, with $J_0$ the Gram matrix of the
frame of $W_0$, $c\in\pcMat_{n\times b}$ and $z\in\pcMat_b$ with
\[
z+z^*=-c^*J_0c ,
\]
$n(c,z)$ fixes $e$, maps $w\in W_0$ to $w+e\,c^*J_0w$ and $f$ to $f-c+e\,z$ (in the frames: the $k$-th vector
$f_k$ goes to $f_k-\sum_lc_{lk}w_l+\sum_ie_iz_{ik}$, $(w_l)$ the frame of $W_0$). The relation $z+z^*=-c^*J_0c$ is
exactly the condition that $n(c,z)$ is an isometry ($\langle e_i,f_j\rangle=\delta_{ij}$, the forms being
conjugate-linear in the first variable), and $n(c,z)n(c',z')=n(c+c',z+z'-c^*J_0c')$, so
$n(c,z)^{-1}=n(-c,z^*)$. The centre is $Z_D=\{n(0,z):z+z^*=0\}$: \emph{its elements are parametrised by the
skew-hermitian $b\times b$ matrices $z$}, and $z\mapsto n(0,z)$ is a homomorphism from the additive group.
For an automorphic function $f$ on $\auU(X)$ and a hermitian $b\times b$ matrix $\beta$ over $E$ the
\emph{centre coefficient} is
\begin{equation}\label{pcu:eq-psibeta}
\pcuF_\beta f(p)=\int_{Z_D(E)\backslash Z_D(\auA)}f(n(0,z)p)\,\psi_\beta(z)^{-1}\,dz ,\qquad
\psi_\beta(z)=\psi\bigl(-\tfrac12\jmath\operatorname{tr}(z\beta)\bigr),
\end{equation}
probability measure, $\psi$ the standard character and $\jmath=-1/\sqrt{-d}$ the constant of the Weil data
({\cite[\S5.1.1]{P49}}). For $z$ skew-hermitian and $\beta$ hermitian, $\operatorname{tr}(z\beta)$ is purely imaginary,
so $-\frac12\jmath\operatorname{tr}(z\beta)$ lies in $\auA_{E^+}$. We use the same formula for $\psi_\beta(z)$ when
$z+z^*=-c^*J_0c$: since $\psi$ is trivial on $\jmath\auA_{E^+}$ ($\operatorname{Tr}_{E/E^+}\jmath=0$), only the
skew-hermitian part $z+\frac12c^*J_0c$ of $z$ contributes. Writing $z=\jmath h$ with $h$ hermitian,
$\psi_\beta(z)=\psi(\operatorname{tr}(h\beta)/2d)$; the trace pairing on hermitian matrices being non-degenerate, the
characters of $Z_D(\auA)$ trivial on $Z_D(E)$ are exactly the $\psi_\beta$, $\beta$ hermitian over $E$. This is the
character by which the centre acts on the terms (Lemma~\ref{pcu:laws}): the term of $x$ is multiplied by
$\psi_{x^*J_Vx}$, and $\pcuF_\beta$ keeps the $x$ with Gram matrix $x^*J_Vx=\beta$.

A pair $(\beta,q)$, $\beta$ hermitian over $E$, $q\in\auGL_b(\auA_E)$, is a \emph{test pair for $U$} if
either (a) $\beta=J_U$ is the form of $U$ and $q=s_0\,x_U\,Y$, where $s_0$ is the finite idele with
$s_0\widehat\auO_E=\alpha\widehat\auO_E$ fixed in {\cite[\S5.1.1]{P49}} (denoted $a$ there; $\alpha$ is the idele by
which the Weil data scale $\psi$; for the set of test pairs only the ideal $s_0\widehat\auO_E$ matters: replacing $s_0$
by $s_0\varepsilon$ with $\varepsilon$ a unit idele replaces $Y$ by $\varepsilon Y$, and $\varepsilon Yk=Y(\varepsilon k)$
with $\varepsilon k\in\auGL_b(\widehat\auO_E)$), $x_U$ the adelic frame of $U$, and the Gram matrix of the columns
of $x_UY_f$ for the form of $U$ lies in $\auGL_b(\widehat\auO_E)$, i.e.\ these columns span a lattice self-dual
for $J_U$ (a condition invariant under $Y\mapsto Yk$, $k\in\auGL_b(\widehat\auO_E)$, and $Y\mapsto\varepsilon Y$); or (b) $\beta$ is invertible and not of the form $a^*\mathrm{diag}(t')a$ at some infinite
place.

\emph{Coordinates.} At the finite places, matrices of $\auHom(X,V)$ are written in the coordinates
$\tilde x=s_0^{-1}x$ of Lemma~\ref{pgx:dict} (on both components at a split place). In them the self-dual
lattice $\alpha\auHom(L_X,L_V)$, whose vacuum is the datum at the finite places outside $S\cup\bar S$
(the \emph{standard vacuum}), is the lattice of integral matrices, at $w\in S$ the data of
Definition~\ref{pgx:conv} are functions of $\tilde x$ (the paragraph after Lemma~\ref{pcx:mod}),
and ``integral'' always refers to these coordinates. By Lemma~\ref{pgx:dict}(d), a point
$x=s_0y$ of $\auHom(\ell,V)(\auA_f)$ has coordinate $\tilde x=y$.

Write $z_0\in\pcMat_{b\times(2b+n)}(\auO_S)$ for the block $f^*$ of the inverse of the frame of $D$
read at $S$ (the matrix whose rows are the coordinate functionals of the $\ell'$-part), so that
$\Phi^{\mathrm{row}}_{z_0}$ and $\Psi^{\mathrm{col}}_{z_0}$ are the standard data of the frame of $D$.

\subsubsection{Smooth Witt data}\label{pcu:ss-witt}

\begin{lemma}[Smooth local Witt theorem for spans]\label{pcu:witt}
Let $J$ be a hermitian $m\times m$ complex matrix and $x$ an $m\times b$ matrix with $\det(x^*Jx)\ne
0$. There is a smooth map $k$ from $\pcMat_{m\times b}(\C)$ to the unitary group $\auU(J)$ such that
for all $y$ near $x$ there is $c\in\pcMat_b(\C)$ with $k(y)\,x\,c=y$.
\end{lemma}
\begin{proof}
Put $w=(x^*Jx)^{-1}x^*$, so $ywJx=y$ for every $y$. For $a\in\pcMat_m(\C)$ let $L(a)=(axw-xwa^*)J$.
Then $L(a)^*J+JL(a)=0$, i.e.\ $L(a)$ is in the Lie algebra of $\auU(J)$, and $(a-L(a))x=x(wa^*Jx)$
because $L(a)x=axwJx-xwa^*Jx=ax-xwa^*Jx$. Let $T(a)=\exp(L(a))(a-L(a))$. Since $L(1)=0$, $T(1)=1$,
and the derivative of $T$ at $1$ is $(\mathrm{id}-L)+L=\mathrm{id}$. By the inverse function
theorem there is a smooth $A$ on $\pcMat_m(\C)$ with $T(A(z))=z$ for $z$ near $1$ (extend a local
inverse smoothly). For $y$ put $z(y)=1-(x-y)wJ$, which is near $1$ for $y$ near $x$, and
$k(y)=\exp(L(A(z(y))))\in\auU(J)$, smooth in $y$. For $y$ near $x$, with $a=A(z(y))$,
\[
k(y)\,x\,(wa^*Jx)=\exp(L(a))(a-L(a))x=T(a)x=z(y)x=x-(x-y)wJx=y. \qedhere
\]
\end{proof}

\begin{lemma}[Witt data with a bump]\label{pcu:wittbump}
For $J=\mathrm{diag}(t'\oplus s)$ and $x_0$ as above there are: a compact neighbourhood $N$ of
$\mathrm{diag}(t')$ in $\auGL_b(\C)$ and a smooth $\iota$ on $\pcMat_b(\C)$ with $\iota(A)=A^{-1}$ for $A\in N$; a smooth $k\colon\pcMat_{(b+n)\times b}(\C)\to\auU(J)$; and
a smooth compactly supported $\sigma\ge0$ on $\mathrm{Fr}$ with $\sigma(x_0,\dots,x_0)=1$, such that
whenever $\sigma(y)\ne0$, for every $\tau$: $y_\tau^*Jy_\tau\in N$ (so
$\iota(y_\tau^*Jy_\tau)=(y_\tau^*Jy_\tau)^{-1}$), and $k(y_\tau)x_0c=y_\tau$ for some $c$.
\end{lemma}
\begin{proof}
Let $N$ be a compact neighbourhood of $\mathrm{diag}(t')$ in the open set $\auGL_b(\C)$, $\chi$ a smooth
function with compact support in $\auGL_b(\C)$ equal to $1$ on $N$, and $\iota(A)=\chi(A)\operatorname{adj}(A)/\det A$
($0$ off the support of $\chi$); it is smooth on $\pcMat_b(\C)$ and inverts on $N$. Lemma~\ref{pcu:witt} gives $k$.
The conditions $y^*Jy\in N$ and $k(y)x_0c=y$ hold on a neighbourhood of $x_0$ (the first because
$y\mapsto y^*Jy$ is continuous and $x_0^*Jx_0=\mathrm{diag}(t')$ is interior to $N$), and $\sigma$ is a
product over $\tau$ of a bump supported in it.
\end{proof}

For $x\in\pcMat_{(b+n)\times b}(\C)$ let $\mathrm{pr}(x)=x\,\iota(x^*Jx)\,x^*Jx_0$ (where $\iota$ inverts,
the $J$-orthogonal projection of $x_0$ to the span of $x$); $\mathrm{pr}$ is smooth on all of
$\pcMat_{(b+n)\times b}(\C)$. If $C$ is invertible and $\iota$ inverts the Gram matrices of $x$ and $xC$, then
$\mathrm{pr}(xC)=xC\,C^{-1}(x^*Jx)^{-1}C^{*-1}C^*x^*Jx_0=\mathrm{pr}(x)$.

\begin{lemma}\label{pcu:frame}
If $\sigma(\mathrm{pr}\circ x)\ne0$ then for every $\tau$, $x_\tau=k(\mathrm{pr}(x_\tau))x_0c$ with
$c^*\mathrm{diag}(t')c=x_\tau^*Jx_\tau$; in particular $x_\tau^*Jx_\tau$ has the form
$c^*\mathrm{diag}(t')c$.
\end{lemma}
\begin{proof}
By Lemma~\ref{pcu:wittbump} applied to $y=\mathrm{pr}\circ x$, $y_\tau=k(y_\tau)x_0c_1$ and $y_\tau$ has
invertible Gram matrix; since $y_\tau=x_\tau A$ with $A=\iota(\cdot)x_\tau^*Jx_0$ and
$\det(y_\tau^*Jy_\tau)\ne0$, $A$ is invertible and $x_\tau=k(y_\tau)x_0c_1A^{-1}$. Put $c=c_1A^{-1}$;
as $k(y_\tau)$ is unitary, $x_\tau^*Jx_\tau=c^*x_0^*Jx_0c=c^*\mathrm{diag}(t')c$.
\end{proof}

\subsubsection{The two real vectors and the identity of periods}\label{pcu:ss-period}

Let $\mu$ be the weight ($\mu_\tau\ge1$). Let $w$ be the Fock vector of $(W_0,U)$ and $z$ that of
$(W_0,W_0)$ given in Proposition~\ref{pcx:P-top}. Let $\omega^\flat$ be the real Weil representation
of $\auU(W_0)_\R\times\auU(T)_\R$ on the Fock space of $\auHom(W_0,T)_\R$ twisted by
$\prod_\tau\det(h_\tau)^{j-\mu_\tau}$ on $\auU(T)_\R$, and $w\boxtimes z$ the external product (a vector
of $\auHom(W_0,U\oplus W_0)_\R$ with finite moments, its moments being bounded by products of those of
$w$, $z$ and the vacuum, since it is an iterated integral of translated vacua,
Lemma~\ref{pcx:superpos}).

\emph{The vector $\Lambda$ for $(X,U)$.} By Lemma~\ref{pcu:bump} with
$\tau(B)=\sigma((B_\tau;0)_\tau)\,\bigl|\prod_\tau\det(B_\tau)^{j+\mu_\tau}\bigr|^2$ ($\tau(1)=1$) there are
$\kappa$ and $\Lambda$ with finite moments, zero term $w$ along $D^{\mathrm s}$, and
$\mathrm{sl}_D(B)\Lambda=\sigma((B;0))\,\bigl|\prod_\tau\det B_\tau^{j+\mu_\tau}\bigr|^2\kappa\,w$.

\emph{The vector $\Phi$ for $(X,T)$.} For $y,x\in\mathrm{Fr}$ put $x'_\tau=k(y_\tau)^{-1}x_\tau$,
$(x'_\tau)_1$ its upper $b\times b$ block, and
\[
\pcuS(y,x)=\sigma(y)\sigma(x')\,\kappa\prod_\tau\overline{\det(x'_\tau)_1}^{\,2\mu_\tau}\;
\omega^\flat(1,k\circ y)(w\boxtimes z).
\]
The function $x\mapsto\pcuS(\mathrm{pr}\circ x,x)$ is smooth and compactly supported (on its support
$x_\tau=k(y_\tau)x'_\tau$ with $y$, $x'$ in the support of $\sigma$), of the form
$\tau(x)\omega^\flat(u(x))v$ with $\tau$, $u$ smooth (the twist is a polynomial in the entries of
$k(y_\tau)$ and their conjugates). By (M1) there is $\Phi$ with finite moments and
$\mathrm{sl}_D(x)\Phi=\pcuS(\mathrm{pr}\circ x,x)$.

\begin{lemma}[Support of the slices]\label{pcu:support}
If $\mathrm{sl}_D(x)\Phi\ne0$ then $x_\tau^*Jx_\tau$ has the form $c^*\mathrm{diag}(t')c$ for every
$\tau$.
\end{lemma}
\begin{proof}
Then $\sigma(\mathrm{pr}\circ x)\ne0$; apply Lemma~\ref{pcu:frame}.
\end{proof}

For a finite datum $F$ of a pair $(W_0,V)$, an adelic matrix $x$, a point $(g,h)$ and characters
$\alpha,\beta$, write $\vartheta^{\alpha,\beta}_{F,x,(g,h)}(\phi)=\theta\bigl(\psrho(x)\omega_{\alpha,
\beta}(g,h)(\phi\otimes F)\bigr)$ for the theta functional on real vectors $\phi$.

\begin{lemma}[Product formula]\label{pcu:product}
For a finite datum $f$ of $(W_0,U)$, a datum $F$ of $(W_0,W_0)$, real vectors $\phi$, $v_1$ and
characters $\alpha,\beta$,
\[
\vartheta^{\alpha,\alpha\beta}_{f\boxtimes F,(x;0),(g,y\oplus l)}(\phi\boxtimes v_1)=
\vartheta^{\alpha,\beta}_{f,x,(g,y)}(\phi)\;\theta^{W_0,W_0}_{\alpha,\alpha}(g,l;v_1\otimes F).
\]
\end{lemma}
\begin{proof}
$\auHom(W_0,U\oplus W_0)=\auHom(W_0,U)\oplus\auHom(W_0,W_0)$ orthogonally; the Weil representation,
the Heisenberg operators and $\theta$ are tensor products along this decomposition (the character
$\alpha\beta$ on $\auU(U\oplus W_0)$ restricted to $\auU(U)\times\auU(W_0)$ is $\alpha\beta\otimes
\alpha\beta$, and the factor $\beta$ is assigned to the first factor, $\alpha$ to the second, by the
convention \eqref{pcx:eq-twist}), and the translation is $0$ on the second factor.
\end{proof}

Let $\gamma=\nu^{-1}\upsilon^{j}$ (the character of \S\ref{pcx:ss-arch} for the rank $n$ of $W_0$),
$I_0(\nu)=\Theta^{\mathrm L}_{W_0,W_0}(z,\Psi^{\mathrm{col}}_1)(1)$, and for $h\in\auU(T)_\R$ let
$\rho(h)=\sigma\bigl(\mathrm{pr}\circ(h_\tau^{-1}x_0)_\tau\bigr)\ge0$; $\rho$ is continuous and
$\rho(1)=1$. For $B\in\auGL_b(\auA_E)$ let $\zeta_B$ be the idele equal to $\det B$ at the infinite
places and $1$ at the finite ones, and $\mathrm N(\zeta)=\prod_\tau|\zeta_\tau|^2$.

\begin{lemma}[Identity of periods]\label{pcu:period}
For every $h\in\auU(T)_\R$ there is $p\in\auU(U)_\R\times\auU(W_0)_\R$ such that for all
$B\in\auGL_b(\auA_E)$, characters $\beta$, finite data $f$ of $(W_0,U)$, adelic $x$, and $g$, $y$,
\[
\int_{[\auU(W_0)]}\vartheta^{\gamma,\gamma\beta}_{f\boxtimes F_1,(x;0),(g,y\oplus u)}
\Bigl(\omega^\flat(1,ph)\,\mathrm{sl}_D\bigl((ph)^{-1}(B;0)\bigr)\Phi\Bigr)du
=\frac{\rho(h)\,I_0(\nu)}{\nu_{\mathrm{bc}}(\zeta_B)\,\mathrm N(\zeta_B)^{j}}\;
\vartheta^{\gamma,\beta}_{f,x,(g,y)}\bigl(\mathrm{sl}_D(B)\Lambda\bigr),
\]
where $F_1$ is the finite datum $\Psi^{\mathrm{col}}_1=\chi_S^{-1}\circ\det$ at $S$ and the vacuum
elsewhere, and $\nu_{\mathrm{bc}}(t)=\nu(t/\bar t)$.
\end{lemma}
\begin{proof}
Let $x_{1,\tau}=h_\tau^{-1}x_0$, so $x_{1,\tau}^*Jx_{1,\tau}=\mathrm{diag}(t')$ and $\iota$ inverts it;
let $y=\mathrm{pr}\circ x_1$, so $\rho(h)=\sigma(y)$.

If $\sigma(y)\ne0$, Lemma~\ref{pcu:frame} gives $x_{1,\tau}=k(y_\tau)x_0u_\tau$ with $u_\tau^*\mathrm{diag}
(t')u_\tau=\mathrm{diag}(t')$, i.e.\ $u\in\auU(U)_\R$. Then $h_\tau k(y_\tau)\mathrm{diag}(u_\tau,1)$ fixes
$x_0$, hence (being unitary) is $\mathrm{diag}(1,q_\tau)$ with $q_\tau\in\auU(W_0)_\R$. Put
$p=(u,q^{-1})$, so $p h=k\circ y^{-1}$ pointwise. If $\sigma(y)=0$ take $p=1$.

At each $\tau$, $(ph)^{-1}(B;0)=h^{-1}x_0u^{-1}B_\tau=x_{1,\tau}C_\tau$ with $C_\tau=u_\tau^{-1}B_\tau$
invertible, and the slice of $\Phi$ there is $\pcuS(y',x_1C)$ with $y'=\mathrm{pr}\circ(x_1C)$.
\emph{If this slice is non-zero, then $y'=y$.} Indeed then $\sigma(x'')\ne0$ for $x''_\tau=k(y'_\tau)^{-1}x_{1,\tau}C_\tau$,
so by Lemma~\ref{pcu:wittbump} the Gram matrix $x''^*_\tau Jx''_\tau=C_\tau^*\mathrm{diag}(t')C_\tau$ ($k(y'_\tau)$ is
unitary) lies in $N$; as $\mathrm{diag}(t')\in N$ too, $\iota$ inverts the Gram matrices of $x_1$ and of $x_1C$,
and $y'=\mathrm{pr}\circ(x_1C)=\mathrm{pr}\circ x_1=y$.
If $\sigma(y)=0$ the slice therefore vanishes (it contains the factor $\sigma(y')=\sigma(y)$ when non-zero), and so does the right
side ($\rho(h)=0$). Otherwise $k(y_\tau)^{-1}x_{1,\tau}C_\tau=x_0u_\tau C_\tau=(B_\tau;0)$. If $\sigma((B;0))=0$, the right
side vanishes ($\mathrm{sl}_D(B)\Lambda$ contains the factor $\sigma((B;0))$), and so does the left: a non-zero slice
would have $y'=y$, hence contain the factor $\sigma(k(y)^{-1}x_1C)=\sigma((B;0))=0$. If $\sigma((B;0))\ne0$, then
$C_\tau^*\mathrm{diag}(t')C_\tau=(B_\tau;0)^*J(B_\tau;0)\in N$ by Lemma~\ref{pcu:wittbump}, so again $y'=y$ and the
slice is $\pcuS(y,x_1C)$. It remains to treat this last case, $\sigma(y)\sigma((B;0))\ne0$, where
\[
\omega^\flat(1,ph)\pcuS(y,x_1C)=\sigma(y)\,\sigma((B;0))\,\kappa\prod_\tau\overline{\det B_\tau}^{\,2\mu_\tau}
\,(w\boxtimes z),
\]
the operators $\omega^\flat(1,ph)$ and $\omega^\flat(1,k\circ y)$ cancelling. By
Lemma~\ref{pcu:product} with $\alpha=\gamma$, the integrand is this scalar times
$\vartheta^{\gamma,\beta}_{f,x,(g,y)}(w)\,\theta^{W_0,W_0}_{\gamma,\gamma}(g,u;z\otimes F_1)$; the integral
over $u$ of the last factor is $\Theta^{\mathrm L}_{W_0,W_0}(z,\Psi^{\mathrm{col}}_1)(g)=I_0(\nu)$ by
Lemma~\ref{pcx:const}. On the right, $\mathrm{sl}_D(B)\Lambda=\sigma((B;0))|\prod\det B_\tau^{j+\mu_\tau}|^2
\kappa w$. It remains to see $\prod_\tau\bar\zeta_\tau^{2\mu_\tau}=|\prod_\tau\zeta_\tau^{j+\mu_\tau}|^2/
(\nu_{\mathrm{bc}}(\zeta)\mathrm N(\zeta)^j)$ for $\zeta=\zeta_B$. With $P=\prod\zeta_\tau$ and
$Q=\prod\zeta_\tau^{\mu_\tau}$ one has $\nu_{\mathrm{bc}}(\zeta)=Q^2/|Q|^2$ (weight $\mu$ at the
infinite places, \S\ref{pcx:ss-statement}), $\mathrm N(\zeta)^j=|P|^{2j}$, and
$|P^jQ|^2|Q|^2/(Q^2|P|^{2j})=|Q|^4/Q^2=\bar Q^2$.
\end{proof}

\subsubsection{Centre coefficients at test pairs}\label{pcu:ss-coeff}

Put $G=\auU(T)(\auA)$, $H=\auU(U)(\auA)\times\auU(W_0)(\auA)\subset G$ (block diagonal), and
$\Gamma_2=\auU(W_0)(E)$. The unfolding of the theta kernel along $D$ uses three facts, stated here;
(U1), (U2) are proved in \S\ref{pcu:ss-open} and (U3) below.

\begin{enumerate}
\item[(U1)] \emph{Unfolding.} Let $V$ be a space of rank $a$, $F=[\phi,\Phi]$ a datum of
$\auHom(X,V)$ with $\Phi$ one of the standard data $\Phi^{\mathrm{row}}_{z_0}$ (for $V=T$) or
$(\Psi^{\mathrm{col}}_{z_0})^{\mathsf t}$ at $\bar S$ (for $V=U$), and $p=n(c,z)m(q,g)$. The kernel
$\theta(\omega(p,h)F)$ is $\sum_{x\in\pcMat_{a\times b}(E)}\mathrm{term}_x(p,h)$, where
$\mathrm{term}_x(p,h)$ is the partial theta sum over the $\xi\in\auHom(X,V)(E)$ whose restriction to
$\ell$ is $x$. Let $\beta$ be invertible. If there is no $x$ with $x^*J_Vx=\beta$, or if
$\mathrm{term}_x(p,\cdot)\equiv0$ for all such $x$, then $\pcuF_\beta\int_{[\auU(V)]}\theta(\omega(\cdot,h)F)dh$
vanishes at $p$. Otherwise the $x$ with $x^*J_Vx=\beta$ form one $\auU(V)(E)$-orbit, and if
$\mathrm{term}_{x}(p,\cdot)$ is absolutely integrable on $\mathrm{Stab}(x)(E)\backslash\auU(V)(\auA)$ (condition
\eqref{pcu:eq-unfoldabs}) then $\pcuF_\beta\int_{[\auU(V)]}\theta(\omega(\cdot,h)F)dh$ at $p$ equals
$\int_{\mathrm{Stab}(x)(E)\backslash\auU(V)(\auA)}\mathrm{term}_x(p,h)\,dh$. Condition \eqref{pcu:eq-unfoldabs}
holds in the two situations of Lemma~\ref{pcu:integrable} (the orbit of $x_0$ on the large side and the
orbit of $1$ on the small side, at a test pair of type (a)), which are the ones used in
Lemma~\ref{pcu:coeff} and in type (a) of the proof of Proposition~\ref{pct:vantop}.
\item[(U2)] \emph{Term formula.} With $x'=x_V^{-1}x$ (standard coordinates) and $x'q=s_0Y$,
\[
\mathrm{term}_x(p,h)=\psi_{x^*J_Vx}(z)\,\chi_V(\det q)^{-1}
|\det q|^{a/2}\,\mathbf 1_{\mathrm{int}}(y)\,\varepsilon_V\;
\vartheta_{F^{(y)},\,x'c^*J_{W_0},\,(g,h)}\bigl(\mathrm{sl}_D(h^{-1}Y)\phi\bigr),
\]
where $y$ is the finite part of $h^{-1}Y$, $\mathbf 1_{\mathrm{int}}(y)$ is $1$ if $y$ is integral and
$0$ otherwise, $\varepsilon_V\ne0$ is the constant in the slices of the vacuum, $F^{(y)}$ is the
finite datum of $(W_0,V)$ equal at $S$ (or $\bar S$) to the slice $y'\mapsto\Phi((y\ y'))$ of $\Phi$
at $y$ and to the vacuum elsewhere, $\chi_V(u)=\alpha_V(u/\bar u)\upsilon(u)^{-a}$ for the source
character $\alpha_V$, and the characters of $\vartheta$ are those of the pair $(X,V)$ with the target
character divided by $\upsilon^b$.
\item[(U3)] \emph{Stages.} Let $\Theta\colon\Gamma_2\backslash G\to\C$ be integrable, and suppose
there are a measurable $\rho'\ge0$ on $G$, left $H$-invariant, non-zero on a neighbourhood of $1$, and
$c_0\in\C$ with $\int_{\Gamma_2\backslash H}\Theta(\xi g)\,d\xi=c_0\rho'(g)$ for almost all $g$. Then
$\int_{\Gamma_2\backslash G}\Theta=R\,c_0$ with $R=\int_{H\backslash G}\rho'\in(0,\infty]$, where if
$R=\infty$ then $c_0=0$ and the identity reads $\int\Theta=0$.
\end{enumerate}
\begin{proof}[Proof of (U3)]
By the quotient integral formula, $\int_{\Gamma_2\backslash G}\Theta=\int_{H\backslash G}
\int_{\Gamma_2\backslash H}\Theta(\xi g)d\xi\,dg=c_0\int_{H\backslash G}\rho'$. $R>0$ because
$\rho'$ is non-zero near $1$ and $H\backslash G$ carries a quasi-invariant measure positive on open
sets; if $R=\infty$, integrability of $\Theta$ forces $c_0=0$.
\end{proof}

\begin{lemma}[Slices of a datum at a split place]\label{pcu:W5slice}
Let $v=w\bar w$ be a place of $E^+$ split in $E$, $V$ a space of rank $a$, and realise $\psHS_v(\auHom(X,V))$ on
$\auS(\auHom(X_w,V_w))$ as in \textup{(W5)}: $F$ is a function of $a\times(2b+n)$ matrices
$\xi=(\xi_\ell\ \xi_{\ell'}\ \xi_{W_0})$, written in the frame $(e,f,\text{frame of }W_0)$ of $D$, with
$\psrho(a_w,0)$ acting by translation and $\psrho(0,a_{\bar w})$ by multiplication by $[\xi,a_{\bar w}]$. Then the
slice of $F$ along $D$ at a point $u=(u_w,u_{\bar w})\in\auHom(\ell,V)_v=\pcMat_{a\times b}(E_w)\times\pcMat_{a\times
b}(E_{\bar w})$ is, up to a constant $c_v>0$ depending only on Haar measures, the function of
$\xi_{W_0}\in\auHom(W_{0,w},V_w)$ (the \textup{(W5)} model of $\auHom(W_0,V)_v$)
\[
F[u](\xi_{W_0})=c_v\int_{\pcMat_{a\times b}(E_w)}F(u_w\ t\ \xi_{W_0})\,[t,u_{\bar w}]\,dt ,
\]
where $t$ runs through the $\ell'$-columns and $[t,u_{\bar w}]$ is the pairing $[\,,]$ between $\auHom(\ell'_w,V_w)$
and $\auHom(\ell_{\bar w},V_{\bar w})$. If $F$ is the indicator of the integral $\ell'$-columns times a function
$F_1(\xi_\ell,\xi_{W_0})$, and the integral matrices of $\auHom(\ell'_w,V_w)$ and $\auHom(\ell_{\bar w},V_{\bar w})$
are dual lattices for $[\,,]$, then $F[u](\xi_{W_0})=c'_v\,\mathbf 1_{\mathrm{int}}(u_{\bar w})\,F_1(u_w,\xi_{W_0})$,
$c'_v>0$.
\end{lemma}
\begin{proof}
In the mixed model of $D$ at $v$ the Schr\"odinger variable is $u\in\auHom(\ell,V)_v=I^*_v$ (maps vanishing on
$\ell'\oplus W_0$), $I_v$ (maps vanishing on $\ell\oplus W_0$) acts by multiplication, and $\auHom(W_0,V)_v$ on the
values (\S\ref{pcu:ss-open}). Compare with the \textup{(W5)} model, which is the Schr\"odinger model of the
polarisation $\auHom(X_w,V_w)\oplus\auHom(X_{\bar w},V_{\bar w})$. The two polarisations agree on
$\auHom(\ell_w,V_w)$ (translated in both), on $\auHom(\ell'_{\bar w},V_{\bar w})$ (multiplication in both), and on
$\auHom(W_0,V)_v$, which is the orthogonal summand on which both are the \textup{(W5)} model of $(W_0,V)$; they differ
exactly on $\auHom(\ell'_w,V_w)$ (translated in \textup{(W5)}, in $I_v$ for $D$) and $\auHom(\ell_{\bar w},V_{\bar w})$
(multiplication in \textup{(W5)}, in $I^*_v$ for $D$). These two are in duality for $[\,,]$, since
$\langle e_i,f_j\rangle=\delta_{ij}$ and the form pairs the $w$-component with the $\bar w$-component. The
intertwiner between the Schr\"odinger models of the two transversal Lagrangians of the symplectic space
$\auHom(\ell'_w,V_w)\oplus\auHom(\ell_{\bar w},V_{\bar w})$ is the partial Fourier transform
$\varphi\mapsto c_v\int\varphi(t)[t,\cdot\,]\,dt$ (Weil's intertwining operator, unique up to a scalar by (P1); one
checks directly that it carries translation by $t_0$ to multiplication by $[\cdot,t_0]=[t_0,\cdot\,]^{-1}$ and multiplication by
$[\cdot,y_0]$ to translation by $y_0$, which are the operators of (P2) in the new model). Tensoring with the identity on the other factors gives
the intertwiner from the \textup{(W5)} model to the mixed model, and the slice at $u$ is its value at $u$; this is the
formula. For the second statement, $\int_{\mathrm{int}}[t,u_{\bar w}]\,dt$ is the volume of the integral matrices if
$u_{\bar w}$ lies in the dual lattice, which is the lattice of integral matrices by hypothesis, and $0$ otherwise.
\end{proof}

\begin{lemma}[Finite slices at test pairs of type (a)]\label{pcu:finslice}
Let $(J_U,q)$ be a test pair of type (a), $q=s_0x_UY$, and $p=n(c,z)m(q,g)$; let $F=[\phi,\Phi^{\mathrm{row}}_{z_0}]_S$
($V=T$), $K_f$ the stabiliser in $\auU(T)(\auA_f)$ of the standard lattice $L_T=L_U\perp L_{W_0}$, $H_f=\auU(U)(\auA_f)
\times\auU(W_0)(\auA_f)$, and let $M_0\subset U(\auA_f)$ be the lattice spanned by the columns of $x_0Y_f$ (self-dual for the
form of $U$, by the definition of a test pair). The slice of $\Phi=\Phi^{\mathrm{row}}_{z_0}=\epsilon\cdot(\text{standard datum of
the frame of }D)$ entering (U2) is computed as follows. In the frame of $D$ at $S$ the rows of $z_0$ are the coordinate functionals of
$\ell'$, so for $\xi=(\xi_\ell\ \xi_{\ell'}\ \xi_{W_0})$ the matrix $\binom{z_0}{\xi}$ is $\begin{pmatrix}0&1_b&0\\\xi_\ell&\xi_{\ell'}&\xi_{W_0}
\end{pmatrix}$; it lies in $\auGL_{2b+n}(\auO_S)$ if and only if $\xi_{\ell'}$ is integral and $(\xi_\ell\ \xi_{W_0})\in\auGL_{b+n}(\auO_S)$, and
then its determinant is $\pm\det(\xi_\ell\ \xi_{W_0})$ (expansion along the rows of $z_0$), with $\chi_S(\pm1)=\nu(1)=1$. Hence
$\Phi(\xi)=\epsilon\,\mathbf 1_{\mathrm{int}}(\xi_{\ell'})\,\Phi_1(\xi_\ell,\xi_{W_0})$ with $\Phi_1(y,y')=\chi_S^{-1}(\det(y\ y'))$ if the square
matrix $(y\ y')$ (columns of $\ell\oplus W_0$) lies in $\auGL_{b+n}(\auO_S)$ and $0$ otherwise. At each place $v=w\bar w$ with $w\in S$
the frame of $D$ is integral with unimodular Gram matrix, so $\alpha\auHom(L_X,L_T)_v$ is the sum of its intersections
with the summands of the decomposition along $\ell,\ell',W_0$; hence (Lemma~\ref{pgx:dict}(b)) the integral $\ell'_w$-columns
and the integral matrices of $\auHom(\ell_{\bar w},T_{\bar w})$, both in the coordinates $\tilde x$, are dual lattices, and Lemma~\ref{pcu:W5slice} gives the slice at $u=(u_w,u_{\bar w})$:
$c'_v\,\mathbf 1_{\mathrm{int}}(u_{\bar w})\,\epsilon\,\Phi_1(u_w,\cdot)$. The factor $\mathbf 1_{\mathrm{int}}(u_{\bar w})$ is the integrality
condition at $\bar S$ recorded in $\mathbf 1_{\mathrm{int}}(y)$ in (U2), and $c'_v$ is part of $\varepsilon_T$; so in (U2) the datum
$F^{(y)}$ at $S$ is $y'\mapsto\Phi((y\ y'))=\epsilon\,\Phi_1(y,y')$, $y$ the $S$-component of the slice point.
\begin{enumerate}
\item[\textup{(a)}] For $\kappa\in K_f$, $\omega_{\gamma_X,\gamma_X}(1,\kappa)F=F$; hence $\mathrm{term}_{x_0}(p,h\kappa)=\mathrm{term}_{x_0}(p,h)$.
\item[\textup{(b)}] If $\mathrm{term}_{x_0}(p,h)\ne0$, then the finite part of $h^{-1}x_0Y$ is integral and $h_f\in H_fK_f$.
\item[\textup{(c)}] Let $h_f=(y_f,y'_f)\in H_f$ and $x_f=y_f^{-1}Y_f$. If $\mathrm{term}_{x_0}(p,h)\ne0$ then $x_f$ is integral,
$\det x_f\in\auO_S^\times$, and the finite datum $F^{(y)}$ of (U2) at $S$ is
$\epsilon\,\chi_S^{-1}(\det x_f)\,f_U\otimes F_1$, where $f_U$ is the vacuum datum (indicator of the integral matrices) of
$\auHom(W_0,U)$ at $S$ and $F_1=\Psi^{\mathrm{col}}_{1_n}$ is the datum of $\auHom(W_0,W_0)$ at $S$ of
Definition~\ref{pcx:def-std} (with $z=1_n$), supported on $\auGL_n(\auO_S)$.
\end{enumerate}
\end{lemma}
\begin{proof}
(a) $\mathrm{term}_{x_0}(p,h\kappa)=\theta_{W_0}((\omega(p,h)\omega(1,\kappa)F)[x_0])$, so the second claim follows from the
first. Here $\omega=\omega_{\gamma_X,\gamma_X}$ with $\gamma_X=\nu^{-1}\upsilon^{(2b+n)/2}=\gamma\,\upsilon^{b}$ (the character attached to the rank $2b+n$ of $X$; $\gamma=\nu^{-1}\upsilon^j$ as before Lemma~\ref{pcu:period}, the shift by $\upsilon^b$ being the one in (U2)), so $\omega(1,\kappa)=\gamma_X(\det\kappa)\,
\omega_{\auHom(X,T)}(1,\kappa)$, and $\det\kappa=(\det\kappa_v)_v$ with $\det\kappa_v$ a norm-one unit at every finite place $v$ of
$E^+$ ($\kappa_v$ stabilises the lattice $L_{T,v}$). Two facts on the characters are used: $\upsilon$ is unramified at every finite
place, i.e.\ trivial on $\widehat\auO_E^\times$ ({\cite[\S5.1.1]{P49}}); and $\nu$ is unramified outside $\pcm$, i.e.\ trivial on
$\auU(1)(\auO_{E^+_v})$ for $v$ not below a prime of $S$ (\S\ref{pcx:ss-statement}). We compute $\omega(1,\kappa_v)$ place by place.

\emph{$v$ not below $S$.} The datum at $v$ is the vacuum of the self-dual lattice $\alpha\auHom(L_X,L_T)_v$ (the integral
matrices in the coordinates $\tilde x$), which $\kappa_v$
stabilises; by the last assertion of (W4) (for the splitting by $\upsilon$ of {\cite[Definition~2.1]{P49}}),
$\omega_{\auHom(X,T),v}(1,\kappa_v)$ fixes it. The character gives $(\gamma_X)_v(\det\kappa_v)=\nu_v(\det\kappa_v)^{-1}
\upsilon_v(\det\kappa_v)^{(2b+n)/2}=1$, both factors being trivial on norm-one units by the two facts above.

\emph{$v=w\bar w$ with $w\in S$.} By {\cite[Proposition~2.17]{P49}} (first formula, $g=1$, source rank $2b+n$),
$(\omega_{\auHom(X,T),v}(1,\kappa_v)\Phi)(\xi)=\upsilon_w(\det\kappa_w)^{-(2b+n)}|\det\kappa_w|_w^{-(2b+n)/2}\Phi(\kappa_w^{-1}\xi)$,
where $\kappa_w\in\auGL_{b+n}(\auO_{E_w})$ is the component of $\kappa_v$ on $T_w$; the first two factors are $1$
($\upsilon$ unramified, $\det\kappa_w$ a unit). Since $\binom{z_0}{\kappa_w^{-1}\xi}=\mathrm{diag}(1,\kappa_w^{-1})\binom{z_0}\xi$,
$\Phi(\kappa_w^{-1}\xi)=\chi_S(\det\kappa_w)\Phi(\xi)$ (the support condition is unchanged). The character: under
$\auU(1)(E^+_v)=\{(t,t^{-1})\}\cong E_w^\times$ the element $\det\kappa_v$ is $(t,t^{-1})$ with $t=\det\kappa_w$, and
$\nu(t,t^{-1})=\nu(t/\bar t)=\chi_S(t)$ (for $t\in E_w^\times$, $t/\bar t=(t,t^{-1})$, as in the proof of
{\cite[Proposition~2.17]{P49}}), while $\upsilon_v(\det\kappa_v)=\upsilon_w(t)\upsilon_{\bar w}(t^{-1})=1$. So $(\gamma_X)_v(\det\kappa_v)=
\chi_S(\det\kappa_w)^{-1}$, and the total factor at $v$ is $\chi_S(\det\kappa_w)^{-1}\chi_S(\det\kappa_w)=1$.

The real component of $F$ is not acted on ($\kappa$ is finite). Hence $\omega(1,\kappa)F=F$.

(b) By (U2), $\mathrm{term}_{x_0}(p,h)\ne0$ forces $\mathbf 1_{\mathrm{int}}(y)=1$ for the finite part $y=h_f^{-1}x_0Y_f$; its
columns span $h_f^{-1}M_0$, so $h_f^{-1}M_0\subset L_T$, and $h_f^{-1}M_0$ is unimodular ($h_f$ is an isometry). A
unimodular sublattice $M$ of a unimodular lattice $L$ is an orthogonal summand: for $l\in L$ the form
$m\mapsto\langle m,l\rangle$ is integral on $M$, hence equals $\langle\cdot,m_l\rangle$ for some $m_l\in M$, and
$l-m_l\perp M$; so $L=M\perp(L\cap M^\perp)$ with $L\cap M^\perp$ unimodular. Thus $L_T=h_f^{-1}M_0\perp N_1$ and
$h_fL_T=M_0\perp h_fN_1$, with $h_fN_1$ a unimodular lattice of $M_0^\perp=W_0(\auA_f)$. Self-dual hermitian lattices of
the same rank in the same space are isometric ($E/E^+$ is unramified at every finite place, so they have orthonormal
bases, Jacobowitz~\cite{Jacobowitz}, as in Step~2 of the proof of Proposition~\ref{pcu:selfdual}); so there are
$u_1\in\auU(U)(\auA_f)$ with $u_1M_0=L_U$ and $u_2\in\auU(W_0)(\auA_f)$ with $u_2h_fN_1=L_{W_0}$. Then
$\kappa=(u_1,u_2)h_f$ stabilises $L_T$, i.e.\ $\kappa\in K_f$, and $h_f=(u_1,u_2)^{-1}\kappa\in H_fK_f$.

(c) Here $y=h_f^{-1}x_0Y_f=x_0x_f=\binom{x_f}0$. By (b) $x_f$ is integral. The Gram matrix of the columns of $x_0x_f$
equals that of the columns of $x_0Y_f$ ($y_f$ is an isometry), which lies in $\auGL_b(\widehat\auO_E)$ by the definition
of a test pair; it is $x_f^*G_0x_f$ with $G_0$ the Gram matrix of the integral basis, a unit matrix. Hence
$\det x_f\cdot\overline{\det x_f}$ is a unit at every finite place of $E^+$; as $\det x_f$ is integral (on both components at a
split place), $\det x_f$ is a unit. For $y'=\binom{y'_1}{y'_2}$ ($y'_1$ of size $b\times n$, $y'_2$ of size $n\times n$),
$(y\ y')=\begin{pmatrix}x_f&y'_1\\0&y'_2\end{pmatrix}$ has determinant $\det x_f\det y'_2$, and lies in $\auGL_{b+n}(\auO_S)$
if and only if $y'_1$, $y'_2$ are integral and $\det y'_2$ is a unit (a product of two integers is a unit only if both
are). Hence $F^{(y)}(y')=\epsilon\chi_S^{-1}(\det x_f)\,\mathbf 1_{\mathrm{int}}(y'_1)\,\chi_S^{-1}(\det y'_2)
\mathbf 1_{\auGL_n(\auO_S)}(y'_2)$, which is the claim.
\end{proof}

\begin{lemma}[Coefficients at test pairs]\label{pcu:coeff}
There is $K>0$ such that for every tame $\nu$ of weight $\mu$ there is $C_\nu\in\C$ with
$|C_\nu|=K|I_0(\nu)|$ and, for every test pair $(\beta,q)$ for $U$, all $c,z,g$ and $p=n(c,z)m(q,g)$,
\[
\pcuF_\beta\,\Theta^{\mathrm L}_{X,T}(\Phi,\Phi^{\mathrm{row}}_{z_0})(p)=C_\nu\,
\pcuF_\beta\,\Theta^{\mathrm S}_{X,U}(\Lambda,\Psi^{\mathrm{col}}_{z_0})(p).
\]
\end{lemma}
\begin{proof}
\emph{Case (b).} On the small side ($V=U$, $a=b$) an $x$ with $x^*J_Ux=\beta$ is invertible and
exhibits $\beta$ as $a^*\mathrm{diag}(t')a$ at every infinite place; there is none, and the coefficient
is $0$ by (U1). On the large side, by (U2) the term of $x$ is a multiple of the slice of $\Phi$ at
$h^{-1}Y$, whose Gram matrix is $q'^*\beta q'$ ($q=s_0q'$); by Lemma~\ref{pcu:support} it vanishes
unless $\beta$ is represented by $t'$ at every infinite place, which it is not. Both sides are $0$.

\emph{Case (a): reduction to one term.} Now $\beta=J_U$. On the large side the $x\in\pcMat_{(b+n)\times
b}(E)$ with $x^*J_Tx=J_U$ are the isometric embeddings of $U$ in $T=U\oplus W_0$; by Witt's theorem
they form one orbit, that of $x_0$, with stabiliser $\Gamma_2=\auU(W_0)(E)$. On the small side they
are the elements of $\auU(U)(E)$, one orbit with trivial stabiliser. By (U1), whose condition
\eqref{pcu:eq-unfoldabs} holds on both sides by Lemma~\ref{pcu:integrable},
\[
\pcuF_\beta\Theta^{\mathrm L}(p)=\int_{\Gamma_2\backslash G}\mathrm{term}^{\mathrm L}_{x_0}(p,h)dh,\qquad
\pcuF_\beta\Theta^{\mathrm S}(p)=\int_{\auU(U)(\auA)}\mathrm{term}^{\mathrm S}_1(p,y)dy .
\]

\emph{Inner integral.} Write $q=s_0x_UY$ with the columns of $x_UY_f$ spanning a self-dual lattice; let
$K_f$ be the stabiliser of the standard lattice of $T$ in $\auU(T)(\auA_f)$, and for $g'\in G$ let
$\rho'(g')=\rho(g'_\infty)$ if $g'_f\in H_fK_f$ and $0$ otherwise; $\rho'$ is left $H$-invariant
(Lemma~\ref{pcu:period} gives $\rho$ as a function of the span of $h^{-1}x_0$, which is unchanged by
left multiplication by $H_\infty$) and non-zero near $1$. We claim: for $g'=(h_\infty,u_f)$ there is
$k'\in H$ with
\[
\int_{\auU(U)(\auA)}\int_{[\auU(W_0)]}\mathrm{term}^{\mathrm L}_{x_0}(p,(y,y')k'g')\,dy'\,dy=
\rho'(g')\,\kappa_0\,\epsilon\,\nu_{\mathrm{bc}}(d_0)\,I_0(\nu)\int_{\auU(U)(\auA)}
\mathrm{term}^{\mathrm S}_1(p,y)dy,
\]
with $d_0=\det(s_0x_U)$, $\kappa_0=\varepsilon_T\varepsilon_U^{-1}\mathrm N((d_0)_\infty)^{j}$ (with $\mathrm N$ as defined before Lemma~\ref{pcu:period}, applied to the archimedean components), and $\epsilon$ the
constant of modulus $1$ with $\Phi^{\mathrm{row}}_{z_0}=\epsilon\cdot(\text{standard datum of the frame
of }D)$. If $u_f\notin H_fK_f$, then $(y,y')g'$ has finite part in $(y_f,y'_f)u_f\notin H_fK_f$, so the
term vanishes by Lemma~\ref{pcu:finslice}(b), and both sides are $0$ ($k'=1$). If $u_f=(u_1,u_2)\kappa$ with $\kappa\in K_f$, let $(\delta_1,\delta_2)=p$ be the element
given by Lemma~\ref{pcu:period} for $h_\infty$ and $k'=((\delta_1,u_1^{-1}),(\delta_2,u_2^{-1}))$
(adjusted so that $k'g'=(\delta h_\infty,\kappa)$). The term is right $K_f$-invariant (Lemma~\ref{pcu:finslice}(a)), so $\kappa$ drops. Fix $y\in\auU(U)(\auA)$ and put
$B=y^{-1}Y$, $x_f$ its finite part. By Lemma~\ref{pcu:finslice}(c) on the large side (and (U2)
on the small side), both terms vanish unless $x_f$ is integral, and then $\det x_f$ is a unit and the large finite
slice is $\epsilon\chi_S^{-1}(\det x_f)$ times the vacuum of $(W_0,U)$ tensored with $F_1$; on the small side the constant $\chi_S^{-1}(\det\bar x_f)$ times
the vacuum. The archimedean point $\delta h_\infty$ is moved into the real slice by (M3), and the
integral over $[\auU(W_0)]$ is the left side of Lemma~\ref{pcu:period}; replacing it by the right side
and comparing scalars gives the claim. The scalar comparison uses: $\det B=\zeta_B\cdot\det x_f$ with
$|\det y_\tau|=1$ at every infinite place and $\det x_f$ a unit, so $\mathrm N(\zeta_B)=\prod_\tau|\det Y_\tau|^2$; $\chi_S^{-1}(\bar\epsilon')=
\chi_S^{-1}(\epsilon')\nu_{\mathrm{bc}}(\epsilon')$ for units $\epsilon'$ at $S$ (definition of
$\chi_S$); and the factor $\nu(\det y)^2$ between the characters $\nu\upsilon^j$ and $\nu^{-1}
\upsilon^j$ of $\auU(U)$ on the two sides; $j=n/2$ makes the powers of $\upsilon$ agree.

\emph{Conclusion.} By (U3) applied to $\Theta(h)=\mathrm{term}^{\mathrm L}_{x_0}(p,h)$ (integrable by
Lemma~\ref{pcu:integrable}(a)), the large coefficient is $R\,\kappa_0\epsilon\nu_{\mathrm{bc}}(d_0)I_0(\nu)$ times the small one,
where $R=\int_{H\backslash G}\rho'$ depends only on $\sigma$ and the measures. If $R<\infty$ put
$K=R|\kappa_0|$ and $C_\nu=R\kappa_0\epsilon\nu_{\mathrm{bc}}(d_0)I_0(\nu)$ ($|\epsilon|=|\nu_{\mathrm{bc}}|=1$).
If $R=\infty$ then (U3) gives $0$ on the left and $c_0=0$, i.e.\ $I_0(\nu)\cdot(\text{small coefficient})
=0$, so the identity holds with $K=|\kappa_0|$ and $C_\nu=\kappa_0\epsilon\nu_{\mathrm{bc}}(d_0)I_0(\nu)$.
\end{proof}
\subsubsection{From test pairs to all standard data}\label{pcu:ss-std-data}

Fix $\nu$ and $C=C_\nu$. The difference
$\mathcal D(\Psi)=\mathcal T(\Psi)-C\,\Theta^{\mathrm S}_{X,U}(\Lambda,\Psi)$, where $\mathcal T$ is the
linear map of Lemma~\ref{pcx:linfactor} attached to $\Phi'\mapsto\Theta^{\mathrm L}_{X,T}(\Phi,\Phi')$
(Lemma~\ref{pcx:Sact}(a)), satisfies, by Lemma~\ref{pcx:section-laws} and Lemma~\ref{pcx:Sact}:
(E1) on the space $\mathcal M$ of \S\ref{pcx:ss-match} it is linear with values automorphic functions
of moderate growth (Corollary~\ref{pcx:conv}); (E2) $\mathcal D(\Psi(g\,\cdot))(xg)=\eta(\det g)^{-b}
\mathcal D(\Psi)(x)$ for $g$ at $S$; (E3) its values are right invariant under the stabiliser of the
standard lattice away from $S$; (E4) $\mathcal D(\Psi^{\mathrm{col}}_{z'})=\Theta^{\mathrm L}(\Phi,
\Phi^{\mathrm{row}}_{z'})-C\,\Theta^{\mathrm S}(\Lambda,\Psi^{\mathrm{col}}_{z'})$ for integral $z'$
(Lemmas~\ref{pcx:std-matched}, \ref{pcx:section-laws}(b)). We call it the \emph{equivariant lift}.

\begin{proposition}[Self-dual lattices suffice]\label{pcu:selfdual}
Let $\mathcal D$ satisfy (E1)--(E4) and suppose $\pcuF_\beta\mathcal D(\Psi^{\mathrm{col}}_{z_0})
(n(c,z)m(q,g))=0$ for every test pair $(\beta,q)$ for $U$ and all $c,z,g$. Then $\mathcal D(\Psi)=0$ for
every locally constant $\Psi$ having a matched $\Phi$ for $\chi_S$; in particular
$\mathcal D(\Psi^{\mathrm{col}}_{z'})=0$ for all integral $z'$.
\end{proposition}
\begin{proof}
\emph{Step 1 (reduction to coefficients at unimodular indices).} By Lemma~\ref{pcu:gauss} below it
suffices that $\pcuF_\beta\mathcal D(\Psi^{\mathrm{col}}_{z_0})((1,g_f))=0$ for all $g_f$
trivial at $S$ and all hermitian $\beta$ with $s_0^{-1}\bar s_0^{-1}
\beta$ (the normalised index) in $\auGL_b(\auO_S)$.

\emph{Step 2 (from unimodular indices to test pairs).} For such $\beta$ the statement of
Lemma~\ref{pcu:unimod} reduces the vanishing at $(1,g_f)$ to the vanishing at points
$n(l,z)m(s_0q,g)$ of $P_D$ at which the Gram matrix $q^*\beta q$ is unimodular at all finite places,
i.e.\ $q k$ has orthonormal columns for $\beta$ at the finite places, $(qk)^*\beta(qk)=1$, for some $k\in
\auGL_b(\widehat\auO_E)$ (a unimodular hermitian lattice over $\widehat\auO_E$ has an orthonormal
basis, because $E/E^+$ is unramified at every finite place ({\cite[\S14]{P46}}), so that every unit of
$\widehat\auO_{E^+}$ is a local norm; proof: the reduction of the lattice is a non-degenerate hermitian space over the quadratic extension of residue fields, so it has a vector $\bar x$ with $\bar h(\bar x,\bar x)\ne0$ (otherwise polarisation and the surjectivity of the residue trace give $\bar h=0$); a lift $x$ has $h(x,x)$ a unit of the base ring, hence the norm of a unit since the extension is unramified, so after scaling $h(x,x)=1$; then the lattice is $\mathcal O x\perp(x^\perp\cap\text{lattice})$ with unimodular second summand, and one inducts on the rank; this is the argument at a place inert in $E$; at a place split in $E$, $\auO_{E,v}=\auO_v\times\auO_v$ and a unimodular hermitian lattice is a pair of lattices $(L',L'')$ in perfect duality under the induced bilinear pairing, so a basis of $L'$ together with its dual basis of $L''$ is an orthonormal basis; see also Jacobowitz~\cite{Jacobowitz}).

\emph{Step 3 (Landherr).} Let $(\beta,q,k)$ be as in Step 2. If $\beta$ is not represented by $t'$ at
some infinite place, $(\beta,s_0q)$ is a test pair of type (b). Otherwise $\beta$ has signature $t'$
at every infinite place and an orthonormal lattice at every finite place, so by
Lemma~\ref{pcx:landherr} there are $c\in\auGL_b(E)$ and $w\in\auU(U)(\auA)$ with $\beta=c^*J_Uc$ and
$x_\beta=c^{-1}x_Uw$, where $x_\beta=qk$ (at the finite places). The coefficient is unchanged by the
rational point $m(c,1)\in P_D(E)$: $\pcuF_{c^*J_Uc}f(p)=\pcuF_{J_U}f(m(c,1)p)$ for automorphic $f$, and
$m(c,1)n(l,z)m(s_0q,g)=n(l',z')m(s_0cq,g)$. Now $cq=x_UY$ with $Y=x_U^{-1}cq$, and the Gram matrix of the columns of $x_UY_f=cq$ for $J_U$ is
$(cq)^*J_U(cq)=k^{*-1}\bigl((cqk)^*J_U(cqk)\bigr)k^{-1}=k^{*-1}k^{-1}$ (since
$(cqk)^*J_U(cqk)=(qk)^*\beta(qk)=1$, as $\beta=c^*J_Uc$), which lies in $\auGL_b(\widehat\auO_E)$ as $k$ does; so $(J_U,s_0cq)$ is a test pair of type (a). In both cases the hypothesis gives $0$.
\end{proof}

\begin{lemma}[Gauss sums and weak approximation]\label{pcu:gauss}
Under (E1)--(E4), if $\pcuF_\beta\mathcal D(\Psi^{\mathrm{col}}_{z_0})(1,g_f)=0$ for every $g_f$ trivial at
$S$ and every $\beta$ with unimodular normalised index, then $\mathcal D(\Psi)=0$ for every locally
constant $\Psi$ with a matched function.
\end{lemma}
\begin{proof}
Let $\mathcal Q=\mathcal M\cap\ker\mathcal D$; by (E2) it is stable under $\Psi\mapsto\Psi(g\,\cdot)$. Let
$\Psi_0=\Psi^{\mathrm{col}}_{z_0}$ and, for $w\in\pcMat_b(F_S)$, $\Psi_w(y)=\Psi_0(y+L\,w\,(z_0y))$ with $L$
the $\ell$-block of the frame of $D$ at $S$ (the shear by $n(0,w)$ at $S$). For $f\in\mathcal M$: an
automorphic function vanishes if it vanishes at all $(1,g_f)$ with $g_f$ trivial at $S$ (weak
approximation at $S$ and (E3), since $\mathcal D f$ is continuous and left $\auU(X)(E)$-invariant), and
a function on $Z_D(E)\backslash Z_D(\auA)$ vanishes if all its Fourier coefficients do, the characters of
this compact group being the $\psi_\beta$, $\beta$ hermitian over $E$ (\S\ref{pcu:ss-top-setting}). For fixed
$x=(1,g_f)$ and $\beta$, the functional $\Lambda_{\beta,x}\colon f\mapsto\pcuF_\beta\mathcal Df(x)$ is linear
on $\mathcal M$ and satisfies $\Lambda_{\beta,x}(\Psi_w)=\psi_S(\operatorname{tr}(\tau_\beta w))
\Lambda_{\beta,x}(\Psi_0)$ by (E2) (the shear is the action of $n(0,w)$ at $S$, which multiplies the
coefficient by the character), where here and below $\psi_S(t)$ stands for $\psi_S(\jmath s_0\bar s_0t)$, a character of
$F_S$ trivial exactly on $\auO_S$ (as $s_0\widehat\auO_E=\alpha\widehat\auO_E$ and $\jmath\alpha\bar\alpha\widehat\auO_E$ is the dual of
$\widehat\auO_E$ for $\psi$, {\cite[\S5.1.1]{P49}}), and $\tau_\beta=-s_0^{-1}\bar s_0^{-1}\beta$ is, up to sign, the
normalised index (at $v=w\bar w$, $w\in S$, $Z_D(E^+_v)\cong\pcMat_b(E_w)$ through the $w$-component, and
$\psi_{\beta,v}(z)=\psi_w(-\jmath\operatorname{tr}(z_w\beta))$ by \eqref{pcu:eq-psibeta}; only whether $\det\tau_\beta$ is a
unit matters below). So: every $f\in\mathcal M$ killed
by all functionals $\Lambda$ with $\Lambda(\Psi_w)=\psi_S(\operatorname{tr}(\tau w))\Lambda(\Psi_0)$ for all
$w$ and $\Lambda(\Psi_0)=0$ when $\tau\in\auGL_b(\auO_S)$ (hypothesis) lies in $\mathcal Q$. Let $\pi\colon
\auO_S\to\pcR$ be reduction modulo the product of the primes of $S$ and $\eta$ the character of
$\pcR^\times$ with $\eta\circ\pi=\chi_S^{-1}$ on units, extended by $0$. Consider
$f=\sum_{r\in\pcMat_b(\pcR)}\eta(\det r)\Psi_{\varpi^{-1}\tilde r}$, $\tilde r$ a lift of $r$. For a
functional as above, $\Lambda(\Psi_w)=\Lambda(\Psi_0)$ for integral $w$ (as $\Psi_w=\Psi_0$ then), so
$\tau$ is integral if $\Lambda(\Psi_0)\ne0$; and
$\Lambda(f)=\Lambda(\Psi_0)\sum_r\eta(\det r)\psi_S(\varpi^{-1}\operatorname{tr}(\tau\tilde r))$, a
Gauss sum on $\pcMat_b(\pcR)$ which vanishes when $\det\tau$ is not a unit ($\eta$ primitive, as
$\chi_S$ is tame with non-trivial tame part at each $w\in S$); when $\det\tau$ is a unit,
$\Lambda(\Psi_0)=0$. So $f\in\mathcal Q$, hence $f(g\,\cdot)\in\mathcal Q$ for $g=-\varpi^{-1}1_N$. A
direct computation shows that every non-zero term of
$f(g\,\cdot)$ at $y$ equals $\Psi^{\mathrm{col}}_{z'_0}(y)$ for the complementary block $z'_0$ of the
inverse frame, so $f(g\,\cdot)=n(\cdot)\Psi^{\mathrm{col}}_{z'_0}$ with $n(y)\in\Z_{\ge0}$ the number of
non-vanishing terms; $n$ is locally constant, and $n>0$ at the point $x_0$ of
Lemma~\ref{pcx:generate} for $s=z'_0$ (the term $r=1$ is $1$ there). The proof of
Lemma~\ref{pcx:generate} goes through verbatim with $p=\Psi^{\mathrm{col}}_{s}$ replaced by $np$:
the claim becomes $\Psi(k^{-1}x_0)\,n(dky)p(dky)=n(dky)|p(dky)|\Psi(y)$, the average
$I(y)=\int n(dky)|p(dky)|\,dk$ is still constant on the support of $\Psi$ by right invariance of
$dk$, and $I(x_0)>0$ by continuity. Hence every locally constant $\Psi$ with a matched function lies
in $\mathcal Q$.
\end{proof}

\begin{lemma}[Unimodular Gram points suffice]\label{pcu:unimod}
Let $f=\mathcal D(\Psi^{\mathrm{col}}_{z_0})$ and $\beta$ with unimodular normalised index. If
$\pcuF_\beta f$ vanishes at every point $n(l,z)m(s_0q,g)\in P_D(\auA)$ whose Gram matrix $q^*\beta q$
is unimodular at all finite places, then $\pcuF_\beta f(1,g_f)=0$ for all $g_f$ trivial at $S$.
\end{lemma}
\noindent The proof is in \S\ref{pcu:ss-unimod}.

\begin{proof}[Proof of Proposition~\ref{pcx:P-top}]
Let $\Phi$, $\Lambda$ be the vectors of \S\ref{pcu:ss-period}, $x=\Phi$, $x'=\Lambda$; they have finite
moments and the zero term of $x'$ along $D^{\mathrm s}$ (the $b$ planes) is $w$. Let $K$, $C_\nu$ be as
in Lemma~\ref{pcu:coeff}. For each $\nu$, the difference $\mathcal D$ built with $C_\nu$ satisfies the
hypothesis of Proposition~\ref{pcu:selfdual}: by (E4) and the linearity of $\pcuF_\beta$ (both
integrals are automorphic of moderate growth), $\pcuF_\beta\mathcal D(\Psi^{\mathrm{col}}_{z_0})=
\pcuF_\beta\Theta^{\mathrm L}(\Phi,\Phi^{\mathrm{row}}_{z_0})-C_\nu\pcuF_\beta\Theta^{\mathrm S}(\Lambda,
\Psi^{\mathrm{col}}_{z_0})$, which vanishes at test pairs by Lemma~\ref{pcu:coeff}. Hence
$\Theta^{\mathrm L}_{X,T}(\Phi,\Phi^{\mathrm{row}}_{z'})=C_\nu\Theta^{\mathrm S}_{X,U}(\Lambda,
\Psi^{\mathrm{col}}_{z'})$ for every integral $z'$: this is (b) of $\mathrm{RC}$ with $c=C_\nu$. For (a):
$|c|^2=K^2|I_0(\nu)|^2=K^2L_*(\pcm,\nu,j)\,\pcxstr_{W_0}(\nu;z)$ by Definition~\ref{pcx:def-ratio}
($\operatorname{rk}W_0=2j$), and $L_*(\pcm,\nu,\operatorname{rk}T-\operatorname{rk}X/2)=L_*(\pcm,\nu,j)$;
so (a) holds with the constant $K^{-2}$.
\end{proof}

\begin{remark}[Hypotheses of Proposition~\ref{pcx:P-top}]\label{pcu:rem-top}
No sign condition is used: the slices of $\Phi$ vanish outside the real class of $U$ by construction
(Lemma~\ref{pcu:support}), which is what case (b) of Lemma~\ref{pcu:coeff} needs. The only hypotheses
are that $W_0$ has even rank and splits off an adapted plane (for Lemma~\ref{pcx:const}), and
$\mu_\tau\ge1$.
\end{remark}

\subsubsection{Unfolding along $P_D$: proofs of (U1) and (U2)}\label{pcu:ss-open}

\emph{Adelic slices.} For a datum $f$ of $\auHom(X,V)$ (all places) and $x\in\pcMat_{a\times b}(\auA_E)$
let $f[x]$ be the datum of $\auHom(W_0,V)$ obtained by restricting the mixed model of $D$ at every
place (Schr\"odinger in the $\auHom(\ell,V)$ variable at every place, as in (W2), (W2$'$) of
{\cite[\S5.1.1]{P49}}) to the fibre over $x$: $f[x]=(\psrho(x\,e^*)f)[0]$, $[0]$ the zero term. The
theta distribution of $\auHom(X,V)$ is $\theta(f)=\sum_{x\in\pcMat_{a\times b}(E)}\theta_{W_0}(f[x])$:
in the mixed model the rational points of $\auHom(X,V)$ decompose as $\auHom(\ell,V)(E)\oplus
\auHom(W_0,V)(E)\oplus\auHom(\ell',V)(E)$, the Schr\"odinger variable is summed over the first and
the remaining sum is $\theta_{W_0}$ by Poisson summation in $\auHom(\ell',V)$ (absolutely convergent
for Schwartz data). Put $\mathrm{term}_x(p,h)=\theta_{W_0}((\omega(p,h)f)[x])$.

\begin{lemma}[Levi and radical laws]\label{pcu:laws}
Let $\chi_V(u)=\alpha(u/\bar u)\upsilon(u)^{-a}$ for the source character $\alpha$. Then
\begin{align*}
(\omega_{\alpha,\beta\upsilon^b}(m(q,g),h)f)[x]&=\chi_V(\det q)^{-1}|\det q|^{a/2}\,
\omega_{\alpha,\beta}(g,h)\bigl(f[h^{-1}xq]\bigr),\\
(\omega(n(c,z),1)f)[x]&=\psi_{x^*J_Vx}(z)\,\psrho(x c^*J_{W_0})
\bigl(f[x]\bigr),\qquad\psi_{x^*J_Vx}(z)=\psi\bigl(-\tfrac12\jmath\operatorname{tr}(z\,x^*J_Vx)\bigr),
\end{align*}
with $\psi_\beta$ as in \eqref{pcu:eq-psibeta}.
\end{lemma}
\begin{proof}
These are the Schr\"odinger-model formulas (W2), (W2$'$) for the Siegel-type parabolic of the
polarisation $\auHom(\ell,V)\oplus\auHom(\ell',V)$, tensored with the action on $\auHom(W_0,V)$. For
$m(q,g)$: together with $h$ it acts on the Schr\"odinger variable $x\in\auHom(\ell,V)$ by
$x\mapsto hxq^{-1}$, so the new function at $x$ is the old one at $h^{-1}xq$, with the factor of (W2) for $\mathrm{GL}(\auHom(\ell,V))$, which is
$\upsilon$ of its determinant $\det(q)^{a}\det(h)^{-b}$ times $|\det q|^{a/2}$; the factor
$\upsilon(\det h)^{-b}$ is absorbed by the shift $\beta\upsilon^b\to\beta$ in the target character
and the factor $\alpha(\det q/\overline{\det q})$ comes from the $\auU(X)$-character restricted along
$\ell\oplus\ell'$, giving $\chi_V(\det q)^{-1}$; $g$ acts on $\auHom(W_0,V)$.

\emph{The radical.} Write $\mathsf W=\auHom(X,V)=I\oplus I^*\oplus\mathsf W'$ with $I^*=\auHom(\ell,V)$ (maps vanishing on
$\ell'\oplus W_0$), $I$ the maps vanishing on $\ell\oplus W_0$, $\mathsf W'=\auHom(W_0,V)$; the form of $\mathsf W$ is
$\langle w_1,w_2\rangle=\operatorname{tr}(w_1^\dagger w_2)$, and in the frames $w^\dagger=G_X^{-1}M^*J_V$ for a map with matrix $M$, $G_X$
the Gram matrix of the frame of $X$. The mixed model is $\auS(I^*;\psHS(\mathsf W'))$ with $\psrho(u_0)$ ($u_0\in I^*$) the
translation, $\psrho(i)$ ($i\in I$) the multiplication by $[u,i]$, and $\mathsf W'$ acting on the values, as in (P2) and
{\cite[Lemma~2.3]{P49}}. The element $(g,1)$ acts on $\mathsf W$ by $w\mapsto wg^{-1}$. Let $u\in I^*$ with $u(e)=x$. Since
$n(c,z)^{-1}=n(-c,z^*)$, $u\,n(c,z)^{-1}=u+d(u)+b(u)$ with $d(u)\in\mathsf W'$, $d(u)(w)=-x\,c^*J_0w$, and $b(u)\in I$,
$b(u)(f)=x\,z^*$; also $\mathsf W'\ni w'\mapsto w'n(c,z)^{-1}=w'+C(w')$ with $C(w')\in I$, and $I$ is fixed.

(i) \emph{The centre, $c=0$.} Then $d=0$, $z^*=-z$, and $n(0,z)$ acts as the element $n(b)$ of the Siegel parabolic of
$I\oplus I^*$, tensored with the identity on $\mathsf W'$ ({\cite[Lemma~2.5]{P49}}). By {\cite[Lemma~2.4]{P49}},
$\omega(n(b))\varphi(u)=\beta(u,bu)^{-1}\varphi(u)$, where $\beta(u,bu)=\psi(\frac12\jmath\langle bu,u\rangle)$. In the frames,
$G_X^{-1}$ exchanges the $e$- and $f$-blocks, the matrix of $u$ is $(x\ 0\ 0)$ and that of $bu$ is $(0\ {-xz}\ 0)$, so
$\langle bu,u\rangle=\operatorname{tr}(G_X^{-1}(bu)^*J_Vu)=\operatorname{tr}(-z^*x^*J_Vx)=\operatorname{tr}(z\,x^*J_Vx)$. Hence
$\omega(n(0,z),1)$ is multiplication by $\psi(-\frac12\jmath\operatorname{tr}(z\,x^*J_Vx))=\psi_{x^*J_Vx}(z)$.

(ii) \emph{The element $n(c,z_2)$, $z_2=-\frac12c^*J_0c$.} Here $z_2$ is hermitian. We claim that $A\colon F\mapsto
(u\mapsto\psrho(-d(u))F(u))=(u\mapsto\psrho(xc^*J_0)F(u))$ satisfies $A\psrho(y)A^{-1}=\psrho(n(c,z_2)y)$ for all $y\in\mathsf W$; we
use $\psrho(y)^{-1}=\psrho(-y)$ and $\beta(y,y')=\psi(\frac12\jmath\langle y',y\rangle)$, and that $\psi(\jmath r)=1$ for $r\in\auA_{E^+}$. For
$y=i\in I$ both sides are multiplication by $[u,i]$. For $y=w'\in\mathsf W'$, $A\psrho(w')A^{-1}F(u)=[-d(u),w']\psrho(w')F(u)$, while
$\psrho(w'+C(w'))=\psrho(C(w'))\psrho(w')$ ($\langle w',C(w')\rangle=0$) multiplies by $[u,C(w')]$ and then applies $\psrho(w')$; and
$[u,C(w')]=[d(u),w']^{-1}$, because $n(c,z_2)$ preserves $[\,,]$ and $u\perp w'$ (expand $[u+d(u)+b(u),w'+C(w')]=1$ using
$I\perp\mathsf W'$, $I$ isotropic). For $y=u_0\in I^*$, with $d_0=d(u_0)$, $b_0=b(u_0)$: on one side
$A\psrho(u_0)A^{-1}F(u)=\psrho(-d(u))\psrho(d(u)+d_0)F(u+u_0)=\beta(d(u),d_0)^{-1}\psrho(d_0)F(u+u_0)$; on the other,
$\psrho(u_0+d_0+b_0)F(u)=\beta(b_0,u_0)^{-1}[u,b_0]\psrho(d_0)F(u+u_0)$. Now $\beta(b_0,u_0)=\psi(\frac12\jmath\operatorname{tr}(x_0^*J_Vx_0z_2))=1$
($x_0=u_0(e)$, and the trace of a product of two hermitian matrices is real); with $P=x^*J_Vx_0$, $M=c^*J_0c$ one finds
$\langle u,b_0\rangle=\operatorname{tr}(Pz_2)=-\frac12\operatorname{tr}(PM)$, so $[u,b_0]=\psi(\frac12\jmath(\overline{\langle u,b_0\rangle}-\langle u,b_0\rangle))
=\psi(-\jmath\langle u,b_0\rangle)=\psi(\frac12\jmath\operatorname{tr}(PM))$, and $\langle d_0,d(u)\rangle=\operatorname{tr}(x_0^*J_Vx\,M)=
\operatorname{tr}(P^*M)=\overline{\operatorname{tr}(PM)}$, so $\beta(d(u),d_0)^{-1}=\psi(-\frac12\jmath\,\overline{\operatorname{tr}(PM)})=
\psi(\frac12\jmath\operatorname{tr}(PM))$. The two sides agree, which proves the claim. By (P1), $\omega(n(c,z_2),1)$ is a scalar
multiple of $A$; the scalar is $1$, because by (W2) (with $\chi_I=1$ on unipotent elements, and $n(c,z_2)$ acting trivially on
$I^\perp/I=\mathsf W'$) $\omega(n(c,z_2),1)$ commutes with evaluation at $u=0$, as $A$ does.

(iii) \emph{General $n(c,z)$.} $n(c,z)=n(0,z_1)\,n(c,z_2)$ with $z_1=z+\frac12c^*J_0c$ skew-hermitian, by the composition law. By
(i) and (ii), $\omega(n(c,z),1)F(u)=\psi_{x^*J_Vx}(z_1)\,\psrho(xc^*J_0)F(u)$, and $\psi_{x^*J_Vx}(z_1)=\psi_{x^*J_Vx}(z)$ since $\psi$ is
trivial on $\jmath\auA_{E^+}$ (\S\ref{pcu:ss-top-setting}). The characters $\alpha$, $\beta$ are trivial on unipotent elements. This is
the second formula.
\end{proof}

\begin{corollary}[Local form]\label{pcu:laws-local}
The same formulas hold at each place $v$ of $E^+$ separately: for a datum $f_v$ of $\auHom(X,V)_v$ in the mixed model of $D$
at $v$ and $x\in\pcMat_{a\times b}(E_v)$,
$(\omega_v(m(q,g),h)f_v)[x]=\chi_{V,v}(\det q)^{-1}|\det q|_v^{a/2}\,\omega_v(g,h)(f_v[h^{-1}xq])$ and
$(\omega_v(n(c,z),1)f_v)[x]=\psi_v(-\frac12\jmath\operatorname{tr}(z\,x^*J_Vx))\,\psrho_v(xc^*J_{W_0})(f_v[x])$, with the local components of the
characters. At an infinite place $\tau$, $f_\tau[x]=\mathrm{sl}_D(x)f_\tau$.
\end{corollary}
\begin{proof}
The proof of Lemma~\ref{pcu:laws} is place by place: (W2), (W2$'$) are identities between local operators, the adelic
mixed model, the slices $f[x]$ and the Weil representation are restricted tensor products of the local ones, and the
characters are products of their local components. Each step of that proof is therefore an identity at the place $v$,
and the adelic formulas are the tensor products of the local ones. At $\tau$ the mixed model is that of
\S\ref{pcu:ss-mixed} and the slice is $\mathrm{sl}_D$ by definition.
\end{proof}

\begin{lemma}[Surjectivity of the standard data]\label{pcu:rowsurj}
\textup{(a)} If $\Phi^{\mathrm{row}}_z(x)\ne0$ ($z\in\pcMat_{b\times N}(\auO_S)$, $x\in\pcMat_{a\times N}(F_S)$, $N=a+b$), then at
every place of $S$ the matrix $x$ has rank $a$, i.e.\ $x\colon F_w^N\to F_w^a$ is surjective. \textup{(b)} The same holds
for $\Psi^{\mathrm{col}}_{1_n}$ on $\pcMat_{n\times n}(F_S)$: it is supported on $\auGL_n(\auO_S)$. \textup{(c)} If
$(\Psi^{\mathrm{col}}_z)^{\mathsf t}(x)=\Psi^{\mathrm{col}}_z(x^{\mathsf t})\ne0$ ($x\in\pcMat_{b\times N}$), then $x$ has rank $b$ at
every place, i.e.\ is surjective. Hence, as $S\cap\bar S=\emptyset$ (beginning of \S\ref{pct:sec}),
Proposition~\ref{pcx:decay}(a) applies with $P=S$ to the data $[y,\Phi^{\mathrm{row}}_z]_S$ of $\auHom(X,T)$ and
$[y,\Psi^{\mathrm{col}}_{1_n}]_S$ of $\auHom(W_0,W_0)$, and with $P=\bar S$ to $[y,(\Psi^{\mathrm{col}}_z)^{\mathsf t}]_{\bar S}$ of
$\auHom(X,U)$.
\end{lemma}
\begin{proof}
(a) By Definition~\ref{pcx:def-std}, $\binom zx\in\auGL_N(\auO_S)$; the $a$ rows of $x$ are rows of an invertible matrix,
hence linearly independent over each $E_w$, $w\in S$. (b) By Definition~\ref{pcx:def-std} with $z=1_n$. (c) By
Definition~\ref{pcx:def-std}, $zx^{\mathsf t}\in\auGL_b(\auO_S)$, so $x^{\mathsf t}$, hence $x$, has rank $b$.
\end{proof}

\begin{proof}[Proof of (U2)]
Apply Lemma~\ref{pcu:laws} to $p=n(c,z)m(q,g)$: \[\mathrm{term}_x(p,h)=\psi_{x^*J_Vx}(z)\chi_V(\det q)^{-1}|\det q|^{a/2}\theta_{W_0}(\psrho(x'c^*J_{W_0})\omega(g,h)f[h^{-1}xq]).\]
With $f=[\phi,\Phi]$ and $x'q=s_0Y$, the slice $f[h^{-1}s_0Y]$ is the tensor product of the real
slice $\mathrm{sl}_D(h^{-1}Y)\phi$ (the factor $s_0$ is trivial at infinity), the slice of $\Phi$ at
the $S$-component $y$ of $h^{-1}Y$ (by definition of the Schr\"odinger model at $S$), and at the
other finite places the slice of the standard vacuum, which is $\varepsilon_V\mathbf 1_{\mathrm{int}}$
times the vacuum of $\auHom(W_0,V)$ (the vacuum of the self-dual lattice $\alpha\auHom(L_X,L_V)$ is the
indicator of the lattice, which factors along the unimodular splitting; $\varepsilon_V$ is the
product of the local volume normalisations, independent of $x$). Here the slice point is the finite part of
$h^{-1}s_0Y$, whose coordinate is $y=h^{-1}Y$ by Lemma~\ref{pgx:dict}(d); so its integrality, and the
slice of $\Phi$ at $S$, are read on $y$, as in (U2). This is (U2).
\end{proof}

\begin{proof}[Proof of (U1)]
Fix $p=n(c,z)m(q,g)$. Let $I\subset\auHom(X,V)$ be the isotropic subspace of the maps vanishing on
$\ell^\perp=\ell\oplus W_0$, so that $I^\perp$ is the space of maps vanishing on $\ell$,
$I^\perp/I=\auHom(W_0,V)$, and $w\mapsto w|_\ell$ identifies $\auHom(X,V)(E)/I^\perp(E)$ with
$\pcMat_{a\times b}(E)$. We claim $\theta_{W_0}(F[x])=c_x(F)$ in the notation of {\cite[Lemma~8.3]{P49}},
i.e.\ $F[x]=\psZ_I\psrho(w)F$ for any $w$ with $w|_\ell=x$. Let $I^*$ be the maps vanishing on
$\ell'\oplus W_0$; it is isotropic, in duality with $I$, and $\auHom(X,V)=(I\oplus I^*)\perp\auHom(W_0,V)$
with $I^*=\auHom(\ell,V)$. The mixed model of $D$ is the model $\psHS(I^*)\hat\otimes\psHS(\auHom(W_0,V))$
in which $\psrho(I^*)$ translates the variable $u\in I^*(\auA)$, $\psrho(I)$ multiplies by the characters
$u\mapsto\psi(\langle u,i\rangle)$, and $\psrho(\auHom(W_0,V))$ acts on the values. The map
$F\mapsto F[0]$ (evaluation at $u=0$) satisfies $(\psrho(i)F)[0]=F[0]$ for $i\in I(\auA)$ and
$(\psrho(w)F)[0]=\psrho(w)(F[0])$ for $w\in\auHom(W_0,V)(\auA)$, i.e.\ the intertwining properties of
$\psZ_I$ in (W2). These determine the map up to a scalar: a continuous map with them factors through the
quotient of $\psHS(I^*)$ by the span of the $(\psi(\langle\cdot,i\rangle)-1)\varphi$, whose dual consists of
the distributions $T$ with $(\psi(\langle\cdot,i\rangle)-1)T=0$ for all $i$; such $T$ is supported at $u=0$, and
at the archimedean places no derivative of $\delta_0$ survives (if $T=\sum_{|\alpha|\le m}c_\alpha
\partial^\alpha\delta_0$ with some $c_\alpha\ne0$, $|\alpha|=m\ge1$, multiplying by $\psi(\langle\cdot,i\rangle)-1$,
which vanishes to first order at $0$ with arbitrary linear part, gives a non-zero distribution of order
$m-1$), so $T$ is a multiple of $\delta_0$. Hence $F[0]=c\,\psZ_IF$, and $c=1$ because both satisfy the
expansion of $\theta$: $\theta(F)=\sum_{u\in I^*(E)}\theta_{W_0}((\psrho(u)F)[0])$ (Poisson summation in the
mixed model, as recalled before Lemma~\ref{pcu:laws}) and (W3) with $W(E)/I^\perp(E)=I^*(E)$, and
$\theta\ne0$. As $F[x]=(\psrho(xe^*)F)[0]$ with $xe^*\in I^*(\auA)$ restricting to $x$ on $\ell$, this is the
claim; the expansion $\theta=\sum_xc_x$ is (W3).

\emph{Step 1: absolute convergence and moderate growth of the absolute series.} For every finite datum
$f_0$ there is $Y$, depending on $f_0$ and on $p$, with
\begin{equation}\label{pcu:eq-absmod}
\sum_{x}\bigl|c_x\bigl(\omega(p,h)(\psrho_\infty(z')\Psi\otimes f_0)\bigr)\bigr|\le Y\,\mathrm{ht}(h)^{4s}\,
\|\Psi\|_{2s}\qquad(h\in\auU(V)(\auA),\ z'\in\auHom(X,V)_\infty,\ \Psi\in\pzF_{2s}),
\end{equation}
$s$ as in (R2) for $\auHom(X,V)$. Indeed, for $\gamma\in\auU(V)(E)$, $c_x\circ\omega(1,\gamma)=c_{\gamma^{-1}x}$
up to a factor of absolute value $1$ ((W1), (W2): $\gamma$ stabilises $I$, and the characters are trivial on
rational points), so the left side is unchanged when $h$ is replaced by $\gamma h$; by
Lemma~\ref{pcx:reduce} we may assume $h_f$ in a fixed compact set and $\|h_\infty\|\le C\,\mathrm{ht}(h)$.
Then $\omega(p,h)(\psrho_\infty(z')\Psi\otimes f_0)$ is a phase times $\psrho_\infty(z'')\omega_\infty
(p_\infty,h_\infty)\Psi\otimes\omega_f(p_f,h_f)f_0$ ((W1)), the finite vector lies up to phases in a finite
set, and {\cite[Lemma~8.3(c)]{P49}} followed by {\cite[Lemma~8.2(a)]{P49}} with index $2s$ (``Other
spaces'' paragraph before {\cite[Lemma~8.3]{P49}}) gives \eqref{pcu:eq-absmod} with
$Y=C'(2DC\|p_\infty\|)^{4s}$. In particular $\sum_x|\mathrm{term}_x(p,h)|<\infty$ for every $h$, and
$|\mathrm{term}_x(n(0,z)p,h)|=|\mathrm{term}_x(p,h)|$ by Lemma~\ref{pcu:laws}.

\emph{Step 2: the coefficient as an integral of the $\beta$-part.} By Proposition~\ref{pcx:decay}(a) (its
hypothesis holds for the standard data, Lemma~\ref{pcu:rowsurj}) $|\theta(\omega(n(0,z)p,h)F)|\le
C\,\mathrm{ht}(h)^{-N_0}$ uniformly for $z$ in a compact fundamental set of $Z_D(E)\backslash Z_D(\auA)$,
integrable on $[\auU(V)]$ by (R1); so Fubini applies to the double integral over $Z_D(E)\backslash
Z_D(\auA)\times[\auU(V)]$. For fixed $h$, Step~1 and Lemma~\ref{pcu:laws} with $c=0$, $q=1$
($\mathrm{term}_x(n(0,z)p,h)=\psi_{x^*J_Vx}(z)\mathrm{term}_x(p,h)$, $\psi_\beta$ as in \eqref{pcu:eq-psibeta}) allow integrating term by term against $\psi_\beta(z)^{-1}$ over the compact
group, which keeps exactly the $x$ with $x^*J_Vx=\beta$. Hence
\[
\pcuF_\beta\int_{[\auU(V)]}\theta(\omega(\cdot,h)F)\,dh\,(p)=\int_{[\auU(V)]}\theta_\beta(h)\,dh,\qquad
\theta_\beta(h)=\sum_{x^*J_Vx=\beta}\mathrm{term}_x(p,h),
\]
the integral on the right absolutely convergent ($|\theta_\beta(h)|\le\sup_z|\theta(\omega(n(0,z)p,h)F)|$).

\emph{Step 3: unfolding.} $\mathrm{term}_{\gamma x}(p,\gamma h)=\mathrm{term}_x(p,h)$ for $\gamma\in
\auU(V)(E)$. If there is no $x$ with $x^*J_Vx=\beta$, or all their terms vanish identically, then
$\theta_\beta\equiv0$ and the coefficient is $0$ by Step~2. Otherwise, $\beta$ being invertible, the $x$
with $x^*J_Vx=\beta$ form one $\auU(V)(E)$-orbit (Witt's theorem: an isometric embedding of the
nondegenerate space $(E^b,\beta)$ into $V$ is unique up to $\auU(V)(E)$). Suppose that
\begin{equation}\label{pcu:eq-unfoldabs}
\int_{\mathrm{Stab}(x_0)(E)\backslash\auU(V)(\auA)}|\mathrm{term}_{x_0}(p,h)|\,dh<\infty
\qquad\Bigl(\text{equivalently }\int_{[\auU(V)]}\sum_{x^*J_Vx=\beta}|\mathrm{term}_x(p,h)|\,dh<\infty\Bigr)
\end{equation}
for one (hence every) $x_0$ of the orbit; the two integrals are equal by Tonelli (non-negative terms,
$\sum_{x}|\mathrm{term}_x(p,h)|=\sum_{\gamma\in\mathrm{Stab}(x_0)(E)\backslash\auU(V)(E)}|\mathrm{term}_{x_0}(p,\gamma h)|$).
Then Fubini gives $\int_{[\auU(V)]}\theta_\beta=\int_{\mathrm{Stab}(x_0)(E)\backslash\auU(V)(\auA)}
\mathrm{term}_{x_0}(p,h)\,dh$, which is (U1).
\end{proof}

\begin{remark}[What (U1) does not assert]\label{pcu:rem-U1}
Condition \eqref{pcu:eq-unfoldabs} is not automatic. For the orbit of $x_0$ on the large side the
stabiliser is $\auU(W_0)$, and by Lemma~\ref{pcn:slicechar} the inner integral over $[\auU(W_0)]$ is a theta
integral for the pair $(W_0,W_0)$ of equal ranks $n$; $W_0\supset\pcxH$ is isotropic, so Weil's
criterion for absolute convergence of such integrals for arbitrary Schwartz--Bruhat data (which needs the
rank of the source to be smaller than the rank minus the Witt index of the target) fails. What makes
\eqref{pcu:eq-unfoldabs} true below is the finite datum $\Phi^{\mathrm{row}}_{z_0}$: its slices along $x_0$
contain the factor $\chi_S^{-1}\circ\det$ supported on invertible matrices, to which
Proposition~\ref{pcx:decay}(a) applies. Moderate growth of $\sum_x|\mathrm{term}_x|$ (Step~1) and decay of the
signed sum $\theta_\beta$ (Step~2) do not by themselves give \eqref{pcu:eq-unfoldabs}.
\end{remark}

\begin{lemma}[Polynomial volume of the orbits]\label{pcu:orbitvol}
Let $m\ge b$, $J$ a hermitian $m\times m$ matrix $\mathrm{diag}(\pm1)$ and $\beta_0$ an invertible hermitian
$b\times b$ matrix, at each infinite place $\tau$. Let $\mathrm{Fr}=\prod_\tau\pcMat_{m\times b}(\C)$ with
Lebesgue measure and euclidean norm $|\cdot|$, $Q(x)=(x_\tau^*Jx_\tau)_\tau$, and $\Omega=Q^{-1}(\beta_0)$.
Suppose $\Omega\ne\emptyset$. Then $\Omega$ is a closed submanifold and a single orbit of $G=\prod_\tau\auU(J)$
acting by left multiplication; the orbit map identifies $G/G_{x_1}$ ($x_1\in\Omega$) with $\Omega$, and for
every $K>\dim_\R\mathrm{Fr}$ the $G$-invariant measure $\mu$ on $\Omega$ satisfies
$\int_\Omega(1+|x|)^{-K}d\mu(x)<\infty$.
\end{lemma}
\begin{proof}
On the open set $\mathrm{Fr}^\circ$ where every $Q(x)_\tau$ is invertible, $Q$ is a submersion onto an open
subset of $\mathcal H=\prod_\tau\mathrm{Herm}_b(\C)$: $dQ_x(\xi)=\xi^*Jx+x^*J\xi$, and $\xi=\frac12x\,Q(x)^{-1}\eta$
gives $\eta$. So $\Omega$ is a closed submanifold, and for each $\beta\in Q(\mathrm{Fr}^\circ)$ the fibre carries
the Leray measure $\mu_\beta=dx/dQ$ (the quotient of the Lebesgue measures), with
$\int_{\mathrm{Fr}^\circ}f\,dx=\int_{\mathcal H}\int_{Q^{-1}(\beta)}f\,d\mu_\beta\,d\beta$ for measurable $f\ge0$.
Left multiplication by $h\in G$ preserves $Q$ and $dx$ ($|\det h_\tau|=1$), hence each $\mu_\beta$. By Witt's
theorem for hermitian forms over $\C$, $\Omega$ is one $G$-orbit; a transitive smooth action of a Lie group
with countably many components identifies $G/G_{x_1}$ with $\Omega$, and $G$, $G_{x_1}$ (a product of
unitary groups) are unimodular, so the invariant measure on $\Omega$ is unique up to a scalar and $\mu$ is a
multiple of $\mu_{\beta_0}$. For $c\in\prod_\tau\auGL_b(\C)$, right multiplication $x\mapsto xc$ maps
$Q^{-1}(\beta)$ to $Q^{-1}(c^*\beta c)$ and multiplies $dx$ and $d\beta$ by positive constants depending
continuously on $c$, so $\int_{Q^{-1}(c^*\beta_0c)}f\,d\mu_{c^*\beta_0c}=J(c)\int_\Omega f(xc)\,d\mu_{\beta_0}(x)$ with
$J$ continuous and positive. The map $c\mapsto c^*\beta_0c$ is a submersion at $c=1$ (derivative
$\eta\mapsto\eta^*\beta_0+\beta_0\eta$, onto with $\eta=\frac12\beta_0^{-1}\zeta$), so there are an open
neighbourhood $N$ of $\beta_0$ and a continuous $c\colon N\to\prod\auGL_b(\C)$ with $c(\beta)^*\beta_0c(\beta)=\beta$
and $\|c(\beta)\|\le2$. Put $I=\int_\Omega(1+|x|)^{-K}d\mu_{\beta_0}$. As $1+|xc|\le2(1+|x|)$ for $\|c\|\le2$,
\[
\infty>\int_{\mathrm{Fr}}(1+|x|)^{-K}dx\ge\int_N J(c(\beta))\int_\Omega(1+|xc(\beta)|)^{-K}d\mu_{\beta_0}(x)\,d\beta
\ge2^{-K}I\int_NJ(c(\beta))\,d\beta,
\]
and the last integral is positive; so $I<\infty$.
\end{proof}

\begin{lemma}[Short representatives]\label{pcu:section}
In the setting of \S\ref{pcu:ss-top-setting} ($J=\mathrm{diag}(t'\oplus s)$, $x_0=\binom{1_b}0$, $J_U=\mathrm{diag}(t')$),
for every $x\in\mathrm{Fr}$ with $x_\tau^*Jx_\tau=J_U$ for all $\tau$ there is $h\in\auU(T)_\R$ with
$h^{-1}x_0=x$ and $\|h\|\le(1+\sqrt n)(1+|x|)$ (operator norm, maximum over $\tau$).
\end{lemma}
\begin{proof}
Fix $\tau$ and drop it. Let $P=xJ_Ux^*J$ and $Q=1-P$; then $x^*JQ=x^*J-J_U^2x^*J=0$, so $Q$ maps into
$x^\perp=\ker(x^*J)$, $Q$ is the identity on $x^\perp$, and $\|Q\|\le1+|x|^2$. Let $z_1$ be an $(n+b)\times n$
matrix whose columns are a euclidean orthonormal basis of $x^\perp$ and $G=z_1^*Jz_1$. For $w\in x^\perp$
put $w'=QJw\in x^\perp$; since $Q^*=1-JxJ_Ux^*$ and $x^*Jw=0$, $w'^*Jw=w^*JQ^*Jw=w^*w$, and $|w'|\le\|Q\||w|$.
So for every $\xi\in\C^n$ there is $\xi'$ with $|\xi'^*G\xi|\ge|\xi||\xi'|/\|Q\|$, i.e.\ $\|G^{-1}\|\le\|Q\|$. Write
$G=u^*Du$ with $u$ unitary and $D$ real diagonal; $T=\mathrm{span}(x)\perp x^\perp$ and Sylvester's law give
$x^\perp$ the signs of $W_0$, so after permuting, $z=z_1u^*|D|^{-1/2}$ satisfies $z^*Jz=\mathrm{diag}(s)$ and
$|z|\le\sqrt n\,\|Q\|^{1/2}\le\sqrt n(1+|x|)$. The matrix $k=(x\ z)$ has $k^*Jk=J$, so $k\in\auU(J)$, $kx_0=x$;
put $h=k^{-1}=Jk^*J$, so $\|h\|=\|k\|\le|x|+|z|\le(1+\sqrt n)(1+|x|)$.
\end{proof}

\begin{lemma}[Moments under the Weil representation]\label{pcu:moments}
Let $\mathcal F$ be the Fock space of a real hermitian space $\C^m$, realised as in
Lemma~\ref{pcx:superpos}, and let $k$ be a unitary operator on $\mathcal F$ with $k\psrho(w)k^{-1}=
\psrho(\sigma w)$ for a real symplectic map $\sigma$ of $\C^m$ (for instance $\omega_\infty(g,h)$, by (W1)).
For every $N\ge0$ there is $C_N$, depending only on $N$, $m$ and the normalisations, with
$\|kv\|_N\le C_N\|\sigma\|^{2N+2m}\|v\|_N$ for every $v$ with $\|v\|_N<\infty$.
\end{lemma}
\begin{proof}
By Stone--von Neumann a unitary operator with this intertwining property is determined by $\sigma$ up to a
scalar of absolute value $1$; write $M(\sigma)$ for any choice. Let $E=k^{-1}e$ and $E(w)=\langle\psrho(w)e,E\rangle$
its function in the realisation of Lemma~\ref{pcx:superpos}. \emph{Reduction.} We first show $|E(w)|\le
\exp(-|w|^2/(1+\|\sigma\|^2))$. Write the
polar decomposition $\sigma^{-1}=\kappa P$ with $P=(\sigma^{-\mathsf t}\sigma^{-1})^{1/2}$ positive symmetric and
$\kappa$ orthogonal; both are symplectic, so $\kappa$ commutes with the complex structure. The vacuum is, up
to scalars, the only vector killed by the annihilation operators of the complex structure, a family that
$M(\kappa)$ permutes; so $M(\kappa)e\in\C e$, $M(\kappa)^{-1}\psrho(w)e\in\C\,\psrho(\kappa^{-1}w)e$, and
$|E(w)|=|\langle\psrho(\kappa^{-1}w)e,M(P)e\rangle|$, with $|\kappa^{-1}w|=|w|$. $P=\exp X$ with $X$ symmetric and
anticommuting with the complex structure $i$, so $i$ exchanges the eigenspaces of $X$ for $\mu$ and $-\mu$ and
there is a euclidean orthonormal basis $(e_j,ie_j)$ of $\C^m$ with $Pe_j=\lambda_je_j$, $P(ie_j)=\lambda_j^{-1}ie_j$,
$\lambda_j\ge1$, $\lambda_j\le\|P\|=\|\sigma^{-1}\|=\|\sigma\|$ (for a symplectic $\sigma$, $\sigma^{-1}=
\mathsf J^{-1}\sigma^{\mathsf t}\mathsf J$ with $\mathsf J$ the orthogonal matrix of the symplectic form in a basis $(e_j,ie_j)$). A unitary change of basis is again of the form $M(\kappa')$ and
does not change the moduli, so we may assume this basis is the standard one. Then $\mathcal F$ is the tensor
product of the Fock spaces of the lines $\C e_j$, $e$ and $M(P)$ factor (by the same uniqueness), and in the
Schr\"odinger model of each line, with $\psrho(p,q)f(t)=e^{2\pi iqt}f(t-p)$ up to a phase, the vacuum is the Gaussian
$\gamma(t)=2^{1/4}e^{-\pi t^2}$ and $M(P_j)f(t)=\lambda_j^{-1/2}f(t/\lambda_j)$. \emph{The normalisation.} The commutation
relation $\psrho(p,0)\psrho(0,q)=e^{-2\pi ipq}\psrho(0,q)\psrho(p,0)$ of this model, compared with $\psrho(w)\psrho(y)=
e^{\kappa(w,y)}\psrho(y)\psrho(w)$, $\kappa(w,y)=2i\Iml\operatorname{tr}(y^*w)$ (Lemma~\ref{pcx:superpos}), forces the
coordinate of $(p,q)$ on the line $\C e_j$ to be $\sqrt\pi(p+iq)$, so $|w|^2=\pi(p^2+q^2)$. For a normalised Gaussian
$\gamma_a(t)=(2a)^{1/4}e^{-\pi at^2}$ ($a>0$), completing the square gives $|\langle\psrho(p,q)\gamma_a,\gamma_a\rangle|=
\exp(-\frac\pi2(ap^2+q^2/a))$, while in the model of Lemma~\ref{pcx:superpos} $|\langle\psrho(w)e,e\rangle|=e(w)=
e^{-|w|^2/2}=\exp(-\frac\pi2(p^2+q^2))$; the unitary isomorphism carrying $e$ to $\gamma_a$ (up to a phase) preserves these
moduli, so $a=1$. Then a Gaussian integral gives
\[
|\langle\psrho(p,q)\gamma,M(P_j)\gamma\rangle|=\Bigl(\frac{2\lambda_j}{1+\lambda_j^2}\Bigr)^{1/2}
\exp\Bigl(-\pi\,\frac{p^2+\lambda_j^2q^2}{1+\lambda_j^2}\Bigr)\le\exp\Bigl(-\pi\,\frac{p^2+q^2}{1+\lambda_j^2}\Bigr)
=\exp\Bigl(-\frac{|w_j|^2}{1+\lambda_j^2}\Bigr)
\]
($w_j$ the component of $w$ on $\C e_j$; the exponent is obtained by writing $(t-p)^2+t^2/\lambda_j^2=(1+\lambda_j^{-2})(t-t_0)^2+
p^2/(1+\lambda_j^2)$, and the prefactor is $\le1$). Taking the product over $j$, with $\lambda_j\le\|\sigma\|$, gives the bound on
$|E(w)|$. Hence $\|E\|_N\le\int(1+|w|)^Ne^{-|w|^2/(1+\|\sigma\|^2)}dw\le C'_N\|\sigma\|^{N+2m}$ ($\|\sigma\|\ge1$).
\emph{The estimate.} As $\|E\|_0<\infty$, $E=\int E(w)\psrho(w)e\,dw$ in $\mathcal F$ (Lemma~\ref{pcx:superpos}(b)).
Then $(kv)(x)=\langle\psrho(x)e,kv\rangle=\langle k^{-1}\psrho(x)e,v\rangle=\langle\psrho(\sigma^{-1}x)E,v\rangle$, and
$\psrho(\sigma^{-1}x)\psrho(w)$ is a scalar of absolute value $1$ times $\psrho(\sigma^{-1}x+w)$; so
$|(kv)(x)|\le\int|E(w)||v(\sigma^{-1}x+w)|dw$. Substituting $u=\sigma^{-1}x+w$ ($|\det\sigma|=1$) and using
$1+|\sigma(u-w)|\le\|\sigma\|(1+|u|)(1+|w|)$,
$\|kv\|_N\le\|\sigma\|^N\|E\|_N\|v\|_N\le C'_N\|\sigma\|^{2N+2m}\|v\|_N$.
\end{proof}

\begin{lemma}[Integrability at test pairs of type (a)]\label{pcu:integrable}
Let $(J_U,q)$ be a test pair of type (a) for $U$ and $p=n(c,z)m(q,g)$.
\begin{enumerate}
\item[\textup{(a)}] Let $V=T$ and $F=[\phi,\Phi^{\mathrm{row}}_{z_0}]_S$ with $\phi$ a Fock vector of $\auHom(X,T)_\R$
with finite moments. Then $\int_{\Gamma_2\backslash G}|\mathrm{term}_{x_0}(p,h)|\,dh<\infty$ ($G=\auU(T)(\auA)$,
$\Gamma_2=\auU(W_0)(E)$, the stabiliser of $x_0$).
\item[\textup{(b)}] Let $V=U$ and $F=[\lambda,(\Psi^{\mathrm{col}}_{z_0})^{\mathsf t}]_{\bar S}$ with $\lambda$ a Fock
vector of $\auHom(X,U)_\R$ with finite moments. Then $\int_{\auU(U)(\auA)}|\mathrm{term}_1(p,y)|\,dy<\infty$.
\end{enumerate}
In particular \eqref{pcu:eq-unfoldabs} holds for $\beta=J_U$ in Lemma~\ref{pcu:coeff} (with $\phi=\Phi$,
$\lambda=\Lambda$ of \S\ref{pcu:ss-period}) and in type (a) of the proof of Proposition~\ref{pct:vantop}
(with $\phi=v$).
\end{lemma}
\begin{proof}
Write $q=s_0x_UY$ as in the proof of Lemma~\ref{pcu:coeff}, so that, in the notation of (U2), $Y$ is the
$b\times b$ block and $h^{-1}(Y;0)=h^{-1}x_0Y$. Constants $C$ may depend on $p$, $\phi$, $\lambda$.

(a) \emph{Integration in stages.} With $H_2=\auU(W_0)(\auA)$, Tonelli and the quotient integral formula
give $\int_{\Gamma_2\backslash G}|\mathrm{term}_{x_0}(p,h)|dh=\int_{H_2\backslash G}j(h)\,dh$ with
$j(h)=\int_{[\auU(W_0)]}|\mathrm{term}_{x_0}(p,s'h)|ds'$, a measurable function of $H_2h$ with values in
$[0,\infty]$. By Lemma~\ref{pcu:finslice}(a), (b), $|\mathrm{term}_{x_0}(p,\cdot)|$ is
right invariant under $K_f$, and $\mathrm{term}_{x_0}(p,\cdot)$ vanishes at $h=(h_\infty,u_f)$ unless $u_f\in H_fK_f$
($H=\auU(U)(\auA)\times H_2$). If $u_f=(u_1,u_2)\kappa$ with $\kappa\in K_f$, left invariance of $j$ under $H_2$
and right invariance under $K_f$ give $j(h)=j(h_0)$ with $h_0=(k_\infty,(y_f,1))$, $y_f=u_1$, for any
$k_\infty\in H_{2,\infty}h_\infty$. In
$H_2\backslash G=(H_{2,\infty}\backslash G_\infty)\times(H_{2,f}\backslash G_f)$ the archimedean coset of $h_0$
is determined by $x=k_\infty^{-1}x_0\in\Omega=\{x:x_\tau^*Jx_\tau=J_U\}$, and by Lemma~\ref{pcu:section} we may
take $\|k_\infty\|\le C(1+|x|)$.

\emph{The term at $s'h_0$.} Let $s'\in H_2$. The finite part of $(s'h_0)^{-1}x_0Y$ is $x_0y_f^{-1}Y_f$, and
its archimedean part is $xY_\infty$ ($s'$ fixes $x_0$). By (U2) and Lemma~\ref{pcu:finslice}(c) (with $h_f=(y_f,s'_f)$; the finite slice does not depend on $s'_f$),
$\mathrm{term}_{x_0}(p,s'h_0)=0$ unless $x_f=y_f^{-1}Y_f$ is integral, and then
\[
|\mathrm{term}_{x_0}(p,s'h_0)|=C_0\,\bigl|\theta\bigl(\psrho(t)\,\omega(g,s'h_0)(\mathrm{sl}_D(xY_\infty)\phi\otimes
f_U\otimes F_1)\bigr)\bigr|,
\]
with $C_0$ independent of $s'$, $x$, $y_f$, $t=x'c^*J_{W_0}$ fixed, $f_U$ the vacuum datum of $(W_0,U)$ and $F_1$
the datum $\Psi^{\mathrm{col}}_1$ of $(W_0,W_0)$ at $S$ (vacuum elsewhere). The set $\mathcal Y$ of $y_f\in\auU(U)(\auA_f)$
with $y_f^{-1}Y_f$ integral is compact ($y_f^{-1}\in\pcMat_b(\widehat\auO_E)Y_f^{-1}$ is bounded, and unitary
groups are closed). Put $A_x=\omega_\infty(1,k_\infty)\,\mathrm{sl}_D(xY_\infty)\phi$, a vector of
$\auHom(W_0,T)_\R$. By Lemma~\ref{pcu:moments} and Lemma~\ref{pcu:slicedecay}, for every $M'$ and $K$,
\[
\|A_x\|_{M'}\le C\|k_\infty\|^{2M'+2m}\|\mathrm{sl}_D(xY_\infty)\phi\|_{M'}\le C(1+|x|)^{2M'+2m}(1+|xY_\infty|)^{-K}
\|\phi\|_{M'+K}\le C'(1+|x|)^{2M'+2m-K}
\]
($m$ the complex dimension of $\auHom(W_0,T)_\R$; $|xY_\infty|\ge|x|/\|Y_\infty^{-1}\|$). Since $s'$ and
$(1,(y_f,1))$ are block diagonal for $T=U\oplus W_0$, $\omega(g,s'h_0)(\mathrm{sl}_D(xY_\infty)\phi\otimes f_U\otimes
F_1)=\omega(g,s')\bigl(A_x\otimes\omega_f(1,y_f)f_U\otimes F_1\bigr)$. Expand $A_x=\int A_x(w)\psrho(w_1)e\boxtimes
\psrho(w_2)e\,dw$ along $\auHom(W_0,T)=\auHom(W_0,U)\oplus\auHom(W_0,W_0)$ (Lemma~\ref{pcx:superpos}(b); the
integral converges in the space of data of class $s$, on which $\theta\circ\psrho(t)\circ\omega(g,s')$ is
continuous) and apply Lemma~\ref{pcu:product} to each coherent state:
\[
|\mathrm{term}_{x_0}(p,s'h_0)|\le C_0\int|A_x(w)|\,\bigl|\theta\bigl(\psrho(t)\omega(g,1)(\psrho(w_1)e\otimes
\omega_f(1,y_f)f_U)\bigr)\bigr|\,\bigl|\theta^{W_0,W_0}(g,s';\psrho(w_2)e\otimes F_1)\bigr|\,dw .
\]
The first theta factor is bounded uniformly in $w_1$ and $y_f\in\mathcal Y$: write $\psrho(t)\omega(g,1)=
\omega(g,1)\psrho(g^{-1}t)$; the finite data $\psrho_f((g^{-1}t)_f)\omega_f(1,y_f)f_U$ lie, up to factors of absolute
value $1$, in a finite set (finite data are fixed by an open subgroup, and $\mathcal Y$ is compact), and
Lemma~\ref{pcx:thetagrowth} is uniform in the archimedean translation. For the second, $F_1$ is supported on
$\auGL_n(\auO_S)$, i.e.\ on surjective matrices (Lemma~\ref{pcu:rowsurj}(b)), so Proposition~\ref{pcx:decay}(a) gives, for $N_0$ with
$\int_{[\auU(W_0)]}\mathrm{ht}^{-N_0}<\infty$ (R1), constants $B,M$ with $|\theta^{W_0,W_0}(g,s';\psrho(w_2)e\otimes
F_1)|\le C\|\psrho(w_2)e\|_M\mathrm{ht}(s')^{-N_0}\mathrm{ht}(g)^B\le C(1+|w_2|)^M\mathrm{ht}(s')^{-N_0}$. So, with
$M'=\max(M,2s)$,
\[
j(h_0)\le C\,\mathbf 1_{\mathcal Y}(y_f)\,\|A_x\|_{M'}\int_{[\auU(W_0)]}\mathrm{ht}(s')^{-N_0}ds'\le C''\,
\mathbf 1_{\mathcal Y}(y_f)\,(1+|x|)^{2M'+2m-K}.
\]

\emph{Conclusion.} The finite cosets with $j\ne0$ lie in the image of the compact set
$\{(y_f,1):y_f\in\mathcal Y\}K_f$, of finite measure. On the archimedean factor, $H_{2,\infty}\backslash G_\infty$ is identified with $\Omega$ by
$h\mapsto h^{-1}x_0$ (Lemma~\ref{pcu:orbitvol} with $m=b+n$, $\beta_0=J_U$), the quotient measure being a
multiple of $\mu$; choosing $K>2M'+2m+\dim_\R\mathrm{Fr}$, Lemma~\ref{pcu:orbitvol} gives
$\int_\Omega(1+|x|)^{2M'+2m-K}d\mu<\infty$. Hence $\int_{H_2\backslash G}j<\infty$.

(b) The stabiliser of $1$ is trivial. By (U2) (factor $\mathbf 1_{\mathrm{int}}$), $\mathrm{term}_1(p,y)=0$ unless $x_f=y_f^{-1}Y_f$ is
integral, i.e.\ $y_f\in\mathcal Y$, and then
$|\mathrm{term}_1(p,y)|=C_0|\theta(\psrho(t)\omega(g,y)(\mathrm{sl}_D(y_\infty^{-1}Y_\infty)\lambda\otimes f^{(x_f)}))|$, where
$f^{(x_f)}$ is the finite datum of (U2); as the datum at $\bar S$ is locally constant with compact support and $x_f$
runs over a compact set, $f^{(x_f)}$ takes only finitely many values. Writing $\psrho(t)\omega(g,y)=\omega(g,y)\psrho((g,y)^{-1}t)$ and using
Lemma~\ref{pcx:thetagrowth} (uniform in the archimedean translation; the finite data form a finite set up to
phases since $y_f\in\mathcal Y$), $|\mathrm{term}_1(p,y)|\le C\|\mathrm{sl}_D(y_\infty^{-1}Y_\infty)\lambda\|_s
\,\mathrm{ht}(y)^{A}$. Here $\|\cdot\|_s\le c_1\|\cdot\|_{2s}$ (Lemma~\ref{pcx:superpos}(b)) and
Lemma~\ref{pcu:slicedecay} give $\|\mathrm{sl}_D(y_\infty^{-1}Y_\infty)\lambda\|_s\le C(1+|y_\infty^{-1}|)^{-K}
\|\lambda\|_{2s+K}$; and $\mathrm{ht}(y)\le C(1+\|y_\infty\|)$ for $y_f\in\mathcal Y$ (the lattice at $y$ lies in a
fixed lattice $N^{-1}\auO_E^b$, whose non-zero vectors have archimedean length $\ge c>0$, and
$|(y^{-1}\gamma)_\infty|\ge|\gamma_\infty|/\|y_\infty\|$), with $\|y_\infty\|=\|y_\infty^{-1}\|\le|y_\infty^{-1}|$
($y_\infty^{-1}=J_Uy_\infty^*J_U$). So $\int_{\auU(U)(\auA)}|\mathrm{term}_1|\le C\,\mathrm{vol}(\mathcal Y)
\int_{\auU(U)_\R}(1+|y|)^{A-K}dy$, finite for $K$ large by Lemma~\ref{pcu:orbitvol} with $m=b$,
$\beta_0=J_U$, $\Omega=\auU(U)_\R$ (inversion preserves Haar measure).
\end{proof}

\subsubsection{Singular forms are governed by unimodular coefficients: proof of
Lemma~\ref{pcu:unimod}}\label{pcu:ss-unimod}

This is an adelic unitary form of a theorem of Freitag (\emph{Singular modular forms and theta
relations}, Lecture Notes in Math.\ 1487, Springer 1991; see also J.-S.~Li, Amer.\ J.\ Math.\ 119
(1997)); we give the proof. Throughout, $\Sigma$ is a set of finite places, $\Sigma^c=\Sigma\cup\bar
\Sigma$, $\mathcal O_\Sigma$ the adeles integral at $\Sigma^c$, $K_\Sigma$ the integral points of
$\auU(X)$ supported on $\Sigma^c$, $\delta=s_0\bar s_0$ (a finite idele of $E^+$), and for $\gamma\in\pcMat_b(\auA_E)$
\emph{Siegel's bracket} $\beta[\gamma]=\delta^{-1}\gamma^*\beta\gamma$. At a finite place $v$ the form $\psi$ of
{\cite[\S5.1.1]{P49}} is $\psi_\Q\circ\operatorname{Tr}_{E/\Q}$, and $\jmath\alpha\bar\alpha\auO_{E,v}$ is the dual $\mathfrak d_{E,v}^{-1}$ of
$\auO_{E,v}$ for it (this is the choice of $\alpha$ in {\cite[\S5.1.1]{P49}}: $\widehat\auO_E$ is its own dual for
$(z,u)\mapsto\psi(\jmath\alpha\bar\alpha\bar zu)$; it holds at every finite place, including those above $2$, $d$
and $q_1$, so no further fractional ideal enters; note that on $\auA_{E^+}$ one has
$\psi(t)=\psi_\Q(\operatorname{Tr}_{E^+/\Q}(2t))$, which cancels the factor $\frac12$ in \eqref{pcu:eq-psibeta}; and
$s_0\bar s_0/\alpha\bar\alpha$ is a unit idele, so $\delta\widehat\auO_{E^+}=\alpha\bar\alpha\widehat\auO_{E^+}$); as $E/E^+$ is unramified at $v$, there is a unit $\theta_v\in\auO_{E,v}^\times$ with
$\bar\theta_v=-\theta_v$ ($\theta_v=\zeta-\bar\zeta$ for a unit $\zeta$ generating the residue extension if $v$ is inert,
$\theta_v=(1,-1)$ if $v$ splits), so the integral skew-hermitian matrices at $v$ are $\theta_v$ times the integral
hermitian ones. A point $g$ \emph{acts on the
frame of $D$ by $\gamma$} if $g\,e=e\,\gamma$ ($e$ the frame of $\ell$, adelically).

\emph{Setting.} $D_1$ gives an adapted splitting $X=(\langle u\rangle\oplus\langle u'\rangle)\perp X'$
(one plane, the splitting $A$), and $X'=(\ell\oplus\ell')\perp W_2$ (the splitting $B$) with
$\ell$ the frame of $D$; in Lemma~\ref{pcu:unimod} this is $X=\pcxH\oplus(\pcxH^b\oplus W)$ with
the plane of $D_1$ (by (M4) and Lemma~\ref{pcx:adapted} the frames can be chosen so). For a row
$a\in\pcMat_{1\times b}(\auA_E)$ let $n_+(a)\in\auU(X)(\auA)$ be the unipotent element $u\mapsto u-e\,a^*$
(with the forced action on $u'$ and $\ell'$, identity on $W_2$), and $n_-(a)$ the same with $u,u'$
exchanged. For $f$ automorphic and $\beta$ with $\det\beta\ne0$ put
\[
f^+(g)=\int_{\pcMat_{1\times b}(E)\backslash\pcMat_{1\times b}(\auA_E)}\pcuF_\beta f(n_-(b)g)\,db,\qquad
\varphi^+_g(a)=f^+(n_+(a)g),
\]
and $f^-,\varphi^-_g$ the same with $u,u'$ exchanged.

\begin{lemma}[Partial means]\label{pcu:pm}
Let $f$ be automorphic and $\det\beta\ne0$, and put $\chi_\beta(a,b)=\psi_\beta(b^*a-a^*b)=\psi(-\jmath\,a\beta b^*)$.
\textup{(a)} $n_+(a)n_-(b)=n(0,b^*a-a^*b)\,n_-(b)n_+(a)$, so
$\pcuF_\beta f(n_-(b)n_+(a)g)=\chi_\beta(a,b)^{-1}\,\pcuF_\beta f(n_+(a)n_-(b)g)$. \textup{(b)} $f^+$ is continuous and
$f^+(n_-(b)g)=f^+(g)$; $\varphi^+_{n_+(a)g}(x)=\varphi^+_g(x+a)$; $\varphi^-_{n_+(a)g}(b)=\chi_\beta(a,b)^{-1}
\varphi^-_g(b)$. \textup{(c)} If $\varphi^+_{g'}=c\,\varphi^+_g$ then $\pcuF_\beta f(n_-(b)n_+(a)g')=c\,
\pcuF_\beta f(n_-(b)n_+(a)g)$ for all $a,b$, and $\varphi^-_{g'}=c\,\varphi^-_g$. In particular
$\varphi^+_g=0$ implies $\pcuF_\beta f(g)=0$ and $\varphi^-_g=0$.
\end{lemma}
\begin{proof}
(a) The isometry conditions force $n_+(a)$: $u\mapsto u-ea^*$, $f_k\mapsto f_k+a_ku'$, and the
identity on $u'$, $e$, $W_2$; and $n_-(b)$: $u'\mapsto u'-eb^*$, $f_k\mapsto f_k+b_ku$, the identity on $u$, $e$, $W_2$.
Both products agree on $u,u',e,W_2$, and $n_+(a)n_-(b)f_k-n_-(b)n_+(a)f_k=\sum_ie_i(\bar b_ia_k-\bar a_ib_k)$, so
they differ by the element $n(0,z)$, $z=b^*a-a^*b$ (skew-hermitian), of the centre $Z_D$; it commutes with $n_\pm$, and
$\pcuF_\beta f(n(0,z)p)=\psi_\beta(z)\pcuF_\beta f(p)$ (substitute $z'+z$ in \eqref{pcu:eq-psibeta}). For the
formula for $\chi_\beta$: $\operatorname{tr}(z\beta)=s-\bar s$ with $s=a\beta b^*$, $\bar\jmath=-\jmath$, so
$-\frac12\jmath(s-\bar s)=-\frac12(\jmath s+\overline{\jmath s})$, and $\psi(\bar t)=\psi(t)$. (b) $b\mapsto\pcuF_\beta f(n_-(b)g)$ is continuous
and invariant under rational $b$ ($n_-(b)$ rational fixes the frame of $D$, so is in $P_D(E)$ with
trivial Levi part on $\ell$ and $\pcuF_\beta f$ is left invariant under it); the mean over the compact
quotient is translation invariant. The second identity is the definition; the third follows from
(a) and the first. (c) For rational $a_0$, by (a) and the left invariance under $n_+(a_0)$,
$\varphi^+_g(a_0+a)$ is the Fourier coefficient at the character $b\mapsto\chi_\beta(a_0,b)$ of the
periodic function $b\mapsto\pcuF_\beta f(n_-(b)n_+(a)g)$. As $\beta$ is invertible, $a_0\beta$ runs through
$\pcMat_{1\times b}(E)$, so $a_0\mapsto\chi_\beta(a_0,\cdot)=\psi(-\jmath\,a_0\beta\,\overline{\cdot}^{\mathsf t})$ runs through all characters of $\pcMat_{1\times b}(E)\backslash
\pcMat_{1\times b}(\auA_E)$, so a continuous periodic function is determined by these coefficients
(Fourier uniqueness on the compact group). This gives the first claim; exchanging the order of the
shifts by (a) (the factor does not depend on the point) and averaging over $a$ gives the second.
\end{proof}

\begin{lemma}[Singularity]\label{pcu:sing}
Let $\mathcal U$ be a non-empty set of finite places. If $f^+(\varsigma g)=f^+(g)$ for every
$\varsigma\in N_A(E_{\mathcal U})$ (the shears $\varsigma=1+uc$, $cu=0$, at the places of $\mathcal U$) and every $g$ fixing the frame of $D$ off
$\Sigma^c$, then $f^+(rg)=f^+(g)$ for every $r$ in the radical $N_A(\auA)$ of the stabiliser of $u$
whose components off $\Sigma^c$ are trivial, for such $g$.
\end{lemma}
\begin{proof}
Call $g$ \emph{admissible} if it fixes the frame of $D$ off $\Sigma^c$. $N_A$ is a two-step unipotent
group: its centre is $Z_A\cong\{z\in\auA_E:z+\bar z=0\}$ (the elements $u'\mapsto u'+zu$) and
$N_A/Z_A\cong X'$ (the elements $n_A(c)$, $c\in X'$). The function $f^+$ itself need not be left invariant
under $N_A(E)$; we pass to its mean over $Z_A$.

\emph{The group $Z'$.} Let $I_1=\langle u\rangle\oplus\ell$ and let $Z'$ be the group of the isometries
$1+\nu$ of $X$ with $\nu=0$ on $I_1^\perp=I_1\oplus W_2$ and $\nu(X)\subseteq I_1$. The product of two such
$\nu$ is $0$, so $Z'$ is abelian and $(1+\nu)(1+\nu')=1+\nu+\nu'$; it contains $Z_D$, the elements $n_-(b)$
and $Z_A$, and it is the direct product of these three subgroups (write $\nu(f_k)=\sum_ie_iz_{ik}+b_ku$,
$\nu(u')=-eb^*+\zeta u$). Let $\psi'(1+\nu)=\psi_\beta(z)$, a character of $Z'(\auA)$ trivial on
$Z'(E)$, on the $n_-(b)$ and on $Z_A$, and put
\[
f^0(g)=\int_{Z'(E)\backslash Z'(\auA)}f(z'g)\,\psi'(z')^{-1}\,dz'=\int_{Z_A(E)\backslash Z_A(\auA)}f^+(\zeta g)\,d\zeta
\]
(probability measures; the two expressions agree by the definitions of $\pcuF_\beta$ and $f^+$ and Fubini's
theorem on the compact quotient). It is continuous. In particular $f^+$ is left invariant under $Z_A(E)$
(the elements of $Z_A$ commute with $Z_D$ and with the $n_-(b)$).

\emph{$f^0$ is left invariant under $N_A(E)$.} Let $\gamma\in N_A(E)$. It stabilises $I_1$ (it fixes $u$
and maps $e_k$ to $e_k-\langle c,e_k\rangle u$), hence $I_1^\perp$, so it normalises $Z'$; and
$\gamma^{-1}(1+\nu)\gamma=1+\gamma^{-1}\nu\gamma$ with $(\gamma^{-1}\nu\gamma)(f_k)=\gamma^{-1}\nu(f_k)$ (as
$\gamma f_k-f_k\in\langle u\rangle\subseteq\ker\nu$), whose components along $e$ are those of $\nu(f_k)$ (as
$\gamma^{-1}$ changes elements of $I_1$ by multiples of $u$). So $\psi'(\gamma^{-1}z'\gamma)=\psi'(z')$, and the
substitution $z'\mapsto\gamma z'\gamma^{-1}$, which preserves $Z'(E)$ and the measure, gives
$f^0(\gamma g)=\int f(\gamma\,\gamma^{-1}z'\gamma\,g)\psi'(z')^{-1}dz'=\int f(z'g)\psi'(z')^{-1}dz'=f^0(g)$, by the
left invariance of $f$ under $\gamma$.

\emph{Consequences of the hypothesis.} Let $g$ be admissible. For $\zeta\in Z_A(\auA)$ the point $\zeta g$ is
admissible ($Z_A$ fixes $e$), and $Z_A(E_{\mathcal U})\subseteq N_A(E_{\mathcal U})$; so the continuous function
$\zeta\mapsto f^+(\zeta g)$ on $Z_A(\auA)$ is invariant under $Z_A(E)$ and, by hypothesis, under
$Z_A(E_{\mathcal U})$. By strong approximation for $\mathbb G_a$ ($\mathcal U\ne\emptyset$, so $E$ is dense in the
adeles away from $\mathcal U$ and $E+E_{\mathcal U}$ is dense in $\auA_E$; the same holds for the vector spaces $Z_A$
and $X'$) it is constant. Hence
\[
f^+(g)=f^0(g)\qquad\text{for every admissible }g .
\]
Moreover, for $\varsigma\in N_A(E_{\mathcal U})$, as $Z_A$ is central in $N_A$ and $\zeta g$ is admissible,
$f^0(\varsigma g)=\int f^+(\varsigma\zeta g)\,d\zeta=\int f^+(\zeta g)\,d\zeta=f^0(g)$.

\emph{Conclusion.} Let $r\in N_A(\auA)$ be trivial off $\Sigma^c$ and $g$ admissible; then $rg$ is admissible.
The set $N_A(E)\cdot N_A(E_{\mathcal U})$ is dense in $N_A(\auA)$: for $c\in X'(E)$, $c_{\mathcal U}\in X'(E_{\mathcal U})$,
$n_A(c)n_A(c_{\mathcal U})=n_A(c+c_{\mathcal U})z$ with $z\in Z_A(\auA)$, and one approximates first the image of $r$
in $X'(\auA)$ and then, multiplying by elements of $Z_A(E)$ and $Z_A(E_{\mathcal U})$, the central component. Choose
$\gamma_i\in N_A(E)$, $\varsigma_i\in N_A(E_{\mathcal U})$ with $\gamma_i\varsigma_i\to r$. By the two invariances of
$f^0$ and its continuity, $f^0(rg)=\lim f^0(\gamma_i\varsigma_ig)=\lim f^0(\varsigma_ig)=f^0(g)$. As $g$ and $rg$
are admissible, $f^+(rg)=f^0(rg)=f^0(g)=f^+(g)$.
\end{proof}

\begin{definition}[Straight points and dual index]\label{pcu:def-straight}
A \emph{straight point} is $P=(g,g_0,t,t',\gamma)$ with $g$ fixing the frame of $D$ off $\Sigma^c$ and
equal to $g_0$ on $\Sigma^c$, where $g_0$ acts on $(u,e)$ by $t\oplus\gamma$ and on $(u',e)$ by
$t'\oplus\gamma$, with $t'^*t=1$. Its \emph{dual index} $\mathrm{di}(P)$ is the index in
$\widehat\auO_{E,\Sigma^c}$ of $\{x:\ x\det(\beta[\gamma])^{-1}\in\widehat\auO_{E,\Sigma^c}\}$ (the
denominator on $\Sigma^c$ of $\det\beta[\gamma]^{-1}$), a positive integer when $\beta[\gamma]$ is
integral on $\Sigma^c$.
\end{definition}

\begin{lemma}\label{pcu:bracket-int}
Let $f$ be automorphic and right $K_\Sigma$-invariant. If $\pcuF_\beta f(Y)\ne0$ for a $Y$ equal on
$\Sigma^c$ to a point acting on the frame of $D$ by $\gamma$, then $\beta[\gamma]$ is integral on
$\Sigma^c$. If moreover $\varphi^+_P\ne0$ for a straight point $P$, then $\beta[\gamma_P]$ is integral on
$\Sigma^c$.
\end{lemma}
\begin{proof}
For integral $M$, the element $n(0,M)$ taken on $\Sigma^c$ is in $K_\Sigma$; moving it across $Y$ turns it
into $n(0,\gamma M\gamma^*)$ on $\Sigma^c$, which multiplies $\pcuF_\beta f(Y)$ by $\psi_\beta(\gamma M\gamma^*)=
\psi(-\frac12\jmath\operatorname{tr}(Mh))$ restricted to $\Sigma^c$, $h=\gamma^*\beta\gamma$ (see \eqref{pcu:eq-psibeta}). Right
invariance makes this $1$ for all integral skew-hermitian $M$. Write $M=\theta_vH$ at $v\in\Sigma^c$, $H$ integral
hermitian, and $c=\jmath\theta_v\in E^+_v$. The diagonal entries give
$\psi(-\frac12cH_{ii}h_{ii})=\psi_\Q(\operatorname{Tr}_{E^+_v/\Q_p}(-cH_{ii}h_{ii}))=1$ for all $H_{ii}\in\auO_{E^+,v}$, and the
pairs of off-diagonal entries give $\psi(-\frac12c\operatorname{Tr}_{E/E^+}(H_{ij}h_{ji}))=\psi(-cH_{ij}h_{ji})=1$ for all
$H_{ij}\in\auO_{E,v}$. Hence $c\,h_{ji}\in\mathfrak d_{E,v}^{-1}$ for all $i,j$ (for $i=j$: $ch_{ii}\in
\mathfrak d_{E^+,v}^{-1}\subseteq\mathfrak d_{E,v}^{-1}$), i.e.\ $h\in c^{-1}\jmath\alpha\bar\alpha\,
\pcMat_b(\auO_{E,v})=\delta\,\pcMat_b(\auO_{E,v})$ ($\theta_v$ a unit): $\beta[\gamma]=\delta^{-1}h$ is integral on $\Sigma^c$.
For the second claim, $\varphi^+_P(a)\ne0$ gives by Lemma~\ref{pcu:pm}(b) and the definition some $b$
with $\pcuF_\beta f(n_-(b)n_+(a)g)\ne0$, and $n_\pm$ fix the frame of $D$, so the point acts by
$\gamma_P$ on $\Sigma^c$.
\end{proof}

\begin{lemma}[Iwasawa decomposition for the flag]\label{pcu:iwasawa}
\textup{(a)} For every $g\in\auU(X)(\auA)$ there is $\kappa$, integral at all finite places, such that
$g\kappa$ stabilises the flag $\langle u\rangle\subset\langle u\rangle\oplus\ell$, i.e.\ acts on $(u,e)$ by a
block upper triangular matrix $\binom{\alpha\ \beta'}{0\ \gamma}$. \textup{(b)} If $g'$ acts on $(u,e)$ by
such a matrix, then $g'=r\,m$ with $r\in N_A(\auA)$ (radical of the stabiliser of $\langle u\rangle$)
and $m$ in the Levi $\auGL_1\times\auU(X')$ acting on $u$ by $\alpha$, on $u'$ by $\bar\alpha^{-1}$, and on $e$
by $\gamma$.
\end{lemma}
\begin{proof}
(a) At each finite place $v$, let $\Lambda$ be the unimodular lattice of $X_v$ and
$Y_1=g_v^{-1}\langle u\rangle\subset Y_2=g_v^{-1}(\langle u\rangle\oplus\ell)$ (isotropic). The
sublattices $Y_i\cap\Lambda$ are primitive and isotropic; in a unimodular hermitian lattice over the
local ring $\auO_{E,v}$ every primitive isotropic flag extends to a hyperbolic basis adapted to it
(lift a basis of $Y_1\cap\Lambda$, then of $Y_2\cap\Lambda$ modulo it, pick dual vectors in $\Lambda$ by
unimodularity, and orthogonalise; Jacobowitz~\cite{Jacobowitz}, also at split places; here again
$E/E^+$ is unramified at $v$), so there is
$\kappa_v\in\auU(\Lambda)$ carrying the standard flag lattices $\auO u\subset\auO u\oplus\auO e$ to
$Y_i\cap\Lambda$. For almost all $v$, $g_v\in\auU(\Lambda)$ and $\kappa_v=g_v^{-1}$ works up to the standard
flag; put $\kappa_\infty=1$ (at the infinite places there is no condition: take $\kappa_\infty$ with
$g_\infty\kappa_\infty$ in the parabolic, by the real Iwasawa decomposition, the maximal compact acting
transitively on isotropic flags of the given type, and note $\kappa_\infty$ is integral vacuously).
Then $g\kappa$ maps $\langle u\rangle$ to itself and $\langle u\rangle\oplus\ell$ to itself.
(b) $g'$ stabilises $\langle u\rangle$, hence lies in $P_A=M_AN_A$ with $M_A$ the stabiliser of $u$ and
$u'$ lines and of $X'$. Let $m\in M_A$ act by $\alpha$ on $u$ and by the component $g'_{X'}$ (the map induced
on $\langle u\rangle^\perp/\langle u\rangle\cong X'$) on $X'$; then $r=g'm^{-1}$ acts trivially on
$\langle u\rangle$ and on $\langle u\rangle^\perp/\langle u\rangle$, so $r\in N_A$. Since $g'$ maps $e$ into
$e\gamma+u\beta'$, the induced map on $X'$ sends $e$ to $e\gamma$.
\end{proof}

\begin{lemma}[Straightening]\label{pcu:straighten}
Under the hypotheses of Lemma~\ref{pcu:unimod}, let $X_0$ fix the frame of $D$ off $\Sigma^c$, equal
$g$ on $\Sigma^c$, with $f^+(X_0)\ne0$, and let $g\kappa$ ($\kappa\in K_\Sigma$-integral) act on the
frame $(u,e)$ by $\binom{\alpha\ \beta'}{0\ \gamma}$ (such $\kappa$ exists by Lemma~\ref{pcu:iwasawa}). Then there is a straight point $P$ with $\gamma_P=\gamma$
and $\varphi^+_P\ne0$.
\end{lemma}
\begin{proof}
By Lemma~\ref{pcu:iwasawa}(b) write $g\kappa=r\,m$ with $r\in N_A$ and $m$ in the Levi of $A$ ($m$ acts by $\alpha$ on $u$ and
preserves $X'$, acting on $e$ by $\gamma$). Put $P_0=r^{-1}_{\Sigma^c}X_0\kappa_{\Sigma^c}$ (parts on
$\Sigma^c$): it equals $m$ on $\Sigma^c$, fixes the frame off $\Sigma^c$, and $f^+(P_0)=f^+(X_0\kappa_{\Sigma^c})
=f^+(X_0)\ne0$ by Lemma~\ref{pcu:sing} and right $K_\Sigma$-invariance. Then $P=(P_0,m,\alpha,
\bar\alpha^{-1},\gamma)$ is straight and $\varphi^+_P(0)=f^+(P_0)\ne0$.
\end{proof}

\begin{lemma}[Index decreases]\label{pcu:index}
Let $h$ act on $(u,e)$ by $\binom{t\ 0}{x\ \gamma}$, with $\gamma$ and $\beta[\gamma]$ invertible and
$\gamma^{-1}x$ not integral on $\Sigma^c$, and let $h\kappa$ ($\kappa$ integral) act by $\binom{\alpha\ \beta'}
{0\ \gamma'}$ with $\beta[\gamma']$ integral on $\Sigma^c$. Then $\mathrm{di}(\gamma')<\mathrm{di}(\gamma)$.
\end{lemma}
\begin{proof}
$\kappa^{-1}$ acts on the frame by an integral matrix $a$ (the frame is integral and unimodular), and
comparing, $x=\gamma'a_{21}$, $\gamma=\gamma'a_{22}$; hence $\gamma^{-1}x=a_{22}^{-1}a_{21}$ with $a_{21}$
integral, so $\det a_{22}$ is not a unit on $\Sigma^c$. Since $\beta[\gamma]=a_{22}^*\beta[\gamma']a_{22}$ and
$\beta[\gamma']$ is integral, $\det\beta[\gamma]^{-1}=\det\beta[\gamma']^{-1}\,\mathrm N(\det a_{22})^{-1}$ with
$\mathrm N(\det a_{22})=\det a_{22}\overline{\det a_{22}}$ integral and not a unit on $\Sigma^c$; as
$\det\beta[\gamma']$ is integral, the denominator of $\det\beta[\gamma]^{-1}$ is that of
$\det\beta[\gamma']^{-1}$ multiplied by the index of $\mathrm N(\det a_{22})\widehat\auO_{E,\Sigma^c}$, which is
$>1$.
\end{proof}

\begin{proof}[Proof of Lemma~\ref{pcu:unimod}]
(We prove it in the generality above, with $\Sigma=$ the places outside $S\cup\bar S$ used in
Proposition~\ref{pcu:selfdual}: there $f=\mathcal D(\Psi^{\mathrm{col}}_{z_0})$ is right invariant under
$K_\Sigma$ by (E3), and the shear invariance of $f^\pm$ at a non-empty set of places $\mathcal U$ (namely $\bar S$)
is the lemma proved after this one.)

Suppose $\pcuF_\beta f(g)\ne0$ for some $g$ fixing the frame off $\Sigma^c$. By Lemma~\ref{pcu:pm}(c),
$\varphi^+_g\ne0$, i.e.\ $f^+(n_+(a)g)\ne0$ for some $a$; by Lemma~\ref{pcu:straighten} there is a straight
point with $\varphi^+\ne0$. Choose one, $P$, with least $\mathrm{di}(P)$ (finite by
Lemma~\ref{pcu:bracket-int}); write $\gamma=\gamma_P$, $t,t'$ its scalars.

\emph{Supports.} If $\varphi^+_P(a)\ne0$ then $\gamma^{-1}a^*t$ is integral on $\Sigma^c$: otherwise the
point $n_+(a)g_0$ acts on $(u,e)$ by $\binom{t\ \ 0}{-a^*t\ \gamma}$, straightening $n_+(a)P$
(Lemma~\ref{pcu:straighten}, legitimate since $f^+(n_+(a)P)=\varphi^+_P(a)\ne0$) gives $Q$ with
$\varphi^+_Q\ne0$, and Lemma~\ref{pcu:index} with Lemma~\ref{pcu:bracket-int} gives
$\mathrm{di}(Q)<\mathrm{di}(P)$, a contradiction. The same holds for $\varphi^-$ with $t'$ (apply the
argument to the straight point with $u,u'$ exchanged, which has the same $\gamma$, and note that
$\varphi^+_{P}=0$ would force $\varphi^-_P=0$ by Lemma~\ref{pcu:pm}(c)).

\emph{Periods.} Let $x$ be a column with $\beta[\gamma]x$ integral on $\Sigma^c$ and
$a=\mathbf 1_{\Sigma^c}(\gamma x t'^*)^*$. For $b$ in the support of $\varphi^-_P$, $\chi_\beta(a,b)=1$:
indeed, with $c'=\gamma^{-1}b^*t'$ (integral on $\Sigma^c$ by Supports) and $t'^*t=1$, on $\Sigma^c$
$a\beta b^*=t'x^*\gamma^*\beta\gamma c't'^{-1}=\delta\,(\beta[\gamma]x)^*c'\in\delta\widehat\auO_E$, so
$\jmath a\beta b^*\in\jmath\alpha\bar\alpha\widehat\auO_E$, on which $\psi$ is trivial, and $a=0$ off $\Sigma^c$. By Lemma~\ref{pcu:pm}(b),
$\varphi^-_{n_+(a)P}=\varphi^-_P$, hence by Lemma~\ref{pcu:pm}(c) (with $u,u'$ exchanged)
$\varphi^+_{n_+(a)P}=\varphi^+_P$, i.e.\ $a$ is a period of $\varphi^+_P$. If $\varphi^+_P(y)\ne0$ then
$\varphi^+_P(y+a)\ne0$, and by Supports $\gamma^{-1}(y+a)^*t$ and $\gamma^{-1}y^*t$ are integral, so
$\gamma^{-1}a^*t=x$ (using $t'^*t=1$) is integral on $\Sigma^c$. Thus $\beta[\gamma]x$ integral implies $x$
integral: $\beta[\gamma]^{-1}$ is integral on $\Sigma^c$. With Lemma~\ref{pcu:bracket-int}, $\beta[\gamma]$ is
unimodular on $\Sigma^c$.

\emph{Contradiction.} By Lemma~\ref{pcu:pm}(c) and $\varphi^+_P\ne0$ there is a point $Y=n_-(b)n_+(a)P$,
equal on $\Sigma^c$ to a point acting on the frame of $D$ by $\gamma$ and fixing it elsewhere, with
$\pcuF_\beta f(Y)\ne0$. This contradicts the hypothesis of Lemma~\ref{pcu:unimod} (in the form: vanishing
at points acting by $\gamma$ on $\Sigma^c$, trivially elsewhere, with $\beta[\gamma]$ unimodular on $\Sigma^c$;
these are points $n(l,z)m(s_0q,g)$ of the statement: $Y$ stabilises $\ell$, so lies in $P_D(\auA)$, and its Levi
component on $\ell$ is $s_0q$ with $q=s_0^{-1}\gamma$ on $\Sigma^c$ and $q=s_0^{-1}$ at the finite places of $S\cup\bar S$;
then $q^*\beta q=\beta[\gamma]$ on $\Sigma^c$ and $q^*\beta q=s_0^{-1}\bar s_0^{-1}\beta$, the normalised index, at $S\cup\bar S$,
both unimodular).
\end{proof}

\begin{proof}[Shear invariance of $f^\pm$ for $f=\mathcal D(\Psi^{\mathrm{col}}_{z_0})$]
Let $\varsigma$ be a shear at $\bar S$ ($1+uc$ there, $cu=0$, trivial off $S\cup\bar S$) and $g$ fixing the frame
of $D$ off $\Sigma^c$. By (E2), $\mathcal D(\Psi')(xg)$ with $\Psi'=\Psi_0\circ(1+k\,c'\,\cdot)$
($k$, $c'$ the $S$-components of $(\varsigma g)^{-1}u$ and $c\,\varsigma g$) equals $\mathcal D(\Psi_0)$ at $\varsigma^{-1}$-translated
points, so $f^+(\varsigma g)$ is the value at $\Psi'$ of the linear form $\Lambda=$ ``$f^+$ at $\varsigma g$'' (linear on
$\mathcal M$ by Lemma~\ref{pcu:pm}(b) and (E1)), while $f^+(g)=\Lambda(\Psi_0)$ after moving $\varsigma$. The form
$\Lambda$ is invariant under the shears $y\mapsto y+k\,w\,z_0y$ ($w$ any row), because these come from the
shifts $n_-(\cdot)$ at $S$, which fix $f^+$ (Lemma~\ref{pcu:pm}(b)); here $z_0 k=0$ and $c'k=0$. It remains:
\emph{if a linear form $\Lambda$ on a space of Schwartz--Bruhat functions is invariant under all
$y\mapsto y+k\,w\,z_0y$, then $\Lambda(\Psi_0\circ(1+kr\,\cdot))=\Lambda(\Psi_0)$ for every row $r$ with
$rk=0$.} Average over integral $t$: $\Lambda$ is unchanged by replacing $\Phi$ with
$\bar\Phi(y)=\int\Phi(y+k\,r\,t\,z_0y)dt$ (a finite combination, by local constancy), and
$\bar{\Psi_0\circ(1+kr\cdot)}=\bar\Psi_0$: if $\Psi_0(y+k\rho)\ne0$ for some $\rho$, then $z_0y$ has an
integral left inverse $g'$ and $(1+kr)(y+krtz_0y)=y+kr(t+t_0)z_0y$ with $t_0=yg'$ integral, so the two
averages agree by translation invariance of $dt$; otherwise both vanish.
\end{proof}

\subsubsection{Prescribed slices: proof of (M1), and the decay of slices}\label{pcu:ss-M1}

Write the real Fock spaces as $\mathcal F'$ (of $\auHom(X,T)_\R\cong\C^{\iota'}$) and $\mathcal F$ (of
$\auHom(X_1,T)_\R\cong\C^{\iota}$), with moments $\|v\|_N=\int(1+|y|)^N|v(y)|dy$, Heisenberg operators
$\psrho$ and vacua $e$. The mixed model of $D$ is characterised by: there are linear maps $L,K\colon
\C^p\to\C^{\iota'}$, $I\colon\C^\iota\to\C^{\iota'}$ (here $p$ is the complex dimension of
$\auHom(\ell,T)_\R$) with $\C^{\iota'}=L\C^p\oplus I\C^\iota\oplus K\C^p$, and bounded maps
$\mathrm{sl}(x)\colon\mathcal F'\to\mathcal F$, such that
\begin{equation}\label{pcu:eq-mixed}
\mathrm{sl}(x)\,\psrho(La'+Iw+Kb)v=\beta(2x+a',b)\,\psrho(w)\,\mathrm{sl}(x+a')v,\qquad
\mathrm{sl}(x)e=e^{-|x|^2}e,
\end{equation}
$\beta(x,b)=\exp(i\,\Im\,b^*x)$ (the Schr\"odinger variable is translated by $L$, multiplied by a
character by $K$, and the rest acts on the values). Let $y\mapsto(A y,Wy,By)$ be the inverse of
$(a,w,b)\mapsto La+Iw+Kb$.

\begin{lemma}[Slices decay]\label{pcu:slicedecay}
For all $N,M$ there is $C$ with $(1+|x|)^M\|\mathrm{sl}(x)v\|_N\le C\|v\|_{N+M}$ for all $x$, $v$.
In particular $\int\|\mathrm{sl}(x)v\|_0\,dx\le C'\|v\|_{2p+1}$.
\end{lemma}
\begin{proof}
For $v=\psrho(y)e$, \eqref{pcu:eq-mixed} gives $\mathrm{sl}(x)v=(\text{phase})\,e^{-|x+Ay|^2}\psrho(Wy)e$.
From $1+|x|\le(1+|x+Ay|)(1+|Ay|)$ and $(1+t)^Me^{-t^2}\le e^{M^2}$,
$(1+|x|)^Me^{-|x+Ay|^2}\le e^{M^2}(1+\|A\||y|)^M$; and $\|\psrho(w)e\|_N\le(1+|w|)^N\|e\|_N$ (translate
the weight). Hence $(1+|x|)^M\|\mathrm{sl}(x)\psrho(y)e\|_N\le C(1+|y|)^{N+M}$, and
Lemma~\ref{pcx:superpos} ($v=\int v(y)\psrho(y)e\,dy$, $\mathrm{sl}(x)$ bounded) gives the bound. The
second statement follows with $M=2p+1$, since $\int(1+|x|)^{-2p-1}dx<\infty$ on $\C^p=\R^{2p}$.
\end{proof}
This is the estimate used in (M2): the integral $\int\mathrm{sl}_D(x)\Lambda\,dx$ converges
absolutely in $\mathcal F$ as soon as $\|\Lambda\|_{2p+1}<\infty$.

\begin{proof}[Proof of (M1)]
Let $\tau$, $u$, $v$ be as in (M1); write $\xi=v$, and for $g$ in the real unitary group of (M1) let $S_g$ be the real
symplectic map of $\auHom(X_1,T)_\R=\C^\iota$ by which $g$ acts, so that $\omega(g)\psrho(w)\omega(g)^{-1}=\psrho(S_gw)$
((W1)). Write $S_x=S_{u(x)}$; on the compact set $\mathrm{supp}\,\tau$, $x\mapsto S_x$ is smooth and
$B=\sup\|S_x\|$, $B'=\sup\|S_x^{-1}\|$ are finite. No compactness of the unitary groups is used. Functions of vectors
are taken in the realisation of Lemma~\ref{pcx:superpos}: $F(w)=\langle\psrho(w)e,F\rangle$.

\emph{Step 1: Gaussian vectors.} Put $E_g=\omega(g)e$. (i) $|E_g(w)|\le\exp(-|w|^2/(1+\|S_g^{-1}\|^2))$: this is the
bound proved in the first part (``Reduction'') of the proof of Lemma~\ref{pcu:moments}, applied to the unitary operator
$k=\omega(g)^{-1}$, for which $k\psrho(w)k^{-1}=\psrho(S_g^{-1}w)$ and $k^{-1}e=E_g$. (ii) In a Schr\"odinger model of
$\C^\iota$ in which $e$ is the Gaussian $\gamma_{Z_0}(t)=\exp(\boldsymbol\pi i\,t^{\mathsf t}Z_0t)$, $Z_0=i1$, the operator
$\omega(g)$ is, up to a scalar, the metaplectic operator of $S_g$, and the metaplectic operators carry $\gamma_{Z_0}$ to a
multiple of $\gamma_{Z}$ with $Z$ the image of $Z_0$ under the linear fractional action on the Siegel upper half space of
$S_g$ or of $(S_g^{\mathsf t})^{-1}$, according to the normalisation (Folland~\cite[Chapter~4, Thm.~4.65]{FollandPS}: the metaplectic operator of $A$ carries $\gamma_Z$ to a multiple of $\gamma_{\alpha((A^{\mathsf t})^{-1})Z}$). So $E_g=c(g)\gamma_{Z(g)}$ with $Z(g)$ smooth in $g$; and $c(g)=\langle
E_g,e\rangle/\langle\gamma_{Z(g)},\gamma_{Z_0}\rangle$ is smooth, because the denominator is an explicit Gaussian integral,
smooth and nowhere zero in $Z$, and $g\mapsto\langle\omega(g)e,e\rangle$ is smooth, $e$ being a Schwartz function and hence a
smooth vector (Folland~\cite[Chapter~4, Thm.~4.45]{FollandPS}). Computing the Gaussian integral
$\langle\psrho(w)\gamma_{Z_0},\gamma_Z\rangle$ gives $E_g(w)=c'(g)\exp Q_g(w)$ with $Q_g$ a quadratic form in the real
coordinates of $w$, with complex coefficients, and $c'$, $Q_g$ smooth in $g$. Hence every derivative of order $|\alpha|$ in
$g$ of $E_g(w)$, and in $w$ of order $|\alpha'|$, is $E_g(w)$ times a polynomial of degree $\le2|\alpha|+|\alpha'|$ in $w$
whose coefficients are smooth in $g$.

\emph{Step 2: the functions of the prescribed slices.} By Lemma~\ref{pcx:superpos} and boundedness of $\omega(g)$,
$\omega(g)\xi=\int\xi(y)\psrho(S_gy)E_g\,dy$. The Heisenberg relations give
$\langle\psrho(w)e,\psrho(a)E\rangle=e^{i\vartheta(a,w)}E(w-a)$ with $\vartheta$ a real quadratic polynomial in $(a,w)$
(the cocycle of the Heisenberg group). So the function of $G(x)=\tau(x)\omega(u(x))\xi$ is
\[
g_x(w)=\tau(x)\int\xi(y)\,e^{i\vartheta(S_xy,w)}\,E_{u(x)}(w-S_xy)\,dy .
\]
Put $\varrho=1+|y|+|w|$. By Step 1 and the chain rule, each $x$-derivative of the integrand brings a factor $O(\varrho^2)$,
uniformly on $\mathrm{supp}\,\tau$, so with $c=(1+B'^2)^{-1}$
\[
\bigl|\partial_x^\alpha\bigl[\tau(x)e^{i\vartheta(S_xy,w)}E_{u(x)}(w-S_xy)\bigr]\bigr|\le C_\alpha\,\varrho^{2|\alpha|}e^{-c|w-S_xy|^2}.
\]
Since $\varrho\le(1+|w-S_xy|)(1+B)(1+|y|)$, this is $\le C'_\alpha(1+|y|)^{2|\alpha|}(1+|w-S_xy|)^{2|\alpha|}e^{-c|w-S_xy|^2}$,
which is integrable against $|\xi(y)|$; so differentiation under the integral is legitimate, and, using
$1+|w|\le(1+|w-S_xy|)(1+B|y|)$ and integrating first in $w$,
\begin{equation}\label{pcu:eq-M1bound}
\int\!\!\int(1+|w|)^N\,|\partial_x^\alpha g_x(w)|\,dw\,dx\le C_{N,\alpha}\,\|\xi\|_{N+2|\alpha|}<\infty
\end{equation}
for all $N$, $\alpha$ (the $x$-integral runs over the compact set $\mathrm{supp}\,\tau$).

\emph{Step 3: construction.} Put $h_w(x)=e^{|x|^2}g_x(w)$, smooth with compact support in $x$;
$\hat h_w(b)=\int h_w(x)e^{-2\pi i\langle x,b\rangle}dx$ (Fourier transform in $x\in\C^p=\R^{2p}$); and
$k(b)=K(-\pi i\,b)$, so that by \eqref{pcu:eq-mixed} $\mathrm{sl}(x)\psrho(k(b))=e^{2\pi i\langle x,b\rangle}\mathrm{sl}(x)$.
Define
\[
\Lambda=\int\!\!\int\hat h_w(b)\,\psrho(k(b))\,\psrho(Iw)e\,db\,dw .
\]
\emph{Moments.} $\|\psrho(k(b))\psrho(Iw)e\|_N\le(1+\|k\||b|)^N(1+\|I\||w|)^N\|e\|_N$ (translation of the weight).
Integrating by parts in $x$, $(1+|b|)^M|\hat h_w(b)|\le C_M\sum_{|\alpha|\le M}\int|\partial_x^\alpha h_w(x)|dx\le
C'_M\sum_{|\alpha|\le M}\int|\partial_x^\alpha g_x(w)|dx$ ($e^{|x|^2}$ and its derivatives are bounded on
$\mathrm{supp}\,\tau$). Taking $M=N+2p+1$, $\int(1+\|k\||b|)^N(1+|b|)^{-M}db<\infty$, and \eqref{pcu:eq-M1bound} gives
$\iint|\hat h_w(b)|\,\|\psrho(k(b))\psrho(Iw)e\|_N\,db\,dw<\infty$. So the integral converges absolutely for every
$\|\cdot\|_N$, and $\|\Lambda\|_N<\infty$ for all $N$.

\emph{Slices.} The Hilbert norm is $\le\|\cdot\|_0$ (Lemma~\ref{pcx:superpos}), so the integral converges absolutely in
$\mathcal F'$ and the bounded operator $\mathrm{sl}(x)$ passes inside. By \eqref{pcu:eq-mixed} with $a'=0$, $b=0$,
$\mathrm{sl}(x)\psrho(Iw)e=\psrho(w)\mathrm{sl}(x)e=e^{-|x|^2}\psrho(w)e$; and by Fourier inversion ($h_w$ smooth with
compact support), $\int\hat h_w(b)e^{2\pi i\langle x,b\rangle}db=h_w(x)$. Hence
\[
\mathrm{sl}(x)\Lambda=\int h_w(x)e^{-|x|^2}\psrho(w)e\,dw=\int g_x(w)\psrho(w)e\,dw=G(x)=\tau(x)\,\omega(u(x))\xi,
\]
the last but one equality by Lemma~\ref{pcx:superpos}, as $\|G(x)\|_0<\infty$ by \eqref{pcu:eq-M1bound} with $\alpha=0$
(pointwise in $x$ by the same bound before integrating in $x$).
\end{proof}

\subsubsection{Exchange on standard data: proof of Proposition~\ref{pcx:P-std}}\label{pcu:ss-std}

\begin{proof}[Proof of Proposition~\ref{pcx:P-std} from Propositions~\ref{pcx:P-slice},
\ref{pcx:P-top} and~\ref{pcx:P-descent}]
Let $W$ (rank $N=b+r$, even, $q$ negative signs), $U$ (rank $b$), $V$ (rank $r$) with
$q_U+q=b+q_V$, $W_0=\pcxH\oplus W$ (adapted, Lemma~\ref{pcx:adapted}), $z$ for $(W_0,W_0)$ and $y$ for
$(W,U)$ be given.

\emph{Step 1.} Proposition~\ref{pcx:P-slice} for the adapted splitting $W_0=W\oplus\pcxH$ gives $w$ for
$(W_0,U)$ with finite moments and zero term $y$.

\emph{Step 2.} Proposition~\ref{pcx:P-top} for $W_0$, $U$, $z$, $w$ gives $X=\pcxH^b\oplus W_0=
\pcxH^{b+1}\oplus W$, $T=U\oplus W_0$, and $x$, $x'$ with zero term $w$ along the $b$ planes and
$\mathrm{RC}(X,T,U;x,x';\nu\mapsto\pcxstr_{W_0}(\nu;z))$. Here $\operatorname{rk}T=b+N+2$ and $T$ has
$q_U+q+1$ negative signs.

\emph{Step 3: the tower.} For $i=0,\dots,b+1$ let $X_i=\pcxH^{i}\oplus W$ (so $X_{b+1}=X$, $X_0=W$,
each $X_{i+1}=X_i\oplus\pcxH$ adapted and of even rank). Let $T_{b+1}=T$ and, for $0\le i\le b$,
$T_i=V\oplus\pcxH^i$ (Lemma~\ref{pcx:adapted}), of rank $r+2i$ with $q_V+i$ negative signs. Then
$\operatorname{rk}T_{i+1}=\operatorname{rk}T_i+2$ and $q_{T_{i+1}}=q_{T_i}+1$ for all $i$: for $i<b$
this is clear, and for $i=b$, $\operatorname{rk}T=b+N+2=r+2b+2$ and $q_T=q_U+q+1=q_V+b+1$ by the
hypothesis $q_U+q=b+q_V$.

\emph{Step 4: descent.} Apply Proposition~\ref{pcx:P-descent} $b+1$ times, to $W'=X_{i+1}=X_i\oplus
\pcxH$, $T=T_{i+1}$, $T_0=T_i$, $T'=U$, $i=b,b-1,\dots,0$: each application keeps $\ell$ and replaces
the second vector by its zero term along the plane removed. After $b+1$ steps one gets $v$ for
$(W,V)$ with $\mathrm{RC}(W,V,U;v,y';\ell)$, where $y'$ is the iterated zero term of $x'$ along the
$b+1$ planes. By (M4) this is the zero term of $x'$ along $X=\pcxH^b\oplus W_0$ followed by the zero
term along $W_0=\pcxH\oplus W$, i.e.\ the zero term of $w$, which is $y$ (Steps 1--2). So
$\mathrm{RC}(W,V,U;v,y;\nu\mapsto\pcxstr_{W_0}(\nu;z))$, which is Proposition~\ref{pcx:P-std}.
\end{proof}

\begin{remark}[Consistency conditions on Proposition~\ref{pcx:P-descent}]\label{pcu:rem-descent}
The proof needs from Proposition~\ref{pcx:P-descent}, as stated in \S\ref{pcx:ss-arch}: (i) it is
applied with the zero term taken along the splitting with the same orientation as the one in which
Proposition~\ref{pcx:P-top} and Proposition~\ref{pcx:P-slice} give zero terms (along $D^{\mathrm s}$
for the $b$ planes); (ii) the sign condition ``with the signs
required'' means $q_{T_0}=q_T-1$; (iii) the exponent in $\mathrm{RC}$(a) changes from
$\operatorname{rk}T-\operatorname{rk}W'/2=r-N/2+i+1$ to $r-N/2+i$, and the constant $C$ is allowed
to change while $\ell$ is kept; (iv) $T_0$ is an arbitrary space of rank $\operatorname{rk}T-2$ with
$q_T-1$ negative signs (not necessarily $T$ minus a plane); removal of $p\ge1$ planes follows by
induction from one plane, as in Step 4. Proposition~\ref{pcx:P-descent} is proved for $T=T_0\oplus
\pcxH$ (adapted), as in the set-up \S\ref{pct:ss-setup}, and Step~3 works with this form: take $T_b=U\oplus W$ (rank
$b+N=r+2b$, $q_U+q=q_V+b$ negative signs), so $T=U\oplus W_0=T_b\oplus\pcxH$; descend once to
$(X_b,T_b)$; then $T_b$ and $V\oplus\pcxH^b$ have the same rank and signs, and Lemma~\ref{pcx:move}
(with $W_1=W=X_b$, $k=1$) transfers $\mathrm{RC}(X_b,T_b,U;\cdot)$ to
$\mathrm{RC}(X_b,V\oplus\pcxH^b,U;\cdot)$, multiplying $c$ by a number of modulus $1$ independent of
the data and keeping $C$; the remaining $b$ descents are of the form $T_{i+1}=T_i\oplus\pcxH$.
\end{remark}

\subsubsection{Dependencies of Theorem~\ref{pcx:X}}\label{pcu:ss-final}

In the deductions of \S\ref{pcx:ss-assembly} and \S\ref{pcx:ss-equal} the statements of \S\ref{pcx:ss-arch} are applied
as follows: (1) in the proof of Proposition~\ref{pcx:P-matches},
Proposition~\ref{pcx:P-std} is applied to $(W,U^-,V)$, whose sign condition $q_{U^-}+q=b+q_V$ is
equivalent to $q=q_U+q_V$; the moments of $y^{\mathrm{ad}}$ are those of $y$; (2) in the proof of
Theorem~\ref{pcx:X}, Proposition~\ref{pcx:P-matches} is applied to $(W,V',V)$ with $y=v$ a vector of
$(V',W)$, $N=a+b=2b+2e$ is even, $q=q_V+q_{V'}$, and $a-N/2=e$; (3) Lemma~\ref{pcx:twice} applies
Proposition~\ref{pcx:P-matches} to $(J,0,J)$ and $(J,J,0)$, which satisfy its hypotheses (for $b=0$
Proposition~\ref{pcx:P-top} is the identity on $X=T=W_0$ and the descent is from $(W_0,W_0)$ to
$(J,J)$). By Corollary~\ref{pcx:conv}(iii) both integrals of \eqref{pcx:eq-X} converge absolutely. The
statements on which Theorem~\ref{pcx:X} depends:

\medskip
{\small\noindent\begin{tabular}{@{}p{0.36\textwidth}p{0.6\textwidth}@{}}
\toprule
Statement & Where proved\\ \midrule
Theorem~\ref{pcx:X} & \S\ref{pcx:ss-assembly}, from \ref{pcx:P-matches}, \ref{pcx:P-equal},
\ref{pcx:Lpos}, \ref{pcx:mod}, \eqref{pcx:eq-twist}\\
Convergence (Cor.~\ref{pcx:conv}) & \S\ref{pcx:ss-conv}, from Prop.~\ref{pcx:decay}, (R1), (R2), (R3)\\
Lemma~\ref{pcx:Lpos}, Cor.~\ref{pcx:cne0} & \S\ref{pcx:ss-statement}; Hecke's non-vanishing on
$\Rel s=1$\\
Lemmas~\ref{pcx:std-matched}, \ref{pcx:mod} & \S\ref{pcx:ss-elem}\\
Prop.~\ref{pcx:P-matches} & \S\ref{pcx:ss-assembly}, from \ref{pcx:P-std}, \ref{pcx:P-match},
\ref{pcx:P-flip}\\
Prop.~\ref{pcx:P-match} & \S\ref{pcx:ss-match} (with Cor.~\ref{pcx:conv})\\
Prop.~\ref{pcx:P-flip} & \S\ref{pcx:ss-flip}, from {\cite[Prop.~7.17]{P49}}\\
Prop.~\ref{pcx:P-std} & \S\ref{pcu:ss-std}, from \ref{pcx:P-slice}, \ref{pcx:P-top},
\ref{pcx:P-descent}, \ref{pcx:adapted}, (M4)\\
Prop.~\ref{pcx:P-slice} & \S\ref{pcu:ss-slice}, from (M1), (M2)\\
Prop.~\ref{pcx:P-top} & \S\S\ref{pcu:ss-witt}--\ref{pcu:ss-std-data}, from \ref{pcx:const},
\ref{pcx:landherr}, \ref{pcx:generate}, \ref{pcx:linfactor}, \ref{pcx:section-laws}, \ref{pcx:Sact},
(M1)--(M3), (U1)--(U3), Lemmas~\ref{pcu:unimod} (\S\ref{pcu:ss-unimod}), \ref{pcu:integrable}\\
Prop.~\ref{pcx:P-descent} & \S\ref{pct:ss-descent} (Proposition~\ref{pct:descent})\\
Prop.~\ref{pcx:P-equal} & \S\ref{pcx:ss-equal}, from \ref{pcx:P-matches}, \ref{pcx:P-nonvan},
\ref{pcx:P-plane}, \ref{pcx:P-signs}, \ref{pcx:adapted}\\
Prop.~\ref{pcx:P-nonvan} & \S\ref{pct:ss-nonvan}, with Lemma~\ref{pct:negcoeff} $=$ Theorem~\ref{pcn:negcoeff} (local vanishing, Lemma~\ref{pcn:local}; no unfolding)\\
Prop.~\ref{pcx:P-plane}, Lemma~\ref{pcx:const} & \S\S\ref{pct:ss-plane}, \ref{pct:ss-const}\\
Prop.~\ref{pcx:P-signs} & \S\ref{pcx:ss-signs}, from \ref{pcx:move}, \ref{pcx:const}\\
Lemmas~\ref{pcx:adapted}, \ref{pcx:landherr}, \ref{pcx:move} & \S\S\ref{pcx:ss-arch},
\ref{pcx:ss-signs}; Landherr, Hasse norm theorem, Jacobowitz\\
Prop.~\ref{pcx:decay} & \S\ref{pcx:ss-conv}; (R1) reduction theory, (R2) moderate growth\\
(M1)--(M4), Lemma~\ref{pcu:mixedmodel} & \S\ref{pcu:ss-mixed}, \S\ref{pcu:ss-M1}; Bargmann transform (Folland)\\
\bottomrule
\end{tabular}}
\medskip

Thus Theorem~\ref{pcx:X}, with absolutely convergent integrals, rests on the proofs listed and on
the published results listed in \S\S\ref{pct:ss-status}, \ref{pcn:ss-dep}.



% ===== C-exchange-tower.tex =====
% proofs/C-exchange-tower.tex
% The plane constant term of theta integrals on standard data and what rests on it:
% Lemma pcx:const, Propositions pcx:P-plane, pcx:P-descent, pcx:P-nonvan of C-exchange.tex.
% Labels of this file start with pct:. Macros: those of C-exchange.tex and C-determinant.tex.
\providecommand{\pctN}{\mathsf N}
\providecommand{\pctP}{\mathsf P}
\providecommand{\pctL}{\mathcal L}
\providecommand{\pctZ}{\mathcal Z}
\providecommand{\pctK}{\mathsf K}
\providecommand{\pctsl}{\operatorname{sl}}
\providecommand{\pctcone}{\operatorname{cone}}

\subsection{The plane constant term}\label{pct:sec}

This section proves Lemma~\ref{pcx:const} and Propositions~\ref{pcx:P-plane},
\ref{pcx:P-descent} and (modulo one input from the large triple, see \S\ref{pct:ss-nonvan})
\ref{pcx:P-nonvan}. All four rest on one computation, the constant term of a theta integral
$\Theta^{\mathrm L}$ along the isotropic line of an added hyperbolic plane
(Theorem~\ref{pct:plane}). \emph{Convention:} in this section a Fock vector $v$ \emph{has finite
moments} if $\|v\|_N=\int(1+|y|)^N|v(y)|\,dy<\infty$ for every $N\ge0$ (moments of all orders; the norms
of Lemma~\ref{pcx:superpos}); by Lemma~\ref{pcx:superpos}(b) such a $v$ lies in $\pzF_{s'}$ for every
integer $s'\ge0$. Notation is that of \S\ref{pcx:ss-arch}: $\nu$ is a tame character of
weight $\mu$ (all $\mu_\tau\ge1$) for a modulus $\pcm\ne(1)$, $S\ne\emptyset$ is the set of primes
of $\pcm$ (so $S\cap\bar S=\emptyset$), $F=F_S$, $O=\auO_S$, $\chi=\chi_S$, $Q=\prod_{w\in S}q_w$,
and for a space $W$ of even rank $N$
\[
\gamma_W=\nu^{-1}\upsilon^{N/2},\qquad
\Theta^{\mathrm L}_{W,T}(v,\Phi)(g)=\int_{[\auU(T)]}\theta^{W,T}_{\gamma_W,\gamma_W}(g,h;[v,\Phi]_S)\,dh .
\]
We write $\tilde\nu(x)=\nu(x/\bar x)$ for the base change of $\nu$, a unitary Hecke character of
$\auA_E^\times/E^\times$, $|\cdot|$ for the idele norm of $E$, and
\[
L^\sharp(x)=|L^S(x,\nu)|\,Q^{x-1/2}\qquad(x\ge1).
\]
Then $L_*(\pcm,\nu,e+1)=L^\sharp(x_e)^2L_*(\pcm,\nu,e)$ for every $e\in\Z$, where $x_e=e+1$ if
$e\ge0$ and $x_e=-e$ if $e<0$: for $e\ge0$ this is $Q^{(e+1)^2-e^2}|L^S(e+1,\nu)|^2$, and for
$e=-(m+1)$ with $m\ge0$, inverting the case $e=m$ through $L_*(-e')=L_*(e')^{-1}$ gives
$L_*(-m)=L^\sharp(m+1)^2L_*(-m-1)$, which is the claim since $x_e=m+1$. By Lemma~\ref{pcx:Lpos}, $L^\sharp(x)>0$ for $x\ge1$. Facts
(W1)--(W5) are those of {\cite[Lemma~5.2]{P49}}.

The base change $\tilde\nu$ has archimedean component $\tilde\nu_\infty(a)=\prod_\tau(a_\tau/|a_\tau|)^{2\mu_\tau}$
on $(E\otimes\R)^\times=\prod_\tau\C^\times$; at a place $v=w\bar w$ with $w\in S$, $\tilde\nu_w$ and
$\tilde\nu_{\bar w}$ are ramified (their restrictions to units are $\chi_w$ and $\chi_w^{-1}$ composed
with conjugation), and $\tilde\nu$ is unramified outside $S\cup\bar S\cup\infty$. Hence the Euler
factors of $L(s,\tilde\nu)$ at $S\cup\bar S$ are $1$ and $L^S(s,\nu)=L(s,\tilde\nu)$.

\subsubsection{Set-up: the plane, its parabolic and the line terms}\label{pct:ss-setup}

Let $W'=W\oplus\pcxH_A$ and $V'=V\oplus\pcxH_X$ be adapted splittings (Lemma~\ref{pcx:adapted}):
$\pcxH_A$ has the isotropic basis $e_A,f_A$ with $\langle e_A,f_A\rangle=1$, the lattice of $W'$ is
$L_W\oplus\widehat\auO_Ee_A\oplus\widehat\auO_Ef_A$, and likewise $e_X,f_X$ for $V'$. Put
$N=\operatorname{rk}W$ (always even below), $r=\operatorname{rk}V$, $b=N-r$ and
\[
d=r-N/2\in\Z .
\]
Then $\operatorname{rk}W'=N+2$, $\operatorname{rk}V'=r+2$, $(N+2)-(r+2)=b$, and
$\gamma_{W'}=\gamma_W\upsilon$.

\emph{Parabolics.} $\pctP_A\subset\auU(W')$ is the stabiliser of the line $Ee_A$, $\pctN_A$ its
unipotent radical, and for $u\in\auA_E^\times$, $g\in\auU(W)(\auA)$ let $m_A(u,g)$ be the Levi
element acting by $g$ on $W$, by $e_A\mapsto ue_A$ and $f_A\mapsto\bar u^{-1}f_A$. For $V'$ we use the
parametrisation $\ell_X(a,h)$: $h$ on $V$, $e_X\mapsto a^{-1}e_X$, $f_X\mapsto\bar af_X$. The group
$\pctN_A(E)\backslash\pctN_A(\auA)$ is compact; $dn$ is its probability measure. For a continuous
function $F$ on $\auU(W')(\auA)$, left invariant under $\pctN_A(E)$, its \emph{constant term} is
$F_A(x)=\int_{\pctN_A(E)\backslash\pctN_A(\auA)}F(nx)\,dn$. The group $\pctN_X(\auA)$ has
coordinates $(c,t)$, $c\in V(\auA)$, $t$ in the $\auA_{E^+}$-line $\{t\in\auA_E:t+\bar t=0\}$ (a
central extension of $V(\auA)$ by this line), and $\ell_X(a,h)$ acts by conjugation through
$c\mapsto hc\,\bar a^{-1}$, $t\mapsto(a\bar a)^{-1}t$; as $|\det h|=1$, this gives
$\int f(\ell_X(a,h)n\ell_X(a,h)^{-1})\,dn=|a|^{r+1}\int f(n)\,dn$ for a Haar measure $dn$ of
$\pctN_X(\auA)$.

\emph{The plane datum.} Let $\mathrm{pr}\in\pcMat_{N\times(N+2)}(O)$ be the matrix, in the frames, of
the orthogonal projection $W'\to W$; it is integral because the splitting is adapted. For
$z\in\pcMat_{b\times N}(O)$ put
\[
\Phi'_z=\Phi^{\mathrm{row}}_{z\,\mathrm{pr}}\qquad\text{on }\pcMat_{(r+2)\times(N+2)}(F)
\]
(Definition~\ref{pcx:def-std} for $(W',V')$; $z\,\mathrm{pr}\in\pcMat_{b\times(N+2)}(O)$). For $b=0$,
$\Phi'_\emptyset=\Phi^{\mathrm{row}}_\emptyset=\chi^{-1}\circ\det$ on $\auGL_{N+2}(O)$, which is
$\Psi^{\mathrm{col}}_1$ of Definition~\ref{pcx:def-ratio}.

\emph{Coordinates on $\auHom(W',V')$.} Write a map $\phi\colon W'\to V'$ as a $3\times3$ block matrix
with rows indexed by $(V,e_X,f_X)$ and columns by $(W,e_A,f_A)$. In the form
$\langle\phi,\psi\rangle=\operatorname{tr}(\psi^*\phi)$ the block $(y,a)$ pairs non-trivially only
with the block $(y^\vee,a^\vee)$, where $e^\vee=f$, $f^\vee=e$, $V^\vee=V$, $W^\vee=W$. Put
\[
\lambda_{\mathrm t}=E_{e_X,e_A},\qquad\lambda_{\mathrm f}=E_{f_X,f_A},\qquad
I_0=(V,f_A)\oplus(e_X,W)\oplus(e_X,f_A),
\]
where $E_{y,a}$ is the map $a\mapsto y$ (zero on the other basis vectors and on $W$), and $(y,a)$
denotes the corresponding block. Then $I_0$ is isotropic, $\Pi=E\lambda_{\mathrm t}\oplus
E\lambda_{\mathrm f}$ is a hyperbolic plane orthogonal to $I_0$, and
$I_0^\perp/I_0=\Pi\oplus\auHom(W,V)$. Also $I_A=(V',f_A)=(V,f_A)\oplus(e_X,f_A)\oplus(f_X,f_A)$
(the maps vanishing on $W\oplus Ee_A$) is isotropic and $I_A^\perp/I_A=\auHom(W,V')$, and
$I_A\cong V'$ by $x\mapsto(w'\mapsto\langle w',e_A\rangle x)$.

\begin{definition}[Slices and line terms]\label{pct:def-line}
For $x\in V'(\auA)$ let $x\otimes f_A^\vee\in\auHom(W',V')(\auA)$ be $w'\mapsto\langle w',f_A\rangle x$
(this is $x$ placed in the column $e_A$), and define the \emph{slice}
$\pctsl_A(x)=\psZ_{I_A}\circ\psrho(x\otimes f_A^\vee)\colon\psHS(\auHom(W',V'))\to
\psHS(\auHom(W,V'))$ (zero terms as in (W2)). For $\alpha\in\auA_E$ the two \emph{line terms}
$\pctL_{\mathrm t}(\alpha),\pctL_{\mathrm f}(\alpha)\colon\psHS(\auHom(W',V'))\to\psHS(\auHom(W,V))$
are
\[
\pctL_{\mathrm t}(\alpha)=\psZ_{I_0+E\lambda_{\mathrm f}}\circ\psrho(\alpha\lambda_{\mathrm t}),\qquad
\pctL_{\mathrm f}(\alpha)=\psZ_{I_0+E\lambda_{\mathrm t}}\circ\psrho(\alpha\lambda_{\mathrm f}).
\]
By the transitivity of zero terms in (W2), $\pctL_{\mathrm t}(\alpha)=\psZ_{(e_X,W)}\circ
\pctsl_A(\alpha e_X)$ (first the zero term along $I_A$, then along $\auHom(W,Ee_X)\subset
\auHom(W,V')$), and $\pctL_{\mathrm f}(\alpha)$ is the zero term along $\auHom(W',Ee_X)$ followed by
the zero term along $(V,f_A)$, after $\psrho(\alpha\lambda_{\mathrm f})$. All these operators are
tensor products of local operators; the archimedean factors, on Fock spaces, are written
$\pctsl_{A,\infty}(x)$ and $\pctL^\infty_\beta(a)$ ($\beta\in\{\mathrm t,\mathrm f\}$,
$a\in(E\otimes\R)$).
\end{definition}

\begin{lemma}[Poisson's formula along the line]\label{pct:poisson}
For every datum $f$ of $\auHom(W',V')$,
$\sum_{\xi\in E}\theta_{W,V}(\pctL_{\mathrm t}(\xi)f)=\sum_{\xi\in E}\theta_{W,V}(\pctL_{\mathrm f}(\xi)f)$,
both series converging absolutely.
\end{lemma}
\begin{proof}
Put $F_0=\psZ_{I_0}f\in\psHS(\Pi\oplus\auHom(W,V))$. Since $\lambda_{\mathrm t},\lambda_{\mathrm f}\in
I_0^\perp$, (W2) gives $\pctL_{\mathrm t}(\xi)f=\psZ_{E\lambda_{\mathrm f}}\psrho(\xi\lambda_{\mathrm t})F_0$
and $\pctL_{\mathrm f}(\xi)f=\psZ_{E\lambda_{\mathrm t}}\psrho(\xi\lambda_{\mathrm f})F_0$, the zero
terms being taken in $\Pi\oplus\auHom(W,V)$. Apply (W3) in this space, once with the isotropic line
$I=E\lambda_{\mathrm f}$ (then $W/I^\perp\cong E\lambda_{\mathrm t}$) and once with $I=E\lambda_{\mathrm t}$:
both sums equal $\theta_{\Pi\oplus\auHom(W,V)}(F_0)$, with absolute convergence.
\end{proof}

\begin{lemma}[Levi elements]\label{pct:levi}
Let $u,c\in\auA_E^\times$, $g\in\auU(W)(\auA)$, $h\in\auU(V)(\auA)$, and assume $u=1$ or $r=N$.
Put $x=cu$, $P_{\mathrm t}(x)=x$, $P_{\mathrm f}(x)=\bar x^{-1}$. Then for every datum $f$ and
$\alpha\in\auA_E$,
\[
\pctL_\beta(\alpha)\,\omega_{\gamma_{W'},\gamma_{W'}}\bigl(m_A(u,g),\ell_X(c,h)\bigr)f
=\frac{|u|^{r+1}\,\tilde\nu(x)}{|x|^{N/2+[\beta=\mathrm f]}}\;
\omega_{\gamma_W,\gamma_W}(g,h)\,\pctL_\beta\bigl(\alpha P_\beta(x)\bigr)f .
\]
\end{lemma}
\begin{proof}
We use (W2) in the following explicit form for a line split off one of the two spaces. For
characters $a,b$ put $\chi^{b}_k(x)=b(x/\bar x)\upsilon(x)^{-k}$ ($x\in\auA_E^\times$; the character of
a space of rank $k$). (\emph{Target line}) If $T$ has rank $k$, $V'=V\oplus\pcxH_X$ and $\psZ$ is the
zero term along $\auHom(T,Ee_X)$, then for $q\in\auU(T)(\auA)$
\[
\psZ\,\omega_{a,\,b\upsilon}(q,\ell_X(c,h))=\chi^{a}_k(c)^{-1}|c|^{-k/2}\,\omega_{a,b}(q,h)\,\psZ ;
\]
(\emph{source line}) if $T'$ has rank $k$ and $\psZ$ is the zero term along the maps $W'\to T'$
vanishing on $W\oplus Ee_A$, then $\psZ\,\omega_{a\upsilon,\,b}(m_A(u,g),q)=\chi^{b}_k(u)^{-1}|u|^{k/2}\,
\omega_{a,b}(g,q)\,\psZ$. Both are (W2): the element stabilises the isotropic subspace $I$ and induces
$(q,h)$, resp.\ $(g,q)$, on $I^\perp/I$; its determinant on $I(\auA)$ is $c^{-k}$ times a determinant of
$q$, resp.\ $\bar u^{k}$ times a determinant of $q$, and $\chi_I=\upsilon(\det)^{-1}|\det|^{1/2}$
combines with the change of the splitting characters ($\det\ell_X(c,h)=(\bar c/c)\det h$,
$\det m_A(u,g)=(u/\bar u)\det g$, the extra $\upsilon$ being evaluated there) to the stated scalar
(the first character of the splitting is that of the target group). Moreover, by (W1),
$\psrho(\alpha\lambda_\beta)\omega(p)=\omega(p)\psrho(p^{-1}\alpha\lambda_\beta)$ with
$p^{-1}\lambda_{\mathrm t}=cu\,\lambda_{\mathrm t}$ and $p^{-1}\lambda_{\mathrm f}=\overline{cu}^{-1}
\lambda_{\mathrm f}$, which turns $\alpha$ into $\alpha P_\beta(x)$.

$\beta=\mathrm t$: $\pctL_{\mathrm t}=\psZ_{(e_X,W)}\circ\pctsl_A$; first the source line on
$\auHom(W',V')$ ($k=r+2$, $a=\gamma_W$, $b=\gamma_{W'}$), then the target line on $\auHom(W,V')$ ($k=N$,
$a=\gamma_W$, $b=\gamma_W$). The scalar is $\chi^{\gamma_{W'}}_{r+2}(u)^{-1}|u|^{(r+2)/2}\cdot
\chi^{\gamma_W}_N(c)^{-1}|c|^{-N/2}$. Now $\chi^{\gamma_W}_N(x)=\nu(x/\bar x)^{-1}\upsilon(x/\bar x)^{N/2}
\upsilon(x)^{-N}=\tilde\nu(x)^{-1}$, because the character $\upsilon$ of $\auU(1)$ is normalised by
$\upsilon(x/\bar x)=\upsilon(x)^2$ (well defined since $\upsilon^2$ is trivial on $\auA_{E^+}^\times$;
$N$ is even); likewise
$\chi^{\gamma_{W'}}_{N+2}=\tilde\nu^{-1}$. If $u=1$ the first factor is $1$; if $r=N$ it is
$\tilde\nu(u)|u|^{(N+2)/2}=\tilde\nu(u)|u|^{r+1-N/2}$. In both cases the scalar is
$\tilde\nu(cu)|u|^{r+1}|cu|^{-N/2}$.

$\beta=\mathrm f$: first the target line on $\auHom(W',V')$ ($k=N+2$, $a=\gamma_{W'}$), then the source
line on $\auHom(W',V)\to\auHom(W,V)$ ($k=r$); the scalar is $\tilde\nu(c)|c|^{-(N+2)/2}\cdot
\chi^{\gamma_W}_r(u)^{-1}|u|^{r/2}$, which for $u=1$ or $r=N$ equals $\tilde\nu(cu)|u|^{r+1}
|cu|^{-N/2-1}$ ($r/2=r+1-N/2-1$ when $r=N$).

The hypothesis is used exactly to identify $\chi_{r+2}$, resp.\ $\chi_r$, at $u$ with $\tilde\nu^{-1}$ and
the power of $|u|$; for $r\ne N$, $u\ne1$ the formula is false.
\end{proof}

\begin{remark}
The restriction $u=1$ or $r=N$ is essential: for $r\ne N$ the constant term at $m_A(u,g)$ is not
independent of $u$. Only the cases used below occur: $u=1$ (Propositions~\ref{pcx:P-plane},
\ref{pcx:P-descent}, \ref{pcx:P-nonvan}) and $r=N$ (Lemma~\ref{pcx:const}).
\end{remark}

\subsubsection{Cone integrals and the statement}\label{pct:ss-statement}

Let $\pctK'=\pctK'_\infty\pctK'_f$ be the maximal compact subgroup of $\auU(V')(\auA)$ ($\pctK'_f$ the
stabiliser of the lattice, $\pctK'_\infty$ maximal compact in $\auU(V')(E^+\otimes\R)$), with
probability Haar measure $d\kappa$, and let $da$ be a Haar measure of $(E\otimes\R)^\times=
\prod_\tau\C^\times$. Write $\omega'_\infty$ for the archimedean component of
$\omega_{\gamma_{W'},\gamma_{W'}}$ and $\|\cdot\|_M$ for the moments of \S\ref{pcx:ss-conv}.

\begin{lemma}[Slices decay]\label{pct:slicedecay}
For all $N_0,M$ there is $C$ such that for every Fock vector $w$ of $\auHom(W',V')\otimes\R$,
$a\in E\otimes\R$, $x\in V'(E\otimes\R)$, $\kappa\in\pctK'_\infty$ and $\beta\in\{\mathrm t,\mathrm f\}$:
$(1+|a|_\infty)^M\|\pctL^\infty_\beta(a)w\|_{N_0}\le C\|w\|_{N_0+M}$,
$\|\pctsl_{A,\infty}(x)w\|_{N_0}\le C\|w\|_{N_0}$ and $\|\omega'_\infty(1,\kappa)w\|_{M}=\|w\|_M$, where
$|a|_\infty=\max_\tau|a_\tau|$.
\end{lemma}
\begin{proof}
By Lemma~\ref{pcx:superpos}, $w=\int w(y)\psrho_\infty(y)e\,dy$, so it suffices to bound the operators
on the coherent states $\psrho_\infty(y)e$ by $C(1+|y|)^{N_0+M}$, resp.\ $C(1+|y|)^{N_0}$, in the norm
$\|\cdot\|_{N_0}$. Let $I$ be an isotropic subspace of a real hermitian space $Y=I\oplus I'\oplus
Y_I$ ($I'$ dual to $I$). In the mixed model $\psHS_\infty(Y)=\mathcal S(I')\otimes\psHS_\infty(Y_I)$
the vacuum is a Gaussian $\varphi_0\otimes e_I$, $\psrho(y_1)$ for $y_1\in I^\perp$ acts by a phase
times $\psrho(\bar y_1)$ on the second factor, and $\psrho(y_2)$, $y_2\in I'$, translates the first
factor; $\psZ_I$ is evaluation of the first factor at $0$. Writing $y=y_1+y_2$, one gets
$\psZ_I\psrho(y)e=\epsilon\,\varphi_0(y_2)\psrho(\bar y_1)e_I$ with $|\epsilon|=1$, and
$\|\psrho(\bar y_1)e_I\|_{N_0}\le C(1+|y_1|)^{N_0}$. Applied to $\pctL^\infty_\beta(a)=
\psZ_{I_\beta}\psrho(a\lambda_\beta)$, the $I'$-component of $y+a\lambda_\beta$ contains $a$ in the
coordinate dual to $\lambda_{\beta^\vee}$, so $|\varphi_0(y_2)|\le Ce^{-c|a-y'|^2}$ with $y'$ a
coordinate of $y$, and $e^{-c|a-y'|^2}\le C_M(1+|a|)^{-M}(1+|y|)^M$. The slice is the case
$I=I_A$ without the factor in $a$. Finally $\pctK'_\infty$ fixes the line of $e$ (it lies in the
maximal compact subgroup of the symplectic group of $\auHom(W',V')\otimes\R$, which acts on $e$
by a character), so $\omega'_\infty(1,\kappa)\psrho(y)e=\epsilon'\psrho(\kappa y)e$ with $|\epsilon'|=1$
and $|\kappa y|=|y|$; this gives the equality of moments.
\end{proof}

\begin{definition}[Cone integrals]\label{pct:def-cone}
For a Fock vector $v$ of $\auHom(W',V')\otimes\R$ with finite moments and $x>0$ put
\[
\pctcone_x(v)=\int_{\pctK'_\infty\times(E\otimes\R)^\times}|a|^x\,\tilde\nu_\infty(a)\,
\pctL^\infty_{\mathrm t}(a)\,\omega'_\infty(1,\kappa)v\;d\kappa\,da,
\]
and $\pctcone^\vee_x(v)$ the same with $\pctL^\infty_{\mathrm f}$. (Here $\tilde\nu_\infty(a)=
\prod_\tau(a_\tau/|a_\tau|)^{2\mu_\tau}$ depends only on $\mu$.)
\end{definition}

By Lemma~\ref{pct:slicedecay} the integrand has $\|\cdot\|_{N_0}$-norm at most
$C|a|^x(1+|a|_\infty)^{-M}\|v\|_{N_0+M}$, and $\int|a|^x(1+|a|_\infty)^{-M}da<\infty$ for $x>0$ and
$M>2x[E^+:\Q]$ (in polar coordinates $|a|^x=\prod_\tau r_\tau^{2x}$ and $da=\prod_\tau
r_\tau^{-1}dr_\tau\,d\theta_\tau$). So the cone integrals converge absolutely in every norm
$\|\cdot\|_{N_0}$ and have finite moments.

For a datum $f$ of $\auHom(W',V')$, $g\in\auU(W)(\auA)$, $h\in\auU(V)(\auA)$, $\kappa\in\pctK'$
and $\alpha\in\auA_E$ put
\[
\Lambda_\beta(f;\alpha;g,h,\kappa)=\theta_{W,V}\bigl(\omega_{\gamma_W,\gamma_W}(g,h)\,\pctL_\beta(\alpha)\,
\omega_{\gamma_{W'},\gamma_{W'}}(1,\kappa)f\bigr),
\]
and, for $u\in\auA_E^\times$ and real $x$, the \emph{zeta integral}
\[
\pctZ(f,g,u;x)=\int_{E^\times\backslash\auA_E^\times}\int_{[\auU(V)]}\int_{\pctK'}\tilde\nu(au)|au|^x
\sum_{\xi\in E^\times}\Lambda_{\mathrm t}(f;\xi au;g,h,\kappa)\,d\kappa\,dh\,da,
\]
whenever it converges absolutely. The substitution $a\mapsto au^{-1}$ shows
$\pctZ(f,g,u;x)=\pctZ(f,g,1;x)$.

\begin{theorem}[The plane constant term]\label{pct:plane}
Let $N$ be even and $d=r-N/2$. There are constants $C_{\mathrm t},C_{\mathrm f}>0$ and an idele
$\alpha_0$, depending only on the splittings, on $\mu$ and on the Weil data, such that for every
tame $\nu$ of weight $\mu$ there are $\varepsilon_{\mathrm t},\varepsilon_{\mathrm f}\in\C$ with
$|\varepsilon_\beta|=C_\beta$ and the following holds for every Fock vector $v$ of
$\auHom(W',V')\otimes\R$ with finite moments, every $z\in\pcMat_{b\times N}(O)$, every
$g\in\auU(W)(\auA)$ and every $u\in\auA_E^\times$ with $u=1$ or $r=N$. Let
$f=[v,\Phi'_z]_S$ and $P(x)=|\alpha_0|^xQ^{x-1/2}$.
\begin{enumerate}
\item[(a)] There is $c_X\ne0$, depending only on Haar measures, such that the function
$\pctZ_{v,z,g}(x)=c_X\pctZ(f,g,u;x)$ is defined by an absolutely convergent integral for every
real $x$, does not depend on $u$, and is real-analytic on $\R$.
\item[(b)] $\bigl(\Theta^{\mathrm L}_{W',V'}(v,\Phi'_z)\bigr)_A(m_A(u,g))=\pctZ_{v,z,g}(d+1)$.
\item[(c)] For $x\ge1$:
$\pctZ_{v,z,g}(x)=\varepsilon_{\mathrm t}P(x)L^S(x,\nu)\,\allowbreak \Theta^{\mathrm L}_{W,V}(\pctcone_xv,
\Phi^{\mathrm{row}}_z)(g)$ and
$\pctZ_{v,z,g}(1-x)=\varepsilon_{\mathrm f}P(x)L^S(x,\nu)\,\allowbreak \Theta^{\mathrm L}_{W,V}(\pctcone^\vee_xv,
\Phi^{\mathrm{row}}_z)(g)$.
\end{enumerate}
In particular there are $C>0$ and, for each $\nu$, $\kappa_\nu\in\C$ with $|\kappa_\nu|=C\,
L^\sharp(d+1)$ if $d\ge0$, $|\kappa_\nu|=C\,L^\sharp(-d)$ if $d<0$, such that, for all $v,z,g$ and
$u$ as above,
\begin{equation}\label{pct:eq-plane}
\bigl(\Theta^{\mathrm L}_{W',V'}(v,\Phi'_z)\bigr)_A(m_A(u,g))=\begin{cases}
\kappa_\nu\,\Theta^{\mathrm L}_{W,V}(\pctcone_{d+1}v,\Phi^{\mathrm{row}}_z)(g),&d\ge0,\\
\kappa_\nu\,\Theta^{\mathrm L}_{W,V}(\pctcone^\vee_{-d}v,\Phi^{\mathrm{row}}_z)(g),&d<0.
\end{cases}
\end{equation}
If $r=N$, the left side does not depend on $u$.
\end{theorem}

The last two assertions follow from (a)--(c): for $d\ge0$ take $x=d+1\ge1$ in the first formula
of (c), for $d<0$ take $x=-d\ge1$ in the second, so that $1-x=d+1$; then $\kappa_\nu=
\varepsilon_\beta P(x)L^S(x,\nu)$ has modulus $C_\beta|\alpha_0|^xL^\sharp(x)$, and $x$ is fixed by
the spaces. Independence of $u$ for $r=N$ is (a)+(b). Note that only values of $L^S$ at points
$\ge1$ enter, and no functional equation is used. The rest of this subsection and the next prove
(a)--(c).

\subsubsection{Proof of Theorem~\ref{pct:plane}}\label{pct:ss-proof}

\emph{Step 1: expansion of the constant term.} For $y\in V'(E)$ and a datum $F$ of
$\auHom(W',V')$ put $T_y(F)=\int_{\pctN_A(E)\backslash\pctN_A(\auA)}\theta_{W,V'}(\pctsl_A(y)
\omega(n,1)F)\,dn$, with $\omega=\omega_{\gamma_{W'},\gamma_{W'}}$.

\begin{lemma}\label{pct:expand}
Let $x'\in\auU(W')(\auA)$ and $f=[v,\Phi'_z]_S$ as in Theorem~\ref{pct:plane}, where the Fock vector $v$
has moments of all orders ($\|v\|_N<\infty$ for every $N$). Then
$\bigl(\Theta^{\mathrm L}_{W',V'}(v,\Phi'_z)\bigr)_A(x')=\int_{[\auU(V')]}\sum_{y\in V'(E)}
T_y(\omega(x',h)f)\,dh$, and $T_y=0$ unless $\langle y,y\rangle=0$.
\end{lemma}
\begin{proof}
By (W3) for the isotropic subspace $I_A$ (whose dual is the column $e_A$, identified with $V'$ by
$y\mapsto y\otimes f_A^\vee$), $\theta(F)=\sum_{y\in V'(E)}\theta_{W,V'}(\pctsl_A(y)F)$, the series
converging absolutely. Let $\mathsf N_c\subset\pctN_A(\auA)$ be a compact set of representatives of
$\pctN_A(E)\backslash\pctN_A(\auA)$. For fixed $x',h$, {\cite[Lemma~8.3(b)]{P49}} (the elements
$(nx',h)$, $n\in\mathsf N_c$, have finite parts in a compact set and bounded archimedean parts; the
archimedean vector $v$ has moments of all orders by hypothesis, so it lies in every $\pzF_{s'}$
(Lemma~\ref{pcx:superpos}(b)), hence is a smooth vector (Lemma~\ref{pcn:smooth}), and
{\cite[Lemma~8.3(b)]{P49}} applies to it directly) bounds
$\sum_y|\theta_{W,V'}(\pctsl_A(y)\omega(nx',h)f)|$ uniformly for $n\in\mathsf N_c$, and Fubini's theorem
gives
$\int\theta(\omega(nx',h)f)\,dn=\sum_yT_y(\omega(x',h)f)$ for each $h$. The kernel
$\theta^{W',V'}_{\gamma_{W'},\gamma_{W'}}(nx',h;f)$ is bounded by $C\,\mathrm{ht}(nx')^B
\mathrm{ht}(h)^{-N_0}$ (Proposition~\ref{pcx:decay}(a): $\Phi'_z$ is supported on matrices of full
rank $r+2$ by Lemma~\ref{pcx:std-matched}); $\mathrm{ht}$ is invariant under rational points and
$\pctN_A(E)\backslash\pctN_A(\auA)x'$ is compact, so Fubini applies on
$(\pctN_A(E)\backslash\pctN_A(\auA))\times[\auU(V')]$ by (R1). This gives the formula. The centre of
$\pctN_A(\auA)$ consists of the elements $\phi\mapsto\phi+t\langle\cdot,e_A\rangle e_A$-type
transvections ($t+\bar t=0$), which act on $\pctsl_A(y)$ by the character
$t\mapsto\psi(t\langle y,y\rangle)$ ((W1) and the commutation relation of the Heisenberg group);
its mean over $E\backslash\auA$ vanishes unless $\langle y,y\rangle=0$.
\end{proof}

\emph{Step 2: the term $y=0$ vanishes.} The element $m_A(u,g)$ stabilises $I_A$, so by (W2)
$T_0(\omega(m_A(u,g),h)f)$ is a multiple of $\theta_{W,V'}(\omega(g,h)\pctsl_A(0)f)$ (the
radical acts trivially on the slice at $0$). At a place $w\in S$, in the model (W5), the zero term
along $I_A$ is the composite of an average over translations by $I_{A,w}$ and of the average over
the characters $x\mapsto[x,a]$, $a\in I_{A,\bar w}$; the latter restricts a function to the
annihilator of $I_{A,\bar w}$, which consists of the matrices $x$ with $xe_A=0$. Now
$\Phi'_z(x)\ne0$ forces $\binom{z\,\mathrm{pr}}{x}\in\auGL_{N+2}(O)$, while $z\,\mathrm{pr}\,e_A=0$ and
$xe_A=0$; so $\Phi'_z$ vanishes on that annihilator and $\pctsl_A(0)f=0$.

\emph{Step 3: unfolding along the isotropic lines.}
\begin{lemma}\label{pct:unfold}
Let $x'=m_A(u,g)$ with $u=1$ or $r=N$, and put $\Phi_{x'}(h)=\sum_{\xi\in E^\times}T_{\xi e_X}(\omega(x',h)f)$.
\begin{enumerate}
\item[(i)] For $\xi\ne0$, $T_{\xi e_X}(F)=\theta_{W,V}(\pctL_{\mathrm t}(\xi)F)$.
\item[(ii)] $\Phi_{x'}$ is left invariant under $\pctN_X(\auA)$ and under the rational points
$\pctP_X(E)$ of the stabiliser of $Ee_X$, and $\sum_{y\ne0,\ \langle y,y\rangle=0}T_y(\omega(x',h)f)=
\sum_{\delta\in\pctP_X(E)\backslash\auU(V')(E)}\Phi_{x'}(\delta h)$.
\item[(iii)] If the integral $\pctZ(f,g,u;d+1)$ converges absolutely, then
$\bigl(\Theta^{\mathrm L}_{W',V'}(v,\Phi'_z)\bigr)_A(x')=c_X\pctZ(f,g,u;d+1)$, with $c_X\ne0$
depending only on Haar measures.
\end{enumerate}
\end{lemma}
\begin{proof}
(i) $\pctN_A(\auA)$ modulo its centre is $W(\auA)$, acting on $\pctsl_A(\xi e_X)F$ through
$\psrho$ of the subspace $\auHom(W,E\xi e_X)=(e_X,W)$ (by (W1)); the centre acts trivially as
$\langle\xi e_X,\xi e_X\rangle=0$. By the first formula of (W2), the mean over
$(e_X,W)(E)\backslash(e_X,W)(\auA)$ of $\theta_{W,V'}\circ\psrho$ is $\theta_{W,V}\circ\psZ_{(e_X,W)}$,
and $\psZ_{(e_X,W)}\pctsl_A(\xi e_X)=\pctL_{\mathrm t}(\xi)$ (Definition~\ref{pct:def-line}).

(ii) An element of $\pctN_X(\auA)$ fixes $e_X$, stabilises $I_0+E\lambda_{\mathrm f}$ and induces the
identity on $\auHom(W,V)$; by (W1), (W2) it commutes with $\pctL_{\mathrm t}(\xi)$ (the character
$\chi_I$ being trivial on unipotent elements). For $\delta\in\auU(V')(E)$, (W1) and the
invariance $\theta\circ\omega(\delta)=\theta$ give $T_{\delta y}(\omega(1,\delta)F)=T_y(F)$; an element
of $\pctP_X(E)$ permutes the vectors $\xi e_X$, $\xi\in E^\times$. By Witt's theorem the nonzero
isotropic vectors of $V'(E)$ are the $\delta^{-1}\xi e_X$, $\delta\in\auU(V')(E)$, $\xi\in E^\times$,
the pair $(\pctP_X(E)\delta,\xi)$ being determined up to the action of $\pctP_X(E)$ on $\xi$; grouping
the vectors by the line they span gives the identity.

(iii) By Lemma~\ref{pct:levi} with $c=a$ (and (W1) to move $\omega(1,\kappa)$),
$|a|^{r+1}\Phi_{x'}(\ell_X(a,h)\kappa)=\sum_{\xi}\tilde\nu(au)|au|^{d+1}\allowbreak \Lambda_{\mathrm t}(f;\xi au;g,h,\kappa)$.
By the Iwasawa decomposition $\auU(V')(\auA)=\pctN_X(\auA)\,\ell_X(\auA_E^\times\times\auU(V)(\auA))
\,\pctK'$ and the modulus computed at the end of \S\ref{pct:ss-setup}, for a measurable function $\Phi$ on
$\auU(V')(\auA)$, left invariant under $\pctN_X(\auA)$ and $\pctP_X(E)$,
\[
\int_{\pctP_X(E)\backslash\auU(V')(\auA)}\Phi=c_X\int_{E^\times\backslash\auA_E^\times}\int_{[\auU(V)]}
\int_{\pctK'}|a|^{r+1}\Phi(\ell_X(a,h)\kappa)\,d\kappa\,dh\,da,
\]
whenever the right side converges absolutely; here $\pctP_X(E)=\pctN_X(E)\ell_X(E^\times\times
\auU(V)(E))$ and $\pctN_X(E)\backslash\pctN_X(\auA)$ has volume $1$ (the standard integration
formula for a parabolic subgroup; e.g.\ Moeglin--Waldspurger, \emph{Spectral decomposition and
Eisenstein series}). Applied to $|\Phi_{x'}|$ (whose right side is
bounded by the absolutely convergent $\pctZ$ integrand) and then to $\Phi_{x'}$, and combined with
(ii) and the unfolding $\int_{[\auU(V')]}\sum_{\delta\in\pctP_X(E)\backslash\auU(V')(E)}\Phi(\delta h)\,dh=
\int_{\pctP_X(E)\backslash\auU(V')(\auA)}\Phi$ (legitimate since $|\Phi|$ is integrable on the right),
Lemma~\ref{pct:expand} and Step~2 give the claim.
\end{proof}

\emph{Step 4: the line terms of the plane datum factor.} Let $\varpi\in O$ be the product of
uniformisers, $\psi$ the standard additive character, and $\alpha_0$ the finite idele denoted $\alpha$ in
{\cite[\S5.1.1]{P49}} (we keep the name $\alpha_0$, as $\alpha$ is a variable idele below). This is the choice the theorem requires: the data $f=[v,\Phi'_z]_S$ are data of
\S\ref{pc:ss-data}, whose components off $S$ are the vacua of the lattices $\alpha_0\auHom(L_{W'},L_{V'})_v$
(self-dual because $\widehat\auO_E$ is self-dual for $\psi(\jmath\alpha_0\aubar\alpha_0\,\aubar zu)$,
Lemma~\ref{pgx:vac}), and whose components at $w\in S$ are supported on $a_w\pcMat(\auO_w)$ with
$a_w\auO_w=\alpha_{0,w}\auO_w$ (Definition~\ref{pgx:conv}). (Self-duality alone would fix only the products
$\alpha_{0,w}\overline{\alpha_{0,\bar w}}$; what is used below is that the scaling idele is the one defining the
data, and this holds by definition.) The components of $\alpha_0$ are units outside the primes dividing the
relative different of $E/K$ (proof of Proposition~\ref{pgv:alpha}(a)), and $\alpha_0$ depends only on the
Weil data. At $w\in S$ we use the coordinates of Definition~\ref{pgx:conv}
described after Lemma~\ref{pcx:mod}, in which the vacuum lattice is $\pcMat(\auO_w)$ and the
component of $\alpha_0$ at $w$ becomes the unit $\alpha_{0,w}/a_w$, its \emph{unit component} (equal to
$\alpha_{0,w}$ where $\alpha_{0,w}$ is a unit, as then $a_w=1$). For $\alpha\in\auA_E$ write $\alpha_S=(\alpha_w)_{w\in S}\in F$
and $\alpha_{\bar S}^\flat=(\overline{\alpha_{\bar w}})_{w\in S}\in F$, and put
\[
\lambda_\varpi(\alpha)=\begin{cases}\chi(\alpha_S)^{-1}\chi(\varpi\alpha^\flat_{\bar S}),&\alpha_f\text{ integral
off }S\cup\bar S,\ \alpha_S\in O^\times,\ \varpi\alpha^\flat_{\bar S}\in O^\times,\\0,&\text{otherwise},
\end{cases}
\]
and $\lambda_1$ the same with $\varpi$ replaced by $1$.

\begin{lemma}[Gauss sums]\label{pct:gauss}
Let $w\in S$, $\chi_w$ the tame part (a character of $O_w^\times$ trivial on $1+\varpi_wO_w$,
non-trivial), and $\psi_w$ any character of $E_w$ of conductor $O_w$ (it will be applied to
$\lambda_w^{\pm1}$ of Lemma~\ref{pgx:dict}(b)). There is $g_w\in\C$ with
$|g_w|^2=q_w^{-1}$ such that for all $c\in E_w$,
$\int_{O_w^\times}\chi_w^{-1}(t)\psi_w(ct)\,dt=g_w\,\chi_w(\varpi_wc)$ if $\varpi_wc\in O_w^\times$,
and $0$ otherwise ($dt$ of total mass $1$ on $O_w$).
\end{lemma}
\begin{proof}
If $c\in O_w$ the integrand's $\psi_w$-factor is $1$ and $\int_{O_w^\times}\chi_w^{-1}=0$. If
$v(c)\le-2$, replacing $t$ by $t(1+\varpi_ws)$, $s\in O_w$, does not change $\chi_w^{-1}(t)$ and
multiplies $\psi_w(ct)$ by $\psi_w(c\varpi_wst)$; averaging over $s$ gives $0$ as $c\varpi_wt\notin
O_w$. If $c=\varpi_w^{-1}c_0$, $c_0\in O_w^\times$, the integral is $q_w^{-1}\sum_{t\in\auF_w^\times}
\chi_w^{-1}(t)\psi_w(\varpi_w^{-1}c_0t)=q_w^{-1}\chi_w(c_0)G$ with the Gauss sum $G$ of $\chi_w^{-1}$
and the non-trivial character $t\mapsto\psi_w(\varpi_w^{-1}t)$ of $\auF_w$, $|G|^2=q_w$. Put
$g_w=q_w^{-1}G$.
\end{proof}

\begin{lemma}[Factorisation]\label{pct:factor}
There are $l_{\mathrm t}=l_{\mathrm t}(\nu),l_{\mathrm f}=l_{\mathrm f}(\nu)\in\C^\times$ whose moduli
$|l_{\mathrm t}|,|l_{\mathrm f}|$ depend only on the splittings and the Weil data, and for each $\nu$ a number $g_S$ with $|g_S|^2=Q^{-1}$, such that for all $v$, all
$z\in\pcMat_{b\times N}(O)$ and all $\alpha\in\auA_E$,
\[
\pctL_\beta(\alpha_0\alpha)[v,\Phi'_z]_S=l_\beta\,g_S\,\lambda_\varpi(\alpha)\,
[\pctL^\infty_\beta(\alpha_\infty)v,\Phi^{\mathrm{row}}_z]_S .
\]
Moreover the finite points $\pctK'_f$ fix $[v,\Phi'_z]_S$ in $\omega_{\gamma_{W'},\gamma_{W'}}$.
\end{lemma}
\begin{proof}
All operators are tensor products over the places, and so is the datum. At $\infty$ the factor is
$\pctL^\infty_\beta(\alpha_\infty)v$ by definition.

\emph{Places $v\notin S\cup\bar S\cup\infty$.} The datum is the vacuum $\psvac_L$ of the self-dual
lattice $L=\alpha_0\auHom(L_{W'},L_{V'})_v$, for which $\lambda_{\mathrm t},\lambda_{\mathrm f}$ and the blocks
of $I_\beta$ are spanned by vectors of $\auHom(L_{W'},L_{V'})_v$ (the splittings are adapted). Let $I=I_\beta$ and
$y=\alpha_0\alpha\lambda_\beta$. For $i\in I\cap L$, $\psZ_I\psrho(i)=\psZ_I$ and $\psrho(i)\psvac_L=
\psvac_L$ (with the normalisation $\psrho^\circ$ of (W4)), while
$\psrho(i)\psrho(y)=\psi([i,y])\psrho(y)\psrho(i)$; hence $\psZ_I\psrho(y)\psvac_L=0$ unless
$\psi([i,y])=1$ for all $i\in I\cap L$, i.e.\ unless $y\in L+I^\perp$, i.e.\ (as $\lambda_\beta$ pairs
with the primitive lattice vector $\lambda_{\beta^\vee}\in I$, and $\alpha_0$ makes $L$ self-dual)
unless $\alpha_v$ is integral. If it is, $\psZ_I\psrho(y)\psvac_L$ is a root of unity times
$\psZ_I\psvac_L$, the vacuum of the self-dual lattice of $I^\perp/I$; the root of unity depends only
on the Weil data and goes into $l_\beta$.

\emph{Places $w\in S$} (and their conjugates, which carry the same local factor $v=w\bar w$). In
the model (W5) data are functions on $\pcMat_{(r+2)\times(N+2)}(E_w)$ (matrices in the frames of
$W'$, $V'$). \emph{Change of coordinates.} Let $P=(P_W\,|\,P_{e}\,|\,P_f)\in\auGL_{N+2}(O_w)$ be the
matrix whose columns are the coordinates of the frame of $W$, $e_A$, $f_A$ in the frame of $W'$, and
$R=P^{-1}$, with block rows $R_W=\mathrm{pr}$, $R_e$, $R_f$; similarly $G\in\auGL_{r+2}(O_w)$ for
$(V,e_X,f_X)$ in $V'$. Both are integral with integral inverse because the splittings are adapted.
Write $x=G\tilde xR$; then $\tilde x$ is the block matrix with rows $(V,e_X,f_X)$ and columns
$(W,e_A,f_A)$ used above, the substitution preserves the measure (unimodular), and
\[
\begin{pmatrix}z\,\mathrm{pr}\\ x\end{pmatrix}=\begin{pmatrix}1&0\\0&G\end{pmatrix}
\begin{pmatrix}z\,\mathrm{pr}\,P\\ \tilde x\end{pmatrix}R,\qquad z\,\mathrm{pr}\,P=(z\,|\,0\,|\,0),
\]
so $\Phi'_z(x)=\epsilon_w\,\chi_w^{-1}\bigl(\det\binom{(z\,|\,0\,|\,0)}{\tilde x}\bigr)$ (with value $0$
when the matrix is not in $\auGL_{N+2}(O_w)$), where $\epsilon_w=\chi_w^{-1}(\det G\det R)$ is a constant of
modulus one (it depends on the frames and, through $\chi=\chi_S$, on $\nu$); it goes into $l_\beta$. The Heisenberg operators and the
isotropic subspaces of Definition~\ref{pct:def-line} are, in the coordinates $\tilde x$, the
translations and characters in the blocks indicated. As in Step~2, the zero term along an isotropic $I$ is $\Phi\mapsto(x\mapsto
\int_{I_w}\Phi(x+i)\,di)$ on the annihilator of $I_{\bar w}$, which is the set of matrices vanishing
in the blocks dual to those of $I$; $\psrho(y)$ is translation by $y_w$ times the character
$[\,\cdot\,,y_{\bar w}]$. For $\beta=\mathrm t$, $I=I_{\mathrm t}=(V,f_A)\oplus(e_X,W)\oplus(e_X,f_A)
\oplus(f_X,f_A)$, the dual blocks are $(V,e_A),(f_X,W),(f_X,e_A),(e_X,e_A)$, and the translation puts
$\alpha_w$ (times the unit component of $\alpha_0$) in the block $(e_X,e_A)$ and multiplies by
$\lambda_w(\epsilon'_w\overline{\alpha_{\bar w}}\,t)^{\pm1}$, $t$ the entry $(f_X,f_A)$, where
$\lambda_w(t)=\psi_w(\jmath(a\bar a)_wt)$ is the character of conductor $O_w$ of
Lemma~\ref{pgx:dict}(b) (the Heisenberg pairing in the coordinates $\tilde x=a^{-1}x$) and
$\epsilon'_w$ is the conjugate of the unit component of $\alpha_0$ at $\bar w$; replacing $c$ by
$\pm\epsilon'_wc$ in Lemma~\ref{pct:gauss} (applied to $\lambda_w^{\pm1}$) multiplies the result by
$\chi_w(\pm\epsilon'_w)$, of modulus one, which goes into $l_\beta$; the $\chi_w$-values of these unit components, which arise when they are removed from the
determinant, have modulus one and go into $l_\beta$. The resulting matrix is
\[
M=\begin{pmatrix}z&0&0\\ y&0&p\\ q&\alpha_w&s\\0&0&t\end{pmatrix}
\]
with $y\in\pcMat_{r\times N}$ the variable of $\auHom(W,V)$ and $p,q,s,t$ integrated over $O_w$
(the measure of $O_w$ being $1$). Expanding along the column $e_A$, $\det M=\pm\alpha_wt\det\binom zy$,
and $M\in\auGL_{N+2}(O_w)$ iff all entries are integral and $\alpha_w,t,\det\binom zy$ are units. So
the integrand is $\chi_w^{-1}(\pm1)\chi_w^{-1}(\alpha_w)\mathbf 1_{O_w^\times}(\alpha_w)\cdot
\Phi^{\mathrm{row}}_z(y)\cdot\chi_w^{-1}(t)\mathbf 1_{O_w^\times}(t)\lambda_w(\epsilon'_w\overline{\alpha_{\bar w}}t)^{\pm1}\cdot
\mathbf 1_{O_w}(p,q,s)$, and Lemma~\ref{pct:gauss} gives the factor
$g_w\chi_w(\alpha_w)^{-1}\chi_w(\varpi_w\overline{\alpha_{\bar w}})$ times $\Phi^{\mathrm{row}}_z(y)$ when
$\alpha_w$ and $\varpi_w\overline{\alpha_{\bar w}}$ are units, and $0$ otherwise. For $\beta=\mathrm f$
the same computation with $I_{\mathrm f}$ puts $t$ in the corner $(e_X,e_A)$ and $\alpha_w$ in the
corner $(f_X,f_A)$; the determinant and the result are the same. Put $g_S=\prod_{w\in S}g_w$.

\emph{Invariance under $\pctK'_f$.} At $w\in S$, $k\in\auGL_{r+2}(O_w)$ acts by $\Phi\mapsto
\gamma_{W'}(\det k)\,\Phi(k^{-1}\cdot)$ (W5), and $\Phi'_z(k^{-1}x)=\chi_w(\det k)\Phi'_z(x)$ (the
row datum is matched, Lemma~\ref{pcx:std-matched}, and Lemma~\ref{pcx:rightinv}), while
$\gamma_{W'}(\det k)=\chi_w(\det k)^{-1}$ on units ($\upsilon$ being trivial on the determinants of
integral points); elsewhere $\pctK'_f$ fixes the vacuum by (W4).
\end{proof}

\emph{Step 5: Tate's integral.} The function $\lambda_\varpi$ satisfies $\lambda_\varpi(\alpha)=
\lambda_1(\alpha c)$ for the finite idele $c$ equal to $\bar\varpi_w$ at the places $\bar w\in\bar S$
and to $1$ elsewhere, and $|c^{-1}|=Q$. On the units $\widehat\auO_E^\times$ of the finite ideles,
$\lambda_1(u)\tilde\nu_f(u)=1$ (at $w\in S$: $\tilde\nu_w(u_w)\tilde\nu_{\bar w}(u_{\bar w})=
\chi_w(u_w)\chi_w(\overline{u_{\bar w}})^{-1}$; elsewhere $\tilde\nu$ is unramified), and $\lambda_1(a)=0$
unless $a$ is integral and a unit at $S\cup\bar S$.

\begin{lemma}\label{pct:fintate}
Let $\rho$ be a Haar measure on $\auA_{E,f}^\times$. For $x>1$,
$\int_{\auA_{E,f}^\times}\lambda_1(a)\tilde\nu_f(a)|a|^x\,d\rho(a)=L^S(x,\nu)\,\rho(\widehat\auO_E^\times)$,
absolutely convergent.
\end{lemma}
\begin{proof}
Write $a=u\varpi^m$ with $u$ a unit and $\varpi^m=\prod_{v\notin S\cup\bar S}\varpi_v^{m_v}$, $m_v\ge0$
almost all $0$. The integrand is $\tilde\nu(\varpi^m)|\varpi^m|^x$, so the integral is
$\rho(\widehat\auO_E^\times)\prod_{v\notin S\cup\bar S}(1-\tilde\nu(\varpi_v)q_v^{-x})^{-1}$.
\end{proof}

\begin{lemma}[Mellin transforms]\label{pct:mellin}
Let $T$, $\mathsf n>0$ be measurable on a measure space with $\int|T|\mathsf n^R<\infty$ and
$\int|T|\mathsf n^{1-R}<\infty$ for every $R>1$. Then $\int|T|\mathsf n^x<\infty$ for every real $x$
and $x\mapsto\int T\mathsf n^x$ is real-analytic on $\R$.
\end{lemma}
\begin{proof}
$\mathsf n^x\le\mathsf n^{1-R}+\mathsf n^R$ for $x\in[1-R,R]$. Near $x_0$, with $0<\delta<1$,
$|\log\mathsf n|^k\mathsf n^{x_0}\le k!\delta^{-k}(\mathsf n^{x_0+\delta}+\mathsf n^{x_0-\delta})$, so the
series $\sum_k\frac{(x-x_0)^k}{k!}\int T(\log\mathsf n)^k\mathsf n^{x_0}$ converges absolutely for
$|x-x_0|<\delta$ and equals $\int T\mathsf n^x$.
\end{proof}

\begin{lemma}[Poisson on the line, in the variable $a$]\label{pct:poisson-a}
For $f=[v,\Phi'_z]_S$, $a\in\auA_E^\times$ and $(g,h,\kappa)$:
$|a|\allowbreak \sum_{\xi\in E^\times}\allowbreak \Lambda_{\mathrm t}(f;\xi a;g,h,\kappa)=\sum_{\xi\in E^\times}
\allowbreak \Lambda_{\mathrm f}(f;\xi\bar a^{-1};g,h,\kappa)$.
\end{lemma}
\begin{proof}
By Lemma~\ref{pct:factor}, $\Lambda_\beta(f;0;\cdot)=0$, since $\lambda_\varpi(0)=0$. Apply
Lemma~\ref{pct:poisson} to $F=\omega(m_A(1,g),\ell_X(a,h))\omega(1,\kappa)f$ and use
Lemma~\ref{pct:levi} with $u=1$, $c=a$: the terms become $\tilde\nu(a)|a|^{-N/2}\Lambda_{\mathrm t}(\xi a)$
and $\tilde\nu(a)|a|^{-N/2-1}\Lambda_{\mathrm f}(\xi\bar a^{-1})$; cancel $\tilde\nu(a)|a|^{-N/2-1}$.
\end{proof}

\begin{proof}[Proof of Theorem~\ref{pct:plane}]
Let $K_0=c^{-1}\alpha_0\in\auA_E^\times$ and $\mathcal K(g,h;w)=\theta^{W,V}_{\gamma_W,\gamma_W}(g,h;
[w,\Phi^{\mathrm{row}}_z]_S)$. By Lemma~\ref{pct:factor} (the finite part of $\kappa$ acting trivially),
for $\alpha\in\auA_E^\times$,
\begin{multline*}
\tilde\nu(\alpha K_0)|\alpha K_0|^x\Lambda_\beta(f;\alpha K_0;g,h,\kappa)\\
=\Bigl[\tilde\nu_\infty(\alpha_\infty)|\alpha_\infty|^x\mathcal K\bigl(g,h;\pctL^\infty_\beta(\alpha_\infty)
\omega'_\infty(1,\kappa_\infty)v\bigr)\Bigr]\cdot\Bigl[l_\beta g_S\tilde\nu(K_0)|K_0|^x\,
\lambda_1(\alpha_f)\tilde\nu_f(\alpha_f)|\alpha_f|^x\Bigr].
\end{multline*}
\emph{Convergence for $x>1$, side $\mathrm t$.} In $\pctZ(f,g,u;x)$ substitute $\xi au=\alpha K_0$;
the sum over $E^\times$ and the integral over $E^\times\backslash\auA_E^\times$ combine into an integral
over $\auA_E^\times=(E\otimes\R)^\times\times\auA_{E,f}^\times$. By Tonelli the integral of the absolute
values is the product of $\int_{\auA_{E,f}^\times}|\lambda_1\tilde\nu_f||\cdot|^x\,d\rho<\infty$
(Lemma~\ref{pct:fintate}) and of $\int|\alpha_\infty|^x|\mathcal K(g,h;\pctL^\infty_{\mathrm t}(\alpha_\infty)
\omega'_\infty(1,\kappa)v)|\,d\alpha_\infty\,dh\,d\kappa$, which is finite: by
Proposition~\ref{pcx:decay}(a) ($\Phi^{\mathrm{row}}_z$ is supported on matrices of rank $r$),
$|\mathcal K(g,h;w)|\le C\,\mathrm{ht}(g)^B\mathrm{ht}(h)^{-N_0}\|w\|_{M'}$, and Lemma~\ref{pct:slicedecay}
and (R1) bound the rest. Hence the integral converges absolutely and, by Fubini and
Lemma~\ref{pct:fintate},
\[
\pctZ(f,g,u;x)=l_{\mathrm t}g_S\tilde\nu(K_0)|K_0|^x\rho(\widehat\auO_E^\times)L^S(x,\nu)\int_{[\auU(V)]}
\mathcal K(g,h;\pctcone_xv)\,dh,
\]
where the cone integral was taken inside $\mathcal K$ (a continuous linear functional of $w$ for
the norm $\|\cdot\|_{M'}$, Proposition~\ref{pcx:decay}(a)). The last integral is
$\Theta^{\mathrm L}_{W,V}(\pctcone_xv,\Phi^{\mathrm{row}}_z)(g)$.

\emph{Side $\mathrm f$.} By Lemma~\ref{pct:poisson-a}, the integrand of $\pctZ(f,g,u;1-x)$ is
$\tilde\nu(au)|au|^{-x}\sum_\xi\Lambda_{\mathrm f}(f;\xi P_{\mathrm f}(au);\cdot)$, and
$\tilde\nu(P_{\mathrm f}(y))=\tilde\nu(y)$, $|P_{\mathrm f}(y)|=|y|^{-1}$. The substitution
$y'=P_{\mathrm f}(au)$ preserves the Haar measure of $E^\times\backslash\auA_E^\times$, and the same
computation with $\beta=\mathrm f$ shows absolute convergence for $x>1$ and
$\pctZ(f,g,u;1-x)=l_{\mathrm f}g_S\tilde\nu(K_0)|K_0|^x\rho(\widehat\auO_E^\times)L^S(x,\nu)\,
\Theta^{\mathrm L}_{W,V}(\pctcone^\vee_xv,\Phi^{\mathrm{row}}_z)(g)$.

\emph{(a).} The integrand of $\pctZ$ is $T\cdot|au|^x$ with $T$ independent of $x$; by the two
previous paragraphs Lemma~\ref{pct:mellin} applies. Independence of $u$ was noted before the
theorem.

\emph{(b).} Lemma~\ref{pct:unfold}(iii), with the absolute convergence at $d+1$ from (a).

\emph{(c).} For $x>1$ these are the two formulas just proved, with
$\varepsilon_\beta=c_Xl_\beta g_S\sqrt Q\,\tilde\nu(K_0)\rho(\widehat\auO_E^\times)$: indeed
$|K_0|^x=Q^x|\alpha_0|^x$ and $|g_S|=Q^{-1/2}$, so $|\varepsilon_\beta|=|c_Xl_\beta|\rho(\widehat
\auO_E^\times)=:C_\beta$ does not depend on $\nu$, and $c_X\,l_\beta g_S\tilde\nu(K_0)|K_0|^x=
\varepsilon_\beta P(x)$. At $x=1$ both sides are continuous on $[1,\infty)$: the left side by (a),
$L^S(x,\nu)=L(x,\tilde\nu)$ because $L(s,\tilde\nu)$ is entire ($\tilde\nu$ is a non-trivial unitary
Hecke character, Lemma~\ref{pcx:Lpos}; Hecke, see Tate's thesis, Main Theorem~4.4.1), and the
theta integral of the cone integral because it is, by the computation above, a Mellin transform
$\int\mathcal K(g,h;\cdots)|\alpha_\infty|^x$ converging absolutely for all $x>0$,
hence continuous on $(0,\infty)$ (the proof of Lemma~\ref{pct:mellin} applies on every compact
subinterval).
\end{proof}

\subsubsection{Strong approximation and the constancy in equal rank}\label{pct:ss-const}

Let $Y=Y_0\oplus\pcxH$ be an adapted splitting, $\pcxH=Ee\oplus Ef$, $\langle e,f\rangle=1$, the form
linear in the first variable. For $c\in Y_0$ and $\beta\in E$ with $\beta+\bar\beta=-\langle c,c\rangle$
(over $E$ or over $\auA_E$) let $n(c,\beta)$ and $n^-(c,\beta)$ be the isometries
\[
n(c,\beta)\colon e\mapsto e,\ y\mapsto y-\langle y,c\rangle e,\ f\mapsto f+c+\beta e;\qquad
n^-(c,\beta)\colon f\mapsto f,\ y\mapsto y-\langle y,c\rangle f,\ e\mapsto e+c+\beta f
\]
($y\in Y_0$); one checks directly that they preserve the form. The $n(c,\beta)$ form the unipotent
radical $\pctN$ of the stabiliser $\pctP$ of $Ee$, the $n^-(c,\beta)$ that of the opposite parabolic,
and $m(t,\delta)$ ($t\in\auA_E^\times$, $\delta\in\auU(Y_0)$) acts by $e\mapsto te$, $f\mapsto\bar t^{-1}f$,
$y\mapsto\delta y$; every element of $\pctP$ is $n\,m(t,\delta)$, uniquely. For a set $R$ of finite
places of $E$ and $x\in\auU(Y)(\auA)$ let $\ell_R(x)$ be the element with the components of $x$ at the
places of $E^+$ below $R\cup\bar R$ and $1$ elsewhere (also at $\infty$); $\ell_R$ is a homomorphism,
and $\auU(Y)(F_R)=\ell_R(\auU(Y)(\auA))$. Let $\pctK=\prod_v\pctK_v$ be the standard maximal compact
subgroup (lattice stabilisers at finite $v$) and $\pctK^S=\{k\in\prod_{v<\infty}\pctK_v:\ell_S(k)=1\}$.

\begin{lemma}[Quasi-reflections from radicals]\label{pct:qref}
Let $v\in Y_0(E)$, $\sigma\in E^\times$ with $\sigma+\bar\sigma=\langle v,v\rangle$, and
$\rho(y)=y-v\sigma^{-1}\langle y,v\rangle$ the quasi-reflection. Let $j\in E^\times$ with $\bar j=-j$,
$\beta_1=-\sigma-j^{-1}$, $\lambda=(j\bar\sigma)^{-1}$. Then
\[
u=n^-(v,\beta_1)\,n(0,j)\,n^-(-\lambda v,\lambda\bar\beta_1)\in\pctP(E),\qquad u=n\,m(\lambda,\rho)\
\text{ for some }n\in\pctN(E).
\]
\end{lemma}
\begin{proof}
The parameters are admissible: $\beta_1+\bar\beta_1=-(\sigma+\bar\sigma)-(j^{-1}+\bar j^{-1})=
-\langle v,v\rangle$, $j+\bar j=0$, and for $\beta_3=\lambda\bar\beta_1$, $c_3=-\lambda v$ one needs
$\beta_3+\bar\beta_3=-\langle c_3,c_3\rangle$, i.e.\
$\lambda\bar\beta_1+\bar\lambda\beta_1=\lambda\bar\lambda(\beta_1+\bar\beta_1)$, i.e.\
$\bar\beta_1\bar\lambda^{-1}+\beta_1\lambda^{-1}=\beta_1+\bar\beta_1$. Here $\lambda^{-1}=j\bar\sigma=
1-j\bar\beta_1$ (as $1-j\bar\beta_1=1+j\bar\sigma+j\bar j^{-1}$ and $j\bar j^{-1}=-1$) and
$\bar\lambda^{-1}=1-\bar j\beta_1=1+j\beta_1$; so the left side is $\bar\beta_1(1+j\beta_1)+\beta_1(1-j\bar\beta_1)=\beta_1+\bar\beta_1$. Now compute. With
$c_3=-\lambda v$, $\beta_3=\lambda\bar\beta_1$: $n^-(c_3,\beta_3)e=e+c_3+\beta_3f$; then $n(0,j)$ gives
$(1+j\beta_3)e+c_3+\beta_3f$, and $1+j\beta_3=\lambda$ (as $\lambda(1-j\bar\beta_1)=1$); then
$n^-(v,\beta_1)$ gives $\lambda(e+v+\beta_1f)+c_3-\langle c_3,v\rangle f+\beta_3f=\lambda e+(\lambda\beta_1+
\lambda\langle v,v\rangle+\lambda\bar\beta_1)f=\lambda e$. So $ue=\lambda e$ and $u\in\pctP(E)$. For
$y\in Y_0$: $n^-(c_3,\beta_3)y=y+\bar\lambda\langle y,v\rangle f$, then $n(0,j)$ gives
$y+\bar\lambda\langle y,v\rangle(f+je)$, and the $Y_0$-component after $n^-(v,\beta_1)$ is
$y+j\bar\lambda\langle y,v\rangle v$. Since $j\bar\lambda=j(\bar j\sigma)^{-1}=-\sigma^{-1}$, this is
$\rho(y)$. Hence the Levi component of $u$ is $m(\lambda,\rho)$.
\end{proof}

\begin{lemma}[Constancy by strong approximation]\label{pct:SA}
Let $S$ be a non-empty finite set of finite places of $E$, and $D\colon\auU(Y)(\auA)\to\C$ continuous
with: (i) $D(\gamma x)=D(x)$ for $\gamma\in\auU(Y)(E)$; (ii) $D(xs)=D(x)$ for $s\in\auU(Y)(F_S)$;
(iii) $D(xk)=D(x)$ for $k\in\pctK^S$; (iv) if $D(nx)=D(x)$ for all $n\in\pctN(\auA)$ and all $x$,
then $D(m(t,y))$ depends only on $y$. Then $D$ is constant.
\end{lemma}
\begin{proof}
\emph{(1)} If $a,b\in\auU(Y)(\auA)$ agree at all places not below $S\cup\bar S$, then $D(a)=D(b)$:
$a=b\,\ell_S(b^{-1}a)$ and (ii).

\emph{(2)} $D$ is left invariant under $\pctN(\auA)$ and $\pctN^-(\auA)$. The parameters $(c,\beta)$
form a closed subset of $Y_0(\auA)\times\auA_E$, and by strong approximation for the additive
group ($E+\prod_{w\in S\cup\bar S}E_w$ is dense in $\auA_E$, as $S\ne\emptyset$; Cassels--Fr\"ohlich,
\emph{Algebraic Number Theory}, strong approximation) applied to the coordinates of $c$ and to
the solutions $\beta$, the admissible parameters of the form
$(c_E+c_S,\beta_E+\beta_S)$, with $(c_E,\beta_E)$ rational and $(c_S,\beta_S)$ supported at $S\cup\bar S$,
are dense. Indeed, let $(c,\beta)$ be admissible, $\beta+\bar\beta=-\langle c,c\rangle$. Choose $c_E\in Y_0(E)$
and $c_S$ supported at $S\cup\bar S$ with $c'=c_E+c_S$ close to $c$. Then
$\delta=\langle c',c'\rangle-\langle c_E,c_E\rangle=\langle c_E,c_S\rangle+\langle c_S,c_E\rangle+\langle c_S,c_S\rangle\in\auA_{E^+}$
is supported at the places below $S$; put $\beta'_E=-\frac12\langle c_E,c_E\rangle\in E^+$ and
$\beta'_S=-\frac12\delta$, so that $(c',\beta'_E+\beta'_S)$ is admissible. The difference
$\beta-\beta'_E-\beta'_S$ is $\tau+\eta$ with $\tau=-\frac12(\langle c,c\rangle-\langle c',c'\rangle)$, small,
and $\eta$ of trace zero. The trace-zero elements form the line $j\auA_{E^+}$ for a fixed $j\in E^\times$
with $\bar j=-j$, and by strong approximation for $\auA_{E^+}$ (the set of places below $S$ is non-empty)
there are $\eta_E\in jE^+$ and $\eta_S\in j\auA_{E^+}$ supported below $S$ with $\eta-\eta_E-\eta_S$ small.
Then $(c',\beta'_E+\eta_E+\beta'_S+\eta_S)$ is admissible, of the required form, and close to $(c,\beta)$. For those, $n(c_E+c_S,\beta_E+\beta_S)x$ agrees with $n(c_E,\beta_E)x$ off $S\cup\bar S$, so
(1) and (i) give $D(n(\cdot)x)=D(x)$; continuity extends this to all of $\pctN(\auA)$. The same for
$\pctN^-$.

\emph{(3)} By (2) and (iv) there is $D_0$ with $D(m(t,y))=D_0(y)$.

\emph{(4)} Let $R$ be a set of finite places. The set $\Gamma_R$ of $\delta\in\auU(Y_0)(E)$ for which
there is $t\in\auA_E^\times$ with $D(\ell_R(m(t,\delta))x)=D(x)$ for all $x$ is a subgroup ($\ell_R$
and $m$ are homomorphisms). It contains the quasi-reflections: with $u=n_1n_2n_3=n\,m(\lambda,\rho)$
from Lemma~\ref{pct:qref}, $\ell_R(m(\lambda,\rho))=\ell_R(n)^{-1}\ell_R(n_1)\ell_R(n_2)\ell_R(n_3)$, a
product of elements of $\pctN(\auA)$ and $\pctN^-(\auA)$ (cutting the parameters to $R$ gives
admissible adelic parameters), under which $D$ is left invariant by (2). As $\auU(Y_0)(E)$ is
generated by quasi-reflections (Witt's theorem for hermitian forms, proved by quasi-reflections;
see e.g.\ Hahn--O'Meara, \emph{The Classical Groups and K-Theory}),
$\Gamma_R=\auU(Y_0)(E)$.

\emph{(5)} For finite $R$, all $t$ and all $y\in\auU(Y_0)(\auA)$: $D(\ell_R(m(t,y)))=D(1)$. By weak
approximation for $\auU(Y_0)$ (a rational variety by the Cayley transform; Platonov--Rapinchuk,
\emph{Algebraic Groups and Number Theory}, strong approximation) the elements $\gamma k$ with
$\gamma\in\auU(Y_0)(E)$ and $\ell_R(k)=1$ are dense, and $y\mapsto D(\ell_R(m(t,y)))$ is continuous, so
it suffices to take $y=\gamma k$, for which $\ell_R(m(t,y))=\ell_R(m(t,\gamma))=m(\ell_R(t),\ell_R(\gamma))$.
By (3) its value is $D_0(\ell_R(\gamma))=D(\ell_R(m(a,\gamma)))$ for the $a$ given by (4), which is
$D(1)$ (take $x=1$ in (4)).

\emph{(6)} By weak approximation for $\auU(Y)$ at the infinite places, $\auU(Y)(E)\,\auU(Y)(\auA_f)$ is
dense; by (i) and continuity it suffices to show $D(k)=D(1)$ for $k\in\auU(Y)(\auA_f)$. Let $T$ be the
cofinite set of finite places with $k_v\in\pctK_v$, and use the Iwasawa decomposition
$\auU(Y)(\auA)=\pctP(\auA)\pctK$ (valid for the lattice stabilisers since the splitting is adapted:
the lattice contains the hyperbolic pair $e,f$): $k\kappa=n\,m(t,y)$ with $\kappa\in\pctK$. The element
$b=k^{-1}\ell_{T^c}(k\kappa)$ has components $\kappa_v$ at $v\in T^c$, $k_v^{-1}\in\pctK_v$ at $v\in T$ and
$1$ at $\infty$; so $b=s\,k'$ with $s=\ell_S(b)\in\auU(Y)(F_S)$ and $k'\in\pctK^S$. By (ii), (iii),
$D(k)=D(kb)=D(\ell_{T^c}(k\kappa))=D(\ell_{T^c}(n)\,\ell_{T^c}(m(t,y)))=D(\ell_{T^c}(m(t,y)))=D(1)$, by
(2) and (5) with $R=T^c$ finite.
\end{proof}

\begin{proof}[Proof of Lemma~\ref{pcx:const}]
Let $Y$ have rank $2k$, $Y=Y_0\oplus\pcxH$ adapted, $v$ with finite moments, and
$F(x)=\Theta^{\mathrm L}_{Y,Y}(v,\Phi)(x)$ with $\Phi=\Psi^{\mathrm{col}}_1=\chi^{-1}\circ\det$ on
$\auGL_{2k}(O)$ (zero elsewhere), which is the row datum $\Phi^{\mathrm{row}}_\emptyset$ for $b=0$. We
check the hypotheses of Lemma~\ref{pct:SA} for $D=F$.

(i) and continuity: the kernel is left invariant in $x$ under $\auU(Y)(E)$ by (W1) ($\gamma_Y$ is
trivial on rational determinants), and the integral over $[\auU(Y)]$ converges uniformly for $x$ in
compact sets (Corollary~\ref{pcx:conv}(i): $\Phi$ is supported on invertible matrices).

(ii): Let $g\in\auGL_{2k}(F)=\auU(Y)(F_S)$ (read in the frames, $S$ consisting of split places).
In the notation of \S\ref{pcx:ss-match} with $b=0$, $a=2k$: $\mathrm{Inj}_0$ is a point, and
$\mathsf T_1=1$. For $\Phi_g=\xi(\det g)^{-1}\Phi(\cdot\,g^{-1})$ one computes, substituting
$r'=rxg^{-1}\in\auGL_{2k}(O)$, $\mathsf S_{\Phi}(x)=\xi(\det x)^{-1}\int_{\auGL_{2k}(O)}\xi(\det r')
\chi^{-1}(\det r')\,dr'=\xi(\det x)^{-1}$ and $\mathsf S_{\Phi_g}(x)=\xi(\det g)^{-1}\mathsf S_\Phi(xg^{-1})=
\xi(\det x)^{-1}$ for $x\in\auGL_{2k}(F)$ (using $\xi=\chi$ on $O^\times$). Since $k^{[a]}=k^{-1}$ for
$b=0$, both $(\Phi,1)$ and $(\Phi_g,1)$ are same-section pairs (Definition~\ref{pcx:def-section}).
Lemma~\ref{pcx:linfactor}, applied to $\Lambda=\Theta^{\mathrm L}_{Y,Y}(v,\cdot)$ (whose hypothesis is
Lemma~\ref{pcx:Sact}(a)), gives $\Lambda(\Phi_g)=\mathcal T(1)=\Lambda(\Phi)$. With
Lemma~\ref{pcx:Sact}(b) for $b=0$: $F(xg)=\Lambda(\Phi_g)(xg)=\Lambda(\Phi)(x)=F(x)$.

(iii): For $k\in\pctK^S$, $\omega_{\gamma_Y,\gamma_Y}(k,1)[v,\Phi]_S=[v,\Phi]_S$: the components of $k$
at $S\cup\bar S$ and $\infty$ are $1$, the vacuum is fixed by (W4) at the other finite places, and
$\gamma_Y(\det k)=1$ ($\nu$ and $\upsilon$ are trivial on the determinants of integral points off
$S\cup\bar S$). Hence $F(xk)=F(x)$.

(iv): If $F$ is left invariant under $\pctN(\auA)$, then $F=F_A$ (the constant term along the line
$Ee$), and by the last assertion of Theorem~\ref{pct:plane} applied with $W'=V'=Y$, $W=V=Y_0$,
$b=0$ ($r=N=2k-2$, $d=k-1\ge0$), $F(m(t,y))=F_A(m(t,y))$ does not depend on $t$.

By Lemma~\ref{pct:SA} ($S\ne\emptyset$ as $\pcm\ne(1)$), $F$ is constant.
\end{proof}

\subsubsection{The image of the cone integral}\label{pct:ss-cone}

Keep the set-up of \S\ref{pct:ss-setup}. At the infinite places write $Z_X$ for the archimedean
zero term along $(e_X,W)=\auHom(W,Ee_X)\subset\auHom(W,V')$, from Fock vectors of
$\auHom(W,V')\otimes\R$ to those of $\auHom(W,V)\otimes\R$; thus $\pctL^\infty_{\mathrm t}(a)=
Z_X\circ\pctsl_{A,\infty}(ae_X)$. The following two facts about the real mixed models are used.

\begin{lemma}[Real mixed models]\label{pct:realmixed}
\textup{(a)} For every Fock vector $v$ of $\auHom(W,V)\otimes\R$ with finite moments there is a Fock
vector $v_B$ of $\auHom(W,V')\otimes\R$ with finite moments and $Z_Xv_B=v$.
\textup{(b)} Let $\beta$ be a smooth compactly supported function on $V'(E\otimes\R)$, $x\mapsto k(x)$ a
smooth map to $\auU(V')(E^+\otimes\R)$, and $w$ a Fock vector of $\auHom(W,V')\otimes\R$ with finite
moments. There is a Fock vector $v'$ of $\auHom(W',V')\otimes\R$ with finite moments such that
$\pctsl_{A,\infty}(x)v'=\beta(x)\,\omega_\infty(1,k(x))w$ for all $x$.
\end{lemma}
\begin{proof}
We use the mixed model and its properties (M1)--(M4) of \S\ref{pcu:ss-mixed} (with (M1) proved in
\S\ref{pcu:ss-M1}, and Lemma~\ref{pcu:slicedecay}).

(b) This is (M1) for the adapted splitting $A$ of the source space $W'$ (with $p=1$, $T=V'$): the
slice $\mathrm{sl}_A(x)$ there is the value of the mixed model at the value $x$ on $e_A$, i.e.\
$\pctsl_{A,\infty}(x)$ of Definition~\ref{pct:def-line} (translation by $x$ in the column $e_A$
followed by the zero term along the column $f_A$, which is evaluation of the Schr\"odinger variable,
\S\ref{pcu:ss-mixed}); take $\tau=\beta$ and $u(x)=(1,k(x))$.

(a) $Z_X$ is a zero term along a line of the \emph{target}. The flip $\mathrm{Fl}$ of
\S\ref{pcx:ss-flip} (fact (W6)), at the infinite places, is induced by the measure-preserving
isometry $w\mapsto w^\dagger$ of $\auHom(W,V')_\R$ onto $\auHom(V'^-,W)_\R$, so it preserves the moments
$\|\cdot\|_M$ and intertwines the Heisenberg operators; it carries $(e_X,W)=\auHom(W,Ee_X)$ onto the
maps $V'^-\to W$ vanishing on $(Ee_X)^{\perp}$, i.e.\ the isotropic subspace attached to the adapted
splitting $V'^-=V^-\oplus\pcxH_X$ of the \emph{source} $V'^-$, and hence carries $Z_X$ to the zero term
$\mathrm{sl}(0)$ along that splitting (both are characterised by (W2) through the same isotropic
subspace). So it suffices to show that $\mathrm{sl}(0)$ is onto on vectors with finite moments, for a
source splitting. Given $v$ with finite moments, let $\tau$ be smooth with compact support and
$\tau(0)=1$; by (M1) with $u\equiv1$ there is $\Lambda$ with finite moments and
$\mathrm{sl}(x)\Lambda=\tau(x)v$; then $\mathrm{sl}(0)\Lambda=v$. Put $v_B=\mathrm{Fl}^{-1}\Lambda$.
\end{proof}

\begin{lemma}[Equivariance of the real zero term]\label{pct:zeroterm-equiv}
Let $p\in\auU(V')(E^+\otimes\R)$ with $pe_X=ae_X$, $a\in(E\otimes\R)^\times$. There is
$h\in\auU(V)(E^+\otimes\R)$ such that $Z_X\,\omega_\infty(1,p)=\tilde\nu_\infty(a)^{-1}|a|^{N/2}\,
\omega_\infty(1,h)\,Z_X$ on Fock vectors with finite moments, the representations being the
archimedean components of $\omega_{\gamma_W,\gamma_{W'}}$ on $\auHom(W,V')$ and of
$\omega_{\gamma_W,\gamma_W}$ on $\auHom(W,V)$.
\end{lemma}
\begin{proof}
$p$ stabilises $Ee_X$, hence $(e_X,W)$, and induces on $(e_X,W)^\perp/(e_X,W)=\auHom(W,V)$ the
action of an element $h$ of $\auU(V)$ (the image of $p$ in the Levi quotient of the stabiliser of
$Ee_X$). The archimedean form of (W2) gives $Z_X\omega(1,p)=\chi(p)\omega(1,h)Z_X$ with
$|\chi(p)|=|\det p|_{(e_X,W)}|^{1/2}=|a|^{N/2}$ ($p$ acts on $\auHom(W,Ee_X)$ by the scalar $a$). The
unitary part of $\chi(p)$, together with the change of the archimedean splitting characters from
$\gamma_{W'}$ to $\gamma_W$ on the target, is the character $\prod_\tau(a_\tau/|a_\tau|)^{-2\mu_\tau}=
\tilde\nu_\infty(a)^{-1}$ of the weight (the determinant of $p$ is $(a/\bar a)\det h$).
\end{proof}

\begin{lemma}[Smooth Witt]\label{pct:smoothwitt}
There are a neighbourhood $\mathcal U$ of $e_X$ in $V'(E\otimes\R)$ and a smooth map
$k\colon\mathcal U\to\auU(V')(E^+\otimes\R)$ with $k(x)e_X=x$ for every isotropic $x\in\mathcal U$.
\end{lemma}
\begin{proof}
Place by place. The set $\mathcal C$ of non-zero isotropic vectors of $V'_\tau\cong\C^{p,q}$ is a
smooth submanifold, and $\auU(p,q)$ acts transitively on it (Witt's theorem over $\C/\R$). The orbit
map $\pi\colon k\mapsto ke_X$ is a smooth surjective equivariant map of homogeneous spaces, hence a
submersion, and has a smooth local section $s$ near $e_X$ with $s(e_X)=1$. Let
$P\colon\mathcal U\to\mathcal C$ be a smooth retraction of a neighbourhood of $e_X$ onto $\mathcal C$
(a tubular neighbourhood), and put $k=s\circ P$; for isotropic $x$, $P(x)=x$.
\end{proof}

\begin{proposition}[The image of the cone integral]\label{pct:coneimage}
Let $N$ be even. There is a function $B>0$ on $\R$, real-analytic,
such that every Fock vector $v$ of $\auHom(W,V)\otimes\R$ with finite moments has a Fock vector $v'$
of $\auHom(W',V')\otimes\R$ with finite moments satisfying, for all $x>0$, all tame $\nu$ of weight
$\mu$, all $z\in\pcMat_{b\times N}(O)$ and all $g$,
\[
\Theta^{\mathrm L}_{W,V}(\pctcone_xv',\Phi^{\mathrm{row}}_z)(g)=B(x)\,\Theta^{\mathrm L}_{W,V}(v,
\Phi^{\mathrm{row}}_z)(g).
\]
In particular, for fixed $x>0$ there is $v'$ with $\Theta^{\mathrm L}_{W,V}(\pctcone_xv',\cdot)=
\Theta^{\mathrm L}_{W,V}(v,\cdot)$ on all these data (replace $v$ by $B(x)^{-1}v$).
\end{proposition}
\begin{proof}
Let $\mathcal U$, $k$ be as in Lemma~\ref{pct:smoothwitt}, and $\beta\ge0$ smooth with compact support in
$\mathcal U$, $\beta(e_X)=1$, vanishing unless $\frac12|e_{X,\tau}|^2\le|x_\tau|^2\le2|e_{X,\tau}|^2$ for
every $\tau$. Lemma~\ref{pct:realmixed}(b) needs $k$ on all of $V'(E\otimes\R)$; we extend it. Since
$k(e_X)=1$, after shrinking $\mathcal U$ we may assume $k(\mathcal U)\subset\exp(\mathcal B)$, where $\mathcal B$
is a ball about $0$ in the Lie algebra $\mathfrak g$ of $\auU(V')(E^+\otimes\R)$ on which $\exp$ is a
diffeomorphism onto an open set, so $k=\exp\circ X$ with $X\colon\mathcal U\to\mathcal B$ smooth. Take
$\beta$ as above with support in this smaller $\mathcal U$, a smooth $\phi$ with compact support in $\mathcal U$
and $\phi=1$ on $\mathrm{supp}\,\beta$, and replace $k$ by $\exp(\phi X)$ (extended by $1$ outside
$\mathcal U$): it is smooth on $V'(E\otimes\R)$ and equal to $k$ on $\mathrm{supp}\,\beta$, which is all
that enters below. By Lemma~\ref{pct:realmixed}(a) write $v=Z_Xv_B$, and by (b) choose $v'$ with
$\pctsl_{A,\infty}(y)v'=\beta(y)\omega_\infty(1,k(y))v_B$. Let $\kappa\in\pctK'_\infty$, $a\in(E\otimes\R)^\times$,
$y=\kappa^{-1}ae_X$ (isotropic). By (W1), (W2) at $\infty$ (the slice is equivariant for the action
of $\auU(V')$ on the target), $\pctsl_{A,\infty}(ae_X)\omega'_\infty(1,\kappa)v'=\omega_\infty(1,\kappa)
\pctsl_{A,\infty}(y)v'=\beta(y)\,\omega_\infty(1,\kappa k(y))v_B$. If $\beta(y)\ne0$, then
$\kappa k(y)e_X=\kappa y=ae_X$, and Lemma~\ref{pct:zeroterm-equiv} gives $h=h(\kappa,a)$ with
\[
|a|^x\tilde\nu_\infty(a)\pctL^\infty_{\mathrm t}(a)\omega'_\infty(1,\kappa)v'=\varrho(\kappa,a)|a|^x\,
\omega_\infty(1,h)v,\qquad\varrho(\kappa,a)=\beta(\kappa^{-1}ae_X)|a|^{N/2}\ge0 .
\]
The functional $w\mapsto\Theta^{\mathrm L}_{W,V}(w,\Phi^{\mathrm{row}}_z)(g)$ is continuous for some
moment norm (Proposition~\ref{pcx:decay}(a)) and invariant under $w\mapsto\omega_\infty(1,h)w$, $h\in
\auU(V)(E^+\otimes\R)$ (right translation of the variable of integration over $[\auU(V)]$). Hence
$\Theta^{\mathrm L}_{W,V}(\pctcone_xv',\Phi^{\mathrm{row}}_z)(g)=B(x)\Theta^{\mathrm L}_{W,V}(v,
\Phi^{\mathrm{row}}_z)(g)$ with $B(x)=\int\varrho(\kappa,a)|a|^x\,d\kappa\,da$. On the support of
$\varrho$, $|a_\tau|$ lies in $[2^{-1/2},2^{1/2}]$ (as $\kappa$ preserves lengths), so $\varrho$ is continuous
with compact support in $\pctK'_\infty\times(E\otimes\R)^\times$; thus $B$ is real-analytic
(differentiate under the integral) and $B(x)>0$ since $\varrho\ge0$ and $\varrho(1,1)=\beta(e_X)=1>0$.
\end{proof}

\subsubsection{Adding a plane: proof of Proposition~\ref{pcx:P-plane}}\label{pct:ss-plane}

\begin{proof}[Proof of Proposition~\ref{pcx:P-plane}]
Let $W'=W\oplus\pcxH$ be adapted, $\operatorname{rk}W=N=2d$, and $v$ of $(W,W)$ with finite moments.
Apply \S\ref{pct:ss-setup} with $V'=W'$, $V=W$, $X=A$, $b=0$: then $r=N$ and the integer of
Theorem~\ref{pct:plane} is $r-N/2=d\ge0$. Proposition~\ref{pct:coneimage} with $x=d+1$ gives $v'$
with finite moments and $\Theta^{\mathrm L}_{W,W}(\pctcone_{d+1}v',\Phi^{\mathrm{row}}_\emptyset)=
\Theta^{\mathrm L}_{W,W}(v,\Phi^{\mathrm{row}}_\emptyset)$ for all $\nu$. Let $F=\Theta^{\mathrm L}_{W',W'}
(v',\Psi^{\mathrm{col}}_1)$; $\Psi^{\mathrm{col}}_1=\Phi'_\emptyset$. By Lemma~\ref{pcx:const} (proved in
\S\ref{pct:ss-const}), $F$ is constant, so $F(1)=F_A(m_A(1,1))$, and by \eqref{pct:eq-plane}
\[
F(1)=\kappa_\nu\,\Theta^{\mathrm L}_{W,W}(\pctcone_{d+1}v',\Psi^{\mathrm{col}}_1)(1)=\kappa_\nu\,
\Theta^{\mathrm L}_{W,W}(v,\Psi^{\mathrm{col}}_1)(1),\qquad|\kappa_\nu|=C\,L^\sharp(d+1).
\]
Since $L_*(\pcm,\nu,d+1)=L^\sharp(d+1)^2L_*(\pcm,\nu,d)$,
$\pcxstr_{W'}(\nu;v')=|F(1)|^2/L_*(\pcm,\nu,d+1)=C^2\,|\Theta^{\mathrm L}_{W,W}(v,\Psi^{\mathrm{col}}_1)(1)|^2/\allowbreak 
L_*(\pcm,\nu,d)=C^2\,\pcxstr_W(\nu;v)$, with $C^2>0$ independent of $\nu$.
\end{proof}

\subsubsection{Removing a plane: proof of Proposition~\ref{pcx:P-descent}}\label{pct:ss-descent}

The precise statement proved (the parenthetical remark of Proposition~\ref{pcx:P-descent} summarises
it) is:

\begin{proposition}\label{pct:descent}
Let $W'=W\oplus\pcxH_A$ be an adapted splitting, $N=\operatorname{rk}W$ even, $T'$ a space of rank
$b$, $T_0$ of rank $r$ and $T$ of rank $r+2$, with $N=b+r$ and $q_T=q_{T_0}+1$, and $d=r-N/2$. If
$\mathrm{RC}(W',T,T';x,x';\ell)$ holds, then $\mathrm{RC}(W,T_0,T';v,y;\ell)$ holds with $y=Z_Ax'$
the zero term of $x'$ along the line $Ee_A$ (defined below) and $v=\pctcone_{d+1}(x)$ if $d\ge0$,
$v=\pctcone^\vee_{-d}(x)$ if $d<0$, the cone integrals being formed with an adapted splitting
$V'=T_0\oplus\pcxH_X$ having the signs of $T$. Only values $L^S(z,\nu)$ with $z\ge1$ enter; no
functional equation is used.
\end{proposition}

For a space $U$ of rank $b$ let $I^U_A\subset\auHom(W',U)$ be the maps vanishing on $W\oplus Ee_A$
(isotropic, $\cong U$, with $I^{U\perp}_A/I^U_A=\auHom(W,U)$), $Z_A$ the archimedean zero term along
it, and recall $\Theta^{\mathrm S}_{X,U}(y,\Psi)=\int_{[\auU(U)]}\theta^{X,U}_{\upsilon^b,\nu\upsilon^{n/2}}
(\cdot,h;[y,\Psi^{\mathsf t}]_{\bar S})\,dh$ for $X$ of rank $n$ (Definition~\ref{pcx:def-large}).

\begin{lemma}[The constant term on the small side]\label{pct:small}
There is $\kappa_A\ne0$, depending only on $A$ and $U$, such that for all $y'$ of $(W',U)$ with
finite moments, $z\in\pcMat_{b\times N}(O)$ and $g\in\auU(W)(\auA)$,
\[
\bigl(\Theta^{\mathrm S}_{W',U}(y',\Psi^{\mathrm{col}}_{z\,\mathrm{pr}})\bigr)_A(m_A(1,g))=\kappa_A\,
\Theta^{\mathrm S}_{W,U}(Z_Ay',\Psi^{\mathrm{col}}_z)(g).
\]
\end{lemma}
\begin{proof}
As in Lemma~\ref{pct:expand} (with Proposition~\ref{pcx:decay}(b) for the bound, the finite datum
being supported on injective matrices), the constant term is $\int_{[\auU(U)]}\sum_{u\in U(E)}
T_u\,dh$, with $T_u$ the mean over $\pctN_A$ of the term of $u$ in the expansion (W3) along $I^U_A$,
and $T_u=0$ for anisotropic $u$.

\emph{$u=0$.} As $m_A(1,g)$ fixes $f_A$ and stabilises $I^U_A$, (W2) gives
$T_0=\theta_{W,U}(\omega(g,h)\psZ_{I^U_A}f')$ for $f'=[y',\Psi^{\mathrm{col}\,\mathsf t}_{z\,\mathrm{pr}}]_{\bar S}$,
the characters $(\upsilon^b,\nu\upsilon^{N/2+1})$ of $(W',U)$ becoming those $(\upsilon^b,\nu\upsilon^{N/2})$
of $(W,U)$ because $\chi_{I^U_A}(m_A(1,g),h)=\upsilon(\det h)^{-1}$. At $\bar w\in\bar S$, in the model (W5),
$\Psi^{\mathrm{col}}_{z\,\mathrm{pr}}(Y)=\chi^{-1}(\det(z\,\mathrm{pr}\,Y))$ (for integral $Y$) depends only on
the rows of $Y$ along $W$; the zero term restricts to the matrices vanishing in the row dual to $I^U_A$
and integrates the row of $I^U_A$ over $O$, which gives $\Psi^{\mathrm{col}}_z$ times $1$. At the other
finite places the zero term of the vacuum is a non-zero multiple of the vacuum (W4), and at $\infty$
it is $Z_A$. So $\psZ_{I^U_A}f'=\kappa_A[Z_Ay',\Psi^{\mathrm{col}\,\mathsf t}_z]_{\bar S}$.

\emph{Isotropic $u\ne0$.} As in Lemma~\ref{pct:unfold}(i), $T_u$ is $\theta$ of the zero term along
$\auHom(W,Eu)$ of the slice at $u$ of $\omega(m_A(1,g),h)f'$; the slice is $\omega(g,h)$ applied to a
datum whose component at $\bar S$ is supported, like $\Psi^{\mathrm{col}\,\mathsf t}_z$, on matrices of full
rank $b$. The functional $\theta\circ\psZ_{\auHom(W,Eu)}\circ\omega(g,h)$ is invariant under the
translations by $\auHom(W,E\,h^{-1}u)(\auA)$, an isotropic subspace whose elements pair
non-trivially at $\bar S$ with every full-rank matrix (as $h^{-1}u\ne0$), so it kills such data by
Lemma~\ref{pcx:meanzero}. Hence $T_u=0$, and the lemma follows.
\end{proof}

\begin{proof}[Proof of Proposition~\ref{pct:descent} (hence of Proposition~\ref{pcx:P-descent})]
By Lemma~\ref{pcx:adapted} there is an adapted splitting $V'=T_0\oplus\pcxH_X$ with the signs of
$T$; by Lemma~\ref{pcx:move} (with $W_1=W=W'$, where $k=1$ may be taken) and the definition of RC,
$\mathrm{RC}(W',V',T';x,x';\ell)$ holds with $c$ replaced by $u^{-1}c$, $|u|=1$. Let $C_0$ and $c$
be as in Definition~\ref{pcx:def-RC} for this identity:
$C_0|c|^2=L_*(\pcm,\nu,d+1)\ell(\nu)$ (as $(r+2)-(N+2)/2=d+1$). For $z\in\pcMat_{b\times N}(O)$
apply RC(b) to the integral matrix $z\,\mathrm{pr}$ ($\Phi^{\mathrm{row}}_{z\,\mathrm{pr}}=\Phi'_z$) and take
constant terms along $A$ at $m_A(1,g)$. By Theorem~\ref{pct:plane} with $u=1$ and Lemma~\ref{pct:small},
\[
\kappa_\nu\,\Theta^{\mathrm L}_{W,T_0}(v,\Phi^{\mathrm{row}}_z)(g)=c\,\kappa_A\,\Theta^{\mathrm S}_{W,T'}
(Z_Ax',\Psi^{\mathrm{col}}_z)(g),
\]
with $v$ as in the statement and $|\kappa_\nu|=C\,L^\sharp(x_d)$, $x_d=d+1$ or $-d$; $\kappa_\nu\ne0$ by
Lemma~\ref{pcx:Lpos}. Put $c_0=c\kappa_A/\kappa_\nu$ and $C'=C_0C^2/|\kappa_A|^2$. Then
$C'|c_0|^2=C_0|c|^2C^2/|\kappa_\nu|^2=L_*(\pcm,\nu,d+1)\ell(\nu)/L^\sharp(x_d)^2=L_*(\pcm,\nu,d)\ell(\nu)$, by
$L_*(e+1)=L^\sharp(e+1)^2L_*(e)$ ($e=d\ge0$) and its form $L_*(-z)=L^\sharp(z+1)^2L_*(-z-1)$ for
$d=-(z+1)<0$. The vectors $v$ and $y=Z_Ax'$ have finite moments (\S\ref{pct:ss-statement} and
Lemma~\ref{pct:slicedecay}). This is $\mathrm{RC}(W,T_0,T';v,y;\ell)$.
\end{proof}

\begin{remark}[Orientation of the zero term]\label{pct:rem-orient}
Proposition~\ref{pct:descent} holds for \emph{every} adapted splitting $A$ of a plane off $W'$, the
zero term $Z_A$ being taken along $I^U_A$ (maps vanishing on $W\oplus Ee_A$), i.e.\ in the orientation of
the isotropic frame $(e_A,f_A)$ of $A$; the space $T_0$ is arbitrary of rank $\operatorname{rk}T-2$ with
$q_{T_0}=q_T-1$, and the exponent of RC(a) passes from $\operatorname{rk}T-(N+2)/2=d+1$ to $d$, with
$\ell$ kept and the constant changed. Applied to the splitting $D^{\mathrm s}$ with the isotropic frames
exchanged (as in Remark~\ref{pcu:rem-descent}) it gives $y=\mathrm{sl}_{D^{\mathrm s}}(0)x'$:
Proposition~\ref{pct:descent} is stated for an arbitrary adapted splitting and is applied with that
splitting, so no conversion is needed.
\end{remark}

\subsubsection{Non-vanishing: proof of Proposition~\ref{pcx:P-nonvan}}\label{pct:ss-nonvan}

For spaces $W$ (even rank $N$), $V$ (rank $r$), $b=N-r$, say that $\mathrm{Van}_\nu(W,V)$ holds if
$\Theta^{\mathrm L}_{W,V}(v,\Phi^{\mathrm{row}}_z)(g)=0$ for all $v$ with finite moments, all
$z\in\pcMat_{b\times N}(O)$ and all $g$.

\begin{proposition}[Vanishing descends]\label{pct:vandesc}
Let $k\in\Z_{\ge0}$ with $k\le\mu_\tau$ for all $\tau$, let $W'=W\oplus
\pcxH_A$, $V'=V\oplus\pcxH_X$ be adapted, $N$ even and $d=r-N/2\ge-k$. Then
$\mathrm{Van}_\nu(W',V')\Rightarrow\mathrm{Van}_\nu(W,V)$. The same holds for $W'$, $V'$ of the
ranks and signs of $W\oplus\pcxH^j$, $V\oplus\pcxH^j$, $j\ge1$.
\end{proposition}
\begin{proof}
Let $v,z,g$ be given and $I=\Theta^{\mathrm L}_{W,V}(v,\Phi^{\mathrm{row}}_z)(g)$. Take $B$ and $v'$ of
Proposition~\ref{pct:coneimage}, and $\pctZ=\pctZ_{v',z,g}$ of Theorem~\ref{pct:plane} (with $u=1$).
By (c), for $x\ge1$, $\pctZ(x)=\varepsilon_{\mathrm t}P(x)L(x,\tilde\nu)B(x)I$. The function
$L(s,\tilde\nu)$ is entire, and $L(j,\tilde\nu)\ne0$ for every integer $j>-k$: for $j\ge1$ by
Lemma~\ref{pcx:Lpos}; for $-k<j\le0$ by the functional equation $\Lambda(s,\tilde\nu)=
\epsilon\Lambda(1-s,\tilde\nu^{-1})$ of the completed $L$-function $\Lambda(s,\tilde\nu)=\prod_\tau
\Gamma_\C(s+\mu_\tau)L(s,\tilde\nu)$ (Tate's thesis, Main Theorem~4.4.1; the archimedean factor of
$a\mapsto(a/\bar a)^{\mu_\tau}$ is $\Gamma_\C(s+\mu_\tau)$), since $\Gamma_\C(j+\mu_\tau)$ and
$\Gamma_\C(1-j+\mu_\tau)$ are finite and non-zero and $L(1-j,\tilde\nu^{-1})\ne0$ ($1-j\ge1$). Both
$\pctZ$ and $\varepsilon_{\mathrm t}PL(\cdot,\tilde\nu)BI$ are real-analytic on $\R$ and agree on
$[1,\infty)$, hence everywhere. At $x=d+1>-k$, Theorem~\ref{pct:plane}(b) and the hypothesis give
$\pctZ(d+1)=(\Theta^{\mathrm L}_{W',V'}(v',\Phi'_z))_A(m_A(1,g))=0$, while
$\varepsilon_{\mathrm t}P(d+1)L(d+1,\tilde\nu)B(d+1)\ne0$. So $I=0$. For $j$ planes, induct, replacing
$W'$, $V'$ by spaces with the same signs (Lemma~\ref{pcx:move}, $\mathrm{Van}$ depending only on
signs) and using that $d$ only increases going up.
\end{proof}

The second input is a special case (period $0$) of the large-triple identity of
Proposition~\ref{pcx:P-top}:

\begin{proposition}[Vanishing at the top]\label{pct:vantop}
Let $J$ be positive definite of rank $2k$, $Y=\pcxH\oplus J$, $X=\pcxH^{2k}\oplus Y$ (adapted),
$T=J^-\oplus Y$. If $\Theta^{\mathrm L}_{Y,Y}(w,\Psi^{\mathrm{col}}_1)\equiv0$ for all $w$ with finite
moments, then $\mathrm{Van}_\nu(X,T)$.
\end{proposition}
In the notation of \S\S\ref{pcu:ss-top-setting}--\ref{pcu:ss-std-data} take $W_0=Y$ (rank $n=2j$,
$j=k+1$, one negative sign, the plane split off adaptedly), $U=J^-$ ($b=2k$, signs $t'=-1$), so that
$X=\pcxH^{b}\oplus W_0$ and $T=U\oplus W_0$ are those of Proposition~\ref{pct:vantop}. The
proof of Proposition~\ref{pcx:P-top} in \S\ref{pcu:sec} is written for the particular vector $\Phi$ of
\S\ref{pcu:ss-period}, whose slices are supported on the real class of $U$ (Lemma~\ref{pcu:support});
for an arbitrary $v$ that support property is replaced, in case (b) of Lemma~\ref{pcu:coeff}, by the
following vanishing, which uses the signs.

\begin{lemma}[Coefficients at non-negative indices]\label{pct:negcoeff}
In the situation just described, let $v$ be any Fock vector of $\auHom(X,T)_\R$ with finite moments
and $\beta$ an invertible hermitian $b\times b$ matrix over $E$ that is not negative definite at some
infinite place. Then $\pcuF_\beta\Theta^{\mathrm L}_{X,T}(v,\Phi^{\mathrm{row}}_{z_0})\equiv0$ on
$P_D(\auA)$.
\end{lemma}
\noindent Lemma~\ref{pct:negcoeff} is Theorem~\ref{pcn:negcoeff}, proved in \S\ref{pcn:sec}. It uses
only that the signs of $U$ are $-1$, that $W_0$ has a positive sign, and that
$\mu_\tau+p_0-\frac n2\ge1$ ($p_0=2k+1$ positive signs of $W_0$, $n=2k+2$), i.e.\
$k+1<2k+1+\mu_\tau$, which holds. The proof does not unfold: it regards the coefficient as a continuous
functional of the archimedean datum at a place $\tau$ where $\beta$ is not negative definite, invariant under
$\auU(T_\tau)$ and $\beta$-equivariant under the centre of the unipotent radical; such a functional is carried by
the orbit of index $\beta$ and descends to a functional on the fibre invariant under the stabiliser of a point, and a
compact subgroup of that stabiliser admits no non-zero invariant vector, by computing the action of the rotations of
a negative line orthogonal to the point on polynomial Fock vectors (Lemmas~\ref{pcn:rot}--\ref{pcn:nofix},
\ref{pcn:local}).

\begin{proof}[Proof of Proposition~\ref{pct:vantop}]
Fix $\nu$ and $v$ with finite moments, and let $\mathcal D=\mathcal T$ be the linear map of
Lemma~\ref{pcx:linfactor} attached to $\Phi'\mapsto\Theta^{\mathrm L}_{X,T}(v,\Phi')$ (this is the
difference of \S\ref{pcu:ss-std-data} with $C=0$). It satisfies (E1)--(E4) of \S\ref{pcu:ss-std-data}
(their proofs do not use the particular vector), with
$\mathcal D(\Psi^{\mathrm{col}}_{z'})=\Theta^{\mathrm L}_{X,T}(v,\Phi^{\mathrm{row}}_{z'})$. By
Proposition~\ref{pcu:selfdual} it suffices to show $\pcuF_\beta\Theta^{\mathrm L}_{X,T}(v,
\Phi^{\mathrm{row}}_{z_0})(p)=0$ at every test pair $(\beta,q)$ for $U$.

Type (b): $\beta$ is not of the form $a^*\mathrm{diag}(t')a=-a^*a$ at some infinite place, i.e.\ not
negative definite there; Lemma~\ref{pct:negcoeff}.

Type (a): $\beta=J_U$. By (U1) and Witt's theorem the coefficient is
$\int_{\auU(W_0)(E)\backslash\auU(T)(\auA)}\mathrm{term}_{x_0}(p,h)\,dh$ (the stabiliser of $x_0$ is
$\auU(W_0)$), absolutely convergent by Lemma~\ref{pcu:integrable}(a) (with $\phi=v$, which has finite
moments of all orders). By the quotient integral formula it is
$\int_{\auU(W_0)(\auA)\backslash\auU(T)(\auA)}\int_{[\auU(W_0)]}\mathrm{term}_{x_0}(p,lh)\,dl\,dh$, and we show
the inner integral is $0$ for every $h$. By (U2), $\mathrm{term}_{x_0}(p,lh)$ is a scalar (independent
of $l$, since $l$ fixes $x_0$ and $|\det|=1$) times $\vartheta_{F^{(y)},x'c^*J_{W_0},(g,lh)}(\phi)$ with
$\phi=\mathrm{sl}_D(h^{-1}Y)v$, a vector of $\auHom(W_0,T)_\R=\auHom(W_0,U)_\R\oplus\auHom(W_0,W_0)_\R$
with finite moments, and $F^{(y)}=f\boxtimes F_1$, $F_1=[\,\cdot\,,\Psi^{\mathrm{col}}_1]_S$, the slice
of $\Phi^{\mathrm{row}}_{z_0}$ being $\chi^{-1}\circ\det$ in the $W_0$-block (as in the proof of
Lemma~\ref{pcu:period}). Expand $\phi=\int\phi(y)\psrho_\infty(y)e\,dy$ (Lemma~\ref{pcx:superpos}); a
coherent state is $\psrho(y_1)e\boxtimes\psrho(y_2)e$ along the orthogonal decomposition, so by
Lemma~\ref{pcu:product} the integrand is
$\int\phi(y)\,\vartheta_{f,\dots}(\psrho(y_1)e)\,\theta^{W_0,W_0}_{\gamma,\gamma}(g,l;\psrho(y_2)e\otimes
F_1)\,dy$. The bound of Proposition~\ref{pcx:decay}(a) for the second factor
($\le C(1+|y_2|)^M\mathrm{ht}(l)^{-N_0}\mathrm{ht}(g)^B$) and $\|\phi\|_M<\infty$ allow Fubini in
$(y,l)$; the integral over $l$ of the second factor is
$\Theta^{\mathrm L}_{W_0,W_0}(\psrho(y_2)e,\Psi^{\mathrm{col}}_1)(g)=0$ by hypothesis (coherent states
have finite moments). So the coefficient vanishes, and Proposition~\ref{pcu:selfdual} gives
$\Theta^{\mathrm L}_{X,T}(v,\Phi^{\mathrm{row}}_{z'})=0$ for all integral $z'$, i.e.\
$\mathrm{Van}_\nu(X,T)$.
\end{proof}

\begin{proof}[Proof of Proposition~\ref{pcx:P-nonvan}]
Suppose $\pcxstr_Y(\nu;z)=0$ for all $z$. By Lemma~\ref{pcx:const}, $\Theta^{\mathrm L}_{Y,Y}(z,
\Psi^{\mathrm{col}}_1)\equiv0$, so $\mathrm{Van}_\nu(X,T)$ by Proposition~\ref{pct:vantop}. Here $X$ has
rank $6k+2$ and $2k+1$ negative signs, $T$ rank $4k+2$ and $2k+1$ negative signs; they have the ranks
and signs of $J\oplus\pcxH^{2k+1}$ and $0\oplus\pcxH^{2k+1}$. Proposition~\ref{pct:vandesc} with
$j=2k+1$ (the integer $d$ is $\ge-k$ at every level, $-k$ at the bottom) gives
$\mathrm{Van}_\nu(J,0)$. But for $(J,0)$, $b=2k$, $z=1_{2k}$, $v$ the vacuum and $g=1$, the kernel
on $\auHom(J,0)=0$ is $\Phi^{\mathrm{row}}_{1}(\emptyset)=\chi^{-1}(\det1)=1$, so the integral is the
volume of $[\auU(0)]$, non-zero. Contradiction.
\end{proof}

\subsubsection{The inputs of Theorem~\ref{pcx:X}}\label{pct:ss-status}
This section proves Theorem~\ref{pct:plane}; Lemmas~\ref{pct:poisson}--\ref{pct:poisson-a},
\ref{pct:qref}, \ref{pct:SA}, \ref{pct:realmixed}, \ref{pct:zeroterm-equiv}, \ref{pct:smoothwitt},
\ref{pct:small}; Lemma~\ref{pcx:const}; Propositions~\ref{pct:coneimage}, \ref{pcx:P-plane},
\ref{pct:descent} (= \ref{pcx:P-descent}), \ref{pct:vandesc}, \ref{pct:vantop}, hence
Proposition~\ref{pcx:P-nonvan}. Lemma~\ref{pct:negcoeff} is proved in \S\ref{pcn:sec}
(Theorem~\ref{pcn:negcoeff}), without using (U1); type (a) of Proposition~\ref{pct:vantop} uses Lemma~\ref{pcu:integrable}.

With \S\ref{pcu:ss-final} and \S\ref{pcn:ss-dep}, Theorem~\ref{pcx:X} therefore rests only on the
following published results (with their hypotheses verified where used): the facts (W1)--(W6) on the Weil
representation (Weil, Kudla, Harris--Kudla--Sweet; {\cite[Lemma~5.2]{P49}}, {\cite[Proposition~7.17]{P49}});
Stone--von Neumann and Weil's intertwining operator (M2); reduction theory (R1) (Borel) and moderate
growth of theta functions (R2) (Lemma~\ref{pcx:thetagrowth}, from Godement's form of Borel's reduction theorem); Landherr's theorem, the Hasse norm theorem and Jacobowitz's
classification of unimodular hermitian lattices (Lemmas~\ref{pcx:landherr}, \ref{pcu:selfdual}); Witt's
theorem and generation of unitary groups by quasi-reflections (Hahn--O'Meara); strong approximation
for the additive group (Cassels--Fr\"ohlich); weak approximation for unitary groups
(Platonov--Rapinchuk); the Iwasawa decomposition and the integration formula for parabolic
subgroups (Moeglin--Waldspurger); Hecke's theorems for the non-trivial unitary character
$\tilde\nu$ --- $L(s,\tilde\nu)$ entire, non-vanishing on $\Rel s=1$, and Tate's functional equation
(Tate's thesis, Main Theorem~4.4.1); classical Gauss sums (Lemma~\ref{pct:gauss}, proved); the existence
of smooth local sections of submersions (Lemma~\ref{pct:smoothwitt}); Witt's theorem over $\C$, the
Fock-model action of the maximal compact subgroup (Folland~\cite[Prop.~4.39]{FollandPS},
Kashiwara--Vergne~\cite{KashiwaraVergne}), the unitarity of the Bargmann transform (Lemma~\ref{pcu:mixedmodel}), Stone's theorem and the inverse function theorem. The density of
polynomial (finite Hermite) vectors among smooth vectors, the equivalence of the norms $\|\cdot\|_s$ with
$\|(\mathsf N+1)^s\cdot\|$, the bound of moments by these norms, the smoothness of the action on smooth vectors, the
mixed model on smooth vectors, invariant distributions and adapted charts are proved in
\S\S\ref{pcn:ss-smooth}, \ref{pcn:ss-local} (Lemmas~\ref{pcn:smooth}, \ref{pcn:smoothaction}, \ref{pcn:mixed},
\ref{pcn:local}; see also Reed--Simon~\cite[Appendix to \S V.3, Lemma~2 and Thm.~V.13]{ReedSimonI} for Hermite expansions).



% ===== C-exchange-neg.tex =====
% proofs/C-exchange-neg.tex
% Proof of Lemma pct:negcoeff of C-exchange-tower.tex: centre coefficients of the theta integral
% vanish at indices that are not negative definite at some real place.
% Labels of this file start with pcn:. Macros: those of C-exchange.tex, C-exchange-top.tex and
% C-exchange-tower.tex.
\providecommand{\pcnK}{\mathsf K}
\providecommand{\pcnC}{\mathsf C}
\providecommand{\pcnS}{\mathsf S}

\subsection{Centre coefficients at indices that are not negative definite}\label{pcn:sec}

This section proves Lemma~\ref{pct:negcoeff}, the last input of Proposition~\ref{pct:vantop} and hence of
Proposition~\ref{pcx:P-nonvan}. The argument: the coefficient, as a function of the archimedean
component of the datum at a real place $\tau$ where the index is not negative definite, is a continuous linear
functional which is invariant under $\auU(T_\tau)$ (right invariance of the measure on $[\auU(T)]$) and transforms by
the character of index $\beta$ under the centre of the unipotent radical. Such a functional is carried by the orbit
$\{x:x^*Jx=\beta\}$ at $\tau$ and descends to a functional on the fibre invariant under the stabiliser of a point
$x$; a compact subgroup of that stabiliser contains, for every point of the Fock space, a circle of rotations of a
negative line; these act on polynomial Fock vectors by a substitution times a positive power of the rotation
parameter, which forces every invariant functional to vanish. No unfolding of the global integral is needed. We first
fix the statement and the conventions, then do the computation at one real place (\S\S\ref{pcn:ss-fock},
\ref{pcn:ss-local}), and then the adelic argument (\S\ref{pcn:ss-proof}).

\subsubsection{Statement and conventions}\label{pcn:ss-statement}

We are in the situation of Proposition~\ref{pct:vantop}, in the notation of
\S\S\ref{pcu:ss-top-setting}--\ref{pcu:ss-coeff}: $J$ is positive definite of rank $2k$ ($k\ge0$),
$W_0=Y=\pcxH\oplus J$ has rank $n=2k+2$ and signature $(p_0,q_0)=(2k+1,1)$ at every infinite place,
$U=J^-$ has rank $b=2k$ and signs $t'=-1$, $X=\pcxH^b\oplus W_0$ with the adapted splitting
$D\colon X=(\ell\oplus\ell')\perp W_0$, and $T=U\oplus W_0$, of rank $b+n$ and signature
\[
(p',q')=(p_0,\,b+q_0)=(2k+1,\,2k+1)
\]
at every infinite place. The weight $\mu$ of $\nu$ has $\mu_\tau\ge1$ for all $\tau$. For a space $W$ of
even rank $N$, $\gamma_W=\nu^{-1}\upsilon^{N/2}$; the theta integral is
$\Theta^{\mathrm L}_{X,T}(v,\Phi)(g)=\int_{[\auU(T)]}\theta^{X,T}_{\gamma_X,\gamma_X}(g,h;[v,\Phi]_S)\,dh$.
The centre coefficient $\pcuF_\beta$ is that of \S\ref{pcu:ss-top-setting}: $z$ runs through the skew-hermitian
matrices, $\beta$ is hermitian, and $\psi_\beta(z)=\psi(-\frac12\jmath\operatorname{tr}(z\beta))$ \eqref{pcu:eq-psibeta}; the
coefficient $\pcuF_\beta$ keeps the terms of the $x$ with Gram matrix $x^*J_Tx=\beta$.

\begin{theorem}[$=$ Lemma~\ref{pct:negcoeff}]\label{pcn:negcoeff}
Let $v$ be a Fock vector of $\auHom(X,T)_\R$ with finite moments, and let $\beta$ be an invertible
hermitian $b\times b$ matrix over $E$ such that, for some infinite place $\tau$, $\tau(\beta)$ is not
negative definite. Then $\pcuF_\beta\Theta^{\mathrm L}_{X,T}(v,\Phi^{\mathrm{row}}_{z_0})(p)=0$ for every
$p\in P_D(\auA)$.
\end{theorem}

The only properties of the situation that the proof uses are:
\begin{enumerate}
\item[(H1)] every sign of $U$ is $-1$ at $\tau$, so that $T_\tau$ has exactly $b$ more negative signs
than $W_{0,\tau}$;
\item[(H2)] $W_0$ has a positive sign at $\tau$ ($p_0\ge1$);
\item[(H3)] the integer $N_\tau=\mu_\tau+p_0-\frac n2$ is $\ge1$. Here $N_\tau=\mu_\tau+k\ge1$.
\end{enumerate}
In terms of the weight $\mu^\flat_\tau=\frac n2-\mu_\tau$ of the character $\gamma_{W_0}$ by which
$\auU(T)$ acts on the slices (Lemma~\ref{pcn:slicechar}), (H3) reads $\mu^\flat_\tau<p_0$, i.e.\
$k+1<2k+1+\mu_\tau$. No condition of the form $2q_{W_0}\le n$ is needed (it holds anyway: $q_0=1$).

\begin{lemma}[The character on the slices]\label{pcn:slicechar}
For $x\in\pcMat_{(b+n)\times b}(E)$ let $\mathrm{term}_x(p,h)=\theta_{W_0}((\omega(p,h)F)[x])$ be the
term of (U1), $F=[v,\Phi^{\mathrm{row}}_{z_0}]_S$, $\omega=\omega_{\gamma_X,\gamma_X}$. Let
$\pcnS_x\subset\auU(T)(\auA)$ be the pointwise stabiliser of the columns of $x$. For $s\in\pcnS_x$,
\[
\mathrm{term}_x(p,sh)=\theta_{W_0}\bigl(\omega^\flat(1,s)\,\Xi_h\bigr),\qquad\Xi_h=(\omega(p,h)F)[x],
\]
where $\omega^\flat=\omega_{\gamma_X\upsilon^{-b},\,\gamma_X\upsilon^{-b}}$ is the Weil representation
of the pair $(W_0,T)$ whose $\auU(T)$-character is $\gamma_X\upsilon^{-b}=\gamma_{W_0}$.
\end{lemma}
\begin{proof}
$\omega(p,1)$ and $\omega(1,s)$ commute, so $(\omega(p,sh)F)[x]=(\omega(1,s)\omega(p,h)F)[x]$. The Levi law
of Lemma~\ref{pcu:laws} with $q=1$, $g=1$ gives $(\omega(1,s)f)[x]=\omega^\flat(1,s)(f[s^{-1}x])$, and
$s^{-1}x=x$. The character: in the Schr\"odinger model of the variable $\auHom(\ell,T)\cong T^b$, the
group $\auU(T)$ acts linearly with determinant $\det(h)^b$, and the factor $\upsilon(\det h)^{b}$ of
(W2) for this action is absorbed by the \emph{target} character, which becomes $\gamma_X\upsilon^{-b}$.
(This is the target-character shift of Lemma~\ref{pcu:laws}. Since $\Theta^{\mathrm L}$ uses the
same character $\gamma_X$ on both groups, the scalar $\chi_V$ of Lemma~\ref{pcu:laws} is the same with
either character.) Finally
$\gamma_X\upsilon^{-b}=\nu^{-1}\upsilon^{(2b+n)/2-b}=\nu^{-1}\upsilon^{n/2}=\gamma_{W_0}$.
\end{proof}

\emph{Convention on $\upsilon$.} All characters above are evaluated on determinants, which lie in the
norm-one ideles $\auU(1)(\auA)$; $\upsilon$ there means the restriction $\upsilon|_{\auU(1)}$, as in
$\upsilon_k=(\upsilon|_{\auU(1)})^{\lfloor k/2\rfloor}$ of {\cite[\S5.1.1]{P49}} (for even rank $N$,
$\upsilon_N=\upsilon^{N/2}$ on $\auU(1)$, which is the $\upsilon^{N/2}$ in $\gamma_W$). The identity
$\upsilon(x/\bar x)=\upsilon(x)^2$ ($x\in\auA_E^\times$), used in \S\ref{pct:sec} to pass between a
character of $\auU(1)$ and one of $\auA_E^\times$, is exact: $\upsilon(\bar x)=\upsilon(x\bar x)
\upsilon(x)^{-1}$, and $\upsilon(x\bar x)=\omega_{E/E^+}(\mathrm N x)=1$ because the quadratic character
of $E/E^+$ is trivial on norms. So no factor is affected by this normalisation. At an infinite place
$\upsilon_\tau|_{\auU(1)}(z)=z$ ({\cite[Theorem~3.29]{P49}}) and $\nu_\tau|_{\auU(1)}(z)=z^{\mu_\tau}$ (the
base change $\tilde\nu(x)=\nu(x/\bar x)$ has $\tilde\nu_\tau(a)=(a/|a|)^{2\mu_\tau}$), so
$\gamma_{W_0,\tau}(z)=z^{n/2-\mu_\tau}=z^{\mu^\flat_\tau}$.

\subsubsection{Smooth vectors of Fock spaces}\label{pcn:ss-smooth}

This subsection proves the facts on smooth vectors used below (and in \S\ref{pct:sec}); everything is
elementary once the Fock model is written in monomials. Let $W_\infty$ be a real symplectic space of
dimension $2D$ with a positive complex structure, realised as in the paragraph preceding
{\cite[Lemma~8.2]{P49}} (which applies to every pair of spaces, see the first paragraph of the proof of
Lemma~\ref{pcx:thetagrowth}): $\pzF$ is the Fock space of entire functions on $\C^D$ with
$\|F\|^2=\boldsymbol\pi^{-D}\int|F(z)|^2e^{-|z|^2}dz$, $\psrho=\psrho_\infty$ is given by {\cite[(41)]{P49}},
$\mathsf B=(\mathsf e_k,i\mathsf e_k)_{k\le D}$ is the real orthonormal basis used there, $z_k=z\cdot\mathsf e_k$,
and for $w\in W_\infty$, $\partial_w$ is the generator of the strongly continuous unitary one-parameter group
$s\mapsto\psrho(sw)$ (a closed skew-adjoint operator, by Stone's theorem). On a polynomial $P$, by
{\cite[(41)]{P49}} (proof of {\cite[Lemma~8.2(b)]{P49}}),
\begin{equation}\label{pcn:eq-dw}
\partial_wP=(z\cdot w)P-\mathsf D_{\bar w}P,\qquad a^+_kP:=z_kP=\tfrac12(\partial_{\mathsf e_k}-i\partial_{i\mathsf e_k})P,
\qquad a^-_kP:=\mathsf D_{\bar{\mathsf e}_k}P=-\tfrac12(\partial_{\mathsf e_k}+i\partial_{i\mathsf e_k})P,
\end{equation}
so on polynomials $\partial_w=\sum_k(c_k(w)a^+_k+c'_k(w)a^-_k)$ with $|c_k(w)|,|c'_k(w)|\le|w|$, real linear in $w$.
The inner product of $\pzF$ is the Fischer inner product ({\cite[\S3.1.4]{P49}}): $a^-_k=\partial/\partial z_k$
is adjoint to $a^+_k$ on polynomials, the monomials $z^\alpha$ are orthogonal with $\|z^\alpha\|^2=\alpha!$,
and they form an orthogonal basis of $\pzF$ (the Taylor series of $F\in\pzF$ converges to $F$ in $\pzF$).
Write $v=\sum_\alpha c_\alpha z^\alpha$ and put
\[
p_k(v)^2=\sum_\alpha(1+|\alpha|)^{2k}\alpha!\,|c_\alpha|^2\in[0,\infty],\qquad
\mathcal S(\pzF)=\{v\in\pzF:\ p_k(v)<\infty\text{ for all }k\ge0\},
\]
a Fr\'echet space (an intersection of weighted $\ell^2$ spaces). Note $p_k(v)=\|(\mathsf N+1)^kv\|$ for the
number operator $\mathsf N=\sum_ka^+_ka^-_k$, $\mathsf Nz^\alpha=|\alpha|z^\alpha$. Recall
$\|\Psi\|_s=\max\|\partial_{w_1}\cdots\partial_{w_j}\Psi\|$ over $j\le2s$, $w_i\in\mathsf B$
(before {\cite[Lemma~8.2]{P49}}). If $\pzF=\pzF'\hat\otimes\pzF''$ is the Fock space of an orthogonal sum
$W'\oplus W''$ (monomials in the two sets of variables), then $p_k(v'\otimes v'')\le p_k(v')p_k(v'')$, since
$1+|\alpha'|+|\alpha''|\le(1+|\alpha'|)(1+|\alpha''|)$.

\begin{lemma}[Smooth vectors]\label{pcn:smooth}
\textup{(a)} Every $v\in\mathcal S(\pzF)$ lies in the domain of every product $\partial_{w_1}\cdots\partial_{w_j}$
(each factor applied to a vector of the domain of the next), the product acts on the Taylor series
term by term through \eqref{pcn:eq-dw}, $w\mapsto\partial_wv$ is real linear, $\partial_w\mathcal S(\pzF)\subset
\mathcal S(\pzF)$, and
\[
\|\partial_{w_1}\cdots\partial_{w_j}v\|\le(2D)^jj^{j/2}|w_1|\cdots|w_j|\;p_{\lceil j/2\rceil}(v).
\]
\textup{(b)} Let $\mathsf B'$ be any real basis of $W_\infty$ and $v\in\pzF$ a vector in the domain of every
product $\partial_{u_1}\cdots\partial_{u_j}$ with $u_i\in\mathsf B'$, $j\le2s$. Then $p_s(v)\le C(\mathsf B',s)\max
\|\partial_{u_1}\cdots\partial_{u_j}v\|$ (maximum over these products). In particular $\mathcal S(\pzF)$ is the
space $\pzF^{\mathrm{sm}}$ of vectors in the domains of all products of the $\partial_w$, it equals
$\bigcap_s\pzF_s$, the norms $\|\cdot\|_s$ define its topology
($(D+1)^{-s}p_s\le\|\cdot\|_s\le(2D)^{2s}(2s)^{s}p_s$), and the polynomials are dense in it.
\textup{(c)} In the realisation of Lemma~\ref{pcx:superpos}, $(1+|x|)^{2s}|v(x)|\le(D+1)^{2s}\|v\|_s$ for all $x$,
hence $\|v\|_M\le C_{M,s}\|v\|_s$ whenever $2s>M+2D$, with $C_{M,s}=(D+1)^{2s}\int(1+|x|)^{M-2s}dx$.
\textup{(d)} Let $k$ be a unitary operator on $\pzF$ with $k\psrho(w)k^{-1}=\psrho(\sigma w)$ for a real
symplectic map $\sigma$ (for instance $\omega_\infty(g,h)$, by (W1)). Then $k\,\mathcal S(\pzF)=\mathcal S(\pzF)$,
$\partial_wkv=k\,\partial_{\sigma^{-1}w}v$ for $v\in\mathcal S(\pzF)$, and $\|kv\|_s\le(2D\|\sigma^{-1}\|)^{2s}\|v\|_s$.
The Heisenberg operators $\psrho(y)$ preserve $\mathcal S(\pzF)$, with $\|\psrho(y)v\|_s\le(1+c_\rho|y|)^{2s}\|v\|_s$.
\end{lemma}
\begin{proof}
(a) A product of $j$ operators $a^\pm_{k}$ maps $z^\alpha$ to a multiple of $z^{\alpha+\delta}$, where $\delta$
depends only on the product, and $\|a^+_kz^\alpha\|^2=(\alpha_k+1)\alpha!$, $\|a^-_kz^\alpha\|^2=\alpha_k\alpha!$;
so it multiplies the norm of $z^\alpha$ by at most $\prod_{i<j}(|\alpha|+i+1)^{1/2}\le j^{j/2}(1+|\alpha|)^{j/2}$, and
distinct $\alpha$ go to orthogonal vectors. Expanding $\partial_{w_1}\cdots\partial_{w_j}$ on polynomials into
$(2D)^j$ such products with coefficients of modulus $\le|w_1|\cdots|w_j|$, the termwise operator $T$ given by
\eqref{pcn:eq-dw} satisfies $\|Tv\|\le(2D)^jj^{j/2}|w_1|\cdots|w_j|p_{\lceil j/2\rceil}(v)$ and, similarly,
$p_k(Tv)\le C\,p_{k+j}(v)$. For $j=1$: if $v_N$ is the truncation of the Taylor series of $v\in\mathcal S(\pzF)$ to
degrees $\le N$, then $v_N\to v$ and $\partial_wv_N=Tv_N\to Tv$ in $\pzF$; as $\partial_w$ is closed, $v$ is in
its domain and $\partial_wv=Tv\in\mathcal S(\pzF)$. Induction on $j$ gives the rest; linearity in $w$ holds
because $T$ is linear in $w$.

(b) For a polynomial $q$ and $v$ as in (b), skew-adjointness gives $\langle v,\partial_{u_1}\cdots\partial_{u_j}q
\rangle=(-1)^j\langle\partial_{u_j}\cdots\partial_{u_1}v,q\rangle$ (one factor at a time; polynomials lie in every
domain by (a)). Expand $(\mathsf N+1)^s=(1+\sum_ka^+_ka^-_k)^s$ into $(D+1)^s$ products of at most $2s$
operators $a^\pm_k$, each a combination, through \eqref{pcn:eq-dw} and the linearity of $w\mapsto\partial_w$ on
polynomials, of products of at most $2s$ operators $\partial_u$, $u\in\mathsf B'$, with coefficients bounded in
terms of $\mathsf B'$ and $s$ (for $\mathsf B'=\mathsf B$ the coefficients have total modulus $\le1$ for each
product of $a^\pm$'s). Hence $\langle v,(\mathsf N+1)^sq\rangle=\langle\tilde v,q\rangle$ for a vector $\tilde v$
with $\|\tilde v\|\le C(\mathsf B',s)\max\|\partial_{u_1}\cdots\partial_{u_j}v\|$ (for $\mathsf B'=\mathsf B$,
$C=(D+1)^s$). With $q=z^\alpha$: $(1+|\alpha|)^s\alpha!\,c_\alpha=\langle\tilde v,z^\alpha\rangle$, and Parseval
gives $p_s(v)^2=\sum_\alpha|\langle\tilde v,z^\alpha\rangle|^2/\alpha!=\|\tilde v\|^2$. This is the first claim.
With (a) it gives $\pzF^{\mathrm{sm}}=\mathcal S(\pzF)$ and the two inequalities between $\|\cdot\|_s$ and
$p_s$. A $\|\cdot\|_s$-Cauchy sequence of smooth vectors is $p_s$-Cauchy, so its limit has $p_s<\infty$; hence
$\bigcap_s\pzF_s=\mathcal S(\pzF)$. The truncations $v_N$ converge to $v$ in every $p_k$.

(c) The isomorphism with the model of Lemma~\ref{pcx:superpos} intertwines the Heisenberg operators and
carries vacuum to vacuum, so it preserves the operators $\partial_w$ and the norms $\|\cdot\|_s$, and $v(x)=
\langle\psrho(x)e,v\rangle$ corresponds to $\langle\psrho(x)1,\Psi\rangle$, which by {\cite[(41)]{P49}} (as in the
proof of {\cite[Lemma~8.2(c)]{P49}}) has absolute value $e^{-|x|^2/2}|\Psi(\bar x)|$. By \eqref{pcn:eq-dw},
$z^\beta\Psi$ is a product of $|\beta|$ operators $a^+_k$ applied to $\Psi$, each a combination of two
$\partial_w$ with coefficients $\frac12$, so $\|z^\beta\Psi\|\le\|\Psi\|_s$ for $|\beta|\le2s$; with $|F(z)|\le
\|F\|e^{|z|^2/2}$ (reproducing kernel) and the multinomial expansion of $(1+\sum_k|z_k|)^{2s}\ge(1+|z|)^{2s}$,
whose coefficients add up to $(D+1)^{2s}$, this gives the pointwise bound (this is the proof of
{\cite[Lemma~8.2(b)]{P49}}, which works for every $s$), and integrating gives the moment bound.

(d) Conjugating $s\mapsto\psrho(s\sigma^{-1}w)$ by $k$ gives $s\mapsto\psrho(sw)$, so $k$ carries the domain of
$\partial_{\sigma^{-1}w}$ onto that of $\partial_w$ and $\partial_wk=k\,\partial_{\sigma^{-1}w}$ there. For $v\in\mathcal
S(\pzF)$, induction with (a) shows that $kv$ is in the domain of every product of the $\partial_w$, with
$\partial_{w_1}\cdots\partial_{w_j}kv=k\,\partial_{\sigma^{-1}w_1}\cdots\partial_{\sigma^{-1}w_j}v$; by (b), $kv\in\mathcal
S(\pzF)$, and writing $\sigma^{-1}w_i$ in the basis $\mathsf B$ (coordinates of modulus $\le\|\sigma^{-1}\|$) and
using linearity gives the bound, as in the proof of {\cite[Lemma~8.2(a)]{P49}}. The same applies to $k^{-1}$. The
statement on $\psrho(y)$ is Lemma~\ref{pcx:superpos}(a) (with $\bigcap_s\pzF_s=\mathcal S(\pzF)$).
\end{proof}

\begin{lemma}[Smooth action of a Lie group]\label{pcn:smoothaction}
Let $G$ be a Lie group, $\omega$ a strongly continuous unitary representation of $G$ on $\pzF$ with
$\omega(g)\psrho(w)\omega(g)^{-1}=\psrho(\sigma_gw)$ for a continuous homomorphism $g\mapsto\sigma_g$ into the real
symplectic group, and suppose that for every $X$ in the Lie algebra the generator $A_X$ of $t\mapsto\omega(\exp tX)$
has the polynomials in its domain and acts on them by a polynomial $\mathsf Q_X$ of degree $\le2$ in the
$a^\pm_k$ (for the Weil representation of a unitary group this is {\cite[Lemma~3.6]{P49}}, which holds by
{\cite[Theorem~3.29]{P49}}). Then:
\begin{enumerate}
\item[\textup{(i)}] $\mathcal S(\pzF)$ lies in the domain of $A_X$, and $A_X=\mathsf Q_X$ there (term by term), a
continuous operator of $\mathcal S(\pzF)$; the map $G\times\mathcal S(\pzF)\to\mathcal S(\pzF)$, $(g,v)\mapsto
\omega(g)v$, is continuous, and the operators $\omega(g)$, $g$ in a compact set, are equicontinuous.
\item[\textup{(ii)}] If $\psi\colon G\to\mathcal S(\pzF)$ is smooth, so is $g\mapsto\omega(g)\psi(g)$, and for the
left-invariant vector field $\tilde X$ of $X$, $\tilde X(\omega\psi)(g)=\omega(g)\bigl(\mathsf Q_X\psi(g)+(\tilde X\psi)(g)
\bigr)$. Consequently, for every compact $K\subset G$, $k$ and $r$ there are $C,k'$ with
$\sup_K p_k(\tilde X_1\cdots\tilde X_{r'}(\omega\psi))\le C\max\sup_K p_{k'}(\tilde X_{i_1}\cdots\tilde X_{i_{r''}}\psi)$
($r'\le r$, the maximum over $r''\le r'$ and over a fixed basis of the Lie algebra).
\end{enumerate}
\end{lemma}
\begin{proof}
(i) $\mathsf Q_X$ acts termwise continuously on $\mathcal S(\pzF)$ ($p_k(\mathsf Q_Xv)\le Cp_{k+1}(v)$, as in
Lemma~\ref{pcn:smooth}(a)); since $A_X$ is closed, truncation as in Lemma~\ref{pcn:smooth}(a) gives the first claim.
By Lemma~\ref{pcn:smooth}(d), $\partial_{w_1}\cdots\partial_{w_j}\omega(g)v=\omega(g)\sum_Vc_V(g)Vv$, a sum over
products $V$ of $j$ operators $\partial_u$, $u\in\mathsf B$, with coefficients $c_V$ polynomial in the entries of
$\sigma_g^{-1}$, hence continuous in $g$ and bounded on compact sets; with the strong continuity of $\omega$ on
$\pzF$ this gives the continuity and the equicontinuity.

(ii) First, for $u\in\mathcal S(\pzF)$, $t^{-1}(\omega(e^{tX})u-u)\to\mathsf Q_Xu$ in $\mathcal S(\pzF)$. Indeed, for a
product $W$ of $j$ operators $\partial_w$, write as above $W\omega(e^{tX})=\omega(e^{tX})W_t$ on $\mathcal S(\pzF)$,
$W_t=\sum_Vc_V(t)V$ with $c_V$ smooth and $W_0=W$; then
\[
W\,\frac{\omega(e^{tX})u-u}t=\sum_V\frac{c_V(t)-c_V(0)}t\,\omega(e^{tX})Vu+\frac{\omega(e^{tX})Wu-Wu}t
\longrightarrow W'u+\mathsf Q_XWu\quad\text{in }\pzF,
\]
$W'=\sum_Vc'_V(0)V$. On polynomials $P,q$: $\langle W\omega(e^{tX})P,q\rangle=\langle\omega(e^{tX})P,W^\dagger q\rangle$
($W^\dagger$ the reversed product times $(-1)^j$), whose derivative at $0$ is $\langle\mathsf Q_XP,W^\dagger q\rangle=
\langle W\mathsf Q_XP,q\rangle$, while the derivative of $\langle\omega(e^{tX})W_tP,q\rangle$ is $\langle\mathsf Q_XWP+
W'P,q\rangle$. So $W\mathsf Q_X=\mathsf Q_XW+W'$ on polynomials, hence on $\mathcal S(\pzF)$ (both sides are
continuous $\mathcal S(\pzF)\to\pzF$, Lemma~\ref{pcn:smooth}(a), and polynomials are dense). Thus the limit is
$W\mathsf Q_Xu$, for every $W$, which is the claim by Lemma~\ref{pcn:smooth}(b). Now
\[
\frac{\omega(ge^{tX})\psi(ge^{tX})-\omega(g)\psi(g)}t=\omega(g)\Bigl[\omega(e^{tX})\frac{\psi(ge^{tX})-\psi(g)}t+
\frac{\omega(e^{tX})\psi(g)-\psi(g)}t\Bigr];
\]
the second term tends to $\mathsf Q_X\psi(g)$, and the first to $(\tilde X\psi)(g)$ by the equicontinuity and
continuity in (i). So $\tilde X(\omega\psi)=\omega\cdot\psi_1$ with $\psi_1=\mathsf Q_X\psi+\tilde X\psi$ continuous, of the
same form; by induction all iterated derivatives along the left-invariant fields of a basis exist, are given by
the stated formula and are continuous. The left-invariant fields form a global frame, so in local coordinates all
partial derivatives exist and are continuous, which (with the mean value inequality for each seminorm) means that
$\omega\psi$ is smooth. The estimate follows from the formula, the continuity of $\mathsf Q_X$ and (i).
\end{proof}

\subsubsection{Rotations of a negative line in the Fock model}\label{pcn:ss-fock}

Fix an infinite place $\tau$. Choose orthonormal bases of $T_\tau=T\otimes_{E,\tau}\C$ and
$W_{0,\tau}$, with signs $s'_a$ ($a$ running through the $b+n$ rows) and $\varepsilon_j$ ($j$ running
through the $n$ columns); the first $b$ rows are those of $U$, with $s'_a=-1$ (H1). An element of
$\auHom(W_0,T)_\tau$ is a matrix $X=(X_{aj})\in\pcMat_{(b+n)\times n}(\C)$, the form is
$\operatorname{tr}(X^\dagger X)=\sum_{a,j}s'_a\varepsilon_j|X_{aj}|^2$, and $h\in\auU(T_\tau)$ acts by
$X\mapsto hX$. As in {\cite[\S3.1.4]{P49}}, the Fock space is realised on functions of $X$, the
holomorphic coordinates being
\[
Z_{aj}=\overline{X_{aj}}\ \text{ if }s'_a\varepsilon_j=1,\qquad Z_{aj}=X_{aj}\ \text{ if }s'_a\varepsilon_j=-1,
\]
the vacuum $e$ is a Gaussian function of $\sum_{a,j}|X_{aj}|^2$, and the polynomial vectors are the
functions $P(Z)e$, $P$ a polynomial. The Fock space of $\auHom(W_0,T)_\R$ is the tensor product of
those at the infinite places; its polynomial vectors are $P(Z)e$ with $P$ a polynomial in the
variables $Z_{w,aj}$ of all infinite places $w$. A polynomial vector determines its polynomial (the map
$P\mapsto P(Z)e$ is injective, $e$ having no zeros and $Z$ being a real-linear coordinate system).
In this realisation the polynomial vectors $P(Z)e$ are the polynomial vectors (polynomials) of the model of
\S\ref{pcn:ss-smooth} ({\cite[\S3.1.4]{P49}}), and the smooth vectors are the space $\mathcal S(\pzF)$ of
Lemma~\ref{pcn:smooth}. Write $\pcnK_\tau=\{c\in\auU(T_\tau):
c\,\mathrm{diag}(s')=\mathrm{diag}(s')c\}\cong\auU(p')\times\auU(q')$ for the maximal compact
subgroup; $c\in\pcnK_\tau$ is unitary for the standard form, and $c=\mathrm{diag}(c_+,c_-)$ along the
positive and negative rows.

\begin{lemma}[Compact action on Fock functions]\label{pcn:Kaction}
Let $\omega^\flat_\tau(1,\cdot)$ be the action of $\auU(T_\tau)$ (placed at $\tau$) in the
representation $\omega^\flat$ of Lemma~\ref{pcn:slicechar}. For $c\in\pcnK_\tau$ and every
polynomial $P$,
\[
\omega^\flat_\tau(1,c)\bigl(P(Z)e\bigr)=\chi_\tau(c)\,P\bigl(Z(c^{-1}X)\bigr)\,e,\qquad
\chi_\tau(c)=\det(c_+)^{\mu^\flat_\tau-q_0}\det(c_-)^{\mu^\flat_\tau-p_0},
\]
where $Z(c^{-1}X)$ means the coordinates $Z$ evaluated at $c^{-1}X$ (the variables of the other
infinite places are unchanged). In particular the space $\mathcal P_d$ of polynomial vectors $P(Z)e$
with $P$ homogeneous of degree $d$ is finite-dimensional and $\pcnK_\tau$-stable.
\end{lemma}
\begin{proof}
An element $c\in\pcnK_\tau$ preserves the form and the standard euclidean structure of
$\auHom(W_0,T)_\tau$, hence the complex structure defining the Fock model; such a symplectic
transformation acts on polynomial vectors by linear substitution of the variables times a
character: this is the statement ``$K\times K'$ acts by linear substitutions'' in the proof of
{\cite[Lemma~3.6]{P49}}, applicable since (V) holds by {\cite[Theorem~3.29]{P49}} (for the standard form of the
action of $\auU(m)\subset\mathrm{Sp}(2m,\R)$ in the Fock model see Folland~\cite[Ch.~4, Prop.~4.39, p.~184]{FollandPS},
and Kashiwara--Vergne~\cite{KashiwaraVergne}). As $\pcnK_\tau$ is connected and polynomial vectors are
dense (Lemma~\ref{pcn:smooth}(b)), this determines the action: $f\mapsto\chi(c)f(c^{-1}\cdot)$ for a
character $\chi$ of $\pcnK_\tau$. (The proof of {\cite[Lemma~3.6]{P49}} uses about the pair
$(Y_j,Y_{j+1})$ only (V) and the shape of the twisting characters, so it applies verbatim to the
pair $(W_0,T)$ with the character $\gamma_{W_0}$ on $\auU(T)$; the linear-substitution form is the
statement that $\pcnK_\tau$ preserves the positive complex structure, so that the metaplectic
action of $c$ is $f\mapsto f(c^{-1}\cdot)$ up to a scalar character of the metaplectic cover of
$\auU(p')\times\auU(q')$.) As $e$ is a
function of $|X|^2$ and $|c^{-1}X|=|X|$, $f=P(Z)e$ goes to $\chi(c)P(Z(c^{-1}X))e$. The map
$X\mapsto c^{-1}X$ is complex linear in the coordinates $Z$: for $s'_a\varepsilon_j=1$,
$Z_{aj}(c^{-1}X)=\sum_{a'}\overline{(c^{-1})_{aa'}}Z_{a'j}$, and the sum runs over $a'$ with
$s'_{a'}=s'_a$ because $c^{-1}$ is block diagonal; likewise without conjugation if
$s'_a\varepsilon_j=-1$. So homogeneous polynomials of degree $d$ go to homogeneous polynomials of
degree $d$.

The character $\chi$ is the eigencharacter of the vacuum, $\chi(c)=\det(c_+)^{a_+}\det(c_-)^{a_-}$
for integers $a_\pm$, which we read off from the derived action on $e$ ($P=1$). By (V) of
{\cite[Lemma~3.6]{P49}}, which holds by {\cite[Theorem~3.29]{P49}}, the untwisted operator
$\omega_{\auHom(W_0,T)}(1,h)$ at $\tau$ is the oscillator action times $\det(h)^{-n/2}$, and the
Weyl-ordered oscillator operator of the diagonal generator of the row $c$ is
$-\sum_ja^-_{jc}a^+_{jc}+\frac n2$ in the notation of {\cite[(12)]{P49}} (see the proof of
{\cite[Lemma~3.6]{P49}}). On $e=1$ the operator $a^-_{jc}a^+_{jc}$ gives $1$ if $\eta_{jc}=
-\varepsilon_js'_c=1$ ($a^+=z$, $a^-=\partial$) and $0$ otherwise ($a^+=-\partial$). The twist of
$\omega^\flat$ on $\auU(T)$ is $\gamma_{W_0,\tau}(\det h)=\det(h)^{\mu^\flat_\tau}$. So the weight of
$e$ on the row $c$ is $-\#\{j:\varepsilon_j\ne s'_c\}+\frac n2-\frac n2+\mu^\flat_\tau$, that is
$\mu^\flat_\tau-q_0$ for a positive row and $\mu^\flat_\tau-p_0$ for a negative row.
\end{proof}

\begin{lemma}[Moving a positive vector into the span]\label{pcn:witt}
Let $x\in\pcMat_{(b+n)\times b}(\C)$ have Gram matrix $x^*\,\mathrm{diag}(s')\,x$ which is not negative
semidefinite, and let $i$ be a row with $s'_i=1$ (it exists by (H2)). There are $u\in\auU(T_\tau)$
and $w\in\C^b$ with $uxw=e_i$, the $i$-th basis vector.
\end{lemma}
\begin{proof}
There is $w$ with $w^*x^*\mathrm{diag}(s')xw>0$; rescaling, the vector $xw$ has length $1$, as does
$e_i$. Both span non-degenerate lines of the same length, so by Witt's theorem for the hermitian space
$T_\tau$ there is an isometry $u$ with $u(xw)=e_i$.
\end{proof}

\begin{lemma}[A negative line orthogonal to the frame]\label{pcn:line}
Let $x''\in\pcMat_{(b+n)\times b}(\C)$ and $w\in\C^b$ with $x''w=e_i$, $s'_i=1$, and let $y_j\in\C^{b+n}$
be given for the $q_0$ columns $j$ with $\varepsilon_j=-1$. There is $r\in\C^{b+n}$ with $r_a=0$ whenever
$s'_a=1$, $r^{\mathsf t}x''=0$, $r^{\mathsf t}y_j=0$ for all $j$ with $\varepsilon_j=-1$, and
$r^{\mathsf t}\bar r=1$.
\end{lemma}
\begin{proof}
Consider the linear map $f\colon\C^{b+n}\to\C^b\times\C^{q_0}\times\C^{p'}$,
$f(r)=(r^{\mathsf t}x'',(r^{\mathsf t}y_j)_j,(r_a)_{s'_a=1})$. By (H1), $q'=b+q_0$, so source and
target both have dimension $b+q_0+p'=b+n$. If $\ker f=0$, $f$ is bijective, and there is $r$ with
$f(r)=(0,0,(1)_a)$; then $r_i=1$, but $r_i=r^{\mathsf t}e_i=r^{\mathsf t}x''w=0$, a contradiction. So
there is $r\ne0$ in $\ker f$; normalise $r^{\mathsf t}\bar r=\sum|r_a|^2=1$.
\end{proof}

\begin{lemma}[Rotations of the line]\label{pcn:rot}
Let $r$ be as in Lemma~\ref{pcn:line}, $\Pi=\bar r\,r^{\mathsf t}$ and, for $|\eta|=1$,
$c_\eta=1+(\bar\eta-1)\Pi$. Then $c_\eta\in\pcnK_\tau$, $c_\eta x''=x''$, and for every polynomial $P$
\[
\omega^\flat_\tau(1,c_\eta)\bigl(P(Z)e\bigr)=\eta^{N_\tau}\,(P\circ\sigma_\eta)(Z)\,e ,
\]
where $N_\tau=\mu_\tau+p_0-\frac n2$ and $\sigma_\eta$ is the linear substitution fixing the variables
of the other places and the variables of the positive rows, and sending, for a negative row $a$,
\[
Z_{aj}\mapsto Z_{aj}+(\bar\eta-1)\,r_a\sum_{a'}\bar r_{a'}Z_{a'j}\quad(\varepsilon_j=-1),\qquad
Z_{aj}\mapsto Z_{aj}+(\eta-1)\,\bar r_a\sum_{a'}r_{a'}Z_{a'j}\quad(\varepsilon_j=1),
\]
the sums running over the negative rows $a'$.
\end{lemma}
\begin{proof}
$\Pi$ is hermitian, $\Pi^2=\Pi$ (as $r^{\mathsf t}\bar r=1$), and $\mathrm{diag}(s')\Pi=\Pi\,
\mathrm{diag}(s')=-\Pi$ ($r$ is supported on negative rows). For a hermitian $M$ with $M\Pi=\Pi M=
\epsilon\Pi$, $c_\eta^*Mc_\eta=M+\epsilon\bigl((\eta-1)+(\bar\eta-1)+(\eta-1)(\bar\eta-1)\bigr)\Pi=M$,
since $\eta\bar\eta=1$. With $M=1$ and $M=\mathrm{diag}(s')$, $c_\eta$ is unitary for both forms, so
$c_\eta\in\pcnK_\tau$; and $c_\eta x''=x''+(\bar\eta-1)\bar r(r^{\mathsf t}x'')=x''$. Its positive block
is $1$ and its negative block is $1+(\bar\eta-1)\bar r_-r_-^{\mathsf t}$, of determinant
$1+(\bar\eta-1)r^{\mathsf t}\bar r=\bar\eta$. By Lemma~\ref{pcn:Kaction},
$\chi_\tau(c_\eta)=\bar\eta^{\,\mu^\flat_\tau-p_0}=\eta^{p_0-n/2+\mu_\tau}=\eta^{N_\tau}$. Finally
$c_\eta^{-1}=c_{\bar\eta}$, so $(c_\eta^{-1}X)_{aj}=X_{aj}+(\eta-1)\bar r_a\sum_{a'}r_{a'}X_{a'j}$;
this changes only negative rows $a$, and involves only negative rows $a'$. If $\varepsilon_j=1$ then
$Z_{a'j}=X_{a'j}$ for all negative $a'$ and this is the second formula; if $\varepsilon_j=-1$ then
$Z_{a'j}=\overline{X_{a'j}}$ and conjugating gives the first.
\end{proof}

\begin{lemma}[Fixed polynomial vectors vanish on a hyperplane section]\label{pcn:fixed}
Assume (H3). Let $r$ be as in Lemma~\ref{pcn:line} and $P$ a polynomial such that
$\omega^\flat_\tau(1,c_\eta)(P(Z)e)=P(Z)e$ for all $|\eta|=1$. Then $P(z)=0$ at every point $z$ (a value
of all the variables) with $\sum_{a'}\bar r_{a'}z_{a'j}=0$ for every $j$ with $\varepsilon_j=-1$.
\end{lemma}
\begin{proof}
By Lemma~\ref{pcn:rot} and injectivity of $P\mapsto P(Z)e$, $\eta^{N_\tau}P\circ\sigma_\eta=P$ as
polynomials, for $|\eta|=1$. At such a point $z$, the first formula of Lemma~\ref{pcn:rot} leaves the
variables $z_{aj}$ ($\varepsilon_j=-1$) unchanged, so $\sigma_\eta(z)=z+(\eta-1)z'$, where $z'_{aj}=
\bar r_a\sum_{a'}r_{a'}z_{a'j}$ for negative rows $a$ and $\varepsilon_j=1$, and $z'=0$ otherwise. The
function $g(\eta)=\eta^{N_\tau}P(z+(\eta-1)z')$ is a polynomial in $\eta$, equal to $P(z)$ on the unit
circle, hence identically $P(z)$; and $g(0)=0$ because $N_\tau\ge1$. So $P(z)=0$.
\end{proof}

\begin{lemma}[No fixed polynomial vectors]\label{pcn:nofix}
Assume (H1)--(H3). Let $x''$, $w$, $i$ be as in Lemma~\ref{pcn:line} and
$\pcnK'=\{c\in\pcnK_\tau:cx''=x''\}$. If a polynomial vector $P(Z)e$ is fixed by
$\omega^\flat_\tau(1,c)$ for all $c\in\pcnK'$, then $P=0$.
\end{lemma}
\begin{proof}
Fix a point $z$. Apply Lemma~\ref{pcn:line} with $y_j=(\overline{z_{\tau,aj}})_a$; then $r^{\mathsf t}y_j=0$
gives, after conjugation, $\sum_a\bar r_az_{\tau,aj}=0$ for $\varepsilon_j=-1$. By
Lemma~\ref{pcn:rot} all $c_\eta$ lie in $\pcnK'$, so they fix $P(Z)e$, and Lemma~\ref{pcn:fixed} gives
$P(z)=0$. As $z$ is arbitrary, $P=0$.
\end{proof}

\begin{lemma}[No invariant functionals]\label{pcn:nofunct}
Assume (H1)--(H3) and let $\pcnK'$ be as in Lemma~\ref{pcn:nofix}. Let $\mathcal S$ be either the space
$\mathcal S(\pzF)$ of Lemma~\ref{pcn:smooth} for the Fock space of $\auHom(W_0,T)_\R$ (all infinite places), or the
same space for the Fock space of $\auHom(W_0,T)_\tau$ alone. Every continuous linear functional
$l$ on $\mathcal S$ with $l\circ\omega^\flat_\tau(1,c)=l$ for all $c\in\pcnK'$ is zero.
\end{lemma}
\begin{proof}
$\pcnK'$ is a closed subgroup of the compact group $\pcnK_\tau$, hence compact; let $dc$ be its
probability Haar measure. Let $f\in\mathcal S$ be a polynomial vector, homogeneous of degree $d$ in the
variables of $\tau$ (any polynomial vector is a finite sum of such vectors times polynomials in the variables of
the other places, on which $\omega^\flat_\tau(1,c)$ acts trivially). By Lemma~\ref{pcn:Kaction}, $c\mapsto
\omega^\flat_\tau(1,c)f$ is a continuous (indeed polynomial in the entries of $c$, $\bar c$) map into
a finite-dimensional space of such vectors, so $m=\int_{\pcnK'}\omega^\flat_\tau(1,c)f\,dc$ lies in it,
and $m$ is fixed by $\pcnK'$ (invariance of $dc$). By Lemma~\ref{pcn:nofix} (applied to $m$, resp.\ to $m\otimes
e^\tau$ with $e^\tau$ the vacuum of the other infinite places in the local case), $m=0$. As $l$ restricted
to a finite-dimensional space is continuous, $l(m)=\int l(\omega^\flat_\tau(1,c)f)
\,dc=l(f)$. So $l$ vanishes on every polynomial vector. Polynomial vectors are dense in $\mathcal S$
(Lemma~\ref{pcn:smooth}(b)). By continuity, $l=0$.
\end{proof}

\subsubsection{A local vanishing statement at the place $\tau$}\label{pcn:ss-local}

Keep the place $\tau$ of \S\ref{pcn:ss-fock}. Write $G_\tau=\auU(T_\tau)$, $V_\ell=\auHom(\ell,T)_\tau\cong
\pcMat_{(b+n)\times b}(\C)\cong\C^{p}$ ($p=(b+n)b$; a real vector space of dimension $2p$ with euclidean norm
$|\cdot|$), and $Q(x)=x^*J_Tx\in\mathrm{Herm}_b(\C)$ for $x\in V_\ell$. Let $\pzF_\tau$ and $\pzF^\flat_\tau$ be the
Fock spaces of $\auHom(X,T)_\tau$ and $\auHom(W_0,T)_\tau$, and $\mathcal S_\tau=\mathcal S(\pzF_\tau)$, $\mathcal
S_\flat=\mathcal S(\pzF^\flat_\tau)$ their spaces of smooth vectors (Lemma~\ref{pcn:smooth}), with the seminorms
$p_k$. On $\pzF^\flat_\tau$, $G_\tau$ acts by $\omega^\flat_\tau(1,\cdot)$ (Lemma~\ref{pcn:slicechar}); by
Lemma~\ref{pcn:smoothaction} (with {\cite[Lemma~3.6]{P49}} and {\cite[Theorem~3.29]{P49}}) this action on $\mathcal S_\flat$ is
continuous and smooth in the sense of that lemma. Let $\mathfrak z_\tau$ be the centre $Z_D$ at $\tau$, written
additively: the space of skew-hermitian $b\times b$ complex matrices (\S\ref{pcu:ss-top-setting}). By
Lemma~\ref{pcu:laws}, $z\in\mathfrak z_\tau$ acts on the slice at $x$ through the character
$e(z,Q(x))$, where $e(z,y)=\psi_\tau(-\frac12\jmath\operatorname{tr}(zy))$ for $y\in\mathrm{Herm}_b(\C)$ (with $\jmath$ read
through $\tau$, $\tau(\jmath)=i/\sqrt d$). Writing $z=iH$ with $H$ hermitian,
$-\frac12\jmath\operatorname{tr}(zy)=\operatorname{tr}(Hy)/(2\sqrt d)$ is real, and $\psi_\tau(t)=e^{-4\pi it}$ for real $t$
($\psi_\tau=\psi_{\Q,\infty}\circ\operatorname{Tr}_{\C/\R}$). As $(H,y)\mapsto\operatorname{tr}(Hy)$ is a positive definite real
inner product on $\mathrm{Herm}_b(\C)$, the pairing $(z,y)\mapsto e(z,y)$ is non-degenerate: every unitary character of
$\mathrm{Herm}_b(\C)$ is $e(z,\cdot)$ for exactly one $z\in\mathfrak z_\tau$, and the real-linear forms
$y\mapsto-\frac12\jmath\operatorname{tr}(zy)$, $z\in\mathfrak z_\tau$, span the dual of $\mathrm{Herm}_b(\C)$.

\emph{The mixed model at $\tau$.} By \S\ref{pcu:ss-mixed} at the single place $\tau$, there is an isomorphism
$\mathsf M\colon\pzF_\tau\to L^2(V_\ell;\pzF^\flat_\tau)$, $(\mathsf Mv)(x)=\mathrm{sl}_D(x)v$, compatible with the
Heisenberg operators as in \eqref{pcu:eq-mixed}; it is unitary up to a constant $c_{\mathsf M}>0$,
$\|\mathsf Mv\|^2=c_{\mathsf M}\|v\|^2$, by Lemma~\ref{pcu:mixedmodel} (at the single place $\tau$). In the notation of \eqref{pcu:eq-mixed}, taking two of
$a',w,b$ equal to $0$,
\begin{equation}\label{pcn:eq-heis}
\mathsf M\psrho(La')\mathsf M^{-1}f(x)=f(x+a'),\quad\mathsf M\psrho(Kb)\mathsf M^{-1}f(x)=e^{2i\,\Im\,b^*x}f(x),\quad
\mathsf M\psrho(Iw)\mathsf M^{-1}f(x)=\psrho^\flat(w)f(x),
\end{equation}
and $\mathsf Me=e^{-|x|^2}e^\flat$. By the local form of Lemma~\ref{pcu:laws} (Corollary~\ref{pcu:laws-local}) and
by Lemma~\ref{pcn:slicechar} together with (M3) for the character, for $v\in\pzF_\tau$, $z\in\mathfrak z_\tau$, $h\in G_\tau$:
\begin{equation}\label{pcn:eq-mixedlaws}
\mathsf M(\omega_\tau(n(0,z),1)v)(x)=e(z,Q(x))\,(\mathsf Mv)(x),\qquad
\mathsf M(\omega_\tau(1,h)v)(x)=\omega^\flat_\tau(1,h)\bigl((\mathsf Mv)(h^{-1}x)\bigr).
\end{equation}
Let $\mathcal T$ be the space of functions $V_\ell\to\mathcal S_\flat$ which are finite sums of elements of
$C_c^\infty(V_\ell;\mathcal S_\flat)$ and of functions $g\otimes u$, $g$ a Schwartz function on $V_\ell$, $u\in\mathcal
S_\flat$, with the seminorms $q_r(f)=\max_{|\beta|\le r}\sup_x(1+|x|)^{r+2p}p_r(\partial^\beta f(x))$
($\partial^\beta$ the real partial derivatives on $V_\ell$). For a compact $B\subset V_\ell$, $C_c^\infty(B;\mathcal S_\flat)$
(functions supported in $B$) is a Fr\'echet space for the $q_r$, and $C_c^\infty(V_\ell;\mathcal S_\flat)$ is their union.

\begin{lemma}[Smooth vectors in the mixed model]\label{pcn:mixed}
\textup{(a)} For $f\in\mathcal T$, $\iota(f)=\mathsf M^{-1}f$ lies in $\mathcal S_\tau$, and for every $k$ there are
$C,r$ with $p_k(\iota f)\le C\,q_r(f)$.
\textup{(b)} $\iota(C_c^\infty(V_\ell;\mathcal S_\flat))$ is dense in $\mathcal S_\tau$.
\textup{(c)} For $f\in C_c^\infty(V_\ell;\mathcal S_\flat)$, $z\in\mathfrak z_\tau$, $h\in G_\tau$:
$\omega_\tau(n(0,z),1)\iota(f)=\iota(e'_zf)$ with $e'_z(x)=e(z,Q(x))$, and
$\omega_\tau(1,h)\iota(f)=\iota(h\cdot f)$ with $(h\cdot f)(x)=\omega^\flat_\tau(1,h)f(h^{-1}x)$; both $e'_zf$ and
$h\cdot f$ lie in $C_c^\infty(V_\ell;\mathcal S_\flat)$.
\end{lemma}
\begin{proof}
(a) Let $D_{a'}$ be the derivative in the direction $a'$ and $m_b$ the multiplication by $2i\,\Im\,b^*x$. For
$f\in\mathcal T$, $t^{-1}(f(\cdot+ta')-f)\to D_{a'}f$ in $L^2(V_\ell;\pzF^\flat_\tau)$ (mean value inequality and
dominated convergence, the difference quotients being bounded by $\sup_{|s|\le1}\|D_{a'}f(x+sa')\|\le
C(1+|x|)^{-2p}q_1(f)$); $t^{-1}(e^{2it\Im b^*x}-1)f\to m_bf$ likewise; and $t^{-1}(\psrho^\flat(tw)f(x)-f(x))\to
\partial^\flat_wf(x)$ for each $x$, with $\|t^{-1}(\psrho^\flat(tw)u-u)\|\le\|\partial^\flat_wu\|$ ($u\in\mathcal S_\flat$;
write $\psrho^\flat(tw)u-u=\int_0^t\psrho^\flat(sw)\partial^\flat_wu\,ds$), so again in $L^2$. By \eqref{pcn:eq-heis},
$\iota(f)$ is in the domains of $\partial_{La'}$, $\partial_{Kb}$, $\partial_{Iw}$, and $\mathsf M\partial_{La'}\iota(f)=
D_{a'}f$, $\mathsf M\partial_{Kb}\iota(f)=m_bf$, $\mathsf M\partial_{Iw}\iota(f)=\partial^\flat_wf$ (applied to the values;
Lemma~\ref{pcn:smooth}(a)). These three operators map $\mathcal T$ into itself. Let $\mathsf B'$ be a real basis of
$\auHom(X,T)_\tau$ made of vectors $La'$, $Kb$, $Iw$. By induction, $\iota(f)$ is in the domain of every product
$\partial_{u_1}\cdots\partial_{u_j}$, $u_i\in\mathsf B'$, and $\mathsf M\partial_{u_1}\cdots\partial_{u_j}\iota(f)=P f$ for
a product $P$ of $j$ of the operators $D_{a'}$, $m_b$, $\partial^\flat_w$. By Lemma~\ref{pcn:smooth}(a) and the
definition of $q_r$, $\|Pf\|_{L^2}^2\le\int C(1+|x|)^{2j}\,p_j(\text{derivatives of order}\le j\text{ of }f(x))^2dx
\le C'q_{2j}(f)^2$. Lemma~\ref{pcn:smooth}(b) for the basis $\mathsf B'$ gives $\iota(f)\in\mathcal S_\tau$ and the bound.

(b) By Lemma~\ref{pcn:smooth}(b) the polynomial vectors are dense in $\mathcal S_\tau$; each is a combination of
products of operators $a^+_k$ applied to the vacuum $e$, and each $a^+_k$ is a combination of operators
$\partial_u$, $u\in\mathsf B'$ (by \eqref{pcn:eq-dw} and linearity). As $\mathsf Me=e^{-|x|^2}e^\flat$ lies in
$\mathcal S(V_\ell)\otimes\mathcal S_\flat$, which the three operators of (a) preserve, $\mathsf M$ maps every polynomial
vector into the algebraic tensor product $\mathcal S(V_\ell)\otimes\mathcal S_\flat$. For $g\in\mathcal S(V_\ell)$,
$u\in\mathcal S_\flat$ and a cut-off $\chi$ equal to $1$ near $0$, $q_r((1-\chi(\cdot/R))g\otimes u)\to0$ as
$R\to\infty$, so by (a) $\iota(\chi(\cdot/R)g\otimes u)\to\iota(g\otimes u)$ in $\mathcal S_\tau$. Hence the closure of
$\iota(C_c^\infty(V_\ell;\mathcal S_\flat))$ contains every polynomial vector.

(c) The identities follow from \eqref{pcn:eq-mixedlaws} and the injectivity of $\mathsf M$. $e'_zf$ is smooth with
the same support; $h\cdot f$ is smooth with compact support because $(x,u)\mapsto\omega^\flat_\tau(1,h)u$ is
continuous linear in $u$ (Lemma~\ref{pcn:smooth}(d)).
\end{proof}

\begin{lemma}[Local vanishing]\label{pcn:local}
Assume (H1)--(H3). Let $\beta_\tau\in\mathrm{Herm}_b(\C)$ be invertible and not negative definite, and let
$\lambda$ be a continuous linear functional on $\mathcal S_\tau$ such that, for all $v\in\mathcal S_\tau$,
$z\in\mathfrak z_\tau$ and $h\in G_\tau$,
\[
\text{(i)}\ \ \lambda(\omega_\tau(n(0,z),1)v)=e(z,\beta_\tau)\,\lambda(v),
\qquad\text{(ii)}\ \ \lambda(\omega_\tau(1,h)v)=\lambda(v).
\]
Then $\lambda=0$.
\end{lemma}
\begin{proof}
Put $\Lambda=\lambda\circ\iota$ on $C_c^\infty(V_\ell;\mathcal S_\flat)$; by Lemma~\ref{pcn:mixed}(a) its restriction to
each $C_c^\infty(B;\mathcal S_\flat)$ is continuous, and by Lemma~\ref{pcn:mixed}(b) it suffices to show $\Lambda=0$.
Write $e_z(x)=e(z,Q(x)-\beta_\tau)=e(z,Q(x))\,e(z,\beta_\tau)^{-1}$, a smooth function of modulus $1$ on
$V_\ell$; by (i) and Lemma~\ref{pcn:mixed}(c), $\Lambda((e_z-1)g)=0$ for all $g\in C_c^\infty(V_\ell;\mathcal S_\flat)$ and
all $z$. Let $\Omega=Q^{-1}(\beta_\tau)\subset V_\ell$.

\emph{Step A: $\Lambda$ kills the functions vanishing on $\Omega$.} (1) Let $x_0\notin\Omega$. By the
non-degeneracy of the pairing there is $z$ with $e_z(x_0)\ne1$, hence $e_z\ne1$ on a neighbourhood $U_{x_0}$. If
$g\in C^\infty_c(V_\ell;\mathcal S_\flat)$ has support in $U_{x_0}$, then $g=(e_z-1)g'$ with $g'=g/(e_z-1)\in
C_c^\infty(V_\ell;\mathcal S_\flat)$, so $\Lambda(g)=0$. (2) For $g\in C_c^\infty(B;\mathcal S_\flat)$ and $z$ fixed,
$s^{-1}(e_{sz}-1)g\to\vartheta_zg$ in $C_c^\infty(B;\mathcal S_\flat)$ as $s\to0$, where $\vartheta_z(x)=2\pi i\,\jmath
\operatorname{tr}(z(Q(x)-\beta_\tau))$ (from $e(z,y)=\exp(2\pi i\jmath\operatorname{tr}(zy))$; every derivative of $s^{-1}(e^{s\vartheta}-1)-
\vartheta$, $\vartheta$ a polynomial, is bounded on $B$ by $|s|$ times a constant). By continuity
$\Lambda(\vartheta_zg)=0$; the real-linear forms $y\mapsto-\frac12\jmath\operatorname{tr}(zy)$ span the dual
of $\mathrm{Herm}_b(\C)$, so $\Lambda(t_kg)=0$ for each real coordinate $t_k$ ($1\le k\le b^2$) of $Q-\beta_\tau$ in a
fixed basis of $\mathrm{Herm}_b(\C)$. (3) \emph{Charts adapted to $\Omega$.} Let $x_1\in\Omega$. The differential
$dQ_{x_1}(\xi)=\xi^*J_Tx_1+x_1^*J_T\xi$ is onto $\mathrm{Herm}_b(\C)$: given $\eta$, $\xi=\frac12x_1\beta_\tau^{-1}\eta$
gives $\eta$ (as $x_1^*J_Tx_1=\beta_\tau$). Let $K_1=\ker dQ_{x_1}$ (real dimension $2p-b^2$) with an orthogonal
projection $\pi_1\colon V_\ell\to K_1$. The map $\Xi(x)=(\pi_1(x-x_1),Q(x)-\beta_\tau)\in K_1\times\mathrm{Herm}_b(\C)$
has bijective differential at $x_1$, so by the inverse function theorem it is a diffeomorphism from an open
neighbourhood $N_{x_1}$ of $x_1$ onto a product $B_1\times B_2$ of open balls centred at $0$; in the coordinates
$(y,t)=\Xi(x)$, $\Omega\cap N_{x_1}=\{t=0\}$. If $g\in C_c^\infty(V_\ell;\mathcal S_\flat)$ has support in $N_{x_1}$ and
$g|_\Omega=0$, write $\tilde g=g\circ\Xi^{-1}$; for $(y,t)\in B_1\times B_2$, $\tilde g(y,t)=\tilde g(y,t)-\tilde g(y,0)=
\sum_kt_k\tilde g_k(y,t)$ with $\tilde g_k(y,t)=\int_0^1\partial_{t_k}\tilde g(y,st)\,ds$ (Hadamard; $B_2$ is convex,
and differentiation under the integral sign is legitimate for the continuous seminorms $p_r$). With a cut-off
$\chi\in C_c^\infty(N_{x_1})$ equal to $1$ on the support of $g$, $g=\sum_kt_k\,g_k$ with $g_k=\chi\cdot(\tilde g_k\circ\Xi)
\in C_c^\infty(V_\ell;\mathcal S_\flat)$, so $\Lambda(g)=0$ by (2). (4) A general $g\in C_c^\infty(V_\ell;\mathcal S_\flat)$ with
$g|_\Omega=0$ is $\sum_i\rho_ig$ for a finite smooth partition of unity on its support subordinate to sets $U_{x_0}$
and $N_{x_1}$; each $\rho_ig$ vanishes on $\Omega$, so $\Lambda(g)=0$. In particular, if $\Omega=\emptyset$ then
$\Lambda=0$ and we are done. Assume henceforth $\Omega\ne\emptyset$.

\emph{Step B: the induced functional on $\Omega$.} By Lemma~\ref{pcu:orbitvol} (at the single place $\tau$, with
$m=b+n$, $\beta_0=\beta_\tau$), $\Omega$ is a closed submanifold and one orbit of $G_\tau$; fix $x_1\in\Omega$ and let
$S_1=\mathrm{Stab}_{G_\tau}(x_1)$, the unitary group of the orthogonal complement of the span of $x_1$, so $G_\tau$
and $S_1$ are unimodular and the orbit map identifies $G_\tau/S_1$ with $\Omega$ (so it is a submersion with smooth
local sections). Cover $\Omega$ by charts $N_{x_i}$ as in Step~A(3), with coordinates $(y,t)$, and let $(\rho_i)$ be a
locally finite smooth partition of unity on $\Omega$ with $\mathrm{supp}\,\rho_i\cap\Omega$ compact in
$N_{x_i}\cap\Omega$; fix $\chi_i\in C_c^\infty(B_{2,i})$ equal to $1$ near $0$. For $f\in C_c^\infty(\Omega;\mathcal S_\flat)$ put
\[
\mathrm{ext}(f)=\sum_i\bigl((y,t)\mapsto\chi_i(t)\,(\rho_if)(\Xi_i^{-1}(y,0))\bigr)\circ\Xi_i ,
\]
each term extended by $0$ outside $N_{x_i}$ (it has compact support in $N_{x_i}$: $\mathrm{supp}(\rho_i)\cap\Omega$ is a
compact subset of $\Xi_i^{-1}(B_{1,i}\times\{0\})$, so its $y$-projection is compact in $B_{1,i}$). The sum is finite on each compact set, so
$\mathrm{ext}$ is a continuous linear map $C_c^\infty(A;\mathcal S_\flat)\to C_c^\infty(V_\ell;\mathcal S_\flat)$ for every
compact $A\subset\Omega$, and $\mathrm{ext}(f)|_\Omega=\sum_i\rho_if=f$. Put $E(f)=\Lambda(\mathrm{ext}(f))$. By Step~A,
$\Lambda(g)=E(g|_\Omega)$ for every $g\in C_c^\infty(V_\ell;\mathcal S_\flat)$ ($g-\mathrm{ext}(g|_\Omega)$ vanishes on
$\Omega$), and $E(f)=\Lambda(g)$ for every compactly supported smooth extension $g$ of $f$. The group $G_\tau$ acts on
$C_c^\infty(\Omega;\mathcal S_\flat)$ by $(h\cdot f)(x)=\omega^\flat_\tau(1,h)f(h^{-1}x)$; $h\cdot\mathrm{ext}(f)$ (the action of
Lemma~\ref{pcn:mixed}(c)) is a compactly supported smooth extension of $h\cdot f$, so by (ii) and
Lemma~\ref{pcn:mixed}(c), $E(h\cdot f)=\Lambda(h\cdot\mathrm{ext}(f))=\lambda(\omega_\tau(1,h)\iota(\mathrm{ext}f))=E(f)$.

\emph{Step C: descent to the fibre.} Let $dg$, $ds$ be Haar measures on $G_\tau$, $S_1$. For $F\in C_c^\infty(G_\tau;\mathcal
S_\flat)$ put $f_F(gx_1)=\int_{S_1}\omega^\flat_\tau(1,gs)F(gs)\,ds$. By Lemma~\ref{pcn:smoothaction}(ii), $g\mapsto
\omega^\flat_\tau(1,g)F(g)$ is smooth with compact support, so the integral (over a compact part of $S_1$ for $g$ in a
compact set) defines a smooth right $S_1$-invariant function of $g$, hence (local sections) an element of
$C_c^\infty(\Omega;\mathcal S_\flat)$, and by the estimate of Lemma~\ref{pcn:smoothaction}(ii) $F\mapsto f_F$ is continuous
from $C_c^\infty(K;\mathcal S_\flat)$ to $C_c^\infty(Kx_1;\mathcal S_\flat)$ for each compact $K\subset G_\tau$ (topologies of
uniform convergence of all derivatives in all $p_k$). It is onto: given $f$, let $K_1$ be a compact neighbourhood of
$\mathrm{supp}f$; choose $a\in C_c^\infty(G_\tau)$, $a\ge0$, with $A(g)=\int_{S_1}a(gs)\,ds>0$ for $gx_1\in K_1$ (cover
the compact set $K_1$ by finitely many open sets $U_jx_1$, $U_j\subset G_\tau$ open relatively compact, and take
$a=\sum a_j$ with $a_j\ge0$ positive on $U_j$), and put $F(g)=a(g)A(g)^{-1}\omega^\flat_\tau(1,g^{-1})f(gx_1)$ where
$gx_1\in K_1$, $F=0$ elsewhere; $F$ is smooth with compact support (Lemma~\ref{pcn:smoothaction}(ii) applied after
the inversion $g\mapsto g^{-1}$) and $f_F=f$, as $A(gs)=A(g)$. Let $\tilde E(F)=E(f_F)$, a linear functional continuous on
each $C_c^\infty(K;\mathcal S_\flat)$. For $(L_hF)(g)=F(h^{-1}g)$, the substitution $g\mapsto hg$ gives $f_{L_hF}=h\cdot
f_F$, so $\tilde E\circ L_h=\tilde E$ by Step~B. For $s\in S_1$ and $(R_sF)(g)=\omega^\flat_\tau(1,s)F(gs)$, right
invariance of $ds$ gives $f_{R_sF}=f_F$, so $\tilde E\circ R_s=\tilde E$.

\emph{Haar measure.} Fix $a\in C_c^\infty(G_\tau)$ with $\int a=1$. For $F\in C_c^\infty(G_\tau;\mathcal S_\flat)$ the map
$h\mapsto L_hF$ is continuous into $C_c^\infty(K';\mathcal S_\flat)$ for a compact $K'$ (left translation commutes with
left-invariant vector fields, and compactly supported continuous functions are uniformly continuous), so
$a*F=\int a(h)L_hF\,dh$ exists in that Fr\'echet space, $\tilde E(a*F)=\int a(h)\tilde E(L_hF)dh=\tilde E(F)$, and
evaluating at $g$, $(a*F)(g)=\int a(h)F(h^{-1}g)dh=\int a(gk^{-1})F(k)\,dk$ ($k=h^{-1}g$; $G_\tau$ is unimodular). The
map $k\mapsto a(\cdot\,k^{-1})\otimes F(k)$ is likewise continuous into some $C_c^\infty(K'';\mathcal S_\flat)$, and its
integral has the same values, so equals $a*F$. Hence
\begin{equation}\label{pcn:eq-haar}
\tilde E(F)=\int_{G_\tau}\tilde E\bigl(a(\cdot\,k^{-1})\otimes F(k)\bigr)\,dk .
\end{equation}
For $\varphi\in\mathcal S_\flat$, $T_\varphi(b)=\tilde E(b\otimes\varphi)$ is a continuous linear functional on $C_c^\infty(G_\tau)$
with $T_\varphi(b(h^{-1}\cdot))=T_\varphi(b)$. Applying \eqref{pcn:eq-haar} to $F=b\otimes\varphi$ gives $T_\varphi(b)=
\int b(k)\phi(k)\,dk$ with $\phi(k)=T_\varphi(a(\cdot\,k^{-1}))$, a continuous function. Left invariance gives
$\int b(k)(\phi(hk)-\phi(k))\,dk=0$ for all $b$, so $\phi(hk)=\phi(k)$ and $\phi$ is a constant $\epsilon(\varphi)$:
$T_\varphi(b)=\epsilon(\varphi)\int b$ (a left-invariant distribution is a multiple of Haar measure). $\epsilon$ is linear,
and continuous ($\epsilon(\varphi)=\tilde E(a\otimes\varphi)$). Then \eqref{pcn:eq-haar} and $\int a(gk^{-1})dg=\int a$ give
$\tilde E(F)=\int_{G_\tau}\epsilon(F(k))\,dk$. Finally $R_s(b\otimes\varphi)=b(\cdot\,s)\otimes\omega^\flat_\tau(1,s)\varphi$ and
$\int b(gs)\,dg=\int b$, so $\tilde E\circ R_s=\tilde E$ gives
\[
\epsilon\circ\omega^\flat_\tau(1,s)=\epsilon\qquad(s\in S_1).
\]

\emph{Step D: no invariant functional on the fibre.} The Gram matrix $\beta_\tau$ of $x_1$ is invertible and not
negative definite, hence not negative semidefinite. Let $u$, $w$, $i$ be as in Lemma~\ref{pcn:witt} for $x_1$,
$x''=ux_1$ and $\pcnK'=\{c\in\pcnK_\tau:cx''=x''\}$; then $u^{-1}\pcnK'u\subset S_1$. The functional
$\epsilon'=\epsilon\circ\omega^\flat_\tau(1,u)^{-1}$ is continuous on $\mathcal S_\flat$ (Lemma~\ref{pcn:smooth}(d)), and for
$c\in\pcnK'$, $\epsilon'\circ\omega^\flat_\tau(1,c)=\epsilon\circ\omega^\flat_\tau(1,u^{-1}cu)\circ\omega^\flat_\tau(1,u)^{-1}=\epsilon'$.
By Lemma~\ref{pcn:nofunct} (the local case), $\epsilon'=0$, so $\epsilon=0$.

\emph{Conclusion.} $\tilde E(F)=\int\epsilon(F(k))dk=0$ for all $F$. As $F\mapsto f_F$ is onto, $E=0$; by Step~B,
$\Lambda(g)=E(g|_\Omega)=0$ for every $g\in C_c^\infty(V_\ell;\mathcal S_\flat)$; so $\lambda=0$.
\end{proof}

\subsubsection{Proof of Theorem~\ref{pcn:negcoeff}}\label{pcn:ss-proof}

\emph{Step 0: reduction to coherent states.} Let $\Phi=\Phi^{\mathrm{row}}_{z_0}$. For a Fock vector
$w$ with finite moments, Proposition~\ref{pcx:decay}(a), applied with $P=S$ to the pair $(X,T)$ (its
hypothesis holds: by Lemma~\ref{pcu:rowsurj}(a) every $x$ with $\Phi(x)\ne0$ is surjective at every place of $S$,
and $S\cap\bar S=\emptyset$), gives $|\theta(\omega(p,h)[w,\Phi]_S)|\le C\|w\|_M\,\mathrm{ht}(h)^{-N_0}$, with $C$
uniform for $p$ in a compact set. By Lemma~\ref{pcx:superpos},
$\theta(\omega(p,h)[w,\Phi]_S)=\int w(y)\,\theta(\omega(p,h)[\psrho(y)e,\Phi]_S)\,dy$, and since
$\|\psrho(y)e\|_M\le(1+|y|)^M\|e\|_M$ and $\int|w(y)|(1+|y|)^Mdy=\|w\|_M<\infty$, Fubini (with (R1) for the
integral over $[\auU(T)]$ and compactness of $Z_D(E)\backslash Z_D(\auA)$) gives
\[
\pcuF_\beta\Theta^{\mathrm L}_{X,T}(w,\Phi)(p)=\int w(y)\,\pcuF_\beta\Theta^{\mathrm L}_{X,T}(\psrho(y)e,\Phi)(p)\,dy .
\]
So it suffices to prove the theorem for $v=\psrho(y)e$. Up to a factor of modulus $1$ this vector is a product
$v=v_\tau\otimes v^\tau$, $v^\tau=\bigotimes_{w\ne\tau}v_w$, of the vectors $v_w=\psrho_w(y_w)e_w$, which lie in the spaces
$\mathcal S(\pzF_w)$ of Lemma~\ref{pcn:smooth} (the vacuum is a polynomial vector, and Lemma~\ref{pcn:smooth}(d)).

\emph{Step 1: the coefficient as a continuous functional at $\tau$.} Let $\pzF_\infty$ be the archimedean
Fock space of $\auHom(X,T)_\R$, of complex dimension $D_\infty$, $\pzF_\tau$ its factor at $\tau$, and $\mathcal S_\tau=
\mathcal S(\pzF_\tau)$. Write $p=p_\infty p_f$ and $p_\infty=(p_w)_w$, and put $u^\tau=\bigotimes_{w\ne\tau}\omega_w(p_w,1)v_w$,
a vector of $\mathcal S(\bigotimes_{w\ne\tau}\pzF_w)$ (Lemma~\ref{pcn:smooth}(d)). For $\varphi\in\mathcal S_\tau$ let
\[
\lambda(\varphi)=\pcuF_\beta\Theta^{\mathrm L}_{X,T}\bigl(w_\varphi,\Phi\bigr)(p),\qquad
w_\varphi=\omega_\infty(p_\infty,1)^{-1}(\varphi\otimes u^\tau),
\]
so that $\lambda(\omega_\tau(p_\tau,1)v_\tau)$ is the coefficient to be computed ($w_\varphi=v$ for this $\varphi$). Fix $M$ as
in Step~0 and an integer $s$ with $2s>M+2D_\infty$. By Lemma~\ref{pcn:smooth}(c), (d) and the tensor product
inequality for $p_s$,
\[
\|w_\varphi\|_M\le C_{M,s}\|w_\varphi\|_s\le C\|\varphi\otimes u^\tau\|_s\le C'p_s(\varphi\otimes u^\tau)\le C'p_s(u^\tau)\,p_s(\varphi),
\]
with $C$ depending on $\|\sigma_{p_\infty}\|$ only. By Step~0's bound,
\[
|\lambda(\varphi)|\le\sup_{z}\int_{[\auU(T)]}\bigl|\theta\bigl(\omega(n(0,z)p,h)[w_\varphi,\Phi]_S\bigr)\bigr|\,dh
\le C''\|w_\varphi\|_M\le C'''p_s(\varphi)
\]
($z$ in a compact fundamental domain of $Z_D(E)\backslash Z_D(\auA)$, so $n(0,z)p$ stays in a compact set; (R1) for
$\int\mathrm{ht}^{-N_0}<\infty$; probability measure on the centre). So $\lambda$ is a well-defined continuous linear
functional on $\mathcal S_\tau$, every integral involved converging absolutely.

\emph{Step 2: the two invariance properties.} (ii) Let $h_0\in G_\tau$, placed at $\tau$. As $\omega_\tau(1,h_0)$
commutes with $\omega_\infty(p_\infty,1)$ and with the other places,
$\theta(\omega(n(0,z)p,h)[w_{\omega_\tau(1,h_0)\varphi},\Phi]_S)=\theta(\omega(n(0,z)p,hh_0)[w_\varphi,\Phi]_S)$. The integral
over $[\auU(T)]$ of an absolutely integrable function is invariant under right translation by $h_0$ (the quotient
measure of the unimodular group $\auU(T)(\auA)$ by the discrete subgroup $\auU(T)(E)$ is right invariant), so,
integrating then over $z$ (Fubini, by the bounds of Steps~0--1), $\lambda(\omega_\tau(1,h_0)\varphi)=\lambda(\varphi)$.
(i) Let $z_\tau\in\mathfrak z_\tau$ (the centre at $\tau$, \S\ref{pcn:ss-local}) and $n_\tau=n(0,z_\tau)$, placed at $\tau$. Since
$\omega(p,h)[w,\Phi]_S=\omega(1,h)[\omega_\infty(p_\infty,1)w,\omega_f(p_f)\Phi]_S$, we get, with
$\Phi'=\omega_f(p_f)\Phi$,
(here $[\cdot,\Phi']$ denotes the datum whose finite part is $\omega_f(p_f)$ applied to the finite part of $[\cdot,\Phi]_S$)
\[
\theta\bigl(\omega(n(0,z)p,h)[w_{\omega_\tau(n_\tau,1)\varphi},\Phi]_S\bigr)
=\theta\bigl(\omega(n(0,z)n_\tau,h)[\varphi\otimes u^\tau,\Phi']\bigr)
=\theta\bigl(\omega(n(0,z+z_\tau)p,h)[w_\varphi,\Phi]_S\bigr),
\]
because $n(0,\cdot)$ is a homomorphism from the additive group into $Z_D(\auA)$ and $n_\tau$ commutes with the
finite components. Integrating over $h$ and then against $\psi_\beta(z)^{-1}dz$ over the compact group
$Z_D(E)\backslash Z_D(\auA)$ and substituting $z\mapsto z-z_\tau$ (invariance of its Haar measure; the coset of $z+z_\tau$
is well defined since $z_\tau\in Z_D(\auA)$) gives
$\lambda(\omega_\tau(n_\tau,1)\varphi)=\psi_{\beta,\tau}(z_\tau)\lambda(\varphi)$ with $\psi_{\beta,\tau}(z_\tau)=
\psi_\tau(-\frac12\jmath\operatorname{tr}(z_\tau\tau(\beta)))=e(z_\tau,\tau(\beta))$.

\emph{Step 3: conclusion.} Let $\tau$ be the infinite place of the theorem, where $\tau(\beta)$ is invertible and not
negative definite; (H1)--(H3) hold (\S\ref{pcn:ss-statement}). By Steps~1--2, $\lambda$ satisfies the hypotheses of
Lemma~\ref{pcn:local} with $\beta_\tau=\tau(\beta)$ (on $\mathcal S_\tau$ itself; no change of model is needed), so $\lambda=0$; in particular
$\pcuF_\beta\Theta^{\mathrm L}_{X,T}(v,\Phi)(p)=\lambda(\omega_\tau(p_\tau,1)v_\tau)=0$. With Step~0 this proves the theorem
for every $v$ with finite moments. \qed

\begin{remark}[What the argument does not use]\label{pcn:rem}
The proof uses neither the unfolding (U1) nor any integrability of a term over a stabiliser quotient: the only
global inputs are the absolute convergence of the theta integral and of its coefficient
(Proposition~\ref{pcx:decay}(a), (R1)) and the right invariance of the measure on $[\auU(T)]$. The orbit
$\{x:x^*J_Tx=\beta\}$ and its stabiliser enter only locally at $\tau$, through Lemma~\ref{pcn:local}. The
finite part of $\Phi$ plays no role beyond the hypothesis of Proposition~\ref{pcx:decay}(a), and nothing is assumed
at the infinite places other than $\tau$.
\end{remark}

\subsubsection{Dependencies}\label{pcn:ss-dep}
With Theorem~\ref{pcn:negcoeff}, Lemma~\ref{pct:negcoeff} is proved, and with it
Proposition~\ref{pct:vantop} and Proposition~\ref{pcx:P-nonvan}. The proof uses results proved
in the paper (Lemma~\ref{pcu:laws} in its local form, Corollary~\ref{pcu:laws-local}; Lemmas~\ref{pcu:rowsurj}, \ref{pcu:orbitvol}, \ref{pcx:superpos} and~{\cite[Lemma~8.2]{P49}}; Proposition~\ref{pcx:decay}; {\cite[Lemma~3.6]{P49}} with
{\cite[Theorem~3.29]{P49}}; the mixed model of \S\ref{pcu:ss-mixed}), the facts on smooth vectors, mixed models, invariant
distributions and adapted charts proved in \S\S\ref{pcn:ss-smooth}, \ref{pcn:ss-local}, and the following published
results: Witt's theorem for hermitian spaces over $\C$; the Fock-model formula for the action of the
maximal compact subgroup (background reference: Folland~\cite[Prop.~4.39]{FollandPS}, Kashiwara--Vergne~\cite{KashiwaraVergne}; the proof rests on {\cite[Lemma~3.6]{P49}}); Stone's theorem on
one-parameter unitary groups; the inverse function theorem; reduction theory (R1). It does not use (U1). Together with
\S\ref{pct:ss-status} and \S\ref{pcu:ss-final}, Theorem~\ref{pcx:X} depends only on published results (those
listed in \S\ref{pct:ss-status}, and the ones just listed).



\section{The triangle of lifts}\label{C1:g3}
% ===== C-triangle.tex =====
% proofs/C-triangle.tex
% Archimedean types at the nodes of the triangle (Lemma pc:arch for the pairs of the triangle) and
% the analytic steps (ii), (iii) of Theorem pc:tri. Labels of this file start with ptr:.
% Only amsmath/amssymb/amsthm are used. Conventions and macros: proofs/C-determinant.tex.
\providecommand{\auU}{\mathrm U}
\providecommand{\auPl}{\Phi}
\providecommand{\auHom}{\operatorname{Hom}}
\providecommand{\auNm}{\operatorname{Nm}}
\providecommand{\auGL}{\mathrm{GL}}
\providecommand{\auA}{\mathbb A}
\providecommand{\aubar}[1]{\overline{#1}}
\providecommand{\psrho}{\varrho}
\providecommand{\pzF}{\mathfrak F}
\providecommand{\pcF}{\mathcal F}
\providecommand{\pcG}{\mathsf G}
\providecommand{\pcm}{\mathfrak m}
\providecommand{\pcR}{R}
\providecommand{\pcMat}{\mathrm M}
\providecommand{\pdsm}{m}
\providecommand{\pdsK}{K}
\providecommand{\ptrnot}[1]{\par\smallskip{\raggedright\noindent\textsc{Not proved here.} [#1]\par}\smallskip}
\providecommand{\ptrwt}{\operatorname{wt}}
\providecommand{\ptrht}{\operatorname{ht}}
\providecommand{\ptrsgn}{\operatorname{sgn}}

\subsection{The archimedean types at the nodes of the triangle}\label{ptr:sec}

This section treats the archimedean part of Lemma~\ref{pc:arch} for the pairs of the triangle of
\S\ref{pc:ss-triangle}, and the steps (ii) (convergence) and (iii) (archimedean multiplicity one)
of the deduction of Theorem~\ref{pc:tri} described at the end of \S\ref{pc:ss-exchange}. The
notation is that of \S\ref{pc:ss-triangle}, of {\cite[\S3.1]{P49}} and of {\cite[\S8.1]{P49}}; inner
products are conjugate-linear in the first variable.

\emph{Contents of this section.}
\begin{enumerate}
\item[(1)] The Harish-Chandra parameters of the types $e_{i,n}$ at the nodes, as defined in
\S\ref{pc:ss-triangle}, \emph{match} (in the sense of {\cite[(13)]{P49}}) along every
ascending pair $(\pcG_{i,n},\pcG_{i,n+1})$ and along every descending pair
$(\pcG_{i+1,n},\pcG_{i,n+1})$, for the characters of the lifts $a_{i,n}$ and for the characters of
the descents determined in Definition~\ref{ptr:desc-char}; the latter are the characters forced by
the hypotheses of the exchange identity (Lemma~\ref{ptr:chars}). This is checked from the
parameters themselves (Proposition~\ref{ptr:match}).
\item[(2)] Consequently (A1) (integrability, cuspidality, rapid decay) and (A2) (the lifts and the
transposed lifts with the joint harmonic preserve the types: the \emph{edge} property) hold at
all nodes and along all these pairs, and the joint harmonic exists and is unique
(Theorem~\ref{ptr:edge}).
\item[(3)] Convergence of all integrals that occur in the squares (Proposition~\ref{ptr:conv}).
\item[(4)] $Z_\infty\neq0$ is reduced to the positivity (F) of the Fock coefficients of a termwise
harmonic polynomial at positive definite elements (Proposition~\ref{ptr:Z}); (F) is
Theorem~\ref{ptq:F}.
\item[(5)] Archimedean multiplicity one (step (iii)) is reduced to the local smoothing statement
(H) (Theorem~\ref{ptr:square}); (H) is Theorem~\ref{ptq:H}.
\end{enumerate}
Step (v) of \S\ref{pc:ss-exchange} (the adjoint law, Theorem~\ref{pc:tri}(c)) is treated in
\S\ref{ptv:sec}.

\subsubsection{Parameters at the nodes}\label{ptr:ss-param}

\paragraph*{Notation}
Recall from \S\ref{pc:ss-triangle}: $\pcG_{i,n}$ $(i,n\ge0,\ i+n\le m)$ has rank $N_{i,n}=r+nt$, a
stem of rank $r$ and $n$ positive definite blocks of rank $t$ numbered $i+1,\dots,i+n$. The stem
coordinates are numbered $p=1,\dots,r$; we put $a(p)=(p-k)^+$. The sign of the stem coordinate $p$
is $\epsilon_p^{(i,n)}=(-1)^{E_{i,n}(a(p))}$ with
\begin{equation}\label{ptr:eq-E}
E_{i,n}(a)=i+a+\min\bigl((a-i)^+,n\bigr)=\begin{cases}i+a,&a\le i,\\ 2a,&i<a\le i+n,\\
i+a+n,&a>i+n.\end{cases}
\end{equation}
The stem of the parameter consists of the numbers of $\varsigma_\tau$ (the case $j=0$ of
{\cite[(74)]{P49}}): the $k-1$ numbers $\tfrac k2-1,\dots,1-\tfrac k2$ (these are the coordinates
$p<k$) and the numbers $-\kappa(\tau)-i'\Delta$, $1\le i'\le m$ (the coordinate $p=k-1+i'$, so that
$a(p)=i'-1$). The parameter $\Lambda^{(i,n)}_\tau$ of $e_{i,n}$ is: the stem numbers with the signs
$\epsilon^{(i,n)}$, and for each block $i+l$ $(1\le l\le n)$ the $t$ numbers
$-\kappa_{i+l}(\tau)+\tfrac{t-1}2-b$ $(0\le b<t)$ with sign $+$; all numbers lowered by
$\tfrac12$ if $n$ is odd. Here $\kappa_l=\kappa+(l+\tfrac12)\Delta$.

\begin{lemma}[Consistency with the lower chain and with the first column]\label{ptr:cons}
\textup{(a)} $\Lambda^{(0,n)}_\tau=\Lambda^{(n)}_\tau$ of {\cite[(74)]{P49}}, with the same signs.
\textup{(b)} $\epsilon^{(i,0)}=(-1)^i\epsilon^{(0,0)}$, and $\Lambda^{(i,0)}_\tau$ and $\Lambda^{(0,0)}_\tau$
define the same discrete series of the common group $\auU(\pcG_{i,0})=\auU(Y_0)$.
\textup{(c)} The number of negative signs of $\pcG_{i,n}$ is at most $r$; the ranks $N_{i,n}$ are
odd exactly when $n$ is even (so $r$ and $t$ are odd).
\end{lemma}
\begin{proof}
(a) For $i=0$, \eqref{ptr:eq-E} gives the sign $+$ when $a\le n$ and $(-1)^{a+n}$ when $a>n$. For
the number $-\kappa-i'\Delta$ ($a=i'-1$) the sign is $-$ exactly when $i'\ge n+2$ and
$i'\equiv n\pmod2$, which is the rule of {\cite[(74)]{P49}} for $j=n$; the small numbers ($a=0$)
have sign $+$; the blocks and the half-shift are the same. (b) For $n=0$, $E_{i,0}(a)=i+a$ in all
cases, while $E_{0,0}(a)=a$. The group of a hermitian form and of its negative is the same, the
compact Cartan subgroup and the root system are the same, and the compact roots are those
$e_a-e_b$ with equal signs, which does not change when all signs change; the discrete series with
Harish-Chandra parameter $\lambda$ is characterised by $\lambda$ and the set of compact roots
(Harish-Chandra; \cite[Theorem~9.20, p.~310]{Knapp}, stated for linear connected semisimple groups; it
applies to $\auU(p,q)=Z\cdot\mathrm{SU}(p,q)$, $Z\cong\auU(1)$ the compact centre,
$Z\cap\mathrm{SU}(p,q)$ finite, since a discrete series of $\auU(p,q)$ restricts irreducibly to one of
$\mathrm{SU}(p,q)$ and $Z$ acts by a character, which is determined by $\lambda$ as $Z$ is connected;
see the paragraph ``On (D3)'' in {\cite[\S3.1.3]{P49}}). (c) The negative signs occur only on
the stem. For (c), second part: $N_j=r+jt$ is odd exactly when $j$ is even (proof of
{\cite[Lemma~3.6]{P49}}); hence $r$ is odd and $t$ is odd.
\end{proof}

\paragraph*{The derived constants for arbitrary characters}
For automorphic characters $\alpha,\beta$ of $\auU(1)$ write $\ptrwt\alpha=(\ptrwt_\tau\alpha)_\tau$
for their weights, normalised as in {\cite[\S3.1]{P49}} ($\upsilon_\tau$ has weight $1$,
$\upsilon_k=\upsilon^{\lfloor k/2\rfloor}$, $\mu_j$ has weight $-\kappa_j$).

\begin{lemma}\label{ptr:const}
Let $S,T$ be spaces of ranks $n<n'$ and $\alpha,\beta$ characters. Under \textup{(V)}
({\cite[Lemma~3.6]{P49}}, proved in {\cite[Theorem~3.29]{P49}}) the derived actions of
$\omega_{\alpha,\beta,\infty}(\cdot,1)$ and $\omega_{\alpha,\beta,\infty}(1,\cdot)$ at $\tau$ on
polynomial vectors are {\cite[(12)]{P49}} with $\gamma=\tau_0+\tfrac{n'}2$, $\gamma'=\tau'_0+\tfrac n2$,
\[
\tau_0=\tfrac{n'}2-\ptrwt_\tau\alpha,\qquad\tau'_0=\ptrwt_\tau\beta-\tfrac n2 .
\]
\end{lemma}
\begin{proof}
The proof of {\cite[Lemma~3.6]{P49}} uses only (V) and the weights of the two characters: the factor
$\alpha(\det g)^{-1}\beta(\det h)$ contributes $\det(g)^{-\ptrwt\alpha}\det(h)^{\ptrwt\beta}$ at
$\tau$, (V) contributes $\det(g)^{n'/2}\det(h)^{-n/2}$, and the Weyl-ordered compact operators carry
the constants $\tfrac{n'}2$, $\tfrac n2$. (For $\omega_{\mu_{j+1}}=\omega_{\mu_{j+1}\upsilon_{n'},
\mu_{j+1}\upsilon_n}$ this gives back $\tau_0=\kappa_{j+1}+\tfrac{n'}2-\lfloor\tfrac{n'}2\rfloor$,
$\tau'_0=-\kappa_{j+1}+\lfloor\tfrac n2\rfloor-\tfrac n2$.)
\end{proof}

By the computation in the proof of {\cite[Lemma~3.7]{P49}} (which uses only the definitions of
$\xi,\xi'$), for a row parameter $\lambda$ and a column parameter $\lambda'$ the conditions
{\cite[(13)]{P49}} are equivalent to the following \emph{block rule}: there is a set $B$ of $t$
columns and a bijection $r\mapsto c(r)$ of the rows onto the other columns such that
\begin{equation}\label{ptr:eq-rule}
\begin{gathered}
\lambda'_{c(r)}=-\lambda_r+\tau_0+\tau'_0,\qquad|\lambda'_{c(r)}-\tau'_0|>\tfrac t2,\qquad
\varepsilon'_{c(r)}=\varepsilon_r\,\ptrsgn(\lambda'_{c(r)}-\tau'_0),\\
\{\lambda'_c:c\in B\}=\{\tau'_0+\tfrac{t-1}2-b:0\le b<t\},\qquad\varepsilon'\text{ constant on }B.
\end{gathered}
\end{equation}
Indeed $\xi_r=\lambda_r-\tau_0$, $\xi'_c=\lambda'_c-\tau'_0$, and $\eta_{rc(r)}=-\varepsilon_r
\varepsilon'_{c(r)}=\ptrsgn\xi_r=-\ptrsgn\xi'_{c(r)}$. In words: the column parameter is the
negative of the row parameter shifted by $\tau_0+\tau'_0$, plus a string of $t$ numbers of one
sign centred at $\tau'_0$; the signs are kept above the centre and changed below it.

\paragraph*{The two kinds of pairs}
An \emph{ascending pair} is $(\pcG_{i,n},\pcG_{i,n+1})$ with the characters of $a_{i,n}$:
$\alpha=\mu_{i+n+1}\upsilon_{r+(n+1)t}$, $\beta=\mu_{i+n+1}\upsilon_{r+nt}$; rows are the
coordinates of $\pcG_{i,n}$, with $\lambda=-\Lambda^{(i,n)}_\tau$; columns those of $\pcG_{i,n+1}$,
with $\lambda'=\Lambda^{(i,n+1)}_\tau$. Thus, as in {\cite[\S3.1]{P49}}, the lift with the joint
harmonic is to take forms of type $e_{i,n}$ to forms of type $e_{i,n+1}$.

A \emph{descending pair} is $(\pcG_{i+1,n},\pcG_{i,n+1})$, $i+n+1\le m$, rows the coordinates of
$\pcG_{i+1,n}$ with $\lambda=\Lambda^{(i+1,n)}_\tau$, columns those of $\pcG_{i,n+1}$ with
$\lambda'=-\Lambda^{(i,n+1)}_\tau$. Then (as in {\cite[Proposition~8.4(c)]{P49}}) the transposed lift
with the joint harmonic is to take forms of type $e_{i,n+1}$ (the conjugate of the type of the
column parameter $-\Lambda^{(i,n+1)}$) to forms of type $e_{i+1,n}$ (the type of the row parameter):
this is the shape of the descent $d_{i,n}\colon\pcF_{i,n+1}\to\pcF_{i+1,n}$.

\begin{definition}[Characters of the descents]\label{ptr:desc-char}
For $i+n+1\le m$ let $\alpha^\downarrow_{i,n},\beta^\downarrow_{i,n}$ be automorphic characters of
$\auU(1)$ with
\[
\ptrwt_\tau\alpha^\downarrow_{i,n}=\kappa_{i+1}(\tau)+\Bigl\lceil\frac{r+(n+1)t}2\Bigr\rceil,\qquad
\ptrwt_\tau\beta^\downarrow_{i,n}=\kappa_{i+1}(\tau)+\Bigl\lceil\frac{r+nt}2\Bigr\rceil
\]
(their finite parts are fixed in Lemma~\ref{ptr:chars}). The descent kernel is
$\theta_{\alpha^\downarrow_{i,n},\beta^\downarrow_{i,n}}$ for $\auHom(\pcG_{i+1,n},\pcG_{i,n+1})$.
\end{definition}

\begin{proposition}[Matching at the nodes]\label{ptr:match}
Under \textup{(V)}, for every $\tau$:
\textup{(a)} every ascending pair satisfies the block rule \eqref{ptr:eq-rule}, with
$\tau_0=\kappa_{i+n+1}+\tfrac12[n\text{ odd}]$, $\tau'_0=-\kappa_{i+n+1}-\tfrac12[n\text{ even}]$, the
string $B$ being block $i+n+1$;
\textup{(b)} every descending pair, with the characters of Definition~\ref{ptr:desc-char},
satisfies \eqref{ptr:eq-rule}, with $\tau_0=-\kappa_{i+1}-\tfrac12[n\text{ odd}]$,
$\tau'_0=\kappa_{i+1}+\tfrac12[n\text{ even}]$, the string $B$ being block $i+1$;
\textup{(c)} conversely, the weights of Definition~\ref{ptr:desc-char} are the only ones for which
\textup{(b)} holds with the string $B$ equal to block $i+1$ and the old coordinates paired with the
same coordinates.
\end{proposition}
\begin{proof}
Throughout, $[P]$ is $1$ if $P$ holds and $0$ otherwise, and we use Lemma~\ref{ptr:cons}(c):
$N_{i,n}=r+nt$ is odd iff $n$ is even, so $\tfrac{N_{i,n}}2-\lfloor\tfrac{N_{i,n}}2\rfloor=\tfrac12
[n\text{ even}]$ and $\lceil\tfrac{N_{i,n}}2\rceil-\tfrac{N_{i,n}}2=\tfrac12[n\text{ even}]$.

(a) \emph{Constants.} By Lemma~\ref{ptr:const} with $n'=N_{i,n+1}$, $\ptrwt\alpha=-\kappa_{i+n+1}+
\lfloor n'/2\rfloor$: $\tau_0=\kappa_{i+n+1}+\tfrac12[n+1\text{ even}]$; with $\ptrwt\beta=
-\kappa_{i+n+1}+\lfloor N_{i,n}/2\rfloor$: $\tau'_0=-\kappa_{i+n+1}-\tfrac12[n\text{ even}]$. So
$\tau_0+\tau'_0=\tfrac12[n\text{ odd}]-\tfrac12[n\text{ even}]$.
\emph{Old numbers.} With $\lambda=-\Lambda^{(i,n)}$ the first condition of \eqref{ptr:eq-rule}
reads $\Lambda^{(i,n+1)}_{c(r)}=\Lambda^{(i,n)}_r+\tau_0+\tau'_0$. Pair each stem coordinate and
each coordinate of the blocks $i+1,\dots,i+n$ of $\pcG_{i,n}$ with the same coordinate of
$\pcG_{i,n+1}$. Both carry the same unshifted number, lowered by $\tfrac12$ according to the parity
of $n$, resp.\ $n+1$; the difference is $-\tfrac12$ if $n$ is even and $+\tfrac12$ if $n$ is odd,
which is $\tau_0+\tau'_0$.
\emph{New string.} Block $i+n+1$ of $\pcG_{i,n+1}$ carries $-\kappa_{i+n+1}+\tfrac{t-1}2-b$
lowered by $\tfrac12[n+1\text{ odd}]=\tfrac12[n\text{ even}]$: the string centred at $\tau'_0$, with
the constant sign $+$.
\emph{Distances.} The block $i+l$ $(l\le n)$ is centred at $-\kappa_{i+l}\pm\tfrac12$, at distance at
least $(n+1-l)\Delta-1\ge\Delta-1$ from $\tau'_0$, so its numbers are at distance $\ge\Delta-\tfrac{t+1}2>
\tfrac t2$. The stem numbers are $\tfrac k2-1,\dots$ (at distance $>\kappa_{i+n+1}-k>\tfrac t2$) and
$-\kappa-i'\Delta$, at distance $|(i+n+\tfrac32-i')\Delta|-\tfrac12\ge\tfrac12\Delta-\tfrac12>\tfrac t2$,
since $\Delta-t\ge r^8+2$ ({\cite[\S8.1]{P49}}).
\emph{Signs.} The blocks lie above the centre (their centres $-\kappa_{i+l}$ exceed
$-\kappa_{i+n+1}$) and keep the sign $+$; the small stem numbers keep their sign $(-1)^i$. The number $-\kappa-i'\Delta$,
$a=i'-1$, lies above the centre iff $i'<i+n+\tfrac32$, i.e.\ $a\le i+n$. By \eqref{ptr:eq-E}:
for $a\le i$, $E_{i,n}(a)=E_{i,n+1}(a)=i+a$; for $i<a\le i+n$ both are $2a$; for $a=i+n+1$,
$E_{i,n}(a)=2i+2n+1$ is odd and $E_{i,n+1}(a)=2a$ is even; for $a>i+n+1$, $E_{i,n+1}(a)=E_{i,n}(a)+1$.
So the sign is kept exactly when $a\le i+n$, which is the rule. (The small numbers have $a=0\le i$.)

(b) \emph{Constants.} Here $n'=N_{i,n+1}$ (columns) and the rows have rank $N_{i+1,n}$. By
Lemma~\ref{ptr:const}, $\tau_0=\tfrac{n'}2-\kappa_{i+1}-\lceil\tfrac{n'}2\rceil=-\kappa_{i+1}-\tfrac12
[n+1\text{ even}]=-\kappa_{i+1}-\tfrac12[n\text{ odd}]$, and $\tau'_0=\kappa_{i+1}+\lceil
\tfrac{N_{i+1,n}}2\rceil-\tfrac{N_{i+1,n}}2=\kappa_{i+1}+\tfrac12[n\text{ even}]$. So
$\tau_0+\tau'_0=\tfrac12[n\text{ even}]-\tfrac12[n\text{ odd}]$.
\emph{Old numbers.} With $\lambda=\Lambda^{(i+1,n)}$, $\lambda'=-\Lambda^{(i,n+1)}$ the first condition
reads $\Lambda^{(i,n+1)}_{c(r)}=\Lambda^{(i+1,n)}_r-(\tau_0+\tau'_0)$. Pair the stem coordinates, and
the blocks $i+2,\dots,i+n+1$ of $\pcG_{i+1,n}$ with the same blocks of $\pcG_{i,n+1}$. Both carry
the same unshifted number, lowered according to $n$, resp.\ $n+1$; the difference
$\Lambda^{(i+1,n)}-\Lambda^{(i,n+1)}$ is $-\tfrac12[n\text{ odd}]+\tfrac12[n\text{ even}]=\tau_0+\tau'_0$.
\emph{New string.} Block $i+1$ of $\pcG_{i,n+1}$ carries $-\kappa_{i+1}+\tfrac{t-1}2-b-\tfrac12
[n\text{ even}]$; its negatives, the entries of $\lambda'$, are $\kappa_{i+1}+\tfrac12[n\text{ even}]
-\tfrac{t-1}2+b$, the string centred at $\tau'_0$, with the constant sign $+$.
\emph{Distances.} The entries of $\lambda'$ on old coordinates are the negatives of those of
$\Lambda^{(i,n+1)}$: the blocks $i+l$, $l\ge2$, centred at $\kappa_{i+l}\pm\tfrac12$, at distance
$\ge\Delta-1$ from $\tau'_0$; the small stem numbers, at distance $>\kappa_{i+1}-k$; the numbers
$\kappa+i'\Delta\pm\tfrac12$, at distance $\ge\tfrac12\Delta-\tfrac12$. All exceed $\tfrac t2$.
\emph{Signs.} The blocks $i+l$, $l\ge2$, lie above $\tau'_0$ and keep $+$. The small stem
numbers (negated) lie below and change sign: indeed $E_{i+1,n}(0)=i+1$, $E_{i,n+1}(0)=i$. The number
$\kappa+i'\Delta$ ($a=i'-1$) lies above $\kappa_{i+1}+\tfrac12[\,]$ iff $i'\ge i+2$, i.e.\ $a\ge i+1$.
By \eqref{ptr:eq-E}: for $a\le i$, $E_{i+1,n}(a)=i+1+a$ and $E_{i,n+1}(a)=i+a$ differ by one; for
$a=i+1$, $E_{i+1,n}(a)=2i+2$ and $E_{i,n+1}(a)=2a$ are both even; for $i+1<a\le i+n+1$ both are $2a$;
for $a>i+n+1$ both are $i+a+n+1$. So the sign is changed exactly when $a\le i$, i.e.\ below the
centre, which is the rule.

(c) The block rule determines $\tau'_0$ (the centre of the new string) and $\tau_0+\tau'_0$ (the
shift), hence $\tau_0,\tau'_0$, hence by Lemma~\ref{ptr:const} the weights.
\end{proof}

\begin{remark}\label{ptr:rem-desc}
For $n=0$ the descending pair is $(\pcG_{i+1,0},\pcG_{i,1})=(\pcG_{i,0}^-,\pcG_{i,1})$, and by (W7)
(Lemma~\ref{pc:SW}) the complex conjugate of the kernel of $a_{i,0}$ is a kernel for this pair,
with parameters negated: $(-\Lambda^{(i,0)},\Lambda^{(i,1)})\mapsto(\Lambda^{(i,0)},-\Lambda^{(i,1)})$,
and $\Lambda^{(i,0)}$, $\Lambda^{(i+1,0)}$ define the same representation (Lemma~\ref{ptr:cons}(b)).
This is the descending pair of Proposition~\ref{ptr:match}(b) for $n=0$, consistent with
$d_{i,0}=a_{i,0}^*$ (Theorem~\ref{pc:tri}(a)); by (c) the weights of the conjugate characters are
those of Definition~\ref{ptr:desc-char}.
\end{remark}

\begin{lemma}[Gap and regularity]\label{ptr:gap}
Every $\Lambda^{(i,n)}_\tau$ consists of $N_{i,n}$ distinct numbers in $\tfrac{N_{i,n}-1}2+\Z$, two
numbers with different signs differ by at least $\tfrac12(\Delta-t)-1\ge\tfrac12r^8$, and
$N_{i,n}\le N_m$. Hence the conclusions of {\cite[Theorem~3.17]{P49}}, {\cite[Lemma~3.25]{P49}},
{\cite[Theorem~3.2]{P49}} and {\cite[Corollary~3.3]{P49}} hold for the types $e_{i,n}$.
\end{lemma}
\begin{proof}
The numbers are those of the stem of $\varsigma$ and of $n$ strings, all shifted by the same
$-\tfrac12[n\text{ odd}]$. The parity: for $n=0$ they lie in $\tfrac{r-1}2+\Z$ by
{\cite[(74)]{P49}} ($j=0$); every string adds $t$ (odd) to the rank and the shift by $\tfrac12$
when $n$ changes parity restores $\tfrac{N-1}2+\Z$. Distinctness and the gap: by \eqref{ptr:eq-E}
negative signs occur only on stem numbers (the small numbers when $i$ is odd, and some of the
numbers $-\kappa-i'\Delta$, shifted); a small number is at distance $\ge\kappa-k$ from every
non-small number, and a number $-\kappa-i'\Delta$ is at distance
$\ge\min(\tfrac12\Delta-\tfrac{t-1}2-\tfrac12,\ \Delta)$ from the other non-small numbers; both
bounds are at least $\tfrac12(\Delta-t)-1$ as in {\cite[\S8.1]{P49}}; distinct strings are $\Delta$ apart, and within
a string the numbers are consecutive. As $N_{i,n}\le r+mt=N_m$, the numerical hypotheses of the four
cited statements (gap $\ge N^2+1$, resp.\ $\ge B+N^2+1$, $\ge5N^3$) hold by the inequality
$\tfrac12r^8\ge500(N_j+N_{j+1})^3+10N_j^5+x^{25}$ of {\cite[\S3.1]{P49}}; these statements use the
parameter only through its regularity, the gap and the rank ({\cite[Remark~3.1]{P49}}, proof of
{\cite[Lemma~3.25]{P49}}).
\end{proof}

\subsubsection{Joint harmonics and the edge property at the nodes}\label{ptr:ss-edge}

For a pair $(S,T)$ (ascending or descending) with row parameter $\lambda$ and column parameter
$\lambda'$ at every $\tau$, let $e_S=d\,\overline{c}$ be the type on $S$ whose conjugate has parameter
$\lambda$ (so $e_S=e_{i,n}$ for an ascending pair and $e_S=\overline{e_{i+1,n}}$ for a descending one)
and $e_T$ the type with parameter $\lambda'$ ($e_{i,n+1}$, resp.\ $\overline{e_{i,n+1}}$). Put
$\mathsf P=\int e_S(x)\,\omega_\infty(x^{-1},1)\,dx$ and $\mathsf P'=\int e_T(x)\,\omega_\infty(1,x)\,dx$
on the Fock space, as in {\cite[Theorem~3.15]{P49}}.

\begin{theorem}[Edges]\label{ptr:edge}
Assume \textup{(V)} and \textup{(D1)--(D3)} of {\cite[\S3.1.3]{P49}}. For every ascending and every
descending pair of the triangle:
\begin{enumerate}
\item[(a)] a joint harmonic $P_{S,T}=\bigotimes_\tau\Phi_\tau$ exists, is unique up to a scalar, and
no monomial of it contains two variables of one column from rows of opposite signs, or of one row
from columns of opposite signs;
\item[(b)] $\mathsf P P_{S,T}=\mathsf P'P_{S,T}$;
\item[(c)] every form of type $e_S$ or $\overline{e_S}$, $e_T$ or $\overline{e_T}$ is a cusp form,
rapidly decreasing on Siegel sets, and the forms of a given level form a finite-dimensional space;
\item[(d)] for every finite datum $f$ and every form $y$ of type $e_S$, the lift
$h\mapsto\int_{[\auU(S)]}\theta(\omega(g,h)(P_{S,T}\otimes f))\,y(g)\,dg$ converges absolutely and is a
form of type $e_T$; for every form $z$ of type $\overline{e_T}$ the transposed lift
$g\mapsto\int_{[\auU(T)]}\theta(\omega(g,h)(P_{S,T}\otimes f))\,z(h)\,dh$ converges absolutely and is a
form of type $\overline{e_S}$.
\end{enumerate}
Here $\omega$ is $\omega_{\alpha,\beta}$ with the characters of $a_{i,n}$, resp.\ of
Definition~\ref{ptr:desc-char}. In particular \textup{(A1)} and \textup{(A2)} of Lemma~\ref{pc:arch}
hold for the pairs $(\pcG_{i,0},\pcG_{i,1})$, and the maps $a_{i,n}$, $d_{i,n}$ of
Theorem~\ref{pc:tri} (the latter formed with any finite datum) are well defined between the spaces
$\pcF_{i,n}$.
\end{theorem}
\begin{proof}
(a) By Proposition~\ref{ptr:match} the parameters match at every $\tau$; {\cite[Propositions~3.11 and~3.13]{P49}} are stated for an arbitrary matching pair with distinct entries, and the
passage to the tensor product over $\tau$ is as in the proof of {\cite[Theorem~3.15]{P49}}.

(b) The proof of {\cite[Theorem~3.15]{P49}} uses: {\cite[Lemma~3.6]{P49}} (here Lemma~\ref{ptr:const}), the
matching (Proposition~\ref{ptr:match}), the conclusion of {\cite[Theorem~3.2]{P49}} for the two types
(Lemma~\ref{ptr:gap}), {\cite[Lemma~3.5]{P49}} (stated for any regular $\lambda$ with integrable
$e_\lambda$), {\cite[Lemma~3.14]{P49}} (stated for any matching pair), and (D1)--(D3). It applies
verbatim. (For a descending pair the row and column representations are $\pi_{\Lambda^{(i+1,n)}}$
and $\pi_{-\Lambda^{(i,n+1)}}$, and $\mathsf P$, $\mathsf P'$ are the projectors $\sigma(e_\lambda)$,
$\sigma'(e_{\lambda'})$ exactly as there.)

(c) {\cite[Proposition~8.4(a)]{P49}} for $G_j$ uses only (A1), (A2) for the type $e_j$, i.e.\ the
conclusions of {\cite[Theorem~3.2]{P49}} and {\cite[Corollary~3.3]{P49}}; these hold for every $e_{i,n}$ by
Lemma~\ref{ptr:gap}. The groups $\auU(\pcG_{i,n})$ with $n=m-i$ may be compact (no negative sign);
then the argument is the one given there for $j=m$.

(d) The proof of {\cite[Proposition~8.4(b)]{P49}}, (c) uses: the growth {\cite[Lemma~8.2]{P49}} (which is
stated for an arbitrary pair of spaces and characters: the norm $\|\cdot\|_s$ and the bound
$\|\omega_\infty(g,h)\|_s\le(2D\|g\|\|h\|)^{2s}$ depend only on $D=fNN'$), reduction theory,
statement (c) above, the conclusion of {\cite[Theorem~3.2]{P49}}, and the identity (b). It applies
verbatim, with $\Phi_j$ replaced by $P_{S,T}$, and gives the two assertions. The last sentence:
$a_{i,n}$ is the lift of (d) for the ascending pair with $f=\delta_\iota$ at $\pcm_\vartheta$ and
the vacuum elsewhere, so it maps $\pcF_{i,n}=\pcF(\pcG_{i,n},e_{i,n})$ into forms of type
$e_{i,n+1}$; and the descent along a descending pair is a transposed lift of (d) with
$e_S=\overline{e_{i+1,n}}$, $e_T=\overline{e_{i,n+1}}$, so it maps forms of type
$\overline{e_T}=e_{i,n+1}$ to forms of type $\overline{e_S}=e_{i+1,n}$. Rapid decay in (A1) of
Lemma~\ref{pc:arch} is (c), and in the form used below it reads: on a Siegel set
$\mathfrak S_\infty\mathsf cK'$ every such form is bounded by $C_N\|g_\infty\|^{-N}$ for every $N$.
This completes the proof. (The steps of the proofs of {\cite[Theorem~3.15]{P49}} and {\cite[Proposition~8.4]{P49}} referred to in (a)--(d)
are also written out one by one, for an arbitrary pair of standard spaces satisfying the block rule, in
Steps 0--4 of the proof of Theorem~\ref{pgv:G3}, which use neither the present theorem nor anything deduced from it; for the present pairs the excess is $t$ and the
properties used there are Proposition~\ref{ptr:match} and Lemma~\ref{ptr:gap}.)
\end{proof}

\begin{remark}[What the previous text assumed]\label{ptr:rem-prev}
In \S\ref{pc:ss-triangle} the parameter at the node $(i,n)$ is defined directly, and
Lemma~\ref{pc:arch} asserts the edge property for $(\pcG_{i,0},\pcG_{i,1})$ only. What the squares
need (step (iii) of \S\ref{pc:ss-exchange}; Theorem~\ref{ptr:square} below) is the edge property
for all four pairs of every square: for the square at $(i,n)$, with $U_1=\pcG_{i,n+1}$,
$U_2=\pcG_{i+1,n+1}$, $A=\pcG_{i,n+2}$, $A'=\pcG_{i+1,n}$, the pairs $(U_1,A)$ and $(A',U_2)$ are
ascending (at $(i,n+1)$ and $(i+1,n)$) and $(U_2,A)$, $(A',U_1)$ are descending (at $(i,n+1)$ and
$(i,n)$). All are covered by Theorem~\ref{ptr:edge}. The sign identity used in
\S\ref{pc:ss-exchange} for the numbers of negative signs is a consequence of \eqref{ptr:eq-E} and
was checked there.
\end{remark}

\subsubsection{The characters of a square}\label{ptr:ss-chars}

Fix a square at $(i,n)$, $i+n+2\le m$: $U_1=\pcG_{i,n+1}$, $U_2=\pcG_{i+1,n+1}$, $A=\pcG_{i,n+2}$,
$A'=\pcG_{i+1,n}$, $W=U_1\oplus U_2$ of rank $N=2M$, $M=r+(n+1)t$. For an $m$-fold twist $\vartheta$
let $\nu=(\nu_{n+1}\vartheta^{\,n+1})^{-1}$, of weight $(n+1)\Delta$ (\S\ref{pc:ss-exchange}), and
$\gamma=\nu^{-1}\upsilon^{M}$, so that $\ptrwt_\tau\gamma=M-(n+1)\Delta$.

In step (i) of \S\ref{pc:ss-exchange} the kernel of $(W,A)$ restricted to $\auU(U_1)\times\auU(U_2)$
is the product of the kernels of $(U_1,A)$ and $(U_2,A)$, and the kernel of $(A',W)$ is the product
of those of $(A',U_1)$ and $(A',U_2)$; for this the target characters on $\auU(A)$ of the first two
must multiply to $\gamma$, and the source characters on $\auU(A')$ of the last two must multiply to
$\gamma$ (Theorem~\ref{pcx:X}: $\auU(V)$ acts through $\gamma(\det h)$ on the left side and $\auU(V')$
through $\gamma(\det h')^{-1}$ on the right side), the Weil representations being multiplicative for
orthogonal sums. The pair $(U_1,A)$ is ascending at $(i,n+1)$, with target character
$\mu_{i+n+2}\upsilon_M$; $(U_2,A)$ is descending at $(i,n+1)$; $(A',U_2)$ is ascending at $(i+1,n)$; $(A',U_1)$ is descending at $(i,n)$. In the conventions of
\S\ref{pc:ss-data} the target character of $a_{i,n+1}$ and the source character of $a_{i+1,n}$ are
both $\mu_{i+n+2}\upsilon_{r+(n+1)t}=\mu_{i+n+2}\upsilon_M$.

\begin{lemma}[The descents' characters are those of the squares]\label{ptr:chars}
Put $\beta^\downarrow_{i,n+1}=\gamma\,(\mu_{i+n+2}\upsilon_M)^{-1}$ and
$\alpha^\downarrow_{i,n}=\gamma\,(\mu_{i+n+2}\upsilon_M)^{-1}$ (for the twist $\vartheta$ of the
square). Their weights are those of Definition~\ref{ptr:desc-char}. Consequently the characters
required by the exchange identity for the four kernels of a square make the four pairs edges (Theorem~\ref{ptr:edge}).
\end{lemma}
\begin{proof}
$\ptrwt(\mu_{i+n+2}\upsilon_M)=-\kappa_{i+n+2}+\lfloor M/2\rfloor$, so both characters have weight
$M-(n+1)\Delta+\kappa_{i+n+2}-\lfloor M/2\rfloor=\kappa_{i+1}+\lceil M/2\rceil$, because
$\kappa_{i+n+2}-(n+1)\Delta=\kappa_{i+1}$. Definition~\ref{ptr:desc-char} asks
$\ptrwt\beta^\downarrow_{i,n+1}=\kappa_{i+1}+\lceil\tfrac{r+(n+1)t}2\rceil$ and
$\ptrwt\alpha^\downarrow_{i,n}=\kappa_{i+1}+\lceil\tfrac{r+(n+1)t}2\rceil$: the same. The last
sentence is Proposition~\ref{ptr:match}(c).
\end{proof}
The characters of $d_{i,n}$ are fixed by the squares, consistently, and $\beta^\downarrow_{i,0}$
agrees with the character given by $d_{i,0}=a_{i,0}^*$ (Remark~\ref{ptr:rem-desc},
Proposition~\ref{ptv:chars}). The square at $(i,n)$ also fixes, by (E2) at $U_1$ and $U_2$,
$\beta^\downarrow_{i,n}=\gamma\,(\text{source character of }a_{i,n+1})^{-1}$ and
$\alpha^\downarrow_{i,n+1}=\gamma\,(\text{target character of }a_{i+1,n})^{-1}$ (weights checked in
Corollary~\ref{ptr:tri-b}); in particular $\alpha^\downarrow_{i,m-i-1}$ is fixed by the square at
$(i,m-i-2)$ when $m-i\ge2$, and the characters of $d_{m-1,0}$ are those of (W7)
(Proposition~\ref{ptv:chars}). They depend on $\vartheta$ through $\gamma$; the tame
part of $\gamma$ at the primes of $\pcm_\vartheta$ is the character $\chi$ of the square, as required
in step (i).

\subsubsection{Convergence in the squares}\label{ptr:ss-conv}

\begin{proposition}[Step (ii)]\label{ptr:conv}
Let a square be as in \S\ref{ptr:ss-chars}, $\pcm=\pcm_\vartheta$, $y_1$ a form of type
$e_{i,n+1}$ on $U_1$ and $y_2$ a form of type $\overline{e_{i+1,n+1}}$ on $U_2$, $\xi\in\pcMat_{a\times
N}(\pcR_\pcm)$ of rank $a=\dim A$ at every prime of $\pcm$ and $\eta\in\pcMat_{N\times b}(\pcR_\pcm)$ of rank
$b=\dim A'$, and let $\chi_2$ be a unitary character of $\auU(1)$. There are $M_0,C$ such that
for every Fock vector $x$ of $\auHom(W,A)\otimes\R$ and $v$ of $\auHom(A',W)\otimes\R$ with finite
moments of order $M_0$:
\begin{enumerate}
\item[(a)] the function $(g_1,g_2,h)\mapsto y_1(g_1)y_2(g_2)\chi_2(\det g_2)\,\theta^{W,A}(g_1\oplus g_2,h;
[x,\delta_\xi])$ is integrable on $[\auU(U_1)]\times[\auU(U_2)]\times[\auU(A)]$, with integral of absolute
value at most $C\|x\|_{M_0}$;
\item[(b)] the function $(g_1,g_2,h')\mapsto y_1(g_1)y_2(g_2)\chi_2(\det g_2)\,\theta^{A',W}(h',g_1\oplus
g_2;[v,\delta_\eta])$ is integrable on $[\auU(U_1)]\times[\auU(U_2)]\times[\auU(A')]$, with the same bound
in $\|v\|_{M_0}$.
\end{enumerate}
(The characters of the kernels are any; $\|\cdot\|_M$ is the moment of Proposition~\ref{pcx:decay}.)
In particular all iterated integrals in steps (i), (iii) of \S\ref{pc:ss-exchange} converge
absolutely and may be computed in any order.
\end{proposition}
\begin{proof}
\emph{Heights of block matrices.} With the heights of Definition~\ref{pcx:def-ht}, for $g_l\in
\auU(U_l)(\auA)$ the lattice at $g=g_1\oplus g_2$ is $\Lambda_{g_1}\oplus\Lambda_{g_2}$ (the frame of
$W$ is the union of the frames) and $|(\gamma_1,\gamma_2)|_g^2=|\gamma_1|_{g_1}^2+|\gamma_2|_{g_2}^2$; so
$\lambda_1(g)=\min(\lambda_1(g_1),\lambda_1(g_2))$ and $\ptrht(g)=\max(\ptrht g_1,\ptrht g_2)\le
\ptrht(g_1)\ptrht(g_2)$.

\emph{Heights on Siegel sets.} Let $\mathfrak S_\infty\mathsf cK'$ be as in the proof of
{\cite[Proposition~8.4(b)]{P49}} for $U_l$. For $g=g_\infty c k'$ in it and $0\ne\gamma\in\Lambda_g$,
$\gamma$ lies in the lattice $(ck')L$ of $S(E)$, which belongs to a finite set of lattices as $c$
runs over the finite set $\mathsf c$ and $k'$ over $K'$; so $\|\gamma\|\ge c_0$ for a fixed euclidean
norm and some $0<c_0\le1$. Since $\|\gamma\|=\|g_\infty g_\infty^{-1}\gamma\|\le\|g_\infty\|\,
\|g_\infty^{-1}\gamma\|$ (operator norm; the norm of $g_\infty=(g_\tau)_\tau$ is the maximum over the
archimedean factors), $|\gamma|_g=\|g_\infty^{-1}\gamma\|\ge c_0\|g_\infty\|^{-1}$. Hence
$\ptrht(g)=\max(1,\lambda_1(g)^{-1})\le c_0^{-1}\max(1,\|g_\infty\|)$ on the Siegel set. By
Theorem~\ref{ptr:edge}(c), $|y_l(g)|\le C_{N'}\|g_\infty\|^{-N'}$ there for every $N'\ge0$; using
$N'=0$ where $\|g_\infty\|\le1$ and $N'=B$ where $\|g_\infty\|\ge1$,
$|y_l(g)|\ptrht(g)^B\le c_0^{-B}\max(C_0,C_B)$ on the Siegel set, so it is bounded there, hence on $\auU(U_l)(\auA)$ (left
$\auU(U_l)(E^+)$-invariance), for every $B$. As $[\auU(U_l)]$ has finite volume,
$|y_l|\ptrht^B$ is integrable.

(a) By Proposition~\ref{pcx:decay}(a) for $(S_1,S_2)=(W,A)$ and $P$ the primes of $\pcm$ (the point
datum $\delta_\xi$ is supported on a matrix of rank $a$, i.e.\ surjective, at every place of $P$),
for $N_0$ large there are $B,M_0,C$ with $|\theta^{W,A}(g,h;[x,\delta_\xi])|\le C\|x\|_{M_0}
\ptrht(g)^B\ptrht(h)^{-N_0}$. With the two preceding paragraphs the integrand is bounded by
$C\|x\|_{M_0}\,|y_1|\ptrht^B(g_1)\cdot|y_2|\ptrht^B(g_2)\cdot\ptrht(h)^{-N_0}$, a product of integrable
functions by (R1) of \S\ref{pcx:ss-conv}. (b) The same with Proposition~\ref{pcx:decay}(b) for
$(S_1,S_2)=(A',W)$: $\delta_\eta$ is supported on a matrix of rank $b$, i.e.\ injective. The last
sentence follows by Fubini's theorem; the path integrals of steps (i), (iii) (lift to $A$, then
pairing on $[\auU(A)]$; transposed lift to $A'$, then lift to $U_2$) are iterated integrals of these
functions, by the factorisation of the kernels in step (i).
\end{proof}

\subsubsection{Archimedean multiplicity one}\label{ptr:ss-mult}

Let $(S,T)$ be a pair with matching parameters as in Theorem~\ref{ptr:edge} (or any pair for which
the conclusions (a)--(d) of that theorem hold), $P=P_{S,T}$ its joint harmonic, $\mathsf P$,
$\mathsf P'$ as there, and $\pzF_s$ the Banach space of {\cite[\S8.1]{P49}}. We use the following
local statement.

\begin{enumerate}
\item[(H)] For every polynomial vector $Q$ of the Fock space of $\auHom(S,T)\otimes\R$ there is a
continuous compactly supported function $\alpha_Q$ on $\auU(T)_\infty$ with
$\mathsf PQ=\int\alpha_Q(x)\,\omega_\infty(1,x)\,\mathsf PP\,dx$ (Bochner integral in $\pzF_s$).
\end{enumerate}
\noindent (H) holds for every ascending and descending pair of the triangle: this is Theorem~\ref{ptq:H}
(proved by Howe's argument, Proposition~\ref{ptq:howe}, modulo (V) and the published results
(D1)--(D5)).

\begin{lemma}\label{ptr:conv-lift}
Let $(S,T)$ be a pair for which \textup{(H)} holds (for instance any pair of the triangle, by
Theorem~\ref{ptq:H}). For every polynomial $Q$, every finite datum $f$ and every form $y$ of type
$e_S$, the lift $F_Q$ of $y$ with $Q\otimes f$ (absolutely convergent as in
Theorem~\ref{ptr:edge}(d)) is $F_Q(h)=\int\alpha_Q(x)F_P(hx)\,dx$, with $\alpha_Q$ independent of
$f$, $y$, the characters and the spaces of the same signatures.
\end{lemma}
\begin{proof}
In the proof of {\cite[Proposition~8.4(b)]{P49}}, $F_\Psi(h)=L_h(\Psi)$ with $L_h$ continuous on
$\pzF_s$, $L_h(\Psi)=L_h(\mathsf P\Psi)$ and $L_h(\omega_\infty(1,x)\Psi)=L_{hx}(\Psi)$. Apply $L_h$ to
(H): $F_Q(h)=L_h(\mathsf PQ)=\int\alpha_Q(x)L_{hx}(\mathsf PP)\,dx=\int\alpha_Q(x)F_P(hx)\,dx$.
\end{proof}

\begin{lemma}\label{ptr:avg}
Let $e$ be an integrable type on a space $T$ with parameter of the class of Lemma~\ref{ptr:gap},
$\alpha_1,\alpha_2\in C_c(\auU(T)_\infty)$. There is $r=r(\alpha_1,\alpha_2)\in\C$ such that for
every form $F_1$ of type $e$ and every form $F_2$ of type $\overline e$ (both rapidly decreasing),
$\int_{[\auU(T)]}(\alpha_1*F_1)(\alpha_2*F_2)=r\int_{[\auU(T)]}F_1F_2$, where $(\alpha*F)(h)=\int\alpha(x)
F(hx)dx$.
\end{lemma}
\begin{proof}
$\overline{F_2}$ is of type $e$. By \S\ref{pc:ss-data} (consequence of Schur orthogonality),
$F_1=\Phi_1(v_0)$, $\overline{F_2}=\Phi_2(v_0)$ for intertwiners $\Phi_l\colon\pi\to L^2$, and
$\langle\Phi_2(u),\Phi_1(v)\rangle=\langle\overline{F_2},F_1\rangle\langle u,v\rangle$. Hence
$B(x,x'):=\int F_1(hx)F_2(hx')dh=\langle\rho(x')\overline{F_2},\rho(x)F_1\rangle=\langle\pi(x')v_0,
\pi(x)v_0\rangle\int F_1F_2$, and by Fubini (the integrand is bounded by
$|\alpha_1(x)\alpha_2(x')|\,|F_1(hx)F_2(hx')|$, integrable since the forms are bounded and in $L^2$)
the claim holds with $r=\iint\alpha_1(x)\alpha_2(x')\langle\pi(x')v_0,\pi(x)v_0\rangle dx\,dx'$.
\end{proof}

\begin{lemma}[Density of polynomials for moment norms]\label{ptr:dense}
In the realisation of Lemma~\ref{pcx:superpos}, every Fock vector $x$ with $\|x\|_{M'}<\infty$ for
all $M'$ is, for each $M$, a limit in $\|\cdot\|_M$ of polynomial vectors. Here the polynomial vectors are the images of
those of {\cite[(41)]{P49}} under the isomorphism of Lemma~\ref{pcx:superpos}; it carries the vacuum
to the vacuum and intertwines the Heisenberg action, so they are the vectors $Pe$ of the realisation.
\end{lemma}
\begin{proof}
By Lemma~\ref{pcx:superpos}, $x=\int x(z)\rho(z)e\,dz$; since $\|\rho(z)e\|_M\le C(1+|z|)^M$ and
$z\mapsto\rho(z)e$ is continuous for $\|\cdot\|_M$, the truncations to $|z|\le R$ converge to $x$
(since $\|\rho(z)e\|_M\le C(1+|z|)^M$, the tail is at most $C\int_{|z|>R}(1+|z|)^M|x(z)|\,dz$, which
tends to $0$ because $\|x\|_M<\infty$, Lemma~\ref{pcx:superpos}(b)), and each truncation is a limit of finite
combinations of vectors $\rho(z)e$. Each $\rho(z)e$ is $\exp(\ell_z)\,e$ with $\ell_z$ affine in the
coordinates in which the polynomial vectors are the $Pe$ (the vacuum is a Gaussian $e$); the Taylor
truncations $\sum_{k\le K}\ell_z^k/k!\cdot e$ are polynomial vectors, bounded by
$e^{|\ell_z|}|e|$, and $(1+|w|)^Me^{|\ell_z(w)|}|e(w)|$ is integrable; so they converge in
$\|\cdot\|_M$ by dominated convergence.
\end{proof}

\subsubsection{The square law}\label{ptr:ss-square}

We state the passage from the exchange identity to one square in a form that serves both the
triangle (Theorem~\ref{pc:tri}(b)) and the grids (Proposition~\ref{pgr:square}).

\begin{definition}[Edge square]\label{ptr:def-square}
An \emph{edge square} of excess $d\ge0$ and weight $\mu$ consists of spaces $U_1,U_2,A,A'$ of
ranks $n_1,n_2,a,b$ with $a=b+2d$, $N=n_1+n_2=a+b$, $W=U_1\oplus U_2$, such that, with $q$, $q_A$,
$q_{A'}$ the numbers of negative signs of $W,A,A'$: $q=q_A+q_{A'}$, $2q+2\le N\le2\mu_\tau+2q$; of
types $e_1,e_2,e,e'$ on $U_1,U_2,A,A'$; and of a family $I$ of data $(\pcm,\nu,\chi,\text{characters})$,
$\pcm\ne(1)$ a modulus, $\nu$ a tame character for $\pcm$ of weight $\mu$ with tame character $\chi$,
$\gamma=\nu^{-1}\upsilon^{N/2}$, and characters for the four kernels such that:
\begin{enumerate}
\item[(E1)] (edges) the four pairs are edges in the sense of Theorem~\ref{ptr:edge}(a)--(d) with
polynomials $P_1,P_2,P'_1,P'_2$: $(U_1,A)$ from $e_1$ to $e$, $(U_2,A)$ from $\overline{e_2}$ to
$\overline e$, $(A',U_1)$ from $\overline{e'}$ to $\overline{e_1}$, $(A',U_2)$ from $e'$ to $e_2$;
the polynomials and types do not depend on the member of $I$, only the characters do, with fixed
weights; and \textup{(H)} holds for $(U_1,A)$ and $(U_2,A)$;
\item[(E2)] (characters) the target characters on $A$ of the kernels of $(U_1,A)$, $(U_2,A)$
multiply to $\gamma$; the source characters on $A'$ of $(A',U_1)$, $(A',U_2)$ multiply to $\gamma$;
and for $l=1,2$ the source character on $U_l$ of $(U_l,A)$ times the target character on $U_l$ of
$(A',U_l)$ is $\gamma$.
\end{enumerate}
For $\xi\in\pcMat_{a\times N}(\pcR_\pcm)$ with blocks $\xi_l$ (columns of $U_l$) and
$\eta\in\pcMat_{N\times b}(\pcR_\pcm)$ with blocks $\eta_l$ (rows of $U_l$), forms $y_1$ of type $e_1$
and $y_2$ of type $\overline{e_2}$, put
\[
L(\xi)=\int_{[\auU(A)]}\Lambda^{(1)}_{\xi_1}y_1\cdot\Lambda^{(2)}_{\xi_2}y_2,\qquad
L'(\eta)=\int_{[\auU(U_2)]}y_2\cdot\Lambda'^{(2)}_{\eta_2}T'^{(1)}_{\eta_1}y_1,
\]
where $\Lambda^{(l)}_\zeta$ is the lift along $(U_l,A)$ with $P_l$ and the point datum $\delta_\zeta$,
$T'^{(1)}_\zeta$ the transposed lift along $(A',U_1)$ and $\Lambda'^{(2)}_\zeta$ the lift along
$(A',U_2)$. (By (E1) and Theorem~\ref{ptr:edge}(d) these are forms of the right types, and the
integrals converge.)
\end{definition}

\begin{theorem}[Square law]\label{ptr:square}
For an edge square there are $\kappa\in\C$ and $C>0$, depending only on the archimedean data, and
for each member of $I$ a number $c$ with $|c|^2=C\,L_*(\pcm,\nu,d)\ne0$, such that for all forms
$y_1,y_2$ the functions $L$ (restricted to the $\xi$ of rank $a$) and $L'$ (restricted to the $\eta$
of rank $b$) satisfy
\begin{equation}\label{ptr:eq-square}
\sum_\eta\Psi(\eta)L'(\eta)=\frac{\kappa}{c}\sum_\xi\Phi(\xi)L(\xi)
\end{equation}
for all $\Phi,\Psi$ matched for $\chi$; moreover $L(h\xi)=\chi(\det h)\,L(\xi)$ and
$L'(\eta h')=\chi(\det h')\,L'(\eta)$ for $h\in\auGL_a(\pcR_\pcm)$, $h'\in\auGL_b(\pcR_\pcm)$. Thus
$L,L'$ satisfy the square law of \S\ref{pc:ss-exchange} with constant $\kappa/c$ (matched functions
vanish off full-rank matrices).
\end{theorem}
\begin{proof}
\emph{Transformation laws.} For $h\in\auGL_a(\pcR_\pcm)$ choose $k\in K_A$ with reduction $h$. By
Lemma~\ref{pc:level}, $\omega(1,k^{-1})$ sends the datum $[P_l,\delta_{h\xi_l}]$ to the target character
of $(U_l,A)$ at $\det k$, inverted, times $[P_l,\delta_{\xi_l}]$; translating the lifted forms by
$k^{-1}$ (the pairing on $[\auU(A)]$ is invariant) gives $L(h\xi)=c_1(\det k)^{-1}c_2(\det k)^{-1}L(\xi)$
with $c_1c_2=\gamma$ the product of the two target characters (E2), and $\gamma^{-1}$ restricted to
$K_A$ at the primes of $\pcm$ is $\chi\circ\det$ by the definition of $\chi$ (the factor
$\upsilon^{N/2}$ is unramified). The law of $L'$ is the same with $K_{A'}$ and the source characters.

\emph{Kernels multiply.} $\psrho$, $\theta$ and the Weil representations of
$\auHom(W,A)=\auHom(U_1,A)\oplus\auHom(U_2,A)$ are tensor products, and
$[x_1x_2,\delta_\xi]=[x_1,\delta_{\xi_1}]\otimes[x_2,\delta_{\xi_2}]$ (finite data multiply blockwise;
the vacuum of an orthogonal sum of lattices is the tensor product). So for Fock vectors $x$ of
$\auHom(W,A)\otimes\R$ the function
\[
L_x(\xi)=\int_{[\auU(U_1)]\times[\auU(U_2)]\times[\auU(A)]}y_1(g_1)\,y_2(g_2)\,\lambda(\det g_2)\,
\theta^{W,A}_{\alpha,\gamma}(g_1\oplus g_2,h;[x,\delta_\xi])
\]
($\alpha=\alpha_1$, where $\alpha_l$ is the source character of $(U_l,A)$, and
$\lambda=\alpha_1\alpha_2^{-1}$, unitary; by (E2) also $\lambda=\alpha'_2\alpha'^{-1}_1$ for the target
characters $\alpha'_l=\gamma\alpha_l^{-1}$ of $(A',U_l)$, so that $y_1y_2\lambda(\det g_2)
\theta^{W,A}_{\alpha_1,\gamma}$ and $y_1y_2\lambda(\det g_2)\theta^{A',W}_{\gamma,\alpha'_1}$ are the
products of the two kernels on each side) equals $L(\xi)$ for $x=P_1P_2$, by Fubini; the triple integral converges
absolutely by Proposition~\ref{ptr:conv}(a) (whose proof uses only (E1) and the full rank of $\xi$),
and $|L_x(\xi)|\le C_\xi\|x\|_{M_0}$. In the same way $L'(\eta)$ is the triple integral
$L'_{v}(\eta)$ over $[\auU(U_1)]\times[\auU(U_2)]\times[\auU(A')]$ of $y_1y_2\lambda(\det g_2)
\theta^{A',W}_{\gamma,\alpha'}(h',g_1\oplus g_2;[v,\delta_\eta])$ with $v=P'_1P'_2$ and
$\alpha'=\gamma\alpha^{-1}$ by (E2), absolutely convergent by Proposition~\ref{ptr:conv}(b) (here the
lift along $(A',U_2)$ is moved to the other side of the pairing on $U_2$ by Fubini).

\emph{Multiplicity one.} Fix $y_1,y_2$. For a polynomial $Q=Q_1Q_2$ with $Q_l$ on
$\auHom(U_l,A)\otimes\R$, Lemma~\ref{ptr:conv-lift} and Lemma~\ref{ptr:avg} (with $F_1=\Lambda^{(1)}
_{\xi_1}y_1$ of type $e$ and $F_2=\Lambda^{(2)}_{\xi_2}y_2$ of type $\overline e$) give
$L_{Q}(\xi)=r(\alpha_{Q_1},\alpha_{Q_2})L(\xi)$, with $r$ independent of $\xi$, $y_1,y_2$ and the
member of $I$. Every polynomial on $\auHom(W,A)\otimes\R$ is a sum of such products; so for every
polynomial $Q$ there is $\kappa(Q)$ with $L_Q=\kappa(Q)L$ on all cases (all $\xi$ of rank $a$, all
$y_1,y_2$, all members of $I$). The map $x\mapsto(L_x(\text{case}))_{\text{cases}}$ is, case by case,
continuous for $\|\cdot\|_{M_0}$, and the set of functions on the cases that are multiples of $L$ is a
line, closed for pointwise convergence. By Lemma~\ref{ptr:dense}, for every $x$ with all moments
finite there is $\kappa(x)$ with $L_x=\kappa(x)L$.

\emph{Exchange.} Apply Theorem~\ref{pcx:X} to $(W,V,V')=(W,A,A')$ with the vector $v=P'_1P'_2$ (a
polynomial, so all its moments are finite): it gives $C>0$ and $x$ with finite moments, and for each
member of $I$ a number $c$ with $|c|^2=CL_*(\pcm,\nu,d)\ne0$ (Corollary~\ref{pcx:cne0}; the excess called $e$ there is $d$ here), such that
for matched $\Phi,\Psi$ and all $g$,
$\int_{[\auU(A)]}\theta^{W,A}_{\alpha,\gamma}(g,h;[x,\Phi])dh=c\int_{[\auU(A')]}\theta^{A',W}_{\gamma,
\alpha'}(h',g;[v,\Psi])dh'$. Write $[x,\Phi]=\sum_\xi\Phi(\xi)[x,\delta_\xi]$ (a finite sum, $\Phi$
vanishing off rank $a$), likewise for $\Psi$; multiply by $y_1(g_1)y_2(g_2)\lambda(\det g_2)$, put
$g=g_1\oplus g_2$, and integrate over $[\auU(U_1)]\times[\auU(U_2)]$; by the absolute convergence just
proved the integrals may be exchanged with the finite sums, and
$\sum_\xi\Phi(\xi)L_x(\xi)=c\sum_\eta\Psi(\eta)L'_v(\eta)$. With $L_x=\kappa L$, $\kappa=\kappa(x)$,
and $L'_v=L'$ this is \eqref{ptr:eq-square}.
\end{proof}

\begin{corollary}[Theorem~\ref{pc:tri}(b)]\label{ptr:tri-b}
Assume \textup{(V)} ({\cite[Theorem~3.29]{P49}}) and \textup{(D1)--(D5)}. Then the squares of the triangle at $(i,n)$, $i+n+2\le m$, with the family of $m$-fold
twists, are edge squares of excess $t$ and weight $(n+1)\Delta$, and Theorem~\ref{pc:tri}(b) holds,
with $\kappa^\sharp_{i,n}=\kappa$, $c_{i,n}(\vartheta)=c$.
\end{corollary}
\begin{proof}
The rank and sign conditions are verified in \S\ref{pc:ss-exchange}; (E1) is
Theorem~\ref{ptr:edge} with Remark~\ref{ptr:rem-prev}; (E2) holds for the weights by
Lemma~\ref{ptr:chars} and the analogous computation at $U_1$, $U_2$: the weights of the source
character of $a_{i,n+1}$ and of $\beta^\downarrow_{i,n}$ add up to
$-\kappa_{i+n+2}+\lfloor\tfrac{M+t}2\rfloor+\kappa_{i+1}+\lceil\tfrac{M-t}2\rceil=M-(n+1)\Delta=\ptrwt
\gamma$, and those of $\alpha^\downarrow_{i,n+1}$ and of the target character of $a_{i+1,n}$ add up to
$\kappa_{i+1}+\lceil\tfrac{M+t}2\rceil-\kappa_{i+n+2}+\lfloor\tfrac{M-t}2\rfloor=\ptrwt\gamma$
($M\pm t$ have the same parity). The descent characters are then fixed by (E2), consistently in all
squares (Proposition~\ref{ptv:chars}); (H) for the ascending and descending pairs is
Theorem~\ref{ptq:H}. Theorem~\ref{ptr:square} gives the square law for $L,L'$ with constant
$\kappa/c$, for the splitting of Definition~\ref{pcx:def-matched} (the first $b$ indices of the frame of
$W=U_1\oplus U_2$). The splitting $(\mathsf B,\mathsf A)=(\mathsf N_0,\mathsf N_1\sqcup\mathsf T')$ of
Lemma~\ref{pc:next} is obtained from it by the permutation that moves the $r+nt$ indices of $\mathsf N_0$
in front of the $r+(n+1)t$ indices of $\mathsf N_1$, of sign $(-1)^{(r+nt)(r+(n+1)t)}=1$ because $r+nt$ and
$r+(n+1)t$ have opposite parities ($t$ is odd, Lemma~\ref{ptr:cons}(c)); so by Lemma~\ref{ptw:split} the law holds
for this splitting with the same constant, and step (iv) of \S\ref{pc:ss-exchange}
(Lemma~\ref{pc:next}) turns it into (b).
\end{proof}

\begin{remark}[The grid squares]\label{ptr:rem-grid}
Proposition~\ref{pgr:square} is Theorem~\ref{ptr:square} for the square $(V',V;W_1,W_2)$ of a grid,
$U_l=W_l$, $A=V$, $A'=V'$, with excess $d=t$ and $\nu=(\mu^0_{a+1}\mu^0_{b+1}\vartheta^n)^{-1}$; its
functions $L_{f,l}$, $L'_{f,l}$ are obtained from $L,L'$ above by replacing the pairing with $y_2$ by
the linear functional $l$, which changes nothing in the proof (Fubini is applied with
$l$ applied after the integrals, $l$ being a linear functional on a finite-dimensional space of forms).
The rank, sign and weight conditions are verified in Lemma~\ref{pgv:esq}; (E1) requires the edge
property along the four edges of a grid square; it is checked in Proposition~\ref{pgv:match} and
Theorem~\ref{pgv:G3}, clause by clause in Lemma~\ref{pgv:E1g} (for edges read with the groups
exchanged use Theorem~\ref{ptw:Hc} and Proposition~\ref{ptw:revedge}); (E2) is the relation
$\chi_{f_1}=\chi_{e_2}$, $\chi_{f_2}=\chi_{e_1}$ of the edge characters there together with the
product condition at $V$, $V'$. The transformation laws of Theorem~\ref{ptr:square} are the
invariance used in Proposition~\ref{pgr:turned}, and the identity \eqref{ptr:eq-square} is the input
of Lemma~\ref{pgr:sumturn}.
\end{remark}

\subsubsection{The zeta integral does not vanish}\label{ptr:ss-Z}

Let $(S,T)$ be one of the pairs above, $\tau$ fixed, $G=\auU(S)_\tau$ with row parameter $\lambda$,
$\sigma=\omega_\infty(\cdot,1)$ at $\tau$, $\Phi=\Phi_\tau$ the local joint harmonic, $c_\lambda$ the
coefficient of a unit highest weight vector of the lowest $K$-type $\Lambda$ of $\pi_\lambda$,
$e_\lambda=d_\lambda\overline{c_\lambda}$ and $\mathsf P=\sigma(e_\lambda)$.

\begin{enumerate}
\item[(F)] For every positive definite $w\in G$, $\langle\Phi,\sigma(w)\Phi\rangle>0$.
\end{enumerate}
\noindent (F) holds: this is Theorem~\ref{ptq:F} (proved in full; the vacuum coefficient is Lemma~\ref{ptq:vac}).

\begin{proposition}\label{ptr:Z}
Then $Z_\infty\neq0$ for the pair $(S,T)$; precisely
$Z_\infty=\prod_\tau d_{\lambda_\tau}^{-1}\langle\Phi_\tau,\mathsf P_\tau\Phi_\tau\rangle$ and each factor is
positive.
\end{proposition}
\begin{proof}
$Z_\infty$ is the integral of $\langle P,\sigma(g)P\rangle\langle v_0,\pi_S(g)v_0\rangle$; as
$\overline{\pi_S}\cong\pi_\lambda$ with $\overline{v_0}$ the highest vector (proof of
{\cite[Theorem~3.15]{P49}}), $\langle v_0,\pi_S(g)v_0\rangle=\overline{c_\lambda(g)}$, and the integral is a
product over $\tau$ of $\int\overline{c_\lambda}\langle\Phi,\sigma(g)\Phi\rangle=d_\lambda^{-1}
\langle\Phi,\mathsf P\Phi\rangle$ (Fubini; $c_\lambda\in L^1$ by Lemma~\ref{ptr:gap}). Fix $\tau$.

\emph{Two averages over $K$.} $\Phi$ is a highest weight vector of weight $\Lambda$ for $K$
({\cite[Definition~3.9]{P49}}, {\cite[Lemma~3.10(iii)]{P49}}), so it generates an irreducible $K$-module
$V$ of highest weight $\Lambda$. The function $k\mapsto\langle\Phi,\sigma(k)\Phi\rangle/\|\Phi\|^2$
extends holomorphically to $K_\C$, equals $1$ at $1$, and is multiplied by $\Lambda(l)\Lambda(u)$
under $x\mapsto lxu$ ($l$ lower, $u$ upper triangular in each block, for the order of decreasing
$\lambda$); $\pdsm_\Lambda$ of {\cite[(57)]{P51}} has the same properties, the Gauss decomposition is
dense, so $\langle\Phi,\sigma(k)\Phi\rangle=\|\Phi\|^2\pdsm_\Lambda(k)$, and likewise
$\langle v,kv\rangle=\pdsm_\Lambda(k)$ for the unit highest vector $v$ of the lowest $K$-type. By Schur
orthogonality on $K$ (mass $1$), with $d_\Lambda=\dim V=(\int|\pdsm_\Lambda|^2)^{-1}$, for $u\in K$:
\begin{equation}\label{ptr:eq-Kavg}
\int_K\pdsm_\Lambda(k^{-1}u)\,\sigma(k)\Phi\,dk=d_\Lambda^{-1}\sigma(u)\Phi .
\end{equation}
(Identify $V$ with the lowest $K$-type by $\Phi\mapsto\|\Phi\|v$; then $\pdsm_\Lambda(k^{-1}u)=
\langle kv,uv\rangle$ and $\int\langle kv,uv\rangle kv\,dk=d_\Lambda^{-1}uv$.)

\emph{Positivity of the averaged integrand.} Since $e_\lambda\in L^1$ and
$|\langle\Phi,\sigma(g)\Phi\rangle|\le\|\Phi\|^2$, $A(g)=\overline{c_\lambda(g)}\langle\Phi,\sigma(g)\Phi
\rangle$ is integrable and, by bi-invariance of Haar measure and Fubini,
$\int_GA=\int_G\bar A(g)\,dg$, $\bar A(g)=\iint_{K\times K}A(k_1gk_2)\,dk_1dk_2$. Fix $g=\kappa z$
($\kappa\in K$, $z$ positive definite, {\cite[\S5.2]{P51}}). Then $k_1\kappa zk_2=(k_1\kappa k_2)
(k_2^{-1}zk_2)$ with $k_2^{-1}zk_2$ positive definite, and {\cite[(59)]{P51}} ({\cite[Theorem~3.17]{P49}},
applicable by Lemma~\ref{ptr:gap}) with the substitution $u\mapsto k_2^{-1}u$ gives
\[
\overline{c_\lambda(k_1\kappa zk_2)}=d_\Lambda\int_K\frac{\overline{\pdsm_\Lambda(k_1\kappa u)}\,
\pdsm_\Lambda(k_2^{-1}u)}{\pdsm_{-\mu}(u^{-1}zu)}\,du ,
\]
{\emergencystretch=6em the denominator being positive ({\cite[Lemma~5.5]{P51}}). Insert this in $\bar A(g)$, with
$\langle\Phi,\sigma(k_1\kappa zk_2)\Phi\rangle=\langle\sigma(k_1^{-1})\Phi,\allowbreak \sigma(\kappa z)\sigma(k_2)
\Phi\rangle$; all integrands are continuous on compact sets, so Fubini applies. By
\eqref{ptr:eq-Kavg}, $\int\pdsm_\Lambda(k_2^{-1}u)\sigma(k_2)\Phi\,dk_2=d_\Lambda^{-1}\sigma(u)\Phi$ and,
with $k=k_1^{-1}$, $\int\overline{\pdsm_\Lambda(k_1\kappa u)}\langle\sigma(k_1^{-1})\Phi,\cdot\rangle
dk_1=\langle d_\Lambda^{-1}\sigma(\kappa u)\Phi,\cdot\rangle$. Hence
\[
\bar A(\kappa z)=d_\Lambda^{-1}\int_K\frac{\langle\sigma(\kappa u)\Phi,\sigma(\kappa z)\sigma(u)\Phi
\rangle}{\pdsm_{-\mu}(u^{-1}zu)}du=d_\Lambda^{-1}\int_K\frac{\langle\Phi,\sigma(u^{-1}zu)\Phi\rangle}
{\pdsm_{-\mu}(u^{-1}zu)}\,du>0
\]
by (F) (Theorem~\ref{ptq:F}). So $\langle\Phi,\mathsf P\Phi\rangle=d_\lambda\int_G\bar A>0$.\par}
\end{proof}




% ===== C-triangle2.tex =====
% proofs/C-triangle2.tex
% Continuation of proofs/C-triangle.tex (labels ptr:): the positivity (F) of Fock coefficients and
% the smoothing lemma (H) (Howe's argument). Labels of this file start with ptq:.
% Only amsmath/amssymb/amsthm are used. Conventions and macros: proofs/C-triangle.tex.
\providecommand{\auU}{\mathrm U}
\providecommand{\auHom}{\operatorname{Hom}}
\providecommand{\pzF}{\mathfrak F}
\providecommand{\pcG}{\mathsf G}
\providecommand{\ptqsgn}{\operatorname{sgn}}

\subsection{Positivity of Fock coefficients and the smoothing lemma}\label{ptq:sec}

This section continues \S\ref{ptr:sec} and proves the two local statements used there: the
positivity (F) of \S\ref{ptr:ss-Z} (Theorem~\ref{ptq:F}) and the smoothing statement (H) of
\S\ref{ptr:ss-mult} (Theorem~\ref{ptq:H}). Notation is that of \S\ref{ptr:sec} and of
{\cite[\S\S3.1.4--3.1.6]{P49}}; inner products are conjugate-linear in the first variable.
Throughout, $(S,T)$ is one of the ascending or descending pairs of the triangle
(\S\ref{ptr:ss-param}), $\tau$ is fixed, $R$ indexes the rows (an orthogonal basis of
$S\otimes_{E,\tau}\C$, signs $\varepsilon_r$), $C$ the columns (signs $\varepsilon'_c$),
$\lambda,\lambda'$ are the row and column parameters, which match in the sense of
\eqref{ptr:eq-rule} (Proposition~\ref{ptr:match}), $G=\auU(S)_\tau$, $G'=\auU(T)_\tau$,
$\sigma=\omega_\infty(\cdot,1)$, $\sigma'=\omega_\infty(1,\cdot)$ at $\tau$, with derived actions on
polynomial vectors given by {\cite[(12)]{P49}} (Lemma~\ref{ptr:const}, using {\cite[Theorem~3.29]{P49}}),
and $\mathsf P=\sigma(e_\lambda)$, $\mathsf P'=\sigma'(e_{\lambda'})$, where for $f\in L^1$ we write
$\sigma(f)=\int f(x)\sigma(x)\,dx$. By Lemma~\ref{ptr:gap} and {\cite[Theorem~3.2]{P49}}, $e_\lambda$,
$e_{\lambda'}$ and all their left and right derivatives are integrable.

We prove (F) (the vacuum coefficient is treated in Lemma~\ref{ptq:vac}) and (H). Besides
(D1)--(D3) (already used in \S\ref{ptr:ss-edge}), the proof of (H) uses two published results: the theorem of Harish-Chandra
that a $K$-finite, $\mathfrak Z$-finite function is reproduced by convolution with a smooth
compactly supported function, and Schmid's theorem on the $K$-types of discrete series.

\subsubsection{Matrix coefficients between harmonic polynomials}\label{ptq:ss-coef}

In this subsection $I$ is a finite set, $\eta\colon I\to\{\pm1\}$, $\pzF$ the Fock space in
variables $z_i$ ($i\in I$) with the Fischer inner product ($z_i$ adjoint to $\partial_i$, $\langle
1,1\rangle=1$), $a^\pm_i$ as in {\cite[\S3.1.4]{P49}} ($a^+_i=z_i$, $a^-_i=\partial_i$ if $\eta_i=1$;
$a^+_i=-\partial_i$, $a^-_i=z_i$ if $\eta_i=-1$), so that $(a^+_i)^*=\eta_ia^-_i$ on polynomials.
The space $W=\C^I$ carries the hermitian form $\sum_i\eta_i\bar w_iw'_i$ and its unitary group
$\auU(\eta)$; $\varrho$ is the Heisenberg representation on $\pzF$, normalised so that on
polynomial vectors
\begin{equation}\label{ptq:eq-heis}
d\varrho(w)=c_0\sum_i\bigl(w_i\,a^+_i-\eta_i\,\bar w_i\,a^-_i\bigr)\qquad(w\in W)
\end{equation}
for a nonzero real constant $c_0$ (this is the Fock model of {\cite[\S3.1.4]{P49}}: $z_i$ is the
coordinate or its conjugate according to $\eta_i$; see Remark~\ref{ptq:rem-conv} for the other
normalisation). Let $u\in\auU(\eta)$ and let $\omega(u)$ be a unitary operator on $\pzF$ with
\begin{equation}\label{ptq:eq-cov}
\omega(u)\varrho(w)\omega(u)^{-1}=\varrho(uw)\qquad(w\in W),
\end{equation}
which maps smooth vectors of $\varrho$ to smooth vectors. For polynomials $Q,P$ write
$\langle Q,\omega(u)P\rangle$ for the coefficient between the polynomial vectors.

\begin{lemma}[Moving the ladder operators]\label{ptq:ladder}
For polynomials $P,Q$ and $j\in I$:
\[
\langle Q,\omega(u)a^+_jP\rangle=\sum_iu_{ij}\eta_i\langle a^-_iQ,\omega(u)P\rangle,\qquad
\eta_j\langle Q,\omega(u)a^-_jP\rangle=\sum_i\overline{u_{ij}}\,\langle a^+_iQ,\omega(u)P\rangle .
\]
\end{lemma}
\begin{proof}
By \eqref{ptq:eq-cov}, $\omega(u)d\varrho(w)P=d\varrho(uw)\omega(u)P$, and $d\varrho(uw)$ is
skew-adjoint (it is $c_0\sum_i((uw)_ia^+_i-((uw)_ia^+_i)^*)$). Hence
$\langle Q,\omega(u)d\varrho(w)P\rangle=-\langle d\varrho(uw)Q,\omega(u)P\rangle$. Expanding both sides
with \eqref{ptq:eq-heis} and using conjugate-linearity in the first variable,
\[
\sum_jw_j\langle Q,\omega(u)a^+_jP\rangle-\sum_j\eta_j\bar w_j\langle Q,\omega(u)a^-_jP\rangle
=-\sum_i\overline{(uw)_i}\langle a^+_iQ,\omega(u)P\rangle+\sum_i\eta_i(uw)_i\langle a^-_iQ,\omega(u)P\rangle .
\]
Both sides are sums of a complex-linear and a complex-antilinear function of $w\in\C^I$; such a
decomposition is unique. Comparing the linear parts gives the first identity, the antilinear parts
the second.
\end{proof}

Let $P_\pm$ be the projections onto the coordinates with $\eta_i=\pm1$, $\bar u$ the entrywise
conjugate, and
\[
C(u)=P_+uP_++P_-\bar uP_- .
\]
Its diagonal blocks are the block $u_{++}$ of $u$ and the conjugate of the block $u_{--}$; both are
invertible for $u\in\auU(\eta)$ (from $u^*\mathrm{diag}(\eta)u=\mathrm{diag}(\eta)$,
$u_{++}^*u_{++}=1+u_{-+}^*u_{-+}$ and $u_{--}^*u_{--}=1+u_{+-}^*u_{+-}$), so $C(u)$ is invertible.
Fix a map $d$ from $I$ to a set of \emph{labels}, and assume that $u$ is \emph{block diagonal}:
$u_{ij}=0$ when $d(i)\ne d(j)$. Call a polynomial $P$ \emph{harmonic} if $\partial_i\partial_jP=0$
whenever $d(i)=d(j)$ and $\eta_i\ne\eta_j$.

\begin{lemma}\label{ptq:pderiv}
\textup{(a)} If $\partial_kQ=0$ for all $k$ with $d(k)=d(i)$, $\eta_k\ne\eta_i$, then
$\langle Q,\omega(u)\partial_iP\rangle=\sum_k\overline{C(u)_{ki}}\,\langle z_kQ,\omega(u)P\rangle$.

\textup{(b)} If $\partial_iQ=0$ for all $i$ with $d(i)=d(j)$, $\eta_i\ne\eta_j$, then
$\langle z_jQ,\omega(u)P\rangle=\sum_i\overline{A_{ij}}\,\langle Q,\omega(u)\partial_iP\rangle$ with
$A=C(u)^{-1}$.
\end{lemma}
\begin{proof}
(a) If $\eta_i=1$, then $\partial_i=a^-_i$ and the second identity of Lemma~\ref{ptq:ladder} gives
$\sum_k\overline{u_{ki}}\langle a^+_kQ,\omega(u)P\rangle$. For $\eta_k=1$, $a^+_k=z_k$ and
$C(u)_{ki}=u_{ki}$. For $\eta_k=-1$, $a^+_kQ=-\partial_kQ$, and the term vanishes: either
$d(k)\ne d(i)$ and $u_{ki}=0$, or $d(k)=d(i)$, $\eta_k\ne\eta_i$ and $\partial_kQ=0$; and
$C(u)_{ki}=0$. If $\eta_i=-1$, then $\partial_i=-a^+_i$ and the first identity gives
$-\sum_ku_{ki}\eta_k\langle a^-_kQ,\omega(u)P\rangle$. For $\eta_k=-1$, $a^-_k=z_k$, the coefficient is
$u_{ki}=\overline{C(u)_{ki}}$; for $\eta_k=1$, $a^-_kQ=\partial_kQ$ and the term vanishes as before,
while $C(u)_{ki}=0$.

(b) Put $w_k=\langle z_kQ,\omega(u)P\rangle$ and let $b=\{i:d(i)=d(j),\ \eta_i=\eta_j\}$. For $i\in b$
the hypothesis of (a) holds (if $d(k)=d(i)$ and $\eta_k\ne\eta_i$ then $d(k)=d(j)$, $\eta_k\ne\eta_j$),
so $(C(u)^*w)_i=\langle Q,\omega(u)\partial_iP\rangle$. The matrix $C(u)$ is block diagonal for the
partition of $I$ by the pair $(d,\eta)$ ($C(u)_{ki}=0$ if $\eta_k\ne\eta_i$, and $u_{ki}=0$ if
$d(k)\ne d(i)$), and so are $C(u)^*$ and $(C(u)^*)^{-1}=A^*$; $b$ is one of the blocks. Hence
$w_j=(A^*C(u)^*w)_j=\sum_{i\in b}(A^*)_{ji}(C(u)^*w)_i$, and $(A^*)_{ji}=\overline{A_{ij}}$ vanishes for
$i\notin b$.
\end{proof}

\begin{lemma}[Harmonic induction]\label{ptq:induct}
Let $\mathcal S$ be a space of polynomials containing $1$ and containing $z_jP$ whenever $P\in\mathcal S$
and $\partial_iP=0$ for all $i$ with $d(i)=d(j)$, $\eta_i\ne\eta_j$. Then $\mathcal S$ contains every
harmonic polynomial.
\end{lemma}
\begin{proof}
The operators $\partial_i\partial_j$ are homogeneous, so the homogeneous components of a harmonic
polynomial are harmonic; it suffices to treat a harmonic $P$ homogeneous of degree $n$, by induction
on $n$. For $n=0$, $P\in\C\cdot1\subset\mathcal S$. For $n\ge1$, Euler's identity gives
$nP=\sum_jz_j\,\partial_jP$; each $\partial_jP$ is harmonic, homogeneous of degree $n-1$, hence in
$\mathcal S$, and $\partial_i\partial_jP=0$ whenever $d(i)=d(j)$, $\eta_i\ne\eta_j$. So
$z_j\partial_jP\in\mathcal S$ and $P\in\mathcal S$.
\end{proof}

For a matrix $X$ let $L_X$ be the substitution $(L_XQ)(z)=Q(zX)$, that is,
$z_j\mapsto\sum_iX_{ij}z_i$. Then $L_XL_Y=L_{XY}$, and $\langle L_XQ,P\rangle=\langle Q,L_{X^*}P\rangle$:
the Fischer space of polynomials of degree $n$ is the $n$-th symmetric power of the space of linear
forms with the induced inner product, $L_X$ is the symmetric power of its action on linear forms,
and on linear forms $\langle L_Xz_j,z_k\rangle=\overline{X_{kj}}=\langle z_j,L_{X^*}z_k\rangle$.

\begin{proposition}[Coefficients between harmonic polynomials]\label{ptq:coef}
Let $u\in\auU(\eta)$ be block diagonal and $c_u=\langle1,\omega(u)1\rangle$. For harmonic $P,Q$,
\[
\langle Q,\omega(u)P\rangle=c_u\,\langle L_{A}Q,P\rangle,\qquad A=C(u)^{-1}.
\]
\end{proposition}
\begin{proof}
\emph{Step 1: $Q=1$.} Let $\mathcal S=\{P:\langle1,\omega(u)P\rangle=c_u\langle1,P\rangle\}$; it
contains $1$. Let $R\in\mathcal S$ with $\partial_iR=0$ for $d(i)=d(j)$, $\eta_i\ne\eta_j$. The element
$u^{-1}=\mathrm{diag}(\eta)u^*\mathrm{diag}(\eta)$ is block diagonal, and $\omega(u)^*=\omega(u)^{-1}$
satisfies \eqref{ptq:eq-cov} for $u^{-1}$. By Lemma~\ref{ptq:pderiv}(b) for $u^{-1}$ and $P=1$,
$\langle z_jR,\omega(u)^{-1}1\rangle=\sum_i\overline{A'_{ij}}\langle R,\omega(u)^{-1}\partial_i1\rangle=0$. Hence
$\langle1,\omega(u)z_jR\rangle=\overline{\langle z_jR,\omega(u)^{-1}1\rangle}=0=c_u\langle\partial_j1,R\rangle
=c_u\langle1,z_jR\rangle$, so $z_jR\in\mathcal S$. By Lemma~\ref{ptq:induct}, every harmonic $P$
lies in $\mathcal S$.

\emph{Step 2.} Let $\mathcal S'$ be the space of $Q$ such that the identity holds for all harmonic
$P$; $1\in\mathcal S'$ by Step 1. Let $R\in\mathcal S'$ with $\partial_iR=0$ for $d(i)=d(j)$,
$\eta_i\ne\eta_j$, and $P$ harmonic. The derivatives $\partial_iP$ are harmonic. By
Lemma~\ref{ptq:pderiv}(b) and the induction hypothesis,
\[
\langle z_jR,\omega(u)P\rangle=\sum_i\overline{A_{ij}}\,c_u\langle L_AR,\partial_iP\rangle
=c_u\sum_i\overline{A_{ij}}\langle z_iL_AR,P\rangle=c_u\Bigl\langle\sum_iA_{ij}z_i\,L_AR,P\Bigr\rangle
=c_u\langle L_A(z_jR),P\rangle .
\]
So $z_jR\in\mathcal S'$, and Lemma~\ref{ptq:induct} gives the claim for all harmonic $Q$.
\end{proof}

\begin{remark}\label{ptq:rem-conv}
If the Fock model is normalised with $\bar w_ia^+_i-\eta_iw_ia^-_i$ in \eqref{ptq:eq-heis}, then
Lemma~\ref{ptq:ladder} and Proposition~\ref{ptq:coef} hold with $u$ replaced by $\bar u$ in
the coefficients (the proof is the same, with linear and antilinear parts exchanged). Below $u$ is
hermitian positive definite, and so is $\bar u$; the conclusions do not change.
\end{remark}

\subsubsection{The positivity (F)}\label{ptq:ss-F}

We use the following lemma on the vacuum coefficient.

\begin{lemma}[The vacuum coefficient does not vanish]\label{ptq:vac}
In the setting of \S\ref{ptq:ss-coef}, let $u\in\auU(\eta)$ be block diagonal and $\omega(u)$ as
there. Then $c_u=\langle1,\omega(u)1\rangle\ne0$.
\end{lemma}
\begin{proof}
Let $I_\pm=\{i:\eta_i=\pm1\}$, $W_\pm=\{w\in W:w_i=0\text{ for }i\notin I_\pm\}$, $\mathcal P_\pm$ the
polynomials in the variables $z_i$, $i\in I_\pm$, and $H_\pm$ the closure of $\mathcal P_\pm$ in $\pzF$.
Recall that $\varrho$ is a Heisenberg representation: $\varrho(x)\varrho(y)=\beta(x,y)\varrho(x+y)$ with
$|\beta|=1$, and $\varrho$ is irreducible ({\cite[\S2.1]{P49}}).

\emph{Step 1: $\varrho(w)1\in H_\pm$ for $w\in W_\pm$.} By \eqref{ptq:eq-heis}, for $w\in W_+$,
$d\varrho(w)=c_0\sum_{i\in I_+}(w_iz_i-\bar w_i\partial_i)$, and for $w\in W_-$,
$d\varrho(w)=c_0\sum_{i\in I_-}(\bar w_iz_i-w_i\partial_i)$; so $d\varrho(w)$ maps $\mathcal P_\pm$ into
itself. In the Fischer norm, for $P$ of degree $\le d$, $\|z_iP\|^2=\|P\|^2+\|\partial_iP\|^2$ and
$\sum_i\|\partial_iP\|^2\le d\|P\|^2$ (both by $[\partial_i,z_i]=1$, $z_i^*=\partial_i$ and Euler's identity
on homogeneous components, which are orthogonal); hence $\|z_iP\|,\|\partial_iP\|\le\sqrt{d+1}\,\|P\|$ and
$\|d\varrho(w)P\|\le M\sqrt{d+1}\,\|P\|$ with $M=2|c_0|\sum_i|w_i|$. Thus $\|d\varrho(w)^n1\|\le
M^n\sqrt{n!}$. Let $A$ be the skew-adjoint generator of the unitary group $s\mapsto\varrho(sw)$ (a group because $\beta(sw,s'w)=1$, $\beta$ being the exponential of an alternating form); it extends
$d\varrho(w)$ on polynomial vectors, which are smooth, so $1\in\bigcap_nD(A^n)$ with $A^n1=d\varrho(w)^n1$.
Since $\sum_n|s|^nM^n\sqrt{n!}/n!<\infty$ for every $s$, the spectral theorem gives
$\varrho(sw)1=\sum_n s^nA^n1/n!$, a norm-convergent series of elements of $\mathcal P_\pm$. Take $s=1$.

\emph{Step 2.} Suppose $c_u=0$. Polynomials in $\mathcal P_+$ and in $\mathcal P_-$ are harmonic
(each involves only variables with one value of $\eta$), so by Proposition~\ref{ptq:coef}
$\langle Q,\omega(u)P\rangle=0$ for $Q\in\mathcal P_-$, $P\in\mathcal P_+$, and by continuity
$\langle x',\omega(u)x\rangle=0$ for $x'\in H_-$, $x\in H_+$. By Step 1, for $w\in W_+$, $w'\in W_-$,
\[
0=\langle\varrho(w')1,\omega(u)\varrho(w)1\rangle=\langle1,\varrho(w')^{-1}\omega(u)\varrho(w)1\rangle
=\langle1,\omega(u)\varrho(u^{-1}w')^{-1}\varrho(w)1\rangle ,
\]
by \eqref{ptq:eq-cov}; and $\varrho(u^{-1}w')^{-1}\varrho(w)$ is a scalar of modulus one times
$\varrho(w-u^{-1}w')$ (Heisenberg relations; $\varrho(y)^{-1}$ is a scalar times $\varrho(-y)$). Hence
$\langle\omega(u)^{-1}1,\varrho(y)1\rangle=0$ for every $y\in W_++u^{-1}W_-$.

\emph{Step 3: $W_++u^{-1}W_-=W$.} The element $u^{-1}$ lies in $\auU(\eta)$, so its block
$(u^{-1})_{--}$ is invertible (as shown before Lemma~\ref{ptq:pderiv}). If $w'\in W_-$ and
$u^{-1}w'\in W_+$, then $(u^{-1})_{--}w'=0$, so $w'=0$. Thus $W_+\cap u^{-1}W_-=0$, and the complex
dimensions $|I_+|+|I_-|=|I|$ add up.

\emph{Step 4.} By Steps 2 and 3, $\omega(u)^{-1}1$ is orthogonal to $\varrho(y)1$ for all $y\in W$. The
closed span of the $\varrho(y)1$ is $\varrho(W)$-stable (Heisenberg relations) and non-zero, hence is
$\pzF$ by irreducibility. So $\omega(u)^{-1}1=0$, which is absurd since $\omega(u)$ is unitary. Hence
$c_u\ne0$.
\end{proof}

\begin{theorem}[(F)]\label{ptq:F}
Let $(S,T)$ be an ascending or descending pair of the triangle, $\Phi=\Phi_\tau$ the local joint
harmonic of Theorem~\ref{ptr:edge}, and $w\in G$ positive definite. Then
$\langle\Phi,\sigma(w)\Phi\rangle>0$. More generally $\langle P,\sigma(w)P\rangle>0$ for every
polynomial $P\ne0$ none of whose monomials contains two variables of one column from rows of
opposite signs.
\end{theorem}
\begin{proof}
Write $w=\exp X$ with $X$ in the $-1$ eigenspace $\mathfrak p$ of the Cartan involution (hermitian,
off-diagonal for the signs); $\operatorname{tr}X=0$, so $\det w=1$, and $w_t=\exp tX$ ($0\le t\le1$) are
positive definite.

Let $u_t=\iota(w_t,1)$, $x\mapsto xw_t^{-1}$ on $W=\auHom(S,T)_\tau$.

\emph{The element $u=u_1$.} In the coordinates $x_{rc}$, $u$ acts on each column separately (it
mixes the row index only), so it is block diagonal for the labels $d(r,c)=c$; it is hermitian
positive definite for the positive majorant (the right multiplication by the positive definite
$w^{-1}$), and lies in $\auU(\eta)$, $\eta_{rc}=-\varepsilon_r\varepsilon'_c$. A polynomial is harmonic
for $d$ if $\partial_{rc}\partial_{r'c}P=0$ when $\eta_{rc}\ne\eta_{r'c}$, i.e.\ $\varepsilon_r\ne
\varepsilon_{r'}$; this holds if no monomial of $P$ contains $z_{rc}z_{r'c}$ with $\varepsilon_r\ne
\varepsilon_{r'}$, in particular for $\Phi$ (Theorem~\ref{ptr:edge}(a)). The operator $\sigma(w)$
satisfies \eqref{ptq:eq-cov} for $u$ by (W1$_\infty$) of {\cite[\S3.2.7]{P49}}.

\emph{The constant.} $c(t)=\langle1,\sigma(w_t)1\rangle$ is continuous (strong continuity) and $c(0)=1$.
It is real: let $J=\mathrm{diag}(\varepsilon)\in K$; since $w_t$ is hermitian and $w_t^*Jw_t=J$,
$w_t^{-1}=Jw_tJ$. The group $K$ acts on polynomial vectors by linear substitutions times a character
(proof of {\cite[Lemma~3.6]{P49}}), so $\sigma(J)1=\chi1$ with $|\chi|=1$, and
$\overline{c(t)}=\langle1,\sigma(w_t^{-1})1\rangle=\langle\sigma(J)^{-1}1,\sigma(w_t)\sigma(J)^{-1}1\rangle=
|\chi|^{-2}c(t)=c(t)$. It never vanishes: everything said above about $u$ holds for $u_t$ (block
diagonal for the labels $d(r,c)=c$, in $\auU(\eta)$, and \eqref{ptq:eq-cov} holds for $\sigma(w_t)$ and
$u_t$ by (W1$_\infty$)), so Lemma~\ref{ptq:vac} gives $c(t)\ne0$. Hence $c(1)=c_u>0$. Moreover
$C(u)$ is positive definite: its blocks are the principal submatrix $u_{++}$ of the positive definite
$u$ and the conjugate of $u_{--}$.

\emph{Positivity.} $C=C(u)$ is positive definite, hence so is $A=C^{-1}$; write $A=BB^*$ with $B$
invertible. By Proposition~\ref{ptq:coef} with $Q=P$,
\[
\langle P,\sigma(w)P\rangle=c_u\langle L_BL_{B^*}P,P\rangle=c_u\langle L_{B^*}P,L_{B^*}P\rangle=c_u\,
\|L_{B^*}P\|^2>0,
\]
since $L_{B^*}$ is invertible (with inverse $L_{(B^*)^{-1}}$) and $P\ne0$.
\end{proof}

\begin{remark}\label{ptq:rem-F}
Theorem~\ref{ptq:F} is (F) of \S\ref{ptr:ss-Z}. Hence Proposition~\ref{ptr:Z} holds
unconditionally: $Z_\infty\ne0$ for every pair of the triangle, with each local factor positive.
The same argument gives $\langle P,\sigma'(w)P\rangle>0$ for positive definite $w\in G'$ and $P\ne0$
none of whose monomials contains two variables of one row from columns of opposite signs (now
$u=\iota(1,w)$ is block diagonal for the labels $d(r,c)=r$).
\end{remark}

\subsubsection{The smoothing lemma (H)}\label{ptq:ss-H}

We use, besides (D1)--(D3) of {\cite[\S3.1.3]{P49}}, the following published results.
\begin{enumerate}
\item[(D4)] \emph{(Schmid \cite[Lemma~(2.22), Lemma~(3.8), Corollary~(3.9), Corollary~(4.19) and the proof of Theorem~(1.3) on p.~560]{Schmid75}. Also a consequence of Blattner's
formula, proved by Hecht and Schmid. Schmid's positive system is $\{\alpha:\langle\alpha,\lambda\rangle>0\}$, as here; the discrete series concerned are not holomorphic in general, see the paragraph ``On (D3)'' in {\cite[\S3.1.3]{P49}}.)} Every $K$-type of $\pi_\lambda$ has highest weight of the form
$\Lambda^K(\lambda)+\sum_\beta n_\beta\beta$, $n_\beta\in\Z_{\ge0}$, the sum over the positive
noncompact roots $\beta$ (positive meaning $\langle\beta,\lambda\rangle>0$).
\item[(D5)] \emph{(Harish-Chandra \cite[Theorem~1, p.~18]{HC66}; used in \cite[Chapter~I, proof of Lemma~14, p.~15]{HarishChandraAF}; see also \cite[\S2.2, p.~193]{BorelJacquet}.)} Let $f$ be a smooth
function on a connected reductive group, $K$-finite on the left and on the right and
$\mathfrak Z(\mathfrak g)$-finite. Then there is $\beta\in C_c^\infty$ with $\beta*f=f$, where
$(\beta*f)(g)=\int\beta(y)f(y^{-1}g)\,dy$. The published statement is the right-convolution
form: for $f$ smooth, $\mathfrak Z(\mathfrak g)$-finite and right $K$-finite there is
$\alpha\in C_c^\infty$ with $f*\alpha=f$, $(f*\alpha)(g)=\int f(gy^{-1})\alpha(y)\,dy$. The left
form follows: $f^\vee(g)=f(g^{-1})$ is right $K$-finite (as $f$ is left $K$-finite) and
$\mathfrak Z(\mathfrak g)$-finite (as $\mathfrak Z(\mathfrak g)$ is stable under the principal
antiautomorphism of $U(\mathfrak g)$, which intertwines the right and left regular actions under
$f\mapsto f^\vee$). Take $\alpha$ with $f^\vee*\alpha=f^\vee$ and put $\beta(y)=\alpha(y^{-1})$.
Substituting $y\mapsto y^{-1}$ (the group is unimodular),
$(\beta*f)(g)=\int\alpha(y)f(yg)\,dy=\int f^\vee(g^{-1}y^{-1})\alpha(y)\,dy=(f^\vee*\alpha)(g^{-1})=f(g)$.
\end{enumerate}
(D5) applies to $f=e_{\lambda'}$ on $G'$: $e_{\lambda'}=d_{\lambda'}\overline{c_{\lambda'}}$ is a
coefficient of the $K'$-finite vector $v'$ of the irreducible unitary $\pi_{\lambda'}$, so it is
smooth, bi-$K'$-finite, and $\mathfrak Z(\mathfrak g')$ acts on it through the infinitesimal
character. Let $H_\lambda\subset\pzF_\tau$ be the closed span of the images of the $G$-equivariant
isometries $\pi_\lambda\to\sigma$; it is $G$-stable, its orthogonal projection $p$ commutes with
$\sigma(G)$, and $H_\lambda\cong\pi_\lambda\hat\otimes M$ for a Hilbert space $M$, so the $K$-types of
$H_\lambda$ are those of $\pi_\lambda$.

\paragraph*{Four properties of $\mathsf P$}

\begin{lemma}\label{ptq:P}
\textup{(a)} \emph{(Integration by parts.)} For a polynomial $Q$ and $Y$ in $\mathfrak g_\C$,
$\mathsf P(d\sigma(Y)Q)=-\sigma(R_{Y}e_\lambda)Q$, where $R_Y$ is the right derivative extended
complex-linearly. Consequently $\mathsf P(E_{ba}Q)=0$ if $\varepsilon_a=\varepsilon_b$ and
$\lambda_a>\lambda_b$, and $\mathsf P(E_{aa}Q)=\Lambda_a\mathsf PQ$, $\Lambda=\Lambda^K(\lambda)$.

\textup{(b)} \emph{(Image.)} For every $x\in\pzF_\tau$, $\mathsf Px\in H_\lambda$; and $\mathsf P$ is
self-adjoint.

\textup{(c)} \emph{(Off the cone.)} Let $Q$ be a polynomial with $E_{aa}Q=\nu_aQ$ for all $a$ and
$E_{ab}Q=0$ whenever $\varepsilon_a=\varepsilon_b$, $\lambda_a>\lambda_b$, where $\nu-\Lambda$ is not a
sum of roots $e_a-e_b$ with $\lambda_a>\lambda_b$ with nonnegative coefficients. Then
$\mathsf P\,\sigma(y)Q=0$ for all $y\in G$, and $\mathsf P(E(X)Q)=0$ for every $X\in\mathfrak{gl}_n(\C)$,
$E(X)$ the row action {\cite[(12)]{P49}}.

\textup{(d)} \emph{(Fixing.)} There is $\beta_0\in C_c^\infty(G')$ with
$\sigma'(\beta_0)\mathsf P\Phi=\mathsf P\Phi$.
\end{lemma}
\begin{proof}
(a) For real $Y$: since $Q$ is a smooth vector with derivative $d\sigma(Y)Q$,
$\mathsf Pd\sigma(Y)Q=\frac d{dt}\big|_0\mathsf P\sigma(\exp tY)Q$, and
$\mathsf P\sigma(\exp tY)Q=\int e_\lambda(g\exp(-tY))\sigma(g)Q\,dg$ (right invariance of Haar
measure). Now $e_\lambda(g\exp(-tY))-e_\lambda(g)=-\int_0^t(R_Ye_\lambda)(g\exp(-sY))ds$, so by Fubini
(as $R_Ye_\lambda\in L^1$ and $\sigma$ is unitary)
$\mathsf P\sigma(\exp tY)Q-\mathsf PQ=-\int_0^t\int(R_Ye_\lambda)(g)\sigma(g\exp sY)Q\,dg\,ds$; the inner
integral is continuous in $s$ by dominated convergence, which gives the formula at $t=0$. Both sides
are complex-linear in $Y$. Now $e_\lambda(g)=d_\lambda\langle\pi_\lambda(g)v,v\rangle$ ($v$ the unit highest
vector of the lowest $K$-type, $c_\lambda(g)=\langle v,\pi_\lambda(g)v\rangle$), so for real $Y$,
$R_Ye_\lambda(g)=d_\lambda\langle\pi_\lambda(g)d\pi_\lambda(Y)v,v\rangle$, and since this is antilinear in
$d\pi_\lambda(Y)v$, for complex $Y$ it is $d_\lambda\langle\pi_\lambda(g)d\pi_\lambda(\bar Y)v,v\rangle$, where
$\bar Y=-JY^*J$ ($J=\mathrm{diag}(\varepsilon)$) is the conjugation of $\mathfrak g_\C$ with respect to
$\mathfrak g$. For $\varepsilon_a=\varepsilon_b$, $\bar E_{ba}=-E_{ab}$ and $\bar E_{aa}=-E_{aa}$. As
$d\pi_\lambda(E_{ab})v=0$ for $\lambda_a>\lambda_b$ and $d\pi_\lambda(E_{aa})v=\Lambda_av$ (with $\Lambda_a$ real),
$R_{E_{ba}}e_\lambda=0$ and $R_{E_{aa}}e_\lambda=-\Lambda_ae_\lambda$; with $d\sigma(E_{ab})=E_{ab}$ on
polynomials this gives the two consequences.

(b) $e_\lambda(x^{-1})=\overline{e_\lambda(x)}$ (proof of {\cite[Theorem~3.15]{P49}}), so $\sigma(e_\lambda)^*=
\sigma(e_\lambda)$. Let $y=\mathsf Px$. For $Y\in\mathfrak g$, $\sigma(\exp tY)y=\int e_\lambda(\exp(-tY)g)
\sigma(g)x\,dg$ and $e_\lambda(\exp(-tY)g)=d_\lambda\langle\pi_\lambda(g)v,\pi_\lambda(\exp tY)v\rangle$. As the left
derivatives of $e_\lambda$ are integrable, $y$ is a smooth vector and, for $Y,Y'\in\mathfrak g_\C$,
$d\sigma(Y)d\sigma(Y')y=d_\lambda\int\langle\pi_\lambda(g)v,d\pi_\lambda(Y)d\pi_\lambda(Y')v\rangle\sigma(g)x\,dg$
(the expression is linear in the second slot, so complex-linear in $Y,Y'$). Hence $y$ has
$T$-weight $\Lambda$, is killed by the $E_{ab}$ with $\varepsilon_a=\varepsilon_b$, $\lambda_a>\lambda_b$,
and $d\sigma(\Omega)y=\kappa_\lambda y$, because $v$ has these properties in $\pi_\lambda$. By the proof of
{\cite[Lemma~3.5]{P49}} (which uses only these three properties, $e_\lambda\in L^1$ and (D1)--(D3)), the
closed span of $\sigma(G)y$ is a direct integral of copies of $\pi_\lambda$, hence is contained in
$H_\lambda$.

(c) The polynomial vectors of bounded degree form a finite-dimensional space stable under $K$
(which acts by linear substitutions, proof of {\cite[Lemma~3.6]{P49}}), so $Q$ is $K$-finite; being a
highest weight vector of weight $\nu$ for $\mathfrak k_\C$ in a finite-dimensional unitary
representation, it generates an irreducible $K$-module of highest weight $\nu$. The projection $p$
commutes with $K$ and with the derived action on smooth vectors, so $pQ$ is zero or generates a
$K$-type of $H_\lambda$ of highest weight $\nu$. By (D4) the highest weights of the $K$-types of
$H_\lambda$ lie in $\Lambda+\sum\Z_{\ge0}\beta$, contained in $\Lambda$ plus the cone of the hypothesis;
so $pQ=0$. Let $y\in G$ and $x=\mathsf P\sigma(y)Q$. By (b), $x\in H_\lambda$ and $\mathsf Px\in H_\lambda$,
and
\[
\|x\|^2=\langle\mathsf P\sigma(y)Q,x\rangle=\langle\sigma(y)Q,\mathsf Px\rangle=\langle p\,\sigma(y)Q,\mathsf
Px\rangle=\langle\sigma(y)pQ,\mathsf Px\rangle=0 .
\]
Differentiating $t\mapsto\mathsf P\sigma(\exp tY)Q=0$ at $0$ gives $\mathsf P(E(Y)Q)=0$ for
$Y\in\mathfrak g$, hence for $Y\in\mathfrak g_\C=\mathfrak{gl}_n(\C)$ by linearity.

(d) By (D5) choose $\beta_0\in C_c^\infty(G')$ with $\beta_0*e_{\lambda'}=e_{\lambda'}$. Since
$\overline{c_{\lambda'}(y^{-1}g)}=c_{\lambda'}(g^{-1}y)$, this reads $\int\beta_0(y)c_{\lambda'}(hy)\,dy=
c_{\lambda'}(h)$ for all $h$. Let $x=\mathsf P\Phi$. In the proof of {\cite[Theorem~3.15]{P49}} (valid for the
present pair, Theorem~\ref{ptr:edge}(b)) it is shown that $x$ is a smooth vector for $\sigma'$ of
$T'$-weight $\Lambda'$, killed by the compact raising operators of the columns, with
$d\sigma'(\Omega')x=\kappa_{\lambda'}x$; so {\cite[Lemma~3.5]{P49}} gives $\langle x,\sigma'(h)x\rangle=
\|x\|^2c_{\lambda'}(h)$. Then for $g\in G'$,
\[
\langle\sigma'(g)x,\sigma'(\beta_0)x\rangle=\int\beta_0(y)\langle x,\sigma'(g^{-1}y)x\rangle dy=\|x\|^2\!\int
\beta_0(y)c_{\lambda'}(g^{-1}y)dy=\|x\|^2c_{\lambda'}(g^{-1})=\langle\sigma'(g)x,x\rangle .
\]
So $\sigma'(\beta_0)x-x$, which lies in the closed span of $\sigma'(G')x$, is orthogonal to it, and
is $0$.
\end{proof}

\paragraph*{Fischer's argument}

For a polynomial space we write $\mathcal D_j$ for the polynomials of degree $\le j$. The operators
$E_{ab}$, $E'_{cd}$ with equal signs preserve degree, and on polynomials $E_{ab}^*=E_{ba}$,
$E'^*_{cd}=E'_{dc}$ for the Fischer product (from $(a^+_i)^*=\eta_ia^-_i$, $s_{ab}=1$; the shifts are
real). For $\varepsilon_a\ne\varepsilon_b$, $\Delta_{ab}=\sum_{c:\eta_{ac}=-1}\partial_{ac}\partial_{bc}$ as in
{\cite[\S3.1.4]{P49}}; note that $\eta_{bc}=-\eta_{ac}$.

\begin{lemma}\label{ptq:fischer}
Let $\mathcal G$ be a space of polynomials containing $\mathcal D_{j-1}$, and $Q\in\mathcal D_j$
orthogonal to $\mathcal G\cap\mathcal D_j$. Let $\varepsilon_a\ne\varepsilon_b$ and suppose that
$E_{ba}\Delta_{ab}Q\in\mathcal G$. Then $\Delta_{ab}Q=0$. The same holds for columns, with
$\Delta'_{cd}$ and $E'_{dc}$.
\end{lemma}
\begin{proof}
Put $Q''=\Delta_{ab}Q\in\mathcal D_{j-2}$. By {\cite[(12)]{P49}}, $\pm E_{ba}=\sum_ca^+_{bc}a^-_{ac}$; for
$\eta_{ac}=-1$ (so $\eta_{bc}=1$) the term is $z_{bc}z_{ac}$, and for $\eta_{ac}=1$ it is
$-\partial_{bc}\partial_{ac}$. Hence $\sum_{c:\eta_{ac}=-1}z_{ac}z_{bc}Q''=\pm E_{ba}Q''+
\sum_{c:\eta_{ac}=1}\partial_{ac}\partial_{bc}Q''$. The first term is in $\mathcal G\cap\mathcal D_j$; the
second is in $\mathcal D_{j-4}\subset\mathcal G$. By adjointness,
\[
\|Q''\|^2=\langle\Delta_{ab}Q,Q''\rangle=\Bigl\langle Q,\sum_{c:\eta_{ac}=-1}z_{ac}z_{bc}Q''\Bigr\rangle=0 .
\]
For columns the computation is the same with rows and columns exchanged.
\end{proof}

\begin{lemma}\label{ptq:tensor}
Let $Q$ satisfy $E_{aa}Q=\Lambda_aQ$ for all $a$ and $E_{ab}Q=0$ for $\varepsilon_a=\varepsilon_b$,
$\lambda_a>\lambda_b$. Suppose that, for $\varepsilon_a\ne\varepsilon_b$, $\Delta_{ab}Q=0$ whenever
$\Delta_{ab}Q$ has these two properties for the weight $\Lambda+e_a-e_b$. Then $\Delta_{ab}Q=0$ for all
$\varepsilon_a\ne\varepsilon_b$.
\end{lemma}
\begin{proof}
From the canonical commutation relations, for $\varepsilon_x=\varepsilon_z$ and $\varepsilon_a\ne\varepsilon_b$,
$[E_{xz},\Delta_{ab}]=\delta_{za}\Delta_{xb}-\delta_{bx}\Delta_{az}$ (the $\Delta_{ab}$ transform like
the matrix units). Argue by induction on $\lambda_b-\lambda_a$. $\Delta_{ab}Q$ has $T$-weight
$\Lambda+e_a-e_b$ by the case $x=z$. Let $\varepsilon_x=\varepsilon_z$, $\lambda_x>\lambda_z$. Then
$E_{xz}\Delta_{ab}Q=\delta_{za}\Delta_{xb}Q-\delta_{bx}\Delta_{az}Q$; in the first term $\lambda_x>\lambda_a$,
so $\lambda_b-\lambda_x<\lambda_b-\lambda_a$ (and $\varepsilon_x=\varepsilon_a\ne\varepsilon_b$); in the second
$\lambda_z<\lambda_b$, so $\lambda_z-\lambda_a<\lambda_b-\lambda_a$. By induction both vanish, so
$\Delta_{ab}Q$ has the two properties, and is $0$ by hypothesis.
\end{proof}

\paragraph*{Howe's argument}

\begin{proposition}[Howe's argument]\label{ptq:howe}
Let $\Phi$ be the joint harmonic of the matching pair $(\lambda,\lambda')$ ({\cite[\S3.1.5]{P49}}) and
$\mathcal G$ a space of polynomials such that: \textup{(i)} $E'_{cd}\mathcal G\subset\mathcal G$ for all
$c,d$; \textup{(ii)} $\Phi\in\mathcal G$; \textup{(iii)} $E(X)Q\in\mathcal G$ for every $X$ and every $Q$ as in
Lemma~\ref{ptq:P}(c); \textup{(iv)} $E_{ba}\mathcal G'\subset\mathcal G$ for all polynomials $\mathcal G'$,
when $\varepsilon_a=\varepsilon_b$, $\lambda_a>\lambda_b$; \textup{(v)} $(E_{aa}-\Lambda_a)Q\in\mathcal G$ for
every polynomial $Q$. Then $\mathcal G$ contains all polynomials.
\end{proposition}
\begin{proof}
We show $\mathcal D_j\subset\mathcal G$ by induction on $j$ ($\mathcal D_{-1}=0$). Suppose
$\mathcal D_{j-1}\subset\mathcal G$ and $\mathcal D_j\not\subset\mathcal G$. Let $\mathcal N$ be the
orthogonal complement of $\mathcal G\cap\mathcal D_j$ in $\mathcal D_j$; $\mathcal N\ne0$.

\emph{$\mathcal N$ is stable under the compact column operators.} If $\varepsilon'_c=\varepsilon'_d$, then
$E'_{dc}$ preserves $\mathcal G\cap\mathcal D_j$ by (i) and degree, so its adjoint $E'_{cd}$ preserves
$\mathcal N$. The $E'_{cc}$ are self-adjoint and commute, so $\mathcal N$ is a finite-dimensional
module for the compact column algebra on which the Cartan subalgebra acts semisimply; it contains
a nonzero vector $Q$ that is an $E'_{cc}$-eigenvector for all $c$ and is killed by every $E'_{cd}$ with
$\varepsilon'_c=\varepsilon'_d$, $\lambda'_c>\lambda'_d$ (a weight vector whose weight is maximal for a linear
order compatible with these roots).

\emph{Row weight and raising.} By (v), $\mathcal G$ is $E_{aa}$-stable, so $\mathcal N$ is, and
$(E_{aa}-\Lambda_a)Q\in\mathcal N\cap\mathcal G=0$. If $\varepsilon_a=\varepsilon_b$, $\lambda_a>\lambda_b$, then
by (iv) $E_{ba}\mathcal D_j\subset\mathcal G\cap\mathcal D_j$, so $E_{ab}Q$, which lies in $\mathcal D_j$,
is orthogonal to $\mathcal D_j$, and $E_{ab}Q=0$.

\emph{Column Laplacians.} For $\varepsilon'_c\ne\varepsilon'_d$, $E'_{dc}\Delta'_{cd}Q\in\mathcal G$ by (i)
(as $\Delta'_{cd}Q\in\mathcal D_{j-2}\subset\mathcal G$), so $\Delta'_{cd}Q=0$ by
Lemma~\ref{ptq:fischer}.

\emph{Row Laplacians.} We apply Lemma~\ref{ptq:tensor}: let $\varepsilon_a\ne\varepsilon_b$ and suppose
$\Delta_{ab}Q$ has $T$-weight $\Lambda+e_a-e_b$ and is killed by the compact raising operators. If
$e_a-e_b$ is not a nonnegative combination of roots $e_x-e_z$ with $\lambda_x>\lambda_z$, then (iii)
gives $E_{ba}\Delta_{ab}Q\in\mathcal G$, and $\Delta_{ab}Q=0$ by Lemma~\ref{ptq:fischer}. Otherwise,
pairing with $\lambda$ shows $\lambda_a>\lambda_b$. All row parameters satisfy $|\xi_r|>\tfrac t2$
(every row is matched, {\cite[(13)]{P49}}), and $\xi_r-\lambda_r$ is constant; so either
$\xi_a>\tfrac t2$, or $\xi_b<\xi_a<-\tfrac t2$. By {\cite[Lemma~3.10(i)]{P49}}, $s\,w_r\ge0$ for
$s=\ptqsgn\xi_r$; the polynomial $Q$ is killed by the row raising operators and by the $\Delta'_{cd}$
and has $E_{aa}Q=\Lambda_aQ$, so {\cite[Lemma~3.12(c)]{P49}} gives: in the first case
$\partial_{ac}Q=0$ whenever $\eta_{ac}=-1$; in the second $\partial_{bc}Q=0$ whenever $\eta_{bc}=1$. Every
term $\partial_{ac}\partial_{bc}Q$ of $\Delta_{ab}Q$ has $\eta_{ac}=-1$, $\eta_{bc}=1$, so
$\Delta_{ab}Q=0$. By Lemma~\ref{ptq:tensor}, all row Laplacians kill $Q$.

\emph{Conclusion.} $Q$ is a pre-joint harmonic ({\cite[Definition~3.9]{P49}}), so $Q=c\Phi$ by
{\cite[Proposition~3.13]{P49}}. By (ii), $Q\in\mathcal G\cap\mathcal N=0$, a contradiction.
\end{proof}

\begin{theorem}[(H)]\label{ptq:H}
Let $(S,T)$ be an ascending or descending pair of the triangle. For every polynomial $Q$ at $\tau$
there is $\beta_Q\in C_c^\infty(G')$ with $\mathsf PQ=\sigma'(\beta_Q)\mathsf P\Phi$. Consequently
\textup{(H)} of \S\ref{ptr:ss-mult} holds: for every polynomial vector $Q$ of the Fock space of
$\auHom(S,T)\otimes\R$ there is $\alpha_Q\in C_c(\auU(T)_\infty)$ with $\mathsf PQ=\int\alpha_Q(x)
\omega_\infty(1,x)\mathsf PP_{S,T}\,dx$.
\end{theorem}
\begin{proof}
Let $\mathcal G$ be the space of polynomials $Q$ such that $\mathsf PQ=\sigma'(\beta)\mathsf P\Phi$ for some
$\beta\in C_c^\infty(G')$. We check the hypotheses of Proposition~\ref{ptq:howe}.

(i) Let $Q\in\mathcal G$, $\mathsf PQ=\sigma'(\beta)\mathsf P\Phi$, and $Y\in\mathfrak g'$. As $\mathsf P$ is
bounded and commutes with $\sigma'$, and $Q$ is a smooth vector for $\sigma'$ with derivative
$d\sigma'(Y)Q=E'(Y)Q$, we get $\mathsf PE'(Y)Q=\frac d{dt}|_0\sigma'(\exp tY)\mathsf PQ$. And
$\sigma'(\exp tY)\sigma'(\beta)\mathsf P\Phi=\sigma'(\beta(\exp(-tY)\,\cdot))\mathsf P\Phi$, whose derivative at
$0$ is $\sigma'(\beta_Y)\mathsf P\Phi$ with $\beta_Y(x)=\frac d{dt}|_0\beta(\exp(-tY)x)\in C_c^\infty$
(differentiation under the integral, $\beta$ being smooth with compact support and $\sigma'$ unitary).
So $E'(Y)Q\in\mathcal G$ for $Y\in\mathfrak g'$, hence for $Y\in\mathfrak g'_\C$ by linearity.

(ii) Lemma~\ref{ptq:P}(d).

(iii), (iv), (v): for such polynomials $\mathsf P$ of them is $0=\sigma'(0)\mathsf P\Phi$, by
Lemma~\ref{ptq:P}(c) and (a) respectively.

So $\mathcal G$ contains all polynomials. For the second assertion: $\mathsf P=\bigotimes_\tau\mathsf
P_\tau$, $P_{S,T}=\bigotimes_\tau\Phi_\tau$, $\omega_\infty(1,x)=\bigotimes_\tau\sigma'_\tau(x_\tau)$, and every
polynomial vector is a finite sum of products $\bigotimes_\tau Q_\tau$; for a product take
$\alpha_Q(x)=\prod_\tau\beta_{Q_\tau}(x_\tau)$, and sum. The identity holds in the Hilbert space; both sides
lie in $\pzF_s$: $\mathsf P$ maps $\pzF_s$ into itself (the integral defining it converges absolutely in
$\pzF_s$ by {\cite[Lemma~8.2(a)]{P49}} and (A1), as in the proof of {\cite[Proposition~8.4(b)]{P49}}), so
$\mathsf PQ,\mathsf PP_{S,T}\in\pzF_s$, and the Bochner integral of the continuous compactly supported
$\pzF_s$-valued function $x\mapsto\alpha_Q(x)\omega_\infty(1,x)\mathsf PP_{S,T}$ converges in $\pzF_s$ by
{\cite[Lemma~8.2(a)]{P49}}; and $\pzF_s$ injects continuously into the Hilbert space, so the identity
holds in $\pzF_s$.
\end{proof}

\begin{remark}\label{ptq:rem-H}
Theorem~\ref{ptq:H} is (H) of \S\ref{ptr:ss-mult}, for all ascending and descending pairs, hence for
the pairs $(U_1,A)$, $(U_2,A)$ of every square of the triangle (Remark~\ref{ptr:rem-prev}). So
Lemma~\ref{ptr:conv-lift}, Theorem~\ref{ptr:square} and Corollary~\ref{ptr:tri-b} hold without the
hypothesis (H). The proof uses only that the parameters match with distinct entries, the
integrability of $e_\lambda$, $e_{\lambda'}$ and their derivatives, (V) and (D1)--(D5); it applies
verbatim to any pair satisfying these (for instance the edges of the grids,
Theorem~\ref{pgv:G3}(e)); for edges read with the groups exchanged use Theorem~\ref{ptw:Hc} and
Proposition~\ref{ptw:revedge}.
\end{remark}

\subsubsection{Dependencies of \texorpdfstring{{\cite[Propositions~14.35 and~14.37]{P46}}}{[Row D, Propositions~14.35 and~14.37]}}\label{ptq:ss-status}

Theorem~\ref{ptq:F} uses Lemma~\ref{ptq:vac} (proved here), (W1$_\infty$) and {\cite[Theorem~3.29]{P49}}; Lemma~\ref{ptq:P}, Proposition~\ref{ptq:howe} and
Theorem~\ref{ptq:H} use (V) and (D1)--(D5); Theorem~\ref{ptr:square} and Corollary~\ref{ptr:tri-b} use
(D1)--(D5) and Theorem~\ref{pcx:X}. The finite parts of the descent characters and the comparison with
\S\ref{pc:ss-triangle} are treated in \S\ref{ptv:sec} (Propositions~\ref{ptv:chars}, \ref{ptv:arch},
\ref{ptv:tri}); the adjoint law is proved there from (RS), which is Theorem~\ref{ptw:RS}; the edge
hypothesis (E1) for the grid squares (Remark~\ref{ptr:rem-grid}) is Lemma~\ref{pgv:E1g}. Hence
{\cite[Propositions~14.35 and~14.37]{P46}} rest, besides the proofs written here, on Theorem~\ref{pcx:X}
and the published results (D1)--(D5), (W1$_\infty$), (W2$_\infty$), (M), (S).




% ===== C-triangle3.tex =====
% proofs/C-triangle3.tex
% Continuation of proofs/C-triangle.tex (ptr:) and proofs/C-triangle2.tex (ptq:): the embedding of
% F_s in the Fock space, the characters of the descents, the reduction of the adjoint law
% (Theorem pc:tri(c)) to the square law of the reversed squares, and the final check of
% Lemma pc:arch and Theorem pc:tri. Labels of this file start with ptv:.
% Only amsmath/amssymb/amsthm are used. Conventions and macros: proofs/C-triangle.tex.
\providecommand{\auU}{\mathrm U}
\providecommand{\auHom}{\operatorname{Hom}}
\providecommand{\auGL}{\mathrm{GL}}
\providecommand{\auNm}{\operatorname{Nm}}
\providecommand{\auA}{\mathbb A}
\providecommand{\auO}{\mathcal O}
\providecommand{\auF}{\mathbb F}
\providecommand{\aubar}[1]{\overline{#1}}
\providecommand{\psrho}{\varrho}
\providecommand{\psHS}{\mathcal H}
\providecommand{\psvac}{\mathbf v}
\providecommand{\pzF}{\mathfrak F}
\providecommand{\pcF}{\mathcal F}
\providecommand{\pcG}{\mathsf G}
\providecommand{\pcm}{\mathfrak m}
\providecommand{\pcR}{R}
\providecommand{\pcMat}{\mathrm M}
\providecommand{\ptrwt}{\operatorname{wt}}
\providecommand{\ptvtr}{\operatorname{tr}}
\providecommand{\ptvF}{\mathsf F}

\subsection{The adjoint law and the clauses of the triangle}\label{ptv:sec}

This section continues \S\ref{ptr:sec} and \S\ref{ptq:sec}. Notation is that of
\S\ref{pc:ss-data}, \S\ref{pc:ss-triangle}, \S\ref{pc:ss-exchange} and \S\ref{ptr:sec}; inner
products are conjugate-linear in the first variable; for a function $f$ on a group, $\aubar f$ is
its complex conjugate.

\emph{Outline.}
\begin{enumerate}
\item[(1)] The Banach space $\pzF_s$ of {\cite[\S8.1]{P49}} injects continuously into the Fock
space (Lemma~\ref{ptv:embed}); this is the fact used in the last step of the proof of
Theorem~\ref{ptq:H}.
\item[(2)] The characters of the descents are fixed consistently by the squares and by
Theorem~\ref{pc:tri}(a) (Proposition~\ref{ptv:chars}): the compatibility of finite parts
stated after Lemma~\ref{ptr:chars} holds, as an identity of automorphic characters.
\item[(3)] The adjoint law, Theorem~\ref{pc:tri}(c), is reduced to the square law for the
\emph{reversed squares} (statement (RS) of \S\ref{ptv:ss-RS}), by complete arguments at the finite
places: reversal of a theta kernel modulo an ideal (Proposition~\ref{ptv:rev}), translation by
central elements (Lemma~\ref{ptv:central}), and a Gauss sum identity for the extended data
(Proposition~\ref{ptv:fourier}); (RS) itself is Theorem~\ref{ptw:RS}.
\item[(4)] Lemma~\ref{pc:arch} is proved for all its pairs, and every clause of
Theorem~\ref{pc:tri} is proved, clause (c) through (RS) (\S\ref{ptv:ss-check}).
\end{enumerate}
The edge hypothesis (E1) for the grid squares (Remark~\ref{ptr:rem-grid}) is Lemma~\ref{pgv:E1g}.

\subsubsection{The space $\pzF_s$ in the Fock space}\label{ptv:ss-embed}

Recall from {\cite[\S8.1]{P49}}: $\pzF$ is the Fock space (a Hilbert space), $\pzF^{\mathrm{sm}}$ its
space of smooth vectors, $\partial_w$ ($w\in W_\infty$) the derivative at $s=0$ of $\psrho_\infty(sw)$
on $\pzF^{\mathrm{sm}}$, $\|\Psi\|_s$ the maximum of the norms $\|\partial_{w_1}\cdots\partial_{w_n}\Psi\|$
over $0\le n\le 2s$ and $w_i$ in a fixed real basis $\mathsf B$, and $\pzF_s$ the completion of
$\pzF^{\mathrm{sm}}$ for $\|\cdot\|_s$. The definition asserts that this completion is ``contained
in $\pzF$''; this is the content of the following lemma.

\begin{lemma}\label{ptv:embed}
The identity of $\pzF^{\mathrm{sm}}$ extends to a continuous linear map $\iota\colon\pzF_s\to\pzF$ with
$\|\iota x\|\le\|x\|_s$, and $\iota$ is injective. Consequently a Bochner integral in $\pzF_s$ of a
continuous compactly supported (or $\|\cdot\|_s$-integrable) function is mapped by $\iota$ to the
Bochner integral in $\pzF$ of the same function, and two elements of $\pzF_s$ with the same image
in $\pzF$ are equal.
\end{lemma}
\begin{proof}
The term $n=0$ of the maximum gives $\|\Psi\|\le\|\Psi\|_s$ on $\pzF^{\mathrm{sm}}$, so the identity is
uniformly continuous and extends to $\iota$ with the stated bound. Let $x\in\pzF_s$ with $\iota x=0$
and let $x_k\in\pzF^{\mathrm{sm}}$ converge to $x$ in $\|\cdot\|_s$; so $x_k\to0$ in $\pzF$. Fix a word
$D=\partial_{w_1}\cdots\partial_{w_n}$ with $n\le2s$, $w_i\in\mathsf B$. The sequence $Dx_k$ is Cauchy in
$\pzF$, with a limit $y$. For $w\in W_\infty$ the operators $\psrho_\infty(sw)$, $s\in\R$, form a
strongly continuous one-parameter unitary group (on a line the cocycle of the Heisenberg group is
trivial), so $\psrho_\infty(sw)^*=\psrho_\infty(-sw)$ and, for $a,b\in\pzF^{\mathrm{sm}}$,
$\langle\partial_wa,b\rangle=-\langle a,\partial_wb\rangle$ (differentiate
$\langle\psrho_\infty(sw)a,b\rangle=\langle a,\psrho_\infty(-sw)b\rangle$ at $s=0$); and
$\pzF^{\mathrm{sm}}$ is stable under the $\partial_w$. Hence for every $b\in\pzF^{\mathrm{sm}}$
\[
\langle y,b\rangle=\lim_k\langle Dx_k,b\rangle=\lim_k(-1)^n\langle x_k,\partial_{w_n}\cdots\partial_{w_1}b
\rangle=0 .
\]
As $\pzF^{\mathrm{sm}}$ is dense in $\pzF$, $y=0$. So $\|Dx_k\|\to0$ for each of the finitely many
words $D$, i.e.\ $\|x_k\|_s\to0$, and $x=\lim x_k=0$. The statements on Bochner integrals follow
because a continuous linear map commutes with Bochner integrals.
\end{proof}

\begin{remark}\label{ptv:rem-embed}
This is exactly what the last step of the proof of Theorem~\ref{ptq:H} uses: the two sides of the
identity there are elements of $\pzF_s$ (the Bochner integral of $x\mapsto\alpha_Q(x)\omega_\infty(1,x)
\mathsf PP_{S,T}$ converges in $\pzF_s$ by {\cite[Lemma~8.2(a)]{P49}}; $\mathsf PQ\in\pzF_s$ for the
same reason, $\mathsf P$ being a Bochner integral of $e_S(x)\omega_\infty(x^{-1},1)Q$), their images in
$\pzF$ coincide, and Lemma~\ref{ptv:embed} gives equality in $\pzF_s$.
\end{remark}

\subsubsection{The characters of the descents}\label{ptv:ss-chars}

Recall (\S\ref{pc:ss-data}) that $\theta_{\alpha,\beta}$ for $\auHom(S,T)$ carries
$\alpha(\det g)^{-1}\beta(\det h)$, $g\in\auU(S)$, $h\in\auU(T)$: $\alpha$ is the \emph{source}
character and $\beta$ the \emph{target} character. Write $M_k=r+kt$ for the rank of the nodes in
column $k$. The lift $a_{i,n}$ has source character $\mu_{i+n+1}\upsilon_{M_{n+1}}$ and target
character $\mu_{i+n+1}\upsilon_{M_n}$. The descent $d_{i,n}$ is a transposed lift for
$\auHom(\pcG_{i+1,n},\pcG_{i,n+1})$ with source character $\alpha^\downarrow_{i,n}$ (on
$\pcG_{i+1,n}$) and target character $\beta^\downarrow_{i,n}$ (on $\pcG_{i,n+1}$). For a square at
$(i,n)$ and a twist $\vartheta$ put $\gamma_n=\nu_{n+1}\vartheta^{\,n+1}\upsilon^{M_{n+1}}$; this is
$\gamma=\nu^{-1}\upsilon^{N/2}$ of \S\ref{ptr:ss-chars} ($N=2M_{n+1}$). Recall
$\mu_j=\mu^\flat\nu_1^{\,j}\vartheta^j$ (\S\ref{pc:ss-triangle}), so
$\mu_{j+1}=\mu_j\,\nu_1\vartheta$, and $\upsilon_k=\upsilon^{\lfloor k/2\rfloor}$.

In the square at $(i,n)$ (with $U_1=\pcG_{i,n+1}$, $U_2=\pcG_{i+1,n+1}$, $A=\pcG_{i,n+2}$,
$A'=\pcG_{i+1,n}$), condition (E2) of Definition~\ref{ptr:def-square} consists of four equations:
\begin{align*}
&\text{at }A:&\mu_{i+n+2}\upsilon_{M_{n+1}}\cdot\beta^\downarrow_{i,n+1}&=\gamma_n,&
&\text{at }A':&\alpha^\downarrow_{i,n}\cdot\mu_{i+n+2}\upsilon_{M_{n+1}}&=\gamma_n,\\
&\text{at }U_1:&\mu_{i+n+2}\upsilon_{M_{n+2}}\cdot\beta^\downarrow_{i,n}&=\gamma_n,&
&\text{at }U_2:&\alpha^\downarrow_{i,n+1}\cdot\mu_{i+n+2}\upsilon_{M_n}&=\gamma_n .
\end{align*}
(At $A$: target characters of $a_{i,n+1}$ and of $d_{i,n+1}$; at $A'$: source characters of
$d_{i,n}$ and of $a_{i+1,n}$; at $U_1$: source character of $a_{i,n+1}$ times target character of
$d_{i,n}$; at $U_2$: source character of $d_{i,n+1}$ times target character of $a_{i+1,n}$.) So each
square fixes all four characters of $d_{i,n}$ and $d_{i,n+1}$, and the characters of $d_{i,n+1}$
($n+1\ge1$) are fixed twice, by the squares at $(i,n)$ and at $(i,n+1)$; those of $d_{i,0}$ are
fixed by the square at $(i,0)$ and, for $\beta^\downarrow_{i,0}$, also by
Theorem~\ref{pc:tri}(a).

\begin{proposition}\label{ptv:chars}
\textup{(a)} For $i+n+2\le m$ define $\alpha^\downarrow_{i,n},\beta^\downarrow_{i,n},
\alpha^\downarrow_{i,n+1},\beta^\downarrow_{i,n+1}$ by the four equations of the square at $(i,n)$.
If also $i+n+3\le m$, the characters of $d_{i,n+1}$ so obtained coincide with those defined by the
square at $(i,n+1)$:
\[
\begin{gathered}\beta^\downarrow_{i,n+1}=\gamma_n\,(\mu_{i+n+2}\upsilon_{M_{n+1}})^{-1}=\gamma_{n+1}\,
(\mu_{i+n+3}\upsilon_{M_{n+3}})^{-1},\\
\alpha^\downarrow_{i,n+1}=\gamma_n\,(\mu_{i+n+2}\upsilon_{M_n})^{-1}=\gamma_{n+1}\,
(\mu_{i+n+3}\upsilon_{M_{n+2}})^{-1}.\end{gathered}
\]
So for each $\vartheta$ there is one family of characters $(\alpha^\downarrow_{i,n},
\beta^\downarrow_{i,n})_{i+n+1\le m}$ satisfying \textup{(E2)} in every square of the triangle (the
characters of $d_{i,m-i-1}$, which lies in no square as the lower edge, are fixed by the square at
$(i,m-i-2)$ when $m-i\ge2$; for $i=m-1$ the characters of $d_{m-1,0}$ are those given by
\textup{(W7)}, Remark~\ref{ptw:rem-W7}, and \textup{(b)} identifies its target character). Their weights are those of Definition~\ref{ptr:desc-char}, and their
tame parts at the primes of $\pcm_\vartheta$ are those of $\mu_{i+1}^{-1}$.

\textup{(b)} The target character of the descent $d_{i,0}$ defined by Theorem~\ref{pc:tri}(a)
(through \textup{(W7)}, Remark~\ref{ptr:rem-desc}) is $\beta^\downarrow_{i,0}$ of \textup{(a)}.
\end{proposition}
\begin{proof}
(a) In the quotient of the two expressions for $\beta^\downarrow_{i,n+1}$,
\[
\frac{\gamma_{n+1}}{\gamma_n}\cdot\frac{\mu_{i+n+2}}{\mu_{i+n+3}}\cdot\frac{\upsilon_{M_{n+1}}}
{\upsilon_{M_{n+3}}}=\nu_1\vartheta\,\upsilon^{t}\cdot(\nu_1\vartheta)^{-1}\cdot
\upsilon^{\lfloor M_{n+1}/2\rfloor-\lfloor M_{n+3}/2\rfloor}=\upsilon^{t-t}=1,
\]
because $\gamma_{n+1}/\gamma_n=\nu_1\vartheta\,\upsilon^{M_{n+2}-M_{n+1}}$, $\nu_{n+2}=\nu_{n+1}\nu_1$,
and $M_{n+3}=M_{n+1}+2t$, so $\lfloor M_{n+3}/2\rfloor=\lfloor M_{n+1}/2\rfloor+t$. The same computation
with $M_n,M_{n+2}$ in place of $M_{n+1},M_{n+3}$ gives the second identity. The weights:
$\ptrwt\mu_j=-\kappa_j$, $\ptrwt\nu_k=-k\Delta$, $\vartheta$ has weight $0$ and $\kappa_{j}+\Delta=
\kappa_{j+1}$; so $\ptrwt\beta^\downarrow_{i,n}=-(n+1)\Delta+M_{n+1}+\kappa_{i+n+2}-\lfloor M_{n+2}/2
\rfloor=\kappa_{i+1}+M_{n+1}-\lfloor(M_{n+1}+t)/2\rfloor$, and as $M_{n+1}+t$ and $M_n=M_{n+1}-t$ have
the same parity, $M_{n+1}-\lfloor(M_{n+1}+t)/2\rfloor=\lceil(M_{n+1}-t)/2\rceil=\lceil M_n/2\rceil$, which
is Definition~\ref{ptr:desc-char}; for $\alpha^\downarrow_{i,n}$ the computation is that of
Lemma~\ref{ptr:chars}. The tame parts: $\upsilon$ is unramified at all finite places, and
$\gamma_n\mu_{i+n+2}^{-1}=\nu_{n+1}\vartheta^{n+1}\mu_{i+n+2}^{-1}\upsilon^{M_{n+1}}=\mu_{i+1}^{-1}
\upsilon^{M_{n+1}}$, because $\mu_{i+n+2}=\mu_{i+1}(\nu_1\vartheta)^{n+1}$; so every character above
is $\mu_{i+1}^{-1}$ times a power of $\upsilon$.

(b) By (a), $\beta^\downarrow_{i,0}=\gamma_0(\mu_{i+2}\upsilon_{M_2})^{-1}=\mu_{i+1}^{-1}\upsilon^{k_0}$
for an integer $k_0$. On the other hand, by (W7) (Lemma~\ref{pc:SW}) the complex conjugate of the
kernel of $a_{i,0}$ is a kernel for $\auHom(\pcG_{i,0}^-,\pcG_{i,1})=\auHom(\pcG_{i+1,0},\pcG_{i,1})$
whose target character is the conjugate $\mu_{i+1}^{-1}\upsilon_{M_0}^{-1}$ of the target character
of $a_{i,0}$ times a power of $\upsilon$ (the character $\lambda$ of (W7) is a product of powers of
$\upsilon\circ\det$); so it is $\mu_{i+1}^{-1}\upsilon^{k_1}$ for an integer $k_1$ (explicitly $k_1=\lceil r/2\rceil$,
Remark~\ref{ptw:rem-W7}). Its weight is that of
Definition~\ref{ptr:desc-char} by Remark~\ref{ptr:rem-desc} (Proposition~\ref{ptr:match}(c)), and
so is the weight of $\beta^\downarrow_{i,0}$ by (a). As $\upsilon$ has weight $1$ at every
$\tau$, $k_0=k_1$, and the two characters are equal.
\end{proof}

\begin{remark}\label{ptv:rem-chars}
A priori the characters on $U_1$ and $U_2$ of the two kernels at $A$ need not agree, and the
discrepancy $\nu_0$ is carried by a factor $\nu_0(\det g_2)$ in the integrand; in
Theorem~\ref{ptr:square} this is the factor $\lambda(\det g_2)$. Proposition~\ref{ptv:chars} shows
that in the triangle no choice is involved: the characters of every descent are determined by any
one of the conditions that involve it, and all these determinations agree. With it,
Corollary~\ref{ptr:tri-b} needs no further input concerning finite parts.
\end{remark}

\subsubsection{Adjoints of lifts and translations by central elements}\label{ptv:ss-adj}

Throughout this subsection $(S,T)$ is a pair of spaces of ranks $N\le N'$, $\pcm$ a modulus,
$\alpha,\beta$ characters, $P$ a Fock vector with finite moments of $\auHom(S,T)\otimes\R$ and $\Psi$ a
function on $\pcMat_{N'\times N}(\pcR_\pcm)$ \emph{supported on injective matrices} (at every prime of
$\pcm$). We assume that $(S,T)$ with $(\alpha,\beta,P)$ is an edge in the sense of
Theorem~\ref{ptr:edge}(c), (d), from a type $e_S$ to a type $e_T$: the lift
$\Theta=\Theta_{\alpha,\beta}(P,\Psi)$ maps $\pcF(S,e_S)$ to $\pcF(T,e_T)$, the transposed lift
$\Theta^{\mathsf t}=\Theta^{\mathsf t}_{\alpha,\beta}(P,\Psi)$ maps forms of type $\aubar{e_T}$ to forms
of type $\aubar{e_S}$, and all these forms are bounded and rapidly decreasing.

\begin{lemma}[Adjoints]\label{ptv:adj}
For $z\in\pcF(T,e_T)$ put $\Theta^*z(g)=\int_{[\auU(T)]}\aubar{\theta_{\alpha,\beta}(g,h;P,\Psi)}\,z(h)
\,dh$. Then $\Theta^*z=\aubar{\Theta^{\mathsf t}\aubar z}\in\pcF(S,e_S)$, and
$\langle\Theta y,z\rangle=\langle y,\Theta^*z\rangle$ for all $y\in\pcF(S,e_S)$, the double integral
being absolutely convergent.
\end{lemma}
\begin{proof}
$\aubar z$ is a form of type $\aubar{e_T}$, so $\Theta^{\mathsf t}\aubar z$ is a form of type
$\aubar{e_S}$ and its conjugate, which is $\Theta^*z$, lies in $\pcF(S,e_S)$ (it is bounded,
continuous, invariant under an open subgroup because the datum is, and reproduced by $e_S$). By
Proposition~\ref{pcx:decay}(b) ($\Psi$ is supported on injective matrices), for $N_0$ given there
are $B,C$ with $|\theta_{\alpha,\beta}(g,h;P,\Psi)|\le C\,\mathrm{ht}(g)^{-N_0}\mathrm{ht}(h)^{B}$. As in
the proof of Proposition~\ref{ptr:conv}, $|z|\,\mathrm{ht}^B$ is integrable on $[\auU(T)]$, $y$ is
bounded and $\mathrm{ht}^{-N_0}$ is integrable on $[\auU(S)]$ for $N_0$ large ((R1) of
\S\ref{pcx:ss-conv}). So $\iint|\theta(g,h)\,y(g)\,z(h)|\,dg\,dh<\infty$, and Fubini's theorem gives
$\langle\Theta y,z\rangle=\iint\aubar{\theta(g,h)}\,\aubar{y(g)}\,z(h)\,dg\,dh=\langle y,\Theta^*z\rangle$.
\end{proof}

The same computation applies to a transposed lift: if $\Theta^{\mathsf t}$ maps $\pcF(T,\aubar{e_T}')$
to $\pcF(S,\aubar{e_S}')$ for types $e'$ and its datum is supported on injective matrices, its adjoint
is $(\Theta^{\mathsf t})^*y(h)=\int_{[\auU(S)]}\aubar{\theta(g,h)}\,y(g)\,dg$, with
$\langle\Theta^{\mathsf t}x,y\rangle=\langle x,(\Theta^{\mathsf t})^*y\rangle$.

\begin{lemma}[Central translations]\label{ptv:central}
Let $\zeta\in\auU(1)(\auA_{E^+,f})$ and write $\zeta$ also for the scalar element of $\auU(S)(\auA_f)$
and of $\auU(T)(\auA_f)$; it is central. There is $\eta\in\C$, $|\eta|=1$, depending on $\zeta$,
$\alpha$, $\beta$, $S$ and $T$ but not on $y$, with $\Theta(\psrho(\zeta)y)=\eta\,\psrho(\zeta)\Theta(y)$ for all
$y$, where $\psrho$ is right translation; the same holds for transposed lifts.
\end{lemma}
\begin{proof}
The pair $(\zeta,\zeta)$ acts on $\auHom(S,T)(\auA)$ by $x\mapsto\zeta x\zeta^{-1}=x$. By (W1) the operator
$\omega_{\auHom(S,T)}(\zeta,\zeta)$ commutes with all $\psrho(x)$, hence is a scalar, of modulus one as
$\omega$ is unitary; so $\omega_{\alpha,\beta}(\zeta,\zeta)=\eta_0$ with $|\eta_0|=1$, and as
$(\zeta,\zeta)$ is central, $\theta_{\alpha,\beta}(g\zeta,h\zeta;P,\Psi)=\eta_0\,\theta_{\alpha,\beta}
(g,h;P,\Psi)$. Since $\zeta$ is central, $g\mapsto g\zeta^{-1}$ preserves $[\auU(S)]$ and its measure,
and
\[
\Theta(\psrho(\zeta)y)(h)=\int\theta(g,h)\,y(g\zeta)\,dg=\int\theta(g\zeta^{-1},h)\,y(g)\,dg=
\aubar{\eta_0}\int\theta(g,h\zeta)\,y(g)\,dg=\aubar{\eta_0}\,\Theta(y)(h\zeta),
\]
using $\theta(g\zeta^{-1},h)=\aubar{\eta_0}\theta(g,h\zeta)$. Put $\eta=\aubar{\eta_0}$. For transposed
lifts exchange the roles of $g$ and $h$.
\end{proof}

\subsubsection{Reversal of a theta kernel modulo an ideal}\label{ptv:ss-rev}

The adjoint of a descent is a lift whose kernel is the complex conjugate of a kernel of
$\auHom(\pcG_{i+1,n},\pcG_{i,n+1})$, viewed as a function on $\auU(\pcG_{i,n+1})\times
\auU(\pcG_{i+1,n})$. We write it as a kernel of the reversed space $\auHom(\pcG_{i,n+1},\pcG_{i+1,n})$.

\begin{lemma}[Invariant pairings]\label{ptv:pairing}
Let $F$ be a finite field, $p,q\ge1$, and $e\colon\pcMat_{p\times q}(F)\times\pcMat_{q\times p}(F)\to
\C^\times$ bi-additive, non-degenerate (for each $\rho\ne0$, $e(\rho,\cdot)\ne1$, and for each $v\ne0$,
$e(\cdot,v)\ne1$) and invariant: $e(g\rho h^{-1},hvg^{-1})=e(\rho,v)$ for $g\in\auGL_p(F)$,
$h\in\auGL_q(F)$. Then there is a non-trivial character $\psi$ of $F$ with
$e(\rho,v)=\psi(\ptvtr(\rho v))$.
\end{lemma}
\begin{proof}
Fix a non-trivial character $\psi_1$ of $F$. Every character of $\pcMat_{p\times q}(F)$ is
$\rho\mapsto\psi_1(\ptvtr(\rho X))$ for a unique $X\in\pcMat_{q\times p}(F)$ (the pairing
$(\rho,X)\mapsto\psi_1(\ptvtr\rho X)$ is non-degenerate, and both groups have the same order). So
$e(\rho,v)=\psi_1(\ptvtr(\rho\,\mu(v)))$ for a map $\mu$ of $\pcMat_{q\times p}(F)$ to itself, additive by
bi-additivity, injective by non-degeneracy. Invariance gives $\mu(hvg^{-1})=h\mu(v)g^{-1}$. The
invertible matrices span $\pcMat_k(F)$ additively for every $k\ge1$ and every $F$ (for $k=1$ because
$1$ generates $\mathbb F_p\subset F$ and $F^\times$ spans $F$; for $k\ge2$ every elementary matrix
$E_{ij}$ is a difference of two invertible matrices); so $\mu(XvY)=X\mu(v)Y$ for all $X\in\pcMat_q(F)$,
$Y\in\pcMat_p(F)$. With $v=E_{ij}$ and $X=E_{ii}$, $Y=E_{jj}$, $\mu(E_{ij})=E_{ii}\mu(E_{ij})E_{jj}$ is a
multiple $c_{ij}E_{ij}$; with $X=E_{ki}$, $Y=E_{jl}$, $c_{kl}=c_{ij}=:c$, and with $X=a1$ ($a\in F$),
$\mu(av)=a\mu(v)$; so $\mu(v)=cv$, $c\ne0$. Put $\psi(x)=\psi_1(cx)$.
\end{proof}

For spaces $S,T$ let $S^-$ be $S$ with the negated form. By (W7) there is a conjugate-linear isometry
$\mathsf C\colon\psHS(\auHom(S,T))\to\psHS(\auHom(S^-,T))$ with $\mathsf C\psrho(x)=\psrho^-(x)\mathsf C$
and $\theta^-\circ\mathsf C=\aubar\theta$; by {\cite[Proposition~7.17]{P49}}, applied to
$\auHom(S^-,T)$, the map $w\mapsto w^\dagger$ identifies $\auHom(S^-,T)'$ with $\auHom(T,S)$ and gives
$\mathrm{Fl}$ with $\mathrm{Fl}\,\psrho^-(x)=\psrho'(x^\dagger)\mathrm{Fl}$ and
$\theta'\circ\mathrm{Fl}=c_\theta\theta^-$. Put $\mathcal A=\mathrm{Fl}\circ\mathsf C$ and
$\phi(x)=x^\dagger$, where $\dagger$ is the adjoint for the forms of $S$ and $T$ (the adjoint for
$S^-$ and $T$ followed by the identification is $x\mapsto-x^\dagger$; the sign is absorbed in
$\phi$ and changes nothing below). Then $\mathcal A$ is conjugate-linear, isometric (the local maps are
normalised to be isometric below; the normalisation constant of $\mathrm{Fl}$ is absorbed in
$c_\theta$), and
\begin{equation}\label{ptv:eq-A}
\mathcal A\,\psrho(x)=\psrho'(\phi(x))\,\mathcal A,\qquad\theta'\circ\mathcal A=c_\theta\,\aubar\theta,
\qquad\mathcal A\,\omega_{\auHom(S,T)}(g,h)=\lambda_S(\det g)\lambda_T(\det h)^{-1}\,
\omega_{\auHom(T,S)}(h,g)\,\mathcal A,
\end{equation}
where $\lambda_S,\lambda_T$ are the products of the character of (W7) with the character
$\upsilon(\det g)^{N'}\upsilon(\det h)^{-N}$ of {\cite[Proposition~7.17]{P49}} ($\det$ the
determinants on $S$ and $T$); they are powers of $\upsilon$ depending only on $N,N'$, and in fact both
are trivial by Lemma~\ref{ptw:W7}. We keep them in the formulas below and use $\lambda_S=\lambda_T=1$
from \S\ref{ptv:ss-RS} on. As in \S\ref{pcx:ss-flip} we take
$\mathrm{Fl}$, $\mathsf C$, hence $\mathcal A$, as tensor products of local maps $\mathcal A_v$, each
conjugate-linear, isometric and intertwining the local Heisenberg operators through $\phi$, with
the normalisation at the finite places fixed below; then $c_\theta\ne0$ depends only on $S$, $T$
and the frames.

Let $\zeta_\alpha=\aubar\alpha/\alpha\in\auU(1)(\auA_{E^+,f})$, with $\alpha$ the finite idele of
\S\ref{pc:ss-data} (the vacuum of the datum is that of the lattice $\alpha\auHom(L_S,L_T)$).

\begin{proposition}[Reversal]\label{ptv:rev}
For every modulus $\pcm$ there are a character $\psi_\pcm$ of the additive group of $\pcR_\pcm$,
non-trivial on every factor $\auF_w$, a scalar $\zeta_\pcm\in\auU(1)(\auA_{E^+,f})$ equal to
$\zeta_\alpha$ outside the primes of $\pcm$, both independent of $S,T$, and for all $\alpha,\beta$ a
number $c=c(S,T,\pcm,\alpha)$ with $|c|^2=C_{S,T}\,(\#\pcR_\pcm)^{-NN'}$, $C_{S,T}>0$ depending only on
$S,T$, such that for every Fock vector $P$, every $\Psi$ on $\pcMat_{N'\times N}(\pcR_\pcm)$ and all
$g\in\auU(S)(\auA)$, $h\in\auU(T)(\auA)$
\[
\aubar{\theta_{\alpha,\beta}(g,h;P,\Psi)}=c\,\theta_{\beta\lambda_T,\,\alpha\lambda_S}(h,g\zeta_\pcm;
P^\flat,\Psi^\flat),\qquad\Psi^\flat(\varrho)=\sum_{v\in\pcMat_{N'\times N}(\pcR_\pcm)}\aubar{\Psi(v)}\,
\psi_\pcm(\ptvtr(\varrho v)),
\]
where the right side is a kernel of $\auHom(T,S)$, $\varrho\in\pcMat_{N\times N'}(\pcR_\pcm)$, and
$P^\flat=\mathcal A_\infty P$, a Fock vector with the same moments as $P$.
\end{proposition}
\begin{proof}
\emph{Places outside $\pcm$.} Let $v$ be a finite place, $L=\alpha\auHom(L_S,L_T)_v$. By (W4) the
vacuum $\psvac_L$ spans the line of vectors on which $\psrho(l)$, $l\in L$, acts by $\epsilon_L$; by
\eqref{ptv:eq-A}, $\mathcal A_v\psvac_L$ spans the line on which $\psrho'(l')$, $l'\in\phi(L)$, acts by
$\aubar{\epsilon_L}\circ\phi^{-1}$. Now $\phi(\alpha x)=\aubar\alpha\phi(x)$ and
$\auHom(L_S,L_T)^\dagger=\auHom(L_T,L_S)$ (the lattices are unimodular), so $\phi(L)=\zeta_\alpha L'$,
$L'=\alpha\auHom(L_T,L_S)_v$, and the scalar $\zeta_\alpha$ acts on $\auHom(T,S)$ as the element
$(1,\zeta_\alpha)$ of $\auU(T)\times\auU(S)$ (by $x'\mapsto\zeta_\alpha x'$). By (W4), $\omega'(1,\zeta_\alpha)
\psvac_{L'}$ spans the line of $\zeta_\alpha L'$; we normalise $\mathcal A_v$ by
$\mathcal A_v\psvac_L=\omega'_v(1,\zeta_{\alpha,v})\psvac_{L'}$ (both are unit vectors). At almost all
$v$, $\zeta_{\alpha,v}$ is a unit and $\omega'_v(1,\zeta_{\alpha,v})\psvac_{L'}=\psvac_{L'}$ (W4), so the
local maps form a restricted tensor product.

\emph{Places of $\pcm$.} Let $v=w\aubar w$, $w\mid\pcm$, $q=\#\auF_w$, $\varpi$ a uniformiser. In the
model (W5), $\psHS_v(\auHom(S,T))$ is the space of Schwartz--Bruhat functions on
$X=\pcMat_{N'\times N}(E_w)$; $\psrho(a)$, $a\in X$, is translation by $a$, and $\psrho(a')$, $a'$ in
the $\aubar w$-component $X^*$, is multiplication by $x\mapsto[x,a']$, a character which, for $a'$
integral, is trivial on $M_w=a_w\pcMat(\auO_w)$ exactly when $a'\in M_{\aubar w}$, the dual lattice
(Lemma~\ref{pgx:vac}; $M_w=\pcMat(\auO_w)$ and $M_{\aubar w}=\pcMat(\auO_{\aubar w})$ when $\alpha$ is a unit at
$w,\aubar w$; in general, with the data of Definition~\ref{pgx:conv}, read $\pcMat(\auO_w)$ as $M_w$ and
$\pcMat(\auO_{\aubar w})$ as $M_{\aubar w}$ in the rest of this paragraph), and which identifies
$\varpi^{-1}\pcMat(\auO_{\aubar w})/\pcMat(\auO_{\aubar w})$ with the dual of
$\pcMat(\auO_w)/\varpi\pcMat(\auO_w)$ ({\cite[(63)]{P49}}: $\lambda_w$ is trivial on $\auO_w$ and not on
$\varpi^{-1}\auO_w$). So the space $D$ of the functions $\Psi^\natural$ ($\Psi$ on
$\pcMat_{N'\times N}(\auF_w)$) is the space of vectors fixed by $\psrho(\Lambda)$,
$\Lambda=\varpi\pcMat(\auO_w)\oplus\pcMat(\auO_{\aubar w})$, and the operators $\psrho(a)$,
$a\in\pcMat(\auO_w)$, and $\psrho(a')$, $a'\in\varpi^{-1}\pcMat(\auO_{\aubar w})$, act on $D$ by
translations $\tau_a$ and by all the characters $m_b$ of $\pcMat_{N'\times N}(\auF_w)$. The adjoint
exchanges the components at $w$ and $\aubar w$ (it involves the conjugation of $E_v$), so $\phi$ maps
$X$ onto the $\aubar w$-component $X'^*$ of $\auHom(T,S)_v$ and $X^*$ onto $X'=\pcMat_{N\times N'}(E_w)$,
and it is equivariant: $\phi(hxg^{-1})=g\phi(x)h^{-1}$ for unitary $g,h$. The only lattices of $X'$ (of
$X'^*$) stable under $\auGL_N(\auO_w)\times\auGL_{N'}(\auO_w)$ are the $\varpi^k\pcMat(\auO)$ (such a
lattice is an $\auO$-module stable under left and right multiplication by matrix rings, hence a
module over $\pcMat_N(\auO)\otimes\pcMat_{N'}(\auO)^{\mathrm{op}}$ inside $\pcMat_{N\times N'}(E_w)$); and
$\phi$ is an isometry for $[\,,]$. Hence $\phi(\Lambda)=\varpi^{k}\pcMat(\auO_w)\oplus\varpi^{1-k}
\pcMat(\auO_{\aubar w})$ for an integer $k$ (depending only on $\alpha_w$, the frames being unimodular),
and $\mathcal A_vD$ is the space of functions on $X'$ that are $\varpi^k\pcMat(\auO_w)$-periodic and
supported on $\varpi^{k-1}\pcMat(\auO_w)$. Let $\zeta_{\pcm,v}\in\auU(1)(E^+_v)$ be the scalar whose
$w$-component moves $\varpi^{k-1}\pcMat(\auO_w)$ to $\pcMat(\auO_w)$ under the action of $(1,\zeta)$ in
(W5); then $T_v=\omega'_v(1,\zeta_{\pcm,v})^{-1}\mathcal A_v$ maps $D$ bijectively onto the space $D'$ of the
functions $\Psi'^\natural$, $\Psi'$ on $\pcMat_{N\times N'}(\auF_w)$ (as $(1,\zeta)$ normalises the
Heisenberg group, $T_v$ intertwines $\psrho$ and $\psrho'\circ\phi_\zeta$ with $\phi_\zeta$ the
composite of $\phi$ and the action of $(1,\zeta)^{-1}$). Through $T_v$ the translations $\tau_a$ on $D$
become multiplications by characters $\rho\mapsto e(\rho,a)$ on $D'$, and the multiplications $m_b$
become translations; $e$ is bi-additive, non-degenerate ($T_v$ is bijective) and invariant under
$\auGL_N(\auF_w)\times\auGL_{N'}(\auF_w)$ (by the equivariance of $\phi$, of the reduction maps of
$K_S$, $K_T$ (Lemma~\ref{pc:level}) and the commutation of $(1,\zeta)$ with $\omega'$). By
Lemma~\ref{ptv:pairing}, $e(\rho,a)=\psi_w(\ptvtr(\rho a))^{-1}$ for a non-trivial character $\psi_w$
of $\auF_w$. The map $\ptvF\colon\Psi\mapsto\Psi^\flat$ of the statement (at $w$) satisfies
$\ptvF(\tau_a\Psi)(\rho)=\psi_w(\ptvtr\rho a)^{-1}\ptvF(\Psi)(\rho)$, the same rule. For the
multiplications: if $T_vm_b=\tau_{b'}T_v$, the commutation relation $\tau_am_b=m_b(a)m_b\tau_a$ and the
conjugate-linearity of $T_v$ give $e(b',a)=m_b(a)$ for all $a$, which determines $b'$; the same
computation for $\ptvF$ gives the same $b'$. So $U=\ptvF^{-1}T_v$ is a linear map of $D$ commuting with
all $\tau_a$ and all $m_b$. The operators commuting with all $m_b$ are multiplications by functions,
and such a multiplication commutes with all translations only if it is constant; so $U$ is a scalar
$\varepsilon_v$, i.e.\ $T_v(\Psi^\natural)=\varepsilon_v(\Psi^\flat)^\natural$. Norms: in the
normalisation of \S\ref{pc:ss-data}, $\|\Psi^\natural\|^2=q^{-NN'}\sum|\Psi|^2$, and by finite
Plancherel $\sum_\rho|\Psi^\flat(\rho)|^2=q^{NN'}\sum_v|\Psi(v)|^2$; $\mathcal A_v$ and $\omega'_v$ preserve
norms; so $|\varepsilon_v|^2=q^{-NN'}$. We normalise $\mathcal A_v$ at $w$ by $\mathcal A_v\psvac=
\omega'_v(1,\zeta_{\pcm,v})u'$, where $u'$ is the unit vector proportional to
$(\mathbf 1^\flat)^\natural=q^{NN'}\delta_0^\natural$ (the characteristic function of $\varpi\pcMat(\auO_w)$, i.e.\ the vacuum of
$\varpi\pcMat(\auO_w)\oplus\varpi^{-1}\pcMat(\auO_{\aubar w})$): the vacuum is $\mathbf 1^\natural$ up to its norm, its image
lies on the line of $\omega'_v(1,\zeta_{\pcm,v})(\mathbf 1^\flat)^\natural$ because $T_v(\mathbf 1^\natural)=\varepsilon_v(\mathbf 1^\flat)^\natural$, so this is possible; this
changes $\varepsilon_v$ by a number of modulus one. Put $\psi_\pcm=\prod_w\psi_w$ on
$\pcR_\pcm=\prod_w\auF_w$; it is non-trivial on every factor.

\emph{Globally.} Put $\zeta_\pcm=\prod_v\zeta_{\pcm,v}$ with $\zeta_{\pcm,v}=\zeta_{\alpha,v}$ for $v$ not
below $\pcm$. By the local statements, $\mathcal A[P,\Psi]=\varepsilon\,\omega'(1,\zeta_\pcm)
[P^\flat,\Psi^\flat]$ with $\varepsilon=\prod_{w\mid\pcm}\varepsilon_w$, $|\varepsilon|^2=(\#\pcR_\pcm)^{-NN'}$.
By \eqref{ptv:eq-A},
\begin{align*}
\aubar{\theta_{\alpha,\beta}(g,h;P,\Psi)}&=\alpha(\det g)\beta(\det h)^{-1}\,\aubar{\theta(\omega(g,h)
[P,\Psi])}=c_\theta^{-1}\alpha(\det g)\beta(\det h)^{-1}\theta'(\mathcal A\,\omega(g,h)[P,\Psi])\\
&=c_\theta^{-1}(\alpha\lambda_S)(\det g)(\beta\lambda_T)(\det h)^{-1}\,\varepsilon\,\theta'\bigl(
\omega'(h,g)\,\omega'(1,\zeta_\pcm)[P^\flat,\Psi^\flat]\bigr)\\
&=c_\theta^{-1}\varepsilon\,(\alpha\lambda_S)(\det\zeta_\pcm)^{-1}\,\theta_{\beta\lambda_T,\alpha\lambda_S}
(h,g\zeta_\pcm;P^\flat,\Psi^\flat).
\end{align*}
So $c=c_\theta^{-1}\varepsilon\,(\alpha\lambda_S)(\det\zeta_\pcm)^{-1}$ and $|c|^2=|c_\theta|^{-2}
(\#\pcR_\pcm)^{-NN'}$. The archimedean map $\mathcal A_\infty$ is induced by the real-linear isometry
$\phi$ followed by complex conjugation, which preserve the moments (\S\ref{pcx:ss-flip}). The
integer $k$ and the scalar $\zeta_{\pcm}$ depend only on $\alpha$ and $\lambda_w$, not on $S$, $T$.
Neither does $\psi_w$: for $a\in X$ the operator $\mathcal A_v\psrho(a)\mathcal A_v^{-1}=\psrho'(\phi(a))$
is the multiplication by $x'\mapsto[x',\phi(a)]'$ on the functions on $X'$, and by the definition of the
bracket ({\cite[\S5.1.1]{P49}}: $[x,y]=\psi(\tfrac12\jmath(\langle y,x\rangle-\langle x,y\rangle))$ with
$\langle w_1,w_2\rangle=\ptvtr(w_1^\dagger w_2)$, for every pair of spaces) and $\phi(a)^\dagger=a$,
$[x',\phi(a)]'=\psi_v\bigl(\tfrac12\jmath(\ptvtr(ax')-\aubar{\ptvtr(ax')})\bigr)$ is the value at
$\ptvtr(ax')\in E_w$ of a character of $E_w$ that depends only on $\psi_v$ and $\jmath$; conjugation by
$\omega'_v(1,\zeta_{\pcm,v})$ replaces $x'$ by a fixed scalar multiple, and the reductions modulo $\pcm$ are
taken in the coordinates $a_w^{-1}x$, $a_w^{-1}x'$ for all spaces. So $e(\rho,a)$ is the value at
$\ptvtr(\rho a)$ of one character of $\auF_w$, the same for all $S$, $T$.
\end{proof}

\begin{corollary}[Adjoints of descents]\label{ptv:desc-adj}
Let $(S,T)=(\pcG_{i+1,n},\pcG_{i,n+1})$ and let $d=\Theta^{\mathsf t}_{\alpha,\beta}(P,\varepsilon)\colon
\pcF_{i,n+1}\to\pcF_{i+1,n}$ be a descent, with $\varepsilon$ supported on injective matrices. Its
adjoint $d^*\colon\pcF_{i+1,n}\to\pcF_{i,n+1}$, $d^*y(h)=\int_{[\auU(S)]}\aubar{\theta(g,h)}y(g)\,dg$, is
\[
d^*=c\cdot R\circ\psrho(\zeta_\pcm^{-1}),\qquad R=\Theta^{\mathsf t}_{\beta\lambda_T,\alpha\lambda_S}
(P^\flat,\varepsilon^\flat),
\]
where $R$, the \emph{reversed lift}, is the transposed lift of the space $\auHom(T,S)$ (integration
over $[\auU(S)]$), and $c$ is the number of Proposition~\ref{ptv:rev}.
\end{corollary}
\begin{proof}
By Proposition~\ref{ptv:rev}, $d^*y(h)=c\int\theta_{\beta\lambda_T,\alpha\lambda_S}(h,g\zeta_\pcm)y(g)\,dg$,
and the substitution $g\mapsto g\zeta_\pcm^{-1}$ (central, measure preserving) gives
$c\int\theta(h,g)\,y(g\zeta_\pcm^{-1})\,dg$. That $d^*$ is the adjoint of $d$ and takes values in
$\pcF_{i,n+1}$ is the remark after Lemma~\ref{ptv:adj} with Theorem~\ref{ptr:edge}(d) for the
descending pair (the conjugate of the lift of $\aubar y$ along it).
\end{proof}

\subsubsection{A Gauss sum identity for extended data}\label{ptv:ss-fourier}

In this subsection $\pcR$ is a finite product of finite fields $\auF_w$, $\psi$ a character of
$(\pcR,+)$ non-trivial on every factor, and $\chi$ a character of $\pcR^\times$ non-trivial on every
factor $\auF_w^\times$, extended by $0$ to non-units. For finite sets $\mathsf I,\mathsf J$ and a
function $f$ on $\pcMat_{\mathsf I\times\mathsf J}(\pcR)$ put
$\ptvF f(\rho)=\sum_vf(v)\psi(\ptvtr(\rho v))$, $\rho\in\pcMat_{\mathsf J\times\mathsf I}(\pcR)$. Since
$(\rho,v)\mapsto\psi(\ptvtr\rho v)$ is a non-degenerate pairing, $\sum_s\psi(\ptvtr(Cs))=
\#\pcMat\cdot[C=0]$ and $\sum_\rho|\ptvF f(\rho)|^2=\#\pcMat_{\mathsf I\times\mathsf J}(\pcR)\sum_v|f(v)|^2$.

\begin{lemma}[Gauss sum of $\chi\circ\det$]\label{ptv:gauss}
Let $\mathsf T$ be a finite set and $G=\sum_{d\in\auGL_{\mathsf T}(\pcR)}\chi(\det d)\psi(\ptvtr d)$. Then
$|G|^2=(\#\pcR)^{|\mathsf T|^2}$, and every function $S$ on $\pcMat_{\mathsf T}(\pcR)$ with
$S(gd)=\chi(\det g)S(d)$ for all $g\in\auGL_{\mathsf T}(\pcR)$ satisfies
$\ptvF S(y)=G\,\chi(\det y)^{-1}S(1)$ for all $y$ (with $\chi(\det y)^{-1}=0$ if $\det y$ is not a
unit).
\end{lemma}
\begin{proof}
If $d$ is not invertible, some component $d_w$ is singular; choose $u\ne0$ with ${}^tu\,d_w=0$ and,
for $t\in\auF_w^\times$, $v$ with ${}^tuv=t-1$; then $g_w=1+v\,{}^tu$ has $g_wd_w=d_w$ and $\det g_w=t$.
With $g=g_w$ at $w$ and $1$ elsewhere, $S(d)=\chi_w(t)S(d)$ for all $t$, so $S(d)=0$. For invertible $d$,
$S(d)=\chi(\det d)S(1)$. Hence $\ptvF S(y)=S(1)\sum_{d\in\auGL}\chi(\det d)\psi(\ptvtr(yd))$. If $y$ is
invertible, the substitution $d=y^{-1}d'$ gives $\chi(\det y)^{-1}G$. If $y_w$ is singular, choose
$g_w$ as above with $y_wg_w=y_w$ (right kernel) and $\det g_w=t$; the substitution $d\mapsto gd$
multiplies the sum by $\chi_w(t)$, so the sum vanishes. For $|G|$: take $S_0=\chi\circ\det$ on
$\auGL$, $0$ elsewhere; Plancherel gives $|G|^2\#\auGL=(\#\pcR)^{|\mathsf T|^2}\#\auGL$.
\end{proof}

\paragraph*{The extension of a datum}
Let $\mathsf N$, $\mathsf T$ be finite sets, $\mathsf M=\mathsf N\sqcup\mathsf T$ (in the triangle:
$\mathsf N$ indexes the nodes of column $n$, $\mathsf M$ those of column $n+1$, $\mathsf T$ the new
block). For $p$ on $\pcMat_{\mathsf M\times\mathsf N}(\pcR)$ its extension $p^+$ is the function on
$\pcMat_{(\mathsf M\sqcup\mathsf T)\times(\mathsf N\sqcup\mathsf T)}(\pcR)$ with
\[
p^+\begin{pmatrix}\eta&s\\0&1\end{pmatrix}=\frac{p(-\eta)}{\#\pcMat_{\mathsf M\times\mathsf T}(\pcR)}
\qquad(\eta\in\pcMat_{\mathsf M\times\mathsf N},\ s\in\pcMat_{\mathsf M\times\mathsf T}),
\]
and $p^+=0$ at all other matrices; this is the function $p^+$ of \S\ref{pc:ss-exchange} (Lemma~\ref{pc:next}),
for $\mathsf N_1=\mathsf M$, $\mathsf N_0=\mathsf N$, $\mathsf T'=\mathsf T$. Note
$\aubar{p^+}=(\aubar p)^+$. If $p$ is supported on injective matrices, so is $p^+$.

\begin{lemma}\label{ptv:F-ext}
For $\rho=\begin{pmatrix}A&B\\C&D\end{pmatrix}$ with $A\in\pcMat_{\mathsf N\times\mathsf M}$,
$B\in\pcMat_{\mathsf N\times\mathsf T}$, $C\in\pcMat_{\mathsf T\times\mathsf M}$, $D\in\pcMat_{\mathsf T}$:
$\ptvF(p^+)(\rho)=[C=0]\,\ptvF p(-A)\,\psi(\ptvtr D)$.
\end{lemma}
\begin{proof}
$\ptvtr\bigl(\rho\bigl(\begin{smallmatrix}\eta&s\\0&1\end{smallmatrix}\bigr)\bigr)=\ptvtr(A\eta)+
\ptvtr(Cs)+\ptvtr D$; sum over $s$ (giving $\#\pcMat\cdot[C=0]$) and over $\eta$ (giving
$\sum_\eta p(-\eta)\psi(\ptvtr A\eta)=\ptvF p(-A)$).
\end{proof}

\paragraph*{The reversed square law}
Let $\mathsf N,\mathsf T,\mathsf M$ be as above, $\mathsf T'$ a copy of $\mathsf T$, $\mathsf M'=\mathsf M
\sqcup\mathsf T'$. Put $\mathsf U=\mathsf N\sqcup\mathsf M'=\mathsf N\sqcup\mathsf M\sqcup\mathsf T'$,
$\mathsf B=\mathsf M$ (the middle block) and $\mathsf A=\mathsf N\sqcup\mathsf T'$, so that
$\mathsf U=\mathsf B\sqcup\mathsf A$ and $|\mathsf A|=|\mathsf B|$; we identify $\mathsf A$ with
$\mathsf M=\mathsf N\sqcup\mathsf T$ through $\mathsf T'\cong\mathsf T$. Let $L$ on
$\pcMat_{\mathsf M\times\mathsf U}(\pcR)$ and $L'$ on $\pcMat_{\mathsf U\times\mathsf M}(\pcR)$ satisfy the
square law of \S\ref{pc:ss-exchange} for $(\mathsf B,\mathsf A)$ with constant $c$, and assume that
$L(kx)=\chi(\det k)L(x)$ for all $k\in\auGL_{\mathsf M}(\pcR)$ and \emph{all} $x$. Let
$\iota\in\pcMat_{\mathsf M\times\mathsf N}$ and $\iota'\in\pcMat_{\mathsf M'\times\mathsf M}$ be the
inclusion matrices; for $y\in\pcMat_{\mathsf N\times\mathsf M}$ write $\binom y{\iota'}\in
\pcMat_{\mathsf U\times\mathsf M}$, and for $x\in\pcMat_{\mathsf M\times\mathsf M'}$ write $(\iota\ x)\in
\pcMat_{\mathsf M\times\mathsf U}$.

\begin{proposition}\label{ptv:fourier}
There is $\gamma$ with $|\gamma|^2=(\#\pcR)^{|\mathsf T|(|\mathsf T|+2|\mathsf N|)}$, independent of
$L,L',c,p$, such that for every $p$ on $\pcMat_{\mathsf M\times\mathsf N}(\pcR)$
\[
\gamma\sum_{y\in\pcMat_{\mathsf N\times\mathsf M}}\ptvF(\aubar p)(y)\,L'\tbinom y{\iota'}=c\sum_{x\in
\pcMat_{\mathsf M\times\mathsf M'}}\ptvF(\aubar{p^+})(x)\,L(\iota\ x).
\]
\end{proposition}
\begin{proof}
\emph{Step 1.} For $q$ on $\pcMat_{\mathsf N\times\mathsf M}$ we show
$\sum_yq(y)L'\binom y{\iota'}=c\sum_xq^+(x)L(\iota\ x)$, where $q^+$ is the extension for the sets
$(\mathsf N,\mathsf M,\mathsf T)$ in the roles of $(\mathsf M,\mathsf N,\mathsf T)$:
$q^+\bigl(\begin{smallmatrix}\eta&s\\0&1\end{smallmatrix}\bigr)=q(-\eta)/\#\pcMat_{\mathsf N\times
\mathsf T}$, $\eta\in\pcMat_{\mathsf N\times\mathsf M}$, $s\in\pcMat_{\mathsf N\times\mathsf T'}$, rows
$\mathsf N\sqcup\mathsf T$ ($=\mathsf M$), columns $\mathsf M\sqcup\mathsf T'$. For $y\in\pcMat_{\mathsf N
\times\mathsf M}$ let $w\in\pcMat_{\mathsf A\times\mathsf B}$ have rows $-y$ (on $\mathsf N$) and $0$ (on
$\mathsf T'$). In the decomposition $\mathsf U=\mathsf B\sqcup\mathsf A$, $\binom y{\iota'}$ has the
block $1$ on the rows $\mathsf B$ and $-w$ on the rows $\mathsf A$, i.e.\ it is $\binom1{-w}$; and
$(w\ 1)\in\pcMat_{\mathsf A\times\mathsf U}$, read with rows $\mathsf M$, is $(\iota\ x_y)$ with
$x_y=\bigl(\begin{smallmatrix}-y&0\\0&1\end{smallmatrix}\bigr)$. By Lemma~\ref{pc:pointlaw}, as
$\#\auGL_{\mathsf A}=\#\auGL_{\mathsf B}$, $L'\binom y{\iota'}=c\,L(\iota\ x_y)$. For
$s\in\pcMat_{\mathsf N\times\mathsf T'}$ the matrix $u_s=\bigl(\begin{smallmatrix}1&s\\0&1\end{smallmatrix}
\bigr)\in\auGL_{\mathsf M}$ (rows and columns $\mathsf N,\mathsf T$) has determinant $1$, $u_s\iota=\iota$
and $u_s x_y=\bigl(\begin{smallmatrix}-y&s\\0&1\end{smallmatrix}\bigr)$; so $L(\iota\ x)$ is the same for
all $x=\bigl(\begin{smallmatrix}-y&s\\0&1\end{smallmatrix}\bigr)$, and summing against $q^+$ gives the
claim.

\emph{Step 2.} Apply Step 1 to $q=\ptvF(\aubar p)$. It remains to show
$\sum_x\ptvF(\aubar{p^+})(x)L(\iota\ x)=\gamma\sum_xq^+(x)L(\iota\ x)$ with $\gamma=G\cdot\#\pcMat_{\mathsf N
\times\mathsf T}(\pcR)$, $G$ the Gauss sum of Lemma~\ref{ptv:gauss} for $\mathsf T$; then $|\gamma|^2=
(\#\pcR)^{|\mathsf T|^2+2|\mathsf N||\mathsf T|}$. By Lemma~\ref{ptv:F-ext} applied to $\aubar p$ (note
$\aubar{p^+}=(\aubar p)^+$; here $x$ plays the role of $\rho$, with rows $\mathsf N\sqcup\mathsf T$ and
columns $\mathsf M\sqcup\mathsf T'$),
\[
\sum_x\ptvF(\aubar{p^+})(x)L(\iota\ x)=\sum_{A,B}q(-A)\sum_D\psi(\ptvtr D)\,K(A,B,D),\qquad
K(A,B,D)=L\Bigl(\iota\ \begin{pmatrix}A&B\\0&D\end{pmatrix}\Bigr).
\]
For $g\in\auGL_{\mathsf T}$, $k=\mathrm{diag}(1,g)\in\auGL_{\mathsf M}$ fixes $\iota$ and multiplies the
last block row by $g$, so $K(A,B,gD)=\chi(\det g)K(A,B,D)$; by Lemma~\ref{ptv:gauss} at $y=1$,
$\sum_D\psi(\ptvtr D)K(A,B,D)=G\,K(A,B,1)$. Hence the sum is $G\sum_{A,B}q(-A)K(A,B,1)=G\,\#\pcMat_{\mathsf
N\times\mathsf T}\sum_xq^+(x)L(\iota\ x)$.
\end{proof}

\subsubsection{The reversed squares and the adjoint law}\label{ptv:ss-RS}

Fix $i+n+2\le m$ and a twist $\vartheta$, with $\pcm=\pcm_\vartheta$, $\pcR=\pcR_\vartheta$,
$\psi=\psi_\pcm$, $\zeta=\zeta_\pcm$ of Proposition~\ref{ptv:rev}, and $\chi=\chi_\vartheta$ the tame
character of $(\nu_{n+1}\vartheta^{n+1})^{-1}$, non-trivial at every prime of $\pcm$
(\S\ref{pc:ss-exchange}). The nodes of column $k$ share an index set; let $\mathsf N$, $\mathsf M=
\mathsf N\sqcup\mathsf T$, $\mathsf M'=\mathsf M\sqcup\mathsf T'$ be those of columns $n$, $n+1$, $n+2$,
the last block being $\mathsf T$, resp.\ $\mathsf T'$; the inclusion matrices $\iota$, $\iota'$ of
\S\ref{ptv:ss-fourier} are the point data of $a_{i+1,n}$ and $a_{i,n+1}$. For a matrix $\xi$ write
$a[\xi]$ for the lift along the same ascending pair as $a$ with the point datum $\delta_\xi$, so
$a_{i+1,n}=a_{i+1,n}[\iota]$, $a_{i,n+1}=a_{i,n+1}[\iota']$.

Let $d_{i,k}=\Theta^{\mathsf t}(P_{i,k},\varepsilon_{i,k})$ ($k=n,n+1$) be the descents of
\S\ref{pc:ss-exchange}(iv): $P_{i,k}$ is the joint harmonic of the descending pair at $(i,k)$,
the characters are those of Proposition~\ref{ptv:chars}, $\varepsilon_{i,0}$ is the datum of
Theorem~\ref{pc:tri}(a) and $\varepsilon_{i,k+1}=\varepsilon_{i,k}^+$. By induction from
$\varepsilon_{i,0}$ (which is $\delta_\iota$ up to a constant: by (W7) in the model (W5) the map
$\mathsf C$ is complex conjugation of functions at the primes of $\pcm$, up to a constant), all
$\varepsilon_{i,k}$ are supported on injective matrices. Let $R_{k}[\varrho]$ be the reversed lift of
Corollary~\ref{ptv:desc-adj} for $d_{i,k}$ with the point datum $\delta_\varrho$ in place of
$\varepsilon_{i,k}^\flat$, and $R_k[f]=\sum_\varrho f(\varrho)R_k[\varrho]$. For $y\in\pcF_{i+1,n}$ and
$z\in\pcF_{i,n+2}$ define, for $\xi=(\xi_1\ \xi_2)\in\pcMat_{\mathsf M\times(\mathsf N\sqcup\mathsf M')}
(\pcR)$ and $\eta=\binom{\eta_1}{\eta_2}\in\pcMat_{(\mathsf N\sqcup\mathsf M')\times\mathsf M}(\pcR)$,
\[
\mathfrak L(\xi)=\bigl\langle z,\,R_{n+1}[\xi_2]\,a_{i+1,n}[\xi_1]\,y\bigr\rangle,\qquad
\mathfrak L'(\eta)=\bigl\langle z,\,a_{i,n+1}[\eta_2]\,R_n[\eta_1]\,y\bigr\rangle .
\]
These are the two composites around the \emph{reversed square} with corners
$U_1=\pcG_{i+1,n}$ (source), $U_2=\pcG_{i,n+2}$ (target), $A=\pcG_{i+1,n+1}$, $A'=\pcG_{i,n+1}$, all
maps being well defined between the spaces $\pcF$ by Theorem~\ref{ptr:edge}(d) and
Corollary~\ref{ptv:desc-adj}.

\begin{enumerate}
\item[(RS)] There is a function $\kappa^{\mathrm{rev}}_{i,n}$ on $m$-fold twists with modulus $1$ up to norm
powers (Definition~\ref{pc:modulus}) such that for every $\vartheta$ and all $y,z$:
$\mathfrak L(kx)=\chi(\det k)\mathfrak L(x)$ for $k\in\auGL_{\mathsf M}(\pcR)$ and all $x$, and
$\mathfrak L,\mathfrak L'$ satisfy the square law of \S\ref{pc:ss-exchange} for
$(\mathsf B,\mathsf A)=(\mathsf M,\mathsf N\sqcup\mathsf T')$ with constant $\kappa^{\mathrm{rev}}_{i,n}
(\vartheta)$.
\end{enumerate}

\noindent\emph{(RS) is proved in Theorem~\ref{ptw:RS}}: the reversed square is an edge square of excess
$e=0$ and weight $(n+1)\Delta$ (so $L_*(\pcm,\nu,0)=1$ and $|c|^2=C$), with $y_1=y$, $y_2=\aubar z$.
The rank and sign conditions are those of the square at $(i,n)$ rearranged; (E1) for the reversed
pairs $(A',U_1)$, $(U_2,A)$, including (H) for $(U_2,A)$, is Proposition~\ref{ptw:revedge} (through
Howe's argument with the groups exchanged, Theorem~\ref{ptw:Hc}, and the preservation of polynomial
vectors by $\mathcal A_\infty$, Lemma~\ref{ptw:poly}); (E2) holds because $\lambda_S=\lambda_T=1$
(Lemma~\ref{ptw:W7}), so the four conditions are those of Proposition~\ref{ptv:chars}.

\begin{theorem}[Theorem~\ref{pc:tri}(c)]\label{ptv:adjoint}
The adjoint law \textup{(T4)} holds for the triangle
of \S\ref{pc:ss-triangle}, with functions $\kappa_{i,n}$ of modulus $1$ up to norm powers.
\end{theorem}
\begin{proof}
We use \textup{(RS)} for all $i+n+2\le m$, which is Theorem~\ref{ptw:RS}.
\emph{The adjoint of $a_{i,n+1}$.} By Lemma~\ref{ptv:adj} for the ascending pair at $(i,n+1)$ (the
point datum $\delta_{\iota'}$ is supported on an injective matrix), every $z\in\pcF_{i,n+2}$ has
$x=a_{i,n+1}^*z\in\pcF_{i,n+1}$ with $\langle a_{i,n+1}v,z\rangle=\langle v,x\rangle$ for all $v$.

\emph{Reduction.} For $v'\in\pcF_{i+1,n}$, with $d^*$ as in Corollary~\ref{ptv:desc-adj},
$\langle d_{i,n}x,v'\rangle=\langle x,d^*_{i,n}v'\rangle=\langle z,a_{i,n+1}d^*_{i,n}v'\rangle$ and
$\langle d_{i,n+1}z,a_{i+1,n}v'\rangle=\langle z,d^*_{i,n+1}a_{i+1,n}v'\rangle$. So (T4) with constant
$\kappa_{i,n}$ follows from
\begin{equation}\label{ptv:eq-adj}
a_{i,n+1}\,d^*_{i,n}=\kappa_{i,n}(\vartheta)\,d^*_{i,n+1}\,a_{i+1,n}\qquad\text{on }\pcF_{i+1,n}.
\end{equation}

\emph{The reversed commutation.} Let $p=\varepsilon_{i,n}$, a function on $\pcMat_{\mathsf M\times
\mathsf N}(\pcR)$, so $\varepsilon_{i,n+1}=p^+$, $\varepsilon_{i,n}^\flat=\ptvF(\aubar p)$ and
$\varepsilon_{i,n+1}^\flat=\ptvF(\aubar{p^+})$ (Proposition~\ref{ptv:rev}, with $\psi=\psi_\pcm$). By
linearity of lifts in the datum,
\[
\langle z,a_{i,n+1}R_n[\ptvF\aubar p]y\rangle=\sum_{y'}\ptvF(\aubar p)(y')\,\mathfrak L'\tbinom{y'}{\iota'},
\qquad\langle z,R_{n+1}[\ptvF\aubar{p^+}]a_{i+1,n}y\rangle=\sum_x\ptvF(\aubar{p^+})(x)\,\mathfrak L(\iota\ x).
\]
By (RS) and Proposition~\ref{ptv:fourier}, $\gamma\langle z,a_{i,n+1}R_n[\ptvF\aubar p]y\rangle=
\kappa^{\mathrm{rev}}_{i,n}\langle z,R_{n+1}[\ptvF\aubar{p^+}]a_{i+1,n}y\rangle$ for all $z\in\pcF_{i,n+2}$.
Both vectors lie in $\pcF_{i,n+2}$, on which the Petersson product is definite; so
\[
a_{i,n+1}R_n=\frac{\kappa^{\mathrm{rev}}_{i,n}}{\gamma}\,R_{n+1}\,a_{i+1,n}\qquad\text{on }\pcF_{i+1,n},
\]
$R_k$ now denoting the reversed lift with the datum $\varepsilon_{i,k}^\flat$.

\emph{Conclusion.} By Corollary~\ref{ptv:desc-adj}, $d^*_{i,k}=c_k\,R_k\,\psrho(\zeta^{-1})$, with the same
central $\zeta$ for $k=n,n+1$ (it does not depend on the spaces). By Lemma~\ref{ptv:central},
$a_{i+1,n}\psrho(\zeta^{-1})=\eta\,\psrho(\zeta^{-1})a_{i+1,n}$ with $|\eta|=1$. Hence
\[
a_{i,n+1}d^*_{i,n}=c_n\frac{\kappa^{\mathrm{rev}}_{i,n}}{\gamma}R_{n+1}a_{i+1,n}\psrho(\zeta^{-1})=
\frac{c_n\,\eta\,\kappa^{\mathrm{rev}}_{i,n}}{\gamma\,c_{n+1}}\,d^*_{i,n+1}a_{i+1,n},
\]
which is \eqref{ptv:eq-adj} with $\kappa_{i,n}=c_n\eta\kappa^{\mathrm{rev}}_{i,n}/(\gamma c_{n+1})$. By
Proposition~\ref{ptv:rev}, $|c_n|^2=C\,(\#\pcR_\vartheta)^{-M_nM_{n+1}}$ and $|c_{n+1}|^2=C'(\#\pcR_\vartheta)^{
-M_{n+1}M_{n+2}}$; by Proposition~\ref{ptv:fourier}, $|\gamma|^2=(\#\pcR_\vartheta)^{t(t+2M_n)}$; and
$\#\pcR_\vartheta=M^f\auNm Q_\vartheta$ (\S\ref{pc:ss-data}). So $|\kappa_{i,n}|^2$ is a constant times
$|\kappa^{\mathrm{rev}}_{i,n}|^2$ times an integral power of $\auNm Q_\vartheta$, and $\kappa_{i,n}$ has
modulus $1$ up to norm powers by Lemma~\ref{pc:mod}(a).
\end{proof}

\subsubsection{Final check of Lemma~\ref{pc:arch} and Theorem~\ref{pc:tri}}\label{ptv:ss-check}

\begin{proposition}[Lemma~\ref{pc:arch}]\label{ptv:arch}
Lemma~\ref{pc:arch} holds for all its pairs, modulo the inputs of Theorems~\ref{ptr:edge},
\ref{ptq:F} and {\cite[Proposition~4.8]{P54}} \textup{((V), (D1)--(D3))}.
\end{proposition}
\begin{proof}
\emph{The pairs $(\pcG_{i,0},\pcG_{i,1})$.} (A1) and (A2) are Theorem~\ref{ptr:edge}(c), (d) (the last
paragraph of its statement). (A3): Proposition~\ref{ptr:Z} with Theorem~\ref{ptq:F}
(Remark~\ref{ptq:rem-F}) gives $Z_\infty=\prod_\tau d_{\lambda_\tau}^{-1}\langle\Phi_\tau,\mathsf P_\tau
\Phi_\tau\rangle$ with every factor positive, and the integral converges absolutely because
$|\langle P,\omega_\infty(g,1)P\rangle|\le\|P\|^2$ and the coefficient is integrable
(Lemma~\ref{ptr:gap}); it involves only archimedean data, at which the twists are trivial.

\emph{The pairs $(W_j,V_j)$.} (A1): for $W_j$ this is {\cite[Lemma~4.3]{P54}}; $[\auU(V_j)]$ is compact
(the group $\auU(V_j)_\infty$ is compact), so the type $e_{V_j}$ is integrable and its forms are bounded
on a compact quotient, which is the decay statement there. (A3): write $\sigma=\omega_\infty(\cdot,1)$
and $\mathsf P=\int e_{W_j}(x)\sigma(x^{-1})\,dx$. As in the first paragraph of the proof of
Proposition~\ref{ptr:Z} (with $\lambda=-\varsigma_\tau$, $\langle v_0,\pi_{W_j}(g)v_0\rangle=
\aubar{c_\lambda(g)}$, proof of {\cite[Proposition~4.8]{P54}}, Step 1),
$Z_\infty=\prod_\tau d_{\lambda_\tau}^{-1}\langle\Phi_\tau,\sigma_\tau(e_{\lambda_\tau})\Phi_\tau\rangle$, the
integral being absolutely convergent by {\cite[Lemma~4.3]{P54}}. Step 1 of the proof of
{\cite[Proposition~4.8]{P54}} shows $\sigma_\tau(e_{\lambda_\tau})\Phi_\tau=\Phi_\tau$ for the joint harmonic
$\Phi_\tau\ne0$. Hence $Z_\infty=\prod_\tau d_{\lambda_\tau}^{-1}\|\Phi_\tau\|^2>0$, independent of the
twist. (A2): the lift $h\mapsto\int\theta(g,h;P,\Psi)\varphi(g)\,dg$ converges absolutely and defines a
bounded function right invariant under an open subgroup, by the argument of the proof of
{\cite[Proposition~8.4(b)]{P49}} ({\cite[Lemma~8.2]{P49}} and rapid decay, {\cite[Lemma~4.3]{P54}}), which uses
the row side only through $\mathsf PP=P$ (Step 1 of the proof of {\cite[Proposition~4.8]{P54}}). It is
reproduced by $e_{V_j}$: writing the lift as $L_h(P)$ with $L_h$ continuous on $\pzF_s$ and
$L_h(\omega_\infty(1,x)\Psi)=L_{hx}(\Psi)$ (proof of Lemma~\ref{ptr:conv-lift}), right convolution by
$e_{V_j}$ is $L_h(\int e_{V_j}(x)\omega_\infty(1,x)P\,dx)$, and $\int e_{V_j}(x)\omega_\infty(1,x)P\,dx=P$
by Schur orthogonality on the compact group $\auU(V_j)_\infty$, because $P$ is a highest weight vector
of the irreducible representation of highest weight $\Lambda'=\lambda'-\rho'$ (Step 1 of the proof of
{\cite[Proposition~4.8]{P54}}) and $e_{V_j}$ is, by definition, the type of that representation
(Proposition~\ref{ptw:A2} and Remark~\ref{ptw:rem-A2}).
\end{proof}

\begin{proposition}[The clauses of Theorem~\ref{pc:tri}]\label{ptv:tri}
\textup{(i)} The maps $a^\vartheta_{i,n}$ are well defined between the spaces $\pcF_{i,n}$
(Theorem~\ref{ptr:edge}(d)) and have adjoints (Lemma~\ref{ptv:adj}).
\textup{(ii)} With $d_{i,0}$ the transposed lift along $(\pcG_{i+1,0},\pcG_{i,1})$ whose kernel is
the complex conjugate of that of $a_{i,0}$ (it is a kernel for this pair by \textup{(W7)}, with the
characters of Proposition~\ref{ptv:chars}(b)), $d_{i,0}=a_{i,0}^*$ under $\pcF_{i+1,0}=\pcF_{i,0}$: this
is Lemma~\ref{ptv:adj} for $a_{i,0}$, because $\pcG_{i+1,0}=\pcG_{i,0}^-$ has the same group, quotient
and measure. This is \textup{(a)}.
\textup{(iii)} \textup{(b)} is Corollary~\ref{ptr:tri-b} (with Theorem~\ref{ptq:H} for (H) and
Proposition~\ref{ptv:chars} for the finite parts of the characters), modulo the published results
(D1)--(D5) and Theorem~\ref{pcx:X}.
\textup{(iv)} \textup{(c)} is Theorem~\ref{ptv:adjoint} (which uses \textup{(RS)}, Theorem~\ref{ptw:RS}).
\end{proposition}
\begin{proof}
Each item is the cited statement; for (ii), $a_{i,0}^*z(g)=\int\aubar{\theta(g,h)}z(h)\,dh$ and by
(W7) $\aubar{\theta^{\pcG_{i,0},\pcG_{i,1}}_{\alpha,\beta}(g,h;F)}=\theta^{\pcG_{i,0}^-,\pcG_{i,1}}
_{\alpha',\beta'}(g,h;\mathsf CF)$ with $\alpha',\beta'$ the conjugate characters times the powers of
$\upsilon$ of (W7), and $\mathsf C[P,\delta_\iota]$ is a datum $[P^-,c\,\delta_\iota]$, $|c|=1$ (at the
primes of the modulus, in the model (W5), $\mathsf C$ is complex conjugation up to a constant, since
$[x,a']^-=\aubar{[x,a']}$; elsewhere it maps the vacuum to a multiple of the vacuum, (W4)).
\end{proof}




% ===== C-triangle4.tex =====
% proofs/C-triangle4.tex
% Continuation of proofs/C-triangle.tex (ptr:), C-triangle2.tex (ptq:), C-triangle3.tex (ptv:):
% Howe's argument with rows and columns exchanged, the smoothing lemma (H) for the reversed pairs,
% the characters of the reversed kernels, and the reversed square law (RS) of \S\ref{ptv:ss-RS};
% then the edge hypothesis (E1) for the grid squares and the check of (A2) for the pairs (W_j,V_j).
% Labels of this file start with ptw:. Only amsmath/amssymb/amsthm are used.
% Conventions and macros: proofs/C-triangle.tex, C-triangle3.tex.
\providecommand{\auU}{\mathrm U}
\providecommand{\auHom}{\operatorname{Hom}}
\providecommand{\auGL}{\mathrm{GL}}
\providecommand{\auNm}{\operatorname{Nm}}
\providecommand{\auA}{\mathbb A}
\providecommand{\auO}{\mathcal O}
\providecommand{\aubar}[1]{\overline{#1}}
\providecommand{\psrho}{\varrho}
\providecommand{\psHS}{\mathcal H}
\providecommand{\pzF}{\mathfrak F}
\providecommand{\pcF}{\mathcal F}
\providecommand{\pcG}{\mathsf G}
\providecommand{\pcm}{\mathfrak m}
\providecommand{\pcR}{R}
\providecommand{\pcMat}{\mathrm M}
\providecommand{\pop}{\mathsf p}
\providecommand{\ptrwt}{\operatorname{wt}}
\providecommand{\ptwsgn}{\operatorname{sgn}}

\subsection{The reversed square law, the grid edges, and the pairs $(W_j,V_j)$}\label{ptw:sec}

This section continues \S\ref{ptr:sec}, \S\ref{ptq:sec} and \S\ref{ptv:sec}, whose notation it
uses; inner products are conjugate-linear in the first variable.

\emph{Contents of this section.}
\begin{enumerate}
\item[(1)] Howe's argument and the uniqueness of the joint harmonic hold with the roles of rows and
columns exchanged (Propositions~\ref{ptw:unique} and~\ref{ptw:howe}); the only point that does not
transpose word for word, the string $B$ of columns that are not matched, is handled by the
observation that a column of $B$ carries no variable of any polynomial that is highest of the
lowest weight for the columns and harmonic for the rows (Lemma~\ref{ptw:B}). Consequently the
smoothing statement holds with the groups exchanged, for every ascending and descending pair
(Theorem~\ref{ptw:Hc}).
\item[(2)] At the infinite places the reversal map $\mathcal A_\infty$ of \S\ref{ptv:ss-rev} maps
polynomial vectors onto polynomial vectors (Lemma~\ref{ptw:poly}), and the characters $\lambda_S$,
$\lambda_T$ of \eqref{ptv:eq-A} are trivial; equivalently, the character of (W7) is
$\upsilon(\det g)^{-N'}\upsilon(\det h)^{N}$ (Lemma~\ref{ptw:W7}).
\item[(3)] Hence the reversed pairs are edges in the sense used by Theorem~\ref{ptr:square},
including (H) (Proposition~\ref{ptw:revedge}), the reversed squares satisfy (E2), and the reversed
square law (RS) of \S\ref{ptv:ss-RS} holds (Theorem~\ref{ptw:RS}). With Theorem~\ref{ptv:adjoint}
this proves the adjoint law, Theorem~\ref{pc:tri}(c), modulo the same published results
(D1)--(D5) as Theorem~\ref{pc:tri}(b).
\end{enumerate}
Finally, \S\ref{ptw:ss-A2} identifies the type on $V_j$ used in the check of (A2) for the pairs
$(W_j,V_j)$ in Proposition~\ref{ptv:arch}, and \S\ref{ptw:ss-grid} records where the edge hypothesis
(E1) for the grid squares is proved.

\subsubsection{Rows and columns exchanged}\label{ptw:ss-sym}

We work at one place $\tau$, in the polynomial model of {\cite[\S3.1.4]{P49}}, for a matching pair
$(\lambda,\lambda')$ with distinct entries in the sense of {\cite[(13)]{P49}}: rows $R$ ($n$ of
them), columns $C$ ($n'=n+t$), every row $r$ matched with a column $c(r)$, and the string
$B=C\smallsetminus c(R)$ of $t$ columns with $\{\xi'_c:c\in B\}=\{\tfrac{t-1}2-b:0\le b<t\}$ and
$\varepsilon'$ constant on $B$. (By Proposition~\ref{ptr:match} and Lemma~\ref{ptr:gap} this is the
situation at every $\tau$ for every ascending and every descending pair of the triangle, the
derived actions being {\cite[(12)]{P49}} by Lemma~\ref{ptr:const}.) Notation is that of
{\cite[\S3.1.5]{P49}}: $\Lambda=\Lambda^K(\lambda)$, $\Lambda'=\Lambda^K(\lambda')$, $v_a$, $v'_c$ the
weights of the vacuum, $w_r=\Lambda_r-v_r$, $w'_c=\Lambda'_c-v'_c$.

\paragraph*{The column operators in the form of the row operators}
Put $\eta'_{rc}=-\eta_{rc}=\varepsilon_r\varepsilon'_c$. For $\varepsilon'_c=\varepsilon'_d$ we have
$\eta_{rc}=\eta_{rd}$ for all $r$, and {\cite[(12)]{P49}} gives, using $\partial_{rc}z_{rd}=
z_{rd}\partial_{rc}+\delta_{cd}$,
\[
E'_{cd}=-\sum_{r:\eta_{rc}=1}\partial_{rc}z_{rd}+\sum_{r:\eta_{rc}=-1}z_{rc}\partial_{rd}+\gamma'\delta_{cd}
=D'^{+}_{c\leftarrow d}-D'^{-}_{d\leftarrow c}+\delta_{cd}v'_c,\qquad
D'^{s}_{c\leftarrow d}=\sum_{r:\eta'_{rd}=s}z_{rc}\partial_{rd},
\]
with $v'_c=\gamma'-\#\{r:\eta'_{rc}=-1\}$. For $\varepsilon_a=\varepsilon_b$, the proof of
{\cite[Proposition~3.11]{P49}} writes $E_{ab}=D^+_{a\leftarrow b}-D^-_{b\leftarrow a}+\delta_{ab}v_a$,
$D^s_{a\leftarrow b}=\sum_{c:\eta_{bc}=s}z_{ac}\partial_{bc}$, $v_a=\gamma-\#\{c:\eta_{ac}=-1\}$. For
$\varepsilon'_c\ne\varepsilon'_d$ the degree-lowering part of $E'_{cd}$ is
$-\sum_{r:\eta'_{rc}=-1}\partial_{rc}\partial_{rd}$, so up to sign
$\Delta'_{cd}=\sum_{r:\eta'_{rc}=-1}\partial_{rc}\partial_{rd}$; and for rows $a,b$ and a column $c$ with
$\eta'_{ac}=1=-\eta'_{bc}$ (so $\varepsilon_a\ne\varepsilon_b$),
$\sum_{d:\varepsilon'_d=\varepsilon'_c}\partial_{ad}\partial_{bd}=\Delta_{ab}$, because $\eta_{ad}=-1$ means
$\varepsilon'_d=\varepsilon_a=\varepsilon'_c$. Thus the substitution
\[
(R,C,\varepsilon,\varepsilon',\eta,\lambda,\lambda',E,\Delta,D,v)\longmapsto
(C,R,\varepsilon',\varepsilon,\eta',\lambda',\lambda,E',\Delta',D',v')
\]
carries the formulas for the row operators into those for the column operators and conversely,
and the Fischer product is unchanged. {\cite[Lemma~3.12]{P49}}, Lemma~\ref{ptq:fischer} and
Lemma~\ref{ptq:tensor} are statements about these formulas and the Fischer product only (their
proofs use nothing else: in {\cite[Lemma~3.12]{P49}} the identities $E_{ab}=D^+-D^-+\delta v$, the
adjointness of $z_i$ and $\partial_i$, and the vanishing of $\sum_b\partial_{bc}\partial_{bd}$ through
$\Delta'_{cd}$; in Lemma~\ref{ptq:tensor} the commutation relations of the $E$ with the $\Delta$). So
they hold with rows and columns exchanged. We record the column form of {\cite[Lemma~3.12]{P49}}.

\begin{lemma}[Column polarisations]\label{ptw:polar}
Let $P$ be a polynomial killed by the column raising operators $E'_{cd}$ ($\varepsilon'_c=\varepsilon'_d$,
$\lambda'_c>\lambda'_d$) and by the row Laplacians $\Delta_{ab}$ ($\varepsilon_a\ne\varepsilon_b$). Let
$\varepsilon'_c=\varepsilon'_d$ and $s=\pm1$.
\textup{(b)} If $s\lambda'_c>s\lambda'_d$, then $D'^s_{c\leftarrow d}P=0$.
\textup{(c)} If $E'_{cc}P=\Lambda'_cP$ and $k=s\,w'_c\ge0$, then $D'^s_{c\leftarrow c}P=kP$ and
$\partial_{rc}P=0$ for every $r$ with $\eta'_{rc}=-s$.
\end{lemma}
\begin{proof}
{\cite[Lemma~3.12(b)]{P49}}, (c) after the substitution above.
\end{proof}

\begin{lemma}[Column weights]\label{ptw:relw}
For a matched column $c=c(r)$, $w'_c\,\ptwsgn\xi'_c\ge0$; for $c\in B$, $w'_c=0$.
\end{lemma}
\begin{proof}
By {\cite[Lemma~3.10(ii)]{P49}}, $w'_{c(r)}=-w_r$, and $\xi'_{c(r)}=-\xi_r$ by {\cite[(13)]{P49}}; so
$w'_{c(r)}\ptwsgn\xi'_{c(r)}=w_r\ptwsgn\xi_r\ge0$ by {\cite[Lemma~3.10(i)]{P49}}. The second assertion is
{\cite[Lemma~3.10(ii)]{P49}}.
\end{proof}

\begin{lemma}[The string carries no variable]\label{ptw:B}
Let $P$ be killed by the column raising operators and by the row Laplacians, with
$E'_{cc}P=\Lambda'_cP$ for a column $c\in B$. Then $\partial_{rc}P=0$ for every row $r$: no monomial of
$P$ contains a variable of the column $c$.
\end{lemma}
\begin{proof}
By Lemma~\ref{ptw:relw}, $w'_c=0$, so $k=s\,w'_c=0\ge0$ for both $s=1$ and $s=-1$.
Lemma~\ref{ptw:polar}(c) with $s=1$ gives $\partial_{rc}P=0$ when $\eta'_{rc}=-1$, and with $s=-1$ when
$\eta'_{rc}=1$.
\end{proof}

The uniqueness of {\cite[Proposition~3.13]{P49}} omits the column weight condition. Howe's argument
with the groups exchanged produces polynomials for which the \emph{row} weight is not known, so we
need the uniqueness with that condition omitted instead. Its proof is that of
{\cite[Proposition~3.13]{P49}} with rows and columns exchanged, except that the columns of $B$ have
to be treated separately.

\begin{proposition}[Uniqueness, rows and columns exchanged]\label{ptw:unique}
Call a polynomial $F\ne0$ a \emph{column pre-joint harmonic} if it satisfies all conditions of
{\cite[Definition~3.9]{P49}} except possibly $E_{aa}F=\Lambda_aF$. Two column pre-joint harmonics
are proportional. In particular every column pre-joint harmonic is a multiple of the joint harmonic.
\end{proposition}
\begin{proof}
Let $x_0$ be the matrix with $z_{rc}(x_0)=1$ if $c=c(r)$ and $0$ otherwise, as in the proof of
{\cite[Proposition~3.13]{P49}}. As there, it suffices to show that a column pre-joint harmonic $F$ has
$F(x_0)\ne0$: for two of them $P,Q$, the polynomial $P-\frac{P(x_0)}{Q(x_0)}Q$ satisfies the same
linear conditions and vanishes at $x_0$, so it is $0$.

Suppose $F(x_0)=0$. Write $F(x_0+u)=\sum_mc_mu^m$, choose a monomial $u^m$ of least positive degree
$\delta$ with $c_m\ne0$ and a variable $i$ occurring in it; then $\partial_iF\ne0$. Let $V=\sum_k
z_{a_k}\partial_{b_k}$ be a vector field such that exactly one term, $z_{a_1}\partial_{b_1}$ with
$b_1=i$, has $z_{a_1}(x_0)\ne0$. In $VF$ the term $z_{a_1}(x_0)\partial_iF$ contributes a non-zero
coefficient of $u^{m-e_i}$, of degree $\delta-1$, and every other term $z_{a_k}\partial_{b_k}F$ has
only terms of degree $\ge\delta$ ($z_{a_k}$ vanishes at $x_0$ and $\partial_{b_k}F$ has no terms of
degree $<\delta-1$). As $F$ has no term of degree $\le\delta-1$, $VF$ is not a multiple of $F$. We
reach a contradiction by producing such a $V$ with $VF\in\C F$.

Let $i=(r,c)$ and $s=\eta'_{rc}$.

\emph{$c\in B$.} By Lemma~\ref{ptw:B} (applicable: $F$ is killed by the column raising operators and
the row Laplacians and $E'_{cc}F=\Lambda'_cF$), $\partial_{rc}F=0$, a contradiction.

\emph{$c=c(r_1)$ is matched.} Then $|\xi'_c|>\tfrac t2$. If $-s\,\xi'_c>\tfrac t2$, then
$\ptwsgn\xi'_c=-s$, so $(-s)w'_c\ge0$ by Lemma~\ref{ptw:relw}, and Lemma~\ref{ptw:polar}(c) with $-s$
gives $\partial_{r'c}F=0$ whenever $\eta'_{r'c}=s$, in particular $\partial_iF=0$: a contradiction. So
$s\,\xi'_c>\tfrac t2$. Then $\ptwsgn\xi_{r_1}=-\ptwsgn\xi'_c=-s$, so by {\cite[(13)]{P49}}
$\eta_{r_1c}=-s$, i.e.\ $\eta'_{r_1c}=s=\eta'_{rc}$, and hence $\varepsilon_{r_1}=\varepsilon_r$. Three
cases.

($\alpha$) $r=r_1$. By Lemma~\ref{ptw:polar}(c) with $s$ ($s\,w'_c\ge0$ by
Lemma~\ref{ptw:relw}), $F$ is an eigenvector of $D'^s_{c\leftarrow c}=\sum_{r':\eta'_{r'c}=s}
z_{r'c}\partial_{r'c}$. At $x_0$, $z_{r'c}(x_0)\ne0$ only for $c=c(r')$, i.e.\ $r'=r_1=r$; so the only
term not vanishing at $x_0$ is $z_{rc}\partial_{rc}=z_{rc}\partial_i$. Contradiction.

($\beta$) $r\ne r_1$ and $s\lambda_{r_1}<s\lambda_r$ (the entries of $\lambda$ are distinct). Then
$(-s)\lambda_{r_1}>(-s)\lambda_r$ and $\varepsilon_{r_1}=\varepsilon_r$, so by {\cite[Lemma~3.12(b)]{P49}} (the
row raising operators and the column Laplacians kill $F$) the field
$D^{-s}_{r_1\leftarrow r}=\sum_{d:\eta_{rd}=-s}z_{r_1d}\partial_{rd}$ kills $F$. Here $z_{r_1d}(x_0)\ne0$
only for $d=c(r_1)=c$, and the term $d=c$ occurs because $\eta_{rc}=-\eta'_{rc}=-s$; it is
$z_{r_1c}\partial_i$. Contradiction.

($\gamma$) $r\ne r_1$ and $s\lambda_{r_1}>s\lambda_r$. Since $\xi-\lambda$ is constant on the rows,
$s\xi_r<s\xi_{r_1}=-s\xi'_c<-\tfrac t2$; so $\ptwsgn\xi_r=-s$, and with $c_2=c(r)\ne c$,
{\cite[(13)]{P49}} gives $\eta_{rc_2}=-s$, i.e.\ $\eta'_{rc_2}=s=\eta'_{rc}$, so
$\varepsilon'_{c_2}=\varepsilon'_c$. Moreover $s\xi'_{c_2}=-s\xi_r>-s\xi_{r_1}=s\xi'_c$, so
$s\lambda'_{c_2}>s\lambda'_c$. By Lemma~\ref{ptw:polar}(b), $D'^s_{c_2\leftarrow c}=\sum_{r':\eta'_{r'c}=s}
z_{r'c_2}\partial_{r'c}$ kills $F$. Here $z_{r'c_2}(x_0)\ne0$ only for $c_2=c(r')$, i.e.\ $r'=r$, and
that term occurs ($\eta'_{rc}=s$); it is $z_{rc_2}\partial_i$. Contradiction.

So $F(x_0)\ne0$. The joint harmonic of {\cite[Proposition~3.11]{P49}} is a column pre-joint harmonic,
whence the last assertion.
\end{proof}

\begin{proposition}[Howe's argument, groups exchanged]\label{ptw:howe}
Let $\Phi$ be the joint harmonic of $(\lambda,\lambda')$ and $\mathcal G$ a space of polynomials such
that: \textup{(i)} $E_{ab}\mathcal G\subset\mathcal G$ for all rows $a,b$; \textup{(ii)}
$\Phi\in\mathcal G$; \textup{(iii)} $E'(X)Q\in\mathcal G$ for every $X\in\mathfrak{gl}_{n'}(\C)$ and
every polynomial $Q$ with $E'_{cc}Q=\nu'_cQ$ for all $c$ and $E'_{cd}Q=0$ whenever
$\varepsilon'_c=\varepsilon'_d$, $\lambda'_c>\lambda'_d$, where $\nu'-\Lambda'$ is not a non-negative
combination of roots $e_c-e_d$ with $\lambda'_c>\lambda'_d$; here $E'(X)$ is the column action;
\textup{(iv)} $E'_{dc}Q\in\mathcal G$ for every polynomial $Q$, when $\varepsilon'_c=\varepsilon'_d$,
$\lambda'_c>\lambda'_d$; \textup{(v)} $(E'_{cc}-\Lambda'_c)Q\in\mathcal G$ for every polynomial $Q$
and every $c$. Then $\mathcal G$ contains all polynomials.
\end{proposition}
\begin{proof}
We show $\mathcal D_j\subset\mathcal G$ by induction on $j$, $\mathcal D_{-1}=0$. Suppose
$\mathcal D_{j-1}\subset\mathcal G$, $\mathcal D_j\not\subset\mathcal G$, and let $\mathcal N\ne0$ be the
orthogonal complement of $\mathcal G\cap\mathcal D_j$ in $\mathcal D_j$ for the Fischer product.

\emph{A highest vector for the rows.} The row operators preserve degree when
$\varepsilon_a=\varepsilon_b$, and $E_{ab}^*=E_{ba}$ for the Fischer product (\S\ref{ptq:ss-H}). If
$\varepsilon_a=\varepsilon_b$, $E_{ba}$ preserves $\mathcal G\cap\mathcal D_j$ by (i), so $E_{ab}$ preserves
$\mathcal N$. The $E_{aa}$ are self-adjoint and commute; so $\mathcal N$ contains $Q\ne0$ with
$E_{aa}Q=\nu_aQ$ for all $a$ (some $\nu$) and $E_{ab}Q=0$ whenever $\varepsilon_a=\varepsilon_b$,
$\lambda_a>\lambda_b$ (a weight vector of maximal weight for a linear order compatible with these
roots).

\emph{Column weight and column raising.} By (v), $\mathcal G$ is $E'_{cc}$-stable, hence so is
$\mathcal N$ ($E'_{cc}$ is self-adjoint and preserves degree), and $(E'_{cc}-\Lambda'_c)Q\in\mathcal N\cap
\mathcal G=0$. If $\varepsilon'_c=\varepsilon'_d$ and $\lambda'_c>\lambda'_d$, then by (iv)
$E'_{dc}\mathcal D_j\subset\mathcal G\cap\mathcal D_j$, so $E'_{cd}Q\in\mathcal D_j$ is orthogonal to
$\mathcal D_j$, and $E'_{cd}Q=0$.

\emph{Row Laplacians.} For $\varepsilon_a\ne\varepsilon_b$, $\Delta_{ab}Q\in\mathcal D_{j-2}\subset\mathcal G$,
so $E_{ba}\Delta_{ab}Q\in\mathcal G$ by (i), and $\Delta_{ab}Q=0$ by Lemma~\ref{ptq:fischer}.

\emph{Column Laplacians.} We apply Lemma~\ref{ptq:tensor} for the columns (valid by
\S\ref{ptw:ss-sym}): let $\varepsilon'_c\ne\varepsilon'_d$ and suppose that $\Delta'_{cd}Q$ has column
weight $\Lambda'+e_c-e_d$ and is killed by the column raising operators. If $e_c-e_d$ is not a
non-negative combination of roots $e_x-e_z$ with $\lambda'_x>\lambda'_z$, then (iii) gives
$E'_{dc}\Delta'_{cd}Q\in\mathcal G$, and $\Delta'_{cd}Q=0$ by Lemma~\ref{ptq:fischer} for columns.
Otherwise, pairing with $\lambda'$ gives $\lambda'_c>\lambda'_d$. Now $Q$ is killed by the column raising
operators and by the row Laplacians, and $E'_{xx}Q=\Lambda'_xQ$ for all $x$. If $c\in B$ or $d\in B$,
Lemma~\ref{ptw:B} shows that $Q$ contains no variable of that column, so every term
$\partial_{rc}\partial_{rd}Q$ of $\Delta'_{cd}Q$ vanishes. Otherwise $c,d$ are matched, $|\xi'_c|,|\xi'_d|>
\tfrac t2$, and $\xi'_c>\xi'_d$ ($\xi'-\lambda'$ is constant); so either $\xi'_c>\tfrac t2$, or
$\xi'_d<\xi'_c<-\tfrac t2$. In the first case $w'_c\ge0$ (Lemma~\ref{ptw:relw}) and
Lemma~\ref{ptw:polar}(c) with $s=1$ gives $\partial_{rc}Q=0$ whenever $\eta'_{rc}=-1$; in the second
$-w'_d\ge0$ and Lemma~\ref{ptw:polar}(c) with $s=-1$ gives $\partial_{rd}Q=0$ whenever $\eta'_{rd}=1$.
Every term of $\Delta'_{cd}=\sum_{r:\eta'_{rc}=-1}\partial_{rc}\partial_{rd}$ has $\eta'_{rc}=-1$ and
$\eta'_{rd}=1$ ($\varepsilon'_c\ne\varepsilon'_d$), so $\Delta'_{cd}Q=0$. By Lemma~\ref{ptq:tensor} for
the columns, all column Laplacians kill $Q$.

\emph{Conclusion.} $Q$ is a column pre-joint harmonic, so $Q=c\Phi$ by
Proposition~\ref{ptw:unique}. By (ii), $Q\in\mathcal G\cap\mathcal N=0$, a contradiction.
\end{proof}

\begin{remark}\label{ptw:rem-howe}
In the proof of Proposition~\ref{ptq:howe} the step ``row Laplacians'' uses that every row is
matched; in the exchanged argument the corresponding step meets the columns of $B$, and
Lemma~\ref{ptw:B} disposes of them. What Howe's argument really needs from the parameters, in either
direction, is: for two indices $x,y$ of the projected group with opposite signs and $\lambda_x>\lambda_y$
(resp.\ $\lambda'_x>\lambda'_y$), either $x$ carries no variable of one of the two kinds or $y$ carries
no variable of the other kind. For matched indices this follows from {\cite[Lemma~3.12(c)]{P49}} and
the sign of $w$; for the unmatched string it is the statement $w'=0$ of {\cite[Lemma~3.10(ii)]{P49}},
through Lemma~\ref{ptw:B}.
\end{remark}

\subsubsection{The smoothing lemma with the groups exchanged}\label{ptw:ss-H}

Let $(S,T)$ be an ascending or descending pair of the triangle, $\tau$ fixed, $G=\auU(S)_\tau$,
$G'=\auU(T)_\tau$, $\sigma=\omega_\infty(\cdot,1)$, $\sigma'=\omega_\infty(1,\cdot)$ at $\tau$, $\Phi$ the
joint harmonic, and $\mathsf P=\sigma(e_\lambda)$, $\mathsf P'=\sigma'(e_{\lambda'})$ the two projections
(proof of {\cite[Theorem~3.15]{P49}}, valid for the pair by Theorem~\ref{ptr:edge}(b)). Let
$H_{\lambda'}\subset\pzF_\tau$ be the closed span of the images of the $G'$-equivariant isometries
$\pi_{\lambda'}\to\sigma'$. For $\beta\in C_c^\infty(G)$ write $\sigma(\beta)=\int\beta(x)\sigma(x)\,dx$.

\begin{lemma}\label{ptw:P}
\textup{(a)} For a polynomial $Q$, $\mathsf P'(E'_{dc}Q)=0$ if $\varepsilon'_c=\varepsilon'_d$ and
$\lambda'_c>\lambda'_d$, and $\mathsf P'(E'_{cc}Q)=\Lambda'_c\mathsf P'Q$.
\textup{(b)} $\mathsf P'$ is self-adjoint and $\mathsf P'x\in H_{\lambda'}$ for every $x$.
\textup{(c)} If $Q$ is a polynomial with $E'_{cc}Q=\nu'_cQ$ for all $c$ and $E'_{cd}Q=0$ whenever
$\varepsilon'_c=\varepsilon'_d$, $\lambda'_c>\lambda'_d$, where $\nu'-\Lambda'$ is not a non-negative
combination of roots $e_c-e_d$ with $\lambda'_c>\lambda'_d$, then $\mathsf P'(E'(X)Q)=0$ for every
$X\in\mathfrak{gl}_{n'}(\C)$.
\textup{(d)} There is $\beta_0\in C_c^\infty(G)$ with $\sigma(\beta_0)\mathsf P'\Phi=\mathsf P'\Phi$.
\end{lemma}
\begin{proof}
(a)--(c) are Lemma~\ref{ptq:P}(a)--(c) for $(G',\sigma',\lambda',\mathsf P')$ in place of
$(G,\sigma,\lambda,\mathsf P)$. Their proofs use only: that the operator is $\sigma(e_\lambda)$ for a
unitary representation $\sigma$ whose derived action on polynomial vectors is the row action
{\cite[(12)]{P49}} (Lemma~\ref{ptr:const}), with $e_\lambda=d_\lambda\aubar{c_\lambda}$ the type of the
discrete series $\pi_\lambda$ and $c_\lambda$ the coefficient of the highest vector of its lowest
$K$-type; the integrability of $e_\lambda$ and of its derivatives (Lemma~\ref{ptr:gap}); Fischer
adjointness; that polynomial vectors of bounded degree form a finite-dimensional space stable
under $K$; {\cite[Lemma~3.5]{P49}} and (D1)--(D4). All of these hold for $(G',\sigma',\lambda')$ with the
column action (the parameter $\lambda'$ is of the class of Lemma~\ref{ptr:gap}, and
{\cite[Lemma~3.5]{P49}} and (D1)--(D4) are stated for an arbitrary unitary group and regular parameter).

(d) By (D5) applied to $e_\lambda$ on $G$ (smooth, bi-$K$-finite and $\mathfrak Z(\mathfrak g)$-finite,
as in \S\ref{ptq:ss-H}) choose $\beta_0\in C_c^\infty(G)$ with $\beta_0*e_\lambda=e_\lambda$; as
$\aubar{c_\lambda(y^{-1}g)}=c_\lambda(g^{-1}y)$, this reads $\int\beta_0(y)c_\lambda(hy)\,dy=c_\lambda(h)$ for
all $h$. Let $x=\mathsf P'\Phi$; by Theorem~\ref{ptr:edge}(b), $x=\mathsf P\Phi$. In the proof of
Lemma~\ref{ptq:P}(b) it is shown that $\mathsf P\Phi$ is a smooth vector for $\sigma$ of $T$-weight
$\Lambda$, killed by the compact raising operators of the rows, with $d\sigma(\Omega)x=\kappa_\lambda x$.
So {\cite[Lemma~3.5]{P49}} gives $\langle x,\sigma(h)x\rangle=\|x\|^2c_\lambda(h)$, and for $g\in G$
\[
\langle\sigma(g)x,\sigma(\beta_0)x\rangle=\int\beta_0(y)\langle x,\sigma(g^{-1}y)x\rangle dy=\|x\|^2\!\int
\beta_0(y)c_\lambda(g^{-1}y)dy=\|x\|^2c_\lambda(g^{-1})=\langle\sigma(g)x,x\rangle .
\]
So $\sigma(\beta_0)x-x$ lies in the closed span of $\sigma(G)x$ and is orthogonal to it; it is $0$.
\end{proof}

\begin{theorem}[(H) with the groups exchanged]\label{ptw:Hc}
Let $(S,T)$ be an ascending or descending pair of the triangle. For every polynomial $Q$ at $\tau$
there is $\beta_Q\in C_c^\infty(G)$ with $\mathsf P'Q=\sigma(\beta_Q)\mathsf P'\Phi$. Consequently, for
every polynomial vector $Q$ of the Fock space of $\auHom(S,T)\otimes\R$ there is a continuous
compactly supported $\beta_Q$ on $\auU(S)_\infty$ with
\begin{equation}\label{ptw:eq-Hc}
\mathsf P'Q=\int_{\auU(S)_\infty}\beta_Q(x)\,\omega_\infty(x,1)\,\mathsf P'P_{S,T}\,dx,
\end{equation}
the integral being a Bochner integral in $\pzF_s$; here $\mathsf P'=\int e_T(x)\omega_\infty(1,x)\,dx$
is the operator of \S\ref{ptr:ss-edge}.
\end{theorem}
\begin{proof}
Let $\mathcal G$ be the space of polynomials $Q$ with $\mathsf P'Q=\sigma(\beta)\mathsf P'\Phi$ for some
$\beta\in C_c^\infty(G)$. We check the hypotheses of Proposition~\ref{ptw:howe}.

(i) Let $Q\in\mathcal G$, $\mathsf P'Q=\sigma(\beta)\mathsf P'\Phi$, and $Y\in\mathfrak g$. As $\mathsf P'$ is
bounded and commutes with $\sigma$, and $Q$ is a smooth vector for $\sigma$ with derivative
$d\sigma(Y)Q=E(Y)Q$ (Lemma~\ref{ptr:const}), $\mathsf P'E(Y)Q=\frac d{dt}|_0\sigma(\exp tY)\mathsf P'Q$; and
$\sigma(\exp tY)\sigma(\beta)\mathsf P'\Phi=\sigma(\beta(\exp(-tY)\,\cdot))\mathsf P'\Phi$, whose derivative
at $0$ is $\sigma(\beta_Y)\mathsf P'\Phi$ with $\beta_Y(x)=\frac d{dt}|_0\beta(\exp(-tY)x)\in C_c^\infty(G)$.
So $E(Y)Q\in\mathcal G$ for $Y\in\mathfrak g$, hence for $Y\in\mathfrak g_\C=\mathfrak{gl}_n(\C)$, in
particular for the $E_{ab}$.

(ii) Lemma~\ref{ptw:P}(d). (iii) Lemma~\ref{ptw:P}(c): $\mathsf P'$ of such a polynomial is
$0=\sigma(0)\mathsf P'\Phi$. (iv), (v) Lemma~\ref{ptw:P}(a), in the same way.

So $\mathcal G$ contains all polynomials. The global statement follows from the local one exactly
as in the proof of Theorem~\ref{ptq:H}: $\mathsf P'=\bigotimes_\tau\mathsf P'_\tau$, $P_{S,T}=
\bigotimes_\tau\Phi_\tau$, $\omega_\infty(x,1)=\bigotimes_\tau\sigma_\tau(x_\tau)$, and every polynomial
vector is a finite sum of products $\bigotimes_\tau Q_\tau$, for which $\beta_Q=\prod_\tau
\beta_{Q_\tau}$. The identity holds in the Hilbert space; both sides lie in $\pzF_s$ (the right side
is a Bochner integral of a continuous compactly supported $\pzF_s$-valued function, by
{\cite[Lemma~8.2(a)]{P49}}; the left side because $\mathsf P'$ is the Bochner integral of
$e_T(x)\omega_\infty(1,x)Q$, convergent in $\pzF_s$ by {\cite[Lemma~8.2(a)]{P49}} and the integrability
of $e_T$ times a power of $\|x\|$, Lemma~\ref{ptr:gap}); so it holds in $\pzF_s$ by
Lemma~\ref{ptv:embed}.
\end{proof}

\subsubsection{The reversal at the infinite places and the characters of (W7)}\label{ptw:ss-rev}

Let $(S,T)$ be spaces of ranks $N\le N'$, and $\mathcal A=\mathrm{Fl}\circ\mathsf C$, $\phi(x)=x^\dagger$ as in
\S\ref{ptv:ss-rev}, with $\mathcal A=\bigotimes_v\mathcal A_v$. Write $W_\infty=\auHom(S,T)\otimes\R$,
$W'_\infty=\auHom(T,S)\otimes\R$, $\pzF$, $\pzF'$ their Fock spaces ({\cite[\S8.1]{P49}}), and
$K\times K'$ the maximal compact subgroup of $\auU(S)_\infty\times\auU(T)_\infty$ used there.

\paragraph*{The complex structures}
In the realisation {\cite[(41)]{P49}} of {\cite[\S8.1]{P49}} there are a real-linear isomorphism
$u\colon W_\infty\to\C^D$ and the standard hermitian form $h(z,z')=\sum_kz_k\aubar{z'_k}$ with
$(\psrho_\infty(y)F)(z)=e^{h(z,u(y))-|u(y)|^2/2}F(z-u(y))$ (in {\cite[(41)]{P49}}, $u(y)=\aubar y$ and
$z\cdot y=h(z,\aubar y)$); the complex structure $J=u^{-1}\circ i\circ u$ of $W_\infty$ is the positive
complex structure commuting with $K\times K'$ of {\cite[\S8.1]{P49}}. From the formula,
$\psrho_\infty(x)\psrho_\infty(y)=e^{iA(x,y)}\psrho_\infty(y)\psrho_\infty(x)$ with
$A(x,y)=2\Iml h(u(y),u(x))$ (the exponents of the two products differ by $-h(u(x),u(y))+h(u(y),u(x))$),
so $A(w,Jw)=2|u(w)|^2>0$ for $w\ne0$. The same holds for $W'_\infty$ with $u',J',A'$.

\begin{lemma}\label{ptw:J}
\textup{(a)} A complex structure $J_1$ on $W_\infty$ commuting with $K\times K'$ and with
$A(w,J_1w)>0$ for all $w\ne0$ equals $J$.
\textup{(b)} $A'(\phi x,\phi y)=-A(x,y)$, and $\phi\circ J=-J'\circ\phi$.
\end{lemma}
\begin{proof}
(a) At each $\tau$, in orthogonal bases of $S_\tau$ and $T_\tau$, $W_\tau=\auHom_\C(S_\tau,T_\tau)$ is
the sum of the four blocks $\auHom(S_{\tau,\pm},T_{\tau,\pm})$ ($S_{\tau,\pm}$ the spans of the basis
vectors of sign $\pm$), stable under $K_\tau\times K'_\tau=\auU(S_{\tau,+})\times\auU(S_{\tau,-})\times
\auU(T_{\tau,+})\times\auU(T_{\tau,-})$, which acts on a non-zero block $\auHom(S_{\tau,a},T_{\tau,b})$
through $x\mapsto h_bxg_a^{-1}$. This is an irreducible complex representation, and it is not
isomorphic to its complex conjugate (the centres of $\auU(S_{\tau,a})$ and $\auU(T_{\tau,b})$ act by
$\zeta^{-1}\zeta'$, resp.\ $\zeta\zeta'^{-1}$); so it is irreducible as a real representation and its
real commutant is $\C$ (the complex structure $i$ of the block and the scalars). Two distinct
non-zero blocks, also for different $\tau$, are non-isomorphic real representations (some central
factor acts non-trivially on one and trivially on the other). Hence $J_1$ preserves every block and
acts on it by an element $a+bi$ of the commutant with $(a+bi)^2=-1$, i.e.\ by $\epsilon i$,
$\epsilon=\pm1$. The same holds for $J$, with signs $\epsilon_0$. On a block, $A(w,J_1w)=\epsilon
A(w,iw)$ and $A(w,Jw)=\epsilon_0A(w,iw)$ are both positive for $w\ne0$; so $\epsilon=\epsilon_0$ and
$J_1=J$.

(b) By \eqref{ptv:eq-A}, $\mathcal A_\infty$ is conjugate-linear with $\mathcal A_\infty\psrho_\infty(x)=
\psrho'_\infty(\phi x)\mathcal A_\infty$. Applying $\mathcal A_\infty$ to $\psrho(x)\psrho(y)=e^{iA(x,y)}
\psrho(y)\psrho(x)$ gives $\psrho'(\phi x)\psrho'(\phi y)=e^{-iA(x,y)}\psrho'(\phi y)\psrho'(\phi x)$, so
$e^{iA'(\phi x,\phi y)}=e^{-iA(x,y)}$ for all real $x,y$; replacing $x$ by $sx$, $s\in\R$, gives
$A'(\phi x,\phi y)=-A(x,y)$. Put $J_1=-\phi^{-1}J'\phi$. It is a complex structure; it commutes with
$K\times K'$, because $\phi(hxg^{-1})=g\phi(x)h^{-1}$ for unitary $g,h$ (\S\ref{ptv:ss-rev}) and $J'$
commutes with $K'\times K$; and $A(w,J_1w)=-A(w,\phi^{-1}J'\phi w)=A'(\phi w,J'\phi w)>0$ for $w\ne0$.
By (a), $J_1=J$, i.e.\ $\phi J=-J'\phi$.
\end{proof}

\begin{lemma}[$\mathcal A_\infty$ preserves polynomial vectors]\label{ptw:poly}
$\mathcal A_\infty$ maps smooth vectors to smooth vectors, with $\mathcal A_\infty\partial_y=
\partial'_{\phi y}\mathcal A_\infty$, maps the vacuum to a multiple of the vacuum, and maps the space of
polynomial vectors of degree $\le d$ onto the corresponding space of $\pzF'$, for every $d$. The
same holds for $\mathcal A_\infty^{-1}$.
\end{lemma}
\begin{proof}
For a smooth vector $v$, $s\mapsto\psrho'(s\phi y)\mathcal A_\infty v=\mathcal A_\infty\psrho(sy)v$ is
differentiable with derivative $\mathcal A_\infty\partial_yv$ at $0$ ($\mathcal A_\infty$ is a bounded
real-linear map); this gives the first assertion. From the formula for $\psrho_\infty$,
$\partial_yF=h(z,u(y))F-D_{u(y)}F$, where $D_{u}$ is the derivative in the direction $u$ (complex-linear
in $u$, $F$ being entire), and $u(Jy)=iu(y)$; hence
\[
\partial_y+i\partial_{Jy}=2\,h(z,u(y)),\qquad\partial_y-i\partial_{Jy}=-2D_{u(y)}
\]
(multiplication by a linear form, resp.\ a derivation), and likewise on $\pzF'$. By
conjugate-linearity and Lemma~\ref{ptw:J}(b),
\[
\mathcal A_\infty(\partial_y\pm i\partial_{Jy})=(\partial'_{\phi y}\mp i\partial'_{\phi Jy})\mathcal A_\infty
=(\partial'_{\phi y}\pm i\partial'_{J'\phi y})\mathcal A_\infty .
\]
So $\mathcal A_\infty$ carries the operators $D_{u(y)}$ to the operators $D'_{u'(\phi y)}$ and the
multiplications by the linear forms $h(z,u(y))$ to the multiplications by $h(z',u'(\phi y))$; as $y$
runs over $W_\infty$ these run over all directional derivatives and all linear forms. The vacuum $1$
is killed by all $D_u$; so $\mathcal A_\infty1$ is a smooth vector (an entire function) killed by all
$D'_{u'}$, i.e.\ a constant $c$, and $c\ne0$ since $\mathcal A_\infty$ is injective. The polynomial
vectors of degree $\le d$ are the span of the products of at most $d$ multiplications by linear forms
applied to $1$; so $\mathcal A_\infty$ maps them into the polynomial vectors of degree $\le d$ of
$\pzF'$. The same argument for $\mathcal A_\infty^{-1}$ (which intertwines $\psrho'$ and
$\psrho\circ\phi^{-1}$) gives the reverse inclusion.
\end{proof}

\begin{lemma}[The exponents in (W7)]\label{ptw:W7}
Let $S,T$ have ranks $N,N'$. The character $\lambda$ of \textup{(W7)} (Lemma~\ref{pc:SW}) is
$\lambda(g,h)=\upsilon(\det g)^{-N'}\upsilon(\det h)^{N}$. Consequently the characters $\lambda_S,\lambda_T$
of \eqref{ptv:eq-A} are trivial: $\mathcal A\,\omega_{\auHom(S,T)}(g,h)=\omega_{\auHom(T,S)}(h,g)\,\mathcal A$,
and Proposition~\ref{ptv:rev} and Corollary~\ref{ptv:desc-adj} hold with $\lambda_S=\lambda_T=1$:
\[
\aubar{\theta_{\alpha,\beta}(g,h;P,\Psi)}=c\,\theta_{\beta,\alpha}(h,g\zeta_\pcm;P^\flat,\Psi^\flat).
\]
\end{lemma}
\begin{proof}
By (W7), $\lambda(g,h)=\upsilon(\det g)^a\upsilon(\det h)^b$ for integers $a,b$; we determine them at an
infinite place $\tau$. Let $W=\auHom(S,T)$ and $W^-=\auHom(S^-,T)$; at $\tau$ both have the same
underlying real space, symplectic forms of opposite signs, the same symplectic group and hence the
same metaplectic double cover $\mathrm{Mp}$, and the same image of $\auU(S)_\tau\times\auU(T)_\tau$. Let
$\omega^{\mathrm{Mp}}$, $\omega^{\mathrm{Mp},-}$ be the metaplectic representations on the Fock spaces
of $\psrho_\tau$, $\psrho^-_\tau$ ({\cite[\S3.2.7]{P49}}). The local component $\mathsf C_\tau$ of $\mathsf C$
is conjugate-linear with $\mathsf C_\tau\psrho_\tau(x)=\psrho^-_\tau(x)\mathsf C_\tau$, so
$\tilde u\mapsto\mathsf C_\tau\omega^{\mathrm{Mp}}(\tilde u)\mathsf C_\tau^{-1}$ is a strongly continuous
unitary representation of $\mathrm{Mp}$ with $\mathsf C_\tau\omega^{\mathrm{Mp}}(\tilde u)
\mathsf C_\tau^{-1}\psrho^-(w)(\cdots)^{-1}=\psrho^-(uw)$. By the irreducibility of $\psrho^-_\tau$ it
differs from $\omega^{\mathrm{Mp},-}$ by a continuous character of $\mathrm{Mp}$; as $\mathrm{Mp}$ is
connected with perfect Lie algebra $\mathfrak{sp}$, this character is trivial. By (V) for $W$ and
for $W^-$ ({\cite[Theorem~3.29]{P49}}, whose proof uses only that $\omega$ is the Weil representation of
$\auU(W)$ defined with $\upsilon$, together with (W1$_\infty$), (W2$_\infty$), (M), and (S) when $W$
is definite), for $u=\iota(g,h)$ with a lift $\tilde u$,
$\omega_\tau(u)=d(\tilde u)\,\omega^{\mathrm{Mp}}(\tilde u)$ and $\omega^-_\tau(u)=d(\tilde u)\,
\omega^{\mathrm{Mp},-}(\tilde u)$ with the same genuine function $d(\tilde u)=\det(g)^{N'/2}
\det(h)^{-N/2}$, $|d|=1$. Hence
\[
\mathsf C_\tau\omega_\tau(u)\mathsf C_\tau^{-1}=\aubar{d(\tilde u)}\,\omega^{\mathrm{Mp},-}(\tilde u)=
d(\tilde u)^{-2}\omega^-_\tau(u)=\det(g)^{-N'}\det(h)^{N}\,\omega^-_\tau(u),
\]
and since $\upsilon_\tau(z)=z$ on $\auU(1)$ ({\cite[Theorem~3.29]{P49}}), $a=-N'$, $b=N$ ($\det$ maps
$\auU(S)_\tau$ onto $\auU(1)_\tau$). By \eqref{ptv:eq-A}, $\lambda_S(\det g)\lambda_T(\det h)^{-1}$ is
$\lambda(g,h)$ times the factor $\upsilon(\det g)^{N'}\upsilon(\det h)^{-N}$ of {\cite[Proposition~7.17]{P49}}
(\S\ref{pcx:ss-flip}), i.e.\ $1$; as $\det$ maps $\auU(S)(\auA)$ and $\auU(T)(\auA)$ onto
$\auU(1)(\auA_{E^+})$, $\lambda_S=\lambda_T=1$. The rest is Proposition~\ref{ptv:rev} and
Corollary~\ref{ptv:desc-adj} with this value.
\end{proof}

\begin{remark}\label{ptw:rem-W7}
For the descent $d_{i,0}$ this gives the exponent needed in Proposition~\ref{ptv:chars}(b): the
conjugate of the kernel of $a_{i,0}$ (characters $\alpha=\mu_{i+1}\upsilon_{M_1}$,
$\beta=\mu_{i+1}\upsilon_{M_0}$) is a kernel of $\auHom(\pcG_{i,0}^-,\pcG_{i,1})$ with characters
$\alpha^{-1}\upsilon^{M_1}$ and $\beta^{-1}\upsilon^{M_0}$, since $\aubar{\theta_{\alpha,\beta}(g,h;F)}=
\alpha(\det g)\beta(\det h)^{-1}\lambda(g,h)\theta^-(\omega^-(g,h)\mathsf CF)$. Its target character is
$\mu_{i+1}^{-1}\upsilon^{M_0-\lfloor M_0/2\rfloor}=\mu_{i+1}^{-1}\upsilon^{\lceil r/2\rceil}$, of weight
$\kappa_{i+1}+\lceil r/2\rceil$, as Definition~\ref{ptr:desc-char} requires.
\end{remark}

\paragraph*{The reversed pairs}
Let $(S,T)=(\pcG_{i+1,k},\pcG_{i,k+1})$ be a descending pair, $i+k+1\le m$, with the characters
$\alpha=\alpha^\downarrow_{i,k}$, $\beta=\beta^\downarrow_{i,k}$ of Proposition~\ref{ptv:chars} (for a
twist $\vartheta$), the types $e_S=\aubar{e_{i+1,k}}$, $e_T=\aubar{e_{i,k+1}}$ and the joint harmonic
$P=P_{S,T}$ of \S\ref{ptr:ss-edge}. The \emph{reversed pair} is $(T,S)$, i.e.\ the space
$\auHom(T,S)$, with source character $\beta$ on $\auU(T)$, target character $\alpha$ on $\auU(S)$, the
Weil representation $\omega^{\mathrm{rev}}=\omega_{\beta,\alpha}$ of $\auHom(T,S)$, the polynomial
$P^\flat=\mathcal A_\infty P$ (a polynomial vector by Lemma~\ref{ptw:poly}), source type $e_T$ and
target type $e_S$. Put $\mathsf P^{\mathrm{rev}}=\int_{\auU(T)_\infty}e_T(x)\,\omega^{\mathrm{rev}}_\infty
(x^{-1},1)\,dx$, the operator $\mathsf P$ of \S\ref{ptr:ss-edge} for this pair.

\begin{proposition}[Reversed pairs are edges]\label{ptw:revedge}
\textup{(c)} Every form of type $e_T$, $\aubar{e_T}$, $e_S$ or $\aubar{e_S}$ is a cusp form, rapidly
decreasing on Siegel sets.
\textup{(d)} For every finite datum $f$ and every form $y$ of type $e_T$, the lift
$g\mapsto\int_{[\auU(T)]}\theta^{\mathrm{rev}}(\omega^{\mathrm{rev}}(h,g)(P^\flat\otimes f))\,y(h)\,dh$
converges absolutely and is a form of type $e_S$; for every form $z$ of type $\aubar{e_S}$ the
transposed lift $h\mapsto\int_{[\auU(S)]}\theta^{\mathrm{rev}}(\omega^{\mathrm{rev}}(h,g)(P^\flat\otimes
f))\,z(g)\,dg$ converges absolutely and is a form of type $\aubar{e_T}$.
\textup{(H)} For every polynomial vector $Q$ of the Fock space of $\auHom(T,S)\otimes\R$ there is
$\alpha_Q\in C_c(\auU(S)_\infty)$ with $\mathsf P^{\mathrm{rev}}Q=\int\alpha_Q(x)\,
\omega^{\mathrm{rev}}_\infty(1,x)\,\mathsf P^{\mathrm{rev}}P^\flat\,dx$ (Bochner integral in $\pzF'_s$).
\textup{(L)} Lemma~\ref{ptr:conv-lift} holds for the reversed pair: for every polynomial vector $Q$,
finite datum $f$ and form $y$ of type $e_T$, the lift $F_Q$ of $y$ with $Q\otimes f$ satisfies
$F_Q(g)=\int\alpha_Q(x)F_{P^\flat}(gx)\,dx$.
\end{proposition}
\begin{proof}
(c) is Theorem~\ref{ptr:edge}(c). By Lemma~\ref{ptw:W7} and the first lines of the computation
``Globally'' in the proof of Proposition~\ref{ptv:rev} (valid for every datum $F$), for every
$F\in\psHS(\auHom(S,T))$
\[
\aubar{\theta_{\alpha,\beta}(g,h;F)}=c_\theta^{-1}\,\theta_{\beta,\alpha}(h,g;\mathcal AF),
\]
and $\mathcal A$ is a bijection. (d) Let $F'=P^\flat\otimes f$ and $F=\mathcal A^{-1}F'=P\otimes f_1$
($f_1=\mathcal A_f^{-1}f$). Then $\int\theta_{\beta,\alpha}(h,g;F')y(h)\,dh=c_\theta\,\aubar{\int
\theta_{\alpha,\beta}(g,h;F)\,\aubar{y(h)}\,dh}$. As $\aubar y$ has type $\aubar{e_T}$, the integral on
the right is the transposed lift along the descending pair, absolutely convergent and a form of type
$\aubar{e_S}$ by Theorem~\ref{ptr:edge}(d); its conjugate is a form of type $e_S$. Likewise the
transposed reversed lift of $z$ is $c_\theta$ times the conjugate of the lift of $\aubar z$ (type
$e_S$) along the descending pair, a form of type $e_T$, whose conjugate has type $\aubar{e_T}$.

(H) At the infinite places, Lemma~\ref{ptw:W7} gives $\mathcal A_\infty\omega_\infty(g,h)=
\omega^{\mathrm{rev}}_\infty(h,g)\mathcal A_\infty$ (the archimedean factors of the characters move
through the conjugate-linear map as in the proof of Proposition~\ref{ptv:rev}). With
$e_T(x^{-1})=\aubar{e_T(x)}$ and the unimodularity of $\auU(T)_\infty$,
\[
\mathcal A_\infty\mathsf P'v=\int\aubar{e_T(x)}\,\omega^{\mathrm{rev}}_\infty(x,1)\mathcal A_\infty v\,dx
=\int e_T(x)\,\omega^{\mathrm{rev}}_\infty(x^{-1},1)\mathcal A_\infty v\,dx=\mathsf P^{\mathrm{rev}}
\mathcal A_\infty v,
\]
$\mathsf P'=\int e_T(x)\omega_\infty(1,x)dx$ being the column projection of the descending pair
($\mathcal A_\infty$ is continuous and real-linear, so it commutes with Bochner integrals up to
conjugation of the scalar factor). Given $Q$, the vector $Q_1=\mathcal A_\infty^{-1}Q$ is a polynomial
vector (Lemma~\ref{ptw:poly}); by Theorem~\ref{ptw:Hc}, $\mathsf P'Q_1=\int\beta(x)\omega_\infty(x,1)
\mathsf P'P\,dx$. Applying $\mathcal A_\infty$,
\[
\mathsf P^{\mathrm{rev}}Q=\int\aubar{\beta(x)}\,\omega^{\mathrm{rev}}_\infty(1,x)\,\mathsf P^{\mathrm{rev}}
P^\flat\,dx,
\]
an identity in the Hilbert space; both sides lie in $\pzF'_s$ (as in the proof of
Theorem~\ref{ptw:Hc}), so it holds there by Lemma~\ref{ptv:embed}. Take $\alpha_Q=\aubar\beta$.

(L) The proof of Lemma~\ref{ptr:conv-lift} uses the functional $L_g(\Psi)=\int_{[\auU(T)]}
\theta^{\mathrm{rev}}(\omega^{\mathrm{rev}}(h,g)(\Psi\otimes f))y(h)\,dh$, continuous on $\pzF'_s$ with
$L_g(\omega^{\mathrm{rev}}_\infty(1,x)\Psi)=L_{gx}(\Psi)$ and $L_g(\mathsf P^{\mathrm{rev}}\Psi)=L_g(\Psi)$.
These are established in the proof of {\cite[Proposition~8.4(b)]{P49}} from {\cite[Lemma~8.2]{P49}}
(stated for an arbitrary pair of spaces), reduction theory, the rapid decay of $y$ (by (c)) and
$y=R(e_T)y$; no property of the parameters is used. With (H) the proof of
Lemma~\ref{ptr:conv-lift} applies verbatim.
\end{proof}

\subsubsection{The reversed square law}\label{ptw:ss-RS}

\begin{lemma}[Change of splitting]\label{ptw:split}
In the square law of \S\ref{pc:ss-exchange} let $\mathsf N=\mathsf B\sqcup\mathsf A=\mathsf B'\sqcup
\mathsf A'$ be two splittings with $|\mathsf B|=|\mathsf B'|$. If $L,L'$ satisfy the square law with
constant $c$ for $(\mathsf B,\mathsf A)$, they satisfy it with constant $\chi(\epsilon)c$ for
$(\mathsf B',\mathsf A')$, for a sign $\epsilon\in\{\pm1\}\subset\pcR^\times$ depending only on the
two splittings (the functions on matrices being read through bijections of $\mathsf A$ with
$\mathsf A'$ and of $\mathsf B$ with $\mathsf B'$); in particular $|\chi(\epsilon)c|=|c|$.
\end{lemma}
\begin{proof}
Let $\sigma$ be the permutation matrix of $\mathsf N$ carrying $\mathsf B$ onto $\mathsf B'$ and $\mathsf A$
onto $\mathsf A'$ through the given bijections, $\epsilon=\det\sigma$. For $k\in\auGL_{\mathsf N}$ put
$\tilde k=k\sigma$: the columns $\mathsf B$ of $\tilde k$ are the columns $\mathsf B'$ of $k$, and the rows
$\mathsf A$ of $\tilde k^{-1}=\sigma^{-1}k^{-1}$ are the rows $\mathsf A'$ of $k^{-1}$. So if $(\Phi,\Psi)$ is
matched for $(\mathsf B',\mathsf A')$, then $\Phi(x_{\tilde k})=\chi(\det k)\Psi(y_{\tilde k})=
\chi(\epsilon)^{-1}\chi(\det\tilde k)\Psi(y_{\tilde k})$, i.e.\ $(\Phi,\chi(\epsilon)^{-1}\Psi)$ is matched for
$(\mathsf B,\mathsf A)$, and $\chi(\epsilon)^{-1}\sum\Psi L'=c\sum\Phi L$. The transformation laws do not
involve the splitting. As $\epsilon^2=1$, $\chi(\epsilon)=\pm1$.
\end{proof}

\begin{theorem}[(RS)]\label{ptw:RS}
Let $i+n+2\le m$. The reversed square with $U_1=\pcG_{i+1,n}$, $U_2=\pcG_{i,n+2}$, $A=\pcG_{i+1,n+1}$,
$A'=\pcG_{i,n+1}$, types $e_1=e_{i+1,n}$, $e_2=e_{i,n+2}$, $e=e_{i+1,n+1}$, $e'=e_{i,n+1}$, the
ascending pairs $(U_1,A)$, $(A',U_2)$ (the lifts $a_{i+1,n}$, $a_{i,n+1}$ with their joint harmonics),
the reversed pairs $(U_2,A)$, $(A',U_1)$ of the descending pairs at $(i,n+1)$ and $(i,n)$ (with
$P^\flat_{i,n+1}$, $P^\flat_{i,n}$ and the characters of Proposition~\ref{ptv:chars}), and the family of
$m$-fold twists $\vartheta$ with $\pcm=\pcm_\vartheta$, $\nu=(\nu_{n+1}\vartheta^{n+1})^{-1}$,
$\chi=\chi_\vartheta$, has excess $0$ and weight $(n+1)\Delta$ and satisfies Definition~\ref{ptr:def-square} with (E1) for the
reversed pairs replaced by (c), (d), (H) and (L) of Proposition~\ref{ptw:revedge}. The conclusion of
Theorem~\ref{ptr:square} holds for it (see the proof). Consequently
\textup{(RS)} of \S\ref{ptv:ss-RS} holds, with $\kappa^{\mathrm{rev}}_{i,n}(\vartheta)=\pm\kappa/c(\vartheta)$,
$|c(\vartheta)|^2=C$; so $|\kappa^{\mathrm{rev}}_{i,n}|^2=|\kappa|^2/C$ is constant and
$\kappa^{\mathrm{rev}}_{i,n}$ has modulus $1$ up to norm powers.
\end{theorem}
\begin{proof}
\emph{Ranks and signs.} $a=\operatorname{rk}A=M_{n+1}=\operatorname{rk}A'=b$, so $e=0$, and
$N=M_n+M_{n+2}=2M_{n+1}=a+b$. Write $q(X)$ for the number of negative signs. The sign identity of
\S\ref{pc:ss-exchange} for the square at $(i,n)$ is $q(\pcG_{i,n+1})+q(\pcG_{i+1,n+1})=
q(\pcG_{i,n+2})+q(\pcG_{i+1,n})$; rearranged it is $q(U_1)+q(U_2)=q(A)+q(A')$, the identity required
here. Moreover the number $q=q(U_1)+q(U_2)$ equals the number $q(\pcG_{i,n+1})+q(\pcG_{i+1,n+1})$ of
the square at $(i,n)$, and $N$ is the same; so the inequalities $2q+2\le N\le2\mu_\tau+2q$, with
$\mu=(n+1)\Delta$ the weight of $\nu$, are those verified in \S\ref{pc:ss-exchange} for that square
(Corollary~\ref{ptr:tri-b}).

\emph{(E1).} The pairs $(U_1,A)=(\pcG_{i+1,n},\pcG_{i+1,n+1})$ and $(A',U_2)=(\pcG_{i,n+1},
\pcG_{i,n+2})$ are ascending, from $e_1$ to $e$ and from $e'$ to $e_2$, and are edges by
Theorem~\ref{ptr:edge}; (H) for $(U_1,A)$ is Theorem~\ref{ptq:H}. The pair $(U_2,A)$ is the reversed
pair of the descending pair $(\pcG_{i+1,n+1},\pcG_{i,n+2})$, for which $e_T=\aubar{e_{i,n+2}}=
\aubar{e_2}$, $e_S=\aubar{e_{i+1,n+1}}=\aubar e$: it goes from $\aubar{e_2}$ to $\aubar e$, as
Definition~\ref{ptr:def-square} requires; and $(A',U_1)$ is the reversed pair of
$(\pcG_{i+1,n},\pcG_{i,n+1})$, from $\aubar{e_{i,n+1}}=\aubar{e'}$ to $\aubar{e_{i+1,n}}=\aubar{e_1}$.
By Proposition~\ref{ptw:revedge}, (c), (d) hold for them and (H) for $(U_2,A)$. The polynomials and
types do not depend on $\vartheta$; the characters do, with fixed weights.

\emph{What the proof of Theorem~\ref{ptr:square} uses of (E1).} The transformation laws use only
Lemma~\ref{pc:level} and (E2); the factorisation of the kernels uses only the tensor structure of
the Weil representations; Proposition~\ref{ptr:conv} uses the rapid decay of $y_1,y_2$ ((c)) and
finite moments of the Fock vectors; the multiplicity-one step uses Lemma~\ref{ptr:conv-lift} for
$(U_1,A)$ and $(U_2,A)$, Lemma~\ref{ptr:avg} for the forms $\Lambda^{(1)}_{\xi_1}y_1$ of type $e$ and
$\Lambda^{(2)}_{\xi_2}y_2$ of type $\aubar e$ ((d); $e=e_{i+1,n+1}$ is integrable of the class of
Lemma~\ref{ptr:gap}), and Lemma~\ref{ptr:dense}; the exchange step uses Theorem~\ref{pcx:X} with the
polynomial $v=P'_1P'_2$. Parts (a), (b) of Theorem~\ref{ptr:edge} enter only through (d) and
Lemma~\ref{ptr:conv-lift}, which for the reversed pairs are Proposition~\ref{ptw:revedge}(d), (L).
So the proof applies.

\emph{(E2).} By Lemma~\ref{ptw:W7} the kernel of the reversed pair of the descent $d_{i,k}$ has
source character $\beta^\downarrow_{i,k}$ (on $\pcG_{i,k+1}$) and target character
$\alpha^\downarrow_{i,k}$ (on $\pcG_{i+1,k}$). The four conditions, with $\gamma=\nu^{-1}\upsilon^{N/2}=
\gamma_n$ of \S\ref{ptv:ss-chars}, are:
\begin{align*}
&\text{at }A=\pcG_{i+1,n+1}:&\mu_{i+n+2}\upsilon_{M_n}\cdot\alpha^\downarrow_{i,n+1}&=\gamma_n,&
&\text{at }A'=\pcG_{i,n+1}:&\beta^\downarrow_{i,n}\cdot\mu_{i+n+2}\upsilon_{M_{n+2}}&=\gamma_n,\\
&\text{at }U_1=\pcG_{i+1,n}:&\mu_{i+n+2}\upsilon_{M_{n+1}}\cdot\alpha^\downarrow_{i,n}&=\gamma_n,&
&\text{at }U_2=\pcG_{i,n+2}:&\beta^\downarrow_{i,n+1}\cdot\mu_{i+n+2}\upsilon_{M_{n+1}}&=\gamma_n
\end{align*}
(at $A$: target characters of $a_{i+1,n}$ and of the reversed $d_{i,n+1}$; at $A'$: source
characters of the reversed $d_{i,n}$ and of $a_{i,n+1}$; at $U_1$: source character of $a_{i+1,n}$
times target character of the reversed $d_{i,n}$; at $U_2$: source character of the reversed
$d_{i,n+1}$ times target character of $a_{i,n+1}$). These are exactly the four equations of the
square at $(i,n)$ in \S\ref{ptv:ss-chars} (at $U_2$, $U_1$, $A'$, $A$ there, respectively), which hold
by Proposition~\ref{ptv:chars}.

\emph{Conclusion.} Theorem~\ref{ptr:square} gives $\kappa$, $C>0$ and, for each $\vartheta$, $c$ with
$|c|^2=C\,L_*(\pcm,\nu,0)=C$ (Corollary~\ref{pcx:cne0}), such that $L,L'$ satisfy the square law with
constant $\kappa/c$, for the splitting used in Theorem~\ref{pcx:X}. With $y_1=y$, $y_2=\aubar z$ we
have $L(\xi)=\mathfrak L(\xi)$ for $\xi$ of rank $a$ and $L'(\eta)=\mathfrak L'(\eta)$ for $\eta$ of rank
$b$, the only matrices on which matched functions are supported: indeed $\Lambda^{(1)}_{\xi_1}=a_{i+1,n}[\xi_1]$,
$R_{n+1}[\xi_2]$ is the transposed lift along the reversed pair $(U_2,A)$, so by Fubini (absolute
convergence, Proposition~\ref{ptr:conv}(a))
\[
\mathfrak L(\xi)=\int_{[\auU(U_2)]}\aubar{z(h)}\int_{[\auU(A)]}\theta^{\mathrm{rev}}(h,g;\xi_2)\,
(a_{i+1,n}[\xi_1]y)(g)\,dg\,dh=\int_{[\auU(A)]}\Lambda^{(1)}_{\xi_1}y\cdot\Lambda^{(2)}_{\xi_2}\aubar z
=L(\xi),
\]
and $\mathfrak L'(\eta)=\int_{[\auU(U_2)]}\aubar z\cdot a_{i,n+1}[\eta_2]R_n[\eta_1]y=L'(\eta)$ with
$T'^{(1)}_{\eta_1}=R_n[\eta_1]$, $\Lambda'^{(2)}_{\eta_2}=a_{i,n+1}[\eta_2]$ (Fubini, by
Proposition~\ref{ptr:conv}(b)). The law $\mathfrak L(h\xi)=\chi(\det h)\mathfrak L(\xi)$ for
$h\in\auGL_{\mathsf M}(\pcR)$ and \emph{all} $\xi$ is proved directly as the transformation law in the
proof of Theorem~\ref{ptr:square}: choose $k\in K_A$ reducing to $h$; by Lemma~\ref{pc:level},
$\omega(1,k^{-1})$ on $\auHom(U_1,A)$ and on $\auHom(U_2,A)$ send the point data
at $h\xi_1$, $h\xi_2$ to the target characters at $\det k$, inverted, times the data at $\xi_1,\xi_2$;
so $a_{i+1,n}[h\xi_1]y=c_1(\det k)^{-1}\psrho(k^{-1})a_{i+1,n}[\xi_1]y$, and the substitution
$g\mapsto gk$ in the integral over $[\auU(A)]$ defining $R_{n+1}[h\xi_2]$ gives
$R_{n+1}[h\xi_2]\psrho(k^{-1})=c_2(\det k)^{-1}R_{n+1}[\xi_2]$; with $c_1c_2=\gamma$ ((E2) at $A$) and
$\gamma^{-1}|_{K_A}=\chi\circ\det$ at the primes of $\pcm$, $\mathfrak L(h\xi)=\chi(\det h)\mathfrak L(\xi)$.
By Lemma~\ref{ptw:split} the
law holds for the splitting $(\mathsf B,\mathsf A)=(\mathsf M,\mathsf N\sqcup\mathsf T')$ of (RS) with
constant $\pm\kappa/c$. This is (RS).
\end{proof}

\begin{corollary}[The adjoint law]\label{ptw:adjoint}
Theorem~\ref{pc:tri}\textup{(c)} holds: the adjoint law \textup{(T4)} holds for the triangle, with
functions $\kappa_{i,n}$ of modulus $1$ up to norm powers. It depends on the same inputs as
Theorem~\ref{pc:tri}\textup{(b)}: (V) ({\cite[Theorem~3.29]{P49}}), Theorem~\ref{pcx:X}, the published results
(D1)--(D5), and (W7) (Lemmas~\ref{pc:SW} and~\ref{ptw:W7}).
\end{corollary}
\begin{proof}
Theorem~\ref{ptv:adjoint} with Theorem~\ref{ptw:RS}. In the constant of Theorem~\ref{ptv:adjoint},
$|c_n|^2$, $|c_{n+1}|^2$, $|\gamma|^2$ are constants times integral powers of $\auNm Q_\vartheta$ and do not
vanish, and $|\kappa^{\mathrm{rev}}_{i,n}|^2$ is constant; Lemma~\ref{pc:mod}(a).
\end{proof}

\subsubsection{The pairs $(W_j,V_j)$: the type on $V_j$}\label{ptw:ss-A2}

The proof of (A2) for $(W_j,V_j)$ in Proposition~\ref{ptv:arch} uses that $e_{V_j}$ is the type of
the representation of highest weight $\Lambda'$ generated by the joint harmonic $P_j=P^W_j$, as in
the definition of $\Theta_{j,\vartheta}$ in \S\ref{pc:sec}. This subsection proves the facts about
that type which (A2) requires.

Let $\sigma'=\omega_{\alpha,\beta,\infty}(1,\cdot)$ on the Fock space of $\auHom(W_j,V_j)\otimes\R$, with
the characters of the pair, $G'=\auU(V_j)_\infty$ (compact, connected) with Haar measure of total
mass $1$, and let $\rho_j$ be the representation of $G'$ on the span of $\sigma'(G')P_j$. For a
finite-dimensional irreducible unitary representation $\pi$ of $G'$ with a unit highest weight
vector $v$ put $e_\pi=d_\pi\,\aubar{c_\pi}$, $c_\pi(x)=\langle v,\pi(x)v\rangle$ (the convention
$e=d\,\aubar c$ of \S\ref{pc:ss-data}).

\begin{proposition}\label{ptw:A2}
\textup{(a)} $\rho_j$ is irreducible, of highest weight $\Lambda'=\lambda'-\rho'$ at every $\tau$, and
$P_j$ is a highest weight vector of it.
\textup{(b)} $\int_{G'}e_{\rho_j}(x)\,\sigma'(x)P_j\,dx=P_j$.
\textup{(c)} For every $\pcm$, $\Psi$, characters as in \textup{(A2)} and form $y$ of type $e_{W_j}$,
the lift $F(h)=\int_{[\auU(W_j)]}\theta(g,h;P_j,\Psi)\,y(g)\,dg$ converges absolutely, is bounded,
right invariant under an open subgroup of $\auU(V_j)(\auA_f)$, and $R(e_{\rho_j})F=F$.
Hence \textup{(A2)} of Lemma~\ref{pc:arch} holds for $(W_j,V_j)$ with
$e_{V_j}=e_{\rho_j}$, and with no other type of the form $e_\pi$ unless $e_\pi=e_{\rho_j}$.
\end{proposition}
\begin{proof}
(a) By the proof of {\cite[Lemma~3.6]{P49}}, $K'=G'$ acts on polynomial vectors by linear substitutions
times a character, so $\sigma'(G')P_j$ spans a finite-dimensional unitary representation. By Step~1
of the proof of {\cite[Proposition~4.8]{P54}}, $P_j$ (the tensor product of the $\pop_{\xi_0,\tau}$,
{\cite[Lemma~4.7]{P54}}) has $T'$-weight $\Lambda'$ and is killed by the $E'_{cd}$ with
$\lambda'_c>\lambda'_d$ (all roots are compact), i.e.\ it is a highest weight vector. A
finite-dimensional representation of a compact connected group generated by a highest weight vector
$v$ of weight $\Lambda'$ is irreducible: decompose $v=\sum_iv_i$ along a decomposition into
irreducibles; each non-zero $v_i$ is a highest weight vector of weight $\Lambda'$, so the summands
with $v_i\ne0$ are isomorphic to the irreducible representation $\pi_{\Lambda'}$; inside the isotypic
part $\pi_{\Lambda'}\otimes M$ the highest weight vectors are $v^+\otimes M$, so $v=v^+\otimes m$ and $v$
generates $\pi_{\Lambda'}\otimes m$.

(b) Put $\hat P=P_j/\|P_j\|$. By Schur orthogonality, for every $u$ in the space of $\rho_j$,
$\langle u,\int e_{\rho_j}(x)\rho_j(x)\hat P\,dx\rangle=d\int\aubar{\langle\hat P,\rho_j(x)\hat P\rangle}
\langle u,\rho_j(x)\hat P\rangle\,dx=\langle u,\hat P\rangle$; both vectors lie in the space of $\rho_j$, so
they are equal.

(c) Convergence, boundedness and invariance are shown in the proof of Proposition~\ref{ptv:arch}, which uses of the present proposition only parts (a) and (b), proved above without it. With $L_h$ as
in the proof of Lemma~\ref{ptr:conv-lift} ($L_h$ continuous on $\pzF_s$,
$L_h(\sigma'(x)\Psi)=L_{hx}(\Psi)$), $R(e_{\rho_j})F(h)=\int e_{\rho_j}(x)L_h(\sigma'(x)P_j)\,dx=
L_h(\int e_{\rho_j}(x)\sigma'(x)P_j\,dx)=L_h(P_j)=F(h)$ by (b). The last assertion: if $R(e_\pi)F=F$ for
all such $F$ and some $F\ne0$, then $\pi\cong\rho_j$ (the $G'$-span of the right translates of $F$ is
a quotient of $\rho_j$ by the formula $F(hx)=L_h(\sigma'(x)P_j)$), and the highest weight vectors
agree up to a scalar, so $e_\pi=e_{\rho_j}$.
\end{proof}

\begin{remark}[The scalar functions on $V_j$]\label{ptw:rem-A2}
The scalar functions attached to $\mathcal A(V_j;\vartheta)$ in \S\ref{pc:sec} are
$h\mapsto\langle\xi_0,\pi^\natural(h_\infty)^{-1}\varphi(h_f)\rangle$, with $\xi_0$ as in
{\cite[Lemma~4.7]{P54}} ($\pop_{\xi_0}=P_j$). This is consistent with Proposition~\ref{ptw:A2}: by
{\cite[(4)]{P54}} with $\xi=\pi^\natural(h_\infty)\xi_0$ and $\pop_{\pi^\natural(h_\infty)\xi_0}=
\sigma'(h_\infty)\pop_{\xi_0}$, the scalar lift with $P_j$ is
$h\mapsto\langle\pi^\natural(h_\infty)\xi_0,(\Theta'y)(h_f)\rangle$ (up to the character $\chi\circ\det$ of
Remark~\ref{pc:rem-norm}). Whether $\xi_0$ is the highest or the lowest weight vector of
$W_{\lambda_V}$ for the standard action depends on whether $\pi^\natural$ is the standard action or its
composite with complex conjugation (compare {\cite[Remark~3.16]{P49}}); nothing depends on it, since
Lemma~\ref{pc:algpet} holds for every unit vector in place of $v^+$, and (A2), (A3) and the norms
use only Proposition~\ref{ptw:A2}.
\end{remark}

\subsubsection{The edge hypothesis for the grid squares}\label{ptw:ss-grid}

The edge hypothesis (E1) for the squares of the grids (Remark~\ref{ptr:rem-grid},
Proposition~\ref{pgr:square}) is proved in Lemma~\ref{pgv:E1g} from the block rule of
Proposition~\ref{pgv:match}; every edge of a grid square is read in the direction of
Theorem~\ref{pgv:G3}(d), so Theorems~\ref{ptr:edge} and~\ref{ptq:H} suffice there and the exchanged
forms of this section are not needed.



\section{Lifts through the second grid}\label{C1:g4}
% ===== C-grid.tex =====
% proofs/C-grid.tex
% Theorem pc:auxcore (integral and short lifts through the second grid): architecture, the
% combinatorics of paths through two grids, transfer of periods along turns, and the deduction of
% pc:auxcore (a), (b), (c) from the intermediate propositions. Labels start with pgr:.
% Conventions and macros: those of proofs/C-determinant.tex.
\providecommand{\auI}{\mathcal I}
\providecommand{\auAut}{\operatorname{Aut}}
\providecommand{\auNm}{\operatorname{Nm}}
\providecommand{\auU}{\mathrm U}
\providecommand{\auGL}{\mathrm{GL}}
\providecommand{\auAz}{\mathcal A_0}
\providecommand{\auAm}{\mathcal A_m}
\providecommand{\auO}{\mathcal O}
\providecommand{\pcFF}{\mathfrak F}
\providecommand{\pcm}{\mathfrak m}
\providecommand{\pcb}{\mathfrak b}
\providecommand{\pgrY}{\widetilde Y}
\providecommand{\pgupt}{\mathsf u}
\providecommand{\pcMat}{\mathrm M}
\providecommand{\aubar}[1]{\overline{#1}}
\providecommand{\pgrG}{\Gamma}
\providecommand{\pgrfirst}{\Gamma_{\mathrm{first}}}
\providecommand{\pgrlast}{\Gamma_{\mathrm{last}}}
\providecommand{\pgrsq}{\operatorname{Sq}}
\providecommand{\pgrE}{\operatorname{Ed}}
\providecommand{\pgrTr}{\mathrel{\vdash}}

\subsection{Lifts through the second grid: Theorem~\ref{pc:auxcore}}\label{pgr:sec}

This section concerns Theorem~\ref{pc:auxcore}, from which {\cite[Proposition~14.38]{P46}} is deduced in
\S\ref{pc:ss-aux}. The notation is that of \S\ref{pc:sec}. \S\ref{pgr:ss-arch} fixes the objects
and states the intermediate propositions exactly; \S\ref{pgr:ss-paths} proves the combinatorics of
paths through the two grids; \S\ref{pgr:ss-transfer} proves that the turns compose to a transfer
from the first path to the last; \S\ref{pgr:ss-deduce} deduces Theorem~\ref{pc:auxcore} from the
intermediate propositions.

\subsubsection{Architecture}\label{pgr:ss-arch}

\paragraph*{The two grids}
For $0\le a,b\le m$ there are two spaces with frames (spaces in the sense of \S\ref{pc:ss-data}),
defined completely in \S\ref{pgv:ss-def} (Definition~\ref{pgv:def-spaces}):
\begin{itemize}
\item $Y^a_b$, \emph{the first grid}: a stem of rank $r$, $a$ blocks of rank $t$ of the first kind and
$b$ blocks of rank $t$ of the second kind; $Y^0_b=Y_b$ with the lattice of {\cite[\S5.1.3]{P49}};
\item $\pgrY^b_a$, \emph{the second grid}: a stem of rank $r+1$, $a$ blocks of rank $t+1$ and $b$
blocks of rank $t-1$; $\pgrY^0_0=Y_0\oplus E\tilde e$ with $\langle\tilde e,\tilde e\rangle=1$.
\end{itemize}
The signs are \eqref{pgv:eq-sign}: they depend only on the second index of the node $(a,b)$.
$Y^m_m$, $\pgrY^m_m$ and $\pgrY^m_0$ are positive definite, $\pgrY^l_l=Y^l_l\oplus E\tilde e_l$ with
$\langle\tilde e_l,\tilde e_l\rangle=(-1)^l$ as spaces with frames, and $\pgrY:=\pgrY^m_0$ is the space
of \S\ref{pc:ss-aux} (rank $\tilde n=r+1+m(t-1)$). A half $h_1$ of a two-grid path runs through the
nodes $(a,b)\mapsto\pgrY^b_a$. Each node carries a type $e_p$ (Harish-Chandra parameter $-\lambda_p$,
\eqref{pgv:eq-param}); the forms moved at a node are the forms of type $e_p$: in the second grid these
are forms, in the first grid \emph{conjugate forms}.

For an $m$-fold twist $\vartheta$ let $\pcm=\pcm(\vartheta)=\pcb\prod_{w\in Q_\vartheta}w$, with
$\pcb=(\pi_M)^b$, and $R_\pcm=\auO_E/\pcm$; here the exponent $b$ is the global parameter of {\cite[\S14.3]{P46}}, not the second grid index. This modulus is distinct from the modulus
$\pcm_\vartheta=(\pi_M)\prod_{w\in Q_\vartheta}w$ of \S\ref{pc:ss-aux}, which appears below only
through the finite ring $R_{\pcm_\vartheta}=\auO_E/\pcm_\vartheta$ (a product of finite fields) in
the index $[K_m:J^\circ_\vartheta]=\#\auGL_{r+mt}(R_{\pcm_\vartheta})$ of the principal congruence subgroup
$J^\circ_\vartheta=J_{Y_m}(\pcm_\vartheta)$ (Lemma~\ref{pgr:jcirc}). An \emph{edge} is a pair of adjacent nodes in one grid,
oriented towards the larger index; it is \emph{vertical} if it changes the second index in the first
grid or the first index in the second grid, and \emph{horizontal} otherwise. Its characters
$\alpha_e,\beta_e$, its edge character $\chi_e$ of $R_\pcm^\times$, its joint harmonic $P_e$ and its
\emph{moves} (the lift $\Lambda_{S_e}$ with the point datum $\delta_{S_e}$ when the first index grows,
the transposed lift $T_{S_e}$ when the second index grows, $S_e\in\pcMat_{q\times p}(R_\pcm)$) are
those of Definition~\ref{pgv:def-edges}; $\chi_e$ depends only on the kind and the index of $e$.

\paragraph*{Paths and turns}
A \emph{grid path} from $(a_0,b_0)$ to $(a,b)$ is a finite sequence of nodes starting at
$(a_0,b_0)$ and ending at $(a,b)$, each step being \emph{down}, $(a',b'+1)\to(a',b')$, or
\emph{right}, $(a',b')\to(a'+1,b')$. A \emph{two-grid path} $\pgrG=(l,h_0,h_1)$ consists of an index
$0\le l\le m$ (the \emph{diagonal index}) and two grid paths $h_0,h_1$ from $(0,m)$ to $(l,l)$, the
\emph{halves}: $h_0$ lives in the first grid (node $(a,b)$ is $Y^a_b$) and $h_1$ in the second
(node $(a,b)$ is $\pgrY^b_a$). Read as one path from $\pgrY^m_0$ to $Y^0_m$, $\pgrG$ goes along $h_1$
backwards in the second grid to $\pgrY^l_l$, crosses to $Y^l_l$, and goes along $h_0$ backwards to
$Y^0_m$. The \emph{first path} $\pgrfirst=(0,c,c)$ has both halves equal to the column
$c\colon(0,m)\to(0,m-1)\to\dots\to(0,0)$; the \emph{last path} $\pgrlast=(m,\rho,\rho)$ has both halves
equal to the row $\rho\colon(0,m)\to(1,m)\to\dots\to(m,m)$.

A \emph{corner} of a grid path is a node $(a,b)$ reached by a step down from $(a,b+1)$ and left by a
step right to $(a+1,b)$. For two-grid paths $\pgrG,\pgrG'$ and $s=(i,j)$ we write
$\pgrG\pgrTr_s\pgrG'$ (\emph{a turn at $s$}) in two cases:
\begin{itemize}
\item[(off)] for some $z\in\{0,1\}$ the half $h_z$ of $\pgrG$ has a corner at $(a,b)$, $\pgrG'$
equals $\pgrG$ except that in $h_z$ the two steps $(a,b+1)\to(a,b)\to(a+1,b)$ are replaced by
$(a,b+1)\to(a+1,b+1)\to(a+1,b)$, and $s=(a,b)$ if $z=0$, $s=(b,a)$ if $z=1$;
\item[(diag)] $\pgrG=(j,h_0,h_1)$ where both halves end with the step $(j,j+1)\to(j,j)$, and
$\pgrG'=(j+1,h'_0,h'_1)$ where $h'_z$ is $h_z$ with this last step replaced by
$(j,j+1)\to(j+1,j+1)$; here $s=(j,j)$.
\end{itemize}
The set of \emph{edges} $\pgrE(\pgrG)$ is the set of edges of $h_0$ in the first grid together with
those of $h_1$ in the second.

\paragraph*{Path data and periods}
An \emph{edge tuple} for $\vartheta$ is a family $S=(S_e)$ of matrices over $R_\pcm$, one for each
edge. A \emph{node family} is a family $g=(g_p)$ with $g_p\in\auGL_{p}(R_\pcm)$; it is
\emph{admissible} if $g_{Y^0_m}=1$, $g_{\pgrY^m_0}\equiv1\bmod\pcb$, and $g$ at $\pgrY^a_a$ is
$g_{Y^a_a}\oplus1$ for every $a$ (\S\ref{pgv:ss-def}). It acts by
$(gS)_e=g_q^{-1}S_eg_p$ for $e=(p,q)$, and we put $\chi_w(g)=\prod_{e\in w}\chi_e(\det g_q/\det g_p)$
for a set $w$ of edges. A complex number is a \emph{sum of at most $N$ roots of unity} if it is
$\sum_{k<N'}\zeta_k$ with $N'\le N$ and roots of unity $\zeta_k$.

\begin{definition}[Path data]\label{pgr:def-datum}
Let $w$ be a finite set of edges and $N\ge0$. A function $\Psi$ of edge tuples (for $\vartheta$) is a
\emph{path datum on $w$ with bound $N$} if: it depends only on $(S_e)_{e\in w}$;
$\chi_w(g)\Psi(gS)=\Psi(S)$ for every admissible $g$ and every $S$; $\Psi(S)\ne0$ implies that
$S_e$ has a left inverse for every vertical $e\in w$; and every value $\Psi(S)$ is a sum of at most
$N$ roots of unity.
\end{definition}

For a two-grid path $\pgrG$ with diagonal index $l$, a half $h$ in the grid $z$ and an edge tuple
$S$, let $h_*(S)$ be the composite of the moves of $S$ along $h$; it maps conjugate forms on $Y_m$
(if $z=0$), respectively forms on $\pgrY^m_0$ (if $z=1$), to conjugate forms on $Y^l_l$,
respectively forms on $\pgrY^l_l$. A \emph{period triple} is $\xi=(\vartheta,y',\varphi)$ with
$\vartheta$ an $m$-fold twist, $y'$ a conjugate form on $Y_m$ and $\varphi$ a form on $\pgrY^m_0$.
For a real point $g$ of $\auU(\pgrY^l_l)$ the \emph{diagonal period} of a conjugate form $y''$ on
$Y^l_l$ and a form $\varphi''$ on $\pgrY^l_l$ is $\mathsf P_{l,g}(y'',\varphi'')=\int_{[\auU(Y^l_l)]}
y''(x)\,\varphi''((x\oplus\pgupt^l)g)\,dx$, with $\pgupt=\pgupt_\pcm=\aubar u_0/u_0$, $u_0=\alpha
\varpi_\pcm$ (Definition~\ref{pgv:def-period}). For a family $\Psi=(\Psi_\vartheta)$
the \emph{period of $\Psi$ on $\pgrG$ at $g$} is
\[
I_{\pgrG}(\Psi;g)(\xi)=\sum_{S\in\mathcal S_\vartheta(\pgrG)}\Psi_\vartheta(S)\,\mathsf P_{l,g}\bigl(h_{0*}(S)y',
h_{1*}(S)\varphi\bigr),
\]
where $\mathcal S_\vartheta(\pgrG)$ is the finite set of edge tuples that vanish off $\pgrE(\pgrG)$.

\begin{definition}[Proportionality, transfer]\label{pgr:def-transfer}
Let $\Xi$ be a complex function of period triples. $\Psi$ is \emph{proportional to $\Xi$ on
$\pgrG$}, written $\Psi\propto_{\pgrG}\Xi$, if for every real point $g$ of $\auU(\pgrY^l_l)$ there is
$\beta(g)\in\C$ with $I_{\pgrG}(\Psi;g)(\xi)=\beta(g)\Xi(\xi)$ for all period triples $\xi$. For two-grid
paths $\pgrG',\pgrG$, a function $k\ge1$ with integer values on the $m$-fold twists and a function
$c$ with values in $\C^\times$, we say that \emph{$(\pgrG',\pgrG)$ transfers periods with bound $k$
and constant $c$} if for every function $N$ and every family $\Psi'$ such that each
$\Psi'_\vartheta$ is a path datum on $\pgrE(\pgrG')$ with bound $N(\vartheta)$, there is a family
$\Psi$ such that each $\Psi_\vartheta$ is a path datum on $\pgrE(\pgrG)$ with bound
$k(\vartheta)N(\vartheta)$, and $\Psi\propto_{\pgrG}\Xi$ implies
$\Psi'\propto_{\pgrG'}\Xi/c$ (the function $(\vartheta,y',\varphi)\mapsto\Xi(\vartheta,y',\varphi)/
c(\vartheta)$) for every $\Xi$.
\end{definition}

Put $V(\vartheta)=(\#R_{\pcm(\vartheta)})^{2n_m}$ (a bound for the number of column vectors over
$R_\pcm$ at a node) and $\mathcal E(\vartheta)=(\#R_{\pcm(\vartheta)})^{(2n_m)^2}$ (a bound for the
number of matrices of an edge). For $0\le i,j<m$ let $\rho'_{ij}$ be the ramified ratio
\eqref{pc:eq-rhoprime}.

\paragraph*{The intermediate statements}
The proof of Theorem~\ref{pc:auxcore} consists of the following statements. Those of
\S\S\ref{pgr:ss-paths}--\ref{pgr:ss-deduce} are proved below; the others are stated here exactly and
proved in the later sections cited after each of them (see \S\ref{pgr:ss-status}).

\begin{proposition}[First path]\label{pgr:first}
There are a family $\Psi^0=(\Psi^0_\vartheta)$, each $\Psi^0_\vartheta$ a path datum on
$\pgrE(\pgrfirst)$ with bound $\mathcal E(\vartheta)^{2m}$, a real point $g_0$ of
$\auU(\pgrY^0_0)$, a number $\alpha_{\mathrm{first}}\in\C$ and a surjective linear map
$\mathsf r\colon\pcFF\to\auAz$, $\pcFF$ the space of forms of type $e_{(1,0,m)}$ on $\pgrY^m_0=\pgrY$
fixed by $J(\pcb)$,
such that for every period triple $(\vartheta,y',\varphi)$ with $\varphi\in\pcFF$
\[
\int_{[\auU(Y_m)]}y'\cdot X_{0,m}^\vartheta\,(\mathsf r\varphi)=\alpha_{\mathrm{first}}\cdot
I_{\pgrfirst}(\Psi^0;g_0)(\vartheta,y',\varphi),
\]
where $X^\vartheta_{0,m}=w^\vartheta_{m-1}a_{0,m-1}\circ\dots\circ w^\vartheta_0a_{0,0}$ is the chain of
\S\ref{pc:ss-triangle}, so that $X^\vartheta_{0,m}|_{\auAz}=\iota_\vartheta$ by \eqref{pc:eq-iota}.
\end{proposition}
\noindent Proposition~\ref{pgr:first} is proved in \S\ref{pgw:ss-first} (Theorem~\ref{pgw:first}) from
three inputs stated there: (I) the injectivity of the untwisted chain of the first column of the
second grid, which is Theorem~\ref{pgy:I}; (V) the deep-level vanishing, which is
Theorem~\ref{pgq:V}; and (P) the compatibility of data under the change of modulus $\pcb\to\pcm$,
which is Proposition~\ref{pgy:P}.

\begin{proposition}[One turn]\label{pgr:turn}
For $0\le i,j<m$ there is $c_{ij}$, with values in $\C^\times$ on the $m$-fold twists and of modulus
$\rho'_{ij}$ up to norm powers (Definition~\ref{pc:modulus}), such that every pair
$(\pgrG',\pgrG)$ with $\pgrG'\pgrTr_{(i,j)}\pgrG$ transfers periods with bound $V$ and constant
$c_{ij}$.
\end{proposition}
\noindent Proposition~\ref{pgr:turn} is proved in \S\ref{pgr:ss-off} from
Propositions~\ref{pgr:square} (the square law on a square of a grid) and~\ref{pgr:turndiag} (the
turn on the diagonal), which are stated there.

\begin{proposition}[Last path]\label{pgr:last}
Let $\Xi_{\mathrm{last}}(\Psi)(\vartheta,y',\varphi)=\Lambda(\Psi_\vartheta;y',\varphi)/
\#\auGL_{r+mt}(R_{\pcm_\vartheta})$, where $\Lambda$ is the sum \eqref{pgw:eq-Lambda}: over
$S\in\mathcal S_\vartheta(\pgrlast)$, $\Psi_\vartheta(S)$ times the class integral over
$\Gamma\backslash\auU(Y^m_m)(\auA_f)$, for the measure giving the level of $Y^m_m$ the volume $1$, of
the branching pairing of the algebraic lifts (each normalised by the level of its source node) of
$y'$ and $\varphi$ along the two halves of $\pgrlast$ with the point data $S$, the second evaluated
at $h\oplus\pgupt_\pcm^m$. Then $\Psi\propto_{\pgrlast}\Xi_{\mathrm{last}}(\Psi)$ for every family
$\Psi$.
\end{proposition}
\noindent Proposition~\ref{pgr:last} is proved in \S\ref{pgw:ss-last} (Theorem~\ref{pgw:last}).

\begin{proposition}[Denominators and conjugates on the last path]\label{pgr:den}
There is an integer $N_0\ge1$ with $N_0\le\exp(x^{99}b^2)$ such that, with
$L(\vartheta)=V(\vartheta)^{m^2}\mathcal E(\vartheta)^{2m}$:
\begin{enumerate}
\item[(a)] if $\Psi_\vartheta$ is a path datum on $\pgrE(\pgrlast)$, the pair $(y',\varphi)$ is
fixed by the levels of the two outer nodes, $\iota\in\auI$ with $\ell(\iota)\nmid\auNm Q_\vartheta$,
and both $y'$ and $\varphi$, read as algebraic forms, have gauge at most $1$ at $\iota$ at every
point, then $N_0\Lambda(\Psi_\vartheta;y',\varphi)$ is integral at $\iota$;
\item[(b)] if $\vartheta=1$, $\Psi_1$ is a path datum on $\pgrE(\pgrlast)$ with bound $L(1)$, the
pair is fixed by the outer levels and $\|\sigma(y'(g))\|^2\le B_0$, $\|\sigma(\varphi(g))\|^2\le
B_1$ for all $g$, then $|\sigma(\Lambda(\Psi_1;y',\varphi))|\le\exp(x^{200}b^2)\sqrt{B_0B_1}$.
\end{enumerate}
\end{proposition}
\noindent Here $\Lambda$ is the sum \eqref{pgw:eq-Lambda}, formed with the integral choices of the
kernels and of the branching form of Definition~\ref{pgz:def-choice} (Theorem~\ref{pgw:last} holds for
them, with another constant $\beta(g)$). Part (a) is proved in \S\ref{pgz:ss-a}
(Theorem~\ref{pgz:a}), with $N_0=\mathsf M\,\mathsf d^{2m}$; the class sums in $\Lambda$ carry no
denominators (Proposition~\ref{pgz:expand}, Remark~\ref{pgz:rem-den}). Part (b) is proved in
Theorem~\ref{pgz:b} from the size bounds (B1)--(B3) of \S\ref{pgz:ss-b}, which are proved in
\S\ref{pgt:sec}.

\begin{proposition}[Values read by delta forms]\label{pgr:delta}
Let $\lambda'$ be the weight of the conjugate forms on $Y_m$ dual to the weight of
$\auAm(\vartheta)$, $\langle\ ,\ \rangle_{\mathrm{du}}$ the integral duality pairing between them,
and $h\ge1$ its hook number. For a finite adelic point $x_f$ of $\auU(Y_m)$ and a vector $v$ of the
dual weight let $\delta_{x_f,v}$ be the conjugate delta form at $x_f$ with value $v$. Then:
\begin{enumerate}
\item[(a)] $[K_m:J^\circ_\vartheta]\int_{[\auU(Y_m)]}\delta_{x_f,v}\cdot X^\vartheta_{0,m}y=
\langle v,\iota_\vartheta y(x_f)\rangle_{\mathrm{du}}/\dim\lambda'$ for $y\in\auAz$;
\item[(b)] $\delta_{x_f,v}$ is fixed by the outer level; if $x_f$ is integral above $\ell(\iota)$ (in the sense of {\cite[Corollary~6.4]{P49}}) and $v$
is integral, $\delta_{x_f,v}$ has gauge at most $1$ at $\iota$; and
$\|\sigma(\delta_{x_f,v}(g))\|^2\le\|\sigma(v)\|^2$;
\item[(c)] a vector $u$ with $\langle w,u\rangle_{\mathrm{du}}\in h\cdot\mathcal O_\iota$ for all
integral $w$ is integral at $\iota$; and if $|\sigma\langle v,u\rangle_{\mathrm{du}}|\le K\|\sigma v\|$
for all $v$ then $\|\sigma u\|^2\le C_{\mathrm{du}}K^2$, with $h\,C_{\mathrm{du}}^{1/2}\exp(x^{99}b^2)
\le\exp(x^{201}b^2-x^{200}b^2)$.
\end{enumerate}
\end{proposition}
\noindent Proposition~\ref{pgr:delta} is proved in \S\ref{pgs:ss-delta} (Proposition~\ref{pgs:delta}),
with $C_{\mathrm{du}}=1$.

\begin{proposition}[Short forms of deep level]\label{pgr:short}
$\pcFF$ is spanned by forms that are integral at every $\iota\in\auI$ and whose values have all
conjugates of squared norm at most $\exp(x^{99}b^2)$.
\end{proposition}

\subsubsection{Paths through two grids}\label{pgr:ss-paths}

For a grid path $h$ let $\pgrsq(h)$ be the set of unit squares, named by their corner with least
indices, between the first column and the steps down of $h$: a step down from $(a,b+1)$ to $(a,b)$
contributes the squares $(i,b)$, $i<a$. For a two-grid path $\pgrG=(l,h_0,h_1)$ put
\[
\pgrsq(\pgrG)=\{(i,j):i,j<l\}\cup\pgrsq(h_0)\cup\{(j,i):(i,j)\in\pgrsq(h_1)\}.
\]
Then $\pgrsq(\pgrfirst)=\emptyset$ (no square lies left of the first column) and
$\pgrsq(\pgrlast)=\{0,\dots,m-1\}^2$ (the halves have no step down and $l=m$).

\begin{lemma}\label{pgr:sqturn}
If $\pgrG\pgrTr_s\pgrG'$, then $s\notin\pgrsq(\pgrG)$, $\pgrsq(\pgrG')=\pgrsq(\pgrG)\cup\{s\}$, and
$s\in\{0,\dots,m-1\}^2$.
\end{lemma}
\begin{proof}
(off) Let $h_z$ have the corner $(a,b)$, so $h_z=q\cdot(a,b+1)\to(a,b)\to(a+1,b)\cdot r$ with $q$ a
path from $(0,m)$ to $(a,b+1)$ and $r$ one from $(a+1,b)$ to $(l,l)$. As $r$ only goes down and
right, $a+1\le l\le b$. A step down from $(a',b'+1)$ to $(a',b')$ contributes squares in the row
$b'$ only, and each row is descended at most once by a path; so the squares of $h_z$ in row $b$
are exactly the $(i,b)$, $i<a$, contributed by the step $(a,b+1)\to(a,b)$, and
$(a,b)\notin\pgrsq(h_z)$. After the turn the step down in row $b$ is $(a+1,b+1)\to(a+1,b)$,
contributing $(i,b)$, $i\le a$, and the other steps down are unchanged: exactly one new square,
$(a,b)$. Next, every square $(i,j)$ of a half of $\pgrG$ satisfies $j\ge l$ (a path from $(0,m)$ to
$(l,l)$ descends only rows $\ge l$) and $i<l$ (the first index along the path is at most $l$, and
a square has first index less than that of its step). Since $a<l$, the pair $(b,a)$ is not a square
of either half; since $b\ge l$, neither $(a,b)$ nor $(b,a)$ lies in $\{(i,j):i,j<l\}$. This gives
$s\notin\pgrsq(\pgrG)$ and $\pgrsq(\pgrG')=\pgrsq(\pgrG)\cup\{s\}$ in both cases $z=0,1$ (for $z=1$
after the transposition $(a,b)\mapsto(b,a)$ in the definition of $\pgrsq(\pgrG)$). Finally $a<l\le
b<m$ since the step from $(a,b+1)$ exists and $b+1\le m$.

(diag) Here $\pgrG=(j,h_0,h_1)$ with both halves ending with $(j,j+1)\to(j,j)$. The square $(j,j)$
is not in $\{i,i'<j\}$, and not a square of a half nor its mirror image, since the squares of a half
of $\pgrG$ have first index $<j$. After the turn, $l=j+1$, so $\{i,i'<j+1\}$ gains the squares
$(j,j)$, $(i,j)$ and $(j,i)$ for $i<j$; on the other hand the last step down of each half, which
contributed the squares $(i,j)$, $i<j$, is removed. So $\pgrsq(h'_0)=\pgrsq(h_0)\smallsetminus\{(i,j):
i<j\}$, $\pgrsq(h'_1)=\pgrsq(h_1)\smallsetminus\{(i,j):i<j\}$, and
\[
\pgrsq(\pgrG')=\{i,i'<j\}\cup\{(j,j)\}\cup\{(i,j),(j,i):i<j\}\cup\bigl(\pgrsq(h_0)\cup\pgrsq(h_1)^{\mathrm t}
\smallsetminus\{(i,j),(j,i):i<j\}\bigr)=\pgrsq(\pgrG)\cup\{(j,j)\},
\]
where $^{\mathrm t}$ denotes the mirror image. As the step from $(j,j+1)$ exists, $j+1\le m$.
\end{proof}

\begin{lemma}\label{pgr:exturn}
Every two-grid path $\pgrG\ne\pgrlast$ admits a turn $\pgrG\pgrTr_s\pgrG'$.
\end{lemma}
\begin{proof}
A grid path from $(a_0,b_0)$ to $(a,b)$ either has a corner, or ends with a step down, or has no
step down (then $b=b_0$): by induction on the length, appending a step right to a path ending with a
step down creates a corner, and appending a step right to a path without step down keeps it so. If
a half of $\pgrG=(l,h_0,h_1)$ has a corner, (off) applies. Otherwise, if $l=m$, both halves go from
$(0,m)$ to $(m,m)$ without a step down, so they are the row and $\pgrG=\pgrlast$. If $l<m$, each half
must descend from row $m$ to row $l$, so it is not without step down and, having no corner, ends with
a step down $(l,l+1)\to(l,l)$; then (diag) applies.
\end{proof}

\begin{proposition}[Turn induction]\label{pgr:induction}
Let $\mathsf M(\pgrG,T)$ be a property of two-grid paths $\pgrG$ and finite sets $T$ of pairs. If
$\mathsf M(\pgrlast,\emptyset)$ holds, and $\mathsf M(\pgrG',T)$ implies
$\mathsf M(\pgrG,T\cup\{s\})$ whenever $\pgrG\pgrTr_s\pgrG'$ and $s\notin T$, then
$\mathsf M(\pgrfirst,\{0,\dots,m-1\}^2)$ holds.
\end{proposition}
\begin{proof}
Put $D(\pgrG)=\{0,\dots,m-1\}^2\smallsetminus\pgrsq(\pgrG)$. We show $\mathsf M(\pgrG,D(\pgrG))$
for every $\pgrG$ by induction on $\#D(\pgrG)$; for $\pgrfirst$ it is the claim. If
$\pgrG=\pgrlast$, $D(\pgrG)=\emptyset$. Otherwise take a turn $\pgrG\pgrTr_s\pgrG'$
(Lemma~\ref{pgr:exturn}). By Lemma~\ref{pgr:sqturn}, $s\in\{0,\dots,m-1\}^2\smallsetminus
\pgrsq(\pgrG)$ and $D(\pgrG)=D(\pgrG')\cup\{s\}$ with $s\notin D(\pgrG')$, so
$\#D(\pgrG')<\#D(\pgrG)$, $\mathsf M(\pgrG',D(\pgrG'))$ holds by induction, and the hypothesis gives
$\mathsf M(\pgrG,D(\pgrG))$.
\end{proof}

\subsubsection{Transfer of periods from the first path to the last}\label{pgr:ss-transfer}

\begin{lemma}\label{pgr:transfer-alg}
\begin{enumerate}
\item[(a)] $(\pgrG,\pgrG)$ transfers periods with bound $1$ and constant $1$.
\item[(b)] If $(\pgrG',\pgrG)$ transfers periods with bound $k$ and constant $c$, it does so with
every bound $k'\ge k$.
\item[(c)] If $(\pgrG_1,\pgrG_2)$ transfers periods with bound $k$ and constant $c$, and
$(\pgrG_2,\pgrG_3)$ with bound $k'$ and constant $c'$, then $(\pgrG_1,\pgrG_3)$ transfers periods
with bound $kk'$ and constant $cc'$.
\end{enumerate}
\end{lemma}
\begin{proof}
(a) Take $\Psi=\Psi'$. (b) A path datum with bound $kN$ has bound $k'N$. (c) Let $\Psi^1$ be path
data on $\pgrE(\pgrG_1)$ with bounds $N$. The first transfer gives $\Psi^2$ on $\pgrE(\pgrG_2)$ with
bounds $kN$; the second, applied with the bounds $kN$, gives $\Psi^3$ on $\pgrE(\pgrG_3)$ with
bounds $k'kN$. If $\Psi^3\propto_{\pgrG_3}\Xi$, then $\Psi^2\propto_{\pgrG_2}\Xi/c'$, hence
$\Psi^1\propto_{\pgrG_1}(\Xi/c')/c=\Xi/(cc')$.
\end{proof}

\begin{proposition}[First to last]\label{pgr:firstlast}
With the constants $c_{ij}$ of Proposition~\ref{pgr:turn}, $c=\prod_{0\le i,j<m}c_{ij}$ has modulus
$\prod_{0\le i,j<m}\rho'_{ij}$ up to norm powers, and $(\pgrfirst,\pgrlast)$ transfers periods with
bound $V^{m^2}$ and constant $c$.
\end{proposition}
\begin{proof}
The modulus follows from Lemma~\ref{pc:mod}(a). Apply Proposition~\ref{pgr:induction} to the
property ``$(\pgrG,\pgrlast)$ transfers periods with bound $V^{\#T}$ and constant
$\prod_{s\in T}c_s$''. For $(\pgrlast,\emptyset)$ it is Lemma~\ref{pgr:transfer-alg}(a). If
$\pgrG\pgrTr_s\pgrG'$, $s\notin T$, and the property holds for $(\pgrG',T)$, then $s\in\{0,\dots,
m-1\}^2$ by Lemma~\ref{pgr:sqturn}, $(\pgrG,\pgrG')$ transfers with bound $V$ and constant $c_s$ by
Proposition~\ref{pgr:turn}, and Lemma~\ref{pgr:transfer-alg}(c) gives the property for
$(\pgrG,T\cup\{s\})$, with bound $V\cdot V^{\#T}=V^{\#(T\cup\{s\})}$.
\end{proof}

\subsubsection{Proof of Theorem~\ref{pc:auxcore} from the intermediate statements}\label{pgr:ss-deduce}

Theorem~\ref{pc:auxcore} is to be read with $\pcFF$ the space of forms of level $J(\pcb)$ on
$\pgrY$, of the weight $\tilde\lambda$ of the outer node of the second grid, and with ``$\varphi$
integral at $\iota$'' meaning that $\varphi$, read as an algebraic form, has gauge at most $1$ at
$\iota$ at every point.

\begin{lemma}[The two levels on $Y_m$]\label{pgr:jcirc}
Let $\vartheta$ be an $m$-fold twist, $J_\vartheta$ and $\chi_\vartheta$ the level and the character of
$\auAm(\vartheta)$ ({\cite[\S14.3]{P46}}, {\cite[Lemma~8.5]{P49}}), and
$J^\circ_\vartheta=J_{Y_m}(\pcm_\vartheta)$ the kernel of the reduction $K_m\to\auGL_{r+mt}
(R_{\pcm_\vartheta})$ (\S\ref{pc:ss-data}).
\begin{enumerate}
\item[(a)] $[K_m:J^\circ_\vartheta]=\#\auGL_{r+mt}(R_{\pcm_\vartheta})$.
\item[(b)] $J^\circ_\vartheta\subseteq J_\vartheta$, $\chi_\vartheta$ is trivial on $J^\circ_\vartheta$,
and every element of $\auAm(\vartheta)$, in particular $\iota_\vartheta y$ for $y\in\auAz$, is right
invariant under $J^\circ_\vartheta$. The group $J^\circ_\vartheta$ is free.
\item[(c)] At $w\in Q_\vartheta$ the reduction of $J_\vartheta$ is the block upper triangular
subgroup of $\auGL_{r+mt}(\auF_w)$ with diagonal blocks of sizes $r,t,\dots,t$, so that
$[J_\vartheta:J^\circ_\vartheta]=C\prod_{w\in Q_\vartheta}g_{w,r}\,g_{w,t}^{\,m}\,q_w^{e_m}$ with
$g_{w,N}=\#\auGL_N(\auF_w)$, $e_m=mrt+\frac12m(m-1)t^2$ and $C\ge1$ an integer independent of
$\vartheta$. In particular $J_\vartheta\ne J^\circ_\vartheta$ as soon as $Q_\vartheta\ne\emptyset$.
\end{enumerate}
\end{lemma}
\begin{proof}
(a) The reduction $K_m\to\auGL_{r+mt}(R_{\pcm_\vartheta})$ is onto (\S\ref{pc:ss-data}), with kernel
$J^\circ_\vartheta$. (b) An element of $J^\circ_\vartheta$ is $\equiv1$ modulo $\pi_M$, hence at the
chosen places, and its reduction at $w\in Q_\vartheta$ is the identity, which is block upper
triangular with trivial diagonal blocks; so it lies in $J_\vartheta$, and $\chi_\vartheta$, a product
of tame characters of the determinants of the diagonal blocks ({\cite[Lemma~8.5]{P49}}), is $1$ on
it. A form $\psi$ of level $(J_\vartheta,\chi_\vartheta)$ satisfies $\psi(gk)=\chi_\vartheta(k)\psi(g)$
for $k\in J_\vartheta$, hence $\psi(gk)=\psi(g)$ for $k\in J^\circ_\vartheta$. A subgroup of a free level is free.
(c) The level $J_\vartheta$ consists of the elements of $K_m$ that are $\equiv1$ at the chosen places
and block upper triangular modulo $Q_\vartheta$ ({\cite[\S14.3]{P46}}); {\cite[Lemma~8.5]{P49}}
shows that $\iota_\vartheta z$ is equivariant under this group: a $k\in K_m$ whose reduction
$\bar k_w$ at $w\in Q_\vartheta$ preserves the flag of frames of $Y_0\subset Y_1\subset\dots\subset
Y_m$ satisfies $\bar k_{i+1}\iota_i=\iota_i\bar k_i$ for its upper left blocks, which is the
condition of {\cite[Lemma~5.5]{P49}}. By (a) and the Chinese remainder theorem, $[J_\vartheta:
J^\circ_\vartheta]$ is the order of the reduction of $J_\vartheta$ in $\prod_{w\mid\pcm_\vartheta}
\auGL_{r+mt}(\auF_w)$; the factor at the chosen places does not depend on $\vartheta$, and at
$w\in Q_\vartheta$ it is the order $g_{w,r}g_{w,t}^{\,m}q_w^{e_m}$ of the parabolic subgroup, $e_m$
being the dimension of its unipotent radical.
\end{proof}
\noindent Lemma~\ref{pgr:jcirc}(c) is the index used in Lemma~\ref{pc:index}, where the inner product
of $\auAm(\vartheta)$ (normalised by $J_\vartheta$) enters the chain determinant. Below, the delta forms
and the integrals over $[\auU(Y_m)]$ are normalised by $J^\circ_\vartheta$; since
$[K_m:J_\vartheta]/[K_m:J^\circ_\vartheta]=C^{-1}\prod_w(g_{w,r}g_{w,t}^mq_w^{e_m})^{-1}$ is not a norm
power, the two normalisations must not be interchanged.

\begin{proposition}\label{pgr:lastchain}
There are a function $c$
with values in $\C^\times$, of modulus $\prod_{i,j<m}\rho'_{ij}$ up to norm powers, a family $\Psi$
with $\Psi_\vartheta$ a path datum on $\pgrE(\pgrlast)$ with bound $L(\vartheta)$, a number
$\beta\in\C$ and the surjection $\mathsf r$ of Proposition~\ref{pgr:first}, such that for every
period triple with $\varphi\in\pcFF$
\[
[K_m:J^\circ_\vartheta]\int_{[\auU(Y_m)]}y'\cdot X^\vartheta_{0,m}(\mathsf r\varphi)=\beta\,
\frac{\Lambda(\Psi_\vartheta;y',\varphi)}{c(\vartheta)} .
\]
\end{proposition}
\begin{proof}
Let $\Psi^0,g_0,\alpha_{\mathrm{first}},\mathsf r$ be as in Proposition~\ref{pgr:first}, with bounds
$\mathcal E^{2m}$. By Proposition~\ref{pgr:firstlast} there is $\Psi$ on $\pgrE(\pgrlast)$ with
bounds $V^{m^2}\mathcal E^{2m}=L$, and, since $\Psi\propto_{\pgrlast}\Xi_{\mathrm{last}}(\Psi)$
(Proposition~\ref{pgr:last}), $\Psi^0\propto_{\pgrfirst}\Xi_{\mathrm{last}}(\Psi)/c$. At $g_0$ this
gives $\beta_0$ with $I_{\pgrfirst}(\Psi^0;g_0)(\xi)=\beta_0\Lambda(\Psi_\vartheta;y',\varphi)/
(\#\auGL_{r+mt}(R_{\pcm_\vartheta})\,c(\vartheta))$, and $[K_m:J^\circ_\vartheta]=
\#\auGL_{r+mt}(R_{\pcm_\vartheta})$ by Lemma~\ref{pgr:jcirc}(a). With
Proposition~\ref{pgr:first}, $\beta=\alpha_{\mathrm{first}}\beta_0$ serves.
\end{proof}

\begin{proof}[Proof of Theorem~\ref{pc:auxcore}\textup{(a)}, \textup{(b)} from
Propositions~\ref{pgr:first}--\ref{pgr:delta}]
Let $c,\Psi,\beta,\mathsf r$ be as in Proposition~\ref{pgr:lastchain}, and $N_0$ as in
Proposition~\ref{pgr:den}. By Proposition~\ref{pgr:delta}(a) applied to $y=\mathsf r\varphi$, for
every $x_f$ and $v$,
\begin{equation}\label{pgr:eq-value}
\Bigl\langle v,\frac{c(\vartheta)}{c(1)}\iota_\vartheta(\mathsf r\varphi)(x_f)\Bigr\rangle_{\mathrm{du}}=
\beta'\,\Lambda(\Psi_\vartheta;\delta_{x_f,v},\varphi),\qquad\beta'=\frac{\dim\lambda'\cdot\beta}{c(1)}.
\end{equation}
Choose $s\in\C^\times$ and $\varepsilon\in\{0,1\}$ with $s\beta'=h\,\varepsilon N_0$: if
$\beta'=0$ take $s=1$, $\varepsilon=0$; otherwise $s=hN_0/\beta'$, $\varepsilon=1$. Then
\[
\Bigl\langle v,\frac{c(\vartheta)}{c(1)}\iota_\vartheta(s\,\mathsf r\varphi)(x_f)\Bigr\rangle_{\mathrm{du}}
=h\,\varepsilon\,N_0\Lambda(\Psi_\vartheta;\delta_{x_f,v},\varphi).
\]

(a) Let $\ell(\iota)\nmid\auNm Q_\vartheta$ and $\varphi$ be integral at $\iota$. Integrality of an
algebraic form at $\iota$ is, by definition ({\cite[\S6.1]{P49}}, (N2)), the bound $\le1$ on the
gauge at $(\iota,x_f)$ of its value at $x_f$, for every $x_f\in\auU(Y_m)(\auA_f)$. It suffices to check
this at points $x_f$ integral above $\ell=\ell(\iota)$: for any $x_f$, {\cite[Corollary~6.4(a)]{P49}}
(with $d=1$) gives $\gamma\in\Gamma$ with $\gamma x_f$ integral above $\ell$, and for an algebraic
form $\psi$ one has $\psi(\gamma x_f)=\gamma\psi(x_f)$, so by (G1) of {\cite[\S6.1]{P49}} the gauge
of $\psi(x_f)$ at $(\iota,x_f)$ equals the gauge of $\psi(\gamma x_f)$ at $(\iota,\gamma x_f)$. For such
$x_f$ and integral $v$, $\delta_{x_f,v}$ is fixed by the outer level and has gauge at most $1$
(Proposition~\ref{pgr:delta}(b)); $\varphi\in\pcFF$ is fixed by the level of $\pgrY^m_0$. So
$N_0\Lambda(\Psi_\vartheta;\delta_{x_f,v},\varphi)$ is integral at $\iota$ by
Proposition~\ref{pgr:den}(a), the right side lies in $h\mathcal O_\iota$, and
Proposition~\ref{pgr:delta}(c) shows that $(c(\vartheta)/c(1))\iota_\vartheta(s\mathsf r\varphi)(x_f)$
is integral at $\iota$.

(b) Let $\vartheta=1$ and $\|\sigma(\varphi(g))\|^2\le B$ for all $g$. By
Proposition~\ref{pgr:delta}(b) the delta form satisfies $\|\sigma\delta_{x_f,v}(g)\|^2\le
\|\sigma v\|^2$, and by Proposition~\ref{pgr:den}(b)
$|\sigma\Lambda(\Psi_1;\delta_{x_f,v},\varphi)|\le\exp(x^{200}b^2)\|\sigma v\|\sqrt B$. As
$|\sigma\varepsilon|\le1$, $N_0\le\exp(x^{99}b^2)$ and $h$ is a positive integer,
\[
\bigl|\sigma\langle v,\iota_1(s\mathsf r\varphi)(x_f)\rangle_{\mathrm{du}}\bigr|\le h\exp(x^{99}b^2)
\exp(x^{200}b^2)\sqrt B\,\|\sigma v\| .
\]
By Proposition~\ref{pgr:delta}(c) with $K=h\exp((x^{99}+x^{200})b^2)\sqrt B$,
\[
\|\sigma(\iota_1(s\mathsf r\varphi)(x_f))\|^2\le C_{\mathrm{du}}h^2\exp(2(x^{99}+x^{200})b^2)B\le
\exp(2x^{201}b^2)\,B .
\]
Since $x_f$ was an arbitrary finite adelic point and algebraic forms are functions on
$\auU(Y_m)(\auA_f)$, this is the bound of Theorem~\ref{pc:auxcore}(b) for all $g$. Finally $c$ has modulus $\prod\rho'_{ij}$ up to norm powers by Proposition~\ref{pgr:lastchain}, and
$\mathsf r$ is onto.
\end{proof}

\begin{proof}[Proof of Theorem~\ref{pc:auxcore}\textup{(c)}]
This is Proposition~\ref{pgr:short}, proved in \S\ref{pgs:ss-short} (Proposition~\ref{pgs:short}).
\end{proof}

\subsubsection{Turns off the diagonal}\label{pgr:ss-off}

Throughout, $\vartheta$ is an $m$-fold twist, $R=R_{\pcm(\vartheta)}$, and for finite sets
$\mathsf X,\mathsf Y$ we write $\mathrm M_{\mathsf X\times\mathsf Y}$ for the matrices over $R$ with
rows indexed by $\mathsf X$ and columns by $\mathsf Y$.

\paragraph*{Matched functions}
Let $\mathsf N$ be a finite set with a bijection $\varepsilon\colon\mathsf B\sqcup\mathsf A\to
\mathsf N$, and $\chi$ a character of $R^\times$ (extended by $0$). For $\gamma\in\auGL_{\mathsf N}
(R)$ let $y_\gamma\in\mathrm M_{\mathsf N\times\mathsf B}$ be its columns $\varepsilon(\mathsf B)$
and $x_\gamma\in\mathrm M_{\mathsf A\times\mathsf N}$ the rows $\varepsilon(\mathsf A)$ of
$\gamma^{-1}$; matched functions and the square law are as in \S\ref{pc:ss-exchange}.

\begin{lemma}\label{pgr:extend}
A matrix $y\in\mathrm M_{\mathsf N\times\mathsf B}$ with a left inverse is $y_\gamma$ for some
$\gamma\in\auGL_{\mathsf N}(R)$.
\end{lemma}
\begin{proof}
$R=\prod_w\auO_E/w^{k_w}$ is a finite product of local rings, and it suffices to treat one factor
$R_w$. If $Ly=1$, then $R_w^{\mathsf N}=y(R_w^{\mathsf B})\oplus\ker L$, and $\ker L$ is a direct
summand of a free module over a local ring, hence free, of rank $\#\mathsf N-\#\mathsf B=\#\mathsf A$
(compare ranks modulo the maximal ideal). A basis of $\ker L$, indexed by $\mathsf A$, together with
the columns of $y$ gives $\gamma$.
\end{proof}

\begin{lemma}\label{pgr:match}
\begin{enumerate}
\item[(a)] If $x_\gamma=x_{\gamma'}$, there is $h\in\auGL_{\mathsf B}(R)$ with $y_{\gamma'}=y_\gamma h$
and $\det\gamma'=\det\gamma\det h$.
\item[(b)] Let $\Psi$ on $\mathrm M_{\mathsf N\times\mathsf B}$ vanish off the matrices $y_\gamma$ and
satisfy $\chi(\det h)\Psi(yh)=\Psi(y)$ for $h\in\auGL_{\mathsf B}(R)$. Then
$\Phi=\mathrm{match}_\chi\Psi$, defined by $\Phi(x_\gamma)=\chi(\det\gamma)\Psi(y_\gamma)$ and $\Phi=0$
off the $x_\gamma$, is well defined, and $\Phi,\Psi$ are matched. Conversely, if $\Phi,\Psi$ are
matched then $\Phi=\mathrm{match}_\chi\Psi$.
\item[(c)] If $\Phi,\Psi$ are matched, $G\in\auGL_{\mathsf N}(R)$ and $k\in\auGL_{\mathsf A}(R)$,
then $x\mapsto\chi(\det G)\Phi(xG)$ and $y\mapsto\Psi(G^{-1}y)$ are matched, and
$\Phi(kx)=\chi(\det k)^{-1}\Phi(x)$.
\end{enumerate}
\end{lemma}
\begin{proof}
(a) Order $\mathsf N$ as $\mathsf B\sqcup\mathsf A$ by $\varepsilon$ and put $u=\gamma^{-1}\gamma'$.
The rows $\mathsf A$ of $u^{-1}=\gamma'^{-1}\gamma$ are $x_{\gamma'}\gamma=x_\gamma\gamma=(0\ 1)$, so
$u^{-1}$, hence $u$, is block upper triangular with lower right block $1$:
$u=\bigl(\begin{smallmatrix}h&*\\0&1\end{smallmatrix}\bigr)$, $h$ invertible, $\det u=\det h$. Then
$y_{\gamma'}=\gamma u\binom10=\gamma\binom h0=y_\gamma h$. (b) Well defined by (a) and the law of
$\Psi$. The support conditions and the identity $\Phi(x_\gamma)=\chi(\det\gamma)\Psi(y_\gamma)$ are
the definition. A function matched with $\Psi$ vanishes off the $x_\gamma$ and has the same values
there. (c) $x_\gamma G=x_{G^{-1}\gamma}$, $y_{G^{-1}\gamma}=G^{-1}y_\gamma$ and
$\det(G^{-1}\gamma)=\det\gamma/\det G$; the support conditions transform alike. For the last claim
let $\gamma'=\gamma\,\mathrm{diag}(1,k^{-1})$: then $x_{\gamma'}=kx_\gamma$, $y_{\gamma'}=y_\gamma$
and $\det\gamma'=\det\gamma/\det k$.
\end{proof}

\begin{lemma}\label{pgr:compl}
Let $\mathsf N=\mathsf N_1\sqcup\mathsf N_2$ and $\gamma\in\auGL_{\mathsf N}(R)$. If the rows
$\mathsf N_2$ of $y_\gamma$ form a matrix with a left inverse, then the columns $\mathsf N_1$ of
$x_\gamma$ form a matrix with a left inverse; likewise with $\mathsf N_1,\mathsf N_2$ exchanged.
\end{lemma}
\begin{proof}
Let $L\in\mathrm M_{\mathsf B\times\mathsf N_2}$ with $L\,(y_\gamma)_{\mathsf N_2}=1$, and let $P\in
\mathrm M_{(\mathsf B\sqcup\mathsf A)\times\mathsf N}$ have the rows $(0\ L)$ (columns $\mathsf N_1$,
$\mathsf N_2$) and $x_\gamma$. Then $P\gamma=\bigl(\begin{smallmatrix}1&*\\0&1\end{smallmatrix}
\bigr)$ (as $x_\gamma y_\gamma=0$ and $x_\gamma\gamma_{\mathsf A}=1$), so $P$ is invertible. The
block $\mathsf N_1\times\mathsf N_1$ of $P^{-1}P=1$ reads $P^{-1}_{\mathsf N_1,\mathsf A}\,
(x_\gamma)_{\mathsf N_1}=1$, because the columns $\mathsf N_1$ of the rows $\mathsf B$ of $P$
vanish.
\end{proof}

\paragraph*{A square and its turn}
Fix $z\in\{0,1\}$ and $a\le b<m$. In the grid $z$ let $V'=(a,b)$, $W_1=(a,b+1)$, $W_2=(a+1,b)$,
$V=(a+1,b+1)$; the \emph{old edges} are $e_1=(V',W_1)$, $e_2=(V',W_2)$ and the \emph{new edges}
$f_1=(W_1,V)$, $f_2=(W_2,V)$. Let $\mathsf B,\mathsf A,\mathsf N_1,\mathsf N_2$ be the frames of
$V',V,W_1,W_2$ and $\mathsf N=\mathsf N_1\sqcup\mathsf N_2$; the ranks satisfy
$\#\mathsf B+\#\mathsf A=\#\mathsf N$ (both are twice the rank of $V'$ plus $2t$ in the first grid,
and likewise in the second), and we fix a bijection $\varepsilon\colon\mathsf B\sqcup\mathsf A\to
\mathsf N$. For an edge tuple $S$ put $y(S)=\binom{S_{e_1}}{S_{e_2}}\in\mathrm M_{\mathsf N\times
\mathsf B}$, $x(S)=(S_{f_1}\ S_{f_2})\in\mathrm M_{\mathsf A\times\mathsf N}$, and let $S[y]$,
$S[x]$ be $S$ with the old, respectively new, edges replaced by the blocks of $y$, respectively $x$.
Put $\chi=(\chi_{e_1}\chi_{e_2})^{-1}$.

The edge characters satisfy $\chi_{f_1}=\chi_{e_2}$ and $\chi_{f_2}=\chi_{e_1}$: $e_1$ and $f_2$
change the second index at $b$, $e_2$ and $f_1$ change the first index at $a$, and the character of
an edge depends only on whether it is vertical and on the index that changes. In the grid $z=0$ the
edges $e_1,f_2$ are vertical and $e_2,f_1$ horizontal; in the grid $z=1$ it is the reverse. If
$a<b$, neither $V'$ nor $V$ is an outer node or a diagonal node, so admissible families are
unconstrained at $V'$ and $V$.

From here to Lemma~\ref{pgr:sumturn} assume $a<b$.
Let $w$ be a set of edges none of which has $V'$ as an end or $V$ as its target.

\begin{proposition}[Turning a path datum]\label{pgr:turned}
Let $\Psi'$ be a path datum on $w\cup\{e_1,e_2\}$ with bound $N$. For an edge tuple $S$ let
$\Psi'_S(y)=\Psi'(S[y])$. Then $\Psi(S)=(\mathrm{match}_\chi\Psi'_S)(x(S))$ defines a path datum on
$w\cup\{f_1,f_2\}$ with bound $N$, and for every $S$ the functions $x\mapsto\Psi(S[x])$ and
$\Psi'_S$ are matched for $\chi$.
\end{proposition}
\begin{proof}
\emph{The match exists.} If $\Psi'_S(y)\ne0$, the block of $y$ on the vertical old edge has a left
inverse, hence so has $y$, and $y=y_\gamma$ by Lemma~\ref{pgr:extend}. For
$k\in\auGL_{\mathsf B}(R)$ let $g'$ be the node family equal to $k$ at $V'$ and to $1$ elsewhere; it
is admissible, it maps $S[y]$ to $S[yk]$ (the edges of $w$ do not meet $V'$), and
$\chi_{w\cup\{e_1,e_2\}}(g')=\chi_{e_1}(\det k)^{-1}\chi_{e_2}(\det k)^{-1}=\chi(\det k)$. The
invariance of $\Psi'$ gives $\chi(\det k)\Psi'_S(yk)=\Psi'_S(y)$. By Lemma~\ref{pgr:match}(b),
$\mathrm{match}_\chi\Psi'_S$ is defined and matched with $\Psi'_S$; since $\Psi'_S$ depends on $S$
only through $(S_e)_{e\in w}$, so does the match, and $\Psi(S[x])=(\mathrm{match}_\chi\Psi'_S)(x)$.
In particular $\Psi$ depends only on the new edges and on $w$.

\emph{Invariance.} Let $g$ be admissible, $G=\mathrm{diag}(g_{W_1},g_{W_2})\in\auGL_{\mathsf N}(R)$,
and let $g''$ be $g$ with $1$ at $V'$ (admissible). Then $g''(S[y])=(gS)[G^{-1}y]$, and the
invariance of $\Psi'$ gives $\Psi'_S(y)=\mathsf a\,\Psi'_{gS}(G^{-1}y)$ with
$\mathsf a=\chi_{e_1}(\det g_{W_1})\chi_{e_2}(\det g_{W_2})\chi_w(g)$. By Lemma~\ref{pgr:match}(c) and
the uniqueness in (b), $(\mathrm{match}_\chi\Psi'_S)(x)=\mathsf a\,\chi(\det G)\,(\mathrm{match}_\chi
\Psi'_{gS})(xG)$. Moreover $x(gS)=g_V^{-1}x(S)G$, so by the last claim of
Lemma~\ref{pgr:match}(c), $\Psi(gS)=\chi(\det g_V)(\mathrm{match}_\chi\Psi'_{gS})(x(S)G)$. Hence
$\Psi(S)=\mathsf a\,\chi(\det G)\chi(\det g_V)^{-1}\Psi(gS)$, and
\[
\mathsf a\,\chi(\det G)\chi(\det g_V)^{-1}=\chi_{e_2}(\det g_{W_1})^{-1}\chi_{e_1}(\det g_{W_2})^{-1}
\chi_{e_1}\chi_{e_2}(\det g_V)\,\chi_w(g)=\chi_{w\cup\{f_1,f_2\}}(g)
\]
by $\chi_{f_1}=\chi_{e_2}$, $\chi_{f_2}=\chi_{e_1}$.

\emph{Left inverses.} Let $\Psi(S)\ne0$. Then $x(S)=x_\gamma$ with $\Psi'_S(y_\gamma)\ne0$. The
edges of $w$ carry the same matrices in $S$ and $S[y_\gamma]$, so their vertical ones are split
injective. If $f_1$ is vertical, so is $e_2$; then the rows $\mathsf N_2$ of $y_\gamma$ are split
injective, and the columns $\mathsf N_1$ of $x_\gamma$, which form $S_{f_1}$, are split injective
by Lemma~\ref{pgr:compl}. The case of $f_2$ is symmetric.

\emph{Values.} $\Psi(S)=\chi(\det\gamma)\Psi'_S(y_\gamma)$ is a root of unity times a sum of at most
$N$ roots of unity.
\end{proof}

\begin{lemma}[Sums over a turned square]\label{pgr:sumturn}
In the situation of Proposition~\ref{pgr:turned}, let $J$ be a set, $C\in\C$, and for $i\in J$ let
$L_i$ on $\mathrm M_{\mathsf A\times\mathsf N}$ and $L'_i$ on $\mathrm M_{\mathsf N\times\mathsf B}$
satisfy the square law for $\chi$ with constant $C$. Let $i(S)\in J$ depend only on
$(S_e)_{e\in w}$. Then
\[
\sum_{S\in\mathcal S(w\cup\{e_1,e_2\})}\Psi'(S)\,L'_{i(S)}(y(S))=C\sum_{S\in\mathcal S(w\cup\{f_1,
f_2\})}\Psi(S)\,L_{i(S)}(x(S)),
\]
$\mathcal S(u)$ being the set of edge tuples vanishing off $u$.
\end{lemma}
\begin{proof}
Since $e_1,e_2\notin w$, $(S,y)\mapsto S[y]$ is a bijection from $\mathcal S(w)\times\mathrm
M_{\mathsf N\times\mathsf B}$ onto $\mathcal S(w\cup\{e_1,e_2\})$, and likewise for the new edges.
For fixed $S\in\mathcal S(w)$ the square law applied to the matched pair of
Proposition~\ref{pgr:turned} gives $\sum_y\Psi'_S(y)L'_{i(S)}(y)=C\sum_x\Psi(S[x])L_{i(S)}(x)$; sum
over $S$.
\end{proof}

\paragraph*{The square law on a square of a grid}
For $f$ a form at $W_1$ (a conjugate form if $z=0$) and $l$ a linear functional on the forms at
$W_2$, let
\[
L_{f,l}(x)=l\bigl(T_{x_{\mathsf N_2}}\Lambda_{x_{\mathsf N_1}}f\bigr),\qquad
L'_{f,l}(y)=l\bigl(\Lambda_{y_{\mathsf N_2}}T_{y_{\mathsf N_1}}f\bigr),
\]
where $\Lambda_M$ and $T_M$ denote the moves (lift and transposed lift) with the point datum $M$
along the edge concerned: $L_{f,l}$ goes through $V$, $L'_{f,l}$ through $V'$.

\begin{proposition}[Square law on a grid square]\label{pgr:square}
For $z\in\{0,1\}$ and $a,b<m$ there are $\kappa\in\C$ and a function $c^z_{ab}$ with values in
$\C^\times$, of modulus $\rho'_{ab}$ if $z=0$ and $\rho'_{ba}$ if $z=1$ up to norm powers, such that
for every $\vartheta$, every $f$ and every $l$ the functions $L_{f,l}$, $L'_{f,l}$ satisfy the square
law for $\chi$ with constant $\kappa/c^z_{ab}(\vartheta)$.
\end{proposition}
\noindent Proposition~\ref{pgr:square} is proved in \S\ref{pgv:ss-square} (Theorem~\ref{pgv:square})
from Theorem~\ref{ptr:square}; the hypothesis (E1) of an edge square holds for the four pairs of a
grid square by Lemma~\ref{pgv:E1g}. The character is $\nu=(\mu^0_{a+1}\mu^0_{b+1}\vartheta^n)^{-1}$,
$n=b+1$ if $z=0$ and $n=a+1$ if $z=1$.

\begin{proposition}[A turn off the diagonal]\label{pgr:turnoff}
Let $z\in\{0,1\}$, $a<b<m$, and let $(\pgrG',\pgrG)$ be a turn
of type (off) at the corner $(a,b)$ of the half $h_z$. Then $(\pgrG',\pgrG)$ transfers periods with
bound $1$ and constant $c^z_{ab}$.
\end{proposition}
\begin{proof}
Write $h_z=q\cdot(V'$-corner$)\cdot r$ as in the definition of (off), and let $w$ be the set of the
edges of $q$, of $r$ and of the other half. The edges of $q$ have both ends in the rectangle
$\{a'\le a,\ b'\ge b+1\}$, those of $r$ in $\{a'\ge a+1,\ b'\le b\}$, and those of the other half in
the other grid; so no edge of $w$ meets $V'$ or ends at $V$, and $\pgrE(\pgrG')=w\cup\{e_1,e_2\}$,
$\pgrE(\pgrG)=w\cup\{f_1,f_2\}$. Let $\Psi'$ be path data on $\pgrE(\pgrG')$ with bounds $N$, and let
$\Psi$ be given by Proposition~\ref{pgr:turned}; it has bounds $N$.

Let $(\vartheta,y',\varphi)$ be a period triple and $g$ a real point. Let $F$ be the starting form of
the half $h_z$ ($y'$ if $z=0$, $\varphi$ if $z=1$). For $S$ put $f_S=q_*(S)F$, a form at $W_1$, and
let $l_S$ be the functional $u\mapsto\mathsf P_{l,g}(r_*(S)u,\,h_{1*}(S)\varphi)$ if $z=0$, and
$u\mapsto\mathsf P_{l,g}(h_{0*}(S)y',\,r_*(S)u)$ if $z=1$. Both depend on $S$ only through $w$. As
$h_{z*}$ for $\pgrG'$ is $r_*\circ\Lambda_{S_{e_2}}\circ T_{S_{e_1}}\circ q_*$ and for $\pgrG$ it is
$r_*\circ T_{S_{f_2}}\circ\Lambda_{S_{f_1}}\circ q_*$,
\[
I_{\pgrG'}(\Psi';g)(\xi)=\sum_S\Psi'_\vartheta(S)\,L'_{f_S,l_S}(y(S)),\qquad
I_{\pgrG}(\Psi;g)(\xi)=\sum_S\Psi_\vartheta(S)\,L_{f_S,l_S}(x(S)).
\]
By Proposition~\ref{pgr:square} and Lemma~\ref{pgr:sumturn} (with $J$ the pairs $(f,l)$),
$I_{\pgrG'}(\Psi';g)(\xi)=\kappa\,c^z_{ab}(\vartheta)^{-1}I_{\pgrG}(\Psi;g)(\xi)$. If
$I_{\pgrG}(\Psi;g)=\beta(g)\Xi$, then $I_{\pgrG'}(\Psi';g)=\kappa\beta(g)\,\Xi/c^z_{ab}$.
\end{proof}

\begin{proposition}[A turn on the diagonal]\label{pgr:turndiag}
For $0\le j<m$ there is a function $c_{jj}$ with values in $\C^\times$, of modulus $\rho'_{jj}$ up to
norm powers, such that every turn $\pgrG'\pgrTr_{(j,j)}\pgrG$ transfers periods with bound $V$ and
constant $c_{jj}$.
\end{proposition}
\noindent Proposition~\ref{pgr:turndiag} is proved in \S\ref{pgu:sec} (Theorem~\ref{pgu:turndiag})
from Proposition~\ref{pgr:square} and the properties (G1)--(G4) of the grid data listed in
\S\ref{pgu:ss-data}, which are proved in \S\ref{pgv:sec}, with the line reversal of
Lemma~\ref{pgx:line}; $c_{jj}=c^0_{jj}/c'_{\mathrm{line}}$.

\begin{proof}[Proof of Proposition~\ref{pgr:turn} from Propositions~\ref{pgr:square}
and~\ref{pgr:turndiag}]
Let $\pgrG'\pgrTr_{(i,j)}\pgrG$. In case (off) with $z=0$ at the corner $(a,b)$ we have $(i,j)=(a,b)$
and $a<b$ (proof of Lemma~\ref{pgr:sqturn}); with $z=1$, $(i,j)=(b,a)$ and $a<b$; in case (diag),
$i=j$. So for $i<j$ the turn is (off) with $z=0$ at $(i,j)$, and we take $c_{ij}=c^0_{ij}$, of
modulus $\rho'_{ij}$; for $i>j$ it is (off) with $z=1$ at $(j,i)$, and $c_{ij}=c^1_{ji}$, of modulus
$\rho'_{ij}$; for $i=j$ we take $c_{jj}$ of Proposition~\ref{pgr:turndiag}. By
Proposition~\ref{pgr:turnoff} and Lemma~\ref{pgr:transfer-alg}(b) (as $V\ge1$) the off-diagonal
turns transfer periods with bound $V$.
\end{proof}

\subsubsection{Where the propositions of this section are proved}\label{pgr:ss-status}
Theorem~\ref{pc:auxcore}(a),(b) is deduced above from Propositions~\ref{pgr:first}--\ref{pgr:delta}, and
Theorem~\ref{pc:auxcore}(c) is Proposition~\ref{pgr:short}. These propositions are proved as follows:
Proposition~\ref{pgr:first} in \S\ref{pgw:ss-first} (Theorem~\ref{pgw:first}), with its inputs in
\S\S\ref{pgy:sec}, \ref{pgq:sec}--\ref{pgm:sec}; Proposition~\ref{pgr:square} in \S\ref{pgv:ss-square}
(Theorem~\ref{pgv:square}, with Lemma~\ref{pgv:E1g}); Proposition~\ref{pgr:turndiag} in \S\ref{pgu:sec}
(Theorem~\ref{pgu:turndiag}), using \S\ref{pgv:sec} and Lemma~\ref{pgx:line}; Proposition~\ref{pgr:last} in
\S\ref{pgw:ss-last} (Theorem~\ref{pgw:last}); Proposition~\ref{pgr:den} in \S\S\ref{pgz:ss-a},
\ref{pgz:ss-b}, \ref{pgt:sec} (Theorems~\ref{pgz:a}, \ref{pgz:b}); Propositions~\ref{pgr:delta} and
\ref{pgr:short} in \S\ref{pgs:sec}. Besides these proofs, Theorem~\ref{pc:auxcore} uses Theorem~\ref{pcx:X}
and the published results named in \S\S\ref{ptr:sec}--\ref{ptw:sec} and \S\S\ref{pgq:sec}--\ref{pgm:sec}.




% ===== C-grid2.tex =====
% proofs/C-grid2.tex
% Continuation of proofs/C-grid.tex (labels pgr:): proofs of the open propositions pgr:delta and
% pgr:short (= pc:auxcore(c)). Labels start with pgs:. Conventions and macros: those of
% proofs/C-grid.tex and proofs/C-determinant.tex.
\providecommand{\auI}{\mathcal I}
\providecommand{\auAut}{\operatorname{Aut}}
\providecommand{\auGL}{\mathrm{GL}}
\providecommand{\auU}{\mathrm U}
\providecommand{\auO}{\mathcal O}
\providecommand{\pcFF}{\mathfrak F}
\providecommand{\pcb}{\mathfrak b}
\providecommand{\auA}{\mathbb A}
\providecommand{\pgsT}{\mathbb T}
\providecommand{\pgsC}{\mathsf C}
\providecommand{\pgsev}{\mathrm{ev}}
\providecommand{\pgsdu}{\mathrm{du}}

\subsection{Lifts through the second grid: Weyl modules, delta forms and short forms}\label{pgs:sec}

This section continues \S\ref{pgr:sec} and uses its notation. It proves
Proposition~\ref{pgr:delta} (values read by delta forms; \S\ref{pgs:ss-delta}), after the facts on
integral Weyl modules that it needs (\S\ref{pgs:ss-weyl}), and Proposition~\ref{pgr:short}, that
is Theorem~\ref{pc:auxcore}(c) (\S\ref{pgs:ss-short}).

\subsubsection{Integral Weyl modules and their duality}\label{pgs:ss-weyl}

All statements are multiplicative over the infinite places $\tau$, and we write them for one
group $\auGL_N$; for a weight $\lambda=(\kappa,e)$ of $\prod_\tau\auGL_N$ one takes tensor products
over $\tau$ (and $\det^{e_\tau}$ at $\tau$). Let $\mu$ be a diagram with columns $\mathfrak c$ of
heights $1\le h_{\mathfrak c}<N$, given by the numbers $\kappa(h)$ of columns of height $h$, and
$|\mu|=\sum_{\mathfrak c}h_{\mathfrak c}$ its number of boxes. We use the model of
{\cite[\S5.2.6]{P49}}: the tensor space $\pgsT_\mu(R)=\bigotimes_{\mathfrak c}\Lambda^{h_{\mathfrak
c}}R^N$ with the orthonormal basis $(e_U)$ indexed by the fillings $U$, the evaluation tensors
$\pgsev_y=\bigotimes_{\mathfrak c}y_1\wedge\dots\wedge y_{h_{\mathfrak c}}$ ($y_i$ the columns of
$y\in\mathrm M_N(R)$), their span $\pgsC_\mu\subseteq\pgsT_\mu(\C)$, the products of minors
$b_U(y)=(e_U,\pgsev_y)$, and the isomorphism $\mathsf m\colon\pgsC_\mu\to W_\mu$,
$\mathsf m(u)(y)=(u,\pgsev_y)$, onto the span $W_\mu$ of the products of minors; $\Sigma_\mu$ is the
$\Z$-span of the $b_U$, with the $\Z$-basis $(b_T)$ of standard products of minors, and
$\Sigma_\mu^\#$ its dual lattice for the invariant form.

\paragraph*{The coordinates of the paper}
The weight $\lambda$ is realised ({\cite[\S14.3]{P46}}) on $W_\lambda$ with the basis $(e_T)$ dual to
the standard products of minors, with integral lattice $W_\lambda(\Z)=\bigoplus\Z e_T$, and with an
invariant hermitian form whose Gram matrix $G_\lambda$ in this basis is a positive definite integer
matrix. We identify $W_\lambda$ with $W_\mu$ (twisted by $\det^{e}$) through the invariant form, so
that $e_T$ is the element of $W_\mu$ with $\langle e_T,b_{T'}\rangle=\delta_{TT'}$. By
{\cite[(51)]{P49}}, $\langle\mathsf m u,b_{U}\rangle=(u,e_U)$ for $u\in\pgsC_\mu$; hence
\begin{equation}\label{pgs:eq-lattice}
W_\lambda(\Z)=\Sigma_\mu^\#=\mathsf m\bigl(\pgsC_\mu\cap\pgsT_\mu(\Z)\bigr),\qquad
\pgsev(y):=\sum_Tb_T(y)\,e_T=\mathsf m(\pgsev_y),
\end{equation}
the second because $\langle\mathsf m\,\pgsev_y,b_T\rangle=(\pgsev_y,e_T)=b_T(y)$. The invariant form
of $W_\lambda$ is the standard form of $\pgsT_\mu$ restricted to $\pgsC_\mu$
({\cite[\S5.2.6]{P49}}); in particular
\begin{equation}\label{pgs:eq-evnorm}
\|\pgsev(y)\|^2=\prod_{\mathfrak c}\det\bigl(y_{[h_{\mathfrak c}]}^*\,y_{[h_{\mathfrak c}]}\bigr)
\le\prod_{\mathfrak c}\ \prod_{i\le h_{\mathfrak c}}\|y_i\|^2 ,
\end{equation}
where $y_{[h]}$ consists of the first $h$ columns of $y$; the inequality is Hadamard's. A vector of
$W_\lambda$ is \emph{integral at $\iota$} if its coordinates in $(e_T)$ lie in $\auO_\iota$.

\begin{lemma}[Gram matrix and Fischer form]\label{pgs:gram}
Let $F_\mu=(\langle b_T,b_{T'}\rangle_F)_{T,T'}$ be the matrix of the Fischer form of
{\cite[Lemma~5.22]{P49}} on the standard products of minors, and $H_\mu$ the hook number of that lemma.
Then $F_\mu$ has integer entries and $G_\lambda F_\mu=F_\mu G_\lambda=H_\mu\cdot1$.
\end{lemma}
\begin{proof}
The entries are integers since the $b_T$ have integer coefficients. By {\cite[Lemma~5.22]{P49}}, the
invariant form of $W_\mu$ is $H_\mu^{-1}\langle\ ,\,\rangle_F$, so the Gram matrix of $(b_T)$ is
$H_\mu^{-1}F_\mu$. The basis $(e_T)$ is dual to $(b_T)$ for this real symmetric form, so its Gram
matrix is the inverse, $G_\lambda=H_\mu F_\mu^{-1}$.
\end{proof}

\begin{lemma}[Duality by complementary minors]\label{pgs:star}
Let $\kappa'$ be the diagram with $\kappa'(h)=\kappa(N-h)$ for $1\le h<N$, $\mu'$ the corresponding
diagram, and $d=\sum_{1\le h<N}\kappa(h)$. There is an integer matrix $J$ from coordinates of
weight $\kappa$ to coordinates of weight $\kappa'$ (in the bases $(e_T)$) which has an integral
inverse, satisfies ${}^tJ\,G_{\mu'}J=G_\mu$, and $S_{\kappa'}({}^tg^{-1})\,J=\det(g)^{-d}\,J\,
S_\kappa(g)$ for $g\in\auGL_N(\C)$, where $S$ denotes the Schur matrices.
\end{lemma}
\begin{proof}
For a subset $U\subseteq[N]$ of $h$ elements let $\star e_U=\epsilon_Ue_{U^c}$ with the sign
$\epsilon_U$ of the permutation $(U,U^c)$ of $[N]$; then $\star\colon\Lambda^hR^N\to\Lambda^{N-h}R^N$
is a signed bijection of bases, characterised by $(\star\omega,\beta)\,e_{[N]}=\omega\wedge\beta$
for $\omega\in\Lambda^h$, $\beta\in\Lambda^{N-h}$ (as $e_U\wedge e_{U^c}=\epsilon_Ue_{[N]}$). For
invertible $g$, $(\star\Lambda^hg\,\omega,\beta)e_{[N]}=g\omega\wedge g(g^{-1}\beta)=\det(g)\,
\omega\wedge g^{-1}\beta=\det(g)\,(\Lambda^{N-h}({}^tg^{-1})\star\omega,\beta)e_{[N]}$, that is
(Jacobi's identity for complementary minors)
$\star\circ\Lambda^hg=\det(g)\,\Lambda^{N-h}({}^tg^{-1})\circ\star$. Let $P=\bigotimes_{\mathfrak c}
\star\colon\pgsT_\mu\to\pgsT_{\mu'}$, the columns of height $h$ going to columns of height $N-h$; $P$
is a signed permutation of the orthonormal bases, so it is orthogonal and preserves the integral
tensors, and $P\circ(\bigotimes\Lambda^{h_{\mathfrak c}}g)=\det(g)^{d}(\bigotimes\Lambda^{N-h_{\mathfrak
c}}{}^tg^{-1})\circ P$. The subspace $\pgsC_\mu$ is the $\auGL_N$-submodule generated by
$\pgsev_1$ (as $g\,\pgsev_y=\pgsev_{gy}$), and $P\pgsev_1=\pm\pgsev_{w_0}$ with $w_0$ the matrix
of the reversal of $[N]$ (each $\star e_{[h]}=\pm e_{\{h+1,\dots,N\}}$, the exterior product of the first $N-h$ columns $e_N,\dots,e_{h+1}$ of $w_0$); so $P$ maps $\pgsC_\mu$ onto
the submodule generated by $\pgsev_{w_0}$, which is $\pgsC_{\mu'}$ ($\dim\pgsC_\mu=\dim\pgsC_{\mu'}$,
both irreducible). As $P$ and $P^{-1}={}^tP$ preserve integral tensors, $P$ maps
$\pgsC_\mu\cap\pgsT_\mu(\Z)$ onto $\pgsC_{\mu'}\cap\pgsT_{\mu'}(\Z)$. Let $J$ be the matrix of
$\mathsf m'\circ P\circ\mathsf m^{-1}$ in the bases $(e_T)$, $(e_{T'})$. By \eqref{pgs:eq-lattice}
it maps $W_\mu(\Z)$ onto $W_{\mu'}(\Z)$, so it and its inverse are integral; it is isometric because
$P$ is orthogonal and the invariant forms are the restrictions of the standard forms; and the
intertwining relation is that of $P$ (for the action of $g$ on $W_\mu$ by $S_\kappa(g)$, the
transpose conventions of {\cite[\S6.1]{P49}} being the same on both sides).
\end{proof}

\begin{lemma}[The duality pairing]\label{pgs:pairing}
Let $\lambda=(\kappa,e)$ and $\lambda'=(\kappa',e')$ be weights of $\prod_\tau\auGL_N$ whose doubled
Harish-Chandra parameters are $l$ and $-l$. Then $\kappa'(h)=\kappa(N-h)$ for $1\le h<N$ and
$e+e'=-d$ at every $\tau$. Put $A=G_{\lambda'}J$, with $J$ the tensor product over $\tau$ of the
matrices of Lemma~\ref{pgs:star}, and $\langle v',v\rangle_{\pgsdu}={}^tv'Av$ for coordinate vectors
$v'$ of weight $\lambda'$ and $v$ of weight $\lambda$. Then:
\begin{enumerate}
\item[(P1)] $\langle\ ,\,\rangle_\pgsdu$ is invariant: $\langle\rho'(g)v',\rho(g)v\rangle_\pgsdu=
\langle v',v\rangle_\pgsdu$ for $g\in\prod_\tau\auGL_N(\C)$;
\item[(P2)] there is an integer matrix $C$ with $CA=h\cdot1$, where $h=H_{\lambda'}=\prod_\tau
H_{\mu'_\tau}$ is the hook number of $\lambda'$;
\item[(P3)] for every $v$ of weight $\lambda$ the vector $v^\flat=\overline{Jv}$ of weight $\lambda'$
satisfies $\|v^\flat\|=\|v\|$ and $\langle a,v\rangle_\pgsdu=\langle v^\flat,a\rangle$ for all $a$,
where $\langle p,q\rangle={}^t\bar pG_{\lambda'}q$ is the invariant hermitian form of weight
$\lambda'$.
\end{enumerate}
\end{lemma}
\begin{proof}
With $\mu_p=e+\sum_{h>p}\kappa(h)$ the highest weight, the doubled parameter is
$(2\mu_p+N-1-2p)_{p<N}$; its negative, reordered decreasingly, is the doubled parameter of the
highest weight $\mu'_p=-\mu_{N-1-p}$. Since $\mu_p-\mu_{p+1}=\kappa(p+1)$, this gives
$\kappa'(h)=\mu'_{h-1}-\mu'_h=\mu_{N-h-1}-\mu_{N-h}=\kappa(N-h)$ and $e'=\mu'_{N-1}=-\mu_0=-e-d$.

(P1) The invariance of the hermitian form $G=G_{\lambda'}$ says ${}^t\overline{\rho'(g)}G\rho'(g)=G$
for $g$ in the compact group $\prod_\tau\auU(N)$; since $\rho'$ has integral (hence real)
coefficients in the entries of $g$ and in $\det g^{\pm1}$, $\overline{\rho'(g)}=\rho'(\bar g)=
\rho'({}^tg^{-1})$, so ${}^t\rho'(g)\,G\,\rho'({}^tg^{-1})=G$ on $\prod\auU(N)$, hence on
$\prod\auGL_N(\C)$ (Zariski density of the unitary group). By Lemma~\ref{pgs:star} and $e+e'=-d$,
$J\rho(g)=J\det(g)^eS_\kappa(g)=\det(g)^{e+d}S_{\kappa'}({}^tg^{-1})J=\rho'({}^tg^{-1})J$. So
${}^t(\rho'(g)v')GJ\rho(g)v={}^tv'\,{}^t\rho'(g)G\rho'({}^tg^{-1})Jv={}^tv'GJv$.

(P2) By Lemma~\ref{pgs:gram} for $\lambda'$, $F_{\lambda'}G_{\lambda'}=h\cdot1$; with $J^{-1}$
integral (Lemma~\ref{pgs:star}), $C=J^{-1}F_{\lambda'}$ gives $CA=J^{-1}F_{\lambda'}G_{\lambda'}J=h$.

(P3) $\langle a,v\rangle_\pgsdu={}^taGJv={}^t(Jv)Ga$ ($G=G_{\lambda'}$ symmetric), which is
$\langle\overline{Jv},a\rangle$. And $\|\overline{Jv}\|^2={}^t\overline v\,{}^tJG_{\lambda'}J\,v=
\|v\|^2$ by the isometry of $J$ (Lemma~\ref{pgs:star}; $J$ is real).
\end{proof}

\begin{remark}
The pairing of Proposition~\ref{pgr:delta} is this $A=G_{\lambda'}J$, and it
measures integrality in the coordinates $(e_T)$; by \eqref{pgs:eq-lattice} this is the lattice
$W_\lambda(\Z)=\Sigma_\mu^\#$, as the plan of \S\ref{pgr:sec} required. Lemma~\ref{pgs:gram} is the
precise form of the inclusion $H_\mu\Sigma_\mu\subseteq\Sigma_\mu^\#$ of {\cite[Lemma~5.22]{P49}} that is
used: in the coordinates $(e_T)$ the vector $b_{T'}$ has coordinates $G_\lambda^{-1}$ times the
unit vector, and $H_\mu G_\lambda^{-1}=F_\mu$ is integral.
\end{remark}

\subsubsection{Values read by delta forms: proof of Proposition~\ref{pgr:delta}}\label{pgs:ss-delta}

\paragraph*{The weights}
Let $\lambda=(\kappa,e)$ be the weight of $\auAm(\vartheta)$ on $Y_m$ (of rank $N=r+mt$), with
doubled Harish-Chandra parameter $l$, and $\lambda'=(\kappa',e')$ the weight of the conjugate forms
on $Y_m$, the weight of the outer node $Y^0_m$ of the first grid; its doubled parameter is $-l$
(the type of a conjugate form is the negated type). So Lemma~\ref{pgs:pairing} applies, and we let
$\langle\ ,\,\rangle_\pgsdu$ be the pairing $A=G_{\lambda'}J$ of that lemma and $h=H_{\lambda'}$.
The forms on $Y_m$ of weight $\lambda$ and the conjugate forms of weight $\lambda'$ are read from
algebraic forms $\psi$, $\eta$ (of weights $\lambda$, $\lambda'$) as in Lemma~\ref{pc:algpet}:
\[
f_\psi(u,g)=\langle v'^{+},u^{-1}\psi(g)\rangle_\pgsdu,\qquad f_\eta(u,g)=\langle u^{-1}\eta(g),v^{+}
\rangle_\pgsdu\qquad(u\in\auU(Y_m)_\infty,\ g\in G_f),
\]
with $v^+$, $v'^+$ the highest weight vectors of the parameters $l$, $-l$, normalised by
$\langle v'^+,v^+\rangle_\pgsdu=1$. (The reading of conjugate forms is defined through the pairing in
this way, and $J$ maps the evaluation vector at the frame ordered by $l$, which is $v^+$, to the one at
the reversed frame, which is $v'^+$.) For a finite adelic point $x$
and a vector $v$ of weight $\lambda'$, the \emph{delta form} $\delta_{x,v}$ is the algebraic form of
weight $\lambda'$ and level $J^\circ_\vartheta=J_{Y_m}(\pcm_\vartheta)$ (the principal congruence
subgroup of Lemma~\ref{pgr:jcirc}, not the level $J_\vartheta$ of $\auAm(\vartheta)$) supported on
$\Gamma xJ^\circ_\vartheta$ with $\delta_{x,v}(\gamma xk)=\gamma v$; it is well defined because
$J^\circ_\vartheta$ is free ({\cite[Lemma~14.5]{P46}}: it consists of elements $\equiv1$ modulo the
chosen places, and the rank is $\le2n_m$). We use the same name for the conjugate form $f_{\delta_{x,v}}$.

\begin{lemma}[Size of the hook number]\label{pgs:hooksize}
$h\le L^{3L}$, where $L=\sum_\tau|\mu'_\tau|\le40f\,n_mb\le x^{19}b$, and
$h\exp(x^{99}b^2)\exp(x^{200}b^2)\le\exp(x^{201}b^2)$.
\end{lemma}
\begin{proof}
$h=\prod_\tau H_{\mu'_\tau}\le\prod_\tau L_\tau^{3L_\tau}\le L^{3L}$ by {\cite[Lemma~5.22]{P49}}. The
bound for $L$ is Lemma~\ref{pgs:boxes} below (with $N\le2n_m$), together with $f\le x^{12}$ and
$n_m\le40x^6$ ({\cite[Lemma~8.8(b)]{P49}}) and $1600\le x$. Then $\log L\le\log b+19\log x\le2\log b$,
so $\log h\le3L\log L\le6x^{19}b\log b\le x^{20}b^2$, and $x^{20}+x^{99}+x^{200}\le x^{201}$.
\end{proof}

\begin{proposition}[Proposition~\ref{pgr:delta}]\label{pgs:delta}
Proposition~\ref{pgr:delta} holds, with $C_{\mathrm{du}}=1$. More precisely:
\begin{enumerate}
\item[(a)] for $y\in\auAz$, $[K_m:J^\circ_\vartheta]\int_{[\auU(Y_m)]}\delta_{x,v}\cdot X^\vartheta_{0,m}y=
\langle v,\iota_\vartheta y(x)\rangle_\pgsdu/\dim W_{\lambda}$;
\item[(b)] $\delta_{x,v}$ is fixed by the principal congruence subgroup modulo $\pcm(\vartheta)$ (the
level of the outer node $Y^0_m$); if $x$ is integral above $\ell(\iota)$ and $v$ is integral at
$\iota$, then $\delta_{x,v}$ has gauge at most $1$ at $\iota$ at every point; and
$\|\sigma(\delta_{x,v}(g))\|^2\le\|\sigma(v)\|^2$ for all $g$ and $\sigma\in\auAut(\C)$;
\item[(c)] if $u$ is a vector of weight $\lambda$ with $\langle w,u\rangle_\pgsdu\in h\auO_\iota$ for
all integral $w$, then $u$ is integral at $\iota$; if $\sigma\in\auAut(\C)$, $K\ge0$ and
$|\sigma\langle v,u\rangle_\pgsdu|\le K\|\sigma v\|$ for all $v$, then $\|\sigma u\|^2\le K^2$; and
$h\exp(x^{99}b^2)\le\exp(x^{201}b^2-x^{200}b^2)$.
\end{enumerate}
\end{proposition}
\begin{proof}
(a) Let $\psi=\iota_\vartheta y$, the algebraic form of weight $\lambda$ whose form is
$X^\vartheta_{0,m}y$ \eqref{pc:eq-iota}; it is right invariant under $J^\circ_\vartheta$
(Lemma~\ref{pgr:jcirc}(b)). Let $\eta=\delta_{x,v}$. As in the proof of
Lemma~\ref{pc:algpet}, $[\auU(Y_m)]$ is the disjoint union over the classes
$x'\in\Gamma\backslash G_f/J^\circ_\vartheta$ of the images of $\auU(Y_m)_\infty\times x'J^\circ_\vartheta$, of
volume $[K_m:J^\circ_\vartheta]^{-1}$ (the group $J^\circ_\vartheta$ is free), and the integrand is right
$J^\circ_\vartheta$-invariant. On the class of $x'$ the integral of $f_\eta f_\psi$ is
$[K_m:J^\circ_\vartheta]^{-1}\int_{\auU(Y_m)_\infty}\langle u^{-1}\eta(x'),v^+\rangle_\pgsdu\langle v'^+,
u^{-1}\psi(x')\rangle_\pgsdu\,du$. By Schur orthogonality for the irreducible representation
$W_\lambda$ of the compact group $\auU(Y_m)_\infty$ of mass $1$ and its dual $W_{\lambda'}$, paired
invariantly by (P1),
\[
\int\langle u^{-1}a,v^+\rangle_\pgsdu\langle v'^+,u^{-1}b\rangle_\pgsdu\,du=
\frac{\langle a,b\rangle_\pgsdu\langle v'^+,v^+\rangle_\pgsdu}{\dim W_\lambda}=
\frac{\langle a,b\rangle_\pgsdu}{\dim W_\lambda}.
\]
Only the class of $x$ contributes, and there $\eta(x)=v$.

(b) The group $J^\circ_\vartheta$ is the principal congruence subgroup of $K_m$
modulo $\pcm_\vartheta$, which contains the principal congruence subgroup modulo
$\pcm(\vartheta)=\pcb\prod_{w\in Q}w$, since $\pi_M\mid\pcb$. At the point $\gamma xk$ the value is
$\gamma v$, and by (G1), (G2) of {\cite[\S6.1]{P49}} its gauge at $(\iota,\gamma xk)$ is the gauge of
$v$ at $(\iota,x)$, that is the largest absolute value of a coordinate of $\rho(x_\ell)^{-1}\iota(v)$;
if $x$ is integral above $\ell$, $\rho(x_\ell)^{-1}$ has integral entries (G2), so this is $\le1$
for integral $v$. Finally $\sigma(\gamma v)=\rho(\gamma^\sigma)\sigma(v)$, where $\rho(\gamma^\sigma)$
is the action of $\gamma$ through the embeddings $\sigma\circ\tau\circ c$, again an element of the
compact group $\prod_\tau\auU(N)$, which preserves the invariant form; so $\|\sigma(\gamma v)\|=
\|\sigma v\|$. (Off the class of $x$ the value is $0$.)

(c) Let $C$ be as in (P2) and $c_T$ its rows, integral. Then $h\,u_T=\langle c_T,u\rangle_\pgsdu\in
h\auO_\iota$, so $u_T\in\auO_\iota$. For the second claim let $\zeta=\sigma(u)$ (coordinatewise),
$a=\zeta^\flat$ as in (P3) and $v=\sigma^{-1}(a)$. As $A$ has integer entries,
$\sigma\langle v,u\rangle_\pgsdu=\langle a,\zeta\rangle_\pgsdu=\langle\zeta^\flat,\zeta^\flat\rangle=
\|\zeta\|^2$, while $\|\sigma v\|=\|a\|=\|\zeta\|$. So $\|\zeta\|^2\le K\|\zeta\|$, whence
$\|\zeta\|^2\le K^2$. The last claim is Lemma~\ref{pgs:hooksize}.
\end{proof}

\begin{remark}
In Proposition~\ref{pgr:delta}(c) the constant $C_{\mathrm{du}}$ may therefore be taken equal to $1$,
and the hypothesis $hC_{\mathrm{du}}^{1/2}\exp(x^{99}b^2)\le\exp(x^{201}b^2-x^{200}b^2)$ used in the
proof of Theorem~\ref{pc:auxcore}(b) holds. The weight $\lambda'$ is the weight of the conjugate
forms on $Y_m$ of the last path, and the bound on its boxes is Lemma~\ref{pgs:boxes} applied to the outer node $Y^0_m$ of the first grid.
\end{remark}

\subsubsection{Short forms of deep level: proof of Proposition~\ref{pgr:short}}\label{pgs:ss-short}

\paragraph*{The parameters of the grids}
At a node $(a,b)$ of either grid the type is given (\S\ref{pgr:ss-arch}, extending
\S\ref{pc:ss-triangle}) by its doubled Harish-Chandra parameter at each $\tau$, which consists of:
on the stem, the numbers $2\varsigma'_{\tau,p}$, $p<r$ (and, in the second grid, $p\le r$ for the
extra line), where $\varsigma'_{\tau,p}=k/2-1-p$ for $p<k-1$ and $\varsigma'_{\tau,p}=\Delta/2-
\kappa_{p+2-k}(\tau)$ otherwise; on the $l$-th block of the first kind ($l<a$, rank $t$ or $t+1$),
the numbers $2\kappa_{l+1}(\tau)+t-1-2b'$; on the $l$-th block of the second kind ($l<b$, rank $t$ or
$t-1$), the numbers $t-1-2(\kappa_{l+1}(\tau)+z)-2b'$; here $z\in\{0,1\}$ is the grid, $b'$ runs over
the rank of the block, and all entries are then raised by $1$ if $z=1$, and lowered by $1$ if $z=0$
and $a+b$ is odd; the weight of the forms moved at the node has the negative of this parameter in
the second grid and in the first grid (where conjugate forms are moved). The weight $\tilde\lambda$ of
Theorem~\ref{pc:auxcore} is the weight at the outer node $\pgrY^m_0=\pgrY$ (in the internal
indexing: no block of the first kind, $m$ blocks of rank $t-1$ of the second kind).

\begin{lemma}[Sizes of the weights on the grids]\label{pgs:boxes}
Let $\lambda=(\kappa,e)$ be the weight at a node of either grid (indices $\le m$), of rank $N$.
Then $N\le2n_m$; every entry of its doubled parameter has absolute value at most $20b$; $|e_\tau|\le
10b$; and its diagram at $\tau$ has at most $20Nb$ boxes, so at most $40n_mb$, and at most $40fn_mb
\le x^{19}b$ in total.
\end{lemma}
\begin{proof}
Recall the data used: $t=5r$ and $n_m=k+m(t+1)$ (proof of {\cite[Lemma~8.8(b)]{P49}}); $r=k+2\varrho+1$
with $\varrho=m/2-1$ ({\cite[Lemma~13.4(S7)]{P46}}), so $r+1=k+m$; and $\kappa_j(\tau)=\kappa(\tau)+(j+\frac12)\Delta$
with $\Delta=2r^8$, the parameters of the grids. The rank is at most $r+1+a(t+1)+b(t-1)\le r+1+2mt=k+m+2mt\le2(k+m(t+1))=2n_m$.
For $j\le2m$, $\kappa_j(\tau)=\kappa(\tau)+(j+\frac12)\Delta\le b+(2m+1)\cdot2r^8\le b+x^{42}\le2b$, by
$\kappa(\tau)\le b$ ({\cite[Lemma~13.4(S5)]{P46}}), $m\le2\sqrt x$, $r\le3x^5$, $x\ge(\log
b)^{100}\ge128^{100}$ and $x^{42}\le b$ (S1). The indices occurring are $p+2-k\le r+2-k=m+1$ and
$l+1\le m$. Hence the stem entries satisfy $|2\varsigma'_{\tau,p}|\le\max(k,2\kappa_j(\tau))\le4b$
(as $k\le r\le b$, and, for the entries of the form $\Delta-2\kappa_j(\tau)$, $|\Delta-2\kappa_j(\tau)|\le2\kappa_j(\tau)$
because $\kappa_j(\tau)\ge\Delta/2$, as $\kappa(\tau)\ge1$ by (S5) and $j\ge0$), the entries of the blocks are at most $2\kappa_{l+1}(\tau)+2t+2\le4b+10r+2\le5b$
in absolute value, and the final shift adds at most $1$: all entries are at most $6b\le20b$.

Let $l_0>l_1>\dots>l_{N-1}$ be the entries, so $l_p=2\mu_p+N-1-2p$ with the highest weight
$\mu_0\ge\dots\ge\mu_{N-1}=e_\tau$. The row lengths of the diagram are
$\mu_p-\mu_{N-1}=\frac12(l_p-l_{N-1})-(N-1-p)\le\frac12(|l_p|+|l_{N-1}|)\le20b$, so the diagram has at
most $20Nb$ boxes. Also $2e_\tau+N-1\le2\mu_0+N-1=l_0\le20b$ and $2e_\tau-(N-1)=l_{N-1}\ge-20b$, so
$|e_\tau|\le10b$. Finally $40fn_mb\le40\cdot x^{12}\cdot40x^6\,b\le x^{19}b$ by
{\cite[Lemma~8.8(b)]{P49}} and {\cite[Lemma~13.4(S3)]{P46}}.
\end{proof}

\paragraph*{Short forms span at a free level}
The following lemma is Step~1 of the proof of {\cite[Proposition~14.22]{P46}}, for an arbitrary weight;
the only new point is the determinant part of the weight, which is absorbed by short bases of the
lattices of the powers $g^{e_\tau}$.

\begin{lemma}[Short forms span]\label{pgs:shortspan}
Let $S$ be a positive definite standard space of rank $n\ge1$ (so $\auU(S)(\auA_{E^+,f})=
\auU(E^n)(\auA_{E^+,f})$ via the frame), $\lambda=(\kappa,e)$ a weight, $|\kappa|=\sum_\tau|\mu_\tau|$
its total number of boxes, $J\subseteq K_S$ a free level, $c_0$ a positive integer in the different
of $E$, $C=6c_0(fn)^2$ the constant of {\cite[Lemma~14.21]{P46}}, and $B\ge C^{6(nf+|\kappa|)}$.
Then $\mathcal A_\lambda(J,1)$ is spanned by forms $\varphi$ that are integral at every
$j\colon\C\to\C_\ell$, for every prime $\ell$, and satisfy $\|\sigma(\varphi(g))\|^2\le B$ for all
$g\in G_f$ and $\sigma\in\auAut(\C)$.
\end{lemma}
\begin{proof}
As in Step~1 of the proof of {\cite[Proposition~14.22]{P46}}, since $J$ is free,
$\mathcal A_\lambda(J,1)$ is the direct sum over the classes $x\in\Gamma\backslash G_f/J$ of the
spaces $\{\delta_{x,\xi}:\xi\in W_\lambda\}$, where $\delta_{x,\xi}(\gamma xk)=\gamma\xi$; the form
$\delta_{x,\xi}$ is integral at $j$ if and only if $\xi$ has gauge at most $1$ at $(j,x)$ (by (G1),
(G2)), and $\|\sigma(\delta_{x,\xi}(g))\|=\|\sigma\xi\|$ or $0$ (proof of
Proposition~\ref{pgs:delta}(b)). So it suffices, for fixed $x$, to span $W_\lambda$ by vectors $\xi$
of gauge at most $1$ at $(j,x)$ for all $j$ and with $\|\sigma\xi\|^2\le B$ for all $\sigma$.

By {\cite[Lemma~14.21]{P46}} choose $Y_0\in\auGL_n(E)$ with $x^{-1}Y_0$ integral and, for each $\tau$,
$Y_\tau\in\auGL_n(E)$ with $(x^{e_\tau})^{-1}Y_\tau$ integral ({\cite[Lemma~14.21]{P46}} applied to the
element $x^{e_\tau}$ of $G_f$; $e_\tau$ may be negative), all entries of all of them having all
conjugates of absolute value at most $C$. For a family $\zeta=(\zeta_\tau)$ of $n\times n$ matrices
whose entries are roots of unity put
\[
\xi(\zeta)=\prod_\tau\tau'(\det Y_\tau)\cdot\rho_0(Y_0)\,\pgsev(\zeta),
\]
where $\rho_0$ is the action of weight $(\kappa,0)$ through the embeddings, $\tau'=\tau\circ c$ is the
embedding through which the factor at $\tau$ is acted on, and $\pgsev(\zeta)=\bigotimes_\tau
\pgsev(\zeta_\tau)$ as in \eqref{pgs:eq-lattice}.

\emph{They span.} $\rho_0(Y_0)$ and the scalar are invertible, and the $\pgsev(\zeta)$ span
$W_\lambda$: a linear form vanishing on them is $\zeta\mapsto\sum_Tc_Tb_T(\zeta)$, a polynomial
vanishing at all families with root of unity entries, hence $0$, and the $b_T$ are linearly
independent, so $c=0$.

\emph{Integrality.} The action of $x_\ell^{-1}$ in weight $\lambda$ is $\prod_\tau\tau'(\det
x^{-1})^{e_\tau}$ times $\rho_0(x^{-1})$, and $\rho_0(x^{-1})\rho_0(Y_0)=\rho_0(x^{-1}Y_0)$, $\det(
x^{-1})^{e_\tau}\det Y_\tau=\det((x^{e_\tau})^{-1}Y_\tau)$, all components at the primes above $\ell$
taken through $j$. So $\rho(x_\ell)^{-1}j(\xi(\zeta))$ is a product of determinants of integral
matrices times $\rho_0(x^{-1}Y_0)$, a matrix whose entries are polynomials with integer coefficients
in the entries of the integral matrix $x^{-1}Y_0$ (Schur matrices; {\cite[\S6.1]{P49}} (N1)), applied
to $\pgsev(\zeta)$, whose coordinates $b_T(\zeta)$ are integer polynomials in roots of unity. All
coordinates are integral: the gauge is at most $1$.

\emph{Size.} Fix $\sigma$. The coordinates of $\sigma(\xi(\zeta))$ are those of
$\prod_\tau\det(Y'_\tau)\cdot\rho_0(Y'_0)\pgsev(\zeta')$, where $Y'_\tau=\sigma\tau'(Y_\tau)$, $Y'_0$ is
the family $\sigma\tau'(Y_0)$, and $\zeta'=\sigma(\zeta)$ has root of unity entries. By (N1),
$\rho_0(Y'_0)\pgsev(\zeta')=\bigotimes_\tau\pgsev(z_\tau)$ with $z_\tau=\tilde Y_\tau\zeta'_\tau$,
$\tilde Y_\tau$ being $\sigma\tau'(Y_0)$ or a transpose of it. Each entry of $z_\tau$ is a sum of $n$
products of a number of absolute value $\le C$ and a root of unity, so each column of $z_\tau$ has
squared norm at most $n(nC)^2\le C^5$ (as $n\le C$). By \eqref{pgs:eq-evnorm},
$\|\pgsev(z_\tau)\|^2\le C^{5|\mu_\tau|}$. By Hadamard's inequality $|\det Y'_\tau|^2\le(nC^2)^n\le
C^{3n}$. As the invariant form is the tensor product over $\tau$,
\[
\|\sigma\xi(\zeta)\|^2\le\prod_\tau C^{3n}\cdot\prod_\tau C^{5|\mu_\tau|}\le C^{3nf+5|\kappa|}\le B.
\qedhere
\]
\end{proof}

\begin{proposition}[Proposition~\ref{pgr:short}; Theorem~\ref{pc:auxcore}(c)]\label{pgs:short}
$\pcFF=\mathcal A_{\tilde\lambda}(J(\pcb),1)$ is spanned by forms that are integral at every
$j\colon\C\to\C_\ell$ for every prime $\ell$ (in particular at every $\iota\in\auI$) and whose values
have all conjugates of squared norm at most $\exp(x^{99}b^2)$.
\end{proposition}
\begin{proof}
Apply Lemma~\ref{pgs:shortspan} with $S=\pgrY$, $n=\tilde n=n_m-m\le n_m$, $\lambda=\tilde\lambda$,
$J=J(\pcb)$ (free by {\cite[Lemma~14.5]{P46}}, as recalled in \S\ref{pc:ss-aux}) and $c_0=q_1d$ (which
lies in the different of $E\subseteq\Q(\zeta_{q_1d})$, as in the proof of {\cite[Proposition~14.22]{P46}}).
By {\cite[Lemma~8.8(b)]{P49}}, $C=6q_1d(f\tilde n)^2\le x^{102}$; by Lemma~\ref{pgs:boxes},
$|\tilde\kappa|\le x^{19}b$; and $\tilde nf\le40x^{18}$. Hence
\[
\log C^{6(\tilde nf+|\tilde\kappa|)}\le6\cdot102\,(40x^{18}+x^{19}b)\log x\le x^{21}b\log b\le
x^{99}b^2,
\]
and $B=\exp(x^{99}b^2)$ is admissible.
\end{proof}

\begin{remark}
In the printed Theorem~\ref{pc:auxcore}, ``integral at $\iota$'' for $\varphi\in\pcFF$ means gauge at
most $1$ at $\iota$ at every point, which is the notion used here. For the number of boxes of
$\tilde\lambda$ at each place we use the explicit bound $x^{19}b$ (Lemma~\ref{pgs:boxes}).
\end{remark}




% ===== C-grid3.tex =====
% proofs/C-grid3.tex
% Continuation of proofs/C-grid.tex (pgr:) and proofs/C-grid2.tex (pgs:): the turn on the
% diagonal, Proposition pgr:turndiag. Labels start with pgu:. Conventions and macros: those of
% proofs/C-grid.tex and proofs/C-determinant.tex.
\providecommand{\auI}{\mathcal I}
\providecommand{\auGL}{\mathrm{GL}}
\providecommand{\auU}{\mathrm U}
\providecommand{\auO}{\mathcal O}
\providecommand{\auA}{\mathbb A}
\providecommand{\auNm}{\operatorname{Nm}}
\providecommand{\auHom}{\operatorname{Hom}}
\providecommand{\aubar}[1]{\overline{#1}}
\providecommand{\pcm}{\mathfrak m}
\providecommand{\pcMat}{\mathrm M}
\providecommand{\pctr}{\operatorname{tr}}
\providecommand{\psHS}{\mathcal H}
\providecommand{\psrho}{\varrho}
\providecommand{\psvac}{\mathbf v}
\providecommand{\psZ}{Z}
\providecommand{\pgrY}{\widetilde Y}
\providecommand{\pgrG}{\Gamma}
\providecommand{\pgrE}{\operatorname{Ed}}
\providecommand{\pgrTr}{\mathrel{\vdash}}
\providecommand{\pcb}{\mathfrak b}
\providecommand{\pguK}{\mathsf K}
\providecommand{\pguL}{\mathsf L}
\providecommand{\pgupt}{\mathsf u}

\subsection{The turn on the diagonal}\label{pgu:sec}

This section continues \S\ref{pgr:sec} and uses its notation. It proves
Proposition~\ref{pgr:turndiag} from Proposition~\ref{pgr:square} (the square law on a square of
a grid, at the square $(j,j)$ of the first grid), in the form of Theorem~\ref{pgu:turndiag}
below. The argument has five ingredients: the theta kernel of an edge ending at a diagonal node of
the second grid splits as a product of a kernel of the first grid and a kernel of a line
(\S\ref{pgu:ss-sum}); the kernel of a line can be reversed, at the cost of a finite Fourier
transform (\S\ref{pgu:ss-line}); two seesaw identities move periods between the diagonal nodes and
the node $S$ common to the two paths (\S\ref{pgu:ss-seesaw}); a path datum is turned across the
square after a Fourier transform in the column of the line (\S\ref{pgu:ss-turned}); and these are
assembled in \S\ref{pgu:ss-proof}.

\subsubsection{The data at a diagonal turn}\label{pgu:ss-data}

Fix $0\le j<m$, an $m$-fold twist $\vartheta$, $\pcm=\pcm(\vartheta)$ and $R=R_\pcm$. We use the
following nodes.
\begin{itemize}
\item In the first grid: $V'=Y^j_j$ (removed by the turn), $W_1$ the node $(j,j+1)$ (kept),
$V=Y^{j+1}_{j+1}$ (added), and $S_1$ the node $(j+1,j)$ (not on either path).
\item In the second grid: $\widetilde V'=\pgrY^j_j=V'\oplus L$ (removed), $S$ the node at the place
$(j,j+1)$ of the half $h_1$ (kept), and $\widetilde V=\pgrY^{j+1}_{j+1}=V\oplus L'$ (added). Here
$L=E\tilde e_j$ and $L'=E\tilde e_{j+1}$ are lines with $\langle\tilde e_l,\tilde e_l\rangle=
\epsilon_l$, $\epsilon_{l}=-\epsilon_{l+1}$, and the frames of $\widetilde V'$, $\widetilde V$ are
those of $V'$, $V$ followed by $\tilde e_j$, $\tilde e_{j+1}$.
\end{itemize}
We write $\mathsf B,\mathsf A,\mathsf N_1,\mathsf N_2$ for the frames of $V',V,W_1,S$, $\bullet$ for
the index of the line, and $\mathsf N=\mathsf N_1\sqcup\mathsf N_2$.

\paragraph*{Facts on the grid data}
The proof uses the following properties of the grid data. They are proved in \S\ref{pgv:sec} from
the complete definition of the grids given there: (G1) is Proposition~\ref{pgv:G1}, (G2) is
Proposition~\ref{pgv:G2} (with $\alpha_L=\alpha_{V'S_1}$, $\beta^{\mathrm c}=1$,
$\alpha^{\mathrm r}=\upsilon$, $\beta_{L'}=\beta_{S_1V}$), (G3) is Theorem~\ref{pgv:G3}, (G4) is
Proposition~\ref{pgv:G4}.
\begin{enumerate}
\item[(G1)] $S$ and $S_1$ have the same rank and signs, and the frames identify them as spaces
(\S\ref{pc:ss-data}); we identify them, so that $\mathsf N_2$ is also the frame of $S_1$, and
$\#\mathsf B+\#\mathsf A=\#\mathsf N$ (both sides equal $2(r+(2j+1)t)$). (The types at $S$ and $S_1$
differ; only spaces, groups, lattices and frames are identified, which is all that is used:
Proposition~\ref{pgx:G1}.)
\item[(G2)] (Splitting characters.) Write the kernel of an edge $(p,q)$ with point datum $X$ and
Fock vector $P$ as $\theta_{(p,q)}(g_p,g_q;P,\delta_X)=\theta^{p,q}_{\alpha_{pq},\beta_{pq}}(g_p,
g_q;P,\delta_X)$ (notation of \S\ref{pc:ss-data}). There are automorphic characters
$\alpha_L,\beta^{\mathrm c}$ (column line) and $\alpha^{\mathrm r},\beta_{L'}$ (row line) of
$\auU(1)$, unramified outside $\pcm$ and tamely ramified at $\pcm$, such that
\[
\alpha_{\widetilde V'S}(\det(h\oplus k))=\alpha_{V'S_1}(\det h)\,\alpha_L(k),\qquad
\beta_{\widetilde V'S}=\beta_{V'S_1}\,\beta^{\mathrm c},
\]
\[
\alpha_{S\widetilde V}=\alpha_{S_1V}\,\alpha^{\mathrm r},\qquad
\beta_{S\widetilde V}(\det(h\oplus k))=\beta_{S_1V}(\det h)\,\beta_{L'}(k),
\]
and $\alpha^{\mathrm r}=(\beta^{\mathrm c})^{-1}\upsilon$. The tame characters of the edges satisfy
$\chi_{(\widetilde V',S)}=\chi_{(V',S_1)}$ and $\chi_{(S,\widetilde V)}=\chi_{(S_1,V)}$; moreover
$\chi_{(W_1,V)}=\chi_{(V',S_1)}$ and $\chi_{(S_1,V)}=\chi_{(V',W_1)}$ (opposite edges of a square
have equal characters, as in \S\ref{pgr:ss-off}).
\item[(G3)] (Edges.) The pairs of the grids are edges in the sense of
Theorem~\ref{ptr:edge}\textup{(a)--(d)}, with joint harmonics $P_{pq}$, and the local statement
\textup{(H)} of \S\ref{ptr:ss-mult} holds for them. This is the input that
Proposition~\ref{pgr:square} also uses (step (iii) of \S\ref{pc:ss-exchange}).
\item[(G4)] (Diagonal period.) There is $\pgupt=\pgupt_\pcm\in\auU(1)(\auA_{E^+,f})$, depending only
on $\pcm$, such that the diagonal period on a path with diagonal index $l$ is
\[
\mathsf P_{l,g}(y'',\varphi'')=\int_{[\auU(Y^l_l)]}y''(x)\,\varphi''\bigl((x\oplus\pgupt^l)g\bigr)\,dx ,
\]
$g$ a real point of $\auU(\pgrY^l_l)$, $x\oplus\pgupt^l$ acting on $Y^l_l\oplus E\tilde e_l$. The
element $\pgupt$ is $\aubar u_0/u_0$, $u_0=\alpha\varpi_\pcm$ (Definition~\ref{pgv:def-period}); it is
exactly the element of Lemma~\ref{pgx:line}.
\end{enumerate}
All forms on the grids are rapidly decreasing on the automorphic quotients (Theorem~\ref{ptr:edge}(c)
for the types of (G3)).

The edges of the turn are $e_1=(V',W_1)$, $e_2=(\widetilde V',S)$ (old) and $f_1=(W_1,V)$,
$f_2=(S,\widetilde V)$ (new). In the first grid $e_1$ is vertical and $f_1$ horizontal; in the
second grid $e_2$ is horizontal and $f_2$ vertical. For an edge tuple $S$ (we write $\mathbf S$ for
edge tuples in this section, to avoid a clash with the node $S$) we cut the matrices of the edges
at the line:
\[
\mathbf S_{e_2}=(y_2\ |\ v),\quad y_2\in\pcMat_{\mathsf N_2\times\mathsf B},\ v\in
\pcMat_{\mathsf N_2\times\bullet};\qquad
\mathbf S_{f_2}=\binom{x_2}{v'},\quad x_2\in\pcMat_{\mathsf A\times\mathsf N_2},\ v'\in
\pcMat_{\bullet\times\mathsf N_2},
\]
and put $y(\mathbf S)=\binom{\mathbf S_{e_1}}{y_2}\in\pcMat_{\mathsf N\times\mathsf B}$,
$v(\mathbf S)=v$, $x(\mathbf S)=(\mathbf S_{f_1}\ x_2)\in\pcMat_{\mathsf A\times\mathsf N}$,
$v'(\mathbf S)=v'$. Conversely $\mathbf S[y,v]$ and $\mathbf S[x,v']$ denote $\mathbf S$ with the
old, respectively new, edges replaced by the blocks of $(y,v)$, respectively $(x,v')$.

\subsubsection{Orthogonal sums and the splitting of kernels}\label{pgu:ss-sum}

\begin{lemma}[Orthogonal sums]\label{pgu:sum}
Let $W=W_1\perp W_2$ be hermitian spaces with self-dual lattices $L_1,L_2$ (for $[\,,]$), and
identify $\psHS(W)=\psHS(W_1)\otimes\psHS(W_2)$ (restricted tensor products; $\psrho(w_1+w_2)=
\psrho(w_1)\otimes\psrho(w_2)$). Then $\omega_W(g_1,g_2)=\omega_{W_1}(g_1)\otimes\omega_{W_2}(g_2)$
for $g_l\in\auU(W_l)(\auA)$, and $\theta_W=c\,\theta_{W_1}\otimes\theta_{W_2}$ with $c\ne0$.
\end{lemma}
\begin{proof}
Since $\beta_W(x_1+x_2,y_1+y_2)=\beta_{W_1}(x_1,y_1)\beta_{W_2}(x_2,y_2)$, the tensor product is a
model of $\psrho_W$. The functional $\theta_{W_1}\otimes\theta_{W_2}$ is invariant under
$\psrho(W(E))$, so it is $c^{-1}\theta_W$ by (W1). By (W1) for $W,W_1,W_2$ the operator
$\kappa(g)=(\omega_{W_1}(g_1)\otimes\omega_{W_2}(g_2))^{-1}\omega_W(g_1,g_2)$ commutes with all
$\psrho(w)$, hence is a scalar, and $\kappa$ is a continuous character of $\auU(W_1)(\auA)\times
\auU(W_2)(\auA)$. It suffices to show $\kappa(g_1,1)=1$ (and symmetrically). We follow the proof of
{\cite[Proposition~7.17]{P49}}, with $\kappa_0=\kappa(\cdot,1)$ in place of $\kappa_0$ there.
\emph{Rational points:} $\theta_W\circ\omega_W(\gamma)=\theta_W$ and the same for $W_1$ give
$\kappa_0(\gamma)=1$. \emph{Stabiliser of $L_1$:} $L=L_1\oplus L_2$ is self-dual, its vacuum is
$\psvac_{L_1}\otimes\psvac_{L_2}$ (by (W4), the line of the character is unique), and both
$\omega_W(k,1)$ and $\omega_{W_1}(k)\otimes1$ fix it for $k$ in the stabiliser of $L_1$ (last
assertion of (W4)); so $\kappa_0(k)=1$. \emph{Parabolic elements:} let $I\subset W_{1,v}$ be
isotropic and $p\in P_I(E^+_v)$ acting trivially on $I^\perp/I$ (in $W_1$). Then $I$ is isotropic in
$W_v$, $I^{\perp_W}/I=(I^{\perp_{W_1}}/I)\perp W_2$, and $(p,1)$ acts trivially on it; by the uniqueness
of the zero term ({\cite[\S7.5]{P49}}), $\psZ^W_I=c'\,\psZ^{W_1}_I\otimes1$, and (W2$'$) for $W$ and $W_1$
gives $\chi_I(p)\psZ^W_I=\psZ^W_I\omega_W(p,1)=\kappa_0(p)\chi_I(p)\psZ^W_I$, the character $\chi_I(p)$
being the same for $W$ and $W_1$ (it depends on $\det p|_I$ only). So $\kappa_0(p)=1$. The remaining
steps (all finite places, by the Iwasawa decomposition; the infinite places, by density of the
rational points) are those of the proof of {\cite[Proposition~7.17]{P49}}, word for word.
\end{proof}

\begin{corollary}[Splitting at the diagonal nodes]\label{pgu:split}
Let $\pguK^{\mathrm c}_{v,B}(k,g)=\theta^{L,S}_{\alpha_L,\beta^{\mathrm c}}(k,g;B,\delta_v)$ for
$v\in\pcMat_{\mathsf N_2\times\bullet}$ and a polynomial $B$ on $\auHom(L,S)\otimes\R$, and
$\pguK^{\mathrm r}_{v',B'}(g,k)=\theta^{S,L'}_{\alpha^{\mathrm r},\beta_{L'}}(g,k;B',\delta_{v'})$ for
$v'\in\pcMat_{\bullet\times\mathsf N_2}$ and a polynomial $B'$ on $\auHom(S,L')\otimes\R$. Then
there are constants $c_1,c_2\ne0$, independent of $\vartheta$, such that for polynomials $A$ on
$\auHom(V',S)\otimes\R$, $A'$ on $\auHom(S,V)\otimes\R$, and all $h,k,g$,
\begin{align*}
\theta_{e_2}\bigl(h\oplus k,g;A\otimes B,\delta_{(y_2|v)}\bigr)&=c_1\,\theta_{(V',S_1)}(h,g;A,\delta_{y_2})\,
\pguK^{\mathrm c}_{v,B}(k,g),\\
\theta_{f_2}\bigl(g,h\oplus k;A'\otimes B',\delta_{\binom{x_2}{v'}}\bigr)&=c_2\,\theta_{(S_1,V)}(g,h;A',\delta_{x_2})\,
\pguK^{\mathrm r}_{v',B'}(g,k).
\end{align*}
\end{corollary}
\begin{proof}
$\auHom(\widetilde V',S)=\auHom(V',S)\perp\auHom(L,S)$, with the self-dual lattices
$\alpha\auHom(L_{V'},L_S)$ and $\alpha\auHom(L_L,L_S)$, and $(h\oplus k,g)$ acts as $(h,g)$ on the
first summand and $(k,g)$ on the second. By Lemma~\ref{pgu:sum}, $\omega(h\oplus k,g)=\omega(h,g)
\otimes\omega(k,g)$ and $\theta=c\,\theta\otimes\theta$. The datum factors:
$A\otimes B$ at the infinite places, the vacuum $\psvac^\pcm=\psvac\otimes\psvac$ outside
$S(\pcm)$, and at $v\in S(\pcm)$ the function $\delta^\natural_{(y_2|v)}(x_1|x_2)=
\delta^\natural_{y_2}(x_1)\delta^\natural_v(x_2)$ (the norms at $S(\pcm)$ are $L^2$ norms of product
measures, so the normalisations factor up to a constant independent of $\vartheta$). The characters
factor by (G2). This gives the first identity; the second is the same with
$\auHom(S,\widetilde V)=\auHom(S,V)\perp\auHom(S,L')$.
\end{proof}

\subsubsection{Reversing a line}\label{pgu:ss-line}

This subsection treats the case where $\alpha$ is a unit at the primes of $\pcm$. The general case is
Lemma~\ref{pgx:line} (\S\ref{pgx:ss-line}), which has the statement below with the data of
Definition~\ref{pgx:conv} and with $\pgupt$ exactly the element of Definition~\ref{pgv:def-period}; it is
the form used in the rest of the paper, and its proof follows the one given here. In this subsection we
assume that $\alpha$ is a unit at the primes of $\pcm$, so that at $v\in S(\pcm)$ the self-dual
lattices $\alpha\auHom(L_S,L_T)_v$ are the lattices of integral matrices. Let $n=\#\mathsf
N_2$ be the rank of $S$, and $L^-$ the line $L$ with the opposite form; it is the line $L'$ (same
frame vector, sign $-\epsilon_j=\epsilon_{j+1}$), and we identify them.

\begin{lemma}[Line reversal when $\alpha$ is a unit at the primes of $\pcm$; the general case is Lemma~\ref{pgx:line}]\label{pgu:line}
Assume that $\alpha$ is a unit at the primes of $\pcm$. There are an element $\pgupt=\pgupt_\pcm$ of $\auU(1)(\auA_{E^+,f})$ depending only on $\pcm$, a
linear bijection $\tau$ from the polynomials on $\auHom(L,S)\otimes\R$ to those on
$\auHom(S,L')\otimes\R$, independent of $\vartheta$, a primitive additive character $\psi_\pcm$ of
$R$, and $c_{\mathrm{line}}=c_{\mathrm{line}}(\vartheta)\in\C^\times$ with
\[
|c_{\mathrm{line}}|^2=C_0\,(\auNm\pcm)^{-n},\qquad C_0>0\text{ independent of }\vartheta,
\]
such that for every $k\in\auU(L)(\auA_f)$, $g\in\auU(S)(\auA)$, $v$ and $B$
\[
\pguK^{\mathrm c}_{v,B}(k,g)=c_{\mathrm{line}}\,\zeta(k)
\sum_{v'\in\pcMat_{\bullet\times\mathsf N_2}(R)}\psi_\pcm(v'v)\,\pguK^{\mathrm r}_{v',\tau B}(g,k\pgupt)
\]
where $\zeta(k)$ is a number of modulus one depending only on $k$, on the characters and on $\pcm$.
Here $\pgupt$ is the element constructed in the proof (it is $k_0$ outside $S(\pcm)$ and has
$w$-component $\varpi^{-k_w}$ at $v\in S(\pcm)$); it is in general not the element $\pgupt$ of
(G4), and the case of (G4) is treated by Lemma~\ref{pgx:line}.
\end{lemma}

\begin{proof}
\emph{The flip.} Let $\mathrm{Fl}\colon\psHS(\auHom(L,S))\to\psHS(\auHom(S,L^-))$ be the isomorphism
of {\cite[Proposition~7.17]{P49}} ((W6)): $\mathrm{Fl}\,\psrho(w)=\psrho(w^\dagger)\mathrm{Fl}$,
$\theta\circ\mathrm{Fl}=c_\theta\theta$, and $\mathrm{Fl}\,\omega(k,g)=\upsilon(\det k)^n\upsilon(\det
g)^{-1}\omega(g,k)\mathrm{Fl}$. With the twists of \S\ref{pc:ss-data},
\[
\mathrm{Fl}\,\omega^{L,S}_{\alpha_L,\beta^{\mathrm c}}(k,g)=\omega^{S,L^-}_{(\beta^{\mathrm c})^{-1}
\upsilon,\ \alpha_L^{-1}\upsilon^n}(g,k)\,\mathrm{Fl}
=\alpha_L(k)^{-1}\upsilon(k)^n\beta_{L'}(k)^{-1}\,
\omega^{S,L'}_{\alpha^{\mathrm r},\beta_{L'}}(g,k)\,\mathrm{Fl},
\]
using $\alpha^{\mathrm r}=(\beta^{\mathrm c})^{-1}\upsilon$ (G2), $\omega_{a,b}(g,h)=a(\det g)^{-1}b(\det h)\,\omega(g,h)$, and
$\beta^{\mathrm c}(g)\upsilon(g)^{-1}=\alpha^{\mathrm r}(g)^{-1}$. (Write
$\zeta_1(k)=\alpha_L(k)^{-1}\upsilon(k)^n\beta_{L'}(k)^{-1}$, as in Lemma~\ref{pgx:line}; it is a character of $\auU(1)$; at the fixed points $k$ used below it is a number of
modulus one; this is the factor $\zeta(k)$ of the statement, together with the value at
$\pgupt$ in the last step.) Hence
\[
\pguK^{\mathrm c}_{v,B}(k,g)=c_\theta^{-1}\theta\bigl(\mathrm{Fl}\,\omega^{L,S}(k,g)[B,\delta_v]\bigr)
=\zeta_1(k)\,c_\theta^{-1}\,\theta\bigl(\omega^{S,L'}_{\alpha^{\mathrm r},\beta_{L'}}(g,k)\,
\mathrm{Fl}[B,\delta_v]\bigr),
\]
and it remains to compute $\mathrm{Fl}[B,\delta_v]$.

\emph{Local factors.} The Heisenberg representations of $\auHom(L,S)_v$ and $\auHom(S,L^-)_v$ are
irreducible, and $w\mapsto w^\dagger$ preserves $\beta$ locally; so there are unitary local
intertwiners $\mathrm{Fl}_v$, unique up to factors of modulus one, and by the uniqueness of
$\mathrm{Fl}$, $\mathrm{Fl}=c_0\bigotimes_v\mathrm{Fl}_v$, where $|c_0|$ does not depend on the
choice of the unitary $\mathrm{Fl}_v$, hence not on $\vartheta$.

\emph{Infinite places.} In the Fock model, $\psHS_\infty(W)$ consists of the functions on
$W\otimes\R$ holomorphic for the complex structure of $W$ (multiplication by $i$ on the coordinates
of positive sign, by $-i$ on those of negative sign) and square integrable for the Gaussian; the
vacuum is $1$, and $\psrho_\infty$ is defined from $\beta$ and the complex structure. The map
$\dagger$ is $\R$-linear, preserves $\beta$, and carries the complex structure of $\auHom(L,S)$ to
that of $\auHom(S,L^-)$: it is antilinear, and the signs of the coordinates are negated (the form
of $L^-$ is that of $L$ negated). Hence $F\mapsto F\circ\dagger^{-1}$ is a unitary intertwiner, and
we take it as $\mathrm{Fl}_\infty$. It maps polynomials of degree $d$ to polynomials of degree $d$;
put $\tau B=B\circ\dagger^{-1}$ (in coordinates, a transposition of the variables with signs). It
does not depend on $\vartheta$.

\emph{Finite places outside $S(\pcm)$.} $\mathrm{Fl}_v$ maps the vacuum of the self-dual lattice
$M_v=\alpha\auHom(L_L,L_S)_v$ to a vector on which $\psrho(l^\dagger)$, $l\in M_v$, acts by
$\epsilon_{M_v}(l)$; as in the proof of {\cite[Proposition~7.17]{P49}} (stabiliser of $L$), this is the
character of the self-dual lattice $M_v^\dagger=\aubar\alpha\auHom(L_S,L_L)_v$, so by (W4)
$\mathrm{Fl}_v\psvac_{M_v}=c_v\psvac_{M^\dagger_v}$ with $|c_v|=1$. Let $k_0\in\auU(L')(\auA_f)$ be
the scalar $\aubar\alpha/\alpha$ on the line (it has norm one). Then $M^\dagger_v=k_{0,v}
\,\alpha\auHom(L_S,L_{L'})_v$, and by (W4) $\psvac_{M^\dagger_v}=c'_v\,\omega(1,k_{0,v})
\psvac_{\alpha\auHom(L_S,L_{L'})_v}$ with $|c'_v|=1$ ($\omega$ is unitary). At the places where
$\alpha$ is a unit both constants are $1$ for a suitable choice of $\mathrm{Fl}_v$.

\emph{Places of $S(\pcm)$.} Let $v=w\aubar w\in S(\pcm)$ with $w\mid\pcm$, $\varpi$ a uniformiser at
$w$ and $k_w=\operatorname{ord}_w\pcm$. In the model (W5), $\psHS_v(\auHom(L,S))$ is the space of
Schwartz--Bruhat functions of the columns $x\in E_w^n$, translations by $\auHom(L_w,S_w)$ act by
translation and those by $\auHom(L_{\aubar w},S_{\aubar w})$ by multiplication with characters;
similarly $\psHS_v(\auHom(S,L'))$ consists of functions of the rows $\rho\in E_w^n$. The map
$\dagger$ is antilinear, so it exchanges the components at $w$ and $\aubar w$: it carries
translations of the first model into multiplications of the second and conversely. An operator with
these commutation relations is, up to a scalar (irreducibility), a Fourier transform: there is
$\kappa_w\in E_w^\times$ with
\[
\mathrm{Fl}_v\Psi(\rho)=\int_{E_w^n}\Psi(x)\,\psi_w(\kappa_w\rho x)\,dx ,
\]
$\psi_w$ the local component of $\psi$ and $dx$ the measure giving $\auO_w^n$ volume $1$, up to a
factor of modulus one. As $\mathrm{Fl}_v$ maps the vacuum $1_{\auO_w^n}$ of the integral lattice to
a multiple of the vacuum $1_{\auO_w^n}$ of the target (as in the previous paragraph, $\alpha$ being a
unit at $w$), and the transform of $1_{\auO_w^n}$ is the indicator of the dual lattice of
$\auO_w^n$ for $(\rho,x)\mapsto\psi_w(\kappa_w\rho x)$, this dual lattice is $\auO_w^n$; so $dx$ is
self-dual, the formula is unitary (Plancherel), and the factor has modulus one. Now let
$\tilde v\in\auO_w^n$ lift $v$. Then
\[
\mathrm{Fl}_v\delta_v^\natural(\rho)=\int_{\tilde v+\varpi^{k_w}\auO_w^n}\psi_w(\kappa_w\rho x)\,dx
=q_w^{-k_wn}\,\psi_w(\kappa_w\rho\tilde v)\,1_{\varpi^{-k_w}\auO_w^n}(\rho),
\]
$q_w$ the order of the residue field. Let $\pgupt_v\in\auU(L')(E^+_v)\cong E_w^\times$ have
$w$-component $\varpi^{k_w}$. By (W5) and unitarity, $(\omega(1,\pgupt_v)F)(\rho)=\eta\,
q_w^{k_wn/2}F(\varpi^{-k_w}\rho)$ with $|\eta|=1$ (the factor $|\det|^{-n/2}$ makes the operator
unitary). Hence
\[
\omega(1,\pgupt_v)\,\mathrm{Fl}_v\delta_v^\natural(\rho)=\eta\,q_w^{-k_wn/2}\,\psi_w(\kappa_w
\varpi^{-k_w}\rho\tilde v)\,1_{\auO_w^n}(\rho)=\eta\,q_w^{-k_wn/2}\sum_{v'}\psi_{w}^{\pcm}(v'v)\,
\delta^\natural_{v'}(\rho),
\]
where $\psi^\pcm_w(r)=\psi_w(\kappa_w\varpi^{-k_w}\tilde r)$ is a well defined character of
$\auO_w/\varpi^{k_w}$ ($\kappa_w$ is a unit for $\psi_w$ in the sense that $\auO_w$ is self-dual for
$\psi_w(\kappa_w\,\cdot\,)$), primitive. Put $\psi_\pcm=\prod_{w\mid\pcm}\psi^\pcm_w$ on
$R=\prod_w\auO_w/\varpi^{k_w}$.

\emph{Assembling.} Let $\pgupt$ be $k_0^{-1}$ outside $S(\pcm)$ (at the places where $\alpha$ is
not a unit) and $\pgupt_v$ at $v\in S(\pcm)$; it depends only on $\pcm$. The local results
give
\[
\mathrm{Fl}[B,\delta_v]=c\;\omega^{S,L'}(1,\pgupt^{-1})\Bigl[\tau B,\sum_{v'}\psi_\pcm(v'v)\delta_{v'}\Bigr],
\qquad|c|=|c_0|\,\prod_{w\mid\pcm}q_w^{-k_wn/2}=|c_0|\,(\auNm\pcm)^{-n/2},
\]
the operators $\omega(1,\cdot)$ being taken without the determinant twists (which only contribute
a factor of modulus one, depending on $\pcm$ and the characters). Inserting this in the formula of
the first step and using $\omega(g,k)\omega(1,\pgupt^{-1})=\omega(g,k\pgupt^{-1})$, we get the
identity with $k\pgupt^{-1}$ in place of $k\pgupt$; replacing $\pgupt$ by $\pgupt^{-1}$ (both depend
only on $\pcm$) gives the statement, with $c_{\mathrm{line}}=c/c_\theta$ and
$C_0=|c_0/c_\theta|^2$.
\end{proof}

\subsubsection{The seesaw}\label{pgu:ss-seesaw}

A \emph{moderate weight} from a space $S_0$ to a space $T'$ is a pair $(N,j)$ of a continuous
function $N$ on $\auU(S_0)(\auA)$, left invariant under $\auU(S_0)(E^+)$ and bounded by a power of
the height (Definition~\ref{pcx:def-ht}), and a continuous map $j\colon\auU(S_0)(\auA)\to
\auU(T')(\auA)$ compatible with the rational points and with $\mathrm{ht}(j(x))\le C\,\mathrm{ht}(x)^C$.
For forms $y$ on $S_0$, $F$ on $T'$ and a real point $r$ of $\auU(T')$ put
\[
\langle y,F\rangle_{N,j,r}=\int_{[\auU(S_0)]}N(x)\,y(x)\,F(j(x)r)\,dx ;
\]
it converges absolutely when $y$ is rapidly decreasing and $F$ is bounded.

\begin{lemma}[Seesaw]\label{pgu:seesaw}
Let $(S_0,T)$ and $(T',S')$ be pairs of spaces that are edges (\textup{(G3)}) with joint harmonics
$\Phi$ and $\Phi'$ and fixed characters, for a family of moduli. Let $X$ be a set, and for each
$x\in X$ let $W_x$ be a finite set and, for $i\in W_x$, let $u_{x,i}\in\C$, $y_{x,i}$ a form on
$S_0$ of the source type of $\Phi$, $f_{x,i}$ a form on $S'$ of the conjugate of the target type of
$\Phi'$, $(N_{x,i},j_{x,i})$ a moderate weight from $S_0$ to $T'$, and $\Psi'_{x,i}$ a finite
datum for $(T',S')$ modulo $\pcm_x$. Let $\Xi$ be a function on $X$. Call a polynomial $p$ on
$\auHom(T',S')\otimes\R$ \emph{matched} if there are a polynomial $Q$ on $\auHom(S_0,T)\otimes\R$,
finite data $\Psi_{x,i}$ for $(S_0,T)$ modulo $\pcm_x$ and moderate weights $(M_{x,i},j'_{x,i})$ from
$S'$ to $T$ such that
\begin{enumerate}
\item[(a)] for every real point $r'$ of $\auU(T)$ there is $\beta(r')$ with
$\sum_{i\in W_x}u_{x,i}\langle f_{x,i},\Theta(\Phi,\Psi_{x,i})y_{x,i}\rangle_{M_{x,i},j'_{x,i},r'}=
\beta(r')\,\Xi(x)$ for all $x$;
\item[(b)] for all $x,i$ and all $a\in\auU(S_0)(\auA)$, $b\in\auU(S')(\auA)$,
$N_{x,i}(a)\,\allowbreak \theta^{T',S'}(j_{x,i}(a),b;p,\Psi'_{x,i})=M_{x,i}(b)\,\allowbreak \theta^{S_0,T}(a,j'_{x,i}(b);Q,
\Psi_{x,i})$.
\end{enumerate}
If the matched polynomials span all polynomials, then for every real point $r$ of $\auU(T')$ there
is $\beta\in\C$ with
\[
\sum_{i\in W_x}u_{x,i}\,\bigl\langle y_{x,i},\Theta^{\mathsf t}(\Phi',\Psi'_{x,i})f_{x,i}
\bigr\rangle_{N_{x,i},j_{x,i},r}=\beta\,\Xi(x)\qquad\text{for all }x\in X .
\]
\end{lemma}

\begin{proof}
For a Fock vector $\mathsf v$ of $\auHom(T',S')\otimes\R$ with finite moments put
\[
\Lambda_{x,i}(\mathsf v)=\int_{[\auU(S_0)]}\int_{[\auU(S')]}N_{x,i}(a)\,y_{x,i}(a)\,
\theta^{T',S'}(j_{x,i}(a),b;\mathsf v,\Psi'_{x,i})\,f_{x,i}(b)\,db\,da ,\qquad
\Lambda_x=\sum_{i\in W_x}u_{x,i}\Lambda_{x,i}.
\]
\emph{Convergence and continuity.} By (R2) of \S\ref{pcx:ss-conv} there are $M,A,Y$ with
$|\theta^{T',S'}(a',b;\mathsf v,\Psi')|\le Y\|\mathsf v\|_M(\mathrm{ht}(a')\mathrm{ht}(b))^A$, where
$\|\cdot\|_M$ is the moment norm of Proposition~\ref{pcx:decay}. The forms $y_{x,i}$ and $f_{x,i}$
are rapidly decreasing and $N_{x,i}$, $j_{x,i}$ are moderate, so the integrand is bounded by
$\|\mathsf v\|_M$ times an integrable function ((R1) of \S\ref{pcx:ss-conv}), and $\Lambda_{x,i}$ is a
continuous linear functional for $\|\cdot\|_M$; the integrals may be taken in any order.

\emph{Matched polynomials.} Let $p$ be matched, with $Q,\Psi,M$. By Lemma~\ref{ptr:conv-lift}
(this uses (H), included in (G3)) there is $\alpha_Q\in C_c(\auU(T)_\infty)$, independent of $x,i$,
with $\Theta(Q,\Psi_{x,i})y_{x,i}(h)=\int\alpha_Q(r')\,\Theta(\Phi,\Psi_{x,i})y_{x,i}(hr')\,dr'$. By
Fubini, (b), and the definition of the lift,
\[
\Lambda_{x,i}(p)=\int_{[\auU(S')]}M_{x,i}(b)f_{x,i}(b)\,\Theta(Q,\Psi_{x,i})y_{x,i}(j'_{x,i}(b))\,db=
\int\alpha_Q(r')\,\langle f_{x,i},\Theta(\Phi,\Psi_{x,i})y_{x,i}\rangle_{M_{x,i},j'_{x,i},r'}\,dr',
\]
the exchange of integrals being justified by the boundedness of the lifted forms and the rapid
decrease of $f_{x,i}$. Multiplying by $u_{x,i}$, summing and using (a),
$\Lambda_x(p)=\bigl(\int\alpha_Q\beta\bigr)\,\Xi(x)$ for all $x$. So the vector $(\Lambda_x(p))_{x\in
X}\in\C^X$ lies on the line $\C\,\Xi$.

\emph{All Fock vectors.} The set of $\mathsf v$ with $(\Lambda_x(\mathsf v))_x\in\C\,\Xi$ is a
linear subspace, closed for $\|\cdot\|_M$ (each $\Lambda_x$ is continuous, and a line is closed in
$\C^X$ with the product topology). It contains the matched polynomials, hence all polynomials, hence
(Lemma~\ref{ptr:dense}) every Fock vector with finite moments of all orders.

\emph{Descent.} Let $r$ be a real point of $\auU(T')$ and $\mathsf v_r=\omega_\infty(r,1)\Phi'$
(without determinant twists), which has finite moments of all orders. Translating the transposed
lift by $r$ moves the Fock vector: $\Theta^{\mathsf t}(\Phi',\Psi')f(a'r)=\int\theta^{T',S'}(a'r,b;
\Phi',\Psi')f(b)\,db=\zeta_r\int\theta^{T',S'}(a',b;\mathsf v_r,\Psi')f(b)\,db$, with $\zeta_r$ the
value at $r$ of the archimedean component of the determinant twist (of modulus one, independent of
$x,i$, since all characters of the family have the same weights). Hence
$\langle y_{x,i},\Theta^{\mathsf t}(\Phi',\Psi'_{x,i})f_{x,i}\rangle_{N_{x,i},j_{x,i},r}=\zeta_r
\Lambda_{x,i}(\mathsf v_r)$, and the claim follows from the previous step applied to $\mathsf v_r$.
\end{proof}

\subsubsection{Turning a path datum with a line}\label{pgu:ss-turned}

We keep the notation of \S\ref{pgu:ss-data}. Let $w$ be a set of edges none of which has $V'$ or
$\widetilde V'$ as an end or $V$ or $\widetilde V$ as its target, let $\chi=(\chi_{e_1}
\chi_{e_2})^{-1}$, and fix a bijection $\varepsilon\colon\mathsf B\sqcup\mathsf A\to\mathsf N$ and an
additive character $\psi$ of $R$. Matched functions are those of \S\ref{pc:ss-exchange} for $\chi$
and $\varepsilon$ (as in \S\ref{pgr:ss-off}). For a path datum $\Psi'$ on $w\cup\{e_1,e_2\}$ and an
edge tuple $\mathbf S$ put
\[
\widehat\Psi'_{\mathbf S}(v',y)=\sum_{v\in\pcMat_{\mathsf N_2\times\bullet}(R)}\psi(v'v)\,
\Psi'(\mathbf S[y,v])\qquad(v'\in\pcMat_{\bullet\times\mathsf N_2},\ y\in\pcMat_{\mathsf N\times
\mathsf B}),
\]
the Fourier transform in the column of the line; it depends on $\mathbf S$ only through
$(\mathbf S_e)_{e\in w}$.

\begin{proposition}[Turning with a line]\label{pgu:turned}
Let $\Psi'$ be a path datum on $w\cup\{e_1,e_2\}$ with bound $N$. Then
\[
\Psi(\mathbf S)=\bigl(\mathrm{match}_\chi\widehat\Psi'_{\mathbf S}(v'(\mathbf S),\cdot\,)\bigr)(x(\mathbf S))
\]
defines a path datum on $w\cup\{f_1,f_2\}$ with bound $\#\pcMat_{\mathsf N_2\times\bullet}(R)\cdot N
\le V(\vartheta)N$, and for every $\mathbf S$ and $v'$ the functions $x\mapsto\Psi(\mathbf S[x,v'])$
and $\widehat\Psi'_{\mathbf S}(v',\cdot)$ are matched.
\end{proposition}

\begin{proof}
\emph{Transformation law of the transform.} Let $g$ be an admissible node family and
$k\in\auGL_{\mathsf B}(R)$. Let $g'$ be $g$ with $k$ at $V'$ and $k\oplus1$ at $\widetilde V'$; it
is admissible (these nodes are neither outer nor touched by the other condition except through
$k\oplus1$). Put $G=\mathrm{diag}(g_{W_1},g_S)\in\auGL_{\mathsf N}(R)$. Then
$g'(\mathbf S[y,v])=(g\mathbf S)[G^{-1}yk,\ g_S^{-1}v]$, the edges of $w$ not meeting $V',\widetilde
V'$. The invariance of $\Psi'$ under $g'$ and the substitution $v\mapsto g_Sv$ in the sum give
\begin{equation}\label{pgu:eq-law}
\widehat\Psi'_{\mathbf S}(v',y)=\mathsf a(g)\,\chi(\det k)\,\widehat\Psi'_{g\mathbf S}(v'g_S,\,
G^{-1}yk),\qquad\mathsf a(g)=\chi_{e_1}(\det g_{W_1})\chi_{e_2}(\det g_S)\chi_w(g),
\end{equation}
since $\psi(v'g_S\cdot g_S^{-1}v)=\psi(v'v)$ and $\chi_{w\cup\{e_1,e_2\}}(g')=\chi(\det k)\mathsf
a(g)$ ($\det(k\oplus1)=\det k$).

\emph{The match exists.} If $\widehat\Psi'_{\mathbf S}(v',y)\ne0$ then $\Psi'(\mathbf S[y,v])\ne0$
for some $v$, so $\mathbf S_{e_1}$ (the rows $\mathsf N_1$ of $y$; $e_1$ is vertical) has a left
inverse, hence so has $y$, and $y=y_\gamma$ by Lemma~\ref{pgr:extend}; and \eqref{pgu:eq-law} with
$g=1$ gives $\chi(\det k)\widehat\Psi'_{\mathbf S}(v',yk)=\widehat\Psi'_{\mathbf S}(v',y)$. By
Lemma~\ref{pgr:match}(b) the match exists and is matched with $\widehat\Psi'_{\mathbf S}(v',\cdot)$.
Since $v'(\mathbf S[x,v'])=v'$, $x(\mathbf S[x,v'])=x$ and the transform depends only on
$(\mathbf S_e)_{e\in w}$, the last assertion follows, and $\Psi$ depends only on the edges of
$w\cup\{f_1,f_2\}$.

\emph{Invariance.} For admissible $g$: $v'(g\mathbf S)=v'(\mathbf S)g_S$, because $g$ is
$g_V\oplus1$ at $\widetilde V$, and $x(g\mathbf S)=g_V^{-1}x(\mathbf S)G$. Let
$H=\widehat\Psi'_{g\mathbf S}(v'(\mathbf S)g_S,\cdot)$ and $\Phi=\mathrm{match}_\chi H$. By
\eqref{pgu:eq-law} with $k=1$, $\widehat\Psi'_{\mathbf S}(v'(\mathbf S),y)=\mathsf a(g)H(G^{-1}y)$, so by
Lemma~\ref{pgr:match}(c) and uniqueness, $\Psi(\mathbf S)=\mathsf a(g)\chi(\det G)\Phi(x(\mathbf S)G)$.
On the other hand $\Psi(g\mathbf S)=\Phi(g_V^{-1}x(\mathbf S)G)=\chi(\det g_V)\Phi(x(\mathbf S)G)$
(last claim of Lemma~\ref{pgr:match}(c)). Hence $\Psi(\mathbf S)=\mathsf a(g)\chi(\det G)\chi(\det
g_V)^{-1}\Psi(g\mathbf S)$, and exactly as at the end of the invariance step of the proof of
Proposition~\ref{pgr:turned}, using $\chi_{f_1}=\chi_{e_2}$ and $\chi_{f_2}=\chi_{e_1}$ (G2) and
$\det(g_V\oplus1)=\det g_V$, the factor is $\chi_{w\cup\{f_1,f_2\}}(g)$.

\emph{Left inverses.} Let $\Psi(\mathbf S)\ne0$. Then $x(\mathbf S)=x_\gamma$ with
$\widehat\Psi'_{\mathbf S}(v'(\mathbf S),y_\gamma)\ne0$, so $\Psi'(\mathbf S[y_\gamma,v])\ne0$ for some $v$:
the vertical edges of $w$ carry split injective matrices, and so does $e_2$ if it is vertical
(it is not: $e_2$ is horizontal). The new vertical edge is $f_2$; its opposite old edge $e_1$ is
vertical, so the rows $\mathsf N_1$ of $y_\gamma$ form a split injective matrix and, by
Lemma~\ref{pgr:compl}, so do the columns $\mathsf N_2$ of $x_\gamma$, i.e.\ $x_2$. If $Lx_2=1$, then
$(L\ 0)\binom{x_2}{v'}=1$, so $\mathbf S_{f_2}$ is split injective.

\emph{Values.} $\Psi(\mathbf S)=\chi(\det\gamma)\sum_v\psi(v'v)\Psi'(\mathbf S[y_\gamma,v])$, a root of
unity times a sum of $\#\pcMat_{\mathsf N_2\times\bullet}(R)$ terms, each a root of unity times a
sum of at most $N$ roots of unity. Finally $\#\pcMat_{\mathsf N_2\times\bullet}(R)=(\#R)^n\le
(\#R)^{2n_m}=V(\vartheta)$.
\end{proof}

\begin{lemma}[Sums over the turned square with a line]\label{pgu:sumturn}
In the situation of Proposition~\ref{pgu:turned}, let $J$ be a set, $C\in\C$, and for $i\in J$ let
$L_i$ on $\pcMat_{\mathsf A\times\mathsf N}$ and $L'_i$ on $\pcMat_{\mathsf N\times\mathsf B}$
satisfy the square law for $\chi$ with constant $C$. Let $i(\mathbf S,v')\in J$ depend only on
$(\mathbf S_e)_{e\in w}$ and $v'$. Then
\[
\sum_{\mathbf S\in\mathcal S(w\cup\{e_1,e_2\})}\Psi'(\mathbf S)\sum_{v'}\psi(v'v(\mathbf S))\,
L'_{i(\mathbf S,v')}(y(\mathbf S))=C\sum_{\mathbf S\in\mathcal S(w\cup\{f_1,f_2\})}\Psi(\mathbf S)\,
L_{i(\mathbf S,v'(\mathbf S))}(x(\mathbf S)).
\]
\end{lemma}

\begin{proof}
$(\mathbf S,y,v)\mapsto\mathbf S[y,v]$ is a bijection from $\mathcal S(w)\times\pcMat_{\mathsf
N\times\mathsf B}\times\pcMat_{\mathsf N_2\times\bullet}$ onto $\mathcal S(w\cup\{e_1,e_2\})$, and
$(\mathbf S,x,v')\mapsto\mathbf S[x,v']$ one onto $\mathcal S(w\cup\{f_1,f_2\})$. For fixed
$\mathbf S\in\mathcal S(w)$ the left side is $\sum_{v'}\sum_y\widehat\Psi'_{\mathbf S}(v',y)
L'_{i(\mathbf S,v')}(y)$, which by the square law applied to the matched pair of
Proposition~\ref{pgu:turned} is $C\sum_{v'}\sum_x\Psi(\mathbf S[x,v'])L_{i(\mathbf S,v')}(x)$; sum
over $\mathbf S$.
\end{proof}

\subsubsection{Proof of Proposition~\ref{pgr:turndiag}}\label{pgu:ss-proof}

Let $\pgrG'\pgrTr_{(j,j)}\pgrG$ be a turn of type (diag): $\pgrG'=(j,h'_0,h'_1)$ where $h'_z=q_z\cdot
\bigl((j,j+1)\to(j,j)\bigr)$, and $\pgrG=(j+1,h_0,h_1)$ with $h_z=q_z\cdot\bigl((j,j+1)\to(j+1,j+1)
\bigr)$, $q_z$ a grid path from $(0,m)$ to $(j,j+1)$. Let $w$ be the set of edges of $q_0$ (first
grid) and $q_1$ (second grid). The nodes of $q_z$ lie in $\{(a,b):a\le j,\ b\ge j+1\}$, so no edge of
$w$ meets $V',\widetilde V'$ or ends at $V,\widetilde V$, and $\pgrE(\pgrG')=w\cup\{e_1,e_2\}$,
$\pgrE(\pgrG)=w\cup\{f_1,f_2\}$. For a period triple $x=(\vartheta,y',\varphi)$ and an edge tuple
$\mathbf S$ let $f_{\mathbf S}=q_{0*}(\mathbf S)y'$ (a conjugate form at $W_1$) and
$\varphi_{\mathbf S}=q_{1*}(\mathbf S)\varphi$ (a form at $S$); they depend on $\mathbf S$ only through
$w$. We write $T_M$, $\Lambda_M$ for the transposed lift and the lift with the point datum $\delta_M$
and the joint harmonic of the edge concerned.

\paragraph*{Column and row periods}
For a real point $z$ of $\auU(S)$, a polynomial $B$ on $\auHom(L,S)\otimes\R$, a polynomial $B'$ on
$\auHom(S,L')\otimes\R$ and families $\Psi'$, $\Psi$ of path data on $\pgrE(\pgrG')$,
$\pgrE(\pgrG)$, put
\begin{align*}
I^{\mathrm c}(\Psi';z,B)(x)&=\sum_{\mathbf S\in\mathcal S_\vartheta(\pgrG')}\Psi'_\vartheta(\mathbf S)
\int_{[\auU(S)]}\pguK^{\mathrm c}_{v(\mathbf S),B}(\pgupt^j,b)\,\varphi_{\mathbf S}(b)\,
\bigl(\Lambda_{y_2}T_{\mathbf S_{e_1}}f_{\mathbf S}\bigr)(bz)\,db,\\
I^{\mathrm r}(\Psi;z,B')(x)&=\sum_{\mathbf S\in\mathcal S_\vartheta(\pgrG)}\Psi_\vartheta(\mathbf S)
\int_{[\auU(S)]}\pguK^{\mathrm r}_{v'(\mathbf S),B'}(b,\pgupt^{j+1})\,\varphi_{\mathbf S}(b)\,
\bigl(T_{x_2}\Lambda_{\mathbf S_{f_1}}f_{\mathbf S}\bigr)(bz)\,db,
\end{align*}
where $\Lambda_{y_2}$ is the lift along $(V',S_1)$, $T_{x_2}$ the transposed lift along $(S_1,V)$,
and the forms on $S_1$ are read on $S$ by (G1). The integrals converge absolutely (rapid decrease of
$\varphi_{\mathbf S}$, boundedness of the other form, (R2) of \S\ref{pcx:ss-conv} for the line
kernel).

\begin{lemma}[First seesaw]\label{pgu:see1}
Let $\Xi$ be a function of period triples. If for every $z$ and $B$ there is $\beta(z,B)$ with
$I^{\mathrm c}(\Psi';z,B)=\beta(z,B)\,\Xi$, then $\Psi'\propto_{\pgrG'}\Xi$.
\end{lemma}
\begin{proof}
Apply Lemma~\ref{pgu:seesaw} with $(S_0,T)=(V',S_1)$, $(T',S')=(\widetilde V',S)$, $W_x=\mathcal
S_\vartheta(\pgrG')$, $u_{x,\mathbf S}=\Psi'_\vartheta(\mathbf S)$, $y_{x,\mathbf S}=T_{\mathbf
S_{e_1}}f_{\mathbf S}$, $f_{x,\mathbf S}=\varphi_{\mathbf S}$, $N=1$, $j(a)=a\oplus\pgupt^j$,
$\Psi'_{x,\mathbf S}=\delta_{\mathbf S_{e_2}}$, $\Phi'$ the harmonic of $e_2$. The conclusion of the
lemma at a real point $r$ of $\auU(\widetilde V')$ is $I_{\pgrG'}(\Psi';r)=\beta\,\Xi$ by (G4) and the
definition of $T_{\mathbf S_{e_2}}$ (the half $h'_1$ ends with the transposed lift along $e_2$). The
polynomials on $\auHom(\widetilde V',S)\otimes\R=\auHom(V',S)\otimes\R\oplus\auHom(L,S)\otimes\R$
are spanned by the products $A\otimes B$. Each product is matched: take $Q=A$, $\Psi_{x,\mathbf
S}=\delta_{y_2(\mathbf S)}$, $M_{x,\mathbf S}(b)=c_1\pguK^{\mathrm c}_{v(\mathbf S),B}(\pgupt^j,b)$ and
$j'$ the identification of (G1); then (b) is Corollary~\ref{pgu:split}, and the sum in (a) at
$r'=z$ is $c_1I^{\mathrm c}(\Psi';z,B)$, a multiple of $\Xi$ by hypothesis.
\end{proof}

\begin{lemma}[Second seesaw]\label{pgu:see2}
If $\Psi\propto_{\pgrG}\Xi$, then for every $z$ and $B'$ there is $\beta$ with $I^{\mathrm
r}(\Psi;z,B')=\beta\,\Xi$.
\end{lemma}
\begin{proof}
Apply Lemma~\ref{pgu:seesaw} with $(S_0,T)=(S,\widetilde V)$, $(T',S')=(S_1,V)$, $W_x=\mathcal
S_\vartheta(\pgrG)$, $u_{x,\mathbf S}=\Psi_\vartheta(\mathbf S)$, $y_{x,\mathbf S}=\varphi_{\mathbf S}$,
$f_{x,\mathbf S}=\Lambda_{\mathbf S_{f_1}}f_{\mathbf S}$, $N_{x,\mathbf S}(b)=c_2\pguK^{\mathrm
r}_{v'(\mathbf S),B'}(b,\pgupt^{j+1})$, $j$ the identification of (G1), $\Psi'_{x,\mathbf
S}=\delta_{x_2(\mathbf S)}$, $\Phi'$ the harmonic of $(S_1,V)$, and $r=z$; the conclusion is the
claim (up to the factor $c_2$). Every polynomial $A'$ on $\auHom(S_1,V)\otimes\R$ is matched: take
$Q=A'\otimes B'$, $\Psi_{x,\mathbf S}=\delta_{\mathbf S_{f_2}}$, $M=1$ and $j'(a)=a\oplus\pgupt^{j+1}$.
Then (b) is Corollary~\ref{pgu:split}, and the sum in (a) at a real point $r'$ of
$\auU(\widetilde V)$ is $I_{\pgrG}(\Psi;r')$ by (G4), a multiple of $\Xi$ by hypothesis.
\end{proof}

\begin{theorem}[Proposition~\ref{pgr:turndiag}]\label{pgu:turndiag}
For $0\le j<m$ put
$c_{jj}=c^0_{jj}/c'_{\mathrm{line}}$, with $c^0_{jj}$ of Proposition~\ref{pgr:square} ($z=0$,
$a=b=j$) and $c'_{\mathrm{line}}$ of Lemma~\ref{pgx:line}. Then $c_{jj}$ has modulus $\rho'_{jj}$ up
to norm powers, and every turn $\pgrG'\pgrTr_{(j,j)}\pgrG$ transfers periods with bound $V$ and
constant $c_{jj}$.
\end{theorem}

\begin{proof}
We use Proposition~\ref{pgr:square}, which is Theorem~\ref{pgv:square}, and the facts (G1)--(G4), which are
Propositions~\ref{pgv:G1}, \ref{pgv:G2}, \ref{pgv:G4} and Theorem~\ref{pgv:G3}.
\emph{Modulus.} $|c^0_{jj}|^2=C(\auNm Q_\vartheta)^e\rho'^2_{jj}$ (with $e$ the norm-power exponent of Definition~\ref{pc:modulus}) and $|c'_{\mathrm{line}}|^{-2}=
C_0^{-1}(\auNm\pcm)^n$ with $n=r+(2j+1)t$; as $\auNm\pcm=\auNm\pcb\,\auNm Q_\vartheta$ with $\pcb=(\pi_M)^b$,
$|c_{jj}|^2=C\,C_0^{-1}(\auNm\pcb)^{n}(\auNm Q_\vartheta)^{e+n}\rho'^2_{jj}$.

\emph{The datum.} Let $\Psi'$ be a family of path data on $\pgrE(\pgrG')$ with bounds $N$. For each
$\vartheta$ let $\Psi_\vartheta$ be the datum of Proposition~\ref{pgu:turned} (with $w$ as above,
$\psi=\psi_\pcm$ of Lemma~\ref{pgx:line}, and the bijection $\varepsilon$ of the square law of
Proposition~\ref{pgr:square}); it is a path datum on $\pgrE(\pgrG)$ with bound $V(\vartheta)N(\vartheta)$.

\emph{The identity.} We show $I^{\mathrm c}(\Psi';z,B)(x)=c'_{\mathrm{line}}(\vartheta)\,\kappa\,
c^0_{jj}(\vartheta)^{-1}\,I^{\mathrm r}(\Psi;z,\tau B)(x)$. By Lemma~\ref{pgx:line} at $k=\pgupt^j$,
\[
\begin{gathered}I^{\mathrm c}(\Psi';z,B)(x)=c'_{\mathrm{line}}\sum_{\mathbf S}\Psi'(\mathbf S)\sum_{v'}\psi_\pcm(v'v(\mathbf
S))\,l_{\mathbf S,v'}\bigl(\Lambda_{y_2}T_{\mathbf S_{e_1}}f_{\mathbf S}\bigr),\\
l_{\mathbf S,v'}(F)=\int_{[\auU(S)]}\pguK^{\mathrm r}_{v',\tau B}(b,\pgupt^{j+1})\varphi_{\mathbf S}(b)F(bz)\,db\end{gathered}
\]
(a finite sum of absolutely convergent integrals). Here $l_{\mathbf S,v'}$ is a linear functional on
the conjugate forms at $S_1$, depending on $\mathbf S$ only through $w$, and
$l(\Lambda_{y_2}T_{\mathbf S_{e_1}}f)=L'_{f,l}(y(\mathbf S))$ in the notation of \S\ref{pgr:ss-off}
for the square $V'=(j,j)$, $W_1=(j,j+1)$, $W_2=S_1=(j+1,j)$, $V=(j+1,j+1)$ of the first grid.
By Proposition~\ref{pgr:square} ($z=0$, $a=b=j$), the functions $L_{f,l},L'_{f,l}$ satisfy the square
law for $(\chi_{e_1}\chi_{(V',S_1)})^{-1}=\chi$ (G2) with constant $\kappa/c^0_{jj}(\vartheta)$. By
Lemma~\ref{pgu:sumturn} with $J$ the pairs $(f,l)$ and $i(\mathbf S,v')=(f_{\mathbf S},l_{\mathbf
S,v'})$,
\[
I^{\mathrm c}(\Psi';z,B)(x)=c'_{\mathrm{line}}\frac{\kappa}{c^0_{jj}}\sum_{\mathbf S\in\mathcal
S_\vartheta(\pgrG)}\Psi(\mathbf S)\,L_{f_{\mathbf S},l_{\mathbf S,v'(\mathbf S)}}(x(\mathbf S)),
\]
and $L_{f,l}(x(\mathbf S))=l(T_{x_2}\Lambda_{\mathbf S_{f_1}}f)$, so the right side is
$c'_{\mathrm{line}}\kappa\,c^0_{jj}{}^{-1}I^{\mathrm r}(\Psi;z,\tau B)(x)$.

\emph{Transfer.} Let $\Psi\propto_{\pgrG}\Xi$. By Lemma~\ref{pgu:see2}, $I^{\mathrm r}(\Psi;z,\tau B)
=\beta(z,B)\Xi$ for every $z,B$. By the identity, $I^{\mathrm c}(\Psi';z,B)=\kappa\beta(z,B)\,
\Xi/c_{jj}$, where $\Xi/c_{jj}$ is the function $x\mapsto\Xi(x)/c_{jj}(\vartheta)$ and $\kappa\beta$
does not depend on $x$. By Lemma~\ref{pgu:see1} applied to $\Xi/c_{jj}$, $\Psi'\propto_{\pgrG'}
\Xi/c_{jj}$.
\end{proof}




\section{The grids and the first path}\label{C1:g5}
% ===== C-grid4.tex =====
% proofs/C-grid4.tex
% Continuation of proofs/C-grid.tex (pgr:), C-grid2.tex (pgs:), C-grid3.tex (pgu:): the complete
% definition of the two grids, the facts (G1)--(G4) of \S pgu:ss-data, and Proposition pgr:square.
% Labels start with pgv:. Conventions and macros: those of proofs/C-grid.tex and C-determinant.tex.
\providecommand{\auI}{\mathcal I}
\providecommand{\auGL}{\mathrm{GL}}
\providecommand{\auU}{\mathrm U}
\providecommand{\auO}{\mathcal O}
\providecommand{\auA}{\mathbb A}
\providecommand{\auNm}{\operatorname{Nm}}
\providecommand{\aubar}[1]{\overline{#1}}
\providecommand{\pcm}{\mathfrak m}
\providecommand{\pcb}{\mathfrak b}
\providecommand{\pcMat}{\mathrm M}
\providecommand{\pgrY}{\widetilde Y}
\providecommand{\pcnot}[1]{\par\smallskip{\raggedright\noindent\textsc{Not proved here.} [#1]\par}\smallskip}
\providecommand{\pgupt}{\mathsf u}
\providecommand{\pgvsgn}{\operatorname{sgn}}
\providecommand{\pgvwt}{\operatorname{wt}}
\providecommand{\pzF}{\mathfrak F}
\providecommand{\auPl}{\Phi}

\subsection{The grids: definition, the facts (G1)--(G4), and the square law}\label{pgv:sec}

This section continues \S\S\ref{pgr:sec}, \ref{pgs:sec}, \ref{pgu:sec} and uses their notation.
\S\ref{pgr:ss-arch} describes the two grids briefly, and \S\ref{pgu:ss-data} isolates four
properties (G1)--(G4) of them on which the diagonal turn (Theorem~\ref{pgu:turndiag}) rests. Here
\S\ref{pgv:ss-def} gives the complete definition (spaces, frames, signs, parameters, characters of
every edge, the moves, the levels and the diagonal period), \S\ref{pgv:ss-G} proves (G1), (G2), (G4)
and discusses the unit property of the idele $\alpha$ at the primes of $\pcm$ assumed in
\S\ref{pgu:ss-line}, \S\ref{pgv:ss-G3} proves (G3) (every pair of the grids is an edge, and (H)
holds), and \S\ref{pgv:ss-square} deduces Proposition~\ref{pgr:square} from Theorem~\ref{ptr:square}.


Throughout, $r$, $k=r+1-m$, $t=5r$, $\Delta=2r^8$, $\kappa$, $\kappa_l=\kappa+(l+\frac12)\Delta$ and
the characters $\mu^0_l$ are those of the data ({\cite[\S14.3]{P46}}); $r$ and $t$ are odd, $m$ is
even, so $k$ is even; $\kappa_l(\tau)$ is an integer ($\Delta$ is even). We write $\upsilon$ also for
the restriction of $\upsilon$ to $\auU(1)$, of weight $1$ and unramified at every finite place, and
$\pgvwt$ for weights of characters of $\auU(1)$, normalised as in \S\ref{ptr:ss-param}, so that
$\pgvwt\mu^0_l=-\kappa_l$ and $\pgvwt\vartheta=0$ for a twist. For an $m$-fold twist $\vartheta$,
$\pcm=\pcm(\vartheta)=\pcb\prod_{w\in Q_\vartheta}w$ with $\pcb=(\pi_M)^b$, and $R=\auO_E/\pcm$, as
in \S\ref{pgr:ss-arch}. Since the primes of $\pi_M$ are prime to their conjugates and to $Q_\vartheta$
({\cite[\S14.3]{P46}}: $Q\cap\aubar Q=\emptyset$, $\auNm Q$ prime to $M$), $\pcm$ is a modulus in the
sense of \S\ref{pc:ss-data}.

\subsubsection{Definition of the grids}\label{pgv:ss-def}

\paragraph*{Nodes, frames, ranks and signs}
A \emph{node} is a triple $p=(z,a,b)$ with $z\in\{0,1\}$ (the grid) and $0\le a,b\le m$. Its
\emph{frame} $\mathsf F_p$ is the disjoint union of
\begin{itemize}
\item the \emph{stem}, with coordinates numbered $0\le\varrho<r+z$;
\item $a$ \emph{blocks of the first kind}, numbered $0\le l<a$, each with $t+z$ coordinates
(positions $0\le\beta<t+z$);
\item $b$ \emph{blocks of the second kind}, numbered $0\le l<b$, each with $t-z$ coordinates
(positions $0\le\beta<t-z$).
\end{itemize}
So the rank is
\begin{equation}\label{pgv:eq-rank}
N_p=N^z_{a,b}=r+z+a(t+z)+b(t-z).
\end{equation}
For a stem coordinate $\varrho$ put $a(\varrho)=\max(\varrho+1-k,0)$; as $r+1=k+m$, $a(\varrho)=0$ for
$\varrho\le k-1$ and $a(\varrho)=\varrho+1-k$ runs through $1,\dots,m-1+z$ for $k\le\varrho<r+z$. The
\emph{signs} at $p$ are
\begin{equation}\label{pgv:eq-sign}
\epsilon_p(\varrho)=(-1)^{a(\varrho)+\min(a(\varrho),b)}\ \ (\varrho\text{ in the stem}),\qquad
\epsilon_p=+1\text{ on the blocks}.
\end{equation}
Thus $\epsilon_p(\varrho)=+1$ if $a(\varrho)\le b$ and $\epsilon_p(\varrho)=(-1)^{a(\varrho)+b}$ if
$a(\varrho)>b$; the signs depend on $z$ and $b$ only (and $z$ only through the length of the stem).
For $z=0$, $a=0$ these are the signs $\epsilon^{(0,b)}$ of $Y_b$ (\eqref{ptr:eq-E} with $i=0$, the
stem being numbered there from $1$).

\begin{definition}[The spaces]\label{pgv:def-spaces}
The space at the node $p$ is the following space with frame $\mathsf F_p$ (spaces in the sense of
\S\ref{pc:ss-data}):
\begin{itemize}
\item $p=(0,0,b)$: the space $Y_b$ with its frame and lattice ({\cite[\S5.1.3]{P49}});
\item $p=(1,0,0)$: $Y_0\oplus E\tilde e$ with $\langle\tilde e,\tilde e\rangle=1$, the frame of $Y_0$
followed by $\tilde e$ as the stem coordinate $\varrho=r$;
\item every other node: $E^{N_p}$ with the standard lattice and the diagonal hermitian form with the
entries $\epsilon_p$ in the frame.
\end{itemize}
The node $(0,a,b)$ is written $Y^a_b$ (\emph{first grid}) and the node $(1,a,b)$ is written
$\pgrY^b_a$ (\emph{second grid}): \emph{$\pgrY^b_a$ has $a$ blocks of rank $t+1$ and $b$ blocks of
rank $t-1$}. So $Y^0_b=Y_b$, $\pgrY:=\pgrY^m_0$ (no block of rank $t+1$, $m$ blocks of rank $t-1$,
rank $\tilde n=r+1+m(t-1)$) is the space of \S\ref{pc:ss-aux}, and a half $h_1$ of a two-grid path
runs through the nodes $(1,a,b)=\pgrY^b_a$, as in \S\ref{pgr:ss-arch}.
\end{definition}

\begin{remark}[Indexing]\label{pgv:rem-index}
In the second grid the upper index counts the blocks of rank $t-1$ and the lower index those of rank
$t+1$: $\pgrY^b_a$ has $a$ blocks of rank $t+1$ and $b$ of rank $t-1$. With the opposite convention
$\pgrY^m_0$ would have rank $r+1+m(t+1)$, not $\tilde n$, and the node $(j,j+1)$ of $h_1$ would not
have the rank of $Y^{j+1}_j$ (compare Proposition~\ref{pgv:G1}).
\end{remark}

\paragraph*{Parameters and types}
For a node $p=(z,a,b)$ and $\tau\in\auPl$ let $s_p=\frac12$ if $z=0$ and $a+b$ is odd, $s_p=0$
otherwise, and let $\lambda_{p,\tau}$ be the vector on $\mathsf F_p$ with entries
\begin{equation}\label{pgv:eq-param}
\begin{aligned}
&\varsigma'_{\tau,\varrho}+\tfrac z2-s_p&&\text{(stem coordinate }\varrho),\\
&\kappa_{l+1}(\tau)+\tfrac{t+z-1}2-\beta-s_p&&\text{(position }\beta\text{ of the block }l\text{ of the first kind}),\\
&-\kappa_{l+1}(\tau)+\tfrac{t-z-1}2-\beta-s_p&&\text{(position }\beta\text{ of the block }l\text{ of the second kind}),
\end{aligned}
\end{equation}
where $\varsigma'_{\tau,\varrho}=\frac k2-1-\varrho$ for $\varrho\le k-2$ and
$\varsigma'_{\tau,\varrho}=-\kappa(\tau)-(\varrho+2-k)\Delta=\frac\Delta2-\kappa_{\varrho+2-k}(\tau)$
for $\varrho\ge k-1$. (This is the description of \S\ref{pgs:ss-short}, written without doubling.)
The \emph{type at $p$}, $e_p$, is the type of $\auU(\mathsf F_p)_\infty$ whose Harish-Chandra
parameter at $\tau$ is $-\lambda_{p,\tau}$, with the signs $\epsilon_p$. The forms \emph{moved} at
$p$ are the forms of type $e_p$, of any level: in the first grid these are the conjugates of forms
of the type with parameter $\lambda_p$ (the \emph{conjugate forms} of \S\ref{pgr:ss-arch}); in the
second grid the parameter of the space $\pgrY^b_a$ is $-\lambda_{(1,a,b)}$ and they are forms. For
$p=(0,0,n)$, $\lambda_p=\Lambda^{(0,n)}$ of \S\ref{ptr:ss-param}, so $e_{(0,0,n)}=\overline{e_{0,n}}$.

\paragraph*{Edges and their characters}
An \emph{edge} is a pair $(p,q)$ of nodes of one grid $z$ with $q=(z,a+1,b)$ (a \emph{right} edge;
its \emph{index} is $l=a$) or $q=(z,a,b+1)$ (an \emph{up} edge; index $l=b$), where $p=(z,a,b)$. Its
\emph{kind} is $c=1$ (\emph{vertical}) if it is an up edge of the first grid or a right edge of the
second grid, and $c=0$ (\emph{horizontal}) otherwise; this agrees with \S\ref{pgr:ss-arch}. Note
$N_q-N_p=d_{pq}$ with $d_{pq}=t+z$ for right edges and $t-z$ for up edges.

\begin{definition}[Characters, kernels, moves]\label{pgv:def-edges}
Let $\vartheta$ be an $m$-fold twist and $(p,q)$ an edge of kind $c$ and index $l$. Put
$\nu_{pq}=\mu^0_{l+1}\vartheta^{c(l+1)}$ and
\[
\alpha_{pq}=\nu_{pq}\,\upsilon^{\lfloor(N_q+1-c)/2\rfloor},\qquad
\beta_{pq}=\nu_{pq}\,\upsilon^{\lfloor(N_p+1-c)/2\rfloor}
\]
(so the power of $\upsilon$ is $\lfloor N/2\rfloor$ for vertical and $\lceil N/2\rceil$ for horizontal
edges). The \emph{kernel} of the edge with Fock vector $P$ and finite datum $\Psi$ on
$\pcMat_{q\times p}(R)$ (rows indexed by $\mathsf F_q$, columns by $\mathsf F_p$) is
$\theta_{\alpha_{pq},\beta_{pq}}(g_p,g_q;P,\Psi)$ for $\auHom(p,q)$, in the notation of
\S\ref{pc:ss-data}, with the modulus $\pcm(\vartheta)$; recall $\omega_{\alpha,\beta}(g,h)=
\alpha(\det g)^{-1}\beta(\det h)\omega(g,h)$, $g$ on the source $p$, $h$ on the target $q$. The
\emph{edge character} $\chi_{pq}$ is the character $x\mapsto\prod_{w\mid\pcm}\psi_w(x\bmod w)$ of
$R^\times$, $\psi_w$ the tame part of $\nu_{pq}$ at $w$ (trivial at the $w\mid\pcm$ where $\nu_{pq}$
is unramified). Let $P_{pq}$ be the joint harmonic of the edge (Theorem~\ref{pgv:G3}(a)). For
$X\in\pcMat_{q\times p}(R)$ the \emph{move} along the edge is
\[
\Lambda_X=\Theta_{\alpha_{pq},\beta_{pq}}(P_{pq},\delta_X)\colon\{\text{forms of type }e_p\}\to\{\text{forms
of type }e_q\}\quad\text{if }(p,q)\text{ is a right edge},
\]
\[
T_X=\Theta^{\mathsf t}_{\alpha_{pq},\beta_{pq}}(P_{pq},\delta_X)\colon\{\text{forms of type }e_q\}\to
\{\text{forms of type }e_p\}\quad\text{if }(p,q)\text{ is an up edge}.
\]
That these maps are defined and take values as stated is Theorem~\ref{pgv:G3}(d). The definition is thus made in two stages. The characters and the kernels are defined here outright, and the proof of Theorem~\ref{pgv:G3} uses nothing else of this definition; the polynomial $P_{pq}$ and the moves are defined once parts (a) and (d) of that theorem are proved, and they are not used before.
\end{definition}

\noindent So $\pgvwt\alpha_{pq}=-\kappa_{l+1}+\lfloor\frac{N_q+1-c}2\rfloor$ and $\pgvwt\beta_{pq}=
-\kappa_{l+1}+\lfloor\frac{N_p+1-c}2\rfloor$. For the edge $(0,0,n)\to(0,0,n+1)$ (vertical, index $n$)
the characters are $\mu_{n+1}\upsilon_{r+(n+1)t}$, $\mu_{n+1}\upsilon_{r+nt}$, those of $a_{0,n}$ in
\S\ref{pc:ss-triangle}. The edge character depends only on the kind and the index, and $\chi_{pq}$ is
also the character of $R^\times$ through which $\alpha_{pq}$ and $\beta_{pq}$ act on $\det$ of
elements of $K$ (Lemma~\ref{pc:level}; $\upsilon$ is unramified).

\paragraph*{Levels, admissible families and the diagonal period}
The \emph{level} of the node $p$ is the subgroup of $\auGL_{\mathsf F_p}(R)$ equal to $\{1\}$ if
$p=(0,0,b)$ and to the kernel of $\auGL_{\mathsf F_p}(R)\to\auGL_{\mathsf F_p}(\auO_E/\pcb)$
otherwise. For $0\le l\le m$ we identify $\pgrY^l_l$ with $Y^l_l\oplus E\tilde e_l$,
$\langle\tilde e_l,\tilde e_l\rangle=(-1)^l$, as spaces with frames: for $l=0$ this is
Definition~\ref{pgv:def-spaces}, and for $l\ge1$ we fix a bijection
$\mathsf F_{(0,l,l)}\sqcup\{\bullet\}\to\mathsf F_{(1,l,l)}$ preserving signs, $\bullet$ having the
sign $(-1)^l$ (it exists by Lemma~\ref{pgv:count}(c)); both spaces are then standard diagonal spaces
with the same signs, and the bijection is an isometry carrying frame to frame and lattice to lattice.
A node family $g=(g_p)$ is \emph{admissible} if $g_{(0,0,m)}=1$, $g_{(1,0,m)}\equiv1\bmod\pcb$, and
$g_{(1,a,a)}=g_{(0,a,a)}\oplus1$ for all $a$ under this identification. (This makes precise the
definition in \S\ref{pgr:ss-arch}: the level of $Y^0_m=Y_m$ reduces to $\{1\}$ modulo $\pcm$, that of
$\pgrY^m_0$ is $J(\pcb)$.)

\begin{definition}[The diagonal period]\label{pgv:def-period}
Let $\varpi_\pcm\in\auA_{E,f}^\times$ be a finite idele that is $1$ at the primes not dividing
$\pcm$ and has order $\operatorname{ord}_w\pcm$ at each $w\mid\pcm$ (it generates $\pcm\widehat
\auO_E$), and put $u_0=\alpha\varpi_\pcm$ and $\pgupt=\pgupt_\pcm=\aubar u_0/u_0$, an element of
$\auU(1)(\auA_{E^+,f})$ depending only on $\pcm$ (and the fixed idele $\alpha$ of
{\cite[\S5.1.1]{P49}}). For $0\le l\le m$, a real point $g$ of $\auU(\pgrY^l_l)$, a form $y''$ of type
$e_{(0,l,l)}$ and a form $\varphi''$ of type $e_{(1,l,l)}$, the diagonal period is
\[
\mathsf P_{l,g}(y'',\varphi'')=\int_{[\auU(Y^l_l)]}y''(x)\,\varphi''\bigl((x\oplus\pgupt^l)\,g\bigr)\,dx ,
\]
$x\oplus\pgupt^l$ acting on $Y^l_l\oplus E\tilde e_l$ (the line component has trivial archimedean
part).
\end{definition}

\noindent With these definitions the periods $I_\Gamma(\Psi;g)$, path data and transfers of
\S\ref{pgr:ss-arch} are defined without further choice. The integral converges absolutely: $y''$ and
$\varphi''$ are bounded (Theorem~\ref{pgv:G3}(c)) and $[\auU(Y^l_l)]$ has finite volume.

\subsubsection{The facts (G1), (G2), (G4); the idele \texorpdfstring{$\alpha$}{alpha} at \texorpdfstring{$\pcm$}{m}}\label{pgv:ss-G}

\begin{lemma}[Ranks and signs]\label{pgv:count}
Let $p=(z,a,b)$ be a node.
\begin{enumerate}
\item[(a)] $N^1_{a,b}$ is even; $N^0_{a,b}$ is odd if and only if $a+b$ is even.
\item[(b)] The number of negative signs at $p$ is $\lfloor(m-b+z)/2\rfloor$.
\item[(c)] For $0\le l\le m$: $N^1_{l,l}=N^0_{l,l}+1$, and the number of negative signs at $(1,l,l)$
is that at $(0,l,l)$ plus $1$ if $l$ is odd and plus $0$ if $l$ is even.
\item[(d)] For $0\le j<m$: $N^1_{j,j+1}=N^0_{j+1,j}=r+(2j+1)t$, and $(1,j,j+1)$, $(0,j+1,j)$ have the
same number $\lfloor(m-j)/2\rfloor$ of negative signs.
\end{enumerate}
\end{lemma}
\begin{proof}
(a) $N^1_{a,b}=r+1+a(t+1)+b(t-1)$ with $r+1$, $t\pm1$ even; $N^0_{a,b}=r+(a+b)t$ with $r,t$ odd.
(b) Only stem coordinates can be negative, and the coordinate with $a(\varrho)=a'$ is negative iff
$a'>b$ and $a'+b$ is odd. The values $a'=a(\varrho)\ge1$ are $1,\dots,m-1+z$, each once; so the count
is the number of $a'\in(b,m-1+z]$ with $a'\equiv b+1\pmod 2$, which is $\lfloor(m-1+z-b+1)/2\rfloor$.
(c) The ranks: $r+1+l(t+1)+l(t-1)=r+2lt+1$. By (b) the counts are $\lfloor(m-l+1)/2\rfloor$ and
$\lfloor(m-l)/2\rfloor$; as $m$ is even they are equal if $l$ is even and differ by $1$ if $l$ is
odd. (d) $r+1+j(t+1)+(j+1)(t-1)=r+(2j+1)t=r+(j+1)t+jt$; by (b) the counts are
$\lfloor(m-(j+1)+1)/2\rfloor$ and $\lfloor(m-j)/2\rfloor$.
\end{proof}

\noindent Recall the nodes of a diagonal turn at $(j,j)$ (\S\ref{pgu:ss-data}): $V'=Y^j_j=(0,j,j)$,
$W_1=(0,j,j+1)$, $V=Y^{j+1}_{j+1}$, $S_1=(0,j+1,j)$ in the first grid; $\widetilde V'=\pgrY^j_j=
(1,j,j)$, $S=(1,j,j+1)$ (the node at the place $(j,j+1)$ of $h_1$, i.e.\ $\pgrY^{j+1}_j$),
$\widetilde V=\pgrY^{j+1}_{j+1}$ in the second; $L=E\tilde e_j$, $L'=E\tilde e_{j+1}$.

\begin{proposition}[(G1)]\label{pgv:G1}
For $0\le j<m$ there is a bijection $\sigma\colon\mathsf F_S\to\mathsf F_{S_1}$ preserving the signs,
and for any such $\sigma$ the map $E^{\mathsf F_S}\to E^{\mathsf F_{S_1}}$ it induces is an isometry
of spaces with frames (frame to frame, standard lattice to standard lattice). Moreover
$N_{V'}+N_V=N_{W_1}+N_S=2(r+(2j+1)t)$.
\end{proposition}
\begin{proof}
$S$ and $S_1$ are not among the exceptional nodes of Definition~\ref{pgv:def-spaces} (the first index
of $S_1$ and the second index of $S$ are $\ge1$), so both are standard diagonal spaces. By
Lemma~\ref{pgv:count}(d) they have the same rank and the same number of negative signs, so a
sign-preserving bijection of the frames exists, and it carries the diagonal form with entries
$\epsilon_S$ to the one with entries $\epsilon_{S_1}$ and the standard lattice to the standard
lattice. The ranks: $N_{V'}+N_V=(r+2jt)+(r+2(j+1)t)$ and $N_{W_1}=N_S=r+(2j+1)t$.
\end{proof}

\noindent The types at $S$ and $S_1$ are \emph{not} identified by $\sigma$ (for instance the first
kind blocks of $S$ have $t+1$ entries and those of $S_1$ have $t$); (G1) identifies spaces, groups,
lattices and frames only, which is all that \S\ref{pgu:sec} uses.

\begin{proposition}[(G2)]\label{pgv:G2}
For $0\le j<m$ and every $m$-fold twist, with the characters of Definition~\ref{pgv:def-edges},
\[
\alpha_{\widetilde V'S}=\alpha_{V'S_1},\quad\beta_{\widetilde V'S}=\beta_{V'S_1},\quad
\alpha_{S\widetilde V}=\alpha_{S_1V}\,\upsilon,\quad\beta_{S\widetilde V}=\beta_{S_1V}.
\]
Hence (G2) of \S\ref{pgu:ss-data} holds with $\alpha_L=\alpha_{V'S_1}=\mu^0_{j+1}\upsilon^{N_S/2}$,
$\beta^{\mathrm c}=1$, $\alpha^{\mathrm r}=\upsilon$ and $\beta_{L'}=\beta_{S_1V}=\mu^0_{j+1}
\vartheta^{j+1}\upsilon^{\lfloor N_S/2\rfloor}$, all unramified outside $\pcm$ and tamely ramified at
$\pcm$; in particular $\alpha^{\mathrm r}=(\beta^{\mathrm c})^{-1}\upsilon$. Moreover
$\chi_{(\widetilde V',S)}=\chi_{(V',S_1)}=\chi_{(W_1,V)}$ and $\chi_{(S,\widetilde V)}=\chi_{(S_1,V)}=
\chi_{(V',W_1)}$; and for every square $V'=(z,a,b)$, $W_1=(z,a,b+1)$, $W_2=(z,a+1,b)$,
$V=(z,a+1,b+1)$ of a grid, $\chi_{(W_1,V)}=\chi_{(V',W_2)}$ and $\chi_{(W_2,V)}=\chi_{(V',W_1)}$.
\end{proposition}
\begin{proof}
The edges $(\widetilde V',S)=((1,j,j),(1,j,j+1))$ and $(V',S_1)=((0,j,j),(0,j+1,j))$ are horizontal
(an up edge of the second grid, a right edge of the first) with index $j$; so $\nu=\mu^0_{j+1}$ for
both, and the powers of $\upsilon$ are $\lceil N/2\rceil$. For $\alpha$ the relevant ranks are those
of the targets, $N_S=N_{S_1}$ (Lemma~\ref{pgv:count}(d)): equal. For $\beta$ they are
$N_{\widetilde V'}=N_{V'}+1$ and $N_{V'}=r+2jt$, which is odd; so $\lceil(N_{V'}+1)/2\rceil=
\lceil N_{V'}/2\rceil$: equal. The edges $(S,\widetilde V)=((1,j,j+1),(1,j+1,j+1))$ and
$(S_1,V)=((0,j+1,j),(0,j+1,j+1))$ are vertical with index $j$, so $\nu=\mu^0_{j+1}\vartheta^{j+1}$ and
the powers are $\lfloor N/2\rfloor$. For $\beta$ the ranks of the sources agree. For $\alpha$ they are
$N_{\widetilde V}=N_V+1$ and $N_V=r+2(j+1)t$, odd, so $\lfloor(N_V+1)/2\rfloor=\lfloor N_V/2\rfloor+1$:
$\alpha_{S\widetilde V}=\alpha_{S_1V}\upsilon$.

Now $\det(h\oplus k)=\det h\cdot k$ for $h\in\auU(V')$, $k\in\auU(L)=\auU(1)$, so
$\alpha_{\widetilde V'S}(\det(h\oplus k))=\alpha_{V'S_1}(\det h)\alpha_L(k)$ with $\alpha_L=
\alpha_{V'S_1}$, and $\beta_{\widetilde V'S}=\beta_{V'S_1}\beta^{\mathrm c}$ with $\beta^{\mathrm c}=1$;
similarly $\alpha_{S\widetilde V}=\alpha_{S_1V}\alpha^{\mathrm r}$ with $\alpha^{\mathrm r}=\upsilon$,
and $\beta_{S\widetilde V}(\det(h\oplus k))=\beta_{S_1V}(\det h)\beta_{L'}(k)$ with $\beta_{L'}=
\beta_{S_1V}$. These are the identities of (G2). $N_S$ is even ($r$ and $(2j+1)t$ are odd). The
characters $\mu^0_{j+1}$ are tamely ramified at the chosen places and unramified elsewhere,
$\vartheta$ is tamely ramified at $Q_\vartheta$ and unramified elsewhere, $\upsilon$ is unramified at
all finite places: so all four characters are unramified outside $\pcm$ and tamely ramified at
$\pcm$.

The edge character depends only on $\nu_{pq}$, i.e.\ on the kind and the index. $(\widetilde V',S)$,
$(V',S_1)$ and $(W_1,V)=((0,j,j+1),(0,j+1,j+1))$ are horizontal of index $j$; $(S,\widetilde V)$,
$(S_1,V)$ and $(V',W_1)=((0,j,j),(0,j,j+1))$ are vertical of index $j$. In a square of the grid $z$,
$(W_1,V)$ and $(V',W_2)$ are right edges of index $a$, and $(W_2,V)$, $(V',W_1)$ are up edges of
index $b$, of the same kind within each pair.
\end{proof}

\begin{proposition}[(G4)]\label{pgv:G4}
The element $\pgupt_\pcm$ of Definition~\ref{pgv:def-period} lies in $\auU(1)(\auA_{E^+,f})$,
depends only on $\pcm$, and the diagonal period is the integral displayed in (G4) of
\S\ref{pgu:ss-data}. At $v\in S(\pcm)$, $v=w\aubar w$ with $w\mid\pcm$, under
$\auU(1)(E^+_v)\cong E_w^\times$ (projection to the $w$-component) $\pgupt_v$ is
$(\alpha_w\varpi_w^{k_w})^{-1}\aubar{\alpha}_{w}$, where $\varpi_w^{k_w}$ is the $w$-component of
$\varpi_\pcm$, $k_w=\operatorname{ord}_w\pcm$, and $\aubar\alpha_w$ is the $w$-component of
$\aubar\alpha$; outside $S(\pcm)$ it is $\aubar\alpha/\alpha$.
\end{proposition}
\begin{proof}
$\aubar u_0/u_0$ has norm $1$, so it is a finite adelic point of $\auU(1)$; it is defined from $\pcm$
and the fixed $\alpha$ only. (G4) is Definition~\ref{pgv:def-period}. At $v=w\aubar w\in S(\pcm)$ the
component of $\varpi_\pcm$ at $\aubar w$ is $1$, because $\aubar w\nmid\pcm$ ($\pcm$ is prime to
$\aubar\pcm$); the $w$-component of $\aubar u_0$ is $\aubar\alpha_w\cdot\aubar{\varpi_\pcm}{}_{,w}=
\aubar\alpha_w$ (the $w$-component of the conjugate is the conjugate of the $\aubar w$-component).
Outside $S(\pcm)$, $\varpi_\pcm=1$.
\end{proof}

\noindent By Lemma~\ref{pgx:line}, $\pgupt_\pcm$ is exactly the element of the line reversal.

\paragraph*{The idele $\alpha$ at the primes of $\pcm$}
\S\ref{pgu:ss-line} assumes that $\alpha$ is a unit at the primes of $\pcm$. By {\cite[\S5.1.1]{P49}},
$\alpha=1/P'(\zeta_{q_1})$ at the primes dividing $\omega=(1+\sqrt{-d})/2$ and $\alpha=1$ elsewhere,
$P$ the minimal polynomial of $\zeta_{q_1}$ over $K$; and $P'(\zeta_{q_1})$ generates the relative
different $\mathfrak d_{E/K}$, which divides $q_1$ (proof of {\cite[Lemma~7.11]{P49}}).

\begin{proposition}\label{pgv:alpha}
\begin{enumerate}
\item[(a)] $\alpha$ is a unit at every prime $w$ of $E$ with $w\nmid q_1$; in particular at every
chosen place, and at every prime of $\pcm(\vartheta)$ if $\auNm Q_\vartheta$ is prime to $q_1$.
\item[(b)] Let $p\mid q_1$ be a prime that ramifies in $\Q(\zeta_{q_1})$ (every odd $p\mid q_1$, and
$p=2$ if $4\mid q_1$). Then $p$ splits in $K$, and every prime $w$ of $E$ above $p$ is split over $E^+$
($w\ne\aubar w$) and ramified over $K$; $\alpha$ is not a unit at the primes $w\mid p$ that divide
$\omega$.
\end{enumerate}
\end{proposition}
\begin{proof}
(a) At $w\nmid\omega$, $\alpha_w=1$. At $w\mid\omega$ with $w\nmid q_1$, $w\nmid\mathfrak d_{E/K}$, so
$P'(\zeta_{q_1})$ is a unit at $w$. A chosen place lies above $p_0$ or $p_1$, and $p_i\equiv1\pmod{q_1}$
({\cite[\S14.3]{P46}}: $p_i\equiv1\pmod{q_2^2}$ and $q_1\mid q_2$), so $p_i\nmid q_1$.

(b) If $p$ is odd: $\omega\aubar\omega=(1+d)/4$ is divisible by every prime dividing $q_1$
({\cite[\S7]{P49}}), so $d\equiv-1\pmod p$, $p\nmid d$, $-d$ is a non-zero square modulo $p$, and $p$
splits in $K$. The prime $2$ splits in $K$ ({\cite[\S7.1]{P49}}). Let $D\subseteq\mathrm{Gal}(E/\Q)$ be
the decomposition group of $w$; its image in $\mathrm{Gal}(K/\Q)$ is the decomposition group of $p$ in
$K$, which is trivial. Complex conjugation restricts to the non-trivial element of $\mathrm{Gal}(K/\Q)$,
so it is not in $D$, i.e.\ $w\ne\aubar w$. As $E\supseteq\Q(\zeta_{q_1})$ is ramified at $p$ and $K$ is
not, $w$ is ramified over $K$, so $w\mid\mathfrak d_{E/K}$ and $P'(\zeta_{q_1})$ is not a unit at $w$.
\end{proof}

\noindent So the assumption of \S\ref{pgu:ss-line} holds at the chosen places, but it does \emph{not}
follow from the definitions for the primes of $Q_\vartheta$ above $q_1$: the definition of a twist
({\cite[\S14.3]{P46}}) requires only $\auNm Q$ prime to $MN$ and $Q\cap\aubar Q=\emptyset$, and by
(b) the primes above $q_1$ meet the second condition. With the data of Definition~\ref{pgx:conv} the line reversal holds without the assumption
(Lemma~\ref{pgx:line}), so the definition of a twist is not changed. The rest of this section does
not use the assumption.

\subsubsection{(G3): the pairs of the grids are edges, and (H) holds}\label{pgv:ss-G3}

\paragraph*{The block rule for an arbitrary excess}
Let $(S,T)$ be spaces of ranks $n<n'$, $d=n'-n$, with characters $\alpha,\beta$ of integral
weights, and at $\tau$ let $\lambda$ (on the frame $\mathsf R$ of $S$, signs $\varepsilon$) and
$\lambda'$ (on the frame $\mathsf C$ of $T$, signs $\varepsilon'$) be parameters with distinct
entries, the entries of $\lambda$ in $\frac{n-1}2+\Z$ and those of $\lambda'$ in $\frac{n'-1}2+\Z$. Put
$\tau_0=\frac{n'}2-\pgvwt_\tau\alpha$ and $\tau'_0=\pgvwt_\tau\beta-\frac n2$ (Lemma~\ref{ptr:const}).
We say that $(\lambda,\lambda')$ satisfies the \emph{block rule of excess $d$} if there are a set
$\mathsf B\subseteq\mathsf C$ of $d$ columns and a bijection $\varrho\mapsto c(\varrho)$ of $\mathsf R$ onto
$\mathsf C\smallsetminus\mathsf B$ with
\begin{equation}\label{pgv:eq-rule}
\begin{gathered}
\lambda'_{c(\varrho)}=-\lambda_\varrho+\tau_0+\tau'_0,\qquad|\lambda'_{c(\varrho)}-\tau'_0|>\tfrac d2,\qquad
\varepsilon'_{c(\varrho)}=\varepsilon_\varrho\,\pgvsgn(\lambda'_{c(\varrho)}-\tau'_0),\\
\{\lambda'_c:c\in\mathsf B\}=\{\tau'_0+\tfrac{d-1}2-\beta:0\le\beta<d\},\qquad\varepsilon'\text{ constant on }
\mathsf B.
\end{gathered}
\end{equation}
For $d=t$ this is \eqref{ptr:eq-rule}.

\begin{lemma}\label{pgv:rule}
In this situation put $\xi_\varrho=\lambda_\varrho-\tau_0$, $\xi'_c=\lambda'_c-\tau'_0$. Then all
$\xi_\varrho,\xi'_c$ lie in $\frac{d-1}2+\Z$, and the block rule of excess $d$ is the matching
condition {\cite[(13)]{P49}} with $t$ replaced by $d$. The statements and proofs of
{\cite[Lemma~3.10]{P49}}, {\cite[Propositions~3.11 and~3.13]{P49}}, {\cite[Lemmas~3.12 and~3.14]{P49}} and of the local identity in the proof of {\cite[Theorem~3.15]{P49}} hold for such
$(\lambda,\lambda')$ with $t$ replaced by $d$.
\end{lemma}
\begin{proof}
The weights are integers, so $\tau_0\in\frac{n'}2+\Z$, $\tau'_0\in\frac n2+\Z$, $\xi_\varrho\in
\frac{n-1-n'}2+\Z=\frac{d-1}2+\Z$ (as $-d-1\equiv d-1\pmod 2$) and $\xi'_c\in\frac{n'-1-n}2+\Z=\frac{d-1}2+\Z$.
The equivalence with {\cite[(13)]{P49}} is the computation recalled after Lemma~\ref{ptr:const}:
$\eta_{\varrho c(\varrho)}=-\varepsilon_\varrho\varepsilon'_{c(\varrho)}$, and
$\eta_{\varrho c(\varrho)}=\pgvsgn\xi_\varrho=-\pgvsgn\xi'_{c(\varrho)}$ is the sign condition. In the cited
statements and proofs $t$ occurs only as the number $n'-n$ of columns of $T$ not matched with rows:
as the size of $\mathsf B$ and of its string $\{\frac{t-1}2-\beta\}$ ({\cite[Lemma~3.10(ii)]{P49}},
{\cite[Proposition~3.11]{P49}}: columns in $\mathsf B$ carry no variable), as the half-integer class
$\frac{t-1}2+\Z$ of the $\xi$ and the bound $|\xi|\ge\frac{t+1}2$ that it implies with $|\xi|>\frac t2$
({\cite[Lemma~3.10(i)]{P49}}, {\cite[Proposition~3.13]{P49}}), and in {\cite[Lemma~3.14]{P49}} through
$|\lambda'|^2=|\lambda|^2-2s_0\sum\lambda+ns_0^2+t\xi_0^2+\frac{t(t^2-1)}{12}$ and $n'=n+t$, i.e.\ again
as $|\mathsf B|=n'-n$. With the first assertion, each of these uses holds with $d$ in place of $t$; no
other property of $t$ is used (the parity of $t$ enters only through the class of the $\xi$).
\end{proof}

\paragraph*{The rule at every edge of the grids}
For an edge $(p,q)$ of the grid $z$ at $\tau$ we take as row parameter (on $\mathsf F_p$) and column
parameter (on $\mathsf F_q$)
\[
(\lambda,\lambda')=(\lambda_{p,\tau},-\lambda_{q,\tau})\ \text{for a right edge},\qquad
(\lambda,\lambda')=(-\lambda_{p,\tau},\lambda_{q,\tau})\ \text{for an up edge},
\]
with the signs $\epsilon_p,\epsilon_q$ and the characters $\alpha_{pq},\beta_{pq}$. In the language of
\S\ref{ptr:ss-edge} (type $e_S$ on $S=p$ with conjugate of parameter $\lambda$, type $e_T$ on $T=q$ with
parameter $\lambda'$): for a right edge $e_S=e_p$, $e_T=e_q$, so the lift goes from forms of type $e_p$
to forms of type $e_q$; for an up edge $e_S=\overline{e_p}$, $e_T=\overline{e_q}$, so the transposed
lift goes from forms of type $\overline{e_T}=e_q$ to forms of type $\overline{e_S}=e_p$. These are the
directions of the moves of Definition~\ref{pgv:def-edges}.

\begin{proposition}[Matching along the edges]\label{pgv:match}
For every edge $(p,q)$ of either grid, every $\tau$ and every $m$-fold twist, the pair
$(\lambda,\lambda')$ just defined satisfies the block rule of excess $d_{pq}=N_q-N_p$, with $c$ the
identity on the stem and on the blocks of $p$ and $\mathsf B$ the new block of $q$ (the first kind
block $a$ for a right edge, the second kind block $b$ for an up edge). The constants are:
\begin{center}
\begin{tabular}{lllll}
\hline
edge&$d$&$\tau_0$&$\tau'_0$&$\tau_0+\tau'_0$\\
\hline
right, $z=0$&$t$&$\kappa_{a+1}-\frac12[a+b\text{ odd}]$&$-\kappa_{a+1}+\frac12[a+b\text{ even}]$&$s_q-s_p$\\
up, $z=0$&$t$&$\kappa_{b+1}+\frac12[a+b\text{ odd}]$&$-\kappa_{b+1}-\frac12[a+b\text{ even}]$&$s_p-s_q$\\
right, $z=1$&$t+1$&$\kappa_{a+1}$&$-\kappa_{a+1}$&$0$\\
up, $z=1$&$t-1$&$\kappa_{b+1}$&$-\kappa_{b+1}$&$0$\\
\hline
\end{tabular}
\end{center}
\end{proposition}
\begin{proof}
Write $\lambda^\circ_p$ for $\lambda_p+s_p$ (the entries of \eqref{pgv:eq-param} without the shift);
on the common coordinates of $p$ and $q$ (stem and old blocks) $\lambda^\circ_q=\lambda^\circ_p$, since the
entries \eqref{pgv:eq-param} do not depend on $a,b$ except through $s$. We use $\kappa_{l+1}\ge
\frac32\Delta$, $\kappa_{l'+1}-\kappa_{l+1}\ge\Delta$ for $l'>l$, $k\le r$, and $\Delta=2r^8>4t+4r$.

\emph{Constants.} By Lemma~\ref{pgv:count}(a) and the weights after Definition~\ref{pgv:def-edges},
$\tau_0=\kappa_{l+1}+\frac{N_q}2-\lfloor\frac{N_q+1-c}2\rfloor$ and $\tau'_0=-\kappa_{l+1}+
\lfloor\frac{N_p+1-c}2\rfloor-\frac{N_p}2$. For $z=1$ all ranks are even and these are $\pm\kappa_{l+1}$.
For $z=0$, a right edge is horizontal ($c=0$, ceilings) and an up edge vertical ($c=1$, floors), and
$N^0_{a,b}$ is odd iff $a+b$ is even; this gives the table. In the last column we used
$s_p=\frac12[a+b\text{ odd}]$, $s_q=\frac12[a+b\text{ even}]$ for $z=0$.

\emph{First condition.} Right edge: it reads $-\lambda_q=-\lambda_p+\tau_0+\tau'_0$ on the common
coordinates, i.e.\ $\lambda^\circ_q-s_q=\lambda^\circ_p-s_p-(s_q-s_p)$: true. Up edge: it reads
$\lambda_q=\lambda_p+\tau_0+\tau'_0$, i.e.\ $\lambda^\circ_q-s_q=\lambda^\circ_p-s_p+(s_p-s_q)$: true. (For
$z=1$, $s=0$.)

\emph{The new block.} Right edge, $z=0$: the block $a$ of the first kind of $q$ carries
$\kappa_{a+1}+\frac{t-1}2-\beta-s_q$, $0\le\beta<t$; the entries of $\lambda'=-\lambda_q$ there are
$-\kappa_{a+1}+s_q-\frac{t-1}2+\beta$, i.e.\ the set $\{\tau'_0+\frac{t-1}2-\beta\}$ as $\tau'_0=
-\kappa_{a+1}+s_q$. Right edge, $z=1$: the block has $t+1$ entries $\kappa_{a+1}+\frac t2-\beta$; their
negatives form $\{\tau'_0+\frac t2-\beta:0\le\beta\le t\}$, the string of excess $t+1$. Up edge, $z=0$:
the block $b$ of the second kind of $q$ carries $-\kappa_{b+1}+\frac{t-1}2-\beta-s_q=\tau'_0+\frac{t-1}2-
\beta$. Up edge, $z=1$: it carries the $t-1$ entries $-\kappa_{b+1}+\frac{t-2}2-\beta=\tau'_0+\frac{(t-1)-1}2-
\beta$. In all cases the signs on the new block are $+1$.

\emph{Distances and signs, right edges.} On a common coordinate $\varrho$,
$\lambda'_{c(\varrho)}-\tau'_0=-\lambda_q(\varrho)+\kappa_{a+1}-s_q=\kappa_{a+1}-\lambda^\circ_p(\varrho)$
(for both $z$). Every entry of $\lambda^\circ_p$ is at most $\kappa_{a+1}-\Delta+\frac t2$: the stem
entries are $\le\frac k2$, the first kind blocks $l<a$ are $\le\kappa_a+\frac t2$, the second kind
blocks are negative. So $\lambda'_{c(\varrho)}-\tau'_0\ge\Delta-\frac t2>\frac d2$ and its sign is $+$:
the rule asks $\epsilon_q(\varrho)=\epsilon_p(\varrho)$. This holds: the stem signs \eqref{pgv:eq-sign}
depend only on $z$ and on the second index, which a right edge does not change, and the block signs
are $+1$.

\emph{Distances and signs, up edges.} Here $\lambda'_{c(\varrho)}-\tau'_0=\lambda_q(\varrho)+\kappa_{b+1}+
s_q=\lambda^\circ_p(\varrho)+\kappa_{b+1}$ (for both $z$). For a first kind block this is
$\ge2\kappa_1-\frac t2>0$, far from $0$. For a second kind block $l<b$ it is
$\kappa_{b+1}-\kappa_{l+1}+\frac{t-z-1}2-\beta\ge\Delta-\frac t2$. For a stem coordinate with
$\varrho\le k-2$ it is $\frac k2-1-\varrho+\frac z2+\kappa_{b+1}\ge\frac32\Delta-r$. For a stem coordinate
with $\varrho\ge k-1$ put $i=\varrho+2-k\ge1$, so that $a(\varrho)=i-1$; then it is
$-\kappa-i\Delta+\frac z2+\kappa+(b+\frac32)\Delta=(b+\frac32-i)\Delta+\frac z2$, of absolute value
$\ge\frac\Delta2-\frac12>\frac d2$, positive iff $i\le b+1$, i.e.\ iff $a(\varrho)\le b$. So the rule asks:
the blocks keep their sign $+1$, and a stem coordinate keeps its sign iff $a(\varrho)\le b$. By
\eqref{pgv:eq-sign}, with $a'=a(\varrho)$: if $a'\le b$ both exponents are $2a'$; if $a'\ge b+1$ the
exponent at $q$ is $a'+b+1$, one more than at $p$. This is the rule.
\end{proof}

\begin{lemma}[Gap and regularity]\label{pgv:gap}
For every node $p$ and $\tau$, $\lambda_{p,\tau}$ consists of $N_p$ distinct numbers in
$\frac{N_p-1}2+\Z$; two of them at coordinates of different signs differ by at least $\Delta-t\ge
\frac12(\Delta-t)-1$; and $N_p\le2n_m$. Hence the conclusions of {\cite[Theorem~3.17]{P49}},
{\cite[Lemma~3.25]{P49}}, {\cite[Theorem~3.2]{P49}} and {\cite[Corollary~3.3]{P49}} hold for the types $e_p$ and
$\overline{e_p}$.
\end{lemma}
\begin{proof}
\emph{Classes.} For $z=0$ the unshifted entries are integers ($k$ even, $t$ odd), and $N^0_{a,b}-1$ is
even iff $a+b$ is even, which is when $s_p=0$. For $z=1$ every entry lies in $\frac12+\Z$ (the stem
$+\frac12$; blocks of even lengths $t\pm1$), and $N^1_{a,b}$ is even.

\emph{Distinctness.} The stem entries are strictly decreasing in $\varrho$; each block is a string of
consecutive numbers; the centres $\pm\kappa_{l+1}$ of the blocks are $\ge\Delta$ apart from each other
and $\ge\frac32\Delta-r$ from the small stem entries ($|\cdot|\le\frac k2$); and the large stem entries
$-\kappa-i\Delta$ are at distance $|i-l-\frac32|\Delta\ge\frac\Delta2$ from the centres $-\kappa_{l+1}$ and
$>\Delta$ from the centres $+\kappa_{l+1}$. All these exceed the half-lengths $\le\frac{t+1}2$ plus $1$.

\emph{Gap.} A negative coordinate is a stem coordinate with $a(\varrho)=i-1>b$, $i=\varrho+2-k$; so its
entry is $-\kappa-i\Delta$ (shifted) with $i\ge b+2$. Its distance to another stem entry is $\ge\Delta$
(large ones) or $\ge\kappa+\Delta$ (small ones); to the second kind blocks $l<b$ it is
$\ge(i-l-\frac32)\Delta-\frac t2\ge\frac32\Delta-\frac t2$; to the first kind blocks it is
$>2\kappa+\Delta$. So any two entries of different signs differ by at least $\Delta-t$.

\emph{Rank.} $N_p\le r+1+m(t+1)+m(t-1)\le2(k+m(t+1))=2n_m$, as $r+1=k+m$ (Lemma~\ref{pgs:boxes}).

\emph{Consequences.} As in the proof of Lemma~\ref{ptr:gap}, the four cited statements use the
parameter only through its regularity, the gap and the rank, and their numerical hypotheses hold if
$\frac12r^8\ge500(N+N')^3+10N^5+x^{25}$ for the ranks $N,N'\le2n_m$ involved. Now $n_m=k+m(t+1)\le
r+2\sqrt x(5r+1)\le13\sqrt x\,r\le13r^{11/10}$ (as $m\le2\sqrt x$ and $x^5\le r$), so the right side is
at most $500\cdot64\cdot13^3r^{33/10}+320\cdot13^5r^{11/2}+r^5\le2\cdot10^8r^{11/2}\le\frac12r^8$, since
$r\ge x^5$ and $x\ge128^{100}$ (\S\ref{pgs:ss-short}, proof of Lemma~\ref{pgs:boxes}).
\end{proof}

\begin{theorem}[(G3)]\label{pgv:G3}
Assume \textup{(V)} and \textup{(D1)--(D5)} ({\cite[\S3.1.3]{P49}}, \S\ref{ptq:ss-H}). For every edge
$(p,q)$ of the grids, with the parameters above and the characters $\alpha_{pq},\beta_{pq}$ of any
$m$-fold twist:
\begin{enumerate}
\item[(a)] a joint harmonic $P_{pq}=\bigotimes_\tau\Phi_\tau$ exists, is unique up to a scalar,
does not depend on the twist, and no monomial of it contains two variables of one column from rows
of opposite signs, or of one row from columns of opposite signs;
\item[(b)] $\mathsf PP_{pq}=\mathsf P'P_{pq}$, with $\mathsf P,\mathsf P'$ as in \S\ref{ptr:ss-edge};
\item[(c)] every form of type $e_p$, $\overline{e_p}$, $e_q$ or $\overline{e_q}$ is a cusp form, rapidly
decreasing on Siegel sets, and the forms of a given level form a finite-dimensional space;
\item[(d)] for every finite datum the lift and the transposed lift of Theorem~\ref{ptr:edge}(d)
converge absolutely; in particular the moves of Definition~\ref{pgv:def-edges} are defined, $\Lambda_X$
maps forms of type $e_p$ to forms of type $e_q$ and $T_X$ maps forms of type $e_q$ to forms of type
$e_p$;
\item[(e)] the statement \textup{(H)} of \S\ref{ptr:ss-mult} holds for $(p,q)$.
\end{enumerate}
This is (G3) of \S\ref{pgu:ss-data}.
\end{theorem}
\begin{proof}
Fix an edge $(p,q)$ and write $S=E^{\mathsf F_p}$, $T=E^{\mathsf F_q}$, of ranks $n=N_p<n'=N_q$,
$d=d_{pq}=n'-n$, $\alpha=\alpha_{pq}$, $\beta=\beta_{pq}$, and $(\lambda,\lambda')$ for the row and column
parameters at $\tau$ fixed before Proposition~\ref{pgv:match}; $e_S$, $e_T$ are the types of
\S\ref{ptr:ss-edge} ($e_S=e_p$, $e_T=e_q$ for a right edge, $e_S=\overline{e_p}$, $e_T=\overline{e_q}$
for an up edge: in both cases the conjugate of $e_S$ has parameter $\lambda$ and $e_T$ has parameter
$\lambda'$), and $\sigma=\omega_{\alpha,\beta,\infty}(\cdot,1)$, $\sigma'=\omega_{\alpha,\beta,\infty}(1,\cdot)$
at $\tau$. The argument is parallel to that of Theorem~\ref{ptr:edge}: both carry over the proofs of
{\cite[Theorem~3.15]{P49}} and of {\cite[Proposition~8.4]{P49}} for the pair $(Y_j,Y_{j+1})$. It is given here in full and does not use Theorem~\ref{ptr:edge}: we go through the
steps of those two proofs and say for each what it uses.

\emph{Step 0: the derived action.} {\cite[Theorem~3.29]{P49}} is stated for $W=\auHom(Y_j,Y_{j+1})$ at $\tau$,
but Steps 1--3 of its proof concern $W=\auHom(V,V')$ for arbitrary nondegenerate hermitian spaces
$V,V'$ over $E_\tau$: they use only (W1$_\infty$), (W2$_\infty$), (M) and (S), which are stated for every
nondegenerate hermitian space, the identity $\det_W\iota(g,h)=\det(h)^n\det(g)^{-n'}$, and that
$W\oplus W'$ is indefinite. So (V) holds for $(S,T)$. The computation of $d\omega^{\mathrm{Mp}}$ in the
proof of {\cite[Lemma~3.6]{P49}} is made for arbitrary signs $\varepsilon,\varepsilon'$ and ranks (it uses
only $\eta_{rc}\eta_{r'c}=s_{rr'}$ and its analogue for columns). Hence Lemma~\ref{ptr:const} holds for
$(S,T)$: the derived actions are {\cite[(12)]{P49}} with $\gamma=\tau_0+\frac{n'}2$,
$\gamma'=\tau'_0+\frac n2$, and $\tau_0,\tau'_0$ are the numbers of Proposition~\ref{pgv:match}. They
depend on $\alpha,\beta$ only through the weights, which do not depend on the twist
($\pgvwt\vartheta=0$).

\emph{Step 1: proof of} (a). By Proposition~\ref{pgv:match} the pair $(\lambda,\lambda')$ satisfies the
block rule of excess $d$, and by Lemma~\ref{pgv:gap} its entries are distinct and lie in
$\frac{n-1}2+\Z$, $\frac{n'-1}2+\Z$. By Lemma~\ref{pgv:rule} this is the matching condition
{\cite[(13)]{P49}} with $d$ in place of $t$, and {\cite[Lemma~3.10]{P49}} holds with $d$ in place of $t$:
in its proof, $\xi\in\frac{d-1}2+\Z$ and $|\xi_k|>\frac d2$ give $|\xi_k|\ge\frac{d+1}2$, at most
$\xi_r-\frac{d+1}2$ indices have $0<\xi_k<\xi_r$, $|\sum_{c\in\mathsf B}\varepsilon'_c|=d$, so
$2w_r\ge2\xi_r-1-(2\xi_r-d-1)-d=0$; and for $c\in\mathsf B$, $2\xi'_c+\sum_{c'\in\mathsf B}\pgvsgn(\xi'_{c'}-
\xi'_c)=0$ because $\mathsf B$ is a string of $d$ consecutive numbers centred at $0$. The statement and
the proof of {\cite[Proposition~3.11]{P49}} do not contain $t$; they use the matching and
{\cite[Lemma~3.10]{P49}} only, and give a joint harmonic $\Phi_\tau$ with the property of the monomials.
{\cite[Lemma~3.12]{P49}} does not involve the matching. In the proof of {\cite[Proposition~3.13]{P49}},
$t$ occurs in the dichotomy $-s\xi_r>\frac t2$ or $s\xi_r>\frac t2$ and in the third case, where
$s\lambda'_{c(r)}>s\lambda'_c$ gives $s\xi'_c<s\xi'_{c(r)}=-s\xi_r<-\frac t2$, so that $c\notin\mathsf B$;
both hold with $d$, since $|\xi'_c|\le\frac{d-1}2$ on $\mathsf B$. So $\Phi_\tau$ is unique up to a
scalar. The passage to $P_{pq}=\bigotimes_\tau\Phi_\tau$ is the first paragraph of the proof of
{\cite[Theorem~3.15]{P49}}, which uses only that the conditions of {\cite[Definition~3.9]{P49}} are
conditions on the separate factors. The operators {\cite[(12)]{P49}}, hence the joint harmonic, depend
only on the signs, on $\lambda,\lambda'$ and on $\gamma,\gamma'$; by Step~0 they do not depend on the
twist.

\emph{Step 2: proof of} (b). As in the proof of {\cite[Theorem~3.15]{P49}} both sides are tensor products over
$\tau$, and it suffices to prove $\mathsf P\Phi=\mathsf P'\Phi$ at one $\tau$, where
$\mathsf P=\int e_S(x)\sigma(x^{-1})\,dx=\sigma(e_\lambda)$ and $\mathsf P'=\sigma'(e_{\lambda'})$ by the
description of $e_S$, $e_T$ above and $e(x^{-1})=\overline{e(x)}$. By Lemma~\ref{pgv:gap} the
conclusion of {\cite[Theorem~3.2]{P49}} holds for $e_\lambda$ and $e_{\lambda'}$: they and their derivatives
are integrable, so $\mathsf P$, $\mathsf P'$ are bounded commuting operators. Let $u=\mathsf P\Phi$.
Since $\mathsf P$ is bounded and commutes with $\sigma'$, $u$ is smooth for $\sigma'$, has $T'$-weight
$\Lambda^K(\lambda')$ and is killed by the compact raising operators of the columns. The first identity
of {\cite[Lemma~3.14]{P49}} is an identity between the operators {\cite[(12)]{P49}} for arbitrary
$n,n',\gamma,\gamma'$; its second assertion uses the matching through
$\lambda'=\{-\lambda_\varrho+s_0\}\cup\{\xi_0+\frac{d-1}2-\beta:0\le\beta<d\}$, so that
$|\lambda'|^2=|\lambda|^2-2s_0\sum\lambda+ns_0^2+d\xi_0^2+\frac{d(d^2-1)}{12}$, and $n'=n+d$; the final
reduction is an identity of polynomials in $n$, $d$, $s_0$, $\xi_0$, valid for every $d$. So
$d\sigma'(\Omega')u=\mathsf P\,d\sigma(\Omega)\Phi+c\,\mathsf P\Phi=\kappa_{\lambda'}u$, the integration
by parts being justified by the integrability of the derivatives of $e_\lambda$. {\cite[Lemma~3.5]{P49}}
is stated for every group $\auU(p,q)$ (the compact case included, {\cite[Lemma~3.4]{P49}}), every regular
parameter and every unitary representation; applied to $\sigma'$, $\lambda'$ and $u$ it gives
$\mathsf P'u=u$. Symmetrically $\mathsf P\mathsf P'\Phi=\mathsf P'\Phi$, and
$\mathsf P\Phi=\mathsf P'\mathsf P\Phi=\mathsf P\mathsf P'\Phi=\mathsf P'\Phi$.

\emph{Step 3: proof of} (c). This is the proof of {\cite[Proposition~8.4(a)]{P49}} with $G_j$ replaced by
$\auU(S)$ or $\auU(T)$ and $e_j$ by the type concerned. It uses of the type only (A1) and (A2), that is,
the conclusions of {\cite[Theorem~3.2]{P49}} and {\cite[Corollary~3.3]{P49}}, which hold by Lemma~\ref{pgv:gap},
and the general theorems quoted there. If the node has no negative sign (the nodes with $b\ge m-1$ in
the first grid and $b=m$ in the second, Lemma~\ref{pgv:count}(b)) the archimedean group is compact, and
the argument is the one given there for $j=m$; if the group over $E^+$ has no proper parabolic subgroup,
the cusp condition is empty and the rest is unchanged.

\emph{Step 4: proof of} (d). This is the proof of {\cite[Proposition~8.4(b)]{P49}}, (c) with
$(G_j,G_{j+1},\Phi_j,\omega_{\mu_{j+1}},e_j,e_{j+1})$ replaced by $(\auU(S),\auU(T),P_{pq},
\omega_{\alpha,\beta},e_S,e_T)$. It uses: {\cite[Lemma~8.2]{P49}}, whose statement and proof involve the
pair of spaces only through $D=fnn'$ and the facts (W0), (W1), (W4); the inequality $2s=2D+2\le B=x^{25}$,
which holds here because $f\le x^{12}$ and $n,n'\le2n_m\le26r^{11/10}$ (proof of Lemma~\ref{pgv:gap}), so
$D\le676\,x^{12}r^{11/5}\le x^{24}$; reduction theory for $\auU(S)$, resp.\ $\auU(T)$; the bound of
(c) on Siegel sets; (A1) for $e_S$ and $e_T$ with the weight $\|x\|^B$ (Lemma~\ref{pgv:gap}); the
identity $\mathsf PP_{pq}=\mathsf P'P_{pq}$ of (b); and the finite-dimensionality in (c) for the
truncation argument. Each of these has been established for $(S,T)$, and the argument is otherwise
word for word the same. It gives: the lift maps forms of type $e_S$ to forms of type $e_T$, and the
transposed lift maps forms of type $\overline{e_T}$ to forms of type $\overline{e_S}$. For a right edge
this is $e_p\to e_q$ for the lift; for an up edge the transposed lift maps forms of type $e_q$ to forms
of type $e_p$. These are the moves of Definition~\ref{pgv:def-edges}.

\emph{Step 5: proof of} (e). This is the proof of Theorem~\ref{ptq:H} for $(S,T)$. Its ingredients are:
Lemma~\ref{ptq:P}(a)--(d), whose proofs use the integrability of $e_\lambda$, $e_{\lambda'}$ and of
their derivatives (Lemma~\ref{pgv:gap}), {\cite[Lemma~3.5]{P49}}, the facts established for $u=\mathsf
P\Phi$ in Step~2, (D4) and (D5); Lemmas~\ref{ptq:fischer} and~\ref{ptq:tensor}, which are identities
between the operators {\cite[(12)]{P49}}; and Proposition~\ref{ptq:howe}, in whose proof the excess
occurs once, in the step on the row Laplacians: there $\lambda_a>\lambda_b$ gives $\xi_a>\xi_b$, and as
all $|\xi_\varrho|>\frac d2$, either $\xi_a>\frac d2$ or $\xi_b<\xi_a<-\frac d2$; the rest of that step
uses {\cite[Lemma~3.10(i)]{P49}} and {\cite[Lemma~3.12(c)]{P49}}, and the conclusion uses
{\cite[Proposition~3.13]{P49}}, all available by Step~1. The passage from one $\tau$ to the product and to
$\pzF_s$ at the end of the proof of Theorem~\ref{ptq:H} uses {\cite[Lemma~8.2(a)]{P49}} and (A1), as in
Step~4.
\end{proof}

\subsubsection{The square law on a grid square: Proposition~\ref{pgr:square}}\label{pgv:ss-square}

Fix $z\in\{0,1\}$ and $0\le a,b<m$, and the square $V'=(z,a,b)$, $W_1=(z,a,b+1)$, $W_2=(z,a+1,b)$,
$V=(z,a+1,b+1)$ of \S\ref{pgr:ss-off}, with $e_1=(V',W_1)$, $e_2=(V',W_2)$, $f_1=(W_1,V)$,
$f_2=(W_2,V)$. For an $m$-fold twist $\vartheta$ put $n=b+1$ if $z=0$, $n=a+1$ if $z=1$, and
\[
\nu=\nu^z_{ab}(\vartheta)=\bigl(\mu^0_{a+1}\mu^0_{b+1}\vartheta^{\,n}\bigr)^{-1},\qquad
\chi=(\chi_{e_1}\chi_{e_2})^{-1},\qquad M_0=\tfrac12\bigl(N_{W_1}+N_{W_2}\bigr),\qquad\gamma=\nu^{-1}\upsilon^{M_0}.
\]
By \eqref{pgv:eq-rank}, $N_{W_2}-N_{W_1}=2z$: in the first grid $N_{W_1}=N_{W_2}=M_0$, and in the second
grid $N_{W_2}=N_{W_1}+2$, all ranks are even (Lemma~\ref{pgv:count}(a)) and $M_0=N_{W_1}+1$ is an integer.
\begin{lemma}[(E1) for the grid squares]\label{pgv:E1g}
Assume \textup{(V)} and \textup{(D1)--(D5)}, as in Theorem~\ref{pgv:G3}. For every such square and every
$m$-fold twist, the four pairs $(U_1,A)=(W_1,V)$, $(U_2,A)=(W_2,V)$, $(A',U_1)=(V',W_1)$,
$(A',U_2)=(V',W_2)$, with the characters of Definition~\ref{pgv:def-edges} and the types
$e_1=e_{W_1}$, $e_2=e_{W_2}$, $e=e_V$, $e'=e_{V'}$, satisfy \textup{(E1)} of
Definition~\ref{ptr:def-square}, with the polynomials $P_{f_1},P_{f_2},P_{e_1},P_{e_2}$. We refer to
this statement as \textup{(E1g)}.
\end{lemma}
\begin{proof}
We check the clauses of (E1) one by one. Recall (\S\ref{pgr:ss-paths}) that a step which raises the
first index is a right edge and one which raises the second index is an up edge, and that for a right
edge $(p,q)$ Theorem~\ref{pgv:G3} is applied with $e_S=e_p$, $e_T=e_q$, for an up edge with
$e_S=\overline{e_p}$, $e_T=\overline{e_q}$ (the convention fixed before Proposition~\ref{pgv:match}).

\emph{The four pairs are edges, in the required directions.} $V'=(z,a,b)$, $W_1=(z,a,b+1)$,
$W_2=(z,a+1,b)$, $V=(z,a+1,b+1)$. So $f_1=(W_1,V)$ and $e_2=(V',W_2)$ raise the first index (right
edges), and $f_2=(W_2,V)$ and $e_1=(V',W_1)$ raise the second (up edges). By Theorem~\ref{pgv:G3}(a)--(d)
each of them satisfies the conclusions (a)--(d) of Theorem~\ref{ptr:edge} with $(S,T)$ the pair and the
types just recalled. Hence: $(U_1,A)=f_1$ is an edge from $e_S=e_{W_1}=e_1$ to $e_T=e_V=e$;
$(U_2,A)=f_2$ is an edge from $\overline{e_{W_2}}=\overline{e_2}$ to $\overline{e_V}=\overline e$;
$(A',U_1)=e_1$ is an edge from $\overline{e_{V'}}=\overline{e'}$ to $\overline{e_{W_1}}=\overline{e_1}$; and
$(A',U_2)=e_2$ is an edge from $e_{V'}=e'$ to $e_{W_2}=e_2$. These are exactly the four directions
required in (E1). In each case Theorem~\ref{ptr:edge}(d) gives both the lift from $e_S$ to $e_T$ and the
transposed lift from $\overline{e_T}$ to $\overline{e_S}$, which is what Definition~\ref{ptr:def-square}
uses.

\emph{Independence of the member of $I$.} The members of $I$ are indexed by the $m$-fold twists. The
types $e_p$ are defined by the node parameters \eqref{pgv:eq-param}, which do not involve $\vartheta$;
the joint harmonics $P_{pq}$ do not depend on the twist (Theorem~\ref{pgv:G3}(a)); the characters
$\alpha_{pq},\beta_{pq}$ of Definition~\ref{pgv:def-edges} depend on $\vartheta$ only through the factor
$\vartheta^{c(l+1)}$, of weight $0$, so their weights (computed after Definition~\ref{pgv:def-edges}) are
fixed.

\emph{(H) for $(U_1,A)$ and $(U_2,A)$.} This is Theorem~\ref{pgv:G3}(e) for the edges $f_1$ and $f_2$,
read with the same parameters (and hence the same projections $\mathsf P$, $\mathsf P'$) as above.

Every clause of (E1) is thus a conclusion of Theorem~\ref{pgv:G3}, which assumes only (V) and
(D1)--(D5). No edge needs to be read with the groups exchanged.
\end{proof}

\begin{lemma}[Grid squares are edge squares]\label{pgv:esq}
Assume \textup{(V)} and \textup{(D1)--(D5)}. The square, with the family $I$ of the data $(\pcm(\vartheta),\nu^z_{ab}(\vartheta),
\chi,\text{the characters of Definition~\ref{pgv:def-edges}})$ over the $m$-fold twists $\vartheta$, is
an edge square of excess $t$ and weight $\kappa_{a+1}+\kappa_{b+1}$ (Definition~\ref{ptr:def-square}),
and $\chi$ is the tame character of $\nu$.
\end{lemma}
\begin{proof}
\emph{Ranks.} $N_V=N_{V'}+2t$ (both indices grow by one; $(t+z)+(t-z)=2t$), and
$N_{W_1}+N_{W_2}=N_V+N_{V'}=2M_0$, so $\dim A=\dim A'+2t$ and $N=2M_0$.
\emph{Signs.} By Lemma~\ref{pgv:count}(b), with $\nu_-(b')=\lfloor(m-b'+z)/2\rfloor$: $q=\nu_-(b+1)+
\nu_-(b)$ (second indices of $W_1$, $W_2$) and $q_A+q_{A'}=\nu_-(b+1)+\nu_-(b)$ (second indices of $V$,
$V'$). Moreover $q\le m+1$, so $2q+2\le2m+4\le2r\le N$; and $N\le4n_m\le52r^{11/10}<3\Delta\le
2(\kappa_{a+1}+\kappa_{b+1})$ (proof of Lemma~\ref{pgv:gap}).
\emph{The data.} $\pcm(\vartheta)\ne(1)$ is a modulus; $\mu^0_l$ is tamely ramified at the chosen places
and unramified elsewhere, $\vartheta$ tamely ramified at $Q_\vartheta$ and unramified elsewhere, so $\nu$
is a tame character for $\pcm(\vartheta)$, of weight $\kappa_{a+1}+\kappa_{b+1}$. In the grid $z=0$,
$e_1$ is an up edge of index $b$, vertical, so $\nu_{e_1}=\mu^0_{b+1}\vartheta^{b+1}$, and $e_2$ is a
right edge of index $a$, horizontal, $\nu_{e_2}=\mu^0_{a+1}$; in the grid $z=1$, $\nu_{e_1}=\mu^0_{b+1}$
and $\nu_{e_2}=\mu^0_{a+1}\vartheta^{a+1}$. In both cases $\nu_{e_1}\nu_{e_2}=\nu^{-1}$, so the tame
character of $\nu$ is $(\chi_{e_1}\chi_{e_2})^{-1}=\chi$. By Proposition~\ref{pgv:G2} (opposite edges)
also $\nu_{f_1}\nu_{f_2}=\nu^{-1}$.
\emph{(E1)} is Lemma~\ref{pgv:E1g}.
\emph{(E2).} In each of the three products of (E2) the two edges are one right and one up edge of
the same grid, hence of the two kinds $c,c'$ with $\{c,c'\}=\{0,1\}$, and their $\nu$'s multiply to
$\nu^{-1}$ (the pairs are $\{f_1,f_2\}$ at $A$, $\{e_1,e_2\}$ at $A'$, and $\{f_1,e_1\}$ at $U_1$,
$\{f_2,e_2\}$ at $U_2$, each consisting of an edge of index $a$ and of the corresponding edge of index
$b$). It remains to add the powers of $\upsilon$. At $A=V$ the target characters are $\beta_{f_1}$,
$\beta_{f_2}$, with powers computed from the ranks of the sources $W_1$, $W_2$: for $z=0$,
$N_{W_1}=N_{W_2}=M_0$ and $\lfloor\frac{M_0+1-c}2\rfloor+
\lfloor\frac{M_0+1-c'}2\rfloor=\lceil\frac{M_0}2\rceil+\lfloor\frac{M_0}2\rfloor=M_0$; for $z=1$ both ranks
are even and the sum is $\frac{N_{W_1}}2+\frac{N_{W_2}}2=M_0$. At $A'=V'$ the source
characters are $\alpha_{e_1},\alpha_{e_2}$, with powers computed from the targets $W_1$, $W_2$: again
$M_0$. At $U_l=W_l$ the characters are $\alpha_{f_l}$ (power from $N_V$) and $\beta_{e_l}$ (power from
$N_{V'}$), of the two kinds; as $N_V=N_{V'}+2t$ have the same parity and $N_V+N_{V'}=2M_0$, the sum of a
floor and a ceiling (or of two exact halves) of $\frac{N_V}2$, $\frac{N_{V'}}2$ is $M_0$. So each product
is $\nu^{-1}\upsilon^{M_0}=\gamma$.
\end{proof}

\begin{theorem}[Proposition~\ref{pgr:square}]\label{pgv:square}
Assume \textup{(V)} and \textup{(D1)--(D5)} (the inputs of Theorems~\ref{ptr:square} and~\ref{pgv:G3}).
For $z\in\{0,1\}$ and $0\le a,b<m$ there are $\kappa_{\square}\in\C$ and a function
$c^z_{ab}$ on the $m$-fold twists with values in $\C^\times$, of modulus $\rho'_{ab}$ if $z=0$ and
$\rho'_{ba}$ if $z=1$ up to norm powers, such that for every $\vartheta$, every form $f$ of type
$e_{W_1}$ and every linear functional $l$ on the forms of type $e_{W_2}$, the functions $L_{f,l}$,
$L'_{f,l}$ of \S\ref{pgr:ss-off} satisfy the square law for $\chi$ with constant
$\kappa_{\square}/c^z_{ab}(\vartheta)$.
\end{theorem}
\begin{proof}
By Lemma~\ref{pgv:esq} and Theorem~\ref{ptr:square} there are $\kappa_{\square}\in\C$ and $C>0$ depending only on
the archimedean data, and for each $\vartheta$ a number $c=c^z_{ab}(\vartheta)$ with
$|c|^2=C\,L_*(\pcm(\vartheta),\nu,t)\ne0$, such that for all forms $y_1$ of type $e_1=e_{W_1}$ and $y_2$
of type $\overline{e_{W_2}}$ the functions $L,L'$ of Definition~\ref{ptr:def-square} satisfy the square
law for $\chi$ with constant $\kappa_{\square}/c$.

\emph{From $y_2$ to $l$.} Let $\mathcal F_0$ be the space of forms of type $e_{W_2}$ fixed by
$J_{W_2}(\pcm)$. For $X$ a matrix over $R$ and $k'\in J_T(\pcm)$, Lemma~\ref{pc:level} gives
$\omega_{\alpha,\beta}(1,k')[P,\delta_X]=[P,\delta_X]$ (the reduction of $k'$ is $1$), so every lift
with a point datum modulo $\pcm$ is fixed by $J(\pcm)$ of its target, and likewise every transposed lift
(with $k\in J_S(\pcm)$). Hence $T_{x_{\mathsf N_2}}\Lambda_{x_{\mathsf N_1}}f$ and
$\Lambda_{y_{\mathsf N_2}}T_{y_{\mathsf N_1}}f$ lie in $\mathcal F_0$ for all $x,y$. By
Theorem~\ref{pgv:G3}(c), $\mathcal F_0$ is finite-dimensional, and the Petersson product is positive
definite on it; so there is $w\in\mathcal F_0$ with $l(u)=\langle w,u\rangle=\int_{[\auU(W_2)]}
\aubar w\,u$ for $u\in\mathcal F_0$, and $y_2=\aubar w$ is a form of type $\overline{e_{W_2}}$. Since only
the restriction of $l$ to $\mathcal F_0$ enters $L_{f,l},L'_{f,l}$, we may replace $l$ by
$u\mapsto\int y_2u$.

\emph{Identification.} With $y_1=f$: $L'_{f,l}(y)=\int_{[\auU(W_2)]}y_2\cdot\Lambda_{y_{\mathsf N_2}}
T_{y_{\mathsf N_1}}f$ is $L'(\eta)$ for $\eta=y$ (the lift along $(A',U_2)=e_2$ after the transposed lift
along $(A',U_1)=e_1$). And $L_{f,l}(x)=\int_{[\auU(W_2)]}y_2\cdot T_{x_{\mathsf N_2}}F$ with
$F=\Lambda_{x_{\mathsf N_1}}f$ on $V$; the transposed lift along $f_2$ and the lift of $y_2$ along $f_2$
have the same kernel $\theta_{f_2}(g,h)$, and the double integral of
$y_2(g)\theta_{f_2}(g,h)F(h)$ over $[\auU(W_2)]\times[\auU(V)]$ converges absolutely
(Proposition~\ref{ptr:conv}, whose proof uses only (E1), Lemma~\ref{pgv:E1g}), so by Fubini $L_{f,l}(x)=
\int_{[\auU(V)]}\Lambda_{x_{\mathsf N_1}}f\cdot\Lambda_{x_{\mathsf N_2}}y_2=L(\xi)$ for $\xi=x$ (columns
$\mathsf N_1$ for $U_1$, $\mathsf N_2$ for $U_2$). So $L_{f,l},L'_{f,l}$ satisfy the square law for $\chi$
with constant $\kappa_{\square}/c$.

\emph{Modulus.} By definition (after \eqref{pc:eq-rho}) $L_*(\pcm,\nu,t)=(\prod_{w\mid\pcm}q_w)^{t^2}\prod_{z'=1}^t|L^\pcm(z',\nu)|^2$, with
$\prod_{w\mid\pcm(\vartheta)}q_w=\auNm(\pi_M)\,\auNm Q_\vartheta=M^f\auNm Q_\vartheta$ (the ideal $(\pi_M)$ is prime, and only the primes dividing $\pcm(\vartheta)$ count, not their multiplicity $b$). The primes of $\pcm(\vartheta)$ are
those of $\pcm_\vartheta$, so $L^{\pcm(\vartheta)}$ is the $L$-function of \eqref{pc:eq-rhoprime} with the
convention recalled after \eqref{pc:eq-rho}. For $z=0$, $\nu=(\mu^0_{a+1}\mu^0_{b+1}\vartheta^{b+1})^{-1}$ and
\[
|c^0_{ab}(\vartheta)|^2=C\,M^{ft^2}\prod_{z'=1}^t|L(z',(\mu^0_{a+1}\mu^0_{b+1})^{-1})|^2\cdot
(\auNm Q_\vartheta)^{t^2}\,\rho'_{ab}(\vartheta)^2 ;
\]
the product of the $L$-values does not depend on $\vartheta$ and is not zero, because $c^0_{ab}(1)\ne0$.
So $c^0_{ab}$ has modulus $\rho'_{ab}$ up to norm powers (Definition~\ref{pc:modulus}, with exponent $t^2$). For
$z=1$, $\nu=(\mu^0_{b+1}\mu^0_{a+1}\vartheta^{a+1})^{-1}$, and the same computation gives the modulus
$\rho'_{ba}$.
\end{proof}

\begin{remark}[(E1g) and Theorem~\ref{pgv:G3}]\label{pgv:rem-E1}
Lemma~\ref{pgv:E1g} shows that Theorem~\ref{pgv:G3} supplies every clause of (E1) for the grid squares:
$(W_1,V)$ is a right edge, lifting $e_{W_1}$ to $e_V$; $(W_2,V)$ is an up edge, lifting
$\overline{e_{W_2}}$ to $\overline{e_V}$; $(V',W_1)$ is an up edge, lifting $\overline{e_{V'}}$ to
$\overline{e_{W_1}}$; $(V',W_2)$ is a right edge, lifting $e_{V'}$ to $e_{W_2}$. The reduction of (E1) to
the block rule recorded in \S\ref{ptw:ss-grid} is thereby completed: the block rule is
Proposition~\ref{pgv:match}.
\end{remark}




% ===== C-grid5.tex =====
% proofs/C-grid5.tex
% Continuation of proofs/C-grid.tex (pgr:), C-grid2.tex (pgs:), C-grid3.tex (pgu:), C-grid4.tex
% (pgv:): the first path (Proposition pgr:first) and the last path (Proposition pgr:last).
% Labels start with pgw:. Conventions and macros: those of proofs/C-grid.tex and C-determinant.tex.
\providecommand{\auI}{\mathcal I}
\providecommand{\auGL}{\mathrm{GL}}
\providecommand{\auU}{\mathrm U}
\providecommand{\auO}{\mathcal O}
\providecommand{\auA}{\mathbb A}
\providecommand{\auNm}{\operatorname{Nm}}
\providecommand{\auHom}{\operatorname{Hom}}
\providecommand{\aubar}[1]{\overline{#1}}
\providecommand{\pcm}{\mathfrak m}
\providecommand{\pcb}{\mathfrak b}
\providecommand{\pcMat}{\mathrm M}
\providecommand{\pcFF}{\mathfrak F}
\providecommand{\pcF}{\mathcal F}
\providecommand{\auPl}{\Phi}
\providecommand{\psvac}{\mathbf v}
\providecommand{\pgrY}{\widetilde Y}
\providecommand{\pgrG}{\Gamma}
\providecommand{\pgrfirst}{\Gamma_{\mathrm{first}}}
\providecommand{\pgrlast}{\Gamma_{\mathrm{last}}}
\providecommand{\pgrE}{\operatorname{Ed}}
\providecommand{\pgupt}{\mathsf u}
\providecommand{\pgwK}{\mathcal K}
\providecommand{\pgwT}{\mathcal T}
\providecommand{\pgwP}{\mathcal P}
\providecommand{\pgwev}{\operatorname{ev}}
\providecommand{\pgwint}[1]{\int^{#1}}

\subsection{The first and the last path}\label{pgw:sec}

This section continues \S\S\ref{pgr:sec}, \ref{pgs:sec}, \ref{pgu:sec}, \ref{pgv:sec} and uses their
notation; the grids, their characters, moves, levels and the diagonal period are those of
\S\ref{pgv:ss-def}. \S\ref{pgw:ss-first} reduces Proposition~\ref{pgr:first} (Theorem~\ref{pgw:first})
to three precisely stated inputs: the injectivity of the untwisted chain of the second grid, a
deep-level vanishing theorem, and a compatibility of data under change of modulus.
\S\ref{pgw:ss-last} proves Proposition~\ref{pgr:last} in full, in the form of
Theorem~\ref{pgw:last}, which makes the normalisation of $\Lambda$ explicit.

\subsubsection{The first path}\label{pgw:ss-first}

\paragraph*{The columns}
The first path $\pgrfirst$ runs down the first columns of the two grids: the nodes $(0,0,b)=Y_b$
and $(1,0,b)=\pgrY^b_0$, $0\le b\le m$, joined by the up edges $e^z_b=((z,0,b),(z,0,b+1))$, $b<m$.
In the first grid these edges are vertical, with the characters of $a_{0,b}$
(\S\ref{pgv:ss-def}); in the second grid they are horizontal, so their characters do not involve
$\vartheta$, and their tame parts are supported at the chosen places. The moves along them are
transposed lifts $T_X$ (Definition~\ref{pgv:def-edges}). For a node family $g$ and a set $w$ of
edges we use $\chi_w(g)$ and the action $(gS)_e=g_q^{-1}S_eg_p$ of \S\ref{pgr:ss-arch}. Let
$\pcm_0=\pcm_\vartheta=(\pi_M)\prod_{w\in Q_\vartheta}w$; it has the same primes as $\pcm=\pcb
\prod_{w\in Q_\vartheta}w$. For $z\in\{0,1\}$ let $\pcm^z$ be $\pcm_0$ if $z=0$ and $\pcb$ if
$z=1$, and let $\iota$ be the inclusion matrix of frames.

\emph{The chains.} By uniqueness of the joint harmonic (Theorem~\ref{ptr:edge}(a), (G3)), the joint
harmonic $P_{e^0_b}$ is a non-zero multiple $c_bP_{0,b}$ of the harmonic of $a_{0,b}$, with $c_b$
independent of $\vartheta$; replacing $P_{e^0_b}$ by $P_{0,b}$ changes the moves by these constants,
and we absorb them below into the constant $\alpha_{\mathrm{first}}$ (they do not depend on $\vartheta$, $S$ or the
forms). Let $a^{\mathsf t}_{0,b}$ be the transposed lift with the datum of $a_{0,b}$ (modulus $\pcm_0$,
point datum $\delta_\iota$), and
\[
X^{\vartheta,\mathsf t}_{0,m}=w_0a^{\mathsf t}_{0,0}\circ w_1a^{\mathsf t}_{0,1}\circ\dots\circ
w_{m-1}a^{\mathsf t}_{0,m-1}
\]
(conjugate forms on $Y_m$ to conjugate forms on $Y_0$). In the second grid let $T^\flat_b$ be the
transposed lift along $e^1_b$ with the modulus $\pcb$ and the point datum $\delta_\iota$ on
$\pcMat(\auO_E/\pcb)$, and $\Theta^\flat=T^\flat_0\circ\dots\circ T^\flat_{m-1}$; it does not depend on
$\vartheta$. Let $a^\flat_b$ be the lift with the same datum, so that $a^\flat=a^\flat_{m-1}\circ\dots
\circ a^\flat_0$ maps forms of type $\overline{e_{(1,0,0)}}$ on $\pgrY^0_0$ to forms of type
$\overline{e_{(1,0,m)}}$ on $\pgrY^m_0$ (Theorem~\ref{ptr:edge}(d) for the pairs of the grids, (G3)).

\emph{Levels.} By Lemma~\ref{pc:level}: a pair $(k,k')\in K_p\times K_q$ with $k\equiv1$ and
$k'\equiv1$ modulo $\pcm^z$ fixes the data $[P,\delta_\iota]$ of the edge $e^z_b$ for the modulus
$\pcm^z$ (the tame parts of the characters are trivial on such elements, and $\bar k'^{-1}\iota\bar
k=\iota$). Hence $X^{\vartheta,\mathsf t}_{0,m}y'$ is fixed by $J_{Y_0}(\pcm_0)$, $\Theta^\flat\varphi$
by $J_{\pgrY^0_0}(\pcb)$, and $a^\flat y$ by $J_{\pgrY^m_0}(\pcb)$, for all forms $y',\varphi,y$.

\emph{Transposition.} For an edge $(p,q)$ with datum $F$, a form $w$ at $p$ and $x$ at $q$ of the
conjugate types,
\begin{equation}\label{pgw:eq-transp}
\int_{[G_q]}x(h)\,\bigl(\Theta(F)w\bigr)(h)\,dh=\int_{[G_p]}\bigl(\Theta^{\mathsf t}(F)x\bigr)(g)\,w(g)\,dg,
\end{equation}
both sides being $\iint\theta(\omega(g,h)F)w(g)x(h)\,dg\,dh$ by Fubini: theta functions of Schwartz
data have moderate growth, and the forms are rapidly decreasing (Theorem~\ref{ptr:edge}(c), (G3)),
so the double integral converges absolutely.

\paragraph*{Three inputs}
\begin{enumerate}
\item[(I)] (Injectivity of the second column.) $a^\flat$ is injective on the forms of type
$\overline{e_{(1,0,0)}}$ on $\pgrY^0_0$ fixed by $J_{\pgrY^0_0}(\pcb)$.
\item[(V)] (Deep-level vanishing.) There is a real point $g_0$ of $\auU(\pgrY^0_0)$ such that every
$y\in\auAz$ with $\int_{[\auU(Y_0)]}\aubar{y(x)}\,\psi((x\oplus1)g_0)\,dx=0$ for all forms $\psi$ of
type $e_{(1,0,0)}$ on $\pgrY^0_0$ fixed by $J_{\pgrY^0_0}(\pcb)$ is zero.
\item[(P)] (Change of modulus.) For an edge of the second column and $\Psi$ on
$\pcMat(\auO_E/\pcb)$, the data $[P,\Psi]$ for the modulus $\pcb$ and $[P,\Psi\circ\mathrm{red}]$ for
the modulus $\pcm$ coincide, $\mathrm{red}$ being reduction modulo $\pcb$.
\end{enumerate}
\noindent (I) is Theorem~\ref{pgy:I} and (P) is Proposition~\ref{pgy:P} (for the data of
Definition~\ref{pgx:conv}, without any assumption on $\alpha$). (V) is Theorem~\ref{pgq:V}, with $g_0=\omega$
the point of input (S) there; its inputs (C), (Q), (S) are Theorems~\ref{pgk:C}, \ref{pgk:Q}
and~\ref{pgl:S}.

\paragraph*{The restriction map}
Let $\pcFF$ be the forms of type $e_{(1,0,m)}$ on $\pgrY=\pgrY^m_0$ fixed by $J(\pcb)$, and $\pcFF_0$
those of type $e_{(1,0,0)}$ on $\pgrY^0_0$ fixed by $J(\pcb)$; both are finite-dimensional
(Theorem~\ref{ptr:edge}(c)). For a form $y$ on $Y_0$ of type $e_0$ and $\psi\in\pcFF_0$ the diagonal
period at $l=0$ is $\mathsf P_{0,g}(\aubar y,\psi)=\int_{[\auU(Y_0)]}\aubar{y(x)}\psi((x\oplus1)g)\,dx$
(Definition~\ref{pgv:def-period}, $\pgupt^0=1$).

\begin{lemma}\label{pgw:onto}
 $\Theta^\flat(\pcFF)=\pcFF_0$.
\end{lemma}
\begin{proof}
$\Theta^\flat(\pcFF)\subseteq\pcFF_0$: the levels give the invariance under $J(\pcb)$, and the type of a transposed lift is given by Theorem~\ref{ptr:edge}(d). Let $\psi\in\pcFF_0$ be orthogonal to
$\Theta^\flat(\pcFF)$ for the Petersson product. Put $U=a^\flat(\aubar\psi)$; it is fixed by
$J(\pcb)$, so $\aubar U\in\pcFF$. By \eqref{pgw:eq-transp} applied $m$ times,
$\langle U,U\rangle=\int\aubar U\,U=\int\Theta^\flat(\aubar U)\cdot\aubar\psi=\aubar{\langle\Theta^\flat
(\aubar U),\psi\rangle}=0$, with $\langle\varphi',\varphi\rangle=\int\aubar{\varphi'}\varphi$. So $U=0$ and
$\psi=0$ by (I) (Theorem~\ref{pgy:I}). As $\pcFF_0$ is finite-dimensional, $\Theta^\flat(\pcFF)=\pcFF_0$.
\end{proof}

\begin{lemma}\label{pgw:res}
Let $g$ be a real point of $\auU(\pgrY^0_0)$. There is a linear map $\mathsf r_g\colon\pcFF\to\auAz$
such that for every $m$-fold twist $\vartheta$, every $\varphi\in\pcFF$ and every form $y$ of type
$e_0$ on $Y_0$ fixed by $J_{Y_0}(\pcm_\vartheta)$,
$\langle y,\mathsf r_g\varphi\rangle=\mathsf P_{0,g}(\aubar y,\Theta^\flat\varphi)$. For the point
$g=g_0$ of \textup{(V)}, $\mathsf r_{g_0}$ is onto.
\end{lemma}
\begin{proof}
$\auAz=\pcF(Y_0,e_0)^{J_M}$ is finite-dimensional, so there is a unique $\mathsf r_g\varphi\in\auAz$
with $\langle y,\mathsf r_g\varphi\rangle=\mathsf P_{0,g}(\aubar y,\Theta^\flat\varphi)$ for $y\in\auAz$
(the right side is conjugate linear in $y$). Let now $y$ be fixed by $J(\pcm_\vartheta)$ and $u\in
J_M$. Let $v\in J_M$ be $u_w^{-1}$ at the chosen places $w$ and $1$ elsewhere. Then $v$ is trivial at
the primes of $Q_\vartheta$ and at the chosen places lies in $J_M$, so $v\in J(\pcm_\vartheta)$ fixes
$y$; and $uv$ is $1$ at the chosen places, so $uv\in J(\pcb)$ and $(uv)\oplus1\in J_{\pgrY^0_0}(\pcb)$
fixes $\Theta^\flat\varphi$. Since $g$ is real it commutes with finite points, and the substitution
$x\mapsto x(uv)^{-1}$ gives
\[
\mathsf P_{0,g}(\overline{\rho(u)y},\Theta^\flat\varphi)=\mathsf P_{0,g}(\overline{\rho(uv)y},
\rho((uv)\oplus1)\Theta^\flat\varphi)=\mathsf P_{0,g}(\aubar y,\Theta^\flat\varphi);
\]
also $\langle\rho(u)y,\mathsf r_g\varphi\rangle=\langle y,\mathsf r_g\varphi\rangle$ since
$\mathsf r_g\varphi$ is fixed by $J_M$. Both sides of the claimed identity are therefore unchanged when
$y$ is replaced by its mean $\hat y=\int_{J_M}\rho(u)y\,du\in\auAz$ (volume $1$), and for $\hat y$ the
identity is the definition. If $y\in\auAz$ is orthogonal to the image of
$\mathsf r_{g_0}$, then $\mathsf P_{0,g_0}(\aubar y,\Theta^\flat\varphi)=0$ for all $\varphi\in\pcFF$,
hence $\mathsf P_{0,g_0}(\aubar y,\psi)=0$ for all $\psi\in\pcFF_0$ by Lemma~\ref{pgw:onto}, and $y=0$
by (V) (Theorem~\ref{pgq:V}).
\end{proof}

\begin{lemma}\label{pgw:chain}
For every real point $g$, every $\vartheta$, every conjugate form $y'$ on $Y_m$ and every
$\varphi\in\pcFF$,
$\int_{[\auU(Y_m)]}y'\cdot X^\vartheta_{0,m}(\mathsf r_g\varphi)=\mathsf P_{0,g}\bigl(X^{\vartheta,
\mathsf t}_{0,m}y',\Theta^\flat\varphi\bigr)$.
\end{lemma}
\begin{proof}
By \eqref{pgw:eq-transp} applied to the $m$ edges, the left side is $\int_{[\auU(Y_0)]}(X^{\vartheta,
\mathsf t}_{0,m}y')\cdot\mathsf r_g\varphi$. The form $y=\overline{X^{\vartheta,\mathsf t}_{0,m}y'}$ is
of type $e_0$ and fixed by $J(\pcm_\vartheta)$, so this is $\langle y,\mathsf r_g\varphi\rangle=
\mathsf P_{0,g}(\aubar y,\Theta^\flat\varphi)$ by Lemma~\ref{pgw:res}.
\end{proof}

\paragraph*{The path datum}
For an edge tuple $S$ let
\[
\Psi^1_\vartheta(S)=\prod_{z\in\{0,1\}}\prod_{b<m}\mathbf 1\bigl[S_{e^z_b}\equiv\iota\bmod\pcm^z\bigr]
\in\{0,1\}.
\]
Let $\mathcal C$ be the set of families $x=(x_{z,b})_{z\in\{0,1\},\,b<m}$ of left cosets
$x_{z,b}\in\auGL_{\mathsf F_{(z,0,b)}}(R)/H_{z,b}$, $H_{z,b}$ the level of the node $(z,0,b)$
(\S\ref{pgv:ss-def}: $H_{0,b}=1$, $H_{1,b}=\ker(\to\auGL(\auO_E/\pcb))$). Fix representatives
$\gamma_{z,b}\in x_{z,b}$ and let $g_x$ be the node family equal to $\gamma_{z,b}$ at $(z,0,b)$ for
$b<m$, except $g_x=\gamma_{0,0}\oplus1$ at $(1,0,0)$, and to $1$ at all other nodes. Put
\[
\Psi^0_\vartheta(S)=\sum_{x\in\mathcal C}\chi_{\pgrE(\pgrfirst)}(g_x)\,\Psi^1_\vartheta(g_xS).
\]

\begin{lemma}\label{pgw:datum}
\begin{enumerate}
\item[(a)] If $h$ is a node family with $h_{(z,0,b)}\in H_{z,b}$ for all $b\le m$ \textup{(}with
$H_{z,m}$ the levels of the outer nodes\textup{)}, then $\chi_{\pgrE(\pgrfirst)}(h)=1$ and
$\Psi^1_\vartheta(hS)=\Psi^1_\vartheta(S)$.
\item[(b)] $\Psi^0_\vartheta$ is a path datum on $\pgrE(\pgrfirst)$ with bound $\#\mathcal C\le
\mathcal E(\vartheta)^{2m}$.
\end{enumerate}
\end{lemma}
\begin{proof}
(a) At the nodes of the first column $h=1$. At $(1,0,b)$, $h\equiv1\bmod\pcb$, and the characters of
the second column are trivial on such determinants (their tame parts live at the chosen places).
Moreover $h_q^{-1}S_eh_p\equiv S_e$ modulo $\pcm^z$.

(b) $\Psi^0_\vartheta(S)$ depends only on the $S_e$, $e\in\pgrE(\pgrfirst)$, and is a sum of
$\#\mathcal C$ terms each $0$ or a root of unity; $\#\mathcal C\le\prod_{z,b}\#\auGL_{N_{(z,0,b)}}(R)
\le\mathcal E^{2m}$. \emph{Left inverses.} The vertical edges of $\pgrfirst$ are the $e^0_b$. If
$\Psi^0_\vartheta(S)\ne0$, some $\Psi^1_\vartheta(g_xS)\ne0$, so $Y=(g_xS)_{e^0_b}\equiv\iota\bmod
\pcm_0$; then $\iota^{\mathrm t}Y\equiv1$ modulo $\pcm_0$, hence is invertible modulo $\pcm$ (same
primes), $Y$ has a left inverse, and so has $S_{e^0_b}=g_{x,q}Yg_{x,p}^{-1}$. \emph{Invariance.} Let
$g$ be admissible. Left multiplication by $g_{(z,0,b)}$ permutes the cosets; let $\sigma(x)$ be the
image of $x$. Then $(g_xgS)$ is $((gg_x)S)$ with $(gg_x)_p=g_pg_{x,p}$, and $gg_x=g_{\sigma x}h$ with
$h_p=g_{\sigma x,p}^{-1}g_pg_{x,p}\in H_p$ at the nodes $(z,0,b)$, $b<m$; at the outer nodes
$h_p=g_p\in H_p$ (admissibility); and at $(1,0,0)$, as $H_{0,0}=1$ forces $g_{\sigma x,(0,0,0)}=
g_{(0,0,0)}g_{x,(0,0,0)}$, admissibility ($g_{(1,0,0)}=g_{(0,0,0)}\oplus1$) gives $h_{(1,0,0)}=1$. Since
$((g_{\sigma x}h)S)=h(g_{\sigma x}S)$ and $\chi$ is multiplicative, (a) gives
$\chi(g)\chi(g_x)\Psi^1((gg_x)S)=\chi(g_{\sigma x})\Psi^1(g_{\sigma x}S)$. Summing over $x$:
$\chi_{\pgrE(\pgrfirst)}(g)\Psi^0_\vartheta(gS)=\Psi^0_\vartheta(S)$.
\end{proof}

\begin{lemma}\label{pgw:transl}
Let $g$ be a node family with $g=1$ at $(0,0,m)$ and $(1,0,m)$ and $g_{(1,0,0)}=g_{(0,0,0)}\oplus1$.
For every real point $g_\infty$, every conjugate form $y'$ on $Y_m$, every form $\varphi$ on
$\pgrY^m_0$ and every $S$, with $h_z=$ the column of the grid $z$,
\[
\mathsf P_{0,g_\infty}\bigl(h_{0*}(gS)y',h_{1*}(gS)\varphi\bigr)=\chi_{\pgrE(\pgrfirst)}(g)\,
\mathsf P_{0,g_\infty}\bigl(h_{0*}(S)y',h_{1*}(S)\varphi\bigr).
\]
\end{lemma}
\begin{proof}
Choose $k_p\in K_p$ reducing to $g_p$ at each node $(z,0,b)$, with $k=1$ at the outer nodes and
$k_{(1,0,0)}=k_{(0,0,0)}\oplus1$. For an edge $e=(p,q)$ of a column, $X=S_e$ and $Y=(gS)_e=
g_q^{-1}Xg_p$, Lemma~\ref{pc:level} gives $\omega(k_p,k_q)[P,\delta_Y]=\chi_e(\det g_q/\det g_p)
[P,\delta_X]$ (as $\bar k_qY\bar k_p^{-1}=X$, and both characters of the edge act through
$\chi_e$ on determinants of elements of $K$), so $\theta(gk_p,hk_q;P,\delta_Y)=\chi_e(\det g_q/\det
g_p)\theta(g,h;P,\delta_X)$ and, substituting $h\mapsto hk_q$ in the transposed lift,
$\rho(k_p)T_Y(F)=\chi_e(\det g_q/\det g_p)\,T_X(\rho(k_q)F)$. By induction down the column of the grid
$z$, $\rho(k_{(z,0,0)})h_{z*}(gS)F=\chi_{w_z}(g)h_{z*}(S)F$, $w_z$ the edges of the column. Finally
$\mathsf P_{0,g_\infty}(\rho(k)y'',\rho(k\oplus1)\varphi'')=\mathsf P_{0,g_\infty}(y'',\varphi'')$ for
finite $k$ (substitute $x\mapsto xk^{-1}$; $g_\infty$ commutes with $k\oplus1$).
\end{proof}

\begin{theorem}[Proposition~\ref{pgr:first}]\label{pgw:first}
With $g_0$ as in \textup{(V)}, $\mathsf r=
\mathsf r_{g_0}$ and $\Psi^0=(\Psi^0_\vartheta)$, there is $\alpha_{\mathrm{first}}\in\C$ such that for every period
triple $(\vartheta,y',\varphi)$ with $\varphi\in\pcFF$
\[
\int_{[\auU(Y_m)]}y'\cdot X^\vartheta_{0,m}(\mathsf r\varphi)=\alpha_{\mathrm{first}}\,I_{\pgrfirst}(\Psi^0;g_0)
(\vartheta,y',\varphi);
\]
$\Psi^0_\vartheta$ is a path datum on $\pgrE(\pgrfirst)$ with bound $\mathcal E(\vartheta)^{2m}$, and
$\mathsf r$ is onto $\auAz$.
\end{theorem}
\begin{proof}
The inputs (I), (V), (P) hold: Theorem~\ref{pgy:I}, Theorem~\ref{pgq:V} (with $g_0=\omega$) and
Proposition~\ref{pgy:P}. $\mathsf r$ is onto by Lemma~\ref{pgw:res}, and $\Psi^0_\vartheta$ is a path datum by
Lemma~\ref{pgw:datum}. By Lemma~\ref{pgw:chain} the left side is $\mathsf P_{0,g_0}(X^{\vartheta,
\mathsf t}_{0,m}y',\Theta^\flat\varphi)$. \emph{Expansion.} For the edge $e^0_b$ the datum
$\delta_\iota$ modulo $\pcm_0$ and the function $X\mapsto\mathbf 1[X\equiv\iota\bmod\pcm_0]$ modulo
$\pcm$ give the same function $\Psi^\natural$ (Proposition~\ref{pgy:P} with $\pcm'=\pcm\subseteq\pcm_0$, as $\pcb\subseteq(\pi_M)$, the generators $a_w$ being the same), so
$a^{\mathsf t}_{0,b}=\sum_{X\equiv\iota\,(\pcm_0)}T_X$ up to the constant $c_b$; for $e^1_b$, (P) gives
$T^\flat_b=\sum_{X\equiv\iota\,(\pcb)}T_X$, the sums over $X\in\pcMat(R)$. Expanding the composites
(edge tuples on $\pgrE(\pgrfirst)$ are the pairs of tuples on the two columns),
\[
\mathsf P_{0,g_0}(X^{\vartheta,\mathsf t}_{0,m}y',\Theta^\flat\varphi)=W(\vartheta)\sum_{S\in\mathcal S_\vartheta
(\pgrfirst)}\Psi^1_\vartheta(S)\,\mathsf P_{0,g_0}\bigl(h_{0*}(S)y',h_{1*}(S)\varphi\bigr),\qquad
W(\vartheta)=\prod_{b<m}\frac{w_b^\vartheta}{c_b}.
\]
\emph{Symmetrisation.} Each $g_x$ satisfies the hypotheses of Lemma~\ref{pgw:transl}, and $S\mapsto
g_xS$ is a bijection of $\mathcal S_\vartheta(\pgrfirst)$; so
$\sum_S\chi(g_x)\Psi^1(g_xS)\mathsf P(S)=\sum_S\Psi^1(g_xS)\mathsf P(g_xS)=\sum_S\Psi^1(S)\mathsf P(S)$,
writing $\mathsf P(S)$ for the period in the display. Summing over $x$,
$I_{\pgrfirst}(\Psi^0;g_0)(\vartheta,y',\varphi)=\#\mathcal C\cdot\sum_S\Psi^1_\vartheta(S)\mathsf P(S)$.
\emph{Constants.} $\#\mathcal C=\prod_{b<m}\#\auGL_{N_{(0,0,b)}}(R)\,\#\auGL_{N_{(1,0,b)}}(\auO_E/\pcb)$
and $w^\vartheta_b=\#\auGL_{N_{(0,0,b)}}(\auO_E/\pcm_0)$. As in the proof of Theorem~\ref{pgw:last},
the Chinese remainder theorem gives $\#\auGL_N(R)=\kappa_N\#\auGL_N(\auO_E/\pcm_0)$ with $\kappa_N=
\#\auGL_N(\auO_E/\pcb)/\#\auGL_N(\auO_E/(\pi_M))$ independent of $\vartheta$. Hence
$W(\vartheta)/\#\mathcal C=\alpha_{\mathrm{first}}$ does not depend on $\vartheta$, and the theorem follows.
\end{proof}

\subsubsection{The last path}\label{pgw:ss-last}

Throughout this subsection $\vartheta$ is an $m$-fold twist, $\pcm=\pcm(\vartheta)$, and $p,q$ denote
nodes $(z,a,m)$ of the row $b=m$ of either grid. By \eqref{pgv:eq-sign}, $\epsilon_p(\varrho)=+1$
whenever $a(\varrho)\le b$; for $b=m$ this holds for every stem coordinate, since
$a(\varrho)\le m-1+z\le m$. Hence \emph{every node of the row $m$ is a positive definite space},
and so is every space on the last path $\pgrlast$, whose halves run through the nodes
$(0,a,m)=Y^a_m$ and $(1,a,m)=\pgrY^m_a$, $0\le a\le m$, along right edges.

For a definite node $p$ write $G_p=\auU(p)$, $G_{p,\infty}$ for its group of real points (a compact
group, of total mass $1$), $G_{p,f}$ for its finite adelic points, $\Gamma_p=G_p(E^+)$, and $K_p$
for the stabiliser of the lattice. Since $G_{p,\infty}$ is compact, $\Gamma_p$ projects injectively
onto a discrete subgroup of $G_{p,f}$ (an element with trivial finite components is $1$), which we
denote again by $\Gamma_p$; the stabilisers $\Gamma_p\cap xJx^{-1}$ of open compact subgroups $J$ are
finite. For an open compact subgroup $J\subseteq G_{p,f}$ and a function $F$ on
$\Gamma_p\backslash G_{p,f}$ that is right invariant under an open subgroup $J'\subseteq J$ put
\[
\pgwint{J}_{\Gamma_p\backslash G_{p,f}}F(x)\,dx=\sum_{x\in\Gamma_p\backslash G_{p,f}/J'}\frac{F(x)}{[J:J']\,
\#(\Gamma_p\cap xJ'x^{-1})},
\]
the integral for the Haar measure giving $J$ the volume $1$ (a finite sum, as $\Gamma_p\backslash
G_{p,f}$ is compact). It does not depend on $J'$, and $\int^{J}=[K_p:J]\int^{K_p}$ for $J\subseteq
K_p$. With the measures of \S\ref{pc:ss-data} (mass $1$ on $G_{p,\infty}$, volume $1$ on $K_p$),
Fubini's theorem gives, for continuous $F$ on $[G_p]=\Gamma_p\backslash(G_{p,\infty}\times G_{p,f})$ right invariant under an open subgroup $J'$ of $G_{p,f}$ (the group $\Gamma_p\cap xJ'x^{-1}$ acts freely on $G_{p,\infty}\times xJ'$, so the fibre over $x$ has mass $\mathrm{vol}(J')/\#(\Gamma_p\cap xJ'x^{-1})$),
\begin{equation}\label{pgw:eq-unfold}
\int_{[G_p]}F=\pgwint{K_p}_{\Gamma_p\backslash G_{p,f}}\int_{G_{p,\infty}}F(ux)\,du\,dx .
\end{equation}
The \emph{level} $J_p\subseteq K_p$ of the node $p$ is the preimage of the level of $p$ in
$\auGL_{\mathsf F_p}(R)$ (\S\ref{pgv:ss-def}): $J_p=J_p(\pcm)$ for $p=(0,0,m)$ (the node $Y_m$) and
$J_p=J_p(\pcb)$ otherwise. Thus $[K_p:J_p]=\#\auGL_{N_p}(\auO_E/\pcb)$ does not depend on
$\vartheta$ for $p\ne(0,0,m)$.

\paragraph*{Vector forms}
Let $p$ be definite and $(W,\sigma)$ a finite-dimensional complex representation of $G_{p,\infty}$;
an element $\gamma\in\Gamma_p$ acts by $\sigma(\gamma)=\sigma(\gamma_\infty)$. A \emph{$W$-valued form}
is a function $\psi\colon G_{p,f}\to W$, right invariant under an open subgroup, with
$\psi(\gamma x)=\sigma(\gamma)\psi(x)$ for $\gamma\in\Gamma_p$. For $\mu\in W^*$ put
\[
f_{\mu,\psi}(g)=\mu\bigl(\sigma(g_\infty)^{-1}\psi(g_f)\bigr)\qquad(g\in G_p(\auA)),
\]
a function on $[G_p]$: $\sigma(\gamma_\infty g_\infty)^{-1}\psi(\gamma g_f)=\sigma(g_\infty)^{-1}\psi(g_f)$.
(For a weight $\lambda$ and $\mu=\langle v^+,\cdot\,\rangle$ this is the function $f_\varphi$ of
Lemma~\ref{pc:algpet}; we use an abstract representation and a linear form so as not to depend on
the conventions for $W_\lambda$.)

\begin{lemma}[Forms of a type on a definite space]\label{pgw:vec}
Let $p$ be definite, $e$ a type at $p$ with representation $\pi$ and unit vector $v_0$, and let
$W=\pi^*$ with the contragredient action $\sigma(u)\lambda=\lambda\circ\pi(u)^{-1}$ and
$\mu_0=\pgwev_{v_0}\in W^*$. Then $\psi\mapsto f_{\mu_0,\psi}$ is a bijection from the $W$-valued
forms onto $\pcF(p,e)$, and $f_{\mu_0,\psi}$ is fixed by an open subgroup $J$ if and only if $\psi$ is.
\end{lemma}
\begin{proof}
By \S\ref{pc:ss-data} every $\varphi\in\pcF(p,e)$ is $\Phi(v_0)$ for a map $\Phi\colon\pi\to
L^2([G_p])$ intertwining $\pi$ with right translation; replacing $\Phi$ by its composite with the
projection to the forms fixed by the open subgroup $J$ fixing $\varphi$ (which commutes with right
translation by $G_{p,\infty}$) we may assume that the $\Phi(v)$ are fixed by $J$, and, as $\pi$ is
finite-dimensional and $\varphi$ continuous, that they are continuous (they are linear combinations
of right translates of $\varphi$ by elements of $G_{p,\infty}$, by irreducibility). Put
$\Psi(g)(v)=\Phi(v)(g)$, so $\Psi(g)\in\pi^*$, $\Psi(gu)=\Psi(g)\circ\pi(u)$ for $u\in
G_{p,\infty}$ and $\Psi(\gamma g)=\Psi(g)$. Let $\psi(x)=\Psi(x)$ for $x\in G_{p,f}$ (viewed in
$G_p(\auA)$ with trivial archimedean component). Then $\Psi(g)=\Psi(g_fg_\infty)=\sigma(g_\infty)^{-1}
\psi(g_f)$, hence $\varphi(g)=\Psi(g)(v_0)=f_{\mu_0,\psi}(g)$. For $\gamma\in\Gamma_p$, written
$\gamma=\gamma_\infty\gamma_f$ in $G_p(\auA)$, and $x\in G_{p,f}$: $\Psi(\gamma_fx)=\Psi(\gamma x\,
\gamma_\infty^{-1})=\Psi(\gamma x)\circ\pi(\gamma_\infty)^{-1}=\Psi(x)\circ\pi(\gamma_\infty)^{-1}$, that
is, $\psi(\gamma x)=\sigma(\gamma)\psi(x)$ in the notation above; and $\psi$ is fixed by $J$. Uniqueness: $f_{\mu_0,\psi}(ux)=\psi(x)(\pi(u)v_0)$
determines $\psi(x)$ on the span of the $\pi(u)v_0$, which is $\pi$. Conversely, for a $W$-valued
form $\psi$ the map $\Phi(v)(g)=\psi(g_f)(\pi(g_\infty)v)$ intertwines $\pi$ with right translation,
its values are bounded continuous functions fixed by the subgroup fixing $\psi$, and
$f_{\mu_0,\psi}=\Phi(v_0)$; by \S\ref{pc:ss-data} (right convolution with $e$ is the orthogonal
projection onto the image of $\auHom_{G_{p,\infty}}(\pi,L^2)\otimes v_0$) it lies in $\pcF(p,e)$.
\end{proof}

\paragraph*{One lift between definite spaces}
Let $(p,q)$ be a right edge of the row $m$ (of either grid), with characters
$\alpha=\alpha_{pq}$, $\beta=\beta_{pq}$ and joint harmonic $P=P_{pq}$ (Definition~\ref{pgv:def-edges}).
Then $\omega_{\alpha,\beta}=\omega_{\alpha,\beta,\infty}\otimes\omega_{\alpha,\beta,f}$, and the
archimedean factor does not depend on $\vartheta$, which is trivial at the infinite places. Let
$\pgwP$ be the space of polynomials in the Fock model of $\auHom(p,q)\otimes\R$; it is stable under
$\omega_{\alpha,\beta,\infty}(G_{p,\infty}\times G_{q,\infty})$, which preserves the degree, and
$P\in\pgwP$. For a finite datum $\Phi_f$ (a vector of the restricted tensor product of the local
Schwartz spaces at the finite places) and $Q\in\pgwP$ the theta functional $\theta(Q\otimes\Phi_f)$ is
defined and linear in $Q$; for a finite-dimensional subspace of $\pgwP$ it is a linear form, and
$\theta(\omega_{\alpha,\beta}(\gamma,\gamma')F)=\theta(F)$ for rational $\gamma,\gamma'$, because
$\alpha,\beta$ are automorphic.

\begin{lemma}[Compact Howe duality]\label{pgw:howe}
$\pgwP=\bigoplus_\nu\pgwP_\nu$, where each $\pgwP_\nu$ is stable under $G_{p,\infty}\times
G_{q,\infty}$ and isomorphic to $A_\nu\otimes B_\nu$ with $A_\nu$, $B_\nu$ irreducible
finite-dimensional representations of $G_{p,\infty}$, $G_{q,\infty}$; the $A_\nu$ are pairwise
inequivalent, and so are the $B_\nu$.
\end{lemma}
\begin{proof}
In the Fock model, $\omega_{\alpha,\beta,\infty}(g,h)$ acts on $\pgwP$ by $f(z)\mapsto c(g,h)\,f(h^{-1}zg)$ on $\prod_\tau\pcMat_{N_q\times N_p}(\C)$, with $c$ a character of $G_{p,\infty}\times G_{q,\infty}$ (the product of $\alpha_\infty(\det g)^{-1}\beta_\infty(\det h)$ and the determinant factor of the Fock model; the splitting characters make this a genuine representation). Up to the twist by $c$, this is the decomposition of the polynomials on $\pcMat_{N_q\times N_p}(\C)$
under $\auGL_{N_q}\times\auGL_{N_p}$ into the modules $W_\nu(\C^{N_q})^*\otimes W_\nu(\C^{N_p})^*$,
each once, $\nu$ running over the diagrams with at most $\min(N_p,N_q)$ rows, applied at each
$\tau\in\auPl$ and restricted to the unitary groups, as in {\cite[Lemma~5.14]{P49}}
(\cite{HoweTheta}, \cite[Theorem~5.6.7]{GoodmanWallach}, \cite{KashiwaraVergne}); distinct families
$(\nu_\tau)_\tau$ give inequivalent representations of $\prod_\tau\auU(N_p)$ and of
$\prod_\tau\auU(N_q)$, because a diagram with at most $\min(N_p,N_q)$ rows is determined by the
highest weight of $W_\nu(\C^N)$ for $N\ge\min(N_p,N_q)$. Twisting by the character $c$ changes $A_\nu,B_\nu$ by fixed characters and preserves
irreducibility and pairwise inequivalence.
\end{proof}

\begin{lemma}[One lift]\label{pgw:one}
Let $(W,\sigma)$ be an irreducible representation of $G_{p,\infty}$, or $W=0$, and $\mu\in W^*$ (if $W=0$ take $W'=0$). There are an
irreducible representation $(W',\sigma')$ of $G_{q,\infty}$ (possibly $W'=0$), a vector $\mu'\in
W'^*$, and for every finite datum $\Phi_f$ a linear map $\pgwK(\Phi_f)\colon W\to W'$, linear in
$\Phi_f$, all depending only on $(p,q)$, $W$, $\mu$ and not on $\vartheta$, such that for every
$m$-fold twist $\vartheta$ (with $\omega=\omega_{\alpha,\beta}$ formed with $\vartheta$):
\begin{enumerate}
\item[(a)] $\pgwK(\omega_f(\gamma x,\gamma'h)\Phi_f)=\sigma'(\gamma')\,\pgwK(\omega_f(x,h)\Phi_f)\,
\sigma(\gamma)^{-1}$ for $x\in G_{p,f}$, $h\in G_{q,f}$, $\gamma\in\Gamma_p$, $\gamma'\in\Gamma_q$;
\item[(b)] for every $W$-valued form $\psi$ and every finite datum $\Phi_f$,
\[
\int_{[G_p]}\theta\bigl(\omega(g,h)(P\otimes\Phi_f)\bigr)\,f_{\mu,\psi}(g)\,dg=f_{\mu',\psi'}(h),\qquad
\psi'(h)=\pgwint{K_p}_{\Gamma_p\backslash G_{p,f}}\pgwK\bigl(\omega_f(x,h)\Phi_f\bigr)\psi(x)\,dx,
\]
and $\psi'$ is a $W'$-valued form.
\end{enumerate}
\end{lemma}
\begin{proof}
\emph{The data.} If no $A_\nu$ of Lemma~\ref{pgw:howe} is isomorphic to $W$, put $W'=0$. Otherwise
let $\nu_0$ be the unique $\nu$ with $A_\nu\cong W$, $T=B_{\nu_0}$ with its action $\tau$, and fix an
isomorphism $j\colon W\otimes T\to\pgwP_{\nu_0}$ of $G_{p,\infty}\times G_{q,\infty}$-modules. Let
$W'=T^*$ with $\sigma'(h)\lambda=\lambda\circ\tau(h)^{-1}$; it is irreducible. For $Q\in\pgwP$ and
$w\in W$ put
\[
Q_w=\int_{G_{p,\infty}}\mu\bigl(\sigma(u)^{-1}w\bigr)\,\omega_\infty(u,1)Q\,du ,
\]
an integral in the finite-dimensional space of polynomials of degree at most $\deg Q$. The
substitution $u\mapsto u_0u$ gives $Q_{\sigma(u_0)w}=\omega_\infty(u_0,1)Q_w$, so $w\mapsto Q_w$ is
$G_{p,\infty}$-equivariant. Writing $Q=\sum_\nu Q^\nu$ with $Q^\nu\in\pgwP_\nu$, the map $w\mapsto
(Q^\nu)_w$ is an equivariant map $W\to\pgwP_\nu\cong A_\nu^{\oplus\dim B_\nu}$, hence zero unless
$A_\nu\cong W$. So $Q_w=0$ for all $w$ if $W'=0$; otherwise $w\mapsto Q_w$ is an equivariant map
$W\to W\otimes T$, and by Schur's lemma there is a unique $t(Q)\in T$ with $Q_w=j(w\otimes t(Q))$ for
all $w$; $t$ is linear. Since $\omega_\infty(1,h)$ commutes with $\omega_\infty(u,1)$ and $j$ is
$G_{q,\infty}$-equivariant, $t(\omega_\infty(1,h)Q)=\tau(h)t(Q)$. Put $\mu'=\pgwev_{t(P)}\in W'^*$
and define $\pgwK(\Phi_f)\in\auHom(W,W')$ by
\[
\bigl(\pgwK(\Phi_f)w\bigr)(s)=\theta\bigl(j(w\otimes s)\otimes\Phi_f\bigr)\qquad(w\in W,\ s\in T).
\]
None of these depends on $\vartheta$.

(a) For rational $\gamma,\gamma'$, $\omega(\gamma,\gamma')=\omega_\infty(\gamma,\gamma')\otimes
\omega_f(\gamma,\gamma')$ and $\theta\circ\omega(\gamma,\gamma')=\theta$, so
\[
\theta\bigl(j(w\otimes s)\otimes\omega_f(\gamma x,\gamma'h)\Phi_f\bigr)=\theta\bigl(\omega_\infty(\gamma,
\gamma')^{-1}j(w\otimes s)\otimes\omega_f(x,h)\Phi_f\bigr)=\theta\bigl(j(\sigma(\gamma)^{-1}w\otimes
\tau(\gamma')^{-1}s)\otimes\omega_f(x,h)\Phi_f\bigr),
\]
which is $\bigl(\sigma'(\gamma')\pgwK(\omega_f(x,h)\Phi_f)\sigma(\gamma)^{-1}w\bigr)(s)$.

(b) By (a) with $\gamma'=1$, $x\mapsto\pgwK(\omega_f(x,h)\Phi_f)\psi(x)$ is left
$\Gamma_p$-invariant; it is right invariant under the open subgroup fixing $\psi$ and fixing
$\Phi_f$ under $\omega_f(\cdot,1)$, so $\psi'$ is defined, and by (a) with $\gamma=1$,
$\psi'(\gamma'h)=\sigma'(\gamma')\psi'(h)$; $\psi'$ is fixed by the open subgroup of $G_{q,f}$ fixing
$\Phi_f$ under $\omega_f(1,\cdot)$. By \eqref{pgw:eq-unfold} the left side of (b) is
\[
\pgwint{K_p}_{\Gamma_p\backslash G_{p,f}}\int_{G_{p,\infty}}\theta\bigl(\omega_\infty(u,1)Q\otimes
\omega_f(x,h_f)\Phi_f\bigr)\,\mu\bigl(\sigma(u)^{-1}\psi(x)\bigr)\,du\,dx,\qquad Q=\omega_\infty(1,
h_\infty)P,
\]
and the inner integral is $\theta(Q_{\psi(x)}\otimes\omega_f(x,h_f)\Phi_f)$ by linearity of $\theta$
on a finite-dimensional space. If $W'=0$ this vanishes, and so does $f_{\mu',\psi'}$. Otherwise
$Q_{\psi(x)}=j(\psi(x)\otimes\tau(h_\infty)t(P))$, and the inner integral is
\[
\bigl(\pgwK(\omega_f(x,h_f)\Phi_f)\psi(x)\bigr)\bigl(\tau(h_\infty)t(P)\bigr)=\mu'\bigl(\sigma'(h_\infty)^{-1}
\pgwK(\omega_f(x,h_f)\Phi_f)\psi(x)\bigr).
\]
Integrating over $x$ (a finite sum) gives $f_{\mu',\psi'}(h)$.
\end{proof}

\begin{remark}\label{pgw:rem-jets}
$\pgwK$ is the kernel of jets of {\cite[(34)]{P49}}: by {\cite[Lemma~5.14]{P49}} one may take
$j(w\otimes s)$ to be the polynomial $z\mapsto s(S_{\nu_0}(z)w)$, $S_{\nu_0}$ the Schur functor, and
then by {\cite[(33)]{P49}} $\pgwK(\Phi_f)$ is the matrix of the Taylor coefficients of the theta
function of $\Phi_f$ along the entries of $S_{\nu_0}$. A different choice of $j$ multiplies
$\pgwK$ by a non-zero constant and $\mu'$ by its inverse (Schur's lemma for the irreducible module
$W\otimes T$). Only the existence of $\pgwK$ with (a), (b) is used below; the integrality of the
kernels (needed for Proposition~\ref{pgr:den}) concerns the explicit choice of
Definition~\ref{pgz:def-choice}(a).
\end{remark}

\paragraph*{The branching pairing}
Let $Y=Y^m_m$ and $\pgrY^m_m=Y\oplus E\tilde e_m$, identified as spaces with frames
(\S\ref{pgv:ss-def}); both are definite, of ranks $n$ and $n+1$. For $x\in G_Y(\auA)$ and
$u\in\auU(1)(\auA_{E^+})$ we write $x\oplus u\in G_{\pgrY^m_m}(\auA)$.

\begin{lemma}[Branching]\label{pgw:branch}
Let $(W_0,\sigma_0)$, $(W_1,\sigma_1)$ be irreducible representations of $G_{Y,\infty}$ and
$G_{\pgrY^m_m,\infty}$, and $\mu_z\in W_z^*$.
\begin{enumerate}
\item[(a)] The bilinear forms $B$ on $W_0\times W_1$ with $B(\sigma_0(x)w_0,\sigma_1(x\oplus1)w_1)=
B(w_0,w_1)$ for $x\in G_{Y,\infty}$ form a space of dimension at most $1$.
\item[(b)] Fix such a form $B$, non-zero if one exists. For every real point $g$ of
$G_{\pgrY^m_m}$ there is $c(g)\in\C$ such that for all $W_0$-valued forms $\eta$ on $G_Y$, all
$W_1$-valued forms $\varphi$ on $G_{\pgrY^m_m}$ and all $u\in\auU(1)(\auA_{E^+,f})$
\[
\int_{[G_Y]}f_{\mu_0,\eta}(x)\,f_{\mu_1,\varphi}\bigl((x\oplus u)g\bigr)\,dx=c(g)\,
\pgwint{K_Y}_{\Gamma_Y\backslash G_{Y,f}}B\bigl(\eta(h),\varphi(h\oplus u)\bigr)\,dh .
\]
\end{enumerate}
\end{lemma}
\begin{proof}
(a) $G_{\pgrY^m_m,\infty}=\prod_\tau\auU(n+1)$ and $W_1=\bigotimes_\tau W_{1,\tau}$. The restriction
of an irreducible representation of $\auU(n+1)$ to $\auU(n)$ (embedded as $x\mapsto x\oplus1$) is
multiplicity free (the branching law of Weyl, \cite[Theorem~8.1.1]{GoodmanWallach}). Hence the restriction of $W_1$ to $\prod_\tau\auU(n)$ is a direct sum of pairwise
inequivalent irreducible representations $\bigotimes_\tau V_\tau$, and an invariant form $B$ is an
element of $\auHom_{G_{Y,\infty}}(W_1,W_0^*)$, a space of dimension at most $1$ since $W_0^*$ is
irreducible.

(b) Put $b_g(w_0,w_1)=\int_{G_{Y,\infty}}\mu_0(\sigma_0(x)^{-1}w_0)\,\mu_1\bigl(\sigma_1(g)^{-1}
\sigma_1(x\oplus1)^{-1}w_1\bigr)dx$. The substitution $x\mapsto yx$ shows that $b_g$ is invariant in
the sense of (a), so $b_g=c(g)B$ for some $c(g)$ (and $b_g=0$ if no non-zero invariant form exists).
The function $h\mapsto B(\eta(h),\varphi(h\oplus u))$ is left $\Gamma_Y$-invariant, because
$\varphi((\gamma h)\oplus u)=\varphi((\gamma\oplus1)(h\oplus u))=\sigma_1(\gamma\oplus1)\varphi(h\oplus u)$
and $(\gamma\oplus1)_\infty=\gamma_\infty\oplus1$; it is right invariant under an open subgroup. For
$x=x_\infty x_f$ the point $(x\oplus u)g$ has archimedean component $(x_\infty\oplus1)g$ and finite
component $x_f\oplus u$, so by \eqref{pgw:eq-unfold} the left side is
\[
\pgwint{K_Y}_{\Gamma_Y\backslash G_{Y,f}}\int_{G_{Y,\infty}}\mu_0\bigl(\sigma_0(x)^{-1}\eta(h)\bigr)\,
\mu_1\bigl(\sigma_1(g)^{-1}\sigma_1(x\oplus1)^{-1}\varphi(h\oplus u)\bigr)\,dx\,dh=\pgwint{K_Y}c(g)\,
B\bigl(\eta(h),\varphi(h\oplus u)\bigr)\,dh. \qedhere
\]
\end{proof}

\paragraph*{The algebraic lifts along the last path and the sum $\Lambda$}
For $z\in\{0,1\}$ and $1\le a\le m$ let $e^z_a$ be the edge $((z,a-1,m),(z,a,m))$; these are the
edges of the half $h_z$ of $\pgrlast$. Let $\pi^z_0$ be the representation of the type at the outer
node $(z,0,m)$, with unit vector $v^z_0$, $W^z_0=(\pi^z_0)^*$ and $\mu^z_0=\pgwev_{v^z_0}$
(Lemma~\ref{pgw:vec}); applying Lemma~\ref{pgw:one} successively to $e^z_1,\dots,e^z_m$ defines
irreducible representations $W^z_a$ of $G_{(z,a,m),\infty}$, vectors $\mu^z_a\in(W^z_a)^*$ and
kernels $\pgwK^z_a$ for $e^z_a$. For a matrix $X\in\pcMat_{q\times p}(R)$ on the edge $e^z_a=(p,q)$ let
$[\delta_X]_f=\delta_X^\natural\otimes\psvac^\pcm$ be the finite part of the datum of
\S\ref{pc:ss-data}, and define the \emph{algebraic lift}
\[
(\pgwT^z_{a,X}\psi)(h)=\pgwint{J_p}_{\Gamma_p\backslash G_{p,f}}\pgwK^z_a\bigl(\omega_f(x,h)[\delta_X]_f
\bigr)\psi(x)\,dx ,
\]
for the measure giving the level $J_p$ of the node $p$ the volume $1$. For an edge tuple $S$ put
$\pgwT^z_S=\pgwT^z_{m,S_{e^z_m}}\circ\dots\circ\pgwT^z_{1,S_{e^z_1}}$. Fix a form $B$ as in
Lemma~\ref{pgw:branch}(a) for $(W^0_m,W^1_m)$, non-zero if possible. For a function $\Psi$ of edge
tuples, a form $y'$ of type $e_{(0,0,m)}$ (a conjugate form on $Y_m$) and a form $\varphi$ of type
$e_{(1,0,m)}$ on $\pgrY^m_0$, with $y'=f_{\mu^0_0,\psi_{y'}}$ and $\varphi=f_{\mu^1_0,\psi_\varphi}$
(Lemma~\ref{pgw:vec}), put
\begin{equation}\label{pgw:eq-Lambda}
\Lambda(\Psi;y',\varphi)=\sum_{S\in\mathcal S_\vartheta(\pgrlast)}\Psi(S)\,
\pgwint{J_Y}_{\Gamma_Y\backslash G_{Y,f}}B\bigl((\pgwT^0_S\psi_{y'})(h),\,(\pgwT^1_S\psi_\varphi)
(h\oplus\pgupt_\pcm^m)\bigr)\,dh ,
\end{equation}
where $Y=Y^m_m$, $J_Y=J_Y(\pcb)$ is its level and $\pgupt_\pcm$ is the element of
Definition~\ref{pgv:def-period}. This is the precise form of the sum $\Lambda$ of
Proposition~\ref{pgr:last}: the pairing $B$ is the branching pairing, the $\pgwT$ are the algebraic
lifts (with the kernels of jets, Remark~\ref{pgw:rem-jets}), and each class sum is normalised by the
level of the node over which it is taken. Replacing $B$ or the kernels $\pgwK^z_a$ by non-zero
multiples multiplies $\Lambda$ by a constant independent of $\vartheta$, $\Psi$, $y'$, $\varphi$.

\begin{theorem}[Proposition~\ref{pgr:last}]\label{pgw:last}
Let $\Xi_{\mathrm{last}}(\Psi)(\vartheta,y',\varphi)=\Lambda(\Psi_\vartheta;y',\varphi)/
\#\auGL_{r+mt}(R_{\pcm_\vartheta})$, with $\Lambda$ as in \eqref{pgw:eq-Lambda} and
$\pcm_\vartheta=(\pi_M)\prod_{w\in Q_\vartheta}w$. Then $\Psi\propto_{\pgrlast}\Xi_{\mathrm{last}}(\Psi)$
for every family $\Psi=(\Psi_\vartheta)$ of functions of edge tuples.
\end{theorem}
\begin{proof}
Fix a real point $g$ of $G_{\pgrY^m_m}$. Let $\vartheta$, $y'$, $\varphi$, $S$ be given and write
$X_a=S_{e^z_a}$ when the half is clear. For $p=(z,a-1,m)$ the algebraic lift with the measure
normalised by $K_p$ is $[K_p:J_p]^{-1}\pgwT^z_{a,X}$. By Lemma~\ref{pgw:one}(b), applied $m$ times
in each half (the move along $e^z_a$ is the lift with the datum $P_{e^z_a}\otimes[\delta_{X_a}]_f$,
Definition~\ref{pgv:def-edges}),
\[
h_{z*}(S)F=f_{\mu^z_m,\,\psi^z_S},\qquad\psi^z_S=\Bigl(\prod_{a=0}^{m-1}[K_{(z,a,m)}:J_{(z,a,m)}]
\Bigr)^{-1}\pgwT^z_S\psi_F ,
\]
for $F=y'$ ($z=0$) and $F=\varphi$ ($z=1$). By Definition~\ref{pgv:def-period} (that is, (G4)) and
Lemma~\ref{pgw:branch}(b) with $u=\pgupt_\pcm^m$,
\[
\begin{aligned}\mathsf P_{m,g}\bigl(h_{0*}(S)y',h_{1*}(S)\varphi\bigr)&=c(g)\pgwint{K_Y}_{\Gamma_Y\backslash G_{Y,f}}
B\bigl(\psi^0_S(h),\psi^1_S(h\oplus\pgupt_\pcm^m)\bigr)\,dh\\&=\frac{c(g)}{D(\vartheta)}\pgwint{J_Y}
B\bigl((\pgwT^0_S\psi_{y'})(h),(\pgwT^1_S\psi_\varphi)(h\oplus\pgupt_\pcm^m)\bigr)\,dh,\end{aligned}
\]
with $D(\vartheta)=[K_Y:J_Y]\prod_{z}\prod_{a<m}[K_{(z,a,m)}:J_{(z,a,m)}]$. Multiplying by
$\Psi_\vartheta(S)$ and summing over $S\in\mathcal S_\vartheta(\pgrlast)$ gives
$I_{\pgrlast}(\Psi;g)(\vartheta,y',\varphi)=c(g)\Lambda(\Psi_\vartheta;y',\varphi)/D(\vartheta)$.

All the indices in $D(\vartheta)$ are those of $J(\pcb)$, independent of $\vartheta$, except
$[K_{Y_m}:J_{Y_m}(\pcm)]=\#\auGL_N(\auO_E/\pcm)$ for the node $(0,0,m)=Y_m$, $N=r+mt$. The primes of
$\pi_M$ are distinct from those of $Q_\vartheta$ and $(\pi_M)$ is squarefree (\S\ref{pc:ss-data}:
$R_\vartheta$ is a product of fields), so by the Chinese remainder theorem
$\#\auGL_N(\auO_E/\pcm)=\#\auGL_N(\auO_E/\pcb)\prod_{w\in Q_\vartheta}\#\auGL_N(\auO_E/w)=\kappa_0\,
\#\auGL_N(R_{\pcm_\vartheta})$ with $\kappa_0=\#\auGL_N(\auO_E/\pcb)/\#\auGL_N(\auO_E/(\pi_M))$. Hence
$D(\vartheta)=D_0\,\#\auGL_N(R_{\pcm_\vartheta})$ with $D_0>0$ independent of $\vartheta$, and
$I_{\pgrlast}(\Psi;g)(x)=\beta(g)\,\Xi_{\mathrm{last}}(\Psi)(x)$ for all period triples $x$, with
$\beta(g)=c(g)/D_0$.
\end{proof}

\noindent The proof uses only the definitions of \S\ref{pgv:ss-def} (spaces, moves, levels, the
diagonal period (G4)), the compact Howe duality and Weyl's branching law; in particular it does not
use (G3) beyond the definition of the joint harmonic $P_{pq}$ as a polynomial. (That the moves take
forms of type $e_p$ to forms of type $e_q$, which (G3) asserts, is recovered here: $W^z_{a}$ is then
the representation of the type at $(z,a,m)$, but this is not needed.)




% ===== C-grid6.tex =====
% proofs/C-grid6.tex
% Continuation of proofs/C-grid.tex (pgr:), C-grid2.tex (pgs:), C-grid3.tex (pgu:), C-grid4.tex
% (pgv:): the line reversal without the unit assumption on alpha, the caveat on (G1), and
% Proposition pgr:den. Labels start with pgx:. Conventions and macros: those of proofs/C-grid3.tex.
\providecommand{\auI}{\mathcal I}
\providecommand{\auGL}{\mathrm{GL}}
\providecommand{\auU}{\mathrm U}
\providecommand{\auO}{\mathcal O}
\providecommand{\auA}{\mathbb A}
\providecommand{\auNm}{\operatorname{Nm}}
\providecommand{\auHom}{\operatorname{Hom}}
\providecommand{\aubar}[1]{\overline{#1}}
\providecommand{\pcm}{\mathfrak m}
\providecommand{\pcMat}{\mathrm M}
\providecommand{\pcR}{\mathsf R}
\providecommand{\psHS}{\mathcal H}
\providecommand{\psrho}{\varrho}
\providecommand{\psvac}{\mathbf v}
\providecommand{\pgrY}{\widetilde Y}
\providecommand{\pgrTr}{\mathrel{\vdash}}
\providecommand{\pcnot}[1]{\par\smallskip{\raggedright\noindent\textsc{Not proved here.} [#1]\par}\smallskip}
\providecommand{\pguK}{\mathsf K}
\providecommand{\pgupt}{\mathsf u}

\subsection{The line reversal at all primes of the modulus; the identification (G1)}\label{pgx:sec}

This section continues \S\S\ref{pgu:sec}, \ref{pgv:sec} and uses their notation. Lemma~\ref{pgu:line}
assumes that the finite idele $\alpha$ of {\cite[\S5.1.1]{P49}} is a unit at the primes of $\pcm$, and
Proposition~\ref{pgv:alpha}(b) shows that this can fail at the primes of $Q_\vartheta$ above $q_1$.
Here no such assumption is made: the data at the primes of $\pcm$ are placed on the self-dual lattice
$\alpha\auHom$ (Definition~\ref{pgx:conv}), and the reversal of a line is proved with the element
$\pgupt=\pgupt_\pcm=\aubar u_0/u_0$, $u_0=\alpha\varpi_\pcm$, of Definition~\ref{pgv:def-period}
(Lemma~\ref{pgx:line}, which contains Lemma~\ref{pgu:line}). \S\ref{pgx:ss-later} lists the uses of
the line reversal in \S\ref{pgu:sec}, \S\ref{pgx:ss-other} the other places that use the
data at the primes of a modulus, and \S\ref{pgx:ss-G1} shows that the seesaw steps use (G1) only as an
identification of spaces.

\subsubsection{The datum at the primes of the modulus}\label{pgx:ss-conv}

Let $\pcm$ be a modulus (\S\ref{pc:ss-data}), $v=w\aubar w\in S(\pcm)$ with $w\mid\pcm$,
$k_w=\operatorname{ord}_w\pcm$, and let $S,T$ be spaces of ranks $N,N'$. By (W5),
$\psHS_v(\auHom(S,T))$ is the space of Schwartz--Bruhat functions on $X_w=\auHom(S_w,T_w)=
\pcMat_{N'\times N}(E_w)$ (matrices in the frames), $\psrho(a)$ is translation by $a$ for $a$ in the
$w$-component and multiplication by $x\mapsto[x,a']$ for $a'$ in the $\aubar w$-component $X_{\aubar w}$.

\begin{lemma}[The vacuum at a split place]\label{pgx:vac}
The self-dual lattice $M_v=\alpha\auHom(L_S,L_T)_v$ is $M_w\oplus M_{\aubar w}$ with
$M_w=\alpha_w\pcMat_{N'\times N}(\auO_w)$ and $M_{\aubar w}$ the dual of $M_w$ for the pairing
$(x,a')\mapsto[x,a']$ of $X_w\times X_{\aubar w}$. In the model (W5) the vacuum $\psvac_{M_v}$ is,
up to a scalar, the indicator function $\mathbf 1_{M_w}$.
\end{lemma}
\begin{proof}
The lattice $L_S\otimes\auO_v$ is $L_{S,w}\oplus L_{S,\aubar w}$, and similarly for $T$, so $M_v$ is
the direct sum of its components $M_w=\alpha_w\auHom(L_{S,w},L_{T,w})$, which is
$\alpha_w\pcMat_{N'\times N}(\auO_w)$ in the frames, and $M_{\aubar w}$. The hermitian form of
$\auHom(S,T)$ pairs the $w$-component with the $\aubar w$-component and vanishes on each of them; so
$[\,,]$ is trivial on $X_w\times X_w$ and on $X_{\aubar w}\times X_{\aubar w}$, and the self-duality
of $M_v$ says exactly that $M_{\aubar w}$ is the dual of $M_w$ under $[\,,]$. The function
$\mathbf 1_{M_w}$ is invariant under the translations by $M_w$ and under the multiplications by
$[\,\cdot\,,a']$, $a'\in M_{\aubar w}$ (these are trivial on $M_w$). On $M_v$ the character
$\epsilon_{M_v}$ of (W4) is trivial on $M_w$ and on $M_{\aubar w}$ (since $\beta$ is trivial on each
component), so by (W4) $\psvac_{M_v}$ is the unique line with these invariances, which is
$\C\mathbf 1_{M_w}$.
\end{proof}

\begin{definition}[Data at the primes of a modulus]\label{pgx:conv}
For every $w\mid\pcm$ let $a_w$ be the generator of the ideal $\alpha_w\auO_w$ fixed in
{\cite[\S5.1.1]{P49}} ($a_w=\alpha_w$ if $\alpha_w$ is not a unit, $a_w=1$ otherwise). For a function $\Psi$ on
$\pcMat_{N'\times N}(\pcR_\pcm)$ let $\Psi^\natural$ be the function on $\prod_{w\mid\pcm}
\pcMat_{N'\times N}(E_w)$ equal to $\Psi(a^{-1}x\bmod\pcm)$ at $x\in a\,\pcMat_{N'\times N}
(\widehat\auO)$ (that is, $x_w\in a_w\pcMat(\auO_w)$ for all $w\mid\pcm$) and to $0$ elsewhere, and
use at $v\in S(\pcm)$ the $L^2$-norm for the measure giving $M_w=a_w\pcMat(\auO_w)$ the volume $1$.
The datum $[P,\Psi]$, the kernels, the lifts and the transposed lifts are then defined by the
formulas of \S\ref{pc:ss-data}.
\end{definition}

\noindent At the primes where $\alpha_w$ is a unit $a_w=1$; then Definition~\ref{pgx:conv}
is literally the definition of \S\ref{pc:ss-data}. The two definitions differ only at the primes of
$\pcm$ where $\alpha$ is not a unit, which by Proposition~\ref{pgv:alpha}(a) lie above $q_1$; at
such a prime a function supported on the integral matrices would not be the vacuum
(Lemma~\ref{pgx:vac}); for this reason \S\ref{pc:ss-data} uses Definition~\ref{pgx:conv}. With it,
$\mathbf 1^\natural=\mathbf 1_{M_w}$ is the vacuum at every $w\mid\pcm$, and
$\|\Psi^\natural\|^2=(\#\pcMat_{N'\times N}(\pcR_\pcm))^{-1}\sum|\Psi|^2$ exactly as before. Another
choice of the generators, $a'_w=a_w\varepsilon_w$ with units $\varepsilon_w$, replaces $\Psi^\natural$ by
$(\Psi\circ\bar\varepsilon^{-1})^\natural$, $\bar\varepsilon$ the image of $\varepsilon$ in
$\pcR_\pcm^\times$ acting by multiplication. Lemma~\ref{pc:level} holds without change: for $k\in K_S$,
$k'\in K_T$ the matrix $k'^{-1}xk$ lies in $a\pcMat(\widehat\auO)$ iff $x$ does, and
$a^{-1}k'^{-1}xk=k'^{-1}(a^{-1}x)k$ since $a$ is a scalar; the proof of that lemma then applies word
for word. The same remark shows that Corollary~\ref{pgu:split} holds with Definition~\ref{pgx:conv}:
$\delta^\natural_{(y_2|v)}(x_1|x_2)=\delta^\natural_{y_2}(x_1)\delta^\natural_v(x_2)$, the lattices
$a_w\pcMat(\auO_w)$ being products of the corresponding lattices of the two summands and the
measures product measures.

\paragraph*{The coordinates at split places}
The idele $\alpha$ enters the data at a split place in three ways: through the frames (coordinates
$x$), through the coordinates $T$ of (N5) ({\cite[\S6.1]{P49}}), and through the coordinates $a^{-1}x$
of Definition~\ref{pgx:conv}. The following lemma fixes one coordinate and gives the dictionary; all
data at split places in the proof of the Covolume Theorem are written in it.

\begin{lemma}[The normalised coordinates]\label{pgx:dict}
Let $a$ be the idele of {\cite[\S5.1.1]{P49}}, $S,T$ spaces of ranks $N,N'$, $v=w\aubar w$ a place of
$E^+$ split in $E$, and $X_w,X_{\aubar w}$ the components of $\auHom(S,T)_v$, written as matrices $x$
in the frames. Put $\tilde x=a_w^{-1}x$ on $X_w$ and $\tilde x=a_{\aubar w}^{-1}x$ on $X_{\aubar w}$, and
more generally $\tilde x=a^{-1}x$ on $\auHom(S,T)(\auA_f)$.
\begin{enumerate}
\item[(a)] $\alpha_w$ is not a unit if and only if $w$ lies above a prime $p\mid q_1$ and $w\mid\omega$;
then $w\ne\aubar w$, $\aubar w\nmid\omega$ and $\alpha_{\aubar w}=a_{\aubar w}=1$. Every place of $E^+$
above $p\mid q_1$ is $w\aubar w$ with exactly one of $w,\aubar w$ dividing $\omega$. No chosen place
and no prime above $N$ lies above $q_1$. A prime of $Q_\vartheta$ may lie above $q_1$ and divide
$\omega$; such primes are the only primes of a modulus $\pcm_\vartheta$, of $S(\pcm_\vartheta)$, or of
the moduli of the grids at which $a_w\ne1$.
\item[(b)] The self-dual lattice $\alpha\auHom(L_S,L_T)$ is $a\auHom(L_S,L_T)$; at $v$ its components
are $M_w=a_w\pcMat(\auO_w)$ and $M_{\aubar w}=a_{\aubar w}\pcMat(\auO_{\aubar w})$, i.e.\ the integral
matrices in the coordinates $\tilde x$, and they are dual to each other for $[\,,]$. For $x\in X_w$,
$y\in X_{\aubar w}$ one has $[x,y]=\lambda_w(\operatorname{tr}({}^t\tilde y\tilde x))^{\pm1}$ with
$\lambda_w(t)=\psi_w(\jmath(a\aubar a)_wt)$, trivial on $\auO_w$ and not on $\varpi_w^{-1}\auO_w$
($E_{\aubar w}$ identified with $E_w$ by conjugation). In the frame coordinates the dual of
$\pcMat(\auO_w)$ is $\aubar{a_w}\,a_{\aubar w}\pcMat(\auO_{\aubar w})$, which is not
$\pcMat(\auO_{\aubar w})$ when $a_w\ne1$ or $a_{\aubar w}\ne1$.
\item[(c)] On $X_w$ the coordinate $T$ of (N5) and the coordinate $a_w^{-1}x$ of
Definition~\ref{pgx:conv} are both $\tilde x$. Hence for $w\mid\pcm$:
$\Psi^\natural=\Psi(\tilde x\bmod\pcm)\,\mathbf 1[\tilde x\ \text{integral}]$; the datum
$\delta_\iota^\natural$ is $\mathbf 1_{a_w(\iota+\varpi\pcMat(\auO_w))}$ in the frames, which is the
datum of {\cite[\S5.1.3]{P49}} at a level place; and a datum defined as a function of $T$ by (N5) (as in
{\cite[\S\S2.5, 3.2 and~4.1.3]{P54}}) is the same function of $\tilde x$, so it is
of the form $\Psi^\natural$ exactly when, as a function of $T$, it is supported on the integral
matrices and depends only on $T$ modulo $\pcm$.
\item[(d)] For $u\in\auHom(S,T)(\auA_f)$, $u\in\alpha\auHom(L_S,L_T)$ if and only if $\tilde u$ is
integral. If $u=s_0y$ with $s_0=a$, then $\tilde u=y$; for $h\in\auU(T)(\auA_f)$,
$\widetilde{h^{-1}u}=h^{-1}y$.
\item[(e)] Let $u$ be a finite idele of $E$ with $u\aubar u=1$, acting as a central element of
$\auU(T)(\auA_f)$. Then $\omega(1,u)$ carries the vacuum of a self-dual lattice $\Lambda$ of
$\auHom(S,T)(\auA_f)$ to a multiple of modulus one of the vacuum of $u\Lambda$, and for every datum $f$,
all characters $\alpha',\beta'$ and $g\in\auU(S)(\auA)$,
\[
\int_{[\auU(T)]}\theta\bigl(\omega_{\alpha',\beta'}(g,h)\,\omega(1,u)f\bigr)\,dh=\beta'(\det u)^{-1}
\int_{[\auU(T)]}\theta\bigl(\omega_{\alpha',\beta'}(g,h)f\bigr)\,dh
\]
whenever either side converges absolutely.
\item[(f)] The map $w\mapsto w^\dagger$ carries $\alpha\auHom(L_S,L_T)$ onto $\aubar\alpha\auHom(L_T,L_S)=
u\,\alpha\auHom(L_T,L_S)$, $u=\aubar a/a$ (a finite idele with $u\aubar u=1$, central in $\auU(S)$,
equal to $1$ at the places not above $q_1$). At $v$, in the coordinates $\tilde x$ of $\auHom(S,T)$ and
$\tilde X$ of $\auHom(T,S^-)$ (frames of $S$ and $T$), it is $\tilde X=-\nu^{-1}\,{}^t\tilde x$ after
identifying $E_{\aubar w}$ with $E_w$ by conjugation, where $\nu=\aubar{a_{\aubar w}}/a_w$ is the
conjugate of $u_{\aubar w}^{-1}$.
\end{enumerate}
\end{lemma}
\begin{proof}
(a) If $w\nmid\omega$ then $\alpha_w=1$. If $w\mid\omega$ and $w\nmid q_1$, $\alpha_w$ is a unit
(Proposition~\ref{pgv:alpha}(a)). If $w\mid\omega$ and $w$ lies above $p\mid q_1$: every prime $p\mid q_1$
ramifies in $\Q(\zeta_{q_1})$ ($p$ odd, or $p=2$ and $4\mid q_1$, {\cite[\S13.1]{P46}}), so $\alpha_w$
is not a unit and $w\ne\aubar w$ (Proposition~\ref{pgv:alpha}(b)). Since $\omega+\aubar\omega=1$, $\omega$
and $\aubar\omega$ are coprime; $p$ divides $\omega\aubar\omega=(1+d)/4$ and splits in $K$ as
$\mathfrak p\aubar{\mathfrak p}$, so exactly one of $\mathfrak p,\aubar{\mathfrak p}$ divides $\omega$,
the primes of $E$ above it divide $\omega$, and their conjugates (above the other one) do not. Hence
$\aubar w\nmid\omega$ and $\alpha_{\aubar w}=1$. The chosen places lie above $p_0,p_1$ and the primes
dividing $\pi_N$ above $N$; these are $\equiv1\pmod{q_2^2}$ and $q_1\mid q_2$, so they are prime to
$q_1$. The definition of a twist ({\cite[\S14.3]{P46}}) asks only that $\auNm Q$ be prime to $MN$ and
$Q\cap\aubar Q=\emptyset$, which such primes satisfy (Proposition~\ref{pgv:alpha}(b)). The moduli in
question are divisible only by chosen places and primes of $Q_\vartheta$, which gives the last claim.

(b) $a/\alpha$ is a unit idele, so $\alpha\auHom(L_S,L_T)=a\auHom(L_S,L_T)$, whose components at $v$
are $a_w\pcMat(\auO_w)$ and $a_{\aubar w}\pcMat(\auO_{\aubar w})$; their duality is Lemma~\ref{pgx:vac}.
For the pairing: $[x,y]=\psi(\jmath\operatorname{tr}(y^\dagger x))^{\pm1}$ (the sign depending on the
order of the two components), and with $x=a_w\tilde x$, $y=a_{\aubar w}\tilde y$,
$\operatorname{tr}(y^\dagger x)=a_w\aubar{a_{\aubar w}}\operatorname{tr}(\tilde y^\dagger\tilde x)$, the
frames having unimodular Gram matrices at $v$ (which are absorbed in the coordinates, as in
(N5)); so $[x,y]=\lambda_w(\operatorname{tr}({}^t\tilde y\tilde x))^{\pm1}$. The character
$t\mapsto\psi_w(\jmath\alpha_w\aubar{\alpha_{\aubar w}}t)$ is trivial exactly on $\auO_w$ by the
self-duality of $\widehat\auO_E$ for $\psi(\jmath\alpha\aubar\alpha\,\cdot)$ ({\cite[\S5.1.1]{P49}}), and
$a_w\aubar{a_{\aubar w}}/\alpha_w\aubar{\alpha_{\aubar w}}$ is a unit. For the last assertion, for
$c\in E_w^\times$ one has $[cx,y]=[x,\aubar cy]$ ($\aubar c\in E_{\aubar w}$), so the dual of
$\pcMat(\auO_w)=a_w^{-1}M_w$ is $\aubar{a_w}M_{\aubar w}=\aubar{a_w}a_{\aubar w}\pcMat(\auO_{\aubar w})$,
and $\aubar{a_w}a_{\aubar w}$ is a unit only if $a_w=a_{\aubar w}=1$ (by (a) at most one of them is
$\ne1$, and that one is not a unit).

(c) By definition (N5) uses $T=a_w^{-1}x_w$, and Definition~\ref{pgx:conv} uses $a_w^{-1}x$ with the
same $a_w$. The formula for $\Psi^\natural$ is Definition~\ref{pgx:conv} rewritten;
$\delta_\iota^\natural(x)=1$ if and only if $a_w^{-1}x\in\iota+\varpi\pcMat(\auO_w)$. The last
assertion is a restatement.

(d) $\alpha\auHom(L_S,L_T)=a\auHom(L_S,L_T)$ by (b), and $a$ is a scalar commuting with the action of
$\auU(T)(\auA_f)$.

(e) $\omega(1,u)\psrho(l)\omega(1,u)^{-1}=\psrho(ul)$ (W1), so $\omega(1,u)\psvac_\Lambda$ is fixed by
$\psrho(u\Lambda)$ up to the character $\epsilon_{u\Lambda}$, hence is a multiple of
$\psvac_{u\Lambda}$ by (W4), of modulus one as $\omega(1,u)$ is unitary. Since $u$ is central,
$\omega_{\alpha',\beta'}(g,h)\omega(1,u)=\beta'(\det u)^{-1}\omega_{\alpha',\beta'}(g,hu)$, and the
measure on $[\auU(T)]$ is invariant under right translation by $u$.

(f) $(cx)^\dagger=\aubar cx^\dagger$ for $c\in\auA_{E,f}$, and $w\mapsto w^\dagger$ carries
$\auHom(L_S,L_T)$ onto $\auHom(L_T,L_S)$ (unimodular lattices); so $\alpha\auHom(L_S,L_T)$ goes to
$\aubar\alpha\auHom(L_T,L_S)=\aubar a\auHom(L_T,L_S)=u\,a\auHom(L_T,L_S)$, and $u$ acts on
$\auHom(T,S^-)$ through the centre of $\auU(S^-)=\auU(S)$. By (a), $u_v=1$ unless $v$ lies above
$q_1$. In the frames $x^\dagger=-{}^t\aubar x$ at $v$ (as in \S\ref{pcx:ss-flip}, the Gram matrices
being absorbed), and $x=a_w\tilde x$, $x^\dagger=a_{\aubar w}\tilde X$ give
$\tilde X=-a_{\aubar w}^{-1}\aubar{a_w}\,{}^t\aubar{\tilde x}$, which is the formula after the
identification by conjugation.
\end{proof}

\subsubsection{The line reversal}\label{pgx:ss-line}

We keep the notation of \S\ref{pgu:ss-line}: $n=\#\mathsf N_2$ is the rank of $S$, $L=E\tilde e_j$,
$L'=E\tilde e_{j+1}$ is $L$ with the opposite form, and $\pguK^{\mathrm c}_{v,B}$,
$\pguK^{\mathrm r}_{v',B'}$ are the line kernels of Corollary~\ref{pgu:split}, now with the data of
Definition~\ref{pgx:conv}. No assumption on $\alpha$ at the primes of $\pcm$ is made.

\begin{lemma}[Line reversal]\label{pgx:line}
Let $\pgupt=\pgupt_\pcm=\aubar u_0/u_0\in\auU(1)(\auA_{E^+,f})$, $u_0=\alpha\varpi_\pcm$, be the
element of Definition~\ref{pgv:def-period}. There are a linear bijection $\tau$ from the polynomials
on $\auHom(L,S)\otimes\R$ to those on $\auHom(S,L')\otimes\R$, independent of $\vartheta$ and
preserving degrees, a primitive additive character $\psi_\pcm$ of $R=\auO_E/\pcm$ depending only on
$\pcm$ (and the fixed choices of $\alpha$, $\varpi_\pcm$, $a$), and
$c_{\mathrm{line}}=c_{\mathrm{line}}(\vartheta)\in\C^\times$ with
\[
|c_{\mathrm{line}}|^2=C_0\,(\auNm\pcm)^{-n},\qquad C_0>0\text{ independent of }\vartheta,
\]
such that for every $k\in\auU(L)(\auA_f)$, $g\in\auU(S)(\auA)$, $v\in\pcMat_{\mathsf N_2\times
\bullet}(R)$ and every polynomial $B$
\[
\pguK^{\mathrm c}_{v,B}(k,g)=c_{\mathrm{line}}\,\zeta(k)\sum_{v'\in\pcMat_{\bullet\times\mathsf N_2}(R)}
\psi_\pcm(v'v)\,\pguK^{\mathrm r}_{v',\tau B}(g,k\pgupt),
\]
where $\zeta(k)$ is a number of modulus one depending only on $k$, on $\pcm$ and on the characters of
the two line kernels. In particular, for $0\le j<m$,
$\pguK^{\mathrm c}_{v,B}(\pgupt^j,g)=c'_{\mathrm{line}}\sum_{v'}\psi_\pcm(v'v)\,
\pguK^{\mathrm r}_{v',\tau B}(g,\pgupt^{j+1})$ with $c'_{\mathrm{line}}=c_{\mathrm{line}}
\zeta(\pgupt^j)$, $|c'_{\mathrm{line}}|=|c_{\mathrm{line}}|$.
\end{lemma}

\begin{proof}
\emph{The flip.} Let $\mathrm{Fl}\colon\psHS(\auHom(L,S))\to\psHS(\auHom(S,L'))$ be the
isomorphism of {\cite[Proposition~7.17]{P49}} ((W6)) for the pair $(L,S)$, $L'=L^-$:
$\mathrm{Fl}\,\psrho(w)=\psrho(w^\dagger)\mathrm{Fl}$, $\theta\circ\mathrm{Fl}=c_\theta\theta$ with
$c_\theta\ne0$, and $\mathrm{Fl}\,\omega(k,g)=\upsilon(\det k)^n\upsilon(\det g)^{-1}\omega(g,k)
\mathrm{Fl}$. Exactly as in the first step of the proof of Lemma~\ref{pgu:line} (this step does not
involve $\alpha$), using $\alpha^{\mathrm r}=(\beta^{\mathrm c})^{-1}\upsilon$ (G2),
\[
\mathrm{Fl}\,\omega^{L,S}_{\alpha_L,\beta^{\mathrm c}}(k,g)=\zeta_1(k)\,
\omega^{S,L'}_{\alpha^{\mathrm r},\beta_{L'}}(g,k)\,\mathrm{Fl},\qquad
\zeta_1(k)=\alpha_L(k)^{-1}\upsilon(k)^n\beta_{L'}(k)^{-1},
\]
so that
\begin{equation}\label{pgx:eq-flip}
\pguK^{\mathrm c}_{v,B}(k,g)=c_\theta^{-1}\,\zeta_1(k)\,\theta\bigl(\omega^{S,L'}_{\alpha^{\mathrm r},
\beta_{L'}}(g,k)\,\mathrm{Fl}[B,\delta_v]\bigr).
\end{equation}
It remains to compute $\mathrm{Fl}[B,\delta_v]$. Write $M=\alpha\auHom(L_L,L_S)$ and
$M'=\alpha\auHom(L_S,L_{L'})$ for the self-dual lattices of the source and the target, and
$M^\dagger=\{x^\dagger:x\in M\}$. Since $(\alpha x)^\dagger=\aubar\alpha x^\dagger$ and
$\auHom(L_L,L_S)^\dagger=\auHom(L_S,L_{L'})$ (the lattices are unimodular),
$M^\dagger=\aubar\alpha\auHom(L_S,L_{L'})=\pgupt^\circ M'$, where $\pgupt^\circ=\aubar\alpha/\alpha$
acts on $\auHom(S,L')$ through the line, $x\mapsto\pgupt^\circ x$, i.e.\ as the element
$(1,\pgupt^\circ)$ of $\auU(S)\times\auU(L')$.

\emph{Local factors.} The local Heisenberg representations are irreducible and $w\mapsto
w^\dagger$ preserves $\beta$ locally, so there are unitary local intertwiners $\mathrm{Fl}_v$, unique
up to factors of modulus one. At every finite place $v$, $\mathrm{Fl}_v\psvac_{M_v}$ spans the line
on which $\psrho(l^\dagger)$, $l\in M_v$, acts by $\epsilon_{M_v}(l)$, i.e.\ (W4) the line of
$\psvac_{M^\dagger_v}$; and by (W4), $\omega_v(1,h)\psvac_{M'_v}$ spans the line of
$\psvac_{hM'_v}$ for $h\in\auU(L')(E^+_v)$ ($\omega_v$ the local Weil representation without
determinant twists). We normalise $\mathrm{Fl}_v$, for finite $v\notin S(\pcm)$, by
\begin{equation}\label{pgx:eq-norm}
\mathrm{Fl}_v\psvac_{M_v}=\omega_v(1,\pgupt_v)\,\psvac_{M'_v}\qquad(\text{unit vectors}),
\end{equation}
which is possible because there $\pgupt_v=\pgupt^\circ_v$ (Proposition~\ref{pgv:G4}), so that
$\pgupt_vM'_v=M^\dagger_v$; at $v\in S(\pcm)$ we only require $\mathrm{Fl}_v$ to be unitary (it is
made explicit below). At almost all $v$ (those outside
$S(\pcm)$ where $\alpha_v$ and $\aubar\alpha_v$ are units) $\pgupt_v$ stabilises $M'_v$, and
$\omega_v(1,\pgupt_v)\psvac_{M'_v}=\psvac_{M'_v}$ (last assertion of (W4)); so the $\mathrm{Fl}_v$ form a
restricted tensor product. At the infinite places $\mathrm{Fl}_\infty$ is the map $F\mapsto F\circ
\dagger^{-1}$ of the Fock models, and $\tau B=B\circ\dagger^{-1}$, exactly as in the proof of
Lemma~\ref{pgu:line} (that paragraph does not involve $\alpha$): $\tau$ is a signed transposition of
the variables, independent of $\vartheta$ and degree-preserving. By the uniqueness of $\mathrm{Fl}$
up to a scalar, $\mathrm{Fl}=c_0\bigotimes_v\mathrm{Fl}_v$; as $\bigotimes_v\mathrm{Fl}_v$ is unitary,
$|c_0|$ is the operator norm of $\mathrm{Fl}$ and does not depend on $\pcm$ or $\vartheta$.

\emph{Places of $S(\pcm)$.} Let $v=w\aubar w\in S(\pcm)$, $w\mid\pcm$, $\varpi$ a
uniformiser at $w$ with $\varpi^{k_w}$ the $w$-component of $\varpi_\pcm$, $q$ the order of the residue
field. By Lemma~\ref{pgx:vac}, in the models (W5) the source space $\psHS_v(\auHom(L,S))$ consists of
the Schwartz--Bruhat functions of the columns $x\in E_w^n$, with vacuum $\mathbf 1_{A}$,
$A=M_w=a_w\auO_w^n$, and the target space $\psHS_v(\auHom(S,L'))$ of the functions of the rows
$\rho\in E_w^n$, with vacuum $\mathbf 1_{A'}$, $A'=M'_w=a_w\auO_w^n$. The map $\dagger$ involves the
conjugation of $E_v$, so it carries the $w$-component of the source onto the $\aubar w$-component of
the target and the $\aubar w$-component of the source onto the $w$-component of the target; hence
$\mathrm{Fl}_v$ carries the translations of the source model into multiplications by characters of
the target model and the multiplications into translations. By the irreducibility of the
Heisenberg representation, an operator with these commutation relations is a multiple of a Fourier
transform: there are $\kappa\in E_w^\times$ and $\gamma\in\C^\times$ with
\[
\mathrm{Fl}_v\Phi(\rho)=\gamma\int_{E_w^n}\Phi(x)\,\psi_w(\kappa\rho x)\,dx ,
\]
$dx$ the measure giving $A$ the volume $1$ and $\psi_w$ the component of the additive character $\psi$.
In detail: the relation $\mathrm{Fl}_v\psrho(l)=\psrho(l^\dagger)\mathrm{Fl}_v$ says that
$\mathrm{Fl}_v$ intertwines translation by $x_0$ in the source with multiplication by the character
$\rho\mapsto[x_0^\dagger,\rho]$ of the target (the pairing $[\,,]$ of the target model between its
$\aubar w$- and $w$-components), and multiplication by a character of the source with a translation of
the target. The pairing $(x_0,\rho)\mapsto[x_0^\dagger,\rho]$ is a non-degenerate bi-additive pairing
$E_w^n\times E_w^n\to\C^\times$ of the form $\psi_w(b(x_0,\rho))$ with $b$ $E_w$-bilinear in the
coordinates (the map $\dagger$ is $E$-semilinear on matrices and the conjugation is absorbed in the
identification of $X_{\aubar w}$ with the matrices over $E_w$). Since $x^\dagger$ is the adjoint of $x$ and the pairing of the target is defined through the
adjoint for the forms of $S$ and $L'=L^-$, one has $(x_0^\dagger)^\dagger=-x_0$, the Gram matrices of
$S$ cancelling; hence $b(x_0,\rho)=\kappa\,\rho x_0$ for some $\kappa\in E_w^\times$. The displayed
transform satisfies the same commutation relations, and by the Stone--von~Neumann theorem (irreducibility
of the Heisenberg representation and Schur's lemma) an intertwiner is unique up to a scalar; hence
$\mathrm{Fl}_v$ is a multiple of it.
The $w$-component of $M^\dagger_v$ is $B'=\aubar\alpha_w\auO_w^n$ (the $w$-component of
$\aubar\alpha$ is the conjugate of $\alpha_{\aubar w}$, and that of $\auHom(L_S,L_{L'})$ is
$\auO_w^n$). As $\mathrm{Fl}_v\mathbf 1_A$ is a multiple of $\psvac_{M^\dagger_v}=\mathbf 1_{B'}$
(Lemma~\ref{pgx:vac} for $M^\dagger_v$) and the transform of $\mathbf 1_A$ is the indicator of the
annihilator $A^\perp=\{\rho:\psi_w(\kappa\rho x)=1\ \forall x\in A\}$, we get
\begin{equation}\label{pgx:eq-dual}
A^\perp=B',\qquad\text{i.e.}\quad\psi_w(\kappa a_w\aubar\alpha_w\,\cdot\,)\text{ is trivial on }\auO_w
\text{ and on no larger fractional ideal}.
\end{equation}
(In particular $\psi_w(\kappa\rho x)=1$ for $x\in A$ and $\rho\in B'$.)
Now $\pgupt_v$ has $w$-component $\pgupt_w=\aubar\alpha_w\alpha_w^{-1}\varpi^{-k_w}$
(Proposition~\ref{pgv:G4}), so $\pgupt_wA'=\aubar\alpha_w\varpi^{-k_w}\auO_w^n=\varpi^{-k_w}B'$. We fix $\gamma$ so that $\mathrm{Fl}_v$ is unitary (the transform maps $\mathbf 1_A$ to
$\mathrm{vol}(A)\mathbf 1_{B'}$, and by Plancherel a positive multiple of it is unitary); a change of
$\gamma$ by a number of modulus one only changes $c_0$ by a number of modulus one.

Let $v\in\pcMat_{\mathsf N_2\times\bullet}(R)$ and let $\tilde v\in\auO_w^n$ lift its $w$-component.
By Definition~\ref{pgx:conv}, $\delta^\natural_v$ at $w$ is the indicator of $a_w\tilde v+a_w
\varpi^{k_w}\auO_w^n$, and
\[
\mathrm{Fl}_v\delta^\natural_v(\rho)=\gamma\int_{a_w\tilde v+a_w\varpi^{k_w}\auO_w^n}\psi_w(\kappa\rho x)
\,dx=\gamma\,q^{-k_wn}\,\psi_w(\kappa a_w\rho\tilde v)\,\mathbf 1_{\varpi^{-k_w}B'}(\rho),
\]
since the annihilator of $a_w\varpi^{k_w}\auO_w^n$ is $\varpi^{-k_w}A^\perp=\varpi^{-k_w}B'$. Substitute
$\rho=\pgupt_w\rho_1$. As $\pgupt_wA'=\varpi^{-k_w}B'$, the indicator becomes $\mathbf 1_{A'}(\rho_1)$.
Write $\rho_1=a_w(\tilde v'+\varpi^{k_w}y)$ with $\tilde v'\in\auO_w^n$ (a row) and $y\in\auO_w^n$. Then
\[
\kappa a_w\pgupt_w\rho_1\tilde v=\kappa a_w^2\pgupt_w\,\tilde v'\tilde v+\kappa a_w\aubar\alpha_w
\,\varepsilon_w\,y\tilde v,\qquad\varepsilon_w=a_w/\alpha_w\in\auO_w^\times ,
\]
because $a_w^2\pgupt_w\varpi^{k_w}=a_w\aubar\alpha_w\varepsilon_w$. By \eqref{pgx:eq-dual} the second
term lies in the kernel of $\psi_w$. Put $\psi^\pcm_w(r)=\psi_w(\kappa a_w^2\pgupt_w\tilde r)$ for
$r\in\auO_w/\varpi^{k_w}$ with a lift $\tilde r$; since $\kappa a_w^2\pgupt_w=\varpi^{-k_w}\cdot\kappa a_w
\aubar\alpha_w\varepsilon_w$ and $\psi_w(\kappa a_w\aubar\alpha_w\varepsilon_w\,\cdot\,)$ is trivial on
$\auO_w$ and on no larger fractional ideal (\eqref{pgx:eq-dual}; $\varepsilon_w$ is a unit),
$\psi^\pcm_w$ is a well defined character of $\auO_w/\varpi^{k_w}$, primitive (non-trivial on the ideal
$\varpi^{k_w-1}\auO_w/\varpi^{k_w}$). Also the value does not depend on the lift $\tilde v$. Hence
\[
\mathrm{Fl}_v\delta^\natural_v(\pgupt_w\rho_1)=\gamma\,q^{-k_wn}\sum_{v'}\psi^\pcm_w(v'_wv_w)\,
\delta^\natural_{v'}(\rho_1),
\]
the sum over the rows $v'$ modulo $\varpi^{k_w}$, with $\delta^\natural_{v'}$ the target datum of
Definition~\ref{pgx:conv} (the indicator of $a_w\tilde v'+a_w\varpi^{k_w}\auO_w^n$). By (W5), for
$h\in\auU(L')(E^+_v)$ with $w$-component $h_w$, $\omega_v(1,h)F(\rho)=\eta\,|h_w|_w^{-n/2}F(h_w^{-1}
\rho)$ with $|\eta|=1$ (the factor makes the operator unitary). Therefore
\[
\mathrm{Fl}_v\delta^\natural_v=\gamma_v\,\omega_v(1,\pgupt_v)\Bigl[\sum_{v'}\psi^\pcm_w(v'_wv_w)\,
\delta^\natural_{v'}\Bigr]
\]
for a constant $\gamma_v$ independent of the rows $v$ and $v'$. Its modulus follows from the norms:
$\|\delta^\natural_v\|^2=q^{-k_wn}$ (measure giving $A$ volume $1$), and the bracket has squared norm
$\sum_{v'}\|\delta^\natural_{v'}\|^2=q^{k_wn}\cdot q^{-k_wn}=1$ (the $\delta^\natural_{v'}$ have disjoint
supports; measure giving $A'$ volume $1$); $\mathrm{Fl}_v$ and $\omega_v(1,\pgupt_v)$ are unitary. So
$|\gamma_v|=q^{-k_wn/2}$. Put $\psi_\pcm=\prod_{w\mid\pcm}\psi^\pcm_w$ on $R=\prod_w\auO_w/\varpi^{k_w}$; it
is primitive, and depends only on $\pcm$ (through $\kappa$, $a_w$, $\alpha$, $\varpi_\pcm$).

\emph{Assembling.} The datum $[B,\delta_v]$ is $B\otimes\bigotimes_{w\mid\pcm}\delta^\natural_{v,w}
\otimes\psvac^\pcm$. Outside $S(\pcm)$, by \eqref{pgx:eq-norm}, $\mathrm{Fl}_v$ maps the unit vector
$\psvac_{M_v}$ to $\omega_v(1,\pgupt_v)\psvac_{M'_v}$. At the infinite places $\mathrm{Fl}_\infty B=
\tau B$. At $S(\pcm)$ we have just computed the image. Hence
\[
\mathrm{Fl}[B,\delta_v]=c\;\omega^{S,L'}(1,\pgupt)\Bigl[\tau B,\sum_{v'\in\pcMat_{\bullet\times
\mathsf N_2}(R)}\psi_\pcm(v'v)\,\delta_{v'}\Bigr],\qquad |c|=|c_0|\prod_{w\mid\pcm}q_w^{-k_wn/2}=
|c_0|\,(\auNm\pcm)^{-n/2},
\]
with $\omega^{S,L'}$ the Weil representation without determinant twists ($\pgupt$ has trivial
archimedean component). Inserting this into \eqref{pgx:eq-flip} and using
$\omega^{S,L'}_{\alpha^{\mathrm r},\beta_{L'}}(g,k)\,\omega^{S,L'}(1,\pgupt)=\beta_{L'}(\pgupt)^{-1}
\omega^{S,L'}_{\alpha^{\mathrm r},\beta_{L'}}(g,k\pgupt)$ (the twisted and untwisted operators differ by
the character $\beta_{L'}(\det\cdot)$ on the target group), we obtain the identity of the statement
with $c_{\mathrm{line}}=c/c_\theta$, $\zeta(k)=\zeta_1(k)\beta_{L'}(\pgupt)^{-1}$, and
$C_0=|c_0/c_\theta|^2$, which is independent of $\vartheta$. The values of the characters at $k$ and at
$\pgupt$ have modulus one, and so does $\zeta(k)$. The last assertion is the case $k=\pgupt^j$.
\end{proof}

\begin{remark}\label{pgx:rem-line}
(a) Lemma~\ref{pgx:line} has the statement of Lemma~\ref{pgu:line}, with two differences: the data are
those of Definition~\ref{pgx:conv} (which agree with those of \S\ref{pc:ss-data} where $\alpha$ is a
unit at the primes of $\pcm$), and the element is \emph{exactly} the element $\pgupt_\pcm$ of
Definition~\ref{pgv:def-period}. In both lemmas
the characters of $k$ enter only through the factor $\zeta(k)$ of modulus one.

(b) The role of the lattices at $v\in S(\pcm)$: when $\alpha$ is a unit at $w$, the vacuum of the
source and of the target is the indicator of $\auO_w^n$, the self-dual measure gives $\auO_w^n$
volume $1$ and the dilation needed is by $\varpi^{k_w}$. In general the vacuum lattices are
$a_w\auO_w^n$ (Lemma~\ref{pgx:vac}), the image of
the source vacuum is the indicator of $\aubar\alpha_w\auO_w^n$, and the dilation needed is by
$\pgupt_w=\aubar\alpha_w/(\alpha_w\varpi^{k_w})$, which carries $a_w\auO_w^n$ to
$\varpi^{-k_w}\aubar\alpha_w\auO_w^n$. With the data of \S\ref{pc:ss-data} (support on the integral
matrices) the computation also closes, but with the dilation by $\varpi^{-k_w}\alpha_w\aubar\alpha_w$
instead of $\pgupt_w$; this element differs from $\pgupt_w$ by $\alpha_w^2$, which is not a unit at
the primes of Proposition~\ref{pgv:alpha}(b), and it would not match the diagonal period of
Definition~\ref{pgv:def-period}. For this reason the data are placed on the lattices $a_w\auO_w^n$
and the period is that of Definition~\ref{pgv:def-period}.
\end{remark}

\subsubsection{The line reversal in \S\ref{pgu:sec} for general $\alpha$}\label{pgx:ss-later}

The line reversal and its constant enter \S\ref{pgu:sec} at the following places. In each, Lemma~\ref{pgx:line}
is used, with the data of Definition~\ref{pgx:conv}; it contains Lemma~\ref{pgu:line} as the case where
$\alpha$ is a unit at the primes of $\pcm$.
\begin{itemize}
\item \emph{(G4) and the remark after Proposition~\ref{pgv:G4}.} (G4) asks for the element of the line
reversal; Lemma~\ref{pgx:line} produces exactly $\pgupt_\pcm$ of Definition~\ref{pgv:def-period}, so
(G4) holds as stated.
\item \emph{Corollary~\ref{pgu:split}.} Holds with Definition~\ref{pgx:conv}: the factorisation
of the point data and of the norms at $S(\pcm)$ holds because $a_w$ is a scalar (end of
\S\ref{pgx:ss-conv}); the rest of the proof does not involve the primes of $\pcm$.
\item \emph{Lemma~\ref{pgu:see1}.} Holds with Definition~\ref{pgx:conv}. It uses the line kernel only through the function
$M_{x,\mathbf S}(b)=c_1\pguK^{\mathrm c}_{v(\mathbf S),B}(\pgupt^j,b)$, a continuous function
of $b$, left invariant under the rational points and of moderate growth by (R2) of
\S\ref{pcx:ss-conv}, and through Corollary~\ref{pgu:split}; neither the identity of the line
reversal nor $c_{\mathrm{line}}$ occurs. The point $\pgupt^j$ is that of (G4).
\item \emph{Lemma~\ref{pgu:see2}.} Holds, for the same reason, with
$N_{x,\mathbf S}(b)=c_2\pguK^{\mathrm r}_{v'(\mathbf S),B'}(b,\pgupt^{j+1})$ and $j'(a)=a\oplus
\pgupt^{j+1}$.
\item \emph{Proposition~\ref{pgu:turned} and Lemma~\ref{pgu:sumturn}.} They are statements
about functions on matrices over $R$ for an arbitrary additive character $\psi$ of $R$, and do not
involve kernels.
\item \emph{Theorem~\ref{pgu:turndiag}.} Statement and proof hold with $c'_{\mathrm{line}}$ and
$\psi_\pcm$ those of Lemma~\ref{pgx:line}. The proof uses: the identity at $k=\pgupt^j$ (the last
assertion of Lemma~\ref{pgx:line}); the primitivity of
$\psi_\pcm$ is not even needed (Proposition~\ref{pgu:turned} allows any $\psi$); the modulus
$|c'_{\mathrm{line}}|^2=C_0(\auNm\pcm)^{-n}$, $n=r+(2j+1)t$; and that
$\tau$ is a bijection of polynomials independent of $\vartheta$ (in the transfer step, $B$ runs over
all polynomials, so $\tau B$ does too).
\end{itemize}
So Theorem~\ref{pgu:turndiag} $=$ Proposition~\ref{pgr:turndiag} holds without any condition on
$\alpha$ at the primes of $\pcm$.

\subsubsection{Other uses of the integral matrices at the primes of a modulus}\label{pgx:ss-other}

Definition~\ref{pgx:conv} changes the data only at primes $w\mid\pcm$ with $\alpha_w$ not a unit; by
Proposition~\ref{pgv:alpha} these lie above $q_1$, and they occur only for moduli containing primes of
$Q_\vartheta$ above $q_1$ (the moduli of the grids, and $\pcm_\vartheta$ of \S\ref{pc:ss-data}). We
list the places where the proof of the Covolume Theorem uses the data at such primes.
\begin{enumerate}
\item \S\ref{pc:ss-data}, definition of $\Psi^\natural$ and of the norm: this is
Definition~\ref{pgx:conv}. With data supported on the integral matrices the identity $\mathbf
1^\natural=$ vacuum, used whenever a datum with $\Psi=\mathbf 1$ is identified with the vacuum, would be
false at the primes of Proposition~\ref{pgv:alpha}(b) (Lemma~\ref{pgx:vac}).
\item Lemma~\ref{pc:level}: holds verbatim with Definition~\ref{pgx:conv} (\S\ref{pgx:ss-conv}).
\item Proposition~\ref{pc:rallis}, step ``At $\pcm$'': uses only the support condition for the right
action of $\auGL_N(E_w)$, Lemma~\ref{pc:level} and $\|\Psi^\natural\|^2=\|\Psi\|^2/\#\pcMat(\pcR_\pcm)$;
all three hold with Definition~\ref{pgx:conv} (the lattice $a_w\pcMat(\auO_w)$ is stable under the
right action exactly when $\pcMat(\auO_w)$ is, $a_w$ being a scalar).
\item Lemma~\ref{pc:clean}: uses only that the component at one split place is supported on
injective matrices, which is independent of the lattice.
\item Proposition~\ref{ptv:rev}: its proof at $w\mid\pcm$ describes the data as
the vectors fixed by $\psrho(\Lambda)$, $\Lambda=\varpi\pcMat(\auO_w)\oplus\pcMat(\auO_{\aubar w})$, and
asserts that $\pcMat(\auO_{\aubar w})$ is the dual of $\pcMat(\auO_w)$. By Lemma~\ref{pgx:dict}(b) the dual
of $\pcMat(\auO_w)$ is $\aubar{a_w}a_{\aubar w}\pcMat(\auO_{\aubar w})$, so that sentence is correct
only when $\alpha$ is a unit at $w$ and $\aubar w$. With Definition~\ref{pgx:conv} the proof applies
verbatim after replacing $\pcMat(\auO_w)$ by $M_w=a_w\pcMat(\auO_w)$ and $\pcMat(\auO_{\aubar w})$ by its
dual $M_{\aubar w}$: the argument uses only that these two lattices are mutually dual and stable under
$\auGL_N(\auO_w)\times\auGL_{N'}(\auO_w)$, and the scalar $\zeta_\pcm$ there is constructed from them.
(For the squarefree
moduli of \S\ref{pc:ss-data} this concerns the primes of $Q_\vartheta$ above $q_1$.)
\item The exchange identity (\S\S\ref{pcx:sec}--\ref{pcn:sec}). Its proof works at $w\in S$ in the
coordinates $\tilde x$ of Lemma~\ref{pgx:dict} ($a_w^{-1}x$ on the $w$-component, $a_{\aubar w}^{-1}x$ on
the $\aubar w$-component), in which the vacuum lattice $M_w$ and its dual $M_{\aubar w}$ both become the
integral matrices and $[\,,]$ becomes $\lambda_w(\operatorname{tr}({}^t\tilde y\tilde x))^{\pm1}$
(Lemma~\ref{pgx:dict}(b)). Since $a_w$ is a scalar, this
change of coordinates commutes with the actions $x\mapsto k'^{-1}xk$ of the model (W5) and preserves
the measure normalisations of Definition~\ref{pgx:conv}; it carries the datum $\Psi^\natural$ of
Definition~\ref{pgx:conv} to the function $\Psi(\tilde x\bmod\pcm)$ on the integral matrices, and
matched pairs to matched pairs. Hence Lemma~\ref{pcx:mod} (``$[x,\Phi]=[x,\Phi']_S$'') holds with the
data of Definition~\ref{pgx:conv}, the right side being read in these coordinates; this is stated
after that lemma. In the proof of Lemma~\ref{pct:factor} (Step~4 of the proof of
Theorem~\ref{pct:plane}), the idele $\alpha_0$ by which the Weil data scale $\psi$ is the idele
$\alpha$ of {\cite[\S5.1.1]{P49}} (the scaling that makes the standard lattices self-dual), and in the
coordinates $\tilde x$ its component at $w\in S$ becomes the unit $a_w^{-1}(\alpha_0)_w$; this is the
``unit component of $\alpha_0$'' that appears there, and the additive character to which the Gauss sum
Lemma~\ref{pct:gauss} is applied has conductor $\auO_w$ in these coordinates, by the duality
\eqref{pgx:eq-dual} of the two lattices. The map $w\mapsto w^\dagger$ of \S\ref{pcx:ss-flip} carries
the self-dual lattice $\alpha\auHom$ to $\aubar\alpha\auHom$, which differs from $\alpha\auHom$ at the places
above $q_1$; the comparison is made by the central translation of Lemma~\ref{pgx:dict}(e),(f), as
explained in \S\ref{pcx:ss-flip}. These statements hold with the data of Definition~\ref{pgx:conv}.
\item \S\ref{pgu:ss-line}, Lemma~\ref{pgu:line}: the general case is Lemma~\ref{pgx:line}.
\item (G4), the remark after Proposition~\ref{pgv:G4}, and the discussion after
Proposition~\ref{pgv:alpha}: the data are those of Definition~\ref{pgx:conv}, and the diagonal period is
that of Definition~\ref{pgv:def-period}.
\end{enumerate}
No other step of the proof of the Covolume Theorem uses the integral matrices at a prime of a modulus as a vacuum
lattice.

\subsubsection{The identification (G1) in the seesaw steps}\label{pgx:ss-G1}

\begin{proposition}\label{pgx:G1}
The proofs of Lemmas~\ref{pgu:see1} and~\ref{pgu:see2}, of Corollary~\ref{pgu:split} and of
Theorem~\ref{pgu:turndiag} use the bijection $\sigma$ of Proposition~\ref{pgv:G1} only through the
induced isometry $S\to S_1$ of spaces with frames: the identification of $\auU(S)$ with
$\auU(S_1)$ (adelic points, rational points, the groups $K$ and the reductions modulo $\pcm$), of
the frames $\mathsf N_2$, and of $\auHom(V',S)$ with $\auHom(V',S_1)$ and $\auHom(S,V)$ with
$\auHom(S_1,V)$ (with lattices, Weil representations, theta functionals and Fock models). No type of
$S$ is identified with a type of $S_1$.
\end{proposition}
\begin{proof}
\emph{Corollary~\ref{pgu:split}.} The isometry $\sigma_*\colon S\to S_1$ is $E$-linear, carries the frame
to the frame and $L_S$ to $L_{S_1}$, and preserves the signs; so it induces isometries
$\auHom(V',S)\cong\auHom(V',S_1)$, $\auHom(S,V)\cong\auHom(S_1,V)$ carrying the self-dual lattices
$\alpha\auHom$ to each other, the complex structures at the infinite places to each other, and
intertwining the Weil representations along $\auU(S)\cong\auU(S_1)$. By (W1) the theta functional is
carried to a non-zero multiple of the theta functional (the isometry is rational), and the multiple is
absorbed in $c_1,c_2$. The datum is carried to the datum: the matrices are written in the frames, which
$\sigma$ identifies, and the polynomials $A$, $A'$ are transported. The characters of (G2) are characters
of $\det$, and $\det(\sigma_*g\sigma_*^{-1})=\det g$. Nothing else is used.

\emph{Lemma~\ref{pgu:see1}.} In the application of Lemma~\ref{pgu:seesaw} with $(S_0,T)=(V',S_1)$ and
$(T',S')=(\widetilde V',S)$: the form $f_{x,\mathbf S}=\varphi_{\mathbf S}$ lives on $S'=S$ and must be
of the conjugate of the target type of $\Phi'=P_{e_2}$, which involves only the type at $S$; the form
$y_{x,\mathbf S}$ lives on $V'$ and has the source type of $\Phi=P_{(V',S_1)}$; the lifts
$\Theta(\Phi,\Psi)y$ are forms on $S_1$ of the type at $S_1$ (Theorem~\ref{pgv:G3}(d)), and they enter
only through the pairing $\langle f,F\rangle_{M,j',r'}=\int_{[\auU(S)]}M(b)f(b)F(j'(b)r')\,db$, in which
$F$ is merely a bounded continuous function on $[\auU(S_1)]$ composed with the group isomorphism
$j'$ of (G1); the moderate-weight condition on $(M,j')$ requires $j'$ to be continuous, compatible with
rational points and to satisfy $\mathrm{ht}(j'(x))\le C\,\mathrm{ht}(x)^C$, and $j'$, induced by a
rational isometry of spaces with frames, preserves the heights (computed in the frames). Condition (b)
is Corollary~\ref{pgu:split}. The rapid decrease of $\varphi_{\mathbf S}$ and the boundedness of the
lifts are properties of the types at $S$ and at $S_1$ separately (Theorem~\ref{pgv:G3}(c),(d)).

\emph{Lemma~\ref{pgu:see2}.} Symmetric: $(S_0,T)=(S,\widetilde V)$, $(T',S')=(S_1,V)$; $y=\varphi_{\mathbf
S}$ lives on $S$ and has the source type of $P_{(S,\widetilde V)}$, the transposed lift
$\Theta^{\mathsf t}(\Phi',\Psi')f$ lives on $S_1$ with the type at $S_1$, and they are paired by
$\int_{[\auU(S)]}N(b)\,y(b)\,F(j(b)r)\,db$ with $j$ the group isomorphism of (G1). Again only the groups
and Corollary~\ref{pgu:split} are used.

\emph{Theorem~\ref{pgu:turndiag}.} The functionals $l_{\mathbf S,v'}$ are defined on functions on
$[\auU(S_1)]$ by integration over $[\auU(S)]$ through the group isomorphism, and the square law
(Proposition~\ref{pgr:square}) is applied in the first grid to forms at $S_1$ of the type at $S_1$; the
forms at $S$ (of the type at $S$) enter only through $\varphi_{\mathbf S}$ inside $l_{\mathbf S,v'}$.
\end{proof}

\noindent Thus the proofs of \S\ref{pgu:sec} use (G1) only in the form described in the remark after
Proposition~\ref{pgv:G1}.




% ===== C-grid7.tex =====
% proofs/C-grid7.tex
% Continuation of proofs/C-grid.tex ... C-grid6.tex: the inputs (I) and (P) of Theorem pgw:first.
% Labels start with pgy:. Conventions and macros: those of proofs/C-grid5.tex and C-determinant.tex.
\providecommand{\auGL}{\mathrm{GL}}
\providecommand{\auU}{\mathrm U}
\providecommand{\auO}{\mathcal O}
\providecommand{\auHom}{\operatorname{Hom}}
\providecommand{\aubar}[1]{\overline{#1}}
\providecommand{\pcm}{\mathfrak m}
\providecommand{\pcb}{\mathfrak b}
\providecommand{\pcMat}{\mathrm M}
\providecommand{\pcF}{\mathcal F}
\providecommand{\pcR}{\mathsf R}
\providecommand{\psvac}{\mathbf v}
\providecommand{\pgrY}{\widetilde Y}
\providecommand{\pcnot}[1]{\par\smallskip{\raggedright\noindent\textsc{Not proved here.} [#1]\par}\smallskip}
\providecommand{\pgvwt}{\operatorname{wt}}
\providecommand{\pgyG}{\mathsf G}
\providecommand{\pcRal}{\mathcal R}
\providecommand{\pcmc}{c}
\providecommand{\pcW}{\mathcal W}

\subsection{Two inputs of the first path: injectivity (I) and change of modulus (P)}\label{pgy:sec}

This section continues \S\S\ref{pgv:sec}, \ref{pgw:sec}, \ref{pgx:sec} and uses their notation.
Theorem~\ref{pgw:first} was proved from three inputs (I), (V), (P) stated in \S\ref{pgw:ss-first}.
\S\ref{pgy:ss-tri} constructs the triangle of the second grid (in the sense of
Definition~\ref{pc:triangle}), verifies (T1)--(T3) for it, and deduces (I) from
Proposition~\ref{pc:tri-inj} (Theorem~\ref{pgy:I}). \S\ref{pgy:ss-P} proves (P) for the data of
Definition~\ref{pgx:conv}, with no assumption on the idele $\alpha$ (Proposition~\ref{pgy:P}).
Input (V) is Theorem~\ref{pgq:V}.

Throughout, $r,k,t=5r,\Delta,\kappa_l,\mu^0_l$ are as in \S\ref{pgv:sec}; $\mu^0_j=\mu^\flat\nu_1^{\,j}$
with $\nu_1=\upsilon^{-\Delta}\chi_M$ as recalled in \S\ref{pc:ss-triangle}, so $\pgvwt\mu^0_j=-\kappa_j$
and the tame parts of $\nu_1$ at the chosen places have order $>2m$; $\pcb=(\pi_M)^b$.
Concerning hypotheses: the archimedean statement called (V) in Lemma~\ref{ptr:const} and
Theorem~\ref{pgv:G3} ({\cite[Lemma~3.6]{P49}}, proved in {\cite[Theorem~3.29]{P49}}) is unrelated to the input
(V) of \S\ref{pgw:ss-first}; to avoid confusion we write it $(\mathrm V_\infty)$. Like
Theorem~\ref{pgv:G3}, everything below assumes $(\mathrm V_\infty)$ and (D1)--(D5).

\subsubsection{The triangle of the second grid and the input (I)}\label{pgy:ss-tri}

\paragraph*{Spaces, parameters, types}
For a stem coordinate $\varrho$ ($0\le\varrho\le r$) put $a(\varrho)=\max(\varrho+1-k,0)$, as in
\eqref{pgv:eq-sign}, and for integers $i,n\ge0$
\[
E_{i,n}(a)=i+a+\min\bigl((a-i)^+,n\bigr)\qquad\text{(the exponent of \eqref{ptr:eq-E})}.
\]
For $i+n\le m$ let $\pgyG_{i,n}$ be the following space with a frame consisting of a stem of $r+1$
coordinates and $n$ blocks of $t-1$ coordinates, numbered $i+1,\dots,i+n$; its rank is
$N_n=r+1+n(t-1)$, which is even.
\begin{itemize}
\item $\pgyG_{i,0}$ is $\pgrY^0_0=Y_0\oplus E\tilde e$ (Definition~\ref{pgv:def-spaces}) with its
form multiplied by $(-1)^i$;
\item for $n\ge1$, $\pgyG_{i,n}=E^{N_n}$ with the standard lattice and the diagonal form with entry
$(-1)^{E_{i,n}(a(\varrho))}$ at the stem coordinate $\varrho$ and $+1$ on the blocks.
\end{itemize}
The signs of $\pgyG_{i,0}$ are $(-1)^{i+a(\varrho)}=(-1)^{E_{i,0}(a(\varrho))}$, since those of
$\pgrY^0_0$ are $(-1)^{a(\varrho)}$ by \eqref{pgv:eq-sign}. For $i=0$, $E_{0,n}(a)=a+\min(a,n)$, so by
\eqref{pgv:eq-sign} and Definition~\ref{pgv:def-spaces}
\begin{equation}\label{pgy:eq-node}
\pgyG_{0,n}=\pgrY^n_0\quad\text{(the node $(1,0,n)$ of the second grid)},\qquad0\le n\le m .
\end{equation}
At $\tau\in\auPl$ let $\lambda^{(i,n)}_\tau$ be the vector with entries
$\varsigma'_{\tau,\varrho}+\frac12$ at the stem coordinate $\varrho$ ($\varsigma'$ as in
\eqref{pgv:eq-param}) and $-\kappa_{i+l}(\tau)+\frac{t-2}2-\beta$ at the position $0\le\beta<t-1$ of
the block $i+l$. For $i=0$ this is $\lambda_{(1,0,n),\tau}$ of \eqref{pgv:eq-param} ($z=1$, $s=0$). Let
$f_{i,n}$ be the type of $\auU(\pgyG_{i,n})_\infty$ whose Harish-Chandra parameter is
$\lambda^{(i,n)}$ (with the signs of $\pgyG_{i,n}$); for $i=0$, $f_{0,n}=\overline{e_{(1,0,n)}}$, since
$e_{(1,0,n)}$ has parameter $-\lambda_{(1,0,n)}$. For $n=0$ the parameter does not depend on $i$,
and $\pgyG_{i,0},\pgyG_{0,0}$ have the same group; as in Lemma~\ref{ptr:cons}(b) the types
$f_{i,0}$ and $f_{0,0}$ coincide (a change of all signs does not change the compact roots). Put
\[
\pcF_{i,n}=\pcF(\pgyG_{i,n},f_{i,n}),\qquad V=\pcF_{0,0}^{J(\pcb)},\qquad j_i\colon V\hookrightarrow
\pcF_{i,0}=\pcF_{0,0}
\]
the inclusion, $J(\pcb)=J_{\pgrY^0_0}(\pcb)$. So $V$ is the space of forms of type
$\overline{e_{(1,0,0)}}$ on $\pgrY^0_0$ fixed by $J(\pcb)$ of the input (I).

\paragraph*{Lifts and descents}
For $i+n<m$ let
\[
\alpha_{i,n}=\mu^0_{i+n+1}\upsilon^{N_{n+1}/2},\quad\beta_{i,n}=\mu^0_{i+n+1}\upsilon^{N_n/2},\qquad
a_{i,n}=\Theta_{\alpha_{i,n},\beta_{i,n}}(P_{i,n},\delta_\iota)\colon\pcF_{i,n}\to\pcF_{i,n+1},
\]
with the modulus $\pcb$, the point datum at the inclusion matrix $\iota$ (stem and blocks
$i+1,\dots,i+n$ of $\pgyG_{i,n}$ to the same coordinates of $\pgyG_{i,n+1}$) and the joint harmonic
$P_{i,n}$ of the \emph{ascending pair} $(\pgyG_{i,n},\pgyG_{i,n+1})$, rows with
$\lambda=-\lambda^{(i,n)}$, columns with $\lambda'=\lambda^{(i,n+1)}$. Let $w_n=\#\auGL_{N_n}(\auO_E/\pcb)$.
For $i=0$ these are the characters of the up edge $e^1_n=((1,0,n),(1,0,n+1))$ of
Definition~\ref{pgv:def-edges} (horizontal, index $n$, $\nu=\mu^0_{n+1}$, powers
$\lceil N/2\rceil=N/2$), the joint harmonic is $P_{e^1_n}$ (both are joint harmonics for the same
parameters, unique up to a scalar, Theorem~\ref{pgv:G3}(a)), and so, after normalising
$P_{0,n}=P_{e^1_n}$,
\begin{equation}\label{pgy:eq-chain}
a_{0,n}=a^\flat_n,\qquad X_{0,m}=w_{m-1}\cdots w_0\;a^\flat_{m-1}\circ\dots\circ a^\flat_0=
w_{m-1}\cdots w_0\;a^\flat ,
\end{equation}
with $a^\flat_n$, $a^\flat$ of \S\ref{pgw:ss-first} and $X_{i,n}$ the chains of
Definition~\ref{pc:triangle}.

For $i+n+1\le m$ the \emph{descending pair} is $(\pgyG_{i+1,n},\pgyG_{i,n+1})$, rows with
$\lambda=\lambda^{(i+1,n)}$, columns with $\lambda'=-\lambda^{(i,n+1)}$, and characters
\begin{equation}\label{pgy:eq-desc}
\alpha^\downarrow_{i,n}=(\mu^0_{i+1})^{-1}\upsilon^{N_{n+1}/2+c_{i}},\qquad
\beta^\downarrow_{i,n}=(\mu^0_{i+1})^{-1}\upsilon^{N_n/2+c_i}
\end{equation}
for an integer $c_i$ fixed in Lemma~\ref{pgy:chars}; their weights are $\kappa_{i+1}+N_{n+1}/2+c_i$ and
$\kappa_{i+1}+N_n/2+c_i$. The descent $d_{i,n}\colon\pcF_{i,n+1}\to\pcF_{i+1,n}$ is the transposed lift
along the descending pair with its joint harmonic and a finite datum $\varepsilon_{i,n}$ modulo
$\pcb$: $\varepsilon_{i,0}$ is the datum for which $d_{i,0}=a_{i,0}^*$ (Lemma~\ref{pgy:T2}), and
$\varepsilon_{i,n+1}=\varepsilon_{i,n}^+$ (the extension of Lemma~\ref{pc:next}), as in step (iv) of
\S\ref{pc:ss-exchange}. The triangle is over the one-point set $I=\{1\}$ (nothing depends on a twist).

\paragraph*{Matching}
\begin{lemma}\label{pgy:match}
Let $c_i=0$ in \eqref{pgy:eq-desc}. At every $\tau$, every ascending pair satisfies the block rule
of excess $t-1$ \eqref{pgv:eq-rule}, with $\tau_0=\kappa_{i+n+1}$, $\tau'_0=-\kappa_{i+n+1}$, $c$ the
identity on the stem and the blocks of $\pgyG_{i,n}$, and $\mathsf B$ the block $i+n+1$; every
descending pair satisfies it with $\tau_0=-\kappa_{i+1}$, $\tau'_0=\kappa_{i+1}$, $c$ the identity on
the stem and on the blocks $i+2,\dots,i+n+1$, and $\mathsf B$ the block $i+1$ of $\pgyG_{i,n+1}$.
Conversely these are the only weights of the descent characters for which the descending pair
satisfies the block rule; hence $c_i=0$ is forced.
\end{lemma}
\begin{proof}
All ranks are even. \emph{Constants} (Lemma~\ref{ptr:const}, valid for any ranks and characters):
ascending, $\tau_0=\frac{N_{n+1}}2-\pgvwt\alpha_{i,n}=\kappa_{i+n+1}$, $\tau'_0=\pgvwt\beta_{i,n}-
\frac{N_n}2=-\kappa_{i+n+1}$; descending, $\tau_0=\frac{N_{n+1}}2-\pgvwt\alpha^\downarrow_{i,n}=
-\kappa_{i+1}-c_i$, $\tau'_0=\pgvwt\beta^\downarrow_{i,n}-\frac{N_n}2=\kappa_{i+1}+c_i$. In both cases
$\tau_0+\tau'_0=0$.

\emph{Old coordinates.} The stem entries of $\lambda^{(i,n)}$ depend neither on $i$ nor on $n$; the
block $i+l$ carries the same entries in every $\pgyG_{i',n'}$ that contains it. Ascending: the
condition $\lambda'_{c(\varrho)}=-\lambda_\varrho$ reads $\lambda^{(i,n+1)}=\lambda^{(i,n)}$ on the common
coordinates: true. Descending: it reads $-\lambda^{(i,n+1)}=-\lambda^{(i+1,n)}$ on the common
coordinates (stem, blocks $i+2,\dots,i+n+1$): true.

\emph{New block.} Ascending: block $i+n+1$ of $\pgyG_{i,n+1}$ carries $-\kappa_{i+n+1}+\frac{(t-1)-1}2-
\beta=\tau'_0+\frac{d-1}2-\beta$, $d=t-1$. Descending: block $i+1$ of $\pgyG_{i,n+1}$ carries
$-\kappa_{i+1}+\frac{t-2}2-\beta$; the entries of $\lambda'$ there are $\kappa_{i+1}-\frac{t-2}2+\beta$,
the string $\{\tau'_0+\frac{d-1}2-\beta'\}$ exactly when $c_i=0$. The signs on blocks are $+1$.
This gives the converse: the rule determines $\tau'_0$ (centre of the new string) and
$\tau_0+\tau'_0$, hence the weights.

\emph{Distances and signs, ascending.} On a common coordinate $\lambda'_{c(\varrho)}-\tau'_0=
\lambda^{(i,n)}_\varrho+\kappa_{i+n+1}$. Blocks $i+l$, $l\le n$: $\kappa_{i+n+1}-\kappa_{i+l}+\frac{t-2}2-
\beta\ge\Delta-\frac t2>\frac d2$, positive. Stem, $\varrho\le k-2$ (then $a(\varrho)=0$):
$\frac k2-\frac12-\varrho+\kappa_{i+n+1}\ge\frac32\Delta-r>0$. Stem, $\varrho\ge k-1$, $a=a(\varrho)=
\varrho+1-k$: $-\kappa-(a+1)\Delta+\frac12+\kappa+(i+n+\frac32)\Delta=(i+n+\frac12-a)\Delta+\frac12$, of
absolute value $\ge\frac\Delta2-\frac12>\frac d2$, positive iff $a\le i+n$. So the rule asks: a stem
coordinate keeps its sign iff $a\le i+n$ (for $\varrho\le k-2$, $a=0\le i+n$). By the definition of
$E_{i,n}$: for $a\le i$, $E_{i,n}(a)=E_{i,n+1}(a)=i+a$; for $i<a\le i+n$ both are $2a$; for $a>i+n$,
$E_{i,n+1}(a)=E_{i,n}(a)+1$. This is the rule.

\emph{Distances and signs, descending} ($c_i=0$). On a common coordinate $\lambda'_{c(\varrho)}-
\tau'_0=-\lambda^{(i,n+1)}_\varrho-\kappa_{i+1}$. Blocks $i+l$, $l\ge2$: $\kappa_{i+l}-\kappa_{i+1}-
\frac{t-2}2+\beta\ge\Delta-\frac t2$, positive. Stem, $\varrho\le k-2$: $-\frac k2+\frac12+\varrho-\kappa_{i+1}
<-\frac32\Delta+r$, negative. Stem, $\varrho\ge k-1$: $\kappa+(a+1)\Delta-\frac12-\kappa-(i+\frac32)\Delta=
(a-i-\frac12)\Delta-\frac12$, of absolute value $>\frac d2$, positive iff $a\ge i+1$. So the rule asks: a
stem coordinate changes sign iff $a\le i$. Compare $E_{i+1,n}$ (rows) with $E_{i,n+1}$ (columns): for
$a\le i$ they are $i+1+a$ and $i+a$; for $a=i+1$ both are $2i+2$; for $i+1<a\le i+n+1$ both are $2a$;
for $a>i+n+1$ both are $i+a+n+1$. This is the rule.
\end{proof}

\begin{lemma}[Gap and regularity]\label{pgy:gap}
Every $\lambda^{(i,n)}_\tau$ consists of $N_n$ distinct numbers in $\frac{N_n-1}2+\Z$, two entries at
coordinates of opposite signs differ by at least $\frac12(\Delta-t)-1$, and $N_n\le2n_m$. Hence the
conclusions of {\cite[Theorem~3.17]{P49}}, {\cite[Lemma~3.25]{P49}}, {\cite[Theorem~3.2]{P49}} and
{\cite[Corollary~3.3]{P49}} hold for the types $f_{i,n}$ and $\overline{f_{i,n}}$.
\end{lemma}
\begin{proof}
The entries are in $\frac12+\Z$ ($\varsigma'$ is integral, $k$ even; blocks of even length $t-1$), and
$N_n$ is even. Group the entries into: the small stem entries ($\varrho\le k-2$, all with
$a(\varrho)=0$, hence all of the sign $(-1)^{i}$, of absolute value $\le\frac k2$), each large stem entry
$-\kappa-(a+1)\Delta+\frac12$ ($0\le a\le m$), and each block (entries $-\kappa_{i+l}+\frac{t-2}2-\beta$,
$0\le\beta\le t-2$: centre $-\kappa_{i+l}$, half-length $\frac{t-2}2<\frac t2$). Within a group the sign is constant and the entries are
distinct. Between groups: small stem to anything else $\ge\kappa-\frac k2-\frac t2$; two large stem
entries $\ge\Delta$; a large stem entry to a block centre $-\kappa-(i+l+\frac12)\Delta$ at distance
$|a+1-i-l-\frac12|\Delta\ge\frac\Delta2$, so to the block entries $\ge\frac\Delta2-\frac t2-\frac12$; two
blocks $\ge\Delta-t$. All these exceed $\frac12(\Delta-t)-1$ (and $1$), giving distinctness and the gap.
$N_n\le r+1+m(t-1)\le2n_m$ as in Lemma~\ref{pgv:gap}. The consequences follow exactly as in the proof
of Lemma~\ref{pgv:gap} (the cited statements use the parameter only through regularity, gap and rank;
the gap bound $\frac12(\Delta-t)-1$ is the one used there and in Lemma~\ref{ptr:gap}).
\end{proof}

\begin{proposition}[Edges]\label{pgy:edge}
For every ascending and every descending pair (with $c_i=0$), the conclusions (a)--(d) of
Theorem~\ref{ptr:edge} hold, with $e_S=f_{i,n}$, $e_T=f_{i,n+1}$ for an ascending pair and
$e_S=\overline{f_{i+1,n}}$, $e_T=\overline{f_{i,n+1}}$ for a descending one; \textup{(H)} of
\S\ref{ptr:ss-mult} holds for them; and \textup{(A3)} of Lemma~\ref{pc:arch} holds for
$(\pgyG_{i,0},\pgyG_{i,1})$. In particular $a_{i,n}$ maps $\pcF_{i,n}$ to $\pcF_{i,n+1}$ and $d_{i,n}$ maps
$\pcF_{i,n+1}$ to $\pcF_{i+1,n}$, all forms involved are cusp forms, rapidly decreasing, and those of a
given level form finite-dimensional spaces.
\end{proposition}
\begin{proof}
Let $(S,T)$ be one of the pairs, of ranks $n<n'$ with $n'-n=d=t-1$, with its characters
$\alpha,\beta$ and the row and column parameters $(\lambda,\lambda')$ of Lemma~\ref{pgy:match}. Steps
0--5 of the proof of Theorem~\ref{pgv:G3} are written for a pair of standard spaces $(S,T)$ and use the
following properties only, each of which holds here.
\begin{itemize}
\item[(i)] The characters have integral weights, and $\tau_0=\frac{n'}2-\pgvwt_\tau\alpha$,
$\tau'_0=\pgvwt_\tau\beta-\frac n2$ (Step 0: (V) and Lemma~\ref{ptr:const} hold for every pair of
nondegenerate spaces). Here all ranks are even and the constants are computed in the proof of
Lemma~\ref{pgy:match}.
\item[(ii)] $(\lambda,\lambda')$ satisfies the block rule of excess $d$ (Steps 1, 2, 5, through
Lemma~\ref{pgv:rule}). This is Lemma~\ref{pgy:match}, with $d=t-1$.
\item[(iii)] The entries of $\lambda$ and of $\lambda'$ are distinct and lie in $\frac{n-1}2+\Z$,
$\frac{n'-1}2+\Z$, and the conclusions of {\cite[Theorem~3.17]{P49}}, {\cite[Lemma~3.25]{P49}},
{\cite[Theorem~3.2]{P49}} and {\cite[Corollary~3.3]{P49}} hold for the types of the two nodes and their conjugates
(Steps 1--5). This is Lemma~\ref{pgy:gap}.
\item[(iv)] $n,n'\le2n_m$, so that $2fnn'+2\le x^{25}$ (Step 4). This is Lemma~\ref{pgy:gap}.
\end{itemize}
No other property of the spaces enters; in particular the position of the stem relative to the blocks
is not used. So Steps 0--4 give (a)--(d) of Theorem~\ref{ptr:edge} for $(S,T)$, and Step 5 gives (H).
In Step 3 the nodes without negative sign (if any) have compact archimedean group and are treated as
$j=m$ in {\cite[Proposition~8.4]{P49}}. The directions in Step 4: for an ascending pair the lift goes from
forms of type $e_S=f_{i,n}$ to forms of type $e_T=f_{i,n+1}$; for a descending pair the transposed lift
goes from forms of type $\overline{e_T}=f_{i,n+1}$ to forms of type $\overline{e_S}=f_{i+1,n}$.

For (A3) of Lemma~\ref{pc:arch} we go through the proofs of Theorem~\ref{ptq:F} and
Proposition~\ref{ptr:Z} for $(S,T)=(\pgyG_{i,0},\pgyG_{i,1})$. Theorem~\ref{ptq:F} is proved for every
polynomial $P\ne0$ none of whose monomials contains two variables of one column from rows of opposite
signs; its proof uses the Fock model of $\auHom(S,T)$ at $\tau$ with the signs $\eta_{rc}=-\varepsilon_r
\varepsilon'_c$, (W1$_\infty$), the strong continuity of $\sigma$ and the fact that $K$ acts on polynomial
vectors by linear substitutions times a character (Step 0), and nothing else about the pair. The joint
harmonic has the property of the monomials by (a). So (F) holds. The proof of
Proposition~\ref{ptr:Z} uses: $c_\lambda\in L^1$ and formula {\cite[(59)]{P51}} (here by (iii)); that
$\Phi$ is a highest weight vector of weight $\Lambda^K(\lambda)$ for $K$ ({\cite[Definition~3.9]{P49}} and
{\cite[Lemma~3.10(iii)]{P49}}, which does not involve the excess); Schur orthogonality on $K$; and (F).
Hence $Z_\infty=\prod_\tau d_{\lambda_\tau}^{-1}\langle\Phi_\tau,\mathsf P_\tau\Phi_\tau\rangle$ with every
factor positive, which is (A3), as in the proof of Proposition~\ref{ptv:arch}.
\end{proof}

\paragraph*{Characters of the squares}
\begin{lemma}\label{pgy:chars}
With $c_i=0$, for every square at $(i,n)$, $i+n+2\le m$ \textup{(}$U_1=\pgyG_{i,n+1}$,
$U_2=\pgyG_{i+1,n+1}$, $A=\pgyG_{i,n+2}$, $A'=\pgyG_{i+1,n}$, $M=N_{n+1}$\textup{)}, put
$\nu=\nu_1^{-(n+1)}$ and $\gamma=\nu^{-1}\upsilon^{M}$. Then the four conditions of \textup{(E2)} hold (at $A$, $A'$, $U_1$, $U_2$ respectively, with
the source and target characters as listed in Proposition~\ref{ptv:chars}):
\[
\beta_{i,n+1}\beta^\downarrow_{i,n+1}=\alpha^\downarrow_{i,n}\alpha_{i+1,n}=
\alpha_{i,n+1}\beta^\downarrow_{i,n}=\alpha^\downarrow_{i,n+1}\beta_{i+1,n}=\gamma .
\]
Moreover the target character of the
descent $d_{i,0}=a_{i,0}^*$ obtained from \textup{(W7)} is $\beta^\downarrow_{i,0}$.
\end{lemma}
\begin{proof}
Write $\mu^0_{i+n+2}=\mu^0_{i+1}\nu_1^{\,n+1}$, so $\gamma(\mu^0_{i+n+2})^{-1}=(\mu^0_{i+1})^{-1}
\upsilon^M$. At $A$: $\beta_{i,n+1}\beta^\downarrow_{i,n+1}=\mu^0_{i+n+2}\upsilon^{M/2}\allowbreak (\mu^0_{i+1})^{-1}
\upsilon^{M/2}=\nu_1^{n+1}\upsilon^M=\gamma$. At $A'$: $\alpha^\downarrow_{i,n}\alpha_{i+1,n}=(\mu^0_{i+1})^{-1}\upsilon^{N_{n+1}/2}
\allowbreak \mu^0_{i+n+2}\upsilon^{N_{n+1}/2}=\gamma$. At $U_1$:
$\alpha_{i,n+1}\beta^\downarrow_{i,n}=\mu^0_{i+n+2}\upsilon^{N_{n+2}/2}\allowbreak (\mu^0_{i+1})^{-1}\upsilon^{N_n/2}=
\gamma$, since $N_n+N_{n+2}=2M$. At $U_2$: the source character $\alpha^\downarrow_{i,n+1}=
(\mu^0_{i+1})^{-1}\upsilon^{N_{n+2}/2}$ of $d_{i,n+1}$ times the target character
$\beta_{i+1,n}=\mu^0_{i+n+2}\upsilon^{N_n/2}$ of $a_{i+1,n}$ is again $\gamma$. The last assertion is the argument of Proposition~\ref{ptv:chars}(b): by (W7) the target character of
the conjugate kernel is $(\mu^0_{i+1})^{-1}$ times an integral power of $\upsilon$ (the target
character $\beta_{i,0}=\mu^0_{i+1}\upsilon^{N_0/2}$ conjugated, times powers of $\upsilon$), its weight is
forced by the block rule (Lemma~\ref{pgy:match}, converse part, for $n=0$, which applies because the
conjugate kernel is a kernel of the descending pair with negated parameters, as in
Remark~\ref{ptr:rem-desc}); $\beta^\downarrow_{i,0}$ has the same form and the same weight, and
$\upsilon$ has weight $1$ at every $\tau$, so the two are equal.
\end{proof}

\paragraph*{(T1): the squares}
\begin{proposition}[(T1)]\label{pgy:T1}
For $i+n+2\le m$ there is $\sigma_{i,n}\in\C$ with $w_n\,a_{i+1,n}\circ d_{i,n}=\sigma_{i,n}\,w_{n+2}\,
d_{i,n+1}\circ a_{i,n+1}$ on $\pcF_{i,n+1}$.
\end{proposition}
\begin{proof}
We check that the square at $(i,n)$, with the one-member family $(\pcb,\nu,\chi,\text{the
characters above})$, $\chi$ the tame character of $\nu$, is an edge square of excess $e=t-1$ and
weight $\mu_\tau=(n+1)\Delta$ (Definition~\ref{ptr:def-square}). \emph{Ranks}:
$\dim A=N_{n+2}=N_n+2(t-1)$ and $N_{n+1}+N_{n+1}=N_{n+2}+N_n=N=2M$. \emph{Signs}: all negative signs are
on stems, and comparing the exponents $E$ coordinate by coordinate as in \S\ref{pc:ss-exchange}
(for $a\le i$ both sides have the exponents $\{i+a,i+a+1\}$; for $i<a\le i+n+1$ all four are $2a$;
for $a\ge i+n+2$ both sides have $\{i+a+n+1,i+a+n+2\}$) gives $q=q_A+q_{A'}$. As each stem has $r+1$
coordinates, $q\le2r+2$, so $2q+2\le4r+6\le2(r+1)+2(t-1)\le N$. And $N\le4n_m\le52r^{11/10}<2\Delta\le
2\mu_\tau$ (proof of Lemma~\ref{pgv:gap}). \emph{Data}: $\pcb\ne(1)$ is a modulus; $\nu=\nu_1^{-(n+1)}$
has weight $(n+1)\Delta$, is unramified outside the chosen places and tamely ramified there with tame
part $\chi_M^{-(n+1)}$, non-trivial since $n+1\le m$ and the tame parts of $\chi_M$ have order $>2m$.
\emph{(E1)} is Proposition~\ref{pgy:edge} for the four pairs $(U_1,A)$, $(A',U_2)$ (ascending at
$(i,n+1)$, $(i+1,n)$) and $(U_2,A)$, $(A',U_1)$ (descending at $(i,n+1)$, $(i,n)$), whose polynomials
and types do not depend on anything else. \emph{(E2)} is Lemma~\ref{pgy:chars} (the factor
$\upsilon^{N/2}=\upsilon^M$ is unramified, so the tame part of $\gamma^{-1}$ at $\pcb$ is $\chi$).
Theorem~\ref{ptr:square} (which rests on Theorem~\ref{pcx:X} and Proposition~\ref{ptr:conv}; the
latter uses only (E1) and full rank) gives the square law of \S\ref{pc:ss-exchange} for the functions
$L,L'$ of the square with constant $\kappa/c$. Step (iv) of \S\ref{pc:ss-exchange}, i.e.\
Lemma~\ref{pc:next} with $\mathsf N_0$ the frame of $A'$ (so $\#\auGL_{\mathsf N_0}(\auO_E/\pcb)=w_n$) and
$\mathsf A$ the frame of $A$ ($\#\auGL_{\mathsf A}=w_{n+2}$), with $p=\varepsilon_{i,n}$ and
$p^+=\varepsilon_{i,n+1}$, turns it into the identity, by pairing both sides with arbitrary forms of
$\pcF_{i+1,n}$ (the pairings separate, the spaces of a given level being finite-dimensional and the
Petersson product non-degenerate). So $\sigma_{i,n}=\kappa/c$.
\end{proof}

\paragraph*{(T2), (T3): the first lifts}
\begin{lemma}[(T2)]\label{pgy:T2}
There is a finite datum $\varepsilon_{i,0}$ modulo $\pcb$, with $|\varepsilon_{i,0}|=|\delta_{\iota}|$ up to a
constant of modulus one, such that the transposed lift $d_{i,0}$ along $(\pgyG_{i+1,0},\pgyG_{i,1})$ with it
is the adjoint of $a_{i,0}$ under $\pcF_{i+1,0}=\pcF_{i,0}$. In particular
$\langle z,a_{i,0}j_iy\rangle=\langle d_{i,0}z,j_{i+1}y\rangle$ for $y\in V$, $z\in\pcF_{i,1}$.
\end{lemma}
\begin{proof}
This is item (ii) of Proposition~\ref{ptv:tri}, whose proof applies verbatim: Lemma~\ref{ptv:adj}
(its hypotheses: $\delta_\iota$ is supported on injective matrices, and the edge property is
Proposition~\ref{pgy:edge}) gives $a_{i,0}^*z(g)=\int\aubar{\theta(g,h)}z(h)\,dh$; by (W7) the conjugate
kernel is the kernel for $(\pgyG_{i,0}^-,\pgyG_{i,1})=(\pgyG_{i+1,0},\pgyG_{i,1})$ of the datum
$\mathsf C[P_{i,0},\delta_\iota]=[\aubar{P_{i,0}},c\,\delta_{\iota}]$, $|c|=1$ (at the primes of $\pcb$
the idele $\alpha$ is a unit by Proposition~\ref{pgv:alpha}(a), so $a_w=1$ in
Definition~\ref{pgx:conv}, and in the model (W5) $\mathsf C$ is complex conjugation of functions up to
a constant of modulus one, since $[x,a']^-=\aubar{[x,a']}$, as in the proof of
Proposition~\ref{ptv:tri}; outside $S(\pcb)$ it maps the vacuum to a multiple of the vacuum by (W4));
$\aubar{P_{i,0}}$ is the joint harmonic of the descending pair for
$n=0$ (Remark~\ref{ptr:rem-desc}), and its characters are those of Lemma~\ref{pgy:chars}.
\end{proof}

\begin{lemma}[The series in rank $r+1$]\label{pgy:series}
Let $S$ be a space of rank $n=r+1$, $M'=r+t$, $\pcm$ a modulus and $\alpha$ a character unramified
outside $\pcm$. Then $\sum_{\gamma\in G^\pcm/K^\pcm,\ \gamma\notin K^\pcm}|\pcmc_{M',\alpha}(\gamma)|<\frac12$,
and the conclusions of Lemma~\ref{pc:series}\textup{(c)} hold for $\pcRal_{M',\alpha}$ on $S$.
\end{lemma}
\begin{proof}
Here $h:=M'/2=3r=3n-3$, whereas Lemma~\ref{pc:series}(b) assumed $M'\ge6n$; we redo the count of its
proof with the same notation ($q=q_w\ge5$). \emph{Split $v$}: as there,
$1+\varepsilon_v\le\prod_{i<n}(1-q^{i-h})^{-2}$; here $i-h\le-2n+2$, $\sum_{i<n}q^{i-h}\le\frac54q^{-2n+2}\le
\frac15$, so $1+\varepsilon_v\le\exp(5q^{-2n+2})$ and $\varepsilon_v\le10q^{-2n+2}\le6250\cdot25^{-n}q^{-2}\le
2(4/5)^nq^{-2}$, as $3125\le20^n$ for $n\ge3$. \emph{Non-split $v$}: as there, $\varepsilon_v\le\sum_{k\ge1}
s_k^2q^{-kh}$ with $s_k\le2^{k+n-1}q^{k(n-1)}$, so $\varepsilon_v\le4^{n-1}\sum_{k\ge1}(4q^{2n-2-h})^k=
4^{n-1}\sum_k(4q^{-n+1})^k\le2\cdot4^nq^{-n+1}$ (as $4q^{-n+1}\le\frac12$), and $2\cdot4^nq^{-n+1}=
2q^{-2}(4/5)^n(5/q)^nq^3\le250(4/5)^nq^{-2}$ (as $(5/q)^nq^3\le125$ for $n\ge3$, $q\ge5$).
\emph{All places}: as there, $\sum_v\varepsilon_v\le250\cdot2f(4/5)^n\cdot\frac12=250f(4/5)^n$. With
$n>r\ge x^5$, $f\le x^{12}$ (the sizes used in the proof of Lemma~\ref{pc:series}(c)) and
$x\ge128^{100}$, $250f(4/5)^n<\frac13$, so the sum is at most $e^{1/3}-1<\frac12$. The proof of
Lemma~\ref{pc:series}(c) then applies unchanged.
\end{proof}

\begin{proposition}[First lifts]\label{pgy:first}
For $0\le i<m$ there is $C\ne0$ with
$\langle a_{i,0}y',a_{i,0}y\rangle=C\langle y',\pcRal\,y\rangle$ for all $y,y'\in V$, where
$\pcRal=\pcRal_{r+t,\alpha_{i,0}}$ with the modulus $\pcb$; and $a_{i,0}$ is injective on $V$. Moreover for
$u\in J(\pcb)$ there is $g\in K_{\pgyG_{i,1}}$ with $a_{i,0}\rho(u)=\rho(g)a_{i,0}$.
\end{proposition}
\begin{proof}
Proposition~\ref{pc:rallis} is stated for ranks $(r,r+t)$; its proof uses the ranks only through
(SW) of Lemma~\ref{pc:SW} and through Lemma~\ref{pc:series}. For the ranks $(r+1,r+t)$ of
$(\pgyG_{i,0},\pgyG_{i,1})$ we use Lemma~\ref{pgy:series} in place of Lemma~\ref{pc:series}, and (SW) for
the doubled space $\pcW=\pgyG_{i,0}\oplus\pgyG_{i,0}^-$ of rank $2(r+1)$ and $T=\pgyG_{i,1}$ of rank
$M'=6r$: as recalled in the comments to Lemma~\ref{pc:SW}, Weil's theorem needs the convergence of the
Eisenstein series, $M'>2(r+1)$, and of the theta integral, $M'$ minus the Witt index of $T$ greater
than $r+1$; the Witt index is at most $r+1$ (all negative signs are on the stem), and $6r-(r+1)>r+1$.
(These are condition (B) and Proposition~8 of~\cite{WeilSiegel} in the hermitian case, the hypotheses of \cite[Th\'eor\`eme~5]{WeilSiegel}; here $r\ge1$, so $6r>2(r+1)$.)
Lemma~\ref{pc:clean} is rank-free. (A1), (A2), (A3) hold by Proposition~\ref{pgy:edge}. The point
datum $\delta_\iota$ on $\pcMat_{N_1\times N_0}(\auO_E/\pcb)$ is sharp for $J=J(\pcb)$ (with $A$ the projection
to the first $N_0$ rows, $A\iota=1\in\bar J=\{1\}$, and $\alpha_{i,0}$ trivial on $\det J(\pcb)$, its tame
parts living at the chosen places, i.e.\ at the primes of $\pcb$, and being characters of $\auO_E/\pcb$
that are trivial on $1$). So Proposition~\ref{pc:rallis}, with $\#\pcMat_{(r+t)\times(r+1)}$ in place of
$\#\pcMat_{(r+t)\times r}$, gives the formula and the injectivity on $\pcF(\pgyG_{i,0},f_{i,0})^{J(\pcb)}=V$.
For the last assertion choose $g$ with reduction $\operatorname{diag}(\bar u,1)=1$ modulo $\pcb$, e.g.\
$g=1$ at the primes of $\pcb$, and argue as in the proof of Proposition~\ref{pc:first}: here $\bar u=1$, so
by Lemma~\ref{pc:level} the pair $(u^{-1},1)$ fixes $[P,\delta_\iota]$ and $a_{i,0}\rho(u)=a_{i,0}$, i.e.\
the claim with $g=1$.
\end{proof}

\begin{lemma}[(T3)]\label{pgy:T3}
For $i<m$, $d_{i,0}a_{i,0}j_i(V)\subseteq j_{i+1}(V)$, and $d_{i,0}a_{i,0}j_i$ is injective.
\end{lemma}
\begin{proof}
As in the proof of Corollary~\ref{pc:istriangle}. Let $x\in V$ and $z=d_{i,0}a_{i,0}x\in\pcF_{i,0}$. By
Lemma~\ref{pgy:T2}, $\langle z,f\rangle=\langle a_{i,0}x,a_{i,0}f\rangle$ for all $f\in\pcF_{i,0}$. For
$u\in J(\pcb)$, by Proposition~\ref{pgy:first}, $a_{i,0}\rho(u^{-1})=a_{i,0}$, hence
$\langle\rho(u)z,f\rangle=\langle z,\rho(u^{-1})f\rangle=\langle a_{i,0}x,a_{i,0}f\rangle=\langle z,f\rangle$; so
$z$ is $J(\pcb)$-invariant, $z\in V$ (it lies in $\pcF_{i+1,0}=\pcF_{0,0}$ by Proposition~\ref{pgy:edge}).
If $z=0$ then $\langle a_{i,0}x,a_{i,0}x\rangle=0$ and $x=0$ by Proposition~\ref{pgy:first}.
\end{proof}

\begin{theorem}[Input (I)]\label{pgy:I}
The data $(V,\pcF_{i,n},j_i,w_n,a_{i,n},d_{i,n})$ form a triangle of size $m$ over $\{1\}$
(Definition~\ref{pc:triangle}). Consequently $a^\flat=a^\flat_{m-1}\circ\dots\circ a^\flat_0$ is injective
on the forms of type $\overline{e_{(1,0,0)}}$ on $\pgrY^0_0$ fixed by $J_{\pgrY^0_0}(\pcb)$: input (I) of
\S\ref{pgw:ss-first} holds.
\end{theorem}
\begin{proof}
The $\pcF_{i,n}$ are inner product spaces (Petersson products of bounded forms on finite-volume
quotients), $j_i$ is an isometry, $w_n>0$. (T1) is Proposition~\ref{pgy:T1}, (T2) Lemma~\ref{pgy:T2},
(T3) Lemma~\ref{pgy:T3}. By Proposition~\ref{pc:tri-inj} with $(i,n)=(0,m)$, $X_{0,m}\circ j_0$ is
injective on $V$, and $X_{0,m}=w_{m-1}\cdots w_0\,a^\flat$ by \eqref{pgy:eq-node} and \eqref{pgy:eq-chain}.
\end{proof}

\begin{remark}\label{pgy:rem-I}
Only the existence of the constants $\sigma_{i,n}$ is used, not their size or non-vanishing (the
latter follows from Proposition~\ref{pc:tri-inj} if $V\ne0$). The triangle is untwisted, so no
modulus of a twist and no $L$-value estimate enters.
\end{remark}

\subsubsection{The change of modulus (P)}\label{pgy:ss-P}

We recall the data of \S\ref{pc:ss-data} (Definition~\ref{pgx:conv}) at the primes of a modulus: generators $a_w$ of
$\alpha_w\auO_w$, fixed once for every prime $w$ (independently of the modulus), $a_w=1$ where $\alpha_w$
is a unit, and the vacuum at $v\in S(\pcm)$ normalised as the function $\mathbf 1_{M_w}$,
$M_w=a_w\pcMat(\auO_w)$, of norm $1$ for the measure of Definition~\ref{pgx:conv} (a vacuum by
Lemma~\ref{pgx:vac}). Outside $S(\pcm)$ the vacuum is that of $\alpha\auHom(L_S,L_T)$ as in
\S\ref{pc:ss-data}, with the same normalisation at split places.

\begin{proposition}[Input (P)]\label{pgy:P}
Let $\pcm'\subseteq\pcm$ be moduli, $\mathrm{red}\colon\auO_E/\pcm'\to\auO_E/\pcm$ the reduction, $S,T$
spaces, $P$ a Fock vector and $\Psi$ a function on $\pcMat_{N'\times N}(\auO_E/\pcm)$. Then the datum
$[P,\Psi\circ\mathrm{red}]$ for the modulus $\pcm'$ equals the datum $[P,\Psi]$ for the modulus $\pcm$.
In particular, for every edge of the second column and every $\Psi$ on $\pcMat(\auO_E/\pcb)$, the data
$[P,\Psi]$ for $\pcb$ and $[P,\Psi\circ\mathrm{red}]$ for $\pcm(\vartheta)=\pcb\prod_{w\in Q_\vartheta}w$ coincide, with
no assumption on $\alpha$; so the kernels, lifts and transposed lifts coincide.
\end{proposition}
\begin{proof}
Both data are $P$ at the infinite places and the vacuum of $\alpha\auHom(L_S,L_T)$ at the finite places
outside $S(\pcm')$. At $v\in S(\pcm)$, $v=w\aubar w$, $w\mid\pcm$: the $w$-component of
$(\Psi\circ\mathrm{red})^\natural$ is $x\mapsto\Psi(\mathrm{red}(a^{-1}x\bmod\pcm'))=\Psi(a^{-1}x\bmod\pcm)$ on
$a\pcMat(\widehat\auO)$, which is the $w$-component of $\Psi^\natural$; the measures used for the norms are
the same. (The decomposition into components is the factorisation of $\Psi^\natural$ over the primes of
$\pcm$ given by the Chinese remainder theorem, $\auO_E/\pcm=\prod_{w\mid\pcm}\auO_E/w^{k_w}$; for
non-factorisable $\Psi$ the identity is the same computation on $\prod_{w\mid\pcm}\pcMat(E_w)$.) At
$v\in S(\pcm')\smallsetminus S(\pcm)$, $w\mid\pcm'$, $w\nmid\pcm$: the value of $\Psi\circ\mathrm{red}$
does not depend on the $w$-component, so $(\Psi\circ\mathrm{red})^\natural$ is, in the variable $x_w$, the
indicator of $a_w\pcMat(\auO_w)=M_w$, which is the normalised vacuum at $v$ by Lemma~\ref{pgx:vac} and our
normalisation; this is the component of $[P,\Psi]$ at $v$, where the modulus $\pcm$ puts the vacuum. For
the second-column edges take $\pcm'=\pcm(\vartheta)$, $\pcm=\pcb$ (moduli: \S\ref{pgv:sec}).
\end{proof}

\begin{remark}\label{pgy:rem-P}
\S\ref{pc:ss-data} uses the normalised coordinates $a^{-1}x$ of Definition~\ref{pgx:conv}. Had
$\Psi^\natural$ been supported on the integral matrices in the frame coordinates, (P) would still hold at
the primes where $\alpha_w$ is a unit, but it would fail at a prime $w\mid Q_\vartheta$ where $\alpha_w$ is
not a unit, because $\mathbf 1_{\pcMat(\auO_w)}$ is not the vacuum there (Lemma~\ref{pgx:vac}). The proof
of Theorem~\ref{pgw:first} uses (P) only to write $T^\flat_b=\sum_{X\equiv\iota\,(\pcb)}T_X$, and, for the
first column, to identify $\delta_\iota$ modulo $\pcm_0$ with $\mathbf 1[X\equiv\iota\bmod\pcm_0]$ modulo
$\pcm$; both are cases of Proposition~\ref{pgy:P}.
\end{remark}




\section{The last path}\label{C1:g6}
% ===== C-grid8.tex =====
% proofs/C-grid8.tex
% Continuation of proofs/C-grid.tex (pgr:), C-grid2.tex (pgs:), C-grid4.tex (pgv:), C-grid5.tex
% (pgw:), C-grid6.tex (pgx:): Proposition pgr:den (denominators and conjugates on the last path).
% Labels start with pgz:. Conventions and macros: those of proofs/C-grid5.tex.
\providecommand{\auI}{\mathcal I}
\providecommand{\auGL}{\mathrm{GL}}
\providecommand{\auU}{\mathrm U}
\providecommand{\auO}{\mathcal O}
\providecommand{\auA}{\mathbb A}
\providecommand{\auNm}{\operatorname{Nm}}
\providecommand{\auHom}{\operatorname{Hom}}
\providecommand{\aubar}[1]{\overline{#1}}
\providecommand{\pcm}{\mathfrak m}
\providecommand{\pcb}{\mathfrak b}
\providecommand{\pcMat}{\mathrm M}
\providecommand{\pcR}{\mathsf R}
\providecommand{\pcF}{\mathcal F}
\providecommand{\auPl}{\Phi}
\providecommand{\psvac}{\mathbf v}
\providecommand{\psrho}{\varrho}
\providecommand{\pgrY}{\widetilde Y}
\providecommand{\pgrlast}{\Gamma_{\mathrm{last}}}
\providecommand{\pgrE}{\operatorname{Ed}}
\providecommand{\pcnot}[1]{\par\smallskip{\raggedright\noindent\textsc{Not proved here.} [#1]\par}\smallskip}
\providecommand{\pgupt}{\mathsf u}
\providecommand{\pgwK}{\mathcal K}
\providecommand{\pgwT}{\mathcal T}
\providecommand{\pgwint}[1]{\int^{#1}}
\providecommand{\pgzd}{\mathsf d}
\providecommand{\pgzM}{\mathsf M}

\subsection{Denominators and conjugates on the last path: Proposition~\ref{pgr:den}}\label{pgz:sec}

This section proves part (a) of Proposition~\ref{pgr:den} in full, for the sum $\Lambda$ of
\eqref{pgw:eq-Lambda} with the integral choices of the branching form and of the kernels fixed in
\S\ref{pgz:ss-choice}, and reduces part (b) to three size bounds stated in \S\ref{pgz:ss-b}. The
notation is that of \S\S\ref{pgv:sec}, \ref{pgw:sec}, \ref{pgx:sec}. Throughout, $\vartheta$ is an
$m$-fold twist, $\pcm=\pcm(\vartheta)=\pcb\prod_{w\in Q_\vartheta}w$ with $\pcb=(\pi_M)^b$, and the
data at the primes of $\pcm$ are those of Definition~\ref{pgx:conv}. For $z\in\{0,1\}$ and $0\le a\le
m$ let $p^z_a=(z,a,m)$ be the nodes of the half $h_z$ of $\pgrlast$, $G^z_a=G_{p^z_a}$,
$\Gamma^z_a=\Gamma_{p^z_a}$, $J^z_a=J_{p^z_a}$ its level (\S\ref{pgw:ss-last}: $J^0_0=J(\pcm)$, all
others $J(\pcb)$), $N^z_a$ its rank, and $e^z_a=(p^z_{a-1},p^z_a)$ for $1\le a\le m$. Put
$Y=Y^m_m$, so $p^0_m$ is the node $Y$ and $p^1_m$ the node $\pgrY^m_m=Y\oplus E\tilde e_m$.

The structure of the argument is the following. (i) All levels along $\pgrlast$ are free, and the
invariance of a path datum makes the integrand of $\Lambda$ invariant under these levels in every
variable separately; so $\Lambda$ is a plain sum, over the tuples of classes of $2m+1$ levels and
over the edge tuples, of $\Psi_\vartheta(S)$ times branching pairings of vectors reached by $m$
lifts from each end (Proposition~\ref{pgz:expand}). \emph{No index $[J:J']$ and no order of a
stabiliser occurs}, so these contribute nothing to $N_0$. (ii) One lift with an edge datum, multiplied
by the integer $\pgzd$ of \eqref{pgz:eq-d}, does not increase gauges at $\iota$ when
$\ell(\iota)\nmid\auNm Q_\vartheta$ (Lemma~\ref{pgz:kernel}). (iii) The branching pairing of two
vectors of gauge at most $1$ at $h$ and at $h\oplus\pgupt^m$ is integral after multiplication by the
integer $\pgzM$ of \eqref{pgz:eq-M} (Lemma~\ref{pgz:pair}). Then $N_0=\pgzM\pgzd^{2m}$.

\subsubsection{Integral structures and the integral choices}\label{pgz:ss-choice}

\paragraph*{The representations along the last path}
Let $W^z_a$, $\mu^z_a$, $\pgwK^z_a$ be as in \S\ref{pgw:ss-last} (the paragraph before
\eqref{pgw:eq-Lambda}). If some $W^z_a$ is $0$, then $\pgwT^z_S=0$ for all $S$ and $\Lambda=0$, and
there is nothing to prove; we assume all $W^z_a\ne0$.

\emph{$W^z_a$ is the representation of the type at $p^z_a$}, more precisely $W^z_a\cong\pi_{p^z_a}^*$ with
$\pi_p$ the representation of the type $e_p$, as $W^z_0=(\pi^z_0)^*$ in Lemma~\ref{pgw:vec}. By induction on
$a$; let $(p,q)=e^z_a$, a right edge (the first index grows along the row), and $W^z_{a-1}\cong\pi_p^*$. In
the convention fixed before Proposition~\ref{pgv:match}, $e_p$ is the type whose parameter is the conjugate
of the row parameter $\lambda$, and $e_q$ the type of the column parameter $\lambda'$; so $\pi_p^*\cong
\pi_\lambda$ and $\pi_q\cong\pi_{\lambda'}$. By Theorem~\ref{pgv:G3}(a),(b) the joint harmonic $P=P_{pq}\ne0$
satisfies $\mathsf PP=\mathsf P'P=P$ with $\mathsf P=\omega_\infty(e_\lambda,1)$, $\mathsf P'=\omega_\infty(1,
e_{\lambda'})$, $e_\lambda=d_\lambda\overline{c_\lambda}$; by Schur orthogonality $\mathsf P$ maps into the
$\pi_\lambda$-isotypic part for $G_{p,\infty}$ and $\mathsf P'$ into the $\pi_{\lambda'}$-isotypic part for
$G_{q,\infty}$. In the decomposition of Lemma~\ref{pgw:howe}, $P$ therefore lies in $\pgwP_{\nu_1}$ for the
unique $\nu_1$ with $A_{\nu_1}\cong\pi_\lambda$ (the $A_\nu$ are pairwise inequivalent), and
$B_{\nu_1}\cong\pi_{\lambda'}$. As $W=W^z_{a-1}\cong\pi_p^*\cong\pi_\lambda=A_{\nu_1}$, the index $\nu_0$ of
the proof of Lemma~\ref{pgw:one} is $\nu_1$, and $W^z_a=B_{\nu_1}^*\cong\pi_{\lambda'}^*\cong\pi_q^*$. The bounds
used below (the extreme entries of the highest weight, Lemmas~\ref{pgz:deg}, \ref{pgz:char}) are the
same for $\pi_q$ and $\pi_q^*$. We realise $W^z_a$ as $W_{\lambda^z_a}$ with the basis $(e_T)$ and the lattice
$W_{\lambda^z_a}(\Z)=\bigoplus\Z e_T$ of \S\ref{pgs:ss-weyl}, $\lambda^z_a$ the corresponding weight of
$\prod_\tau\auGL_{N^z_a}$. Its matrix coefficients in the basis $(e_T)$ are polynomials with integer
coefficients in the entries of $g$ and in $\det g^{-1}$. For a prime $\ell$ and $j\colon\C\to\C_\ell$,
the \emph{gauge} at $(j,g)$, $g\in G^z_{a,f}$, of $\xi\in W^z_a$ is the largest absolute value of a
coordinate of $g_\ell^{-1}\cdot j(\xi)$, the action of $g_\ell$ being through the embeddings of
{\cite[\S14.3]{P46}} (as in the definition of integral forms there). A $W^z_a$-valued form is
\emph{integral at $\iota$} if each value $\psi(g)$ has gauge at most $1$ at $(\iota,g)$; this is the
hypothesis of Proposition~\ref{pgr:den}(a) on $y'=f_{\mu^0_0,\psi_{y'}}$ and
$\varphi=f_{\mu^1_0,\psi_\varphi}$.

Two facts on the gauge are used. (G-a) For $k\in G^z_{a,f}$ integral with integral inverse at the
places above $\ell$, $k_\ell$ acts on $W_{\lambda^z_a}(\auO_{\C_\ell})$ by a matrix invertible over
$\auO_{\C_\ell}$ (its entries and those of its inverse are integral polynomials in entries that are
integral). (G-b) For a rational point $\gamma\in\Gamma^z_a$, the archimedean action of $\gamma$ on
$W^z_a$ followed by $j$ equals the action of $\gamma_\ell$ on $W^z_a\otimes_j\C_\ell$: both are given
by the same integral polynomials evaluated at the images of the entries of $\gamma$ under the
embeddings concerned ({\cite[\S14.3]{P46}}). Hence the gauge of $\gamma\xi$ at $(j,\gamma g)$ equals
the gauge of $\xi$ at $(j,g)$.

\begin{definition}[Integral choices]\label{pgz:def-choice}
(a) \emph{Kernels.} For the edge $e^z_a=(p,q)$ let $S_\nu(z)$ be the Schur matrix of the diagram
$\nu=\nu_0$ of the edge (Lemma~\ref{pgw:one}, Remark~\ref{pgw:rem-jets}) in the bases $(e_T)$, an
integer matrix of polynomials in the variables $z_{\tau,c,a'}$, homogeneous of total degree
$e(\nu)=\sum_\tau|\nu_\tau|$, and let $\Theta_f$ be the function of {\cite[(42)]{P49}} attached to a
finite datum $f$ on $\auHom(p,q)$, with the scale $s_0$ of {\cite[(45)]{P49}}. We take
\[
\pgwK^z_a(f)=c_{\mathrm{vac}}^{\,N_pN_q}\,S_\nu(\partial)\bigl[\Theta_f(s_0\,\cdot)\bigr](0),
\]
$c_{\mathrm{vac}}$ the constant of {\cite[(58)]{P49}}. By Remark~\ref{pgw:rem-jets} this is the kernel
of Lemma~\ref{pgw:one} for one choice of $j$, multiplied by a non-zero constant independent of
$\vartheta$.

(b) \emph{Branching form.} By Lemma~\ref{pgw:branch}(a) the invariant forms $B$ on $W^0_m\times
W^1_m$ form a space of dimension at most $1$. In the bases $(e_T)$ the conditions of invariance are
linear equations with rational coefficients (the representation matrices have rational coefficients
and the invariance under the compact group $G_{Y,\infty}$ is equivalent, by Zariski density, to the
polynomial identities $B(\rho_0(g)w_0,\rho_1(g\oplus1)w_1)=B(w_0,w_1)$ on $\prod_\tau\auGL_n$). So
if a non-zero invariant form exists, there is one whose matrix in the bases $(e_T)$ has coprime
rational integer entries; we take $B$ to be this matrix (it is unique up to sign), and $B=0$ if
there is no non-zero invariant form.
\end{definition}

\noindent\emph{The bases $(e_T)$ and the Schur matrices.} The Schur matrix $S_\nu$ of (a) is the matrix,
in the bases $(e_T)$ of \S\ref{pgs:ss-weyl}, of the Schur functor $Z\mapsto S(Z)$ of
{\cite[(38)]{P49}}; these are the coordinates in which {\cite[\S14.3]{P46}} writes the matrix of $g$
as the Schur matrix, so $S_\nu$ is the matrix of {\cite[(34)]{P49}} for the edge. Its entries are
polynomials with integer coefficients in the entries of $Z$: in the basis $(b_T)$ of standard products of
minors, $S(Z)b_{T'}(x)=b_{T'}(Z^{\mathsf t}x)$ is, by the Cauchy--Binet formula and straightening, an integer
combination of the $b_T$ with coefficients integer polynomials in $Z$; and $(e_T)$ is the dual basis for
the invariant form, for which $S(Z)^*=S(Z^*)$, so the matrix of $S(Z)$ in $(e_T)$ is the conjugate
transpose of that of $S(Z^*)$ in $(b_T)$, again with integer polynomial entries in $Z$. The operator norm
{\cite[(38)]{P49}} and the degree $e(\nu)$ do not depend on the basis. These are the only properties
of $S_\nu$ used in Lemmas~\ref{pgz:deg}, \ref{pgz:char} and~\ref{pgt:jet}.

\noindent By the last sentence of the paragraph containing \eqref{pgw:eq-Lambda}, Theorem~\ref{pgw:last}
holds for these choices, with another constant $\beta(g)$. The choice of the matrix $B$ only affects
(b) of Proposition~\ref{pgr:den}; see \S\ref{pgz:ss-b}.

\begin{lemma}[The degree of the kernels]\label{pgz:deg}
For every edge $e^z_a$, $e(\nu)\le40fn_mb\le x^{19}b$.
\end{lemma}
\begin{proof}
By Lemma~\ref{pgw:howe} and the proof of Lemma~\ref{pgw:one}, $W'=W^z_a$ is the dual of
$B_{\nu}$, which is the restriction to $\prod_\tau\auU(N_q)$ of $\bigotimes_\tau W_{\nu_\tau}(
\C^{N_q})^*$ twisted by a character; by Lemma~\ref{pgw:howe}, $\nu_\tau$ has at most $\min(N_p,N_q)$
rows. Along the half $h_z$ the node $p^z_a$ has $a$ blocks of the first kind and $p^z_{a-1}$ has
$a-1$, with the same stem and the same blocks of the second kind (\S\ref{pgw:ss-last}), so
$N_q=N_p+t'$ with $t'\in\{t,t+1\}$ the rank of a block of the first kind; thus $N_p<N_q$ and
$\nu_\tau$ has at most $N_q-1$ rows. So the highest weight of $W^z_a$ at $\tau$ is, up to adding a
constant vector (the twisting character, which is the same for all entries), $\pm\nu_\tau$ padded by
at least one zero, reordered. Adding a constant vector does not change differences of entries, and
since a padded zero entry occurs, the largest minus the smallest entry of $\pm\nu_\tau$ padded is
exactly $\nu_{\tau,1}-0$; hence the longest row $\nu_{\tau,1}$ equals the difference between the
largest and the smallest entry of the highest weight. By the proof of Lemma~\ref{pgs:boxes}, for the weight at the node $q$ this difference
is at most $20b$. Hence $|\nu_\tau|\le N_q\nu_{\tau,1}\le2n_m\cdot20b$, and summing over the $f$
places $\tau$ gives $e(\nu)\le40fn_mb\le x^{19}b$ (Lemma~\ref{pgs:boxes}).
\end{proof}

\subsubsection{The levels are free, and $\Lambda$ is a sum without denominators}\label{pgz:ss-free}

\begin{lemma}[Free levels on the last path]\label{pgz:free}
For every $z$, $a\le m$ and $g\in G^z_{a,f}$, $\Gamma^z_a\cap gJ^z_ag^{-1}=1$. Consequently, for every
open subgroup $J'\subseteq J^z_a$ and every function $F$ on $G^z_{a,f}$ that is left
$\Gamma^z_a$-invariant and right $J'$-invariant,
\[
\pgwint{J^z_a}_{\Gamma^z_a\backslash G^z_{a,f}}F(x)\,dx=\sum_{x\in\Gamma^z_a\backslash G^z_{a,f}/J^z_a}
\int_{J^z_a}F(xk)\,dk ,
\]
$dk$ the Haar probability measure of $J^z_a$; if $F$ is right $J^z_a$-invariant this is
$\sum_{x}F(x)$, a sum over representatives of the classes with coefficient $1$.
\end{lemma}
\begin{proof}
The rank is $N^z_a\le2n_m$ (Lemma~\ref{pgs:boxes}). Both $J(\pcb)$ and $J(\pcm)$ are contained in
$J(\pi_M)$, whose elements are $\equiv1$ modulo every chosen place, because the chosen places divide
$\pi_M$ (proof of {\cite[Lemma~8.8(a)]{P49}}). The space at $p^z_a$ is a standard positive definite
space (Definition~\ref{pgv:def-spaces}: for $a=0$, $z=0$ it is $Y_m$, definite by
\S\ref{pgr:ss-arch}; the other nodes of row $m$ are diagonal with all signs $+1$, by
Lemma~\ref{pgv:count}), so {\cite[Lemma~14.5]{P46}} applies and gives the first assertion. For the
formula: the class $\Gamma xJ^z_a$ is the disjoint union of the classes $\Gamma xkJ'$, $k$ running
over $J^z_a/J'$, which are pairwise distinct since $\Gamma\cap xJ^z_ax^{-1}=1$; each has measure
$[J^z_a:J']^{-1}$ for the measure giving $J^z_a$ the volume $1$ (its stabiliser in $\Gamma$ being
trivial). Hence the integral over $\Gamma\backslash\Gamma xJ^z_a$ is
$[J^z_a:J']^{-1}\sum_{k\in J^z_a/J'}F(xk)=\int_{J^z_a}F(xk)\,dk$.
\end{proof}

\begin{lemma}[Moving the data by the levels]\label{pgz:move}
Let $(p,q)=e^z_a$, $X\in\pcMat_{q\times p}(\pcR_\pcm)$, $k\in K_p$, $k'\in K_q$ with reductions
$\bar k,\bar k'$ modulo $\pcm$, and $x\in G_{p,f}$, $h\in G_{q,f}$. Then
\[
\pgwK^z_a\bigl(\omega_f(xk,hk')[\delta_X]_f\bigr)=\chi_{pq}\bigl(\det\bar k'/\det\bar k\bigr)\,
\pgwK^z_a\bigl(\omega_f(x,h)[\delta_{\bar k'X\bar k^{-1}}]_f\bigr).
\]
\end{lemma}
\begin{proof}
$\omega_f(xk,hk')=\omega_f(x,h)\omega_f(k,k')$. By Lemma~\ref{pc:level}, which holds for the data of
Definition~\ref{pgx:conv} (remark after that definition), $\omega_f(k,k')[\delta_X]_f$ is the datum
of the function $y\mapsto\alpha_\pcm(\det\bar k)^{-1}\beta_\pcm(\det\bar k')\delta_X(\bar k'^{-1}y\bar
k)=\alpha_\pcm(\det\bar k)^{-1}\beta_\pcm(\det\bar k')\delta_{\bar k'X\bar k^{-1}}(y)$; and
$\alpha_{pq}$, $\beta_{pq}$ act on determinants of elements of $K$ through the same character
$\chi_{pq}$ (remark after Definition~\ref{pgv:def-edges}). Apply the linear map $\pgwK^z_a\circ
\omega_f(x,h)$.
\end{proof}

For an edge tuple $S$, points $\mathbf x=(x^z_a)_{z\in\{0,1\},\,0\le a<m}$ with $x^z_a\in G^z_{a,f}$,
and $h\in G_{Y,f}$, put $x^0_m=h$, $x^1_m=h\oplus\pgupt^m$ ($\pgupt=\pgupt_\pcm$ of
Definition~\ref{pgv:def-period}) and define the vectors \emph{reached along $h_z$}
\begin{equation}\label{pgz:eq-reached}
v^z_S(\mathbf x,h)=\pgwK^z_m\bigl(\omega_f(x^z_{m-1},x^z_m)[\delta_{S_{e^z_m}}]_f\bigr)\cdots
\pgwK^z_1\bigl(\omega_f(x^z_0,x^z_1)[\delta_{S_{e^z_1}}]_f\bigr)\,\psi^z(x^z_0)\in W^z_m,
\end{equation}
with $\psi^0=\psi_{y'}$, $\psi^1=\psi_\varphi$.

\begin{proposition}[$\Lambda$ as a sum over classes]\label{pgz:expand}
Let $\Psi_\vartheta$ be a path datum on $\pgrE(\pgrlast)$ with bound $N$, and let $y'$, $\varphi$ be
fixed by the levels $J^0_0=J(\pcm)$ and $J^1_0=J(\pcb)$ of the two outer nodes. Choose representatives
$\mathcal R^z_a$ of $\Gamma^z_a\backslash G^z_{a,f}/J^z_a$ ($a<m$) and $\mathcal R_Y$ of
$\Gamma_Y\backslash G_{Y,f}/J_Y$. Then
\[
\Lambda(\Psi_\vartheta;y',\varphi)=\sum_{\mathbf x\in\prod_{z,a<m}\mathcal R^z_a}\ \sum_{h\in\mathcal
R_Y}\ \sum_{S\in\mathcal S_\vartheta(\pgrlast)}\Psi_\vartheta(S)\,B\bigl(v^0_S(\mathbf x,h),
v^1_S(\mathbf x,h)\bigr).
\]
\end{proposition}
\begin{proof}
\emph{Unfolding.} For fixed $S$, the composite $\pgwT^z_S\psi^z$ evaluated at $x^z_m$ is, by the
definition of the algebraic lifts in \S\ref{pgw:ss-last}, the iterated integral
$\pgwint{J^z_{m-1}}\cdots\pgwint{J^z_0}$ of \eqref{pgz:eq-reached} over $x^z_{m-1},\dots,x^z_0$ (each
integral is a finite sum, so the order is immaterial). As $B$ is bilinear and all sums are finite,
\[
\Lambda=\pgwint{J_Y}_{\Gamma_Y\backslash G_{Y,f}}\ \pgwint{J^\bullet}\ F(\mathbf x,h)\,d\mathbf x\,dh,
\qquad F(\mathbf x,h)=\sum_{S\in\mathcal S_\vartheta(\pgrlast)}\Psi_\vartheta(S)\,B\bigl(v^0_S(\mathbf
x,h),v^1_S(\mathbf x,h)\bigr),
\]
where $\pgwint{J^\bullet}$ is the iterated class integral over the $2m$ variables $x^z_a$ ($a<m$), each
for the measure giving $J^z_a$ the volume $1$. Each function of one variable obtained by fixing the
others is left invariant under the corresponding $\Gamma$ (Lemma~\ref{pgw:one}(a), the $\Gamma$-equivariance
of $\psi^z$, and for $h$ also Lemma~\ref{pgw:branch}(b): $B(\eta(\gamma h),\varphi((\gamma
h)\oplus\pgupt^m))=B(\eta(h),\varphi(h\oplus\pgupt^m))$ by the invariance of $B$ and
$(\gamma h)\oplus\pgupt^m=(\gamma\oplus1)(h\oplus\pgupt^m)$), and right invariant under an open
subgroup.

\emph{Invariance.} Let $k^z_a\in J^z_a$ for $a<m$, $k_Y\in J_Y$, and put $k^0_m=k_Y$,
$k^1_m=k_Y\oplus1$ (an element of $J^1_m=J(\pcb)$ of $\pgrY^m_m$). We claim
\begin{equation}\label{pgz:eq-inv}
F\bigl((x^z_ak^z_a),hk_Y\bigr)=F(\mathbf x,h).
\end{equation}
Note $x^0_mk^0_m=hk_Y$ and $x^1_mk^1_m=(h\oplus\pgupt^m)(k_Y\oplus1)=hk_Y\oplus\pgupt^m$, so the left
side is formed with the points $x^z_ak^z_a$ for all $a\le m$. Let $g$ be the node family with
$g_{p^z_a}=\bar k^z_a$ (reduction modulo $\pcm$) for $a\le m$ and $g_p=1$ at the other nodes. It is
admissible (\S\ref{pgv:ss-def}): $g_{(0,0,m)}=1$ because $k^0_0\in J(\pcm)$; $g_{(1,0,m)}\equiv1
\bmod\pcb$ because $k^1_0\in J(\pcb)$; $g_{(1,m,m)}=g_{(0,m,m)}\oplus1$; and at the nodes $(z,a,a)$
with $a<m$, which are not on $\pgrlast$, both sides are $1$. By the outer invariance
$\psi^z(x^z_0k^z_0)=\psi^z(x^z_0)$, and by Lemma~\ref{pgz:move} applied to each factor of
\eqref{pgz:eq-reached},
\[
v^z_S\bigl((x^{z'}_ak^{z'}_a),hk_Y\bigr)=\prod_{a=1}^m\chi_{e^z_a}\Bigl(\frac{\det g_{p^z_a}}{\det
g_{p^z_{a-1}}}\Bigr)\cdot v^z_{g^{-1}S}(\mathbf x,h),
\]
where $(g^{-1}S)_e=g_qS_eg_p^{-1}$ for $e=(p,q)$. The product of the character factors over both
halves is $\chi_{\pgrE(\pgrlast)}(g)$. So the left side of \eqref{pgz:eq-inv} is
$\sum_S\Psi_\vartheta(S)\chi_{\pgrE(\pgrlast)}(g)B(v^0_{g^{-1}S},v^1_{g^{-1}S})(\mathbf x,h)$. The
map $S\mapsto g^{-1}S$ is a bijection of $\mathcal S_\vartheta(\pgrlast)$; substituting $S=gS'$ and
using $\chi_{\pgrE(\pgrlast)}(g)\Psi_\vartheta(gS')=\Psi_\vartheta(S')$ (Definition~\ref{pgr:def-datum})
gives $F(\mathbf x,h)$.

\emph{Conclusion.} By \eqref{pgz:eq-inv} with all but one of the $k$ equal to $1$, $F$ is right
invariant under $J^z_a$ in the variable $x^z_a$ and under $J_Y$ in $h$, the other variables being
fixed. Applying Lemma~\ref{pgz:free} to each of the $2m+1$ class integrals in turn replaces it by the
sum over the representatives, with coefficient $1$.
\end{proof}

\begin{remark}[The denominators of the class sums]\label{pgz:rem-den}
Proposition~\ref{pgz:expand} settles the first point of Proposition~\ref{pgr:den}: in $\Lambda$ every
class sum is normalised by the level over which it is taken, the levels are free
(Lemma~\ref{pgz:free}), and the joint invariance \eqref{pgz:eq-inv} makes each class sum a sum with
coefficients $1$. The indices $[J:J']$ of auxiliary subgroups fixing the individual terms, and the
orders of stabilisers, cancel or are $1$; they contribute nothing to $N_0$, and no condition relating
them to $\ell(\iota)$ is needed. The indices $[K_p:J_p]$ occur only in the comparison of $\Lambda$ with
the periods (the number $D(\vartheta)$ in the proof of Theorem~\ref{pgw:last}), not in $\Lambda$.
\end{remark}

\begin{lemma}[Counting the terms]\label{pgz:count}
The number of triples $(\mathbf x,h,S)$ in Proposition~\ref{pgz:expand} is
$\#\mathcal S_\vartheta(\pgrlast)\cdot\#(\Gamma_Y\backslash G_{Y,f}/J_Y)\prod_{z,\,a<m}
\#(\Gamma^z_a\backslash G^z_{a,f}/J^z_a)$, and $\#\mathcal S_\vartheta(\pgrlast)=\prod_{z,a}
\#\pcMat_{N^z_a\times N^z_{a-1}}(\pcR_\pcm)$.
\end{lemma}
\begin{proof}
An edge tuple vanishing off $\pgrE(\pgrlast)$ is a choice of one matrix for each of the $2m$ edges
$e^z_a$.
\end{proof}

\subsubsection{One lift does not increase gauges after multiplication by $\pgzd$}\label{pgz:ss-lift}

Put
\begin{equation}\label{pgz:eq-d}
\pgzd=M^{4bfn_m^2}\,(dM^b)^{40fn_mb},
\end{equation}
$d$ the discriminant of $K$ (a positive integer; $40fn_mb$ is an integer).

\begin{lemma}[The edge data are standard]\label{pgz:std}
Let $\iota\in\auI$ with $\ell=\ell(\iota)\nmid\auNm Q_\vartheta$, $(p,q)=e^z_a$, $X\in\pcMat_{q\times
p}(\pcR_\pcm)$, and let $x\in G_{p,f}$, $h\in G_{q,f}$ be integral with integral inverses at the places
above $\ell$. Then $c_{\mathrm{vac}}^{\,N_pN_q}M^{bfN_pN_q}\,\omega_f(x,h)[\delta_X]_f$ is an
$\iota$-unit times a standard datum of level $M^b$ at $\ell$ ({\cite[\S7.4]{P49}}).
\end{lemma}
\begin{proof}
\emph{The datum $[\delta_X]_f$.} Outside $S(\pcm)$ it is the vacuum. At $v=w\aubar w\in S(\pcm)$,
$w\mid\pcm$, it is $\delta_X^\natural$, the indicator of $a_w(\tilde X+\pcm M_w^\circ)$ with
$M_w^\circ=\pcMat_{q\times p}(\auO_w)$ and $\tilde X$ an integral lift of $X$ (Definition~\ref{pgx:conv}).
By Lemma~\ref{pgx:vac} the vacuum at $v$ is the indicator of $M_w=a_wM^\circ_w$. As in the proof of
{\cite[Lemma~4.10]{P54}} (paragraph ``the datum as a combination of translates of the vacuum'', with the
lattice $M_w$ in place of the integral matrices; only the self-duality of $M_v=M_w\oplus M_{\aubar w}$ is
used there), the indicator of $a_w\tilde X+\pcm M_w$ is a root of unity times $\psrho^\circ(a_w\tilde
X)\mathbf 1_{\pcm M_w}$, and $\mathbf 1_{\pcm M_w}=[M_w:\pcm M_w]^{-1}\sum_y\zeta_y\psrho^\circ(y)
\mathbf 1_{M_w}$, with $[M_w:\pcm M_w]$ terms, roots of unity $\zeta_y$, and $y$ in the
$\aubar w$-component with $y\in\aubar\pcm^{-1}M_{\aubar w}$. Taking the product over $v\in S(\pcm)$ and
writing $\psvac=c_{\mathrm{vac}}^{-N_pN_q}\psvac^\circ$ {\cite[(58)]{P49}},
\[
c_{\mathrm{vac}}^{\,N_pN_q}\,\auNm(\pcm)^{N_pN_q}\,[\delta_X]_f=\sum_{x'}\zeta_{x'}\,\psrho^\circ(x')
\psvac^\circ ,
\]
a finite sum with roots of unity $\zeta_{x'}$ and translations $x'$ that vanish outside $S(\pcm)$ and
at $v\in S(\pcm)$ have components $a_w\tilde X\in M_w$ and $y\in\aubar\pcm^{-1}M_{\aubar w}$.

\emph{Level $M^b$ at $\ell$.} Let $w\mid\pcm$ lie above $\ell$. As $\ell\nmid\auNm Q_\vartheta$, $w$ is
not in $Q_\vartheta$, so $w\mid\pi_M$; $\ell\nmid q_1$ (since $\ell\mid M$ and $p_0,p_1\equiv1\pmod
{q_1}$), so $a_w=1$ (Proposition~\ref{pgv:alpha}(a) and the convention after Definition~\ref{pgx:conv}),
$M_w=M^\circ_w$, and $M_{\aubar w}$ is the integral lattice at $\aubar w$. Then $M^b\in\pcm_w\cap
\aubar\pcm_{\aubar w}$ ($\pcm_w=\pi_M^b\auO_w$ and $\aubar{\pi_M}^b\mid M^b$), so $M^bx'$ is integral at
the places above $\ell$. Hence the right side is a standard datum of level $M^b$ at $\ell$, with
coefficients roots of unity.

\emph{Norms.} $\auNm(\pcm)=\auNm(\pcb)\auNm Q_\vartheta$ (the primes of $\pi_M$ are not in $Q_\vartheta$),
$\auNm Q_\vartheta$ is prime to $\ell$, and $\auNm(\pcb)=\auNm_{E/\Q}(\pi_M)^b=M^{bf}$ ($\pi_M\in K$,
$[E:K]=f$). So $c_{\mathrm{vac}}^{N_pN_q}M^{bfN_pN_q}[\delta_X]_f$ is an $\iota$-unit times a standard
datum of level $M^b$.

\emph{The action of $(x,h)$.} By the remark preceding {\cite[Theorem~7.15]{P49}}, $\omega(x,h)$ preserves
the standard data of level $M^b$ at $\ell$; $\omega_f=\omega_{\alpha,\beta,f}$ differs from it by
$\alpha(\det x)^{-1}\beta(\det h)$ and the unramified factors of {\cite[(1)]{P54}}, which are
$\iota$-units for $x,h$ integral with integral inverses above $\ell$ (values of algebraic Hecke
characters at ideles that are units above $\ell$, and powers of residue characteristics different from
$\ell$), exactly as in the proof of {\cite[Lemma~4.10(KS1)]{P54}}.
\end{proof}

\begin{lemma}[One lift]\label{pgz:kernel}
Let $\iota\in\auI$ with $\ell(\iota)\nmid\auNm Q_\vartheta$, $(p,q)=e^z_a$, $X\in\pcMat_{q\times
p}(\pcR_\pcm)$, $x\in G_{p,f}$ and $h\in G_{q,f}$. If $\xi\in W^z_{a-1}$ has gauge at most $1$ at
$(\iota,x)$, then $\pgzd\cdot\pgwK^z_a(\omega_f(x,h)[\delta_X]_f)\,\xi$ has gauge at most $1$ at
$(\iota,h)$.
\end{lemma}
\begin{proof}
\emph{Reduction to integral points.} Unitary groups are rational varieties and satisfy weak
approximation~\cite[\S7.1, Propositions~7.3 and~7.4]{PlatonovRapinchuk}, so there are rational points $\gamma\in\Gamma_p$,
$\gamma'\in\Gamma_q$ such that $k=\gamma^{-1}x$ and $k'=\gamma'^{-1}h$ are integral with integral
inverses at the places above $\ell$ (as in the proof of {\cite[Proposition~5.30]{P49}}). By
Lemma~\ref{pgw:one}(a), $\pgwK^z_a(\omega_f(x,h)\Phi)=\gamma'\circ\pgwK^z_a(\omega_f(k,k')\Phi)\circ
\gamma^{-1}$, the rational points acting on $W^z_{a-1}$, $W^z_a$ by the archimedean action. By (G-b)
of \S\ref{pgz:ss-choice}, with $\eta=\pgwK^z_a(\omega_f(x,h)\Phi)\xi$,
\[
h_\ell^{-1}\iota(\eta)=k'^{-1}_\ell\,\iota\bigl(\pgwK^z_a(\omega_f(k,k')\Phi)\bigr)\,k_\ell\,
\bigl(x_\ell^{-1}\iota(\xi)\bigr).
\]
By (G-a), $k_\ell$ and $k'^{-1}_\ell$ are matrices invertible over $\auO_{\C_\ell}$, and
$x_\ell^{-1}\iota(\xi)$ has integral coordinates. So it suffices to show that the matrix of
$\pgzd\,\pgwK^z_a(\omega_f(k,k')[\delta_X]_f)$ in the bases $(e_T)$ has entries in $\auO_\iota$.

\emph{Integrality.} By Lemma~\ref{pgz:std}, $c_{\mathrm{vac}}^{N_pN_q}\omega_f(k,k')[\delta_X]_f=
u\,M^{-bfN_pN_q}f'$ with an $\iota$-unit $u$ and a standard datum $f'$ of level $M^b$ at $\ell$. By
Definition~\ref{pgz:def-choice}(a) and the homogeneity of the entries of $S_\nu$ (of total degree
$e=e(\nu)$),
\[
\pgwK^z_a(\omega_f(k,k')[\delta_X]_f)=u\,M^{-bfN_pN_q}(dM^b)^{-e}\,S_\nu(\partial)\bigl[\Theta_{f'}
(dM^bs_0\,\cdot)\bigr](0).
\]
The series $\Theta_{f'}(dM^bs_0\,\cdot)$ is $\partial$-integral at $\iota$: for $\ell\ne d$ by the last
assertion of {\cite[Theorem~7.15]{P49}} with $\delta=M^b$ ($\aubar\delta=\delta$) and $c=d$; for $\ell=d$ by
the first assertion of {\cite[Theorem~7.15]{P49}} and the argument of the proof of {\cite[Lemma~4.10(KS1)]{P54}}
for $\ell=d$ (there written for $\delta=M^2$; it uses only the product formula of
{\cite[Theorem~7.15]{P49}}, the $\partial$-integrality of $y$ at every prime, and that of $\theta(d\,\cdot)$).
The entries of $S_\nu$ are polynomials with integer coefficients, so the jets are in $\auO_\iota$.
Finally $N_pN_q\le4n_m^2$ (Lemma~\ref{pgs:boxes}) and $e\le40fn_mb$ (Lemma~\ref{pgz:deg}), so
$M^{bfN_pN_q}(dM^b)^e$ divides $\pgzd$.
\end{proof}

\begin{corollary}[Vectors reached along a half]\label{pgz:reach}
Let $\iota$ be as in Lemma~\ref{pgz:kernel} and let $\psi^z$ be integral at $\iota$. Then for every
$S$, $\mathbf x$ and $h$, $\pgzd^mv^0_S(\mathbf x,h)$ has gauge at most $1$ at $(\iota,h)$ and
$\pgzd^mv^1_S(\mathbf x,h)$ has gauge at most $1$ at $(\iota,h\oplus\pgupt^m)$.
\end{corollary}
\begin{proof}
Induction on the number of factors in \eqref{pgz:eq-reached}: $\psi^z(x^z_0)$ has gauge at most $1$
at $x^z_0$, and Lemma~\ref{pgz:kernel} (gauges scale with the absolute value of the scalar) passes from
$x^z_{a-1}$ to $x^z_a$, the last point being $x^z_m$.
\end{proof}

\subsubsection{The branching pairing at the point $h\oplus\pgupt^m$}\label{pgz:ss-pair}

\begin{lemma}[The idele $\alpha$ and the flip scalar]\label{pgz:flip}
(a) $\alpha^{-1}$ and $q_1\alpha$ are integral finite ideles of $E$.
(b) $\mathsf a=\auNm Q_\vartheta\cdot M^b\cdot q_1\cdot\pgupt$ and $\aubar{\mathsf a}=\auNm Q_\vartheta
\cdot M^b\cdot q_1\cdot\pgupt^{-1}$ are integral finite ideles of $E$.
\end{lemma}
\begin{proof}
(a) Let $P\in\auO_K[x]$ be the minimal polynomial of $\zeta=\zeta_{q_1}$ over $K$ (\S\ref{pgv:ss-G},
before Proposition~\ref{pgv:alpha}). It is a monic divisor of $x^{q_1}-1$, so $x^{q_1}-1=PG$ with $G$
monic; $G$ has coefficients in $\auO_K$, being a monic factor of a monic polynomial with integral
coefficients. Differentiating at $\zeta$: $q_1\zeta^{q_1-1}=P'(\zeta)G(\zeta)$, so $q_1/P'(\zeta)=
\zeta^{-1}G(\zeta)\in\auO_E$. Since $\alpha$ is $1/P'(\zeta)$ or $1$ at every prime, $\alpha^{-1}$ and
$q_1\alpha$ are integral.

(b) $\pgupt=\aubar u_0/u_0$ with $u_0=\alpha\varpi_\pcm$, so $\mathsf a=(\auNm Q_\vartheta M^b
\varpi_\pcm^{-1})\cdot\aubar{\varpi_\pcm}\cdot\aubar{q_1\alpha}\cdot\alpha^{-1}$. The integer
$\auNm Q_\vartheta M^b$ lies in $\pcm$ ($\pi_M^b\mid M^b$, and $\auNm w\in w$ for each $w\in
Q_\vartheta$), and $\varpi_\pcm$ generates $\pcm\widehat\auO_E$; so the first factor is integral. The
others are integral by (a) and since complex conjugation preserves integrality. As $\pgupt\aubar\pgupt
=1$, $\aubar{\mathsf a}=\auNm Q_\vartheta M^bq_1\pgupt^{-1}$, integral as the conjugate of $\mathsf a$.
\end{proof}

\begin{lemma}[The character of the branching form]\label{pgz:char}
Suppose $B\ne0$. There are integers $L_\tau$ $(\tau\in\auPl)$ with $|L_\tau|\le30b$ such that, as
polynomial identities in $g\in\prod_\tau\auGL_n$ and $t\in\prod_\tau\auGL_1$ (over any field of
characteristic $0$, through the integral representation matrices),
\[
B\bigl(\rho_0(g)w_0,\rho_1(g\oplus t)w_1\bigr)=\prod_\tau t_\tau^{L_\tau}\,B(w_0,w_1),
\]
$\rho_0,\rho_1$ the representations of $\prod_\tau\auGL_n$, $\prod_\tau\auGL_{n+1}$ on $W^0_m$,
$W^1_m$ in the bases $(e_T)$.
\end{lemma}
\begin{proof}
For $t$ in the compact group $\prod_\tau\auU(1)$ the form $(w_0,w_1)\mapsto B(w_0,\rho_1(1\oplus t)w_1)$
is invariant in the sense of Lemma~\ref{pgw:branch}(a), since $1\oplus t$ commutes with $g\oplus1$; by
the uniqueness in that lemma it is $\chi(t)B$ for a continuous character $\chi$ of
$\prod_\tau\auU(1)$, so $\chi(t)=\prod_\tau t_\tau^{L_\tau}$ with $L_\tau\in\Z$. Together with the
invariance, $B(\rho_0(g)w_0,\rho_1(g\oplus t)w_1)=\chi(t)B(w_0,w_1)$ on $\prod\auU(n)\times\prod\auU(1)$;
both sides are polynomials in the entries of $g$, $t$ and in $\det g^{-1},t^{-1}$, with rational
coefficients, so by Zariski density the identity holds on $\prod\auGL_n\times\prod\auGL_1$ and then over
every field of characteristic $0$. For the bound, decompose $W^1_m$ into the eigenspaces of
$\rho_1(1\oplus t)$, $t\in\prod\auGL_1$; on each the action is by $\prod_\tau t_\tau^{k_\tau}$, with
$k_\tau$ the last coordinate at $\tau$ of a weight of $W^1_m$. Since $B\ne0$, it is non-zero on some
eigenspace, and there $\chi(t)=\prod t_\tau^{k_\tau}$, so $L_\tau=k_\tau$. The coordinates of the
weights of an irreducible representation lie between the smallest and the largest entry of its highest
weight; by the proof of Lemma~\ref{pgs:boxes} these are at most $|e_\tau|+20b\le30b$ in absolute value.
\end{proof}

Put
\begin{equation}\label{pgz:eq-M}
\pgzM=\prod_{\tau\in\auPl}(M^bq_1)^{m|L_\tau|}\qquad(\pgzM=1\text{ if }B=0).
\end{equation}

\begin{lemma}[Clearing the branching pairing]\label{pgz:pair}
Let $\iota\in\auI$ with $\ell=\ell(\iota)\nmid\auNm Q_\vartheta$, $h\in G_{Y,f}$, and let $v'\in W^0_m$
have gauge at most $1$ at $(\iota,h)$ and $v\in W^1_m$ gauge at most $1$ at $(\iota,h\oplus\pgupt^m)$.
Then $\pgzM\,B(v',v)\in\auO_\iota$.
\end{lemma}
\begin{proof}
We may assume $B\ne0$. Let $w'=h_\ell^{-1}\iota(v')$ and $w=(h\oplus\pgupt^m)_\ell^{-1}\iota(v)$, vectors
with coordinates in $\auO_{\C_\ell}$. Here $(h\oplus\pgupt^m)_\ell=(h_\ell\oplus1)(1\oplus\pgupt^m_\ell)$,
and $\pgupt_\ell$ acts through its images $t_\tau\in\C_\ell^\times$ at the primes $w_\tau$ of
{\cite[\S14.3]{P46}}. By Lemma~\ref{pgz:char} over $\C_\ell$ (applied with $g=h_\ell$, $t=\pgupt^m_\ell$),
\[
\iota(B(v',v))=B\bigl(h_\ell w',(h_\ell\oplus\pgupt^m_\ell)w\bigr)=\prod_\tau t_\tau^{mL_\tau}\cdot
B(w',w).
\]
$B$ has integer entries, so $B(w',w)\in\auO_{\C_\ell}$. For each $\tau$, $(M^bq_1)^{m|L_\tau|}
t_\tau^{mL_\tau}$ is $(\auNm Q_\vartheta)^{-m|L_\tau|}$ times the image of $\mathsf a^{m|L_\tau|}$ or of
$\aubar{\mathsf a}^{m|L_\tau|}$ (according to the sign of $L_\tau$), which is integral by
Lemma~\ref{pgz:flip}(b); and $\auNm Q_\vartheta$ is an $\ell$-adic unit. Multiplying over $\tau$ gives
the claim.
\end{proof}

\subsubsection{Proof of Proposition~\ref{pgr:den}(a)}\label{pgz:ss-a}

\begin{theorem}[Proposition~\ref{pgr:den}(a)]\label{pgz:a}
Let $N_0=\pgzM\,\pgzd^{2m}$, with $\pgzd$, $\pgzM$ of \eqref{pgz:eq-d}, \eqref{pgz:eq-M}. Then
$1\le N_0\le\exp(x^{99}b^2)$, $N_0$ does not depend on $\vartheta$, $\Psi$, $y'$, $\varphi$, and the
assertion (a) of Proposition~\ref{pgr:den} holds for $\Lambda$ of \eqref{pgw:eq-Lambda} with the
choices of Definition~\ref{pgz:def-choice}.
\end{theorem}
\begin{proof}
\emph{Integrality.} Let $\Psi_\vartheta$, $(y',\varphi)$ and $\iota$ be as in (a). By
Proposition~\ref{pgz:expand}, $\Lambda$ is a finite sum of terms $\Psi_\vartheta(S)B(v^0_S,v^1_S)$.
$\Psi_\vartheta(S)$ is a sum of roots of unity, hence in $\auO_\iota$. By Corollary~\ref{pgz:reach},
$\pgzd^mv^0_S$ and $\pgzd^mv^1_S$ have gauge at most $1$ at $h$ and $h\oplus\pgupt^m$; by
Lemma~\ref{pgz:pair}, $\pgzM\pgzd^{2m}B(v^0_S,v^1_S)=\pgzM B(\pgzd^mv^0_S,\pgzd^mv^1_S)\in\auO_\iota$.

\emph{Independence.} $\pgzd$ depends only on $M,b,f,n_m,d$. The integers $L_\tau$ depend only on the
weights at $p^0_m$, $p^1_m$ (Lemma~\ref{pgz:char}), not on $\vartheta$.

\emph{Size.} By {\cite[Lemma~13.4]{P46}}: $\log M\le22\log b$ (S5), $f\le x^{12}$ (S3), $d\le\log x$
((S1)), $q_1\le x^{12}$ (S3), $m\le2\sqrt x$ (D3), and $n_m\le40x^6$ ({\cite[Lemma~8.8(b)]{P49}}). So
\[
\log\pgzd=4bfn_m^2\log M+40fn_mb(\log d+b\log M)\le4\cdot1600\,x^{24}\cdot22\,b\log b+1600\,x^{18}b
(\log x+22\,b\log b)\le x^{26}b^2,
\]
using $\log b\le x^{1/100}$ (S1). By Lemma~\ref{pgz:char}, $\log\pgzM\le f\cdot m\cdot30b\cdot(b\log M+
\log q_1)\le x^{12}\cdot2\sqrt x\cdot30b\cdot23\,b\log b\le x^{14}b^2$. Hence $\log N_0\le
x^{14}b^2+4\sqrt x\,x^{26}b^2\le x^{99}b^2$.
\end{proof}

\subsubsection{Proposition~\ref{pgr:den}(b): reduction to three size bounds}\label{pgz:ss-b}

Let $\vartheta=1$, so $\pcm=\pcb$, $\#\pcR_\pcb=\auNm\pcb=M^{bf}$, and put $G_0=\exp(x^{99}b^2)$.
Part (b) follows from Proposition~\ref{pgz:expand} and the following three bounds, which are proved
in \S\ref{pgt:sec}: (B1) is Proposition~\ref{pgt:B1}, (B2) is Proposition~\ref{pgt:B2} and (B3) is
Corollary~\ref{pgt:B3}.

\begin{enumerate}
\item[(B1)] (One lift.) For every edge $e^z_a=(p,q)$, $X\in\pcMat_{q\times p}(\pcR_\pcb)$, $x\in
G_{p,f}$, $h\in G_{q,f}$, $\sigma\in\auAut(\C)$ and $\xi\in W^z_{a-1}$:
$\|\sigma(\pgwK^z_a(\omega_f(x,h)[\delta_X]_f)\xi)\|\le G_0\|\sigma(\xi)\|$, with the kernels of
Definition~\ref{pgz:def-choice}(a), $\sigma$ acting on coordinates in the bases $(e_T)$ and
$\|\ \|$ the invariant norms.
\item[(B2)] (The branching form.) If a non-zero invariant form exists, there is a non-zero invariant integer matrix $B'$ (in the bases
$(e_T)$) with $|B'(v',v)|\le C\|v'\|\,\|v\|$ for all $v',v$, where $C=(L^L)^3$ and $L\ge1$ is an integer
bound for the number of boxes of the weight at $p^1_m$; by Lemma~\ref{pgs:boxes} one may take
$L=\lfloor x^{19}b\rfloor$.
\item[(B3)] (Classes.) For $z\in\{0,1\}$ and $a\le m$, $\#(\Gamma^z_a\backslash G^z_{a,f}/J^z_a)\le G_0$
when $\vartheta=1$.
\end{enumerate}

\noindent Given (B2), the matrix $B$ of Definition~\ref{pgz:def-choice}(b) satisfies the same bound:
by uniqueness $B'=cB$ with $c$ rational, and as $B$ has coprime integer entries and $B'$ integer entries,
$c$ is a non-zero integer, so $|B(v',v)|\le|B'(v',v)|$.

\begin{theorem}[Proposition~\ref{pgr:den}(b)]\label{pgz:b}
Let $\Psi_1$ be a path datum on $\pgrE(\pgrlast)$ with bound $L(1)$, let
$(y',\varphi)$ be fixed by the outer levels, and $\|\sigma(\psi_{y'}(g))\|^2\le B_0$,
$\|\sigma(\psi_\varphi(g))\|^2\le B_1$ for all $g$. Then $|\sigma(\Lambda(\Psi_1;y',\varphi))|\le
\exp(x^{200}b^2)\sqrt{B_0B_1}$.
\end{theorem}
\begin{proof}
The bounds (B1), (B2), (B3) hold by Propositions~\ref{pgt:B1}, \ref{pgt:B2} (if no non-zero invariant
form exists, $B=0$ and the bound is trivial, Remark~\ref{pgt:rem-B2}) and Corollary~\ref{pgt:B3}.
By Proposition~\ref{pgz:expand}, $\sigma(\Lambda)=\sum\sigma(\Psi_1(S))\,B(\sigma v^0_S,\sigma v^1_S)$, as
$B$ has rational entries. By (B1) and induction along \eqref{pgz:eq-reached},
$\|\sigma v^0_S\|\le G_0^m\sqrt{B_0}$ and $\|\sigma v^1_S\|\le G_0^m\sqrt{B_1}$; by (B2) each pairing is
at most $C\,G_0^{2m}\sqrt{B_0B_1}$; $|\sigma\Psi_1(S)|\le L(1)$. By Lemma~\ref{pgz:count} and (B3) the
number of terms is at most $\#\mathcal S_1(\pgrlast)\,G_0^{2m+1}$, and $\#\mathcal S_1(\pgrlast)\le
M^{bf\cdot2m(2n_m)^2}$ (Lemma~\ref{pgs:boxes}: ranks $\le2n_m$). Also $L(1)=V(1)^{m^2}\mathcal
E(1)^{2m}=M^{bf(2n_mm^2+2m(2n_m)^2)}$. With the bounds of {\cite[Lemma~13.4]{P46}} used in the proof of
Theorem~\ref{pgz:a},
\[
\log\bigl(L(1)\,\#\mathcal S_1(\pgrlast)\bigr)\le bf\log M\,(2n_mm^2+4m(2n_m)^2)\le x^{12}\cdot22\,b\log b
\cdot(320x^7+51200\,x^{12.5})\le x^{28}b^2,
\]
$\log C=3L\log L\le x^{20}b^2$, and $(4m+1)x^{99}b^2\le x^{100}b^2$. The total is at most
$x^{28}b^2+x^{20}b^2+x^{100}b^2\le x^{200}b^2$.
\end{proof}






% ===== C-grid9.tex =====
% proofs/C-grid9.tex
% Input (V) of proofs/C-grid5.tex (deep-level vanishing by Poincare series).
% Labels start with pgq:. Conventions and macros: those of proofs/C-grid5.tex and C-grid4.tex.
\providecommand{\auU}{\mathrm U}
\providecommand{\auO}{\mathcal O}
\providecommand{\auA}{\mathbb A}
\providecommand{\auAz}{\mathcal A_0}
\providecommand{\auPl}{\Phi}
\providecommand{\aubar}[1]{\overline{#1}}
\providecommand{\pcb}{\mathfrak b}
\providecommand{\pgrY}{\widetilde Y}
\providecommand{\pcnot}[1]{\par\smallskip{\raggedright\noindent\textsc{Not proved here.} [#1]\par}\smallskip}
\providecommand{\pgqvol}{\operatorname{vol}}
\providecommand{\pgqcol}{\operatorname{col}}
\providecommand{\pgqSl}{\mathsf{Sl}}

\subsection{Deep-level vanishing}\label{pgq:sec}

This section proves input (V) of \S\ref{pgw:ss-first} (Theorem~\ref{pgq:V}), from three
inputs (C), (Q), (S) stated in \S\ref{pgq:ss-inputs} and proved in \S\S\ref{pgk:sec}, \ref{pgl:sec}. The argument halves every bound of
$|y|$: high in the cusp by the cusp property of $y$ and the slow variation of its reproducing kernel
(\S\ref{pgq:ss-cusp}); elsewhere by a Poincar\'e series of deep level whose kernel is a translated
matrix coefficient of the type of $\pgrY^0_0$. Its period against $y$ vanishes by the hypothesis of
(V), and by unfolding it equals a smoothing of $y$ up to an error that is small because the level
$\pcb=(\pi_M)^b$ is deep (\S\ref{pgq:ss-poinc}); averaging the smoothing against the kernel of $y$
returns $y$ times a product of local pairings, which is not small at the chosen point $g_0$
(\S\ref{pgq:ss-bulk}).

\subsubsection{Setting and statement}\label{pgq:ss-stat}

\emph{Spaces and groups.} Let $V=Y_0$, of rank $n=r$, and $Z=\pgrY^0_0=V\oplus E\tilde e$ with
$\langle\tilde e,\tilde e\rangle=1$ (Definition~\ref{pgv:def-spaces}); the lattice of $Z$ is
$L_Z=L_V\oplus\auO_E\tilde e$. Put $G=\auU(V)$, $H=\auU(Z)$, embed $G\to H$ by $h\mapsto h\oplus1$,
and write $\Gamma_G=G(E^+)$, $\Gamma_H=H(E^+)$, $[G]=\Gamma_G\backslash G(\auA)$. Let $f=\#\auPl$.
For $\tau\in\auPl$ we use on $V_\tau=V\otimes_{E,\tau}\C$ the normalised basis of
{\cite[\S5.1]{P51}}, in which the hermitian form has a hermitian unitary matrix $\mathsf J_\tau$, and
on $Z_\tau$ that basis followed by $\tilde e$ (matrix $\mathsf J_\tau\oplus1$). Norms of vectors are
the Euclidean norms in these bases and $\|g\|$ is the operator norm; thus $\|g\|=\|g^{-1}\|\ge1$ for
$g$ in $G_\tau$ or $H_\tau$ (as $g^{-1}=\mathsf Jg^*\mathsf J$), and $\|r\oplus1\|=\|r\|$. For
$g\in G_\infty=\prod_\tau G_\tau$ put $\|g\|=\prod_\tau\|g_\tau\|$ and $\operatorname{dist}(g,1)=
\max_\tau\|g_\tau-1\|$; likewise on $H_\infty$. We write $G_\infty\subset G(\auA)$ and use that real
points commute with finite ones.

\emph{Measures.} $G_\tau$ carries the Haar measure of \textup{(DS1)(ii)} ({\cite[Proposition~5.2]{P51}}),
$G_\infty$ the product measure, $G(\auA_f)$ any Haar measure, $G(\auA)$ the product, and $[G]$ the
quotient measure, so that
\begin{equation}\label{pgq:eq-unfold}
\int_{[G]}\sum_{\delta\in\Gamma_G}F(\delta h)\,dh=\int_{G(\auA)}F(h)\,dh
\end{equation}
for $F\ge0$ measurable, and for $F\in L^1(G(\auA))$ (the inner sum then converging absolutely for
almost all $h$). Replacing the measure of $[G]$ by a positive multiple changes neither the hypothesis
nor the conclusion of (V). Put $V_\tau=\pgqvol\{r\in G_\tau:\|r-1\|\le1\}>0$.

\emph{Coordinates and the first minimum.} For $x\in G(\auA)$ and $v\in V(E)$ write $x^{-1}v\in
V(\auA)$ in the frame; $v$ lies in the lattice $L_V(x)$ if the finite coordinates of $x^{-1}v$ are
integral, and $|v|_x$ is the Euclidean norm of its archimedean coordinates in $\prod_\tau\C^n$.
Likewise $L_Z(x')$, $|\cdot|_{x'}$ for $x'\in H(\auA)$. The \emph{first minimum} is
\[
\lambda(x)=\min\bigl(1,\ \min\{|v|_x:0\ne v\in L_V(x)\}\bigr)\in(0,1].
\]

\emph{Types.} Let $c=\prod_\tau c_\tau$ and $e=d\,\aubar c$ ($d=\prod_\tau d_\tau$) be the coefficient
and the type of $\auAz$ ({\cite[\S14.3]{P46}}), $c_\tau=\langle v_\tau,\pi_\tau(\cdot)v_\tau\rangle$; and
let $c'=\prod_\tau c'_\tau$, $e'=d'\aubar{c'}$ be those of the type $e_{(1,0,0)}$ of $Z$
(\S\ref{pgv:ss-def}: Harish-Chandra parameter $-\lambda_{(1,0,0),\tau}$, signs $\epsilon_{(1,0,0)}$). The
forms of type $e$ on $G$ ({\cite[Definition~8.1]{P49}}) are bounded, continuous, cuspidal
({\cite[Proposition~8.4(a)]{P49}}). Let $J_G(\mathfrak a)$, $J_H(\mathfrak a)$ be the principal
congruence subgroups of level $\mathfrak a$ in the stabilisers of $\widehat L_V$, $\widehat L_Z$. Then
$J_M=J_G(\pi_M)\supseteq J_G(\pcb)$ and $J_M\supseteq J_G(M)$, as $\pi_M$ divides $\pcb=\pi_M^b$ and
$M=\pi_M\aubar{\pi_M}$ ({\cite[\S14.3]{P46}}); and $h\oplus1\in J_H(\mathfrak a)$ if and only if
$h\in J_G(\mathfrak a)$. Since $\pi_M\in\auO_K$, $|\tau(\pi_M)|^2=M$ for every $\tau$.

\begin{theorem}[Input (V)]\label{pgq:V}
There is a real point
$g_0=\omega$ of $H$ such that every $y\in\auAz$ with
\[
\int_{[G]}\aubar{y(x)}\,\psi\bigl((x\oplus1)\omega\bigr)\,dx=0
\]
for all forms $\psi$ of type $e'$ on $H$ fixed by $J_H(\pcb)$ is zero.
\end{theorem}

\subsubsection{Three inputs}\label{pgq:ss-inputs}

Let $N_d=q_1d$, a natural number in the different of $E$ ({\cite[\S5.1.5]{P54}}), and
$\kappa_{\mathrm c}=8fn\,\mathrm{rd}_E+48n^6f^2N_d+1$.

\begin{enumerate}
\item[(C)] (\emph{Cusp estimate.}) Let $F$ be a form of type $e$ on $G$ fixed by $J_G(M)$, let
$0<\delta\le1$, and let $x\in G(\auA)$ satisfy $\lambda(x)<(\delta/(M\kappa_{\mathrm c}^2))^n$. If
$|F(xu)-F(x)|\le C$ for all $u\in G_\infty$ with $\operatorname{dist}(u,1)\le\delta$, then
$|F(x)|\le C$.
\item[(Q)] (\emph{Kernel forms.}) Let $K\subseteq H(\auA_f)$ be an open subgroup of the stabiliser of
$\widehat L_Z$ and $z\in H(\auA)$. The series $\sum_{\gamma\in\Gamma_H}e'((z\gamma g)_\infty)
\mathbf 1_K((z\gamma g)_f)$ converges absolutely for every $g\in H(\auA)$, and there is a form $Q$ of
type $e'$ on $H$ fixed by $K$ with $\aubar{Q(g)}$ equal to its sum.
\item[(S)] (\emph{Restriction pairing.}) Let $\beta_{\mathsf S}=b^{3m}$ and $N_{\mathsf S}=20(2n+1)^2$.
For every $\tau$ there is $\omega_\tau\in H_\tau$ with $\|\omega_\tau\|^2\le\beta_{\mathsf S}$ and
\[
V_\tau\le\beta_{\mathsf S}^{N_{\mathsf S}}\Bigl|\int_{G_\tau}c_\tau(r)\,\aubar{c'_\tau\bigl(
\omega_\tau^{-1}(r\oplus1)\omega_\tau\bigr)}\,dr\Bigr| .
\]
\end{enumerate}
\noindent Inputs (C) and (Q) are proved in \S\ref{pgk:sec} (Theorems~\ref{pgk:C} and~\ref{pgk:Q}), and
input (S) in \S\ref{pgl:sec} (Theorem~\ref{pgl:S}).

\begin{lemma}[Steepness]\label{pgq:steep}
Let $B=\kappa(\tau)+\Delta/2$ and $\Lambda$, $\Lambda'$ the highest weights of the lowest $K$-types for
the parameters $\varsigma_\tau$ of $V$ and $\lambda'=\lambda_{(1,0,0),\tau}$ of $Z$ \textup{(}signs
$\epsilon_{(0,0,0)}$, $\epsilon_{(1,0,0)}$\textup{)}. On the first $k-1$ coordinates both signs are $+$
and $\Lambda=\Lambda'=0$. On the remaining $m$, resp.\ $m+1$, coordinates, in the order of decreasing
parameter, the signs alternate starting with $+$, and $\Lambda$, resp.\ $\Lambda'$, is strictly
decreasing with values in $[-3B,-B]$.
\end{lemma}
\begin{proof}
For $V$ this is {\cite[Lemma~5.1]{P51}} and the remark after it. For $Z$: by \eqref{pgv:eq-param}
with $z=1$, $a=b=0$, $s_p=0$, the stem coordinate $\varrho\in\{0,\dots,r\}$ carries
$\lambda'_\varrho=\varsigma'_{\tau,\varrho}+\frac12$, that is $\frac k2-\frac12-\varrho$ for $\varrho\le
k-2$ and $-\kappa(\tau)-j\Delta+\frac12$ with $j=\varrho+2-k\in\{1,\dots,m+1\}$ otherwise; by
\eqref{pgv:eq-sign} its sign is $(-1)^{a(\varrho)}=(-1)^{j-1}$ for $\varrho\ge k-1$ and $+$ before.
So the far coordinates, in decreasing order, have signs $+,-,\dots,+$, and $m/2$ of them are
negative. Use $\Lambda'=\lambda'+\rho_{\mathrm n}-\rho_{\mathrm c}$ ({\cite[\S5.1]{P51}}). For $i\le k-2$
there are $i$ entries above $\lambda'_i$, all of sign $+$, and below it $k-2-i$ entries of sign $+$ and
the $m+1$ far entries, $m/2$ of sign $-$. Hence $2\rho_{\mathrm n,i}=m/2$, $2\rho_{\mathrm c,i}=
(k-2-i)+(m/2+1)-i$, and $2\Lambda'_i=(k-1-2i)+m/2-(k-2-2i)-(m/2+1)=0$. On the far coordinates
$|\Lambda'_j-\lambda'_j|\le n+1$, where $\lambda'_j=-\kappa(\tau)-j\Delta+\frac12$; the $\lambda'_j$ decrease in steps $\Delta>2n+2$, so $\Lambda'$ is
strictly decreasing, and $\Lambda'_j\ge-\kappa-(m+1)\Delta-n-\frac12\ge-3B$, $\Lambda'_j\le-\kappa-\Delta+n+\frac32\le-B$ because $m\Delta+r\le
\kappa(\tau)$ ({\cite[Lemma~13.4(S7)(g)]{P46}}), and $n+2\le r^8=\Delta/2$ as $\Delta=2r^8$, $r\ge100$ ({\cite[Lemma~13.4(S7)(b)]{P46}}); this also gives $\Delta>2n+2$.
\end{proof}

\subsubsection{Local estimates}\label{pgq:ss-local}

Let $h_V\ge\Delta$ and $h_Z\ge\Delta-t$ be the gaps of $\varsigma_\tau$ and of $-\lambda_{(1,0,0),\tau}$
({\cite[\S5.1]{P51}}, Lemma~\ref{pgv:gap}), $L_\tau=\max_i|\varsigma_{\tau,i}|\le\kappa(\tau)+m\Delta\le2b$ (this is the $L$ of {\cite[\S5.1]{P51}}, bounded there; it is not the quantity $\max_i|2\varsigma_{\tau,i}|$ of {\cite[Lemma~13.4]{P46}}), $C_{0,\tau}=
3n(L_\tau+2n)$ and $\delta_\tau=1/(4C_{0,\tau})$.

\begin{lemma}\label{pgq:local}
Let $\tau\in\auPl$, $g\in G_\tau$, $g'\in H_\tau$.
\begin{enumerate}
\item[(a)] $|c_\tau(g)|\le(2/\|g\|)^j$ for $0\le j\le h_V$, and $|c'_\tau(g')|\le(2/\|g'\|)^j$ for
$0\le j\le h_Z$.
\item[(b)] For real $s\ge8n^2$, $d_\tau\int_{G_\tau}\|g\|^{-s}dg\le8(5/\delta_\tau)^{2n^2}$.
\item[(c)] For $u\in G_\tau$ with $\|u-1\|\le\frac12$, $|c_\tau(gu)-c_\tau(g)|\le2C_{0,\tau}\|u-1\|$.
\item[(d)] $\aubar{c_\tau(g)}=c_\tau(g^{-1})$, $|c_\tau|\le1$, $|c'_\tau|\le1$.
\end{enumerate}
\end{lemma}
\begin{proof}
(a) {\cite[Theorem~5.6]{P51}}; for $c'_\tau$, {\cite[Lemma~3.25]{P49}} with {\cite[Theorem~3.17]{P49}}, which apply
by Lemma~\ref{pgv:gap}. (b) {\cite[Theorem~5.9]{P51}}, as $s/(s-4n^2)\le2$. (c) Steps~1 and 2 of the
proof of {\cite[Lemma~5.7]{P51}} give $\|\pi_\tau(u)v_\tau-v_\tau\|\le C_{0,\tau}\|\log u\|\le2C_{0,\tau}
\|u-1\|$; and $c_\tau(gu)-c_\tau(g)=\langle v_\tau,\pi_\tau(g)(\pi_\tau(u)v_\tau-v_\tau)\rangle$.
(d) $c_\tau(g^{-1})=\langle\pi_\tau(g)v_\tau,v_\tau\rangle$; unit vectors.
\end{proof}

\begin{lemma}[Balls]\label{pgq:balls}
For real $s\ge8n^2$, $\int_{G_\tau}\|r\|^{-s}dr\le2\cdot16^{n^2}V_\tau$.
\end{lemma}
\begin{proof}
For $T\ge1$ let $S_T=\{r:\|r\|\le T\}$, a compact set. Let $r_1,\dots,r_\nu\in S_T$ with
$\|r_ir_j^{-1}-1\|>1$ for $i\ne j$. As $\|r_ir_j^{-1}-1\|\le\|r_i-r_j\|\,\|r_j^{-1}\|$, $\|r_i-r_j\|>1/T$,
so the balls of radius $1/(2T)$ about the $r_i$ in $\mathrm M_n(\C)\cong\R^{2n^2}$ are disjoint and
lie in the ball of radius $2T$ about $0$; hence $\nu\le(4T^2)^{2n^2}$. Take such a family maximal.
Then every $r\in S_T$ has $\|rr_i^{-1}-1\|\le1$ for some $i$, so $S_T\subseteq\bigcup_iB_1r_i$ with
$B_1=\{\|r-1\|\le1\}$, and $\pgqvol(S_T)\le(4T^2)^{2n^2}V_\tau$ ($G_\tau$ is unimodular). Since
$\|r\|^{-s}=\int_1^\infty sT^{-s-1}\mathbf 1[\|r\|\le T]\,dT$, Fubini gives $\int\|r\|^{-s}dr\le
16^{n^2}V_\tau\int_1^\infty sT^{4n^2-s-1}dT=16^{n^2}V_\tau\,s/(s-4n^2)$.
\end{proof}

\begin{lemma}[Slices]\label{pgq:slice}
For $g\in H_\tau$ and real $s\ge16n^2$, $\int_{G_\tau}\|g(r\oplus1)\|^{-s}dr\le2\cdot16^{n^2}V_\tau$.
\end{lemma}
\begin{proof}
The continuous function $r\mapsto\|g(r\oplus1)\|\ge\|r\|/\|g^{-1}\|$ tends to infinity, so it attains
its minimum $t_0$ at some $r_0$. For every $r$,
\[
\|rr_0^{-1}\|=\|(rr_0^{-1})\oplus1\|=\bigl\|(g(r_0\oplus1))^{-1}g(r\oplus1)\bigr\|\le t_0\,\|g(r\oplus1)\|
\le\|g(r\oplus1)\|^2 .
\]
Hence $\|g(r\oplus1)\|^{-s}\le\|rr_0^{-1}\|^{-s/2}$; integrate, use right invariance and
Lemma~\ref{pgq:balls} with $s/2$.
\end{proof}

\begin{lemma}[Averaging]\label{pgq:schur}
Let $F$ be a bounded measurable function on $G_\infty$ with $F(g)=\int_{G_\infty}e(z)F(gz)\,dz$ for all
$g$. Then $\int_{G_\infty}e(u)F(ur)\,du=c(r)F(1)$ for every $r\in G_\infty$.
\end{lemma}
\begin{proof}
The function $e$ is integrable (Lemma~\ref{pgq:local}(a),(b)), so all integrals below converge
absolutely. By the hypothesis at the point $ur$ and the substitution $z=r^{-1}u^{-1}w$,
\[
\int e(u)F(ur)\,du=\int F(w)\Bigl(\int e(u)\,e(r^{-1}u^{-1}w)\,du\Bigr)dw .
\]
Now $e(u)=d\langle\pi(u)v,v\rangle$ and $e(r^{-1}u^{-1}w)=d\,\aubar{\langle\pi(u)\pi(r)v,\pi(w)v\rangle}=
d\langle\pi(w)v,\pi(u)\pi(r)v\rangle$ ($\pi=\bigotimes\pi_\tau$, $v=\bigotimes v_\tau$). Schur's
orthogonality relations for the square-integrable irreducible representation $\pi$ of formal degree
$d$ (conjugate-linear in the first variable: $\int\langle a,\pi(u)b\rangle\aubar{\langle a',\pi(u)b'
\rangle}du=d^{-1}\langle a,a'\rangle\langle b',b\rangle$; see {\cite[(9)]{P49}}) give
$\int e(u)e(r^{-1}u^{-1}w)\,du=d\langle v,\pi(r)v\rangle\langle\pi(w)v,v\rangle=c(r)e(w)$. So the left
side is $c(r)\int e(w)F(w)\,dw=c(r)F(1)$.
\end{proof}

\begin{lemma}[Variation of the kernel]\label{pgq:var}
Let $\mathsf D(g)=\prod_\tau(2/\|g_\tau\|)^{2A}$ on $G_\infty$, with $4A\le h_V$, $2A\ge8n^2$, and
$T_\tau=d_\tau\int_{G_\tau}(2/\|g\|)^{2A}dg$. Then $|e|\le d\,\mathsf D^2$,
$\int d\,\mathsf D=\prod_\tau T_\tau$, and for $u\in G_\infty$ with $\operatorname{dist}(u,1)\le\delta
\le\frac12$,
\[
\int_{G_\infty}|e(u^{-1}z)-e(z)|\,dz\le2\bigl(2fC_0\delta\bigr)^{1/2}\prod_\tau T_\tau,\qquad C_0=
\max_\tau C_{0,\tau}.
\]
\end{lemma}
\begin{proof}
$|e|\le d\mathsf D^2$ is Lemma~\ref{pgq:local}(a) with $j=4A$, and the integral of $d\mathsf D$ is a
product by Fubini. Let $a_\tau=\aubar{c_\tau(u_\tau^{-1}z_\tau)}$, $b_\tau=\aubar{c_\tau(z_\tau)}$. By
Lemma~\ref{pgq:local}(d),(c), $|a_\tau-b_\tau|=|c_\tau(z_\tau^{-1}u_\tau)-c_\tau(z_\tau^{-1})|\le
2C_0\delta$, and as all $|a_\tau|,|b_\tau|\le1$, telescoping gives $|\prod a_\tau-\prod b_\tau|\le
2fC_0\delta$. Also $|\prod a_\tau|\le\mathsf D(u^{-1}z)^2$ and $|\prod b_\tau|\le\mathsf D(z)^2$. Hence
$|\prod a-\prod b|^2\le2fC_0\delta\,(\mathsf D(u^{-1}z)+\mathsf D(z))^2$, and
$|e(u^{-1}z)-e(z)|\le d(2fC_0\delta)^{1/2}(\mathsf D(u^{-1}z)+\mathsf D(z))$. Integrate, using left
invariance.
\end{proof}

\subsubsection{Sizes}\label{pgq:ss-sizes}

Put
\[
A=4f(n+1)^2,\qquad\mathsf D_0=b^{x^{36}},\qquad\Pi=\prod_\tau\Bigl(T_\tau\,(4\beta_{\mathsf S})^{2A}\,
2\cdot16^{n^2}\,\beta_{\mathsf S}^{N_{\mathsf S}}\Bigr) .
\]

\begin{lemma}\label{pgq:sizes}
\begin{enumerate}
\item[(a)] $4A\le h_V$, $2A\le h_Z$, $A\ge16(n+1)^2$, $A\ge2f(n+1)+2$.
\item[(b)] $\mathsf D_0\ge\max\bigl(10,\ f,\ 20\,\Pi,\ 4(2fC_0)^{1/2}\prod_\tau T_\tau\bigr)$ and
$\mathsf D_0^{2n+2}(M\kappa_{\mathrm c}^2)^n\le M^{b/2}$.
\item[(c)] $\beta_{\mathsf S}\ge10^{12}$, $\beta_{\mathsf S}\ge2n+1+2b$, $\beta_{\mathsf S}\ge(n+1)^4
(6B)^{2m+2}$, and $\min(h_V,h_Z)\ge500(n+1)^3+10^{30}\bigl((2n+1)^2\log\beta_{\mathsf S}\bigr)^3$.
\end{enumerate}
\end{lemma}
\begin{proof}
We use {\cite[Lemma~13.4]{P46}} (here $n=r$; the $n$ of (S7) is $m/2$): $x^5/2\le k\le r=n\le3x^5$, $m\le2\sqrt x$, $x^{11}\le f\le x^{12}$,
$\mathrm{rd}_E\le x^{12}$, $q_1\le x^{12}$, $d\le\log x$, $(\log b)^{100}\le x$, $x\ge128^{100}$,
$b\ge x^{10^{12}}$, $b^{20}\le M\le b^{22}$, $\Delta=2r^8$, $t=5r$, $B\le2b$; moreover $L_\tau\le2b$ ({\cite[\S5.1]{P51}}).
(a) $4A\le16x^{12}(4x^5)^2\le x^{40}/2^{9}\le\Delta-t\le\min(h_Z,h_V)$ (as $r\ge x^5/2$, $h_Z\ge\Delta-t$ by Lemma~\ref{pgv:gap} and $h_V\ge\Delta$); the other two
hold as $f\ge4$. (b) $5/\delta_\tau=60n(L_\tau+2n)\le b^2$, so $T_\tau\le2^{2A}\cdot8b^{4n^2}\le
b^{2A+4n^2+3}$ by Lemma~\ref{pgq:local}(b). With $4\beta_{\mathsf S}\le b^{3m+2}$ and
$2\cdot16^{n^2}\le b^{n^2+1}$,
\[
\log_b\Pi\le f\bigl(2A(3m+3)+5n^2+4+60m(2n+1)^2\bigr)\le f\,m\bigl(48f(n+1)^2+244(n+1)^2\bigr)\le
50f^2m(n+1)^2\le x^{35},
\]
since $50x^{24}\cdot2x^{1/2}\cdot16x^{10}\le x^{35}$. So $20\Pi\le b^{x^{35}+1}\le\mathsf D_0$; the
bound for $4(2fC_0)^{1/2}\prod T_\tau\le b^2\prod T_\tau$ is the same computation, and $10,f\le
\mathsf D_0$. Next $\kappa_{\mathrm c}\le8x^{12}\cdot3x^5\cdot x^{12}+48(3x^5)^6x^{24}x^{13}+1\le x^{70}
\le b$, so $\log_b(\mathsf D_0^{2n+2}(M\kappa_{\mathrm c}^2)^n)\le(2n+2)x^{36}+24n\le x^{42}\le10b\le
\log_b M^{b/2}$. (c) $\log_b\beta_{\mathsf S}=3m$ and $b\ge x^{10^{12}}$ give the first two;
$(n+1)^4(12b)^{2m+2}\le b\cdot b^{1.01(2m+2)}\le b^{3m}$ as $m\ge200$ ({\cite[Lemma~13.4(S7)]{P46}}).
Finally $10^{30}((2n+1)^2\,3m\log b)^3\le10^{35}(37x^{10}\cdot6x^{1/2}x^{1/100})^3\le x^{33}$ and
$500(n+1)^3\le x^{17}$, while $h_Z\ge\Delta-t\ge x^{40}/2^8-15x^5$.
\end{proof}

\noindent From now on $A$, $\mathsf D_0$ are as above, $\omega=(\omega_\tau)$ is the real point of $H$
given by (S), and for a real point $\alpha$ of $H$ we put
\[
\phi_\alpha(g)=\prod_\tau\aubar{c'_\tau\bigl(\alpha_\tau^{-1}g_\tau\omega_\tau\bigr)}=\frac{e'(\alpha^{-1}g
\omega)}{d'}\quad(g\in H_\infty),\qquad\pgqSl(\alpha)=\prod_\tau\bigl(2\|\alpha_\tau\|\|\omega_\tau\|\bigr)^{2A}
\,2\cdot16^{n^2}V_\tau .
\]

\begin{lemma}[Slices of the translated coefficient]\label{pgq:slicecoef}
Let $g\in H_\infty$ and $t>0$ with $t\le\prod_\tau\|g_\tau(r_\tau\oplus1)\|$ for all $r\in G_\infty$.
Then $r\mapsto\phi_\alpha(g(r\oplus1))$ is integrable on $G_\infty$ and
$\int_{G_\infty}|\phi_\alpha(g(r\oplus1))|\,dr\le\pgqSl(\alpha)\,t^{-A}$.
\end{lemma}
\begin{proof}
Write $X_\tau=g_\tau(r_\tau\oplus1)$. As $\|X_\tau\|\le\|\alpha_\tau\|\,\|\alpha_\tau^{-1}X_\tau\omega_\tau\|\,
\|\omega_\tau\|$ ($\|\alpha_\tau\|=\|\alpha_\tau^{-1}\|$, $\|\omega_\tau^{-1}\|=\|\omega_\tau\|$),
Lemma~\ref{pgq:local}(a) with $j=2A\le h_Z$ gives $|c'_\tau(\alpha_\tau^{-1}X_\tau\omega_\tau)|\le
(2\|\alpha_\tau\|\|\omega_\tau\|)^{2A}\|X_\tau\|^{-A}\|X_\tau\|^{-A}$. The product over $\tau$ of the
first factors $\|X_\tau\|^{-A}$ is at most $t^{-A}$; integrate the rest by Fubini and
Lemma~\ref{pgq:slice} ($A\ge16n^2$). The function is continuous, hence measurable.
\end{proof}

\subsubsection{The cusp}\label{pgq:ss-cusp}

\begin{lemma}\label{pgq:cusp}
Let $y\in\auAz$ and $S\ge\sup|y|$. If $\lambda(x)<\theta:=(\mathsf D_0^{-2}/(M
\kappa_{\mathrm c}^2))^n$, then $|y(x)|\le S/2$.
\end{lemma}
\begin{proof}
Let $\delta=\mathsf D_0^{-2}\le\frac12$ and $\operatorname{dist}(u,1)\le\delta$. By the reproducing
property at $xu$ and at $x$, and the substitution $z\mapsto u^{-1}z$,
$y(xu)-y(x)=\int(e(u^{-1}z)-e(z))y(xz)\,dz$; by Lemma~\ref{pgq:var} (its hypotheses hold by
Lemma~\ref{pgq:sizes}(a)) and Lemma~\ref{pgq:sizes}(b),
$|y(xu)-y(x)|\le S\cdot2(2fC_0)^{1/2}\mathsf D_0^{-1}\prod_\tau T_\tau\le S/2$. As $y$ is a form of
type $e$ fixed by $J_M\supseteq J_G(M)$, (C) gives $|y(x)|\le S/2$.
\end{proof}

\subsubsection{Poincar\'e series of deep level}\label{pgq:ss-poinc}

Fix $x\in G(\auA)$ and put $\hat x=x\oplus1$. Let $\mathcal R\subset\Gamma_H$ be a set of
representatives of $\Gamma_H/(\Gamma_G\oplus1)$ containing $1$; then $(\sigma,\delta)\mapsto
\sigma(\delta\oplus1)$ is a bijection $\mathcal R\times\Gamma_G\to\Gamma_H$. For $\sigma\in\Gamma_H$
let $\pgqcol(\sigma)\in\prod_\tau\C^{n+1}\cong\R^{2f(n+1)}$ be the archimedean coordinates of
$\hat x^{-1}\sigma\tilde e$. Let $\mathcal R^\flat$ be the set of $\sigma\in\mathcal R$ for which some
$k\in G(\auA_f)$ has $(\hat x^{-1}\sigma(k\oplus1))_f\in J_H(\pcb)$.

\begin{lemma}[Columns]\label{pgq:cols}
\begin{enumerate}
\item[(a)] For $\sigma\in\mathcal R$ and $h\in G(\auA)$, $g=\hat x^{-1}\sigma(h\oplus1)$ satisfies
$\prod_\tau\|g_\tau\|\ge|\pgqcol(\sigma)|/\sqrt f$; and $|\pgqcol(1)|=\sqrt f$.
\item[(b)] For distinct $\sigma,\sigma'\in\mathcal R^\flat$,
$|\pgqcol(\sigma)-\pgqcol(\sigma')|\ge M^{b/2}\lambda(x)$.
\end{enumerate}
\end{lemma}
\begin{proof}
(a) $g_\infty\tilde e=\hat x_\infty^{-1}\sigma\tilde e$, as $h\oplus1$ fixes $\tilde e$; so
$|\pgqcol(\sigma)|^2=\sum_\tau|g_\tau\tilde e|^2\le\sum_\tau\|g_\tau\|^2\le f\prod_\tau\|g_\tau\|^2$ (all
factors are $\ge1$). For $\sigma=1$ the coordinates are those of $\tilde e$, of norm $1$ at each
$\tau$.

(b) Let $v=\sigma\tilde e$, $v'=\sigma'\tilde e\in Z(E)$. If $v=v'$, then $\sigma^{-1}\sigma'$ fixes
$\tilde e$, hence preserves $V=\tilde e^\perp$, so $\sigma^{-1}\sigma'\in\Gamma_G\oplus1$ and
$\sigma=\sigma'$; thus $v\ne v'$. Choose $k,k'$ as in the definition of $\mathcal R^\flat$ and put
$g=(\hat x^{-1}\sigma(k\oplus1))_f\in J_H(\pcb)$. Then $\hat x_f^{-1}v=g\tilde e\equiv\tilde e$ modulo
$\pcb\widehat L_Z$, and likewise for $v'$, so $\hat x_f^{-1}(v-v')\in\pi_M^b\widehat L_Z$ and
$w=\pi_M^{-b}(v-v')$ is a non-zero element of $L_Z(\hat x)$. Since $|\tau(\pi_M)|=M^{1/2}$,
$|\pgqcol(\sigma)-\pgqcol(\sigma')|=|v-v'|_{\hat x}=M^{b/2}|w|_{\hat x}$. Finally $L_Z(\hat x)=
L_V(x)\oplus\auO_E\tilde e$ orthogonally in the coordinates, so for $w=w_0+a\tilde e\ne0$: if
$w_0\ne0$ then $|w|_{\hat x}\ge|w_0|_x\ge\lambda(x)$, and if $w_0=0$ then $a\in\auO_E\smallsetminus0$
and $\max_\tau|\tau(a)|\ge1$, as $\prod_\tau|\tau(a)|^2=|\mathrm N_{E/\Q}(a)|\ge1$.
\end{proof}

Let $y$ be a bounded measurable function on $[G]$ fixed on the right by $J_G(\pcb)$, $S\ge\sup|y|$,
$\pgqvol=\pgqvol(J_G(\pcb))$, and for a real point $\alpha$ of $H$ and $g\in H(\auA)$ let
$\Phi_\alpha(g)=\phi_\alpha(g_\infty)\mathbf 1_{J_H(\pcb)}(g_f)$. Put
\[
\mathsf P_x(\alpha)=\int_{[G]}\sum_{\gamma\in\Gamma_H}\Phi_\alpha\bigl(\hat x^{-1}\gamma(h\oplus1)\bigr)
\,y(h)\,dh,\qquad\mathsf c_x(\alpha)=\int_{G_\infty}\phi_\alpha(r\oplus1)\,y(xr)\,dr .
\]

\begin{lemma}[Unfolding]\label{pgq:unfold}
Let $\mathsf D\ge\max(10,\sqrt f)$ with $\mathsf D^2\le M^{b/2}\lambda(x)$. Then the series defining
$\mathsf P_x(\alpha)$ converges absolutely for almost all $h$, $\mathsf c_x(\alpha)$ converges
absolutely, and
\[
\bigl|\mathsf P_x(\alpha)-\pgqvol\cdot\mathsf c_x(\alpha)\bigr|\le\pgqvol\cdot\pgqSl(\alpha)\,S\,
\frac{10}{\mathsf D}.
\]
\end{lemma}
\begin{proof}
For $\sigma\in\mathcal R$ put $T_\sigma(h)=\Phi_\alpha(\hat x^{-1}\sigma(h\oplus1))y(h)$ on $G(\auA)$.
\emph{Support.} If $T_\sigma(h)\ne0$ then $(\hat x^{-1}\sigma(h\oplus1))_f\in J_H(\pcb)$, so
$\sigma\in\mathcal R^\flat$. \emph{One coset.} Let $\sigma\in\mathcal R^\flat$, with $k$ as in the
definition, and $g_0=\hat x^{-1}\sigma(k\oplus1)$. For $h\in G(\auA)$, $(g_0(h\oplus1))_f\in J_H(\pcb)$ if
and only if $h_f\in J_G(\pcb)$, as $g_{0,f}\in J_H(\pcb)$. Substituting $h\mapsto kh$,
\[
\int_{G(\auA)}|T_\sigma|\le S\,\pgqvol\int_{G_\infty}|\phi_\alpha(g_{0,\infty}(r\oplus1))|\,dr\le S\,\pgqvol\,
\pgqSl(\alpha)\,t_\sigma^{-A}
\]
by Lemma~\ref{pgq:slicecoef}, where $t_\sigma=\max(1,|\pgqcol(\sigma)|/\mathsf D)$ is admissible by
Lemma~\ref{pgq:cols}(a) ($g_0(r\oplus1)=\hat x^{-1}\sigma((kr)\oplus1)$, and $\sqrt f\le\mathsf D$).
\emph{The sum over cosets.} By Lemma~\ref{pgq:cols}, the points $\pgqcol(\sigma)$, $\sigma\in
\mathcal R^\flat$, are $\mathsf D^2$-separated and $|\pgqcol(1)|\le\mathsf D$; so $|\pgqcol(\sigma)|\ge
\mathsf D^2-\mathsf D\ge\mathsf D^2/2$ for $\sigma\ne1$, and the points $p_\sigma=\pgqcol(\sigma)/(\mathsf
D^2/2)$, $\sigma\in\mathcal R^\flat\smallsetminus1$, satisfy $|p_\sigma|\ge1$, $|p_\sigma-p_{\sigma'}|\ge1$.
By {\cite[Lemma~8.22(b)]{P51}} in $\R^{2f(n+1)}$ with exponent $A\ge2f(n+1)+2$,
\[
\sum_{\sigma\in\mathcal R^\flat\smallsetminus1}\Bigl(\frac{\mathsf D}{|\pgqcol(\sigma)|}\Bigr)^A=\Bigl(
\frac2{\mathsf D}\Bigr)^A\sum|p_\sigma|^{-A}\le\Bigl(\frac{10}{\mathsf D}\Bigr)^A\le\frac{10}{\mathsf D}.
\]
Hence $\sum_{\sigma\in\mathcal R}\int|T_\sigma|<\infty$. \emph{Unfolding.} As $y(\delta h)=y(h)$, the
bijection $\mathcal R\times\Gamma_G\to\Gamma_H$ gives $\sum_\gamma\Phi_\alpha(\hat x^{-1}\gamma(h\oplus1))
y(h)=\sum_{\sigma}\sum_{\delta}T_\sigma(\delta h)$. By \eqref{pgq:eq-unfold} applied to $|T_\sigma|$ and
summed over $\sigma$, the double series converges absolutely for almost all $h$ and is integrable on
$[G]$; by dominated convergence and \eqref{pgq:eq-unfold} again,
$\mathsf P_x(\alpha)=\sum_{\sigma\in\mathcal R^\flat}\int_{G(\auA)}T_\sigma$. \emph{Main term.} For
$\sigma=1$ substitute $h\mapsto xh$: $\int T_1=\int_{G(\auA)}\phi_\alpha(h_\infty\oplus1)
\mathbf 1_{J_G(\pcb)}(h_f)y(xh)\,dh=\pgqvol\cdot\mathsf c_x(\alpha)$, as $y(xh_\infty h_f)=y(xh_\infty)$
for $h_f\in J_G(\pcb)$; $\mathsf c_x(\alpha)$ converges absolutely by Lemma~\ref{pgq:slicecoef} with
$g=1$, $t=1$. The other terms are bounded by $S\,\pgqvol\,\pgqSl(\alpha)(\mathsf D/|\pgqcol(\sigma)|)^A$,
whose sum is at most $S\,\pgqvol\,\pgqSl(\alpha)\,10/\mathsf D$.
\end{proof}

\begin{lemma}[The period vanishes]\label{pgq:zero}
Let $y$ be as above, and orthogonal in the sense of Theorem~\ref{pgq:V} at the
point $\omega$. Then $\mathsf P_x(\alpha)=0$ for every real point $\alpha$ of $H$.
\end{lemma}
\begin{proof}
Let $z=(\alpha^{-1},1)\hat x^{-1}\in H(\auA)$. For $h\in G(\auA)$ and $\gamma\in\Gamma_H$, as $\alpha$
and $\omega$ are real, $z\gamma(h\oplus1)(\omega,1)$ has archimedean part $\alpha^{-1}(\hat x^{-1}\gamma
(h\oplus1))_\infty\omega$ and finite part $(\hat x^{-1}\gamma(h\oplus1))_f$. So
$\Phi_\alpha(\hat x^{-1}\gamma(h\oplus1))=d'^{-1}e'((z\gamma g)_\infty)\mathbf 1_{J_H(\pcb)}((z\gamma g)_f)$
with $g=(h\oplus1)\omega$. By (Q) with $K=J_H(\pcb)$ there is a form $Q$ of type $e'$ fixed by
$J_H(\pcb)$ with $\sum_\gamma\Phi_\alpha(\hat x^{-1}\gamma(h\oplus1))=d'^{-1}\aubar{Q((h\oplus1)\omega)}$.
Hence $\mathsf P_x(\alpha)=d'^{-1}\,\aubar{\int_{[G]}\aubar{y(h)}Q((h\oplus1)\omega)\,dh}=0$.
\end{proof}

\subsubsection{Away from the cusp}\label{pgq:ss-bulk}

\begin{lemma}\label{pgq:bulk}
Let $y\in\auAz$ be orthogonal as in Theorem~\ref{pgq:V} at $\omega$,
$S\ge\sup|y|$, and $\lambda(x)\ge\theta$ \textup{(}Lemma~\ref{pgq:cusp}\textup{)}. Then $|y(x)|\le S/2$.
\end{lemma}
\begin{proof}
$y$ is fixed by $J_M\supseteq J_G(\pcb)$. By Lemma~\ref{pgq:sizes}(b), $\mathsf D=\mathsf D_0$ satisfies
$\mathsf D\ge\max(10,\sqrt f)$ and $\mathsf D^2\le M^{b/2}\theta\le M^{b/2}\lambda(x)$. So for every real
$\alpha$, Lemmas~\ref{pgq:unfold} and \ref{pgq:zero} give
\begin{equation}\label{pgq:eq-conv}
|\mathsf c_x(\alpha)|\le\pgqSl(\alpha)\,S\,\frac{10}{\mathsf D_0}.
\end{equation}
\emph{Translation.} Put $F(g)=y(xg)$ on $G_\infty$; it is bounded and $F(g)=\int e(z)F(gz)\,dz$. For
$u\in G_\infty$, $\phi_{(u\oplus1)\omega}(r\oplus1)=\phi_\omega((u^{-1}r)\oplus1)$, so the substitution
$r\mapsto ur$ gives $\mathsf c_x((u\oplus1)\omega)=\int_{G_\infty}\phi_\omega(r\oplus1)F(ur)\,dr$.
\emph{Averaging.} By
Fubini (the integrand is bounded by $S|e(u)||\phi_\omega(r\oplus1)|$, integrable) and
Lemma~\ref{pgq:schur},
\[
\int_{G_\infty}e(u)\,\mathsf c_x((u\oplus1)\omega)\,du=\int\phi_\omega(r\oplus1)\int e(u)F(ur)\,du\,dr
=\eta\,y(x),\qquad\eta=\int_{G_\infty}\phi_\omega(r\oplus1)c(r)\,dr .
\]
By Fubini, $\eta=\prod_\tau\int_{G_\tau}c_\tau(r)\aubar{c'_\tau(\omega_\tau^{-1}(r\oplus1)\omega_\tau)}dr$,
so $|\eta|\ge\prod_\tau V_\tau\beta_{\mathsf S}^{-N_{\mathsf S}}$ by (S). \emph{Bound.} As
$\|(u_\tau\oplus1)\omega_\tau\|\le\|u_\tau\|\beta_{\mathsf S}^{1/2}$ and $\|\omega_\tau\|\le
\beta_{\mathsf S}^{1/2}$, $\pgqSl((u\oplus1)\omega)\le\prod_\tau(2\|u_\tau\|\beta_{\mathsf S})^{2A}
2\cdot16^{n^2}V_\tau$; with $|e(u)|\le d\prod_\tau(2/\|u_\tau\|)^{4A}$ (Lemma~\ref{pgq:var}) and
\eqref{pgq:eq-conv},
\[
|e(u)\,\mathsf c_x((u\oplus1)\omega)|\le S\frac{10}{\mathsf D_0}\prod_\tau d_\tau\Bigl(\frac2{\|u_\tau\|}
\Bigr)^{2A}(4\beta_{\mathsf S})^{2A}\,2\cdot16^{n^2}V_\tau .
\]
Integrating over $u$, $|\eta|\,|y(x)|\le S\frac{10}{\mathsf D_0}\prod_\tau T_\tau(4\beta_{\mathsf S})^{2A}
2\cdot16^{n^2}V_\tau$, and dividing by the lower bound for $|\eta|$ (the $V_\tau$ cancel),
$|y(x)|\le S\cdot10\,\Pi/\mathsf D_0\le S/2$ by Lemma~\ref{pgq:sizes}(b).
\end{proof}

\subsubsection{Proof of Theorem~\ref{pgq:V}}\label{pgq:ss-proof}

The inputs (C), (Q), (S) hold by Theorems~\ref{pgk:C}, \ref{pgk:Q}, \ref{pgl:S}. Let $\omega$ be
given by (S) and let $y\in\auAz$ be orthogonal as in the theorem. The form $y$ is
bounded; let $S=\sup|y|$. For every $x\in G(\auA)$, $|y(x)|\le S/2$: by Lemma~\ref{pgq:cusp} if
$\lambda(x)<\theta$ and by Lemma~\ref{pgq:bulk} otherwise. Hence $S\le S/2$, $S=0$ and $y=0$.
\qed




% ===== C-grid10.tex =====
% proofs/C-grid10.tex
% Continuation of proofs/C-grid8.tex (pgz:): the size bounds (B1)--(B3) of \S\ref{pgz:ss-b}.
% Labels start with pgt:. Conventions and macros: those of proofs/C-grid8.tex.
\providecommand{\auI}{\mathcal I}
\providecommand{\auGL}{\mathrm{GL}}
\providecommand{\auU}{\mathrm U}
\providecommand{\auO}{\mathcal O}
\providecommand{\auNm}{\operatorname{Nm}}
\providecommand{\pcb}{\mathfrak b}
\providecommand{\auPl}{\Phi}
\providecommand{\pkm}{\mathsf m}
\providecommand{\pkev}{\operatorname{ev}}
\providecommand{\pgsdu}{\mathrm{du}}
\providecommand{\aneu}{\mathrm{eu}_d}
\providecommand{\anrd}{\mathrm{rd}}
\providecommand{\pcnot}[1]{\par\smallskip{\raggedright\noindent\textsc{Not proved here.} [#1]\par}\smallskip}

\subsection{The size bounds on the last path}\label{pgt:sec}

This section proves the bounds of \S\ref{pgz:ss-b} that part (b) of Proposition~\ref{pgr:den} was
reduced to in Theorem~\ref{pgz:b}. The notation is that of \S\ref{pgz:sec}; $\vartheta=1$, so
$\pcm=\pcb=(\pi_M)^b$, and $G_0=\exp(x^{99}b^2)$.

\subsubsection{(B3): the number of classes for ranks up to $2n_m$}\label{pgt:ss-B3}

The proof of {\cite[Proposition~7.20]{P49}} uses the hypothesis $n\le n_m$ only in its last paragraph
(``Sizes''): the counting argument (the two elementary facts, the invariant $\mathsf I(g)$, and the
passage to levels) holds for every rank $n\ge1$, since {\cite[Lemma~14.21]{P46}}, which it uses with
$c_0=q_1d$, holds for every $n\ge1$. We record the statement for all ranks and redo the size
comparison for $n\le2n_m$.

\begin{proposition}[Classes in rank $\le 2n_m$]\label{pgt:classes}
Let $X=E^n$ be the positive definite standard space of rank $n\ge1$ (form $\sum x_i\aubar{x_i}$,
lattice $\widehat\auO_E^{\,n}$), $G=\auU(X)$, $\Gamma=G(E^+)$, $K^\circ$ the stabiliser of
$\widehat\auO_E^{\,n}$ in $G_f$, and for an integer $\mathsf G\ge1$ let $K(\mathsf G)=\{k\in K^\circ:
k\equiv1\bmod\mathsf G\}$. Then, with $C_{\mathrm{sb}}=6q_1d(fn)^2$ and $\mathsf B=n^2C_{\mathrm{sb}}^2
\cdot2f\sqrt d$,
\begin{equation}\label{pgt:eq-classes}
\#\bigl(\Gamma\backslash G_f/J\bigr)\le(2\mathsf B+1)^{4fn^2}\,\mathsf G^{2fn^2}
\end{equation}
for every subgroup $J\supseteq K(\mathsf G)$ of $G_f$. If $n\le2n_m$ and $\mathsf G=M^b$, the right
side is at most $\exp(x^{26}b^2)\le G_0$.
\end{proposition}
\begin{proof}
\emph{The bound \eqref{pgt:eq-classes}.} This is the proof of {\cite[Proposition~7.20]{P49}}, which
we check uses no restriction on $n$. Fact (i) ($\auA_{E,f}=E+\widehat\auO_E$ with a representative of
conjugates at most $\mathsf B_0=2f\sqrt d$) and fact (ii) (at most $(2B+1)^{2f}$ elements of $\auO_E$
with all conjugates at most $B$) do not involve $n$. For $g\in G_f$, {\cite[Lemma~14.21]{P46}} (for
the given $n$, with $c_0=q_1d$, which lies in the different of $E$ as in the proof of
{\cite[Proposition~14.22]{P46}}) gives $Y\in\auGL_n(E)$ with $A=g^{-1}Y$ integral and entries with
conjugates at most $C_{\mathrm{sb}}$; fact (i) gives $X'\in M_n(E)$ with $A^{-1}-X'$ integral and
entries bounded by $\mathsf B_0$. As $g^*g=1$, $Y^*Y=A^*A$ and $Y^*YX'=A^*(1-A(A^{-1}-X'))$ are in
$M_n(\auO_E)$, with conjugates at most $nC_{\mathrm{sb}}^2$ and $n^2C_{\mathrm{sb}}^2\mathsf B_0=\mathsf B$;
by fact (ii) the pair $\mathsf I(g)=(Y^*Y,Y^*YX')$ takes at most $(2\mathsf B+1)^{4fn^2}$ values. If
$\mathsf I(g)=\mathsf I(g')$ (with $X'_g$, $X'_{g'}$ the matrices chosen for $g$, $g'$), then $Y^*Y=Y'^*Y'$ is invertible and the second components give $X'_g=X'_{g'}=:X'$; so $A'X'=1-A'(A'^{-1}-X')$ and $A'(A^{-1}-X')$ are integral. Moreover $\gamma=YY'^{-1}$ satisfies $\gamma^*\gamma=1$, so $\gamma\in\Gamma$,
and $k=g'^{-1}\gamma^{-1}g=A'A^{-1}=A'(A^{-1}-X')+A'X'$ is integral and unitary, so $k\in K^\circ$ and
$g=\gamma g'k$. Hence $\#(\Gamma\backslash G_f/K^\circ)\le(2\mathsf B+1)^{4fn^2}$. Finally
$\Gamma\backslash G_f/(J\cap K^\circ)\to\Gamma\backslash G_f/J$ is onto and the fibres of
$\Gamma\backslash G_f/(J\cap K^\circ)\to\Gamma\backslash G_f/K^\circ$ have at most $[K^\circ:K(\mathsf
G)]\le\#M_n(\auO_E/\mathsf G)=\mathsf G^{2fn^2}$ elements, since $k\mapsto k\bmod\mathsf G$ embeds
$K^\circ/K(\mathsf G)$ in $M_n(\auO_E/\mathsf G)$ and $[E:\Q]=2f$.

\emph{Sizes for $n\le2n_m$.} By {\cite[Lemma~13.4]{P46}}: $f\le q_1\le x^{12}$ (S3), $d\le\log x$
((S1): $\log x\ge10^{14}\cdot2d$), $\log M\le22\log b$ (S5), $\log b\le x^{1/100}$ (S1); and
$n_m\le40x^6$ ({\cite[Lemma~8.8(b)]{P49}}). So $n\le80x^6$ and
\[
C_{\mathrm{sb}}=6q_1d(fn)^2\le6\cdot x^{12}\cdot\log x\cdot x^{24}\cdot6400x^{12}\le x^{50}.
\]
Since $n\le C_{\mathrm{sb}}$ and $\mathsf B_0=2f\sqrt d\le C_{\mathrm{sb}}$, we have $\mathsf B\le
C_{\mathrm{sb}}^5$ and $2\mathsf B+1\le C_{\mathrm{sb}}^6$, so $\log(2\mathsf B+1)\le300\log x$. Also
$fn^2\le x^{12}\cdot6400x^{12}=6400x^{24}$. Hence
\[
4fn^2\log(2\mathsf B+1)\le25600\,x^{24}\cdot300\log x\le x^{26},\qquad
2fn^2\log(M^b)\le12800\,x^{24}\cdot22\,b\log b\le x^{25}b^2,
\]
and the sum is at most $x^{26}b^2\le x^{99}b^2$.
\end{proof}

\begin{corollary}[(B3)]\label{pgt:B3}
Let $\vartheta=1$. For $z\in\{0,1\}$ and $a\le m$, $\#(\Gamma^z_a\backslash G^z_{a,f}/J^z_a)\le
\exp(x^{26}b^2)\le G_0$; likewise $\#(\Gamma_Y\backslash G_{Y,f}/J_Y)\le G_0$.
\end{corollary}
\begin{proof}
The node $p^z_a=(z,a,m)$ has rank $N^z_a\le2n_m$ (Lemma~\ref{pgs:boxes}). Its space is a positive
definite standard space: for $(z,a)\ne(0,0)$ it is $E^{N^z_a}$ with the standard lattice and the
diagonal form with the signs $\epsilon_{p^z_a}$ (Definition~\ref{pgv:def-spaces}; the node $(1,0,0)$
is not on $\pgrlast$ since $m\ge1$), and all these signs are $+1$, since the number of negative signs
at $(z,a,m)$ is $\lfloor z/2\rfloor=0$ (Lemma~\ref{pgv:count}(b)); for $(z,a)=(0,0)$ it is $Y_m$,
which via its frame is the positive definite standard space of rank $n_m-1$ with the lattice
$\widehat\auO_E^{\,n_m-1}$ ({\cite[\S5.1.3]{P49}}, and {\cite[Lemma~8.8(a)]{P49}} for the same identification).
So Proposition~\ref{pgt:classes} applies with $G=G^z_a$, $\Gamma=\Gamma^z_a$, $K^\circ=K_{p^z_a}$.
For $\vartheta=1$ every level on $\pgrlast$ is $J(\pcb)$, the principal congruence subgroup of
$K_{p^z_a}$ of level $\pcb$ (\S\ref{pgw:ss-last}, \S\ref{pgz:sec}). As $\pcb=(\pi_M)^b$ and
$\pi_M\aubar{\pi_M}=M$, $M^b\in\pcb$, so $K(M^b)\subseteq J(\pcb)$, and \eqref{pgt:eq-classes} with
$\mathsf G=M^b$ gives the bound. The node $Y=Y^m_m$ (not the space $Y_m$ above) is $p^0_m$, so the last assertion is the case $a=m$,
$z=0$.
\end{proof}

\begin{remark}\label{pgt:rem-B3}
Corollary~\ref{pgt:B3} is (B3) of \S\ref{pgz:ss-b}. The statement of {\cite[Proposition~7.20]{P49}}
may be widened from $n\le n_m$ to all $n$ for the bound \eqref{pgt:eq-classes}, without change of
proof; only its final size estimate $\exp(x^{30})$ uses $n\le n_m$ and $\mathsf G=MN$.
\end{remark}

\subsubsection{(B2): an integral branching matrix for an arbitrary interlacing pair}\label{pgt:ss-B2}

We use the tensor model of {\cite[\S5.2.6]{P49}} and the coordinates of \S\ref{pgs:ss-weyl}. As
there, all statements are multiplicative over $\tau\in\auPl$; we write the constructions for one
$\tau$ and take tensor products over $\tau$ at the end. Put $N=n+1$ and regard $\auGL_n\subset
\auGL_N$ as $g\mapsto g\oplus1$, the new coordinate being the last one, and $\auGL_n\times\auGL_1
\subset\auGL_N$ as $(g,s)\mapsto g\oplus s$. In the model of {\cite[\S5.2.6]{P49}} columns of height
equal to the rank are allowed; for $\auGL_n$ a column of height $n$ contributes the factor
$\Lambda^n R^n=Re_{[n]}$, on which $g$ acts by $\det g$, and the vector $e_{[n]}$ has norm $1$. So for a
diagram $\nu$ with $c$ columns of height $n$ and $\nu^-$ the diagram of its other columns,
$\mathbb T_\nu=\mathbb T_{\nu^-}\otimes(Re_{[n]})^{\otimes c}$ identifies $\mathsf C_\nu$ with
$\mathsf C_{\nu^-}\otimes\det^c$ isometrically, preserving integral vectors and products of minors
(a product of minors for $\nu$ is one for $\nu^-$ times $\det^c$). Hence the weight $(\nu^-,e+c)$ of
\S\ref{pgs:ss-weyl} may be realised on $\mathsf C_\nu$ twisted by $\det^e$, with its lattice
\eqref{pgs:eq-lattice} and its invariant form.

\paragraph*{Interlacing}
Let $\lambda=(\kappa,e)$ be a weight of $\auGL_N$ at $\tau$, $\mu$ its diagram (columns of heights
$<N$), so the highest weight is $\mu+e$ with $\mu$ padded to $N$ entries. Let $\lambda_0$ be a weight
of $\auGL_n$ whose highest weight $\mu_0'$ interlaces it: $\mu_1+e\ge\mu'_{0,1}\ge\mu_2+e\ge\dots\ge
\mu'_{0,n}\ge\mu_N+e=e$. Then $\nu=\mu_0'-e$ is a diagram with at most $n$ rows and $\mu_{i+1}\le
\nu_i\le\mu_i$, i.e.\ $\mu/\nu$ is a horizontal strip: every column of $\mu$, of height $h$, either
\emph{keeps} its height in $\nu$ or is \emph{lost} and has height $h-1$ in $\nu$ (columns of height $1$
that are lost disappear); the columns of $\nu$ are exactly these. Let $\ell$ be the number of lost
columns. Then $\lambda_0=(\nu^-,e+c)$ with $c$ the number of kept columns of height $n$. Note
$|\nu^-|\le|\nu|\le|\mu|$.

\begin{lemma}[The branching map]\label{pgt:br}
In this situation let $\varepsilon\colon\mathbb T_\nu(\Z^n)\to\mathbb T_\mu(\Z^N)$ send $e_U$ to the
tensor with factor $e_{U_{\mathfrak c}}$ in a kept column and $e_N\wedge e_{U_{\mathfrak c}}$ in a lost
one, and put $i=\pkm\circ\varepsilon\circ(\pkm^\nu|_{\mathsf C_\nu})^{-1}\colon W_\nu\to W_\mu$ and
$\mathrm{br}=H_{\nu}\,i^*\colon W_\mu\to W_\nu$, $H_\nu$ the hook number of {\cite[Lemma~5.22]{P49}} and
$i^*$ the adjoint for the invariant forms. Then:
\begin{enumerate}
\item[(a)] $\mathrm{br}\circ\rho_\mu(g\oplus s)=s^{\ell}\,\rho_\nu(g)\circ\mathrm{br}$ for $g\in\auGL_n$,
$s\in\auGL_1$, over every field of characteristic $0$ (through the integral representation matrices);
\item[(b)] $\|\mathrm{br}\,v\|\le H_\nu\|v\|$ for the invariant norms;
\item[(c)] $\mathrm{br}$ maps $W_\mu(\Z)$ into $W_\nu(\Z)$, i.e.\ its matrix in the bases $(e_T)$ is
integral;
\item[(d)] $\mathrm{br}\ne0$, and $H_\nu\le|\nu|^{3|\nu|}$.
\end{enumerate}
\end{lemma}
\begin{proof}
This is the proof of {\cite[Proposition~5.24]{P49}}, in which the specific structure of the pair
$(\lambda(X_m),\lambda^\circ)$ was used only to take a smaller denominator; we take the full hook
number $H_\nu$ instead, and then no structure is needed.

\emph{The map $i$.} $\varepsilon$ maps the orthonormal basis $(e_U)$ injectively to $\pm$ elements of
the orthonormal basis of $\mathbb T_\mu$ (the sets $U_{\mathfrak c}\subseteq[n]$ do not contain $N$),
so it is isometric and integral, with $\varepsilon^t$ integral of norm $\le1$. For $x\in M_n(R)$,
whose columns are regarded in $R^N$ with last coordinate $0$, and $(g,s)$, one has
$\varepsilon((\Lambda g)u)=s^{-\ell}(\Lambda(g\oplus s))\varepsilon(u)$ factor by factor (a lost column
gets $s e_N\wedge$), so $\varepsilon$ is $\auGL_n$-equivariant; as in {\cite[\S5.2.6]{P49}}, $\pkm$ factors through $\pi_\mu$, and under the
identifications {\cite[(51)]{P49}} $i=\pi_\mu\varepsilon|_{\mathsf C_\nu}$ and
$i^*=\pi_\nu\varepsilon^t|_{\mathsf C_\mu}$. The projections $\pi_\mu,\pi_\nu$ commute with the
actions of the unitary groups (the subspaces $\mathsf C$ are invariant and the forms are), so
$i^*\rho_\mu(g\oplus s)=s^\ell\rho_\nu(g)i^*$ for unitary $g,s$; both sides are polynomial in the
entries of $g,s$ and $\det g^{-1},s^{-1}$ with rational coefficients, so the identity holds on
$\auGL_n\times\auGL_1$ by Zariski density, and then over every field of characteristic $0$. This is
(a). (b): $\|i^*\|\le\|\pi_\nu\|\,\|\varepsilon^t\|\le1$, the invariant forms being the restrictions of
the standard forms ({\cite[\S5.2.6]{P49}}).

\emph{The kernel.} For $x\in M_n(R)$ and $y\in M_N(R)$ let $\mathcal K(x,y)=(\varepsilon\,\pkev^\nu_x,
\pkev_y)$. By {\cite[(51)]{P49}}, as in {\cite[(54)]{P49}},
$\mathcal K(x,y)=(\pkev^\nu_x,\pi_\nu\varepsilon^t\pkev_y)$ and $(\pi_\nu\varepsilon^t\pkev_y,e_{T'})=
i(b^\nu_{T'})(y)$ for every filling $T'$ of $\nu$. All columns of $\mu$ have height $\le n$, so the
last column of $y$ does not occur in $\pkev_y$. Let first the last row of $y$ be $(\mathsf
t,0,\dots,0)$ and $g\in M_n(R)$ the matrix of the first $n$ rows and columns of $y$. In a kept column
of height $h$ the factor of $\mathcal K$ is $\det(x_a\cdot g_b)_{a,b\le h}=\Delta_{[h],[h]}(g^tx)$; in a
lost column of height $h$ it is the determinant with first row $(\mathsf t,0,\dots,0)$ and rows
$(x_a\cdot g_b)_{b\le h}$, $a<h$, that is $\mathsf t\,\Delta_{\{2,\dots,h\},[h-1]}(g^tx)$. So
$\mathcal K(x,y)=\mathsf t^{\,\ell}\mathsf P(g^tx)$ with $\mathsf P=b^\nu_U$ the product of minors of a
filling $U$ of $\nu$. By {\cite[Lemma~5.22 and~(51)]{P49}}, $H_\nu\mathsf P=\pkm^\nu(w_0)$ with
$w_0\in\mathsf C_\nu\cap\mathbb T_\nu(\Z)$, so $H_\nu\mathcal K(x,y)=\mathsf t^{\,\ell}(w_0,
\pkev^\nu_{g^tx})=(\mathsf t^{\,\ell}gw_0,\pkev^\nu_x)$. For general $y$, over $R=\Z[y_{ab}]$, let $y'$
have the columns $y'_1=y_1$, $y'_b=y_{N1}y_b-y_{Nb}y_1$ $(b\ge2)$; its last row is
$(y_{N1},0,\dots,0)$ and $\pkev_{y'}=y_{N1}^{\,\mathsf M}\pkev_y$ with $\mathsf M=\sum_{\mathfrak
c}(h_{\mathfrak c}-1)$. Exactly as in the proof of {\cite[Proposition~5.24]{P49}},
$y_{N1}^{\,\mathsf M}H_\nu\pi_\nu\varepsilon^t\pkev_y=y_{N1}^{\,\ell}g'w_0$ ($g'$ the first $n$ rows and
columns of $y'$), both being vectors of $\mathsf C_\nu\otimes\Q(y)$ with the same pairings against all
$\pkev^\nu_x$; the right side has coordinates in $R$. Hence $y_{N1}^{\,\mathsf M}H_\nu
i(b^\nu_{T'})\in\Z[y_{ab}]$, and since $y_{N1}$ is prime in $\Z[y_{ab}]$ and $H_\nu i(b^\nu_{T'})\in
\Q[y_{ab}]$, $H_\nu i(b^\nu_{T'})\in\Z[y_{ab}]$; by {\cite[(52)]{P49}}, $H_\nu i(b^\nu_{T'})\in\Sigma_\mu$.

\emph{(c).} For $u\in W_\mu(\Z)=\Sigma_\mu^\#$ and every $T'$, $\langle\mathrm{br}\,u,b^\nu_{T'}\rangle=
\langle u,H_\nu i(b^\nu_{T'})\rangle\in\Z$; the coordinate of $\mathrm{br}\,u$ on $e_{T'}$ is this
number (the basis $(e_T)$ being dual to $(b_T)$, \S\ref{pgs:ss-weyl}), so $\mathrm{br}\,u\in
W_\nu(\Z)$. (The products of minors $b^\nu_{T'}$ for $\nu$ and for $\nu^-$ differ by $\det^c$, which
does not affect this.)

\emph{(d).} Let $y\in M_N(\C)$ have first $n$ rows and columns $1_n$ and last row $(1,0,\dots,0)$, so
$\mathsf t=1$, $g=1$ and $\mathcal K(x,y)=\mathsf P(x)$. As $\mathsf P$ is a product of minors of a
generic matrix, it is a non-zero polynomial, so $\mathsf P(x)\ne0$ for some $x$. Then
$(\pkev^\nu_x,i^*\pkev_y)=\mathcal K(x,y)\ne0$ (by the paragraph on the kernel, since $\pkev_y\in\mathsf C_\mu$ by the definition of $\mathsf C_\mu$ and $i^*=\pi_\nu\varepsilon^t|_{\mathsf C_\mu}$), so $i^*\ne0$ and $\mathrm{br}\ne0$. The bound on
$H_\nu$ is {\cite[Lemma~5.22]{P49}}.
\end{proof}

\begin{proposition}[(B2)]\label{pgt:B2}
Let $\lambda^0$, $\lambda^1$ be the weights of $W^0_m$ (rank $n$) and $W^1_m$ (rank $N=n+1$), and
suppose there is a non-zero invariant form on $W^0_m\times W^1_m$ (in the sense of
Lemma~\ref{pgw:branch}(a)). Then there is a non-zero integer matrix $B'$ in the bases $(e_T)$ with
$B'(\rho_0(g)w_0,\rho_1(g\oplus1)w_1)=B'(w_0,w_1)$ for $g\in\prod_\tau\auGL_n$ and
\[
|B'(v',v)|\le L^{3L}\,\|v'\|\,\|v\|\qquad\text{for all }v'\in W^0_m,\ v\in W^1_m\ (\text{over }\C),
\]
for every integer $L$ at least the total number of boxes $\sum_\tau|\mu^1_\tau|$ of $\lambda^1$; in
particular for $L=\lfloor x^{19}b\rfloor$ (Lemma~\ref{pgs:boxes}). The same bound holds after
applying any $\sigma\in\auAut(\C)$ to the coordinates of $v'$, $v$, and it holds for the matrix $B$ of
Definition~\ref{pgz:def-choice}(b). As $L^{3L}=(L^L)^3$, this is (B2) of \S\ref{pgz:ss-b}.
\end{proposition}
\begin{proof}
\emph{Interlacing.} A non-zero invariant form is a non-zero element of
$\auHom_{\prod_\tau\auU(n)}(W^1_m,(W^0_m)^*)$, so $(W^0_m)^*=\bigotimes_\tau(W^0_{m,\tau})^*$ occurs in
the restriction of $W^1_m=\bigotimes_\tau W^1_{m,\tau}$ to $\prod_\tau\auU(n)$; hence, at each $\tau$,
$(W^0_{m,\tau})^*$ occurs in the restriction of $W^1_{m,\tau}$ to $\auU(n)$, and by the branching law
of Weyl (\cite[Theorem~8.1.1]{GoodmanWallach}) its highest weight interlaces that of
$W^1_{m,\tau}$. Let $\lambda_0$ be the weight of $(W^0_m)^*$; by the proof of
Lemma~\ref{pgs:pairing} its doubled parameter is the negative of that of $\lambda^0$. At each $\tau$
we are in the situation of Lemma~\ref{pgt:br} with $\lambda=\lambda^1_\tau$ and $\lambda_0$ at $\tau$;
let $\mathrm{br}=\bigotimes_\tau\mathrm{br}_\tau$, an integral matrix $W^1_m\to W_{\lambda_0}$
(Lemma~\ref{pgt:br}(c), tensor products of integral matrices in the tensor product bases). With the
twists by $\det^{e}$ of $\lambda^1$ and $\det^{e+c}$ of $\lambda_0$ at each $\tau$, Lemma~\ref{pgt:br}(a)
with $s=1$ gives $\mathrm{br}\,\rho_1(g\oplus1)=\rho_{\lambda_0}(g)\,\mathrm{br}$, since $\det(g\oplus1)
=\det g$ and the kept columns of height $n$ contribute $\det(g)^c$.

\emph{The matrix.} Let $A=G_{\lambda^0}J$ be the duality pairing of Lemma~\ref{pgs:pairing} between
weight $\lambda^0$ (in the role of $\lambda'$) and $\lambda_0$ (in the role of $\lambda$), and put
$B'(v',v)=\langle v',\mathrm{br}\,v\rangle_\pgsdu={}^tv'A\,\mathrm{br}\,v$. It is an integer matrix.
By (P1), $B'(\rho_0(g)v',\rho_1(g\oplus1)v)=\langle\rho_0(g)v',\rho_{\lambda_0}(g)\mathrm{br}\,v
\rangle_\pgsdu=B'(v',v)$. It is non-zero since $\mathrm{br}\ne0$ (Lemma~\ref{pgt:br}(d)) and $A$ is
invertible ($J$ is, and $G_{\lambda^0}$ is positive definite).

\emph{The bound.} By (P3), $\langle v',u\rangle_\pgsdu=\langle u^\flat,v'\rangle$ with $\|u^\flat\|=
\|u\|$, so by Cauchy--Schwarz for the invariant hermitian form and Lemma~\ref{pgt:br}(b),
\[
|B'(v',v)|\le\|v'\|\,\|\mathrm{br}\,v\|\le\Bigl(\prod_\tau H_{\nu_\tau}\Bigr)\|v'\|\,\|v\|\le
\prod_\tau L_\tau^{3L_\tau}\,\|v'\|\,\|v\|\le L^{3L}\|v'\|\,\|v\|,
\]
with $L_\tau=|\nu_\tau|\le|\mu^1_\tau|$ and $\sum_\tau L_\tau\le L$ (the norm of a tensor product of
operators being the product of the norms). For $\sigma\in\auAut(\C)$: $B'$ has integer entries, so
$\sigma(B'(v',v))=B'(\sigma v',\sigma v)$ and the bound applies to $\sigma v',\sigma v$. For $B$: by
Lemma~\ref{pgw:branch}(a) $B'=cB$ with $c$ rational, and as $B$ is primitive and $B'$ integral, $c$ is
a non-zero integer; so $|B(v',v)|\le|B'(v',v)|$. Finally $\sum_\tau|\mu^1_\tau|\le40fn_mb\le x^{19}b$ by
Lemma~\ref{pgs:boxes}.
\end{proof}

\begin{remark}\label{pgt:rem-B2}
If there is no non-zero invariant form, $B=0$ by Definition~\ref{pgz:def-choice}(b) and the bound
used in the proof of Theorem~\ref{pgz:b} holds trivially; so Proposition~\ref{pgt:B2} is all that
proof needs. The constant obtained is $L^{3L}=(L^L)^3$, the constant asked for in (B2).
\end{remark}

\subsubsection{(B1): the operator bound for one jet kernel}\label{pgt:ss-B1}

{\cite[Proposition~14.23]{P46}} is stated for the kernels of the chain $X_j\to X_{j+1}$, with their row
lengths, and {\cite[Lemma~4.10(KS2)]{P54}} needs $M^2\in\pcm$, which fails for $\pcm=\pcb=(\pi_M)^b$ when
$b>2$. Their proofs, however, give an operator bound for any jet kernel between standard definite
spaces of increasing ranks with a translate of the vacuum as datum, and in the archimedean bound the
modulus plays no role: an indicator datum is an average of translates of the vacuum with coefficients
of absolute value $1$ (after normalisation). We write this out.

\begin{lemma}[Jets of translates of the vacuum]\label{pgt:jet}
Let $X=E^n$, $X'=E^{n'}$ be positive definite standard spaces, $n<n'$, $W=\auHom(X,X')$,
$D=fnn'$. Let $\nu=(\nu_\tau)$ be a diagram with at most $n$ rows at each $\tau$, $e=\sum_\tau|\nu_\tau|$,
and $S_\nu(z)$, $z\in\prod_\tau M_{n'\times n}(\C)$, the matrix of the Schur functor
$W_\nu(\C^n)\to W_\nu(\C^{n'})$ of {\cite[\S5.1.5]{P49}} (tensor product over $\tau$), in the bases
$(e_T)$. Let $\phi=\psrho^\circ(x)\omega(k)\psvac^\circ$ with $x\in W(\auA_f)$, $k\in\auU(W)(\auA_f)$. Then
for all $\sigma\in\auAut(\C)$ and all $\xi\in W_\nu(\C^n)$,
\[
\bigl\|\sigma\bigl(S_\nu(\partial)[\Theta_\phi(s_0\,\cdot)](0)\,\xi\bigr)\bigr\|\le G_1\,\|\sigma\xi\|,\qquad
G_1=\Omega^{-e/2}(2e)^{e/2}\,4^{D}\exp\bigl(C_0(q_1nn'+1)^2\bigr),
\]
$C_0=C_0(q_0,d)$ the constant of {\cite[Lemma~5.7{\rm(ii)}]{P49}} and $(2e)^{e/2}=1$ if $e=0$.
\end{lemma}
\begin{proof}
This is Steps 2--4 of the proof of {\cite[Proposition~14.23]{P46}} at the scale $s_0$ instead of
$\pi_Ms_0$, with a cruder treatment of the Schur matrix. As there we may assume $\sigma|_K=1$: if
$\sigma'=c\circ\sigma$ ($c$ complex conjugation), then $|\sigma'w|=|\sigma w|$ for numbers and
$\|\sigma'v\|=\|\sigma v\|$ for vectors, the Gram matrices being real.

Let $\eta$ be a covector on $W_\nu(\C^{n'})$ with $\|\sigma\eta\|\le1$ for the dual norm, and
$Q(z)=\eta\cdot S_\nu(z)\xi$, a polynomial homogeneous of degree $e$. By {\cite[Lemma~5.7(i) and~(33)]{P49}} for $\Lambda_\sigma\phi$ at the scale $t_\sigma$ (Step~2 of the proof of
{\cite[Proposition~14.23]{P46}}),
\[
\sigma\bigl(Q(\partial)[\Theta_\phi(s_0\,\cdot)](0)\bigr)=\int Q^\sigma(t_\sigma\bar y)\,e^{-|y|^2/2}\,
\theta\bigl(\psrho_\infty(y)(1\otimes\Lambda_\sigma\phi)\bigr)\,d\mu(y),\qquad d\mu(y)=\pi^{-D}dy,
\]
$Q^\sigma$ obtained by applying $\sigma$ to the coefficients and $\varpi_\sigma$ to the variables. The
entries of $S_\nu$ have integer coefficients, so $Q^\sigma(z)=\sigma\eta\cdot S_\nu(z')\sigma\xi$ with
$z'$ a permutation of the components of $z$, and by {\cite[(38)]{P49}} (at each $\tau$, and the
operator norm of a tensor product being the product of the norms)
\[
|Q^\sigma(z)|\le\|\sigma\xi\|\prod_{\tau,p}s_{\tau,p}^{\,l_{\tau,p}}\le\|\sigma\xi\|\,|z|^{e},
\]
where $s_{\tau,p}\le|z_\tau|_{\mathrm{HS}}\le|z|$ are the singular values and $l_{\tau,p}$ the rows of
$\nu_\tau$, $\sum_{\tau,p}l_{\tau,p}=e$. For $z=t_\sigma\bar y$, $|z|^2=\Omega^{-1}|y|^2$
({\cite[Lemma~5.7]{P49}}). By {\cite[Lemma~5.7(ii)]{P49}} the theta value is at most $\exp(C_0(q_1nn'+1)^2)$.
Using $u^{e/2}e^{-u/4}\le(2e)^{e/2}e^{-e/2}\le(2e)^{e/2}$ for $u\ge0$ (maximum at $u=2e$) with $u=|y|^2$,
\[
\int|y|^{e}e^{-|y|^2/2}\,d\mu(y)\le(2e)^{e/2}\int e^{-|y|^2/4}\,d\mu(y)=(2e)^{e/2}\,4^{D}.
\]
Hence $|\sigma(\eta\cdot S_\nu(\partial)[\Theta_\phi(s_0\,\cdot)](0)\xi)|\le G_1\|\sigma\xi\|$; since $S_\nu$ has
integer coefficients, the left side is $|\sigma\eta\cdot\sigma(S_\nu(\partial)[\cdots](0)\xi)|$, and
taking the supremum over $\eta$ gives the claim.
\end{proof}

\begin{proposition}[(B1)]\label{pgt:B1}
Let $\vartheta=1$, $e^z_a=(p,q)$ an edge of $\pgrlast$, $X\in\pcMat_{q\times p}(\pcR_\pcb)$, $x\in
G_{p,f}$, $h\in G_{q,f}$, $\sigma\in\auAut(\C)$ and $\xi\in W^z_{a-1}$. With the kernels of
Definition~\ref{pgz:def-choice}(a),
\[
\bigl\|\sigma\bigl(\pgwK^z_a(\omega_f(x,h)[\delta_X]_f)\,\xi\bigr)\bigr\|\le\exp(x^{52}b\log b)\,
\|\sigma\xi\|\le G_0\,\|\sigma\xi\| .
\]
This is (B1) of \S\ref{pgz:ss-b}.
\end{proposition}
\begin{proof}
\emph{The spaces.} The nodes $p=p^z_{a-1}$, $q=p^z_a$ carry positive definite standard spaces
(proof of Corollary~\ref{pgt:B3}) of ranks $N_p<N_q\le2n_m$ (Lemma~\ref{pgs:boxes}; the rank grows by
$t+1$ along $h_z$). By Definition~\ref{pgz:def-choice}(a), $\pgwK^z_a(\Phi)=c_{\mathrm{vac}}^{\,N_pN_q}
S_\nu(\partial)[\Theta_\Phi(s_0\,\cdot)](0)$ with $S_\nu$ the Schur matrix of the diagram $\nu$ of the
edge, of degree $e=e(\nu)\le40fn_mb\le x^{19}b$ (Lemma~\ref{pgz:deg}); $\nu_\tau$ has at most $N_p$
rows (Lemma~\ref{pgw:howe} and the proof of Lemma~\ref{pgw:one}: $W^z_{a-1}\cong A_{\nu}$ is, up to a
character, $W_\nu$ for $\auGL_{N_p}$).

\emph{The datum.} As $\vartheta=1$, $Q_\vartheta=\emptyset$ and $S(\pcb)$ consists of the places
$v=w\aubar w$ with $w\mid\pi_M$, where $a_w=1$ (proof of Lemma~\ref{pgz:std}). By the first paragraph of
the proof of Lemma~\ref{pgz:std},
\[
c_{\mathrm{vac}}^{\,N_pN_q}\,[\delta_X]_f=\auNm(\pcb)^{-N_pN_q}\sum_{x'}\zeta_{x'}\,\psrho^\circ(x')
\psvac^\circ ,
\]
a sum of $\prod_{w\mid\pcb}[M_w:\pcb M_w]=\auNm(\pcb)^{N_pN_q}$ terms with roots of unity $\zeta_{x'}$.
By Step~1 of the proof of {\cite[Proposition~14.23]{P46}}, in the form used in the proof of
{\cite[Lemma~4.10(KS2)]{P54}} for twisted Weil representations $\omega_{a,b}$ with unitary Hecke characters
of algebraic weights (which is the case of $\omega_f=\omega_{\alpha,\beta,f}$, as in the proof of
Lemma~\ref{pgz:std}), $\omega_f(x,h)\psrho^\circ(x')\psvac^\circ=\epsilon_{x'}\psrho^\circ(x'')\omega(k)
\psvac^\circ$ with $k\in\auU(W)(\auA_f)$ the image of $(x,h)$, some $x''\in W(\auA_f)$, and
$|\sigma\epsilon_{x'}|=1$. (Here $\epsilon_{x'}$ is a root of unity times the values at $\det x$, $\det h$
of the splitting characters $\alpha,\beta$ of Definition~\ref{pgv:def-edges}, which are products of the
characters $\mu^0_l$ and of powers of $\upsilon|_{\auU(1)}$, all of algebraic weight. The unramified
factors $\eta_w$ of the model at a split place, {\cite[(1)]{P54}}, do not enter $\epsilon_{x'}$ with their
non-unitary part: $|\eta_w(1,h)|=|\det h|_w^{-r/2}$ is the normalisation that makes the substitution
$\Psi\mapsto\Psi(h^{-1}Tg)$ unitary, i.e.\ it belongs to the operator $\omega(k)$, and the remaining
unitary part of $\eta_w$ is a value of a power of $\upsilon$. Finally, for a character $\lambda$ of
$\auU(1)(\auA)$ trivial on $\auU(1)(E^+)$ with algebraic weight and $a\in\auU(1)(\auA_f)$: as
$\auU(1)$ is anisotropic, $a^h=\gamma\kappa$ for some $h\ge1$, $\gamma\in\auU(1)(E^+)$ and $\kappa$ in a
compact open subgroup on which $\lambda$ has finite order, so $\lambda(a)^h$ is a root of unity times
$\prod_\tau\tau(\gamma)^{-k_\tau}$ with $\gamma\aubar\gamma=1$; all conjugates of this number have modulus
$1$, $E$ being a CM field, hence so have those of $\lambda(a)$. This is the argument of Step~1 of the
proof of {\cite[Proposition~14.23]{P46}}.) Hence
\[
\pgwK^z_a(\omega_f(x,h)[\delta_X]_f)=\auNm(\pcb)^{-N_pN_q}\sum_{x'}\zeta_{x'}\epsilon_{x'}\,
S_\nu(\partial)\bigl[\Theta_{\phi_{x'}}(s_0\,\cdot)\bigr](0),\qquad \phi_{x'}=\psrho^\circ(x'')\omega(k)\psvac^\circ ,
\]
and as the matrices $S_\nu(\partial)[\cdots](0)$ are applied to $\xi$ and then conjugated by $\sigma$
(which commutes with the finite sum and the rational scalar $\auNm(\pcb)^{-N_pN_q}$), the triangle
inequality and Lemma~\ref{pgt:jet} give the bound $G_1\|\sigma\xi\|$, an average of
$\auNm(\pcb)^{N_pN_q}$ terms each at most $G_1\|\sigma\xi\|$.

\emph{Sizes.} By {\cite[Lemma~13.4(S6)]{P46}} and (S4), $\Omega^{-1}=e^{-\aneu}\anrd_E/2\pi\le
x^{12}e^{|\aneu|}\le x^{13}$ ((S1): $|\aneu|\le\log x$). With $e\le x^{19}b$ and $\log(2e)\le
1+19\log x+\log b\le2\log b$ ((S1)), $\tfrac e2\log\Omega^{-1}+\tfrac e2\log(2e)\le x^{19}b(7\log x+\log
b)\le2x^{19}b\log b$. Next $D=fN_pN_q\le x^{12}(80x^6)^2\le x^{25}$, so $D\log4\le x^{26}$; and
$C_0\le\log\log b\le x$ ({\cite[Lemma~8.7(c)]{P49}}, (S1)), $q_1N_pN_q+1\le x^{12}\cdot6400x^{12}+1\le x^{25}$,
so $C_0(q_1N_pN_q+1)^2\le x^{51}$. In all, $\log G_1\le2x^{19}b\log b+x^{26}+x^{51}\le x^{52}b\log b\le
x^{99}b^2$.
\end{proof}

\begin{remark}\label{pgt:rem-B1}
The bound does not depend on $\pcm$, $X$ or the points $x,h$; in particular it holds verbatim for
every twist $\vartheta$ (with $\pcm(\vartheta)$ in place of $\pcb$, using the full statement of
Lemma~\ref{pgz:std}), though only $\vartheta=1$ is used. The row lengths of $\nu$ enter only through
their sum $e$; a sharper bound, with $\sum l\log l$ in place of $e\log e$, follows from
{\cite[Lemma~5.8]{P49}} as in Step~3 of the proof of {\cite[Proposition~14.23]{P46}}, but is not needed.
\end{remark}




\section{Kernel forms, the restriction pairing and the coefficient lemmas}\label{C1:g7}
% ===== C-grid11.tex =====
% proofs/C-grid11.tex
% Inputs (C) and (Q) of proofs/C-grid9.tex (Section pgq:ss-inputs). Labels start with pgk:.
% Conventions and macros: those of proofs/C-grid9.tex.
\providecommand{\auU}{\mathrm U}
\providecommand{\auO}{\mathcal O}
\providecommand{\auA}{\mathbb A}
\providecommand{\auAz}{\mathcal A_0}
\providecommand{\auPl}{\Phi}
\providecommand{\aubar}[1]{\overline{#1}}
\providecommand{\pgrY}{\widetilde Y}
\providecommand{\pgkvol}{\operatorname{vol}}
\providecommand{\pgkP}{\mathsf P}

\subsection{The cusp estimate and the kernel forms}\label{pgk:sec}

This section proves the inputs (Q) and (C) of \S\ref{pgq:ss-inputs} (Theorems~\ref{pgk:Q} and
\ref{pgk:C}). Notation is that of \S\ref{pgq:ss-stat}: $V=Y_0$ of rank $n$, $Z=V\oplus E\tilde e$,
$G=\auU(V)$, $H=\auU(Z)$, $\Gamma_G=G(E^+)$, $\Gamma_H=H(E^+)$, $f=\#\auPl$, the normalised bases at
the places $\tau\in\auPl$, operator norms $\|\cdot\|$, $\operatorname{dist}(u,1)=\max_\tau\|u_\tau-1\|$,
the coefficients $c,c'$ and the types $e=d\aubar c$ of $\auAz$ and $e'=d'\aubar{c'}$ of $e_{(1,0,0)}$.
Forms of a type are meant in the sense of {\cite[Definition~8.1]{P49}} (for $H$, with $H$ in place of
$G_j$): bounded continuous functions, left invariant under the rational points, right invariant under
an open subgroup of the finite adelic points, with $R(e)y=y$, $(R(e)y)(g)=\int e(x)y(gx)\,dx$. For
$y\in H_\infty$ (or $G_\infty$) put $\pgkP(y)=\prod_\tau\|y_\tau\|\ge1$; then $\pgkP(y^{-1})=\pgkP(y)$
and $\pgkP(yy')\le\pgkP(y)\pgkP(y')$.

\subsubsection{Kernel forms of a level: input (Q)}\label{pgk:ss-Q}

\begin{remark}[What is already in the paper]\label{pgk:rem-Q}
{\cite[Lemma~4.9(c)]{P51}} bounds, uniformly in $g$, the sum over $\gamma$ of $\pgkP((g^{-1}\gamma x)_\infty)
^{-B'}\mathbf 1_{\Omega_f}((g^{-1}\gamma x)_f)$, with a constant depending on the reduction of $x$, and
{\cite[Proposition~4.11(a)]{P51}} gives the convergence of kernels $\sum_\gamma(e\otimes\varphi')(x^{-1}
\gamma y)$; both are written for a group $\auU(V)$ with the norms adapted to a Witt basis. We need the
bound only for one fixed point $z$ and uniformly in $g$, for the group $H$ with the norms of the
normalised basis; the disjointness argument of the proof of {\cite[Lemma~4.9(c)]{P51}} gives it directly,
and we write it out in Lemma~\ref{pgk:sum}, which avoids Siegel sets and the change of norms. What
remains, that the sum is a form of type $e'$, is the reproducing property $e'*e'=e'$.
\end{remark}

Fix $B'=2(n+1)^2$. By Lemma~\ref{pgq:local}(a) with $j=B'\le h_Z$ (Lemma~\ref{pgq:sizes}(a): $h_Z\ge2A
\ge B'$), for $y\in H_\infty$
\begin{equation}\label{pgk:eq-decay}
|e'(y)|=d'\prod_\tau|c'_\tau(y_\tau)|\le d'\,2^{fB'}\,\pgkP(y)^{-B'} .
\end{equation}
By {\cite[Lemma~4.7(c)]{P51}}, applied to the unitary group $H_\tau$ of rank $n+1$ (its proof uses only
the Cartan decomposition $H_\tau=K_\tau A^+K_\tau$, and the operator norm of the normalised basis is
bi-invariant under the maximal compact subgroup $H_\tau\cap\auU(n+1)$, since the matrix of the form is
hermitian and unitary, {\cite[\S5.1]{P51}}), $I_\tau=\int_{H_\tau}\|h\|^{-B'}dh<\infty$, as
$B'>(n+1)^2$.

\begin{lemma}[A uniform sum]\label{pgk:sum}
Let $K\subseteq H(\auA_f)$ be a compact open subgroup and $x\in H(\auA)$. Put $\mathsf M(y)=
\pgkP(y_\infty)^{-B'}\mathbf 1_K(y_f)$ on $H(\auA)$. There is $C_x<\infty$ with
$\sum_{\gamma\in\Gamma_H}\mathsf M(g^{-1}\gamma x)\le C_x$ for all $g\in H(\auA)$.
\end{lemma}
\begin{proof}
$\Gamma_H$ is discrete in $H(\auA)$, so there is an open neighbourhood $W$ of $1$ with $W\cap\Gamma_H=
\{1\}$. Choose $0<\rho\le\frac14$ and an open subgroup $K''\subseteq K$ such that $Z_{4\rho}\times K''
\subseteq x^{-1}Wx$, where $Z_r=\{u\in H_\infty:\operatorname{dist}(u,1)\le r\}$; this is possible since
these sets form a basis of neighbourhoods of $1$. Let $U=Z_\rho\times K''$. For $u,u'\in Z_\rho$,
$\|u_\tau^{-1}\|=\|u_\tau\|\le2$ and $u'_\tau u_\tau^{-1}-1=(u'_\tau-u_\tau)u_\tau^{-1}$, so
$u'u^{-1}\in Z_{4\rho}$; hence $UU^{-1}\subseteq x^{-1}Wx$.

\emph{Disjointness.} If $g^{-1}\gamma xu=g^{-1}\gamma'xu'$ with $u,u'\in U$, then $\gamma'^{-1}\gamma=
xu'u^{-1}x^{-1}\in W\cap\Gamma_H$, so $\gamma=\gamma'$. Thus the sets $g^{-1}\gamma xU$ are disjoint.

\emph{Comparison.} For $y\in H(\auA)$, $u\in U$: $\pgkP((yu)_\infty)\le2^f\pgkP(y_\infty)$ and
$\mathbf 1_K(y_fu_f)=\mathbf 1_K(y_f)$, so $\mathsf M(y)\le2^{fB'}\mathsf M(yu)$. Averaging over $U$ and
using left invariance of Haar measure,
\[
\pgkvol(U)\sum_\gamma\mathsf M(g^{-1}\gamma x)\le2^{fB'}\sum_\gamma\int_{g^{-1}\gamma xU}\mathsf M\le
2^{fB'}\int_{H(\auA)}\mathsf M=2^{fB'}\pgkvol(K)\prod_\tau I_\tau .
\]
So $C_x=2^{fB'}\pgkvol(U)^{-1}\pgkvol(K)\prod_\tau I_\tau$ works ($\pgkvol(U)>0$, $U$ having non-empty
interior).
\end{proof}

\begin{theorem}[Input (Q)]\label{pgk:Q}
Let $K\subseteq H(\auA_f)$ be an open subgroup of the stabiliser of $\widehat L_Z$ and $z\in H(\auA)$.
For $g\in H(\auA)$ the series
\[
\mathsf F(g)=\sum_{\gamma\in\Gamma_H}e'\bigl((z\gamma g)_\infty\bigr)\,\mathbf 1_K\bigl((z\gamma g)_f\bigr)
\]
converges absolutely, $\sup_g\sum_\gamma|\cdots|<\infty$, and $Q=\aubar{\mathsf F}$ is a form of type
$e'$ on $H$ fixed by $K$. This is \textup{(Q)} of \S\ref{pgq:ss-inputs}.
\end{theorem}
\begin{proof}
$K$ is compact (closed in a compact group) and open. Write $\Phi(y)=e'(y_\infty)\mathbf 1_K(y_f)$.
By \eqref{pgk:eq-decay}, $\pgkP(y^{-1})=\pgkP(y)$ and $K=K^{-1}$,
$|\Phi(y)|\le d'2^{fB'}\mathsf M(y^{-1})$. With $y=z\gamma g$, $y^{-1}=g^{-1}\gamma^{-1}z^{-1}$, and
Lemma~\ref{pgk:sum} with $x=z^{-1}$ (sum over $\gamma^{-1}$) gives
\begin{equation}\label{pgk:eq-unif}
\sum_\gamma|\Phi(z\gamma g)|\le C:=d'2^{fB'}C_{z^{-1}}\qquad\text{for all }g\in H(\auA).
\end{equation}
So $\mathsf F$ is defined, $|\mathsf F|\le C$, and $\mathsf F$ is measurable (a pointwise limit of
finite sums of measurable functions). \emph{Invariance.} For $\delta\in\Gamma_H$, $\mathsf F(\delta g)=
\mathsf F(g)$ (substitute $\gamma\mapsto\gamma\delta^{-1}$); for $k\in K$, $(z\gamma gk)_f=(z\gamma g)_fk$
and $(z\gamma gk)_\infty=(z\gamma g)_\infty$, so $\mathsf F(gk)=\mathsf F(g)$.

\emph{Reproducing property.} Since $e'(h^{-1})=\aubar{e'(h)}$ ($c'_\tau(h^{-1})=\aubar{c'_\tau(h)}$ as in
Lemma~\ref{pgq:local}(d)), $R(e')Q=Q$ is equivalent to
\[
\mathsf F(g)=\int_{H_\infty}e'(h^{-1})\,\mathsf F(gh)\,dh\qquad(g\in H(\auA)).
\]
For each $\gamma$, with $y=(z\gamma g)_\infty$, $\int e'(h^{-1})\Phi(z\gamma gh)\,dh=\mathbf 1_K((z\gamma
g)_f)\int e'(yh)e'(h^{-1})\,dh=\Phi(z\gamma g)$, because $\int e'(yh)e'(h^{-1})dh=(e'*e')(y)=e'(y)$. The
last identity is Schur's orthogonality relation for the discrete series $\pi'=\bigotimes_\tau\pi'_\tau$
of formal degree $d'$ and unit vector $v'$ (with the convention recalled in the proof of
Lemma~\ref{pgq:schur}): $e'(y)=d'\langle\pi'(y)v',v'\rangle$, and
$\int\langle v',\pi'(h)v'\rangle\aubar{\langle\pi'(y^{-1})v',\pi'(h)v'\rangle}dh=d'^{-1}\langle v',\pi'(y^{-1})
v'\rangle$, which is the integral $d'^{-2}\int e'(yh)e'(h^{-1})dh$ written out. The exchange of sum and
integral is justified by Tonelli: $\int|e'(h^{-1})|\sum_\gamma|\Phi(z\gamma gh)|\,dh\le C\|e'\|_{L^1}
<\infty$ by \eqref{pgk:eq-unif} and the integrability of $e'$ \eqref{pgk:eq-decay}.

\emph{Continuity.} By the reproducing property, $\mathsf F(g)=\int e'(h^{-1}g_\infty)\mathsf F(h,g_f)\,dh$
after the substitution $h\mapsto g_\infty^{-1}h$. As $g_\infty\mapsto e'(\,\cdot\,g_\infty)$ is
continuous into $L^1(H_\infty)$ (continuity of translation in $L^1$) and $\mathsf F$ is bounded,
$\mathsf F$ is continuous in $g_\infty$, uniformly in $g_f$; and it is locally constant in $g_f$ by
right $K$-invariance. Hence $\mathsf F$ and $Q$ are continuous, and $Q$ is a form of type $e'$ fixed
by $K$ with $\aubar{Q(g)}=\mathsf F(g)$.
\end{proof}

\subsubsection{Coordinates, lattices and two arithmetic facts}\label{pgk:ss-coord}

\emph{Coordinates.} $V$ is a space in the sense of \S\ref{pc:ss-data}: it has a frame, an $E$-basis
whose $\widehat\auO_E$-span $\widehat L_V$ is unimodular. We identify $V(E)=E^n$ by the frame; the
form has a hermitian Gram matrix $\mathsf H\in M_n(E)$, integral with unit determinant at every finite
place. For $\tau\in\auPl$ let $\mathsf P_\tau\in\mathrm{GL}_n(\C)$ be the matrix whose columns are the
normalised basis in frame coordinates, so $\mathsf P_\tau^*\tau(\mathsf H)\mathsf P_\tau=\mathsf J_\tau$.
For $x\in G(\auA)$ and $v\in V(\auA)$ the \emph{coordinates at $x$} are $\xi_x(v)=x^{-1}v$, read at
$\tau$ in the normalised basis (that is, $\mathsf P_\tau^{-1}x_\tau^{-1}\tau(v)$) and at the finite
places in the frame; this is the reading of \S\ref{pgq:ss-stat}, so that $L_V(x)=\{v\in E^n:(x^{-1}v)_f
\text{ integral}\}$, $|v|_x^2=\sum_\tau|\xi_x(v)_\tau|^2$, and $\lambda(x)$ is the first minimum defined
there. In these coordinates the form is $\mathsf S$, with $\mathsf S_\tau=\mathsf J_\tau$ (hermitian and
unitary, so $\|\mathsf S_\tau\|=1$) and $\mathsf S_f=x_f^*\mathsf Hx_f=\mathsf H$; for rational vectors
$v,w$, $\langle v,w\rangle=\xi_x(v)^*\mathsf S\,\xi_x(w)$ at every place. We extend $\xi_x$ to
matrices whose columns are vectors.

\begin{lemma}[Lattice vectors]\label{pgk:lat}
Let $x\in G(\auA)$ and $v,w\in L_V(x)$.
\begin{enumerate}
\item[(a)] $\langle v,w\rangle\in\auO_E$ and $|\tau\langle v,w\rangle|\le|v|_x|w|_x$ for every $\tau$.
If $|v|_x|w|_x<1$ then $\langle v,w\rangle=0$.
\item[(b)] If $v\ne0$ and $y_1,\dots,y_n\in L_V(x)$ span $V$ with $|y_j|_x\le R$, then $|v|_xR\ge1$.
\item[(c)] $L_V(x)$ is a lattice of rank $2fn$ in $\prod_\tau\C^n\cong\R^{2fn}$ \textup{(}embedded by
the archimedean coordinates\textup{)} of covolume $(2^{-f}|D_E|^{1/2})^n$, $D_E$ the discriminant of
$E$.
\end{enumerate}
\end{lemma}
\begin{proof}
(a) At the finite places $\langle v,w\rangle=\xi_x(v)^*\mathsf H\xi_x(w)$ with integral entries, so
$\langle v,w\rangle\in E\cap\widehat\auO_E=\auO_E$. At $\tau$, $|\tau\langle v,w\rangle|=|\xi_x(v)_\tau^*
\mathsf J_\tau\xi_x(w)_\tau|\le|\xi_x(v)_\tau||\xi_x(w)_\tau|\le|v|_x|w|_x$. A non-zero $a\in\auO_E$
has $\prod_{\tau\in\auPl}|\tau(a)|^2=|\mathrm N_{E/\Q}(a)|\ge1$, so $\max_\tau|\tau(a)|\ge1$.
(b) Otherwise $|v|_x|y_j|_x<1$ for all $j$, so $v$ is orthogonal to all $y_j$ by (a), hence to $V$, and
$v=0$ as the form is non-degenerate.
(c) $L_V(x)=E^n\cap x_f\widehat\auO_E^{\,n}$. Since $x_f$ preserves the unimodular form,
$\det x_f$ is a unit at every finite place, so $x_f\widehat\auO_E^{\,n}$ and $\widehat\auO_E^{\,n}$ have the same
volume in $\auA_{E,f}^n$; hence the covolumes of $L_V(x)$ and of $\auO_E^n$ in $E_\infty^n$ (embedded by
$v\mapsto(\tau(v))_\tau$) agree, both being the covolume of $E^n$ in $\auA_E^n$ divided by that volume. The latter is $(2^{-f}
|D_E|^{1/2})^n$ ({\cite[Lemma~8.32]{P51}}; see also \cite[Chapter~I, (5.1) and Proposition~(5.2)]{Neukirch}). The archimedean coordinate map is $v\mapsto(\mathsf
P_\tau^{-1}x_\tau^{-1}\tau(v))_\tau$, of real Jacobian $\prod_\tau|\det\mathsf P_\tau|^{-2}|\det x_\tau|
^{-2}$. Here $|\det x_\tau|=1$ ($x_\tau$ is unitary for a non-degenerate form), and $|\det\mathsf P_\tau|
^{-2}=|\tau(\det\mathsf H)|$ because $\det\mathsf J_\tau=\pm1$; as $\det\mathsf H\in E^+$ is a unit at
every finite place, $\prod_\tau|\tau(\det\mathsf H)|=|\mathrm N_{E^+/\Q}\det\mathsf H|=1$.
\end{proof}

\begin{lemma}[Rounding]\label{pgk:round}
Let $N_d=q_1d$, a natural number in the different $\mathfrak d_E$ of $E$ ({\cite[\S5.1.5]{P54}}). For
every $c\in\auA_E$ there is $\varrho\in E$ with $c_f-\varrho\in\widehat\auO_E$ and $|c_\tau-\tau(\varrho)|
\le2f^2N_d$ for every $\tau\in\auPl$.
\end{lemma}
\begin{proof}
As $E+\widehat\auO_E=\auA_{E,f}$ (approximation theorem), there is $\varrho_0\in E$ with $c_f-\varrho_0
\in\widehat\auO_E$. Embed $E_\infty=\prod_\tau\C\cong\R^{2f}$ with the inner product $(a,b)=\Rel\sum_\tau
\tau(a)\aubar{\tau(b)}=\frac12\mathrm{Tr}_{E/\Q}(a\bar b)$ on $E$. The dual of $\auO_E$ is
$\{b:\frac12\bar b\in\mathfrak d_E^{-1}\}=2\mathfrak d_E^{-1}$. For $0\ne b\in2\mathfrak d_E^{-1}$,
$N_db/2\in N_d\mathfrak d_E^{-1}\subseteq\auO_E$ is non-zero, so $|N_db/2|^2=\sum_\tau|\tau(N_db/2)|^2
\ge f\ge1$ (arithmetic and geometric means); hence the dual has no non-zero vector shorter than $2/N_d$.
By {\cite[Lemma~5.4]{P49}} with $m=2f$, every point of $E_\infty$ is within $(2f)^2N_d/2$ of $\auO_E$:
there is $o\in\auO_E$ with $|(c_\tau-\tau(\varrho_0)-\tau(o))_\tau|\le2f^2N_d$. Take $\varrho=\varrho_0+o$.
\end{proof}

\subsubsection{A short vector gives a cusp}\label{pgk:ss-cusp}

\begin{lemma}[Minkowski basis]\label{pgk:mink}
Let $x\in G(\auA)$ and $\kappa\ge(fn\,\mathrm{rd}_E)^{1/2}$, where $\mathrm{rd}_E=|D_E|^{1/2f}$. There
is an $E$-basis $y_1,\dots,y_n$ of $V$ in $L_V(x)$ with $\alpha_j=|y_j|_x$ non-decreasing and
$\prod_j\alpha_j\le\kappa^n$.
\end{lemma}
\begin{proof}
Let $D=2fn$ and $\lambda_1\le\dots\le\lambda_D$ the successive minima of $L_V(x)\subset\R^D$ for the
Euclidean unit ball, of volume $\pi^{fn}/(fn)!$. By Minkowski's second theorem
(\cite[Chapter~VIII, \S4.3, Theorem~V]{Cassels}; see also \cite[Theorem~1.3]{HenkMinima}) and Lemma~\ref{pgk:lat}(c),
\[
\lambda_1\cdots\lambda_D\le\frac{2^D\,2^{-fn}|D_E|^{n/2}}{\pi^{fn}/(fn)!}\le\Bigl(\frac2\pi\Bigr)^{fn}(fn)^{fn}
\mathrm{rd}_E^{\,fn}\le(fn\,\mathrm{rd}_E)^{fn},
\]
as $|D_E|^{1/2}=\mathrm{rd}_E^{\,f}$. Exactly as in the proof of {\cite[Lemma~8.34]{P51}}, choose
$\Z$-independent $v_1,\dots,v_D$ with $|v_j|_x=\lambda_j$; an $E$-subspace of dimension $e$ contains at
most $2fe$ $\Z$-independent vectors, so we may pick $E$-independent $y_i$ among $v_1,\dots,
v_{2f(i-1)+1}$ with $|y_i|_x\le\lambda_{2f(i-1)+1}$, $1\le i\le n$. Then $|y_i|_x^{2f}\le\prod_{t=1}^{2f}
\lambda_{2f(i-1)+t}$, so $\prod_i|y_i|_x^{2f}\le(fn\,\mathrm{rd}_E)^{fn}$ and $\prod_i|y_i|_x\le(fn\,
\mathrm{rd}_E)^{n/2}\le\kappa^n$. Reorder by length.
\end{proof}

\begin{lemma}[Cusp frame]\label{pgk:frame}
Let $\kappa\ge\max(1,(fn\,\mathrm{rd}_E)^{1/2})$, $\theta>0$ with $\theta\kappa\le1$, and $x\in G(\auA)$
with $\lambda(x)<\theta^n$. There are $1\le i\le n/2$, a non-zero isotropic subspace $A\subset V$ of
dimension $i$, and lattice vectors in $L_V(x)$: $y_1,\dots,y_i$ spanning $A$, of length at most
$\ell_X$, and $y_1,\dots,y_{n-i}$ spanning $A^\perp$, of length at most $\ell_Y$, with $\ell_X\le\ell_Y$
and $\ell_X\ell_Y<\theta\kappa$.
\end{lemma}
\begin{proof}
As $\lambda(x)<\theta^n\le1$, there is $0\ne\xi\in L_V(x)$ with $|\xi|_x<\theta^n$. Take $y_j$,
$\alpha_j$ from Lemma~\ref{pgk:mink}. By Lemma~\ref{pgk:lat}(b) with $R=\alpha_n$, $|\xi|_x\alpha_n\ge1$, so
\[
\prod_{j<n}\alpha_j\le|\xi|_x\prod_{j\le n}\alpha_j<\theta^n\kappa^n .
\]
If $n=1$ this reads $1<\theta\kappa$, which is false; so $n\ge2$. Put $k=n-1$ and $c=\theta\kappa\le1$.
Then $\prod_{j=1}^k\alpha_j\alpha_{k+1-j}=(\prod_{j\le k}\alpha_j)^2<c^{2n}\le c^k$, so some factor has
$\alpha_j\alpha_{k+1-j}<c$; replacing $j$ by $k+1-j$ if needed we get $i$ with $i\le k+1-i=n-i$ (so
$\alpha_i\le\alpha_{n-i}$ by monotonicity) and $\alpha_i\alpha_{n-i}<c$. Put $\ell_X=\alpha_i$,
$\ell_Y=\alpha_{n-i}$, $A=\langle y_1,\dots,y_i\rangle$, $B=\langle y_1,\dots,y_{n-i}\rangle$. For
$a\le i$ and $b\le n-i$, $|y_a|_x|y_b|_x\le\ell_X\ell_Y<1$, so $\langle y_a,y_b\rangle=0$ by
Lemma~\ref{pgk:lat}(a). Thus $A\subseteq B\subseteq A^\perp$; $A$ is isotropic; and $\dim B=n-i=
\dim A^\perp$, so $B=A^\perp$.
\end{proof}

\subsubsection{The unipotent radical and its reduction}\label{pgk:ss-rad}

Let $A\subset V$ be a non-zero isotropic subspace. Its \emph{unipotent radical} $N_A$ is the algebraic
subgroup of $G$ of the $u$ that fix $A$ pointwise and induce the identity on $A^\perp/A$. As for $G$, we write $N_A(R)$ for the elements of $N_A$ read as matrices over $R=E,\auA_E,\dots$; thus $N_A(E)=N_A(E^+)\subseteq\Gamma_G$ and $N_A(\auA)=N_A(\auA_{E^+})$. In the
coordinates at a point $x$, let $\mathsf a$ (an $n\times i$ matrix) and $\mathsf y$ (an $n\times(n-i)$
matrix) be the coordinates of the columns $X=(y_1,\dots,y_i)$ and $Y=(y_1,\dots,y_{n-i})$ of
Lemma~\ref{pgk:frame}; so $\mathsf a^*\mathsf S\mathsf a=0$, $\mathsf a^*\mathsf S\mathsf y=0$. Let $R$ be
$E$, $\auA_E$ or one of its completions, $\mathsf S$ the form (with $\mathsf a$, $\mathsf y$) read
over $R$. For an $n\times i$ matrix $w$ over $R$ with $\mathsf a^*\mathsf Sw=0$ put
\[
P_w=(w\mathsf a^*-\mathsf aw^*)\mathsf S,\qquad\mathsf E(w)=1+2P_w+2P_w^2=1+2\bigl(w\mathsf a^*-\mathsf a
w^*-\mathsf a\,w^*\mathsf Sw\,\mathsf a^*\bigr)\mathsf S .
\]

\begin{lemma}\label{pgk:param}
\begin{enumerate}
\item[(a)] $P_w$ is skew-adjoint for $\mathsf S$, $P_w^2=-\mathsf a(w^*\mathsf Sw)\mathsf a^*\mathsf S$,
$P_w^3=0$, and $\mathsf E(w)=\exp(2P_w)$ lies in $N_A(R)$ \textup{(}read in coordinates\textup{)}.
\item[(b)] $\mathsf E(w)\mathsf E(w')=\mathsf E\bigl(w+w'-\mathsf a(w^*\mathsf Sw')\bigr)$.
\item[(c)] Every $u\in N_A(R)$ is $\mathsf E(w)$ for some $w=\mathsf yC$, $C$ an $(n-i)\times i$ matrix
over $R$.
\end{enumerate}
\end{lemma}
\begin{proof}
Note $w^*\mathsf S\mathsf a=(\mathsf a^*\mathsf Sw)^*=0$. (a) $\mathsf SP_w=\mathsf S(w\mathsf a^*-\mathsf a
w^*)\mathsf S$ is skew-hermitian. Expanding $P_w^2$, the terms containing $\mathsf a^*\mathsf S\mathsf a$,
$\mathsf a^*\mathsf Sw$ or $w^*\mathsf S\mathsf a$ vanish, leaving $-\mathsf aw^*\mathsf Sw\mathsf a^*\mathsf
S$; and $P_wP_w^2=0$ for the same reason. So $\exp(2P_w)=1+2P_w+2P_w^2$ is unitary for $\mathsf S$ (the
exponential of an $\mathsf S$-skew-adjoint nilpotent matrix). It fixes the columns of $\mathsf a$
($P_w\mathsf a=0$), and for $v$ with $\mathsf a^*\mathsf Sv=0$, $P_wv=-\mathsf a(w^*\mathsf Sv)\in A$ and
$P_w^2v\in\mathsf a\cdot R^i$; so it is the identity on $A^\perp/A$.
(b) For two-step nilpotent matrices, $\exp X\exp Y=\exp(X+Y+\frac12[X,Y])$. Here, expanding as in (a),
$P_wP_{w'}=-\mathsf a(w^*\mathsf Sw')\mathsf a^*\mathsf S$, so $[P_w,P_{w'}]=-\mathsf a(k-k^*)\mathsf a^*
\mathsf S$ with $k=w^*\mathsf Sw'$, and this commutes with $P_w,P_{w'}$. Also $P_{-\mathsf ak}=-\mathsf a
(k-k^*)\mathsf a^*\mathsf S$ (and $\mathsf a^*\mathsf S\mathsf ak=0$). Hence $2P_w+2P_{w'}+2[P_w,P_{w'}]=
2P_{w+w'-\mathsf ak}$.
(c) Let $u\in N_A(R)$ and $X=\log u=(u-1)-\frac12(u-1)^2$; since $(u-1)V\subseteq A^\perp$,
$(u-1)A^\perp\subseteq A$, $(u-1)A=0$, we have $(u-1)^3=0$, $u=\exp X$, and $X$ is $\mathsf S$-skew-adjoint
with the same three properties. Choose $\mathsf b$ ($n\times i$ over $R$) with $\mathsf a^*\mathsf S
\mathsf b=1_i$ (possible: $\mathsf a$ has independent columns and $\mathsf S$ is invertible); then
$R^n=A^\perp\oplus\mathsf bR^i$. Put $w_0=X\mathsf b$, so $\mathsf a^*\mathsf Sw_0=-(X\mathsf a)^*\mathsf
S\mathsf b=0$, and $c=w_0^*\mathsf S\mathsf b=-\mathsf b^*\mathsf SX\mathsf b$, which is skew-hermitian.
We claim $X=P_{w_0}+\mathsf ac\mathsf a^*\mathsf S$. On $\mathsf b$: the right side gives $w_0-\mathsf a
(w_0^*\mathsf S\mathsf b)+\mathsf ac=w_0=X\mathsf b$. On $v\in A^\perp$: write $Xv=\mathsf at$; then
$t=\mathsf b^*\mathsf S\mathsf at=\mathsf b^*\mathsf SXv$, and the right side gives $-\mathsf aw_0^*\mathsf
Sv=-\mathsf a\mathsf b^*X^*\mathsf Sv=\mathsf a\mathsf b^*\mathsf SXv=Xv$. Now $\mathsf ac\mathsf a^*
\mathsf S=P_{\mathsf ac/2}$ (as $c^*=-c$), so $X=P_w$ with $w=w_0+\mathsf ac/2$, and $\mathsf a^*\mathsf S
w=0$. Hence $u=\exp(2P_{w/2})=\mathsf E(w/2)$. The columns of $w/2$ lie in $A^\perp$, which is spanned
by the columns of $\mathsf y$, so $w/2=\mathsf yC$.
\end{proof}

\begin{proposition}[Reduction in the radical]\label{pgk:cover}
Let $M\ge1$ be an integer, $0<\delta\le1$, $\kappa\ge48n^6f^2N_d$, $x\in G(\auA)$, and let $A$,
$\ell_X$, $\ell_Y$ be as in Lemma~\ref{pgk:frame} with $\ell_X\ell_Y\le\delta/(M\kappa)$. For every
$u\in N_A(\auA)$ there is $\gamma\in N_A(E)$ such that $k=x^{-1}\gamma ux$ satisfies
$\operatorname{dist}(k_\infty,1)\le\delta$ and $k_f\in J_G(M)$.
\end{proposition}
\begin{proof}
Work in the coordinates at $x$: $\tilde u=x^{-1}ux\in N_A$ read in coordinates. By
Lemma~\ref{pgk:param}(c), $\tilde u=\mathsf E(M\mathsf yc)$ for an adelic $(n-i)\times i$ matrix $c$.
Let $\varrho(\cdot)$ be the entrywise rounding of Lemma~\ref{pgk:round}, put $d=-M\varrho(c)^*(\mathsf y^*
\mathsf S\mathsf y)c$ (an adelic $i\times i$ matrix) and $Z=-M(Y\varrho(c)-X\varrho(d))$, a matrix over
$E$ whose columns lie in $A^\perp$. Let $\gamma=\mathsf E_{\mathsf H}(Z)$ be the element defined by the
same formula in frame coordinates (with $X$, $\mathsf H$ in place of $\mathsf a$, $\mathsf S$); by
Lemma~\ref{pgk:param}(a) over $R=E$, $\gamma\in N_A(E)$. Since the coordinates at $x$ are obtained by
applying $x^{-1}$ (with $\mathsf S=x^*\mathsf Hx$, $\mathsf a=x^{-1}X$, $\mathsf y=x^{-1}Y$), $x^{-1}\gamma
x=\mathsf E(z)$ with $z=x^{-1}Z=-M(\mathsf y\varrho(c)-\mathsf a\varrho(d))$. By Lemma~\ref{pgk:param}(b),
$k=\mathsf E(z)\mathsf E(M\mathsf yc)=\mathsf E(z+M\mathsf yc-\mathsf az^*\mathsf SM\mathsf yc)$. Using
$\mathsf a^*\mathsf S\mathsf y=0$, $z^*\mathsf SM\mathsf yc=-M^2\varrho(c)^*\mathsf y^*\mathsf S\mathsf yc=
Md$, so
\[
k=\mathsf E(MT),\qquad T=\mathsf y\,(c-\varrho(c))-\mathsf a\,(d-\varrho(d)).
\]
\emph{Finite places.} The entries of $c-\varrho(c)$, $d-\varrho(d)$, $\mathsf a_f$, $\mathsf y_f$ (lattice
vectors) and $\mathsf S_f=\mathsf H$ are integral, so $k_f=1+M\mathsf Z$ with $\mathsf Z$ integral. As
$k_f$ preserves the unimodular form, $\det k_f$ is a unit, so $k_f$ lies in the stabiliser of
$\widehat L_V$ and is $\equiv1$ modulo $M$: $k_f\in J_G(M)$.
\emph{Real places.} A matrix with $p$ columns each of length at most $\ell$ has operator norm at most
$\sqrt p\,\ell$, and a $p\times q$ matrix with entries at most $\beta$ has operator norm at most
$\sqrt{pq}\,\beta$. So $\|\mathsf a_\tau\|\le\sqrt n\ell_X$, $\|\mathsf y_\tau\|\le\sqrt n\ell_Y$,
$\|(c-\varrho(c))_\tau\|,\|(d-\varrho(d))_\tau\|\le2nf^2N_d$, and $\|T_\tau\|\le4n^{3/2}f^2N_d\ell_Y$ (as
$\ell_X\le\ell_Y$). With $q=\|T_\tau\|\|\mathsf a_\tau\|\le4n^2f^2N_d\ell_X\ell_Y$ and $\|\mathsf S_\tau
\|=1$,
\[
\|k_\tau-1\|\le4Mq+2M^2q^2,\qquad Mq\le\frac{4n^2f^2N_d\,\delta}{\kappa}\le\frac\delta{12},
\]
so $\|k_\tau-1\|\le\frac\delta3+\frac{\delta^2}{72}\le\delta$.
\end{proof}

\subsubsection{The cusp estimate: input (C)}\label{pgk:ss-C}

\begin{theorem}[Input (C)]\label{pgk:C}
Let $\kappa_{\mathrm c}=8fn\,\mathrm{rd}_E+48n^6f^2N_d+1$. Let $F$ be a form of type $e$ on $G$ fixed by
$J_G(M)$, let $0<\delta\le1$, and let $x\in G(\auA)$ with $\lambda(x)<(\delta/(M\kappa_{\mathrm c}^2))^n$.
If $|F(xu)-F(x)|\le C$ for all $u\in G_\infty$ with $\operatorname{dist}(u,1)\le\delta$, then $|F(x)|\le
C$. This is \textup{(C)} of \S\ref{pgq:ss-inputs}.
\end{theorem}
\begin{proof}
\emph{The cusp.} Let $\kappa=\kappa_{\mathrm c}$ and $\theta=\delta/(M\kappa^2)$. Then $\kappa\ge1$,
$\kappa\ge8fn\,\mathrm{rd}_E\ge(fn\,\mathrm{rd}_E)^{1/2}$ ($\mathrm{rd}_E\ge1$), $\kappa\ge48n^6f^2N_d$,
and $\theta\kappa=\delta/(M\kappa)\le1$. By Lemma~\ref{pgk:frame} there is a non-zero isotropic $A$ with
$\ell_X\ell_Y<\theta\kappa=\delta/(M\kappa)$, and Proposition~\ref{pgk:cover} applies: every $u\in
N_A(\auA)$ has $\gamma_u\in N_A(E)$ with $x^{-1}\gamma_uux\in\mathcal B:=Z_\delta\times J_G(M)$, where
$Z_\delta=\{z\in G_\infty:\operatorname{dist}(z,1)\le\delta\}$.

\emph{Compactness.} Put $\mathcal K=N_A(\auA)\cap x\mathcal Bx^{-1}$, compact ($N_A(\auA)$ is closed and
$\mathcal B$ compact). Every $u\in N_A(\auA)$ lies in $N_A(E)\mathcal K$. So $N_A(E)\backslash N_A(\auA)$
is compact; give it the quotient of a Haar measure of $N_A(\auA)$, of finite mass $\mu>0$.

\emph{Oscillation.} For $u\in N_A(\auA)$ write $x^{-1}\gamma_uux=k=(k_\infty,k_f)$. As $F$ is left
invariant under $\gamma_u\in\Gamma_G$, right invariant under $J_G(M)$, and real and finite points commute,
$F(ux)=F(\gamma_uux)=F(xk_\infty k_f)=F(xk_\infty)$, so $|F(ux)-F(x)|\le C$.

\emph{The constant term vanishes.} The function $\mathsf F=F$ is bounded, continuous, left invariant
under $N_A(E)\subseteq\Gamma_G$ and satisfies $\mathsf F(g)=\int_{G_\infty}\mathsf F(gh)e(h)\,dh$; and the
real points of $N_A$ at $\tau$ form the unipotent radical of the stabiliser in $G_\tau$ of the non-zero
isotropic subspace $\tau(A)$, a proper parabolic subgroup. The argument of the first paragraph of the
proof of {\cite[Lemma~4.10]{P51}} (``the constant term vanishes'') therefore gives
$\int_{N_A(E)\backslash N_A(\auA)}F(ug)\,du=0$ for all $g$. That argument uses {\cite[Lemma~4.8]{P51}} for
$e_\tau=d_\tau\aubar{c_\tau}$; its hypotheses (the bound and smoothness of {\cite[Lemma~4.7(a)]{P51}},(b)
with $B=5n^3$) hold here by Lemma~\ref{pgq:local}(a) with $j=5n^3\le h_V$ (Lemma~\ref{pgq:sizes}(c)
gives $h_V\ge500(n+1)^3\ge5n^3$) and the proof of {\cite[Lemma~4.7(b)]{P51}}, which uses only that $\pi_\tau$ is
irreducible unitary and admissible ({\cite[Remark~4.3(1)]{P51}},(3)).

\emph{Conclusion.} The function $u\mapsto F(ux)$ is continuous and left $N_A(E)$-invariant, so
\[
\mu\,|F(x)|=\Bigl|\int_{N_A(E)\backslash N_A(\auA)}\bigl(F(x)-F(ux)\bigr)du\Bigr|\le\mu\,C .
\qedhere
\]
\end{proof}




% ===== C-grid12.tex =====
% proofs/C-grid12.tex
% Input (S) of proofs/C-grid9.tex (the restriction pairing), proved by Laplace's method.
% Labels start with pgl:. Conventions and macros: those of proofs/C-grid9.tex and B-period.tex.
\providecommand{\auU}{\mathrm U}
\providecommand{\aubar}[1]{\overline{#1}}
\providecommand{\pgqvol}{\operatorname{vol}}
\providecommand{\pdsK}{\mathbf K}
\providecommand{\pdsg}{\mathfrak g}
\providecommand{\pdsm}{m}
\providecommand{\pbptr}{\operatorname{tr}}
\providecommand{\pglF}{\mathsf F}
\providecommand{\pglLarge}{\mathrm{Large}}
\providecommand{\pglSmall}{\mathrm{Small}}

\subsection{The restriction pairing}\label{pgl:sec}

This section proves input (S) of \S\ref{pgq:ss-inputs} (Theorem~\ref{pgl:S}). The argument is the
Laplace method of {\cite[\S7]{P51}}, with two changes: the compact factor of {\cite[\S7]{P51}} is replaced
by a second noncompact group, and the stationary position has a single vector. Lemmas of
{\cite[\S\S5 and~7]{P51}} and~{\cite[\S3.2]{P49}} are reused in the general form (arbitrary rank and
parameter) stated and proved in \S\ref{pgm:sec}; what differs is written in full. \emph{Dependencies:} Theorem~\ref{pgl:S} is proved in full, modulo the
published results already used in {\cite[\S\S5 and~7]{P51}} and~{\cite[\S3.2]{P49}} (Schwarz's lemma,
the Casimir value {\cite[(60)]{P51}} through {\cite[Lemma~5.7]{P51}}).

\subsubsection{Statement and reductions}\label{pgl:ss-stat}

We use the notation of \S\ref{pgq:ss-stat} and \S\ref{pgq:ss-local}: $n=r$, $m=2\varrho+2$, $k$,
$B=\kappa(\tau)+\Delta/2$, $\beta_{\mathsf S}=b^{3m}$, $N_{\mathsf S}=20(2n+1)^2$, $V_\tau=\pgqvol\{r\in
G_\tau:\|r-1\|\le1\}$, and the coefficients $c_\tau$ (type $e_0$, Harish-Chandra parameter
$\varsigma_\tau$, {\cite[\S5.1]{P51}}) and $c'_\tau$ (type $e_{(1,0,0)}$, Harish-Chandra parameter
$-\lambda'$ with $\lambda'=\lambda_{(1,0,0),\tau}$).

\begin{theorem}[Input (S)]\label{pgl:S}
For every $\tau\in\auPl$ there is $\omega_\tau\in H_\tau$ with $\|\omega_\tau\|^2\le\beta_{\mathsf S}$ and
\[
V_\tau\le\beta_{\mathsf S}^{N_{\mathsf S}}\Bigl|\int_{G_\tau}c_\tau(r)\,\aubar{c'_\tau\bigl(\omega_\tau^{-1}
(r\oplus1)\omega_\tau\bigr)}\,dr\Bigr| .
\]
\end{theorem}

Fix $\tau$ and drop it. \emph{Matrices.} Let $p$ be the unitary matrix of {\cite[\S5.1]{P51}} with
$p^{-1}\mathsf J_\tau p=J:=\operatorname{diag}(\varepsilon_i)$ (signs on $V$), and put $p_Z=p\oplus1$,
$J_Z=J\oplus1$. Conjugation by $p$, $p_Z$ identifies $G_\tau$, $H_\tau$ with $G=\auU(J)$, $H=\auU(J_Z)$,
preserves operator norms, carries $r\oplus1$ to $r\oplus1$, and carries the Haar measure to a Haar
measure. We work in $G$, $H$ from now on; the ball $\{\|r-1\|\le1\}$ is preserved. By
Lemma~\ref{pgm:coords} (with {\cite[Lemma~5.1]{P51}} and~Lemma~\ref{pgq:steep}), the coordinates of $\C^n$ split as $X\sqcup\{1,\dots,m\}$
and those of $\C^{n+1}$ as $X\sqcup\{0,\dots,m\}$ (the same set $X$ of $k-1$ coordinates, of sign $+$,
the extra coordinate of $\C^{n+1}$ lying among the far ones), numbered so that: on $X$ both weights
vanish; for $V$, $\Lambda_1>\dots>\Lambda_m$ lie in $[-3B,-B]$ and the coordinate $j$ has sign
$(-1)^{j-1}$; for $Z$, $\Lambda'_0>\dots>\Lambda'_m$ lie in $[-3B,-B]$ and the coordinate $i$ has sign
$(-1)^i$. Here $\Lambda=\Lambda(\varsigma)$ and $\Lambda'=\Lambda(\lambda')$ are the highest weights of the
lowest $K$-types ({\cite[\S5.1]{P51}}). All weights are integers, and $m$ is even.

\emph{The functions $\pglF_\lambda$.} For a unitary group $\auU(J_0)$ of rank $N$ and a regular
$\lambda\in(\frac{N-1}2+\Z)^N$ let $\pglF_\lambda$ be the function $F$ of {\cite[(14)]{P49}}, with its
$\Lambda=\Lambda(\lambda)$, $\mu=\lambda+\rho$. By {\cite[Theorem~3.17]{P49}} (gap $\ge N^2+1$, which holds here
by Lemma~\ref{pgq:sizes}(c)), $c=\pglF_\varsigma$ on $G$ and $c'=\pglF_{-\lambda'}$ on $H$.

\begin{lemma}[Negation]\label{pgl:neg}
$\Lambda(-\lambda)=-\Lambda(\lambda)$, $\mu(-\lambda)=-\mu(\lambda)$, $-\lambda$ has the same gap as
$\lambda$, and $\pglF_{-\lambda}=\aubar{\pglF_\lambda}$.
\end{lemma}
\begin{proof}
Replacing $\lambda$ by $-\lambda$ exchanges the sets $\{j:\lambda_j<\lambda_i\}$ and $\{j:\lambda_j>
\lambda_i\}$, so $\rho$, $\rho_{\mathrm c}$, $\rho_{\mathrm n}$ change sign; this gives the first two
claims, and the third is clear. For unitary $u$, {\cite[Lemma~3.18(a)]{P49}} gives $\pdsm_{-w}(u)=
\pdsm_w(u^{-1})=\pdsm_w(u^*)=\aubar{\pdsm_w(u)}$. Hence $|\pdsm_{-\Lambda}|=|\pdsm_\Lambda|$ on $K$,
$d_{-\Lambda}=d_\Lambda$, and the numerator of {\cite[(14)]{P49}} for $-\lambda$ is
$\aubar{\pdsm_\Lambda(ku)}\pdsm_\Lambda(u)$, the conjugate of that for $\lambda$. The denominators are
$\pdsm_\mu(w)$ and $\pdsm_{-\mu}(w)$ for positive definite $w=u^{-1}zu\in\auU(J_0)$; they are equal and
positive by {\cite[Lemma~5.5]{P51}} (its statement $\pdsm_{-\mu}(w)=\pdsm_\mu(w)>0$ does not use the
integer $c$).
\end{proof}

So $\aubar{c'}=\pglF_{\lambda'}$ on $H$ and $c=\aubar{\pglF_{-\varsigma}}$ on $G$, and with
\[
\mu_Z:=\operatorname{diag}(\Lambda')\ (\text{weight of }\lambda'),\qquad\mu_V:=\operatorname{diag}(-\Lambda)\
(\text{weight of }-\varsigma),
\]
the integrand of Theorem~\ref{pgl:S} is $\pglF_{\lambda'}(\omega^{-1}(r\oplus1)\omega)\,
\aubar{\pglF_{-\varsigma}(r)}$. The weight $\mu_Z$ is \emph{steep} in the sense of
{\cite[(63)]{P51}} (zero on $X$, strictly decreasing in $[-3B,-B]$ with alternating signs elsewhere),
and $\mu_V$ is $0$ on $X$ and lies in $[B,3B]$ elsewhere.

\subsubsection{The stationary position}\label{pgl:ss-pos}

Write $\mu_i=\Lambda'_i$ ($0\le i\le m$) and $\nu_j=-\Lambda_j$ ($1\le j\le m$): integers with
$\mu_0>\dots>\mu_m$, $\nu_1<\dots<\nu_m$, $|\mu_i|\le3B$, $B\le\nu_j\le3B$, so $1\le2B\le\nu_j-\mu_i\le
6B$. Let $p(t)=\prod_i(t-\mu_i)$ (degree $m+1$), $q(t)=\prod_j(t-\nu_j)$ (degree $m$) and $q_j=q/(t-\nu_j)$.
For a real polynomial $F$ let $[F]\in\R^{n+1}$ vanish on $X$ and have the entry
$[F]_i=F(\mu_i)\,|p'(\mu_i)q(\mu_i)|^{-1/2}$ on the coordinate $i$; then $\mu_Z[F]=[tF]$. For real
vectors vanishing on $X$ put $\mathcal B(x,y)=\sum_i(-1)^ix_iy_i$, the form $x^*J_Zy$.

\begin{lemma}\label{pgl:gram}
\textup{(i)} $\mathcal B([F_1],[F_2])=\sum_iF_1F_2(\mu_i)/(p'(\mu_i)q(\mu_i))$.
\textup{(ii)} $\mathcal B([q],[q])=1$, $\mathcal B([q_j],[q])=0$, and $\mathcal B([q_j],[q_l])=
(-1)^{j-1}\delta_{jl}\,|q'(\nu_j)|/p(\nu_j)$, with $p(\nu_j)>0$.
\textup{(iii)} $|[q]_i|\le(6B)^{m/2}$, and with $s_j=(p(\nu_j)/|q'(\nu_j)|)^{1/2}$, $|s_j[q_j]_i|\le
(6B)^m$.
\end{lemma}
\begin{proof}
(i) Exactly $i$ of the $\mu$'s exceed $\mu_i$, so $p'(\mu_i)$ has sign $(-1)^i$; $q(\mu_i)=
\prod_j(\mu_i-\nu_j)$ has sign $(-1)^m=+1$. So $|p'(\mu_i)q(\mu_i)|=(-1)^ip'(\mu_i)q(\mu_i)$.
(ii) By (i) and Lemma~\ref{pgm:residue} (with $s=m+1$ nodes): $\mathcal B([q],[q])=\sum_iq(\mu_i)/
p'(\mu_i)$ is the coefficient of $t^m$ in $q$, which is $1$; $\mathcal B([q_j],[q])=\sum_iq_j(\mu_i)/
p'(\mu_i)$ is the coefficient of $t^m$ in $q_j$ (degree $m-1$), which is $0$; and $\mathcal B([q_j],
[q_l])=\sum_iq_j(\mu_i)/((\mu_i-\nu_l)p'(\mu_i))=-q_j(\nu_l)/p(\nu_l)$, as $\deg q_j\le m$ and $\nu_l$ is
not a node. This vanishes for $j\ne l$, and for $j=l$ it is $-q'(\nu_j)/p(\nu_j)$. Now $p(\nu_j)>0$ as
$\nu_j>\mu_i$ for all $i$, and $q'(\nu_j)=\prod_{l\ne j}(\nu_j-\nu_l)$ has sign $(-1)^{m-j}=(-1)^j$.
(iii) The nodes are distinct integers, so $|p'(\mu_i)|\ge1$, $|q'(\nu_j)|\ge1$, and $|q_j(\mu_i)|\le
|q(\mu_i)|$ as $|\mu_i-\nu_j|\ge1$. Hence $[q]_i^2\le|q(\mu_i)|\le(6B)^m$, and $(s_j[q_j]_i)^2\le
p(\nu_j)\,q_j(\mu_i)^2/|q(\mu_i)|\le p(\nu_j)|q_j(\mu_i)|\le(6B)^{m+1}(6B)^{m-1}$.
\end{proof}

\begin{proposition}[Stationary position]\label{pgl:stat}
Let $\iota$ be the $(n+1)\times n$ matrix that is the identity from the coordinates $X$ of $\C^n$ to
those of $\C^{n+1}$ and sends the far coordinate vector $e_j$ of $\C^n$ to $w_j:=s_j[q_j]$, and let
$v=[q]$. Then $w:=(\iota\mid v)\in H$, $\|w\|^2\le\beta_{\mathsf S}$, $\iota^\dagger\mu_Z\iota=\mu_V$
where $\iota^\dagger=J\iota^*J_Z$, and $\iota^\dagger\iota=1$, $\iota^\dagger v=0$.
\end{proposition}
\begin{proof}
By Lemma~\ref{pgl:gram}(ii) and disjointness of supports, the Gram matrix $w^*J_Zw$ is
$J\oplus1=J_Z$ (the coordinate $j$ of $V$ has sign $(-1)^{j-1}$); so $w\in H$, and $w^{-1}=J_Zw^*J_Z$
has the rows $\iota^\dagger$ and $v^*J_Z$, which gives $\iota^\dagger\iota=1$, $\iota^\dagger v=0$.
Since $t\,q_j=\nu_jq_j+q$, $\mu_Zw_j=\nu_jw_j+s_jv$, and $\mu_Z$ vanishes on the $X$-coordinates; so
$\iota^\dagger\mu_Z\iota e_j=\nu_je_j$ and $\iota^\dagger\mu_Z\iota e_l=0$ ($l\in X$), i.e.\
$\iota^\dagger\mu_Z\iota=\mu_V$. All entries of $w$ have modulus at most $(6B)^m$
(Lemma~\ref{pgl:gram}(iii)), so $\|w\|^2\le\|w\|_{\mathrm F}^2\le(n+1)^2(6B)^{2m}\le\beta_{\mathsf S}$ by
Lemma~\ref{pgq:sizes}(c).
\end{proof}

We put $\omega=w^{-1}$ (so $\|\omega\|=\|w\|$ and $\omega^{-1}(r\oplus1)\omega=w(r\oplus1)w^{-1}$) and,
for $r\in G$, $g_r=w(r\oplus1)w^{-1}\in H$.

\begin{lemma}[The group of the position]\label{pgl:group}
For $r\in G$: \textup{(a)} $g_r-1=w((r-1)\oplus0)w^{-1}$, $g_r\iota=\iota r$, $g_rv=v$, $r=\iota^\dagger
g_r\iota$; \textup{(b)} $\pbptr(g_r-1)=\pbptr(r-1)$ and $\pbptr(\mu_Z(g_r-1))=\pbptr(\mu_V(r-1))$;
\textup{(c)} if $\|r-1\|\le\frac12$, then $\|g_r-1\|\le\frac12\beta_{\mathsf S}$ and, when also
$\|g_r-1\|\le\frac12$, $\log g_r=w((\log r)\oplus0)w^{-1}$ and $\pbptr(\mu_Z\log g_r)=\pbptr(\mu_V\log r)$;
\textup{(d)} if $u$ is unitary on the coordinates $X$ and $r_u=u\oplus1$, then $g_{r_u}=u\oplus1$;
\textup{(e)} $r\mapsto g_r$ is a homeomorphism of $G$ onto a closed subgroup of $H$, and
$\|g_r\|\le\|w\|^2\|r\|$.
\end{lemma}
\begin{proof}
(a) Immediate from $w(r\oplus1)=(\iota r\mid v)$ and $\iota^\dagger\iota=1$. (b) $\pbptr(D\,w(Y\oplus0)
w^{-1})=\pbptr((Y\oplus0)w^{-1}Dw)$, and the upper left $n\times n$ block of $w^{-1}Dw$ is
$\iota^\dagger D\iota$; take $D=1$, resp.\ $D=\mu_Z$, and use Proposition~\ref{pgl:stat}. (c) The
first bound is $\|w\|\|w^{-1}\|\|r-1\|$. The logarithm series commutes with conjugation by $w$ and
with $Y\mapsto Y\oplus0$ termwise (both series converge, $\|(r-1)^j\oplus0\|=\|(r-1)^j\|$), so
$\log g_r=w((\log r)\oplus0)w^{-1}$ (as $(1+Y)\oplus1$, conjugated, is $g_r$); then argue as in (b).
(d) $\iota$ and $v$ are the identity on $X$ and vanish there off it: $w=1_X\oplus w'$. (e) $g_r$
is continuous in $r$, $r=\iota^\dagger g_r\iota$, and the image is the closed subgroup of $h\in H$
fixing $v$ (an $h$ fixing $v$ preserves $\iota(\C^n)=v^\perp$, and $w^{-1}hw=r\oplus1$ with
$r=\iota^\dagger h\iota\in G$).
\end{proof}

\emph{The integrand.} By \S\ref{pgl:ss-stat}, the integrand of Theorem~\ref{pgl:S} is
\begin{equation}\label{pgl:eq-theta}
\theta(r)=\pglF_{\lambda'}(g_r)\,\aubar{\pglF_{-\varsigma}(r)},\qquad|\theta(r)|=|c'(g_r)|\,|c(r)| .
\end{equation}
The phase difference of {\cite[\S7.3]{P51}} is $\varphi(r)=\pbptr(\mu_Z(g_r-1))-\pbptr(\mu_V(r-1))=0$
by Lemma~\ref{pgl:group}(b): the position is exactly stationary and there is no character.

\subsubsection{Sizes and the coefficient lemmas of \texorpdfstring{{\cite[\S7]{P51}}}{[Moment, \S7]}}\label{pgl:ss-sizes}

Put $D=2n+1$, $\beta=\beta_{\mathsf S}$, $T=D^2\log\beta$, $h_0=\lceil10^{30}T^3\rceil$,
$\rho_0=(200\beta^5)^{-1}$.

\begin{lemma}\label{pgl:sizes}
\textup{(a)} $\beta\ge10^{12}$, $(n+1)^2\le T$, $18D^2\le h_0\le\min(h_V,h_Z)$.
\textup{(b)} For $\pglF=\pglF_{\lambda'}$ on $H$ and $\pglF=\pglF_{-\varsigma}$ on $G$ \textup{(}rank
$N\le n+1$, weight $\Lambda_\pglF$\textup{)}: $\max_i|\Lambda_{\pglF,i}|\le6b$ and $d_{\Lambda_\pglF}\le
e^{2T}$. \textup{(c)} $\rho_0\le1/(4C)$ for $C=3(n+1)(2b+2n+2)$.
\end{lemma}
\begin{proof}
(a) Lemma~\ref{pgq:sizes}(c) gives $\beta\ge10^{12}$ and $\min(h_V,h_Z)\ge500(n+1)^3+10^{30}T^3\ge h_0$;
$\log\beta\ge1$ gives $(n+1)^2\le T$, and $18D^2\le h_0$. (b) $|\Lambda_i|\le3B\le6b$
(Lemma~\ref{pgq:steep}, $B\le2b$ by {\cite[Lemma~13.4]{P46}}). Lemma~\ref{pgm:dLambda} gives
$d_\Lambda\le4(24M_\Lambda N)^{2N^2}$ with $M_\Lambda\le6b$: so $d_\Lambda\le4(144b(n+1))^{2(n+1)^2}\le
b^{4(n+1)^2+1}$ (as $144(n+1)\le b$), and $(4(n+1)^2+1)\log b\le2D^2\cdot3m\log b=2T$.
(c) $4C\le48b^2\le200b^{15m}$.
\end{proof}

\begin{lemma}[Transfer]\label{pgl:transfer}
For each of the two functions $\pglF$ of Lemma~\ref{pgl:sizes}(b), with $\pdsK$, $J_0$, the weight
$\Lambda_\pglF$ and the parameter $\lambda_\pglF$ of its group, the following hold.
\begin{enumerate}
\item[(a)] \textup{(Decay)} $|\pglF(g)|\le(2/\|g\|)^{c}$ for integers $0\le c\le h_0$, and
$|\pglF(g)|^4\le(1+\|y\|_{\mathrm F}^2)^{-h_0}$, $y=g-a(g)$ the noncompact part.
\item[(b)] \textup{(Gaussian shape)} If $\|g-1\|_{\mathrm F}\le\delta_1$ with $\delta_1T\le5\cdot10^{-13}$
and $Z=\log g$, there is $\varrho\in[0,1]$ with $|\pglF(g)e^{-\pbptr(\Lambda_\pglF Z)}-\varrho|\le\varrho/10
+e^{-100T}$.
\item[(c)] \textup{(Left $X$-invariance)} If $k_0\in\pdsK$ is the identity off $X$, $\pglF(k_0g)=\pglF(g)$.
\item[(d)] \textup{(Near $1$)} If $\|g-1\|\le\rho_0$, $|\pglF(g)|\ge\frac12$.
\end{enumerate}
\end{lemma}
\begin{proof}
Both functions are $\pglF_\lambda$ for a regular $\lambda\in(\frac{N-1}2+\Z)^N$ ($N=n$, $\lambda=-\varsigma$;
$N=n+1$, $\lambda=\lambda'$), so the general lemmas of \S\ref{pgm:sec} apply (\S\ref{pgm:ss-effect}).
(a) Lemma~\ref{pgm:decay} with $c\le h_0\le\min(h_V,h_Z)$ (Lemma~\ref{pgl:sizes}(a); Lemma~\ref{pgl:neg}
for the gap of $-\varsigma$); the second bound uses $|\pglF|\le|\det a(g)|^{-h_0/2}$, the intermediate
bound of the proof of {\cite[Lemma~3.25]{P49}}.
(b) Proposition~\ref{pgm:gauss}, whose hypotheses $T\ge N^2$ and $d_{\Lambda_\pglF}\le e^{2T}$ are
Lemma~\ref{pgl:sizes}(a),(b).
(c) Lemma~\ref{pgm:leftX}: $X$ has sign $+$ and $\Lambda'=0$, $-\Lambda=0$ there (Lemma~\ref{pgm:coords}).
(d) $|\pglF|=|c'|$, resp.\ $|c|$ ({\cite[Theorem~3.17]{P49}}, gaps $\ge N^2+1$). Lemma~\ref{pgm:lip} with
$C=3N(L+2N)$: for $c$, $N=n$, $L=L_\tau\le2b$; for $c'$, $N=n+1$ and $L=\max_i|\lambda'_i|\le2b$ (the entries
are $\frac k2-\frac12-i$ ($0\le i\le k-2$) or $-\kappa(\tau)-j\Delta+\frac12$ ($1\le j\le m+1$), Lemma~\ref{pgm:coords}; the first
kind has $|\frac k2-\frac12-i|\le\frac k2\le b$, and the second has
$|\lambda'_i|\le\kappa(\tau)+(m+1)\Delta\le b+(2m+1)\cdot2r^8\le b+x^{42}\le2b$ by $\kappa(\tau)\le b$
({\cite[Lemma~13.4(S5)]{P46}}), $\Delta=2r^8$, $m\le2\sqrt x$, $r\le3x^5$ and $x^{42}\le b$ (S1)). In
both cases $C\le3(n+1)(2b+2n+2)$, and $\rho_0\le1/(4C)$ by Lemma~\ref{pgl:sizes}(c).
\end{proof}

\subsubsection{Rows, localisation and the cone}\label{pgl:ss-cone}

Norms are Frobenius norms in this subsection. For $r\in G$ let $x=g_r-1$, $x'=r-1$, $y$, $y'$ the
noncompact parts of $g_r$, $r$, and $P_{\mathrm{off}}$ the coordinate projection of $\C^{n+1}$ onto the
coordinates outside $X$.

\begin{lemma}[Rows]\label{pgl:rows}
$\|x\|^2=2\|y\|^2-2\Rel\pbptr x$, $\|x'\|^2=2\|y'\|^2-2\Rel\pbptr x'$, and
$\|P_{\mathrm{off}}x\|^2\le6(\|y\|^2+\|y'\|^2)$, $\|xP_{\mathrm{off}}\|^2\le6(\|y\|^2+\|y'\|^2)$.
\end{lemma}
\begin{proof}
The row identity $2\Rel x_{ii}=2\sum_j|y_{ij}|^2-\sum_j|x_{ij}|^2$ of the proof of
{\cite[Lemma~7.25]{P51}} holds in any $\auU(J_0)$; summing gives the two identities. Multiply the row
identity of $g_r$ by $-\mu_{Z,i}\ge0$ and sum: $\sum_i(-\mu_{Z,i})\sum_j|x_{ij}|^2=2\sum_i(-\mu_{Z,i})
\sum_j|y_{ij}|^2+2\Rel\pbptr(\mu_Zx)$; the left side is $\ge B\|P_{\mathrm{off}}x\|^2$ and the first
term on the right $\le6B\|y\|^2$. Multiply that of $r$ by $\mu_{V,i}\in[0,3B]$:
$2\Rel\pbptr(\mu_Vx')\le6B\|y'\|^2$. As $\pbptr(\mu_Zx)=\pbptr(\mu_Vx')$ (Lemma~\ref{pgl:group}(b)),
$B\|P_{\mathrm{off}}x\|^2\le6B(\|y\|^2+\|y'\|^2)$. For the columns apply this to $r^{-1}$:
$g_{r^{-1}}-1=J_Zx^*J_Z$, and noncompact parts of inverses have the same norms.
\end{proof}

\begin{lemma}[Localisation]\label{pgl:loc}
Let $0<\rho_1\le\frac16$ and $\|y\|^2+\|y'\|^2\le\rho_1^2$. There is $t\in G$ with $\theta(t)=
\theta(r)$, $\|g_t-1\|\le4\rho_1$ and $\|t-1\|\le4\rho_1$.
\end{lemma}
\begin{proof}
As in {\cite[Lemma~7.28]{P51}}: the block $\alpha$ of $g_r$ on $X\times X$ satisfies $\|\alpha^*\alpha-1\|
\le\|P_{\mathrm{off}}x\|^2\le6\rho_1^2\le\rho_1$ (Lemma~\ref{pgl:rows}), so $\alpha=k_1z_1$, $k_1$
unitary, $\|z_1-1\|\le\rho_1$. Let $u=k_1^{-1}$ and $t=r_ur$ (Lemma~\ref{pgl:group}(d)); then
$g_t=(u\oplus1)g_r$ has the block $z_1$ on $X$, $\pglF_{\lambda'}(g_t)=\pglF_{\lambda'}(g_r)$ and
$\pglF_{-\varsigma}(t)=\pglF_{-\varsigma}(r)$ by Lemma~\ref{pgl:transfer}(c), so $\theta(t)=\theta(r)$;
left multiplication by the block diagonal unitaries $u\oplus1$ keeps $\|y\|$, $\|y'\|$. Counting rows
off $X$, columns off $X$ and the block on $X$, $\|g_t-1\|^2\le13\rho_1^2$. By the two identities of
Lemma~\ref{pgl:rows} and $\pbptr(g_t-1)=\pbptr(t-1)$, $\|t-1\|^2=\|g_t-1\|^2-2\|y_t\|^2+2\|y'_t\|^2\le
15\rho_1^2$.
\end{proof}

\begin{proposition}[Cone]\label{pgl:cone}
For every $r\in G$, $|\theta(r)|\le2\Rel\theta(r)+12e^{-100T}$.
\end{proposition}
\begin{proof}
Let $E=e^{-100T}$, $\rho_1=10^{-13}/T$. If $\|y\|^2+\|y'\|^2\ge\rho_1^2$, Lemma~\ref{pgl:transfer}(a)
gives $|\theta|^4\le(1+\|y\|^2)^{-h_0}(1+\|y'\|^2)^{-h_0}\le(1+\rho_1^2)^{-h_0}\le E^4$ (as in
{\cite[Proposition~7.29]{P51}}), so $|\theta|\le2\Rel\theta+12E$. Otherwise take $t$ by
Lemma~\ref{pgl:loc}, $\delta_1=4\rho_1$, $Z=\log g_t$, $Z'=\log t$. By Lemma~\ref{pgl:group}(c),
$\pbptr(\mu_ZZ)=\pbptr(\mu_VZ')$, and this is imaginary ($Z'\in\pdsg$ has imaginary diagonal). Hence
$\theta(t)=u_1u_2$ with $u_1=\pglF_{\lambda'}(g_t)e^{-\pbptr(\mu_ZZ)}$, $u_2=\aubar{\pglF_{-\varsigma}(t)
e^{-\pbptr(\mu_VZ')}}$. By Lemma~\ref{pgl:transfer}(b) there are $\varrho,\varrho'\in[0,1]$ with
$|u_1-\varrho|\le\varrho/10+E$, $|u_2-\varrho'|\le\varrho'/10+E$, and Lemma~\ref{pgm:conearith} gives the
claim.
\end{proof}

\subsubsection{Proof of Theorem~\ref{pgl:S}}\label{pgl:ss-proof}

Operator norms again. Let $\pglLarge=\{r:\|g_r\|\le\beta+1,\ \|r\|\le\beta+1\}$ and
$\pglSmall=\{r:\|g_r-1\|\le\rho_0,\ \|r-1\|\le\rho_0\}$; they are compact (Lemma~\ref{pgl:group}(e)).

\begin{lemma}\label{pgl:cover}
$\pgqvol(\pglLarge)\le\beta^{16D^2}\pgqvol(\pglSmall)$, and $\{\|r-1\|\le1\}\subseteq\pglLarge$.
\end{lemma}
\begin{proof}
Lemma~\ref{pgm:cover} with $\Gamma=G$, $\mathcal J(r)=g_r\oplus r\in\mathrm{GL}_D(\C)$, $T'=\beta+1$, $\rho=\rho_0$
(Remark~\ref{pgm:rem-apply}): $\|\mathcal J(r)\|=\max(\|g_r\|,\|r\|)=\|\mathcal J(r)^{-1}\|$, $r\in\pglSmall$
iff $\|\mathcal J(r)-1\|\le\rho_0$, and $(8(\beta+1)^2/\rho_0)^{2D^2}\le\beta^{16D^2}$ (as $1/\rho_0\le200\beta^5$ and $\beta\ge10^{12}$). If $\|r-1\|\le1$ then $\|r\|\le2$ and $\|g_r\|\le1+\|w\|^2\|r-1\|\le1+\beta$
(Lemma~\ref{pgl:group}(a), Proposition~\ref{pgl:stat}).
\end{proof}

\begin{proof}[Proof of Theorem~\ref{pgl:S}]
Take $\omega$ as above (and $p_Z\omega p_Z^{-1}$ in $H_\tau$); $\|\omega\|^2\le\beta$. Let
$\mathcal G(r)=(2/\|r\|)^{9D^2}$ and $E=\beta^{-50D^2}$, so $E^2=e^{-100T}$. \emph{Decay.} By
Lemma~\ref{pgl:transfer}(a) with $c=9D^2,18D^2\le h_0$, $|\theta|\le|c(r)|\le\mathcal G$ and
$|\theta|\le\mathcal G^2$; by Lemma~\ref{pgq:balls} ($9D^2\ge8n^2$), $\int_G\mathcal G\le2^{9D^2}\cdot
2\cdot16^{n^2}V_\tau\le16^{9D^2}V_\tau$, so $\theta$ is integrable. \emph{Cone.} $\Omega:=|\theta|-
2\Rel\theta\le12E^2$ (Proposition~\ref{pgl:cone}) and $\Omega\le3\mathcal G^2$, so $\Omega\le6E\mathcal G$.
\emph{Small.} On $\pglSmall$, $|\theta|\ge\frac14$ by Lemma~\ref{pgl:transfer}(d). \emph{Integration.}
With $I=|\int_G\theta|$, $V=\pgqvol(\pglLarge)\ge V_\tau$, $W=\pgqvol(\pglSmall)$:
$\frac14W\le\int|\theta|\le2I+6E\cdot16^{9D^2}V_\tau\le2I+6E\cdot16^{9D^2}V$. By Lemma~\ref{pgl:cover},
$V\le8\beta^{16D^2}I+24\cdot16^{9D^2}\beta^{-34D^2}V$, and $24\cdot16^{9D^2}\beta^{-34D^2}\le\frac12$; as
$V<\infty$, $V_\tau\le V\le16\beta^{16D^2}I\le\beta^{20D^2}I=\beta_{\mathsf S}^{N_{\mathsf S}}I$.
\end{proof}




% ===== C-grid13.tex =====
% proofs/C-grid13.tex
% The coefficient lemmas of B-period.tex / B-discrete-series.tex in arbitrary rank, as used by
% proofs/C-grid12.tex (pgl:), and the two coordinate facts taken from pgq:steep.
% Labels start with pgm:. Conventions and macros: those of proofs/C-grid12.tex and B-period.tex.
\providecommand{\auU}{\mathrm U}
\providecommand{\aubar}[1]{\overline{#1}}
\providecommand{\pdsK}{\mathbf K}
\providecommand{\pdsg}{\mathfrak g}
\providecommand{\pdsm}{m}
\providecommand{\pbptr}{\operatorname{tr}}
\providecommand{\pbpF}[1]{\|#1\|_{\mathrm F}}
\providecommand{\pbpQ}{\mathcal Q}
\providecommand{\pglF}{\mathsf F}
\providecommand{\pgmvol}{\operatorname{vol}}
\providecommand{\pdsk}{\mathfrak k}
\providecommand{\pdsp}{\mathfrak p}
\providecommand{\pcoH}{\mathsf H}

\subsection{The coefficient lemmas in arbitrary rank}\label{pgm:sec}

Section~\ref{pgl:sec} applies several lemmas of {\cite[\S\S5 and~7]{P51}} to functions other
than the coefficient $c$ for which they are stated: to $\pglF_{-\varsigma}$ on $G$ (rank $n$) and to
$\pglF_{\lambda'}$ on $H$ (rank $n+1$). This section states each of these lemmas for an arbitrary
unitary group $\auU(J_0)$ of rank $N$ and an arbitrary parameter, with the hypotheses that are
actually needed and explicit constants, and proves it by going through its proof and
recording every place where the rank, the parameter or the group enters. It also proves the two
facts on coordinates that \S\ref{pgl:ss-stat} takes from Lemma~\ref{pgq:steep}, and re-checks the
numerical inputs of Theorem~\ref{pgl:S}. Two inputs of \S\ref{pgl:sec} are used in the form proved here: the bound needed from
{\cite[Lemma~3.25]{P49}} is the intermediate bound of its proof (Lemma~\ref{pgm:decay}), and the constant of
{\cite[Lemma~7.21]{P51}} is the one for the rank at hand (Lemma~\ref{pgm:dLambda}); with them
Theorem~\ref{pgl:S} holds as stated (\S\ref{pgm:ss-effect}).

\subsubsection{Setting}\label{pgm:ss-set}

Throughout, $J_0=\operatorname{diag}(\varepsilon_1,\dots,\varepsilon_N)$ with $\varepsilon_i=\pm1$,
$G_0=\auU(J_0)$, $\pdsK_0=G_0\cap\auU(N)$ (the block diagonal elements), $\pdsg_0$ the Lie algebra,
and $\lambda\in(\frac{N-1}2+\Z)^N$ has distinct entries (\emph{regular}). From $\lambda$ we form, as in
{\cite[\S3.2.1]{P49}}, $\rho,\rho_{\mathrm c},\rho_{\mathrm n}$, $\Lambda=\Lambda(\lambda)=\lambda+
\rho_{\mathrm n}-\rho_{\mathrm c}\in\Z^N$, $\mu=\lambda+\rho=\Lambda+2\rho_{\mathrm c}\in\Z^N$, the gap
$h=\min\{|\lambda_i-\lambda_j|:\varepsilon_i\ne\varepsilon_j\}$ (put $h=\infty$ if all signs agree),
$L=\max_i|\lambda_i|$, $M_\Lambda=\max(1,\max_i|\Lambda_i|)$, $d_\Lambda=(\int_{\pdsK_0}
|\pdsm_\Lambda|^2)^{-1}$, and $\pglF_\lambda$, the function $F$ of {\cite[(14)]{P49}} on $G_0$. Norms:
$\|\cdot\|$ operator norm, $\pbpF\cdot$ Frobenius norm, $a(x)=\frac12(x+J_0xJ_0)$, and the noncompact
part of $g\in G_0$ is $y=g-a(g)$. We record the elementary bounds
\begin{equation}\label{pgm:eq-rho}
|\rho_i|\le\tfrac{N-1}2,\quad|2\rho_{\mathrm c,i}|\le N-1,\quad|\rho_{\mathrm n,i}|+|\rho_{\mathrm c,i}|
\le\tfrac{N-1}2,\quad|\Lambda_i|\le L+N,
\end{equation}
each half sum running over at most $N-1$ indices $j\ne i$.

\begin{lemma}[What is stated in general]\label{pgm:general}
The following statements are stated, and proved, for an arbitrary $G_0$ and an arbitrary regular
$\lambda\in(\frac{N-1}2+\Z)^N$, with no condition on the gap: {\cite[Lemma~3.18]{P49}},
{\cite[Corollary~3.19]{P49}}, {\cite[Lemma~3.20]{P49}}, {\cite[Proposition~3.21]{P49}}, {\cite[Lemmas~3.22 and~3.23]{P49}} \textup{(}in particular {\cite[(17)]{P49}}\textup{)}, {\cite[Lemma~5.4]{P51}}, the
first assertion $\pdsm_{-\mu}(w)=\pdsm_\mu(w)>0$ of {\cite[Lemma~5.5]{P51}}, and, for integers
$0\le c\le h$, its second assertion. {\cite[Lemmas~7.16 and~7.18]{P51}},
{\cite[Propositions~7.17 and~7.19]{P51}} are stated in arbitrary dimension, with
constants independent of it.
\end{lemma}
\begin{proof}
For {\cite[\S3.2]{P49}} this is the convention of {\cite[\S3.2.1]{P49}} ($G=\auU(p,q)$ of arbitrary rank $n$,
arbitrary regular $\lambda$); the gap enters {\cite[\S3.2]{P49}} only in {\cite[Lemma~3.25]{P49}} and
{\cite[Theorem~3.17]{P49}}. {\cite[Lemma~5.4]{P51}} and the first assertion of {\cite[Lemma~5.5]{P51}}
use only $g^*J_0g=J_0$ (facts (F2), (F3) of that proof use $w\in G_0$ positive definite and nothing
else); the second assertion uses the gap only through ``$|\lambda_i-\lambda_j|\ge h\ge c$ if
$\varepsilon_i\ne\varepsilon_j$'' and the sign of $\rho_i-\rho_j$, valid for any regular $\lambda$.
The proofs in {\cite[\S7.6]{P51}} begin with ``the dimension is arbitrary; no constant depends on
it'', and indeed use only norm inequalities valid in every dimension.
\end{proof}

\subsubsection{Decay and the bound far from \texorpdfstring{$\pdsK_0$}{K}}\label{pgm:ss-decay}

\begin{lemma}[Decay, general form of {\cite[Lemma~3.25]{P49}} and {\cite[Lemma~7.27]{P51}}]\label{pgm:decay}
For $g\in G_0$ and every integer $0\le c\le h$,
\[
|\pglF_\lambda(g)|\le|\det a(g)|^{-c/2}\le(2/\|g\|)^c,\qquad|\pglF_\lambda(g)|^4\le(1+\pbpF y^2)^{-c},
\]
where $y$ is the noncompact part of $g$.
\end{lemma}
\begin{proof}
This is the proof of {\cite[Theorem~5.6]{P51}}, with $\pglF_\lambda$ in place of $c_\tau$: write $g=kz$,
$D=\det a(z)=|\det a(g)|$; for $u\in\pdsK_0$, $w_u=u^{-1}zu$ is positive definite in $G_0$ with $\det
a(w_u)=D$ ({\cite[Lemma~5.4]{P51}}), so $\pdsm_{-\mu}(w_u)=\pdsm_\mu(w_u)\ge D^{c/2}$ by
{\cite[Lemma~5.5]{P51}} (valid by Lemma~\ref{pgm:general}). Then $|\pglF_\lambda(g)|\le d_\Lambda
D^{-c/2}\int\frac12(|\pdsm_\Lambda(ku)|^2+|\pdsm_\Lambda(u)|^2)du=D^{-c/2}$, by invariance of $du$ and the
definition of $d_\Lambda$. The second inequality is {\cite[Lemma~5.4]{P51}}. For the last, as in
{\cite[Lemma~7.27]{P51}}: the block diagonal part of $g^*J_0g=J_0$ is $a^*a-y^*y=1$ ($J_0$ commutes with
$a=a(g)$ and anticommutes with $y$), so $|\det a|^2=\det(1+y^*y)\ge1+\pbptr(y^*y)=1+\pbpF y^2$, and
$|\pglF_\lambda|^4\le|\det a|^{-2c}\le(1+\pbpF y^2)^{-c}$. Neither $N$ nor $\lambda$ enters beyond
Lemma~\ref{pgm:general}.
\end{proof}

{\cite[Lemma~7.27]{P51}} and Lemma~\ref{pgl:transfer}(a) use the bound $|\det a(g)|^{-c/2}$ of
Lemma~\ref{pgm:decay}.

\subsubsection{The Gaussian shape}\label{pgm:ss-gauss}

\begin{lemma}[General form of {\cite[Lemma~7.20]{P51}}]\label{pgm:reduced}
Let $Z\in\pdsg_0$ with $\pbpF Z\le\frac1{20}$; let $u_-$ \textup(resp.\ $u_+$\textup) keep the entries
$Z_{ij}$ with $\varepsilon_i=\varepsilon_j$ and $\lambda_i<\lambda_j$ \textup(resp.\ $\lambda_i>
\lambda_j$\textup) and vanish elsewhere, and $\hat g=e^{-u_-}e^Ze^{-u_+}$. Then
\[
\pglF_\lambda(e^Z)=d_\Lambda\int_{\pdsK_0}\frac{|\pdsm_\Lambda(k)|^2}{\pdsm_{-\mu}(k^{-1}\hat gk)}\,dk .
\]
\end{lemma}
\begin{proof}
We go through the proof of {\cite[Lemma~7.20]{P51}}. (1) $x=e^Z$ has $\|x-1\|\le e^{1/20}-1<1$, so
$x+x^*=2+(x-1)+(x-1)^*$ is positive definite and $1\in A_x$; with the polar factor $k_x\in A_x$,
{\cite[Proposition~3.21]{P49}} gives $\pglF_\lambda(x)=d_\Lambda I(x;k_x)=d_\Lambda I(x;1)=d_\Lambda
\mathcal J(1,1)$. The proof invokes {\cite[Theorem~3.17]{P49}} at this point only to replace $c$ by
$F$; for $\pglF_\lambda$ this step is the definition, and no gap is needed. (2) Step (iv) of the proof
of {\cite[Proposition~3.21]{P49}} gives both $\partial_1\mathcal J=0$ and $\partial_2\mathcal J=0$ on $U_x$,
for any $\pcoH$ (rank-free). (3) $\|u_\pm\|\le\pbpF Z$, so $e^{tu_+}x^{-1}e^{tu_-}$ is within
$e^{3/20}-1<1$ of $1$ for $t\in[0,1]$, hence accretive; $e^{tu_\pm}$ are block diagonal (only entries
with $\varepsilon_i=\varepsilon_j$), so the path lies in $U_x$ and $\mathcal J(1,1)=\mathcal J(e^{u_-},
e^{u_+})$. (4) $b=e^{u_-}$ has unit diagonal and $b_{cd}=0$ unless $c=d$ or ($\varepsilon_c=
\varepsilon_d$, $\lambda_c<\lambda_d$); for such pairs $\Lambda_c\le\Lambda_d$ by
{\cite[Lemma~3.18(c)]{P49}}, which holds for every regular $\lambda$. {\cite[Lemma~3.18(b)]{P49}} (left
version; $b$ is triangular for the order of decreasing $\lambda$, which is a total order since
$\lambda$ is regular) gives $\pdsm_\Lambda(bk)=\pdsm_\Lambda(k)$; likewise $\pdsm_\Lambda(k^*e^{u_+})=
\aubar{\pdsm_\Lambda(k)}$. (5) $(k^*e^{u_+}x^{-1}e^{u_-}k)^{-1}=k^{-1}\hat gk$ and
{\cite[Lemma~3.18(a)]{P49}}. The rank enters nowhere; the parameter enters only through its
regularity (total order, {\cite[Lemma~3.18(c)]{P49}}).
\end{proof}

\begin{lemma}[General form of {\cite[Lemma~7.21]{P51}}]\label{pgm:dLambda}
$d_\Lambda\le4\,(24M_\Lambda N)^{2N^2}$.
\end{lemma}
\begin{proof}
We go through the proof of {\cite[Lemma~7.21]{P51}}, with $r$ replaced by $N$ and $3b$ by $M_\Lambda$.
Let $\varepsilon_0=1/(6M_\Lambda N)\le\frac16$. By {\cite[(57)]{P51}}, $\pdsm_\Lambda(k)=(\det k)^{-M_\Lambda}
\prod_{t=1-M_\Lambda}^{M_\Lambda}\Delta_{S_t(\Lambda)}(k)$ (factors with $t\le-M_\Lambda$ are $\det k\cdot
(\det k)^{-1}$, those with $t>M_\Lambda$ are $1$): $2M_\Lambda$ minors. If $k\in\pdsK_0$ and
$\|k-1\|\le\varepsilon_0$, then $|\det k|=1$, and $\Delta_S(k)=\det(1+(k-1)_{SS})$ is a product of at
most $N$ factors $1+\xi$ with $|\xi|\le\|(k-1)_{SS}\|\le\varepsilon_0\le\frac12$ (eigenvalues), so
$|\Delta_S(k)|\ge e^{-2N\varepsilon_0}$ and $|\pdsm_\Lambda(k)|\ge e^{-4M_\Lambda N\varepsilon_0}=
e^{-2/3}\ge\frac12$. A maximal family $k_1,\dots,k_\nu\in\pdsK_0$ with $\|k_ik_j^{-1}-1\|>
\varepsilon_0$ ($i\ne j$) has $\|k_i-k_j\|>\varepsilon_0$; the balls of radius $\varepsilon_0/2$ about the
$k_i$ in $M_N(\C)$ (real dimension $2N^2$) are disjoint and lie in the ball of radius $2$, so
$\nu\le(4/\varepsilon_0)^{2N^2}$. By maximality $\pdsK_0$ is covered by the $\nu$ right translates
$\{\|k-1\|\le\varepsilon_0\}k_i$, each of $dk$-mass equal to that of $\{\|k-1\|\le\varepsilon_0\}$; so
this ball has mass $\ge1/\nu$, and $d_\Lambda^{-1}\ge\frac14\nu^{-1}$, i.e.\ $d_\Lambda\le4(4/
\varepsilon_0)^{2N^2}$. The rank enters through the number of factors of a minor ($N$), the
dimension $2N^2$, and $\varepsilon_0$; the parameter only through $M_\Lambda$.
\end{proof}

\begin{proposition}[General form of {\cite[Proposition~7.22]{P51}}]\label{pgm:gauss}
Let $T\ge N^2$ with $d_\Lambda\le e^{2T}$. Let $g\in G_0$ with $\pbpF{g-1}\le\delta_1$, $\delta_1T\le
5\cdot10^{-13}$, and $Z=\log g$. There is $\varrho\in[0,1]$ with
\[
\bigl|\pglF_\lambda(g)\,e^{-\pbptr(\Lambda Z)}-\varrho\bigr|\le\varrho/10+e^{-100T}.
\]
Here $\Lambda$ is regarded as the diagonal matrix $\operatorname{diag}(\Lambda)$; no condition on the
gap is needed.
\end{proposition}
\begin{proof}
We go through the proof of {\cite[Proposition~7.22]{P51}}, naming each place where $N$ or $\lambda$
enters. (1) $\delta=2\delta_1$; as $\pbpF{g-1}\le\frac12$, $\pbpF Z\le-\log(1-\pbpF{g-1})\le\delta$
(Frobenius is submultiplicative), $Z\in\pdsg_0$ by the argument of Step~2 of {\cite[Lemma~5.7]{P51}}
(rank-free), and $\delta T\le10^{-12}$; as $T\ge N^2\ge1$, $\delta\le10^{-12}$. (2)
Lemma~\ref{pgm:reduced}. (3) With $w=-\mu$, $M'=\max|w_i|$: $\pdsm_w(y)=(\det y)^{-M'}\prod_{t=1-M'}^{M'}
\Delta_{S_t(w)}(y)$. For $i\in S_t(w)$, $j\notin S_t(w)$, $\mu_i<\mu_j$, hence $\lambda_i<\lambda_j$
({\cite[Lemma~3.18(c)]{P49}}); so {\cite[Proposition~7.19]{P51}} (arbitrary dimension) applies to each
factor with $o=\lambda$ (injective: regularity). Counting $t$ as stated,
$\pdsm_w(k^{-1}\hat gk)=\exp(-\pbptr(\mu Z)+\frac14\pbpQ_w(Z)+r(k))$ with $|r(k)|\le10^3\delta
(\pbpQ_w(Z)+\pbpQ_w(k))$, where $\pbpQ_w(Y)=\sum_{i,j}|w_i-w_j||Y_{ij}|^2$. (4) \emph{Rank enters.} By
\eqref{pgm:eq-rho}, $|\pbptr(2\rho_{\mathrm c}Z)|\le(N-1)\sum_i|Z_{ii}|\le(N-1)\sqrt N\pbpF Z\le N^2\delta$;
$|w_i-w_j|\le|\Lambda_i-\Lambda_j|+2(N-1)$, and $\sum_{ij}|k_{ij}|^2=N$, so $\pbpQ_w(k)\le\pbpQ_\Lambda(k)+
2N^2$. With $q=\pbpQ_w(Z)/4$: $\pdsm_w(k^{-1}\hat gk)=\exp(q-\pbptr(\Lambda Z)-r'(k))$, $|r'(k)|\le N^2
\delta+10^3\delta(4q+\pbpQ_\Lambda(k)+2N^2)\le10^4\delta\,s(k)$, $s(k)=q+\pbpQ_\Lambda(k)+N^2$. (5) As
$d_\Lambda\int|\pdsm_\Lambda|^2=1$, $\pglF_\lambda(g)e^{-\pbptr(\Lambda Z)}-e^{-q}=d_\Lambda\int
|\pdsm_\Lambda(k)|^2e^{-q}(e^{r'(k)}-1)dk$. (6) {\cite[Lemma~7.18]{P51}} (any dimension), multiplied
over $t$, gives $|\pdsm_\Lambda(k)|^2\le e^{-\pbpQ_\Lambda(k)/2}$. (7) If $s(k)\le5\cdot10^4T$, then
$|r'|\le5\cdot10^8\delta T\le5\cdot10^{-4}$ and $|e^{r'}-1|\le\frac1{10}$. Otherwise $|r'|\le s/4$ (as
$10^4\delta\le\frac14$) and $|\pdsm_\Lambda|^2e^{-q}|e^{r'}-1|\le 2e^{-\pbpQ_\Lambda/2-q+s/4}\le
2e^{-s/4+N^2/2}\le 2e^{-12500T+T/2}\le e^{-5100T}$ (using $|e^{r'}-1|\le e^{|r'|}+1\le2e^{s/4}$); \emph{this is the place where $T\ge N^2$ is
used}. (8) Integrating, the difference is at most $e^{-q}/10+d_\Lambda e^{-5100T}\le e^{-q}/10+
e^{-5098T}$ by the hypothesis $d_\Lambda\le e^{2T}$; take $\varrho=e^{-q}$. The constants are those of
the proof of {\cite[Proposition~7.22]{P51}}, which uses $r^2\le T$ and $d_\Lambda\le e^{2T}$, and these are replaced by the hypotheses.
\end{proof}

\subsubsection{Left invariance and the lower bound near \texorpdfstring{$1$}{1}}\label{pgm:ss-near}

\begin{lemma}[General form of {\cite[Lemma~7.26]{P51}}]\label{pgm:leftX}
Let $X$ be a set of coordinates on which $\varepsilon\equiv+1$ and $\Lambda=0$. If $k_0\in\pdsK_0$ is the
identity on the coordinates outside $X$, then $\pglF_\lambda(k_0g)=\pglF_\lambda(g)$ for $g\in G_0$.
\end{lemma}
\begin{proof}
The proof uses {\cite[(17)]{P49}} (general, Lemma~\ref{pgm:general}), $\Lambda=0$ on $X$, and
that $k_0$ is block diagonal, which holds as $X$ lies in one sign class. For $t\ge1$, $S_t(\Lambda)\cap
X=\emptyset$ and the rows of $k_0Y$ in $S$ are those of $Y$; for $t\le0$, $S_t(\Lambda)\supseteq X$, and
$(k_0)_{ij}=0$ for $i\in S$, $j\notin S$ (as $j\notin X$), so $\Delta_S(k_0Y)=\det k_0\,\Delta_S(Y)$, and
this factor cancels against $(\det k_0Y)^{-1}$ in {\cite[(57)]{P51}}. Rank and parameter do not enter.
\end{proof}

\begin{lemma}[General form of {\cite[Lemma~5.7]{P51}}]\label{pgm:lip}
Let $\pi$ be a discrete series of $G_0$ with Harish-Chandra parameter $\lambda$, $v_0$ a unit vector of
its lowest $\pdsK_0$-type, of highest weight $\Lambda$, and $c_0=\langle v_0,\pi(\cdot)v_0\rangle$. Put
$C=3N(L+2N)$. If $\|g-1\|\le1/(4C)$, then $|c_0(g)|\ge\frac12$. If $h\ge N^2+1$, the same holds for
$\pglF_\lambda$.
\end{lemma}
\begin{proof}
Step~1 of the proof of {\cite[Lemma~5.7]{P51}} uses {\cite[(60)]{P51}} for $\pi$ (published, valid for
any discrete series of any $\auU(p,q)$: the Casimir value depends only on the infinitesimal character
$\lambda$, and the $\pdsk$-Casimir value only on the $\pdsK_0$-type $\Lambda$, a product of two unitary
groups), and the bounds
$|\Lambda_i+\rho_{\mathrm c,i}|=|\lambda_i+\rho_{\mathrm n,i}|\le L+N$, $|\rho|^2\le N^3/4$,
$|\Lambda_a|\le L+N$ (\eqref{pgm:eq-rho}): they give $\sum_{\varepsilon_a\ne\varepsilon_b}
\|d\pi(E_{ab})v_0\|^2\le N(L+2N)^2+N^3/4\le2N(L+2N)^2$, hence $\|d\pi(Y)v_0\|\le\sqrt2N(L+2N)\|Y\|$
on $\pdsp_0$ (Cauchy--Schwarz gives $\|d\pi(Y)v_0\|\le\pbpF Y\,(2N)^{1/2}(L+2N)$, and $\pbpF Y\le
\sqrt N\|Y\|$), and $\|d\pi(X)v_0\|\le N(L+N)\|X\|$ on $\pdsk_0$. So
$\|d\pi(X)v_0\|\le C\|X\|$ on $\pdsg_0$. Step~2 is rank-free: $\|\log g\|\le2\|g-1\|$, $\log g\in\pdsg_0$,
and $|c_0(g)-1|\le\int_0^1\|d\pi(\log g)v_0\|dt\le2C\|g-1\|\le\frac12$. The rank enters through $C$
only, the parameter through $L$ only. If $h\ge N^2+1$, {\cite[Theorem~3.17]{P49}} gives $\pglF_\lambda=
c_0$ for the discrete series with parameter $\lambda$ (whose lowest $\pdsK_0$-type is $\Lambda$, by the
formula for the lowest $K$-type quoted in {\cite[\S5.1]{P51}}).
\end{proof}

\subsubsection{The lemmas for two noncompact groups}\label{pgm:ss-two}

Lemmas~\ref{pgl:rows}, \ref{pgl:loc} and Proposition~\ref{pgl:cone} are proved in \S\ref{pgl:ss-cone}
with complete arguments except for two references, ``the last lines of the proof of
{\cite[Proposition~7.29]{P51}}'' and ``the proof of {\cite[Lemma~7.31]{P51}}'', and Lemma~\ref{pgl:gram}
refers to {\cite[Lemma~7.10]{P51}} with $2n$ nodes. We state and prove the three general facts behind
these references; they involve neither the rank nor the parameter.

\begin{lemma}[Residue sums, any number of nodes]\label{pgm:residue}
Let $\mu_0,\dots,\mu_{s-1}$ be distinct reals, $p(t)=\prod_i(t-\mu_i)$, $F$ a polynomial of degree
$\le s-1$ and $c$ not a node. Then $\sum_iF(\mu_i)/p'(\mu_i)$ is the coefficient of $t^{s-1}$ in $F$, and
$\sum_iF(\mu_i)/((\mu_i-c)p'(\mu_i))=-F(c)/p(c)$.
\end{lemma}
\begin{proof}
The proof of {\cite[Lemma~7.10]{P51}} with $2n$ replaced by $s$: Lagrange interpolation at the $s$
nodes, and $H(t)=(F(t)p(c)-F(c)p(t))/(t-c)$, of degree $\le s-1$ with coefficient $-F(c)$ of $t^{s-1}$.
\end{proof}

\begin{lemma}[The cone inequality]\label{pgm:conearith}
Let $0\le E\le1$, $\varrho,\varrho'\in[0,1]$ and $u_1,u_2\in\C$ with $|u_1-\varrho|\le\varrho/10+E$,
$|u_2-\varrho'|\le\varrho'/10+E$. Then $z=u_1u_2$ satisfies $|z|\le2\Rel z+12E$.
\end{lemma}
\begin{proof}
Put $P=\varrho\varrho'\ge0$. Then $|z-P|\le\varrho|u_2-\varrho'|+\varrho'|u_1-\varrho|+|u_1-\varrho|
|u_2-\varrho'|\le(0.2P+2E)+(0.01P+0.2E+E^2)\le0.21P+3.2E$, using $\varrho,\varrho',E\le1$. Hence
$|z|-2\Rel z\le(P+|z-P|)-2(P-|z-P|)=-P+3|z-P|\le-0.37P+9.6E\le12E$.
\end{proof}

\begin{lemma}[Covering]\label{pgm:cover}
Let $\Gamma$ be a unimodular locally compact group with a Haar measure, $D\ge1$, and $\mathcal J\colon
\Gamma\to\mathrm{GL}_D(\C)$ a continuous injective homomorphism with $\|\mathcal J(x)^{-1}\|=\|\mathcal
J(x)\|$, such that $\{x:\|\mathcal J(x)\|\le T'\}$ is compact for every $T'$. For $T'\ge1$ and
$0<\rho\le1$,
\[
\pgmvol\{x:\|\mathcal J(x)\|\le T'\}\le(8T'^2/\rho)^{2D^2}\,\pgmvol\{x:\|\mathcal J(x)-1\|\le\rho\}.
\]
\end{lemma}
\begin{proof}
Let $x_1,\dots,x_\nu$ lie in $L=\{\|\mathcal J\|\le T'\}$ with $\|\mathcal J(x_ix_j^{-1})-1\|>\rho$ for
$i\ne j$. As $\mathcal J(x_i)-\mathcal J(x_j)=(\mathcal J(x_ix_j^{-1})-1)\mathcal J(x_j)$ and
$\|\mathcal J(x_j)^{-1}\|\le T'$, $\|\mathcal J(x_i)-\mathcal J(x_j)\|>\rho/T'$; the balls of radius
$\rho/(2T')$ about the $\mathcal J(x_i)$ in $M_D(\C)\cong\R^{2D^2}$ are disjoint and lie in the ball of
radius $2T'$, so $\nu\le(4T'^2/\rho)^{2D^2}$. Take such a family maximal and put $S:=\{\|\mathcal
J-1\|\le\rho\}$. If $x\in L$ is not among the $x_i$, it cannot be added, so for some $i$ either
$\|\mathcal J(xx_i^{-1})-1\|\le\rho$ or $\|\mathcal J(x_ix^{-1})-1\|\le\rho$, i.e.\ $xx_i^{-1}\in S$ or
$xx_i^{-1}\in S^{-1}$ (the separation condition need not be symmetric in $i,j$). Hence $L\subseteq
\bigcup_i(S\cup S^{-1})x_i$, and by right invariance and unimodularity ($\pgmvol(S^{-1})=\pgmvol(S)$)
$\pgmvol(L)\le2\nu\pgmvol(S)\le2(4T'^2/\rho)^{2D^2}\pgmvol(S)\le(8T'^2/\rho)^{2D^2}\pgmvol(S)$.
\end{proof}

\begin{remark}[Application in \S\ref{pgl:sec}]\label{pgm:rem-apply}
Lemma~\ref{pgl:cover} is Lemma~\ref{pgm:cover} with $\Gamma=G$, $\mathcal J(r)=g_r\oplus r$, $D=2n+1$,
$T'=\beta+1$, $\rho=\rho_0$: $\mathcal J$ is a continuous injective homomorphism
(Lemma~\ref{pgl:group}(e)), $\|\mathcal J(r)^{\pm1}\|=\max(\|g_r\|,\|r\|)$ as $\|h^{-1}\|=\|h\|$ in $G$
and $H$, $\{\|\mathcal J\|\le T'\}$ is closed in the compact set $\{\|r\|\le T'\}$, and
$(8(\beta+1)^2\cdot200\beta^5)^{2D^2}\le\beta^{16D^2}$ as $\beta\ge10^{12}$. Proposition~\ref{pgl:cone}
ends with Lemma~\ref{pgm:conearith} applied to $u_1,u_2$ (with $E=e^{-100T}$). The proofs of
Lemmas~\ref{pgl:rows} and~\ref{pgl:loc} are complete as stated:
in Lemma~\ref{pgl:rows} the weights used are $-\mu_{Z,i}\in\{0\}\cup[B,3B]$ and $\mu_{V,i}\in\{0\}\cup
[B,3B]$, with zero exactly on $X$ (Lemma~\ref{pgm:coords}); in Lemma~\ref{pgl:loc} the element $r_u=
u\oplus1$ lies in $G$ because $X$ has sign $+$, $g_{r_u}=u\oplus1$ needs that $X$ is the same set in
$\C^n$ and $\C^{n+1}$ (Lemma~\ref{pgm:coords}), and the invariance $\theta(r_ur)=\theta(r)$ is
Lemma~\ref{pgm:leftX} on $H$ and on $G$.
\end{remark}

\subsubsection{The coordinates of \texorpdfstring{$V$}{V} and \texorpdfstring{$Z$}{Z}}\label{pgm:ss-coords}

Fix $\tau$. By \S\ref{pgq:ss-stat}, $V_\tau$ carries the normalised basis of $Y_0$, which after
conjugation by $p$ ({\cite[\S5.1]{P51}}) gives the coordinates $\varrho=0,\dots,r-1$ of $\C^n$ ($n=r$),
attached to the stem coordinates of the frame of $Y_0$; and $Z_\tau$ carries that basis followed by
$\tilde e$, so that after conjugation by $p_Z=p\oplus1$ the coordinates of $\C^{n+1}$ are the same
$\varrho=0,\dots,r-1$ followed by $\varrho=r$, the coordinate of $\tilde e$
(Definition~\ref{pgv:def-spaces}: the frame of $Z=\pgrY^0_0$ is that of $Y_0$ followed by $\tilde e$ as
the stem coordinate $\varrho=r$). The inclusion $r\mapsto r\oplus1$ is the identity on the
coordinates $\varrho<r$. Put $a(\varrho)=\max(\varrho+1-k,0)$.

\begin{lemma}[The two coordinate facts]\label{pgm:coords}
Let $X_V=\{\varrho<r:\Lambda_\varrho=0\}$ and $X_Z=\{\varrho\le r:\Lambda'_\varrho=0\}$, where
$\Lambda=\Lambda(\varsigma_\tau)$ \textup(signs of $V$\textup) and $\Lambda'=\Lambda(\lambda')$,
$\lambda'=\lambda_{(1,0,0),\tau}$ \textup(signs $\epsilon_{(1,0,0)}$\textup). Then:
\begin{enumerate}
\item[(a)] $X_V=X_Z=\{0,1,\dots,k-2\}$, and on these coordinates the signs of $V$ and of $Z$ are $+$;
\item[(b)] at every coordinate $\varrho<r$ the signs of $V$ and $Z$ agree, namely $(-1)^{a(\varrho)}$,
and $\lambda'_\varrho=\varsigma_{\tau,\varrho}+\frac12$;
\item[(c)] the extra coordinate $\varrho=r$ of $Z$ has sign $+$, is not in $X_Z$, and carries the
smallest entry $\lambda'_r=-\kappa(\tau)-(m+1)\Delta+\frac12$ of $\lambda'$ and the smallest value of
$\Lambda'$; in the numbering $0,\dots,m$ of the far coordinates of $Z$ by decreasing $\Lambda'$ used in
\S\ref{pgl:ss-stat}, it is the coordinate $m$. The far coordinates are $\varrho=k-1+i$ \textup($0\le
i\le m$\textup) for $Z$, of sign $(-1)^i$, and $\varrho=k-2+j$ \textup($1\le j\le m$\textup) for $V$, of
sign $(-1)^{j-1}$.
\end{enumerate}
\end{lemma}
\begin{proof}
\emph{Signs.} By \eqref{pgv:eq-sign} for $p=(z,0,0)$, $\epsilon_p(\varrho)=(-1)^{a(\varrho)+\min(a(\varrho),0)}
=(-1)^{a(\varrho)}$ for every stem coordinate, for $z=0$ (the space $Y_0=Y^0_0$, whose signs are
those of {\cite[Lemma~5.1]{P51}}: $(-1)^{\max(0,\varrho+1-k)}$) and for $z=1$ alike. So the signs agree
at $\varrho<r$; they are $+$ for $\varrho\le k-1$; and at $\varrho=r$, $a(r)=r+1-k=m$ (as $r+1=k+m$) is
even, so the sign is $+$.

\emph{Parameters.} By \eqref{pgv:eq-param} with $z=1$, $a=b=0$, $s_p=0$, $\lambda'_\varrho=
\varsigma'_{\tau,\varrho}+\frac12$, and $\varsigma'_{\tau,\varrho}$ is the entry $\varsigma_{\tau,\varrho}$ of
{\cite[Lemma~5.1]{P51}} ($\frac k2-1-\varrho$ for $\varrho\le k-2$, $-\kappa(\tau)-(\varrho+2-k)\Delta$ for
$\varrho\ge k-1$); this gives (b), and $\lambda'_r=-\kappa(\tau)-(m+1)\Delta+\frac12$, the smallest
entry since the entries decrease along $\varrho$ (by $1$ on $\varrho\le k-2$, then from $\lambda'_{k-2}=\frac32-\frac k2$ to
$\lambda'_{k-1}=-\kappa(\tau)-\Delta+\frac12<\frac32-\frac k2$ as $\kappa(\tau)+\Delta>k$, then by $\Delta$).

\emph{Zero sets.} By {\cite[Lemma~5.1]{P51}}, $\Lambda_\varrho=0$ for $\varrho\le k-2$, and by the remark
after it (and Lemma~\ref{pgq:steep}) $\Lambda_\varrho\in[-3B,-B]$, in particular $\ne0$ as $B\ge1$, for
$\varrho\ge k-1$; so $X_V=\{0,\dots,k-2\}$. The computation in the proof of Lemma~\ref{pgq:steep}
(which uses only the signs and parameters just determined) gives $\Lambda'_\varrho=0$ for $\varrho\le k-2$
and $\Lambda'_\varrho\in[-3B,-B]$ for $k-1\le\varrho\le r$; so $X_Z=\{0,\dots,k-2\}=X_V$ as sets of
coordinates of $\C^{n+1}\supset\C^n$, and $r\notin X_Z$. This is (a).

\emph{Order.} On $\varrho\ge k-1$, $\lambda'$ decreases strictly with $\varrho$, in steps $\Delta$, and
$|\Lambda'_\varrho-\lambda'_\varrho|\le|\rho_{\mathrm n,\varrho}|+|\rho_{\mathrm c,\varrho}|\le n/2$
\eqref{pgm:eq-rho}; as $\Delta=2r^8>n+1$, $\Lambda'$ also decreases strictly
with $\varrho$ there, so the order by decreasing $\Lambda'$ is the order of $\varrho$, $\varrho=k-1+i$,
and $\varrho=r$ is $i=m$; its sign is $(-1)^{a(\varrho)}=(-1)^{i}$. The same argument for $V$ (steps
$\Delta$, $|\Lambda_\varrho-\varsigma_{\tau,\varrho}|\le(n-1)/2$) gives $\varrho=k-2+j$, sign
$(-1)^{\varrho+1-k}=(-1)^{j-1}$. This is (c).
\end{proof}

So the two facts that \S\ref{pgl:ss-stat} attributes to Lemma~\ref{pgq:steep} (``$X$ is the same set
of coordinates in $V$ and $Z$'', ``the extra coordinate of $Z$ lies among the far ones'') follow
from the definitions (``on the first $k-1$ coordinates'' in Lemma~\ref{pgq:steep}, the frame being
that of Definition~\ref{pgv:def-spaces}).

\subsubsection{Effect on Theorem~\ref{pgl:S}}\label{pgm:ss-effect}

We list the parameters with which the lemmas above are applied in \S\ref{pgl:sec} and check every
hypothesis. Notation of \S\ref{pgl:ss-sizes}: $D=2n+1$, $\beta=b^{3m}$, $T=D^2\log\beta$,
$h_0=\lceil10^{30}T^3\rceil$, $\rho_0=(200\beta^5)^{-1}$.

\begin{enumerate}
\item[(1)] \emph{The two functions.} $\pglF_{-\varsigma}$ on $G=\auU(J)$: $N=n$, parameter $-\varsigma$,
weight $-\Lambda$; $\pglF_{\lambda'}$ on $H=\auU(J_Z)$: $N=n+1$, parameter $\lambda'$, weight $\Lambda'$.
Both parameters are regular and lie in $(\frac{N-1}2+\Z)^N$: the entries of $\varsigma$ are integers and
$n=r$ is odd; those of $\lambda'$ are in $\frac12+\Z$ ($k$ even) and $(N-1)/2=r/2\in\frac12+\Z$.
Lemma~\ref{pgl:neg} gives the weights and the gaps.
\item[(2)] \emph{Lemma~\ref{pgm:decay}} with $c\le h_0$ requires $h_0\le\min(h_V,h_Z)$:
Lemma~\ref{pgq:sizes}(c) gives $\min(h_V,h_Z)\ge500(n+1)^3+10^{30}T^3\ge h_0$.
\item[(3)] \emph{Proposition~\ref{pgm:gauss}} requires $T\ge N^2$ and $d_\Lambda\le e^{2T}$. First,
$N^2\le(n+1)^2\le D^2\le T$ as $\log\beta\ge1$. Second, $M_\Lambda\le3B\le6b$ (Lemma~\ref{pgm:coords},
Lemma~\ref{pgq:steep}; $B=\kappa(\tau)+\Delta/2\le b+r^8\le2b$ by {\cite[Lemma~13.4(S5)]{P46}}, and
$r^8\le b$ as $r\le3x^5$, $b\ge x^{10^{12}}$), so Lemma~\ref{pgm:dLambda} gives $d_\Lambda\le4(144b
(n+1))^{2(n+1)^2}\le b^{4(n+1)^2+1}$, using $144(n+1)\le b$ ((S7)(b): $8\cdot\frac m2\cdot r\le b$ with
$\frac m2\ge100$); and $(4(n+1)^2+1)\log b\le2D^2\cdot3m\log b=2T$. This is the constant of
Lemma~\ref{pgl:sizes}(b), now justified by Lemma~\ref{pgm:dLambda} instead of {\cite[Lemma~7.21]{P51}}.
\item[(4)] \emph{Lemma~\ref{pgm:leftX}} with $X$ of Lemma~\ref{pgm:coords}: sign $+$, and $\Lambda'=0$,
$-\Lambda=0$ there.
\item[(5)] \emph{Lemma~\ref{pgm:lip}} for $c$ on $G$ ($N=n$, $L=L_\tau\le2b$) and for $c'$ on $H$
($N=n+1$, parameter $-\lambda'$, $L'=\max_i|\lambda'_i|\le\kappa(\tau)+(m+1)\Delta\le b+x^{42}\le2b$,
using $m\le2\sqrt x$, $\Delta=2r^8\le2\cdot3^8x^{40}$). The gaps exceed $N^2+1$ (item (2)), so
$|\pglF|=|c|$, $|c'|$ and the lemma applies to both $\pglF$. Since $C\le3(n+1)(2b+2n+2)$ and
$4C\le48b^2\le200\beta^5$, $\rho_0\le1/(4C)$.
\item[(6)] \emph{Stationary position.} Lemma~\ref{pgl:gram} uses Lemma~\ref{pgm:residue} with $s=m+1$;
Proposition~\ref{pgl:stat} uses $\|w\|^2\le(n+1)^2(6B)^{2m}\le\beta$, which follows from
Lemma~\ref{pgq:sizes}(c) ($\beta\ge(n+1)^4(6B)^{2m+2}$).
\item[(7)] \emph{Final count.} The proof of Theorem~\ref{pgl:S} uses $9D^2,18D^2\le h_0$, $9D^2\ge8n^2$
(Lemma~\ref{pgq:balls}), $\beta\ge10^{12}$ (Lemma~\ref{pgq:sizes}(c)), Lemma~\ref{pgm:cover} with the
constant $\beta^{16D^2}$, and $16\beta^{16D^2}\le\beta^{20D^2}$, all of which hold here.
\end{enumerate}

Hence every lemma used by \S\ref{pgl:sec} holds with the constants used there, and
Theorem~\ref{pgl:S} holds as stated, with $\beta_{\mathsf S}=b^{3m}$ and $N_{\mathsf S}=20(2n+1)^2$.




% ===== C-conj.tex =====
% proofs/C-conj.tex
% The law of Galois conjugation (CL) of the norms of the lifts of characters, used in the proof of
% Theorem psu:arch (= Proposition po:arch). To be \input after C-onelift, C-small, C-small2,
% C-exchanges3, C-periods, C-small3. Labels of this file start with pcj:.
% Only amsmath/amssymb/amsthm are used.
\providecommand{\auO}{\mathcal O}
\providecommand{\auU}{\mathrm U}
\providecommand{\auHom}{\operatorname{Hom}}
\providecommand{\auAut}{\operatorname{Aut}}
\providecommand{\auNm}{\operatorname{Nm}}
\providecommand{\auA}{\mathbb A}
\providecommand{\auPl}{\Phi}
\providecommand{\aubar}[1]{\overline{#1}}
\providecommand{\plone}{\mathbf 1}
\providecommand{\psrho}{\varrho}
\providecommand{\psvac}{\mathbf v}
\providecommand{\psHS}{\mathcal H}
\providecommand{\poA}{\mathcal A}
\providecommand{\pop}{\mathsf p}
% macros of this file (prefix \pcj)
\providecommand{\pcjT}{\operatorname{T}}
\providecommand{\pcjJ}{\mathrm J}
\providecommand{\pcjB}{\mathsf B}
\providecommand{\pcjS}{\mathsf S}
\providecommand{\pcjG}{\mathcal G}
\providecommand{\pcjtr}{\operatorname{Tr}}
\providecommand{\pcjsh}{\psrho^{\sharp}}
\providecommand{\pcjhat}[1]{\widehat{#1}}

\section{The Galois conjugation of the lifts of characters: statement (CL)}\label{pcj:sec}

This section proves the statement (CL) isolated in {\cite[\S7.3.6]{P54}}, which is the only input
of {\cite[Theorem~7.43]{P54}} (that is, of {\cite[Proposition~2.24]{P54}}) not proved there. We use the
notation of {\cite[\S5.1.1]{P49}}, {\cite[\S5.2.1]{P49}}, {\cite[\S5.2.8]{P49}}, {\cite[\S2.1]{P54}},
{\cite[\S2.7]{P54}} (\emph{the lift of a character}) and {\cite[\S3.1.1]{P54}}.

\begin{theorem}[Statement (CL)]\label{pcj:CL}
Let $\sigma\in\auAut(\C/K)$ and $2\le z\le t$. There is $a_\sigma(z)>0$ such that for every critical
pair $p=(\nu,z)$ there is a critical pair $p^\sigma=(\nu^{\sigma},z)$, whose weight $-c^\sigma$ satisfies
$c^\sigma(\sigma\circ\tau)=c(\tau)$ for all $\tau\in\auPl$, with
\[
\bigl|\sigma\bigl(\|\varphi_{p,1}\|^2\bigr)\bigr|=a_\sigma(z)\,\|\varphi_{p^\sigma,1}\|^2 .
\]
\end{theorem}

\emph{Why {\cite[Lemma~5.7]{P49}} does not suffice, and what is done here.} {\cite[Proposition~5.27]{P49}}
gives a $\sigma$-semilinear map $\Lambda_\sigma$ which conjugates jets but commutes with the
Heisenberg translations only up to roots of unity $\zeta_x$; then $\Lambda_\sigma\omega(g)$ and
$\omega(g)\Lambda_\sigma$ differ by an undetermined translation $\psrho_f(y_g)$, and nothing can be
said about the image of the datum of a lift. We show (Proposition~\ref{pcj:Lambda}) that all these
roots of unity are $1$: the jet-conjugating map is unique, and it carries $\psrho_f(x)$ to
$\psrho_f(\pcjB x)$ \emph{exactly}, for the idele $\pcjB=\beta_\sigma$ of {\cite[Theorem~7.9]{P49}}; at
the same time $\sigma$ acts on the values of $\psi_f$ as multiplication by $\pcjB\aubar\pcjB$
(Lemma~\ref{pcj:char}). Writing $\pcjB=ua$ with $u$ of norm one and $a$ a unit equal to $1$ at the
chosen places (Lemma~\ref{pcj:split}), the map $\omega(u)^{-1}\Lambda_\sigma$ fixes the line of the
vacuum, conjugates the Weil operators up to explicit scalars $\sigma(\delta_k)/\delta_k$ with
$|\delta_k|^2\in\Q$ (Lemma~\ref{pcj:vac}), and multiplies the block datum by a constant
(Lemma~\ref{pcj:datum}). Summing over the classes of $\auU(1)$ gives the conjugate of each
coordinate of the lift (Proposition~\ref{pcj:coord}); the scalars cancel between $\sigma$ and
$c\sigma c$ in the norm.

\emph{Dependency order.} Lemma~\ref{pcj:loc} (local constancy) $\to$ Lemma~\ref{pcj:transl} (exact
translates of the theta products) $\to$ Lemma~\ref{pcj:jets} (exact conjugation of their jets)
$\to$ Lemma~\ref{pcj:char} ($\sigma$ on $\psi_f$) $\to$ Proposition~\ref{pcj:Lambda} (the map
$\Lambda$) $\to$ Lemma~\ref{pcj:split} $\to$ Lemma~\ref{pcj:vac} (vacuum and Weil operators) $\to$
Lemma~\ref{pcj:datum} (the block datum) $\to$ Lemmas~\ref{pcj:lift}, \ref{pcj:nusigma}
(the lift as a class sum; conjugate characters) $\to$ Proposition~\ref{pcj:coord} $\to$
Theorem~\ref{pcj:CL}.

\subsection{Setting}\label{pcj:ss-set}

Fix $z\ge2$ and let $W=\auHom(E,V^{(z)})=E^{\mathsf C}$, $\mathsf C=[2z]\times[1]$, with the form
$\sum_i\aubar x_iy_i$; $L=\alpha\pcjhat\auO_E^{\,\mathsf C}$, $\Lambda_L=W(E)\cap L=\mathfrak
a_0^{\mathsf C}$ ({\cite[\S5.2.8]{P49}}). We write $\psrho_f$, $\omega=\omega_f$ for the finite
Heisenberg and Weil operators on $\psHS_f=\psHS_f(W)$, which are unitary for its inner product,
$\beta=\beta_f$, $[x,y]=\beta(x,y)\beta(y,x)^{-1}$, and $\psvac=\psvac_L$; so
$\psrho_f(x)\psrho_f(y)=\beta(x,y)\psrho_f(x+y)$ with $\beta(x,y)=\psi_f(\tfrac12\jmath\langle
y,x\rangle)$, and $\psrho_f(l)\psvac=\epsilon_L(l)\psvac$ for $l\in L$, $\epsilon_L(l)=\psi_f(-\tfrac12
\omega_0\jmath\langle l,l\rangle)$, $\omega_0=(1+\sqrt{-d})/2$ ({\cite[\S5.1.1]{P49}}). We use three
elementary facts: (a) $\psi_f(\aubar y)=\psi_f(y)$, and $\aubar\jmath=-\jmath$; (b) $\psi_f(y)=1$
for $y\in E$ with $\pcjtr_{E/\Q}y=0$, because $\psi_f(y)=\psi_\infty(y)^{-1}=e^{2\boldsymbol\pi i\,
\pcjtr_{E/\Q}y}$; (c) $\psi_f(\jmath\alpha\aubar\alpha t)=1$ for $t\in\pcjhat\auO_E$, by the
self-duality of $\pcjhat\auO_E$ for $(z,u)\mapsto\psi(\jmath\alpha\aubar\alpha\aubar zu)$; since
$\langle l,l'\rangle\in\alpha\aubar\alpha\pcjhat\auO_E$ for $l,l'\in L$, (c) gives
$\psi_f(\jmath\,\mathsf t\langle l,l'\rangle)=1$ for $l,l'\in L$ and $\mathsf t\in\pcjhat\auO_E$.

\emph{Taylor series and jets.} For $f\in\psHS_f$ and $c\in\C^\times$ let $\pcjT_c(f)\in\C[[z]]$
be the Taylor series at $0$ of $z\mapsto\Theta_f(cz)$, in the variables $z_{\tau,\mathsf c}$
($\tau\in\auPl$, $\mathsf c\in\mathsf C$), with $\Theta_f$ as in {\cite[(42)]{P49}}. The map
$f\mapsto\pcjT_c(f)$ is linear and injective, because $f\mapsto\Theta_f$ is injective
({\cite[Lemma~5.26]{P49}}) and an entire function is determined by its Taylor series. For a polynomial
$Q$ the jet is $\pcjJ_{Q,c}(f)=Q(\partial)[\Theta_f(c\,\cdot)](0)=\theta(Q(c\,\cdot)\otimes f)$
({\cite[(43)]{P49}}); if $Q$ is homogeneous of degree $D$, $\pcjJ_{Q,c}(f)=c^D\pcjJ_{Q,1}(f)$.
For $\sigma\in\auAut(\C/K)$ let $\sigma_*$ be the continuous ring endomorphism of $\C[[z]]$ which is
$\sigma$ on the constants and maps $z_{\tau,\mathsf c}$ to $z_{\sigma\tau,\mathsf c}$; here
$\sigma\circ\tau\in\auPl$ because $\sigma$ fixes $K$. Writing $\pcjT_c(f)=\sum_\gamma t_\gamma z^\gamma$
and $Q=\sum_\gamma q_\gamma z^\gamma$, $\pcjJ_{Q,c}(f)=\sum_\gamma\gamma!\,q_\gamma t_\gamma$; as the
renaming of the variables does not change $\gamma!$,
\begin{equation}\label{pcj:eq-jet}
\pcjT_{c'}(f')=\sigma_*\pcjT_c(f)\quad\Longrightarrow\quad\pcjJ_{\sigma_*Q,c'}(f')=\sigma\bigl(\pcjJ_{Q,c}
(f)\bigr)\ \text{ for every polynomial }Q .
\end{equation}

\emph{The data attached to $\sigma$.} Fix $\sigma\in\auAut(\C/K)$ and $\Omega_\sigma$, $\mathfrak c$,
$\mathfrak a=\mathfrak c\aubar{\mathfrak c}^{-1}$, $\beta_\sigma$ as in {\cite[Theorem~7.9]{P49}}; let
$\pcjB$ be $\beta_\sigma$ regarded in $\auA_{E,f}^\times$, so that $\pcjB\pcjhat\auO_E=\mathfrak a
\pcjhat\auO_E$, $\pcjB L=\mathfrak aL$, and $\pcjB\aubar\pcjB$ is a unit ({\cite[Theorem~7.9(a)]{P49}}).
Let $F_0$ be the vector of {\cite[(56)]{P49}} ($\Theta_{F_0}=\Pi$), $F_0^\sigma$ the vector attached
in the proof of {\cite[Proposition~5.27]{P49}} to the lattice $\Lambda^\sigma=\Omega_\sigma\mathfrak a$
and the number $\Omega_\sigma$ ($\Theta_{F_0^\sigma}=\Pi^\sigma$, the product {\cite[(56)]{P49}}
with $\theta^\sigma$ and $\Omega_\sigma$ in place of $\theta$ and $\Omega_0$), and
$t_\sigma=\sigma(\varpi_{\mathrm{vac}}/\mathsf u)\,c_\psi/\Omega_\sigma$, so that $|t_\sigma|^2=
|s_0|^2=\Omega^{-1}$ ({\cite[Proposition~5.27 and~(45)]{P49}}).

\begin{definition}\label{pcj:def-image}
Let $\gamma\in W(E)$ and let $\mathsf m\ge1$ be an integer with $\mathsf m\gamma\in\Lambda_L$. An
\emph{$\mathsf m$-image} of $\gamma$ is a $\gamma'\in W(E)$ with $\gamma'-\pcjB\gamma\in2q_1^2\mathsf m\,
\pcjB L$.
\end{definition}

Since $W(E)$ is dense in $W(\auA_{E,f})$ and $\pcjB L$ is open, $\mathsf m$-images exist for every
admissible $\mathsf m$. If $\gamma'$ and $\delta'$ are $\mathsf m$-images of $\gamma$ and $\delta$,
then $\gamma'+\delta'$ is an $\mathsf m$-image of $\gamma+\delta$; $0$ is an $\mathsf m$-image of $0$;
and if $\mathsf m\mid\mathsf m'$, an $\mathsf m'$-image is an $\mathsf m$-image.

\begin{lemma}[Local constancy]\label{pcj:loc}
For $w\in\psHS_f$ and $x_0\in W(\auA_{E,f})$ there is an open subgroup $U$ of $W(\auA_{E,f})$ with
$\psrho_f(x_0+y)w=\psrho_f(x_0)w$ for all $y\in U$.
\end{lemma}
\begin{proof}
By (W4), $w$ is a finite sum $\sum_jc_j\psrho_f(x_j)\psvac$. For $y\in L$,
$\psrho_f(y)\psrho_f(x_j)\psvac=[y,x_j]\epsilon_L(y)\psrho_f(x_j)\psvac$. The characters
$y\mapsto[y,x_j]$, $y\mapsto\beta(x_0,y)$ of $W(\auA_{E,f})$ are continuous, and $\epsilon_L(y)=1$
for $y\in NL$, $N$ even and large. So there is an open $U\subseteq L$ on which all of them are $1$;
then $\psrho_f(y)w=w$ and $\psrho_f(x_0+y)w=\beta(x_0,y)^{-1}\psrho_f(x_0)\psrho_f(y)w=\psrho_f(x_0)w$.
\end{proof}

\subsection{Exact translates and exact conjugation of jets}\label{pcj:ss-exact}

\begin{lemma}[Exact translates]\label{pcj:transl}
For $\gamma\in W(E)$, with $u_k=\zeta_{q_1}^k/q_1$, $a_{k,i}(z)=\frac{\Omega_0}{c_\psi}u_k^\dagger(z_i)$,
$a^\sigma_{k,i}(z)=\frac{\Omega_\sigma}{c_\psi}u_k^\dagger(z_i)$ and $\mathsf x_{k,i}(\gamma)=
-\pcjtr_{E/K}(\aubar{u_k}\gamma_i)\in K$,
\[
\Theta_{\psrho_f(\gamma)F_0}(z)=\prod_{i,k}\theta_{\mathsf x_{k,i}(\gamma)\Omega_0}\bigl(a_{k,i}(z)\bigr)^{q_1},
\qquad
\Theta_{\psrho_f(\gamma)F_0^\sigma}(z)=\prod_{i,k}\theta^\sigma_{\mathsf x_{k,i}(\gamma)\Omega_\sigma}\bigl(
a^\sigma_{k,i}(z)\bigr)^{q_1},
\]
with no root of unity, the twisted translates being taken for $\Lambda$ and for $\Lambda^\sigma$.
\end{lemma}
\begin{proof}
\emph{The root of unity of {\cite[Lemma~5.26]{P49}} is $1$ for rational $\gamma$.} In the proof of that
lemma, $\theta=\theta\circ\psrho(-\gamma)$ and $\psrho(-\gamma)=\psrho_\infty(-\gamma^\natural)\otimes
\psrho_f(-\gamma)$; so
$\theta(\psrho_\infty(z)(1\otimes\psrho_f(\gamma)f))=\beta_\infty(-\gamma^\natural,z)\,
\beta_f(-\gamma,\gamma)\,\theta(\psrho_\infty(z-\gamma^\natural)(1\otimes f))$, and the root of unity
there is $\beta_f(-\gamma,\gamma)=\psi_f(-\frac12\jmath\langle\gamma,\gamma\rangle)$. Now
$\langle\gamma,\gamma\rangle\in E^+$, so $\pcjtr_{E/K}\langle\gamma,\gamma\rangle\in K\cap\R=\Q$ and
$\pcjtr_{E/\Q}(\jmath\langle\gamma,\gamma\rangle)=\pcjtr_{K/\Q}(\jmath)\pcjtr_{E/K}\langle\gamma,
\gamma\rangle=0$; by fact (b) the root of unity is $1$. Hence, for every $f$,
$\Theta_{\psrho_f(\gamma)f}(z)=e^{\langle z,\gamma^\natural\rangle-|\gamma^\natural|^2/2}\Theta_f(z-
\gamma^\natural)$.

\emph{The product.} Put $\mathsf y=\mathsf x_{k,i}(\gamma)\Omega_0=-a_{k,i}(\gamma^\natural)$. By
{\cite[(30)]{P49}} and $\boldsymbol\pi|\Omega_0|^2/(Ac_\psi^2)=1$ (as in {\cite[\S5.2.8]{P49}}),
$\langle z,\gamma^\natural\rangle=\sum_{i,k}q_1H(a_{k,i}(\gamma^\natural),a_{k,i}(z))=-\sum_{i,k}q_1
\frac{\boldsymbol\pi}A\aubar{\mathsf y}\,a_{k,i}(z)$ and $|\gamma^\natural|^2=\sum_{i,k}q_1\frac{
\boldsymbol\pi}A|\mathsf y|^2$ (the $\mathsf y$ depending on $(i,k)$), while $a_{k,i}(z-\gamma^\natural)=
a_{k,i}(z)+\mathsf y$. So $e^{\langle z,\gamma^\natural\rangle-|\gamma^\natural|^2/2}\Pi(z-\gamma^\natural)
=\prod_{i,k}\bigl[e^{-\frac{\boldsymbol\pi}A(a_{k,i}(z)\aubar{\mathsf y}+\frac12|\mathsf y|^2)}\theta(a_{k,i}
(z)+\mathsf y)\bigr]^{q_1}=\prod_{i,k}\theta_{\mathsf y}(a_{k,i}(z))^{q_1}$ by the definition of the
twisted translate ({\cite[\S5.2.7]{P49}}). For $F_0^\sigma$ the same computation applies: the area of
$\Lambda^\sigma=\Omega_\sigma\mathfrak a$ is $A^\sigma=|\Omega_\sigma|^2\sqrt d/2$ (as $\auNm\mathfrak a=1$),
so $\boldsymbol\pi|\Omega_\sigma|^2/(A^\sigma c_\psi^2)=2\boldsymbol\pi/(\sqrt dc_\psi^2)=1$.
\end{proof}

\begin{lemma}[Exact conjugation of jets of translates]\label{pcj:jets}
Let $\gamma\in W(E)$, $\mathsf m\ge1$ with $\mathsf m\gamma\in\Lambda_L$, and $\gamma'$ an
$\mathsf m$-image of $\gamma$. Then $\pcjT_{t_\sigma}(\psrho_f(\gamma')F_0^\sigma)=\sigma_*\pcjT_{s_0}
(\psrho_f(\gamma)F_0)$.
\end{lemma}
\begin{proof}
By Lemma~\ref{pcj:transl} and~{\cite[(45)]{P49}}, $\pcjT_{s_0}(\psrho_f(\gamma)F_0)=\prod_{i,k}G_{k,i}
(\ell_{k,i})^{q_1}$, where $G_{k,i}$ is the Taylor series of $\theta_{\mathsf x_{k,i}\Omega_0}$ and
$\ell_{k,i}=a_{k,i}(s_0z)=(\varpi_{\mathrm{vac}}/\mathsf u)\,u_k^\dagger(z_i)$; likewise
$\pcjT_{t_\sigma}(\psrho_f(\gamma')F_0^\sigma)=\prod_{i,k}G'_{k,i}(\ell'_{k,i})^{q_1}$ with $G'_{k,i}$ the
Taylor series of $\theta^\sigma_{\mathsf x'_{k,i}\Omega_\sigma}$, $\mathsf x'_{k,i}=\mathsf x_{k,i}
(\gamma')$, and $\ell'_{k,i}=a^\sigma_{k,i}(t_\sigma z)=\sigma(\varpi_{\mathrm{vac}}/\mathsf u)\,
u_k^\dagger(z_i)$. As $\sigma_*$ is a continuous ring homomorphism, $\sigma_*(G(\ell))=(\sigma G)
(\sigma_*\ell)$ for a series $G$ in one variable and a linear form $\ell$. Since $\sigma$ fixes $K$
and $\tau(E)$ is a CM field, $\sigma(\aubar{\tau(u_k)})=\aubar{(\sigma\tau)(u_k)}$; so $\sigma_*\ell_{k,i}
=\sigma(\varpi_{\mathrm{vac}}/\mathsf u)\sum_\tau\aubar{(\sigma\tau)(u_k)}\,z_{\sigma\tau,i}=\ell'_{k,i}$.
It remains to see $\sigma G_{k,i}=G'_{k,i}$, which is {\cite[Theorem~7.9(b)]{P49}} for $\mathsf x=
\mathsf x_{k,i}$, $n=q_1\mathsf m$, $\mathsf x'=\mathsf x'_{k,i}$. Indeed $n\mathsf x=-\pcjtr_{E/K}
(\zeta_{q_1}^{-k}\mathsf m\gamma_i)\in\pcjtr_{E/K}(\mathfrak a_0)\subseteq\auO_K$, as $\mathfrak a_0
\subseteq\mathfrak d_{E/K}^{-1}$ is an $\auO_E$-module; and, $\pcjtr_{E/K}$ being
$\auA_{K,f}$-linear and $\beta_\sigma\in\auA_{K,f}$, $\mathsf x'-\beta_\sigma\mathsf x=-\pcjtr_{E/K}
(\aubar{u_k}(\gamma'_i-\pcjB\gamma_i))$ with $\aubar{u_k}(\gamma'_i-\pcjB\gamma_i)\in2q_1\mathsf m\,
\mathfrak a\,\alpha\pcjhat\auO_E=2q_1\mathsf m\,\mathfrak a\,\pcjhat{\mathfrak a_0}$, whose trace lies in
$2q_1\mathsf m\,\pcjhat{\mathfrak a}=2n\pcjhat{\mathfrak a}$.
\end{proof}

\begin{lemma}[Closeness]\label{pcj:close}
Let $\gamma,\delta\in W(E)$, $\mathsf m\gamma,\mathsf m\delta\in\Lambda_L$, and $\gamma',\delta'$ their
$\mathsf m$-images. Then $\beta(\delta',\gamma')=\psi_f(\frac12\jmath\langle\pcjB\gamma,\pcjB\delta
\rangle)=\psi_f(\frac12\jmath\,\pcjB\aubar\pcjB\langle\gamma,\delta\rangle)$.
\end{lemma}
\begin{proof}
$\langle\gamma',\delta'\rangle-\langle\pcjB\gamma,\pcjB\delta\rangle=\langle\gamma'-\pcjB\gamma,
\delta'\rangle+\langle\pcjB\gamma,\delta'-\pcjB\delta\rangle$. Here $\delta,\gamma\in\mathsf m^{-1}L$, so
$\pcjB\gamma,\delta'\in\pcjB\mathsf m^{-1}L$, and $\gamma'-\pcjB\gamma,\delta'-\pcjB\delta\in2q_1^2
\mathsf m\pcjB L$. As $\langle\pcjB l,\pcjB l'\rangle=\pcjB\aubar\pcjB\langle l,l'\rangle$ with
$\pcjB\aubar\pcjB$ a unit, both terms lie in $2\langle L,L\rangle\pcjhat\auO_E$, and
$\psi_f(\frac12\jmath\cdot2\mathsf t)=1$ for $\mathsf t\in\langle L,L\rangle\pcjhat\auO_E$ by fact (c).
\end{proof}

\subsection{The jet-conjugating map}\label{pcj:ss-Lambda}

Every vector of $\psHS_f$ is a finite combination of vectors $\psrho_f(\gamma)F_0$, $\gamma\in W(E)$
({\cite[\S5.2.8]{P49}}: the $\psrho_f(x)F_0$ span, and each is a root of unity times
$\psrho_f(\gamma)F_0$ for rational $\gamma$ close to $x$).

\begin{lemma}[Definition of $\Lambda$ on translates]\label{pcj:def}
There is a unique map $\Lambda\colon\psHS_f\to\psHS_f$ with $\pcjT_{t_\sigma}(\Lambda f)=\sigma_*
\pcjT_{s_0}(f)$ for all $f$. It is $\sigma$-semilinear and injective, and $\Lambda(\psrho_f(\gamma)
F_0)=\psrho_f(\gamma')F_0^\sigma$ whenever $\gamma'$ is an $\mathsf m$-image of $\gamma\in W(E)$.
\end{lemma}
\begin{proof}
Uniqueness and injectivity follow from the injectivity of $\pcjT_{t_\sigma}$ and $\pcjT_{s_0}$ and
the injectivity of $\sigma_*$. For $f=\sum_jc_j\psrho_f(\gamma_j)F_0$ put $\Lambda f=\sum_j\sigma(c_j)
\psrho_f(\gamma_j')F_0^\sigma$ with $\gamma'_j$ any $\mathsf m_j$-images. By Lemma~\ref{pcj:jets} and the
$\sigma$-semilinearity of $\sigma_*\circ\pcjT_{s_0}$, $\pcjT_{t_\sigma}(\Lambda f)=\sigma_*\pcjT_{s_0}(f)$; so
$\Lambda f$ does not depend on the choices (injectivity of $\pcjT_{t_\sigma}$), $\Lambda$ is well defined
and $\sigma$-semilinear, and has the stated value on translates.
\end{proof}

\begin{lemma}[$\sigma$ on the additive character]\label{pcj:char}
$\sigma(\psi_f(y))=\psi_f(\pcjB\aubar\pcjB\,y)$ for all $y\in\auA_{E,f}$.
\end{lemma}
\begin{proof}
Let $\lambda\in q_1\Lambda_L$, $\delta\in W(E)$, $\mathsf m\ge1$ with $\mathsf m\lambda,\mathsf m\delta
\in\Lambda_L$, and $\lambda',\delta'$ their $\mathsf m$-images. By {\cite[Lemma~7.13]{P49}} (for the frame
$(\zeta_{q_1}^k)$, $\mathsf b=q_1$), $\psrho_f(\lambda)F_0=\epsilon_L(\lambda)F_0$. Apply $\Lambda$, using
Lemma~\ref{pcj:def} with the $\mathsf m$-image $0$ of $0$: $\psrho_f(\lambda')F_0^\sigma=\Lambda(\psrho_f
(\lambda)F_0)=\sigma(\epsilon_L(\lambda))F_0^\sigma$. Now compute $\Lambda v$ for $v=\psrho_f(\delta)
\psrho_f(\lambda)F_0$ in two ways. First $v=\epsilon_L(\lambda)\psrho_f(\delta)F_0$, so $\Lambda v=
\sigma(\epsilon_L(\lambda))\psrho_f(\delta')F_0^\sigma$. Second $v=\beta(\delta,\lambda)\psrho_f(\delta+
\lambda)F_0$ and $\delta'+\lambda'$ is an $\mathsf m$-image of $\delta+\lambda$, so
\[
\Lambda v=\sigma(\beta(\delta,\lambda))\psrho_f(\delta'+\lambda')F_0^\sigma=\sigma(\beta(\delta,\lambda))
\beta(\delta',\lambda')^{-1}\psrho_f(\delta')\psrho_f(\lambda')F_0^\sigma=\sigma(\beta(\delta,\lambda))
\beta(\delta',\lambda')^{-1}\sigma(\epsilon_L(\lambda))\psrho_f(\delta')F_0^\sigma .
\]
Since $F_0^\sigma\ne0$ ($\Pi^\sigma\ne0$) and $\psrho_f(\delta')$ is invertible, $\sigma(\beta(\delta,
\lambda))=\beta(\delta',\lambda')$, which by Lemma~\ref{pcj:close} gives
$\sigma(\psi_f(\frac12\jmath\langle\lambda,\delta\rangle))=\psi_f(\frac12\jmath\pcjB\aubar\pcjB\langle
\lambda,\delta\rangle)$. Fix $\lambda\ne0$; for $y\in E$ the vector $\delta=2y(\jmath\langle\lambda,
\lambda\rangle)^{-1}\lambda$ has $\frac12\jmath\langle\lambda,\delta\rangle=y$. So the assertion holds for
$y\in E$. Both sides are characters of $\auA_{E,f}$ trivial on an open subgroup ($\psi_f$ is), hence
continuous, and $E$ is dense in $\auA_{E,f}$.
\end{proof}

\begin{proposition}[The jet-conjugating map scales translations exactly]\label{pcj:Lambda}
The map $\Lambda$ of Lemma~\ref{pcj:def} is a $\sigma$-semilinear bijection of $\psHS_f$ with
\[
\pcjT_{t_\sigma}(\Lambda f)=\sigma_*\pcjT_{s_0}(f),\qquad\Lambda\,\psrho_f(x)=\psrho_f(\pcjB x)\,
\Lambda\qquad(f\in\psHS_f,\ x\in W(\auA_{E,f})).
\]
\end{proposition}
\begin{proof}
\emph{Rational $x$.} Let $\delta,\gamma\in W(E)$ and $v=\psrho_f(\gamma)F_0$. Let $\mathsf m_0\ge1$
with $\mathsf m_0\delta,\mathsf m_0\gamma\in\Lambda_L$, let $\mathsf m$ be a multiple of $\mathsf m_0$,
and $\delta',\gamma'$ $\mathsf m$-images. As in the proof of Lemma~\ref{pcj:char},
\[
\Lambda\psrho_f(\delta)v=\sigma(\beta(\delta,\gamma))\psrho_f(\delta'+\gamma')F_0^\sigma=\sigma(\beta
(\delta,\gamma))\beta(\delta',\gamma')^{-1}\psrho_f(\delta')\Lambda v=\psrho_f(\delta')\Lambda v,
\]
because $\sigma(\beta(\delta,\gamma))=\psi_f(\frac12\jmath\pcjB\aubar\pcjB\langle\gamma,\delta\rangle)=
\beta(\delta',\gamma')$ by Lemmas~\ref{pcj:char} and~\ref{pcj:close}. Take $\mathsf m=\mathsf m_0\,j!$ and
let $j\to\infty$: then $\delta'\to\pcjB\delta$ in $W(\auA_{E,f})$, the left side does not change, and by
Lemma~\ref{pcj:loc} (at $x_0=\pcjB\delta$, $w=\Lambda v$) $\psrho_f(\delta')\Lambda v=\psrho_f(\pcjB\delta)
\Lambda v$ for $j$ large. By linearity, $\Lambda\psrho_f(\delta)=\psrho_f(\pcjB\delta)\Lambda$ on
$\psHS_f$.

\emph{All $x$.} Let $f\in\psHS_f$ and $\delta_j\in W(E)$, $\delta_j\to x$. By Lemma~\ref{pcj:loc}
(at $x$ and $f$, and at $\pcjB x$ and $\Lambda f$), for $j$ large $\psrho_f(\delta_j)f=\psrho_f(x)f$ and
$\psrho_f(\pcjB\delta_j)\Lambda f=\psrho_f(\pcjB x)\Lambda f$, so $\Lambda\psrho_f(x)f=\Lambda\psrho_f(\delta_j)
f=\psrho_f(\pcjB\delta_j)\Lambda f=\psrho_f(\pcjB x)\Lambda f$.

\emph{Bijectivity.} $\Lambda$ is injective (Lemma~\ref{pcj:def}). Its image is a non-zero subspace
stable under every $\psrho_f(y)=\Lambda\psrho_f(\pcjB^{-1}y)\Lambda^{-1}$ ($\pcjB$ is an idele), hence is
$\psHS_f$ by the irreducibility of the Heisenberg representation.
\end{proof}

\begin{remark}\label{pcj:rem-pk}
Proposition~\ref{pcj:Lambda} sharpens the commutation rules in the proof of
{\cite[Proposition~5.27]{P49}}: the map $\Lambda_\sigma$ defined there satisfies the jet identity of {\cite[Lemma~5.7(i)]{P49}}, hence equals $\Lambda$
by uniqueness, and its roots of unity $\zeta'$ and $\zeta_x$ are all $1$ when the images $\gamma'$ are
chosen as in Definition~\ref{pcj:def-image} (the only source of roots of unity there is
{\cite[Lemma~5.26]{P49}}, whose root of unity is $1$ for rational translations by
Lemma~\ref{pcj:transl}).
\end{remark}

\subsection{The vacuum, the Weil operators and the block datum}\label{pcj:ss-vac}

\begin{lemma}[Splitting the scale]\label{pcj:split}
Let $\mathsf R$ be a finite set of finite places of $E$ with $\mathsf R\cap\aubar{\mathsf R}=\emptyset$.
There are $u\in\auA_{E,f}^\times$ with $u\aubar u=1$ and $a\in\pcjhat\auO_E^\times$ with $a_w=1$ for
$w\in\mathsf R$, such that $\pcjB=ua$. Then $\pcjB\aubar\pcjB=a\aubar a$.
\end{lemma}
\begin{proof}
It suffices to find a unit $a$, equal to $1$ on $\mathsf R$, with $a\aubar a=\pcjB\aubar\pcjB$, and
to put $u=\pcjB a^{-1}$; we do it at each finite place $v$ of $E^+$. If $v$ does not split, $E_v$ is a
field and $\pcjB_v\aubar{\pcjB_v}=\auNm(\pcjB_v)$ is a unit, so $\pcjB_v$ is a unit; put $a_v=\pcjB_v$
($v$ is not below $\mathsf R$, whose places satisfy $w\ne\aubar w$). If $v=w\aubar w$ splits, write
elements of $E_v=E_w\times E_{\aubar w}$ as pairs, with $\aubar{(x,y)}=(\aubar y,\aubar x)$ for the
conjugation isomorphisms $E_{\aubar w}\leftrightarrow E_w$. Then $\pcjB\aubar\pcjB=(\nu_w,\aubar{\nu_w})$
with $\nu_w=\pcjB_w\aubar{\pcjB_{\aubar w}}$ a unit, and $(a_w,a_{\aubar w})=(1,\aubar{\nu_w})$ if
$w\in\mathsf R$, $(\nu_w,1)$ otherwise (at most one of $w,\aubar w$ is in $\mathsf R$; name them so that
$\aubar w\notin\mathsf R$), satisfies $a\aubar a=(a_w\aubar{a_{\aubar w}},a_{\aubar w}\aubar{a_w})=
(\nu_w,\aubar{\nu_w})$.
\end{proof}

We apply Lemma~\ref{pcj:split} with $\mathsf R=\mathsf M$, the chosen places (all of them lie above
split places of $E^+$ and $\mathsf M\cap\aubar{\mathsf M}=\emptyset$, {\cite[\S3.1.1]{P54}}). The norm-one
idele $u$ is a point of $\auU(1)(\auA_{E^+,f})$ and acts on $W$ as a scalar, a central element of
$\auU(W)(\auA_f)$. Put $\Lambda'=\omega(u)^{-1}\Lambda$, and, for $x\in W(\auA_{E,f})$,
\[
\pcjsh(x)=e(x)\psrho_f(x),\qquad e(x)=\psi_f\bigl(\tfrac12\omega_0\jmath\langle x,x\rangle\bigr),
\]
so that $\pcjsh(l)\psvac=\psvac$ for $l\in L$ ($e=\epsilon_L^{-1}$ on $L$). For
$k\in\auU(W)(\auA_f)$ let $I_k=[L:L\cap k^{-1}L]$ and $\pcjS_k=\sum_{q\in L/(L\cap k^{-1}L)}\pcjsh(kq)
\psvac$.

\begin{lemma}[Vacuum and Weil operators]\label{pcj:vac}
\begin{enumerate}
\item[(a)] $\Lambda'$ is a $\sigma$-semilinear bijection with $\Lambda'\psrho_f(x)=\psrho_f(ax)\Lambda'$
and $\Lambda'\pcjsh(x)=\pcjsh(ax)\Lambda'$ for all $x$, and $\sigma(\psi_f(y))=\psi_f(a\aubar ay)$.
\item[(b)] $\Lambda'\psvac=\mathsf c\,\psvac$ for a number $\mathsf c\ne0$.
\item[(c)] $\pcjsh(x)\pcjsh(y)=\psi_f(-\jmath\omega_0\langle x,y\rangle)\,\pcjsh(x+y)$. The terms of
$\pcjS_k$ do not depend on the representatives $q$, $\pcjS_k\ne0$, and $\omega(k)\psvac=\delta_k\pcjS_k$
with $\delta_k\in\C^\times$, $|\delta_k|^2=I_k^{-1}$.
\item[(d)] For $k\in\auU(W)(\auA_f)$ put $\mathsf m_\sigma(k)=\sigma(\delta_k)/\delta_k$. Then
$\Lambda'\omega(k)=\mathsf m_\sigma(k)\,\omega(k)\Lambda'$ and $\Lambda\omega(k)=\mathsf m_\sigma(k)\,
\omega(k)\Lambda$ on $\psHS_f$; $\mathsf m_\sigma$ is a homomorphism, $\mathsf m_\sigma(k)=1$ if $kL=L$,
and, for $\sigma'=c\circ\sigma\circ c$, $\mathsf m_{\sigma'}(k)\,\aubar{\mathsf m_\sigma(k)}=1$.
\end{enumerate}
\end{lemma}
\begin{proof}
(a) By (W1), $\omega(u)^{-1}\psrho_f(\pcjB x)\omega(u)=\psrho_f(u^{-1}\pcjB x)=\psrho_f(ax)$, which with
Proposition~\ref{pcj:Lambda} gives the first law. The last assertion is Lemma~\ref{pcj:char} with
$\pcjB\aubar\pcjB=a\aubar a$; as $\langle ax,ax\rangle=a\aubar a\langle x,x\rangle$, it gives
$\sigma(e(x))=e(ax)$, hence the second law.

(b) For $l\in L$, $a^{-1}l\in L$ and $\psrho_f(l)\Lambda'\psvac=\Lambda'\psrho_f(a^{-1}l)\psvac=\sigma
(\epsilon_L(a^{-1}l))\Lambda'\psvac$, and $\sigma(\epsilon_L(a^{-1}l))=\psi_f(a\aubar a\cdot(-\frac12\omega_0
\jmath)(a\aubar a)^{-1}\langle l,l\rangle)=\epsilon_L(l)$. By (W4) $\Lambda'\psvac$ lies on the line
of $\psvac$, and it is not $0$.

(c) $e(x)e(y)/e(x+y)=\psi_f(-\frac12\omega_0\jmath(\langle x,y\rangle+\langle y,x\rangle))$ and
$\beta(x,y)=\psi_f(\frac12\jmath\langle y,x\rangle)$; with $s=\omega_0\langle x,y\rangle$ and
$1-\omega_0=\aubar{\omega_0}$ the product is $\psi_f(\frac12\jmath(\aubar s-s))$, and $\psi_f(\frac12
\jmath\aubar s)=\psi_f(-\aubar{\frac12\jmath s})=\psi_f(\frac12\jmath s)^{-1}$ by fact (a); this is the
formula. For $l'\in L\cap k^{-1}L$, $\pcjsh(kq)\pcjsh(kl')\psvac=\pcjsh(kq)\psvac$ (as $kl'\in L$) and
$\pcjsh(kq)\pcjsh(kl')=\psi_f(-\jmath\omega_0\langle q,l'\rangle)\pcjsh(k(q+l'))=\pcjsh(k(q+l'))$ by fact
(c); so the terms depend only on $q$ modulo $L\cap k^{-1}L$. For $l\in L$ the same computation gives
$\pcjsh(kl)\pcjsh(kq)\psvac=\pcjsh(k(q+l))\psvac$, and $q\mapsto q+l$ permutes the classes; so
$\pcjsh(kl)\pcjS_k=\pcjS_k$, that is $\psrho_f(kl)\pcjS_k=e(l)^{-1}\pcjS_k=\epsilon_L(l)\pcjS_k$ (note
$e(kl)=e(l)$). Also $\psrho_f(kl)\omega(k)\psvac=\omega(k)\psrho_f(l)\psvac=\epsilon_L(l)\omega(k)\psvac$ by
(W1). As $kL$ is self-dual, (W4) shows that $\omega(k)\psvac$ is a multiple of $\pcjS_k$ once
$\pcjS_k\ne0$. The vectors $\pcjsh(x)\psvac$, $x\in W(\auA_{E,f})/L$, are, up to roots of unity, the
basis of (W4); they are unit vectors (the $\psrho_f(x)$ are unitary) and pairwise orthogonal, being
eigenvectors of the unitary operators $\psrho_f(l)$, $l\in L$, with the characters $l\mapsto[l,x]
\epsilon_L(l)$, which are distinct for distinct classes because $L$ is self-dual. The classes $kq$,
$q\in L/(L\cap k^{-1}L)$, are distinct modulo $L$; so $\pcjS_k$ is a sum of $I_k$ orthonormal vectors,
$\|\pcjS_k\|^2=I_k$, and $1=\|\omega(k)\psvac\|^2=|\delta_k|^2I_k$ (normalising $\|\psvac\|=1$).

(d) By (a), $\Lambda'\pcjsh(kq)\psvac=\pcjsh(akq)\Lambda'\psvac=\mathsf c\,\pcjsh(k(aq))\psvac$, and
$q\mapsto aq$ permutes $L/(L\cap k^{-1}L)$ ($a$ is a unit scalar commuting with $k$); so $\Lambda'
\pcjS_k=\mathsf c\,\pcjS_k$ and $\Lambda'\omega(k)\psvac=\sigma(\delta_k)\mathsf c\,\pcjS_k=\mathsf
m_\sigma(k)\,\omega(k)\Lambda'\psvac$. For $x\in W(\auA_{E,f})$, by (W1) and $e(kx)=e(x)$,
$\omega(k)\pcjsh(x)=\pcjsh(kx)\omega(k)$; hence
\[
\Lambda'\omega(k)\pcjsh(x)\psvac=\Lambda'\pcjsh(kx)\omega(k)\psvac=\pcjsh(akx)\Lambda'\omega(k)\psvac=
\mathsf m_\sigma(k)\pcjsh(akx)\omega(k)\Lambda'\psvac=\mathsf m_\sigma(k)\,\omega(k)\Lambda'\pcjsh(x)\psvac,
\]
and these vectors span $\psHS_f$. As $u$ is central in $\auU(W)(\auA_f)$ and $\omega$ is a
representation, $\Lambda\omega(k)=\omega(u)\Lambda'\omega(k)=\mathsf m_\sigma(k)\omega(k)\Lambda$. The
homomorphism property follows from $\Lambda\omega(kk')=\Lambda\omega(k)\omega(k')$ and $\Lambda\ne0$. If
$kL=L$ then $I_k=1$, $\pcjS_k=\psvac$ and $\omega(k)\psvac=\psvac$ (W4), so $\delta_k=1$. Finally
$\sigma'(x)=\aubar{\sigma(\aubar x)}$ and $\aubar{\delta_k}=I_k^{-1}\delta_k^{-1}$ with $I_k\in\Q$, so
$\sigma'(\delta_k)=\aubar{I_k^{-1}\sigma(\delta_k)^{-1}}$ and $\mathsf m_{\sigma'}(k)=I_k^{-1}/(\delta_k
\aubar{\sigma(\delta_k)})=\aubar{\delta_k}/\aubar{\sigma(\delta_k)}=\aubar{\mathsf m_\sigma(k)}^{-1}$.
\end{proof}

\emph{The block datum.} By (W5), at a place $v=w\aubar w$ of $E^+$ with $w\in\mathsf M$, $\psHS_v$ is
the space of Schwartz--Bruhat functions on $E_w^{2z}$, the vacuum is $\plone[\alpha_w\auO_w^{2z}]$, i.e.\ $\plone[\auO_w^{2z}]$ in the coordinate $T=\tilde x$ of (N5) (Lemma~\ref{pgx:dict}; at a chosen place $\alpha_w$ is a unit and the canonical idele of {\cite[\S5.1.1]{P49}} is trivial there, so $T=x$), in which all functions on $E_w^{2z}$ are written below
(proof of {\cite[Lemma~3.2]{P54}}), $\psrho(y_w,0)$ is the translation by $y_w$ and $\psrho(0,y_{\aubar w})$
the multiplication by $x\mapsto[(x,0),(0,y_{\aubar w})]$ (with $y$ read in the same coordinate). Write $E_{\mathsf M}^{2z}=\bigoplus_{w\in
\mathsf M}E_w^{2z}$. The block datum of {\cite[\S3.1.1]{P54}} for $\vartheta=1$ is $\Phi_{\mathsf M}=
\Psi_{\mathsf M}\otimes\psvac^{\mathsf M}$, $\Psi_{\mathsf M}=\bigotimes_{w\in\mathsf M}\plone[T\equiv e_1
\bmod w]$, with $\psvac^{\mathsf M}$ the vacuum at the other places.

\begin{lemma}[The block datum]\label{pcj:datum}
For every Schwartz--Bruhat function $\Psi$ on $E_{\mathsf M}^{2z}$,
$\Lambda'(\Psi\otimes\psvac^{\mathsf M})=\mathsf c\,(\sigma\circ\Psi)\otimes\psvac^{\mathsf M}$. In
particular $\Lambda\Phi_{\mathsf M}=\mathsf c\,\omega(u)\Phi_{\mathsf M}$.
\end{lemma}
\begin{proof}
Let $y=(y_w,y_{\aubar w})_{w\in\mathsf M}\in W(\auA_{E,f})$ be supported at the places above $\mathsf M$,
and put $Y_1=(y_w,0)_w$, $Y_2=(0,y_{\aubar w})_w$. Then $\psrho_f(y)=\beta(Y_1,Y_2)^{-1}\psrho_f(Y_1)
\psrho_f(Y_2)$, so $\psrho_f(y)\psvac=F_y\otimes\psvac^{\mathsf M}$ with
\[
F_y(x)=\beta(Y_1,Y_2)^{-1}\,[(x,0)+Y_1,Y_2]\,\plone[\auO^{2z}](x+y_w)\qquad(x\in E_{\mathsf M}^{2z}).
\]
Since $a_w=1$ for $w\in\mathsf M$, $a(x,0)=(x,0)$, $aY_1=Y_1$ and $ay=(y_w,a_{\aubar w}y_{\aubar w})$. By
Lemma~\ref{pcj:vac}(a), $\sigma(\beta(U,U'))=\beta(aU,aU')$ and $\sigma([U,U'])=[aU,aU']$ for all $U,U'$
(both are values of $\psi_f$ at multiples of $\langle U',U\rangle$, and $\langle aU',aU\rangle=a\aubar a
\langle U',U\rangle$). Hence $\sigma\circ F_y=F_{ay}$, and by Lemma~\ref{pcj:vac}(a),(b),
$\Lambda'(F_y\otimes\psvac^{\mathsf M})=\psrho_f(ay)\Lambda'\psvac=\mathsf c\,F_{ay}\otimes\psvac^{\mathsf M}=
\mathsf c\,(\sigma\circ F_y)\otimes\psvac^{\mathsf M}$. The functions $F_y$ span the Schwartz--Bruhat
functions: up to non-zero constants they are the products of a character $x\mapsto[(x,0),(0,q)]$ and of
a translate of $\plone[\auO^{2z}]$; these characters are all the continuous characters of
$E_{\mathsf M}^{2z}$ ($[\,,]$ is non-degenerate and $E_w^{2z}$, $E_{\aubar w}^{2z}$ are isotropic and
in duality), and the indicator of a coset of $\varpi^j\auO^{2z}$ inside $\auO^{2z}$ is an average of such
characters times $\plone[\auO^{2z}]$. By $\sigma$-semilinearity of both sides the first assertion
follows. For the second, $\Psi_{\mathsf M}$ takes the values $0,1$, so $\sigma\circ\Psi_{\mathsf M}=
\Psi_{\mathsf M}$ and $\Lambda\Phi_{\mathsf M}=\omega(u)\Lambda'\Phi_{\mathsf M}=\mathsf c\,\omega(u)
\Phi_{\mathsf M}$.
\end{proof}

\subsection{The lifts of characters and their conjugates}\label{pcj:ss-lift}

Let $p=(\nu,z)$ be a critical pair, $\nu$ of weight $-c$, $l_\tau=c(\tau)-z$, $D=\sum_\tau l_\tau$. By
{\cite[\S2.1]{P54}} and {\cite[\S3.1.1]{P54}} (with $\vartheta=1$: $a=1$, $b=\nu\upsilon^z$), the
representation defining $\varphi_{p,1}$ is
\[
\omega_\nu(g,h)=\nu(g)^{-1}\upsilon(g)^{-z}\,\omega(k_{g,h}),\qquad k_{g,h}\colon x\mapsto hxg^{-1},
\]
for $g\in\auU(1)(\auA_{E^+})$, $h\in\auU(V^{(z)})(\auA_{E^+})$, where $\omega$ is the Weil representation
of $\auU(W)(\auA)$ defined with $\upsilon$; we write $\omega_{\nu,f}$, $\nu_f$, $\upsilon_f$ for the finite
parts. Note $k_{g,h}k_{g',h'}=k_{gg',hh'}$, and the scalar $u$ of \S\ref{pcj:ss-vac} is $k_{u^{-1},1}$.
Put $\Phi^\circ=c_{\mathrm{vac}}^{2z}\Phi_{\mathsf M}$. The weight of $\poA(V^{(z)};p,1)$ has one row of
length $l_\tau$ at $\tau$, so $W_\lambda=\bigotimes_\tau\operatorname{Sym}^{l_\tau}\C^{2z}\subseteq
\bigotimes_\tau(\C^{2z})^{\otimes l_\tau}$, with the form for which the pure tensors $e_I$ of basis vectors
are orthonormal ({\cite[\S2.1]{P54}}). For a multi-index $\beta=(\beta_\tau)_\tau$, $\beta_\tau\in
\N^{2z}$ of degree $l_\tau$, let $e_\beta$ be the sum of the $e_I$ of content $\beta$ and $n_\beta=
\langle e_\beta,e_\beta\rangle\in\Z_{>0}$ their number; the $e_\beta$ form an orthogonal basis. By
{\cite[(3)]{P54}}, $\pop_\xi(z)=\langle\xi,\bigotimes_\tau z_\tau^{\otimes l_\tau}\rangle$ in the Fock
variables $z_\tau=(z_{\tau,\mathsf c})_{\mathsf c}$; it is homogeneous of degree $l_\tau$ in $z_\tau$, and
$P_\beta:=\pop_{e_\beta}=n_\beta z^\beta$. For $\sigma\in\auAut(\C/K)$ let $\pi\beta$ be the multi-index
$(\pi\beta)_{\tau'}=\beta_{\sigma^{-1}\tau'}$; then $\sigma_*P_\beta=n_\beta z^{\pi\beta}$, and $n_{\pi\beta}
=n_\beta$.

\begin{lemma}[The lift as a class sum of jets]\label{pcj:lift}
Let $\pcjG$ be a set of representatives of the finite set $\auU(1)(E^+)\backslash\auU(1)(\auA_{E^+,f})/
K^1_M$. There is $w_0>0$, depending only on $z$, $\mathsf M$ and the fixed Haar measure on
$[\auU(1)]$, such that for every critical pair $p=(\nu,z)$, $h\in\auU(V^{(z)})(\auA_{E^+,f})$ and
$\xi\in W_\lambda$,
\[
\langle\xi,\varphi_{p,1}(h)\rangle=w_0\sum_{g\in\pcjG}\pcjJ_{\pop_\xi,\pi_Ms_0}\bigl(\omega_{\nu,f}(g,h)
\Phi^\circ\bigr).
\]
Moreover, for every polynomial $Q$ homogeneous of degree $l_\tau$ in $z_\tau$, every $s\in\auU(1)(E^+)$,
$\kappa\in K^1_M$ and $g$, $\pcjJ_{Q,c}(\omega_{\nu,f}(sg\kappa,h)\Phi^\circ)=\pcjJ_{Q,c}(\omega_{\nu,f}
(g,h)\Phi^\circ)$.
\end{lemma}
\begin{proof}
For $g\in\auU(1)(\auA_{E^+})$ write $g=g_\infty g_f$. The proof of {\cite[Lemma~3.3]{P54}} shows that
$\omega_{\nu,\infty}(g_\infty,1)$ fixes every polynomial homogeneous of degree $l_\tau$ in $z_\tau$ (it
multiplies such a polynomial by the product of $\epsilon^{-l_\tau}$, of the splitting character and of
$\nu_\tau(\epsilon)^{-1}\upsilon_\tau(\epsilon)^{-z}$, which is trivial). Hence, by {\cite[(43)]{P49}},
$\theta(\omega_\nu(g,h)(Q\otimes\Phi))=\theta(Q\otimes\omega_{\nu,f}(g_f,h)\Phi)=\pcjJ_{Q,1}(\omega_{\nu,f}
(g_f,h)\Phi)$. For $s\in\auU(1)(E^+)$, $\nu(s)=\upsilon(s)=1$ and $k_{s,1}\in\auU(W)(E^+)$, so
$\theta\circ\omega_\nu(s,1)=\theta$ by (W1); this gives the invariance under $s$ (for $c=1$, hence for
all $c$ by homogeneity). For $\kappa\in K^1_M$, $\omega_{\nu,f}(\kappa,1)\Phi_{\mathsf M}=\Phi_{\mathsf M}$:
at $w\in\mathsf M$ by {\cite[(1)]{P54}} ($\kappa_w\equiv1\bmod w$, the tame character and $\eta_w$ are
trivial on it, and $x\equiv e_1$ if and only if $x\kappa_w\equiv e_1$), and at the other places because
$\nu$, $\upsilon$ are unramified there and $\omega(k_{\kappa,1})$ fixes the local vacuum, $k_{\kappa,1}$
being a unit scalar (W4). Now {\cite[(4)]{P54}} with $y=\plone$, {\cite[(13)]{P54}} and these
invariances give $\langle\xi,\varphi_{p,1}(h)\rangle=\kappa_p[K^1:K^1_M]\,v_0\sum_{g\in\pcjG}\pcjJ_{\pop_\xi,1}
(\omega_{\nu,f}(g,h)\Phi_{\mathsf M})$, where $v_0$ is the common volume of the sets $\auU(1)(E^+)
\backslash\auU(1)(E^+)\auU(1)(\R)gK^1_M$ (they are translates of one another, $\auU(1)$ being
commutative). Since $\kappa_p=c_{\mathrm{vac}}^{2z}(\pi_Ms_0)^D$ and $\pop_\xi$ is homogeneous of degree
$D$, $\kappa_p\pcjJ_{\pop_\xi,1}(\omega_{\nu,f}(g,h)\Phi_{\mathsf M})=\pcjJ_{\pop_\xi,\pi_Ms_0}(\omega_{\nu,f}
(g,h)\Phi^\circ)$; put $w_0=[K^1:K^1_M]v_0$.
\end{proof}

\begin{lemma}[Conjugate characters]\label{pcj:nusigma}
Let $\chi$ be an automorphic character of $\auU(1)$ of weight $-c$ and $\sigma\in\auAut(\C/K)$. There is a
unique automorphic character $\chi^{(\sigma)}$ of $\auU(1)$ with finite part $\sigma\circ\chi_f$; it has
weight $-c^\sigma$, $c^\sigma(\sigma\tau)=c(\tau)$, it is unitary, and its restriction to the units at a
finite place $w$ is $\sigma$ composed with that of $\chi$. For $\sigma'=c\circ\sigma\circ c$,
$\chi^{(\sigma')}=\chi^{(\sigma)}$. If $(\nu,z)$ is a critical pair, so is $(\nu^{(\sigma)},z)$.
\end{lemma}
\begin{proof}
Write $\chi_\infty(x)=\prod_\tau\tau(x)^{e(\tau)}$ on $\auU(1)(E^+\otimes\R)$, identified with
$\prod_\tau\{|\epsilon|=1\}$ through the $\tau$, where $e(\tau)$ is determined by $c(\tau)$ through the
convention for weights; put $\chi^{(\sigma)}_\infty(x)=\prod_{\tau'}\tau'(x)^{e(\sigma^{-1}\tau')}$ and
$\chi^{(\sigma)}=(\sigma\circ\chi_f)\chi^{(\sigma)}_\infty$. It is a continuous homomorphism ($\chi_f$ is
trivial on an open subgroup, hence so is $\sigma\circ\chi_f$). For $s\in\auU(1)(E^+)$, $\chi_f(s)=
\chi_\infty(s)^{-1}=\prod_\tau\tau(s)^{-e(\tau)}$, so $\sigma(\chi_f(s))=\prod_\tau(\sigma\tau)(s)^{-e(\tau)}
=\chi_\infty^{(\sigma)}(s)^{-1}$: $\chi^{(\sigma)}$ is trivial on $\auU(1)(E^+)$. It is unitary because
$\auU(1)(E^+)\backslash\auU(1)(\auA_{E^+})$ is compact. Uniqueness: two such characters agree on the
subgroup $\auU(1)(E^+)\auU(1)(\auA_{E^+,f})$, which is dense in $\auU(1)(\auA_{E^+})$ because $\auU(1)(E^+)$ is dense in $\auU(1)(E^+\otimes\R)$ (weak approximation for the torus $\auU(1)$, which is rational over $E^+$ by the Cayley parametrisation $y\mapsto y/\aubar y$). For $\sigma'$: $\sigma'(\chi_f(g))=\aubar{\sigma(
\aubar{\chi_f(g)})}=\aubar{\sigma(\chi_f(g))^{-1}}=\sigma(\chi_f(g))$, since $|\chi_f(g)|=1$ and
$|\sigma(\chi_f(g))|=1$ (unitarity of $\chi^{(\sigma)}$). For a critical pair: $z\le c^\sigma\le4b$ since
$c^\sigma$ is a permutation of $c$, and $\nu^{(\sigma)}$ is tamely ramified exactly where $\nu$ is.
\end{proof}

By Lemma~\ref{pcj:nusigma} applied to $\upsilon|_{\auU(1)}$ (of weight $1$),
$|\sigma(\upsilon_f(g))|=1$ for all finite $g$. For $\sigma\in\auAut(\C/K)$ define, on
$\auU(1)(\auA_{E^+,f})$,
\[
\rho_\sigma(g)=\mathsf m_\sigma(k_{g,1})\,\upsilon_f(g)^z\,\sigma(\upsilon_f(g))^{-z},
\]
a character (Lemma~\ref{pcj:vac}(d)). Put $\mathsf l_\sigma=\sigma(c_{\mathrm{vac}})^{2z}c_{\mathrm{vac}}^{-2z}
\mathsf c\ne0$; then $\Lambda\Phi^\circ=\mathsf l_\sigma\,\omega(u)\Phi^\circ$ by Lemma~\ref{pcj:datum}. The
constants $\mathsf l_\sigma$, $u$, $\mathsf c$ depend on $\sigma$ and $z$ only.

\begin{lemma}[The kernels]\label{pcj:kernel}
For every automorphic character $\nu$ of $\auU(1)$, every character $\chi$ of $\auU(1)(\auA_{E^+,f})$
and all finite $g,h$,
\[
\Lambda\bigl(\omega_{\nu,f}(g,h)\Phi^\circ\bigr)=\mathsf l_\sigma\,\mathsf m_\sigma(k_{1,h})\,
\frac{\chi(gu^{-1})\upsilon_f(gu^{-1})^z\mathsf m_\sigma(k_{g,1})}{\sigma(\nu_f(g))\,\sigma(\upsilon_f(g))^{z}}
\,\omega_{\chi,f}(gu^{-1},h)\Phi^\circ ,
\]
where $\omega_{\chi,f}(g,h)=\chi(g)^{-1}\upsilon_f(g)^{-z}\omega(k_{g,h})$.
\end{lemma}
\begin{proof}
{\emergencystretch=6em By $\sigma$-semilinearity, Lemma~\ref{pcj:vac}(d) and $k_{g,h}=k_{g,1}k_{1,h}$,
$\Lambda(\omega_{\nu,f}(g,h)\Phi^\circ)=\sigma(\nu_f(g)\upsilon_f(g)^z)^{-1}\allowbreak \mathsf m_\sigma(k_{g,1})
\allowbreak \mathsf m_\sigma(k_{1,h})\,\allowbreak \omega(k_{g,h})\Lambda\Phi^\circ$, and $\omega(k_{g,h})\Lambda\Phi^\circ=\mathsf
l_\sigma\,\omega(k_{g,h}k_{u^{-1},1})\Phi^\circ=\mathsf l_\sigma\,\omega(k_{gu^{-1},h})\Phi^\circ=\mathsf
l_\sigma\,\chi(gu^{-1})\allowbreak \upsilon_f(gu^{-1})^z\allowbreak \omega_{\chi,f}(gu^{-1},h)\Phi^\circ$.\par}
\end{proof}

\begin{lemma}[The character $\rho_\sigma$]\label{pcj:rho}
$\rho_\sigma$ is trivial on $\auU(1)(E^+)$ and on $K^1$, it has finite order, $|\mathsf m_\sigma(k_{g,1})|
=1$ for all $g$, and $\rho_{\sigma'}=\rho_\sigma$ for $\sigma'=c\circ\sigma\circ c$.
\end{lemma}
\begin{proof}
For $g\in K^1$, $k_{g,1}$ is a unit scalar, so $k_{g,1}L=L$ and $\mathsf m_\sigma(k_{g,1})=1$
(Lemma~\ref{pcj:vac}(d)), and $\upsilon_f(g)=1$ ($\upsilon$ is unramified at every finite place).

Let $s\in\auU(1)(E^+)$. Take any critical pair $p=(\nu,z)$ (for instance a flat one,
{\cite[Lemma~7.27(a)]{P54}}). By {\cite[Proposition~3.4]{P54}}, $\varphi_{p,1}\ne0$; by
Lemma~\ref{pcj:lift} there are $g,h,\beta$ with $\pcjJ_{P_\beta,s_0}(f)\ne0$ for $f=\omega_{\nu,f}(g,h)
\Phi^\circ$. Put $f_s=\omega_{\nu,f}(sg,h)\Phi^\circ$; then $\pcjJ_{P_\beta,s_0}(f_s)=\pcjJ_{P_\beta,s_0}(f)$
(Lemma~\ref{pcj:lift}). By Proposition~\ref{pcj:Lambda} and \eqref{pcj:eq-jet}, $\pcjJ_{\sigma_*P_\beta,
t_\sigma}(\Lambda f_s)=\pcjJ_{\sigma_*P_\beta,t_\sigma}(\Lambda f)=\sigma(\pcjJ_{P_\beta,s_0}(f))\ne0$. Apply
Lemma~\ref{pcj:kernel} with $\chi=\nu^{(\sigma)}_f$ to $f$ and to $f_s$. Since $\sigma_*P_\beta$ is
homogeneous of degree $l_{\sigma^{-1}\tau'}$ in $z_{\tau'}$, which is the degree attached to the weight
$-c^\sigma$ of $\nu^{(\sigma)}$, and $(\nu^{(\sigma)},z)$ is a critical pair (Lemma~\ref{pcj:nusigma}),
Lemma~\ref{pcj:lift} for it gives $\pcjJ_{\sigma_*P_\beta,t_\sigma}(\omega_{\chi,f}(sgu^{-1},h)\Phi^\circ)=
\pcjJ_{\sigma_*P_\beta,t_\sigma}(\omega_{\chi,f}(gu^{-1},h)\Phi^\circ)$. Comparing the two instances of
Lemma~\ref{pcj:kernel}, the ratio of the scalars for $sg$ and for $g$ must be $1$; it is
\[
\frac{\chi(s)\upsilon_f(s)^z\mathsf m_\sigma(k_{s,1})}{\sigma(\nu_f(s))\sigma(\upsilon_f(s))^z}=
\mathsf m_\sigma(k_{s,1})\upsilon_f(s)^z\sigma(\upsilon_f(s))^{-z}=\rho_\sigma(s),
\]
as $\chi(s)=\sigma(\nu_f(s))$. So $\rho_\sigma$ is trivial on $\auU(1)(E^+)K^1$, a subgroup of finite
index; hence $\rho_\sigma$ has finite order and $|\rho_\sigma|=1$, and then $|\mathsf m_\sigma(k_{g,1})|=
|\sigma(\upsilon_f(g))/\upsilon_f(g)|^z=1$. Finally, by Lemma~\ref{pcj:vac}(d),
$\mathsf m_{\sigma'}(k_{g,1})=\aubar{\mathsf m_\sigma(k_{g,1})}^{-1}=\mathsf m_\sigma(k_{g,1})$, and
$\sigma'(\upsilon_f(g))=\sigma(\upsilon_f(g))$ by Lemma~\ref{pcj:nusigma}; so $\rho_{\sigma'}=\rho_\sigma$.
\end{proof}

\begin{definition}\label{pcj:def-psigma}
For a critical pair $p=(\nu,z)$ and $\sigma\in\auAut(\C/K)$ let $\nu^\sigma=\nu^{(\sigma)}\rho_\sigma^{-1}$
(with $\rho_\sigma$ extended by $1$ at the infinite places) and $p^\sigma=(\nu^\sigma,z)$.
\end{definition}

By Lemmas~\ref{pcj:nusigma} and~\ref{pcj:rho}, $\nu^\sigma$ is an automorphic character of weight
$-c^\sigma$ with the same ramification as $\nu^{(\sigma)}$ ($\rho_\sigma$ is unramified), so $p^\sigma$ is
a critical pair, flat if $p$ is; and $p^{\sigma'}=p^\sigma$ for $\sigma'=c\circ\sigma\circ c$.

\begin{proposition}[The conjugates of the coordinates]\label{pcj:coord}
Let $p=(\nu,z)$ be a critical pair, $h\in\auU(V^{(z)})(\auA_{E^+,f})$, and write $a_\beta(h)=\langle
e_\beta,\varphi_{p,1}(h)\rangle$ and $a^\sigma_{\beta'}(h)=\langle e_{\beta'},\varphi_{p^\sigma,1}(h)\rangle$.
Then
\[
\sigma\bigl(a_\beta(h)\bigr)=C_{\sigma,p}\,\mathsf m_\sigma(k_{1,h})\,a^\sigma_{\pi\beta}(h),\qquad
C_{\sigma,p}=\frac{\sigma(w_0)}{w_0}\Bigl(\frac{\sigma(\pi_M)\,t_\sigma}{\pi_Ms_0}\Bigr)^{D}\mathsf
l_\sigma\,\frac{\rho_\sigma(u)}{\sigma(\nu_f(u))\,\upsilon_f(u)^{z}},
\]
and $|C_{\sigma,p}|=|\sigma(w_0)/w_0|\,|\mathsf l_\sigma|$, which does not depend on $p$.
\end{proposition}
\begin{proof}
Lemma~\ref{pcj:kernel} with $\chi=\nu^\sigma_f$, for which $\chi(gu^{-1})=\sigma(\nu_f(g))\rho_\sigma(g)^{-1}
\sigma(\nu_f(u))^{-1}\rho_\sigma(u)$, gives, by the definition of $\rho_\sigma$,
\[
\Lambda\bigl(\omega_{\nu,f}(g,h)\Phi^\circ\bigr)=\mathsf l_\sigma\frac{\rho_\sigma(u)}{\sigma(\nu_f(u))
\upsilon_f(u)^z}\,\mathsf m_\sigma(k_{1,h})\,\omega_{\nu^\sigma,f}(gu^{-1},h)\Phi^\circ .
\]
By Lemma~\ref{pcj:lift}, homogeneity, \eqref{pcj:eq-jet} and Proposition~\ref{pcj:Lambda},
\[
\sigma(a_\beta(h))=\sigma(w_0)\sigma(\pi_M)^D\sum_{g\in\pcjG}\pcjJ_{\sigma_*P_\beta,t_\sigma}\bigl(\Lambda\,
\omega_{\nu,f}(g,h)\Phi^\circ\bigr),
\]
and $\sigma_*P_\beta=n_\beta z^{\pi\beta}=P_{\pi\beta}$, the polynomial $\pop_{e_{\pi\beta}}$ for the weight
of $p^\sigma$, homogeneous of the same total degree $D$. As $g\mapsto gu^{-1}$ permutes the double cosets
and the summands depend only on the double coset (Lemma~\ref{pcj:lift} for $p^\sigma$),
$\sum_g\pcjJ_{P_{\pi\beta},t_\sigma}(\omega_{\nu^\sigma,f}(gu^{-1},h)\Phi^\circ)=(t_\sigma/\pi_Ms_0)^D
a^\sigma_{\pi\beta}(h)/w_0$. This is the formula. For the absolute value: $|t_\sigma|=|s_0|$;
$|\sigma(\pi_M)|^2=M=|\pi_M|^2$ (as used in Step~3 of the proof of {\cite[Proposition~14.23]{P46}} in
{\cite[\S5.1.5]{P49}}, for every $\sigma$); $|\rho_\sigma(u)|=|\upsilon_f(u)|=1$, and $|\sigma(\nu_f(u))|=1$
by Lemma~\ref{pcj:nusigma}.
\end{proof}

\begin{proof}[Proof of Theorem~\ref{pcj:CL}]
Let $\sigma'=c\circ\sigma\circ c\in\auAut(\C/K)$, so that $\sigma(\aubar x)=\aubar{\sigma'(x)}$. For an
embedding $\tau$ of the CM field $E$, $c\circ\tau=\tau\circ c_E$, hence $\sigma'\circ\tau=c\circ(\sigma\tau)
\circ c_E=\sigma\circ\tau$: the permutations $\pi$ for $\sigma$ and $\sigma'$ coincide, and $p^{\sigma'}=
p^\sigma$ (Definition~\ref{pcj:def-psigma}). The norm of {\cite[\S14.3]{P46}} is $\|\Xi\|^2=\sum_{h}
r_h\langle\Xi(h),\Xi(h)\rangle$, $h$ over representatives of $\auU(V^{(z)})(E^+)\backslash\auU(V^{(z)})
(\auA_{E^+,f})/K$ and $r_h=\#(\Gamma\cap hKh^{-1})^{-1}\in\Q_{>0}$ (proof of {\cite[Lemma~2.1]{P54}}); the
level $K$ and its character are the same for $p$ and $p^\sigma$ (they are defined by $\Phi_{\mathsf M}$
and by the action of $\auU(V^{(z)})$, which does not involve $\nu$). With the orthogonal basis $e_\beta$,
$\langle\Xi(h),\Xi(h)\rangle=\sum_\beta|\langle e_\beta,\Xi(h)\rangle|^2/n_\beta$. So, by
Proposition~\ref{pcj:coord} for $\sigma$ and for $\sigma'$,
\[
\sigma\bigl(\|\varphi_{p,1}\|^2\bigr)=\sum_hr_h\sum_\beta\frac{\sigma(a_\beta(h))\,\aubar{\sigma'(a_\beta(h))}}
{n_\beta}=C_{\sigma,p}\aubar{C_{\sigma',p}}\sum_hr_h\,\mathsf m_\sigma(k_{1,h})\aubar{\mathsf m_{\sigma'}(k_{1,h})}
\sum_\beta\frac{|a^\sigma_{\pi\beta}(h)|^2}{n_{\pi\beta}} .
\]
By Lemma~\ref{pcj:vac}(d), $\mathsf m_\sigma(k)\aubar{\mathsf m_{\sigma'}(k)}=1$; and $\beta\mapsto\pi\beta$
is a bijection onto the multi-indices of the weight of $p^\sigma$. Hence $\sigma(\|\varphi_{p,1}\|^2)=
C_{\sigma,p}\aubar{C_{\sigma',p}}\,\|\varphi_{p^\sigma,1}\|^2$, and the theorem holds with
$a_\sigma(z)=|C_{\sigma,p}C_{\sigma',p}|=|\sigma(w_0)\sigma'(w_0)|\,w_0^{-2}\,|\mathsf l_\sigma\mathsf
l_{\sigma'}|>0$, which depends on $\sigma$ and $z$ only. The weight of $p^\sigma$ is $-c^\sigma$ with
$c^\sigma(\sigma\tau)=c(\tau)$ (Lemma~\ref{pcj:nusigma}).
\end{proof}

\begin{remark}\label{pcj:rem-use}
(i) In the proof of {\cite[Theorem~7.43]{P54}}, (CL) is used for a main pair $p$ and the flat pair $q$; $q^\sigma$
is flat, and $p^\sigma$ has $l^\sigma_{\sigma\tau}=l_\tau$, as required there. (ii) We do not decide whether the
unramified character $\rho_\sigma$ is trivial (that is, whether $\nu^\sigma=\nu^{(\sigma)}$): (CL) only
asks for some critical pair with the permuted weight, and Definition~\ref{pcj:def-psigma} avoids computing
the multiplier of the Weil representation on scalars, which would need the explicit splitting. (iii) Nothing in this section uses the size of the constants.
\end{remark}



\appendix

\section{Statements used from \texorpdfstring{\cite{P46}}{[Row D]}, \texorpdfstring{\cite{P49}}{[Theta]}, \texorpdfstring{\cite{P51}}{[Moment]}, \texorpdfstring{\cite{P54}}{[L-values]}, \texorpdfstring{\cite{P55}}{[Galois]}}\label{C1:quoted}

For the convenience of the reader we reproduce here, without proofs, the statements of other papers that are used in this one, in the notation of the paper quoted. Each is headed by its place in that paper.

\begin{quotedthm}{\cite[Lemma~13.4]{P46}}{}
For admissible data the following hold.
\begin{enumerate}
\item[(S1)] $x^{10^{6}}\le x_1$, $x_1^{10^{6}}\le b$,
$4\log x_1\le(\log x_1)^{6}\le\log b$,
$(\log b)^{100}\le x$, $\log b\ge128$, and
$\log x\ge10^{14}(q_0+2d+|\aneu|)$.
\item[(S2)] $x^{11}\le2^{s}\le2^{22}x^{11}$; $l_0(s)\le2592\log x$;
$2m\le r$, so $r/2\le k\le r$.
\item[(S3)] $x^{11}\le T(s)=\#\anPa\le f\le q_1\le x^{12}$, and
$f=\varphi(q_0)\,\#\anPa$; moreover $1\le\#\anPa/\#\anPb\le\log b$.
\item[(S4)] $\anrd_E=\sqrt d\;q_1\prod_{p\mid q_0}p^{-1/(p-1)}$, and
$x^{10}\le\anrd_E\le x^{12}$.
\item[(S5)] $1\le xb/x_1\le\kappa(\tau)\le b\le a\le b^{2}$ for all $\tau$;
$b^{20}\le M\le b^{22}$;
$20\log b\le\log(M\anrd_E/2\pi)\le23\log b$; and
$m\log b\le\log N\le25\,m\log b$.
\item[(S6)] With $\Omega:=e^{\aneu}\cdot2\pi/\anrd_E$: $0<\Omega\le1$ and
$20\log b\le\log(M/\Omega)\le23\log b$.
\item[(S7)] With $n=m/2$ and $\varrho=n-1$ (so that $r=k+2\varrho+1$): $r$ is odd; every prime
dividing $q_1$ divides $q_0$; $q_0^{\,q_0}\le x$; $\varrho\ge10^{5}$;
$k\ge512000\,n^{9}$ and $k\ge1000\,(2\varrho+2)^{9}$; $k\ge2\varrho^{2}+3$; $18\varrho f\le b^{8}$;
$f\le b/1000$; $\anrd_E\le b$; and $N\le b^{60n}$. Moreover, with $\Delta=2r^{8}$:
\begin{itemize}
\item[(a)] $\varrho\ge10^{8}$ (so $n\ge100$); $n^{7}\le k$; $12n^{3}+500\le k$; $k-3n-6>10^{9}n^{3}$;
\item[(b)] $r\ge100$ (so $2r^{8}\ge500r^{3}+10r^{5}$); $3r^{3}f\le r^{7}$; $8nr\le b$; $kf\ge40$; $f\ge4000\,n$;
\item[(c)] $2f\le2^{2\varrho}$ and $2^{-\varrho}\le b^{-6}$;
\item[(d)] $(200\,f\,k\,\anrd_E)^{1000}\le b$ and $(25\,n^{3}f\,\anrd_E)^{1000}\le b$;
\item[(e)] $\sqrt d\le\anrd_E$, and $n,r,f,\sqrt d\le b$; $M\le b^{23}$;
\item[(f)] $\binom{a+k-1}{k-1}^{f}\le b^{2kf}$ (this bounds the dimension $d_\infty=\dim W_{\lambda^X}$ of
{\cite[\S14.3]{P46}}, a product of at most $f$ such binomial coefficients);
\item[(g)] $m\Delta+r\le\kappa(\tau)$ for all $\tau$; $k+2b+2m\Delta\le b^{2}$ (this bounds every
entry of $2\varsigma_\tau$, so that $L_\tau:=\max_i|2\varsigma_{\tau,i}|\le b^{2}$); and
$(r+1)^{4}(b^{2}+r+1+2n^{2})\le b^{5}$.
\end{itemize}
\end{enumerate}
\end{quotedthm}

\begin{quoteddef}{\cite[Definition~14.1]{P46}}{}
A family $y$ \emph{has height at most $c\in\R$} if there are $H\in\C^\times$, $\alpha\in\R$ and
reals $h_\ell\ge0$, one for each prime, such that $|\sigma H|\le e^\alpha$ for all
$\sigma\in\auAut(\C)$, $|H\,y_\iota|_\iota\le\ell(\iota)^{h_{\ell(\iota)}}$ for all $\iota$, and
$\alpha+\sum_{\ell\in T}h_\ell\log\ell\le c$ for every finite set $T$ of primes.
\end{quoteddef}

\begin{quoteddef}{\cite[Definition~14.3]{P46}}{ (Twisted lattice; index; covolume)}
Let $V$ be a finite-dimensional complex inner product space. A \emph{twisted lattice} $\auLat$ on
$V$ consists of a set $\Xi$ of \emph{twists} with a distinguished element $1$; for each
$\xi\in\Xi$ an inner product space $H_\xi$ and an injective linear map
$\Theta_\xi\colon V\to H_\xi$; a relation ``$\xi$ is allowed at $\ell$''; for each $\xi$ and
$\iota$ an $\auO_\iota$-submodule $\Lambda_{\xi,\iota}\subseteq H_\xi$; and a linear isomorphism
$\operatorname{coord}$ of $H_1$ onto a subspace of $\C^{\mathsf C}$, $\mathsf C$ finite, with a
positive definite matrix $G\in M_{\mathsf C}(\Q)$ such that
$\langle\varphi,\varphi'\rangle=\aubar{\operatorname{coord}(\varphi)}^{\,t}G\operatorname{coord}(\varphi')$
on $H_1$. The condition is
\begin{equation*}\tag*{{\cite[(39)]{P46}}}
\langle\varphi,\varphi'\rangle\in\auO_\iota\qquad\text{for }\varphi\in\Lambda_{\xi,\aubar\iota},\
\varphi'\in\Lambda_{\xi,\iota},\ \xi\text{ allowed at }\ell(\iota).
\end{equation*}
\begin{enumerate}
\item[(a)] The module of \emph{vectors of $V$ integral at $\iota$} is
$\auLat_\iota=\{y\in V:\Theta_\xi(y)\in\Lambda_{\xi,\iota}\text{ for every }\xi\text{ allowed at }
\ell(\iota)\}$.
\item[(b)] A finite family $y=(y_1,\dots,y_n)$ in $V$ \emph{has index at most $\delta$} if there
are reals $e_\ell\ge0$ with $\sum_{\ell\in T}e_\ell\log\ell\le\delta$ for all finite $T$ such
that for every $\iota$ there is $u\in M_n(\auO_\iota)$ with $|\det u|_\iota\ge
\ell(\iota)^{-e_{\ell(\iota)}}$ and $\sum_ku_{ik}y_k\in\auLat_\iota$ for all $i$.
\item[(c)] For $\sigma\in\auAut(\C)$ and $y\in V$ let $\|\Theta_1y\|_\sigma^2$ be the value of the
form $G$ at the vector $\sigma(\operatorname{coord}(\Theta_1y))$, where $\sigma$ is applied to
each coordinate. The family $y$ \emph{has covolume at most $c$} if, for every $\sigma$, the Gram
determinant for $G$ of the vectors $\sigma(\operatorname{coord}(\Theta_1y_i))$ is at most $e^c$.
\item[(d)] $Z(\xi)=\det(\Theta_\xi^*\Theta_\xi)>0$.
\end{enumerate}
\end{quoteddef}

\begin{quotedthm}{\cite[Theorem~14.4]{P46}}{ (A twist for each prime)}
Let $\auLat$ be a twisted lattice on $V$ and let $y_1,\dots,y_n,z_1,\dots,z_{n'}$ be a basis of
$V$. For each $\iota\in\auI$ let $\theta_\iota\in\Xi$ be allowed at $\ell(\iota)$. Let
$\gamma\colon\Xi\to\C^\times$ and $c,\delta,b_{\mathrm{short}},b_{\mathrm{den}},b_{\mathrm{vol}}
\in\R$ be such that
\begin{enumerate}
\item[(i)] $y$ has covolume at most $c$ and index at most $\delta$;
\item[(ii)] $\gamma(\xi)\Theta_\xi(z_i)\in\Lambda_{\xi,\iota}$ for all $i$, all $\iota$ and all
$\xi$ allowed at $\ell(\iota)$;
\item[(iii)] $\|\Theta_1z_i\|_\sigma^2\le\exp(b_{\mathrm{short}})$ for all $i$ and $\sigma$;
\item[(iv)] the family $\iota\mapsto|\gamma(\theta_\iota)|^{-2}$ has height at most
$b_{\mathrm{den}}$, and $\iota\mapsto Z(1)/Z(\theta_\iota)$ has height at most $b_{\mathrm{vol}}$.
\end{enumerate}
Then $0\le c+2\delta+n'(b_{\mathrm{short}}+b_{\mathrm{den}})+b_{\mathrm{vol}}$.
\end{quotedthm}

\begin{quotedthm}{\cite[Lemma~14.5]{P46}}{ (Free levels)}
Let $n\le2n_m$ and let $J\subseteq K^\circ$ consist of elements congruent to $1$ modulo every
chosen place. Then $\Gamma\cap gJg^{-1}=1$ for all $g\in G_f$.
\end{quotedthm}

\begin{quoteddef}{\cite[Definition~14.8]{P46}}{ (Prime sums, gain, main amount, remainder)}
Let $y\ge0$ be real and $\sigma\in\auPl$.
\begin{enumerate}
\item[(a)] The \emph{local prime sum} at $p$ is
$s_p(y,\sigma)=v(\pi_p)\sum_{e\in I_p}\bigl(y/p^e-\kappa(u(-p^e)\sigma)\bigr)^+$, and the
\emph{prime sum} is
$\ausfin(y,\sigma)=\sum_{1\le\nu\le a}\Lambda_K(\nu)\bigl(y/\nu-\kappa(u(-\nu)\sigma)\bigr)^+$.
\item[(b)] The \emph{gain} at $p$ is
\[
\augain_p=m\Bigl(\sum_{\tau\in\auPa}s_p(a,\tau)+k\sum_{\sigma\in\auPl}s_p(\kappa(\sigma),\sigma)
\Bigr).
\]
\item[(c)] The size of a derivative of order $j$ is $\ausinf(j)=\log j!+j\log(M/\Omega)$; the
\emph{direction amount} is $\auhdir=\sum_{\sigma\in\auPl}\bigl(\ausinf(\kappa(\sigma))-
\ausfin(\kappa(\sigma),\sigma)\bigr)$; the \emph{main amount} is $\auhmain=m\,k\,\auhdir$.
\item[(d)] The \emph{expected value} is
$\ausmain=a\bigl(\log a-\tfrac1f\sum_{\tau\in\auPl}\log\kappa(\tau)+c_1\bigr)$, and the
\emph{total remainder} is $\aurem=\sum_{\tau\in\auPa}\bigl(\ausfin(a,\tau)-\ausmain\bigr)$.
\end{enumerate}
\end{quoteddef}

\begin{quotedthm}{\cite[Theorem~14.15]{P46}}{ (the Covolume Theorem)}
Let $y_1,\dots,y_n\in\auT$ be linearly independent with $n\ge(1-4/b^2)\audT$, of index at most
$\delta$ and covolume at most $c$ in $\auLat_D$. Then
\[
c+2\delta\ \ge\ \audT\cdot\bigl(\auhmain-3\,f\,a\,(\log b)^4\bigr).
\]
\end{quotedthm}

\begin{quotedthm}{\cite[Lemma~14.21]{P46}}{ (Short bases)}
Let $n\ge1$ and let $c_0$ be a positive integer in the different of $E$. For every
$g\in\auU(E^n)(\auA_{E^+,f})$ there is $Y\in\auGL_n(E)$ with $g^{-1}Y\in M_n(\widehat{\auO}_E)$
all of whose entries have all conjugates of absolute value at most $C_{\mathrm{sb}}=6c_0(fn)^2$.
\end{quotedthm}

\begin{quotedthm}{\cite[Proposition~14.22]{P46}}{ (Short ramified forms span)}
$\auSram$ is spanned by the forms $\varphi\in\auSram$ that are integral at every
$j\colon\C\to\C_\ell$, for every $\ell$, and satisfy $\|\sigma(\varphi(g))\|^2\le e^{\aueu}$ for
all $\sigma\in\auAut(\C)$ and all $g$.
\end{quotedthm}

\begin{quotedthm}{\cite[Proposition~14.23]{P46}}{ (One kernel)}
Let $0\le j<m$, let $\phi$ be a translate of the vacuum and $\sigma\in\auAut(\C)$. Then for all
$g,g'$ and all $\xi$,
\[
\|\sigma(\theta_j(\phi;g,g')\xi)\|\le\exp\bigl(x^{50}+\tfrac12\mathrm{rs}_j\bigr)\,\|\sigma(\xi)\|,
\qquad\mathrm{rs}_j=\sum_{\tau\in\auPl}\ \sum_{\text{rows }i}\ausinf(\text{length of row }i
\text{ at }\tau).
\]
\end{quotedthm}

\begin{quotedthm}{\cite[Proposition~14.27]{P46}}{ (Equivariance of the lower chain)}
For every tame character $\vartheta$ and every $u\in\auG$ there is $\beta_\vartheta(u)$ in the
stabiliser $K^\circ$ of the lattice of $Y_m$ such
that $(\iota_\vartheta(\varrho(u)z))(g)=(\iota_\vartheta z)(g\beta_\vartheta(u))$ for all
$z\in\auAz$ and all $g$.
\end{quotedthm}

\begin{quotedthm}{\cite[Proposition~14.28]{P46}}{ (Passage to the good part)}
Let $z\in\auAz$, $s\in\auG$, and $z'=\#\auG\,\varrho(e_0)\varrho(s)z$.
\begin{enumerate}
\item[(a)] If $c\in\C$ and $c\cdot\iota_\vartheta z$ is integral at $j\colon\C\to\C_\ell$, so is
$c\cdot\iota_\vartheta z'$.
\item[(b)] If $\sigma\in\auAut(\C)$ and all values of $\sigma(\iota_1z)$ have squared norm at most
$e^B$, then $\|\iota_1z'\|_\sigma^2\le\exp(x^{99}+B)$.
\end{enumerate}
\end{quotedthm}

\begin{quotedthm}{\cite[Proposition~14.35]{P46}}{ (Injectivity of the lower chain)}
For every $m$-fold twist $\vartheta$ the map $\iota_\vartheta\colon\auAz\to\auAm(\vartheta)$ is
injective.
\end{quotedthm}

\begin{quotedthm}{\cite[Proposition~14.36]{P46}}{ (Equality of multiplicities)}
For $1\le j\le m$ and every element $h$ of the spherical Hecke algebra away from $M$ (the
compactly supported functions on the finite adelic group outside the chosen places that are
bi-invariant under the stabiliser of the lattice, at all places, split or not),
$\operatorname{tr}(h\,|\,e_0\mathcal A(W_j))=\operatorname{tr}(h\,|\,\auT)$. In particular
$\dim e_0\mathcal A(W_j)=\audT$, and each system of eigenvalues of the Hecke operators at the split
places prime to $M$ has the same multiplicity in $e_0\mathcal A(W_j)$ and in $\auT$. (The
hypothesis of {\cite[Proposition~4.6]{P51}} holds for the data: at every $\tau$, two entries of
$\varsigma_\tau$ on coordinates of different signs {\cite[(43)]{P51}} differ by at least
$\Delta=2r^8\ge500r^3+10r^5$.)
\end{quotedthm}

\begin{quotedthm}{\cite[Proposition~14.37]{P46}}{ (Factorisation of the chain determinant)}
For $1\le j\le m$ and a twist $\vartheta$ let $Z_j(\vartheta)=\det(u^*u)>0$, where
$u=\Theta_{j,\vartheta}$ is the theta lift from $e_0\mathcal A(W_j)$ to algebraic modular forms on
$\auU(V_j)$ with the character $\mu_j^0\vartheta$, normalised as in \S\ref{pc:ss-rallis} (joint
harmonic at infinity; at the chosen places the indicator of the inclusion; at $w\in Q_\vartheta$
the function $x\mapsto\psi_w(\det x^{\mathrm{top}})$; values in the algebraic modular forms of level
$J^V_\vartheta$ with their inner product). There are $e\in\Z$ and the ramified ratios $\rho_{ij}$,
$0<i<j\le m$, of \eqref{pc:eq-rho}, such that
for every $m$-fold twist $\vartheta$
\[
\frac{Z(1)}{Z(\vartheta)}=\Bigl((\auNm Q_\vartheta)^{e/2}\prod_{0<i<j\le m}\rho_{ij}(\vartheta)
\Bigr)^{-\audT}\cdot\prod_{j=1}^m\frac{Z_j(1)}{Z_j(\vartheta^j)} .
\]
\end{quotedthm}

\begin{quotedthm}{\cite[Proposition~14.38]{P46}}{ (Auxiliary forms)}
There are $e'\in\Z$, ramified ratios $\rho'_{ij}$ ($0\le i,j<m$) and
$\gamma\colon\{m\text{-fold twists}\}\to\C^\times$ with $\gamma(1)=1$ and
$|\gamma(\vartheta)|=(\auNm Q_\vartheta)^{e'/2}\prod_{i,j}\rho'_{ij}(\vartheta)$, such that $\auT$
is spanned by forms $y$ with: $\gamma(\vartheta)\Theta_\vartheta y$ is integral at every $\iota$
with $\ell(\iota)\nmid\auNm Q_\vartheta$, for every $m$-fold twist $\vartheta$; and
$\|\Theta_1y\|_\sigma^2\le\exp(2x_1^2b^2)$ for all $\sigma\in\auAut(\C)$.
\end{quotedthm}

\begin{quoteddef}{\cite[Definition~14.39]{P46}}{ (Regular families of twists)}
Put $\aubud=\aueu/(mk)$. A family $\theta\colon\auI\to\{m\text{-fold twists}\}$ is \emph{regular}
if:
\begin{enumerate}
\item[(R1)] $\ell(\iota)\nmid\auNm Q_{\theta_\iota}$ for all $\iota$;
\item[(R2)] for every ramified critical pair $(\nu,z)$, $1\le n\le m$ and sign $\pm$, the family
$\iota\mapsto|L(z,\nu^{-1})|/|L(z,(\nu\theta_\iota^{\pm n})^{-1})|$ has height at most $x_1$;
\item[(R3)] for $1\le j\le m$ there are $\beta_\ell\ge0$ with $\sum_{\ell\in T}\beta_\ell\le
\aubud$ for all finite $T$ such that, for every $\iota$ with prime $\ell$: (i) the twist
$\theta_\iota^{\,j}$ \emph{has denominators at most $\beta_\ell$ at $\iota$ for the $j$-th lift}
in the sense of {\cite[Definition~2.8]{P54}} (the false lifts being those of
{\cite[Definition~2.6]{P54}}: the forms in the image of $e_0$ that are orthogonal to the lifts and
are killed by the Hecke polynomial $P_w(X_w)$ of {\cite[Lemma~2.2(L3)]{P54}} at one split place $w$
off the level; a form is orthogonal to the false lifts if it is so for every such $w$); and (ii)
for every main pair $p$ for $j$ ({\cite[Definition~2.19]{P54}}) and all test forms $\Xi,\Xi'$ at
$\ell$ for $(p,\theta_\iota^{\,j})$ ({\cite[Definition~2.22]{P54}}), the inner product of the orthogonal
projections of $\Xi$ and $\Xi'$ to the line of the normalised lift
$\varphi_{p,\theta_\iota^{\,j}}$ has absolute value at most $e^{\beta_\ell}$ at $\iota$.
\end{enumerate}
\end{quoteddef}

\begin{quotedthm}{\cite[Proposition~14.40]{P46}}{ (Regular families exist)}
For all admissible data there is a regular family of $m$-fold twists.
\end{quotedthm}

\begin{quotedthm}{\cite[Proposition~14.41]{P46}}{ (One lift and Hecke values)}
Let $1\le j\le m$ and $d_j=\dim e_0\mathcal A(W_j)$. There are a constant $C_j\ne0$, for every twist
$\vartheta$ a number $\Delta_j(\vartheta)\ne0$ (the determinant of a \emph{comparison operator}
compressed to the space of lifts) and the positive number $H_j(\vartheta)$ of {\cite[(11)]{P54}}: for the three critical pairs $(\nu_i,z_i)$ of {\cite[\S2.7]{P54}} it is the product of $(\auNm\mathsf S_\vartheta)^{z_i^2}\prod_{1\le z\le z_i}|L^{\mathsf S_\vartheta}(z,(\nu_i\vartheta^{\pm1})^{-1})|^2$, such that
\[
Z_j(\vartheta)\cdot\Delta_j(\vartheta)=C_j\cdot\bigl(H_j(\vartheta)/H_j(1)\bigr)^{d_j},
\]
and
\begin{enumerate}
\item[(a)] $|\sigma\Delta_j(1)|\le\exp(3d_j\cdot f\,m\,b\,(\log b)^2)$ for all
$\sigma\in\auAut(\C)$;
\item[(b)] if $\ell(\iota)\nmid\auNm Q_\vartheta$ and $\vartheta$ has denominators at most $\beta$
at $\iota$ for the $j$-th lift, then $|\Delta_j(\vartheta)|_\iota\le\exp(d_j\,r\,\beta)$.
\end{enumerate}
\end{quotedthm}

\begin{quotedthm}{\cite[Proposition~14.42]{P46}}{ (Heights of twisted Hecke values)}
For a regular family $\theta$ and $1\le j\le m$ the family $\iota\mapsto H_j(1)/H_j(\theta_\iota^
{\,j})$ has height at most $10\,m\,x_1+k\,(6\,\aubud-\auhdir)$.
\end{quotedthm}

\begin{quotedthm}{\cite[Proposition~14.43]{P46}}{ (The cost of one lift)}
For a regular family $\theta$ and $1\le j\le m$ the family $\iota\mapsto Z_j(1)/Z_j(\theta_\iota^
{\,j})$ has height at most $\audT\bigl(f\,a\,(\log b)^4/m-k\,\auhdir\bigr)$.
\end{quotedthm}

\begin{quoteddef}{\cite[Definition~2.1]{P49}}{ (The Weil representation)}
Let $W$ be a non-degenerate hermitian space over $E$. For each place $v$ of $k$ a homomorphism
$\omega_v=\omega_{W,v}$ from $\auU(W)(k_v)$ to the unitary operators of a Heisenberg representation
$\psHS_v(W)$, with $\omega_v(g)\psrho(w)\omega_v(g)^{-1}=\psrho(gw)$, is fixed as follows.
\begin{itemize}
\item[(D1)] $E_v$ a field and $W_v$ split: $\omega_v=\omega_{Y,\upsilon_v}$ of (P3) for a Lagrangian
$Y\subset W_v$ (by {\cite[Lemma~2.2]{P49}} it does not depend on $Y$).
\item[(D2)] $E_v$ a field and $W_v$ not split: the splitting (P4) with $\xi=\upsilon_v$, characterised
intrinsically as follows. The space $W_v\oplus W_v^-$ is split, and $\omega_{W\oplus W^-,v}(g,1)$ is
defined by (D1), independently of the Lagrangian ({\cite[Lemma~2.2]{P49}}). Since $(g,1)$ fixes $W^-$
pointwise, this operator commutes with $1\otimes\psrho(w')$ for $w'\in W^-_v$ on
$\psHS_v(W\oplus W^-)=\psHS_v(W)\widehat\otimes\psHS_v(W^-)$; as $\psrho|_{W^-}$ is irreducible (P1),
the commutant of $1\otimes\psrho(W^-_v)$ is the set of operators $A\otimes1$, so
$\omega_{W\oplus W^-,v}(g,1)=A_g\otimes1$ for a unique bounded $A_g$, and we set $\omega_v(g):=A_g$.
As $A\mapsto A\otimes1$ is an injective, multiplicative, $*$-preserving map, $g\mapsto A_g$ is a
unitary representation, and $A_g\psrho(w)A_g^{-1}=\psrho(gw)$ follows from the same identity for
$(g,1)$ on $w\in W_v$. This is the operator of (P4) \cite[(4.1)--(4.5), Proposition~4.1]{KudlaSplitting},
and it involves no auxiliary choice.
\item[(D3)] $v=w\aubar w$ split in $E$: $W_v=W_w\oplus W_{\aubar w}$ with both summands
$\pwfA$-isotropic, $\auU(W)(k_v)=\auGL(W_w)$ (restriction to $W_w$), and on the model (P2)
$\auS(W_w)$,
\[
(\omega_v(g)\Psi)(x)=\upsilon_w(\det g_w)^{-1}|\det g_w|_w^{-1/2}\,\Psi(g_w^{-1}x).
\]
\end{itemize}
The global representation is $\omega=\bigotimes_v\omega_v$ on $\psHS(W)=\bigotimes'_v\psHS_v(W)$, as
in {\cite[\S5.1.1]{P49}} and (W4).
\end{quoteddef}

\begin{quotedthm}{\cite[Lemma~2.3]{P49}}{ (Distributions supported at $0$)}
Let $v$ be a place, $W_v=X\oplus Y\oplus W'$ with $X,Y$ $\pwfA$-isotropic, $X\oplus Y$ orthogonal
to $W'$ for $\pwfA$, and $(\psHS',\psrho')$ a Heisenberg representation of $W'$ (or $\psHS'=\C$ if
$W'=0$). Realise $\psHS_v(W)$ as $\auS(X)\widehat\otimes\psHS'$, with smooth vectors
$\auS(X;\psHS'^\infty)$, where $X\oplus Y$ acts by \textup{(P2)} on the first factor and $W'$ by
$\psrho'$ on the second. If $T\colon\auS(X;\psHS'^\infty)\to\psHS'^\infty$ is continuous linear,
$T\circ\psrho(y)=T$ for $y\in Y$ and $T\circ\psrho(w')=\psrho'(w')\circ T$ for $w'\in W'$, then
$T=c\,\pwfev$ for a constant $c$, where $\pwfev F=F(0)$.
\end{quotedthm}

\begin{quotedthm}{\cite[Lemma~2.4]{P49}}{ (The Siegel parabolic)}
At every place $v$, $\omega_v(m(a))=\upsilon_v(\det a)^{-1}\pwfL_a$ and $\omega_v(n(b))=\pwfM_b$.
\end{quotedthm}

\begin{quotedthm}{\cite[Lemma~2.5]{P49}}{ (Orthogonal sums and doubling)}
Let $W=W_1\oplus W_2$ be an orthogonal sum of non-degenerate hermitian spaces over $E_v$, and
identify $\psHS_v(W)=\psHS_v(W_1)\widehat\otimes\psHS_v(W_2)$ by an intertwiner (P1). Then
$\omega_{W,v}(g_1,g_2)=\omega_{W_1,v}(g_1)\otimes\omega_{W_2,v}(g_2)$ for $g_i\in\auU(W_i)(k_v)$. In
particular $\omega_{W\oplus W^-,v}(g,1)=\omega_{W,v}(g)\otimes1$.
\end{quotedthm}

\begin{quotedthm}{\cite[Proposition~2.17]{P49}}{ ((W5))}
Let $v=w\aubar w$ be split in $E$. Then $\psHS_v(\auHom(V,V'))$ is realised on
$\auS(\auHom(V_w,V'_w))$, with $\psrho(a_w,0)\Psi(x)=\Psi(x+a_w)$ and $\psrho(0,a_{\aubar w})\Psi(x)=
[x,a_{\aubar w}]\Psi(x)$, and
\[
(\omega_{\auHom(V,V')}(g,h)\Psi)(x)=\upsilon_w(\det h_w)^{-n}\upsilon_w(\det g_w)^{n'}\,
|\det h_w|_w^{-n/2}|\det g_w|_w^{n'/2}\,\Psi(h_w^{-1}xg_w).
\]
If $n,n'$ are even, then $(\omega_\mu(g,h)\Psi)(x)=\chi_w(\det h_w)\chi_w(\det g_w)^{-1}|\det
h_w|_w^{-n/2}|\det g_w|_w^{n'/2}\Psi(h_w^{-1}xg_w)$, where $\chi=\mu\circ(z\mapsto z/\aubar z)$; this
is {\cite[(62)]{P49}}.
\end{quotedthm}

\begin{quotedthm}{\cite[Theorem~3.2]{P49}}{ ((A1))}
Assume that every $\Lambda^{(j)}_\tau$ satisfies \textup{(Dec$_c$)} with $c=B+N_j^2+1$, $B=x^{25}$.
Then for $0\le j\le m$ the function $c_j=\bigotimes_\tau c_{j,\tau}$ is smooth, left and right
$K_{j,\infty}$-finite, an eigenfunction of the centre of the enveloping algebra, and
$\int_{G_{j,\infty}}|\mathsf Xc_j(x)|\,\|x\|^B\,dx<\infty$ for every left- or right-invariant
differential operator $\mathsf X$.
\end{quotedthm}

\begin{quotedthm}{\cite[Corollary~3.3]{P49}}{ ((A2))}
Assume \textup{(Dec$_c$)} with $c=5N_j^3$ for every $\Lambda^{(j)}_\tau$. Let $\mathsf P=\prod_\tau
\mathsf P_\tau$ be a parabolic subgroup of $G_{j,\infty}$ with $\mathsf P_\tau$ proper for at least
one $\tau$, and $\mathsf N$ its unipotent radical. Then $\int_{\mathsf N}|c_j(anb)|dn<\infty$ and
$\int_{\mathsf N}c_j(anb)\,dn=0$ for all $a,b$.
\end{quotedthm}

\begin{quotedthm}{\cite[Lemma~3.4]{P49}}{}
Let $X$ be the Harish-Chandra module of an irreducible unitary representation of $G=\auU(p,q)$ with infinitesimal character $\lambda$ \textup{(}regular, $\lambda\in(\frac{n-1}2+\Z)^n$\textup{)} that contains the $K$-type $\Lambda^K(\lambda)$. Then every $K$-type of $X$ has highest weight $\Lambda^K(\lambda)+\sum n_\beta\beta$ with $n_\beta\in\Z_{\ge0}$ and $\beta$ running over noncompact roots in $\Psi$; in particular $X$ contains no $K$-type $\Lambda^K(\lambda)-A$ as in \textup{(D3)(a)}, and $X$ is the Harish-Chandra module of $\pi_\lambda$.
\end{quotedthm}

\begin{quotedthm}{\cite[Lemma~3.5]{P49}}{ (Rigidity)}
Let $\sigma$ be a unitary representation of $G$ on a separable Hilbert space and $y$ a smooth
vector with: $\sigma(t)y=t^{\Lambda^K(\lambda)}y$ for $t\in T$; $d\sigma(E_{ab})y=0$ for
$\varepsilon_a=\varepsilon_b$, $\lambda_a>\lambda_b$; $d\sigma(\Omega)y=\kappa_\lambda y$. Put
$e_\lambda=d_\lambda\overline{c_\lambda}$. Then $\langle y,\sigma(g)y\rangle=\|y\|^2c_\lambda(g)$ and
$\sigma(e_\lambda)y=y$, provided $e_\lambda\in L^1(G)$.
\end{quotedthm}

\begin{quotedthm}{\cite[Lemma~3.6]{P49}}{ (The derived action)}
After conjugation by an element of $T\times T'$, the derived actions of $\omega_\infty(\cdot,1)$ and
$\omega_\infty(1,\cdot)$ at $\tau$ on polynomial vectors are {\cite[(12)]{P49}}, with
\[
\gamma=\tau_0+\tfrac{n'}2,\quad\gamma'=\tau'_0+\tfrac n2,\qquad
\tau_0=\kappa_{j+1}(\tau)+\tfrac12[j\text{ odd}],\quad\tau'_0=-\kappa_{j+1}(\tau)-\tfrac12[j\text{ even}],
\]
provided the following holds.
\begin{enumerate}
\item[(V)] At $\tau$, the operator $\omega_{\auHom(Y_j,Y_{j+1})}(g,h)$ of {\cite[\S5.1.1]{P49}} is the
genuine oscillator action of the metaplectic covers (Weyl-ordered quadratic operators) times
$\det(g)^{n'/2}\det(h)^{-n/2}$, where $\upsilon_\tau$ restricted to $\auU(1)$ is $z\mapsto z$.
\end{enumerate}
\end{quotedthm}

\begin{quotedthm}{\cite[Lemma~3.7]{P49}}{}
Under {\cite[Lemma~3.6]{P49}}, $(-\Lambda^{(j)}_\tau,\Lambda^{(j+1)}_\tau)$ match.
\end{quotedthm}

\begin{quoteddef}{\cite[Definition~3.9]{P49}}{}
A \emph{joint harmonic} is a polynomial $\Phi\ne0$ with $E_{aa}\Phi=\Lambda_a\Phi$,
$E'_{cc}\Phi=\Lambda'_c\Phi$, $E_{ab}\Phi=0$ ($\varepsilon_a=\varepsilon_b$, $\lambda_a>\lambda_b$),
$E'_{cd}\Phi=0$ ($\varepsilon'_c=\varepsilon'_d$, $\lambda'_c>\lambda'_d$), $\Delta_{ab}\Phi=0$
($\varepsilon_a\ne\varepsilon_b$) and $\Delta'_{cd}\Phi=0$ ($\varepsilon'_c\ne\varepsilon'_d$). A
\emph{pre-joint harmonic} satisfies the same conditions except $E'_{cc}\Phi=\Lambda'_c\Phi$.
\end{quoteddef}

\begin{quotedthm}{\cite[Lemma~3.10]{P49}}{}
\textup{(i)} $w_r\parsgn\xi_r\ge0$ for all $r$. \textup{(ii)} $w_r=-w'_{c(r)}$, and $w'_c=0$ for $c\in B$.
\textup{(iii)} If $\varepsilon_a=\varepsilon_b$ and $\lambda_a>\lambda_b$ then $\Lambda_a\ge\Lambda_b$;
likewise for $\Lambda'$.
\end{quotedthm}

\begin{quotedthm}{\cite[Proposition~3.11]{P49}}{ (Existence)}
For $r\in R$ and an integer $0\le\ell<|w_r|$, call $(r,c(r))$ a pair of \emph{level} $\ell$. For each
$\ell$ and each pair of signs $(\epsilon,\epsilon')$ let $L_{\ell,\epsilon,\epsilon'}$ be the set of
rows $r$ having a pair of level $\ell$ with $\varepsilon_r=\epsilon$, $\varepsilon'_{c(r)}=\epsilon'$, and
$M_{\ell,\epsilon,\epsilon'}=\det(z_{r\,c(r')})_{r,r'\in L_{\ell,\epsilon,\epsilon'}}$. Then
$\Phi=\prod_{\ell,\epsilon,\epsilon'}M_{\ell,\epsilon,\epsilon'}$ is a joint harmonic. Moreover no monomial
of $\Phi$ contains two variables of one column from rows of opposite signs, or of one row from columns
of opposite signs.
\end{quotedthm}

\begin{quotedthm}{\cite[Lemma~3.12]{P49}}{ (Polarisations)}
Let $P$ be killed by the row raising operators and by the $\Delta'_{cd}$. Let $\varepsilon_a=\varepsilon_b$.
\textup{(a)} $\langle D^{-s}_{b\leftarrow a}P,D^{s}_{a\leftarrow b}P\rangle=0$ for $s=\pm1$.
\textup{(b)} If $s\lambda_a>s\lambda_b$ then $D^s_{a\leftarrow b}P=0$.
\textup{(c)} If $E_{aa}P=\Lambda_aP$ and $k=s\,w_a\ge0$, then $D^s_{a\leftarrow a}P=kP$ and
$\partial_{ac}P=0$ whenever $\eta_{ac}=-s$.
\end{quotedthm}

\begin{quotedthm}{\cite[Proposition~3.13]{P49}}{ (Uniqueness)}
Two pre-joint harmonics are proportional. In particular the joint harmonic is unique up to a scalar.
\end{quotedthm}

\begin{quotedthm}{\cite[Lemma~3.14]{P49}}{}
Put $\Omega=\sum E_{ab}E_{ba}$, $\Omega'=\sum E'_{cd}E'_{dc}$, $\mathcal N=\sum_{a,c}a^+_{ac}a^-_{ac}=
\sum_aE_{aa}-n\gamma$ and $s_0=\gamma+\gamma'-\tfrac{n+n'}2$. On polynomials,
$\Omega'=\Omega-2s_0\mathcal N+n^2n'-2\gamma'nn'+n'\gamma'^2-n\gamma^2$. If $(\lambda,\lambda')$ match,
then on every polynomial $\Phi$ with $\sum_aE_{aa}\Phi=(\sum_a\lambda_a)\Phi$ the scalars satisfy
$\kappa_{\lambda'}=\kappa_\lambda-2s_0(\sum_a\lambda_a-n\gamma)+n^2n'-2\gamma'nn'+n'\gamma'^2-n\gamma^2$.
\end{quotedthm}

\begin{quotedthm}{\cite[Theorem~3.15]{P49}}{ ((A3))}
Assume {\cite[Lemma~3.6]{P49}} (that is, (V)), the conclusion of {\cite[Theorem~3.2]{P49}}, and
\textup{(D1)--(D3)}. Let $0\le j<m$. A joint harmonic $\Phi_j=\bigotimes_\tau\Phi_\tau$ exists and is
unique up to a scalar, and
\[
\int_{G_{j,\infty}}e_j(x)\,\omega_\infty(x^{-1},1)\Phi_j\,dx=\int_{G_{j+1,\infty}}e_{j+1}(x)\,\omega_\infty(1,x)
\Phi_j\,dx .
\]
\end{quotedthm}

\begin{quotedthm}{\cite[Theorem~3.17]{P49}}{}
Assume $h\ge n^2+1$. Then $F=c_\lambda$ on $G$; that is, {\cite[(59)]{P51}} holds.
\end{quotedthm}

\begin{quotedthm}{\cite[Lemma~3.18]{P49}}{}
Let $w\in\Z^n$ and let $x$ be invertible.
\begin{enumerate}
\item[(a)] $\pcom_{-w}(x)=\pcom_w(x^{-1})$, and $\overline{\pcom_w(x)}=\pcom_w(x^*)$.
\item[(b)] Let $b$ be invertible, with $\Delta_S(b)=\prod_{a\in S}b_{aa}$ for every $S$. If
$b_{ab}=0$ whenever $w_a<w_b$, then $\pcom_w(xb)=b^w\pcom_w(x)$. If $b_{ab}=0$ whenever
$w_b<w_a$, then $\pcom_w(bx)=b^w\pcom_w(x)$. The condition on the minors holds for diagonal $b$,
for $b=1+sE_{ab}$ \textup($a\ne b$, $s\in\C$\textup), and for $b\in\pcoB\cup\pcoB^-$.
\item[(c)] If $\varepsilon_a=\varepsilon_b$ and $\lambda_a>\lambda_b$, then $\Lambda_a\ge\Lambda_b$.
For all $a,b$: $\mu_a>\mu_b$ if and only if $\lambda_a>\lambda_b$.
\item[(d)] If $Z$ is accretive and $y$ invertible, then $y^*Zy$ and $Z^{-1}$ are accretive, and
every principal minor of $Z$ is nonzero; so $\pcom_w(Z)$ is defined and nonzero.
\end{enumerate}
\end{quotedthm}

\begin{quotedthm}{\cite[Corollary~3.19]{P49}}{}
Let $Y$ be invertible.
\begin{enumerate}
\item[(a)] For $b\in\pcoB$: $\pcom_\Lambda(Yb)=b^\Lambda\pcom_\Lambda(Y)$, $\pcom_\mu(Yb)=
b^\mu\pcom_\mu(Y)$, $\pcom_\Lambda(b^*Y)=\overline{b^\Lambda}\pcom_\Lambda(Y)$ and $\pcom_\mu(b^*Y)
=\overline{b^\mu}\pcom_\mu(Y)$.
\item[(b)] If $\varepsilon_a=\varepsilon_b$ and $\lambda_a>\lambda_b$, then $\pcom_\Lambda(Y(1+sE_{ab}))
=\pcom_\Lambda(Y)=\pcom_\Lambda((1+sE_{ba})Y)$ for $s\in\C$. For $t$ diagonal,
$\pcom_\Lambda(Yt)=t^\Lambda\pcom_\Lambda(Y)=\pcom_\Lambda(tY)$.
\item[(c)] Put $\psi(Y)=\pcom_{-\mu}(Y)^{-1}$ where $\pcom_{-\mu}(Y)\ne0$. If $\lambda_a<\lambda_b$
\textup(any signs\textup), $\psi(Y(1+sE_{ab}))=\psi(Y)$; for $t$ diagonal, $\psi(Yt)=t^\mu\psi(Y)$.
\end{enumerate}
\end{quotedthm}

\begin{quotedthm}{\cite[Lemma~3.20]{P49}}{}
Let $f\colon\pcoH\to\C$ be continuous with $f(yb)=\delta(b)^{-1}f(y)$ for $y\in\pcoH$, $b\in\pcoB$.
Then $\int_Kf(xu)\,du=\int_Kf(u)\,du$ for every $x\in\pcoH$.
\end{quotedthm}

\begin{quotedthm}{\cite[Proposition~3.21]{P49}}{}
For every $x\in\mathrm{GL}_n(\C)$, $a\mapsto I(x;a)$ is constant on $A_x$. Moreover the function
$\mathcal J(p,r)$ defined in the proof is locally constant on the open set $U_x$ defined there; in
particular $\mathcal J(p,r)=\mathcal J(1,1)$ whenever $(p,r)$ is joined to $(1,1)$ by a path in $U_x$.
\end{quotedthm}

\begin{quotedthm}{\cite[Lemma~3.22]{P49}}{}
Let $g_0=k_0z_0\in G$ \textup(polar decomposition\textup) and let $W$ be the open set of
$x\in\mathrm{GL}_n(\C)$ with $k_0^{-1}x$ accretive. Then $\widetilde F(x)=d_\Lambda I(x;k_0)$ is
holomorphic on $W$, and $F=\widetilde F$ on $W\cap G$, a neighbourhood of $g_0$ in $G$. In particular
$F$ is real analytic on $G$, and $F(1)=1$.
\end{quotedthm}

\begin{quotedthm}{\cite[Lemma~3.23]{P49}}{ ($K$-behaviour)}
Let $x\in G$, $a\in A_x$ and $k_1\in K$. Then
\begin{align*}
F(xk_1)&=d_\Lambda\int_K\frac{\pcom_\Lambda(au)\,\pcom_\Lambda(u^*k_1)}{\pcom_{-\mu}(u^{-1}a^{-1}xu)}\,du,
\tag*{{\cite[(16)]{P49}}}\\
F(k_1x)&=d_\Lambda\int_K\frac{\pcom_\Lambda(k_1au)\,\overline{\pcom_\Lambda(u)}}{\pcom_{-\mu}(u^{-1}a^{-1}xu)}
\,du.\tag*{{\cite[(17)]{P49}}}
\end{align*}
Consequently $F(xt)=t^\Lambda F(x)=F(tx)$ for $t\in T$, and if $\varepsilon_a=\varepsilon_b$,
$\lambda_a>\lambda_b$, then $R(E_{ab})F=0$.
\end{quotedthm}

\begin{quotedthm}{\cite[Lemma~3.25]{P49}}{ (Decay)}
For every integer $0\le c\le h$ and $g\in G$, $|F(g)|\le(2/\|g\|)^c$. If $h\ge n^2+1$, then $F\in
L^1(G)\cap L^2(G)\cap L^\infty(G)$.
\end{quotedthm}

\begin{quotedthm}{\cite[Theorem~3.29]{P49}}{}
Assume \textup{(W1$_\infty$)}, \textup{(W2$_\infty$)}, \textup{(M)} for every nondegenerate hermitian
space over $E_\tau$ (in particular for $W$ and, if $W$ is definite, for $W\oplus W'$ as in Step~3),
and, if $W$ is definite, \textup{(S)}. Then \textup{(V$'$)} holds; that is, \textup{(V)} of {\cite[Lemma~3.6]{P49}} holds at
$\tau$, with $\upsilon_\tau|_{\auU(1)}(z)=z$.
\end{quotedthm}

\begin{quotedthm}{\cite[Lemma~5.2]{P49}}{ (Standard facts)}
Let $W$ be a hermitian space over $E$.
\begin{enumerate}
\item[(W1)] $\omega(g)\psrho(w)\omega(g)^{-1}=\psrho(gw)$. There is a functional $\theta=\theta_W$
on $\psHS(W)$, unique up to a scalar, with $\theta\circ\psrho(w)=\theta$ for $w\in W(E)$; it
satisfies $\theta\circ\omega(\gamma)=\theta$ for $\gamma\in\auU(W)(E^+)$.
\item[(W2)] Let $I\subset W$ be an isotropic $E$-subspace, $W_I=I^\perp/I$, and $P_I$ the stabiliser
of $I$ in $\auU(W)$. There is a continuous linear map $\psZ_I\colon\psHS(W)\to\psHS(W_I)$ (the
\emph{zero term}) with $\theta_{W_I}(\psZ_IF)=\int_{I(E)\backslash I(\auA)}\theta(\psrho(i)F)\,di$,
$\psZ_I\psrho(w)=\psrho(\bar w)\psZ_I$ for $w\in I^\perp(\auA)$ with image $\bar w$, and
$\psZ_I\omega(p)=\chi_I(p)\,\omega_{W_I}(\bar p)\psZ_I$ for $p\in P_I(\auA)$, where $\chi_I(p)=
(\upsilon|\cdot|^{-1/2})(\det p|_I)^{-1}=\upsilon(\det p|_I)^{-1}|\det p|_I|^{1/2}$, with
$\det p|_I\in\auA_E^\times$ the determinant on the $\auA_E$-module $I(\auA)$ and $|\cdot|$ the
idele norm of $E$; so $\chi_I=1$ on unipotent elements and on $P_I(E^+)$.
For isotropic $I\subset I'$ one has $\psZ_{I'}=\psZ_{I'/I}\psZ_I$, and $\psZ_I=\psZ_{I,\infty}
\otimes\psZ_{I,f}$.
\item[(W3)] For $F\in\psHS(W)$, $\theta(F)=\sum_{w\in W(E)/I^\perp(E)}\theta_{W_I}(\psZ_I\psrho(w)F)$,
the series converging absolutely. If $W=\auHom(V,V')$, $F=\Psi\otimes f$ with $\Psi$ smooth, and
$\mathcal C$ is a compact set of finite adelic points of $\auU(V)\times\auU(V')$, there are $C,A$ with
$\sum_w|\theta_{W_I}(\psZ_I\psrho(w)\omega(g,h)F)|\le C(\|g_\infty\|\|h_\infty\|)^A$ whenever
$(g_f,h_f)\in\mathcal C$.
\item[(W4)] Let $L\subset W(\auA_f)$ be a lattice that is its own dual for $[\,,]$. The vectors of
$\psHS_f(W)$ on which $\psrho(l)$, $l\in L$, acts by a fixed character form a line. If $\psvac_L$
spans the line of the character $\epsilon_L$ with $\epsilon_L(l)\epsilon_L(l')=\beta(l,l')
\epsilon_L(l+l')$ fixed below, and $\psrho^\circ(x)=\epsilon(x)\psrho(x)$ for a fixed function
$\epsilon$ with values in roots of unity extending $\epsilon_L^{-1}$, then
$\psrho^\circ(l)\psvac_L=\psvac_L$, $\psrho^\circ(x+l)\psvac_L=\zeta\psrho^\circ(x)\psvac_L$ with a
root of unity $\zeta$, the $\psrho^\circ(x)\psvac_L$, $x\in W(\auA_f)/L$, form a basis of
$\psHS_f(W)$, and $\omega(g)\psvac_L\in\C^\times\psvac_{gL}$ for $g\in\auU(W)(\auA_f)$. If
$g\in\auU(W)(\auA_f)$ stabilises $L$ then $\omega(g)\psvac_L=\psvac_L$.
\item[(W5)] At a place $v=w\aubar w$ of $E^+$ split in $E$, $\psHS_v(\auHom(V,V'))$ is the space
of Schwartz--Bruhat functions on $\auHom(V_w,V'_w)$, with $(\omega_\mu(g,h)\Psi)(x)=\mu_w(\det h_w)
\mu_w(\det g_w)^{-1}\upsilon'_w(g_w,h_w)\,\Psi(h_w^{-1}xg_w)$, where, for even ranks $n,n'$, $\mu_w\upsilon'_w$ is given by {\cite[(62)]{P49}}: the factor is
$\chi_w(\det h_w)|\det h_w|_w^{-n/2}\chi_w(\det g_w)^{-1}|\det g_w|_w^{n'/2}$ with $\chi$ the base
change of $\mu$; for arbitrary ranks the untwisted operator is
\[\Psi\mapsto\upsilon_w(\det h_w)^{-n}\upsilon_w(\det g_w)^{n'}|\det h_w|_w^{-n/2}|\det g_w|_w^{n'/2}\Psi(h_w^{-1}xg_w)\]
({\cite[Proposition~2.17]{P49}}), so that all factors multiplying $\Psi(h_w^{-1}xg_w)$ are built from $\mu_w$, $\upsilon_w$ and $|\cdot|_w$ evaluated at $\det g_w,\det h_w$; $\psrho(a_w,0)\Psi(x)=\Psi(x+a_w)$
and $\psrho(0,a_{\bar w})$ is multiplication by the character $x\mapsto[x,a_{\bar w}]$.
\item[(W0), (W6)] Two further facts are used; they are stated in full in {\cite[\S5.2.1]{P49}}.
(W0): $\theta$ is defined on the smooth vectors of the Fock space tensored with $\psHS_f(W)$ and is
continuous in the first factor. (W6): for hermitian spaces $V,V'$ of ranks $n,n'$ the map
$w\mapsto w^\dagger$ identifies $\psHS(\auHom(V,V'))$ with $\psHS(\auHom(V',V^-))$, carrying
$\omega_{\auHom(V,V')}(g,g')$ to $\upsilon(\det g)^{n'}\upsilon(\det g')^{-n}\,
\omega_{\auHom(V',V^-)}(g',g)$. It carries the self-dual lattice $\alpha\auHom(\psM,\psM')$ to
$\aubar\alpha\auHom(\psM',\psM)$, which differs from $\alpha\auHom(\psM',\psM)$ at the places above $q_1$
({\cite[Remark~7.18]{P49}}).
\end{enumerate}
\end{quotedthm}

\begin{quotedthm}{\cite[Lemma~5.4]{P49}}{}
Let $\Lambda$ be a lattice in a euclidean space of dimension $m\ge1$ whose dual lattice has no
non-zero vector shorter than $\rho>0$. Then every point of the space is within $m^2/\rho$ of
$\Lambda$, and the vectors of $\Lambda$ of length at most $3m^2/\rho$ span the space.
\end{quotedthm}

\begin{quotedthm}{\cite[Lemma~5.5]{P49}}{}
Let $k\in K_j$, $k'\in K_{j+1}$ have reductions $\bar k,\bar k'$ at the level places with
$\bar k'\iota_j=\iota_j\bar k$. Then $\omega_{\mu_{j+1}}(k,k')\Upsilon_{j,f}=\prod_w\psi_{j+1,w}
(\det\bar k'_w)\psi_{j+1,w}(\det\bar k_w)^{-1}\cdot\Upsilon_{j,f}$, where $\psi_{j+1,w}$ is the
tame part of $\mu_{j+1}$ at $w$. In particular $J_j\times J_{j+1}$ fixes $\Upsilon_{j,f}$.
\end{quotedthm}

\begin{quotedthm}{\cite[Lemma~5.7]{P49}}{ (Conjugates of jets)}
Let $X,X'$ be standard of ranks $n<n'$, $W=\auHom(X,X')$, and $\sigma\in\auAut(\C)$ with
$\sigma|_K=1$. There are $t_\sigma\in\C$ with $|t_\sigma|^2=\Omega^{-1}$, a permutation
$\varpi_\sigma$ of $\auPl$, and a $\sigma$-semilinear map $\Lambda_\sigma$ of $\psHS_f(W)$ such
that for all $f\in\psHS_f(W)$:
\begin{enumerate}
\item[(i)] the Taylor series of $\Theta_{\Lambda_\sigma f}(t_\sigma\,\cdot)$ is obtained from that
of $\Theta_f(s_0\,\cdot)$ by applying $\sigma$ to the coefficients and $\varpi_\sigma$ to the
variables $z_\tau$;
\item[(ii)] if $f=\psrho^\circ(x)\omega(g)\psvac^\circ$ with $x\in W(\auA_f)$, $g\in\auU(W)
(\auA_f)$, then $|\theta(\psrho_\infty(y)(1\otimes\Lambda_\sigma f))|\le\exp\bigl(C_0(q_1nn'+1)^2
\bigr)$ for all $y$, where $C_0$ depends only on $q_0$ and $d$.
\end{enumerate}
\end{quotedthm}

\begin{quotedthm}{\cite[Lemma~5.8]{P49}}{}
Let $L\ge1$, $u\in\R$, $R\ge0$, let $\rho_i\ge0$ and $l_i\in\N$ ($i$ in a finite set) with
$\sum_i\rho_i\le e^u(2L+2)R$. Then $\prod_i\rho_i^{\,l_i}\le\exp\bigl(\sum_i(l_i/L+\log l_i!+l_iu)
+2LR\bigr)$.
\end{quotedthm}

\begin{quotedthm}{\cite[Lemma~5.14]{P49}}{ (Cauchy decomposition)}
Let $\mathcal P$ be the space of polynomials in $z$. As a module over $\prod_\tau\auU(n)$ acting by
$\omega_{\mu_{j+1},\infty}(g,1)$, $\mathcal P$ is the direct sum of the spaces $\mathcal P_\nu$
spanned by the polynomials $A\cdot S_\nu(z)\xi$, where $\nu=(\nu_\tau)$ runs over the families of
diagrams with at most $n$ rows and $S_\nu$ is the Schur functor of $\nu$; the modules
$\mathcal P_\nu$ are isotypic of pairwise different irreducible types, and $(A,\xi)\mapsto A\cdot
S_\nu(z)\xi$ is injective on $W^*_{\nu}(\C^{n'})\otimes W_\nu(\C^n)$.
\end{quotedthm}

\begin{quotedthm}{\cite[Lemma~5.22]{P49}}{ (The hook number)}
Let $\langle P,Q\rangle_F=(P(\partial)\aubar Q)(0)=\sum_\alpha\alpha!\,P_\alpha\aubar{Q_\alpha}$ be the
Fischer form on polynomials in the entries of $y$, and $H_\mu=\langle f_\mu,f_\mu\rangle_F$. Then
$H_\mu$ is a positive integer, $H_\mu\le|\mu|^{3|\mu|}$ $(0^0=1)$, the invariant form of $W_\mu$ is
$H_\mu^{-1}\langle\ ,\,\rangle_F$, and $H_\mu\Sigma_\mu\subseteq\Sigma_\mu^\#$.
\end{quotedthm}

\begin{quotedthm}{\cite[Proposition~5.24]{P49}}{ ({\cite[Lemma~5.6(c)]{P49}})}
Let $\lambda=\lambda(X_m)$, $n=n_m$, and let $\lambda^{\circ\flat}_\tau$ be $\lambda^\circ_\tau$
without its first $k-1$ rows. Put $c_{\mathrm{br}}=\prod_\tau H_{\lambda^{\circ\flat}_\tau}$ and
$\mathrm{br}=c_{\mathrm{br}}\,i^*$, with $i$ as below. Then {\rm(a), (b), (c)} of
{\cite[Lemma~5.6]{P49}} hold, and $c_{\mathrm{br}}\le L^{3L}$.
\end{quotedthm}

\begin{quotedthm}{\cite[Lemma~5.26]{P49}}{ (Dictionary)}
The map $f\mapsto\Theta_f$ is a bijection from $\psHS_f(W)$ onto the space of entire functions
$\Theta$ on $\C^{\auPl\times\mathsf C}$ for which there are an integer $\mathsf N\ge1$ and a function
$\epsilon$ on $\mathsf N\Lambda_L$, with values in the roots of unity, with
\[
\Theta(z+\lambda^\natural)=\epsilon(\lambda)\,e^{\langle z,\lambda^\natural\rangle+|\lambda^\natural|^2
/2}\,\Theta(z)\qquad(\lambda\in\mathsf N\Lambda_L),\qquad\langle z,w\rangle=\sum\aubar w z .
\]
For $\gamma\in W(E)$, $\Theta_{\psrho_f(\gamma)f}(z)$ is a root of unity times
$e^{\langle z,\gamma^\natural\rangle-|\gamma^\natural|^2/2}\,\Theta_f(z-\gamma^\natural)$.
\end{quotedthm}

\begin{quotedthm}{\cite[Proposition~5.27]{P49}}{ ({\cite[Lemma~5.7]{P49}})}
{\cite[Lemma~5.7]{P49}} holds, with
$t_\sigma=\sigma(\varpi_{\mathrm{vac}}/\mathsf u)\,c_\psi/\Omega_\sigma$, $\varpi_\sigma$ the
permutation $\tau\mapsto\sigma\circ\tau$ of $\auPl$, and $C_0=C_0(q_0,d)$ given at the end of the proof.
\end{quotedthm}

\begin{quotedthm}{\cite[Proposition~5.30]{P49}}{ ({\cite[Lemma~4.13(K1)]{P49}}, (K2))}
{\rm(K1)} and {\rm(K2)} of {\cite[Lemma~4.13]{P49}} hold.
The integrality of the jets in (K1) holds for every $\ell\ne d$, by (S2) and {\cite[Theorem~7.15]{P49}}.
\end{quotedthm}

\begin{quotedthm}{\cite[Corollary~6.4]{P49}}{}
Let $X$ be a standard space, $G=\auU(X)$, $\ell$ a prime.
{\rm(a)} For $g,d\in G_f$ there is $\gamma\in\Gamma$ such that $\gamma g$ and $d\,\gamma g\,d^{-1}$
are integral above $\ell$. In particular every class of $\Gamma\backslash G_f/J$ has a
representative that is integral above $\ell$.
{\rm(b)} If $\varphi$ is a form of level $J\subseteq K^\circ$ and $(x)$ a system of representatives
that are integral above $\ell$, and if $J$ is the full group $\auGL_n(\auO_w)$ at the places above
$\ell$, then $\varphi$ is integral at $\iota$ if and only if all $\iota(\varphi(x))$ have
coordinates in $\pao$.
\end{quotedthm}

\begin{quotedthm}{\cite[Theorem~7.9]{P49}}{ ({\cite[Proposition~5.25{\rm(C3)(ii)}]{P49}})}
Let $\Lambda=\Omega_0\auO_K$ be normalised and $\sigma\in\auAut(\C/K)$. There are $\Omega_\sigma
\in\C^\times$, an ideal $\mathfrak c\ne0$ of $\auO_K$ and a finite idele $\beta_\sigma$ of $K$ such
that, with $\mathfrak a=\mathfrak c\aubar{\mathfrak c}^{-1}$ and $\widehat{\mathfrak a}=\mathfrak
a\widehat\auO_K$:
\begin{enumerate}
\item[(a)] $\Lambda^\sigma=\Omega_\sigma\mathfrak a$, $\beta_\sigma\widehat\auO_K=\widehat{
\mathfrak a}$ (so $\beta_\sigma\aubar{\beta_\sigma}$ is a unit), and $S_\sigma(\mathsf x\Omega_0)
\equiv\mathsf x'\Omega_\sigma$ modulo $\Lambda^\sigma$ whenever $\mathsf x,\mathsf x'\in K$ and
$\mathsf x'-\beta_\sigma\mathsf x\in\widehat{\mathfrak a}$;
\item[(b)] if $\mathsf x\in K$, $n\ge1$, $n\mathsf x\in\auO_K$, and $\mathsf x'\in K$ satisfies
$\mathsf x'-\beta_\sigma\mathsf x\in2n\widehat{\mathfrak a}$, then $\sigma$ maps the Taylor series
at $0$ of $\theta_{\mathsf x\Omega_0}$ to that of $\theta^\sigma_{\mathsf x'\Omega_\sigma}$.
\end{enumerate}
\end{quotedthm}

\begin{quotedthm}{\cite[Lemma~7.11]{P49}}{ (A frame of length prime to $q_1$)}
There is a tight frame of a length $\mathsf h$ prime to $q_1$.
\end{quotedthm}

\begin{quotedthm}{\cite[Lemma~7.13]{P49}}{ (The product and the vacuum have the same character)}
$\Pi^v=\Theta_{F^v}$ for a unique $F^v\in\psHS_f(W)$; $\theta(1\otimes F^v)=1$;
$\psrho_f(l)F^v=\epsilon_L(l)F^v$ for all $l\in\mathsf bL$; and for $\gamma\in W(E)$
\begin{equation*}\tag*{{\cite[(73)]{P49}}}
\Theta_{\psrho_f(\gamma)F^v}(z)=\zeta_\gamma\prod_{i,\mathsf c}\theta_{\mathsf x_{i,\mathsf c}
\Omega_0}\bigl(a_{i,\mathsf c}(z)\bigr)^{\mathsf b},\qquad\mathsf x_{i,\mathsf c}=-\pfT_{E/K}(
\aubar{u_i}\gamma_{\mathsf c}),
\end{equation*}
with a root of unity $\zeta_\gamma$.
\end{quotedthm}

\begin{quotedthm}{\cite[Theorem~7.15]{P49}}{ (Standard data at every prime)}
Let $\ell$ be a prime, $\delta\in\auO_K$, $\delta\ne0$, and $f$ a standard datum of level $\delta$
at $\ell$. There are a resolution of the identity $(u'_j,m_j)_j$ of $E^{\mathsf C}$ with
$\alpha^{-1}u'_j$ integral above $\ell$, and an integral power series $\mathsf H$ in the series
$y(\mu_j(z))$, $\mu_j=(\varpi_{\mathrm{vac}}/\mathsf u)\,u_j'^{\dagger}$, such that
\[
\Theta_f\bigl(\aubar\delta\,s_0\,z\bigr)=\Bigl(\prod_j\theta(\mu_j(z))^{m_j}\Bigr)^{\auNm\delta}
\cdot\mathsf H(z),
\]
and the linear forms $\mu_j$ have integral coefficients. If $\ell\ne d$, the series
$\Theta_f(\aubar\delta s_0\,\cdot)$ is $\partial$-integral, and so is $\Theta_f(c\,\aubar\delta s_0
\,\cdot)$ for every $c\in\pao$.
\end{quotedthm}

\begin{quotedthm}{\cite[Proposition~7.17]{P49}}{ ({\rm(W6)})}
Let $W$ be a hermitian space over $E$ that has a self-dual lattice $L\subset W(\auA_f)$, and let
$W'$ be the $E$-vector space $\aubar W$ $(W$ with the conjugate scalar multiplication$)$ with the
form $\langle w_1,w_2\rangle'=-\langle w_2,w_1\rangle$. Then $\beta'=\beta$ on $W(\auA)=W'(\auA)$,
$\auU(W)=\auU(W')$ as groups acting on this set, and if $\mathrm{Fl}\colon\psHS(W)\to\psHS(W')$ is
an isomorphism with $\mathrm{Fl}\,\psrho(w)=\psrho'(w)\,\mathrm{Fl}$ (it exists and is unique up
to a scalar), then
\[
\mathrm{Fl}\circ\omega_W(g)=\upsilon(\det\nolimits_Wg)^{-1}\cdot\omega_{W'}(g)\circ\mathrm{Fl}
\qquad(g\in\auU(W)(\auA)),
\]
where $\det_W$ is the determinant on the $E$-module $W$, and $\theta_{W'}\circ\mathrm{Fl}$ is a
multiple of $\theta_W$. For $W=\auHom(V,V')$ the map $w\mapsto w^\dagger$ identifies $W'$ with
$\auHom(V',V^-)$, and $\det_W(g,g')=\det(g')^{n}\det(g)^{-n'}$; this is {\rm(W6)} of
{\cite[\S5.2.1]{P49}}.
\end{quotedthm}

\begin{quotedthm}{\cite[Proposition~7.20]{P49}}{}
Let $X=E^n$ be a standard space with $n\le n_m$, $G=\auU(X)$, $\Gamma=G(E^+)$, $K^\circ$ the
stabiliser of $\widehat\auO_E^{\,n}$ in $G_f$, and for an integer $\mathsf G\ge1$ let $K(\mathsf G)$
be the group of the $k\in K^\circ$ with $k\equiv1$ modulo $\mathsf G$. Then
\[
\#\bigl(\Gamma\backslash G_f/K^\circ\bigr)\le(2\mathsf B+1)^{4fn^2},\qquad\mathsf B=n^2C_{\mathrm{sb}}^2
\cdot2f\sqrt d,\quad C_{\mathrm{sb}}=6q_1d(fn)^2,
\]
and $\#(\Gamma\backslash G_f/J)\le(2\mathsf B+1)^{4fn^2}\mathsf G^{2fn^2}$ for every subgroup
$J\supseteq K(\mathsf G)$ of $G_f$. For $\mathsf G=MN$ this is at most $\exp(x^{30})$. In
particular the levels $J_{X_j}$ $(j\le m)$, and the level $J_1$ of the target forms for the trivial
twist, have at most $\exp(x^{30})$ classes.
\end{quotedthm}

\begin{quoteddef}{\cite[Definition~8.1]{P49}}{}
Let $e=e_j$ or $\overline{e_j}$. For a function $F$ on $G_j(\auA)$ put $(R(e)F)(g)=
\int_{G_{j,\infty}}e(x)F(gx)\,dx$ when the integral converges absolutely. A \emph{form of type $e$}
on $G_j$ is a bounded continuous function $y$ on $G_j(\auA)$, left invariant under $G_j(E^+)$, right
invariant under an open subgroup of $G_j(\auA_f)$ (a \emph{level}), with $R(e)y=y$. The
\emph{forms of the archimedean type of the chain} on $G_j$ are the forms of type $e_j$; their
complex conjugates are the forms of type $\overline{e_j}$. The space $\auAz$ of
{\cite[\S14.3]{P46}} is the space of the forms of type $e_0$ and level $J_M$; this is the precise
meaning of ``the fixed vector of the discrete series with parameter $\varsigma_\tau$''. For $j=m$
the forms of type $e_m$ are the functions $g\mapsto l\bigl(g_\infty^{-1}\varphi(g_f)\bigr)$ with
$\varphi$ an algebraic modular form of weight $\lambda^\circ$ and $l$ the linear form given by
$v_m$ (Peter--Weyl); in this way the last lift has values in algebraic modular forms.
\end{quoteddef}

\begin{quotedthm}{\cite[Lemma~8.2]{P49}}{ (Growth)}
Let $s=D+1$.
\begin{enumerate}
\item[(a)] For $g\in G_{j,\infty}$, $h\in G_{j+1,\infty}$ and $\Psi\in\pzF^{\mathrm{sm}}$,
$\|\omega_\infty(g,h)\Psi\|_s\le(2D\|g\|\|h\|)^{2s}\|\Psi\|_s$. Hence $\omega_\infty$ extends to a
strongly continuous representation on $\pzF_s$ with the same bound.
\item[(b)] $(1+|z|)^{2s}\,|\Psi(z)|\,e^{-|z|^2/2}\le(D+1)^{2s}\|\Psi\|_s$ for $\Psi\in\pzF_s$ and
$z\in\C^D$.
\item[(c)] For every $f\in\psHS_f(W)$ there is $C_f$ with $|\theta(\Psi\otimes f)|\le C_f\|\Psi\|_s$
for $\Psi\in\pzF^{\mathrm{sm}}$; so $\theta(\,\cdot\otimes f)$ extends continuously to $\pzF_s$.
\end{enumerate}
\end{quotedthm}

\begin{quotedthm}{\cite[Lemma~8.3]{P49}}{ (Absolute convergence along an isotropic subspace)}
Let $W$ be a hermitian space over $E$, $2D=\dim_\R W_\infty$, $s=D+1$, with $\|\cdot\|_{s'}$ as in the
preceding paragraph. Let $I\subseteq W$ be an isotropic $E$-subspace and, for a datum $F$ and $w\in
W(E)$, $c_w(F)=\theta_{W_I}(\psZ_I\psrho(w)F)$ \textup{(}notation of \textup{(W2))}; $c_w$ depends only on
the class of $w$ in $W(E)/I^\perp(E)$.
\begin{enumerate}
\item[\textup{(a)}] For every $f\in\psHS_f(W)$ there is $C_f$ with $\sum_{w\in W(E)/I^\perp(E)}|c_w(\Psi\otimes
f)|\le C_f\|\Psi\|_{2s}$ for every smooth $\Psi$.
\item[\textup{(b)}] Let $W=\auHom(X,X')$, $f\in\psHS_f(W)$ and $\mathcal C$ a compact subset of
$\auU(X)(\auA_f)\times\auU(X')(\auA_f)$. There is $C$ such that
\[
\sum_{w\in W(E)/I^\perp(E)}\bigl|c_w\bigl(\omega(g,h)(\Psi\otimes f)\bigr)\bigr|\le C\,(2D\|g_\infty\|\|h_\infty\|)^{4s}
\|\Psi\|_{2s}
\]
for all smooth $\Psi$ and all $(g,h)\in\auU(X)(\auA)\times\auU(X')(\auA)$ with $(g_f,h_f)\in\mathcal C$;
the same holds with $\omega$ replaced by $\omega_{\alpha,\beta}$ or by the twisted $\omega_\mu$, which differ
from it by characters of absolute value $1$.
\item[\textup{(c)}] \textup{(Uniformity under archimedean translates.)} For every $f\in\psHS_f(W)$ there is
$C_f$ with
\[
\sum_{w\in W(E)/I^\perp(E)}|c_w(\psrho_\infty(z)\Psi\otimes f)|\le C_f\|\Psi\|_{2s}
\]
for every
$z\in W_\infty$ and every smooth $\Psi$.
\end{enumerate}
\end{quotedthm}

\begin{quotedthm}{\cite[Proposition~8.4]{P49}}{ (The lower chain is well defined)}
Assume {\rm(A1)--(A3)}. Let $0\le j\le m$.
\begin{enumerate}
\item[(a)] Every form $y$ of type $e_j$ or $\overline{e_j}$ on $G_j$ is smooth, $K_{j,\infty}$-finite
and finite under the centre of the enveloping algebra; it is a cusp form; it is rapidly
decreasing on Siegel sets, in the form: for every Siegel set $\mathfrak S_\infty\subset G_{j,\infty}$, finite set $\mathsf c\subset G_j(\auA_f)$, open subgroup $K'$ of $G_j(\auA_f)$ and $N\ge0$ there is $C_N$ with $|y(g)|\le C_N\|g_\infty\|^{-N}$ for $g\in\mathfrak S_\infty\mathsf cK'$; and the forms of a given level form a finite-dimensional space.
\item[(b)] Let $j<m$, $f\in\psHS_f(W)$ and let $y$ be a form of type $e_j$ on $G_j$. Then the
integral $F(h)=\int_{[G_j]}\theta\bigl(\omega_{\mu_{j+1}}(g,h)(\Phi_j\otimes f)\bigr)y(g)\,dg$ converges
absolutely and $F$ is a form of type $e_{j+1}$ on $G_{j+1}$.
\item[(c)] Let $j<m$, $f\in\psHS_f(W)$ and let $z$ be a form of type $\overline{e_{j+1}}$ on
$G_{j+1}$. Then $F^*(g)=\int_{[G_{j+1}]}\theta\bigl(\omega_{\mu_{j+1}}(g,h)(\Phi_j\otimes f)\bigr)z(h)
\,dh$ converges absolutely and $F^*$ is a form of type $\overline{e_j}$ on $G_j$.
\item[(d)] The space $\auAz$ is finite-dimensional and stable under the right translations by
the stabiliser of the lattice at the chosen places, so $\auG$ acts on it; the Poincar\'e map
$\auP$ of {\cite[\S14.3]{P46}} is well defined; and $\iota_\vartheta$ maps $\auAz$ to algebraic
modular forms of weight $\lambda^\circ$ on $\auU(Y_m)$.
\end{enumerate}
In particular the assumption made in {\cite[\S5.1.3]{P49}} holds: the lifts $\Theta^Y_{j,\vartheta}$ and
their transposes $\Theta^{Y,*}_{j,\vartheta}$ preserve the types, and all the forms concerned are
cusp forms and rapidly decreasing.
\end{quotedthm}

\begin{quotedthm}{\cite[Lemma~8.5]{P49}}{ (Level and character on $Y_m$)}
Assume {\rm(A1)--(A3)}. Let $\vartheta$ be a twist, $z\in\auAz$, and let $k\in K_m$ have, at every
level place $w$, a reduction $\bar k_w$ that is block upper triangular for the flag
$\psM_0\subset\psM_1\subset\dots\subset\psM_m$ of the frames, with diagonal blocks
$\bar k_w^{(0)},\dots,\bar k_w^{(m)}$, and $\bar k_w^{(0)}=1$ at the chosen places. Then
$(\iota_\vartheta z)(gk)=\prod_w\prod_{i=1}^m\psi_{i,w}(\det\bar k_w^{(i)})\cdot(\iota_\vartheta z)(g)$,
where $\psi_{i,w}$ is the tame part of $\mu_i$ at $w$. So $\iota_\vartheta z$ has the level
$J_\vartheta$ and the character $\chi_\vartheta$ of {\cite[\S14.3]{P46}}.
\end{quotedthm}

\begin{quotedthm}{\cite[Lemma~8.7]{P49}}{}
\begin{enumerate}
\item[(a)] $C_0(q_0,d)$ is a finite real number that depends only on $q_0$ and $d$: not on the
scales $x,x_1,b$, nor on $s$, $q_1$, $E$ or any other part of admissible data.
\item[(b)] It is the constant $C_0$ of {\cite[Proposition~5.27]{P49}} and of {\cite[Lemma~5.7{\rm(ii)}]{P49}}.
\item[(c)] With the constant $C_P$ of {\cite[(35)]{P46}} in condition {\rm(D2)} of
{\cite[Definition~13.3]{P46}}, which contains $d$ and $C_0(q_0,d)$ as summands (all summands are
non-negative), all admissible data satisfy
$C_0\le10^{-12}\log\log b$ and $d\le10^{-12}\log\log b$.
\end{enumerate}
\end{quotedthm}

\begin{quotedthm}{\cite[Lemma~8.8]{P49}}{}
\begin{enumerate}
\item[(a)] The levels $J_{X_j}$ $(0\le j\le m)$ and $J_1$ contain the principal congruence subgroup
$K(MN)$ of the corresponding standard space, of rank $n_j\le n_m$, respectively $n_m-1$; they are
free.
\item[(b)] $n_m\le40x^6$, $6q_1d(fk)^2\le x^{100}$, and $6q_1d(fn)^2\le x^{102}$ for $n\le n_m$.
\item[(c)] $\log(MN)\le x$.
\end{enumerate}
So the hypotheses of {\cite[Proposition~7.20]{P49}} hold for these levels, and each of them has at
most $\exp(x^{30})$ classes.
\end{quotedthm}

\begin{quotedthm}{\cite[Proposition~4.6]{P51}}{ ({\cite[Proposition~14.36]{P46}}, precise form)}
Assume {\cite[Theorem~4.1]{P51}} and the property \textup{(G)} of {\cite[\S4.1.8]{P51}}. Assume that $f$
is even and that for each of the sign vectors {\cite[(43)]{P51}} and each $\tau$ the entries of
$\varsigma_\tau$ on coordinates of different signs differ by at least $500r^3+10r^5$. (Both hold
for the data: $f=\#(\Z/q_1)^\times$, and these entries differ by at least $\Delta=2r^8\ge500r^3+10r^5$.) Then for
$1\le j\le m$ and every $T\in\mathcal H^M$
\[
\pmtr\bigl(T\mid e_0\pmA(W_j)\bigr)=\pmtr\bigl(T\mid\auT\bigr).
\]
In particular $\dim e_0\pmA(W_j)=\audT$, and every character of $\mathcal H^M$ has the same
multiplicity in $e_0\pmA(W_j)$ and in $\auT$.
\end{quotedthm}

\begin{quotedthm}{\cite[Lemma~4.7]{P51}}{}
Let $e_\tau$ be as in {\cite[(42)]{P51}}, on $G_\tau\cong G_s$.
\begin{enumerate}
\item[(a)] $e_\tau$ is continuous, $e_\tau^*=e_\tau$, $|e_\tau(g)|\le C\|g\|^{-B}$,
$e_\tau\in L^1$, and $e_\tau*e_\tau=e_\tau$.
\item[(b)] There is $\alpha_\tau\in C_c^\infty(G_\tau)$ with $\alpha_\tau*e_\tau=e_\tau$. Hence
$e_\tau=e_\tau*\alpha_\tau^*$, $e_\tau$ is smooth, and $|(De_\tau)(g)|\le C_D\|g\|^{-B}$ for every
left-invariant differential operator $D$.
\item[(c)] $\int_{G_\tau}\|g\|^{-B'}dg<\infty$ for $B'>n^2$.
\end{enumerate}
The same holds for $e$ on $G_\infty$ with $\alpha=\bigotimes\alpha_\tau$, with $\pmP(g)^{-B}$ in place
of $\|g\|^{-B}$.
\end{quotedthm}

\begin{quotedthm}{\cite[Lemma~4.8]{P51}}{ ($e$ is a cusp form)}
Let $N'$ be the unipotent radical of a proper parabolic subgroup of $G_\tau$. Then
$\int_{N'}|e_\tau(x\nu y)|\,d\nu<\infty$ and $\int_{N'}e_\tau(x\nu y)\,d\nu=0$ for all $x,y\in G_\tau$.
\end{quotedthm}

\begin{quotedthm}{\cite[Lemma~4.9]{P51}}{}
Let $x=\omega ak\in\pmS$, $P=P_x$, $N=N_x$.
\begin{enumerate}
\item[(a)] Let $\Omega_f\subset G_f$ be compact. There is $c>0$ such that
$\pmP\bigl((x^{-1}\gamma x)_\infty\bigr)\ge c\,\beta(a)^{f}$ for every $\gamma\in\pmGa\smallsetminus
P(F)$ with $(x^{-1}\gamma x)_f\in\Omega_f$.
\item[(b)] Let $K'\subset G_f$ be compact open. There is $C$ such that $N(\auA_F)\,x\subset
N(F)\,x\,(Z_\delta\times K')$ with $\delta=C\beta(a)^{-1}$, where $Z_\delta=\{z\in G_\infty:
\|z_\tau-1\|\le\delta\text{ for all }\tau\}$.
\item[(c)] Let $B'>n^2$ and $M_{B'}(y)=\pmP(y_\infty)^{-B'}1_{\Omega_f}(y_f)$. There is $C$ such that
$\sum_{\gamma\in\pmGa}M_{B'}(g^{-1}\gamma x)\le C\|a\|^{2fn^2}$ for all $g\in G(\auA_F)$. For $x$ in a compact subset of $G(\auA_F)$, not necessarily in $\pmS$, the series converges
uniformly in $(g,x)$.
\end{enumerate}
\end{quotedthm}

\begin{quotedthm}{\cite[Lemma~4.10]{P51}}{ (Decay in a cusp)}
For every integer $j\ge0$ there is $\varepsilon_0\in(0,1]$ with the following property. Let $N$
be the unipotent radical of a proper parabolic $F$-subgroup of $G$, $x\in G(\auA_F)$,
$K'\subset G_f$ compact open and $\delta>0$ with $N(\auA_F)x\subset N(F)\,x\,(Z_\delta\times K')$.
Let $\mathsf F$ be a continuous function on $G(\auA_F)$ with $|\mathsf F|\le T$, left invariant
under $N(F)$, right invariant under $K'$, and with $\mathsf F(g)=\int_{G_\infty}\mathsf F(gh)e(h)\,dh$
for all $g$. Then $|\mathsf F(x)|\le(\delta/\varepsilon_0)^j\,T$.
\end{quotedthm}

\begin{quotedthm}{\cite[Proposition~4.11]{P51}}{}
Let $\varphi'\in C_c^\infty(G_f)$ and $\mathsf f'=e\otimes\varphi'$.
\begin{enumerate}
\item[(a)] The series $k_{\mathsf f'}(x,y)=\sum_{\gamma\in\pmGa}\mathsf f'(x^{-1}\gamma y)$ converges absolutely and
uniformly on compact sets, and is continuous on $\pmQ\times\pmQ$.
\item[(b)] $\sup_{x\in G(\auA_F)}\sum_{\gamma\in\pmGa_\pmell}|\mathsf f'(x^{-1}\gamma x)|<\infty$.
\item[(c)] There is $C$ such that $\bigl|\sum_{\gamma\in S}\mathsf f'(x^{-1}\gamma x)\bigr|\le C$ for all
$x\in G(\auA_F)$ and all admissible $S$.
\end{enumerate}
\end{quotedthm}

\begin{quotedthm}{\cite[Lemma~5.1]{P51}}{}
The signs of {\cite[(43)]{P51}} on $Y_0$ are $\varepsilon_i=(-1)^{\max(0,\,i+1-k)}$ for the
coordinate numbered $i=0,\dots,r-1$, and the doubled parameter is twice
$\lambda_i=\frac k2-1-i$ for $i<k-1$ and $\lambda_i=-\kappa(\tau)-(i+2-k)\Delta$ for $i\ge k-1$.
Hence the entries $\frac k2-1,\dots,1-\frac k2$ and the entries $-\kappa(\tau)-j\Delta$ with $j$ odd
sit at coordinates of sign $+$, those with $j$ even at coordinates of sign $-$; there are $k+\varrho$
signs $+$ and $\varrho+1$ signs $-$. On the $k-1$ coordinates of the first group $\Lambda=0$, and on
the other $m=2\varrho+2$ coordinates, in the order of decreasing $\lambda$, the signs are $+,-,+,
\dots,-$.
\end{quotedthm}

\begin{quotedthm}{\cite[Proposition~5.2]{P51}}{}
Let $dg$ be the Haar measure on $G$ compatible with the Haar measure of mass $1$ on $\auU(N)$. Then
$\int_G|c_\tau|^2dg=\pbdeg_\tau^{-1}$ with $\pbdeg_\tau=\prod_{i<j}(\lambda_i-\lambda_j)/(j-i)$
(entries in decreasing order). Thus the measures of \textup{(DS1)(ii)} and \textup{(DS1)(iii)} agree.
\end{quotedthm}

\begin{quotedthm}{\cite[Lemma~5.4]{P51}}{}
For $g\in G$: $a(g)$ is invertible, $|\det a(g)|\ge(\|g\|/2)^2$, and $|\det a(kgk')|=|\det a(g)|$ for
$k,k'\in\pdsK$. If $g$ is positive definite, so is $a(g)=\frac12(g+g^{-1})$.
\end{quotedthm}

\begin{quotedthm}{\cite[Lemma~5.5]{P51}}{}
Let $w\in G$ be positive definite and $0\le c\le h$ an integer. Then $\pdsm_{-\mu}(w)=\pdsm_\mu(w)>0$
and $\det a(w)^c\le\pdsm_\mu(w)^2$.
\end{quotedthm}

\begin{quotedthm}{\cite[Theorem~5.6]{P51}}{ (\textup{(DS2)})}
For every $g\in G$ and every integer $0\le c\le h$ (formula {\cite[(59)]{P51}} holds by {\cite[Theorem~3.17]{P49}}),
\[
|c_\tau(g)|\le|\det a(g)|^{-c/2}\le\bigl(2/\|g\|\bigr)^c .
\]
In particular $|c_\tau(g)|\le(2/\|g\|)^{j+\pbeq}$ for $0\le j\le r^7$, $\pbeq=9r^2$, and for $g\in
G_\tau$ with the norm of {\cite[\S3.2.1]{P51}}.
\end{quotedthm}

\begin{quotedthm}{\cite[Lemma~5.7]{P51}}{}
Put $C_0=3N(L+2N)$. If $g\in G$ and $\|g-1\|\le\delta:=1/(4C_0)$, then $|c_\tau(g)|\ge\frac12$.
\end{quotedthm}

\begin{quotedthm}{\cite[Theorem~5.9]{P51}}{ (\textup{(DS3)})}
For every real $\mathsf s>4N^2$, $\pbdeg_\tau\int_G\|g\|^{-\mathsf s}dg\le4(5/\delta)^{2N^2}\mathsf s/
(\mathsf s-4N^2)$. In particular $c_\tau$ is integrable, and
\[
\pbdeg_\tau\int_{G_\tau}\bigl(32/\|g\|\bigr)^{\pbeq}dg\le8\cdot32^{\pbeq}\bigl(60N(L+2N)\bigr)^{2N^2}
\le b^{99r^2}\qquad(\pbeq=9r^2).
\]
\end{quotedthm}

\begin{quotedthm}{\cite[Lemma~7.10]{P51}}{ (Residue sums)}
Let $F$ be a polynomial of degree at most $2n-1$ and $c\in\R$ not a node. Then
$\sum_iF(\mu_i)/p'(\mu_i)$ is the coefficient of $t^{2n-1}$ in $F$, and $\sum_iF(\mu_i)/((\mu_i-c)
p'(\mu_i))=-F(c)/p(c)$.
\end{quotedthm}

\begin{quotedthm}{\cite[Lemma~7.16]{P51}}{ (Third order)}
Let $\|A_i\|\le\rho\le\frac1{10}$ \textup($i=1,2,3$\textup). Then
\[
\det\bigl(1+Q(e^{A_1}e^{A_2}e^{A_3}-1)\bigr)=\exp\Bigl(\pbptr\bigl(Q(A_1+A_2+A_3)\bigr)+
\pbpCT_Q(A_1,A_2,A_3)+r\Bigr),\qquad|r|\le50\rho^3 .
\]
\end{quotedthm}

\begin{quotedthm}{\cite[Proposition~7.17]{P51}}{ (Second order with crossing parts)}
Let $\|A_i\|\le\delta\le\frac1{20}$ and $\|[Q,A_i]\|\le\zeta$ \textup($i=1,2,3$\textup). Then the identity
of {\cite[Lemma~7.16]{P51}} holds with $|r|\le400\,\delta\zeta^2$.
\end{quotedthm}

\begin{quotedthm}{\cite[Lemma~7.18]{P51}}{ (Size of one minor)}
Let $k$ be unitary, $S$ a set of coordinates and $P$ the coordinate projection onto $S$. Then
$|\Delta_S(k)|^2\le\exp(-\|[P,k]\|^2/2)$.
\end{quotedthm}

\begin{quotedthm}{\cite[Proposition~7.19]{P51}}{ (One minor near the identity)}
Let $o$ be an injective real function on the coordinates, $S$ a set of coordinates with $o_i<o_j$
for $i\in S$, $j\notin S$, and $P$ the coordinate projection onto $S$. Let $Z\in\pdsg$ with
$\|Z\|\le\delta\le\frac1{20}$, and let $u_-$ \textup(resp.\ $u_+$\textup) keep the entries $Z_{ij}$ with
$\varepsilon_i=\varepsilon_j$ and $o_i<o_j$ \textup(resp.\ $o_i>o_j$\textup) and vanish elsewhere. Put
$\hat g=e^{-u_-}e^Ze^{-u_+}$. Then for every $k\in\pdsK$,
\[
\Delta_S(k^*\hat gk)=\exp\bigl(\pbptr(PZ)+\tfrac14\|[P,Z]\|^2+r\bigr),\qquad|r|\le10^3\delta\bigl(
\|[P,Z]\|^2+\|[P,k]\|^2\bigr),
\]
and $\det\hat g=e^{\pbptr Z}$.
\end{quotedthm}

\begin{quotedthm}{\cite[Lemma~7.20]{P51}}{ (Removing the unipotent factors)}
Let $Z\in\pdsg$ with $\|Z\|\le\frac1{20}$, and let $u_\pm$ and $\hat g$ be as in
{\cite[Proposition~7.19]{P51}} for $o=\lambda$. Then
\[
c(e^Z)=d_\Lambda\int_{\pdsK}\frac{|\pdsm_\Lambda(k)|^2}{\pdsm_{-(\lambda+\rho)}(k^{-1}\hat gk)}\,dk .
\]
\end{quotedthm}

\begin{quotedthm}{\cite[Lemma~7.21]{P51}}{}
$d_\Lambda\le e^{2T}$.
\end{quotedthm}

\begin{quotedthm}{\cite[Proposition~7.22]{P51}}{ (Gaussian shape, noncompact factor)}
Let $g\in G$ with $\|g-1\|\le\delta_1$, $\delta_1T\le5\cdot10^{-13}$, and $Z=\log g$. There is
$\gamma\in[0,1]$ with $|c(g)e^{-\pbptr(\mu Z)}-\gamma|\le\gamma/10+e^{-100T}$.
\end{quotedthm}

\begin{quotedthm}{\cite[Lemma~7.25]{P51}}{ (Rows and columns)}
For $s\in R_\tau$: $\|x_s\|^2=2\|y_s\|^2-2\Rel\pbptr x_s$, $\|x'_s\|^2=-2\Rel\pbptr x'_s$,
$\|P_{\mathrm{off}}x_s\|^2\le6\|y_s\|^2$ and $\|x_sP_{\mathrm{off}}\|^2\le6\|y_s\|^2$.
\end{quotedthm}

\begin{quotedthm}{\cite[Lemma~7.26]{P51}}{}
If $k_0\in\pdsK$ is the identity on the coordinates outside $X$, then $c(k_0g)=c(g)$ for $g\in G$.
\end{quotedthm}

\begin{quotedthm}{\cite[Lemma~7.27]{P51}}{ (Far from $\pdsK$)}
For $g\in G$ with noncompact part $y$, $|c(g)|^4\le(1+\|y\|^2)^{-h_0}$.
\end{quotedthm}

\begin{quotedthm}{\cite[Lemma~7.28]{P51}}{ (Localisation)}
Let $0<\rho_1\le\frac16$ and $s\in R_\tau$ with $\|y_s\|\le\rho_1$. There is $t\in R_\tau$ with
$\theta(t)=\theta(s)$, $\|x_t\|\le4\rho_1$ and $\|x'_t\|\le4\rho_1$.
\end{quotedthm}

\begin{quotedthm}{\cite[Proposition~7.29]{P51}}{ (The cone)}
For every $s\in R_\tau$, $|\theta(s)|\le2\Rel\theta(s)+12e^{-100T}$.
\end{quotedthm}

\begin{quotedthm}{\cite[Lemma~7.31]{P51}}{}
$\pbvol(\mathrm{Large})\le\beta^{16D^2}\pbvol(\mathrm{Small})$.
\end{quotedthm}

\begin{quotedthm}{\cite[Lemma~8.22]{P51}}{ (Packing)}
\textup{(a)} Let $m\ge1$ and $R\ge0$. At most $(1+2R)^{2mf}$ vectors $v\in\auO_E^m$ satisfy $|\tau
(v_i)|\le R$ for all $i$ and $\tau$.
\textup{(b)} Let $\Sigma\subset\R^D$ be a set with $\|x\|\ge1$ and $\|x-x'\|\ge1$ for all $x\ne x'$ in
$\Sigma$ (Euclidean norm). Then $\sum_{x\in\Sigma}\|x\|^{-s}\le5^s$ for every real $s\ge D+2$.
\end{quotedthm}

\begin{quotedthm}{\cite[Lemma~8.32]{P51}}{ (Covolume of $\auO_E$)}
Let $\auPl$ be a set of representatives of the $f$ pairs of complex conjugate embeddings of the CM
field $E$ ($[E:\Q]=2f$). The image of $\auO_E$ under $a\mapsto(\tau(a))_{\tau\in\auPl}\in\C^f\cong\R^{2f}$
is a lattice of covolume $2^{-f}|D_E|^{1/2}$ for the Lebesgue measure $\prod_\tau d\Rel\,d\Iml$.
\end{quotedthm}

\begin{quotedthm}{\cite[Lemma~8.34]{P51}}{ (Siegel's lemma)}
Let $\Lambda$ be a unimodular lattice of $X_0$ and $V\subseteq X_0$ a subspace of codimension $l\ge1$. There
are $E$-linearly independent $X_1,\dots,X_l\in\Lambda\cap V^\perp$ with $V=\{v:h(X_i,v)=0\ \forall i\}$ and
\[
\prod_{i=1}^l\|X_i\|^{2f}\le(lf\,\pbcrd)^{lf}\,\ppidisc(\Lambda\cap V).
\]
\end{quotedthm}

\begin{quotedthm}{\cite[Lemma~2.1]{P54}}{}
For every $g\in\auU(V)(\auA_{E^+,f})$ the order of $\Gamma\cap gK^\circ g^{-1}$ divides
$n_{\mathrm{Aut}}$, and $n_{\mathrm{Aut}}\le\exp(x^{40})$. If $\Xi,\Xi'\in\poA(V;\vartheta)$ are
integral at $\aubar\iota$ and at $\iota$ respectively, then $n_{\mathrm{Aut}}\langle\Xi,\Xi'\rangle
\in\auO_\iota$.
\end{quotedthm}

\begin{quotedthm}{\cite[Lemma~2.2]{P54}}{ (Standing facts)}
Let $\vartheta$ be a twist.
\begin{enumerate}
\item[(L2)] $\Theta_\vartheta\circ\varrho(s)=\varrho_V(s)\circ\Theta_\vartheta$ for $s\in\auG$.
Hence $\Theta_\vartheta(e_0\poA(W_j))\subseteq\varrho_V(e_0)\poA(V;\vartheta)$.
\item[(L3)] Let $w$ be off the level, with residue field of $q$ elements, $P_w(X)=\sum_{i\ge0}
(-X)^iq^{\binom i2}T_{w,i}$ the operator of (N3) on $\poA(V;\vartheta)$, and $\alpha_w$ the number
with $\omega_\mu(1,d_1)\plone=\alpha_w\,\plone(d_1^{-1}T)$ at $w$ (by {\cite[(1)]{P54}}, $\alpha_w=\eta_w(1,d_1)$ does not depend on $\vartheta$, and
$|\alpha_w|=q^{r/2}$). Put $X_w=(\alpha_wq^{t-1})^{-1}$. Then $P_w(X_w)\,\Theta_\vartheta y=0$ for
every $y\in\poA(W_j)$.
\end{enumerate}
\end{quotedthm}

\begin{quotedthm}{\cite[Proposition~2.15]{P54}}{ (Three exchanges)}
Let $H_j(\vartheta)$ be the number {\cite[(11)]{P54}} below. There are functions $F_i$ $(i\in I)$ on
$\auU(W)(\auA)$, independent of $\vartheta$, a constant $C_{\mathrm x}\ge0$, and for every twist
$\vartheta$ a number $c'_\vartheta\in\C$, such that
\[
c'_\vartheta\,\Xi_{\vartheta,i}=F_i\quad(i\in I),\qquad|c'_\vartheta|^2\,H_j(\vartheta)=
C_{\mathrm x}.
\]
Moreover $yF_iK^\circ(\cdot,h^\circ)$ is integrable on $[\auU(W)]$ for $y\in\poA(W_j)$, and so is
$y\,\Xi_{\vartheta,i}K^\circ(\cdot,h^\circ)$.
\end{quotedthm}

\begin{quotedthm}{\cite[Proposition~2.24]{P54}}{ (The untwisted norm at all conjugates)}
For a main pair $p$ for $j$ the constant family $\|\varphi_{p,1}\|^2$ has height at most
$\aubud-\sum_{\tau\in\auPl}\ausinf(\kappa(\tau))$.
\end{quotedthm}

\begin{quotedthm}{\cite[Lemma~3.2]{P54}}{ (The unramified sum)}
Let $\pexmc=\pexmc_{2z,\alpha}$ be the coefficient of the vacuum of \S\ref{pc:ss-data} for $n=1$,
$M'=2z$ and the modulus $\mathfrak m$, on $G^{\mathfrak m}=\auU(1)(\auA^{\mathsf S}_{E^+,f})$ (the finite
points with trivial components at the places below $\mathsf S$). Then $\sum_{\gamma\in G^{\mathfrak
m}/K^{\mathfrak m}}|\pexmc(\gamma)|<\infty$ and
\[
\sum_{\gamma\in G^{\mathfrak m}/K^{\mathfrak m}}\pexmc(\gamma)=\frac{L^{\mathsf S}\bigl(z,(\nu\vartheta)^{-1}
\bigr)}{\zeta^{\mathsf S}_{E^+}(2z)},
\]
where $L^{\mathsf S}$ is as in {\cite[\S2.7]{P54}} and $\zeta^{\mathsf S}_{E^+}$ is the Dedekind zeta
function of $E^+$ without the Euler factors at the places below $\mathsf S$.
\end{quotedthm}

\begin{quotedthm}{\cite[Lemma~3.3]{P54}}{ (The archimedean factor)}
For $g\in\auU(1)(E^+\otimes\R)$, $\omega_{\alpha,\beta,\infty}(g,1)P=P$. Hence, in
Proposition~\ref{pc:rallis} for the pair $(E,V^{(z)})$, the constant function on $\auU(1)_\infty$ as
the type of the source, and the polynomial $P$, the archimedean integral is
$Z_\infty=\langle P,P\rangle_{\mathrm{Fock}}>0$; it does not depend on $\vartheta$.
\end{quotedthm}

\begin{quotedthm}{\cite[Proposition~3.4]{P54}}{ (Norm of the lift of a character)}
Let $p=(\nu,z)$ be a critical pair with $z\ge2$. There is a real number $A_p>0$, independent of
$\vartheta$, such that for every tame character $\vartheta$ of modulus $Q$
\[
\|\varphi_{p,\vartheta}\|^2=A_p\,(\auNm Q)^{2z-1}\,L^{\mathsf M\cup Q}\bigl(z,(\nu\vartheta)^{-1}\bigr),
\]
and $L^{\mathsf M\cup Q}(z,(\nu\vartheta)^{-1})>0$. Explicitly
\[
A_p=|\kappa_p|^2\,\prod_{w\in\mathsf M}(q_w-1)\cdot\dim W_{\lambda}\cdot\|v^+\|^{-2}\cdot c'\,
\langle P,P\rangle_{\mathrm{Fock}}\,\mathrm{vol}([\auU(V^{(z)})])\,\mathrm{vol}([\auU(1)])\,
\zeta_{E^+}(2z)^{-1},
\]
with $c'\ne0$ the constant of Lemma~\ref{pc:clean}(b) (all factors of $A_p$ are non-zero, and $A_p$ is real and positive, see the proof). In particular, for twists $\vartheta$, $\rho_p(\vartheta)=\|\varphi_{p,\vartheta}\|^2/
\bigl((\auNm Q_\vartheta)^{2z-1}\|\varphi_{p,1}\|^2\bigr)$.
\end{quotedthm}

\begin{quotedthm}{\cite[Lemma~4.3]{P54}}{}
The conclusions of {\cite[Theorem~3.2]{P49}} (with $B=x^{25}$, $N_j$ replaced by $r$) and of
{\cite[Corollary~3.3]{P49}} hold for $c_W=\bigotimes_\tau c_{W,\tau}$. Consequently the conclusions of
{\cite[Proposition~8.4{\rm(a)}]{P49}} hold for forms of type $e_W$ or $\aubar{e_W}$ on $\auU(W)$: they
are cusp forms, rapidly decreasing on Siegel sets, and those of a given level form a
finite-dimensional space.
\end{quotedthm}

\begin{quotedthm}{\cite[Lemma~4.7]{P54}}{}
$\pop_{\xi_0}$ is, up to a non-zero scalar, the joint harmonic $\Phi$ of
{\cite[Proposition~3.11]{P49}} for the matching pair $(\lambda,\lambda')$ of {\cite[Lemma~4.5]{P54}}.
\end{quotedthm}

\begin{quotedthm}{\cite[Proposition~4.8]{P54}}{ ((L4))}
Let $\vartheta$ be a twist and $\Xi\in\poA(V;\vartheta)$. Then $D_\vartheta\Xi\in\poA(W_j)$ (forms of
type $e_W$ and level $J_M$). Consequently statement {\rm(L4)} of {\cite[\S2.1]{P54}} holds: if
$\Xi\in\varrho_V(e_0)\poA(V;\vartheta)$ is orthogonal to $\Theta_\vartheta(e_0\poA(W_j))$, then
$D_\vartheta\Xi=0$.
\end{quotedthm}

\begin{quotedthm}{\cite[Lemma~4.10]{P54}}{ (Kernel scales)}
For all $l,n_N\ge1$ and $e\ge0$ there is $c=c_{l,n_N,e}\in\C^\times$ such that, for all $a,b$,
$\mathfrak m$, $\Psi$ and all $Q$ homogeneous of degree $e$:
\begin{enumerate}
\item[(KS1)] Let $q\ge1$ be an integer with $qM^2\in\mathfrak m$, let $\iota\colon\C\to\C_\ell$ with
$\ell\nmid q$, and suppose that the values of $\Psi$ and the coefficients of $Q$ are integral at
$\iota$ and that $g,h$ are integral with integral inverses at the places above $\ell$. Then
$|c\,\plfK_{a,b}[\Psi,Q](g,h)|_\iota\le1$.
\item[(KS2)] Suppose $l,n_N\le\Delta$, $e\le M$, $M^2\in\mathfrak m$, $|\sigma'(\Psi(x))|\le A$ for all
$x$ and all $\sigma'\in\auAut(\C)$, and $|Q^\sigma(z)|\le B$ for all $z$ with $\max_i|z_i|\le M^4$.
Then $|\sigma(c\,\plfK_{a,b}[\Psi,Q](g,h))|\le e^{x_1}AB$ for all $g,h$.
\end{enumerate}
\end{quotedthm}

\begin{quotedthm}{\cite[Proposition~5.39]{P54}}{ (The series for $M'>2n$)}
Let $M'>2n$. \textup{(a)} $\sum_{\gamma\in G^\psrm/K^\psrm}|\psrmc(\gamma)|<\infty$. \textup{(b)} On the space
$\mathcal B$ of bounded functions on $\auU(S)(\auA)$ that are right invariant under $K^\psrm$, the
series $\psrR=\sum_{\gamma\in G^\psrm/K^\psrm}\psrmc(\gamma)\rho(\gamma)$ converges absolutely in the
supremum norm and maps $\mathcal B$ to itself; it commutes with left translations, with right
translations by $\auU(S)_\infty$ and by the components at $S(\psrm)$, and with $\kappa*$ for every
integrable $\kappa$ on $\auU(S)_\infty$; it maps continuous functions to continuous functions.
\textup{(c)} $\psrR$ is injective on $\mathcal B$.
\end{quotedthm}

\begin{quotedthm}{\cite[Lemma~5.40]{P54}}{ (Weil's formula for a definite target)}
The statement \textup{(SW)} of Lemma~\ref{pc:SW}, including the absolute convergence and the bounds by
a power of the height of $g\in\auU(\mathcal W)(\auA)$, holds when $T$ is positive definite at every infinite place and $M'>2n$: these imply the
hypotheses (i), (ii) of Lemma~\ref{pc:SW}.
\end{quotedthm}

\begin{quotedthm}{\cite[Lemma~5.41]{P54}}{ (Growth of the Siegel Eisenstein series)}
Let $\pcW$ be a space of rank $2n$ with a Lagrangian, $T$ a space of rank $M'>2n$, and $\Phi$ a datum of
$\auHom(\pcW,T)$. There are $C,A$ with $\sum_L|\psZ_L(\omega(x,1)\Phi)|\le C\,\mathrm{ht}(x)^A$ for all
$x\in\auU(\pcW)(\auA)$, the sum over all Lagrangians of $\pcW$.
\end{quotedthm}

\begin{quotedthm}{\cite[Lemma~7.27]{P54}}{}
Let $q=q(\psi,z')$ as in {\cite[Definition~7.26]{P54}}.
\begin{enumerate}
\item[(a)] $q$ is a critical pair; it is flat ($c(\tau)=z'$ for all $\tau$, so $l_\tau=0$), and in the
notation of {\cite[\S3.1.1]{P54}} the character on $\auU(1)$ of its Weil representation is
$\alpha_{q,\vartheta}=\nu_q\vartheta\,\upsilon^{z'}=\psi\vartheta$.
\item[(b)] For every tame $\vartheta$ of modulus $Q$ and every $h$,
$\varphi_{q,\vartheta}(h)=c_{\mathrm{vac}}^{\,2z'}\sum_{c\in\psuC}k_{q,\vartheta}(x_c,h)$.
\item[(c)] For $g\in\auU(1)(\auA_{E^+,f})$ put $G_{q,\vartheta}(g)=\sum_he_h^{-1}\,\aubar{\varphi_{q,
\vartheta}(h)}\,k_{q,\vartheta}(g,h)$, the sum over the classes $h$ of $\auU(V')$ at the level of
$q$. Then $G_{q,\vartheta}(g)=\|\varphi_{q,\vartheta}\|^2/(N_1c_{\mathrm{vac}}^{\,2z'})$ for every $g$.
\end{enumerate}
\end{quotedthm}

\begin{quotedthm}{\cite[Theorem~7.43]{P54}}{ ({\cite[Proposition~2.24]{P54}})}
For a main pair $p=(\nu,z)$ for $j$ the constant family $\|\varphi_{p,1}\|^2$ has height
at most $\aubud-\sum_{\tau\in\auPl}\ausinf(\kappa(\tau))$.
\end{quotedthm}

\begin{quotedthm}{\cite[Proposition~6.7]{P55}}{ ((L))}
Let $w$ be split with $w,\aubar w$ off the level of the pair, let $y_0$ be a form on $\auU(W)$
spherical at $w$ with $P^W_w(X)y_0=p(X)y_0$, $p\in\C[X]$, and let $y=\Theta(\Phi)y_0$ be its theta
lift for a datum $\Phi$ which is the characteristic function of the integral matrices at $w$. Then
$y$ is spherical at $w$ and
\begin{equation*}\tag*{{\cite[(10)]{P55}}}
P_w^V(X)\,y=\prod_{l=0}^{t-1}(1-q^l\alpha_wX)\cdot p\bigl(q^t\alpha_w\beta_w^{-1}X\bigr)\,y .
\end{equation*}
Equivalently, if $y\ne0$ and $\gamma_1,\dots,\gamma_r$ are the roots of $H^W_w(y_0;X)$, the roots
of $H^V_w(y;X)$ are $\alpha_wq^l$ ($0\le l<t$) and $q^t\alpha_w\beta_w^{-1}\gamma_k$; this is
{\cite[(1)]{P55}} with $\lambda_w=\alpha_w$ and rescaling factor $q^t\alpha_w/\beta_w$.
Moreover the lift is $\auG_O$-equivariant as in {\cite[Proposition~14.27]{P46}}.
\end{quotedthm}

\begin{thebibliography}{L-values}
\bibitem[Row D]{P46} D.~A.~Trejo Pizzo, \emph{The sign of the radial intersection form on the idele class group}, preprint, 2026, \url{https://github.com/dtrejopizzo/riemann/tree/main/04-papers/46-row-d-sign-radial-intersection-form}.
\bibitem[Theta]{P49} D.~A.~Trejo Pizzo, \emph{The Short Forms Theorem: theta lifts, complex multiplication and divisibility}, preprint, 2026. % LINK-PLACEHOLDER
\bibitem[Moment]{P51} D.~A.~Trejo Pizzo, \emph{The Spanning Theorem: the second moment of a Poincar\'e map}, preprint, 2026. % LINK-PLACEHOLDER
\bibitem[L-values]{P54} D.~A.~Trejo Pizzo, \emph{The Covolume Theorem: one theta lift, Hecke $L$-values and test forms of least valuation}, preprint, 2026. % LINK-PLACEHOLDER
\bibitem[Galois]{P55} D.~A.~Trejo Pizzo, \emph{The Covolume Theorem: regular families of twists, Selmer groups and Ribet's method}, preprint, 2026. % LINK-PLACEHOLDER
\bibitem[Hodge]{P56} D.~A.~Trejo Pizzo, \emph{Explicit cup exponents for local Galois representations beyond the Fontaine--Laffaille range}, preprint, 2026. % LINK-PLACEHOLDER
\bibitem{HoweTheta} R.~Howe, \emph{$\theta$-series and invariant theory}, in: Automorphic Forms,
Representations and $L$-functions, Proc.\ Sympos.\ Pure Math.\ \textbf{33}, Part~1, Amer.\ Math.\
Soc., 1979, 275--285.
\bibitem{KashiwaraVergne} M.~Kashiwara and M.~Vergne, \emph{On the Segal--Shale--Weil
representations and harmonic polynomials}, Invent.\ Math.\ \textbf{44} (1978), 1--47.
\bibitem{KudlaSplitting} S.~S.~Kudla, \emph{Splitting metaplectic covers of dual reductive pairs},
Israel J.\ Math.\ \textbf{87} (1994), 361--401.
\bibitem{GoodmanWallach} R.~Goodman and N.~R.~Wallach, \emph{Symmetry, Representations, and
Invariants}, Graduate Texts in Mathematics \textbf{255}, Springer, 2009.
\bibitem{PlatonovRapinchuk} V.~Platonov and A.~Rapinchuk, \emph{Algebraic Groups and Number
Theory}, Pure and Applied Mathematics \textbf{139}, Academic Press, 1994.
\bibitem{HarishChandraAF} Harish-Chandra, \emph{Automorphic Forms on Semisimple Lie Groups}, Lecture
Notes in Mathematics 62, Springer, Berlin, 1968.
\bibitem{BorelJacquet} A.~Borel and H.~Jacquet, \emph{Automorphic forms and automorphic
representations}, in: Automorphic Forms, Representations and $L$-functions, Proc.\ Sympos.\ Pure
Math.\ \textbf{33}, Part~1, Amer.\ Math.\ Soc., 1979, 189--202.
\bibitem{Knapp} A.~W.~Knapp, Representation Theory of Semisimple Groups: An Overview Based on Examples, Princeton Univ. Press, 1986.
\bibitem{Schmid75} W.~Schmid, Some properties of square-integrable representations of semisimple Lie groups, Ann. of Math. (2) 102 (1975), 535--564.
\bibitem{FollandPS} G.~B.~Folland, Harmonic Analysis in Phase Space, Ann. of Math. Stud. 122, Princeton Univ. Press, 1989.
\bibitem{BorelFin} A.~Borel, \emph{Some finiteness properties of adele groups over number
fields}, Publ.\ Math.\ Inst.\ Hautes \'Etudes Sci.\ \textbf{16} (1963), 5--30.
\bibitem{Godement63} R.~Godement, \emph{Domaines fondamentaux des groupes arithm\'etiques},
S\'eminaire Bourbaki, vol.~8 (1962--1964), expos\'e 257, 201--225 (Th\'eor\`eme~7, p.~211);
\url{http://www.numdam.org/item/SB_1962-1964__8__201_0}.
\bibitem{HC66} Harish-Chandra, \emph{Discrete series for semisimple Lie groups II. Explicit
determination of the characters}, Acta Math.\ \textbf{116} (1966), 1--111.
\bibitem{Jacobowitz} R.~Jacobowitz, \emph{Hermitian forms over local fields}, Amer.\ J.\ Math.\
\textbf{84} (1962), 441--465.
\bibitem{Landherr} W.~Landherr, \emph{\"Aquivalenz Hermitescher Formen \"uber einem beliebigen
algebraischen Zahlk\"orper}, Abh.\ Math.\ Sem.\ Univ.\ Hamburg \textbf{11} (1935), 245--248.
\bibitem{Neukirch} J.~Neukirch, \emph{Algebraic Number Theory}, Grundlehren Math.\ Wiss.\
\textbf{322}, Springer, 1999.
\bibitem{IchinoSW} A.~Ichino, \emph{A regularized Siegel--Weil formula for unitary groups}, Math.\
Z.\ \textbf{247} (2004), 241--277.
\bibitem{WeilSiegel} A.~Weil, \emph{Sur la formule de Siegel dans la th\'eorie des groupes
classiques}, Acta Math.\ \textbf{113} (1965), 1--87.
\bibitem{Cassels} J.~W.~S.~Cassels, \emph{An Introduction to the Geometry of Numbers}, Grundlehren
der mathematischen Wissenschaften 99, Springer, Berlin, 1959.
\bibitem{HenkMinima} M.~Henk, \emph{Successive minima and lattice points}, Rend.\ Circ.\ Mat.\ Palermo (2) Suppl.\ \textbf{70} (2002), 377--384; arXiv:math/0204158.

% refs-B-size.tex: new references cited by proofs/B-size.tex.
% None: the Ostrowski--Taussky inequality and Hua's integral are proved in the text.
\bibitem{ReedSimonI} M.~Reed and B.~Simon, \emph{Methods of Modern Mathematical Physics I: Functional
Analysis}, revised ed., Academic Press, 1980.


% new bibitems for C-triangle (Knapp already cited elsewhere in the paper)

% New bibitems for C-triangle2.tex (the others already exist: FollandPS in refs-A-coefficient,
% Schmid75 in refs-A-archimedean, HarishChandraAF in refs-A-closure). None new.

% New bibitems for C-triangle3.tex: none (no new reference is cited).

% New bibitems for C-triangle4.tex: none (no new reference is cited).

% refs-C-grid.tex: no new bibitems.

% refs-C-grid2.tex: no new bibitems (pk:hook, pk:eq-dual, au:shortbasis are cited by label).

% refs-C-grid3.tex: no new bibliography items (all citations are internal labels).

% refs-C-grid4.tex: no new bibliography items (all citations are internal labels).

% New bibitems of C-grid5: none.

% refs-C-grid6.tex: no new bibliography items (all citations are internal labels).

% no new references

% refs-C-grid8.tex: no new bibliography items (PlatonovRapinchuk is already in the bibliography of the paper).

% refs-C-grid9.tex: no new bibliography items (C-grid9 cites only labels of the paper).

% refs-C-grid10.tex: no new bibliography items (GoodmanWallach is already in the paper).
\end{thebibliography}

\end{document}
